% Part4_Baryomesonic_era.md — sanitized presentation layer (generated; see source/CONVERSION_LOG.md)
\chapter{\texorpdfstring{$\langle qu(b(i(t|t))ick)\rangle$ — \textbf{Part 4 · Baryomesonic\_era} — Colour-Sealed Composites: Baryons, Mesons, Junctions}{Qubitick - **Part 4 - Baryomesonic _era** - Colour-Sealed Composites: Baryons, Mesons, Junctions}}

\textbf{Dep.} Part 1 (axioms A1–A9; conventions C1–C13; qubit-era definitions D0.1–D1.12; calibration P1–P5; pair canon CT0.\textit{/CT1.1–CT1.11b; ledger dynamics L0–L5); Part 2 (definitions D2.1–D4.3; geometry CT1.12–CT2.11; electroweak EW.}; electromagnetic EM.\textit{; colour CT3.1–CT3.22a; saturation CT4.}); Part 3 (the fermion layer: Def-H.1–Def-H.43, FT.1–FT.155); \textbf{the Census annex} (Appendix A of this volume; the computed cells, cited row-wise). Format: Part 1 §0.4. \textbf{Namespace.} This part owns: Def-BM.1–Def-BM.20; BM.1–BM.152. Tiers 4.0, BM.1–BM.24 in series order. \textbf{DAG law.} Citations resolve to Parts 1–3, Census cells, or earlier tags of this part. Nothing in the canon cites this part. \textbf{0.1 Charter.} The baryomesonic era is the layer of colour-sealed composites: the algebraic sealings of D3.9 (baryon $\varepsilon_{ijk}q^iq^jq^k$; meson $\delta^i{}_j\,q^jr_i$) realized on bound-chain constituents (Def-H.4) at shared junctions. Its foundation stack, by citation: (i) the colour and sealing machinery [CT3.1–CT3.22a; D3.8, D3.9; D2.14; CT4.1–CT4.7]; (ii) the derived sealing grammar and junction structure [FT.29, FT.34(5), FT.19, FT.39]; (iii) the computed baryon-type cells [Census FT.35: the first colour-active, first $D^-$-carrying floor; Census FT.49: the baryon-type sealing phase $\Theta_{\mathrm{int}}\in4\mathbb Z$]; (iv) the channel machinery [FT.109–FT.119: the $\ell$/q-channels with the $\varepsilon$-selected $C_3$ and confinement blindness; FT.141–FT.145: the rung-one channel and the arrow law]; (v) the carrier and parity structure [FT.122–FT.134]; (vi) the sector formalisms [FT.102–FT.108]; (vii) the conservation and symmetry layer [Tier BM.4: the Janus balance and the telescoping endpoint read, the continuum promotion of the hop, the $C_3$ and the slot algebras, the stress-energy tensor of a sealed composite, the canonical-pair obstructions, the differential conservation law]; (viii) the calibration layer [Tier BM.5: the Born width, the coherence volume and energy density, the curvature of a sealed composite, the dimensionless lock identities, the neutron row of the identification layer]; (ix) the flicker layer [Tier BM.6: the topology-blind ensemble of a sealed composite and its derived weights, the path-integrated mass from the curvature law, the binding read, the static force matrix on the two charge classes, the transition with its forced carrier, the family naming rows]; (x) the CPT-sector and charge-restriction layers [Tier BM.7: the EM$^\pm$ sector labels and the EM scar, the mirror involution and its exact phase, the proton and antiproton rows, and the charge-constrained probability space with its three mass routes]; (xi) the additive charge read and the mass-route adjudication [Tier BM.8]; (xii) the mesonic flicker test [Tier BM.9: the contraction grammar, the colour projector, and the protected junction]; (xiii) the metric-energy and vacuum-entropy layer [Tier BM.10: the dispersion read of a sealed composite from the emergent metric with its kinetic energy and zero-point bound, the vacuum-balanced-entropy reframing of the rest mass, the force-structure ladder with its strength bar and the inter-family constraint, and the route-artifact adjudication closing the counting trap]; (xiv) the decay-vertex layer [Tier BM.11: the family transport-profile map discharging the Yukawa-labelled gap, the family mass form with its positivity, splitting law and sum rule, the rung-1 isospin row and the neutrality of the charge vertex, the forced neutral member and the Q-value sign theorem, the angular-momentum budget with its absorption adjudications, the member-wise curvature rows, and the per-tick time evolution of the assembled reaction]; (xv) the electromagnetic self-interaction layer [Tier BM.12: the chart-side EM source read of a sealed composite, the scar's field energy and the mass-splitting theorem, the infinite-range seal asymmetry with its UV/IR finiteness rows, the two-body potential with the binding category bar, and the beta-decay re-adjudication with the endpoint-to-average refutation]; (xvi) the full-scope mass-splitting layer [Tier BM.13: the Lyapunov foundation with its iterated contraction law and class-degenerate zero, the complete-force-operator sector census with its five zeros and one monomial, the metric background's zero contribution and the non-uniform clock with its Mandelstam–Tamm resolvability bound, the two binding zeros and the unconstructed atom, and the exact symbolic ratio with its named order residue]; (xvii) the constituent-transport-splitting layer [Tier BM.14: the continuation-symmetry census of the two vertex types with its three annihilation rows, the observation frames and the charge-blindness of the metric, the constituent transport splitting that fills the residue BM.55(4) names, the six-state quark assembly on the three Platonic mass floors with its branch order and its two inverted hierarchies, the coupling-blindness theorem and the named obstruction to $\alpha_B$, the extended six-term decomposition with its exact critical strength and sign theorem, and the decay re-adjudication with its exact closure margin]; (xviii) the organic-bridge layer [Tier BM.15: the operator lift of the EM correction with its charge-blind linear trace and its quadratic invariant, the derived $\mathbb Z_2$ Wilson sector with the two-torsion and calibration-weight obstructions to the real holonomy lift and its structured obstruction record, the constrained joint entropy as a fiber cardinality with its stationary geometric factor of one and its named missing enumeration, the bridge convergence with the non-derivation theorem reducing Gap 7 to one named scale ratio and two named attachments, the two-bit daughter-floor theorem with the formed-closed dichotomy and its exact margin ladder, and the five-route opening-term census closing the beta-decay energy audit]; (xix) the chiral-assignment layer [Tier BM.16: the left–right assignment rule coupling the four derived gradings and establishing the assignment rule's forced shape, the ray-resolved spin census with the opposite-ray vertex table, the weight–branch–core attachment theorem with the sign-coherent quark mass splitting, the nilpotent chiral vertex operator with its flip, rate and arrow rows, and the vertex read-specification with the exact trichotomy of the beta-decay opening branch]. (xx) the mass-ratio-census layer [Tier BM.17: the Platonic winding census with the ladder law on both priced windings, the Chebyshev tower closing at the dodecahedral top weight, and the intra-generational ratios as exact sine ratios with their minimal polynomials; the inter-generational ratios as the three roots of one cubic form with three channel constants, the cross-channel identity between the lepton channel's heaviest line and the quark channel's cubic constant, and the scale-free hierarchy-order theorem; the six-state ratio census with the splitting order, the doublet ratios, the two inverted hierarchies and the three pairwise non-proportional ladders; the bridge census with its six exact factors, the intra-currency closure and the localized single residue; the inter-family bridge in full with the four-currency census, the exact incommensurability of the two formed ratios, and the three-number localization; the neutron-proton mass ratio with its exact range, its irrational critical locus and its total trichotomy; the Q-value in two formed reads with both margin ladders and the read-relative verdict; and the census adjudication] (xxi) the electromagnetic self-energy correction layer [Tier BM.18: the category adjudication voiding the chart-side distribution factor and deriving the formed exterior value one half on two distribution-free routes; the running coupling as an exact geometric profile on the envelope ratio, the lapse and the curvature lock, vanishing at the member's own Born width and growing with the energy scale; the self-energy as an ensemble average over the flicker with one derived attenuation, the crystallization weight; the branch-blindness theorem refuting a defect-conditioned retention on four grounds; the corrected neutron–proton splitting with its parameter-free bar, its universal ratio bound and its ensemble-robust incommensurability; the Q-value on both formed reads with the uniform closure and the untouched specified branch; and the correction's adjudication] (xxii) the flicker-integral and rest-mass layer [Tier BM.19: the flicker integral replacing the crystallized snapshot, with the derived ensemble distribution, the materialization restriction and the two crystallization weights; the vertex superposition with its exact fibre weights and the two positive parts of the transport splitting; the rung-one $\nu$ line's seal-derived rest mass, its strict positivity, its single-ray occupancy and the four bars against a seesaw-type mechanism; the rest-mass Q-value on the specified and pure-pole reads with the exact critical locus; the relating equation in four exact forms with its flicker weighting and its integer bound; the materialized splitting with its two-sided parameter-free bar and its two tightened bounds; and the tier's adjudication] (xxiii) the six-objective-resolution layer [Tier BM.20: the energy-scale graph of the forward and inverse reads with the derived crossing strength, its exact enclosure and the composition-dependent crossing count; the dimensional bridge, with the $\mathfrak q$-exponent census of the inventory, the seven rows derived or adjudicated, the two bridges separated and the one calibration slot declared; the three-constituent assembly's two formed eigenvalue equations and the tensor-product referent census; the curvature composition's back-reaction fixed point with its three exact solution rows, its pole locus and the non-selection theorem; the general crystallization weight with its measure-free bounds, its threshold family and the refuted defect conditioning; and the kind structure of the arrow, with the spin-parity refutation at the floor and the exchange-structure referent census] (xxiv) the antimatter-and-resonance layer [Tier BM.21: the CP-conjugate chain derived from the three reversals with the T-block on the forward-only tick and the Lyapunov contraction, its exact phase, the mass-degeneracy theorem and the forcing theorem from the ledger's own total-charge-zero balance, with the Dirac parallel as an inventory statement; the CP-pair vacuum fluctuation with the Janus resolvent's own creation weight, its exact maximum one quarter and its two endpoint zeros; the two-chain state space discharging BM.6(2)'s exchange-space specification, the meson as the unique $\delta$-sealed $2$-juncture, the CP-pair saturation of the neutral class at exactly six, the two additive mass routes with their vanishing seal-binding and the named two-chain ensemble residue, and the amplitude/norm reconciliation on the locus distinction; the exchange operator as the flip composed with the CP transport, an unconditional involution with its eigenspace census and its statistics-blind orientation block; the spin–statistics connection $\varepsilon=(-1)^{|k|}=(-1)^{2j}$ from the fourth-root phase alphabet and the winding additivity, closing deficit 11 with the sector choice priced; the resonance spectrum with the transport characteristic polynomial $(z+1)^{2n}/z^{n}$ at every chain count, the degree-16 eigenvalue field, the multiset-orbit law with stabilizer multiplicities, the one exact three-chain collision and the minimal-polynomial degree census; the materialization threshold in its three exact readings with the never-empty locus and the no-forbidden-state theorem; the decay width with its three exact zeros, its bound and its three-tick lifetime floor; and the groupings, the ceilings and the octet adjudication] (xxv) the annihilation-channel and concrete-spin layer [Tier BM.22: the Janus register with its two nilpotent crossing moves, its master generator's decomposition, kernel, spectrum and derived detailed balance; the c-threshold adjudication with the exact subluminal deficit $1-v=m^2/(E(E+p))>0$ and the register boundary that replaces it; the zig-zag as the resolvent's evenness, every leg forward and the record involution fixing each completed word; the pair channel with its charge balance, exact threshold, weight and three CP legs, and the two-photon bar on four grounds; the annihilation channel with its exact direction ratio, detailed-balance indicator and entropy destination, closing residue (R9) at the channel level; the spin–winding relation $\Theta=4j$, the multiplet defect ladder in Chebyshev form with its standing zero and ceiling, and the dimensional bar against fusing the winding with the frame-holonomy angle; the refutation of a spin-gated tunnel as hypothesis (W7) on the four standing grounds, with the bound/transit domain row in its place; the winding-parity identification with the statistics and kind reads separated; the three-chain exchange structure with the $S_n$ character formula $d^{\,c(\sigma)}$, the vanishing carrier wedge against the one-dimensional colour wedge, and the multiplicative exchange law discharging (R7) at $n=3$; and the meson's unconditional boson sign against the baryon's exact odd-count fermion criterion] (xxvi) the closing layer [Tier BM.23: the Haar–Poisson duality with its twirl fixed point, complementary factorization, unique common fixed point and proper-time contraction law; the derived lapse $N(k,E)=\sqrt{1-G/(4\pi k)}$ with its formed domain and its envelope row, closing the lapse gap; the curvature composition selected as $\mathsf g=1/(1+G)$ by Lyapunov contractivity and minimal lapse-domain loss, completing the running coupling; the momentum-space $\Sigma$-resolvent with its form uniqueness, pole reads, on-shell measure and two-body phase space, closing item 17; the multi-chain charge census with the stretched states and the gauged read, closing Q10 and deficit 10; the $S_3$–colour-vertex equivariance with the exchange structure at every $n$ and the statistics per class, closing (R7$'$); the two-chain ensemble measure with the general $w_{\mathrm{mat}}^{(n)}$ law, the sector-mass pincer and the formed meson binding, closing (R8); the seed-tensoring mechanism as a finite-group twirl, closing item 19; the four convention declarations (A-BM.10–A-BM.13), closing item 16; the dodecahedron selection theorem with the five-currency ceiling census, closing deficit 12 as adjudicated; the Jensen residue as an exact rational function with its two-chain companion, closing (R6); the nilpotent supply vertex with the pair amplitude and the two symbolic cross-sections; the house cleaning and the calibration-register consolidation; and the era closure with the balanced ledger] (xxvii) the continuum-limit layer [Tier BM.24: the complete force operator acting over the graph, with the exact block census and the eight-step composed action; the continuum-limit function of content and span with its five closed forms and four formed-domain conditions; the metric–energy–continuum map with the dictionary bijection, the pure-trace linearized solve and the on-shell inverse read; the breaking threshold $E_{\mathrm{crit}}=2r/\mathfrak h$ with the envelope corollary; the locked-sector complement at the breaking locus with the absorption anchors, the one-way membrane and the retention theorem; the universality of the limit on one operator, one arena, one action and three fingerprints; and the derivation of $\mathrm{CT2.6}(\mathrm{ii})\text{--}(\mathrm{iii})$ as theorems — the unbounded ladder and the $\sqrt3$ spectral bound — leaving the extra-axiomatic list empty and closing the foundational era] Its enumerative substrate is the Census block: 25 cells, each carrying its consumption row (Census §Consumption). \textbf{0.2 Standing bars.} Species names are admitted in this part through the identification layer Def-BM.3 alone: a name is a label on a derived structural tuple, it carries no datum, and it occurs in no \textbf{Thm} clause of Parts 1–3, whose own label-discipline rows stand unmodified. No external data outside quarantined audit rows [Census FT.114, FT.120; FT.73(5) pattern]; no measured external constant is identified with any corpus quantity [P5]; topology labels are internal bookkeeping [Def-H.32(4)]; amplitude $\ne$ mass [FT.79(5)]; rung $\ne$ scale [D0.8].

\phantomsection
\section[§Cal-BM — Free internal parameters of the baryomesonic era…]{\texorpdfstring{§Cal-BM — Free internal parameters of the baryomesonic era (P4-class)}{SCal-BM - Free internal parameters of the baryomesonic era (P4-class)}}

None introduced by this part. Inherited and undeclared: the family-channel overall scale $\mu$ [Cal-F; FT.113, FT.119 — every ratio is $\mu$-free]. \textbf{Declared slot, value $\neg$issued (A-BM.9, Tier BM.20):} the bridge's dimensionless calibration ratio $\rho_{\mathrm B}=\hat\mu'^2\mathfrak h^2/(\ln2)^2\in\mathfrak q^0$ [BM.69(2); BM.98; BM.102(5)] — a packaging of Cal-F's $\hat\mu'^2$ and P1's $\mathfrak h$, hence $\neg$a new free parameter and $\neg$a new dimensionful scale; no numerical value attaches to it here or anywhere [P5; BM.69(4)'s four grounds stand: the slot is declared, $\neg$derived]. No second dimensionful scale exists [P1; the derived no-go: Part 1 §E]. Any numerical use of an inherited scale beyond its Part-3 loci requires a logged declaration first.

\textbf{Consolidated calibration register (Tier BM.23, BM.143(4) — a consolidation, $\neg$an addition; the parameter count is unchanged).}

%dsp-level:0
\[
(\mathrm{C\text{-}BM.1})\quad\rho_{\mathrm B}
\]

%dsp-level:0
\[
= \hat\mu'^2\mathfrak h^2/(\ln2)^2\in\mathfrak q^0:\ \ \text{the bridge ratio; slot declared}\ [\mathrm{A\text{-}BM.9}],
\]

%dsp-level:0
\[
\ \text{value}\ \neg\text{issued}\ [\mathrm{BM.143}(4)];
\]

%dsp-level:0
\[
(\mathrm{C\text{-}BM.2})
\]

%dsp-level:0
\[
\quad e_0:\ \ \text{the numerical value of the electric coupling}\ [\mathrm{EM.25}(7);\ \mathrm{BM.88}(8);\ \mathrm{BM.133}(6)];
\]

%dsp-level:0
\[
\ \ \text{value}\ \neg\text{issued};
\]

%dsp-level:0
\[
\ \neg\text{fixed by measurement}\ [\mathrm{P5}];
\]

%dsp-level:0
\[
(\mathrm{C\text{-}BM.3})\quad\mu_\nu,
\]

%dsp-level:0
\[
\ \mu,
\]

%dsp-level:0
\[
\ \rho_\ell,
\]

%dsp-level:0
\[
\ \rho_\nu:\ \ \text{the channel scales, inherited from}\ \mathrm{Cal\text{-}F},
\]

%dsp-level:0
\[
\ \neg\text{declared in this part}\ [\mathrm{FT.145}(3);\ \mathrm{BM.83}(5),(7)].
\]

Each row is a calibration datum of the P4/A-BM.9 species, $\neg$a structural gap; no row carries a value; $\mathfrak h$ remains the single dimensionful scale [P1].

\phantomsection
\section[Tier 4.0 — Composite Vocabulary]{\texorpdfstring{Tier 4.0 — Composite Vocabulary}{Tier 4.0 - Composite Vocabulary}}

\begin{definition}[{\texorpdfstring{Def-BM.1 — Colour-sealed chain composites}{Def-BM.1 - Colour-sealed chain composites}}]\label{tag:Def-BM-1}
\textbf{Def.} \textbf{Constituent} $:=$ a triplet-attached bound chain [D3.8: one fundamental slot index $q^i$ (or dual $r_i$) at the attachment; Def-H.4: a Janus-closed bound chain, terminus locus; $\delta h\ge2$]. \textbf{Junction} $:=$ a shared vertex hosting one elementary invariant contraction of the constituents' slot indices [CT3.17: $\delta^i{}_j$ for a pair, $\varepsilon_{ijk}/\varepsilon^{ijk}$ for a triple; trivalent slot structure CT3.2].

%dsp-level:0
\[
\textbf{baryon:}\quad B(c_1,c_2,c_3)\ :=\ \varepsilon_{ijk}\,c_1^i\,c_2^j\,c_3^k
\]

%dsp-level:0
\[
\qquad [\mathrm{D3.9}\ \text{algebraic form, chain-realized at one junction}];
\]

%dsp-level:0
\[
\textbf{meson:}\quad M(c,\bar c)\ :=\ \delta^i{}_j\,c^j\,\bar c_i
\]

%dsp-level:0
\[
\qquad [\mathrm{D3.9};\ \bar c:\ \text{the conjugate constituent}\ (\omega\text{-reversed chain with dual attachment},\ \mathrm{D2.6},\ \mathrm{D3.8})].
\]

\textbf{Validity clause.} For valid ambient $C$ [D2.11], $\pi_C$ is total on $\mathrm{Colour}(U)$ [CT3.18 $+$ D2.11 applied to the colour ledger; D2.14]: every attachment has its contraction partners inside the composite; a restriction $\Sigma$ with $\partial_{\mathrm{col}}(\Sigma;C)\ne\varnothing$ has $\mathcal C(\Sigma;C)=0$ and its entire colour sector vanishes [D2.14, FT.19: all-or-nothing] — the mask admits only singlets [CT3.18]. \textbf{Read clause.} Each constituent retains its own terminus tuple $(q,\mathbf S,m_{\mathrm{hid}})$ [Def-H.9]; the junction adds the sealing-phase class and the lock content [Census FT.49: $\Theta_{\mathrm{int}}\in4\mathbb Z$ at the baryon-type cell; Census FT.35: junction $D^-$ content]; the card assembly is Def-BM.2.

\end{definition}

\begin{definition}[{\texorpdfstring{Def-BM.2 — The composite card}{Def-BM.2 - The composite card}}]\label{tag:Def-BM-2}
\textbf{Def.} The \textbf{card} of a sealed composite refines the floor-card row template [Census FT.28 — the minimal fermion cell; every sealed-composite card refines its row grammar: registration data $(\sigma_e,\kappa)$, $\Theta_{\mathrm{int}}$, winding, holonomy, aspect rows, mask dimension, rung counts, ternary core read] with:

%dsp-level:0
\[
(1)\ \text{constituent rows: the}\ \Pi_v\ \text{projections of Def-H.9 per constituent};
\]

%dsp-level:0
\[
\qquad (2)\ \text{the junction row: sealing-phase class and lock content}\ [\mathrm{Def\text{-}BM.1}\ \text{read clause}];
\]

%dsp-level:0
\[
(3)\ \text{the sealing-grammar row:}\ N_{\mathrm{strong}}
\]

%dsp-level:0
\[
= \log_2 g\bigl(|\mathrm{Colour}(\Sigma)|\bigr)\ \text{with the grammar counts}\ g\ [\mathrm{CT3.17a}(2)\text{–}(3);\ \mathrm{FT.41}(b)];
\]

%dsp-level:0
\[
(4)\ \text{the sealing-profile row}\ [\mathrm{FT.34}(5)\ \text{partition rule}];
\]

%dsp-level:0
\[
\qquad (5)\ \text{the winding-ledger rows}\ [\text{additive composition}\ \mathrm{FT.34}(3);\ \text{Census FT.43},\ \text{FT.51}].
\]

Card \textbf{values} are Census-cell data at the computed cells; card \textbf{grammar} is this definition. No card row is a species identification [P1, P5].

\end{definition}

\phantomsection
\section[Tier BM.1 — Observability of Sealed Composites]{\texorpdfstring{Tier BM.1 — Observability of Sealed Composites}{Tier BM.1 - Observability of Sealed Composites}}

\begin{theorem}[{\texorpdfstring{BM.1 — Confinement inheritance and the trace-only read}{BM.1 - Confinement inheritance and the trace-only read}}]\label{tag:BM-1}
\textbf{Thm.}

%dsp-level:0
\[
(1)\quad\textbf{sealing is the observability switch:}\ \ \text{a sealed composite is colour-complete on its own support}\ (\partial_{\mathrm{col}}=\varnothing\implies\mathcal C=1)\ \text{and its colour sector survives masking generically}
\]

%dsp-level:0
\[
\qquad\bigl[\mathrm{FT.39}(1):\ \mathcal C=1\Rightarrow F^\Sigma_{\mathrm{strong}}=F_{\mathrm{strong}}|_\Sigma\ne0\ \text{generically, exact exception locus}\ \{f=0\}\cup\{\mathrm{Colour}=\varnothing\}\bigr];
\]

%dsp-level:0
\[
\text{any unsealed constituent read vanishes: one dangling attachment zeroes the whole colour sector}\ [\mathrm{FT.19}];
\]

%dsp-level:0
\[
\ \neg\exists\ \text{triplet pole, physical norm}\ 0\ [\mathrm{CT3.22}];
\]

%dsp-level:0
\[
(2)\quad\textbf{the one-body trace law:}\ \ \text{every Hermitian slot operator on a sealed singlet reads only its trace — baryon:}\ \mathfrak T(H)
\]

%dsp-level:0
\[
= 3a\ \text{for the covariant class}\ H
\]

%dsp-level:0
\[
= aI_3+zP+z^*P^T;
\]

%dsp-level:0
\[
\ \text{meson:}\ 0\ \text{(including the trace part)}\ [\mathrm{FT.117}(1)\text{–}(3)];
\]

%dsp-level:0
\[
\text{hence the}\ C_3\text{-doublet datum}\ (B,\theta)\ \text{is annihilated on every colour-sealed object, and hadron-level triplets cannot carry the family channel}\ [\mathrm{FT.117}(2),(7)];
\]

%dsp-level:0
\[
(3)\quad\textbf{the sealed colour factor:}\ \ \text{singlet-sealed chains keep the undressed arena}\ 9\ \text{— the invariant lines are}\ 1\text{-dimensional}\ [\mathrm{FT.117}(5):\ \dim(\mathbb C^3)^{\otimes3\,\mathfrak{sl}(3)}=\dim(\mathbb C^3\otimes\bar{\mathbb C}^3)^{\mathfrak{sl}(3)}=1;\ \mathrm{Def\text{-}H.42}];
\]

%dsp-level:0
\[
(4)\quad\textbf{the visibility locus:}\ \ \text{the same covariant data are nonzero on the adjoint channel}\ ([H,E_{12}]\ne0),
\]

%dsp-level:0
\[
\ \text{but the adjoint is not a sealed singlet and carries no pole}\ [\mathrm{FT.117}(4);\ \mathrm{CT3.18};\ \mathrm{CT3.22}]\ \text{— the constituent-level channel data}\ (\delta_q=\tfrac2{27}\ [\mathrm{FT.119}])\ \text{are ledger-internal at hadron level}.
\]

\textbf{Req.} D2.11, D2.14, D3.8, D3.9 (Part 2); CT3.17, CT3.18, CT3.22 (Part 2); Def-BM.1; FT.19, FT.39, FT.116(3), FT.117, FT.119, Def-H.9, Def-H.42 (Part 3).

\textbf{Pf.}

%dsp-level:0
\[
(1):\ \mathrm{Def\text{-}BM.1}\ \text{validity clause}:\ \pi_C\ \text{total}\implies\partial_{\mathrm{col}}(\Sigma;C)
\]

%dsp-level:0
\[
= \varnothing\ \text{on the composite's own support}\implies\mathcal C(\Sigma;C)
\]

%dsp-level:0
\[
= 1\ [\mathrm{D2.14}];
\]

%dsp-level:0
\[
\qquad \mathrm{FT.39}(1):\ \text{generic nonvanishing with the exact exception locus}
\]

%dsp-level:0
\[
\mathrm{FT.19}:\ \mathcal C=0\implies F^\Sigma_{\mathrm{strong}}=0\ (\text{one null factor kills the slot product});\qquad\mathrm{CT3.22}:\ \neg\text{triplet pole}\ \wedge\ \text{norm}\ 0
\]

%dsp-level:0
\[
(2)\text{–}(3):\ \mathrm{FT.117}(1)\text{–}(3),
\]

%dsp-level:0
\[
(5)\ \text{verbatim at the composite support}:\ \varepsilon\ \text{spans}\ \wedge^3\mathbb C^3\ (\text{the determinant representation});
\]

%dsp-level:0
\[
\ \text{the one-body sum acts by}\ \mathfrak T(H)\ [\mathrm{CT3.16};\ \mathrm{CT3.7}:\ \mathfrak T(T_a)=0];
\]

%dsp-level:0
\[
\ \text{meson}:\ (H\otimes I-I\otimes H^T)\,
\]

%dsp-level:0
\[
\mathrm{vec}(I)
\]

%dsp-level:0
\[
= 0\ [\text{machine rows re-cited}]
\]

%dsp-level:0
\[
(4):\ \mathrm{FT.117}(4)\ \text{verbatim}:\ \text{adjoint}\notin\operatorname{Im}\mathcal P_M^{(3)}\ [\mathrm{CT3.18}];
\]

%dsp-level:0
\[
\qquad \text{ledger-internality}: (2)\text{'s annihilation on sealed supports}.\ \blacksquare
\]

\textbf{Phys.} Sealing is the observability switch: the colour sector becomes physical exactly when the $\varepsilon/\delta$ grammar completes — and then reads a single number, the trace. Everything the constituent ledger carries (channel angles, doublet data, spectator dressings) stays behind the seal: the hadron sees sums, never slots. This is confinement as an information property, run forward from the subnetwork theorems to the composite level.

\end{theorem}

\begin{theorem}[{\texorpdfstring{BM.2 — Junction phase and winding grammar}{BM.2 - Junction phase and winding grammar}}]\label{tag:BM-2}
\textbf{Thm.}

%dsp-level:0
\[
(1)\quad\textbf{baryon-type sealing phase:}\ \ \text{at the computed baryon-type junction cell,}\ \Theta_{\mathrm{int}}\in4\mathbb Z
\]

%dsp-level:0
\[
\qquad [\text{Census FT.49};\ \text{per materialized closed chain}\ \Theta(c)\in2\mathbb Z,\ \mathrm{FT.22}];
\]

%dsp-level:0
\[
(2)\quad\textbf{additive winding under composition:}\ \ \Gamma
\]

%dsp-level:0
\[
= c_1\cdot P\cdot c_2\cdot\bar P\ \text{materialized with}\ d_{12}
\]

%dsp-level:0
\[
= \Theta(P)\qquad \ge 0:
\]

%dsp-level:0
\[
\quad k(\Gamma)=k_1+k_2+d_{12}
\]

%dsp-level:0
\[
\qquad [\mathrm{FT.34}(3),\ \text{Census FT.50}(2)];
\]

%dsp-level:0
\[
\text{carrier level:}\ \ \Psi(\Gamma)
\]

%dsp-level:0
\[
= \mathrm{Sym}^{k_1+k_2+d_{12}}\subseteq\mathrm{Sym}^{k_1}\otimes\mathrm{Sym}^{k_2}\otimes\mathrm{Sym}^{d_{12}}\ \text{as maximal weight}\ [\mathrm{FT.60}(6),\ \text{under Def-H.38}];
\]

%dsp-level:0
\[
(3)\quad\textbf{sealing partition:}\ \ \text{attachment partitions into contraction groups follow the}\ \mathrm{CT3.17a}\ \text{grammar with counts}\ g(N)
\]

%dsp-level:0
\[
= 1,
\]

%dsp-level:0
\[
1,
\]

%dsp-level:0
\[
3,
\]

%dsp-level:0
\[
10,
\]

%dsp-level:0
\[
25,
\]

%dsp-level:0
\[
105\ (N=2..7)\ [\mathrm{CT3.17a}(2);\ \mathrm{FT.41}(b)];
\]

%dsp-level:0
\[
\ \text{the gluing law of composed floor cells is the constraint-rank law}\ [\text{Census FT.50}];
\]

%dsp-level:0
\[
(4)\quad\textbf{junction lock content:}\ \ \text{a junction at}\ n
\]

%dsp-level:0
\[
\ge 4\ \text{carries}\ D^-\ \text{lock content — retained,}\ \neg\text{accessible}\ [\text{Census FT.35};\ \mathrm{CT4.1}\text{–}\mathrm{CT4.3};\ \mathrm{D4.2};\ \mathrm{D1.11}].
\]

\textbf{Req.} FT.22, FT.34, FT.60, Def-H.38 (Part 3); CT3.17a, CT4.1–CT4.3 (Part 2); D4.2, D1.11 (Parts 1–2); Def-BM.1; Census FT.35, FT.49, FT.50; exact integer arithmetic.

\textbf{Pf.}

%dsp-level:0
\[
(1):\ \text{Census FT.49's card row (computed from FT.22/FT.24's quantization on the tetrahedron-face sealing grammar)};
\]

%dsp-level:0
\[
\qquad \text{no new arithmetic}
\]

%dsp-level:0
\[
(2):\ \mathrm{FT.34}(3)\ \text{verbatim}:\ \Theta(\Gamma)
\]

%dsp-level:0
\[
= \Theta(c_1)+\Theta(P)+\Theta(c_2)+\Theta(\bar P)
\]

%dsp-level:0
\[
= 2(k_1+d_{12}+k_2)\ [\Theta(P)=\Theta(\bar P)=d_{12};\ \mathrm{FT.23}];
\]

%dsp-level:0
\[
\qquad \mathrm{FT.60}(6):\ \text{CG inclusion on the cone}
\]

%dsp-level:0
\[
(3):\ \mathrm{CT3.17a}(2)\ \text{counts re-cited};
\]

%dsp-level:0
\[
\ \text{Census FT.50's junction-saturation and constraint-rank rows}
\]

%dsp-level:0
\[
(4):\ \text{Census FT.35's junction row};\ \mathrm{CT4.1}\text{–}\mathrm{CT4.3}:\ \text{absorption/retention at}\ n\ge4.\ \blacksquare
\]

\textbf{Phys.} The hadron ledger is additive exactly where the fermion ledger is: windings compose, phase classes stack in $4\mathbb Z$ at baryon-type junctions, and the junction is where the first lock content sits — the baryon-type cell is the corpus's first $D^-$-carrying floor. What the junction does not add is new grammar: the sealing partitions are the frozen colour counts, and the gluing is the constraint-rank law the floor cells already proved.

\end{theorem}

\phantomsection
\section[Tier BM.2 — The Family Group and the Invariance Laws of the…]{\texorpdfstring{Tier BM.2 — The Family Group and the Invariance Laws of the Seal}{Tier BM.2 - The Family Group and the Invariance Laws of the Seal}}

\begin{theorem}[{\texorpdfstring{BM.3 — The family group of the seal}{BM.3 - The family group of the seal}}]\label{tag:BM-3}
\textbf{Thm.}

%dsp-level:0
\[
(1)\quad\textbf{the forced group:}\ \ \text{the stabilizer of the baryon seal}\ \varepsilon_{ijk}\ \text{inside the native slot-relabelling group}\ S_3\ \text{is}\ A_3\cong C_3
\]

%dsp-level:0
\[
= S_3\cap SU(3)\ \text{— the unique maximal}\ \varepsilon\text{-preserving subgroup};
\]

%dsp-level:0
\[
\ \text{full}\ S_3\ \text{covariance is two-level-only}\ (\text{every}\ S_3\text{-covariant slot operator has spectrum}\ \{a+3b,a,a\})\ [\mathrm{FT.116}];
\]

%dsp-level:0
\[
(2)\quad\textbf{channel inheritance:}\ \ \text{every triplet-attached constituent carries the democratic}\ C_3\ \text{mechanism with angle}\ \delta_q
\]

%dsp-level:0
\[
= \tfrac2{27}\ [\mathrm{FT.119};\ \mathrm{Def\text{-}H.42}];
\]

%dsp-level:0
\[
\ \text{the triple-angle link}\ 3\delta_q
\]

%dsp-level:0
\[
= \delta_\ell
\]

%dsp-level:0
\[
= \tfrac29\ [\mathrm{FT.119}(4);\ \mathrm{FT.112}];
\]

%dsp-level:0
\[
\ \text{the cross-channel numerator identity}\ 16\delta_\nu
\]

%dsp-level:0
\[
= 9\delta_\ell
\]

%dsp-level:0
\[
= 27\delta_q
\]

%dsp-level:0
\[
= 2\ [\mathrm{FT.143}];
\]

%dsp-level:0
\[
(3)\quad\textbf{exclusions:}\ \ \text{the light-triplet read}\ (A,B)
\]

%dsp-level:0
\[
= (1,\sqrt2)\ \text{is excluded —}\ Q\ \text{is confined to}\ (0.54,0.59)\ \text{across the allowed ratio box while the channel normalization forces}\ Q
\]

%dsp-level:0
\[
= \tfrac23\ \text{phase-independently}\ [\mathrm{FT.120}(6);\ \text{Census row, quarantined external data}];
\]

%dsp-level:0
\[
\ \text{the heavy-triplet proximity rows are audit-only, no identification}\ [\text{Census FT.120};\ \mathrm{P1},\ \mathrm{P5}].
\]

\textbf{Req.} CT3.17 (Part 2); FT.112, FT.116, FT.119, FT.143, Def-H.42 (Part 3); Census FT.120; exact $\mathbb Q(i)$/$\mathbb Q(\sqrt3)$ arithmetic (rows re-cited).

\textbf{Pf.} $(1)$: FT.116 verbatim — the relabelling lattice, the two-level obstruction, the $\varepsilon$-selection [machine rows re-cited]. $(2)$: FT.119(4)/FT.112(3)/FT.143 angle ledgers verbatim; per-constituent application by Def-BM.1's constituent structure (each constituent is a triplet-attached bound chain in FT.119's sense). $(3)$: Census FT.120's exclusion rows (quarantine discipline [FT.73(5) pattern]). $\blacksquare$

\textbf{Phys.} The seal picks its own symmetry: the only slot-relabelling subgroup preserving the baryon contraction is the cyclic three — the same three that carries the channel angles on every constituent. Hadron-level family structure is not imposed on the junction; it is the junction's stabilizer, with its phase fixed to two-thirds and its alternatives excluded by theorem.

\end{theorem}

\begin{theorem}[{\texorpdfstring{BM.4 — The invariance laws of the seal}{BM.4 - The invariance laws of the seal}}]\label{tag:BM-4}
\textbf{Thm.}

%dsp-level:0
\[
(1)\quad\textbf{colour invariance:}\ \ B,
\]

%dsp-level:0
\[
\ M\ \text{are colour singlets}\ [\mathrm{CT3.15}\text{–}\mathrm{CT3.17}:\ \{B,M\}\ \text{singlets};\ \{\delta,\varepsilon,\bar\varepsilon\}\ \text{the invariant grammar}];
\]

%dsp-level:0
\[
\ \text{the mask admits only singlets}\ [\mathrm{CT3.18}];
\]

%dsp-level:0
\[
\ \text{the colour sector is rigid and contributes}\ \neg\mathrm{P4}\ \text{data}\ [\mathrm{CT3.19}];
\]

%dsp-level:0
\[
(2)\quad\textbf{transport compatibility:}\ \ \text{the seed is colour-neutral by theorem}\ [\mathrm{EW.32}:\ T_a\mathsf v=0];
\]

%dsp-level:0
\[
\ \text{the weak action on the extended carrier is seed-factor-only}\ [W_i\mapsto\tfrac12\sigma_i\otimes I_\omega,\ \mathrm{FT.123}(1);\ \text{Def-H.34}]\ \text{and the diagonal}\ \mathrm{Sym}^k\ \text{transport is colour-blind}\ [\mathrm{FT.60};\ \mathrm{Def\text{-}H.38}]\ \implies\ \text{constituent channel data and multiplicity coupling compose with the seal without colour obstruction};
\]

%dsp-level:0
\[
(3)\quad\textbf{the parity side:}\ \ \text{the scalar-CPT action on each constituent's orientation doublet,}\ \mathrm{CPT}|_{\mathcal H_\omega}
\]

%dsp-level:0
\[
= (-1)^{|E|}e^{i\pi\Theta}\,
\]

%dsp-level:0
\[
\mathrm{id},
\]

%dsp-level:0
\[
\ \text{is chirality-even}\ [\mathrm{FT.134}(3)];
\]

%dsp-level:0
\[
\ \text{pole multiplets are}\ m\leftrightarrow-m\ \text{symmetric and the ceilings parity-blind}\ [\mathrm{FT.138}]\ \implies\ \text{both laws pass to sealed composites constituent-wise}\ [\mathrm{Def\text{-}BM.1}\ \text{product structure}].
\]

\textbf{Req.} CT3.15–CT3.19 (Part 2); EW.32 (Part 2); FT.60, FT.123, FT.134, FT.138, Def-H.34, Def-H.38 (Part 3); Def-BM.1.

\textbf{Pf.}

%dsp-level:0
\[
(1):\ \mathrm{CT3.15}\text{–}\mathrm{CT3.17}\ \text{singlet rows re-cited};
\]

%dsp-level:0
\[
\ \mathrm{CT3.18}:\ \operatorname{Im}\mathcal P_M^{(3)}\ \text{admits only singlets};
\]

%dsp-level:0
\[
\ \mathrm{CT3.19}:\ \text{open quantities}\qquad = \varnothing
\]

%dsp-level:0
\[
(2):\ \mathrm{EW.32}:\ T_a\mathsf v=0\ [\text{the seed's colour neutrality}];
\]

%dsp-level:0
\[
\qquad \mathrm{FT.123}(1):\ [W_i\otimes I_\omega,\ \cdot]\ \text{acts on the seed/orientation factors only};
\]

%dsp-level:0
\[
\qquad \mathrm{Def\text{-}BM.1}:\ \text{constituents are independent chains with own carriers}\implies\text{product composition}
\]

%dsp-level:0
\[
(3):\ \mathrm{FT.134}(3),
\]

%dsp-level:0
\[
\ \mathrm{FT.138}\ \text{verbatim per constituent};
\]

%dsp-level:0
\[
\ \text{composite-level passage}: \mathrm{Def\text{-}BM.1}'\text{s product structure}\ (\text{each factor carries its own doublet and}\ \Theta\text{-summand};\ (-1)^{|E|}e^{i\pi\Theta}\ \text{multiplicative over disjoint constituents}).\ \blacksquare
\]

\textbf{Phys.} Both algebras that act on a composite leave the seal intact: colour SU(3) has nothing to do inside a singlet, and weak transport never touches the colour factor — so a baryon's card assembles constituent by constituent, and the parity laws pass through the seal unbroken. The rigidity theorem closes the ledger: the colour side of the hadron grammar contributes no free datum at all.

\end{theorem}

\phantomsection
\section[Tier BM.3 — Consumption and Open Constructions]{\texorpdfstring{Tier BM.3 — Consumption and Open Constructions}{Tier BM.3 - Consumption and Open Constructions}}

\pseudohead{theorem}{\texorpdfstring{BM.5}{BM.5}}{\texorpdfstring{The census-consumption law}{The census-consumption law}}\label{tag:BM-5}
\textbf{Thm.}

%dsp-level:0
\[
(1)\quad\textbf{the substrate:}\ \ \text{the enumerative substrate of this era is the Census block — the}\ 25\ \text{Category-B cells with their configuration spaces, counts, and consumption rows}\ [\text{Census charter and §Consumption}];
\]

%dsp-level:0
\[
\ \text{every construction of this part consumes cells by tag citation; no cell is peripheral};
\]

%dsp-level:0
\[
(2)\quad\textbf{the family table:}
\]

%tbl-level:3
{\scriptsize\setlength{\tabcolsep}{4pt}
\begin{longtable}{>{\raggedright\arraybackslash}p{\dimexpr0.1200\textwidth-2\tabcolsep\relax}|>{\raggedright\arraybackslash}p{\dimexpr0.1127\textwidth-2\tabcolsep\relax}|>{\raggedright\arraybackslash}p{\dimexpr0.2302\textwidth-2\tabcolsep\relax}|>{\raggedright\arraybackslash}p{\dimexpr0.5371\textwidth-2\tabcolsep\relax}}
\toprule
census family & cells & configuration space supplied & consumption at the composite level \\
\addlinespace[2pt]
\midrule
\endfirsthead
\toprule
census family & cells & configuration space supplied & consumption at the composite level \\
\addlinespace[2pt]
\midrule
\endhead
\bottomrule
\endlastfoot
floor cards and ladders & FT.28, FT.32, FT.35, FT.42–FT.45, FT.49–FT.51, FT.63 & the computed fermion floors with full read columns & the cell grammar of sealed composites (Def-BM.2): baryon-type junctions are gluings of these cards; the winding ledgers and duality splits are the multiplet bookkeeping; the $\Sigma_{12}$/tetrahedron-face cells carry the first colour-active and $4\mathbb Z$-phased junctions \\
loop and cover classes & FT.64, FT.65, FT.69, FT.78, FT.96 & admissible-loop lattices, coarse-grained degeneracies, amplitude censuses, $\operatorname{Aut}$-orbits with chirality tags & the enumeration substrate of sealed-object path integrals: sector conditioning $(q,\allowbreak S)$ [Def-H.32(4)] runs over these pools; the chirality-tagged orbits feed the parity reads [BM.4(3)] \\
pole and mass multiplets & FT.76 & defect-angle pole reads under both convention variants, with the ceiling & the pole-side mass cells and the hard range ceiling every composite read inherits [Def-H.39(2)]; the convention variants are priced — the choice converts them into data [BM.6(4)] \\
form-factor tables & FT.98, FT.99 & polygon and Platonic form factors with growth and dominance rows & the size-bound range checks for extended composite skeletons [Def-H.39 ceilings] \\
audit ledgers & FT.81, FT.83, FT.84, FT.114, FT.120 & the read-class inventory, the rung-emergence ledger, the ceiling taxonomy, the quarantined comparison rows & the card grammar and the P1/P5-compliant comparison template; the activation map at $n=\allowbreak 3$/$n\ge\allowbreak 4$; the read-species taxonomy any new read must join \\
topology-blind cells & FT.105 & graph counts and canonical skeletons from first principles, budget-regularized & the regulated topology pool of the composite path integral — the space its fugacity formalism sums over [Def-H.32(1)] \\
\end{longtable}
}

%dsp-level:0
\[
(3)\quad\textbf{the conditioning law:}\ \ \text{sector-conditioned composite sums run on Def-H.32(4)'s labels}\ (q,S)\ \text{over}\ (2)\text{'s pools};
\]

%dsp-level:0
\[
\ \text{the amplitude/norm dichotomy and the amplitude}\leftrightarrow\text{spectral bridge residue are item-13 content}\ [\mathrm{FT.67}\text{–}\mathrm{FT.68};\ \mathrm{FT.80}].
\]

\textbf{Req.} The Census block (all 25 cells); Def-H.32, FT.67, FT.68, FT.80 (Part 3); Def-H.39 (Part 3); Def-BM.2

\textbf{Pf.} $(2)$: the consumption map — each family row references its cells by tag; every Census tag occurs in exactly one family row [exact set arithmetic: $11+5+1+2+5+1=25$]. $(1),(3)$: charter and citation scans. $\blacksquare$

\textbf{Phys.} The census is not an appendix: it is the table of configuration spaces the baryomesonic era enumerates — cards to glue, orbits to condition on, multiplets whose ceilings every composite read inherits. The consumption law states, family by family, which space feeds which construction, and the count closes: twenty-five cells, no residue.

\begin{theorem}[{\texorpdfstring{BM.6 — The construction interfaces (specifications)}{BM.6 - The construction interfaces (specifications)}}]\label{tag:BM-6}
\textbf{Def.} The constructions this era requires and does not yet carry, each stated as a specification (what a construction must supply); none is a declaration, and none is adopted:

%dsp-level:0
\[
(1)\quad\textbf{the multi-chain additive charge read}:\ \text{a read}\ Q\ \text{on sealed composites with}\ (i)\ Q|_{\text{single chain}}
\]

%dsp-level:0
\[
= q_{\mathrm{core}}\ [\mathrm{FT.40};\ \mathrm{EM.9}],
\]

%dsp-level:0
\[
\ (ii)\ \text{additivity under disjoint sealing},
\]

%dsp-level:0
\[
\ (iii)\ \text{observability only on sealed supports}\ [\mathrm{CT3.22}\ \text{discipline}].\ \text{Status: not constructed; the single-chain read is derived; two end registers per chain are counted}\ [\mathrm{FT.41}(a)];
\]

%dsp-level:0
\[
(2)\quad\textbf{the two-chain exchange space}:\ \text{a state space on ordered pairs of bound chains with an exchange involution; the single-chain chiral grade is}\ \neg\text{an exchange sign}\ [\mathrm{FT.125}(5);\ \mathrm{FT.80}(5):\ \neg\exists\ \text{exchange operator in the inventory};\ \mathrm{FT.89}(2)\ \text{kind discipline}];
\]

%dsp-level:0
\[
\ \text{spin–statistics content is external until constructed};
\]

%dsp-level:0
\[
(3)\quad\textbf{the}\ \aleph\text{-domain choice}:\ \text{literal}\ [\mathrm{D3.3}]\ \text{vs chain-local}\ \mathrm{E1}/\mathrm{E2}\ [\mathrm{Def\text{-}H.27}]\ \text{— both path integrals computed and pairwise inequivalent}\ [\mathrm{FT.67},\ \mathrm{FT.74},\ \mathrm{FT.75}];
\]

%dsp-level:0
\[
\ \text{sealed junctions are the natural site of the residual class-(iii) edges}\;
\]

%dsp-level:0
\[
\ \text{the choice requires a logged declaration before use};
\]

%dsp-level:0
\[
(4)\quad\textbf{the defect corner convention}:\ \text{variant A (face winding}\ k_f)\ \text{/ variant B (vertex-figure winding}\ k_{vf})\ \text{— both multiplet tables priced}\ [\text{Census FT.76}(1)\text{–}(2)];
\]

%dsp-level:0
\[
\ \text{a declaration converts them into data immediately};
\]

%dsp-level:0
\[
(5)\quad\textbf{the dynamics rates}\ [\text{Cal-F}\ r]:\ \text{a dynamics construction supplying}\ J^\pm\ \text{rates, or the free datum retained; the tick-rate/field coupling requires emergent time}\ [\mathrm{FT.131}(3)];
\]

%dsp-level:0
\[
(6)\quad\textbf{the momentum-space}\ \Sigma\text{-resolvent}:\ \text{plane-wave boundary-channel states on the coarse-grained channel}\ [\mathrm{Def\text{-}H.13}\ \text{pattern}];
\]

%dsp-level:0
\[
\ \text{constructed at BM.134 (Def-BM.20); consumed by scattering reads of sealed objects};
\]

%dsp-level:0
\[
(7)\quad\textbf{the embedding datum}:\ \text{a circumradius or moment-of-inertia read for rotor-type reads; the ratios are embedding-free and computed}\ [\mathrm{FT.87}(6)];
\]

%dsp-level:0
\[
\ \text{adoption would not lift the}\ 4\hat\mu'^2\ \text{ceiling}\ [\mathrm{FT.76}(3);\ \mathrm{FT.87}(1)];
\]

%dsp-level:0
\[
(8)\quad\textbf{the frame section}:\ e\mapsto\mathcal F_e\ \text{over each skeleton}\ [30\ \text{native frames catalogued},\ \mathrm{FT.93}]\ \text{plus its}\ \mathrm{EW.32}\text{-transport coupling and a global Dirac-type operator on the extended carrier}\ [\mathrm{Def\text{-}H.34}];
\]

%dsp-level:0
\[
\ \text{the weak-chirality pincer on the pair arena is untouched}\ [\mathrm{FT.93}(3)].
\]

\textbf{Req.} The cited tags

\textbf{Pf.} Specification content: each row states its requirement and cites its anchors; no truth claim is made; anchor existence machine-checked [\texttt{checks/\allowbreak{}verify\_\allowbreak{}era\_\allowbreak{}structure.py}]. $\blacksquare$

\textbf{Phys.} The era's debts are itemized as build-specs: what a charge read must satisfy, what an exchange space must carry, which two conventions must be chosen and what each choice costs. Nothing here is adopted; everything here is priced. The rung closes with its obligations named.

\end{theorem}

\phantomsection
\section[Tier BM.4 — Noether's Theorem from the Janus Seal Topology:…]{\texorpdfstring{Tier BM.4 — Noether's Theorem from the Janus Seal Topology: the Discrete Conservation Law, the Continuum Promotion, the Stress-Energy of a Sealed Composite, and the Canonical-Pair Obstructions}{Tier BM.4 - Noether's Theorem from the Janus Seal Topology: the Discrete Conservation Law, the Continuum Promotion, the Stress-Energy of a Sealed Composite, and the Canonical-Pair Obstructions}}

\begin{definition}[{\texorpdfstring{Def-BM.3 — The species-identification layer (declared)}{Def-BM.3 - The species-identification layer (declared)}}]\label{tag:Def-BM-3}
\textbf{Def.} A \textbf{species name} $:=$ a label attached to a derived structural tuple of a sealed composite. The layer is constitutive in exactly one respect — the name$\leftrightarrow$tuple assignment — and non-derived in that same respect: no theorem of Parts 1–4 forces any name onto any tuple, and no name carries a datum into a derivation.

%dsp-level:0
\[
(1)\quad\textbf{the tuple a name attaches to:}
\]

%dsp-level:0
\[
\qquad \bigl(\text{junction grammar}\ [\mathrm{Def\text{-}BM.1}],\ \ \text{sealing-phase class}\ [\mathrm{BM.2}(1)],\ \ \text{charge class}\ [\mathrm{EM.9};\ \text{clause}\ (3)],\ \ m_{\mathrm{hid}}=E_{\mathrm{int}}\ln2/\mathfrak h\ [\mathrm{FT.44}(4)],\ \ w=p/\rho\ [\mathrm{CT4.5};\ \mathrm{BM.10}(3)]\bigr);
\]

%dsp-level:0
\[
(2)\quad\textbf{the declared map (baryon-type, }\varepsilon\text{-sealed 3-junctures):}
\]

%dsp-level:0
\[
\qquad \textbf{neutron}\ :=\ \varepsilon_{ijk}c_1^ic_2^jc_3^k\ \text{in the neutral charge class}\ (q_{\mathrm{net}}=0);
\]

%dsp-level:0
\[
\textbf{proton}\ :=\ \varepsilon_{ijk}c_1^ic_2^jc_3^k\ \text{in the unit charge class}\ (q_{\mathrm{net}}=+1);
\]

%dsp-level:0
\[
\qquad \textbf{antiproton}\ :=\ \omega\text{-image of the proton}\ [\mathrm{D2.6}];
\]

%dsp-level:0
\[
(3)\quad\textbf{the declared map (meson-type, }\delta\text{-sealed 2-junctures):}
\]

%dsp-level:0
\[
\qquad \pi^0\ :=\ \delta^i{}_j\,c^j\bar c_i\ \text{in the neutral class};
\]

%dsp-level:0
\[
\qquad \pi^{\pm}\ :=\ \delta^i{}_j\,c^j\bar c_i\ \text{in the}\ q_{\mathrm{net}}=\pm1\ \text{classes};
\]

%dsp-level:0
\[
(4)\quad\textbf{the charge-class clause (conditional, declared):}
\]

%dsp-level:0
\[
\qquad q_{\mathrm{net}}\ \text{is the additive multi-chain read specified but}\ \neg\text{constructed}\ [\mathrm{BM.6}(1)]:
\]

%dsp-level:0
\[
\text{the class labels}\ 0,
\]

%dsp-level:0
\[
\ \pm1\ \text{of clauses}\ (2)\text{–}(3)\ \text{are the ternary alphabet of a single end read}\ [q_{\mathrm{EM}}\in\{-1,0,+1\}\ [\mathrm{D3.4},\ \mathrm{EM.9},\ \mathrm{P2}]]\ \text{carried to the composite by declaration};
\]

%dsp-level:0
\[
\text{the constituent alphabet remains ternary: the fractional read}\ (\pm\tfrac13,\ \pm\tfrac23)\ \text{is deficit}\ 1'\ \text{content}\ [\mathrm{FT.125}(4)]\ —\ \neg\text{supplied,}\ \neg\text{assumed here};
\]

%dsp-level:0
\[
(5)\quad\textbf{the naming-precedent clause:}
\]

%dsp-level:0
\[
\qquad \mathrm{D3.7}\ \text{names two transit structures (photon}\ :=\ \text{trivial chain}\ \subseteq V_2;\ \text{gluon}\ :=\ \text{trivial chain}\ \cap\,V_3\ne\varnothing)\ \text{and}\ \mathrm{D3.9}\ \text{names the two sealings (baryon, meson)}:
\]

%dsp-level:0
\[
\text{this definition extends that precedent from transit structures and sealings to charge classes of sealed composites};
\]

%dsp-level:0
\[
\qquad \text{a name is a definition-side label,}\ \neg\text{a measurement};
\]

%dsp-level:0
\[
(6)\quad\textbf{the data bar (standing):}
\]

%dsp-level:0
\[
\qquad \neg\exists\ \text{a measured external value in any clause of Parts 1–4}\ [\mathrm{P5}\ \text{stands}];
\]

%dsp-level:0
\[
\qquad \text{external comparison data remain quarantined to the audit rows}\ [\text{Census FT.114},\ \text{FT.120}];
\]

%dsp-level:0
\[
\text{the names of clauses}\ (2)\text{–}(3)\ \text{carry}\ \neg\text{mass value,}\ \neg\text{coupling value,}\ \neg\text{lifetime,}\ \neg\text{radius}:\ \text{every number attached to a named tuple is corpus-derived}\ (m_{\mathrm{hid}},\ \Theta_{\mathrm{int}},\ w,\ E_{\mathrm{int}});
\]

%dsp-level:0
\[
(7)\quad\textbf{scope:}
\]

%dsp-level:0
\[
\qquad \text{species names occur in Part 4 only}\ —\ \text{in this definition, in the}\ \textbf{Phys}\ \text{and}\ \textbf{Rmk}\ \text{prose of this part, and nowhere in a}\ \textbf{Thm}\ \text{clause of Parts 1–3}\ [\text{those tiers' own label-discipline rows stand unmodified}];
\]

%dsp-level:0
\[
\text{extension of the map beyond clauses}\ (2)\text{–}(3)\ \text{requires a logged declaration first}\ [\text{Part 4 §Cal-BM rule}];
\]

%dsp-level:0
\[
\qquad \text{the layer adds}\ \neg\text{dimensionful scale}\ [\mathrm{P1}]\ \text{and}\ \neg\text{P4 parameter}.
\]

\textbf{Rmk (identification discipline).} The layer answers \textit{which name sits on which derived tuple}; it does not answer \textit{which species sits in which gauge representation on which chiral ray} — that assignment rule is application-layer content; its non-derivability here is proven [FT.125(4)]. Naming the neutron is a labelling act on the $\varepsilon$-sealed neutral class; it derives no neutron property, imports no neutron datum, and closes no open row.

\textbf{Req.} D2.6, D3.4, D3.7, D3.9 (Parts 1–2); CT4.5 (Part 2); EM.9 (Part 2); Def-BM.1, BM.2(1), BM.6(1) (this part); FT.44(4), FT.125(4) (Part 3, Census); P1, P2, P5 (Parts 1–2); Census FT.114, FT.120 (quarantine rows)

\textbf{Pf.}

%dsp-level:0
\[
\text{Definition block: the constitutive clause is the name}\leftrightarrow\text{tuple assignment of}\ (2)\text{–}(3);
\]

%dsp-level:0
\[
\ \text{the non-derived clauses are}\ (4)\ \text{(the charge class is declared pending BM.6(1)) and the name}\leftrightarrow\text{tuple direction of}\ (1)\ [\text{no theorem forces it}];
\]

%dsp-level:0
\[
\ \text{the standing bars}\ (6)\text{–}(7)\ \text{are scope clauses cited at use}.\ \blacksquare
\]

\textbf{Phys.} The baryomesonic era names its objects. A name is pinned to a tuple the corpus already derives — the junction grammar, the sealing phase class, the charge class, the informational mass, the equation of state — and the pinning itself is the declared part. The neutron is the $\varepsilon$-sealed three-juncture in the neutral charge class; the proton is the same structure in the unit class; the pion-type triplet is the $\delta$-sealed pair in its three classes. Nothing measured enters: the mass read is $E_{\mathrm{int}}\ln2/\mathfrak h$, the phase class is $4\mathbb Z$, the equation of state is $w=0$, and the charge class is the ternary alphabet carried up by declaration while the additive read is still an open construction. The layer extends a precedent the corpus already set when it named the photon and the gluon as chain classes.

\end{definition}

\begin{theorem}[{\texorpdfstring{BM.7 — The discrete conservation law is Janus topology: no symmetry argument, no action, no limit}{BM.7 - The discrete conservation law is Janus topology: no symmetry argument, no action, no limit}}]\label{tag:BM-7}
\textbf{Thm.}

%dsp-level:0
\[
(1)\quad\textbf{the balance is definitional, and the seal ledger realizes it uniquely:}
\]

%dsp-level:0
\[
\qquad \text{valid universe}\iff\text{network closed}\ \wedge\ \#J^+
\]

%dsp-level:0
\[
= \#J^-\ \wedge\ \in\operatorname{Im}(\text{active mask})\ \wedge\ \neg\exists\,
\]

%dsp-level:0
\[
\text{dangling boundary}\ [\mathrm{D2.11},\ \mathrm{D1.3}];
\]

%dsp-level:0
\[
J^+:\ (\Delta n_{\mathcal R},\delta h,\delta\tau)=(+1,+1,+1),
\]

%dsp-level:0
\[
\qquad J^-:\ (-1,-1,-1)
\]

%dsp-level:0
\[
\qquad [\mathrm{CT0.3}]:
\]

%dsp-level:0
\[
\qquad \text{the seal excess}\ \#J^+-\#J^-=n_{\mathcal R}\ \text{is the ledger's own occupation count};
\]

%dsp-level:0
\[
\text{the D1.2 transition domains force strict alternation}\ (J^-\ \text{undefined at}\ n_{\mathcal R}=0;\ J^+\ \text{undefined at}\ n_{\mathcal R}=1):\ \text{the single-pair history of length}\ n\ \text{is unique}\ \forall n\ [\mathrm{FT.144}(2);\ \text{machine}];
\]

%dsp-level:0
\[
\text{the completed history of length}\ 2m\ \text{is unique and balanced:}\quad(J^+J^-)^m,
\]

%dsp-level:0
\[
\qquad \#J^+=\#J^-=m
\]

%dsp-level:0
\[
\qquad [\mathrm{FT.144}(2),\ \mathrm{FT.156}(5);\ \text{machine}];
\]

%dsp-level:0
\[
\text{the prefix law}\ \#J^+(\text{prefix})\ge\#J^-(\text{prefix})\ \text{holds on every admissible history, with equality}
\]

%dsp-level:0
\[
\iff n_{\mathcal R}
\]

%dsp-level:0
\[
= 0\ [\mathrm{FT.144}(2);\ \text{machine}];
\]

%dsp-level:0
\[
\qquad \text{the pure order reversal of every nonempty completed history leads with}\ J^-:\ \text{inadmissible at its first event}\ [\mathrm{FT.144}(4)(\mathrm{ii});\ \text{machine}];
\]

%dsp-level:0
\[
\textbf{the relaxation census}\ (\text{prefix law alone, alternation dropped}):\quad\#\{\text{completed words of length}\ 2m\}
\]

%dsp-level:0
\[
= \mathrm{Cat}_m
\]

%dsp-level:0
\[
= \frac1{m+1}\binom{2m}m
\]

%dsp-level:0
\[
= 1,
\]

%dsp-level:0
\[
1,
\]

%dsp-level:0
\[
2,
\]

%dsp-level:0
\[
5,
\]

%dsp-level:0
\[
14,\qquad 42
\]

%dsp-level:0
\[
\qquad [\text{machine, exhaustive}\ \le12];
\]

%dsp-level:0
\[
\mathrm{Cat}_m>1
\]

%dsp-level:0
\[
\iff m
\]

%dsp-level:0
\[
\ge 2:\ \text{the D1.2 domains are strictly stronger than the prefix law from}\ 2m
\]

%dsp-level:0
\[
= 4\ \text{on}\ —\ \text{the uniqueness of}\ (1)\ \text{is domain content,}\ \neg\text{balance content};
\]

%dsp-level:0
\[
\textbf{two pairs:}\quad\#\{\text{completed interleavings of total length}\ 2n\}
\]

%dsp-level:0
\[
= \sum_{m_0+m_1=n}\frac{(2n)!}{(2m_0)!\,(2m_1)!}
\]

%dsp-level:0
\[
= 2,
\]

%dsp-level:0
\[
\ 8,
\]

%dsp-level:0
\[
\ 32\quad(2n=2,4,6)
\]

%dsp-level:0
\[
\qquad [\mathrm{FT.144}(2);\ \text{machine: exhaustive}];
\]

%dsp-level:0
\[
\text{balance and the prefix law survive every interleaving}\ [\text{machine}];
\]

%dsp-level:0
\[
\qquad \text{the balanced sector is}\ a_+-a_-=|\mathcal R'|-|\mathcal R|=0\ \text{of the budget identity}\ [\mathrm{FT.108}(4)];
\]

%dsp-level:0
\[
(2)\quad\textbf{the endpoint read is a boundary term, exactly:}
\]

%dsp-level:0
\[
\qquad q_{\mathrm{EM}}(c)=\frac{\sigma_{s(c)}-\sigma_{t(c)}}2\in\{-1,0,+1\}
\]

%dsp-level:0
\[
\qquad [\mathrm{D3.4},\ \mathrm{EM.9};\ \text{all three values attained}];
\]

%dsp-level:0
\[
\textbf{the telescoping identity:}\quad\text{on a materialized chain}\ c=(e_1,\dots,e_k),\ \ t(e_i)=s(e_{i+1})\ [\mathrm{D1.4}]:
\]

%dsp-level:0
\[
\qquad q_{\mathrm{EM}}(c)=\sum_{i=1}^{k}j(e_i),
\]

%dsp-level:0
\[
\qquad j(e):=\frac{\sigma_{s(e)}-\sigma_{t(e)}}2
\]

%dsp-level:0
\[
\qquad [\text{machine: exhaustive}\ k\le7];
\]

%dsp-level:0
\[
\text{interior vertex signs cancel pairwise; the sum is the difference of the two end reads}\ —\ \text{the charge at the start equals the charge at the end};
\]

%dsp-level:0
\[
\textbf{the closed-chain read:}\quad\partial c=0\implies s(c)=t(c)\implies\sigma_{s(c)}=\sigma_{t(c)}\implies q_{\mathrm{EM}}(c)=0\ \text{exactly}
\]

%dsp-level:0
\[
\qquad [\mathrm{D1.4};\ \text{machine: every closed configuration}];
\]

%dsp-level:0
\[
\text{total charge vanishes in every valid universe: each end is sealed against an oppositely oriented end carrying the same sign}\ [\mathrm{EM.11}];
\]

%dsp-level:0
\[
(3)\quad\textbf{the ground is the chain complex, not an extremum:}
\]

%dsp-level:0
\[
\qquad \partial c:=t(c)-s(c)\ \text{(formal difference of vertices)},
\]

%dsp-level:0
\[
\qquad \partial^2=0\ \text{since}\ \partial(\text{vertex})=0
\]

%dsp-level:0
\[
\qquad [\mathrm{D1.4};\ \text{machine}];
\]

%dsp-level:0
\[
\text{closed}\iff\partial c=0\ \wedge\ \#J^+=\#J^-\ [\mathrm{D1.4}];
\]

%dsp-level:0
\[
\qquad \ker\partial=\{\text{closed chains}\},
\]

%dsp-level:0
\[
\qquad \text{every boundary has total multiplicity}\ 0\ [\text{machine}];
\]

%dsp-level:0
\[
\text{conservation is the kernel statement of an integer boundary map: no action functional, no variational derivative and no continuous parameter enters}\ (1)\text{–}(3);
\]

%dsp-level:0
\[
(4)\quad\textbf{the Noether route is unavailable for this read (obstruction):}
\]

%dsp-level:0
\[
\qquad \text{a Noether charge requires a one-parameter group action}\ \mathbb R\to\operatorname{Aut}(\text{the read's value space})\ [\text{math}];
\]

%dsp-level:0
\[
\text{the value space of}\ q_{\mathrm{EM}}\ \text{is the ternary set}\ \{-1,0,+1\},
\]

%dsp-level:0
\[
\ \text{so the action lands in}\ \operatorname{Perm}\{-1,0,+1\}
\]

%dsp-level:0
\[
= S_3\ [\mathrm{EM.9}];
\]

%dsp-level:0
\[
U(1)
\]

%dsp-level:0
\[
= \varinjlim_m\mathbb Z/m
\]

%dsp-level:0
\[
= \mathbb Q/\mathbb Z\ \text{is divisible}\ [\text{math}]:\ \text{its image in a finite group is divisible, hence trivial}\ —\ \text{no transposition of}\ S_3\ \text{has a square root, no}\ 3\text{-cycle has a cube root}\ [\text{machine: the full root census of}\ S_3];
\]

%dsp-level:0
\[
|\operatorname{hom}(\mathbb Z/m,S_3)|
\]

%dsp-level:0
\[
= \#\{g\in S_3:\operatorname{ord}(g)\mid m\}
\]

%dsp-level:0
\[
= 1,
\]

%dsp-level:0
\[
4,
\]

%dsp-level:0
\[
3,
\]

%dsp-level:0
\[
4,
\]

%dsp-level:0
\[
6,
\]

%dsp-level:0
\[
6\ (m=1,2,3,4,6,12)\ [\text{machine}]:\ \text{every finite stage admits nontrivial homs, and the only projectively compatible family is the trivial one};
\]

%dsp-level:0
\[
\therefore\ \neg\exists\,\text{a nontrivial continuous}\ U(1)\ \text{action on the charge alphabet}:\qquad q_{\mathrm{EM}}\ \text{is}\ \neg\text{a Noether charge};
\]

%dsp-level:0
\[
(5)\quad\textbf{priority:}
\]

%dsp-level:0
\[
\qquad (1)\text{–}(3)\ \text{are exact, discrete, integer-valued and hold at every scale, including}\ \mathfrak h\ne0;
\]

%dsp-level:0
\[
\qquad \text{the differential law}\ \partial_\mu J^\mu=0\ \text{is their continuum limit}\ [\mathrm{BM.11}]\ —\ \text{derived from them,}\ \neg\text{converse};
\]

%dsp-level:0
\[
\text{the converse fails:}\ \partial_\mu J^\mu
\]

%dsp-level:0
\[
= 0\ \text{is a density statement and does not force the integer balance at finite cell size}\ [\mathrm{BM.11}(5)].
\]

\textbf{Rmk (label discipline).} $J^{\pm}$, \textit{seal}, \textit{source}, \textit{sink} and \textit{current} name ledger structures [D1.2, D1.3, CT0.3]; the topology labels are internal bookkeeping [Def-H.32(4) pattern]. Where clause (2) is read on a named composite, the name is Def-BM.3's label on the derived tuple; the theorem's content is the tuple.

\textbf{Req.} A4, A6, A7 (Part 1); D1.2, D1.3, D1.4, D2.11, D3.3, D3.4 (Parts 1–2); CT0.3, CT0.9 (Part 1); EM.9, EM.11 (Part 2); P2 (Part 1); Def-BM.1, Def-BM.3 (this part); FT.108(4), FT.144(2), FT.144(4)(ii), FT.156(5) (Part 3); exhaustive ledger-word enumeration, the Catalan census, multinomial interleaving counts, hom-counts and the root census of $S_3$, divisibility of $\mathbb Q/\mathbb Z$ [math; machine: \texttt{checks/\allowbreak{}verify\_\allowbreak{}H\_\allowbreak{}noether.py} §A].

\textbf{Pf.}

%dsp-level:0
\[
(1):\ \mathrm{D2.11}/\mathrm{D1.3}\ \text{verbatim};
\]

%dsp-level:0
\[
\qquad \mathrm{CT0.3}\ \text{increments verbatim}\implies\#J^+-\#J^-=n_{\mathcal R}\ [\text{exact integers}];
\]

%dsp-level:0
\[
\qquad \text{alternation, uniqueness, the prefix law, the reversal inadmissibility}: \mathrm{FT.144}(2),
\]

%dsp-level:0
\[
(4)(\mathrm{ii})\ \text{re-derived by exhaustive enumeration at lengths}\
\]

%dsp-level:0
\[
\le 12\ [\text{machine}]
\]

%dsp-level:0
\[
\text{relaxation census}: \#\{\text{prefix-admissible completed words of length}\ 2m\}
\]

%dsp-level:0
\[
= \mathrm{Cat}_m\ [\text{machine enumeration} = \text{the closed form}\ \tfrac1{m+1}\binom{2m}m];
\]

%dsp-level:0
\[
\qquad \mathrm{Cat}_m=1\iff m\le1\ [\text{exact integers}];
\]

%dsp-level:0
\[
\qquad \text{two pairs}: \binom{2n}{2m_0}\ \text{labelled interleaving patterns summed over}\ m_0+m_1
\]

%dsp-level:0
\[
= n\ [\text{machine}]
\]

%dsp-level:0
\[
= \text{the closed form}
\]

%dsp-level:0
\[
(2):\ \mathrm{EM.9}\ \text{table verbatim}\ (\text{machine: the four sign pairs});
\]

%dsp-level:0
\[
\qquad \text{telescoping}: \sum_{i=1}^k\tfrac12(\sigma_{v_{i-1}}-\sigma_{v_i})
\]

%dsp-level:0
\[
= \tfrac12(\sigma_{v_0}-\sigma_{v_k})\ [\text{pairwise cancellation of the}\ k-1\ \text{interior terms};\ \text{machine: all}\ \sigma\in\{\pm1\}^{k+1},\ k\le7];
\]

%dsp-level:0
\[
\qquad \partial c=0\implies v_0=v_k\implies q_{\mathrm{EM}}=0\ [\text{machine}]
\]

%dsp-level:0
\[
(3):\ \partial c
\]

%dsp-level:0
\[
= t(c)-s(c)\ \text{is a signed pair of vertices};
\]

%dsp-level:0
\[
\ \partial(\text{vertex})
\]

%dsp-level:0
\[
= 0\implies\partial^2c
\]

%dsp-level:0
\[
= 0\ [\mathrm{D1.4};\ \text{machine: total multiplicity}\ 0\ \text{on every enumerated path}];
\]

%dsp-level:0
\[
\qquad \ker\partial=\{s(c)=t(c)\}\ [\text{machine: equality exactly on the closed paths}]
\]

%dsp-level:0
\[
(4):\ \text{a Noether charge is the generator of a one-parameter symmetry action}\ [\text{math}];
\]

%dsp-level:0
\[
\qquad \text{the value space is}\ \{-1,0,+1\}\ [\mathrm{EM.9}]\implies\text{the action factors through}\ S_3;
\]

%dsp-level:0
\[
\qquad U(1)\cong\mathbb Q/\mathbb Z\ \text{divisible}\ [\text{math}];
\]

%dsp-level:0
\[
\ \text{a divisible image in}\ S_3\ \text{consists of elements having}\ m\text{-th roots}\ \forall m:\ \{e\}\ \text{only}\ [\text{machine: the root census}];
\]

%dsp-level:0
\[
\qquad \text{hom-counts}: g^m=e\iff\operatorname{ord}(g)\mid m\ [\text{machine}]
\]

%dsp-level:0
\[
(5):\ \text{conjunction of}\ (1)\text{–}(4);
\]

%dsp-level:0
\[
\ \text{the limit step is BM.11's content, cited as a forward pointer and imported nowhere here}.\ \blacksquare
\]

\textbf{Phys.} Charge conservation in this corpus is a fact about integers, not about symmetries. A valid universe is one in which every crystallization has its dissolution — the definition says so, and the transition domains make the completed history unique: seal, unseal, seal, unseal. The endpoint charge read telescopes: sum the per-edge source along a materialized chain and every interior sign cancels, leaving the two ends; close the chain and the read is exactly zero. The ground is a boundary map with $\partial^2=0$ — the kernel statement of an integer chain complex. And the Noether route is not merely unnecessary here, it is closed: the charge alphabet has three elements, so any continuous charge-rotation action would have to land in a three-element permutation group, and a divisible group has no nontrivial finite image. There is no $U(1)$ to rotate the ternary read, hence no Noether current to build it from. What the continuum limit delivers later is the density form of a law that already holds, exactly, at every cell size.

\end{theorem}

\begin{theorem}[{\texorpdfstring{BM.8 — The continuum promotion: the hop becomes a translation, the $C_3$ becomes $SO(3)$, the tick stays an arrow, and the gauge algebras were never discrete}{BM.8 - The continuum promotion: the hop becomes a translation, the becomes , the tick stays an arrow, and the gauge algebras were never discrete}}]\label{tag:BM-8}
\textbf{Thm.}

%dsp-level:0
\[
(1)\quad\textbf{spatial translations from the discrete hop:}
\]

%dsp-level:0
\[
\qquad \text{on the emergent chart}\ \chi([v])
\]

%dsp-level:0
\[
= (x^0,x^1,x^2,x^3)\ [\mathrm{CT2.6};\ \mathrm{D1.8}:\ x^0:=\tau_n],
\]

%dsp-level:0
\[
\ \text{let}\ T_\ell f(x):
\]

%dsp-level:0
\[
= f(x+\ell)\ \text{be the hop of length}\ \ell\ \text{and}\ D_\ell:
\]

%dsp-level:0
\[
= \tfrac1\ell(T_\ell-1);
\]

%dsp-level:0
\[
\textbf{the exact commutator identity:}\qquad[\hat x,\,D_\ell]=-T_\ell\qquad\text{at every}\ \ell\ne0\qquad[\text{machine: exact on}\ x^k,\ k\le7];
\]

%dsp-level:0
\[
\textbf{the exact binomial expansion:}
\]

%dsp-level:0
\[
\qquad D_\ell\,x^k=\sum_{j=1}^{k}\binom{k}j x^{k-j}\ell^{\,j-1}=k\,x^{k-1}+\ell\cdot(\text{exact rational polynomial})
\]

%dsp-level:0
\[
\qquad [\text{machine}];
\]

%dsp-level:0
\[
\textbf{the Weyl transport:}
\]

%dsp-level:0
\[
\qquad T_\ell\,\hat x=(\hat x+\ell)\,T_\ell
\]

%dsp-level:0
\[
\qquad [\text{machine}];
\]

%dsp-level:0
\[
\qquad \ell\to0:\quad D_\ell\to\partial_x,\quad[\hat x,\partial_x]=-1,\quad\hat p:=-i\partial_x\implies[\hat x,\hat p]=i\ \ (\hbar=1\ [\mathrm{D2.9}]);
\]

%dsp-level:0
\[
\textbf{the dictionary carrier:}
\]

%dsp-level:0
\[
\qquad \text{the three}\ T\text{-odd coefficients}\ w_i\ [\mathrm{D2.1};\ \mathrm{CT1.2}]\ \longmapsto\ N_i=\beta w_i\ \longmapsto\ \bar g_{0i}=\varkappa N_i\ [\mathrm{CT1.12};\ \mathrm{D2.1}];
\]

%dsp-level:0
\[
N_i\ \text{is the Lagrange multiplier of the momentum constraint, whose charge generates spatial diffeomorphisms}\ [\text{math: ADM}]:
\]

%dsp-level:0
\[
\qquad \text{the conserved quantity is the momentum}\ P_i,
\]

%dsp-level:0
\[
\ \text{and}\ P_i\ \text{is}\ T\text{-odd, matching its carrier}\ [w_i\mapsto-w_i\ [\mathrm{FT.156}(3)]];
\]

%dsp-level:0
\[
(2)\quad\textbf{time translations: the tick is a semigroup, and the dynamics strictly contracts}\ —\ \textbf{an obstruction with a conditional charge:}
\]

%dsp-level:0
\[
\qquad \neg\exists\,\tau^{-1}\ [\mathrm{A6};\ \mathrm{FT.156}(5)];
\]

%dsp-level:0
\[
\qquad |\mathfrak L(\tau_{n+1})|\ge|\mathfrak L(\tau_n)|\ [\mathrm{A7};\ \text{machine}],\ \text{strict at every crystallization};
\]

%dsp-level:0
\[
V(\tilde K)\mapsto\bigl(1-\tfrac\gamma4\bigr)^2V(\tilde K),
\]

%dsp-level:0
\[
\qquad \text{deficit}=\tfrac{\gamma(8-\gamma)}{16},
\]

%dsp-level:0
\[
\qquad \gamma=1-4\varepsilon_q\in(0,1]\ [\mathrm{L2},\ \mathrm{D0.3},\ \mathrm{CT0.8};\ \text{machine}];
\]

%dsp-level:0
\[
\bigl(1-\tfrac\gamma4\bigr)^2+\tfrac{\gamma(8-\gamma)}{16}=1\ \text{exactly};
\]

%dsp-level:0
\[
\qquad \tfrac{\gamma(8-\gamma)}{16}=0\iff\gamma\in\{0,8\}\ [\text{machine}];
\]

%dsp-level:0
\[
\qquad \gamma=0\iff\varepsilon_q=\tfrac14\notin[0,\tfrac14)\ [\mathrm{CT0.8}];
\]

%dsp-level:0
\[
\therefore\ \text{on the whole admissible range the tick map strictly contracts}\ V:\quad\neg\exists\,
\]

%dsp-level:0
\[
\text{a nonzero tick-invariant of the pair-state dynamics}\implies\neg\exists\,
\]

%dsp-level:0
\[
\text{an exact time-translation invariance of the microscopic dynamics};
\]

%dsp-level:0
\[
\textbf{the conditional energy read}\ (\text{declared scope, theorem-local, L0's dynamical-scope pattern}):\ \text{under the stationarity hypothesis — the coarse-grained coefficient fields}\ (a_{00},q_{ij},w_i)\ \text{of}\ \mathrm{D2.1}\ \text{are tick-independent across the cells of}\ \mathrm{CT2.6}(\mathrm{ii})\text{–}(\mathrm{iii})\ —\ \text{the continuum limit carries a time-translation one-parameter group whose charge is the energy read};
\]

%dsp-level:0
\[
\text{on the sealed sector the hypothesis is structural rather than imposed: a bound chain's reads are tick-independent}\ [\mathrm{D1.5};\ \text{forward pointer:}\ \mathrm{BM.9}(3)];
\]

%dsp-level:0
\[
\qquad \text{at}\ \gamma\to0\ \text{the contraction}\to1\ \text{and the deficit}\to0\ [\text{exact rationals}];
\]

%dsp-level:0
\[
(3)\quad\textbf{rotations: the}\ C_3\ \textbf{is a rational point of a one-parameter subgroup already present in the register:}
\]

%dsp-level:0
\[
\qquad \mathcal W:=W_1{+}W_2{+}W_3,
\]

%dsp-level:0
\[
\qquad \mathcal W^3=3\mathcal W
\]

%dsp-level:0
\[
\qquad [\mathrm{FT.158}(1);\ \text{machine}];
\]

%dsp-level:0
\[
X:=\hat n\cdot W=\mathcal W/\sqrt3,\quad\hat n=\tfrac{(1,1,1)}{\sqrt3}:
\]

%dsp-level:0
\[
\qquad X^3=X,
\]

%dsp-level:0
\[
\qquad \operatorname{charpoly}(\mathcal W)=x^4-3x^2\implies\operatorname{spec}(X)=\{1,0,-1,0\}
\]

%dsp-level:0
\[
\qquad [\text{machine: Faddeev–LeVerrier, exact traces}];
\]

%dsp-level:0
\[
\textbf{the promotion identity:}
\]

%dsp-level:0
\[
\qquad u(\theta):=1+(\cos\theta-1)X^2-i\sin\theta\,X
\]

%dsp-level:0
\[
\qquad \text{satisfies}
\]

%dsp-level:0
\[
\qquad u(\theta)u(\phi)=u(\theta+\phi)\ \ \forall\theta,\phi
\]

%dsp-level:0
\[
\text{exactly in}\ \mathbb Q[i][c,s,c',s',X]/(c^2{+}s^2{-}1,\ c'^2{+}s'^2{-}1,\ X^3{-}X)\ [\text{machine}]:
\]

%dsp-level:0
\[
\qquad u(0)=1,\quad u(\theta)u(-\theta)=1,\quad u(\theta)^3=u(3\theta)\ [\text{triple-angle polynomials}];
\]

%dsp-level:0
\[
\text{the group law closes in}\ \operatorname{span}\{1,X,X^2\}:\ X^3=X\ \text{is the only relation used}\ [\text{machine: degree census}];
\]

%dsp-level:0
\[
\textbf{the discrete elements are its rational points:}
\]

%dsp-level:0
\[
\qquad \theta=\tfrac{2\pi}3:
\]

%dsp-level:0
\[
\quad U=E_0-\tfrac12\mathcal W^2-\tfrac i2\mathcal W=u\otimes u,
\]

%dsp-level:0
\[
\quad U^3=E_0,
\]

%dsp-level:0
\[
\quad UX_iU^\dagger=X_{i+1}\ (X\in\{W,J,K\})\ [\mathrm{FT.111},\ \mathrm{FT.158}(1);\ \text{machine}];
\]

%dsp-level:0
\[
\theta=\pi:\quad U_{ij}=E_0-(W_i+W_j)^2,\quad U_{ij}^2=E_0,\quad U_{ij}:\ X_i\leftrightarrow X_j,\ X_k\mapsto-X_k\ [\mathrm{FT.156}(4);\ \text{machine}];
\]

%dsp-level:0
\[
C_3\le SO(3)\ \text{is inner in the weak rotation algebra}\ [\mathrm{CT1.4}:\ \operatorname{span}\{W_i\}\cong\mathfrak{su}(2);\ \mathrm{CT1.8}:\ J,K\ \text{vectors}]:
\]

%dsp-level:0
\[
\qquad \text{the promotion is exact,}\ \neg\text{approximate};
\]

%dsp-level:0
\[
\textbf{the charge:}
\]

%dsp-level:0
\[
\qquad \textstyle\sum_iW_i^2=\tfrac32E_0+\tfrac12\Sigma=E_0+\mathbb S=2\,\Pi_{\mathcal R}\ [\mathrm{FT.157}(2);\ \text{machine}]:
\]

%dsp-level:0
\[
\qquad \text{Casimir}\ 2\ \text{on}\ \vee^2\ (j{=}1),\quad 0\ \text{on}\ \wedge^2\ (j{=}0)\ [\text{machine}];
\]

%dsp-level:0
\[
\text{the integer weight lattice}\ \{1,0,-1\}\ \text{is what makes the rational angles admissible:}\ u(\theta+2\pi)
\]

%dsp-level:0
\[
= u(\theta)\ \text{on integer weights}\ [\text{math}];
\]

%dsp-level:0
\[
(4)\quad\textbf{gauge: the slot algebras are natively continuous — no limit is needed for the algebra:}
\]

%dsp-level:0
\[
\qquad [W_i,W_j]=i\epsilon_{ijk}W_k\ \text{on all}\ 9\ \text{pairs}\ [\mathrm{CT1.4};\ \text{machine}];
\]

%dsp-level:0
\[
\mathfrak u(2)=\mathfrak{su}(2)\oplus\mathfrak u(1)_{\mathrm{centre}}:
\]

%dsp-level:0
\[
\quad Y=\tfrac12\sigma_0\ \text{commutes with all}\ W_i,
\]

%dsp-level:0
\[
\quad 3\ \text{traceless}\ +\ 1\ \text{central}=4\ [\mathrm{CT0.13},\ \mathrm{EM.2};\ \text{machine}];
\]

%dsp-level:0
\[
\pi:\mathfrak u(3)\to\mathfrak{su}(3),
\]

%dsp-level:0
\[
\ \ \pi(X)
\]

%dsp-level:0
\[
= X-\tfrac{\mathfrak T(X)}3I_3:\ \ \text{Lie homomorphism on all}\ 81\ \text{basis pairs},
\]

%dsp-level:0
\[
\ \ker\pi
\]

%dsp-level:0
\[
= \operatorname{span}\{I_3\}\ [\mathrm{CT3.7};\ \text{machine}];
\]

%dsp-level:0
\[
\text{the colour block: all}\ 8\ \text{generators traceless},
\]

%dsp-level:0
\[
\ \text{all}\ 64\ \text{commutators in the span},
\]

%dsp-level:0
\[
\ \text{Jacobi on all}\ 512\ \text{triples},
\]

%dsp-level:0
\[
\ [H_1,H_2]
\]

%dsp-level:0
\[
= 0,
\]

%dsp-level:0
\[
\ \det\mathrm K
\]

%dsp-level:0
\[
\ne 0\ [\mathrm{CT3.8},\ \mathrm{FT.82}(3);\ \text{machine: exact rationals}];
\]

%dsp-level:0
\[
\textbf{what the continuum limit supplies is the gauging,}\ \neg\text{the algebra: the parameters become fields on}\ M\ [\mathrm{CT2.6}]\ \text{and the covariant derivative is the chart-side object}\ [\text{math}];
\]

%dsp-level:0
\[
(5)\quad\textbf{the three charge statuses are distinct:}
\]

%dsp-level:0
\[
\qquad \textbf{colour:}\ \text{on a sealed composite every Hermitian slot operator reads only its trace}\ \mathfrak T(H)
\]

%dsp-level:0
\[
= 3a\ [\mathrm{BM.1}(2);\ \text{machine}]\ \text{and the invariant line is}\ 1\text{-dimensional}\ [\mathrm{BM.1}(3);\ \text{machine}:\ 27-\operatorname{rk}=1]:\ \text{the colour charge of a sealed singlet is}\ 0\ \text{exactly};
\]

%dsp-level:0
\[
\textbf{weak isospin:}\ \text{the derived Casimir of}\ (3)\ —\ 2\ \text{on the symmetric sector,}\ 0\ \text{on the antisymmetric one}\ [\mathrm{FT.157}(2)];
\]

%dsp-level:0
\[
\textbf{electric charge:}\ \neg\text{a Noether charge}\ [\mathrm{BM.7}(4)]:\ \text{the ternary read is topological};
\]

%dsp-level:0
\[
\qquad \text{the gauge Noether story of the corpus therefore has two continuous charges and one discrete invariant, not three of a kind}.
\]

\textbf{Rmk (promotion discipline).} \textit{Momentum}, \textit{energy} and \textit{angular momentum} name the charges of the promoted one-parameter groups; the names are physics vocabulary for derived structures [A §0.2]. Clause (2)'s stationarity hypothesis is a theorem-local dynamical scope of the species of L0's crystallization–tilt coupling and CT4.5a's ensemble scope, declared at the point of use; it is not an axiom, and it is not discharged by any tag of Parts 1–3.

\textbf{Req.} A6, A7 (Part 1); D0.3, D1.5, D1.8, D2.1 (Parts 1–2); CT0.8, CT0.13, CT1.2, CT1.4, CT1.8 (Part 1); CT1.12, CT2.6, EM.2 (Part 2; CT2.6's conditions (ii)–(iii) at the point of use — A §0.2a); CT3.7, CT3.8 (Part 2); L2 (Part 1); Def-BM.3, BM.1(2), BM.1(3) (this part); FT.82(3), FT.111, FT.144(2), FT.156(3), FT.156(4), FT.156(5), FT.157(2), FT.158(1) (Part 3); exact $\mathbb Q[x,\ell]$ hop calculus, exact identities in $\mathbb Q[i][c,s,c',s',X]/(c^2{+}s^2{-}1,X^3{-}X)$, Faddeev–LeVerrier characteristic polynomials, exact $\mathbb Q(i)$ conjugation tables, exact $\mathbb Q$ structure constants and Killing determinants, exact rational contraction arithmetic [math; machine: \texttt{checks/\allowbreak{}verify\_\allowbreak{}H\_\allowbreak{}noether.py} §B].

\textbf{Pf.}

%dsp-level:0
\[
(1):\ (T_\ell-1)f\ \text{is divisible by}\ \ell\ \text{in}\ \mathbb Q[x,\ell]\ \text{for every polynomial}\ f\ [\text{machine}];
\]

%dsp-level:0
\[
\qquad \hat x D_\ell f-D_\ell(\hat x f)
\]

%dsp-level:0
\[
= \tfrac1\ell\bigl[x(f(x{+}\ell)-f(x))-((x{+}\ell)f(x{+}\ell)-xf(x))\bigr]
\]

%dsp-level:0
\[
= -\tfrac{\ell f(x+\ell)}\ell
\]

%dsp-level:0
\[
= -T_\ell f\ [\text{exact}]
\]

%dsp-level:0
\[
D_\ell x^k
\]

%dsp-level:0
\[
= \tfrac{(x+\ell)^k-x^k}\ell
\]

%dsp-level:0
\[
= \sum_{j\ge1}\binom kj x^{k-j}\ell^{j-1}\ [\text{binomial theorem; machine coefficient census}];
\]

%dsp-level:0
\[
\qquad \ell\to0\ \text{leaves}\ kx^{k-1}\ \text{and}\ [\hat x,\partial_x]=-1;
\]

%dsp-level:0
\[
\qquad \hat p=-i\partial_x\implies[\hat x,\hat p]=i\ [\hbar=1,\ \mathrm{D2.9}]
\]

%dsp-level:0
\[
\text{dictionary rows verbatim}\ [\mathrm{D2.1},\ \mathrm{CT1.12}];
\]

%dsp-level:0
\[
\qquad N_i\ \text{as the momentum-constraint multiplier and}\ P_i\ \text{as its charge}\ [\text{math: ADM}];
\]

%dsp-level:0
\[
\qquad w_i\mapsto-w_i\ \text{under}\ T\ [\mathrm{FT.156}(3)\ \text{re-cited}]
\]

%dsp-level:0
\[
(2):\ \mathrm{A6}/\mathrm{A7}\ \text{verbatim}\ [\text{machine: monotone and strictly increasing samples}];
\]

%dsp-level:0
\[
\qquad \mathrm{L2}\ \text{contraction re-derived as an exact polynomial identity in}\ \gamma;
\]

%dsp-level:0
\[
\qquad \text{zero set}\ \{0,8\}\ [\text{machine over a rational sweep and by factorization}];
\]

%dsp-level:0
\[
\qquad \gamma\in(0,1]\ \text{on}\ \varepsilon_q\in[0,\tfrac14)\ [\mathrm{D0.3},\ \mathrm{CT0.8}]
\]

%dsp-level:0
\[
\text{strict contraction}\ \forall\gamma
\]

%dsp-level:0
\[
\in (0,1]
\]

%dsp-level:0
\[
\implies V(\mathcal T^n\rho)\to0\ \text{and no nonzero invariant}\ [\mathrm{L1}\ \text{monotone};\ \text{exact rationals}];
\]

%dsp-level:0
\[
\qquad \text{the conditional clause states its hypothesis explicitly and cites no tag for it}
\]

%dsp-level:0
\[
(3):\ \mathcal W^3=3\mathcal W\ [\text{machine}]\implies X^3=X;
\]

%dsp-level:0
\[
\qquad \text{charpoly}\ x^4-3x^2\ [\text{Faddeev–LeVerrier, exact};\ \operatorname{tr}\mathcal W=0,\ \operatorname{tr}\mathcal W^2=6,\ \operatorname{tr}\mathcal W^4=18];
\]

%dsp-level:0
\[
\qquad u(\theta)u(\phi)=u(\theta+\phi)\ [\text{machine in the quotient ring, using the addition formulas and}\ X^4=X^2]
\]

%dsp-level:0
\[
\theta
\]

%dsp-level:0
\[
= \tfrac{2\pi}3:\ (\cos\theta-1)
\]

%dsp-level:0
\[
= -\tfrac32,
\]

%dsp-level:0
\[
\ \sin\theta/\sqrt3
\]

%dsp-level:0
\[
= \tfrac12
\]

%dsp-level:0
\[
\implies u
\]

%dsp-level:0
\[
= E_0-\tfrac12\mathcal W^2-\tfrac i2\mathcal W
\]

%dsp-level:0
\[
= u\otimes u\ [\text{machine entrywise}];
\]

%dsp-level:0
\[
\qquad U^3=E_0,\ UX_iU^\dagger=X_{i+1}\ [\text{machine, all}\ 9];
\]

%dsp-level:0
\[
\qquad \theta=\pi:\ (\cos\pi-1)=-2,\ X=(W_i{+}W_j)/\sqrt2\implies U_{ij}=E_0-(W_i{+}W_j)^2\ [\text{machine}]
\]

%dsp-level:0
\[
\textstyle\sum_iW_i^2=\tfrac32E_0+\tfrac12\Sigma\ [\mathrm{FT.157}(2)\ \text{re-derived};\ \text{machine}];
\]

%dsp-level:0
\[
\qquad \text{Casimir reads on}\ \Pi_{\mathcal R},\ \Pi_{\mathcal V}\ [\text{machine}]
\]

%dsp-level:0
\[
(4):\ \text{all rows re-derived entrywise over}\ \mathbb Q(i)\ \text{and}\ \mathbb Q\ [\text{machine}];
\]

%dsp-level:0
\[
\qquad \text{the homomorphism and kernel rows of}\ \mathrm{CT3.7}\ \text{on the}\ \mathfrak u(3)\ \text{basis}\ \{T_a\}\cup\{I_3\}\ [\text{machine: all}\ 81\ \text{pairs}]
\]

%dsp-level:0
\[
(5):\ \mathrm{BM.1}(2)\text{–}(3)\ \text{re-derived on exact witnesses}\ [\text{machine}:\ \langle\varepsilon|\sum_iH^{(i)}|\varepsilon\rangle=\mathfrak T(H)\langle\varepsilon|\varepsilon\rangle;\ (\ H\otimes I-I\otimes H^{\mathrm T})\operatorname{vec}(I)=0;\ 27-\operatorname{rk}=1];
\]

%dsp-level:0
\[
\qquad \mathrm{BM.7}(4)\ \text{re-cited}.\ \blacksquare
\]

\textbf{Phys.} Three of the four continuum symmetries are already sitting in the register; the fourth is not, and the structure yields why. The hop is a translation with an exact error term: the commutator of position with the finite difference is minus the shift operator at every step size, and the binomial expansion hands over the derivative as its leading term, so the momentum generator and its canonical commutator are the $\ell\to0$ read of an identity that never needed a limit to be true. Rotation is stronger still: because the weak generator satisfies $X^3=X$, the exponential of $-i\theta X$ is a quadratic polynomial in $X$ with $\cos\theta$ and $\sin\theta$ as coefficients, and the group law holds as a ring identity — the $120^\circ$ seal rotation and the $180^\circ$ transpositions are just two rational points of a one-parameter group that was there all along, with integer weights making the periodicity exact. The gauge algebras were never discrete: $\mathfrak{su}(2)$, $\mathfrak{su}(3)$ and the rung-2 centre close as Lie algebras over the rationals, and what the continuum adds is the gauging, not the algebra. Time is the exception, and it is an honest one — the tick has no inverse and the ledger dynamics strictly contracts its Lyapunov function at every admissible threshold, so the microscopic theory has no exact time-translation symmetry and the energy charge exists only under a declared stationarity scope, which the sealed sector satisfies structurally. The charge ledger then splits three ways: colour reads zero on every sealed singlet, weak isospin reads the derived Casimir, and electric charge is not a Noether charge at all.

\end{theorem}

\begin{theorem}[{\texorpdfstring{BM.9 — The stress-energy tensor of a sealed composite: the two reads, their exact sign reconciliation, and the dust source}{BM.9 - The stress-energy tensor of a sealed composite: the two reads, their exact sign reconciliation, and the dust source}}]\label{tag:BM-9}
\textbf{Thm.}

%dsp-level:0
\[
(1)\quad\textbf{the corpus definition and the field equations:}
\]

%dsp-level:0
\[
\qquad T_{\mu\nu}:=-\frac{2}{\sqrt{-g}}\frac{\delta S_{\mathrm{force}}}{\delta g^{\mu\nu}}
\]

%dsp-level:0
\[
\qquad [\mathrm{CT1.22}];
\]

%dsp-level:0
\[
S_{\mathrm{eff}}=S_{\mathrm{geom}}+S_{\mathrm{force}},\qquad S_{\mathrm{geom}}[g]=\frac1{2\kappa^2}\int d^4x\,\sqrt{-g}\,(R-2\Lambda)\qquad[\mathrm{CT1.21}];
\]

%dsp-level:0
\[
\delta S_{\mathrm{eff}}=0\iff G_{\mu\nu}+\Lambda g_{\mu\nu}=\kappa^2T_{\mu\nu},\qquad\kappa^2=8\pi G_N=\frac{2\pi\,\mathfrak h^2}{\ln2}\qquad[\mathrm{CT1.23},\ \mathrm{CT2.9},\ \mathrm{P3}];
\]

%dsp-level:0
\[
(2)\quad\textbf{the canonical read and the exact sign reconciliation:}
\]

%dsp-level:0
\[
\qquad \text{for}\ \mathcal L=-\tfrac12g^{\alpha\beta}\partial_\alpha\phi\,\partial_\beta\phi-\tfrac12m^2\phi^2:
\]

%dsp-level:0
\[
\qquad \frac{\partial\mathcal L}{\partial(\partial^\mu\phi)}=-\partial_\mu\phi\ \ [\text{exact}];
\]

%dsp-level:0
\[
T^{\mathrm{can}}_{\mu\nu}:=\frac{\partial\mathcal L}{\partial(\partial^\mu\phi)}\,\partial_\nu\phi-g_{\mu\nu}\mathcal L=-\partial_\mu\phi\,\partial_\nu\phi-g_{\mu\nu}\mathcal L;
\]

%dsp-level:0
\[
\qquad T^{\mathrm{Hilb}}_{\mu\nu}=\partial_\mu\phi\,\partial_\nu\phi+g_{\mu\nu}\mathcal L;
\]

%dsp-level:0
\[
\boxed{\ T^{\mathrm{can}}_{\mu\nu}=-\,T^{\mathrm{Hilb}}_{\mu\nu}\ }
\]

%dsp-level:0
\[
\qquad \text{exactly, at signature}\ (-,+,+,+)\ [\mathrm{A3},\ \mathrm{C13},\ \mathrm{CT1.12}]\ \text{with this}\ \mathcal L\ \text{sign}
\]

%dsp-level:0
\[
\qquad [\text{machine: all}\ 16\ \text{components on}\ 3\ \text{rational metric witnesses}];
\]

%dsp-level:0
\[
\textbf{the convention is fixed by the energy sign:}\quad\text{at rest}\ T^{\mathrm{Hilb}}_{00}
\]

%dsp-level:0
\[
= \tfrac12\sum_i(\partial_i\phi)^2+\tfrac12m^2\phi^2\ \text{with every coefficient}\ +\tfrac12,
\]

%dsp-level:0
\[
\ \text{while}\ T^{\mathrm{can}}_{00}\ \text{carries}\ -\tfrac12\ [\text{machine: exact coefficient census}];
\]

%dsp-level:0
\[
\therefore\ \mathrm{CT1.22}\text{'s minus is the positive-energy read at}\ (-,+,+,+);
\]

%dsp-level:0
\[
\qquad \text{the variational composition is exact:}\quad\frac{\delta(\sqrt{-g}\,\mathcal L)}{\delta g^{\mu\nu}}
\]

%dsp-level:0
\[
= -\tfrac12\sqrt{-g}\,
\]

%dsp-level:0
\[
\bigl(\partial_\mu\phi\,\partial_\nu\phi+g_{\mu\nu}\mathcal L\bigr)\ [\text{machine}];
\]

%dsp-level:0
\[
(3)\quad\textbf{the Noether identity (translation current):}
\]

%dsp-level:0
\[
\qquad \partial^\mu T^{\mathrm{Hilb}}_{\mu\nu}=(\Box\phi-m^2\phi)\,\partial_\nu\phi
\]

%dsp-level:0
\[
\qquad \text{as an exact jet identity}\ [\text{machine: all}\ 4\ \text{components},\ 3\ \text{metric witnesses}];
\]

%dsp-level:0
\[
\Box\phi-m^2\phi=0\ \text{(the field equation of this}\ \mathcal L)\implies\partial^\mu T^{\mathrm{Hilb}}_{\mu\nu}=0:
\]

%dsp-level:0
\[
\qquad T_{\mu\nu}\ \text{is the Noether current of the translation group promoted in}\ \mathrm{BM.8}(1);
\]

%dsp-level:0
\[
\textbf{the trace identity:}\qquad g^{\mu\nu}T^{\mathrm{Hilb}}_{\mu\nu}=-(\partial\phi)^2-2m^2\phi^2\qquad[\text{machine: exact in}\ 4\ \text{dimensions}];
\]

%dsp-level:0
\[
(4)\quad\textbf{the sealed composite at rest — the neutron}\ [\mathrm{Def\text{-}BM.3}:\ \text{the}\ \varepsilon\text{-sealed 3-juncture in the neutral charge class}]:
\]

%dsp-level:0
\[
\qquad T^{\mathrm{bd}}_{\mu\nu}=(\rho+p)u_\mu u_\nu+p\,g_{\mu\nu}
\]

%dsp-level:0
\[
\qquad [\mathrm{CT1.22}\ \text{bound part}];
\]

%dsp-level:0
\[
u=(1,0,0,0),\quad g_{\mu\nu}=\mathrm{diag}(-1,1,1,1),\quad c=1\ [\mathrm{D2.9}]:
\]

%dsp-level:0
\[
\qquad T_{\mu\nu}=\mathrm{diag}(\rho,\,p,\,p,\,p)
\]

%dsp-level:0
\[
\qquad [\text{machine: exact component identity}];
\]

%dsp-level:0
\[
\textbf{fully sealed}\implies p=0,\quad w=\tfrac p\rho=0
\]

%dsp-level:0
\[
\qquad [\mathrm{CT4.5}:\ T^{(D^-)}_{\mu\nu}=\rho\,u_\mu u_\nu,\ p_{D^-}=0]:
\]

%dsp-level:0
\[
\qquad T_{\mu\nu}=\mathrm{diag}(\rho,0,0,0)\ \ [\text{machine}];
\]

%dsp-level:0
\[
\textbf{static:}\ \ \partial_0T_{\mu\nu}=0;
\]

%dsp-level:0
\[
\qquad \textbf{spherically symmetric:}\ \ \text{the isotropic spatial block}\ p\,
\]

%dsp-level:0
\[
\delta_{ij}\ \text{is invariant under all}\ 48\ \text{signed triad permutations}\ \subset O(3),
\]

%dsp-level:0
\[
\ \text{hence under the promoted}\ SO(3)\ [\mathrm{BM.8}(3);\ \text{machine}];
\]

%dsp-level:0
\[
\textbf{trace and energy conditions:}
\]

%dsp-level:0
\[
\qquad g^{\mu\nu}T_{\mu\nu}=-\rho+3p\overset{p=0}{=}-\rho;
\]

%dsp-level:0
\[
\qquad \text{WEC},\ \text{DEC},\ \text{SEC hold exactly at}\ \rho\ge0,\ p=0\ [\text{machine: rational witnesses}];
\]

%dsp-level:0
\[
\rho\ge0,\qquad\rho=0\iff\neg\exists v:\ n(v)\ge4\qquad[\mathrm{CT4.5};\ \text{machine on enumerated rung profiles}];
\]

%dsp-level:0
\[
(5)\quad\textbf{the source read and its dimensions:}
\]

%dsp-level:0
\[
\qquad \rho=\frac{m_{\mathrm{hid}}c^2}{V}=\frac{m_{\mathrm{hid}}}{V}\ \ [c=1,\ \mathrm{D2.9}],
\]

%dsp-level:0
\[
\qquad m_{\mathrm{hid}}=\frac{E_{\mathrm{int}}\ln2}{\mathfrak h}
\]

%dsp-level:0
\[
\qquad [\mathrm{FT.44}(4),\ \mathrm{FT.42}(1)];
\]

%dsp-level:0
\[
[\mathfrak h]=\mathfrak q^{-1/2},\quad[\text{length}]=\mathfrak q^{-1/2},\quad[\text{energy}]=\mathfrak q^{+1/2},\quad[T_{\mu\nu}]=\mathfrak q^{+2}\qquad[\mathrm{CT2.1},\ \mathrm{D2.9}];
\]

%dsp-level:0
\[
[m_{\mathrm{hid}}]=[E_{\mathrm{int}}]-[\mathfrak h]=0+\tfrac12=\mathfrak q^{+1/2}=[\text{energy}];
\]

%dsp-level:0
\[
\qquad [V]=\mathfrak q^{-3/2};
\]

%dsp-level:0
\[
\qquad [\rho]=\mathfrak q^{1/2-(-3/2)}=\mathfrak q^{+2}=[T_{\mu\nu}]
\]

%dsp-level:0
\[
\qquad [\text{machine: exact exponent arithmetic}];
\]

%dsp-level:0
\[
\textbf{the}\ \ln2\ \textbf{cancellation:}
\]

%dsp-level:0
\[
\qquad \kappa^2\rho=\frac{2\pi\,\mathfrak h^2}{\ln2}\cdot\frac{E_{\mathrm{int}}\ln2}{\mathfrak h\,V}=\frac{2\pi\,\mathfrak h\,E_{\mathrm{int}}}{V}
\]

%dsp-level:0
\[
\qquad [\text{machine: exact, denominators cleared}]:
\]

%dsp-level:0
\[
\text{the}\ \ln2\ \text{of the entropy calibration}\ [\mathrm{CT2.9}]\ \text{and the}\ \ln2\ \text{of the informational mass}\ [\mathrm{FT.44}(4)]\ \text{cancel in the sourced field equation}:G_{\mu\nu}+\Lambda g_{\mu\nu}
\]

%dsp-level:0
\[
= \tfrac{2\pi\mathfrak h E_{\mathrm{int}}}{V}\,
\]

%dsp-level:0
\[
u_\mu u_\nu\ \text{at}\ p\qquad = 0;
\]

%dsp-level:0
\[
(6)\quad\textbf{the width (forward pointer:}\ \mathrm{BM.10}(6)\textbf{):}
\]

%dsp-level:0
\[
\qquad V=(\Delta x)^3,
\]

%dsp-level:0
\[
\qquad \Delta x\ge\frac1{2\,\Delta p}\ \ (\hbar=1):
\]

%dsp-level:0
\[
\qquad \rho=\frac{m_{\mathrm{hid}}}{V}\le8\,m_{\mathrm{hid}}(\Delta p)^3
\]

%dsp-level:0
\[
\qquad [\text{machine: exact coefficients}];
\]

%dsp-level:0
\[
\rho\ \text{is finite for every}\ \Delta p>0:
\]

%dsp-level:0
\[
\qquad T_{\mu\nu}\ \text{is a distribution of nonzero width},\ \neg\text{a}\ \delta\text{-source};
\]

%dsp-level:0
\[
\qquad \text{the width is the scale at which the coarse-grained fields are}\ \mathrm{CT2.6}(\mathrm{ii})\text{-regular}.
\]

\textbf{Rmk (convention record).} Clause (2) prices both sign conventions instead of choosing silently: the canonical read of the classical formula and CT1.22's Hilbert read differ by an exact minus at signature $(-,+,+,+)$ with this $\mathcal L$ sign, and the corpus convention is the one with positive rest energy. Every downstream row — the dust form, the trace, the field equations, the energy conditions — is stated in the CT1.22 convention. No measured value enters any clause [P5; Def-BM.3(6)].

\textbf{Req.} A3 (Part 1); D1.5, D2.9 (Parts 1–2); CT1.21, CT1.22, CT1.23, CT2.1, CT2.9, CT4.5 (Part 2); Def-BM.1, Def-BM.3, BM.8(1), BM.8(3) (this part); FT.42(1) (Part 3); Census FT.44(4); P3, P5 (Parts 1–2); exact multivariate jet-ring polynomial identities, exact rational metric inversion by Cramer, exact coefficient censuses, Fractions [math; machine: \texttt{checks/\allowbreak{}verify\_\allowbreak{}H\_\allowbreak{}noether.py} §C].

\textbf{Pf.}

%dsp-level:0
\[
(1):\ \mathrm{CT1.21}/\mathrm{CT1.22}/\mathrm{CT1.23}\ \text{verbatim};
\]

%dsp-level:0
\[
\qquad \kappa^2=8\pi G_N=8\pi\cdot\tfrac{\mathfrak h^2}{4\ln2}=\tfrac{2\pi\mathfrak h^2}{\ln2}\ [\mathrm{CT2.9};\ \text{machine, denominators cleared}]
\]

%dsp-level:0
\[
(2):\ \partial\mathcal L/\partial(\partial^\mu\phi)
\]

%dsp-level:0
\[
= -g^{\mu\alpha}\partial_\alpha\phi\cdot\tfrac12\cdot2
\]

%dsp-level:0
\[
= -\partial^\mu\phi\ \text{and}\ \partial^\mu\phi\ \text{lowers to}\ \partial_\mu\phi\ \text{in the first slot}\ [\text{exact}];
\]

%dsp-level:0
\[
\qquad T^{\mathrm{can}}+T^{\mathrm{Hilb}}
\]

%dsp-level:0
\[
= 0\ \text{componentwise}\ [\text{machine:}\ 3\ \text{rational metrics, diagonal and fully off-diagonal}];
\]

%dsp-level:0
\[
\qquad \delta\sqrt{-g}
\]

%dsp-level:0
\[
= -\tfrac12\sqrt{-g}\,
\]

%dsp-level:0
\[
g_{\mu\nu}\delta g^{\mu\nu},
\]

%dsp-level:0
\[
\ \delta\mathcal L/\delta g^{\mu\nu}
\]

%dsp-level:0
\[
= -\tfrac12\partial_\mu\phi\partial_\nu\phi\ [\text{math}]\implies\text{the variational composition}\ [\text{machine}]
\]

%dsp-level:0
\[
(3):\ \partial_\mu(d_\mu\phi\,d_\nu\phi)=\Box\phi\,\partial_\nu\phi+\partial^\mu\phi\,\partial_\mu\partial_\nu\phi;
\]

%dsp-level:0
\[
\qquad g_{\mu\nu}\partial^\mu\mathcal L=-\partial_\nu\phi\,(\partial^\alpha\phi)-m^2\phi\,\partial_\nu\phi\ \text{at constant metric};
\]

%dsp-level:0
\[
\qquad \text{the two}\ \partial\partial\phi\ \text{terms cancel by symmetry of}\ \partial_\mu\partial_\nu\ \text{and relabeling}\ [\text{machine: the residual is exactly}\ (\Box\phi-m^2\phi)\partial_\nu\phi]
\]

%dsp-level:0
\[
(4):\ (\rho+p)u_\mu u_\nu+pg_{\mu\nu}\ \text{at}\ u
\]

%dsp-level:0
\[
= (1,0,0,0),
\]

%dsp-level:0
\[
\ g
\]

%dsp-level:0
\[
= \mathrm{diag}(-1,1,1,1):\ (0,0)\text{-entry}\ \rho,
\]

%dsp-level:0
\[
\ (0,i)\text{-entries}\ 0,
\]

%dsp-level:0
\[
\ (i,j)\text{-entry}\ p\delta_{ij}\ [\text{machine}];
\]

%dsp-level:0
\[
\qquad p=0,\ w=0\ [\mathrm{CT4.5}\ \text{re-cited}];
\]

%dsp-level:0
\[
\qquad R\,\mathrm{diag}(\rho,p,p,p)R^{\mathrm T}=\mathrm{diag}(\rho,p,p,p)\ \forall R\in\{\text{signed permutations}\}\ [\text{machine: all}\ 48];
\]

%dsp-level:0
\[
\qquad \text{trace and the three energy conditions at}\ p=0\ [\text{machine}]
\]

%dsp-level:0
\[
(5):\ c=1\ [\mathrm{D2.9}];
\]

%dsp-level:0
\[
\qquad m_{\mathrm{hid}}=E_{\mathrm{int}}\ln2/\mathfrak h\ [\mathrm{FT.44}(4)\ \text{re-cited}];
\]

%dsp-level:0
\[
\qquad \text{exponents}: \tfrac12-3(-\tfrac12)=2\ [\text{machine}];
\]

%dsp-level:0
\[
\qquad \kappa^2\rho:\ (2\pi\mathfrak h^2/\ln2)(E_{\mathrm{int}}\ln2/(\mathfrak h V))
\]

%dsp-level:0
\[
= 2\pi\mathfrak h E_{\mathrm{int}}/V\ [\text{machine: cross-multiplied identity}]
\]

%dsp-level:0
\[
(6):\ \Delta x\,\Delta p=\tfrac12\ \text{on rational witnesses}\ [\mathrm{BM.10}(6),\ \text{forward pointer}];
\]

%dsp-level:0
\[
\qquad (1/(2\Delta p))^3=1/(8(\Delta p)^3)\ [\text{machine}];
\]

%dsp-level:0
\[
\qquad \rho V=m_{\mathrm{hid}}\ \text{with}\ V>0\implies\rho<\infty\ [\text{exact}].\ \blacksquare
\]

\textbf{Phys.} A sealed composite gravitates as dust, and the corpus can now say so with both sign conventions on the table. The classical canonical formula and CT1.22's metric definition differ by an exact minus at the corpus's own signature; the minus is what makes the rest energy positive, so the convention is decided by the energy condition rather than by taste. The same tensor is the translation current: its divergence is the field equation times the field gradient, an identity in the jet variables that vanishes on-shell — Noether's theorem, run in reverse from the symmetry the hop promoted. For the neutron, the fully sealed three-juncture in the neutral class, the bound-part form collapses to $\mathrm{diag}(\rho,0,0,0)$: static, spherically symmetric, pressure-free, exactly the $D^-$ equation of state, and satisfying every energy condition at $\rho\ge0$. Its source is the informational mass over the width volume, and the dimensions balance to $\mathfrak q^{+2}$ without adjustment. The two $\ln2$ factors — one from the horizon-entropy calibration, one from the bit-count mass read — cancel in the sourced field equation, leaving $2\pi\mathfrak h E_{\mathrm{int}}/V$. And because the width is strictly positive, the source is a distribution, not a point: the smoothing scale is where the coarse-graining hypothesis of the continuum limit takes over.

\end{theorem}

\medskip\noindent\textit{\texorpdfstring{BM.10 — The canonical pair: the finite-dimension obstruction, the discrete Weyl pre-limit, and the three exact bounds}{BM.10 - The canonical pair: the finite-dimension obstruction, the discrete Weyl pre-limit, and the three exact bounds}.}\ Structural-constraints record: see Appendix B, Section B.3.5 of the Addendum.\par\medskip

\begin{theorem}[{\texorpdfstring{BM.11 — The continuum limit of the Janus topology: the differential conservation law and its two agreeing routes}{BM.11 - The continuum limit of the Janus topology: the differential conservation law and its two agreeing routes}}]\label{tag:BM-11}
\textbf{Thm.}

%dsp-level:0
\[
(1)\quad\textbf{the discrete divergence theorem carries no correction term:}
\]

%dsp-level:0
\[
\qquad \sum_{j=0}^{n-1}\ell\,(D_\ell f)(x_0+j\ell)=f(x_0+n\ell)-f(x_0)
\]

%dsp-level:0
\[
\qquad \text{EXACTLY on every rational grid}\ [\text{machine:}\ f=x^k,\ k\le5;\ n\le6;\ (x_0,\ell)\in\{(0,1),(\tfrac13,\tfrac25),(-\tfrac74,\tfrac12)\}];
\]

%dsp-level:0
\[
\text{the same total is the exact continuum integral of the derivative:}
\]

%dsp-level:0
\[
\qquad \sum_j\ell\,(D_\ell f)(x_j)=\int_{x_0}^{x_0+n\ell}\!f'(x)\,dx
\]

%dsp-level:0
\[
\qquad [\text{machine: exact equality at every}\ \ell,\ \neg\text{a limit statement}];
\]

%dsp-level:0
\[
\therefore\ \ell\to0\ \text{changes no term of the conservation identity}\ [\mathrm{BM.8}(1);\ \mathrm{CT2.6}];
\]

%dsp-level:0
\[
(2)\quad\textbf{the region law:}
\]

%dsp-level:0
\[
\qquad \text{for a materialized chain}\ c\ \text{and every sub-interval}\ [i,j]:
\]

%dsp-level:0
\[
\qquad \sum_{e\in[i,j]}j(e)=\frac{\sigma_{v_i}-\sigma_{v_j}}2=\text{the boundary read}
\]

%dsp-level:0
\[
\qquad [\text{machine: all}\ (\text{chain},\text{interval})\ \text{pairs},\ k\le5];
\]

%dsp-level:0
\[
\text{the full-chain case is}\ q_{\mathrm{EM}}(c)\ [\mathrm{BM.7}(2)];
\]

%dsp-level:0
\[
\qquad \text{the closed case reads}\ 0\ [\mathrm{D1.4};\ \mathrm{EM.11}:\ U\ \text{valid}\implies Q_{\mathrm{EM}}=\sum_{c\in U}q_{\mathrm{EM}}(c)=0];
\]

%dsp-level:0
\[
\textbf{the two-dimensional form:}
\]

%dsp-level:0
\[
\qquad \sum_{\text{cells}\in R}\bigl(\nabla_{\!\ell}\!\cdot j\bigr)=\text{the boundary flux of}\ R
\]

%dsp-level:0
\[
\qquad [\text{machine: every rectangle of a}\ 6\times5\ \text{rational lattice}];
\]

%dsp-level:0
\[
\textbf{the closed-network form:}
\]

%dsp-level:0
\[
\qquad \text{on a periodic lattice the total divergence is exactly}\ 0\ [\text{machine}]:
\]

%dsp-level:0
\[
\qquad \text{the discrete image of}\ \#J^+=\#J^-\ \text{with no boundary term}\ [\mathrm{D1.3},\ \mathrm{D2.11}];
\]

%dsp-level:0
\[
\text{the seal excess summed over every completed admissible history of length}\
\]

%dsp-level:0
\[
\le 12\ \text{is exactly}\ 0\ [\mathrm{BM.7}(1);\ \text{machine}];
\]

%dsp-level:0
\[
(3)\quad\textbf{the limit and the pointwise law:}
\]

%dsp-level:0
\[
\qquad \text{under}\ \mathrm{CT2.6}(\mathrm{ii})\text{–}(\mathrm{iii})\ [\text{extra-axiomatic, declared at the point of use — A §0.2a}]\ \text{the cells carry unboundedly many real edges and the coefficient fields are of bounded variation}:
\]

%dsp-level:0
\[
\sum_{\text{cells}\in U}(\cdots)\longrightarrow\int_U\!\sqrt{-g}\,\partial_\mu J^\mu\,d^4x,\qquad\text{boundary flux}\longrightarrow\oint_{\partial U}\!\sqrt{-g}\,J^\mu n_\mu\,d^3\Sigma;
\]

%dsp-level:0
\[
\text{the interval-sum moment matrix is nonsingular exactly for polynomial densities of degree}\
\]

%dsp-level:0
\[
\le 4\ [\text{machine: exact determinants}]\implies\text{the interval sums determine the density coefficients uniquely};
\]

%dsp-level:0
\[
\text{since}\ (2)\ \text{holds for every region, the fundamental lemma of the calculus of variations}\ [\text{math}]\ \text{gives}
\]

%dsp-level:0
\[
\qquad \boxed{\ \partial_\mu J^\mu=0\ }
\]

%dsp-level:0
\[
\qquad \text{pointwise};
\]

%dsp-level:0
\[
\textbf{the contrapositive is witnessed:}\quad\text{every nonzero polynomial density of degree}\
\]

%dsp-level:0
\[
\le 3\ \text{has a nonzero interval sum on a rational grid}\ [\text{machine}];
\]

%dsp-level:0
\[
\textbf{the current dictionary:}
\]

%dsp-level:0
\[
\qquad J^0=\text{the seal density}\ (\#J^+-\#J^-\ \text{per cell}=n_{\mathcal R}\ \text{density}\ [\mathrm{CT0.3}]);
\]

%dsp-level:0
\[
\qquad J^i=\text{the hop current on the}\ w_i\ \text{carrier}\ [\mathrm{D2.1};\ \mathrm{CT1.12}];
\]

%dsp-level:0
\[
(4)\quad\textbf{the stress-energy conservation, by two independent routes that agree:}
\]

%dsp-level:0
\[
\qquad \textbf{(i) geometry:}\ \nabla^\mu G_{\mu\nu}
\]

%dsp-level:0
\[
= 0\ \text{on the exact pp-wave witness}\ [\text{machine: nonlinear, all}\ 4\ \text{components, as polynomials in}\ u,v,x,y];
\]

%dsp-level:0
\[
g_{\mu\nu}:\ g_{uu}=H,\ g_{uv}=g_{vu}=-1,\ g_{xx}=g_{yy}=1\implies\det g=-1\ \text{exactly, independent of}\ H\ [\text{machine}];
\]

%dsp-level:0
\[
\qquad \Gamma^\mu_{\mu\lambda}=\partial_\lambda\ln\sqrt{|g|}=0\ [\text{machine}];
\]

%dsp-level:0
\[
\nabla_\lambda g_{\mu\nu}=0\ \text{on all}\ 64\ \text{components}\ [\text{machine}]\implies\nabla_\mu(\Lambda g^{\mu\nu})=0\ \text{on all}\ 4\ [\text{machine}]:
\]

%dsp-level:0
\[
\qquad \mathrm{CT1.23}\ \text{forces}\ \kappa^2\nabla_\mu T^{\mu\nu}=0;
\]

%dsp-level:0
\[
R_{\mu\nu}=-\tfrac12(\partial_x^2+\partial_y^2)H\,\delta^u_\mu\delta^u_\nu,\quad R=0,\quad G_{uu}=R_{uu}\ [\text{machine}];
\]

%dsp-level:0
\[
\qquad \Delta_\perp H=0\implies R_{\mu\nu}=0:\ \text{the vacuum pp-wave}\ [\text{machine}];
\]

%dsp-level:0
\[
\textbf{(ii) Noether:}\ \partial^\mu T_{\mu\nu}
\]

%dsp-level:0
\[
= (\Box\phi-m^2\phi)\partial_\nu\phi\ [\mathrm{BM.9}(3)]\ \text{on an exact on-shell witness}\ (\phi=t,\ m=0:\ \Box\phi=0):
\]

%dsp-level:0
\[
\qquad \partial^\mu T_{\mu\nu}=0\ [\text{machine}];
\]

%dsp-level:0
\[
\textbf{the dust read:}
\]

%dsp-level:0
\[
\qquad \partial_\mu T^{\mu\nu}
\]

%dsp-level:0
\[
= (\partial_0\rho,\,0,\,0,\,0)\ \text{for the static pressureless dust of}\ \mathrm{BM.9}(4):\ \text{the three spatial derivatives drop out at}\ p
\]

%dsp-level:0
\[
= 0\ [\text{machine}];
\]

%dsp-level:0
\[
\text{a static inhomogeneous dust is exactly conserved};
\]

%dsp-level:0
\[
\qquad \text{a time-dependent}\ \rho\ \text{leaves the exact residual}\ \partial_0\rho\ [\text{machine}];
\]

%dsp-level:0
\[
(5)\quad\textbf{the priority clause:}
\]

%dsp-level:0
\[
\qquad (1)\text{–}(3)\ \text{derive the differential law from the integer law}: \partial_\mu J^\mu
\]

%dsp-level:0
\[
= 0\ \text{is the density read of}\ \#J^+
\]

%dsp-level:0
\[
= \#J^-\ [\mathrm{BM.7}(5)];
\]

%dsp-level:0
\[
\text{the converse fails at finite cell size: the integer balance constrains the completed-history census}\ [\mathrm{BM.7}(1):\ \mathrm{Cat}_m,\ \text{the interleaving counts}]\ \text{and the density law does not};
\]

%dsp-level:0
\[
\text{no variational principle is used anywhere in}\ (1)\text{–}(3):\ \text{the input is}\ \mathrm{D1.2}\text{'s transition domains,}\ \mathrm{D1.4}\text{'s boundary map and}\ \mathrm{D2.11}\text{'s validity condition}.
\]

\textbf{Rmk (limit discipline).} Clause (3) is the tier's only use of the extra-axiomatic continuum-limit conditions, inherited at the point of use in the manner of A §0.2a; the discrete rows (1), (2) and (5) are unconditional and hold at every $\mathfrak h\ne0$. The pp-wave witness of clause (4) is a rational-coometric test object, not a claim about the corpus's own emergent metric: it verifies the identity $\nabla^\mu G_{\mu\nu}=0$, which is metric-independent [math].

\textbf{Req.} D1.2, D1.3, D1.4, D2.1, D2.11 (Parts 1–2); CT0.3 (Part 1); CT1.12, CT1.22, CT1.23, CT2.6 (Part 2; CT2.6's conditions (ii)–(iii) at the point of use — A §0.2a); CT4.5, EM.11 (Part 2); BM.7(1), BM.7(2), BM.7(5), BM.8(1), BM.9(3), BM.9(4) (this part); exact rational grid telescoping, exact moment-matrix determinants, exact multivariate pp-wave curvature and covariant-derivative algebra, the fundamental lemma of the calculus of variations [math; machine: \texttt{checks/\allowbreak{}verify\_\allowbreak{}H\_\allowbreak{}noether.py} §E].

\textbf{Pf.}

%dsp-level:0
\[
(1):\ \ell(D_\ell f)(x_j)
\]

%dsp-level:0
\[
= f(x_j+\ell)-f(x_j)\ [\mathrm{BM.8}(1)]\implies\text{the sum telescopes to}\ f(x_0+n\ell)-f(x_0)\ [\text{machine, exact rationals}];
\]

%dsp-level:0
\[
\qquad \int_a^b f'(x)\,dx=f(b)-f(a)\ \text{for}\ f=x^k\ [\text{the exact polynomial antiderivative}\ \tfrac{x^{k+1}}{k+1}];
\]

%dsp-level:0
\[
\qquad \text{the two totals are the same rational number at every}\ \ell\ [\text{machine}]
\]

%dsp-level:0
\[
(2):\ \text{pairwise cancellation of the interior signs}\ [\mathrm{BM.7}(2)\ \text{re-run per interval};\ \text{machine: exhaustive}];
\]

%dsp-level:0
\[
\qquad \text{the}\ 2\mathrm D\ \text{lattice}: \sum_{i,k}[(j^x_{i+1,k}-j^x_{i,k})+(j^y_{i,k+1}-j^y_{i,k})]\ \text{telescopes in each direction separately to the four boundary sums}\ [\text{machine: all rectangles}];
\]

%dsp-level:0
\[
\qquad \text{periodic identification cancels the boundary sums pairwise}\ [\text{machine}]
\]

%dsp-level:0
\[
(3):\ \text{cell}\to\text{chart convergence under}\ \mathrm{CT2.6}(\mathrm{ii})\text{–}(\mathrm{iii})\ [\text{math}];
\]

%dsp-level:0
\[
\qquad \text{moment matrix}: M_{b,m}
\]

%dsp-level:0
\[
= \sum_{j=0}^{b-1}\ell\,
\]

%dsp-level:0
\[
(j\ell)^m,
\]

%dsp-level:0
\[
\ b
\]

%dsp-level:0
\[
= 1..d{+}1:\ \det M
\]

%dsp-level:0
\[
\ne 0\ [\text{machine, exact Fractions}]\implies\text{injectivity of the coefficient map};
\]

%dsp-level:0
\[
\qquad \text{vanishing of all interval sums}\implies\text{all coefficients zero}\implies\partial_\mu J^\mu=0\ [\text{fundamental lemma};\ \text{math}]
\]

%dsp-level:0
\[
(4)\text{(i)}:\ \Gamma^\rho_{\mu\nu}
\]

%dsp-level:0
\[
= \tfrac12g^{\rho\sigma}(\partial_\mu g_{\nu\sigma}+\partial_\nu g_{\mu\sigma}-\partial_\sigma g_{\mu\nu}),
\]

%dsp-level:0
\[
\ R^\rho_{\sigma\mu\nu}
\]

%dsp-level:0
\[
= \partial_\mu\Gamma^\rho_{\nu\sigma}-\partial_\nu\Gamma^\rho_{\mu\sigma}+\Gamma^\rho_{\mu\lambda}\Gamma^\lambda_{\nu\sigma}-\Gamma^\rho_{\nu\lambda}\Gamma^\lambda_{\mu\sigma},
\]

%dsp-level:0
\[
\ R_{\sigma\nu}
\]

%dsp-level:0
\[
= R^\mu_{\sigma\mu\nu}\ [\text{math}];
\]

%dsp-level:0
\[
\qquad \text{all steps exact in}\ \mathbb Q[u,v,x,y]\ [\text{machine}];
\]

%dsp-level:0
\[
\qquad \nabla_\mu G^{\mu\nu}=\partial_\mu G^{\mu\nu}+\Gamma^\mu_{\mu\lambda}G^{\lambda\nu}+\Gamma^\nu_{\mu\lambda}G^{\mu\lambda}=0\ [\text{machine}]
\]

%dsp-level:0
\[
(4)\text{(ii)}:\ \mathrm{BM.9}(3)\text{'s jet identity evaluated at}\ \phi
\]

%dsp-level:0
\[
= t,
\]

%dsp-level:0
\[
\ m
\]

%dsp-level:0
\[
= 0\ [\Box\phi=g^{\mu\nu}\partial_\mu\partial_\nu t=0]:\ \text{eom}
\]

%dsp-level:0
\[
= 0\ \text{and}\ \partial^\mu T_{\mu\nu}
\]

%dsp-level:0
\[
= 0\ [\text{machine}];
\]

%dsp-level:0
\[
\qquad \text{dust}: T^{00}=\rho(x),\ T^{0i}=T^{ij}=0\implies\partial_\mu T^{\mu0}=\partial_0\rho,\ \partial_\mu T^{\mu i}=0\ [\text{machine}]
\]

%dsp-level:0
\[
(5):\ \text{conjunction of}\ (1)\text{–}(3)\ \text{with}\ \mathrm{BM.7}(1)\text{'s census rows}.\ \blacksquare
\]

\textbf{Phys.} The continuum conservation law is a limit of a telescoping sum, and the telescoping is exact — not asymptotic. Sum the per-edge source over any stretch of a materialized chain and the interior signs cancel pairwise, leaving only the two ends; do it on a lattice rectangle and the same cancellation leaves only the boundary flux; close the network and the total is zero, which is the integer balance wearing a density's clothes. Because the interval sums determine a polynomial density uniquely — the moment matrix has an exact nonzero determinant — vanishing on every region forces the density itself to vanish, and $\partial_\mu J^\mu=0$ follows pointwise. The stress-energy tensor is conserved twice over and the two derivations agree: geometrically, because the Einstein tensor is divergence-free (verified nonlinearly on an exact rational pp-wave, where $\det g=-1$ identically and metric compatibility kills the cosmological term), and dynamically, because the jet identity of the translation current vanishes on-shell. Static dust stays conserved however lumpy it is; only $\partial_0\rho$ can spoil it. And the arrow of explanation runs one way: the integers ground the density, never the reverse.

\end{theorem}

\medskip\noindent\textit{\texorpdfstring{BM.12 — The Noether status of the corpus: adjudication and scope}{BM.12 - The Noether status of the corpus: adjudication and scope}.}\ Structural-constraints record: see Appendix B, Section B.1.15 of the Addendum.\par\medskip

\phantomsection
\section[Tier BM.5 — The Calibration Lock: the Born Width, the…]{\texorpdfstring{Tier BM.5 — The Calibration Lock: the Born Width, the Coherence Volume, the Curvature of a Sealed Composite, and the Neutron Row}{Tier BM.5 - The Calibration Lock: the Born Width, the Coherence Volume, the Curvature of a Sealed Composite, and the Neutron Row}}

\begin{definition}[{\texorpdfstring{Def-BM.4 — The Born width and the coherence volume}{Def-BM.4 - The Born width and the coherence volume}}]\label{tag:Def-BM-4}
\textbf{Def.} Let $E_{\mathrm{int}}\in\mathbb Z_{\ge2}$ be the internal directed-edge count of a sealed composite [D0.8; FT.44(4)]: $E_{\mathrm{int}}=2E_{\mathrm{bidir}}+E_{\mathrm{uni}}$; the minimal floor is the digon, $E_{\mathrm{int}}=2$ [FT.28] and $m_{\mathrm{hid}}=E_{\mathrm{int}}\ln2/\mathfrak h$ its informational mass [FT.44(4), FT.42(1)].

%dsp-level:0
\[
(1)\quad\textbf{the Born width:}
\]

%dsp-level:0
\[
\qquad \Delta x_{\mathrm B}\ :=\ \frac1{m_{\mathrm{hid}}}\ =\ \frac{\mathfrak h}{E_{\mathrm{int}}\ln2}
\]

%dsp-level:0
\[
\qquad [\hbar=c=1,\ \mathrm{D2.9};\ \text{the Compton-like length of}\ \mathrm{FT.90}(6)\ \text{re-cited as an existing read}];
\]

%dsp-level:0
\[
\textbf{constitutive clause:}\ \ \Delta x_{\mathrm B}\cdot m_{\mathrm{hid}}=1\ \text{is the definition};
\]

%dsp-level:0
\[
\qquad \textbf{non-derived clause:}\ \ \text{the dimensionless normalization is fixed to}\ 1\ \text{(the Compton normalization),}\ \neg\text{derived};
\]

%dsp-level:0
\[
(2)\quad\textbf{the dimensional-uniqueness clause:}
\]

%dsp-level:0
\[
\qquad \exists!\ \text{free dimensionful scale}\ \mathfrak h\ [\mathrm{P1},\ \mathrm{CT2.11}]:\ \ \forall Q\ \text{dimensionful}:\ Q
\]

%dsp-level:0
\[
= c_Q\,
\]

%dsp-level:0
\[
\mathfrak h^{k_Q},
\]

%dsp-level:0
\[
\ c_Q\ \text{fixed}\ [\mathrm{CT2.11}];
\]

%dsp-level:0
\[
[\text{length}]=\mathfrak q^{-1/2}=[\mathfrak h],
\]

%dsp-level:0
\[
\qquad [m_{\mathrm{hid}}]=\mathfrak q^{+1/2}\ [\mathrm{CT2.1}]\implies\text{every length built from}\ m_{\mathrm{hid}}\ \text{alone is}\ c/m_{\mathrm{hid}}\ \text{with}\ c\ \text{dimensionless}:
\]

%dsp-level:0
\[
\text{the Born width is the unique such length at}\ c=1;
\]

%dsp-level:0
\[
\qquad \text{any other}\ c\ \text{is a P4-class internal choice and is priced, not adopted}\ [\text{clause}\ (4)];
\]

%dsp-level:0
\[
(3)\quad\textbf{the coherence volume:}\qquad V_{\mathrm{eff}}\ :=\ (\Delta x_{\mathrm B})^3\ =\ \frac{\mathfrak h^3}{(E_{\mathrm{int}}\ln2)^3}\qquad[\alpha=1:\ \text{the cube convention}];
\]

%dsp-level:0
\[
(4)\quad\textbf{the geometric-factor clause (priced, not adopted):}
\]

%dsp-level:0
\[
\qquad V=\alpha\,(\Delta x_{\mathrm B})^3\ \text{with}\ \alpha\ \text{dimensionless}:
\]

%dsp-level:0
\[
\qquad \rho(\alpha)=\rho(1)/\alpha,
\]

%dsp-level:0
\[
\qquad \alpha=\tfrac{4\pi}3:\ \rho\mapsto\tfrac3{4\pi}\rho\ \ \text{exactly};
\]

%dsp-level:0
\[
\text{every ratio in which}\ \alpha\ \text{appears in both entries is}\ \alpha\text{-free};
\]

%dsp-level:0
\[
\qquad \alpha\ \text{is}\ \neg\text{a dimensionful datum and}\ \neg\text{a P4 core member};
\]

%dsp-level:0
\[
(5)\quad\textbf{the locus clause:}
\]

%dsp-level:0
\[
\qquad \Delta x_{\mathrm B}\ \text{is a chart-side envelope on}\ M\ [\mathrm{CT2.6}]\ \text{in the sense of BM.10(6)'s width law},
\]

%dsp-level:0
\[
\ \neg\text{a materialized edge length}:
\]

%dsp-level:0
\[
\text{the minimal materialized line element is}\ \widehat{ds}^{\,2}=1\implies ds^2=\mathfrak h^2\ [\mathrm{D1.6};\ \mathrm{CT1.19a}]:
\]

%dsp-level:0
\[
\qquad \Delta x_{\mathrm B}<\mathfrak h\iff E_{\mathrm{int}}\ln2>1\iff E_{\mathrm{int}}\ge2\ [\text{forward pointer:}\ \mathrm{BM.13}(3);\ \text{machine}];
\]

%dsp-level:0
\[
\text{every admissible cell is sub-hop, so the width is a coarse-graining envelope and}\ \neg\text{a ledger distance}\ [\text{locus discipline},\ \mathrm{FT.136}(5)\ \text{pattern}];
\]

%dsp-level:0
\[
(6)\quad\textbf{the scope clause:}
\]

%dsp-level:0
\[
\qquad \mathrm{FT.90}(6)\ \text{refutes a Compton-wavelength relation to the generation threshold}\ (\delta h\ge\lceil\log_2\#\text{classes}\rceil\ \text{depends on the class count alone}):
\]

%dsp-level:0
\[
\text{that refutation is untouched here}\ —\ \text{this definition uses the length as a spread read, and derives no relation between it and any threshold}.
\]

\textbf{Rmk (naming).} \textit{Born width} names the spatial spread of the sealed composite's density read at the Compton normalization; \textit{coherence volume} names its cube. Both are internal bookkeeping on derived quantities [A §0.2]; neither is a measurement, and no external length datum enters [P5; Def-BM.3(6)].

\textbf{Req.} D0.6, D0.8, D2.9 (Parts 1–2); D1.6 (Part 1); CT1.19a, CT2.1, CT2.6, CT2.11 (Part 2); P1, P5 (Parts 1–2); Def-BM.3 (this part); FT.28, FT.42(1), FT.44(4) (Part 3, Census); FT.90(6) (Part 3); BM.10(6) (this part); the sub-hop row cited in clause (5) is a marked forward pointer at the citing site and is imported in no \textbf{Req} line; exact Laurent-monomial arithmetic, exact rational $\ln2$ enclosures [math; machine: \texttt{checks/\allowbreak{}verify\_\allowbreak{}H\_\allowbreak{}calibration\_\allowbreak{}lock.py} §A].

\textbf{Pf.}

%dsp-level:0
\[
\text{Definition block: the constitutive clause is}\ \Delta x_{\mathrm B}\cdot m_{\mathrm{hid}}
\]

%dsp-level:0
\[
= 1;
\]

%dsp-level:0
\[
\ \text{the non-derived clauses are the normalization}\ c
\]

%dsp-level:0
\[
= 1\ \text{of}\ (1)\text{–}(2),
\]

%dsp-level:0
\[
\ \text{the geometric factor}\ \alpha\ \text{of}\ (4),
\]

%dsp-level:0
\[
\ \text{and the locus reading of}\ (5);
\]

%dsp-level:0
\[
\ \text{clause}\ (6)\ \text{is a scope record}.\ \blacksquare
\]

\textbf{Phys.} The sealed composite has one length and it is the reciprocal of its own informational mass: $\mathfrak h$ over the edge count times $\ln2$. That is not a new postulate — the corpus already carried the Compton-like length as an existing read, and with one free scale in the theory there is no other length a mass can build. The cube of that width is the volume the density spreads over, with a dimensionless geometric factor left free and priced rather than chosen. One clause keeps the reading honest: the width is a chart-side envelope, not a hop. Every admissible cell turns out to be \textit{shorter} than a single materialized line element, which is exactly why the width belongs to the continuum chart and not to the ledger.

\end{definition}

\begin{theorem}[{\texorpdfstring{BM.13 — The Born width: the exact substitution, the quantized ladder, and the sub-hop locus}{BM.13 - The Born width: the exact substitution, the quantized ladder, and the sub-hop locus}}]\label{tag:BM-13}
\textbf{Thm.}

%dsp-level:0
\[
(1)\quad\textbf{the exact substitution:}
\]

%dsp-level:0
\[
\qquad \Delta x_{\mathrm B}=\frac1{m_{\mathrm{hid}}}=\frac{\mathfrak h}{E_{\mathrm{int}}\ln2}
\]

%dsp-level:0
\[
\qquad [\mathrm{Def\text{-}BM.4}(1);\ \mathrm{FT.44}(4);\ \text{machine}:\ \Delta x_{\mathrm B}\cdot m_{\mathrm{hid}}=1\ \text{as a Laurent identity}];
\]

%dsp-level:0
\[
\text{the monomial census is exact:}\quad\Delta x_{\mathrm B}=\mathfrak h^1E_{\mathrm{int}}^{-1}(\ln2)^{-1}\pi^0
\]

%dsp-level:0
\[
\qquad [\text{machine: one term, coefficient}\ 1];
\]

%dsp-level:0
\[
[\Delta x_{\mathrm B}]=[\mathfrak h]=\mathfrak q^{-1/2}=[\text{length}]
\]

%dsp-level:0
\[
\qquad [\mathrm{CT2.1};\ \text{machine: exact exponent arithmetic}];
\]

%dsp-level:0
\[
\qquad E_{\mathrm{int}},\ \ln2:\ \text{dimensionless};
\]

%dsp-level:0
\[
(2)\quad\textbf{quantization:}
\]

%dsp-level:0
\[
\qquad \Delta x_{\mathrm B}\in\Bigl\{\frac{\mathfrak h}{n\ln2}:\ n\in\mathbb Z_{\ge2}\Bigr\}:
\]

%dsp-level:0
\[
\qquad \text{a discrete ladder},\ \neg\text{a continuous range}\ [\text{machine: injective in}\ n];
\]

%dsp-level:0
\[
\textbf{the computed cells:}
\]

%dsp-level:0
\[
\quad E_{\mathrm{int}}
\]

%dsp-level:0
\[
= 2\ [\Sigma_{\min};\ \mathrm{FT.28}],
\]

%dsp-level:0
\[
\ 6\ [\Sigma_{12};\ \text{Census FT.35}(2)],
\]

%dsp-level:0
\[
\ 12,
\]

%dsp-level:0
\[
\ 24,
\]

%dsp-level:0
\[
\ 24,
\]

%dsp-level:0
\[
\ 60,
\]

%dsp-level:0
\[
\ 60\ [\text{the five solids};\ \text{Census FT.42}(5)];
\]

%dsp-level:0
\[
\text{the ladder is strictly decreasing with disjoint exact enclosures}\ [\text{machine: rational}\ \ln2\ \text{enclosures of width}\ 2^{-12}];
\]

%dsp-level:0
\[
\qquad \text{the cycle ladder}\ E_{\mathrm{int}}=L\ [\mathrm{FT.32}]\ \text{has no repeated value}\ [\text{machine}];
\]

%dsp-level:0
\[
(3)\quad\textbf{the sub-hop locus:}
\]

%dsp-level:0
\[
\qquad \Delta x_{\mathrm B}<\mathfrak h\iff\frac1{E_{\mathrm{int}}\ln2}<1\iff E_{\mathrm{int}}\ln2>1\iff E_{\mathrm{int}}\ge2
\]

%dsp-level:0
\[
\qquad [\text{machine: exact enclosure}\ \tfrac12<\ln2<1\ \text{drives every step}];
\]

%dsp-level:0
\[
\text{the minimal admissible count is}\ E_{\mathrm{int}}
\]

%dsp-level:0
\[
= 2\ [\mathrm{FT.28}:\ \text{the digon floor};\ \text{a materialized closed chain carries}\ \ge2\ \text{directed edges}]:
\]

%dsp-level:0
\[
\therefore\ \textbf{every admissible sealed composite is sub-hop:}\quad\Delta x_{\mathrm B}\le\frac{\mathfrak h}{2\ln2}<\mathfrak h
\]

%dsp-level:0
\[
\qquad [\text{machine: all computed cells}];
\]

%dsp-level:0
\[
\text{at the inadmissible}\ E_{\mathrm{int}}
\]

%dsp-level:0
\[
= 1\ \text{the order reverses}\ (\Delta x_{\mathrm B}>\mathfrak h):\ \text{the threshold is sharp and sits exactly at the floor}\ [\text{machine}];
\]

%dsp-level:0
\[
\textbf{locus discipline:}\ \ \Delta x_{\mathrm B}\ \text{is a}\ \mathrm{CT2.6}\ \text{chart envelope},
\]

%dsp-level:0
\[
\ \neg\text{a ledger distance}\ [\mathrm{Def\text{-}BM.4}(5);\ \mathrm{D1.6}:\ \text{the minimal materialized element is}\ ds^2=\mathfrak h^2];
\]

%dsp-level:0
\[
(4)\quad\textbf{the Robertson consistency:}
\]

%dsp-level:0
\[
\qquad \Delta x\,\Delta p\ge\tfrac12\ [\mathrm{BM.10}(3),(6);\ \hbar=1]\ \text{at}\ \Delta x=\Delta x_{\mathrm B}\iff\Delta p\ge\frac1{2\Delta x_{\mathrm B}}=\frac{m_{\mathrm{hid}}}2
\]

%dsp-level:0
\[
\qquad [\text{machine}];
\]

%dsp-level:0
\[
\Delta p=\tfrac12m_{\mathrm{hid}}\ \text{saturates Robertson exactly at the Born width}:
\]

%dsp-level:0
\[
\qquad \text{the width is the saturating spread for a momentum spread of half the informational mass};
\]

%dsp-level:0
\[
(5)\quad\textbf{the}\ \alpha\textbf{-variant priced:}
\]

%dsp-level:0
\[
\qquad V=\alpha(\Delta x_{\mathrm B})^3\implies\rho(\alpha)=\rho(1)/\alpha;
\]

%dsp-level:0
\[
\qquad \alpha=\tfrac{4\pi}3\implies\rho\mapsto\tfrac3{4\pi}\rho\ \ \text{exactly}\ [\text{machine: a single Laurent monomial}].
\]

\textbf{Rmk (comparison discipline).} No clause reads $\Delta x_{\mathrm B}$ against an external length. The comparisons of this tag are internal: against the hop $\mathfrak h$ [D0.6], against the minimal materialized element [D1.6], and against the Robertson bound [BM.10(6)]. External comparison remains quarantined [P5; Census FT.114, FT.120].

\textbf{Req.} D0.6, D0.8, D1.6, D2.9 (Parts 1–2); CT1.19a, CT2.1, CT2.6 (Part 2); P1, P5 (Parts 1–2); Def-BM.4 (this part); FT.28, FT.32, FT.42(1), FT.44(4), FT.90(6) (Part 3, Census); Census FT.35(2), FT.42(5) (row-wise); BM.10(3), BM.10(6) (this part); exact Laurent-monomial identities, exact rational $\ln2$ enclosures via $\ln2=\sum_{n\ge1}1/(n2^n)$ with tail $<2^{-N}$ [math; machine: \texttt{checks/\allowbreak{}verify\_\allowbreak{}H\_\allowbreak{}calibration\_\allowbreak{}lock.py} §A].

\textbf{Pf.}

%dsp-level:0
\[
(1):\ m_{\mathrm{hid}}
\]

%dsp-level:0
\[
= E_{\mathrm{int}}\ln2/\mathfrak h\ [\mathrm{FT.44}(4)\ \text{re-cited}]
\]

%dsp-level:0
\[
\implies 1/m_{\mathrm{hid}}
\]

%dsp-level:0
\[
= \mathfrak h/(E_{\mathrm{int}}\ln2)\ [\text{machine: the product is}\ 1\ \text{in}\ \mathbb Q[\mathfrak h,E,\ln2]^{\pm}];
\]

%dsp-level:0
\[
\qquad \text{exponents}: [\mathfrak h]-0-0=\mathfrak q^{-1/2}\ [\mathrm{CT2.1}]
\]

%dsp-level:0
\[
(2):\ \Delta x_{\mathrm B}(n)
\]

%dsp-level:0
\[
= \mathfrak h/(n\ln2):\ n_1
\]

%dsp-level:0
\[
\ne n_2\implies\Delta x_{\mathrm B}(n_1)\ne\Delta x_{\mathrm B}(n_2)\ [\text{machine: injectivity on}\ n\in2..59];
\]

%dsp-level:0
\[
\qquad \text{the five computed cells}: \text{disjoint exact enclosures}\ [\text{machine}:\ \ln2\in(\textstyle\sum_1^{12}\tfrac1{n2^n},\ \cdot+2^{-12})]
\]

%dsp-level:0
\[
(3):\ \tfrac12<\ln2<1\ [\text{machine: the alternating-free series with exact rational partial sums}]
\]

%dsp-level:0
\[
\implies E\ln2>1
\]

%dsp-level:0
\[
\iff E
\]

%dsp-level:0
\[
\ge 2\ [\text{machine: every}\ E\in1..59];
\]

%dsp-level:0
\[
\qquad E_{\mathrm{int}}\ge2\ [\mathrm{FT.28}]\implies\Delta x_{\mathrm B}\le\mathfrak h/(2\ln2)<\mathfrak h
\]

%dsp-level:0
\[
(4):\ \Delta x\,
\]

%dsp-level:0
\[
\Delta p\ge\tfrac12\ [\mathrm{BM.10}(3):\ \text{Robertson};\ \mathrm{BM.10}(6):\ \text{the width law}]\ \text{with}\ \Delta x
\]

%dsp-level:0
\[
= \Delta x_{\mathrm B}
\]

%dsp-level:0
\[
= 1/m_{\mathrm{hid}}\implies\Delta p
\]

%dsp-level:0
\[
\ge m_{\mathrm{hid}}/2\ [\text{machine: exact rationals}];
\]

%dsp-level:0
\[
\qquad \text{equality at}\ \Delta p=m_{\mathrm{hid}}/2\ [\mathrm{BM.10}(3)\ \text{saturation row}]
\]

%dsp-level:0
\[
(5):\ \rho=m_{\mathrm{hid}}/(\alpha\Delta x_{\mathrm B}^3)=\rho(1)/\alpha\ [\text{machine}];
\]

%dsp-level:0
\[
\qquad \alpha=4\pi/3:\ 3/(4\pi)\ \text{is a single Laurent monomial}\ [\text{machine}].\ \blacksquare
\]

\textbf{Phys.} The width of a sealed composite is a reciprocal integer count times the hop, divided by $\ln2$ — one Laurent monomial, no mixture of scales, nothing to tune. The ladder it walks is strictly discrete: two, six, twelve, twenty-four, sixty, each step shorter than the last with exact rational enclosures that never touch. And every rung of it sits \textit{below} a single hop, which is the sharpest locus statement in the tier: the Born width is not a distance the ledger can materialize, it is the envelope the continuum chart assigns to a bound whose internal edges all cancel pairwise. Consistency with the uncertainty bound is exact rather than approximate — half the informational mass is precisely the momentum spread that saturates Robertson at this width.

\end{theorem}

\begin{theorem}[{\texorpdfstring{BM.14 — The coherence volume and the energy density}{BM.14 - The coherence volume and the energy density}}]\label{tag:BM-14}
\textbf{Thm.}

%dsp-level:0
\[
(1)\quad\textbf{the exact volume:}
\]

%dsp-level:0
\[
\qquad V_{\mathrm{eff}}=(\Delta x_{\mathrm B})^3=\frac{\mathfrak h^3}{(E_{\mathrm{int}}\ln2)^3}
\]

%dsp-level:0
\[
\qquad [\mathrm{Def\text{-}BM.4}(3);\ \text{machine: one Laurent monomial}\ \mathfrak h^3E^{-3}(\ln2)^{-3}];
\]

%dsp-level:0
\[
[V_{\mathrm{eff}}]=\mathfrak q^{-3/2}=[\text{length}]^3\qquad[\mathrm{CT2.1};\ \text{machine:}\ 3\cdot(-\tfrac12)=-\tfrac32];
\]

%dsp-level:0
\[
(2)\quad\textbf{the exact energy density:}
\]

%dsp-level:0
\[
\qquad \rho=\frac{m_{\mathrm{hid}}}{V_{\mathrm{eff}}}=\frac{E_{\mathrm{int}}\ln2/\mathfrak h}{\mathfrak h^3/(E_{\mathrm{int}}\ln2)^3}=\frac{(E_{\mathrm{int}}\ln2)^4}{\mathfrak h^4}
\]

%dsp-level:0
\[
\qquad [\text{machine: Laurent identity, denominators cleared}];
\]

%dsp-level:0
\[
\rho\,\mathfrak h^4=(E_{\mathrm{int}}\ln2)^4\ \ \text{exactly};
\]

%dsp-level:0
\[
\qquad [\rho]=[\text{energy}]-[V]=\mathfrak q^{1/2+3/2}=\mathfrak q^{+2}=[T_{\mu\nu}]
\]

%dsp-level:0
\[
\qquad [\mathrm{CT2.1};\ \mathrm{BM.9}(5);\ \text{machine}];
\]

%dsp-level:0
\[
(3)\quad\textbf{the agreement with the general source read:}
\]

%dsp-level:0
\[
\qquad \kappa^2\rho=\frac{2\pi\,\mathfrak h\,E_{\mathrm{int}}}{V}\bigg|_{V=V_{\mathrm{eff}}}=\frac{2\pi\,E_{\mathrm{int}}^4(\ln2)^3}{\mathfrak h^2}
\]

%dsp-level:0
\[
\qquad [\mathrm{BM.9}(5);\ \text{machine: the two forms are the same monomial}];
\]

%dsp-level:0
\[
\kappa^2=8\pi G_N=\frac{2\pi\,\mathfrak h^2}{\ln2},\qquad G_N=\frac{\mathfrak h^2}{4\ln2}\qquad[\mathrm{CT1.23},\ \mathrm{CT2.9},\ \mathrm{P3};\ \text{machine: exact}];
\]

%dsp-level:0
\[
(4)\quad\textbf{the}\ \ln2\ \textbf{ledger:}
\]

%dsp-level:0
\[
\qquad m_{\mathrm{hid}}:\ (\ln2)^1\ [\mathrm{FT.44}(4)];
\]

%dsp-level:0
\[
\qquad V_{\mathrm{eff}}:\ (\ln2)^{-3};
\]

%dsp-level:0
\[
\qquad \rho:\ (\ln2)^4=(\ln2)^1\cdot(\ln2)^3;
\]

%dsp-level:0
\[
\qquad \kappa^2:\ (\ln2)^{-1};
\]

%dsp-level:0
\[
\qquad \kappa^2\rho:\ (\ln2)^3\cdot2\pi;
\]

%dsp-level:0
\[
\text{at a general}\ V\ [\mathrm{BM.9}(5)]\ \text{the}\ \ln2\ \text{of}\ m_{\mathrm{hid}}\ \text{and of}\ \kappa^2\ \text{cancel and no}\ \ln2\ \text{survives};
\]

%dsp-level:0
\[
\qquad \text{at}\ V=V_{\mathrm{eff}}\ \text{the volume contributes}\ (\ln2)^{-3}\ \text{and three powers survive}:
\]

%dsp-level:0
\[
\text{both statements are the same exact identity read at two volumes}\ [\text{machine: the monomial censuses agree}];
\]

%dsp-level:0
\[
(5)\quad\textbf{the}\ \alpha\textbf{-variant:}
\]

%dsp-level:0
\[
\qquad \rho(\alpha)=\frac{(E_{\mathrm{int}}\ln2)^4}{\alpha\,\mathfrak h^4},
\]

%dsp-level:0
\[
\qquad \kappa^2\rho(\alpha)=\frac{2\pi E_{\mathrm{int}}^4(\ln2)^3}{\alpha\,\mathfrak h^2};
\]

%dsp-level:0
\[
\qquad \alpha=\tfrac{4\pi}3\implies\kappa^2\rho\mapsto\frac{3\,E_{\mathrm{int}}^4(\ln2)^3}{2\,\mathfrak h^2}\ \ [\text{machine: the}\ \pi\ \text{cancels exactly}].
\]

\textbf{Rmk (volume convention).} Clause (5) prices the spherical convention instead of adopting it: at $\alpha=4\pi/3$ the $\pi$ of $\kappa^2$ cancels and the curvature scale becomes $\pi$-free. Both conventions are exact; the cube convention of Def-BM.4(3) is the one carried downstream, and every $\alpha$-free ratio is unaffected.

\textbf{Req.} D2.9 (Parts 1–2); CT1.23, CT2.1, CT2.9 (Part 2); P3 (Parts 1–2); Def-BM.4 (this part); FT.42(1), FT.44(4) (Part 3, Census); BM.9(5), BM.13(1) (this part); exact Laurent-monomial arithmetic with denominators cleared [math; machine: \texttt{checks/\allowbreak{}verify\_\allowbreak{}H\_\allowbreak{}calibration\_\allowbreak{}lock.py} §B].

\textbf{Pf.}

%dsp-level:0
\[
(1):\ \text{cubing}\ \mathrm{BM.13}(1)\text{'s monomial}\ [\text{machine}];\qquad3\cdot(-\tfrac12)=-\tfrac32\ [\mathrm{CT2.1}]
\]

%dsp-level:0
\[
(2):\ \rho=\frac{E\ln2}{\mathfrak h}\cdot\frac{(E\ln2)^3}{\mathfrak h^3}=\frac{(E\ln2)^4}{\mathfrak h^4}\ [\text{machine}];
\]

%dsp-level:0
\[
\qquad [\rho]=\mathfrak q^{1/2-(-3/2)}=\mathfrak q^{+2}=[T_{\mu\nu}]\ [\mathrm{CT2.1};\ \mathrm{BM.9}(5)\ \text{row}]
\]

%dsp-level:0
\[
(3):\ \kappa^2\rho=\tfrac{2\pi\mathfrak h^2}{\ln2}\cdot\tfrac{(E\ln2)^4}{\mathfrak h^4}=\tfrac{2\pi E^4(\ln2)^3}{\mathfrak h^2}\ [\text{machine}];
\]

%dsp-level:0
\[
\qquad \mathrm{BM.9}(5):\ \kappa^2\rho
\]

%dsp-level:0
\[
= 2\pi\mathfrak h E/V\ \text{at}\ V
\]

%dsp-level:0
\[
= \mathfrak h^3/(E\ln2)^3\implies\tfrac{2\pi\mathfrak h E(E\ln2)^3}{\mathfrak h^3}
\]

%dsp-level:0
\[
= \tfrac{2\pi E^4(\ln2)^3}{\mathfrak h^2}\ [\text{machine: identical monomials}]
\]

%dsp-level:0
\[
(4):\ \text{monomial censuses}\ [\text{machine}]:\ m_{\mathrm{hid}}\to(0,1,1,0),
\]

%dsp-level:0
\[
\ V_{\mathrm{eff}}\to(3,-3,-3,0),
\]

%dsp-level:0
\[
\ \rho\to(-4,4,4,0),
\]

%dsp-level:0
\[
\ \kappa^2\to(2,0,-1,1),
\]

%dsp-level:0
\[
\ \kappa^2\rho\to(-2,4,3,1)\ \text{in the basis}\ (\mathfrak h,E,\ln2,\pi)
\]

%dsp-level:0
\[
(5):\ \text{division by}\ \alpha\ [\text{machine}];\qquad\alpha=4\pi/3:\ (2\pi)/(4\pi/3)=3/2\ \text{exactly}.\ \blacksquare
\]

\textbf{Phys.} Cube the width and the density follows by one exact division: the fourth power of the information content over the fourth power of the hop. The dimensions land precisely on the stress-energy tensor's own, with no adjustment and no second scale. The same quantity arrives a second way — through the general source read of the stress-energy tier with the volume now supplied by the width — and the two routes give the identical Laurent monomial, which is the first sign that the calibration is closing rather than accumulating. The $\ln2$ bookkeeping is explicit about where its powers live: one from the mass, three from the volume, minus one from the coupling, three surviving in the curvature. And the geometric factor is priced rather than hidden: choose the sphere and the $\pi$ of the coupling cancels outright.

\end{theorem}

\begin{theorem}[{\texorpdfstring{BM.15 — The curvature of a sealed composite: the exact trace of the field equations, and the curvature scale}{BM.15 - The curvature of a sealed composite: the exact trace of the field equations, and the curvature scale}}]\label{tag:BM-15}
\textbf{Thm.} In the orthonormal rest frame $g_{\mu\nu}=\mathrm{diag}(-1,1,1,1)$ [CT1.12]: $\operatorname{sig}=(-,+,+,+)$; [A3, C13], with $T_{\mu\nu}=\mathrm{diag}(\rho,p,p,p)$ [BM.9(4)] and $G_{\mu\nu}+\Lambda g_{\mu\nu}=\kappa^2T_{\mu\nu}$ [CT1.23]:

%dsp-level:0
\[
(1)\quad\textbf{the dust trace:}
\]

%dsp-level:0
\[
\qquad T=g^{\mu\nu}T_{\mu\nu}=-\rho+3p
\]

%dsp-level:0
\[
\qquad [\text{machine: exact component algebra}];
\]

%dsp-level:0
\[
\qquad p=0\ [\mathrm{CT4.5}]\implies T=-\rho;
\]

%dsp-level:0
\[
(2)\quad\textbf{the exact trace of the field equations:}\qquad G^\mu{}_\mu=-R\ \ \text{in}\ 4\ \text{dimensions}\ [\text{math}]\implies-R+4\Lambda=\kappa^2T:
\]

%dsp-level:0
\[
\boxed{\ R=\kappa^2(\rho-3p)+4\Lambda\ }\qquad\text{and at}\ p=0:\qquad R=\kappa^2\rho+4\Lambda\qquad[\text{machine: exact in}\ \mathbb Q[\rho,p,\Lambda,\kappa^2]];
\]

%dsp-level:0
\[
\textbf{the inadmissible form:}\quad R=-\kappa^2\rho\ \text{would require}\ 2\kappa^2\rho-3\kappa^2p+4\Lambda=0\ \text{identically}:
\]

%dsp-level:0
\[
\qquad \textbf{refuted exactly}\ [\text{machine: the residual is a nonzero polynomial}];
\]

%dsp-level:0
\[
\text{the sign of the dust curvature is fixed by CT1.23's own trace}: \ \kappa^2>0,
\]

%dsp-level:0
\[
\ \rho
\]

%dsp-level:0
\[
\ge 0\ [\mathrm{CT4.5}]
\]

%dsp-level:0
\[
\implies R
\]

%dsp-level:0
\[
\ge 4\Lambda,
\]

%dsp-level:0
\[
\ \text{and}\ R
\]

%dsp-level:0
\[
= \kappa^2\rho>0\ \text{at}\ \Lambda\qquad = 0;
\]

%dsp-level:0
\[
(3)\quad\textbf{the component round trip:}
\]

%dsp-level:0
\[
\qquad R_{\mu\nu}=\kappa^2\bigl(T_{\mu\nu}-\tfrac12g_{\mu\nu}T\bigr)+\Lambda g_{\mu\nu}
\]

%dsp-level:0
\[
\qquad [\text{machine: all}\ 4\ \text{components exact}];
\]

%dsp-level:0
\[
p=0:\quad R_{00}=\tfrac12\kappa^2\rho-\Lambda,
\]

%dsp-level:0
\[
\qquad R_{ii}=\tfrac12\kappa^2\rho+\Lambda\ \ (i=1,2,3);
\]

%dsp-level:0
\[
\qquad g^{\mu\nu}R_{\mu\nu}=\kappa^2\rho+4\Lambda=R\ \ [\text{machine: the round trip closes}];
\]

%dsp-level:0
\[
G_{00}=\kappa^2\rho+\Lambda,
\]

%dsp-level:0
\[
\qquad G_{ii}=\kappa^2p-\Lambda\ \overset{p=0}{=}\ -\Lambda;
\]

%dsp-level:0
\[
\qquad G_{\mu\nu}=\kappa^2T_{\mu\nu}-\Lambda g_{\mu\nu}\ \ [\text{machine: all}\ 4\ \text{components}];
\]

%dsp-level:0
\[
(4)\quad\textbf{the curvature scale of the sealed composite:}
\]

%dsp-level:0
\[
\qquad \kappa^2\rho=\frac{2\pi\,E_{\mathrm{int}}^4(\ln2)^3}{\mathfrak h^2}
\]

%dsp-level:0
\[
\qquad [\mathrm{BM.14}(3);\ \text{machine}];
\]

%dsp-level:0
\[
[\kappa^2\rho]=\mathfrak q^{-1}\cdot\mathfrak q^{+2}=\mathfrak q^{+1}=[G_{\mu\nu}]=[\Lambda]=[R]
\]

%dsp-level:0
\[
\qquad [\mathrm{CT2.1};\ \text{machine: exact exponent arithmetic}];
\]

%dsp-level:0
\[
\text{the scale is a pure function of the integer count}\ E_{\mathrm{int}}\ \text{and the two universal factors}\ \mathfrak h,
\]

%dsp-level:0
\[
\ \ln2\ [\mathrm{P1};\ \mathrm{FT.44}(4)]:\ \neg\text{a continuous mass variable enters};
\]

%dsp-level:0
\[
(5)\quad\textbf{the curvature radius:}
\]

%dsp-level:0
\[
\qquad \ell_R:=\frac1{\sqrt{\kappa^2\rho}}:
\]

%dsp-level:0
\[
\qquad \Bigl(\frac{\Delta x_{\mathrm B}}{\ell_R}\Bigr)^{\!2}=(\Delta x_{\mathrm B})^2\,\kappa^2\rho=2\pi\,E_{\mathrm{int}}^2\ln2
\]

%dsp-level:0
\[
\qquad [\text{machine: one Laurent monomial}];
\]

%dsp-level:0
\[
2\pi\ln2>1\ [\text{machine: exact enclosures}\ \pi>\tfrac{223}{71},\ \ln2>\textstyle\sum_1^{12}\tfrac1{n2^n}]\implies\Bigl(\frac{\Delta x_{\mathrm B}}{\ell_R}\Bigr)^{\!2}>1\ \ \forall E_{\mathrm{int}}\ge1:
\]

%dsp-level:0
\[
\qquad \ell_R<\Delta x_{\mathrm B}\ \ \text{on the whole ladder};
\]

%dsp-level:0
\[
\text{the ratio is strictly increasing in}\ E_{\mathrm{int}}\ \text{with disjoint exact enclosures}\ [\text{machine}]:
\]

%dsp-level:0
\[
\qquad \text{the self-curvature of a sealed composite tightens as its information content grows};
\]

%dsp-level:0
\[
(6)\quad\textbf{the}\ \Lambda\ \textbf{clause:}
\]

%dsp-level:0
\[
\qquad \Lambda\ \text{is a P4 core member}\ [\mathrm{P4}]:
\]

%dsp-level:0
\[
\qquad R\ \text{carries exactly one additive free term},\ 4\Lambda,\ \text{with exact coefficient}\ 4\ [\text{machine}];
\]

%dsp-level:0
\[
\text{the locked part is}\ \kappa^2\rho=R-4\Lambda\ \ \text{at}\ p=0:
\]

%dsp-level:0
\[
\qquad \text{the curvature generated by the composite is calibration-locked; the ambient term is not}.
\]

\textbf{Rmk (frame and scope).} Clauses (1)–(3) are component algebra in one orthonormal frame with a diagonal metric; the identities are rational in $(\rho,p,\Lambda,\kappa^2)$ and hold in every frame [math: tensoriality]. Clause (5) compares lengths by squares, so no radical enters any assertion. No solution of the field equations is constructed here: the pp-wave witness of BM.11(4) carries the differential identities, and this tag carries the algebraic dust read.

\textbf{Req.} A3 (Part 1); CT1.12, CT1.23, CT2.1, CT2.9 (Part 2); CT4.5 (Part 2); P1, P4 (Parts 1–2); Def-BM.4 (this part); FT.44(4) (Census); BM.9(4), BM.11(4), BM.13(1), BM.14(3) (this part); exact multivariate polynomial algebra over $\mathbb Q$, Archimedes' rational bounds for $\pi$, exact rational $\ln2$ enclosures [math; machine: \texttt{checks/\allowbreak{}verify\_\allowbreak{}H\_\allowbreak{}calibration\_\allowbreak{}lock.py} §C].

\textbf{Pf.}

%dsp-level:0
\[
(1):\ g^{\mu\nu}T_{\mu\nu}=(-1)\rho+1\cdot p+1\cdot p+1\cdot p=-\rho+3p\ [\text{machine}];\qquad p=0\ [\mathrm{CT4.5}:\ w=p/\rho=0]
\]

%dsp-level:0
\[
(2):\ G_{\mu\nu}=R_{\mu\nu}-\tfrac12g_{\mu\nu}R\implies G^\mu{}_\mu=R-2R=-R\ [\text{math},\ 4\ \text{dimensions}];
\]

%dsp-level:0
\[
\qquad \mathrm{CT1.23}\ \text{traced}:\ -R+4\Lambda=\kappa^2(-\rho+3p)\implies R=\kappa^2(\rho-3p)+4\Lambda\ [\text{machine}]
\]

%dsp-level:0
\[
\text{residual of}\ R
\]

%dsp-level:0
\[
= -\kappa^2\rho:\ \bigl[\kappa^2(\rho-3p)+4\Lambda\bigr]+\kappa^2\rho
\]

%dsp-level:0
\[
= 2\kappa^2\rho-3\kappa^2p+4\Lambda
\]

%dsp-level:0
\[
\ne 0\ \text{as a polynomial}\ [\text{machine}]
\]

%dsp-level:0
\[
(3):\ R_{\mu\nu}
\]

%dsp-level:0
\[
= \kappa^2T_{\mu\nu}+\tfrac12g_{\mu\nu}R-\Lambda g_{\mu\nu}
\]

%dsp-level:0
\[
= \kappa^2(T_{\mu\nu}-\tfrac12g_{\mu\nu}T)+\Lambda g_{\mu\nu}\ [\text{substituting}\ (2);\ \text{machine}];
\]

%dsp-level:0
\[
\qquad G_{\mu\nu}=R_{\mu\nu}-\tfrac12g_{\mu\nu}R=\kappa^2T_{\mu\nu}-\Lambda g_{\mu\nu}\ [\text{machine: all}\ 4];
\]

%dsp-level:0
\[
\qquad \text{the dust values as displayed}\ [\text{machine}]
\]

%dsp-level:0
\[
(4):\ \mathrm{BM.14}(3)\ \text{re-cited};
\]

%dsp-level:0
\[
\qquad \text{exponents}: -1+2=+1\ [\mathrm{CT2.1}:\ [\kappa^2]=\mathfrak q^{-1},\ [T_{\mu\nu}]=\mathfrak q^{+2},\ [G_{\mu\nu}]=[\Lambda]=\mathfrak q^{+1}]
\]

%dsp-level:0
\[
(5):\ (\Delta x_{\mathrm B})^2\kappa^2\rho=\tfrac{\mathfrak h^2}{(E\ln2)^2}\cdot\tfrac{2\pi E^4(\ln2)^3}{\mathfrak h^2}=2\pi E^2\ln2\ [\text{machine}];
\]

%dsp-level:0
\[
\qquad 2\pi\ln2>2\cdot\tfrac{223}{71}\cdot\tfrac{405}{584}>1\ [\text{machine: exact rationals}]\implies\ell_R<\Delta x_{\mathrm B}
\]

%dsp-level:0
\[
(6):\ \text{the coefficient of}\ \Lambda\ \text{in}\ R\ \text{is}\ 4\ [\text{machine: monomial census}];
\]

%dsp-level:0
\[
\qquad \Lambda\in\mathrm{P4}\ \text{core}\ [\mathrm{P4}\ \text{re-cited}];
\]

%dsp-level:0
\[
\qquad R-4\Lambda=\kappa^2\rho\ \text{at}\ p=0\ [\text{machine}].\ \blacksquare
\]

\textbf{Phys.} Trace the field equations and the sign of a dust curvature is not a matter of convention: $R=\kappa^2(\rho-3p)+4\Lambda$, positive for a positive density at vanishing $\Lambda$, and the opposite sign is refuted as a polynomial identity rather than argued away. The component round trip closes — Ricci, scalar, Einstein, back to $\kappa^2T$ — in exact rational arithmetic on all four diagonal entries. Feeding in the width-derived density gives the curvature scale as $2\pi E_{\mathrm{int}}^4(\ln2)^3/\mathfrak h^2$: a pure function of an integer count, dimensionally exactly a curvature. Its radius is always shorter than the Born width, by a factor that grows with the count, so the more information a composite seals, the tighter its own curvature sits inside its spread. One term is not locked, and it is named: the cosmological constant enters the scalar additively with coefficient four, a free internal parameter, and the locked part is what remains after subtracting it.

\end{theorem}

\begin{theorem}[{\texorpdfstring{BM.16 — The calibration lock}{BM.16 - The calibration lock}}]\label{tag:BM-16}
\textbf{Thm.}

%dsp-level:0
\[
(1)\quad\textbf{the lock identity:}
\]

%dsp-level:0
\[
\qquad \boxed{\ (\Delta x_{\mathrm B})^2\cdot\kappa^2\rho\ =\ 2\pi\,E_{\mathrm{int}}^2\ln2\ }
\]

%dsp-level:0
\[
\qquad [\text{machine: one Laurent monomial, coefficient}\ 2];
\]

%dsp-level:0
\[
\mathfrak h\text{-free}\ [\text{machine: the}\ \mathfrak h\text{-exponent is}\ 0];
\]

%dsp-level:0
\[
\qquad \Lambda\text{-free}\ [\text{machine}];
\]

%dsp-level:0
\[
\qquad \text{P4-free}\ [\mathrm{P4}\ \text{core}:\ g,g',h_0,c_W,\Lambda,\varkappa,\beta,\alpha_\Omega\ —\ \neg\text{one occurs}];
\]

%dsp-level:0
\[
\textbf{the three dimensionless reads:}
\]

%dsp-level:0
\[
\qquad \frac{\Delta x_{\mathrm B}}{\mathfrak h}=\frac1{E_{\mathrm{int}}\ln2};
\]

%dsp-level:0
\[
\qquad \kappa^2\rho\,\mathfrak h^2=2\pi E_{\mathrm{int}}^4(\ln2)^3;
\]

%dsp-level:0
\[
\qquad \rho\,\mathfrak h^4=(E_{\mathrm{int}}\ln2)^4
\]

%dsp-level:0
\[
\qquad [\text{machine: all}\ \mathfrak h\text{-free}];
\]

%dsp-level:0
\[
(2)\quad\textbf{the rescaling covariance:}
\]

%dsp-level:0
\[
\qquad \mathfrak h\mapsto\lambda\mathfrak h:
\]

%dsp-level:0
\[
\qquad \Delta x_{\mathrm B}\mapsto\lambda\,\Delta x_{\mathrm B},\quad V_{\mathrm{eff}}\mapsto\lambda^3V_{\mathrm{eff}},\quad m_{\mathrm{hid}}\mapsto\lambda^{-1}m_{\mathrm{hid}},
\]

%dsp-level:0
\[
\rho\mapsto\lambda^{-4}\rho,\qquad\kappa^2\rho\mapsto\lambda^{-2}\kappa^2\rho,\qquad G_N\mapsto\lambda^2G_N,\qquad\kappa^2\mapsto\lambda^2\kappa^2\qquad[\text{machine: exact weight census}];
\]

%dsp-level:0
\[
\text{every dimensionless product of}\ (1)\ \text{is invariant}\ [\text{machine}]:
\]

%dsp-level:0
\[
\qquad \text{rescaling is a symmetry}\ \exists!\ \text{calibration d.o.f.}\ [\mathrm{CT2.11}]\ \text{and the lock rides above it};
\]

%dsp-level:0
\[
(3)\quad\textbf{no continuous mass variable:}
\]

%dsp-level:0
\[
\qquad \text{every read of this tier is a single Laurent monomial in}\ (\mathfrak h,E_{\mathrm{int}},\ln2,\pi)\ [\text{machine: monomial census}]:
\]

%dsp-level:0
\[
\neg\exists\ \text{a sum of scales,}\ \neg\exists\ \text{a tunable mixture};
\]

%dsp-level:0
\[
\qquad m_{\mathrm{hid}}\mathfrak h=E_{\mathrm{int}}\ln2\ \text{is injective in the integer}\ E_{\mathrm{int}}\ [\text{machine}];
\]

%dsp-level:0
\[
\textbf{quantization:}
\]

%dsp-level:0
\[
\qquad m_{\mathrm{hid}}\in\Bigl\{\frac{n\ln2}{\mathfrak h}:n\in\mathbb Z_{\ge2}\Bigr\},
\]

%dsp-level:0
\[
\qquad \Delta x_{\mathrm B}\in\Bigl\{\frac{\mathfrak h}{n\ln2}\Bigr\},
\]

%dsp-level:0
\[
\qquad \kappa^2\rho\in\Bigl\{\frac{2\pi n^4(\ln2)^3}{\mathfrak h^2}\Bigr\}:
\]

%dsp-level:0
\[
\qquad \text{three discrete ladders, strictly monotone}\ [\text{machine}];
\]

%dsp-level:0
\[
(4)\quad\textbf{the conjugate comparison scales}\ [\mathrm{P3}:\ G_N\ \text{is a named combination for external comparison},\ \neg\text{an input}]:
\]

%dsp-level:0
\[
\qquad r_s:=2G_Nm_{\mathrm{hid}}=\frac{\mathfrak h\,E_{\mathrm{int}}}2
\]

%dsp-level:0
\[
\qquad [\text{machine: exact}];
\]

%dsp-level:0
\[
\boxed{\ \Delta x_{\mathrm B}\cdot r_s=2G_N=\frac{\mathfrak h^2}{2\ln2}\ }
\]

%dsp-level:0
\[
\qquad \text{independent of}\ E_{\mathrm{int}}\ [\text{machine: the}\ E\text{-exponent is}\ 0]:
\]

%dsp-level:0
\[
\qquad 2\ln2\cdot(\Delta x_{\mathrm B}r_s)=\mathfrak h^2;
\]

%dsp-level:0
\[
\text{the two length reads are conjugate with a fixed product}: \ E_{\mathrm{int}}\uparrow\implies\Delta x_{\mathrm B}\downarrow,
\]

%dsp-level:0
\[
\ r_s\uparrow\ \text{along one hyperbola}\ [\text{machine}];
\]

%dsp-level:0
\[
(5)\quad\textbf{the crossing gap:}
\]

%dsp-level:0
\[
\qquad \ell_G:=\sqrt{G_N}=\frac{\mathfrak h}{2\sqrt{\ln2}}\ [\mathrm{CT2.9}]:
\]

%dsp-level:0
\[
\qquad \Bigl(\frac{\Delta x_{\mathrm B}}{\ell_G}\Bigr)^{\!2}=\frac{\Delta x_{\mathrm B}^2}{G_N}=\frac4{E_{\mathrm{int}}^2\ln2}
\]

%dsp-level:0
\[
\qquad [\text{machine: squares only, no radical in any assertion}];
\]

%dsp-level:0
\[
\Delta x_{\mathrm B}<\ell_G\iff E_{\mathrm{int}}^2\ln2>4\iff E_{\mathrm{int}}\ge3\qquad[\text{machine: exact enclosures, every}\ E\in1..60];
\]

%dsp-level:0
\[
\frac{r_s}{\Delta x_{\mathrm B}}=\frac{E_{\mathrm{int}}^2\ln2}2
\]

%dsp-level:0
\[
\qquad [\text{machine}]:
\]

%dsp-level:0
\[
\qquad r_s=\Delta x_{\mathrm B}\iff E_{\mathrm{int}}^2\ln2=2\iff E_{\mathrm{int}}=\sqrt{\tfrac2{\ln2}}\in(1,2)\ [\text{machine:}\ 1<\tfrac2{\ln2}<4];
\]

%dsp-level:0
\[
\textbf{the crossing lies strictly below the minimal floor}\ E_{\mathrm{int}}=2\ [\mathrm{FT.28}]:
\]

%dsp-level:0
\[
\qquad \neg\exists\,E_{\mathrm{int}}\in\mathbb Z\ \text{at the crossing};
\]

%dsp-level:0
\[
\qquad r_s>\Delta x_{\mathrm B}\ \ \forall E_{\mathrm{int}}\ge2\ [\text{machine}];
\]

%dsp-level:0
\[
(6)\quad\textbf{the not-locked ledger:}
\]

%dsp-level:0
\[
\qquad \Lambda\ [\mathrm{P4}:\ \text{one additive term}\ 4\Lambda\ \text{in}\ R,\ \mathrm{BM.15}(6)];
\]

%dsp-level:0
\[
\qquad \alpha\ [\mathrm{Def\text{-}BM.4}(4):\ \text{priced},\ \neg\text{adopted}];
\]

%dsp-level:0
\[
\qquad \mathfrak h\ [\mathrm{P1},\ \mathrm{CT2.11}:\ \text{the single free scale}\ —\ \text{absolute values ride on it}];
\]

%dsp-level:0
\[
\textbf{the comparison discipline:}
\]

%dsp-level:0
\[
\quad r_s,\ \ell_G\ \text{are built from}\ G_N,\ \text{a named combination of}\ \mathfrak h\ \text{for external comparison only}\ [\mathrm{P3}]:
\]

%dsp-level:0
\[
\qquad \neg\text{a collapse reading,}\ \neg\text{a horizon reading,}\ \neg\text{an identification is made}\ [\mathrm{P5};\ \mathrm{Def\text{-}BM.3}(6)];
\]

%dsp-level:0
\[
\text{the inequalities of}\ (5)\ \text{are order relations between derived lengths, and no external datum enters them}.
\]

\textbf{Rmk (lock scope).} The lock of clause (1) is a statement about dimensionless products: it says the width and the curvature cannot be varied independently of the integer count. It does not say the corpus predicts a length in metres — the single free scale $\mathfrak h$ fixes the unit [P1, P5], and every external comparison stays quarantined [Census FT.114, FT.120].

\textbf{Req.} D2.9 (Parts 1–2); CT2.1, CT2.9, CT2.11 (Part 2); P1, P3, P4, P5 (Parts 1–2); Def-BM.3(6), Def-BM.4 (this part); FT.28, FT.44(4) (Part 3, Census); Census FT.114, FT.120 (quarantine rows); BM.13(1), BM.14(2), BM.14(3), BM.15(4), BM.15(6) (this part); exact Laurent-monomial arithmetic, exact rational enclosures of $\ln2$ and Archimedes' bounds for $\pi$, square-comparison inequalities [math; machine: \texttt{checks/\allowbreak{}verify\_\allowbreak{}H\_\allowbreak{}calibration\_\allowbreak{}lock.py} §D].

\textbf{Pf.}

%dsp-level:0
\[
(1):\ (\Delta x_{\mathrm B})^2\kappa^2\rho=\tfrac{\mathfrak h^2}{(E\ln2)^2}\cdot\tfrac{2\pi E^4(\ln2)^3}{\mathfrak h^2}=2\pi E^2\ln2\ [\text{machine}];
\]

%dsp-level:0
\[
\qquad \text{the}\ \mathfrak h\text{-exponent census}: 2-2=0\ [\text{machine}];
\]

%dsp-level:0
\[
\qquad \text{the P4 core members occur in no monomial}\ [\text{machine scan of the tier's expressions}]
\]

%dsp-level:0
\[
(2):\ \text{each read is}\ c\,
\]

%dsp-level:0
\[
\mathfrak h^{k}\ \text{with}\ k\in\{1,3,-1,-4,-2,2,2\}\ [\text{machine: weight census}]\implies\mathfrak h\mapsto\lambda\mathfrak h\ \text{acts by}\ \lambda^k\ [\mathrm{CT2.11}\ \text{re-cited}];
\]

%dsp-level:0
\[
\qquad \text{the products of}\ (1)\ \text{have}\ k=0\implies\text{invariant}
\]

%dsp-level:0
\[
(3):\ \text{monomial census}\ [\text{machine}];
\]

%dsp-level:0
\[
\qquad n_1\ln2=n_2\ln2\iff n_1=n_2\ [\text{machine on}\ n\in1..39];
\]

%dsp-level:0
\[
\qquad \text{monotonicity of all three ladders}\ [\text{machine}]
\]

%dsp-level:0
\[
(4):\ r_s=2\cdot\tfrac{\mathfrak h^2}{4\ln2}\cdot\tfrac{E\ln2}{\mathfrak h}=\tfrac{\mathfrak h E}2\ [\text{machine}];
\]

%dsp-level:0
\[
\qquad \Delta x_{\mathrm B}r_s=\tfrac{\mathfrak h}{E\ln2}\cdot\tfrac{\mathfrak h E}2=\tfrac{\mathfrak h^2}{2\ln2}=2G_N\ [\text{machine}];
\]

%dsp-level:0
\[
\qquad \text{the}\ E\text{-exponent is}\ -1+1=0
\]

%dsp-level:0
\[
(5):\ \Delta x_{\mathrm B}^2/G_N=\tfrac{\mathfrak h^2}{(E\ln2)^2}\cdot\tfrac{4\ln2}{\mathfrak h^2}=\tfrac4{E^2\ln2}\ [\text{machine}];
\]

%dsp-level:0
\[
\qquad \tfrac4{E^2\ln2}<1
\]

%dsp-level:0
\[
\iff E^2\ln2>4:\ E
\]

%dsp-level:0
\[
= 2\to4\ln2<4\ (\ln2<1),
\]

%dsp-level:0
\[
\ E
\]

%dsp-level:0
\[
= 3\to9\ln2>4\ (\ln2>\tfrac49)\ [\text{machine: exact enclosures}];
\]

%dsp-level:0
\[
r_s/\Delta x_{\mathrm B}=\tfrac{\mathfrak h E}2\cdot\tfrac{E\ln2}{\mathfrak h}=\tfrac{E^2\ln2}2\ [\text{machine}];
\]

%dsp-level:0
\[
\qquad =1\iff E^2=\tfrac2{\ln2}\in(1,4)\ [\text{machine}]\implies E\in(1,2);
\]

%dsp-level:0
\[
\qquad E\ge2\implies\tfrac{E^2\ln2}2\ge2\ln2>1\ [\ln2>\tfrac12;\ \text{machine}]
\]

%dsp-level:0
\[
(6):\ \mathrm{BM.15}(6),\ \mathrm{Def\text{-}BM.4}(4),\ \mathrm{P1}/\mathrm{CT2.11},\ \mathrm{P3}/\mathrm{P5}\ \text{re-cited}.\ \blacksquare
\]

\textbf{Phys.} Two numbers in this tier carry neither the free scale nor a single internal parameter: the product of the squared Born width with the curvature scale, which is $2\pi$ times the squared edge count times $\ln2$, and the product of the Born width with the comparison radius, which is exactly twice Newton's constant and does not know the edge count at all. That is the lock — width and curvature are conjugate along a hyperbola whose constant is pure calibration, so no sealed composite can slide its size without sliding its field by the reciprocal amount. Rescaling the one free scale moves every dimensionful read by its exact weight and leaves every one of these products standing, which is what makes the lock a statement about the theory rather than about units. The ladder has a hole in it too: the point where the two lengths would coincide sits at an edge count between one and two, strictly below the smallest chain the ledger admits, so no admissible composite is ever at its own comparison radius. What is not locked is listed: the cosmological term, the geometric factor, and the hop itself.

\end{theorem}

\pseudohead{theorem}{\texorpdfstring{BM.17}{BM.17}}{\texorpdfstring{The neutron row of the identification layer}{The neutron row of the identification layer}}\label{tag:BM-17}
\textbf{Thm.}

%dsp-level:0
\[
(1)\quad\textbf{the structure:}
\]

%dsp-level:0
\[
\qquad \Sigma_{12}
\]

%dsp-level:0
\[
= \{p_1,q_1\}\cup\{p_2,q_2\}\ \text{with materialized 2-cycles joined by a bidirectional connector, single observer contact}\ [\text{Census FT.35}];
\]

%dsp-level:0
\[
E_{\mathrm{int}}=6\ (c_1{:}\,2,\ c_2{:}\,2,\ \text{connector}{:}\,2),
\]

%dsp-level:0
\[
\qquad I_{\mathrm{hid}}=6\ln2,
\]

%dsp-level:0
\[
\qquad m_{\mathrm{hid}}=\frac{6\ln2}{\mathfrak h}
\]

%dsp-level:0
\[
\qquad [\text{Census FT.35}(2);\ \text{cited row-wise}];
\]

%dsp-level:0
\[
\mathcal C(\Sigma_{12};C)=1\iff\text{all three attachments seal inside as one}\ \varepsilon\text{-type (baryon-like) triple}
\]

%dsp-level:0
\[
\qquad [\text{Census FT.35}(5);\ \mathrm{D3.9};\ \mathrm{CT3.17};\ \mathrm{FT.34}(5)];
\]

%dsp-level:0
\[
\Theta_{\mathrm{int}}=2k(\Gamma)\in2\mathbb Z,
\]

%dsp-level:0
\[
\qquad \Pi=72\,k(\Gamma)\,(\ln2)^2
\]

%dsp-level:0
\[
\qquad [\text{Census FT.35}(3)];
\]

%dsp-level:0
\[
\qquad \text{the junction vertices are}\ D^-\text{-configured}\ [\text{Census FT.35}(4);\ \mathrm{CT4.1}];
\]

%dsp-level:0
\[
(2)\quad\textbf{the identification:}
\]

%dsp-level:0
\[
\qquad \mathrm{Def\text{-}BM.3}(2)\ \text{labels the}\ \varepsilon\text{-sealed 3-juncture in the neutral charge class}\ q_{\mathrm{net}}
\]

%dsp-level:0
\[
= 0\ \text{as the}\ \textbf{neutron};
\]

%dsp-level:0
\[
\text{the theorem content of every clause below is the tuple}\ \bigl(\varepsilon\text{-junction},\ \Theta_{\mathrm{int}}\in2\mathbb Z,\ q_{\mathrm{net}}=0,\ E_{\mathrm{int}}=6,\ w=0\bigr);
\]

%dsp-level:0
\[
\qquad \text{the name is the layer's label}\ [\mathrm{Def\text{-}BM.3}(1),(7)];
\]

%dsp-level:0
\[
\text{the charge class is}\ \mathrm{Def\text{-}BM.3}(4)\text{'s declared label, pending BM.6(1)'s additive read}\;
\]

%dsp-level:0
\[
\qquad \text{deficit}\ 1'\ \text{is untouched}\ [\mathrm{FT.125}(4)];
\]

%dsp-level:0
\[
(3)\quad\textbf{the derived row at}\ E_{\mathrm{int}}^{(n)}=6\ \text{(all exact):}
\]

%dsp-level:0
\[
\Delta x_{\mathrm B}^{(n)}=\frac{\mathfrak h}{6\ln2}\ (<\mathfrak h\ \text{by the exact factor}\ \tfrac1{6\ln2});
\]

%dsp-level:0
\[
\qquad V_{\mathrm{eff}}^{(n)}=\frac{\mathfrak h^3}{216(\ln2)^3};
\]

%dsp-level:0
\[
\qquad m_{\mathrm{hid}}^{(n)}=\frac{6\ln2}{\mathfrak h};
\]

%dsp-level:0
\[
\rho^{(n)}=\frac{1296(\ln2)^4}{\mathfrak h^4};
\]

%dsp-level:0
\[
\qquad \kappa^2\rho^{(n)}=\frac{2592\,\pi(\ln2)^3}{\mathfrak h^2}
\]

%dsp-level:0
\[
\qquad [2\cdot6^4=2592;\ \text{machine}];
\]

%dsp-level:0
\[
\qquad \ell_R^{(n)}=\frac{\mathfrak h}{6^2\sqrt{2\pi(\ln2)^3}}<\Delta x_{\mathrm B}^{(n)};
\]

%dsp-level:0
\[
r_s^{(n)}=3\,\mathfrak h
\]

%dsp-level:0
\[
\qquad [\text{machine}];
\]

%dsp-level:0
\[
\qquad \Delta x_{\mathrm B}^{(n)}\cdot r_s^{(n)}=\frac{\mathfrak h^2}{2\ln2}=2G_N
\]

%dsp-level:0
\[
\qquad [\text{machine: the}\ E\text{-independent product}];
\]

%dsp-level:0
\[
(\Delta x_{\mathrm B}^{(n)})^2\,\kappa^2\rho^{(n)}=72\,\pi\ln2
\]

%dsp-level:0
\[
\qquad [\text{machine}]:
\]

%dsp-level:0
\[
\qquad \text{the integer}\ 72
\]

%dsp-level:0
\[
= 2\cdot6^2\ \text{is Census FT.35}(3)\text{'s}\ \Pi\text{-coefficient}\ —\ \text{recorded as a cross-check},
\]

%dsp-level:0
\[
\ \neg\text{a derivation of equality}\ [\mathrm{FT.112}\ \mathrm{Rmk}\ \text{pattern}];
\]

%dsp-level:0
\[
\Delta x_{\mathrm B}^{(n)}<\ell_G\iff36\ln2>4\iff\ln2>\tfrac19
\]

%dsp-level:0
\[
\qquad [\text{machine: exact enclosure}];
\]

%dsp-level:0
\[
\qquad w^{(n)}=0,\quad T^{(n)}_{\mu\nu}=\mathrm{diag}(\rho^{(n)},0,0,0)
\]

%dsp-level:0
\[
\qquad [\mathrm{CT4.5};\ \mathrm{BM.9}(4)];
\]

%dsp-level:0
\[
(4)\quad\textbf{the family table (the five computed cells):}
\]

%tbl-level:2
{\footnotesize\setlength{\tabcolsep}{4pt}
\begin{longtable}{>{\raggedright\arraybackslash}p{\dimexpr0.2770\textwidth-2\tabcolsep\relax}|>{\raggedright\arraybackslash}p{\dimexpr0.0529\textwidth-2\tabcolsep\relax}|>{\raggedright\arraybackslash}p{\dimexpr0.1010\textwidth-2\tabcolsep\relax}|>{\raggedright\arraybackslash}p{\dimexpr0.1411\textwidth-2\tabcolsep\relax}|>{\raggedright\arraybackslash}p{\dimexpr0.1330\textwidth-2\tabcolsep\relax}|>{\raggedright\arraybackslash}p{\dimexpr0.1491\textwidth-2\tabcolsep\relax}|>{\raggedright\arraybackslash}p{\dimexpr0.0690\textwidth-2\tabcolsep\relax}|>{\raggedright\arraybackslash}p{\dimexpr0.0770\textwidth-2\tabcolsep\relax}}
\toprule
cell & $E_{\mathrm{int}}$ & $\Delta x_{\mathrm B}/\mathfrak h$ & $V_{\mathrm{eff}}/\mathfrak h^3$ & $\rho\,\allowbreak \mathfrak h^4$ & $\kappa^2\rho\,\allowbreak \mathfrak h^2$ & $r_s/\mathfrak h$ & $\Delta x_{\mathrm B}$ vs $\ell_G$ \\
\addlinespace[2pt]
\midrule
\endfirsthead
\toprule
cell & $E_{\mathrm{int}}$ & $\Delta x_{\mathrm B}/\mathfrak h$ & $V_{\mathrm{eff}}/\mathfrak h^3$ & $\rho\,\allowbreak \mathfrak h^4$ & $\kappa^2\rho\,\allowbreak \mathfrak h^2$ & $r_s/\mathfrak h$ & $\Delta x_{\mathrm B}$ vs $\ell_G$ \\
\addlinespace[2pt]
\midrule
\endhead
\bottomrule
\endlastfoot
$\Sigma_{\min}$ (the digon) [FT.28] & 2 & $\tfrac1{2\ln2}$ & $\tfrac1{8(\ln2)^3}$ & $16(\ln2)^4$ & $32\pi(\ln2)^3$ & 1 & $\Delta x_{\mathrm B}>\allowbreak \ell_G$ \\
$\Sigma_{12}$ (the $\varepsilon$-type junction) [Census FT.35] & 6 & $\tfrac1{6\ln2}$ & $\tfrac1{216(\ln2)^3}$ & $1296(\ln2)^4$ & $2592\pi(\ln2)^3$ & 3 & $\Delta x_{\mathrm B}<\allowbreak \ell_G$ \\
tetrahedron [Census FT.42(5)] & 12 & $\tfrac1{12\ln2}$ & $\tfrac1{1728(\ln2)^3}$ & $20736(\ln2)^4$ & $41472\pi(\ln2)^3$ & 6 & $\Delta x_{\mathrm B}<\allowbreak \ell_G$ \\
cube / octahedron [Census FT.42(5)] & 24 & $\tfrac1{24\ln2}$ & $\tfrac1{13824(\ln2)^3}$ & $331776(\ln2)^4$ & $663552\pi(\ln2)^3$ & 12 & $\Delta x_{\mathrm B}<\allowbreak \ell_G$ \\
dodecahedron / icosahedron [Census FT.42(5)] & 60 & $\tfrac1{60\ln2}$ & $\tfrac1{216000(\ln2)^3}$ & $12960000(\ln2)^4$ & $25920000\pi(\ln2)^3$ & 30 & $\Delta x_{\mathrm B}<\allowbreak \ell_G$ \\
\end{longtable}
}

%dsp-level:0
\[
\text{the width column is strictly decreasing and the curvature column strictly increasing, with disjoint exact enclosures}\ [\text{machine}];
\]

%dsp-level:0
\[
\qquad r_s/\mathfrak h=E_{\mathrm{int}}/2\ \text{across the family}\ [\text{machine}];
\]

%dsp-level:0
\[
\Sigma_{\min}\ \text{is the unique computed cell with}\ \Delta x_{\mathrm B}>\ell_G\ [\mathrm{BM.16}(5);\ \text{machine}];
\]

%dsp-level:0
\[
\qquad \text{the duality-degenerate pairs share}\ E_{\mathrm{int}}\ \text{and hence every read of this tier}\ [\mathrm{FT.44}(1);\ \text{machine}];
\]

%dsp-level:0
\[
(5)\quad\textbf{not a point mass:}
\]

%dsp-level:0
\[
\qquad V_{\mathrm{eff}}^{(n)}>0,\quad\rho^{(n)}<\infty,\quad T^{(n)}_{\mu\nu}\ \text{is a finite-width distribution}\ [\mathrm{BM.9}(6);\ \mathrm{BM.10}(6)];
\]

%dsp-level:0
\[
\text{the width is sub-hop}\ [(3)]\implies\text{it is a chart envelope},
\]

%dsp-level:0
\[
\ \neg\text{a materialized distance}\ [\mathrm{Def\text{-}BM.4}(5);\ \mathrm{D1.6}];
\]

%dsp-level:0
\[
(6)\quad\textbf{the openness clause:}
\]

%dsp-level:0
\[
\qquad E_{\mathrm{int}}^{(n)}
\]

%dsp-level:0
\[
= 6\ \text{is Census-cell data at the computed}\ \varepsilon\text{-type cell}\ [\text{Census FT.35}(2);\ \text{row-wise consumption},\ \mathrm{BM.5}(2)\ \text{pattern}];
\]

%dsp-level:0
\[
\text{a general composite additivity law for}\ E_{\mathrm{int}}\ \text{beyond the computed cells is}\ \neg\text{constructed}\ [\mathrm{Def\text{-}BM.1}\ \text{read clause};\ \mathrm{BM.6}(1),(5)]:
\]

%dsp-level:0
\[
\text{every read of this tier is exact in}\ E_{\mathrm{int}}^{(n)}\ \text{and parametric wherever no cell supplies it}\ —\ \neg\text{a numeric value is invented};
\]

%dsp-level:0
\[
\neg\exists\ \text{a measured value in any clause}\ [\mathrm{P5};\ \mathrm{Def\text{-}BM.3}(6)].
\]

\textbf{Rmk (identification discipline).} Clause (3) is a row of derived exact expressions attached to a Census cell and labelled by Def-BM.3. It derives no neutron property that the tuple does not already carry, imports no neutron datum, and closes no row beyond its own: the additive charge read [BM.6(1)] and the assignment rule [FT.125(4)] stand as they are.

\textbf{Req.} D1.6, D2.9, D3.9 (Parts 1–2); CT3.17, CT4.1 (Part 2); P5 (Parts 1–2); Def-BM.1, Def-BM.3, Def-BM.4, BM.5(2), BM.6(1), BM.6(5) (this part); FT.28, FT.34(5), FT.44(1), FT.125(4) (Part 3); Census FT.35(2), FT.35(3), FT.35(4), FT.35(5), FT.42(5), FT.44(4) (row-wise); BM.9(4), BM.9(6), BM.10(6), BM.13, BM.14, BM.15, BM.16 (this part); exact rational arithmetic, exact $\ln2$ enclosures [math; machine: \texttt{checks/\allowbreak{}verify\_\allowbreak{}H\_\allowbreak{}calibration\_\allowbreak{}lock.py} §E].

\textbf{Pf.}

%dsp-level:0
\[
(1):\ \text{Census FT.35's rows re-cited verbatim}\ (2),
\]

%dsp-level:0
\[
(3),
\]

%dsp-level:0
\[
(4),
\]

%dsp-level:0
\[
(5)\ —\ \text{row-wise consumption},
\]

%dsp-level:0
\[
\ \neg\text{premise import}\ [\mathrm{BM.5}(2);\ \text{A §0.1}]
\]

%dsp-level:0
\[
(2):\ \mathrm{Def\text{-}BM.3}(2)\ \text{re-cited};
\]

%dsp-level:0
\[
\qquad \text{the tuple is the theorem content}\ [\mathrm{Def\text{-}BM.3}(1),(7)];
\]

%dsp-level:0
\[
\qquad \text{the conditional charge class}\ [\mathrm{Def\text{-}BM.3}(4);\ \mathrm{BM.6}(1)]
\]

%dsp-level:0
\[
(3):\ \text{substitution of}\ E_{\mathrm{int}}
\]

%dsp-level:0
\[
= 6\ \text{into BM.13(1), BM.14(1),(2),(3), BM.15(4),(5), BM.16(4)}\ [\text{machine: every coefficient exact}];
\]

%dsp-level:0
\[
\qquad 2\cdot6^4=2592,\ 6^4=1296,\ 6^3=216,\ 2\cdot6^2=72\ [\text{machine}];
\]

%dsp-level:0
\[
\qquad 36\ln2>4\iff\ln2>\tfrac19\ [\text{machine: enclosure}]
\]

%dsp-level:0
\[
(4):\ \text{the same substitutions at}\ E_{\mathrm{int}}\in\{2,6,12,24,60\}\ [\text{Census FT.42}(5)\ \text{card table};\ \mathrm{FT.28}];
\]

%dsp-level:0
\[
\qquad \text{monotonicity and disjointness}\ [\text{machine}];
\]

%dsp-level:0
\[
\qquad r_s/\mathfrak h=E/2\ [\text{machine}];
\]

%dsp-level:0
\[
\qquad \text{duality degeneracy}: E_{\mathrm{int}}=2E\ \text{shared by dual pairs}\ [\mathrm{FT.44}(1)\ \text{re-cited}]
\]

%dsp-level:0
\[
(5):\ \mathrm{BM.9}(6),\ \mathrm{BM.10}(6),\ \mathrm{Def\text{-}BM.4}(5)\ \text{re-cited}
\]

%dsp-level:0
\[
(6):\ \text{openness rows re-cited unmodified};
\]

%dsp-level:0
\[
\qquad \mathrm{P5}\ \text{scan of this tag's text}\ [\text{machine: zero external numeric tokens}].\ \blacksquare
\]

\textbf{Phys.} The corpus has one computed $\varepsilon$-type junction cell, and it carries six internal directed edges: two per cycle plus the two of the connector. Everything this tier derived then evaluates on it in exact rationals — a width of one hop over six $\ln2$, a volume of one hop cubed over $216(\ln2)^3$, a density of $1296(\ln2)^4$ per hop to the fourth, a curvature scale of $2592\pi(\ln2)^3$ per hop squared, and a comparison radius of exactly three hops. The lock reads $72\pi\ln2$ on this cell, and $72$ is the same integer the Census found in its $\Pi$ row — recorded as a coincidence of counts, not promoted to an identity. Def-BM.3's label attaches to that tuple, and the label changes nothing in it. The neutron of this corpus is not a point: it is a finite density over a strictly positive, sub-hop volume, with a curvature radius shorter than its own spread, sitting on a ladder whose rungs are integers and whose crossing point no integer reaches.

\medskip\noindent\textit{\texorpdfstring{BM.18 — The status of the calibration lock: adjudication and scope}{BM.18 - The status of the calibration lock: adjudication and scope}.}\ Structural-constraints record: see Appendix B, Section B.1.16 of the Addendum.\par\medskip

\phantomsection
\section[Tier BM.6 — The Topology-Blind Neutron and the…]{\texorpdfstring{Tier BM.6 — The Topology-Blind Neutron and the Topology-Blind + Static Proton: the Flicker Ensemble, the Path-Integrated Mass, the Binding Read, and the Forced Carrier}{Tier BM.6 - The Topology-Blind Neutron and the Topology-Blind + Static Proton: the Flicker Ensemble, the Path-Integrated Mass, the Binding Read, and the Forced Carrier}}

\begin{definition}[{\texorpdfstring{Def-BM.5 — The flicker ensemble of a sealed composite}{Def-BM.5 - The flicker ensemble of a sealed composite}}]\label{tag:Def-BM-5}
\textbf{Def.} A construction class on Def-H.32's topology-blind path integral (no new axiom; Def-H.32's scopes (S1)–(S4) are inherited at every use).

%dsp-level:0
\[
(1)\quad\textbf{the ensemble:}\qquad\mathcal E(\mathcal X)\ :=\ \bigl\{R\subseteq E_{\mathrm{dir}}(\Gamma):\ \text{the boundary data of}\ \mathcal X\ \text{is preserved on}\ R\bigr\}
\]

%dsp-level:0
\[
\text{the boundary data}\ :
\]

%dsp-level:0
\[
= \ \bigl(\text{charge class}\ [\mathrm{Def\text{-}BM.3}(2)\text{–}(3)],\ \ \Theta_{\mathrm{int}}\ \text{class}\ [\mathrm{BM.2}(1)],\ \ S\ \text{winding class}\ [\mathrm{Def\text{-}H.32}(4):\ \text{the}\ (q,S)\ \text{labels}]\bigr);
\]

%dsp-level:0
\[
\textbf{constitutive clause:}\ \ \text{the ensemble is the sector-conditioned configuration set of}\ \mathrm{Def\text{-}H.32}(1)\text{–}(2)\ \text{restricted to a composite's boundary data};
\]

%dsp-level:0
\[
\qquad \text{the sum below is finite and exact at every budget}\ B\ [\mathrm{Def\text{-}H.32}(1):\ \mathcal G_B=\{\Gamma:\ V(\Gamma)\le B\},\ \mathfrak C_B(\Gamma)=\{R:\ |R|\le B_K\}];
\]

%dsp-level:0
\[
(2)\quad\textbf{the path integral:}
\]

%dsp-level:0
\[
\qquad Z(\mathcal X)\ =\ \sum_{R\in\mathcal E(\mathcal X)}\ \sum_{\sigma}\ w(R,\sigma),
\]

%dsp-level:0
\[
\qquad w(C)=\mathcal J(C)\prod_{e\in E(C)}G(e)\cdot(-1)^{\aleph(\sigma,C)}
\]

%dsp-level:0
\[
\qquad [\mathrm{Def\text{-}H.32}(2);\ \mathrm{A8}];
\]

%dsp-level:0
\[
\text{the two layers (norm and amplitude) are kept separate}\ [\mathrm{CT2.8};\ \mathrm{Def\text{-}H.32}(2)]\ \text{and no layer is substituted for the other};
\]

%dsp-level:0
\[
(3)\quad\textbf{the weights are derived, not declared:}\qquad\text{the per-pair Janus master generator}\ Q=\bigl(\begin{smallmatrix}-x_+&x_+\\ x_-&-x_-\end{smallmatrix}\bigr)\ \text{over}\ \mathbb Z[x_+,x_-]\ \text{has a one-dimensional kernel spanned by}\ (x_-,x_+)\ [\mathrm{FT.108}(4)]:
\]

%dsp-level:0
\[
p\ =\ \frac{x_+}{x_++x_-}\qquad[\text{derived}];\qquad x_+=1-4\varepsilon_q,\quad x_-=4\varepsilon_q\implies p=1-4\varepsilon_q\ \ \text{identically}\qquad[\mathrm{A4},\ \mathrm{D0.4};\ \text{machine}];
\]

%dsp-level:0
\[
\textbf{no P-class datum is added:}\ \ \varepsilon_q\ \text{is the A4 crystallization threshold}\ [\mathrm{D0.4}:\ \text{admissible range}\ [0,\tfrac14)\ [\mathrm{CT0.8}];\ \text{sharp choice}\ \varepsilon_q=0];
\]

%dsp-level:0
\[
(4)\quad\textbf{the information content is distributed:}
\]

%dsp-level:0
\[
\qquad E_{\mathrm{int}}(R)
\]

%dsp-level:0
\[
= |R|
\]

%dsp-level:0
\[
= n_{\mathcal R}\ [\mathrm{Def\text{-}H.32}(1):\ \text{the budget is the ledger's own content measure};\ \mathrm{CT1.29}:\ S_{\mathfrak L}=n_{\mathcal R}\ln2];
\]

%dsp-level:0
\[
\text{the crystallized value}\ E_{\mathrm{int}}
\]

%dsp-level:0
\[
= 6\ [\text{Census FT.35}(2)]\ \text{is one point of the distribution},
\]

%dsp-level:0
\[
\ \neg\text{the ensemble};
\]

%dsp-level:0
\[
\qquad \text{the move system conserves the grade}\ a_+-a_-=|R'|-|R|\ [\mathrm{FT.108}(2)];
\]

%dsp-level:0
\[
(5)\quad\textbf{the no-flicker limit (scope):}
\]

%dsp-level:0
\[
\qquad \varepsilon_q
\]

%dsp-level:0
\[
= 0
\]

%dsp-level:0
\[
\implies x_-
\]

%dsp-level:0
\[
= 0\ \text{exactly}\implies\neg\exists\,
\]

%dsp-level:0
\[
\text{dissolution move}\implies\mathcal E(\mathcal X)\ \text{collapses to its fully crystallized member}\ [\text{machine}];
\]

%dsp-level:0
\[
\text{every}\ \varepsilon_q>0\ \text{gives}\ x_->0\ \text{and a strictly positive flicker weight};
\]

%dsp-level:0
\[
\qquad \text{the}\ B\to\infty\ \text{limit is}\ \neg\text{constructed}\ [\mathrm{Def\text{-}H.32}(1);\ \text{growth data: FT.105}].
\]

\textbf{Rmk (ensemble discipline).} The ensemble is a sum over crystallization patterns at fixed boundary data, not a sum over species: every member carries the composite's charge class, sealing-phase class and winding class. The topology labels are internal bookkeeping [Def-H.32(4) pattern], and no member of any ensemble is identified with an external object [P5; Def-BM.3(6)].

\textbf{Req.} A4, A8 (Part 1); D0.4, D2.9 (Parts 1–2); CT0.8, CT1.29 (Part 1); CT2.8 (Part 2); Def-H.32, FT.105, FT.108(2), FT.108(4), FT.108(5) (Part 3); Def-BM.3, BM.2(1) (this part); Census FT.35(2) (row-wise); exact polynomial kernel arithmetic over $\mathbb Q[x_+,x_-]$, exact rational specialization [math; machine: \texttt{checks/\allowbreak{}verify\_\allowbreak{}H\_\allowbreak{}topo\_\allowbreak{}blind\_\allowbreak{}neutron.py} §A].

\textbf{Pf.}

%dsp-level:0
\[
\text{Definition block: the constitutive clause is}\ (1)\text{'s sector-conditioned configuration set};
\]

%dsp-level:0
\[
\ \text{the inherited clauses are Def-H.32's}\ (1)\text{–}(2)\ \text{and scopes}\ (\mathrm{S1})\text{–}(\mathrm{S4}),
\]

%dsp-level:0
\[
\ \text{already declared};
\]

%dsp-level:0
\[
\ \text{clause}\ (3)\ \text{is derived content recorded at the definition site};
\]

%dsp-level:0
\[
\ (5)\ \text{is a scope clause cited at use}.\ \blacksquare
\]

\textbf{Phys.} A sealed composite is not one topology with a number attached; it is a finite exact sum over crystallization patterns that share its boundary data — the same topology-blind path integral the fermion era built, now conditioned on a composite's charge class, sealing phase and winding class. Nothing is postulated to weight the sum: the Janus move generator has a one-dimensional kernel, its span is the stationary measure, and substituting the crystallization threshold gives the flicker weight as $4\varepsilon_q$ with no new free parameter anywhere. Set that threshold to its sharp value and the ensemble collapses to the single crystallized configuration — the flicker is a consequence of the threshold being nonzero, not an extra ingredient.

\end{definition}

\begin{definition}[{\texorpdfstring{Def-BM.6 — The family naming rows (extension of the identification layer)}{Def-BM.6 - The family naming rows (extension of the identification layer)}}]\label{tag:Def-BM-6}
\textbf{Def.} An extension of Def-BM.3's identification layer to the three derived channel triplets, recorded by a logged declaration as Def-BM.3(7) requires. A name attaches to a \textit{line of a derived spectrum}, ordered by that spectrum; no name attaches to a mass value.

%dsp-level:0
\[
(1)\quad\textbf{the}\ \ell\text{-channel triplet}\ [\mathrm{FT.109}\text{–}\mathrm{FT.113}:\ r_i=1+\sqrt2\cos(\delta_\ell+\tfrac{2\pi i}3),\ \delta_\ell=\tfrac29]:
\]

%dsp-level:0
\[
\qquad \textbf{electron}\ :=\ \text{the}\ i=0\ \text{line},\quad\textbf{muon}\ :=\ \text{the}\ i=1\ \text{line},\quad\textbf{tau}\ :=\ \text{the}\ i=2\ \text{line};
\]

%dsp-level:0
\[
\text{the order rule is the spectrum's own:}\ r_0>r_1>r_2\ \text{at}\ \delta_\ell=\tfrac29\ [\mathrm{FT.113};\ \text{the span is}\ 3477\text{-fold}];
\]

%dsp-level:0
\[
(2)\quad\textbf{the}\ \nu\text{-channel triplet}\ [\mathrm{FT.142}\text{–}\mathrm{FT.144}:\ \delta_\nu=\tfrac18\ \text{on the derived}\ 16\text{-arena}]:
\]

%dsp-level:0
\[
\qquad \textbf{electron neutrino},\ \textbf{muon neutrino},\ \textbf{tau neutrino}\ :=\ \text{its three lines in the same spectral order};
\]

%dsp-level:0
\[
(3)\quad\textbf{the}\ q\text{-channel triplet}\ [\mathrm{FT.118},\ \mathrm{FT.119}:\ \delta_q=\tfrac2{27},\ \text{triplet-attached}]:
\]

%dsp-level:0
\[
\qquad \textbf{the quark triplet}\ :=\ \text{its three lines in the same order};
\]

%dsp-level:0
\[
\qquad \text{the cross-channel identity}\ 16\delta_\nu
\]

%dsp-level:0
\[
= 9\delta_\ell
\]

%dsp-level:0
\[
= 27\delta_q
\]

%dsp-level:0
\[
= 2\ [\mathrm{FT.143}]\ \text{is the derived relation between the three rows};
\]

%dsp-level:0
\[
(4)\quad\textbf{the derived content the names attach to:}
\]

%dsp-level:0
\[
\qquad \sum_ir_i=3=\dim F_{J^\pm}\ [\mathrm{FT.111}(3)];
\]

%dsp-level:0
\[
\qquad \sum_ir_i^2=6;
\]

%dsp-level:0
\[
\qquad Q:=\frac{\sum_ir_i^2}{(\sum_ir_i)^2}=\frac69=\frac23\ \ \textbf{exactly, for every phase}\ [\text{machine};\ \mathrm{FT.120}(6)];
\]

%dsp-level:0
\[
\text{the cosine form is forced by}\ C_3\text{-covariance and the normalization by the seal dimension}\ [\mathrm{FT.111}(2)\text{–}(3);\ \mathrm{FT.158}(4)];
\]

%dsp-level:0
\[
\qquad \text{the angle is a ratio of spatial counts}\ \delta_\ell=(3-1)/(3\cdot3)\ [\mathrm{FT.158}(2)\text{–}(3)];
\]

%dsp-level:0
\[
(5)\quad\textbf{the non-derived clause (the mass form):}
\]

%dsp-level:0
\[
\qquad \text{the bridge from a spectral line to a mass value}\ (\sqrt{m}\propto r\ \text{vs}\ m\propto r\ \text{vs}\ \neg\text{a bridge})\ \text{is item-13 content and is}\ \neg\text{adopted here}\;
\]

%dsp-level:0
\[
\text{the standing bar}\ \text{amplitude}\ne\text{mass}\ [\mathrm{FT.79}(5)]\ \text{is untouched};
\]

%dsp-level:0
\[
\qquad \text{the channel scale}\ \mu\ \text{remains an undeclared P4-class parameter}\ [\text{Cal-F}]:\ \neg\text{a mass value is derived for any named line};
\]

%dsp-level:0
\[
(6)\quad\textbf{the data bar:}
\]

%dsp-level:0
\[
\qquad \neg\exists\ \text{a measured external value in any clause}\ [\mathrm{P5}];
\]

%dsp-level:0
\[
\qquad \text{the ratios}\ (Q,\ \text{the span, the angle identities})\ \text{are corpus-derived and are}\ \neg\text{compared to measurement here};
\]

%dsp-level:0
\[
(7)\quad\textbf{scope:}
\]

%dsp-level:0
\[
\qquad \text{the names occur in Part 4 only}\ [\mathrm{Def\text{-}BM.3}(7)\ \text{extended}];
\]

%dsp-level:0
\[
\qquad \text{further extension requires a further logged declaration}.
\]

\textbf{Rmk (what the naming does and does not do).} Naming the $\ell$-channel lines \textit{electron, muon, tau} records that the corpus derives a three-line $C_3$-covariant spectrum whose Koide ratio is exactly $2/3$ at every phase — the ratio the family is famous for — and that the lines are strictly ordered with a 3477-fold span. It does not derive their masses, it does not adopt a mass form, and it does not compare any ratio to a measured value. The naming is Def-BM.3's labelling act applied to spectral lines instead of junctions.

\textbf{Req.} FT.79(5), FT.109, FT.110, FT.111, FT.112, FT.113, FT.118, FT.119, FT.120(6), FT.142, FT.143, FT.144, FT.158 (Part 3); Def-BM.3 (this part); P4, P5 (Parts 1–2); Cal-F; exact identities in $\mathbb Q[c,s,r]/(c^2{+}s^2{-}1,\ r^2{-}3)$ [math; machine: \texttt{checks/\allowbreak{}verify\_\allowbreak{}H\_\allowbreak{}topo\_\allowbreak{}blind\_\allowbreak{}neutron.py} §F].

\textbf{Pf.}

%dsp-level:0
\[
\text{Definition block: the constitutive clause is the name}\leftrightarrow\text{line attachment by spectral order};
\]

%dsp-level:0
\[
\ \text{the non-derived clause is}\ (5)\ \text{(the mass form is item-13 content,}\ \neg\text{adopted)};
\]

%dsp-level:0
\[
\ \text{clause}\ (4)\text{'s identities are derived}\ [\text{machine:}\ \textstyle\sum_i\cos(\theta+\tfrac{2\pi i}3)=0,\ \sum_i\cos^2=\tfrac32\ \text{in the quotient ring}]\ \text{and are recorded here as the content the names attach to}.\ \blacksquare
\]

\textbf{Phys.} Three channels, three triplets, one mechanism: a $C_3$-covariant seal bilinear whose shape covariance forces, whose normalization the seal dimension fixes, and whose angle is a ratio of spatial counts. The ratio the charged-lepton family is known for falls out as $6/9$ — exactly two thirds, at every phase, with nothing fitted — and the lines are strictly ordered across a 3477-fold span. The names now attach to those lines in spectral order. What the names do not carry is a mass: the bridge from a spectral line to a mass value is the corpus's open item, the amplitude-is-not-mass bar stands, and the overall scale is still an undeclared internal parameter.

\end{definition}

\begin{theorem}[{\texorpdfstring{BM.19 — The flicker ensemble is derived: the stationary measure, the threshold specialization, and the no-flicker limit}{BM.19 - The flicker ensemble is derived: the stationary measure, the threshold specialization, and the no-flicker limit}}]\label{tag:BM-19}
\textbf{Thm.}

%dsp-level:0
\[
(1)\quad\textbf{the measure is forced:}\qquad Q=\begin{pmatrix}-x_+&x_+\\ x_-&-x_-\end{pmatrix}\ \text{over}\ \mathbb Z[x_+,x_-]:\qquad(x_-,\,x_+)\,Q=(0,0)\ \ \text{exactly}\ [\text{machine}];
\]

%dsp-level:0
\[
\neg\exists\ \text{a nonzero constant vector in}\ \ker Q\ [\text{machine: exhaustive on}\ [-3,3]^2]:
\]

%dsp-level:0
\[
\qquad \dim\ker Q=1\ \text{over the fraction field}\implies p=\frac{x_+}{x_++x_-}\ \text{is the unique stationary measure}\ [\mathrm{FT.108}(4)];
\]

%dsp-level:0
\[
(2)\quad\textbf{the specialization is an identity:}
\]

%dsp-level:0
\[
\qquad x_+=1-4\varepsilon_q,\quad x_-=4\varepsilon_q\implies p=\frac{1-4\varepsilon_q}{1-4\varepsilon_q+4\varepsilon_q}=1-4\varepsilon_q
\]

%dsp-level:0
\[
\qquad [\text{machine: exact at six rational thresholds}];
\]

%dsp-level:0
\[
\text{the flicker (dissolution) weight is}\ x_-=4\varepsilon_q:
\]

%dsp-level:0
\[
\qquad \text{admissible on}\ \varepsilon_q\in[0,\tfrac14)\ [\mathrm{CT0.8}]\implies x_-\in[0,1),\ x_+\in(0,1]\ [\text{machine}];
\]

%dsp-level:0
\[
\textbf{no P-class datum enters:}\ \ \varepsilon_q\ \text{is the A4 threshold}\ [\mathrm{D0.4}];
\]

%dsp-level:0
\[
\qquad \neg\text{a new fugacity,}\ \neg\text{a new rate,}\ \neg\text{a P4 addition}\ [\text{Cal-BM stays empty}];
\]

%dsp-level:0
\[
(3)\quad\textbf{the no-flicker limit:}
\]

%dsp-level:0
\[
\qquad \varepsilon_q
\]

%dsp-level:0
\[
= 0\ [\mathrm{D0.4}\ \text{sharp choice}]
\]

%dsp-level:0
\[
\implies x_-
\]

%dsp-level:0
\[
= 0\ \text{exactly}\implies\neg\exists\,
\]

%dsp-level:0
\[
\text{a dissolution move}\implies\mathcal E(\mathcal X)
\]

%dsp-level:0
\[
= \{\text{the crystallized configuration}\}\ [\mathrm{Def\text{-}BM.5}(5);\ \text{machine}];
\]

%dsp-level:0
\[
\varepsilon_q>0\implies x_->0:
\]

%dsp-level:0
\[
\qquad \text{the flicker is strictly present at every nonzero threshold and strictly absent at the sharp one};
\]

%dsp-level:0
\[
(4)\quad\textbf{the ensemble is finite and exact at each budget:}
\]

%dsp-level:0
\[
\qquad |\mathfrak C_B(\Gamma)|<\infty\ [\mathrm{Def\text{-}H.32}(1)];
\]

%dsp-level:0
\[
\qquad \text{the}\ B=4\ \text{sector census is exact integers}\ [\mathrm{FT.106}(3):\ Z^{\mathrm{norm}}(q,S)\ \text{and}\ Z^{\mathrm{amp}}(q,S)\ \text{per sector}];
\]

%dsp-level:0
\[
Z^{\mathrm{norm}}(q,S)
\]

%dsp-level:0
\[
= Z^{\mathrm{norm}}(-q,S)\ \text{unconditionally and}\ Z^{\mathrm{amp}}(q,S)
\]

%dsp-level:0
\[
= Z^{\mathrm{amp}}(-q,S)\ \text{on the cover support}\ [\mathrm{FT.106}(3\mathrm d)];
\]

%dsp-level:0
\[
\qquad Z^{\mathrm{amp}}(q,k)=0\ \text{for}\ k=2\ \text{and odd}\ k\ge3\ [\mathrm{FT.106}(3\mathrm c)];
\]

%dsp-level:0
\[
(5)\quad\textbf{the information content is a distribution:}
\]

%dsp-level:0
\[
\qquad E_{\mathrm{int}}(R)=|R|\ [\mathrm{Def\text{-}BM.5}(4)]:
\]

%dsp-level:0
\[
\qquad \text{the ensemble spreads}\ E_{\mathrm{int}}\ \text{over}\ \{0,\dots,B_K\}\ \text{and the crystallized value is its top point};
\]

%dsp-level:0
\[
a_+-a_-=|R'|-|R|,
\]

%dsp-level:0
\[
\qquad \textstyle\sum_dm_d-\sum_dn_d=V'-V,
\]

%dsp-level:0
\[
\qquad \sum_d d\,(m_d-n_d)=E'_{\mathrm{pot}}-E_{\mathrm{pot}}
\]

%dsp-level:0
\[
\qquad [\mathrm{FT.108}(2):\ \text{the grade map is a homomorphism of the move system}]:
\]

%dsp-level:0
\[
\text{every flicker move changes}\ E_{\mathrm{int}}\ \text{by exactly the}\ J^\pm\ \text{excess it carries}\ —\ \text{the distribution is the move algebra's orbit, not an ansatz};
\]

%dsp-level:0
\[
(6)\quad\textbf{the composite is not a fixed topology:}
\]

%dsp-level:0
\[
\qquad \mathcal E(\mathcal X)\ \text{has more than one member}\iff x_->0\iff\varepsilon_q>0\ [(3)]:
\]

%dsp-level:0
\[
\qquad \text{at every admissible nonzero threshold the neutron and the proton are ensembles}.
\]

\textbf{Req.} A4 (Part 1); D0.4 (Part 1); CT0.8, CT1.29 (Part 1); CT2.8 (Part 2); Def-H.32, FT.105, FT.106(3), FT.108(2), FT.108(4), FT.108(5) (Part 3); Def-BM.5 (this part); P4 (Parts 1–2); §Cal-BM; exact polynomial arithmetic over $\mathbb Q[x_+,x_-]$, exhaustive integer kernel census, exact rational specialization [math; machine: \texttt{checks/\allowbreak{}verify\_\allowbreak{}H\_\allowbreak{}topo\_\allowbreak{}blind\_\allowbreak{}neutron.py} §A].

\textbf{Pf.}

%dsp-level:0
\[
(1):\ (x_-)(-x_+)+(x_+)(x_-)=0\ \text{and}\ (x_-)(x_+)+(x_+)(-x_-)=0\ \text{in}\ \mathbb Q[x_+,x_-]\ [\text{machine}];
\]

%dsp-level:0
\[
\qquad \text{constant vectors}: a(-x_+)+bx_-=0\ \wedge\ ax_++b(-x_-)=0\implies a=b=0\ [\text{machine: exhaustive}]
\]

%dsp-level:0
\[
(2):\ \text{substitution}\ [\text{machine at}\ \varepsilon_q\in\{\tfrac1{64},\tfrac1{16},\tfrac18,\tfrac3{16},\tfrac15,\tfrac{63}{256}\}];
\]

%dsp-level:0
\[
\qquad \varepsilon_q\in[0,\tfrac14)\ [\mathrm{CT0.8}]\implies4\varepsilon_q\in[0,1)\ [\text{machine}]
\]

%dsp-level:0
\[
(3):\ x_-=4\cdot0=0\ [\text{machine}];
\]

%dsp-level:0
\[
\qquad \text{a move system with no}\ J^-\ \text{generator admits no dissolution}\ [\mathrm{D1.2}:\ J^-:n_{\mathcal R}:1\to0];
\]

%dsp-level:0
\[
\qquad \text{the ensemble reduces to the crystallized pattern}\ [\mathrm{Def\text{-}BM.5}(5)]
\]

%dsp-level:0
\[
(4):\ \mathrm{Def\text{-}H.32}(1)\text{'s filtration and}\ \mathrm{FT.106}(3)\text{'s census rows re-cited}\ (\text{exact integers, frozen});
\]

%dsp-level:0
\[
\qquad (5):\ \mathrm{FT.108}(2)\text{'s three grade identities re-cited}\ [\text{machine rows frozen at Tier F.36}]
\]

%dsp-level:0
\[
(6):\ |\mathcal E(\mathcal X)|>1\iff\exists\,\text{a dissolution move}\iff x_->0\ [(2),(3)].\ \blacksquare
\]

\textbf{Phys.} Nothing in the flicker is chosen. The Janus move generator has a kernel, that kernel is a line, and its span is the stationary measure — so the weight of a crystallization against a dissolution is the ratio $x_+/(x_++x_-)$ before anyone names a rate. Substituting the crystallization threshold the axioms already carry turns that ratio into $1-4\varepsilon_q$, an identity rather than a fit, and the flicker weight into $4\varepsilon_q$. At the sharp threshold the dissolution weight is exactly zero and the ensemble collapses to one configuration; at any nonzero threshold the composite is genuinely an ensemble, and the move algebra's grade identities say precisely how each flicker step moves the information content. That is the whole mechanism of "the neutron is not a fixed topology", with no new parameter anywhere in it.

\end{theorem}

\begin{theorem}[{\texorpdfstring{BM.20 — The mass from the curvature law on the path-integrated route}{BM.20 - The mass from the curvature law on the path-integrated route}}]\label{tag:BM-20}
\textbf{Thm.}

%dsp-level:0
\[
(1)\quad\textbf{the mass is the curvature flux:}\qquad R=\kappa^2(\rho-3p)+4\Lambda\ [\mathrm{BM.15}(2)];\qquad p=0,\ \Lambda=0:\quad R=\kappa^2\rho,\quad\rho V=m:
\]

%dsp-level:0
\[
\boxed{\ R\,V=\kappa^2 m\ }\qquad\text{i.e.}\qquad m=\frac1{\kappa^2}\int_\Sigma R\,dV\qquad[\text{machine: exact monomial identity}];
\]

%dsp-level:0
\[
(2)\quad\textbf{the}\ \Lambda\ \textbf{clause:}
\]

%dsp-level:0
\[
\qquad (R-4\Lambda)V=\kappa^2m\ \ \text{identically}\ [\text{machine}]:
\]

%dsp-level:0
\[
\qquad \Lambda\ \text{is the one additive freedom}\ [\mathrm{P4};\ \mathrm{BM.15}(6)];
\]

%dsp-level:0
\[
(3)\quad\textbf{the ensemble form (the path-integrated route):}
\]

%dsp-level:0
\[
\qquad \langle m\rangle_{\mathcal X}=\frac1{\kappa^2}\Bigl\langle\int_\Sigma R\,dV\Bigr\rangle_{\mathcal X}=\sum_{R\in\mathcal E(\mathcal X)}w(R)\,m(R),
\]

%dsp-level:0
\[
\qquad \sum_Rw(R)=1
\]

%dsp-level:0
\[
\qquad [\text{machine: the normalized average equals the averaged flux}];
\]

%dsp-level:0
\[
\text{the mass of a sealed composite is the ensemble average of the curvature flux over its flicker}\ —\ \neg\text{the flux of the crystallized member};
\]

%dsp-level:0
\[
(4)\quad\textbf{the derived sector read:}
\]

%dsp-level:0
\[
\qquad m^2_B(q,S)=\frac{\sum_{\Gamma,R}\nu\,t^{|R|}u^L\#\mathrm{inst}\cdot\tfrac32(2E-|R|)}{\sum_{\Gamma,R}\nu\,t^{|R|}u^L\#\mathrm{inst}}\bigg|_{t=u=1}
\]

%dsp-level:0
\[
\qquad [\mathrm{FT.107}(3);\ \Delta S=\tfrac32n_{\mathcal V}\ [\mathrm{FT.107}(2)]]:
\]

%dsp-level:0
\[
m^2_4(\cdot,0)=\frac{843957}{98956},
\]

%dsp-level:0
\[
\quad m^2_4(\cdot,1)=\frac{1921662}{393679},
\]

%dsp-level:0
\[
\quad m^2_4(\cdot,2)=\frac{91}{23},
\]

%dsp-level:0
\[
\quad m^2_4(\cdot,3)=\frac{13713}{3088}
\]

%dsp-level:0
\[
\qquad [\text{exact rationals};\ \text{machine}];
\]

%dsp-level:0
\[
\text{charge-degenerate in}\ q\ [\mathrm{FT.106}(3\mathrm a),(3\mathrm d)];
\]

%dsp-level:0
\[
\qquad \text{all strictly positive};
\]

%dsp-level:0
\[
\qquad \text{the transit sector}\ (S=0)\ \text{is heaviest}\ [\text{machine}];
\]

%dsp-level:0
\[
(5)\quad\textbf{the bound:}
\]

%dsp-level:0
\[
\qquad 0<m^2_B(q,S)\le\tfrac32\max_\Gamma n_{\mathcal V}\le\tfrac32B(B-1)=18\ \ (B=4);
\]

%dsp-level:0
\[
\qquad \text{the census maximum is}<9
\]

%dsp-level:0
\[
\qquad [\mathrm{FT.107}(4);\ \text{machine}];
\]

%dsp-level:0
\[
(6)\quad\textbf{the curvature scale of the ensemble read:}
\]

%dsp-level:0
\[
\qquad \rho_B=\frac{m_B}{V_{\mathrm{eff}}},
\]

%dsp-level:0
\[
\qquad \kappa^2\rho_B=\frac{2\pi\,\mathfrak h\,m_B}{V_{\mathrm{eff}}}\ \Big/\ \frac{\mathfrak h}{1}=\frac{2\pi\,E^4_{\mathrm{eff}}(\ln2)^3}{\mathfrak h^2}\bigg|_{E_{\mathrm{eff}}:=m_B\mathfrak h/\ln2}
\]

%dsp-level:0
\[
\text{i.e. the sector read enters BM.14's and BM.15's formulas through the effective count}\ E_{\mathrm{eff}}:
\]

%dsp-level:0
\[
= m_B\mathfrak h/\ln2\ [\text{exact rational per sector}]:
\]

%dsp-level:0
\[
\qquad \text{the lock of BM.16(1) applies sector by sector}.
\]

\textbf{Rmk (route discipline).} Clause (1) is the curvature law read as a mass definition — the same law BM.15 derived from CT1.23 — applied configuration by configuration and then averaged, which is what "the path-integrated route" means here. It does not import a mass datum: $m(R)$ at each member is the informational read [FT.44(4)], and the ensemble average is an exact rational combination of integers times $\ln2/\mathfrak h$; the sector read of clause (4) is the deficit-average route [FT.107(3)], and the two coincide where both are defined.

\textbf{Req.} D2.9 (Parts 1–2); CT1.23, CT2.1, CT2.9 (Part 2); P4 (Parts 1–2); Def-H.32, FT.106(3), FT.107(2), FT.107(3), FT.107(4) (Part 3); FT.44(4) (Census); Def-BM.4, Def-BM.5, BM.14(2), BM.14(3), BM.15(2), BM.15(6), BM.16(1), BM.19 (this part); exact Laurent-monomial identities, exact rational ensemble averages [math; machine: \texttt{checks/\allowbreak{}verify\_\allowbreak{}H\_\allowbreak{}topo\_\allowbreak{}blind\_\allowbreak{}neutron.py} §B].

\textbf{Pf.}

%dsp-level:0
\[
(1):\ \mathrm{BM.15}(2)\ \text{at}\ p=0,\ \Lambda=0:\ R=\kappa^2\rho;
\]

%dsp-level:0
\[
\qquad \rho
\]

%dsp-level:0
\[
= m/V\ [\mathrm{BM.14}(2)\ \text{pattern}]
\]

%dsp-level:0
\[
\implies RV
\]

%dsp-level:0
\[
= \kappa^2m\ [\text{machine: substitution of the two defining relations into}\ RV-\kappa^2m\ \text{gives}\ 0\ \text{identically}]
\]

%dsp-level:0
\[
(2):\ R=\kappa^2\rho+4\Lambda\implies(R-4\Lambda)V=\kappa^2\rho V=\kappa^2m\ [\text{machine: the residual polynomial is}\ 0]
\]

%dsp-level:0
\[
(3):\ \text{linearity of the flux in}\ m\ \text{at fixed}\ V,\kappa^2\ [(1)]\implies\langle RV\rangle/\kappa^2=\langle m\rangle;
\]

%dsp-level:0
\[
\qquad \text{normalization}\ \sum w=1\ [\text{machine witness}]
\]

%dsp-level:0
\[
(4)\text{–}(5):\ \mathrm{FT.107}(2)\text{–}(4)\ \text{rows re-cited verbatim}\ (\text{exact rationals, frozen});
\]

%dsp-level:0
\[
\qquad \text{positivity, the ordering and the bound re-verified on the stated fractions}\ [\text{machine}]
\]

%dsp-level:0
\[
(6):\ E_{\mathrm{eff}}:=m_B\mathfrak h/\ln2\ \text{is the count whose informational mass is}\ m_B\ [\mathrm{FT.44}(4)\ \text{inverted}];
\]

%dsp-level:0
\[
\qquad \mathrm{BM.14}(3)/\mathrm{BM.15}(4)\ \text{re-cited at}\ E=E_{\mathrm{eff}}.\ \blacksquare
\]

\textbf{Phys.} Mass is curvature flux. Trace the field equations for a pressureless composite and the scalar curvature times the volume is $\kappa^2$ times the mass, with the cosmological term the only additive ambiguity — so the mass read the corpus derived from the dust stress tensor is the same object as the integral of its own curvature. Run that identity through the flicker and the mass of a composite becomes the ensemble average of the curvature flux, not the flux of its crystallized snapshot. The topology-blind census supplies the average independently, as an exact rational per sector: four fractions at budget four, all positive, all charge-degenerate, with the transit sector heaviest and the whole table under the quadratic budget bound. Both routes meet in one formula, and the calibration lock of the previous tier applies sector by sector through the effective count.

\end{theorem}

\begin{theorem}[{\texorpdfstring{BM.21 — The binding energy: additive in the count, non-additive in the ensemble}{BM.21 - The binding energy: additive in the count, non-additive in the ensemble}}]\label{tag:BM-21}
\textbf{Thm.}

%dsp-level:0
\[
(1)\quad\textbf{the informational count is exactly additive:}
\]

%dsp-level:0
\[
\qquad E_{\mathrm{int}}(\Sigma_{12})=6=(c_1{:}\,2)+(c_2{:}\,2)+(\text{connector}{:}\,2)
\]

%dsp-level:0
\[
\qquad [\text{Census FT.35}(2);\ \text{machine}];
\]

%dsp-level:0
\[
m_{\mathrm{hid}}
\]

%dsp-level:0
\[
= E_{\mathrm{int}}\ln2/\mathfrak h\ [\mathrm{FT.44}(4)]\ \text{is linear in the count}\implies\textbf{the informational binding read vanishes exactly:}
\]

%dsp-level:0
\[
\quad E^{\mathrm{info}}_{\mathrm{bind}}=\bigl(6-(2{+}2{+}2)\bigr)\tfrac{\ln2}{\mathfrak h}=0
\]

%dsp-level:0
\[
\qquad [\text{machine}];
\]

%dsp-level:0
\[
(2)\quad\textbf{the ensemble read is not additive:}
\]

%dsp-level:0
\[
\qquad \mathcal B(S_j;S_1,S_2,S_3)\ :=\ m^2_B(\cdot,S_j)-\sum_{i=1}^{3}m^2_B(\cdot,S_i)
\]

%dsp-level:0
\[
\qquad [\mathrm{BM.20}(4)]:
\]

%dsp-level:0
\[
\mathcal B(1;1,1,1)=-2\,m^2_4(\cdot,1)=-\frac{3843324}{393679};
\]

%dsp-level:0
\[
\qquad \mathcal B(2;1,1,1)=\frac{91}{23}-\frac{5764986}{393679};
\]

%dsp-level:0
\[
\qquad \mathcal B(3;1,1,1)=\frac{13713}{3088}-\frac{5764986}{393679};
\]

%dsp-level:0
\[
\mathcal B(0;1,1,1)=\frac{843957}{98956}-\frac{5764986}{393679};
\]

%dsp-level:0
\[
\qquad \textbf{every one is a nonzero exact rational}\ [\text{machine}];
\]

%dsp-level:0
\[
(3)\quad\textbf{the binding is the flicker's, not the count's:}
\]

%dsp-level:0
\[
\qquad \varepsilon_q
\]

%dsp-level:0
\[
= 0\implies\mathcal E\ \text{collapses}\ [\mathrm{BM.19}(3)]\implies\text{the ensemble average reduces to the crystallized member}\implies\mathcal B\to0\ \text{with}\ (1);
\]

%dsp-level:0
\[
\varepsilon_q>0\implies\mathcal B\ne0\ \text{exactly}\ [(2)]:
\]

%dsp-level:0
\[
\qquad \textbf{a non-flickering composite has no binding read, and a flickering one has a nonzero exact one};
\]

%dsp-level:0
\[
(4)\quad\textbf{the sign structure:}
\]

%dsp-level:0
\[
\qquad \text{all four combinations are bound:}\quad\mathcal B(S_j;1,1,1)<0\ \ \forall S_j\in\{0,1,2,3\}
\]

%dsp-level:0
\[
\qquad [\text{machine: exact rational sign census}];
\]

%dsp-level:0
\[
\text{the exact order of the sector masses is}\ m^2_4(\cdot,0)>m^2_4(\cdot,1)>m^2_4(\cdot,3)>m^2_4(\cdot,2)\ [\text{machine}]:
\]

%dsp-level:0
\[
\qquad \text{three separated constituents outweigh one junction in every sector combination};
\]

%dsp-level:0
\[
\text{the like-sector case is exactly}\ -2\,m^2_4(\cdot,1)\ [\text{machine}];
\]

%dsp-level:0
\[
\qquad |\mathcal B(S_j)|
\]

%dsp-level:0
\[
= 3m^2_4(\cdot,1)-m^2_4(\cdot,S_j)\ \text{is largest for the lightest junction sector}\ (S_j=2)\ \text{and smallest for the heaviest}\ (S_j=0)\ [\text{machine}];
\]

%dsp-level:0
\[
(5)\quad\textbf{the ceiling inheritance:}
\]

%dsp-level:0
\[
\qquad |\mathcal B|
\]

%dsp-level:0
\[
\le 4\cdot\tfrac32B(B-1)
\]

%dsp-level:0
\[
= 72\ \ \text{at}\ B
\]

%dsp-level:0
\[
= 4\ \text{from}\ \mathrm{FT.107}(4)\text{'s bound applied to all four terms}\ [\text{machine: exact bound arithmetic}];
\]

%dsp-level:0
\[
\text{the winding spread at}\ B=4\ \text{is}\ \tfrac{91}{23}\le m^2\le\tfrac{843957}{98956}\ [\mathrm{FT.107}(4)]:
\]

%dsp-level:0
\[
\qquad \text{the binding read stays order-one while the bound grows as}\ \Theta(B^2);
\]

%dsp-level:0
\[
(6)\quad\textbf{the read discipline:}
\]

%dsp-level:0
\[
\qquad \mathcal B\ \text{is a difference of the corpus's own}\ m^2_B\ \text{reads}\ [\mathrm{FT.107}(3)];
\]

%dsp-level:0
\[
\qquad \neg\text{a mass value in MeV,}\ \neg\text{an external datum},\ \neg\text{a comparison}\ [\mathrm{P5}];
\]

%dsp-level:0
\[
\text{the pole-mass side}\ (\text{defect angles, ceilings})\ \text{is a different read and remains item-13 bridge content}\ [\mathrm{FT.76}(3);\ \mathrm{Def\text{-}H.39}].
\]

\textbf{Req.} D2.9 (Parts 1–2); CT2.1 (Part 2); P5 (Parts 1–2); Def-H.32, Def-H.39, FT.76(3), FT.107(2), FT.107(3), FT.107(4) (Part 3); FT.44(4) (Census); Census FT.35(2) (row-wise); Def-BM.5, BM.19(3), BM.20(1), BM.20(4) (this part); exact rational arithmetic [math; machine: \texttt{checks/\allowbreak{}verify\_\allowbreak{}H\_\allowbreak{}topo\_\allowbreak{}blind\_\allowbreak{}neutron.py} §C].

\textbf{Pf.}

%dsp-level:0
\[
(1):\ \text{Census FT.35}(2)\text{'s decomposition row re-cited}\ (2{+}2{+}2=6);
\]

%dsp-level:0
\[
\qquad m_{\mathrm{hid}}\ \text{linear in}\ E_{\mathrm{int}}\ [\mathrm{FT.44}(4)]\implies\text{the difference is}\ 0\cdot\tfrac{\ln2}{\mathfrak h}\ [\text{machine}]
\]

%dsp-level:0
\[
(2):\ \text{exact rational subtraction of}\ \mathrm{FT.107}(3)\text{'s four frozen values}\ [\text{machine}];
\]

%dsp-level:0
\[
\qquad \text{nonvanishing: each numerator is nonzero after clearing denominators}\ [\text{machine}]
\]

%dsp-level:0
\[
(3):\ \mathrm{BM.19}(3)\ \text{re-cited};
\]

%dsp-level:0
\[
\qquad \text{a one-member ensemble has}\ \langle m\rangle=m\ \text{and the count-additivity of}\ (1)\ \text{applies}
\]

%dsp-level:0
\[
(4):\ \mathrm{FT.107}(3)\text{'s four values in exact descending order}\ \tfrac{843957}{98956}>\tfrac{1921662}{393679}>\tfrac{13713}{3088}>\tfrac{91}{23},
\]

%dsp-level:0
\[
\ \text{i.e.}\ S
\]

%dsp-level:0
\[
= 0>1>3>2\ [\text{machine: exact rational comparison}];
\]

%dsp-level:0
\[
\qquad \text{the sign census and the like-sector identity}\ [\text{machine}]
\]

%dsp-level:0
\[
(5):\ 4\cdot\tfrac32\cdot4\cdot3=72\ [\text{machine}];\qquad\mathrm{FT.107}(4)\text{'s spread row re-cited}
\]

%dsp-level:0
\[
(6):\ \mathrm{P5}\ \text{scan of this tag's text}\ [\text{machine: zero external numeric tokens}].\ \blacksquare
\]

\textbf{Phys.} Adding the edges yields six; adding the masses yields nothing. The informational count is exactly additive at the crystallized cell, so the binding read built from it vanishes identically — which is why a binding energy cannot be an integer summation. The flicker ensemble gives a different answer: averaged over the sector censuses, three separated constituents and one junction do not add up, and the mismatch is a nonzero exact rational in every combination of sectors. Turn the threshold to its sharp value and the ensemble collapses, the mismatch collapses with it, and the binding is exactly zero again — so the binding is a property of the flicker itself. Its magnitude stays order-one at budget four against a bound growing quadratically, its sign follows the transit sector's weight, and it remains a read of the corpus's own $m^2_B$ values with no external number anywhere near it.

\end{theorem}

\begin{theorem}[{\texorpdfstring{BM.22 — The static force matrix on the two charge classes}{BM.22 - The static force matrix on the two charge classes}}]\label{tag:BM-22}
\textbf{Thm.} With $h(\mathbf c)=\tfrac12(c_0I_2+\mathbf c\cdot\boldsymbol\sigma)$ and $\hat\kappa$ the ray swap [FT.129(2)], under Def-H.43's declared read promotion:

%dsp-level:0
\[
(1)\quad\textbf{the promoted force matrix:}
\]

%dsp-level:0
\[
\qquad \mathbf F_{\mathrm{ext}}(\mathbf c,q)=h(\mathbf c)\otimes I_2+q\,I_2\otimes\hat\kappa,
\]

%dsp-level:0
\[
\qquad \bigl[h\otimes I_2,\ q\,I_2\otimes\hat\kappa\bigr]=0
\]

%dsp-level:0
\[
\qquad [\mathrm{FT.129}(5);\ \text{machine: all profiles}\times\text{ternary}\ q];
\]

%dsp-level:0
\[
(2)\quad\textbf{the block read:}
\]

%dsp-level:0
\[
\qquad F_{LL}=F_{RR}=h(\mathbf c),
\]

%dsp-level:0
\[
\qquad F_{LR}=F_{RL}=q\,I_2
\]

%dsp-level:0
\[
\qquad [\text{machine: exact block extraction on}\ 4\ \text{profiles}\times3\ \text{charge values}];
\]

%dsp-level:0
\[
\text{the flipping block is the core read itself}:\ \langle L|\mathbf F_{\mathrm{ext}}|R\rangle=q\,I_2\ [\mathrm{FT.129}(5)];
\]

%dsp-level:0
\[
\qquad \text{the flipping set is}\ \{F_{\mathrm{EM}}\ \text{core read}\},\ \text{minimal and unique}\ [\mathrm{FT.129}(5)];
\]

%dsp-level:0
\[
(3)\quad\textbf{the spectrum:}
\]

%dsp-level:0
\[
\qquad \operatorname{spec}=\bigl\{\tfrac12(c_0+s|\mathbf c|)+tq:\ s,t\in\{\pm1\}\bigr\},
\]

%dsp-level:0
\[
\qquad \operatorname{charpoly}=\bigl((y-u)^2-q^2\bigr)\bigl((y+u)^2-q^2\bigr),\quad y=x-\tfrac{c_0}2,\quad u^2=\tfrac{|\mathbf c|^2}4
\]

%dsp-level:0
\[
[\mathrm{FT.130}(5);\ \text{machine: the monic quartic built from the four roots equals the Faddeev–LeVerrier charpoly, and equals the closed form, on every witness}];
\]

%dsp-level:0
\[
(4)\quad\textbf{the two classes:}\qquad\textbf{the chiral-partner splitting is}\ 2|q|\in\{0,2\}\qquad[\mathrm{FT.130}(5);\ \text{machine}]:
\]

%dsp-level:0
\[
\textbf{neutron class}\ (q_{\mathrm{net}}=0\ [\mathrm{Def\text{-}BM.3}(2)]):\quad F_{LR}=0,\quad\text{splitting}\ 0,\quad\text{parity doubling}\ [\mathrm{FT.130}(6)];
\]

%dsp-level:0
\[
\textbf{proton class}\ (q_{\mathrm{net}}=+1\ [\mathrm{Def\text{-}BM.3}(2)]):
\]

%dsp-level:0
\[
\quad F_{LR}=I_2,\quad\text{splitting}\ 2,\quad\text{the doublet is split by exactly two read units};
\]

%dsp-level:0
\[
\textbf{the EM difference between the two classes is exactly the flipping block}\ [(2),(4)]:\ \text{one carries it, the other is annihilated by it};
\]

%dsp-level:0
\[
(5)\quad\textbf{the derived locus and the promotion status:}
\]

%dsp-level:0
\[
\qquad \text{at the derived locus}\ \mathbf F_{\mathrm{ext}}(\mathbf c)=h\otimes I_\omega\ \text{and}\ F_{LR}=0\ \text{exactly}\ [\mathrm{FT.128}(2):\ \text{the no-flip theorem}]:
\]

%dsp-level:0
\[
\text{the flipping block exists only under Def-H.43's declared read promotion}\ [\text{the Def-H.43 declaration}];
\]

%dsp-level:0
\[
\qquad \textbf{status: conditional on that declaration, not on any new one};
\]

%dsp-level:0
\[
(6)\quad\textbf{the invariances:}
\]

%dsp-level:0
\[
\qquad \mathbf F_{\mathrm{ext}}\mapsto\mathbf F_{\mathrm{ext}}+\lambda I_4
\]

%dsp-level:0
\[
\implies F_{LR},
\]

%dsp-level:0
\[
\ \text{Schmidt ranks, commutants and}\ 2|q|\ \text{invariant; the spectrum translates by}\ \lambda
\]

%dsp-level:0
\[
\qquad [\mathrm{FT.130}(7)];
\]

%dsp-level:0
\[
\text{operator-Schmidt rank}=2\iff q\ne0\ \wedge\ \mathbf c\ne0;\quad1\ \text{when}\ h\propto I_2;\quad0\iff\mathbf c=0=q\ [\mathrm{FT.130}(2)];
\]

%dsp-level:0
\[
\qquad \neg\text{an entangling term at any locus}\ [\mathrm{FT.130}(2)];
\]

%dsp-level:0
\[
(7)\quad\textbf{the charge-label status:}
\]

%dsp-level:0
\[
\qquad q_{\mathrm{net}}\ \text{is}\ \mathrm{Def\text{-}BM.3}(4)\text{'s declared class label, pending the additive read}\ [\mathrm{BM.6}(1)]:
\]

%dsp-level:0
\[
\text{clauses}\ (1)\text{–}(4),(6)\ \text{are exact in}\ q\in\{-1,0,+1\}\ [\mathrm{EM.9}];
\]

%dsp-level:0
\[
\qquad \text{the attachment of}\ q=0\ \text{to the neutron class and}\ q=+1\ \text{to the proton class is the declaration}.
\]

\textbf{Req.} CT0.12, CT1.1, CT1.2 (Part 1); EM.9 (Part 2); Def-H.34, Def-H.43, FT.122(6), FT.123, FT.128(2), FT.129(2), FT.129(5), FT.130(2), FT.130(5), FT.130(6), FT.130(7) (Part 3); Def-BM.3 (this part); the Def-H.43 declaration; exact $\mathbb Q(i)$ block algebra, Faddeev–LeVerrier characteristic polynomials over $\mathbb Q$, exact quartic root reconstruction [math; machine: \texttt{checks/\allowbreak{}verify\_\allowbreak{}H\_\allowbreak{}topo\_\allowbreak{}blind\_\allowbreak{}neutron.py} §D].

\textbf{Pf.}

%dsp-level:0
\[
(1)\text{–}(2):\ \text{block extraction of}\ h\otimes I_2+qI_2\otimes\hat\kappa\ \text{with the chiral factor inner}: F_{ab}
\]

%dsp-level:0
\[
= h\,
\]

%dsp-level:0
\[
\delta\text{-block}+q\hat\kappa_{ab}
\]

%dsp-level:0
\[
\implies F_{LL}
\]

%dsp-level:0
\[
= F_{RR}
\]

%dsp-level:0
\[
= h\ (\hat\kappa_{00}=\hat\kappa_{11}=0),
\]

%dsp-level:0
\[
\ F_{LR}
\]

%dsp-level:0
\[
= F_{RL}
\]

%dsp-level:0
\[
= qI_2\ (\hat\kappa_{01}=\hat\kappa_{10}=1)\ [\text{machine}];
\]

%dsp-level:0
\[
\qquad \text{commutation}\ [\text{machine}]
\]

%dsp-level:0
\[
(3):\ \text{the four roots}\ \tfrac12(c_0+s|\mathbf c|)+tq\ \text{build a monic quartic over}\ \mathbb Q\ \text{at witnesses with}\ |\mathbf c|^2\ \text{a perfect square rational}\ (u\in\{\tfrac12,0,1,\tfrac37\});
\]

%dsp-level:0
\[
\ \text{Faddeev–LeVerrier on}\ \mathbf F_{\mathrm{ext}}\ \text{agrees, and the closed form}\ (y^2+u^2-q^2)^2-4u^2y^2\ \text{agrees with both}\ [\text{machine}]
\]

%dsp-level:0
\[
(4):\ 2|q|\ \text{at}\ q\in\{-1,0,+1\}\ [\text{machine}];\qquad\text{the class attachment}\ [\mathrm{Def\text{-}BM.3}(2)\ \text{re-cited}]
\]

%dsp-level:0
\[
(5)\text{–}(6):\ \mathrm{FT.128}(2),\ \mathrm{FT.130}(2),(7)\ \text{rows re-cited}\ (\text{frozen});\qquad(7):\ \mathrm{Def\text{-}BM.3}(4)\ \text{re-cited}.\ \blacksquare
\]

\textbf{Phys.} Put the static force matrix on the two charge classes and they come apart in exactly one place. Both classes share the same internal block $h$; the promoted read adds a chiral-flipping block equal to the core charge times the identity, so the neutral class has no flipping block at all and keeps its parity doubling, while the unit class splits its chiral partners by precisely two read units. That splitting — zero against two — is the whole electromagnetic difference between the neutron row and the proton row in this corpus, and it is a matrix identity, not a coupling constant. The promotion that allows a flipping block is a declared read, and the tier says so; everything downstream of it (the spectrum, the charpoly, the Schmidt ranks, the scalar-shift invariance) is exact. What is still declared rather than derived is which class label sits on which junction: the additive charge read remains an open construction.

\end{theorem}

\begin{theorem}[{\texorpdfstring{BM.23 — The transition: the ladder rule, the forced carrier, and the resolvent rate}{BM.23 - The transition: the ladder rule, the forced carrier, and the resolvent rate}}]\label{tag:BM-23}
\textbf{Thm.}

%dsp-level:0
\[
(1)\quad\textbf{the weak ladder:}
\]

%dsp-level:0
\[
\qquad E_\pm=\tfrac12(\sigma_1\pm i\sigma_2)=\text{the two matrix units in the}\ \sigma_3\ \text{eigenbasis};
\]

%dsp-level:0
\[
\qquad [E_+,E_-]=\sigma_3,
\]

%dsp-level:0
\[
\qquad [\sigma_3,E_\pm]=\pm2E_\pm
\]

%dsp-level:0
\[
\qquad [\mathrm{FT.41}(\mathrm d);\ \text{machine: exact over}\ \mathbb Q(i)];
\]

%dsp-level:0
\[
\text{the charged-current shape}\ W^\pm
\]

%dsp-level:0
\[
= \tfrac1{\sqrt2}(W_1\mp iW_2)\ [\mathrm{D3.10}]\ \text{is the same ladder on the pair register}\ [\mathrm{C8}];
\]

%dsp-level:0
\[
(2)\quad\textbf{the selection rule:}
\]

%dsp-level:0
\[
\qquad q_{\mathrm{core}}=\tfrac12(\sigma_{s}-\sigma_{t})\ [\mathrm{D3.4}]\ \text{and the ladder moves}\ \sigma_3\ \text{by}\ \pm2\ [(1)]\implies\boxed{\ \Delta q=\pm1\ }
\]

%dsp-level:0
\[
\qquad [\text{machine}];
\]

%dsp-level:0
\[
(3)\quad\textbf{the alphabet overflow:}
\]

%dsp-level:0
\[
\qquad \text{shifting the ternary read}\ \{-1,0,+1\}\ [\mathrm{EM.9}]\ \text{by}\ +1\ \text{gives}\ \{0,+1,+2\}\ \text{and by}\ -1\ \text{gives}\ \{-2,-1,0\}:
\]

%dsp-level:0
\[
\qquad \text{both leave the alphabet}\ [\text{machine}];
\]

%dsp-level:0
\[
\therefore\ \text{a charge-changing transition cannot stay inside the composite: it must emit a carrier};
\]

%dsp-level:0
\[
(4)\quad\textbf{the compensation is unique:}
\]

%dsp-level:0
\[
\qquad Q_{\mathrm{EM}}
\]

%dsp-level:0
\[
= \sum_{c\in U}q_{\mathrm{EM}}(c)
\]

%dsp-level:0
\[
= 0\ \text{in every valid universe}\ [\mathrm{EM.11};\ \mathrm{D2.11},\ \mathrm{D1.3}]\implies\text{a composite read changing by}\ \Delta q\ \text{requires a compensating read of exactly}\ -\Delta q:
\]

%dsp-level:0
\[
\#\{q\in\{-1,0,+1\}:\ q=-\Delta q\}=1\qquad[\text{machine}]:\qquad\textbf{the emitted carrier's charge is forced, not chosen};
\]

%dsp-level:0
\[
\text{the compensated total is exactly}\ 0:\ \text{EM.11's balance survives the transition}\ [\text{machine}];
\]

%dsp-level:0
\[
(5)\quad\textbf{the carrier selection theorem:}
\]

%dsp-level:0
\[
\qquad \text{the carrier must satisfy}\ (\mathrm i)\ q
\]

%dsp-level:0
\[
= -1\ \text{for}\ \Delta q
\]

%dsp-level:0
\[
= +1\ [(4)];
\]

%dsp-level:0
\[
\ (\mathrm{ii})\ \text{no colour attachment}\ [\text{else the composite's colour seal breaks}:\ \mathcal C=0\implies\text{the block vanishes}\ [\mathrm{FT.19};\ \mathrm{CT3.18},\ \mathrm{D2.14}]];
\]

%dsp-level:0
\[
\ (\mathrm{iii})\ \text{EM activated}\ (f_v\ne0)\ [\mathrm{D3.3};\ \text{else it cannot carry the read}];
\]

%dsp-level:0
\[
\text{in the channel inventory: the}\ q\text{-channel is triplet-attached}\ [\mathrm{FT.119}]\ \text{(fails ii)};
\]

%dsp-level:0
\[
\qquad \text{the}\ \nu\text{-channel is EM-unactivated}\ [\mathrm{FT.142}(1):\ f_v=0\ \text{reads}]\ \text{(fails iii)};
\]

%dsp-level:0
\[
\therefore\ \textbf{the}\ \ell\text{-channel is the unique carrier class}\ [\mathrm{FT.109}\text{–}\mathrm{FT.113}:\ \text{colour-blind, ternary read active}]:
\]

%dsp-level:0
\[
\qquad \text{the emitted line is Def-BM.6(1)'s electron row};
\]

%dsp-level:0
\[
(6)\quad\textbf{the neutral partner is admissible, not forced:}
\]

%dsp-level:0
\[
\qquad \neg\exists\ \text{a derived law requiring a neutral emission}\ [\text{no lepton-number read exists in the inventory}];
\]

%dsp-level:0
\[
\text{the}\ \nu\text{-channel}\ [\mathrm{FT.142}\text{–}\mathrm{FT.144}]\ \text{is EM-neutral and colour-blind, so it can carry the deficit read}\ \Delta S
\]

%dsp-level:0
\[
= \tfrac32n_{\mathcal V}\ [\mathrm{FT.107}(2)]\ \text{without disturbing the charge balance}:
\]

%dsp-level:0
\[
\textbf{status: admissible carrier for the deficit read; its emission is not a theorem of this corpus};
\]

%dsp-level:0
\[
(7)\quad\textbf{the rate:}
\]

%dsp-level:0
\[
\qquad \sum_n\langle\varnothing|\hat U^n|\varnothing\rangle w^n=\frac1{1-w^2x_+x_-}\ \ \text{in}\ \mathbb Q[w,x_+,x_-]/(w^{12})
\]

%dsp-level:0
\[
\qquad [\mathrm{FT.108}(5);\ \text{machine}];
\]

%dsp-level:0
\[
x_+x_-=4\varepsilon_q(1-4\varepsilon_q)\le\tfrac14,
\]

%dsp-level:0
\[
\qquad \text{equality at}\ \varepsilon_q=\tfrac18
\]

%dsp-level:0
\[
\qquad [\text{machine: exact maximum over the admissible range}]\implies\text{the pole}\ w^2x_+x_-=1\ \text{sits at}\ |w|\ge2;
\]

%dsp-level:0
\[
\text{the per-tick transition weight is an exact rational function of}\ \varepsilon_q\ \text{alone}\ [\mathrm{BM.19}(2)]:
\]

%dsp-level:0
\[
\qquad \Gamma\ \text{is}\ \neg\text{a declared number and}\ \neg\text{a new parameter};
\]

%dsp-level:0
\[
(8)\quad\textbf{the no-flicker theorem:}
\]

%dsp-level:0
\[
\qquad \varepsilon_q=0\implies x_+x_-=0\implies\text{the resolvent is}\ 1\implies\Gamma=0\ \text{exactly}
\]

%dsp-level:0
\[
\qquad [\mathrm{BM.19}(3);\ \text{machine}]:
\]

%dsp-level:0
\[
\textbf{a non-flickering composite never decays};
\]

%dsp-level:0
\[
\qquad \text{the flicker is the mechanism, and the threshold is its only control};
\]

%dsp-level:0
\[
(9)\quad\textbf{the arrow:}
\]

%dsp-level:0
\[
\qquad \text{the transition is forward-only}\ [\mathrm{FT.144}(2):\ \text{the domains force strict alternation;}\ \mathrm{A6}:\ \neg\exists\tau^{-1}]:
\]

%dsp-level:0
\[
\qquad \neg\exists\ \text{an inverse decay in the ledger};
\]

%dsp-level:0
\[
\text{the time coordinate is the ledger order}\ [x^0:=\tau_n,\ \mathrm{D1.8}]:
\]

%dsp-level:0
\[
\qquad \Gamma\ \text{is a weight per tick},\ \neg\text{a rate per external second}\ [\mathrm{P5}].
\]

\textbf{Rmk (provenance of each piece).} Clause (5) is a selection theorem: three conditions on the carrier, two channels eliminated by their own frozen rows, one class left. Clause (6) records honestly what is \textit{not} forced — the neutral partner — because no lepton-number or energy-sharing law exists in the inventory; naming it would be a declaration, and none is issued. Clause (7)'s $\Gamma$ is exact in $\varepsilon_q$ and carries no time unit beyond the tick [D1.8]; no lifetime is derived and none is compared [P5].

\textbf{Rmk (the forcing status of clause (6)).} The neutral emission is conditionally forced: BM.41(2)–(3) derives that every kinematically admissible transition vertex carries a third member, and that its class is unique among the constructed carrier classes — the $\nu$-channel line: a two-body vertex at the equal rest masses of BM.28(1) has no real on-shell solution for any positive member mass ($E_p-m_p=-m_\ell^2/(2m_n)<0$ identically), and the transit (EM-radiation) alternative is $\neg$constructed [Def-BM.8(8)]. Clause (6)'s inventory row is unchanged: $\neg\exists$ a lepton-number read — the forcing is kinematic plus uniqueness, $\neg$a number conservation; its hypothesis is the vertex's kinematic admissibility, which BM.41(4) shows no adopted read satisfies.

\textbf{Req.} A6 (Part 1); D1.2, D1.8, D3.3, D3.4, D3.10 (Parts 1–2); C8, CT0.8, CT3.18 (Part 1, Part 2); D2.11, D2.14 (Part 2); EM.9, EM.11 (Part 2); P5 (Parts 1–2); Def-H.32, FT.19, FT.41(d), FT.107(2), FT.108(3), FT.108(5), FT.109, FT.113, FT.119, FT.142(1), FT.144(2) (Part 3); Def-BM.5, Def-BM.6, BM.19(2), BM.19(3), BM.22 (this part); exact $\mathbb Q(i)$ ladder algebra, exhaustive ternary-alphabet censuses, exact rational resolvent identities and maxima [math; machine: \texttt{checks/\allowbreak{}verify\_\allowbreak{}H\_\allowbreak{}topo\_\allowbreak{}blind\_\allowbreak{}neutron.py} §D–§E].

\textbf{Pf.}

%dsp-level:0
\[
(1):\ E_\pm=\tfrac12(\sigma_1\pm i\sigma_2)\ \text{entrywise}\ [\text{machine}];
\]

%dsp-level:0
\[
\qquad [E_+,E_-]=\sigma_3,\ [\sigma_3,E_\pm]=\pm2E_\pm\ [\text{machine: all commutators exact over}\ \mathbb Q(i)]
\]

%dsp-level:0
\[
(2):\ \text{a}\ \pm2\ \text{move on}\ \sigma_s\ \text{changes}\ \tfrac12(\sigma_s-\sigma_t)\ \text{by}\ \pm1\ [\text{machine}];
\]

%dsp-level:0
\[
\qquad (3):\ \text{the shifted sets}\ \{0,1,2\},\ \{-2,-1,0\}\ \text{against the alphabet}\ \{-1,0,1\}\ [\text{machine}]
\]

%dsp-level:0
\[
(4):\ \mathrm{EM.11}\ \text{re-cited};\qquad\#\{q:q=-\Delta q\}=1\ \text{for}\ \Delta q=\pm1\ [\text{machine}];\qquad\Delta q+(-\Delta q)=0\ [\text{machine}]
\]

%dsp-level:0
\[
(5):\ (\mathrm{ii})\ \mathrm{FT.19}/\mathrm{CT3.18}/\mathrm{D2.14}\ \text{re-cited};\ (\mathrm{iii})\ \mathrm{D3.3}\ \text{re-cited};
\]

%dsp-level:0
\[
\qquad \text{eliminations}: \mathrm{FT.119}\ (\text{triplet-attached}),\ \mathrm{FT.142}(1)\ (\text{EM unactivated})\ \text{re-cited};
\]

%dsp-level:0
\[
\qquad \text{uniqueness}: \text{the channel inventory has three triplets}\ [\mathrm{FT.112},\ \mathrm{FT.119},\ \mathrm{FT.142}]
\]

%dsp-level:0
\[
(6):\ \text{inventory scan: no lepton-number read exists}\ [\mathrm{FT.81}(4)\ \text{pattern}];
\]

%dsp-level:0
\[
\qquad \Delta S=\tfrac32n_{\mathcal V}\ [\mathrm{FT.107}(2)\ \text{re-cited}]
\]

%dsp-level:0
\[
(7):\ \mathrm{FT.108}(5)\text{'s resolvent re-derived coefficient-wise at four rational thresholds}\ [\text{machine}];
\]

%dsp-level:0
\[
\qquad x_+x_-=4\varepsilon_q(1-4\varepsilon_q):\ \text{maximum}\ \tfrac14\ \text{at}\ \varepsilon_q=\tfrac18\ [\text{machine over}\ \varepsilon_q=k/256]
\]

%dsp-level:0
\[
(8):\ \varepsilon_q=0\implies x_+x_-=0\implies1/(1-0)=1\implies\text{no transition weight beyond the identity}\ [\text{machine}];
\]

%dsp-level:0
\[
\qquad (9):\ \mathrm{FT.144}(2),\ \mathrm{A6},\ \mathrm{D1.8}\ \text{re-cited}.\ \blacksquare
\]

\textbf{Phys.} The decay is assembled from three independent necessities and one honest gap. The weak ladder moves the $\sigma_3$ eigenvalue by two units, and the core charge is half a $\sigma$ difference, so any charge-changing transition moves the ternary read by exactly one — which pushes the endpoints straight out of the alphabet, so the change cannot stay inside the composite. Total charge vanishes in every valid universe, and the alphabet contains exactly one value that compensates a unit change, so the emitted carrier's charge is forced. Then the inventory does the rest: the quark-side channel is colour-attached and would break the seal, the neutrino-side channel is electromagnetically inert and could not carry the read, which leaves precisely one class — the $\ell$-channel line the naming layer calls the electron. The neutral partner is admissible and useful for carrying the deficit read, but no law in the corpus forces it, and the tier says so instead of naming a neutrino into the theorem. The rate is a geometric resolvent in the derived weights, its pole safely outside the unit disk, and it vanishes identically at the sharp crystallization threshold: a composite that cannot flicker cannot decay. The whole transition runs forward only, per tick, with no lifetime and no external unit anywhere in it.

\end{theorem}

\medskip\noindent\textit{\texorpdfstring{BM.24 — The status of the topology-blind composite: adjudication and scope}{BM.24 - The status of the topology-blind composite: adjudication and scope}.}\ Structural-constraints record: see Appendix B, Section B.1.17 of the Addendum.\par\medskip

\phantomsection
\section[Tier BM.7 — The Proton in the Positive CPT Sector and Its…]{\texorpdfstring{Tier BM.7 — The Proton in the Positive CPT Sector and Its Mirror: the Sector Labels, the Exact Phase, and the Charge-Blind Geometry}{Tier BM.7 - The Proton in the Positive CPT Sector and Its Mirror: the Sector Labels, the Exact Phase, and the Charge-Blind Geometry}}

\begin{definition}[{\texorpdfstring{Def-BM.7 — The EM$^\pm$ sector labels and the EM scar}{Def-BM.7 - The EM sector labels and the EM scar}}]\label{tag:Def-BM-7}
\textbf{Def.} Part-4 labels for two derived structures: the CPT-orientation classes of a charged sealed composite, and its endpoint charge read. Both labels are naming; both objects are derived.

%dsp-level:0
\[
(1)\quad\textbf{the sectors:}\qquad\mathrm{EM}^+\ :=\ \text{the CPT-orientation class with}\ q_{\mathrm{net}}=+1\ \text{and the unstarred weight}\ w(C);
\]

%dsp-level:0
\[
\mathrm{EM}^-\ :=\ \text{the class with}\ q_{\mathrm{net}}=-1\ \text{and the conjugated weight}\ w(C)^\ast;
\]

%dsp-level:0
\[
\textbf{constitutive clause:}\ \ \text{the two labels name the two orbits of the CPT involution on the charged classes}\ [\mathrm{EM.24}(4);\ \mathrm{BM.25}(4)];
\]

%dsp-level:0
\[
\qquad \textbf{non-derived clause:}\ \ \text{the signs}\ \pm\ \text{in the labels are a convention};
\]

%dsp-level:0
\[
(2)\quad\textbf{the scope clause (the mirror is formal, not dynamical):}
\]

%dsp-level:0
\[
\qquad \text{``backward time''}\ :=\ \Theta_{\mathfrak L}\ [\mathrm{D2.7}:\ \text{ledger order reversed},\ J^+\leftrightarrow J^-,\ G(e)\mapsto G(e)^\ast],
\]

%dsp-level:0
\[
\text{an involution on configuration data bijecting the completed admissible set onto itself}\ [\mathrm{FT.144}(4)(\mathrm i);\ \mathrm{FT.156}(5);\ \text{machine}]\ —\ \neg\text{a dynamical inverse}:
\]

%dsp-level:0
\[
\neg\exists\,\tau^{-1}\ [\mathrm{A6}];
\]

%dsp-level:0
\[
\qquad \text{the pure order reversal (labels unchanged) of every nonempty completed history is inadmissible at its first event}\ [\mathrm{FT.144}(4)(\mathrm{ii});\ \text{machine}];
\]

%dsp-level:0
\[
\therefore\ \mathrm{EM}^-\ \text{is a bookkeeping orbit of the mirror,}\ \neg\text{a universe running backwards};
\]

%dsp-level:0
\[
\qquad \text{no clause of this part asserts a backward evolution};
\]

%dsp-level:0
\[
(3)\quad\textbf{the EM scar:}
\]

%dsp-level:0
\[
\qquad \text{scar}(c)\ :=\ q_{\mathrm{core}}(c)=\tfrac12\bigl(\sigma_{s(c)}-\sigma_{t(c)}\bigr)\in\{-1,0,+1\}
\]

%dsp-level:0
\[
\qquad [\mathrm{D3.4},\ \mathrm{EM.9},\ \mathrm{EM.9a};\ \mathrm{FT.40}]:
\]

%dsp-level:0
\[
\text{the endpoint charge read of a materialized chain}\ —\ \text{registration data on the two ends}\ [\mathrm{D3.3},\ \mathrm{D1.1}],
\]

%dsp-level:0
\[
\ \neg\text{an internal edge}\ [\mathrm{D0.8}];
\]

%dsp-level:0
\[
\textbf{the zero-contribution clause:}\quad\text{the scar enters}\ E_{\mathrm{int}}\ \text{with coefficient}\ 0\ [\mathrm{BM.28}(3);\ \text{machine}:\ \text{Census FT.35}(2)\text{'s decomposition is}\ 2{+}2{+}2,\ \text{all edge counts}];
\]

%dsp-level:0
\[
(4)\quad\textbf{the neutral class:}
\]

%dsp-level:0
\[
\qquad q_{\mathrm{net}}=0\ \text{is the unique fixed point of the charge map}\ q\mapsto-q\ [\mathrm{EM.24}(1);\ \text{machine}]:
\]

%dsp-level:0
\[
\qquad \text{the neutral sector is charge-self-conjugate};
\]

%dsp-level:0
\[
\text{no configuration is fixed pointwise, since}\ \Theta_{\mathfrak L}\ \text{conjugates the weight}\ [\mathrm{CT2.8}(1)]:
\]

%dsp-level:0
\[
\qquad \text{self-conjugacy is a class property,}\ \neg\text{a pointwise one}\ [\mathrm{BM.25}(4)];
\]

%dsp-level:0
\[
(5)\quad\textbf{the data bar:}
\]

%dsp-level:0
\[
\qquad \text{``matter''/``antimatter'' are the labels of}\ (1)\ \text{and carry no external datum}\ [\mathrm{P5}];
\]

%dsp-level:0
\[
\qquad \text{no annihilation channel, no production channel and no cross-section is constructed in this corpus}\ [\mathrm{BM.29}(5)].
\]

\textbf{Rmk (vocabulary).} $\mathrm{EM}^\pm$ are Part-4 labels for CPT-orientation classes; they are not the register's $(S,T)$ sector symbols $\mathcal E_{\pm\pm}$ [CT0.12, CT1.1] and no clause identifies them. The pair-register sectors grade operators; the labels here grade composite configurations.

\textbf{Req.} A6, A8 (Part 1); D0.8, D1.1, D2.7, D3.3, D3.4 (Parts 1–2); CT2.8(1) (Part 2); EM.9, EM.9a, EM.24(1), EM.24(4) (Part 2); P5 (Parts 1–2); FT.40, FT.144(4), FT.156(5) (Part 3); Def-BM.3 (this part); clauses (3)–(4) carry marked forward pointers at their citing sites, imported in no \textbf{Req} line; Census FT.35(2) (row-wise); exhaustive enumeration of ledger records, exact $\mathbb Q(i)$ weight algebra [math; machine: \texttt{checks/\allowbreak{}verify\_\allowbreak{}H\_\allowbreak{}proton\_\allowbreak{}cpt.py} §A, §C].

\textbf{Pf.}

%dsp-level:0
\[
\text{Definition block: the constitutive clause is}\ (1)\text{'s orbit naming};
\]

%dsp-level:0
\[
\ \text{the non-derived clauses are the sign convention of}\ (1),
\]

%dsp-level:0
\[
\ \text{the scope reading of}\ (2)\ \text{and the vocabulary of}\ (5);
\]

%dsp-level:0
\[
\ (3)\text{'s zero-contribution clause and}\ (4)\text{'s fixed-point census are derived and machine-checked}.\ \blacksquare
\]

\textbf{Phys.} Two names are added and nothing else. The positive and negative CPT-orientation sectors label the two orbits that the corpus's own CPT involution makes of the charged classes, and the EM scar is the endpoint charge read the fermion era already derived. The one clause that does real work is the scope clause: the mirror is an involution on configuration data — reverse the ledger order and swap the seal labels — and it is not a way of running the universe backwards, because the tick has no inverse and every reversed history is inadmissible at its first event. So the negative sector is a bookkeeping orbit with equal amplitude, not a backwards-running world. The neutral class is the unique charge fixed point of the mirror, self-conjugate as a class while no single configuration is fixed, since the mirror always conjugates the weight.

\end{definition}

\begin{theorem}[{\texorpdfstring{BM.25 — The CPT involution on the sealed composite: orbits, the exact phase, and the census ratio}{BM.25 - The CPT involution on the sealed composite: orbits, the exact phase, and the census ratio}}]\label{tag:BM-25}
\textbf{Thm.}

%dsp-level:0
\[
(1)\quad\textbf{the three involutions are derived:}
\]

%dsp-level:0
\[
\qquad C:\ \sigma_e\mapsto-\sigma_e\ \forall e\implies q\mapsto-q,\quad f_v\mapsto-f_v,\quad\aleph_v\ \text{invariant}
\]

%dsp-level:0
\[
\qquad [\mathrm{EM.24}(1);\ \mathrm{D3.3}];
\]

%dsp-level:0
\[
P:\ \mathcal P\mapsto E_{\mathcal R}\setminus\mathcal P,\ \ \text{equivalently}\ \alpha_e\mapsto\pi-\alpha_e\qquad[\mathrm{EM.24}(2);\ \mathrm{Def\text{-}EM.21},\ \mathrm{EM.20}(1)];
\]

%dsp-level:0
\[
T:\ \Theta_{\mathfrak L},\qquad w(\Theta_{\mathfrak L}C)=w(C)^\ast\qquad[\mathrm{EM.24}(3);\ \mathrm{D2.7};\ \mathrm{CT2.8}(1)];
\]

%dsp-level:0
\[
\mathrm{CPT}:=C\circ P\circ T,\qquad C^2=P^2=T^2=\mathrm{id}\implies\mathrm{CPT}^2=\mathrm{id}\qquad[\mathrm{EM.24}(4):\ \text{each component derived, no external input}];
\]

%dsp-level:0
\[
(2)\quad\textbf{the mirror on the ledger is an exact involution with an unexpected fixed set:}
\]

%dsp-level:0
\[
\qquad \Theta_{\mathfrak L}^2=\mathrm{id}\ \text{on every admissible single-pair history}\ [\text{machine}];
\]

%dsp-level:0
\[
\Theta_{\mathfrak L}\bigl((J^+J^-)^m\bigr)=(J^+J^-)^m\ \ \forall m:
\]

%dsp-level:0
\[
\qquad \text{reversal composed with the}\ J^+\leftrightarrow J^-\ \text{relabelling fixes every completed word pointwise}\ [\text{machine: all}\ m\le6];
\]

%dsp-level:0
\[
\text{the pure order reversal (labels unchanged) is inadmissible at its first event}\ [\mathrm{FT.144}(4)(\mathrm{ii});\ \text{machine}]:
\]

%dsp-level:0
\[
\qquad \textbf{the formal mirror exists and the dynamical reversal does not};
\]

%dsp-level:0
\[
\Theta_{\mathfrak L}\ \text{bijects the two-pair completed interleavings onto themselves at}\ 2n
\]

%dsp-level:0
\[
= 2,
\]

%dsp-level:0
\[
4,
\]

%dsp-level:0
\[
6\ [\text{machine: set equality and cardinality}];
\]

%dsp-level:0
\[
\qquad \#J^+,\ \#J^-\ \text{are preserved}\implies\text{the Janus balance survives the mirror}\ [\mathrm{D1.3};\ \text{machine}];
\]

%dsp-level:0
\[
(3)\quad\textbf{the weight conjugation and the exact phase:}
\]

%dsp-level:0
\[
\qquad G(e)=\sigma_e\exp\!\bigl(i\alpha_e\widehat{ds}^{\,2}(e)\bigr),\quad\alpha_e\in\{0,\pi\}\ [\mathrm{D3.5},\ \mathrm{CT0.9}]\implies G(e)=\sigma_e\,i^{m_e},\quad m_e\in\mathbb Z:
\]

%dsp-level:0
\[
\text{the phase alphabet is the fourth roots of unity}\ —\ \text{exact in}\ \mathbb Q(i),\ \neg\text{a transcendental};
\]

%dsp-level:0
\[
\qquad \Theta_{\mathfrak L}:\ m_e\mapsto-m_e,\ \sigma_e\ \text{fixed},\ \mathcal J(C)\mapsto\mathcal J(C)^\ast\ [\mathrm{D3.6}];
\]

%dsp-level:0
\[
\boxed{\ w(\Theta_{\mathfrak L}C)=w(C)^\ast\ }
\]

%dsp-level:0
\[
\qquad [\mathrm{CT2.8}(1);\ \text{machine:}\ 24696\ \text{enumerated configurations}\ (\text{sign assignments}\times\text{half-integer}\ \kappa\ \text{ladder}\times\mathcal J\text{-phases})];
\]

%dsp-level:0
\[
w(\mathrm{CPT}\,C)=(-1)^{|E(C)|}\,e^{i\pi\Theta(C)}\,w(C),
\]

%dsp-level:0
\[
\qquad \lambda:=(-1)^{|E|}e^{i\pi\Theta}\in\{+1,-1,+i,-i\},
\]

%dsp-level:0
\[
\qquad |\lambda|^2=1
\]

%dsp-level:0
\[
\qquad [\mathrm{EM.24}(5),\ \mathrm{FT.134}(3);\ \text{machine}];
\]

%dsp-level:0
\[
\bigl|w(\mathrm{CPT}\,C)\bigr|^2=\bigl|w(C)\bigr|^2\ \ \forall C\qquad[\mathrm{EM.24}(5):\ \text{unconditional};\ \text{machine}];
\]

%dsp-level:0
\[
\textbf{the on-shell refinement:}\quad\widehat{ds}^{\,2}(e)\in\mathbb Z\ \forall e\implies\Theta\in\mathbb Z\implies\lambda=(-1)^{|E|+\Theta}\in\{\pm1\}:
\]

%dsp-level:0
\[
\qquad w(\mathrm{CPT}\,C)=w(C)^\ast\iff|E|+\Theta\equiv0\ (\mathrm{mod}\ 2)
\]

%dsp-level:0
\[
\qquad [\mathrm{EM.24}(5);\ \text{machine}];
\]

%dsp-level:0
\[
(4)\quad\textbf{the sector orbit census:}
\]

%dsp-level:0
\[
\qquad \mathrm{CPT}:\ (q,\ \text{orientation})\mapsto(-q,\ \text{flipped})\ \ \text{on}\ \{-1,0,+1\}\times\{w,w^\ast\}:
\]

%dsp-level:0
\[
\qquad \text{an involution with}\ \textbf{no fixed sector}\ [\text{machine}];
\]

%dsp-level:0
\[
\text{the orbit decomposition is exactly three pairs of size two};
\]

%dsp-level:0
\[
\qquad q\mapsto-q\ \text{has exactly one fixed point},\ q=0\ [\text{machine}]:
\]

%dsp-level:0
\[
\textbf{the neutral class is charge-self-conjugate and every charged class is exchanged}\ —\ \mathrm{EM}^+\leftrightarrow\mathrm{EM}^-\ [\mathrm{Def\text{-}BM.7}(1),(4)];
\]

%dsp-level:0
\[
(5)\quad\textbf{the census ratio (the sectors are equal in weight, the neutral is twice either):}
\]

%dsp-level:0
\[
\qquad Z^{\mathrm{norm}}(0,S)
\]

%dsp-level:0
\[
= 2\,
\]

%dsp-level:0
\[
Z^{\mathrm{norm}}(+1,S)\ \ \text{exactly for}\ S
\]

%dsp-level:0
\[
= 0,
\]

%dsp-level:0
\[
1,
\]

%dsp-level:0
\[
2,
\]

%dsp-level:0
\[
3\ [\text{machine}:\ 395824=2\cdot197912,\ 1574716=2\cdot787358,\ 1311552=2\cdot655776,\ 2001024=2\cdot1000512];
\]

%dsp-level:0
\[
Z^{\mathrm{norm}}(+1,S)=Z^{\mathrm{norm}}(-1,S)\ \ \text{exactly}\ [\mathrm{FT.106}(3\mathrm d);\ \text{machine}]:
\]

%dsp-level:0
\[
\qquad \textbf{the two charged sectors have identical weight};
\]

%dsp-level:0
\[
\text{the ratio is the factorization's own:}\ m_q2^{|R|-2}\ \text{with}\ m_q
\]

%dsp-level:0
\[
= (2,1,1)\ \text{at}\ q
\]

%dsp-level:0
\[
= (0,+1,-1)
\]

%dsp-level:0
\[
\implies m_0/m_{\pm1}
\]

%dsp-level:0
\[
= 2\ [\mathrm{FT.106}(3\mathrm a);\ \text{machine}];
\]

%dsp-level:0
\[
Z^{\mathrm{amp}}(+1,S)=Z^{\mathrm{amp}}(-1,S)\ \text{on the cover support},
\]

%dsp-level:0
\[
\qquad Z^{\mathrm{amp}}(q,S)=0\ \ \forall q,\ S\in\{2,3\}
\]

%dsp-level:0
\[
\qquad [\mathrm{FT.106}(3\mathrm c),(3\mathrm d);\ \text{machine}];
\]

%dsp-level:0
\[
(6)\quad\textbf{the baryon-type character:}
\]

%dsp-level:0
\[
\qquad \Theta_{\mathrm{int}}\in4\mathbb Z\ [\text{Census FT.49}]\implies e^{i\pi\Theta}=1;
\]

%dsp-level:0
\[
\qquad \Theta_{\mathrm{int}}\in2\mathbb Z\ [\text{Census FT.35}(3);\ \mathrm{BM.2}(1)]\implies e^{i\pi\Theta}=1\ \ \text{as well}\ [\text{machine}]:
\]

%dsp-level:0
\[
\text{on the}\ \varepsilon\text{-sealed junction class the CPT character is exactly}
\]

%dsp-level:0
\[
\qquad \lambda=(-1)^{|E|}
\]

%dsp-level:0
\[
\qquad [\mathrm{FT.134}(3);\ \text{machine}];
\]

%dsp-level:0
\[
(7)\quad\textbf{the norm pairing:}
\]

%dsp-level:0
\[
\qquad \langle A|A\rangle
\]

%dsp-level:0
\[
= \sum_C|w(C)|^2
\]

%dsp-level:0
\[
\ge 0\ [\mathrm{CT2.8}(3)]\ \text{is CPT-invariant term by term}:\ |w(C)|^2
\]

%dsp-level:0
\[
= |w(\Theta_{\mathfrak L}C)|^2\ [\text{machine on all enumerated pairs}]:
\]

%dsp-level:0
\[
\textbf{the two sectors contribute equally to the norm}\ —\ \text{the amplitude-level statement of}\ \mathrm{EM}^+\leftrightarrow\mathrm{EM}^-.
\]

\textbf{Rmk (locus discipline).} Clause (2)'s fixed set is a statement about \textit{words}: the mirror fixes each completed alternating word because reversing and relabelling commute to the same word. It is not a statement that any dynamical process is reversible — A6 and FT.144(4)(ii) bar that reading, and clause (2) records both rows side by side.

\textbf{Req.} A6, A8 (Part 1); D1.3, D2.7, D3.3, D3.4, D3.5, D3.6 (Parts 1–2); CT0.9 (Part 1); CT2.8(1), CT2.8(3), EM.9, EM.9a, EM.20(1), EM.24(1), EM.24(2), EM.24(3), EM.24(4), EM.24(5), Def-EM.21 (Part 2); Def-BM.7 (this part); FT.40, FT.106(3), FT.134(3), FT.144(4), FT.156(5) (Part 3); Census FT.35(3), FT.49 (row-wise); BM.2(1), BM.7(1) (this part); exact $\mathbb Q(i)$ weight algebra on the fourth-root alphabet, exhaustive ledger-record enumeration, exact integer census arithmetic [math; machine: \texttt{checks/\allowbreak{}verify\_\allowbreak{}H\_\allowbreak{}proton\_\allowbreak{}cpt.py} §A].

\textbf{Pf.}

%dsp-level:0
\[
(1):\ \mathrm{EM.24}(1)\text{–}(4)\ \text{rows re-cited}\ (\text{frozen, machine-verified at Tier F.33});
\]

%dsp-level:0
\[
\qquad \mathrm{CPT}^2=C^2P^2T^2=\mathrm{id}\ [\text{each an involution}]
\]

%dsp-level:0
\[
(2):\ \Theta_{\mathfrak L}(w)
\]

%dsp-level:0
\[
= \text{reverse}(w)\ \text{with labels swapped}: \text{on}\ (J^+J^-)^m\ \text{the reversal gives}\ (J^-J^+)^m\ \text{and the swap returns}\ (J^+J^-)^m\ [\text{machine}];
\]

%dsp-level:0
\[
\qquad \Theta_{\mathfrak L}^2=\mathrm{id}\ [\text{machine}];
\]

%dsp-level:0
\[
\qquad \text{pure reversal}: \text{leads with}\ J^-,\ \text{outside D1.2's domain}\ [\mathrm{FT.144}(4)(\mathrm{ii});\ \text{machine}]
\]

%dsp-level:0
\[
(3):\ \alpha_e\in\{0,\pi\}\ [\mathrm{CT0.9}]\ \wedge\ \widehat{ds}^{\,2}\in\tfrac12\mathbb Z\ [\mathrm{FT.22}\text{–}\mathrm{FT.24}\ \text{ladder}]\implies\alpha_e\widehat{ds}^{\,2}\in\tfrac\pi2\mathbb Z\implies\exp(i\alpha_e\widehat{ds}^{\,2})
\]

%dsp-level:0
\[
= i^{m_e}\ [\text{machine}];
\]

%dsp-level:0
\[
\qquad \Theta_{\mathfrak L}:\ G\mapsto G^\ast\ \text{per edge and}\ \mathcal J\mapsto\mathcal J^\ast\ [\mathrm{D2.7},\ \mathrm{D3.6}]\implies w\mapsto w^\ast\ [\text{machine}]
\]

%dsp-level:0
\[
\lambda=(-1)^{|E|}e^{i\pi\Theta}=i^{2|E|}i^{2\Theta}\in\{i^k\}\ [\text{machine}];
\]

%dsp-level:0
\[
\qquad |\lambda|^2=\lambda\lambda^\ast=1\ [\text{machine}];
\]

%dsp-level:0
\[
\qquad \text{on-shell}\ \Theta\in\mathbb Z\implies i^{2\Theta}=(-1)^\Theta\ [\text{machine}]
\]

%dsp-level:0
\[
(4):\ \text{the map on six sectors, its square, and its fixed-point census}\ [\text{machine}];
\]

%dsp-level:0
\[
\qquad q\mapsto-q\ \text{fixed points}: \{0\}\ [\text{machine}]
\]

%dsp-level:0
\[
(5):\ \mathrm{FT.106}(3)\text{'s census integers re-cited and the ratios verified exactly}\ [\text{machine}];
\]

%dsp-level:0
\[
\qquad m_q=(2,1,1)\ [\mathrm{FT.106}(3\mathrm a)]\implies m_0/m_{\pm1}=2\ \text{and}\ m_{+1}=m_{-1}\ [\text{machine}]
\]

%dsp-level:0
\[
(6):\ e^{i\pi\Theta}=i^{2\Theta}=1\ \text{at}\ \Theta\in2\mathbb Z\ [\text{machine}];\qquad\lambda=(-1)^{|E|}\cdot1\ [\mathrm{FT.134}(3)\ \text{re-cited}]
\]

%dsp-level:0
\[
(7):\ \mathrm{CT2.8}(3)\ \text{re-cited};\qquad|w|^2\ \text{equality term by term from}\ (3)\ [\text{machine}].\ \blacksquare
\]

\textbf{Phys.} The mirror does something sharper than exchange two sectors. On the ledger it is an involution that reverses order and swaps the seal labels — and on a completed history those two acts cancel, so every completed word is fixed pointwise, while the same history read backwards with its labels intact is inadmissible at its first event. The formal mirror exists; the dynamical reversal does not. Weights conjugate exactly, because the phase alphabet is the fourth roots of unity: with $\alpha_e\in\{0,\pi\}$ and a half-integer line element, every edge weight is a sign times a power of $i$, and negating the winding conjugates it. The CPT phase is therefore one of four exact values, its modulus is one, and the amplitude is invariant unconditionally — with the on-shell refinement that the phase is real and the weight literally conjugates exactly when the edge count plus the winding is even. The census closes the picture: the two charged sectors carry identical weight at every winding, the neutral sector exactly twice either, which is the factorization's own $m_q=(2,1,1)$ and not a coincidence. On the baryon-type junction, whose phase class lives in $4\mathbb Z$, the whole CPT character collapses to a sign counting edges.

\end{theorem}

\begin{theorem}[{\texorpdfstring{BM.26 — The proton row: the positive sector}{BM.26 - The proton row: the positive sector}}]\label{tag:BM-26}
\textbf{Thm.}

%dsp-level:0
\[
(1)\quad\textbf{the structure:}
\]

%dsp-level:0
\[
\qquad \text{the}\ \varepsilon\text{-sealed 3-juncture in the unit charge class}\ q_{\mathrm{net}}=+1
\]

%dsp-level:0
\[
\qquad [\mathrm{Def\text{-}BM.3}(2);\ \mathrm{D3.9}:\ B=\varepsilon_{ijk}c_1^ic_2^jc_3^k];
\]

%dsp-level:0
\[
\mathcal C(\Sigma;C)
\]

%dsp-level:0
\[
= 1\iff\text{all three attachments seal inside as one}\ \varepsilon\text{-type triple}\ [\text{Census FT.35}(5);\ \mathrm{CT3.17};\ \mathrm{FT.34}(5)]:
\]

%dsp-level:0
\[
\qquad \text{the colour seal is the observability switch}\ [\mathrm{BM.1}(1)];
\]

%dsp-level:0
\[
(2)\quad\textbf{the sector membership:}
\]

%dsp-level:0
\[
\qquad q_{\mathrm{net}}=+1\ \wedge\ w(C)\implies\ \text{the row lies in}\ \mathrm{EM}^+
\]

%dsp-level:0
\[
\qquad [\mathrm{Def\text{-}BM.7}(1);\ \text{the class label is}\ \mathrm{Def\text{-}BM.3}(4)\text{'s, pending BM.6}(1)];
\]

%dsp-level:0
\[
(3)\quad\textbf{the EM scar:}
\]

%dsp-level:0
\[
\qquad \text{scar}=q_{\mathrm{core}}=\tfrac12(\sigma_s-\sigma_t)\in\{-1,0,+1\}\ [\mathrm{D3.4},\ \mathrm{EM.9}];
\]

%dsp-level:0
\[
\qquad \text{the unit class carries}\ q_{\mathrm{core}}=+1\ \text{at its read end};
\]

%dsp-level:0
\[
\text{the scar is}\ \omega\text{-odd}\ [\mathrm{D2.6};\ \mathrm{FT.129}(3);\ \text{machine}]\ \text{and contributes}\ 0\ \text{to}\ E_{\mathrm{int}}\ [\mathrm{Def\text{-}BM.7}(3);\ \mathrm{BM.28}(3);\ \text{machine}];
\]

%dsp-level:0
\[
(4)\quad\textbf{the information content and the six reads:}
\]

%dsp-level:0
\[
\qquad E_{\mathrm{int}}^{(p)}=6\ (c_1{:}\,2,\ c_2{:}\,2,\ \text{connector}{:}\,2)
\]

%dsp-level:0
\[
\qquad [\text{Census FT.35}(2);\ \text{row-wise}];
\]

%dsp-level:0
\[
\Delta x_{\mathrm B}^{(p)}=\frac{\mathfrak h}{6\ln2},
\]

%dsp-level:0
\[
\qquad V_{\mathrm{eff}}^{(p)}=\frac{\mathfrak h^3}{216(\ln2)^3},
\]

%dsp-level:0
\[
\qquad m_{\mathrm{hid}}^{(p)}=\frac{6\ln2}{\mathfrak h},
\]

%dsp-level:0
\[
\qquad \rho^{(p)}=\frac{1296(\ln2)^4}{\mathfrak h^4};
\]

%dsp-level:0
\[
\kappa^2\rho^{(p)}=\frac{2592\,\pi(\ln2)^3}{\mathfrak h^2},
\]

%dsp-level:0
\[
\qquad r_s^{(p)}=3\,\mathfrak h,
\]

%dsp-level:0
\[
\qquad (\Delta x_{\mathrm B}^{(p)})^2\kappa^2\rho^{(p)}=72\,\pi\ln2
\]

%dsp-level:0
\[
\qquad [\mathrm{BM.13},\ \mathrm{BM.14},\ \mathrm{BM.15}(4),\ \mathrm{BM.16}(1),(4);\ \text{machine}];
\]

%dsp-level:0
\[
(5)\quad\textbf{the EM coupling:}
\]

%dsp-level:0
\[
\qquad F_{LR}=q\,I_2=I_2,
\]

%dsp-level:0
\[
\qquad \text{the chiral-partner splitting}\ 2|q|=2
\]

%dsp-level:0
\[
\qquad [\mathrm{BM.22}(2),(4);\ \mathrm{FT.129}(5),\ \mathrm{FT.130}(5);\ \text{machine}];
\]

%dsp-level:0
\[
\text{the promoted force matrix is}\ \mathbf F_{\mathrm{ext}}
\]

%dsp-level:0
\[
= h\otimes I_2+I_2\otimes\hat\kappa\ \text{with}\ \operatorname{spec}
\]

%dsp-level:0
\[
= \{\tfrac12(c_0+s|\mathbf c|)+t:\ s,t\in\{\pm1\}\}\ [\mathrm{BM.22}(3);\ \text{machine}];
\]

%dsp-level:0
\[
(6)\quad\textbf{the amplitude:}
\]

%dsp-level:0
\[
\qquad |w(\mathrm{CPT}\cdot C^{(p)})|^2=|w(C^{(p)})|^2\ [\mathrm{BM.25}(3)];
\]

%dsp-level:0
\[
\qquad \text{on-shell with}\ |E|+\Theta\equiv0\ (2):\ w(\mathrm{CPT}\cdot C^{(p)})=w(C^{(p)})^\ast;
\]

%dsp-level:0
\[
\text{the CPT image of this row is BM.27's}\ —\ \text{no independent datum is carried by either}.
\]

\textbf{Rmk (route scope of clause (4)).} The six reads of clause (4) are functions of $E_{\mathrm{int}}$ alone, so they are the crystallized-count route's rows: they are identical across the charge classes and clause (4) states exactly that. The ensemble-weighted mass read is route-relative — degenerate on the normalized route and split by the exact factor $\tfrac43$ between the neutral and the unit class on the total-deficit route [BM.28(4)]. Clause (4) is unchanged on its own route; the ensemble route is BM.28(4)'s and BM.30's content.

\textbf{Rmk (identification discipline).} The name attaches by Def-BM.3(2); the theorem content of every clause is the tuple $(\varepsilon\text{-junction},\ q_{\mathrm{net}}=+1,\ E_{\mathrm{int}}=6,\ \Theta_{\mathrm{int}}\in4\mathbb Z,\ 2|q|=2)$. No measured value enters [P5]; the charge class is the declared label pending the additive read [Def-BM.3(4); BM.6(1)].

\textbf{Req.} D3.4, D3.9 (Parts 1–2); CT3.17 (Part 2); EM.9 (Part 2); P5 (Parts 1–2); Def-BM.3, Def-BM.7, BM.1(1), BM.2(1), BM.13, BM.14, BM.15(4), BM.16(1), BM.16(4), BM.22(2), BM.22(3), BM.22(4) (this part); FT.34(5), FT.40, FT.129(3), FT.129(5), FT.130(5) (Part 3); Census FT.35(2), FT.35(3), FT.35(5), FT.49 (row-wise); BM.25(3) (this part); exact $\mathbb Q(i)$ block algebra, exact rational reads [math; machine: \texttt{checks/\allowbreak{}verify\_\allowbreak{}H\_\allowbreak{}proton\_\allowbreak{}cpt.py} §B–§C].

\textbf{Pf.}

%dsp-level:0
\[
(1)\text{–}(2):\ \mathrm{Def\text{-}BM.3}(2)\ \text{and}\ \mathrm{Def\text{-}BM.7}(1)\ \text{re-cited};\qquad\text{Census FT.35}(5)\text{'s sealing row re-cited}
\]

%dsp-level:0
\[
(3):\ \mathrm{D3.4}/\mathrm{EM.9}\ \text{re-cited};
\]

%dsp-level:0
\[
\qquad \omega\text{-oddness}: q_{\mathrm{core}}(\omega c)=\tfrac12(\sigma_t-\sigma_s)=-q_{\mathrm{core}}(c)\ [\mathrm{D2.6};\ \text{machine}];
\]

%dsp-level:0
\[
\qquad \text{the zero contribution}: \mathrm{BM.28}(3)
\]

%dsp-level:0
\[
(4):\ \text{substitution of}\ E_{\mathrm{int}}
\]

%dsp-level:0
\[
= 6\ \text{into BM.13(1), BM.14(1),(2),(3), BM.15(4), BM.16(1),(4)}\ [\text{machine: every coefficient exact}]
\]

%dsp-level:0
\[
(5):\ \mathrm{BM.22}(2)\text{–}(4)\ \text{re-derived at}\ q=+1\ [\text{machine}];\qquad(6):\ \mathrm{BM.25}(3)\ \text{re-cited}.\ \blacksquare
\]

\textbf{Phys.} The proton row is the neutron row with one entry changed. Same junction grammar, same six internal directed edges, same sealing phase class, same width, volume, mass, density, curvature scale, comparison radius and lock — because none of those reads contains a charge symbol. What changes is the flipping block: the promoted static force matrix carries $q I_2$ off-diagonal, so the unit class splits its chiral partners by exactly two read units where the neutral class splits them by zero. The scar is the endpoint read, it is odd under the orientation reversal, and it contributes nothing to the edge count — which is precisely why the geometry is blind to it. The amplitude of the row equals the amplitude of its mirror, unconditionally, and on-shell the mirror's weight is literally the complex conjugate.

\end{theorem}

\begin{theorem}[{\texorpdfstring{BM.27 — The antiproton row: the CPT image, and the resolution of the identity question}{BM.27 - The antiproton row: the CPT image, and the resolution of the identity question}}]\label{tag:BM-27}
\textbf{Thm.}

%dsp-level:0
\[
(1)\quad\textbf{the image:}
\]

%dsp-level:0
\[
\qquad C^{(\bar p)}:=\mathrm{CPT}\cdot C^{(p)}:
\]

%dsp-level:0
\[
\qquad q_{\mathrm{net}}\mapsto-q_{\mathrm{net}}=-1\ [\mathrm{EM.24}(1)],\quad w\mapsto w^\ast\ [\mathrm{CT2.8}(1)],\quad\mathcal P\mapsto E_{\mathcal R}\setminus\mathcal P\ [\mathrm{EM.24}(2)];
\]

%dsp-level:0
\[
\text{the image lies in}\ \mathrm{EM}^-\qquad[\mathrm{Def\text{-}BM.7}(1);\ \mathrm{BM.25}(4):\ \text{the two charged classes form one orbit}];
\]

%dsp-level:0
\[
(2)\quad\textbf{the structure is the same junction:}
\]

%dsp-level:0
\[
\qquad \varepsilon\text{-sealed 3-juncture in the}\ q_{\mathrm{net}}=-1\ \text{class}\ [\mathrm{Def\text{-}BM.3}(2):\ \text{the antiproton row}];
\]

%dsp-level:0
\[
E_{\mathrm{int}}^{(\bar p)}=E_{\mathrm{int}}^{(p)}=6:
\]

%dsp-level:0
\[
\qquad \Theta_{\mathfrak L}\ \text{preserves the record size and exchanges}\ \#J^+\ \text{with}\ \#J^-\ [\mathrm{BM.25}(2);\ \text{machine}]
\]

%dsp-level:0
\[
\implies |R|\ \text{and hence}\ E_{\mathrm{int}}
\]

%dsp-level:0
\[
= |R|\ \text{are invariant}\ [\mathrm{Def\text{-}BM.5}(4);\ \text{machine}];
\]

%dsp-level:0
\[
(3)\quad\textbf{every read is the image's, so every read is identical:}
\]

%dsp-level:0
\[
\qquad \Delta x_{\mathrm B}^{(\bar p)}=\Delta x_{\mathrm B}^{(p)},
\]

%dsp-level:0
\[
\quad V_{\mathrm{eff}}^{(\bar p)}=V_{\mathrm{eff}}^{(p)},
\]

%dsp-level:0
\[
\quad m_{\mathrm{hid}}^{(\bar p)}=m_{\mathrm{hid}}^{(p)},\quad\rho^{(\bar p)}=\rho^{(p)},\quad\kappa^2\rho^{(\bar p)}=\kappa^2\rho^{(p)},
\]

%dsp-level:0
\[
\quad r_s^{(\bar p)}=r_s^{(p)};
\]

%dsp-level:0
\[
\text{the reason is structural, not numerical: none of the six monomials carries a charge exponent}\ [\mathrm{BM.28}(1);\ \text{machine: the}\ q\text{-exponent is}\ 0\ \text{in all six}];
\]

%dsp-level:0
\[
(4)\quad\textbf{the amplitude and the character:}
\]

%dsp-level:0
\[
\qquad |w(C^{(\bar p)})|^2=|w(C^{(p)})|^2\ [\mathrm{EM.24}(5);\ \text{machine}];
\]

%dsp-level:0
\[
\qquad w(C^{(\bar p)})=\lambda\,w(C^{(p)}),\quad\lambda=(-1)^{|E|}e^{i\pi\Theta};
\]

%dsp-level:0
\[
\Theta_{\mathrm{int}}\in4\mathbb Z\ [\text{Census FT.49}]\implies\lambda=(-1)^{|E|}\in\{\pm1\}\qquad[\mathrm{BM.25}(6);\ \text{machine}];
\]

%dsp-level:0
\[
\text{the two rows have equal norm contribution to}\ \langle A|A\rangle
\]

%dsp-level:0
\[
= \sum_C|w(C)|^2\ [\mathrm{CT2.8}(3);\ \mathrm{BM.25}(7);\ \text{machine}];
\]

%dsp-level:0
\[
(5)\quad\textbf{the difference set is exactly two entries:}
\]

%dsp-level:0
\[
\qquad \text{scar}^{(\bar p)}=-\,\text{scar}^{(p)}\ [\omega\text{-odd};\ \text{machine}]\ \text{and the sector orientation}\ w\leftrightarrow w^\ast;
\]

%dsp-level:0
\[
F_{LR}^{(\bar p)}=-I_2=-F_{LR}^{(p)}\ \text{but}\ 2|q^{(\bar p)}|=2=2|q^{(p)}|:
\]

%dsp-level:0
\[
\qquad \textbf{the flipping block is odd and the splitting is even}\ [\mathrm{BM.22}(2),(4);\ \text{machine}];
\]

%dsp-level:0
\[
(6)\quad\textbf{the identity adjudication:}
\]

%dsp-level:0
\[
\qquad C^{(\bar p)}
\]

%dsp-level:0
\[
\ne C^{(p)}\ \text{as configurations}\ [\text{the scar and the orientation differ};\ \mathrm{BM.25}(4):\ \neg\exists\ \text{a fixed charged sector}];
\]

%dsp-level:0
\[
C^{(\bar p)}
\]

%dsp-level:0
\[
= \mathrm{CPT}\cdot C^{(p)}\ \text{with}\ \mathrm{CPT}^2
\]

%dsp-level:0
\[
= \mathrm{id}\ [\mathrm{EM.24}(4)]\implies\textbf{the antiproton row carries no independent datum}: \text{it is one orbit of size two, not two structures};
\]

%dsp-level:0
\[
\text{both statements hold because they are about different objects}\ —\ \text{configurations}\ (6)\text{'s first line, structure}\ (6)\text{'s second};
\]

%dsp-level:0
\[
\qquad \text{no clause asserts an identity of configurations};
\]

%dsp-level:0
\[
(7)\quad\textbf{what is not constructed:}
\]

%dsp-level:0
\[
\qquad \neg\exists\ \text{an annihilation channel},
\]

%dsp-level:0
\[
\ \neg\exists\ \text{a pair-production channel},
\]

%dsp-level:0
\[
\ \neg\exists\ \text{a cross-section between}\ \mathrm{EM}^+\ \text{and}\ \mathrm{EM}^-\ [\mathrm{Def\text{-}BM.7}(5);\ \mathrm{BM.29}(5)]:
\]

%dsp-level:0
\[
\text{the two sectors are amplitude-equal configuration orbits; their interaction is not a theorem of this corpus}.
\]

\textbf{Rmk (route scope of clause (3)).} The six identical reads of clause (3) carry no charge exponent [BM.28(1)], and the equality of the proton and antiproton rows holds on \textbf{every} mass route: the charge census is symmetric, $W(+1)=W(-1)=15$ exactly, and both ensemble layers are $\omega$-symmetric [BM.28(5); FT.106(3d)]. What the census does split is the neutral class against the charged pair, by the exact factor $\tfrac43$ on the total-deficit route [BM.28(4)(ii)] — a splitting that leaves this tag's CPT rows untouched.

\textbf{Rmk (mirror scope).} The image is configuration data under $\Theta_{\mathfrak L}$ composed with the sign and parity maps. It is not a time-reversed process: A6 gives the tick no inverse, and the pure order reversal of a completed history is inadmissible at its first event [BM.25(2)]. Reading clause (1) as a dynamical evolution would contradict FT.144(2), and no clause here does.

\textbf{Req.} A6 (Part 1); D2.6, D2.7 (Parts 1–2); CT2.8(1), CT2.8(3) (Part 2); EM.24(1), EM.24(2), EM.24(4), EM.24(5) (Part 2); P5 (Parts 1–2); Def-BM.3, Def-BM.5(4), Def-BM.7, BM.22(2), BM.22(4), BM.25(2), BM.25(4), BM.25(6), BM.25(7), BM.26 (this part); FT.129(3), FT.130(5), FT.144(2), FT.144(4) (Part 3); Census FT.49 (row-wise); clause (3) carries a marked forward pointer at its citing site, imported in no \textbf{Req} line; exact $\mathbb Q(i)$ arithmetic, exhaustive enumeration [math; machine: \texttt{checks/\allowbreak{}verify\_\allowbreak{}H\_\allowbreak{}proton\_\allowbreak{}cpt.py} §A–§C].

\textbf{Pf.}

%dsp-level:0
\[
(1):\ \mathrm{EM.24}(1),(2)\ \text{and}\ \mathrm{CT2.8}(1)\ \text{composed}\ [\mathrm{EM.24}(4)];\qquad\text{orbit membership}: \mathrm{BM.25}(4)\text{'s census}
\]

%dsp-level:0
\[
(2):\ \text{Census FT.35}(2)\text{'s row re-cited};
\]

%dsp-level:0
\[
\qquad |\Theta_{\mathfrak L}(w)|=|w|\ \text{and}\ \#J^+(\Theta w)=\#J^-(w)\ [\text{machine}]\implies E_{\mathrm{int}}=|R|\ \text{invariant}
\]

%dsp-level:0
\[
(3):\ \text{the six monomials of BM.13(1), BM.14(1),(2), BM.15(4), BM.16(4) have}\ q\text{-exponent}\ 0\ [\text{machine: monomial census}]\implies\text{substituting}\ q
\]

%dsp-level:0
\[
\mapsto -q\ \text{fixes each}
\]

%dsp-level:0
\[
(4):\ \mathrm{EM.24}(5)\text{'s identity re-cited and re-verified}\ [\text{machine}];
\]

%dsp-level:0
\[
\qquad \lambda=(-1)^{|E|}\ \text{at}\ \Theta\in4\mathbb Z\ [\mathrm{BM.25}(6)]
\]

%dsp-level:0
\[
(5):\ q_{\mathrm{core}}\mapsto-q_{\mathrm{core}}\ [\mathrm{D2.6};\ \text{machine}];\qquad F_{LR}=qI_2\mapsto-qI_2\ \text{while}\ 2|q|\mapsto2|q|\ [\text{machine}]
\]

%dsp-level:0
\[
(6):\ \neg\exists\ \text{a fixed charged sector}\ [\mathrm{BM.25}(4);\ \text{machine}]\implies\text{distinct configurations};
\]

%dsp-level:0
\[
\qquad \mathrm{CPT}^2=\mathrm{id}\ [\mathrm{EM.24}(4)]\implies\text{one orbit of size two}
\]

%dsp-level:0
\[
(7):\ \text{inventory scan: no annihilation, production or inter-sector transition operator exists}\ [\mathrm{BM.23}(6)\ \text{pattern}];
\]

%dsp-level:0
\[
\qquad \mathrm{Def\text{-}BM.7}(5)\ \text{re-cited}.\ \blacksquare
\]

\textbf{Phys.} The antiproton is what the mirror does to the proton, and the mirror is cheap: flip every edge sign, complement the alignment pattern, reverse the ledger while swapping the seal labels. Everything the corpus derives about the composite survives that map untouched, because not one of the six geometric reads contains a charge — the width, the volume, the mass, the density, the curvature scale and the comparison radius come out identical, and the amplitude comes out equal in modulus with a phase that is one of four exact roots of unity, collapsing to a bare sign on the baryon-type phase class. Only two entries change: the scar flips sign and the weight conjugates. The flipping block flips with the scar while the splitting, being $2|q|$, does not — so the proton and antiproton split their chiral partners by the same two units in opposite directions. Whether they are "the same object" depends on whether one is counting configurations or structures: as configurations they differ and no charged sector is fixed; as structure they are one orbit of size two, and the second carries no datum the first does not. What the corpus does not have is any process connecting the two sectors — no annihilation, no pair production, no cross-section — and that absence is recorded rather than filled.

\end{theorem}

\begin{theorem}[{\texorpdfstring{BM.28 — The charge-constrained probability space: the census, the three mass routes, and the splitting}{BM.28 - The charge-constrained probability space: the census, the three mass routes, and the splitting}}]\label{tag:BM-28}
\textbf{Thm.}

%dsp-level:0
\[
(1)\quad\textbf{the geometric monomials are charge-blind (the crystallized route):}
\]

%dsp-level:0
\[
\qquad \Delta x_{\mathrm B}=\mathfrak h^1E^{-1}(\ln2)^{-1},
\]

%dsp-level:0
\[
\quad V_{\mathrm{eff}}=\mathfrak h^3E^{-3}(\ln2)^{-3},\quad\rho=\mathfrak h^{-4}E^4(\ln2)^4,\quad\kappa^2\rho=2\pi\,\mathfrak h^{-2}E^4(\ln2)^3;
\]

%dsp-level:0
\[
r_s=\tfrac12\mathfrak h E,
\]

%dsp-level:0
\[
\qquad (\Delta x_{\mathrm B})^2\kappa^2\rho=2\pi E^2\ln2,
\]

%dsp-level:0
\[
\qquad \Delta x_{\mathrm B}r_s=\tfrac{\mathfrak h^2}{2\ln2}
\]

%dsp-level:0
\[
\qquad [\mathrm{BM.13}(1),\ \mathrm{BM.14}(1)\text{–}(3),\ \mathrm{BM.15}(4),\ \mathrm{BM.16}(1),(4)]:
\]

%dsp-level:0
\[
\textbf{the charge exponent is}\ 0\ \textbf{in all seven}\ [\text{machine: monomial census}]\implies\text{at fixed}\ E_{\mathrm{int}}\ \text{these seven reads are identical on every charge class};
\]

%dsp-level:0
\[
\textbf{scope:}\ \text{this clause governs reads that are functions of}\ E_{\mathrm{int}}\ \text{alone}\ —\ \text{the crystallized-count route};
\]

%dsp-level:0
\[
\qquad \text{it says nothing about ensemble-weighted reads}\ [(2)\text{–}(5)];
\]

%dsp-level:0
\[
(2)\quad\textbf{the charge constraint is a filter on the probability space, and the filter differs by class:}
\]

%dsp-level:0
\[
\qquad q_i=\tfrac12(\sigma_{s(c_i)}-\sigma_{t(c_i)})\in\{-1,0,+1\}\ [\mathrm{D3.4},\ \mathrm{EM.9}];
\]

%dsp-level:0
\[
\qquad Q=q_1+q_2+q_3\ [\mathrm{D3.9};\ \text{the additive read of}\ \mathrm{BM.6}(1)\text{'s spec}];
\]

%dsp-level:0
\[
\textbf{the end-pair census:}\quad m_q=(2,1,1)\ \text{at}\ q=(0,+1,-1)
\]

%dsp-level:0
\[
\qquad [\mathrm{EM.9}\text{'s four end-pairs};\ \mathrm{FT.106}(3\mathrm a);\ \text{machine}]:
\]

%dsp-level:0
\[
\qquad \text{the neutral read has two sign realizations, each charged read one};
\]

%dsp-level:0
\[
\textbf{the full census of the 3-constituent composite}\ (4^3=64\ \text{end-sign assignments}):
\]

%dsp-level:0
\[
\qquad \#\{Q=k\}=(1,\,6,\,15,\,\mathbf{20},\,15,\,6,\,1)\ \ \text{for}\ k=-3,\dots,+3
\]

%dsp-level:0
\[
\qquad [\text{machine: exhaustive}];
\]

%dsp-level:0
\[
\text{the generating function}\ (2+z+z^{-1})^3\ \text{reproduces it exactly and sums to}\ 64\ [\text{machine}];
\]

%dsp-level:0
\[
\textbf{the ordered charge-triple census:}\quad\#\{(q_1,q_2,q_3):\ \textstyle\sum q_i=k\}=(1,3,6,\mathbf7,6,3,1):
\]

%dsp-level:0
\[
\qquad \textbf{7 neutral triples against 6 charged}\ [\text{machine}];
\]

%dsp-level:0
\[
\textbf{the probability-space sizes:}
\]

%dsp-level:0
\[
\qquad W(0)=20,
\]

%dsp-level:0
\[
\qquad W(+1)=W(-1)=15,
\]

%dsp-level:0
\[
\qquad \boxed{\ \frac{W(0)}{W(\pm1)}=\frac{20}{15}=\frac43\ }
\]

%dsp-level:0
\[
\qquad [\text{machine: exact}];
\]

%dsp-level:0
\[
\text{one chain restricts}\ 2{:}1\ [m_q];\ \text{three chains under the additive constraint restrict}\ 4{:}3:
\]

%dsp-level:0
\[
\qquad \textbf{the constraint is what changes the ratio};
\]

%dsp-level:0
\[
(3)\quad\textbf{the scar is a boundary condition, not an internal edge:}
\]

%dsp-level:0
\[
\qquad E_{\mathrm{int}}=\#\{\text{internal directed edges}\}\ [\mathrm{D0.8}];
\]

%dsp-level:0
\[
\qquad \text{scar}=q_{\mathrm{core}}\ \text{is endpoint registration data}\ [\mathrm{D1.1},\ \mathrm{D3.3},\ \mathrm{D3.4}];
\]

%dsp-level:0
\[
\text{Census FT.35}(2)\text{'s decomposition}\ 6
\]

%dsp-level:0
\[
= (c_1{:}\,2)+(c_2{:}\,2)+(\text{connector}{:}\,2)\ \text{has}\ \textbf{no endpoint summand}\ [\text{machine}];
\]

%dsp-level:0
\[
\qquad E_{\mathrm{int}}\ \text{is unchanged across all four endpoint sign pairs}\ [\text{machine}];
\]

%dsp-level:0
\[
\therefore\ \text{the scar shifts no crystallized-count read}\ [(1)]\ —\ \textbf{and it does not follow that the scar shifts no read}: \text{the charge constraint acts on the ensemble, not on the count}\ [(2)];
\]

%dsp-level:0
\[
(4)\quad\textbf{the route theorem (the mass read is route-relative):}
\]

%dsp-level:0
\[
\qquad \textbf{(i) the normalized ensemble average}\ \ m^2_B(q,S)
\]

%dsp-level:0
\[
= \dfrac{\sum_R\nu\,t^{|R|}u^L\#\mathrm{inst}\cdot\Delta S(R)}{\sum_R\nu\,t^{|R|}u^L\#\mathrm{inst}}\ \bigg|_{t=u=1}\ [\mathrm{FT.107}(3)]:
\]

%dsp-level:0
\[
\text{a uniform sector factor cancels in a ratio}\ —\ \text{verified exactly for factors}\ 1,
\]

%dsp-level:0
\[
2,
\]

%dsp-level:0
\[
3,
\]

%dsp-level:0
\[
15,
\]

%dsp-level:0
\[
20,
\]

%dsp-level:0
\[
\tfrac73\ [\text{machine}]\implies\textbf{charge-degeneracy on this route is structural},
\]

%dsp-level:0
\[
\ \neg\text{an oversight}:
\]

%dsp-level:0
\[
\Delta S(\Gamma,R)
\]

%dsp-level:0
\[
= \tfrac32n_{\mathcal V}(\Gamma,R)
\]

%dsp-level:0
\[
= \tfrac32(2E-|R|)\ \text{is}\ \sigma\text{-free}\ [\mathrm{FT.107}(2)],
\]

%dsp-level:0
\[
\ \text{so the restriction rescales the measure without reshaping it};
\]

%dsp-level:0
\[
\textbf{(ii) the unnormalized total deficit}\ \ T(q,S)
\]

%dsp-level:0
\[
= \sum_R w(R)\,
\]

%dsp-level:0
\[
\Delta S(R)\ [\text{the}\ I_{\mathrm{hid}}\text{-type read};\ \mathrm{FT.42}(1)\ \text{pattern}]:
\]

%dsp-level:0
\[
\qquad T(0,S)=\tfrac43\,T(\pm1,S)\ \ \text{exactly}\ [\text{machine}]:
\]

%dsp-level:0
\[
\boxed{\ \Delta m\ne0\ \text{on this route},\qquad\frac{\langle m\rangle^{(n)}}{\langle m\rangle^{(p)}}=\frac43,\qquad\text{the neutral class is heavier}\ }
\]

%dsp-level:0
\[
\qquad [\text{machine: exact, ordering derived}];
\]

%dsp-level:0
\[
\textbf{(iii) the amplitude layer}\ [\mathrm{Def\text{-}H.32}(2);\ \mathrm{FT.106}(3)]:
\]

%dsp-level:0
\[
\qquad \frac{Z^{\mathrm{amp}}(0,0)}{Z^{\mathrm{amp}}(\pm1,0)}=\frac{5336}{1196}=\frac{58}{13},
\]

%dsp-level:0
\[
\qquad \frac{Z^{\mathrm{amp}}(0,1)}{Z^{\mathrm{amp}}(\pm1,1)}=\frac{4796}{802}=\frac{2398}{401}
\]

%dsp-level:0
\[
\qquad [\text{machine}]:
\]

%dsp-level:0
\[
\text{a third ratio, winding-dependent,}\ \neg\text{equal to}\ \tfrac43\ \text{or}\ 2;
\]

%dsp-level:0
\[
\qquad \text{its bridge to a mass read is item-13 content and is}\ \neg\text{constructed}\;
\]

%dsp-level:0
\[
(5)\quad\textbf{what the splitting does not touch:}
\]

%dsp-level:0
\[
\qquad W(+1)=W(-1)\ \text{exactly}\ [\text{machine}]\ \text{and}\ Z^{\bullet}(+1,S)=Z^{\bullet}(-1,S)\ [\mathrm{FT.106}(3\mathrm d)]:
\]

%dsp-level:0
\[
\textbf{the proton and the antiproton remain degenerate on all three routes}\ —\ \text{the charge restriction splits the neutral class from the charged pair and never splits the pair}\ [\mathrm{BM.27}(3),(4)];
\]

%dsp-level:0
\[
\text{CPT maps the}\ Q
\]

%dsp-level:0
\[
= +1\ \text{triple set bijectively onto the}\ Q
\]

%dsp-level:0
\[
= -1\ \text{set, and the neutral set onto itself with}\ (0,0,0)\ \text{the}\ \textbf{unique CPT-fixed charge triple}\ [\text{machine}];
\]

%dsp-level:0
\[
(6)\quad\textbf{the identity and difference sets:}
\]

%dsp-level:0
\[
\qquad \textbf{identical on every route:}\ E_{\mathrm{int}},
\]

%dsp-level:0
\[
\ \Delta x_{\mathrm B},
\]

%dsp-level:0
\[
\ V_{\mathrm{eff}},
\]

%dsp-level:0
\[
\ \rho,
\]

%dsp-level:0
\[
\ \kappa^2\rho,
\]

%dsp-level:0
\[
\ r_s,
\]

%dsp-level:0
\[
\ \text{both locks},
\]

%dsp-level:0
\[
\ \Theta_{\mathrm{int}}\ \text{class},
\]

%dsp-level:0
\[
\ |w|^2,
\]

%dsp-level:0
\[
\ 2|q|\ \text{restricted to}\ \{|q|=0\}\ \text{vs}\ \{|q|=1\}\ [(1)];
\]

%dsp-level:0
\[
\textbf{route-dependent:}\ \ \text{the ensemble mass read}\ —\ \text{degenerate on}\ (4)(\mathrm i),
\]

%dsp-level:0
\[
\ \text{split by}\ \tfrac43\ \text{on}\ (4)(\mathrm{ii}),
\]

%dsp-level:0
\[
\ \text{split by}\ \tfrac{58}{13}/\tfrac{2398}{401}\ \text{on}\ (4)(\mathrm{iii});
\]

%dsp-level:0
\[
\textbf{distinguishing on every route:}\ \ F_{LR}
\]

%dsp-level:0
\[
= qI_2,
\]

%dsp-level:0
\[
\ q_{\mathrm{core}},
\]

%dsp-level:0
\[
\ k,
\]

%dsp-level:0
\[
\ \operatorname{sign}(k),
\]

%dsp-level:0
\[
\ \text{the CPT orientation},
\]

%dsp-level:0
\[
\ \text{the sector label},
\]

%dsp-level:0
\[
\ \text{and now}\ W(Q)\ [\text{machine: the sets are disjoint}];
\]

%dsp-level:0
\[
(7)\quad\textbf{no fractional charge exists, structurally:}
\]

%dsp-level:0
\[
\qquad q\in\{-1,0,+1\}\subset\mathbb Z\ [\mathrm{EM.9},\ \mathrm{P2}]\ \text{and}\ Q=\sum_iq_i\implies Q\in\mathbb Z\ \text{with}\ Q\in\{-3,\dots,+3\}\ [\text{machine}]:
\]

%dsp-level:0
\[
\tfrac13,\ \tfrac23\notin\{-3,\dots,+3\}:
\]

%dsp-level:0
\[
\qquad \textbf{fractional charges are impossible under an integer-valued additive read},\ \neg\text{merely absent};
\]

%dsp-level:0
\[
\text{deficit}\ 1'\text{'s charge side}\ [\mathrm{FT.125}(4)]\ \text{is therefore not closable by this read}: \text{a fractional constituent alphabet requires a different read, and none is constructed};
\]

%dsp-level:0
\[
(8)\quad\textbf{the read-choice row:}
\]

%dsp-level:0
\[
\qquad \text{which of}\ (4)\text{'s three routes is ``the mass'' is a read specification},
\]

%dsp-level:0
\[
\ \neg\text{a theorem}\ [\text{Census FT.76's priced-variant pattern}];
\]

%dsp-level:0
\[
\text{all three are exact and machine-verified};
\]

%dsp-level:0
\[
\qquad \text{no route imports a measured value}\ [\mathrm{P5}];
\]

%dsp-level:0
\[
\qquad \neg\exists\ \text{a P4 parameter and}\ \neg\exists\ \text{a new dimensionful scale}\ [\mathrm{P1};\ \text{§Cal-BM stays empty}].
\]

\textbf{Rmk (route discipline).} Clause (1) and clause (4)(i) are not in tension with clause (4)(ii): they are statements about different functionals of the same ensemble. A normalized average is blind to any uniform rescaling of the measure; a total is not. The charge constraint of clause (2) is exactly such a rescaling — uniform within a sector, different between sectors — so it is invisible to the first and visible to the second, with the exact factor $4/3$. Any reading of this tag that quotes one route against another is reading a scope clause as a claim.

\textbf{Req.} D0.8, D1.1, D2.6, D3.3, D3.4, D3.9 (Parts 1–2); CT2.1, CT2.9 (Part 2); EM.9, EM.24(1) (Part 2); P1, P2, P5 (Parts 1–2); Def-BM.4, Def-BM.7, BM.13(1), BM.14(1), BM.14(2), BM.14(3), BM.15(4), BM.16(1), BM.16(4), BM.22(2), BM.22(4), BM.27(3), BM.27(4) (this part); Def-H.32, FT.42(1), FT.84(1), FT.90(5), FT.106(3), FT.107(2), FT.107(3), FT.125(4), FT.129(3) (Part 3); Census FT.35(2) (row-wise); the additive read of BM.6(1)'s specification (constructed in Tier BM.8, cited there as a forward pointer and imported in no \textbf{Req} line); exhaustive end-sign enumeration, generating-function convolution, exact rational ratio arithmetic [math; machine: \texttt{checks/\allowbreak{}verify\_\allowbreak{}H\_\allowbreak{}probability\_\allowbreak{}space\_\allowbreak{}restriction.py} §A–§C].

\textbf{Pf.}

%dsp-level:0
\[
(1):\ \text{the seven monomials re-cited from Tiers BM.5's rows and their exponent vectors read off}\ [\text{machine}];
\]

%dsp-level:0
\[
\qquad q\text{-exponent}\ 0\ \text{in each}\ [\text{machine}]
\]

%dsp-level:0
\[
(2):\ \mathrm{EM.9}\text{'s table}\ \to m_q=(2,1,1)\ [\text{machine}];
\]

%dsp-level:0
\[
\qquad \text{exhaustive enumeration of}\ 4^3
\]

%dsp-level:0
\[
= 64\ \text{end-sign assignments}\ \to\ \#\{Q=k\}
\]

%dsp-level:0
\[
= (1,6,15,20,15,6,1)\ [\text{machine}];
\]

%dsp-level:0
\[
\qquad (2+z+z^{-1})^3\ \text{convolved three times reproduces the same census and sums to}\ 64\ [\text{machine}]
\]

%dsp-level:0
\[
\text{triples}: \#\{\textstyle\sum q_i=k\}=(1,3,6,7,6,3,1),\ \text{summing to}\ 27=3^3\ [\text{machine}];
\]

%dsp-level:0
\[
\qquad W(k)=\sum_{\sum q_i=k}\prod_im_{q_i}\implies W(0)=8+12=20,\ W(\pm1)=12+3=15\ [\text{machine}];
\]

%dsp-level:0
\[
\qquad 20/15=4/3\ [\text{exact}]
\]

%dsp-level:0
\[
(3):\ \mathrm{D0.8}\ \text{counts internal directed edges};\ \mathrm{D3.4}\ \text{reads two endpoint signs};
\]

%dsp-level:0
\[
\qquad \text{Census FT.35}(2)\text{'s three summands are edge counts}\ [\text{machine}];
\]

%dsp-level:0
\[
\qquad \text{invariance across the four sign pairs}\ [\text{machine}]
\]

%dsp-level:0
\[
(4)(\mathrm i):\ \text{write}\ w_q(R)
\]

%dsp-level:0
\[
= c_qv(R)\ \text{with}\ v\ q\text{-free}: \frac{\sum c_qvm}{\sum c_qv}
\]

%dsp-level:0
\[
= \frac{\sum vm}{\sum v}\ \text{for every}\ c_q
\]

%dsp-level:0
\[
\ne 0\ [\text{machine at}\ c\in\{1,2,3,15,20,\tfrac73\}];
\]

%dsp-level:0
\[
\qquad \Delta S=\tfrac32(2E-|R|)\ \text{is}\ \sigma\text{-free}\ [\mathrm{FT.107}(2)]
\]

%dsp-level:0
\[
(4)(\mathrm{ii}):\ T(q,S)
\]

%dsp-level:0
\[
= W(q)\cdot\overline{\Delta S}(S)\ \text{with the same}\ \overline{\Delta S}\ [(\mathrm i)]
\]

%dsp-level:0
\[
\implies T(0)/T(\pm1)
\]

%dsp-level:0
\[
= W(0)/W(\pm1)
\]

%dsp-level:0
\[
= 4/3\ [\text{machine}];
\]

%dsp-level:0
\[
\qquad 4/3>1\implies\text{the neutral class is heavier}\ [\text{machine}]
\]

%dsp-level:0
\[
(4)(\mathrm{iii}):\ \mathrm{FT.106}(3)\text{'s amplitude integers re-cited};
\]

%dsp-level:0
\[
\qquad 5336/1196=58/13\ \text{and}\ 4796/802=2398/401\ [\text{machine: exact reduction}];
\]

%dsp-level:0
\[
\qquad 58/13\ne4/3\ne2\ [\text{machine}]
\]

%dsp-level:0
\[
(5):\ W(+1)=W(-1)=15\ [\text{machine}];
\]

%dsp-level:0
\[
\qquad \mathrm{FT.106}(3\mathrm d)\text{'s}\ \omega\text{-symmetry re-cited};
\]

%dsp-level:0
\[
\qquad \text{CPT}: q_i\mapsto-q_i\ \text{gives}\ Q\mapsto-Q\ \text{and a bijection of triple sets};\ \text{fixed triples}: \{(0,0,0)\}\ [\text{machine}]
\]

%dsp-level:0
\[
(6):\ \text{set arithmetic on the three lists}\ [\text{machine}];
\]

%dsp-level:0
\[
\qquad (7):\ Q\in\mathbb Z\ \text{as a sum of integers, range}\ [-3,3]\ [\text{machine}];
\]

%dsp-level:0
\[
\qquad \tfrac13\notin\mathbb Z\ [\text{exact}]
\]

%dsp-level:0
\[
(8):\ \text{Census FT.76's priced-variant pattern re-cited};
\]

%dsp-level:0
\[
\qquad \mathrm{P1}/\mathrm{P5}\ \text{scan of this tag's text}\ [\text{machine}].\ \blacksquare
\]

\textbf{Phys.} The charge constraint does exactly what a constraint on a sum does: it changes how many microscopic sign assignments land in each class. For one chain the neutral read has two realizations and each charged read one; for three chains bound by $q_1+q_2+q_3=Q$ the full census is $1,6,15,20,15,6,1$ across $Q=-3$ to $+3$, sixty-four assignments in all, and the neutral class ends up with twenty against the unit class's fifteen. That ratio, four thirds, is the whole content of the restriction. Whether it reaches the mass depends on what functional the mass is. A normalized ensemble average cannot see it — a uniform rescaling of the measure cancels between numerator and denominator, which is why the sector mass table is charge-degenerate as a matter of structure rather than luck. A total deficit can see it, and does: on that route the neutral class carries exactly four thirds the information content of either charged class, and is heavier. The amplitude layer carries a third ratio, fifty-eight thirteenths at zero winding and a different one at winding one, and its bridge to a mass read is still an open item. Two things survive every route: the proton and the antiproton stay degenerate, because the census is symmetric under charge conjugation and $(0,0,0)$ is the only self-conjugate triple; and no fractional charge can appear at all, since an integer sum of ternary reads cannot land on a third. What the corpus does not do is pick a route for the reader: all three are exact, all three are priced, and the choice is a read specification.

\end{theorem}

\medskip\noindent\textit{\texorpdfstring{BM.29 — The status of the CPT sectors: adjudication and scope}{BM.29 - The status of the CPT sectors: adjudication and scope}.}\ Structural-constraints record: see Appendix B, Section B.1.18 of the Addendum.\par\medskip

\phantomsection
\section[Tier BM.8 — The Additive Charge Read and the Mass-Route…]{\texorpdfstring{Tier BM.8 — The Additive Charge Read and the Mass-Route Adjudication}{Tier BM.8 - The Additive Charge Read and the Mass-Route Adjudication}}

\begin{definition}[{\texorpdfstring{Def-BM.8 — The additive charge read (the construction BM.6(1) specifies)}{Def-BM.8 - The additive charge read (the construction BM.6(1) specifies)}}]\label{tag:Def-BM-8}
\textbf{Def.} A read $Q$ on sealed composites, constructed to BM.6(1)'s three specification clauses; no new axiom, no new datum.

%dsp-level:0
\[
(1)\quad\textbf{the read:}
\]

%dsp-level:0
\[
\qquad Q(c_1,\dots,c_k)\ :=\ \sum_{i=1}^{k}q_{\mathrm{core}}(c_i),
\]

%dsp-level:0
\[
\qquad q_{\mathrm{core}}(c)=\tfrac12\bigl(\sigma_{s(c)}-\sigma_{t(c)}\bigr)\in\{-1,0,+1\}
\]

%dsp-level:0
\[
\qquad [\mathrm{D3.4},\ \mathrm{EM.9},\ \mathrm{EM.9a},\ \mathrm{FT.40}];
\]

%dsp-level:0
\[
(2)\quad\textbf{spec clause (i) — restriction:}\qquad Q|_{\text{single chain}}=q_{\mathrm{core}}\qquad[\text{immediate at}\ k=1];
\]

%dsp-level:0
\[
(3)\quad\textbf{spec clause (ii) — additivity:}
\]

%dsp-level:0
\[
\qquad Q(\mathcal X\sqcup\mathcal Y)=Q(\mathcal X)+Q(\mathcal Y)\ \ \text{for disjoint sealings}
\]

%dsp-level:0
\[
\qquad [\text{the sum is over disjoint constituent sets}];
\]

%dsp-level:0
\[
(4)\quad\textbf{spec clause (iii) — observability on sealed supports:}
\]

%dsp-level:0
\[
\qquad Q\ \text{is defined only where}\ \mathcal C(\Sigma;C)=1\ [\mathrm{D2.14}:\ \partial_{\mathrm{col}}(\Sigma;C)=\varnothing];
\]

%dsp-level:0
\[
\qquad \text{one dangling attachment zeroes the whole colour sector}\ [\mathrm{FT.19}]\ \text{and the read with it};
\]

%dsp-level:0
\[
(5)\quad\textbf{the derived range and parity:}
\]

%dsp-level:0
\[
\qquad k=3:\quad Q\in\{-3,-2,-1,0,+1,+2,+3\}\subset\mathbb Z
\]

%dsp-level:0
\[
\qquad [\text{machine}];
\]

%dsp-level:0
\[
\qquad k=2\ (\delta\text{-seal}):\ Q\in\{-2,\dots,+2\};
\]

%dsp-level:0
\[
\textbf{integer-valued at every}\ k:\qquad\tfrac13,\ \tfrac23\ \text{are unreachable}\ [\mathrm{BM.28}(7)];
\]

%dsp-level:0
\[
(6)\quad\textbf{the probability-space weights it induces:}
\]

%dsp-level:0
\[
\qquad W(Q):=\#\{(\sigma\text{-assignments}):\ \textstyle\sum_iq_i=Q\}=\bigl[z^{Q}\bigr](2+z+z^{-1})^{k}
\]

%dsp-level:0
\[
\qquad [\mathrm{BM.28}(2);\ \text{machine}]:
\]

%dsp-level:0
\[
k=3:\quad W=(1,6,15,20,15,6,1)\ \text{over}\ Q=-3,\dots,+3,\qquad\sum_QW(Q)=4^3=64\qquad[\text{machine}];
\]

%dsp-level:0
\[
(7)\quad\textbf{the CPT action:}
\]

%dsp-level:0
\[
\qquad \mathrm{CPT}:\ q_i\mapsto-q_i\ [\mathrm{EM.24}(1)]\implies Q\mapsto-Q;
\]

%dsp-level:0
\[
\qquad W(Q)=W(-Q)\ [\text{machine}];
\]

%dsp-level:0
\[
\qquad \text{the unique CPT-fixed triple at}\ k=3\ \text{is}\ (0,0,0)\ [\text{machine}];
\]

%dsp-level:0
\[
(8)\quad\textbf{status and scope:}
\]

%dsp-level:0
\[
\qquad \textbf{constructed}\ —\ \text{BM.6}(1)\text{'s specification is discharged for the sealed classes}\ k\le3;
\]

%dsp-level:0
\[
\qquad \text{the charge-class labels of}\ \mathrm{Def\text{-}BM.3}(4)\ \text{now rest on a derived read rather than on a declaration};
\]

%dsp-level:0
\[
\neg\text{constructed: the coupling of}\ Q\ \text{to a gauge field}\ [\text{deficit}\ 1',\ \mathrm{FT.125}(4);\ \mathrm{FT.82}(7)(\mathrm a)],
\]

%dsp-level:0
\[
\ \text{and any dynamics of}\ Q\ [\mathrm{BM.6}(5)].
\]

\textbf{Rmk (what the construction does not do).} $Q$ is a sum of the ternary end reads. It supplies no fractional constituent alphabet, no gauge coupling, and no mass: the mass routes are BM.28(4)'s and BM.30's, and the read-choice among them is not made here.

\textbf{Req.} D2.14, D3.4, D3.9 (Parts 1–2); EM.9, EM.9a, EM.24(1) (Part 2); P2 (Part 1); Def-BM.3(4), BM.6(1), BM.28(2), BM.28(7) (this part); FT.19, FT.40, FT.82(7), FT.125(4) (Part 3); generating-function convolution, exhaustive end-sign enumeration [math; machine: \texttt{checks/\allowbreak{}verify\_\allowbreak{}H\_\allowbreak{}probability\_\allowbreak{}space\_\allowbreak{}restriction.py} §A, §D].

\textbf{Pf.}

%dsp-level:0
\[
(1)\text{–}(4):\ \text{the three specification clauses of BM.6}(1)\ \text{discharged clause by clause};
\]

%dsp-level:0
\[
\qquad \text{observability}: \mathcal C(\Sigma;C)\ \text{is the}\ \Pi_{\mathcal C}\ \text{projector's support}\ [\mathrm{D2.14}];
\]

%dsp-level:0
\[
\ \mathrm{FT.19}\text{'s all-or-nothing row re-cited}
\]

%dsp-level:0
\[
(5):\ \text{a sum of}\ k\ \text{integers each in}\ \{-1,0,1\}\ \text{lies in}\ \{-k,\dots,k\}\ [\text{machine at}\ k=2,3];
\]

%dsp-level:0
\[
\qquad (6):\ (2+z+z^{-1})^3\ \text{convolved}\ [\text{machine}];
\]

%dsp-level:0
\[
\qquad \textstyle\sum_QW(Q)=(2+1+1)^3=64\ [\text{machine}]
\]

%dsp-level:0
\[
(7):\ q_i
\]

%dsp-level:0
\[
\mapsto -q_i\ \text{componentwise}\ [\mathrm{EM.24}(1)]
\]

%dsp-level:0
\[
\implies Q
\]

%dsp-level:0
\[
\mapsto -Q\ \text{and}\ W(Q)
\]

%dsp-level:0
\[
= W(-Q)\ \text{by the census's symmetry}\ [\text{machine}];
\]

%dsp-level:0
\[
\qquad \text{fixed triples}: -q_i=q_i\ \forall i\iff q_i=0\ \forall i\ [\text{machine}].\ \blacksquare
\]

\textbf{Phys.} The additive charge read the era listed as an open construction is now built, and it costs nothing: sum the ternary endpoint reads of the constituents, restrict to sealed supports, and the three clauses the specification demanded are satisfied by inspection. Two things come out of it that the specification did not ask for. The read is integer-valued at every constituent number, so a fractional charge is not missing from this corpus but unreachable by it. And the read carries its own statistics: the number of sign assignments realizing each net charge is the coefficient of a cubic generating function, giving twenty for the neutral class against fifteen for either unit class, which is the exact measure of how differently the charge constraint filters the two probability spaces.

\end{definition}

\begin{theorem}[{\texorpdfstring{BM.30 — The mass-route adjudication: three exact ratios, one read choice}{BM.30 - The mass-route adjudication: three exact ratios, one read choice}}]\label{tag:BM-30}
\textbf{Thm.}

%dsp-level:0
\[
(1)\quad\textbf{the three routes and their exact charge ratios:}\qquad\begin{array}{l|c|c}\text{route} & \text{functional} & \text{neutral}:\text{charged}\\\hline(\mathrm i)\ \text{normalized ensemble average}\ [\mathrm{FT.107}(3)] & \dfrac{\sum_Rw\,\Delta S}{\sum_Rw} & 1\ \ (\text{degenerate})\\[2mm](\mathrm{ii})\ \text{total deficit}\ [I_{\mathrm{hid}}\text{-type};\ \mathrm{FT.42}(1)] & \sum_Rw\,\Delta S & \tfrac43\\[1mm](\mathrm{iii})\ \text{amplitude layer}\ [\mathrm{Def\text{-}H.32}(2);\ \mathrm{FT.106}(3)] & \sum_Rv_\Gamma(R)\,\Delta S & \tfrac{58}{13}\ (S{=}0),\ \tfrac{2398}{401}\ (S{=}1)\end{array}
\]

%dsp-level:0
\[
\text{all three exact}\ [\text{machine}];\qquad\text{the three ratios are pairwise unequal}\ [\text{machine}];
\]

%dsp-level:0
\[
(2)\quad\textbf{the splitting, where it exists:}
\]

%dsp-level:0
\[
\qquad \text{on route}\ (\mathrm{ii}):
\]

%dsp-level:0
\[
\qquad \frac{\langle m\rangle^{(n)}}{\langle m\rangle^{(p)}}=\frac{W(0)}{W(+1)}=\frac{20}{15}=\frac43,
\]

%dsp-level:0
\[
\qquad \Delta m\ne0,
\]

%dsp-level:0
\[
\qquad \textbf{the neutral class is heavier}\ [\mathrm{BM.28}(4)(\mathrm{ii});\ \text{machine}];
\]

%dsp-level:0
\[
\text{the mechanism is the measure, not the deficit:}\ \Delta S
\]

%dsp-level:0
\[
= \tfrac32(2E-|R|)\ \text{is}\ \sigma\text{-free}\ [\mathrm{FT.107}(2)],
\]

%dsp-level:0
\[
\ \text{so the restriction rescales without reshaping};
\]

%dsp-level:0
\[
(3)\quad\textbf{the degeneracy, where it is structural:}
\]

%dsp-level:0
\[
\qquad \text{on route}\ (\mathrm i)\ \text{any uniform sector factor cancels in the ratio}\ [\text{machine: factors}\ 1,2,3,15,20,\tfrac73]:
\]

%dsp-level:0
\[
\qquad \textbf{charge-degeneracy is a theorem about normalized averages},\ \neg\text{an assumption about composites};
\]

%dsp-level:0
\[
(4)\quad\textbf{what survives every route:}
\]

%dsp-level:0
\[
\qquad W(+1)=W(-1)\ \text{and}\ Z^{\bullet}(+1,S)=Z^{\bullet}(-1,S)\ [\mathrm{BM.28}(5);\ \mathrm{FT.106}(3\mathrm d)]:
\]

%dsp-level:0
\[
\qquad \textbf{the proton and the antiproton are degenerate on all three routes};
\]

%dsp-level:0
\[
\text{the crystallized-count reads are charge-blind}\ [\mathrm{BM.28}(1)]:
\]

%dsp-level:0
\[
\qquad E_{\mathrm{int}},\ \Delta x_{\mathrm B},\ V_{\mathrm{eff}},\ \rho,\ \kappa^2\rho,\ r_s,\ \text{both locks are shared by all three classes};
\]

%dsp-level:0
\[
(5)\quad\textbf{the read-choice row:}
\]

%dsp-level:0
\[
\qquad \neg\exists\ \text{a theorem of Parts 1–4 selecting one route as ``the mass''}\ [\text{inventory: the}\ m_{\mathrm{hid}},\ \text{pole, sector-expectation and}\ m^2_B\ \text{reads are distinct}\ [\mathrm{FT.84}(1);\ \mathrm{FT.87}(8)]];
\]

%dsp-level:0
\[
\text{adoption requires a logged declaration}\ [\text{Census FT.76's priced-variant pattern};\ \mathrm{BM.6}(4)\ \text{precedent}];
\]

%dsp-level:0
\[
\qquad \textbf{no route is adopted here};
\]

%dsp-level:0
\[
(6)\quad\textbf{the ordering is derived and uncompared:}
\]

%dsp-level:0
\[
\qquad \langle m\rangle^{(n)}>\langle m\rangle^{(p)}\ \text{on route}\ (\mathrm{ii})\ [\text{machine:}\ \tfrac43>1];
\]

%dsp-level:0
\[
\qquad \neg\exists\ \text{a measured mass value in any clause and}\ \neg\exists\ \text{a comparison}\ [\mathrm{P5}];
\]

%dsp-level:0
\[
(7)\quad\textbf{the status of the adjudication:}
\]

%dsp-level:0
\[
\qquad \textbf{derived:}\ (1)\text{–}(4),\ (6);
\]

%dsp-level:0
\[
\qquad \textbf{open:}\ \text{the route choice}\ (5),\ \text{the amplitude}\leftrightarrow\text{mass bridge}\,\ \text{the gauge coupling of}\ Q\;
\]

%dsp-level:0
\[
\textbf{closed by this tier:}\ \mathrm{BM.6}(1)\text{'s construction spec}\ [\mathrm{Def\text{-}BM.8}(8)];
\]

%dsp-level:0
\[
\qquad \textbf{unclosable by this read:}\ \text{the fractional constituent alphabet}\ [\mathrm{BM.28}(7);\ \mathrm{Def\text{-}BM.8}(5)].
\]

\textbf{Rmk (no route is privileged).} Clause (2)'s splitting and clause (3)'s degeneracy are both exact and neither is a correction of the other: they are different functionals of one ensemble. A corpus that adopted route (ii) would have a neutron heavier than its proton by the factor four thirds; a corpus that adopted route (i) would have them equal; both statements are machine-verified, and the choice between them is the read specification that clause (5) declines to make.

\textbf{Req.} D2.9 (Parts 1–2); CT2.1 (Part 2); P2, P5 (Parts 1–2); Def-H.32, FT.42(1), FT.84(1), FT.87(8), FT.106(3), FT.107(2), FT.107(3), FT.125(4) (Part 3); Def-BM.8, BM.6(1), BM.6(4), BM.28(1), BM.28(2), BM.28(4), BM.28(5), BM.28(7) (this part); Census FT.76 (the priced-variant pattern); exact rational ratio arithmetic, exhaustive enumeration [math; machine: \texttt{checks/\allowbreak{}verify\_\allowbreak{}H\_\allowbreak{}probability\_\allowbreak{}space\_\allowbreak{}restriction.py} §B–§D].

\textbf{Pf.}

%dsp-level:0
\[
(1):\ \mathrm{BM.28}(4)\text{'s three rows re-cited with their machine values};
\]

%dsp-level:0
\[
\qquad \tfrac{58}{13}\ne\tfrac43\ \text{since}\ 58\cdot3=174\ne13\cdot4=52;
\]

%dsp-level:0
\[
\qquad \tfrac43\ne1\ [\text{machine}]
\]

%dsp-level:0
\[
(2):\ T(q,S)
\]

%dsp-level:0
\[
= W(q)\overline{\Delta S}(S)\ \text{with a common}\ \overline{\Delta S}\ [\mathrm{BM.28}(4)(\mathrm i),(\mathrm{ii})]\implies\text{the ratio is}\ W(0)/W(+1)
\]

%dsp-level:0
\[
= 4/3\ [\text{machine}]
\]

%dsp-level:0
\[
(3):\ \text{the cancellation lemma}\ [\mathrm{BM.28}(4)(\mathrm i);\ \text{machine}]
\]

%dsp-level:0
\[
(4):\ \mathrm{BM.28}(5)\ \text{and}\ \mathrm{BM.28}(1)\ \text{re-cited}
\]

%dsp-level:0
\[
(5):\ \text{inventory scan of the mass reads}\ [\mathrm{FT.84}(1)\text{'s read-class taxonomy};\ \mathrm{FT.87}(8)];
\]

%dsp-level:0
\[
\qquad \text{Census FT.76's variant-pricing pattern re-cited}
\]

%dsp-level:0
\[
(6):\ \tfrac43>1\ [\text{machine}];\qquad\mathrm{P5}\ \text{scan of this tag's text}\ [\text{machine: zero external numeric tokens}]
\]

%dsp-level:0
\[
(7):\ \text{conjunction of}\ (1)\text{–}(6)\ \text{with}\ \mathrm{Def\text{-}BM.8}(8).\ \blacksquare
\]

\textbf{Phys.} One ensemble, three functionals, three answers — and all three are exact. Normalize and the charge constraint disappears, because a uniform rescaling of the measure cancels between numerator and denominator; that is why the sector mass table is charge-degenerate, and it is a structural fact rather than a lucky one. Sum instead of averaging and the constraint is all that matters: twenty sign assignments realize the neutral class against fifteen for either charged class, so the neutral composite carries four thirds the information deficit and weighs more. Take the amplitude layer and a third ratio appears, fifty-eight thirteenths at zero winding and something else at winding one, with its bridge to a mass read still open. What no route disturbs is the mirror: the census is symmetric under charge conjugation, so the proton and antiproton stay degenerate everywhere, and the seven crystallized reads stay blind to charge entirely. The corpus prices the three routes and adopts none, because adopting one is a read specification and not a theorem — and it records the ordering it does derive, neutral heavier, without comparing it to anything.

\end{theorem}

\phantomsection
\section[Tier BM.9 — The Mesonic Flicker Test: the Contraction…]{\texorpdfstring{Tier BM.9 — The Mesonic Flicker Test: the Contraction Grammar, the Colour Projector, and the Protected Junction}{Tier BM.9 - The Mesonic Flicker Test: the Contraction Grammar, the Colour Projector, and the Protected Junction}}

\begin{definition}[{\texorpdfstring{Def-BM.9 — The extended flicker ensemble and the sealing filter}{Def-BM.9 - The extended flicker ensemble and the sealing filter}}]\label{tag:Def-BM-9}
\textbf{Def.} A construction class on Def-BM.5 (no new axiom; Def-H.32's scopes (S1)–(S4) and Def-BM.5's inherited clauses stand).

%dsp-level:0
\[
(1)\quad\textbf{the two restricted ensembles:}
\]

%dsp-level:0
\[
\qquad \mathcal E_\varepsilon(\mathcal X):=\{R\in\mathcal E(\mathcal X):\ \text{the sealing of}\ R\ \text{is}\ \varepsilon_{ijk}\text{-type}\},
\]

%dsp-level:0
\[
\qquad \mathcal E_\delta(\mathcal X):=\{R\in\mathcal E(\mathcal X):\ \delta^i{}_j\text{-type}\}
\]

%dsp-level:0
\[
\qquad [\mathrm{D3.9};\ \mathrm{CT3.17}];
\]

%dsp-level:0
\[
\mathcal E_{\mathrm{ext}}(\mathcal X):=\mathcal E_\varepsilon(\mathcal X)\cup\mathcal E_\delta(\mathcal X);
\]

%dsp-level:0
\[
\qquad \textbf{constitutive clause:}\ \text{the union is taken inside}\ \mathrm{Def\text{-}H.32}(1)\text{'s topology space, which is colour-blind}\ [\mathrm{D2.13}:\ \mathcal M(\Sigma)\ \text{is subject to the rung masks and Janus balance},\ \neg\text{to colour sealing}];
\]

%dsp-level:0
\[
(2)\quad\textbf{the filter clause (the load-bearing one):}\qquad\text{colour sealing is ambient-global}\ [\mathrm{CT3.18}]\ \text{and acts}\ \textbf{downstream}\ \text{as the projector}\ \Pi_{\mathcal C}(\Sigma;C)=\mathcal C(\Sigma;C)\cdot\mathbb I^{(\mathrm{strong})}_\Sigma,\qquad\mathcal C(\Sigma;C)=\begin{cases}1&\partial_{\mathrm{col}}(\Sigma;C)=\varnothing\\0&\text{otherwise}\end{cases}\qquad[\mathrm{D2.14}];
\]

%dsp-level:0
\[
\text{one dangling attachment zeroes the entire colour block}\ [\mathrm{FT.19}:\ \text{all-or-nothing}];
\]

%dsp-level:0
\[
\qquad \mathcal C\ \text{is an idempotent scalar},\ \mathcal C^2=\mathcal C\in\{0,1\}\ [\text{machine}];
\]

%dsp-level:0
\[
\textbf{non-derived clause:}\ \ \text{none}\ —\ (2)\ \text{is}\ \mathrm{D2.14}\text{'s derived projector, cited at the site where the union of}\ (1)\ \text{would otherwise be read as a sum over sealings};
\]

%dsp-level:0
\[
(3)\quad\textbf{the grouping grammar:}
\]

%dsp-level:0
\[
\qquad g(N):=\#\{\text{partitions of}\ N\ \text{attachments into contraction groups of size}\ 2\ (\delta)\ \text{and}\ 3\ (\varepsilon)\}
\]

%dsp-level:0
\[
\qquad [\mathrm{CT3.17a}(2);\ \mathrm{FT.41}(\mathrm b)]:
\]

%dsp-level:0
\[
g(N)=(1,\,1,\,3,\,10,\,25,\,105)\quad\text{for}\ N=2,\dots,7
\]

%dsp-level:0
\[
\qquad [\text{machine: exhaustive partition enumeration reproduces all six values}];
\]

%dsp-level:0
\[
g_\delta(N):=\#\{\text{groupings containing at least one}\ \delta\}=(1,\,\mathbf0,\,3,\,10,\,15,\,105),
\]

%dsp-level:0
\[
\qquad g_\varepsilon(N):=\#\{\text{groupings of}\ 3\text{-blocks only}\}=(0,\,1,\,0,\,0,\,10,\,0);
\]

%dsp-level:0
\[
g_\delta(N)+g_\varepsilon(N)=g(N)\ \ \forall N\ [\text{machine}];\qquad g_\varepsilon(N)\ne0\iff3\mid N\ [\text{machine on the frozen range}];
\]

%dsp-level:0
\[
(4)\quad\textbf{the boundary-data tuple:}
\]

%dsp-level:0
\[
\qquad \text{the labels that condition}\ \mathcal E(\mathcal X)\ \text{are}\ \bigl(Q,\ \Theta_{\mathrm{int}}\ \text{class},\ S,\ \text{attachment count}\ N\bigr)\ [\mathrm{Def\text{-}BM.5}(1);\ \mathrm{Def\text{-}BM.8}]:
\]

%dsp-level:0
\[
\textbf{the attachment count is a boundary datum},
\]

%dsp-level:0
\[
\ \neg\text{a summed-over structure}\ —\ \text{a}\ 3\text{-junction and a}\ 2\text{-junction are different supports, not two members of one ensemble};
\]

%dsp-level:0
\[
(5)\quad\textbf{the scope clause:}
\]

%dsp-level:0
\[
\qquad \text{no clause of this part sums over sealings of a fixed support as if they were alternatives}: \text{the grammar of}\ (3)\ \text{decides which sealings a support admits at all}.
\]

\textbf{Rmk (why the filter, not a cancellation).} The topology-blind path integral is colour-blind by construction — that is what "topology-blind" means here [D2.13's scope note; CT3.22a Rmk]. Colour enters only through $\Pi_{\mathcal C}$, a $0/1$ projector applied after the sum. So the question "do $\delta$-sealed states cancel in the baryon's integral?" has a category answer before it has an arithmetic one: they are not terms that could cancel, they are terms the projector never admits.

\textbf{Req.} D2.13, D2.14, D3.9 (Parts 1–2); CT3.17, CT3.17a(2), CT3.18, CT3.22a (Part 2); Def-H.32 (Part 3); FT.19, FT.41(b) (Part 3); Def-BM.5, Def-BM.8 (this part); exhaustive set-partition enumeration over blocks of size 2 and 3, exact integer arithmetic [math; machine: \texttt{checks/\allowbreak{}verify\_\allowbreak{}H\_\allowbreak{}mesonic\_\allowbreak{}flicker\_\allowbreak{}test.py} §A].

\textbf{Pf.}

%dsp-level:0
\[
\text{Definition block: the constitutive clause is}\ (1)\text{'s union inside the colour-blind topology space};
\]

%dsp-level:0
\[
\ (2)\ \text{is}\ \mathrm{D2.14}\text{'s projector re-cited at the site where the union is formed};
\]

%dsp-level:0
\[
\ (3)\text{'s counts are machine-reproduced from}\ \mathrm{CT3.17a}(2)\text{'s frozen list};
\]

%dsp-level:0
\[
\ (4)\text{–}(5)\ \text{are scope clauses}.\ \blacksquare
\]

\textbf{Phys.} Before any sum is taken, one structural fact settles the question. The topology-blind path integral does not filter by colour — it filters by topology, and colour enters afterwards as a zero-one projector on the support. A sealing is therefore not an alternative term competing inside one integral; it is a property of the support that either passes the projector or annihilates the whole colour block. And the grammar that decides it is small enough to enumerate: partition the attachments into groups of two and three, and count.

\end{definition}

\begin{theorem}[{\texorpdfstring{BM.31 — The boundary-data comparison and the grammar verdict}{BM.31 - The boundary-data comparison and the grammar verdict}}]\label{tag:BM-31}
\textbf{Thm.}

%dsp-level:0
\[
(1)\quad\textbf{the grammar at}\ N=3\ \textbf{is forced:}
\]

%dsp-level:0
\[
\qquad g(3)=1,
\]

%dsp-level:0
\[
\qquad \text{the unique grouping is the}\ \varepsilon\text{-triple},
\]

%dsp-level:0
\[
\qquad g_\delta(3)=\mathbf0
\]

%dsp-level:0
\[
\qquad [\mathrm{Def\text{-}BM.9}(3);\ \text{machine: exhaustive}];
\]

%dsp-level:0
\[
\text{a}\ \delta\ \text{sub-pairing on a}\ 3\text{-attachment support consumes two attachments and leaves exactly one dangling}\ (3-2=1):
\]

%dsp-level:0
\[
\qquad \partial_{\mathrm{col}}\ne\varnothing\implies\mathcal C=0\implies\text{the colour block vanishes}\ [\mathrm{D2.14};\ \mathrm{FT.19}];
\]

%dsp-level:0
\[
\textbf{there is no}\ \delta\text{-sealed member of a}\ 3\text{-junction's ensemble to cancel}: \mathcal E_\delta(\mathcal X_{N=3})=\varnothing;
\]

%dsp-level:0
\[
(2)\quad\textbf{the protection is unique to}\ N=3:
\]

%dsp-level:0
\[
\qquad g_\delta(N)=(1,0,3,10,15,105)\ \text{for}\ N=2,\dots,7\ [\text{machine}]:
\]

%dsp-level:0
\[
\qquad g_\delta(N)>0\ \ \forall N\ge4;
\]

%dsp-level:0
\[
g_\varepsilon(N)\ne0\iff3\mid N:\qquad g_\varepsilon(3)=1,\quad g_\varepsilon(6)=10,\quad g_\varepsilon(N)=0\ \text{at}\ N\in\{2,4,5,7\}\ [\text{machine}];
\]

%dsp-level:0
\[
\text{at}\ N
\]

%dsp-level:0
\[
= 6\ \text{the ensemble genuinely mixes:}\ \tfrac{15}{25}
\]

%dsp-level:0
\[
= \tfrac35\ \text{of the groupings contain a}\ \delta,
\]

%dsp-level:0
\[
\ \tfrac{10}{25}
\]

%dsp-level:0
\[
= \tfrac25\ \text{do not}\ [\text{machine}];
\]

%dsp-level:0
\[
\qquad \text{at}\ N=4\ \text{every grouping does}\ (g_\delta=g=3);
\]

%dsp-level:0
\[
(3)\quad\textbf{the charge classes overlap, so charge does not separate the sealings:}
\]

%dsp-level:0
\[
\qquad Q_{N=3}\in\{-3,\dots,+3\},\quad W_3=(1,6,15,20,15,6,1),\quad\sum W_3=64=4^3\ [\mathrm{Def\text{-}BM.8}(6);\ \text{machine}];
\]

%dsp-level:0
\[
Q_{N=2}\in\{-2,\dots,+2\},\quad W_2=(1,4,6,4,1),\quad\sum W_2=16=4^2\ [\text{machine}];
\]

%dsp-level:0
\[
\text{the overlap is}\ \{-2,-1,0,+1,+2\}:
\]

%dsp-level:0
\[
\qquad \textbf{at}\ Q
\]

%dsp-level:0
\[
= 0\ \text{and}\ Q
\]

%dsp-level:0
\[
= +1\ \text{the two sealings share a charge class}\ —\ \text{the neutron and proton classes are not separated from mesonic ones by charge};
\]

%dsp-level:0
\[
\text{only}\ |Q|=3\ \text{is baryon-exclusive}\ [\text{machine}];
\]

%dsp-level:0
\[
\qquad W_3(Q)
\]

%dsp-level:0
\[
\ne W_2(Q)\ \text{on the whole overlap}\ [\text{machine}]:\ \text{the probability spaces differ even where the classes coincide};
\]

%dsp-level:0
\[
(4)\quad\textbf{the phase class separates only partially:}
\]

%dsp-level:0
\[
\qquad \text{baryon-type}\ \Theta_{\mathrm{int}}=2\Theta_\triangle=4k_\triangle\in4\mathbb Z\ [\text{Census FT.49}(3);\ \mathrm{FT.22}];
\]

%dsp-level:0
\[
\qquad \text{per materialized closed chain}\ \Theta(c)\in2\mathbb Z\ [\mathrm{FT.22}];
\]

%dsp-level:0
\[
\qquad \Theta_{\mathrm{int}}=2k(\Gamma)\in2\mathbb Z\ [\mathrm{FT.34}(3);\ \text{Census FT.35}(3)]:
\]

%dsp-level:0
\[
4\mathbb Z\subsetneq2\mathbb Z\ [\text{machine}]:
\]

%dsp-level:0
\[
\qquad \Theta_{\mathrm{int}}
\]

%dsp-level:0
\[
= 2\ \text{is meson-admissible and baryon-excluded, while}\ \Theta_{\mathrm{int}}
\]

%dsp-level:0
\[
\in 4\mathbb Z\ \text{is admissible to both}\ —\ \textbf{the phase label is a partial separator};
\]

%dsp-level:0
\[
(5)\quad\textbf{the representation content does not separate them either:}
\]

%dsp-level:0
\[
\qquad \dim\bigl((\mathbb C^3)^{\otimes3}\bigr)^{\mathfrak{sl}(3)}=27-26=1,
\]

%dsp-level:0
\[
\qquad \dim\bigl(\mathbb C^3\otimes\bar{\mathbb C}^3\bigr)^{\mathfrak{sl}(3)}=9-8=1
\]

%dsp-level:0
\[
\qquad [\mathrm{BM.1}(3);\ \text{machine: exact rank over}\ \mathbb Q];
\]

%dsp-level:0
\[
\text{both contents carry exactly one singlet line}:
\]

%dsp-level:0
\[
\qquad \textbf{the separator is the attachment count}\ N,
\]

%dsp-level:0
\[
\ \text{carried in the boundary data}\ [\mathrm{Def\text{-}BM.9}(4)],
\]

%dsp-level:0
\[
\ \neg\text{charge},
\]

%dsp-level:0
\[
\ \neg\text{phase alone},
\]

%dsp-level:0
\[
\ \neg\text{representation content};
\]

%dsp-level:0
\[
(6)\quad\textbf{the verdict on the extended ensemble:}
\]

%dsp-level:0
\[
\qquad \mathcal E_{\mathrm{ext}}(\mathcal X_{N=3})=\mathcal E_\varepsilon(\mathcal X_{N=3})\sqcup\varnothing=\mathcal E_\varepsilon(\mathcal X_{N=3})
\]

%dsp-level:0
\[
\qquad [\text{machine: the grammar census}]:
\]

%dsp-level:0
\[
\textbf{the baryon's flicker ensemble is not extended by mesonic members}\ —\ \text{not because they cancel, but because a}\ 3\text{-junction admits no}\ \delta\text{-sealing at all}.
\]

\textbf{Rmk (what would extend it).} Clause (2) is the operative scope: at $N\ge4$ the union is nontrivial and the extended ensemble genuinely mixes sealings. Any composite whose support carries four or more colour attachments — the Census FT.50-type junction gluings [BM.2(3)] — is in that regime, and this tag's protection does not apply to it.

\textbf{Req.} D2.14, D3.9 (Parts 1–2); CT3.17, CT3.17a(2), CT3.18 (Part 2); EM.9 (Part 2); P2 (Part 1); Def-BM.8(6), Def-BM.9 (this part); FT.19, FT.22, FT.34(3), FT.41(b) (Part 3); BM.1(3), BM.2(1), BM.2(3) (this part); Census FT.35(3), FT.49(3) (row-wise); exhaustive set-partition enumeration, exact rational rank over $\mathbb Q$, generating-function convolution [math; machine: \texttt{checks/\allowbreak{}verify\_\allowbreak{}H\_\allowbreak{}mesonic\_\allowbreak{}flicker\_\allowbreak{}test.py} §A–§B].

\textbf{Pf.}

%dsp-level:0
\[
(1):\ \text{partitions of a}\ 3\text{-set into blocks of size}\ \in\{2,3\}:\ \text{the only one is}\ \{1,2,3\}\ [\text{machine}];
\]

%dsp-level:0
\[
\qquad \text{a}\ 2\text{-block leaves}\ 3-2=1\ \text{unpaired}\implies\partial_{\mathrm{col}}\ne\varnothing\ [\mathrm{D2.14}]\implies\mathcal C=0\implies\text{block}\equiv0\ [\mathrm{FT.19}]
\]

%dsp-level:0
\[
(2):\ \text{the censuses}\ g_\delta,\ g_\varepsilon\ \text{by exhaustive enumeration for}\ N=2..7\ [\text{machine}];
\]

%dsp-level:0
\[
\qquad g_\varepsilon(N)\ne0\implies N=3m\ [\text{a partition into}\ 3\text{-blocks requires}\ 3\mid N];
\]

%dsp-level:0
\[
\qquad \tfrac{15}{25}=\tfrac35,\ \tfrac{10}{25}=\tfrac25\ [\text{machine}]
\]

%dsp-level:0
\[
(3):\ (2+z+z^{-1})^3\ \text{and}\ (2+z+z^{-1})^2\ \text{convolved}\ [\mathrm{Def\text{-}BM.8}(6);\ \text{machine}];
\]

%dsp-level:0
\[
\qquad \text{the overlap and the exclusivity of}\ |Q|=3\ [\text{machine: set arithmetic}]
\]

%dsp-level:0
\[
(4):\ \text{Census FT.49}(3)\text{'s}\ 4k_\triangle\ \text{row re-cited};
\]

%dsp-level:0
\[
\qquad \mathrm{FT.22}\text{'s}\ \Theta(c)\in2\mathbb Z\ \text{and}\ \mathrm{FT.34}(3)\text{'s additivity re-cited};
\]

%dsp-level:0
\[
\qquad 4\mathbb Z\subsetneq2\mathbb Z\ \text{with}\ 2\in2\mathbb Z\setminus4\mathbb Z\ [\text{machine}]
\]

%dsp-level:0
\[
(5):\ \text{the common null space of the}\ 8\ \text{traceless generators' actions}: 27-26
\]

%dsp-level:0
\[
= 1\ \text{and}\ 9-8
\]

%dsp-level:0
\[
= 1\ [\text{machine: exact rank over}\ \mathbb Q];
\]

%dsp-level:0
\[
\qquad \mathrm{BM.1}(3)\ \text{re-derived}
\]

%dsp-level:0
\[
(6):\ (1)\text{'s}\ \mathcal E_\delta=\varnothing\ \text{substituted into}\ \mathrm{Def\text{-}BM.9}(1)\text{'s union}.\ \blacksquare
\]

\textbf{Phys.} Three of the four boundary labels fail to tell a baryon from a meson. Their charge classes overlap on everything except $|Q|=3$, so the neutron's neutral class and the proton's unit class are shared with mesonic configurations; their phase classes differ only in half the cases, since $4\mathbb Z$ sits strictly inside $2\mathbb Z$; and their colour representation content is identical in the only respect that matters, one singlet line each. What does separate them is the attachment count, and it separates them absolutely: three attachments admit exactly one contraction grouping, the epsilon triple, and a delta sub-pairing would leave one attachment dangling, which the all-or-nothing law turns into a vanished colour block. So the baryon's flicker ensemble has no mesonic members to cancel. That protection is not general — it is a property of the number three, and from four attachments upward the groupings genuinely mix, fifteen of twenty-five at six.

\end{theorem}

\begin{theorem}[{\texorpdfstring{BM.32 — The mesonic contribution: zero by filtration, and the amplitude layer is orthogonal}{BM.32 - The mesonic contribution: zero by filtration, and the amplitude layer is orthogonal}}]\label{tag:BM-32}
\textbf{Thm.}

%dsp-level:0
\[
(1)\quad\textbf{the contribution vanishes, and the mechanism is the projector:}
\]

%dsp-level:0
\[
\qquad Z^{(n)}_\delta:=\sum_{R\in\mathcal E_\delta(\mathcal X_{N=3})}\sum_\sigma w(R,\sigma)=\sum_{R\in\varnothing}(\cdots)=0
\]

%dsp-level:0
\[
\qquad [\mathrm{BM.31}(1),(6);\ \text{machine}];
\]

%dsp-level:0
\[
Z^{(n)}_{\mathrm{ext}}=Z^{(n)}_\varepsilon+Z^{(n)}_\delta=Z^{(n)}_\varepsilon\qquad[\text{machine}];
\]

%dsp-level:0
\[
\textbf{the zero is a filtration, not a cancellation:}\ \ \neg\exists\ \text{a pair of mesonic terms whose weights sum to zero}\ —\ \neg\exists\ \text{mesonic terms};
\]

%dsp-level:0
\[
(2)\quad\textbf{the projector form of the same statement:}
\]

%dsp-level:0
\[
\qquad \mathcal C(\Sigma;C)\in\{0,1\},\quad\mathcal C^2=\mathcal C\ [\mathrm{D2.14};\ \text{machine}]:
\]

%dsp-level:0
\[
\qquad \Pi_{\mathcal C}\ \text{is an idempotent scalar on the colour slot space};
\]

%dsp-level:0
\[
\text{a}\ \delta\text{-sealed member of a}\ 3\text{-support has}\ \mathcal C=0\implies\Pi_{\mathcal C}=0\implies F^\Sigma_{\mathrm{strong}}=0\ [\mathrm{FT.19}]:
\]

%dsp-level:0
\[
\qquad \text{its weight is multiplied by zero before any sum is formed};
\]

%dsp-level:0
\[
(3)\quad\textbf{both mass routes are untouched:}
\]

%dsp-level:0
\[
\qquad \text{adding zero-weight members leaves a normalized average invariant}\ \bigl(\tfrac{\sum w m}{\sum w}\bigr)\ \textbf{and}\ \text{a total invariant}\ \bigl(\sum wm\bigr)
\]

%dsp-level:0
\[
\qquad [\text{machine: exact on a rational witness ensemble}]:
\]

%dsp-level:0
\[
\langle m\rangle_{\mathrm{ext}}=\langle m\rangle_\varepsilon\ \ \text{on route}\ (\mathrm i),\qquad T_{\mathrm{ext}}=T_\varepsilon\ \ \text{on route}\ (\mathrm{ii})\qquad[\mathrm{BM.28}(4);\ \mathrm{BM.20}(3);\ \text{machine}];
\]

%dsp-level:0
\[
\therefore\ \mathrm{BM.20}\text{'s curvature-flux mass and BM.28's three ratios are}\ \textbf{unchanged}\ \text{by the extension};
\]

%dsp-level:0
\[
(4)\quad\textbf{the amplitude-layer mechanisms are orthogonal to this question:}
\]

%dsp-level:0
\[
\qquad \mathrm{FT.67}(3\mathrm a)\ \text{and}\ \mathrm{FT.106}(2)\ \text{kill}\ \textbf{pure directed cycles}\ (\text{any length, any cell});
\]

%dsp-level:0
\[
\qquad \mathrm{FT.71}\ \text{cancels every odd cycle and every undissolved even cycle};
\]

%dsp-level:0
\[
\qquad \mathrm{FT.69}(3)(\mathrm{iv}),(\mathrm v)\ \text{exclude odd cycles and}\ 4\text{-cycles from covers}:
\]

%dsp-level:0
\[
\text{these are statements about cycles in}\ \Gamma,\ \neg\text{about colour sealings of a support}\ [\mathrm{Def\text{-}BM.9}(2)\ \mathrm{Rmk}];
\]

%dsp-level:0
\[
\qquad \text{they neither produce nor require}\ Z_\delta=0;
\]

%dsp-level:0
\[
(5)\quad\textbf{where mesonic content does survive:}
\]

%dsp-level:0
\[
\qquad Z^{\mathrm{amp}}(q,1)=802\ (q=\pm1),\ 4796\ (q=0):
\]

%dsp-level:0
\[
\qquad \textbf{dissolved digons inside covers}\ [\mathrm{FT.106}(3);\ \text{machine}];
\]

%dsp-level:0
\[
Z^{\mathrm{amp}}(q,S)=0\ \ \forall q,\ S\in\{2,3\}\ [\mathrm{FT.106}(3\mathrm c);\ \text{machine}]:
\]

%dsp-level:0
\[
\qquad \text{a standalone digon}\ —\ \text{a pure directed cycle}\ —\ \text{dies at the amplitude layer};
\]

%dsp-level:0
\[
Z^{\mathrm{amp}}(q,S)=Z^{\mathrm{amp}}(-q,S)\ [\mathrm{FT.106}(3\mathrm d);\ \text{machine}];
\]

%dsp-level:0
\[
\qquad Z^{\mathrm{norm}}(0,S)
\]

%dsp-level:0
\[
= 2Z^{\mathrm{norm}}(\pm1,S)\ \text{at every}\ S\ [\text{machine}]:\ \text{the norm layer keeps its}\ 2{:}1\ \text{ratio where the amplitude layer dies};
\]

%dsp-level:0
\[
(6)\quad\textbf{the two readings reconciled:}
\]

%dsp-level:0
\[
\qquad \text{a meson as a}\ \textbf{sealed composite}\ \text{is excluded from a}\ 3\text{-junction's ensemble by the grammar}\ [(1)];
\]

%dsp-level:0
\[
\qquad \text{mesonic}\ \textbf{chain content}\ \text{enters the amplitude layer only as dissolved digons inside covers}\ [(5)]:
\]

%dsp-level:0
\[
\text{the two statements concern different objects}\ —\ \text{a support's sealing and a cover's cycle structure}\ —\ \text{and neither contradicts the other}.
\]

\textbf{Rmk (no protection claim beyond N = 3).} Clause (3)'s invariance is exact because the added weight is exactly zero. At $N\ge4$ [BM.31(2)] the added weight is nonzero, the average genuinely moves, and no clause of this tag applies.

\textbf{Req.} D2.14 (Parts 1–2); CT2.5, CT2.8, CT3.18 (Part 2); Def-H.32, FT.19, FT.67(2), FT.67(3), FT.69(3), FT.71, FT.106(2), FT.106(3) (Part 3); Def-BM.9, BM.20(3), BM.28(4), BM.31(1), BM.31(2), BM.31(6) (this part); exact rational ensemble algebra, exact integer census arithmetic [math; machine: \texttt{checks/\allowbreak{}verify\_\allowbreak{}H\_\allowbreak{}mesonic\_\allowbreak{}flicker\_\allowbreak{}test.py} §C–§D].

\textbf{Pf.}

%dsp-level:0
\[
(1):\ \mathcal E_\delta=\varnothing\ [\mathrm{BM.31}(1)]\implies\text{the sum is empty}\implies0\ [\text{machine}];\qquad Z_{\mathrm{ext}}=Z_\varepsilon+0\ [\text{machine}]
\]

%dsp-level:0
\[
(2):\ \mathrm{D2.14}\text{'s}\ \mathcal C\ \text{definition re-cited};
\]

%dsp-level:0
\[
\qquad \mathcal C^2=\mathcal C\ \text{at}\ \mathcal C\in\{0,1\}\ [\text{machine}];
\]

%dsp-level:0
\[
\qquad \mathrm{FT.19}\text{'s all-or-nothing row re-cited}
\]

%dsp-level:0
\[
(3):\ \text{adding terms with}\ w
\]

%dsp-level:0
\[
= 0\ \text{changes neither}\ \sum wm\ \text{nor}\ \sum w\implies\text{neither the ratio nor the total}\ [\text{machine: exact Fractions on a}\ 5\text{-member witness ensemble}]
\]

%dsp-level:0
\[
(4):\ \mathrm{FT.67}(3\mathrm a),\ \mathrm{FT.71},\ \mathrm{FT.69}(3)\ \text{rows re-cited};
\]

%dsp-level:0
\[
\qquad \text{their objects are cycles of}\ \Gamma\ [\mathrm{Def\text{-}H.32}(1)\text{'s topology space}],
\]

%dsp-level:0
\[
\ \text{while}\ \mathcal C\ \text{acts on colour attachments}\ [\mathrm{D2.14}]
\]

%dsp-level:0
\[
(5):\ \mathrm{FT.106}(3)\text{'s census integers re-cited and their relations verified exactly}\ [\text{machine}:\ 802,\ 4796,\ \text{the zeros at}\ S\in\{2,3\},\ \text{the}\ \omega\text{-symmetry, the}\ 2{:}1\ \text{norm ratio at all four windings}]
\]

%dsp-level:0
\[
(6):\ (1)\ \text{and}\ (5)\ \text{conjoined with}\ \mathrm{Def\text{-}BM.9}(2)\text{'s locus distinction}.\ \blacksquare
\]

\textbf{Phys.} The answer to the test is zero, and the reason is more interesting than the number. Nothing cancels: there is nothing to cancel. A three-attachment support admits exactly one contraction, so the mesonic subset of its ensemble is empty before any weight is written down, and the colour projector would have annihilated those terms individually anyway — one dangling attachment and the whole block goes. Because the added weight is exactly zero, both mass functionals stay put: the normalized average and the total deficit are each invariant, so the curvature-flux mass and the four-thirds ratio of the previous tier survive the extension untouched. The amplitude-layer cancellations the corpus already proved are about a different object entirely — directed cycles in the skeleton — and they cut the other way: a standalone digon dies there, while dissolved digons inside covers are precisely what survives, with exact census weights. Mesons as sealed composites and mesonic chain content in covers are two different things, and the tier keeps them apart.

\end{theorem}

\medskip\noindent\textit{\texorpdfstring{BM.33 — The status of the mesonic flicker test}{BM.33 - The status of the mesonic flicker test}.}\ Structural-constraints record: see Appendix B, Section B.1.19 of the Addendum.\par\medskip

\phantomsection
\section[Tier BM.10 — The Kinetic Energy from the Metric and the…]{\texorpdfstring{Tier BM.10 — The Kinetic Energy from the Metric and the Vacuum-Balanced Entropy: the Dispersion Read of a Sealed Composite, the Deficit Reframing of the Rest Mass, the Force-Structure Ladder, and the Route-Artifact Adjudication}{Tier BM.10 - The Kinetic Energy from the Metric and the Vacuum-Balanced Entropy: the Dispersion Read of a Sealed Composite, the Deficit Reframing of the Rest Mass, the Force-Structure Ladder, and the Route-Artifact Adjudication}}

\begin{definition}[{\texorpdfstring{Def-BM.10 — The metric-energy read of a sealed composite (the dispersion route)}{Def-BM.10 - The metric-energy read of a sealed composite (the dispersion route)}}]\label{tag:Def-BM-10}
\textbf{Def.} A read constructed on the chart side of a sealed composite from CT1.12's dictionary; no new axiom, no new datum, no new parameter. Local notation (internal bookkeeping): $\gamma_{ij}:=g_{ij}$ (the spatial block), $\gamma^{ij}:=$ the inverse matrix of the $3\times3$ block ($\gamma^{ik}g_{kj}=\delta^i{}_j$), $g^{\mu\nu}:=$ the inverse of the full $4\times4$ metric.

%dsp-level:0
\[
(1)\quad\textbf{locus and hypotheses:}
\]

%dsp-level:0
\[
\qquad \text{the read lives on the emergent chart}\ M\ [\mathrm{CT2.6}],
\]

%dsp-level:0
\[
\ \text{inheriting hypotheses (i)--(iii) at the point of use}\ [\text{A §0.2a}];
\]

%dsp-level:0
\[
\text{the}\ 16\text{-dimensional pair register carries no canonical pair}\ [\mathrm{BM.10}(1)];
\]

%dsp-level:0
\[
\qquad \text{the register-side precursors are the discrete Weyl pairs}\ [\mathrm{BM.10}(2)]\ \text{and the width law}\ [\mathrm{BM.10}(6)];
\]

%dsp-level:0
\[
(2)\quad\textbf{the metric data and the derived ADM split:}
\]

%dsp-level:0
\[
\qquad \text{for the coarse-grained sector coefficients}\ (a_{00},q_{ij},w_i)\ [\mathrm{D2.1}]:
\]

%dsp-level:0
\[
\Omega=e^{\alpha_\Omega a_{00}},\qquad g_{00}=-\Omega^2,\qquad g_{0i}=\Omega^2\varkappa\beta w_i,\qquad g_{ij}=\Omega^2(\delta_{ij}+\varkappa q_{ij})\qquad[\mathrm{CT1.12};\ \mathrm{FT.157}(1)];
\]

%dsp-level:0
\[
\textbf{lapse:}\quad N^2:=\Omega^2+\gamma^{ij}N_iN_j=-\frac{\det g}{\det\gamma}>0;
\]

%dsp-level:0
\[
\qquad \textbf{shift covector:}\quad N_i:=g_{0i}=\Omega^2\varkappa\beta w_i;
\]

%dsp-level:0
\[
\qquad \textbf{shift vector:}\quad N^i:=\gamma^{ij}N_j;
\]

%dsp-level:0
\[
\textbf{the line element (an exact identity):}
\]

%dsp-level:0
\[
\qquad ds^2=-N^2d\tau^2+g_{ij}\bigl(dx^i+N^id\tau\bigr)\bigl(dx^j+N^jd\tau\bigr)=g_{\mu\nu}dx^\mu dx^\nu
\]

%dsp-level:0
\[
\qquad [\text{machine: Schur, denominators cleared}];
\]

%dsp-level:0
\[
(3)\quad\textbf{the momenta:}
\]

%dsp-level:0
\[
\qquad \textbf{the unit-mass coordinate-time read:}\quad\pi_i:=g_{ij}\Bigl(\frac{dx^j}{dx^0}+N^j\Bigr)
\]

%dsp-level:0
\[
\qquad [x^0=\tau_n\ \text{the chart time};\ \mathrm{D1.8};\ \mathrm{FT.159}];
\]

%dsp-level:0
\[
\textbf{the physical spatial momentum:}\qquad p_i:=m\,u^0\,\pi_i,\qquad u^\mu:=\frac{dx^\mu}{d\lambda},\qquad g_{\mu\nu}u^\mu u^\nu=-1;
\]

%dsp-level:0
\[
\text{the shift contribution is}\ g_{ij}N^j=N_i=g_{0i}\ [\text{machine: exact polynomial identity}]:
\]

%dsp-level:0
\[
\qquad \text{the metric's time--space block enters the momentum exactly}\ [\mathrm{FT.157}(1)];
\]

%dsp-level:0
\[
(4)\quad\textbf{the rest input:}
\]

%dsp-level:0
\[
\qquad m:=m_{\mathrm{hid}}=\frac{E_{\mathrm{int}}\ln2}{\mathfrak h}
\]

%dsp-level:0
\[
\qquad [\mathrm{FT.44}(4)\ \text{(Census)}];
\]

%dsp-level:0
\[
\qquad \text{the ensemble route:}\quad m\ \mapsto\ \langle m\rangle\ [\mathrm{BM.20}(3)];
\]

%dsp-level:0
\[
\text{the vacuum-deficit reframing of the rest read is}\ \mathrm{BM.35}\ —\ \text{a marked forward pointer at the citing site, imported in no}\ \textbf{Req}\ \text{line};
\]

%dsp-level:0
\[
(5)\quad\textbf{the mass shell and the energy:}\qquad g^{\mu\nu}p_\mu p_\nu=-m^2\qquad[\text{signature}\ (-,+,+,+);\ \mathrm{CT1.12}];
\]

%dsp-level:0
\[
E:=-p_0\quad\text{is}\ \mathrm{BM.8}(2)\text{'s conditional energy read under its declared stationarity scope (theorem-local, inherited at the point of use)};
\]

%dsp-level:0
\[
\textbf{the Eulerian form:}
\]

%dsp-level:0
\[
\qquad E=N\sqrt{m^2+\gamma^{ij}p_ip_j}-N^ip_i;
\]

%dsp-level:0
\[
\qquad \textbf{the kinetic-energy read:}
\]

%dsp-level:0
\[
\qquad \mathrm{KE}:=E-N\,m\ \ge0\ \ \text{where}\ g_{ij}\succ0\ [\mathrm{CT1.12};\ \text{machine}];
\]

%dsp-level:0
\[
(6)\quad\textbf{status and scope:}
\]

%dsp-level:0
\[
\qquad \textbf{constructed}\ —\ \text{an abbreviation of derived objects:}\ (2)\text{'s split is linear algebra on}\ \mathrm{CT1.12}\text{'s dictionary};
\]

%dsp-level:0
\[
\qquad (3)\text{--}(5)\ \text{are pointwise identities}\ [\text{machine}];
\]

%dsp-level:0
\[
\neg\text{a new parameter}\ [\text{§Cal-BM stays empty}];
\]

%dsp-level:0
\[
\qquad \neg\text{a declaration};
\]

%dsp-level:0
\[
\qquad \text{the names (lapse, shift, kinetic energy) are internal bookkeeping}\ [\text{A §0.2}];
\]

%dsp-level:0
\[
\neg\text{constructed: the back-reaction of the composite's motion on its own sector coefficients}\ [\mathrm{BM.15}\text{'s read is the rest-frame one}]\ \text{and any inter-sector dynamics}\ [\mathrm{BM.29}(5)\ \text{stands}].
\]

\textbf{Rmk (the unit-mass normalization).} The leading kinetic read is $\mathrm{KE}=\tfrac12N\,\gamma^{ij}p_ip_j/m+O(p^4)$: the factor $1/m$ is dimensionally forced, since $[\gamma^{ij}p_ip_j]=\mathfrak q^{+1}=\text{energy}^2$ while $[\mathrm{KE}]=\mathfrak q^{+1/2}=\text{energy}$ [CT2.1; machine: exponent arithmetic]. The unit-mass unit-lapse normalization ($m=1$, $N=1$, $w_i=0$) gives $\mathrm{KE}=\tfrac12\gamma^{ij}p_ip_j$. And $\gamma^{ij}\ne g^{ij}$ where the shift is nonzero: $g^{ij}=\gamma^{ij}-N^iN^j/N^2$ [machine] — the contraction that inverts the spatial momentum is $\gamma^{ij}$; the two coincide exactly in the shift-free gauge.

\textbf{Req.} D1.8, D2.1, D2.9 (Parts 1–2); CT1.12, CT2.1, CT2.6 (Part 2); P1, P5 (Parts 1–2); FT.157(1), FT.159 (Part 3); FT.44(4) (Census); BM.8(2), BM.10(1), BM.10(2), BM.10(6), BM.20(3) (this part); A §0.2a (the continuum-limit hypotheses at the point of use); Schur decomposition, adjugate calculus [math; machine: \texttt{checks/\allowbreak{}verify\_\allowbreak{}H\_\allowbreak{}kinetic\_\allowbreak{}energy\_\allowbreak{}metric.py} §A].

\textbf{Pf.}

%dsp-level:0
\[
\text{Definition block: the constitutive clauses are the displayed identities of}\ (2)\text{--}(5)\ [\text{machine: exact polynomial identities over}\ \mathbb Q[\Omega^2,\varkappa,\beta,w_i,q_{ij},p_\mu,m]\ \text{with denominators cleared by}\ \det\gamma];
\]

%dsp-level:0
\[
\neg\exists\ \text{a non-derived clause: the ADM split is forced by the Schur decomposition of}\ g_{\mu\nu}\ [\text{math}];
\]

%dsp-level:0
\[
\qquad (1),(6)\ \text{are scope records}.\ \blacksquare
\]

\textbf{Phys.} The energy of a moving sealed composite is read off the metric the corpus already derived, not added to it. The dictionary's ten sector coefficients split into lapse, shift, and spatial metric by linear algebra alone — the Schur decomposition of the very matrix CT1.12 builds — and the same three shift coefficients that the weak block carries enter the spatial momentum as its drag term. The mass-shell constraint of the Lorentzian metric converts these data into one exact dispersion identity, whose rest input is the informational mass the calibration tiers locked. Two normalizations — unit mass and coordinate time — are fixed as bookkeeping, and the read is chart-side: the register itself carries no canonical pair, and the definition keeps the obstruction exactly where Tier BM.4 placed it.

\end{definition}

\begin{theorem}[{\texorpdfstring{BM.34 — The kinetic energy from the emergent metric: the ADM split, the mass shell, the dispersion relation, and the zero-point bound}{BM.34 - The kinetic energy from the emergent metric: the ADM split, the mass shell, the dispersion relation, and the zero-point bound}}]\label{tag:BM-34}
\textbf{Thm.}

%dsp-level:0
\[
(1)\quad\textbf{the split is derived:}
\]

%dsp-level:0
\[
\qquad \gamma_{ij}=\Omega^2(\delta_{ij}+\varkappa q_{ij}),\quad N_i=\Omega^2\varkappa\beta w_i,\quad N^i=\gamma^{ij}N_j,\quad N^2=\Omega^2+\gamma^{ij}N_iN_j;
\]

%dsp-level:0
\[
\boxed{\ N^2\det\gamma=-\det g\ }\qquad[\text{Schur; machine: exact polynomial identity over}\ \mathbb Q[\Omega^2,\varkappa,\beta,w_i,q_{ij}]];
\]

%dsp-level:0
\[
\textbf{the inverse metric in ADM form:}
\]

%dsp-level:0
\[
\qquad g^{00}=-\frac1{N^2},
\]

%dsp-level:0
\[
\qquad g^{0i}=\frac{N^i}{N^2},
\]

%dsp-level:0
\[
\qquad g^{ij}=\gamma^{ij}-\frac{N^iN^j}{N^2}
\]

%dsp-level:0
\[
\qquad [\text{machine: cleared adjugate identities}\ +\ 3\ \text{rational witnesses, exact Fractions}];
\]

%dsp-level:0
\[
(2)\quad\textbf{the line element and the vacuum:}
\]

%dsp-level:0
\[
\qquad ds^2=-N^2d\tau^2+g_{ij}\bigl(dx^i+N^id\tau\bigr)\bigl(dx^j+N^jd\tau\bigr)\ \ \text{identically}\ [\text{machine}];
\]

%dsp-level:0
\[
(a_{00},q_{ij},w_i)=0\implies(\Omega,N,N^i,g_{ij})=(1,1,0,\delta_{ij}):
\]

%dsp-level:0
\[
\qquad \textbf{the vacuum's emergent metric is the flat metric}\ [\mathrm{CT1.12};\ \text{machine: exact substitution}];
\]

%dsp-level:0
\[
(3)\quad\textbf{the momentum and its carrier:}
\]

%dsp-level:0
\[
\qquad p_i=m\,u^0\,g_{ij}\Bigl(\frac{dx^j}{dx^0}+N^j\Bigr)\ [\mathrm{Def\text{-}BM.10}(3)];
\]

%dsp-level:0
\[
\qquad g_{ij}N^j=N_i=g_{0i}=\Omega^2\varkappa\beta w_i\ [\text{machine}];
\]

%dsp-level:0
\[
\text{the drag term and the metric's time--space block are the same three coefficients}\ w_i\ —\ \text{the weak aspect, one block, two derived reads}\ [\mathrm{FT.157}(1);\ \mathrm{CT1.12a}(2)];
\]

%dsp-level:0
\[
\textbf{the T-parity of the carrier:}
\]

%dsp-level:0
\[
\quad w\mapsto-w\implies N_i\mapsto-N_i,
\]

%dsp-level:0
\[
\quad N^i\mapsto-N^i,
\]

%dsp-level:0
\[
\quad (N,\gamma_{ij})\ \text{fixed}\ [\mathrm{FT.156}(3);\ \mathrm{FT.159}(3);\ \text{machine}]:
\]

%dsp-level:0
\[
\qquad p_i\ \text{is T-odd, matching its carrier}\ [\mathrm{BM.8}(1)];
\]

%dsp-level:0
\[
(4)\quad\textbf{the mass shell}\iff\textbf{the dispersion relation:}\qquad g^{\mu\nu}p_\mu p_\nu=-m^2\iff\bigl(p_0-N^ip_i\bigr)^2=N^2\bigl(m^2+\gamma^{ij}p_ip_j\bigr)
\]

%dsp-level:0
\[
[\text{machine: the master polynomial identity, cleared by}\ (\det\gamma)^2];
\]

%dsp-level:0
\[
\qquad E:=-p_0,\ \text{the positive-energy branch}:
\]

%dsp-level:0
\[
\qquad \boxed{\ E=N\sqrt{m^2+p^2}-N^ip_i\ },
\]

%dsp-level:0
\[
\qquad p^2:=\gamma^{ij}p_ip_j;
\]

%dsp-level:0
\[
\textbf{the gauges:}\qquad w=0:\quad E^2=N^2(m^2+p^2);\qquad\Omega=1\wedge w=0:\quad\boxed{\ E^2=m^2+p^2\ }\qquad[c=1,\ \mathrm{D2.9}];
\]

%dsp-level:0
\[
\textbf{the dispersion is T-even:}\quad\text{the mass shell is invariant under the simultaneous}\ (N^i,p_i)
\]

%dsp-level:0
\[
\mapsto (-N^i,-p_i)\ [\text{machine: substitution in the master identity}];
\]

%dsp-level:0
\[
\qquad \text{the signature is invariant}\ [\mathrm{FT.159}(3)\ \text{re-cited}];
\]

%dsp-level:0
\[
(5)\quad\textbf{the kinetic energy and the non-relativistic read:}
\]

%dsp-level:0
\[
\qquad \mathrm{KE}:=E-Nm=N\bigl(\sqrt{m^2+p^2}-m\bigr)\ \ge0,
\]

%dsp-level:0
\[
\qquad \mathrm{KE}=0\iff p_i=0;
\]

%dsp-level:0
\[
\text{on the perturbative regime}\quad g_{ij}\succ0\iff\forall a:\ 1+\varkappa\lambda_a(q)>0
\]

%dsp-level:0
\[
\qquad [\mathrm{CT1.12};\ \text{machine: exact rational Sylvester/LDL}^\top\ \text{witnesses}];
\]

%dsp-level:0
\[
\sqrt{1+x}=1+\frac x2-\frac{x^2}8+\frac{x^3}{16}-\frac{5x^4}{128}+\frac{7x^5}{256}+O(x^6)
\]

%dsp-level:0
\[
\qquad [\text{machine: exact rational coefficients;}\ (\cdot)^2\equiv1+x\ \bmod x^6];
\]

%dsp-level:0
\[
\boxed{\ \mathrm{KE}=N\Bigl(\frac{p^2}{2m}-\frac{p^4}{8m^3}+\frac{p^6}{16m^5}\Bigr)+O(p^8)\ },\qquad\text{the leading read:}\quad\mathrm{KE}=\frac{N\,\gamma^{ij}p_ip_j}{2m};
\]

%dsp-level:0
\[
\text{the unit normalization}\ (m{=}1,N{=}1,w{=}0):\quad\mathrm{KE}=\tfrac12\gamma^{ij}p_ip_j\ [\mathrm{Def\text{-}BM.10}\ \mathrm{Rmk}];
\]

%dsp-level:0
\[
\qquad [\mathrm{KE}]=\mathfrak q^{+1/2},\quad[\gamma^{ij}p_ip_j]=\mathfrak q^{+1}\ [\mathrm{CT2.1};\ \text{machine: exponent arithmetic}];
\]

%dsp-level:0
\[
(6)\quad\textbf{the zero-point bound:}
\]

%dsp-level:0
\[
\qquad \Delta x\,\Delta p\ge\tfrac12\ [\mathrm{BM.10}(3),(6);\ \hbar=1];
\]

%dsp-level:0
\[
\qquad \Delta x=\Delta x_{\mathrm B}:\quad\Delta p\ge\frac1{2\Delta x_{\mathrm B}}=\frac{m_{\mathrm{hid}}}2\ [\mathrm{BM.13}(4)];
\]

%dsp-level:0
\[
\text{the isotropic cube convention}\ [\mathrm{Def\text{-}BM.4}(3)],\ \text{the rest read}\ \langle p_i\rangle=0:
\]

%dsp-level:0
\[
\qquad \langle p^2\rangle\ \ge\ 3\Bigl(\frac{m_{\mathrm{hid}}}2\Bigr)^2=\frac34\,m_{\mathrm{hid}}^2
\]

%dsp-level:0
\[
\qquad [\text{machine: exact Fractions}];
\]

$\text{in the vacuum gauge of}\ (4)\ \bigl(E^2=m^2+\langle p^2\rangle\bigr),\ \text{at}\ m=m_{\mathrm{hid}}:\qquad\boxed{\ E^2\ \ge\ \frac74\,m_{\mathrm{hid}}^2\ },\qquad\mathrm{KE}_{\mathrm{NR}}\ \ge\ \frac{\langle p^2\rangle}{2m}\ \ge\ \frac38\,m_{\mathrm{hid}},\qquad E\ \ge\ \frac{\sqrt7}2\,m_{\mathrm{hid}}\ \ \text{exact in}\ \mathbb Q(\sqrt7)$

%dsp-level:0
\[
[\text{machine: the quadratic ring, minimal polynomial}\ x^2-7,\ \text{comparison by squaring}];
\]

%dsp-level:0
\[
\qquad m_{\mathrm{hid}}=E_{\mathrm{int}}\ln2/\mathfrak h\ \text{substituted}:
\]

%dsp-level:0
\[
\quad E^2\ge\frac74\,\frac{(E_{\mathrm{int}}\ln2)^2}{\mathfrak h^2}\ [\text{machine: one Laurent monomial, coefficient}\ \tfrac74];
\]

%dsp-level:0
\[
(7)\quad\textbf{the rest slice and the dust closure:}
\]

%dsp-level:0
\[
\qquad p_i=0:\quad E=Nm;
\]

%dsp-level:0
\[
\qquad \text{the vacuum gauge:}\quad E=m=m_{\mathrm{hid}}=\frac{E_{\mathrm{int}}\ln2}{\mathfrak h}\ [\mathrm{FT.44}(4)];
\]

%dsp-level:0
\[
RV=\varkappa^2m\ [\mathrm{BM.20}(1)],
\]

%dsp-level:0
\[
\qquad \rho=\frac mV,\quad\kappa^2\rho=\frac{2\pi\,\mathfrak h\,E_{\mathrm{int}}}V\ [\mathrm{BM.9}(5);\ \mathrm{BM.15}(2)]:
\]

%dsp-level:0
\[
\qquad \textbf{the informational, the curvature-flux and the dispersion rest reads agree on-shell}\ [\text{machine: Laurent identities}];
\]

%dsp-level:0
\[
\text{KE is the}\ p
\]

%dsp-level:0
\[
\ne 0\ \text{completion the rest-frame dust read does not carry}\ [\mathrm{BM.9}(4):\ T_{\mu\nu}=\mathrm{diag}(\rho,0,0,0)\ \text{at}\ p=0]:
\]

%dsp-level:0
\[
\qquad \text{the kinetic term is the whole composite's transport on the chart}\ [\mathrm{BM.8}(1)].
\]

\textbf{Rmk (chart-side discipline).} The reads of (4)–(6) are continuum-chart objects under CT2.6's hypotheses at the point of use; no canonical pair exists in the finite register [BM.10(1)] and no clause imports one. The pointwise mass-shell identity needs no stationarity; the conserved-charge read $E=-p_0$ inherits BM.8(2)'s theorem-local stationarity hypothesis at the point of use. The bound of (6) applies the Robertson inequality chart-side, at the width law's locus — Def-BM.4(5)'s envelope discipline: the zero-point read is a coarse-graining envelope, never a ledger distance.

\textbf{Req.} A3 (Part 1); D1.8, D2.1, D2.9 (Parts 1–2); CT1.12, CT1.12a(2), CT2.1, CT2.6 (Part 2); P1, P5 (Parts 1–2); FT.44(4) (Census); FT.156(3), FT.157(1), FT.159 (Part 3); Def-BM.4, Def-BM.10, BM.8(1), BM.8(2), BM.9(4), BM.9(5), BM.10(1), BM.10(2), BM.10(3), BM.10(6), BM.13(4), BM.15(2), BM.20(1), BM.20(3) (this part); A §0.2a (the hypotheses at the point of use); Schur decomposition, adjugate calculus, the Sylvester criterion, exact rational $\mathrm{LDL}^\top$ inertia, the binomial series over $\mathbb Q$, the quadratic ring $\mathbb Q(\sqrt7)$, Laurent monomials [math; machine: \texttt{checks/\allowbreak{}verify\_\allowbreak{}H\_\allowbreak{}kinetic\_\allowbreak{}energy\_\allowbreak{}metric.py} §A].

\textbf{Pf.}

%dsp-level:0
\[
(1):\ \det g=\det\gamma\cdot\bigl(g_{00}-g_{0i}\gamma^{ij}g_{0j}\bigr)\ [\text{Schur}];
\]

%dsp-level:0
\[
\qquad g_{00}=-\Omega^2,\ g_{0i}=N_i\implies-\det g=\bigl(\Omega^2+\gamma^{ij}N_iN_j\bigr)\det\gamma=N^2\det\gamma\ [\text{machine}]
\]

%dsp-level:0
\[
\text{the inverse rows in cleared adjugate form}\ (\operatorname{adj}\gamma:=\det\gamma\cdot\gamma^{-1},\ D:=\det\gamma):
\]

%dsp-level:0
\[
\quad (\operatorname{adj}g)_{00}=D;
\]

%dsp-level:0
\[
\quad (\operatorname{adj}g)_{0i}=-\bigl((\operatorname{adj}\gamma)^{ij}N_j\bigr);
\]

%dsp-level:0
\[
\quad D\,
\]

%dsp-level:0
\[
(\operatorname{adj}g)_{ij}
\]

%dsp-level:0
\[
= (\det g)\,
\]

%dsp-level:0
\[
(\operatorname{adj}\gamma)^{ij}+\bigl((\operatorname{adj}\gamma)^{ik}N_k\bigr)\bigl((\operatorname{adj}\gamma)^{j\ell}N_\ell\bigr)\ [\text{machine: three exact polynomial identities}]
\]

%dsp-level:0
\[
\text{the three rational witnesses}\ (\Omega^2,\varkappa,\beta,w_i,q_{ij}\ \text{rational,}\ q\ \text{off-diagonal}):\ \text{every row as exact Fraction equality, including}\ g^{\mu\alpha}g_{\alpha\nu}
\]

%dsp-level:0
\[
= \delta^\mu_\nu\ \text{and}\ \det g\qquad = -N^2\det\gamma
\]

%dsp-level:0
\[
(2):\ \text{expansion in}\ (d\tau,dx^i):\quad g_{00}=-N^2+N_iN^i,\quad g_{0i}=\gamma_{ij}N^j=N_i,\quad g_{ij}=\gamma_{ij}\ [\text{machine}];
\]

%dsp-level:0
\[
\qquad \Omega=e^{\alpha_\Omega\cdot0}=1,\ q=0,\ w=0\implies\text{the flat metric}\ [\text{machine: exact substitution}]
\]

%dsp-level:0
\[
(3):\ \gamma_{ij}N^j=N_i\ [\text{machine: cleared by}\ \det\gamma];
\]

%dsp-level:0
\[
\qquad w
\]

%dsp-level:0
\[
\mapsto -w:\ N_i
\]

%dsp-level:0
\[
= \Omega^2\varkappa\beta w_i
\]

%dsp-level:0
\[
\mapsto -N_i\ \text{and}\ (N,\gamma_{ij})\ \text{carry}\ w\ \text{only through}\ N_iN_j\ \text{and}\ N_i\ [\text{machine:}\ N^2,\gamma_{ij}\ \text{fixed}]
\]

%dsp-level:0
\[
(4):\ g^{\mu\nu}p_\mu p_\nu
\]

%dsp-level:0
\[
= -\frac1{N^2}\bigl(p_0-N^ip_i\bigr)^2+\gamma^{ij}p_ip_j\ [(1)\text{'s inverse rows}]\implies\bigl(p_0-N^ip_i\bigr)^2
\]

%dsp-level:0
\[
= N^2\bigl(m^2+\gamma^{ij}p_ip_j\bigr)\ \text{on shell}\ [\text{machine: cleared by}\ (\det\gamma)^2]
\]

%dsp-level:0
\[
\text{the branch: at}\ p_i=0,\ w=0:\ p_0^2=N^2m^2\implies p_0=\mp Nm;
\]

%dsp-level:0
\[
\qquad E=-p_0=+Nm\ \text{on the positive branch}\ [\text{machine: rational witnesses}];
\]

%dsp-level:0
\[
\qquad \text{the gauges: substitution}\ [\text{machine}];
\]

%dsp-level:0
\[
\qquad \text{T-evenness:}\ (N^i,p_i)\mapsto(-N^i,-p_i)\ \text{fixes}\ (p_0-N^ip_i)^2\ \text{and}\ \gamma^{ij}p_ip_j\ [\text{machine}]
\]

%dsp-level:0
\[
(5):\ \mathrm{KE}=N\bigl(\sqrt{m^2+p^2}-m\bigr)=Nm\bigl(\sqrt{1+p^2/m^2}-1\bigr)\ [\text{exact}];
\]

%dsp-level:0
\[
\qquad \text{the series}\ a=\bigl(1,\tfrac12,-\tfrac18,\tfrac1{16},-\tfrac5{128},\tfrac7{256}\bigr):
\]

%dsp-level:0
\[
\quad (a*a)_n=\delta_{n0}+\delta_{n1}\ \forall n\le5\ [\text{machine: Fractions}]
\]

%dsp-level:0
\[
\text{positivity:}\ g_{ij}\succ0\ [\mathrm{CT1.12}]\implies\gamma^{ij}\succ0
\]

%dsp-level:0
\[
\implies p^2
\]

%dsp-level:0
\[
\ge 0\ \text{with equality iff}\ p_i
\]

%dsp-level:0
\[
= 0\ [\text{machine:}\ \mathrm{LDL}^\top\ \text{on exact rational witnesses}];
\]

%dsp-level:0
\[
\qquad [\mathrm{KE}]=[\gamma^{ij}p_ip_j]-[m]=\mathfrak q^{1-1/2}=\mathfrak q^{+1/2}\ [\mathrm{CT2.1};\ \text{machine}]
\]

%dsp-level:0
\[
(6):\ \Delta p\ge\tfrac1{2\Delta x_{\mathrm B}}=\tfrac{m_{\mathrm{hid}}}2\ [\mathrm{BM.13}(4)\ \text{re-cited}];
\]

%dsp-level:0
\[
\qquad \langle p^2\rangle=\sum_i\bigl((\Delta p_i)^2+\langle p_i\rangle^2\bigr)=\sum_i(\Delta p_i)^2\ge3\bigl(\tfrac{m_{\mathrm{hid}}}2\bigr)^2\ \text{at}\ \langle p_i\rangle=0\ [\text{machine}]
\]

%dsp-level:0
\[
E^2=m^2+\langle p^2\rangle\ \text{in the vacuum gauge of}\ (4):\quad E^2\ge\bigl(1+\tfrac34\bigr)m_{\mathrm{hid}}^2\ \text{at}\ m=m_{\mathrm{hid}}\ [\text{machine}];
\]

%dsp-level:0
\[
\qquad \mathrm{KE}_{\mathrm{NR}}=\tfrac{\langle p^2\rangle}{2m}\ge\tfrac{3/4}{2}\,m_{\mathrm{hid}}=\tfrac38m_{\mathrm{hid}}\ [\text{machine}];
\]

%dsp-level:0
\[
\qquad \bigl(\tfrac{\sqrt7}2\bigr)^2=\tfrac74,\ x^2-7\ \text{irreducible over}\ \mathbb Q\ [\text{machine: the quadratic ring}];
\]

%dsp-level:0
\[
\qquad E^2\ge\tfrac74m_{\mathrm{hid}}^2\wedge E>0\implies E\ge\tfrac{\sqrt7}2m_{\mathrm{hid}}\ [\text{squaring comparison}]
\]

%dsp-level:0
\[
(7):\ p=0\ \text{in}\ (4):\ E^2=N^2m^2\implies E=Nm\ [\text{the positive branch}];
\]

%dsp-level:0
\[
\qquad RV=\varkappa^2m\ \text{and}\ \rho V=m:\ \text{the same monomial}\ [\text{machine}];
\]

%dsp-level:0
\[
\qquad \kappa^2\rho
\]

%dsp-level:0
\[
= \tfrac{2\pi\mathfrak h^2}{\ln2}\cdot\tfrac{E_{\mathrm{int}}\ln2}{\mathfrak h V}
\]

%dsp-level:0
\[
= \tfrac{2\pi\mathfrak h E_{\mathrm{int}}}V\ [\mathrm{BM.9}(5)\ \text{re-verified; machine: denominators cleared}].\ \blacksquare
\]

\textbf{Phys.} Motion enters the corpus through the metric the corpus already built. The dictionary's ten coefficients split into lapse, shift, and spatial metric by the Schur complement alone, and the mass-shell constraint of the Lorentzian metric turns that split into one exact identity: the energy conjugate to ledger time, corrected by the shift's drag, is the lapse times the square root of the rest mass squared plus the spatial momentum squared. In the vacuum gauge that is precisely the dispersion relation, and its expansion gives the kinetic energy with rational coefficients no one fitted — one half, minus one eighth, plus one sixteenth — nonnegative, vanishing only at zero momentum, wherever the dictionary's perturbative positivity holds. Two exact consequences deserve names. The width law converts the uncertainty bound at the Born width into a zero-point bound: a sealed composite spread over its own width carries at least seven quarters of its rest mass squared in energy squared — three eighths of the rest mass in the non-relativistic kinetic read, $\sqrt7/2$ times it in the exact relativistic one — every constant an exact rational over the ledger count. And the rest slice closes a triangle the corpus drew twice before: at zero momentum, the dispersion energy, the informational mass, and the curvature flux are one and the same monomial.

\end{theorem}

\begin{theorem}[{\texorpdfstring{BM.35 — The vacuum-balanced entropy: the rest mass as the deficit against the all-virtual vacuum}{BM.35 - The vacuum-balanced entropy: the rest mass as the deficit against the all-virtual vacuum}}]\label{tag:BM-35}
\textbf{Thm.}

%dsp-level:0
\[
(1)\quad\textbf{the vacuum is the unique entropy maximum:}
\]

%dsp-level:0
\[
\qquad n_{\mathcal R}=0\implies\rho_{pq}=\rho_\circ=\tfrac14I\ \ \forall(p,q),
\]

%dsp-level:0
\[
\qquad S_{\mathrm{raw}}(\rho_\circ)=\ln4=\max_\rho S_{\mathrm{raw}}(\rho),
\]

%dsp-level:0
\[
\qquad \arg\max=\bigl\{\tfrac14I\bigr\}
\]

%dsp-level:0
\[
\qquad [\mathrm{CT0.2}];
\]

%dsp-level:0
\[
\textbf{the machine form:}
\]

%dsp-level:0
\[
\qquad \lambda=\bigl(\tfrac{k_1}8,\dots,\tfrac{k_4}8\bigr),\ \textstyle\sum_ik_i=8:
\]

%dsp-level:0
\[
\qquad S(\lambda)\le\ln4\iff\prod_ik_i^{k_i}\ge256\quad(0^0:=1)
\]

%dsp-level:0
\[
\qquad [\text{integer comparison;}\ \exp\ \text{strictly monotone}];
\]

%dsp-level:0
\[
\text{the exhaustive census of all}\ \tbinom{11}3
\]

%dsp-level:0
\[
= 165\ \text{compositions: the inequality holds at every one, with equality only at}\ (2,2,2,2)\ [\text{machine}]:
\]

%dsp-level:0
\[
\qquad \textbf{maximality and uniqueness without one transcendental evaluated};
\]

%dsp-level:0
\[
(2)\quad\textbf{the vacuum is flat:}
\]

%dsp-level:0
\[
\qquad (a_{00},q_{ij},w_i)=0\implies(N,N^i,g_{ij})=(1,0,\delta_{ij})\ [\mathrm{BM.34}(2)];
\]

%dsp-level:0
\[
\qquad \Theta_R=0\ \text{on geodesic polygons in flat slices}\ [\mathrm{FT.47}(1)];
\]

%dsp-level:0
\[
\textbf{the vacuum is entropy-maximal and defect-free}\ —\ \text{the two reference points against which every deficit read is read};
\]

%dsp-level:0
\[
(3)\quad\textbf{the ledger deficit and the balance identity:}
\]

%dsp-level:0
\[
\qquad \text{one real edge}=1\ \text{bit}=\ln2\ \text{nats}\ [\mathrm{C9};\ \mathrm{CT0.9};\ \mathrm{CT1.29}]:
\]

%dsp-level:0
\[
\qquad S_{\mathfrak L}=n_{\mathcal R}\ln2;
\]

%dsp-level:0
\[
\qquad \text{the vacuum:}\ n_{\mathcal R}=0\implies S_{\mathfrak L}(\mathrm{vac})=0\ [\mathrm{CT0.2};\ \mathrm{D0.1}];
\]

%dsp-level:0
\[
\Delta S_{\mathfrak L}:=S_{\mathfrak L}-S_{\mathfrak L}(\mathrm{vac})=E_{\mathrm{int}}\ln2=I_{\mathrm{hid}}\qquad[\mathrm{FT.42}(1);\ \mathrm{FT.44}(4);\ \text{machine: exact in}\ \mathbb Q\cdot\ln2];
\]

%dsp-level:0
\[
\textbf{the balance identity:}
\]

%dsp-level:0
\[
\qquad \boxed{\ \ln4=2\ln2\ }
\]

%dsp-level:0
\[
\qquad [\text{machine}]:
\]

%dsp-level:0
\[
\qquad \text{each bidirectionally crystallized channel}\ (2\ \text{directed edges}\ [\mathrm{D0.8}])\ \text{books ledger entropy}\ 2\ln2
\]

%dsp-level:0
\[
= \ln4
\]

%dsp-level:0
\[
= \text{the vacuum pair's own entropy};
\]

%dsp-level:0
\[
E_{\mathrm{int}}
\]

%dsp-level:0
\[
= 2E_{\mathrm{bidir}}\ [\text{the bidirectional composites: the digon}\ [\mathrm{FT.28}],\ \text{the polygon ladder}\ [\mathrm{FT.45}(1)],\ \text{the Platonic floors}\ [\mathrm{FT.42}(1)],\ \text{the neutron row}\ [\mathrm{BM.17}]]:
\]

%dsp-level:0
\[
\boxed{\ \Delta S_{\mathfrak L}=E_{\mathrm{bidir}}\ln4=E_{\mathrm{bidir}}\,S_{\mathrm{raw}}(\rho_\circ)\ }
\]

%dsp-level:0
\[
\qquad [\text{machine: exact in}\ \mathbb Q\cdot\ln2]:
\]

%dsp-level:0
\[
\qquad \textbf{crystallization re-books the vacuum's pair entropy into the ledger in full}\ —\ \text{the balance is exact};
\]

%dsp-level:0
\[
(4)\quad\textbf{the reframing identity:}
\]

%dsp-level:0
\[
\qquad m_{\mathrm{rest}}:=\frac{\Delta S_{\mathfrak L}}{\mathfrak h}=\frac{E_{\mathrm{int}}\ln2}{\mathfrak h}=m_{\mathrm{hid}}
\]

%dsp-level:0
\[
\qquad [\mathrm{FT.44}(4);\ \text{machine: Laurent identity}];
\]

%dsp-level:0
\[
[\Delta S_{\mathfrak L}]=\mathfrak q^0,\quad[\mathfrak h]=\mathfrak q^{-1/2}\implies[m_{\mathrm{rest}}]=\mathfrak q^{+1/2}=[\text{energy}]
\]

%dsp-level:0
\[
\qquad [\mathrm{CT2.1};\ \mathrm{D2.9};\ \text{machine: exact exponent arithmetic}];
\]

%dsp-level:0
\[
\textbf{the substance clause:}
\]

%dsp-level:0
\[
\qquad \text{the rest mass is the deficit read of ledger data against the vacuum}\ —\ \neg\text{an independent entity: the ontology adds nothing};
\]

%dsp-level:0
\[
\text{the rest energy:}\quad E_{\mathrm{rest}}=m_{\mathrm{rest}}c^2\ \ \text{with}\ \ c=1\ [\mathrm{D2.9}]:
\]

%dsp-level:0
\[
\qquad \text{the}\ c^2\ \text{is bookkeeping restoration only}\ [\text{machine: the}\ c\text{-exponent is}\ 0];
\]

%dsp-level:0
\[
(5)\quad\textbf{the geometric member of the deficit family:}
\]

%dsp-level:0
\[
\qquad m^2_{\mathrm{pole}}=4\hat\mu'^2\sin^2\Bigl(\frac{\Theta_R}4\Bigr),
\]

%dsp-level:0
\[
\qquad \Theta_R=\textstyle\sum_i\alpha_i-(n-2)\pi
\]

%dsp-level:0
\[
\qquad [\mathrm{Def\text{-}H.37}(2);\ \mathrm{FT.54}(4);\ \mathrm{FT.47}(1)]:
\]

%dsp-level:0
\[
\text{the pole mass is the angle-deficit read against}\ (2)\text{'s flat vacuum};
\]

%dsp-level:0
\[
\qquad \text{the range}\ [0,4\hat\mu'^2]\ [\mathrm{FT.76}(3)];
\]

%dsp-level:0
\[
\qquad \text{the exact witness}\ 2-\sqrt2\in\mathbb Q(\sqrt2)\ [\mathrm{FT.54}(5)];
\]

%dsp-level:0
\[
\neg\exists\ \text{a theorem equating the bit currency}\ (\ln2\ \text{per directed edge})\ \text{and the angle currency}\ (\Theta_R):
\]

%dsp-level:0
\[
\qquad \text{the relating equation is application-layer content};
\]

%dsp-level:0
\[
(6)\quad\textbf{the ensemble refinement and the third currency:}
\]

%dsp-level:0
\[
\qquad \langle m\rangle=\sum_Rw(R)\,m(R)\ [\mathrm{BM.20}(3)];
\]

%dsp-level:0
\[
\qquad \Delta S(\Gamma,R)=\tfrac32\bigl(2E-|R|\bigr)\ \text{nats}\ [\mathrm{FT.107}(2);\ \mathrm{L3}];
\]

%dsp-level:0
\[
\textbf{the three deficit functionals are pairwise distinct:}
\]

%dsp-level:0
\[
\qquad \text{the digon skeleton}\ (2E=2):\quad|R|=2:\ \ \Delta S_{\mathfrak L}=2\ln2\ \ \text{vs}\ \ \Delta S=0;
\]

%dsp-level:0
\[
\qquad R=\varnothing:\ \ \Delta S_{\mathfrak L}=0\ \ \text{vs}\ \ \Delta S=3
\]

%dsp-level:0
\[
\qquad [\text{machine}];
\]

%dsp-level:0
\[
2\ln2\ne3\ \text{by the exact enclosure}\ \tfrac12<\ln2<1\ [\text{machine:}\ \ln2=\textstyle\sum_{n\ge1}\tfrac1{n2^n},\ \text{tail}<2^{-20}];
\]

%dsp-level:0
\[
\qquad \mathrm{BM.20}\ \mathrm{Rmk}\text{'s agreement row stands untouched};
\]

%dsp-level:0
\[
(7)\quad\textbf{the total energy assembly:}
\]

%dsp-level:0
\[
\qquad E_{\mathrm{total}}^2=m_{\mathrm{rest}}^2+p^2\ [\mathrm{BM.34}(4)\ \text{vacuum gauge}];
\]

%dsp-level:0
\[
\qquad E_{\mathrm{total}}=N\sqrt{m_{\mathrm{rest}}^2+p^2}-N^ip_i\ [\mathrm{Def\text{-}BM.10}(5)];
\]

%dsp-level:0
\[
\textbf{NR:}
\]

%dsp-level:0
\[
\qquad \boxed{\ E_{\mathrm{total}}=m_{\mathrm{rest}}+\mathrm{KE}=\frac{\Delta S_{\mathfrak L}}{\mathfrak h}+\frac{N\,\gamma^{ij}p_ip_j}{2m}+O(p^4)\ }
\]

%dsp-level:0
\[
\qquad [\text{machine: exact series assembly}];
\]

%dsp-level:0
\[
(8)\quad\textbf{status and scope:}
\]

%dsp-level:0
\[
\qquad \text{the reframing adds no content: it is an identity of derived reads}\ ((3),(4))\ \text{plus a name}\ [\text{A §0.2 bookkeeping}];
\]

%dsp-level:0
\[
\neg\text{a declaration},\quad\neg\text{a P4 member},\quad\text{§Cal-BM stays empty};
\]

%dsp-level:0
\[
\qquad (3)\text{'s balance identity holds exactly on the bidirectionally crystallized composites};
\]

%dsp-level:0
\[
\text{a unidirectional edge books}\ \ln2\ \text{against the pair freedom}\ \ln4\ (\text{half-balance}):
\]

%dsp-level:0
\[
\qquad \text{the general directed form is}\ (4)\text{'s}\ m_{\mathrm{hid}}\ \text{identity}\ [\mathrm{FT.44}(4)].
\]

\textbf{Rmk (three currencies, one reference point).} The ledger deficit counts the crystallized bits — $\ln2$ per directed edge; the branch-state deficit counts the uncrystallized channels at $\tfrac32$ nats each [FT.107(2)]; the angle deficit counts the holonomy against flatness [FT.54(4)]. All three are reads of structure against the same vacuum — the all-virtual entropy maximum, which is also the flat metric — and no theorem of Parts 1–4 converts one currency into another: item 7 holds the relating-equation question, and this tier does not touch it. The balance identity is the sharpest single row the corpus can state about why the rest read deserves the entropy name: on a bidirectional crystallization the ledger books exactly the entropy the vacuum pair carried, so the mass counts re-booked pair freedom, not destroyed information — $\ln4$ per pair, $2\ln2$ per channel, $E_{\mathrm{int}}\ln2$ in total, one $\mathfrak h$ away from an energy.

\textbf{Req.} C9, D0.1, D0.8, D2.9 (Parts 1–2); CT0.2, CT0.9 (Part 1); CT1.29, CT2.1 (Part 2); P1, P5 (Parts 1–2); Def-H.37(2), FT.47(1), FT.54(4), FT.54(5), FT.107(2), L3 (Part 3); Census FT.28, FT.42(1), FT.44(4), FT.45(1), FT.76(3) (row-wise); Def-BM.4, Def-BM.5, Def-BM.10, BM.17, BM.20(3), BM.34(2), BM.34(4) (this part); exact $\mathbb Q\cdot\ln2$ arithmetic, the exhaustive integer-grid Gibbs census, exact rational $\ln2$ enclosures, Laurent monomials [math; machine: \texttt{checks/\allowbreak{}verify\_\allowbreak{}H\_\allowbreak{}kinetic\_\allowbreak{}energy\_\allowbreak{}metric.py} §B].

\textbf{Pf.}

%dsp-level:0
\[
(1):\ S(\lambda)=-\sum_i\tfrac{k_i}8\ln\tfrac{k_i}8=\ln8-\tfrac18\sum_ik_i\ln k_i;
\]

%dsp-level:0
\[
\qquad S\le\ln4=\ln8-\ln2\iff\textstyle\sum_ik_i\ln k_i\ge8\ln2=\ln256\iff\prod_ik_i^{k_i}\ge256\ [\exp\ \text{strictly monotone}];
\]

%dsp-level:0
\[
\qquad \text{the census over all}\ 165\ \text{compositions of}\ 8\ \text{into}\ 4\ \text{parts: integers only}\ [\text{machine}]
\]

%dsp-level:0
\[
(2):\ \mathrm{BM.34}(2)\ \text{re-cited}\ [\text{machine}];\qquad\Theta_R=\bar KA=0\ \text{at}\ \bar K=0\ [\mathrm{FT.47}(1)\ \text{re-cited}]
\]

%dsp-level:0
\[
(3):\ S_{\mathfrak L}=n_{\mathcal R}\ln2\ [\mathrm{C9};\ \mathrm{CT0.9};\ \mathrm{CT1.29}];
\]

%dsp-level:0
\[
\qquad n_{\mathcal R}=E_{\mathrm{int}}\ [\mathrm{D0.8};\ \mathrm{Def\text{-}BM.5}(4)\ \text{pattern}];
\]

%dsp-level:0
\[
\qquad 2\ln2=\ln4\ \text{and}\ E_{\mathrm{bidir}}\cdot2\ln2=E_{\mathrm{int}}\ln2\ \text{at}\ E_{\mathrm{int}}=2E_{\mathrm{bidir}}\ [\text{machine:}\ \mathbb Q\cdot\ln2]
\]

%dsp-level:0
\[
(4):\ \mathrm{FT.44}(4)\ \text{substituted}\ [\text{machine}];
\]

%dsp-level:0
\[
\qquad \text{the exponents:}\ 0-(-\tfrac12)=\tfrac12\ [\mathrm{CT2.1};\ \text{machine}];
\]

%dsp-level:0
\[
\qquad c=1\ [\mathrm{D2.9}]:\ \text{the restored}\ c^2\ \text{carries}\ \mathfrak q\text{-exponent}\ 0\ [\text{machine}]
\]

%dsp-level:0
\[
(5):\ \mathrm{Def\text{-}H.37}(2)/\mathrm{FT.54}(4)\ \text{re-cited with}\ \mathrm{FT.47}(1)\text{'s}\ \Theta_R;
\]

%dsp-level:0
\[
\qquad \text{the range and the witness re-cited}\ [\mathrm{FT.76}(3);\ \mathrm{FT.54}(5):\ (2-\sqrt2)(2+\sqrt2)=2\ \text{in}\ \mathbb Q(\sqrt2)]
\]

%dsp-level:0
\[
(6):\ \text{the witness evaluations:}\ 2\ln2\ \text{vs}\ \tfrac32(2-2)=0;
\]

%dsp-level:0
\[
\qquad 0\ \text{vs}\ \tfrac32(2-0)=3\ [\text{machine}];
\]

%dsp-level:0
\[
\qquad \text{the enclosure: the partial sum}\ \textstyle\sum_{n=1}^{20}\tfrac1{n2^n}\ \text{with tail}<2^{-20}\implies\tfrac12<\ln2<1\ [\text{machine: exact Fractions}]\implies2\ln2<2<3
\]

%dsp-level:0
\[
(7):\ \mathrm{BM.34}(4)\text{'s gauge row}\ +\ (5)\text{'s series composed}\ [\text{machine}];\qquad(8):\ \text{the scope record}.\ \blacksquare
\]

\textbf{Phys.} The rest mass of the corpus has always been an entropy in disguise, and the disguise comes off here. The vacuum is the unique state of maximal raw information, and it is flat: maximal entropy, zero defect — one reference point in two currencies. Structure is a lowering against that reference, and the ledger books the lowering in bits, one per crystallized directed edge. The balance identity is why the reframing earns its name: a bidirectionally crystallized channel books exactly the entropy its vacuum pair carried — two bits, $\ln4$ — so the ledger deficit is a complete re-booking, not a residue, and division by the single free scale turns it into the informational mass the floors have always carried. The reframing costs nothing: same monomial, same dimensions, no new entity, no declaration. What it clarifies is the family structure of the mass reads — the pole mass reads the angle deficit against the same flat vacuum, the ensemble reads the branch-state deficit against the same all-virtual reference — three currencies, one reference point, no exchange rate derived. The total energy of a moving composite is then the sum of two derived reads: the deficit against the vacuum, at rest; and the metric's kinetic term, in motion.

\end{theorem}

\begin{theorem}[{\texorpdfstring{BM.36 — The force-structure ladder: the mask dimensions, the transport channel, and the inter-family constraint}{BM.36 - The force-structure ladder: the mask dimensions, the transport channel, and the inter-family constraint}}]\label{tag:BM-36}
\textbf{Thm.}

%dsp-level:0
\[
(1)\quad\textbf{the flag and its increments:}
\]

%dsp-level:0
\[
\qquad \dim\mathcal K_n=(6,16,17,25)\ \ \text{strict}\ [\mathrm{CT0.14}];
\]

%dsp-level:0
\[
\qquad (G_5{+}W_3{+}H_1)+(J^+_3{+}J^-_3)+D^+_1=16=C_{\mathrm{cap}}\ [\mathrm{CT1.12a}(1);\ \mathrm{CT0.12}];
\]

%dsp-level:0
\[
\text{rung}\ 1\to2:\quad+1=\dim Z\bigl(\iota_2(\mathfrak u(2))\bigr)=\dim\mathcal K^{(\mathrm{EM})}\ [\mathrm{D2.12};\ \mathrm{CT0.14}];
\]

%dsp-level:0
\[
\qquad \text{rung}\ 2\to3:\quad+8=\dim\mathfrak{su}(3)=\dim\mathcal K^{(\mathrm{strong})}\ [\mathrm{D2.12};\ \mathrm{CT3.8}];
\]

%dsp-level:0
\[
\text{the observable bound:}\quad\dim\mathcal K_\Sigma\le1+5+1+3+8=18\ [\mathrm{FT.26}(3)]
\]

%dsp-level:0
\[
\qquad [\text{machine: exact integer censuses}];
\]

%dsp-level:0
\[
(2)\quad\textbf{the derived structure ratios:}
\]

%dsp-level:0
\[
\qquad \frac{\dim\mathcal K^{(\mathrm{EM})}}{\dim F_W}=\frac13\quad[F_W=F_{W^+}\oplus F_{W^-}\oplus F_Z,\ \dim=3;\ \mathrm{C8}];
\]

%dsp-level:0
\[
\qquad \frac{\dim\mathcal K^{(\mathrm{strong})}}{\dim\mathcal K^{(\mathrm{EM})}}=\frac81=8;
\]

%dsp-level:0
\[
\qquad \frac{\dim\mathcal K^{(\mathrm{EM})}}{\dim\mathcal K_3}=\frac1{25}
\]

%dsp-level:0
\[
\qquad [\text{machine: exact Fractions}];
\]

%dsp-level:0
\[
\textbf{these are the alphabet-capacity ratios of the rung-activation order}\ [\mathrm{CT0.13}]\ —\ \text{the only force-structure ratios the corpus derives};
\]

%dsp-level:0
\[
(3)\quad\textbf{the strength bar:}
\]

%dsp-level:0
\[
\qquad \neg\exists\ \text{a theorem fixing any numerical coupling strength}:
\]

%dsp-level:0
\[
\qquad (g,g')\ \text{are P4-class}\ [\mathrm{EW.33};\ \mathrm{P4}];
\]

%dsp-level:0
\[
\qquad \alpha^{(2)}_{\mathrm{EM}}\ \text{is an operator block read whose numerical value no theorem of the corpus fixes}\ [\mathrm{EM.17};\ \mathrm{EM.25}(8);\ \text{Part 2 §Cal-2's logged row}];
\]

%dsp-level:0
\[
\text{the identification}\ \alpha_X:=\text{a dimension ratio is}\ \neg\text{derived and}\ \neg\text{adopted}:
\]

%dsp-level:0
\[
\qquad \text{reading}\ (2)\text{'s rationals as coupling strengths is a category error of the}\ \pi\text{-against-}1/137\ \text{species}\ [\mathrm{FT.73}(5)\ \text{trap pattern}];
\]

%dsp-level:0
\[
[\text{machine: the text pins}\ —\ \text{the bar on record in this tag; zero numerical strength values in the tier}];
\]

%dsp-level:0
\[
(4)\quad\textbf{the transport channel (the mass-generation locus):}
\]

%dsp-level:0
\[
\qquad \operatorname{Stab}_{\mathfrak g_{\mathrm{EW}}}(h_0\Sigma)
\]

%dsp-level:0
\[
= \mathfrak g_{\mathrm{EW}}:\ \ \neg\text{adjoint breaking}\implies\text{mass generation is a transport effect on a carried seed,}\ \neg\text{a commutator effect on generators}\ [\mathrm{EW.15}];
\]

%dsp-level:0
\[
\text{the seed:}\quad\mathcal H_{\mathsf v}\cong\mathbb C^2,\quad W_i\mapsto\tfrac12\sigma_i,\quad Y\mapsto\tfrac12\sigma_0,\quad T_a\mathsf v=0\ [\mathrm{EW.32}];
\]

%dsp-level:0
\[
\qquad M^2_{AB}=\mu_\star^2\operatorname{Re}\langle\mathcal G_A\mathsf v,\mathcal G_B\mathsf v\rangle,\quad\ker M^2=\text{the stabilizer},\quad h_0=0\implies M^2\equiv0\ [\mathrm{EW.33}];
\]

%dsp-level:0
\[
\textbf{the exact recomputation in the weak frame}\ [\mathrm{EW.20}]:\qquad M^2=\frac{\mu_\star^2}4\begin{pmatrix}g^2&0&0&0\\0&g^2&0&0\\0&0&g^2&-gg'\\0&0&-gg'&g'^2\end{pmatrix}\qquad[\text{machine: Pauli spinors over}\ \mathbb Q(i)\otimes\mathbb Q[g,g']:\ \text{the real parts give the block matrix; the imaginary parts are antisymmetric and cancel in the symmetric form}\ c_Ac_B\ [\mathrm{EW.33}]];
\]

%dsp-level:0
\[
\text{the spectrum}\ \Bigl\{\frac{g^2}4,\ \frac{g^2}4,\ 0,\ \frac{g^2+g'^2}4\Bigr\}\mu_\star^2\ [\text{machine: the charpoly factorization exact in}\ \mathbb Q[g,g',x]];
\]

%dsp-level:0
\[
\qquad \text{the kernel line}\ c\propto(0,0,g',g)\ [\text{machine: exact annihilation}];
\]

%dsp-level:0
\[
\frac{M_Z^2}{M_W^2}=\frac{g^2+g'^2}{g^2}\ \ \text{exact in}\ \mathbb Q(g^2,g'^2)\ [\text{machine: cross-multiplied}];
\]

%dsp-level:0
\[
\qquad e=g\sin\theta_W\ [\mathrm{EW.35}]\ \text{stands untouched};
\]

%dsp-level:0
\[
(5)\quad\textbf{the inter-family constraint theorem:}
\]

%dsp-level:0
\[
\qquad \text{every fermion pole mass takes the form}\quad m^2_{\mathrm{pole}}(c)=\hat\mu^2\,\ell_c^2\,G(c),
\]

%dsp-level:0
\[
\qquad G(c)=\sum_{A,B}c_Ac_B\operatorname{Re}\langle\mathcal G_A\mathsf v,\mathcal G_B\mathsf v\rangle
\]

%dsp-level:0
\[
\qquad [\mathrm{Def\text{-}H.37}(1)]:
\]

%dsp-level:0
\[
\text{an inter-family mass ratio on this channel is a ratio of}\ G(c)\,
\]

%dsp-level:0
\[
\ell_c^2\ \text{values across transport profiles}\ [\text{structural consequence}];
\]

%dsp-level:0
\[
\qquad \textbf{the derived sandwich:}
\]

%dsp-level:0
\[
\qquad 0\ \le\ G(c)\ \le\ \mu_\star^2\,\frac{g^2+g'^2}4\,\|c\|^2
\]

%dsp-level:0
\[
\qquad [\text{machine: the difference form's PSD factorization}];
\]

%dsp-level:0
\[
\textbf{the exact massless line:}\qquad G(c)=0\iff\textstyle\sum_Ac_A\mathcal G_A\in\operatorname{Ann}(\mathsf v)\iff c\parallel(0,0,g',g)\qquad[\mathrm{EW.33};\ \text{machine}];
\]

%dsp-level:0
\[
\textbf{the range ceilings on the computed skeletons:}
\]

%dsp-level:0
\[
\qquad \text{the pole ladder lies in}\ [0,4\hat\mu'^2]\ [\mathrm{FT.76}(3)];
\]

%dsp-level:0
\[
\qquad \neg\exists\ \text{an assignment of defect reads reaching a multi-order family ratio}\ [\mathrm{FT.77},\ \text{the quarantined audit re-cited}];
\]

%dsp-level:0
\[
\text{a cover read spanning}\ 1{:}R\ \text{requires skeleton scale}\ V\ge\tfrac43R\ [\mathrm{FT.79}(6)\ \text{re-cited}]:
\]

%dsp-level:0
\[
\qquad \textbf{the dynamical range must come from the transport data}\ (\ell_c,c_A)\ —\ \neg\text{from the order-one structure of}\ (1)\text{--}(2);
\]

%dsp-level:0
\[
(6)\quad\textbf{the status row (the inter-family ratios):}
\]

%dsp-level:0
\[
\qquad \textbf{constrained,}\ \neg\textbf{derived:}
\]

%dsp-level:0
\[
\qquad \text{the channel}\ ((4)\text{'s transport}),
\]

%dsp-level:0
\[
\ \text{the form (the Gram)},
\]

%dsp-level:0
\[
\ \text{the kernel (the exact massless line)},
\]

%dsp-level:0
\[
\ \text{the sandwich}\ ((5)),
\]

%dsp-level:0
\[
\ \text{the ceilings}\ ((5))\ \text{are derived};
\]

%dsp-level:0
\[
\neg\exists\ \text{a map in Parts 1--4 assigning transport profiles}\ c_A\ \text{to family lines}\ [\text{inventory scan; machine: no clause of}\ \mathrm{Def\text{-}BM.6}\ \text{or of the channel tags outputs a}\ c_A]:
\]

%dsp-level:0
\[
\boxed{\ \textbf{the Yukawa-like coupling}\ —\ \text{the family-profile map and its scale}\ —\ \textbf{is application-layer content}\ —\ \neg\textbf{a structural absence of the base layer}\ };
\]

%dsp-level:0
\[
\text{its construction locus is}\ (4)\text{'s seed transport}\ [\mathrm{EW.32},\ \mathrm{EW.33},\ \mathrm{Def\text{-}H.37}(1)];
\]

%dsp-level:0
\[
\qquad \text{the Koide-form invariant}\ Q
\]

%dsp-level:0
\[
= \tfrac23\ [\mathrm{Def\text{-}BM.6}(4);\ \mathrm{FT.158}(4)]\ \text{stands as a structural invariant of the naming rows};
\]

%dsp-level:0
\[
\qquad \text{its mass-form bridge}\ (\sqrt m\propto r\ \text{vs}\ m\propto r)\ \text{is}\ \neg\text{adopted}\ [\mathrm{Def\text{-}BM.6}(5)].
\]

\textbf{Rmk (what the dimension ratios can and cannot carry).} The ladder $6\subset16\subset17\subset25$ and its ratios $(\tfrac13,\,8,\,\tfrac1{25})$ count activation capacity — how many generator directions the ledger can switch on at each rung. A strength read would require a coupling theorem relating those counts to resolvent weights; no such theorem exists in Parts 1–4, and the couplings the corpus does carry are free positive P4 numbers whose values no clause fixes. The ceilings of (5) sharpen the bar into a quantitative shape: order-one structure ratios cannot carry a multi-order hierarchy on the computed cells, so an inter-family mass generation must come from the transport data — profiles and lengths — whose derivation is application-layer construction, not from the flag.

\textbf{Req.} C8 (Part 1); D2.12 (Part 1); CT0.12, CT0.13, CT0.14 (Part 1); CT1.12a(1), CT3.8, EW.15, EW.20, EW.32, EW.33, EW.35, EM.17, EM.25(8) (Part 2); P4, P5 (Parts 1–2); Def-H.37(1), FT.26(3), FT.73(5), FT.77, FT.79(6), FT.158(4) (Part 3); Census FT.76(3) (row-wise); Def-BM.6 (this part); exact Pauli/Gram algebra over $\mathbb Q(i)\otimes\mathbb Q[g,g']$, charpoly factorization in $\mathbb Q[g,g',x]$, exact Fractions, text pins [math; machine: \texttt{checks/\allowbreak{}verify\_\allowbreak{}H\_\allowbreak{}kinetic\_\allowbreak{}energy\_\allowbreak{}metric.py} §C].

\textbf{Pf.}

%dsp-level:0
\[
(1):\ \mathrm{CT0.14}\text{'s flag re-cited};
\]

%dsp-level:0
\[
\qquad 3+3+5+3+1+1=16=C_{\mathrm{cap}},\quad16+1=17,\quad17+8=25,\quad6<16<17<25,\quad1+5+1+3+8=18\ [\text{machine}]
\]

%dsp-level:0
\[
(2):\ \dim F_W=3\ [\mathrm{C8}];\qquad\tfrac{1}{3},\ \tfrac{8}{1},\ \tfrac{1}{25}\ [\text{machine: exact Fractions}]
\]

%dsp-level:0
\[
(3):\ \mathrm{EW.33}\text{'s P4 row and Part 2 §Cal-2's logged}\ \alpha^{(2)}_{\mathrm{EM}}\ \text{row re-cited}\ [\text{text pin: ``no theorem of the corpus fixes its numerical value''}];
\]

%dsp-level:0
\[
\qquad \mathrm{FT.73}(5)\text{'s trap pattern re-cited}
\]

%dsp-level:0
\[
(4):\ \mathsf v=(0,1)^{\mathsf T}:\quad\sigma_1\mathsf v=(1,0)^{\mathsf T},\ \ \sigma_2\mathsf v=(-i,0)^{\mathsf T},\ \ \sigma_3\mathsf v=(0,-1)^{\mathsf T},\ \ \sigma_0\mathsf v=\mathsf v;
\]

%dsp-level:0
\[
\qquad \mathcal G=(gW_1,gW_2,gW_3,g'Y)\ [\mathrm{EW.33}]
\]

%dsp-level:0
\[
\operatorname{Re}\langle\mathcal G_A\mathsf v,\mathcal G_B\mathsf v\rangle=\tfrac14\,\bigl(\text{the displayed block matrix}\bigr)_{AB};
\]

%dsp-level:0
\[
\qquad \operatorname{Im}\langle\mathcal G_A\mathsf v,
\]

%dsp-level:0
\[
\mathcal G_B\mathsf v\rangle
\]

%dsp-level:0
\[
= -\operatorname{Im}\langle\mathcal G_B\mathsf v,
\]

%dsp-level:0
\[
\mathcal G_A\mathsf v\rangle\ \text{and cancels identically in}\ \textstyle\sum_{A,B}c_Ac_B(\cdot)\ [\mathrm{EW.33};\ \text{machine}];
\]

%dsp-level:0
\[
\qquad \det(xI-\mathrm{Gram})=\bigl(x-\tfrac{g^2}4\bigr)^2\,x\,\bigl(x-\tfrac{g^2+g'^2}4\bigr)\ [\text{machine: expanded and compared in}\ \mathbb Q[g,g',x]]
\]

%dsp-level:0
\[
\mathrm{Gram}\cdot(0,0,g',g)^{\mathsf T}=0\ [\text{machine}];
\]

%dsp-level:0
\[
\qquad \mu_\star^2=c_W\varkappa^2\beta^2h_0^2\big|_{h_0=0}=0\implies M^2\equiv0\ [\text{machine: substitution}];
\]

%dsp-level:0
\[
\qquad \tfrac{(g^2+g'^2)/4}{g^2/4}=\tfrac{g^2+g'^2}{g^2}\ [\text{machine: cross-multiplied}]
\]

%dsp-level:0
\[
(5):\ \mathrm{Def\text{-}H.37}(1)\text{'s form re-cited};
\]

%dsp-level:0
\[
\qquad \tfrac{g^2+g'^2}4I-\mathrm{Gram}:\ \ \text{the spectrum}\ \bigl\{\tfrac{g'^2}4,\tfrac{g'^2}4,\tfrac{g^2+g'^2}4,\tfrac{g^2}4\bigr\}\ [\text{machine: charpoly factorization}]\implies\text{the sandwich by the Rayleigh range};
\]

%dsp-level:0
\[
\qquad \mathrm{FT.76}(3),\ \mathrm{FT.77},\ \mathrm{FT.79}(6)\ \text{re-cited}\ (\neg\text{recomputed; the quarantined rows stay quarantined})
\]

%dsp-level:0
\[
(6):\ \text{the inventory scan of}\ \mathrm{Def\text{-}BM.6}\ \text{and the channel tags}\ \mathrm{FT.109}\text{--}\mathrm{FT.119},
\]

%dsp-level:0
\[
\ \mathrm{FT.142}\text{--}\mathrm{FT.145}:\ \neg\exists\ \text{a clause outputting a transport profile}\ [\text{machine: token scan}];\qquad \ \blacksquare
\]

\textbf{Phys.} The force side of the mass question splits cleanly into what the corpus derives and what it does not. Derived is the structure: the mask flag grows six to sixteen to seventeen to twenty-five by exact increments — one electromagnetic centre, eight gluons — and the only ratios the ladder supplies are capacity ratios: one third, eight, one twenty-fifth. Derived is also the channel: mass can only be a transport defect of a carried seed, because the scalar background breaks nothing adjointly, and the seed's Gram form comes with an exactly computed spectrum — two equal charged directions, one massive neutral direction, one exact kernel direction along which the transported family is massless whatever the couplings do. Not derived is any coupling value: no theorem turns a count into a strength, and the tier records the bar rather than let a rational drift across the gap. And not derived is the family map: nothing in Parts 1–4 assigns transport profiles to family lines, while the ceilings on the computed skeletons say that order-one structure cannot carry a multi-order hierarchy. The inter-family mass ratios are therefore constrained — channel, form, kernel, sandwich, ceiling — and open, with the construction locus named: the seed transport itself.

\end{theorem}

\begin{theorem}[{\texorpdfstring{BM.37 — The mass-route artifacts: the disjoint input alphabets and the closure of the counting trap}{BM.37 - The mass-route artifacts: the disjoint input alphabets and the closure of the counting trap}}]\label{tag:BM-37}
\textbf{Thm.}

%dsp-level:0
\[
(1)\quad\textbf{the three routes, three exact values:}
\]

%dsp-level:0
\[
\qquad (\mathrm i)\ \ \text{the normalized ensemble average}\ [\mathrm{Def\text{-}H.32}(6);\ \mathrm{FT.107}(3)]:\quad\text{neutral}{:}\text{charged}
\]

%dsp-level:0
\[
= 1\ (\text{degenerate})\ —\ \text{a uniform sector factor cancels at}\ \{1,2,3,15,20,\tfrac73\}\ [\mathrm{BM.30}(3);\ \text{machine}];
\]

%dsp-level:0
\[
(\mathrm{ii})\ \ \text{the total deficit}\ [\mathrm{BM.28}(4)(\mathrm{ii})]:
\]

%dsp-level:0
\[
\qquad \frac{W(0)}{W(+1)}=\frac{20}{15}=\frac43
\]

%dsp-level:0
\[
\qquad [\text{machine: the re-enumeration of}\ (2+z+z^{-1})^3\ \text{and the exhaustive}\ 27\text{-triple census}];
\]

%dsp-level:0
\[
(\mathrm{iii})\ \ \text{the amplitude layer}\ [\mathrm{Def\text{-}H.32}(2),(5);\ \mathrm{BM.28}(4)(\mathrm{iii})]:
\]

%dsp-level:0
\[
\qquad \frac{5336}{1196}=\frac{58}{13}\ (S{=}0),
\]

%dsp-level:0
\[
\qquad \frac{4796}{802}=\frac{2398}{401}\ (S{=}1)
\]

%dsp-level:0
\[
\qquad [\mathrm{FT.106}(3)\text{'s frozen integers; machine: exact reductions}];
\]

%dsp-level:0
\[
\text{pairwise unequal:}\qquad1\ \ne\ \frac43\ \ne\ \frac{58}{13}\ \ne\ \frac{2398}{401}\qquad[\text{machine: cross-multiplications in}\ \mathbb Z];
\]

%dsp-level:0
\[
(2)\quad\textbf{the disjointness theorem (the artifact classification):}
\]

%dsp-level:0
\[
\qquad \text{routes}\ (\mathrm i),
\]

%dsp-level:0
\[
(\mathrm{ii})\ \text{are measure functionals of}\ (\Gamma,R):\ \ \text{the weights}\ (\nu,\ t^{|R|},\ u^L,\ \#\mathrm{inst})\ \text{and the deficit}\ \Delta S\ [\mathrm{Def\text{-}H.32}(3),(6);\ \mathrm{FT.107}(2)];
\]

%dsp-level:0
\[
\text{route}\ (\mathrm{iii})\ \text{is the amplitude functional of}\ v_\Gamma(R)\ [\mathrm{Def\text{-}H.32}(2),(5)];
\]

%dsp-level:0
\[
\qquad \text{the dispersion read is the metric functional of}\ (g_{ij},N,N^i,p_i,m)\ [\mathrm{Def\text{-}BM.10}(2)\text{--}(5)];
\]

%dsp-level:0
\[
\textbf{the machine token census on the canonical defining strings:}
\]

%dsp-level:0
\[
\qquad \text{metric-alphabet occurrences inside the defining clauses of routes}\ (\mathrm i)\text{--}(\mathrm{iii}):\ \ \mathbf 0;
\]

%dsp-level:0
\[
\text{measure-alphabet occurrences inside}\ \mathrm{Def\text{-}BM.10}(2),
\]

%dsp-level:0
\[
(3),
\]

%dsp-level:0
\[
(5)\ \text{(the kinematic clauses; clause}\ (4)\text{'s rest input is the shared symbol}\ m):\ \ \mathbf 0;
\]

%dsp-level:0
\[
\boxed{\ \textbf{the input alphabets are disjoint:}\quad\neg\exists\ \text{a route output that is a function of metric data};\qquad\neg\exists\ \text{a dispersion energy that is a function of measure data}\ }
\]

%dsp-level:0
\[
(3)\quad\textbf{the trap and its closure:}
\]

%dsp-level:0
\[
\qquad \textbf{the counting trap}\ :
\]

%dsp-level:0
\[
= \ \text{reading a route ratio}\ \bigl(1,\ \tfrac43,\ \tfrac{58}{13},\ \tfrac{2398}{401}\bigr)\ \text{or the Koide-form invariant}\ Q
\]

%dsp-level:0
\[
= \tfrac23\ [\mathrm{Def\text{-}BM.6}(4);\ \mathrm{FT.158}(4)]\ \text{as a physical mass ratio}\ [\mathrm{FT.89}\text{'s category-trap species}];
\]

%dsp-level:0
\[
\textbf{barred by three independent rows:}
\]

%dsp-level:0
\[
\qquad (\mathrm a)\ \ \text{the disjointness}\ [(2)]:\ \ \text{a route output carries no momentum dependence and cannot transform as a dispersion energy};
\]

%dsp-level:0
\[
(\mathrm b)\ \ \text{amplitude}\ne\text{mass}\ [\mathrm{FT.79}(5)]:\qquad\text{route}\ (\mathrm{iii})\text{'s bridge is item-13 content,}\ \neg\text{constructed}\;
\]

%dsp-level:0
\[
(\mathrm c)\ \ \text{the route choice is a read specification,}\ \neg\text{a theorem}\ [\mathrm{BM.30}(5)]:
\]

%dsp-level:0
\[
\qquad \text{no route is adopted, and none is retracted};
\]

%dsp-level:0
\[
Q = \tfrac23
\]

%dsp-level:0
\[
= \tfrac{3-1}3\ \text{is the doublet/triplet decomposition ratio of the spatial seal triad}\ [\mathrm{FT.158}(5)]:
\]

%dsp-level:0
\[
\qquad \text{a structural invariant of the naming rows, whose mass-form bridge stays}\ \neg\text{adopted}\ [\mathrm{Def\text{-}BM.6}(5)];
\]

%dsp-level:0
\[
(4)\quad\textbf{the completion row:}\qquad\text{the physical mass-energy of a sealed composite is the dispersion read:}
\]

%dsp-level:0
\[
E=N\sqrt{m_{\mathrm{rest}}^2+p^2}-N^ip_i,\qquad m_{\mathrm{rest}}=\frac{\Delta S_{\mathfrak L}}{\mathfrak h}\qquad[\mathrm{BM.34}(4);\ \mathrm{BM.35}(4),(7)];
\]

%dsp-level:0
\[
\text{the routes are static}\ (p\text{-free})\ \text{slices: each evaluates an ensemble with no momentum argument}\ [(2)];
\]

%dsp-level:0
\[
\qquad \textbf{KE}\ —\ \text{the metric's motion term}\ —\ \textbf{is what no route can express};
\]

%dsp-level:0
\[
\text{at}\ p
\]

%dsp-level:0
\[
= 0\ \text{the dispersion rest term coincides on-shell with the curvature-flux and the informational reads}\ [\mathrm{BM.34}(7)]:
\]

%dsp-level:0
\[
\qquad \text{the routes price the ensemble measure's restrictions around the same rest object};
\]

%dsp-level:0
\[
\mathrm{BM.30}(2)\text{'s splitting and}\ \mathrm{BM.30}(3)\text{'s degeneracy stand derived and exact as}\ \textbf{measure rows}\ —\ \text{the closure classifies, it does not retract};
\]

%dsp-level:0
\[
(5)\quad\textbf{the standing rules:}
\]

%dsp-level:0
\[
\qquad \neg\text{a P4 addition}\ [\text{§Cal-BM stays empty}];
\]

%dsp-level:0
\[
\qquad \neg\text{a new dimensionful scale}\ [\mathrm{P1}];
\]

%dsp-level:0
\[
\qquad \neg\text{a measured value and}\ \neg\text{a decimal numeral in this tier's text}\ [\mathrm{P5};\ \text{machine scans}];
\]

%dsp-level:0
\[
\neg\text{a new declaration};
\]

%dsp-level:0
\[
\qquad \text{species names confined to the identification layers}\ [\mathrm{Def\text{-}BM.3}(7),\ \mathrm{Def\text{-}BM.6}(7)];
\]

%dsp-level:0
\[
\qquad \text{topology labels are internal bookkeeping}\ [\mathrm{Def\text{-}H.32}(4)];
\]

%dsp-level:0
\[
(6)\quad\textbf{the adjudication:}
\]

%dsp-level:0
\[
\qquad \boxed{\ \text{the three mass-route ratios are counting and amplitude artifacts}\ —\ \text{exact as measure prices, barred as physical masses}\ };
\]

%dsp-level:0
\[
\text{the counting trap is closed at the category level by the disjoint input alphabets}\ [(2)];
\]

%dsp-level:0
\[
\text{the physical mass-energy read is the metric dispersion route: the rest term is the vacuum-balanced entropy}\ [\mathrm{BM.35}],
\]

%dsp-level:0
\[
\ \text{the motion term is the shift-transported kinetic energy}\ [\mathrm{BM.34}];
\]

%dsp-level:0
\[
\text{the inter-family content is application-layer}\ [\mathrm{BM.36}(6)].
\]

\textbf{Rmk (what the closure does not do).} The closure retracts no route: every route value stands as a derived exact rational, and BM.30(5)'s read-choice row is untouched — adopting any route as "the mass" still requires a logged declaration. What the closure establishes is the classification: the three functionals price measures and amplitudes, and a measure price is not a dispersion energy however clean its rational looks. The Koide-form invariant belongs to the same classification — it is the decomposition ratio the spatial seal triad forces [FT.158(5)] — and converting it into a mass ratio needs exactly the mass-form bridge — application-layer content. The physical reading runs the other way: the rest mass is the vacuum-deficit read, the motion is the metric's kinetic term, and the route ratios stay what they always were — exact prices of how the charge constraint and the sign sum filter one ensemble.

\textbf{Req.} P1, P4, P5 (Parts 1–2); Def-H.32(2), Def-H.32(3), Def-H.32(4), Def-H.32(5), Def-H.32(6), FT.79(5), FT.89, FT.106(3), FT.107(2), FT.107(3), FT.158(4), FT.158(5) (Part 3); Def-BM.3(7), Def-BM.6(4), Def-BM.6(5), Def-BM.6(7), Def-BM.10, BM.28(2), BM.28(4), BM.30(1), BM.30(2), BM.30(3), BM.30(5), BM.34(4), BM.34(7), BM.35(4), BM.35(7), BM.36(6) (this part); exact rational reductions, the exhaustive triple census, generating-function convolution, token censuses [math; machine: \texttt{checks/\allowbreak{}verify\_\allowbreak{}H\_\allowbreak{}kinetic\_\allowbreak{}energy\_\allowbreak{}metric.py} §D].

\textbf{Pf.}

%dsp-level:0
\[
(1):\ (\mathrm i)\ \ \tfrac{\sum_Rc\,v(R)\Delta S(R)}{\sum_Rc\,v(R)}
\]

%dsp-level:0
\[
= \tfrac{\sum_Rv(R)\Delta S(R)}{\sum_Rv(R)}\ \ \forall c
\]

%dsp-level:0
\[
\ne 0\ [\text{machine at}\ c\in\{1,2,3,15,20,\tfrac73\}\ \text{on rational witness data}]
\]

%dsp-level:0
\[
(\mathrm{ii}):\ W(Q)=\bigl[z^Q\bigr](2+z+z^{-1})^3\ \text{re-enumerated}:
\]

%dsp-level:0
\[
\quad W=(1,6,15,20,15,6,1),
\]

%dsp-level:0
\[
\quad W(0)=8+12=20,
\]

%dsp-level:0
\[
\quad W(\pm1)=12+3=15,\quad\tfrac{20}{15}=\tfrac43\ [\text{machine}];
\]

%dsp-level:0
\[
\qquad (\mathrm{iii}):\ \tfrac{5336}{1196}=\tfrac{58}{13}\ (5336\cdot13=1196\cdot58),\ \ \tfrac{4796}{802}=\tfrac{2398}{401}\ [\text{machine: exact}]
\]

%dsp-level:0
\[
\text{the pairwise inequalities:}\ 3\cdot1\ne4\cdot1,
\]

%dsp-level:0
\[
\quad 4\cdot13\ne3\cdot58,
\]

%dsp-level:0
\[
\quad 58\cdot401\ne13\cdot2398,\quad\text{and the remaining pairs}\ [\text{machine: integer cross-multiplications}]
\]

%dsp-level:0
\[
(2):\ \text{the canonical defining strings extracted from the canon files}\ [\mathrm{Def\text{-}H.32}(2),(3),(5),(6);\ \mathrm{FT.107}(2),(3);\ \mathrm{BM.28}(4);\ \mathrm{Def\text{-}BM.10}(2)\text{--}(5)];
\]

%dsp-level:0
\[
\qquad \text{the metric alphabet}\ \{\Omega,\varkappa,\beta,w_i,q_{ij},\gamma,N,N^i,p_i,g_{0i},\text{lapse},\text{shift}\}:\ 0\ \text{occurrences in the route clauses}\ [\text{machine}]
\]

%dsp-level:0
\[
\text{the measure alphabet}\ \{W(Q),\#\mathrm{inst},\nu,v_\Gamma,\Delta S,t^{|R|},u^L,Z^{\bullet},\text{ensemble},\text{flicker}\}:\ 0\ \text{occurrences in}\ \mathrm{Def\text{-}BM.10}(2)\text{--}(5)\ \text{beyond the shared symbol}\ m\ [\text{machine}]
\]

%dsp-level:0
\[
(3):\ (\mathrm a)\ [(2)\ \text{conjoined with}\ \mathrm{Def\text{-}BM.10}(5)\text{'s}\ p\text{-dependence}];
\]

%dsp-level:0
\[
\qquad (\mathrm b)\ \mathrm{FT.79}(5)\ \text{re-cited};
\]

%dsp-level:0
\[
\qquad (\mathrm c)\ \mathrm{BM.30}(5)\ \text{re-cited};
\]

%dsp-level:0
\[
\qquad Q=\tfrac23\ \text{and}\ \tfrac{3-1}3\ [\mathrm{FT.158}(5);\ \mathrm{Def\text{-}BM.6}(4)\ \text{re-cited}]
\]

%dsp-level:0
\[
(4):\ \mathrm{BM.34}(4),\ \mathrm{BM.35}(4),(7)\ \text{conjoined};
\]

%dsp-level:0
\[
\qquad \mathrm{BM.34}(7)\text{'s on-shell agreement row re-cited};
\]

%dsp-level:0
\[
\qquad (5):\ \text{the standing-rule record}\ [\mathrm{BM.33}(5)\ \text{pattern}]\ +\ \text{the machine scans}
\]

%dsp-level:0
\[
(6):\ \text{conjunction of}\ (1)\text{--}(5).\ \blacksquare
\]

\textbf{Phys.} Three exact rationals — one, four thirds, fifty-eight thirteenths — came out of the charge-constrained ensemble, and the temptation is to read them as mass ratios. The corpus closes that trap the only way it closes anything: by classification. Each of the three routes is a functional of the ensemble's combinatorics alone — pattern counts, norm weights, sign-summed amplitudes — and a machine census of the defining clauses shows their input alphabet contains not one metric token, while the dispersion read's definition contains not one measure token beyond the shared rest symbol. Disjoint inputs, disjoint outputs: no route value can transform as a dispersion energy, because no route has ever seen a momentum. What the routes price is the measure — how the charge constraint rescales the probability space, how the sign sum weights the amplitudes — and those prices stand, exact and unretracted. The physical mass-energy is assembled elsewhere: the rest term from the vacuum-balanced entropy, the motion term from the metric's kinetic energy, the two joined by the mass-shell identity. The Koide-form invariant joins the same classification — a decomposition ratio the spatial triad forces, waiting on a mass-form bridge that item 13 holds open. The counting was never wrong; reading the counts as masses was.

\end{theorem}

\medskip\noindent\textit{\texorpdfstring{BM.38 — The status of the metric-energy tier: adjudication and scope}{BM.38 - The status of the metric-energy tier: adjudication and scope}.}\ Structural-constraints record: see Appendix B, Section B.1.20 of the Addendum.\par\medskip

\phantomsection
\section[Tier BM.11 — The Neutron Decay Vertex: the Family Transport…]{\texorpdfstring{Tier BM.11 — The Neutron Decay Vertex: the Family Transport Profiles, the Rung-1 Isospin Row, the Forced Neutral Member, the Angular-Momentum Budget, and the Full Kinematic Balance}{Tier BM.11 - The Neutron Decay Vertex: the Family Transport Profiles, the Rung-1 Isospin Row, the Forced Neutral Member, the Angular-Momentum Budget, and the Full Kinematic Balance}}

\begin{definition}[{\texorpdfstring{Def-BM.11 — The family transport-profile map (the constructed bridge)}{Def-BM.11 - The family transport-profile map (the constructed bridge)}}]\label{tag:Def-BM-11}
\textbf{Def.} A construction on EW.32/EW.33's seed transport and FT.158(1)'s inner $C_3$ rotation; no new axiom, no new dimensionful scale, no Yukawa-type object — none exists in the inventory [EW.15; FT.145(3)]. Local notation: $\mathrm{Gr}_{AB}:=\operatorname{Re}\langle\mathcal G_A\mathsf v,\mathcal G_B\mathsf v\rangle$ (EW.33's Gram data).

%dsp-level:0
\[
(1)\quad\textbf{the seed data:}
\]

%dsp-level:0
\[
\qquad \mathcal G=(gW_1,\ gW_2,\ gW_3,\ g'Y)\ [\mathrm{EW.33}];
\]

%dsp-level:0
\[
\qquad \mathsf v:\ \ W_3\mathsf v=-\tfrac12\mathsf v,\ \ Y\mathsf v=+\tfrac12\mathsf v\ [\mathrm{EW.20}]:
\]

%dsp-level:0
\[
\qquad \textbf{the hypercharge coordinate}\ y_0=\tfrac12\ \textbf{is derived}\ [\mathrm{EW.20}:\ \text{trace normalization},\ \neg\text{assigned}];
\]

%dsp-level:0
\[
(2)\quad\textbf{the family rotation:}
\]

%dsp-level:0
\[
\qquad U=\exp\!\bigl(-\tfrac{2\pi i}3\,\hat n\cdot W\bigr)=E_0-\tfrac12\mathcal W^2-\tfrac i2\mathcal W,\quad\mathcal W=W_1{+}W_2{+}W_3:
\]

%dsp-level:0
\[
\qquad UW_iU^\dagger=W_{i+1},\quad U^3=E_0\ [\mathrm{FT.158}(1);\ \text{machine: exact over}\ \mathbb Q(i),\ \text{via}\ \mathcal W^3=3\mathcal W];
\]

%dsp-level:0
\[
Y\ \text{is the centre of}\ \iota_2(\mathfrak u(2))\implies[Y,W_i]=0\implies UYU^\dagger=Y\ [\mathrm{EW.8};\ \mathrm{D2.12};\ \text{machine}]:
\]

%dsp-level:0
\[
\qquad \textbf{the}\ C_3\ \textbf{cycles the weak triplet coordinates of a transport profile and fixes the hypercharge coordinate};
\]

%dsp-level:0
\[
(3)\quad\textbf{the profile class:}
\]

%dsp-level:0
\[
\qquad \text{every}\ C_3\text{-covariant triplet profile takes the cosine form}\ \lambda_i
\]

%dsp-level:0
\[
= A+B\cos(\theta+\tfrac{2\pi i}3)\ [\mathrm{FT.111}(2)];
\]

%dsp-level:0
\[
\textbf{unit normalization (constitutive convention, the Def-BM.4(1) pattern):}\quad\text{the weak coordinates carry radius}\ 1;
\]

%dsp-level:0
\[
\textbf{the base profile:}\qquad c(0):=\bigl(\cos\delta,\ \cos(\delta+\tfrac{2\pi}3),\ \cos(\delta+\tfrac{4\pi}3),\ y_0\bigr)\in\mathbb R^4;
\]

%dsp-level:0
\[
\textbf{the orbit:}
\]

%dsp-level:0
\[
\qquad c(i):=\operatorname{Ad}(U^i)\,c(0)\ =\ \text{the cyclic shift of the weak triplet,}\ y_0\ \text{fixed}\ [(2);\ \text{machine}]:
\]

%dsp-level:0
\[
\qquad c(i)_3=\cos\bigl(\delta+\tfrac{2\pi(2-i)}3\bigr);
\]

%dsp-level:0
\[
\text{the}\ W_3\ \text{slot cycles through the triplet};
\]

%dsp-level:0
\[
\qquad \text{the Gram-value set}\ \bigl\{M^2\bigl(A_0-B_0\cos(\delta-\tfrac{2\pi j}3)\bigr)\bigr\}_j\ \text{is the orbit's content under}\ (4)\text{'s orbit-label convention};
\]

%dsp-level:0
\[
(4)\quad\textbf{the angle identification (the single declared clause):}
\]

%dsp-level:0
\[
\qquad \delta:=\delta_{\mathrm{channel}}\ —\ \text{per channel:}\quad\ell:\ \tfrac29\ [\mathrm{FT.112}];
\]

%dsp-level:0
\[
\qquad q:\ \tfrac{2}{27}\ [\mathrm{FT.119}(1)];
\]

%dsp-level:0
\[
\qquad \nu:\ \tfrac18\ [\mathrm{FT.142}(5)];
\]

%dsp-level:0
\[
\text{the orbit-label orientation convention}\ (\delta\mapsto-\delta\ \text{in the}\ W_3\ \text{slot})\ \text{is internal bookkeeping};
\]

%dsp-level:0
\[
\qquad \text{the lines attach to}\ \mathrm{Def\text{-}BM.6}\text{'s named rows by its own spectral order};
\]

%dsp-level:0
\[
(5)\quad\textbf{the mass evaluation:}
\]

%dsp-level:0
\[
\qquad m^2(\text{line}\ i)=\hat\mu^2\,\ell_c^2\,G\bigl(c(i)\bigr),
\]

%dsp-level:0
\[
\qquad G(c)=\sum_{A,B}c_Ac_B\,\mathrm{Gr}_{AB}
\]

%dsp-level:0
\[
\qquad [\mathrm{Def\text{-}H.37}(1);\ \mathrm{EW.33}]:
\]

%dsp-level:0
\[
\boxed{\ m^2_i=\hat\mu^2\ell_c^2\Bigl[\Bigl(\frac{3g^2}8+\frac{g'^2}{16}\Bigr)-\frac{gg'}4\,\cos\Bigl(\delta-\frac{2\pi i}3\Bigr)\Bigr]\ }=:M^2\bigl(A_0-B_0\cos\varphi_i\bigr),
\]

%dsp-level:0
\[
\qquad A_0=\frac{3g^2}8+\frac{g'^2}{16},\quad B_0=\frac{gg'}4
\]

%dsp-level:0
\[
\qquad [\text{machine: exact in}\ \mathbb Q[g,g']\otimes\mathbb Q[c,s,r_3]/(c^2{+}s^2{-}1,r_3^2{-}3)];
\]

%dsp-level:0
\[
\textbf{the affine form:}
\]

%dsp-level:0
\[
\qquad r_2\cdot m^2_i=M^2\bigl(A_0r_2+B_0-B_0\,r_i\bigr),
\]

%dsp-level:0
\[
\qquad r_i=1+r_2\cos\Bigl(\delta+\frac{2\pi i}3\Bigr),\quad r_2^2=2
\]

%dsp-level:0
\[
\qquad [\mathrm{FT.113}\ \text{vocabulary rows};\ \text{machine}];
\]

%dsp-level:0
\[
(6)\quad\textbf{interface status:}
\]

%dsp-level:0
\[
\qquad \textbf{derived:}\ \text{the orbit structure}\ [(2)],
\]

%dsp-level:0
\[
\ \text{the cosine shape}\ [\mathrm{FT.111}(2)],
\]

%dsp-level:0
\[
\ \text{the Gram evaluation and the affine form}\ [(5);\ \text{machine}],
\]

%dsp-level:0
\[
\ \text{positivity, the splitting law, the sum rule, the order inversion}\ [\mathrm{BM.39}];
\]

%dsp-level:0
\[
\textbf{non-derived:}\ \text{exactly one angle identification per channel}\ (4)\ \text{and the unit normalization}\ (3);
\]

%dsp-level:0
\[
\qquad \textbf{scales:}\ \hat\mu^2\in\mathrm{Cal\text{-}F}\ (\text{logged});\ \ \ell_c\ \text{is chain data}\ [\mathrm{D1.4}];\ \ (g,g')\ \text{are P4}\ [\mathrm{EW.33}]:
\]

%dsp-level:0
\[
\qquad \neg\text{a new parameter},\ \text{§Cal-BM stays empty};
\]

%dsp-level:0
\[
\text{declared}\ (\text{A-BM.4})\ \text{before use};
\]

%dsp-level:0
\[
(7)\quad\textbf{the locus clause:}
\]

%dsp-level:0
\[
\qquad \text{the construction lives in the pole/transport read species}\ [\mathrm{Def\text{-}H.37}(1)];
\]

%dsp-level:0
\[
\qquad \mathrm{FT.113}/\mathrm{FT.143}(3)\text{'s inertial vocabulary}\ m^2=\mu r^2\ \text{is a distinct read species on the same lines}:
\]

%dsp-level:0
\[
\text{the two species order the lines oppositely}\ [\mathrm{BM.39}(5)];
\]

%dsp-level:0
\[
\qquad \text{the relating equation is item-7 content, untouched}\;
\]

%dsp-level:0
\[
\textbf{BM.36}(6)\textbf{'s gap label is discharged:}\quad\text{the family-profile map exists as this construction};
\]

%dsp-level:0
\[
\qquad \text{the scale data is P4-class calibration}\ (\neg\text{declared}),\ \neg\text{the map}.
\]

\textbf{Rmk (what the construction is).} The map answers BM.36(6)'s named gap with derived material only: the rotation is FT.158(1)'s — the same $C_3$ that forces the channel spectra, inner in the weak algebra; the shape is FT.111(2)'s — every $C_3$-covariant triplet is a cosine triplet; the evaluation is Def-H.37(1)'s Gram transport; the hypercharge coordinate is EW.20's derived $\tfrac12$. The single declared content is that the profile's phase is the channel's own angle — one identification per channel, and the angles are derived rationals. The payoff is structural: the intra-family splitting is exactly $gg'$ — the weak–hypercharge mixed strength — acting on the transport phase, which is the derived form of "the relative force strength acting on the seed's transport profile." No commutator breaking enters anywhere [EW.15], and no Yukawa-type term exists to enter.

\textbf{Req.} D1.4, D2.12 (Parts 1–2); EW.8, EW.20, EW.32, EW.33 (Part 2); P4, P5 (Parts 1–2); Def-H.37(1), FT.111(2), FT.112, FT.113, FT.119(1), FT.142(5), FT.143, FT.145(3), FT.158(1) (Part 3); Def-BM.4(1), Def-BM.6, BM.36(4), BM.36(5), BM.36(6) (this part); Cal-F; the declaration addendum A-BM.4; exact quotient-ring trigonometric identities in $\mathbb Q[g,g',c,s,r_3,r_2]/(c^2{+}s^2{-}1,\ r_3^2{-}3,\ r_2^2{-}2)$, Gram algebra over $\mathbb Q[g,g']$ [math; machine: \texttt{checks/\allowbreak{}verify\_\allowbreak{}H\_\allowbreak{}neutron\_\allowbreak{}decay\_\allowbreak{}vertex.py} §A].

\textbf{Pf.}

%dsp-level:0
\[
\text{Definition block: the constitutive clauses are}\ (3)\text{'s unit normalization and}\ (4)\text{'s angle identification};
\]

%dsp-level:0
\[
\qquad \text{every displayed evaluation is a machine-verified exact identity}\ [\text{§A}];
\]

%dsp-level:0
\[
\qquad (6),(7)\ \text{are the status and locus records}.\ \blacksquare
\]

\textbf{Phys.} The map the corpus named as its open construction is built, and it is made of pieces the derivations already produced. The family rotation is the inner $C_3$ of the weak algebra — the same rotation that forces the cosine shape of every channel spectrum — so a transport profile carried around the orbit is automatically a cosine triplet with the channel's own angle. Feed the triplet through the seed's Gram form and the mass-squared of each line comes out affine in the channel's spectral invariant: a constant set by the couplings, minus the mixed weak–hypercharge strength times the line's cosine. The only thing declared is the identification of the profile's phase with the channel angle; everything else — the orbit, the shape, the evaluation, the positivity, the sum rule — is forced. What the construction does not do is convert currencies: the transport species and the inertial $m=\mu r^2$ vocabulary still wait on the relating equation, and the overall scales remain the undeclared internal numbers they have always been.

\end{definition}

\begin{theorem}[{\texorpdfstring{BM.39 — The family mass form: positivity, the splitting law, the sum rule, the order inversion, and the area ladder}{BM.39 - The family mass form: positivity, the splitting law, the sum rule, the order inversion, and the area ladder}}]\label{tag:BM-39}
\textbf{Thm.}

%dsp-level:0
\[
(1)\quad\textbf{the exact evaluation and positivity:}
\]

%dsp-level:0
\[
\qquad m^2_i=M^2\bigl(A_0-B_0\cos\varphi_i\bigr),
\]

%dsp-level:0
\[
\qquad M^2:=\hat\mu^2\ell_c^2,\quad A_0=\frac{3g^2}8+\frac{g'^2}{16},\quad B_0=\frac{gg'}4
\]

%dsp-level:0
\[
\qquad [\mathrm{Def\text{-}BM.11}(5);\ \text{machine}];
\]

%dsp-level:0
\[
A_0-B_0=\frac38\Bigl(g-\frac{g'}3\Bigr)^2+\frac{g'^2}{48},
\]

%dsp-level:0
\[
\qquad A_0+B_0=\frac38\Bigl(g+\frac{g'}3\Bigr)^2+\frac{g'^2}{48}
\]

%dsp-level:0
\[
\qquad [\text{machine: exact sum-of-squares decompositions}];
\]

%dsp-level:0
\[
\implies M^2(A_0-B_0)\ \le\ m^2_i\ \le\ M^2(A_0+B_0),
\]

%dsp-level:0
\[
\qquad m^2_i>0\ \ \forall i,\ \forall\delta,\ \forall(g,g')\ne(0,0):
\]

%dsp-level:0
\[
\qquad \textbf{every line of every channel is massive under the construction};
\]

%dsp-level:0
\[
\textbf{the kernel line is unreachable:}
\]

%dsp-level:0
\[
\quad c\parallel(0,0,g',g)\ [\mathrm{EW.22}:\ \ker M^2=\operatorname{span}\{Q\},\ Q=W_3+Y;\ \mathrm{FT.12}]\ \text{requires two vanishing cosines}:
\]

%dsp-level:0
\[
\cos\varphi
\]

%dsp-level:0
\[
= 0\wedge\cos(\varphi+\tfrac{2\pi}3)
\]

%dsp-level:0
\[
= 0\ \text{has no solution}\ [\text{machine:}\ c=0\implies s^2=1\implies\cos(\varphi+\tfrac{2\pi}3)=\mp\tfrac{r_3}2,\ r_3^2=3\ne0]:
\]

%dsp-level:0
\[
\qquad \mathrm{BM.36}(5)\text{'s exact massless line carries no family row};
\]

%dsp-level:0
\[
(2)\quad\textbf{the splitting law:}\qquad m^2_i-m^2_j=M^2B_0\bigl(\cos\varphi_j-\cos\varphi_i\bigr)\qquad[\text{machine: exact factorization}]:
\]

%dsp-level:0
\[
\textbf{the intra-family splitting is exactly the mixed weak--hypercharge strength}\ gg'\ \textbf{acting on the transport phase}\ [\mathrm{Def\text{-}BM.11}\ \mathrm{Rmk}];
\]

%dsp-level:0
\[
\text{degenerate}\iff gg'=0\ \vee\ \cos\varphi_i=\cos\varphi_j;
\]

%dsp-level:0
\[
\qquad \textbf{the democratic witness}\ \delta
\]

%dsp-level:0
\[
= 0:\ \ \{\cos\varphi_i\}
\]

%dsp-level:0
\[
= \{1,-\tfrac12,-\tfrac12\}:\ \ \text{two lines degenerate, the third split by}\ \tfrac32M^2B_0\ [\text{machine}];
\]

%dsp-level:0
\[
(3)\quad\textbf{the sum rule:}\qquad\sum_im^2_i=3M^2A_0\qquad[\textstyle\sum_i\cos\varphi_i=0:\ \mathrm{FT.111}(3)/\mathrm{FT.158}(4)\ \text{identity};\ \text{machine}]:
\]

%dsp-level:0
\[
\textbf{phase-independent and channel-independent}\ —\ \text{the seed is channel-blind}\ [\mathrm{EW.32}:\ \neg\text{colour},\ \neg\text{slot index}]:
\]

%dsp-level:0
\[
\qquad \text{all three channels share}\ (A_0,B_0);\ \text{only}\ \delta\in\{\tfrac29,\tfrac{2}{27},\tfrac18\}\ \text{differs};
\]

%dsp-level:0
\[
(4)\quad\textbf{the cross-channel row:}
\]

%dsp-level:0
\[
\qquad 16\,\delta_\nu=9\,\delta_\ell=27\,\delta_q=2,
\]

%dsp-level:0
\[
\qquad 3\delta_q=\delta_\ell
\]

%dsp-level:0
\[
\qquad [\mathrm{FT.143}(5);\ \mathrm{FT.119}(2);\ \text{machine: exact Fractions}]:
\]

%dsp-level:0
\[
\text{the three channels' phase data are exact rational multiples of one seed numerator}\ —\ \text{the inter-channel mass relations are symbolic rows in}\ (A_0,B_0,M^2),
\]

%dsp-level:0
\[
\ \neg\text{numbers}\ [\mathrm{P4};\ \mathrm{P5}];
\]

%dsp-level:0
\[
(5)\quad\textbf{the order inversion:}
\]

%dsp-level:0
\[
\qquad m^2_i\ \text{is affine strictly decreasing in}\ r_i\ [\mathrm{Def\text{-}BM.11}(5)\ \text{affine form};\ B_0>0\ \text{at}\ g,g'>0\ [\mathrm{EW.33}:\ g,g'>0];\ \text{machine: sign census}]:
\]

%dsp-level:0
\[
\textbf{the spectral-maximal line receives the mass-minimum};
\]

%dsp-level:0
\[
\qquad \text{the frozen chains re-cited:}\ \ r_1^2<r_2^2<r_0^2\ [\mathrm{FT.113}\ (\ell);\ \mathrm{FT.143}(1),(3)\ (\nu);\ \text{enclosure-proven, frozen}];
\]

%dsp-level:0
\[
\textbf{the read-species row:}\quad\mathrm{FT.113}\text{'s inertial vocabulary}\ m^2
\]

%dsp-level:0
\[
= \mu r^2\ \text{orders with the spectrum; the transport species orders against it}:
\]

%dsp-level:0
\[
\text{the order is read-relative; neither species is adopted as ``the mass''}\ [\mathrm{Def\text{-}BM.11}(7)];
\]

%dsp-level:0
\[
\qquad Q=\tfrac23\ \text{is the structure invariant of}\ \mathrm{Def\text{-}BM.6}(4),\ \neg\text{a mass claim}\ [\mathrm{FT.143}(4)];
\]

%dsp-level:0
\[
(6)\quad\textbf{the inter-family constraint stands:}
\]

%dsp-level:0
\[
\qquad \mathrm{BM.36}(5)\text{'s rows are untouched by the construction:}\quad0\le G(c)\le\mu_\star^2\tfrac{g^2+g'^2}4\|c\|^2,
\]

%dsp-level:0
\[
\qquad \text{the pole ladder}\subseteq[0,4\hat\mu'^2]\ [\mathrm{FT.76}(3)];
\]

%dsp-level:0
\[
\text{the multi-order hierarchies remain beyond the computed cells}\ [\mathrm{FT.77},\ \mathrm{FT.79}(6)\ \text{re-cited; quarantined rows stay quarantined}]:
\]

%dsp-level:0
\[
\qquad \text{the construction supplies the form,}\ \neg\text{a scale}\ [\mathrm{P4};\ \mathrm{P5}];
\]

%dsp-level:0
\[
(7)\quad\textbf{the area ladder (the holonomy species):}
\]

%dsp-level:0
\[
\qquad m^2_{\mathrm{pole}}=4\hat\mu'^2\sin^2\bigl(\tfrac{\Theta_R}4\bigr),
\]

%dsp-level:0
\[
\qquad \Theta_R=\bar KA=\textstyle\sum_i\alpha_i-(n-2)\pi
\]

%dsp-level:0
\[
\qquad [\mathrm{Def\text{-}H.37}(2);\ \mathrm{FT.54}(4);\ \mathrm{FT.47}(1);\ \text{machine: the exact SU(2) defect re-derived}];
\]

%dsp-level:0
\[
\textbf{the extent-dependent form:}\quad\text{a loop of spatial extent}\ r\ \text{enclosing area}\ A(r)\ \text{in a constant-curvature slice:}
\]

%dsp-level:0
\[
\qquad m^2(r)=4\hat\mu'^2\sin^2\Bigl(\frac{\bar K\,A(r)}4\Bigr)
\]

%dsp-level:0
\[
\qquad [\text{machine: formal identity}];
\]

%dsp-level:0
\[
\textbf{the ratio formula:}
\]

%dsp-level:0
\[
\qquad \frac{m^2_1}{m^2_2}=\frac{\sin^2(\Theta_{R,1}/4)}{\sin^2(\Theta_{R,2}/4)}
\]

%dsp-level:0
\[
\qquad [\text{machine: symbolic}];
\]

%dsp-level:0
\[
\qquad \text{monotone in}\ \Theta_R\in[0,2\pi],\ \ \text{range}\ [0,4\hat\mu'^2]\ [\mathrm{FT.76}(3)\ \text{re-cited}];
\]

%dsp-level:0
\[
\textbf{the winding multiplet:}\quad m^2(m)=4\hat\mu'^2\sin^2\bigl(\tfrac{m\,\Theta_R}4\bigr),\quad m\in\{k,k{-}2,\dots,-k\}
\]

%dsp-level:0
\[
\qquad [\mathrm{FT.60}(5)\ \text{re-cited; Def-H.38 ladder}];
\]

%dsp-level:0
\[
(8)\quad\textbf{the distinction row:}
\]

%dsp-level:0
\[
\qquad \Theta_R\ \text{is the angle defect}\ [\mathrm{FT.47}(1)];
\]

%dsp-level:0
\[
\qquad \text{the winding}\ k=\Theta/2\ \text{is the ledger read}\ [\mathrm{Def\text{-}H.32}(4);\ \mathrm{FT.22}];
\]

%dsp-level:0
\[
\text{the identification}\ \Theta_R=2\pi k\ \text{is ruled out on calibration weight}\ [\mathrm{FT.55}(1)\ \text{re-cited}]:
\]

%dsp-level:0
\[
\qquad \textbf{the two ladders are distinct reads, and the tier keeps them apart};
\]

%dsp-level:0
\[
(9)\quad\textbf{the matching row:}
\]

%dsp-level:0
\[
\qquad \text{the Gram species}\ ((1)\text{--}(5))\ \text{and the holonomy species}\ ((7))\ \text{are related by}\ \mathrm{Def\text{-}H.37}\text{'s leading-order matching}\ \hat\mu'^2\Theta_R^2/4
\]

%dsp-level:0
\[
= \hat\mu^2\ell_c^2G(c)\ [\mathrm{FT.54}(2);\ \mathrm{FT.33}]:
\]

%dsp-level:0
\[
\textbf{recorded, not imposed}\ [\mathrm{Def\text{-}H.37}\ \text{locus clause}];
\]

%dsp-level:0
\[
\qquad \text{the exact-matching derivation remains item-7 content}\ [\text{pointer}].
\]

\textbf{Rmk (two mechanisms, one form.)} Within a channel the families split by the Gram phase — the mixed strength times a cosine; along a winding or area ladder the levels split by the holonomy defect — a sine-squared of the enclosed curvature. The two are the leading-order and exact forms of one transfer [Def-H.37(1)–(2)], and the tier uses each where it is derived: the cosine triplet for the family structure, the defect ladder for the extent structure. Neither carries a scale: $\hat\mu,\hat\mu',\ell_c,(g,g')$ are the logged internal data, and no clause converts them into a number.

\textbf{Req.} P4, P5 (Parts 1–2); EW.22, EW.32, EW.33 (Part 2); Def-H.32(4), Def-H.37, Def-H.38, FT.12, FT.22, FT.33, FT.47(1), FT.54(2), FT.54(4), FT.55(1), FT.60(5), FT.77, FT.79(6), FT.111(3), FT.112, FT.113, FT.119(1), FT.119(2), FT.142(5), FT.143(1), FT.143(3), FT.143(4), FT.143(5), FT.158(4) (Part 3); Census FT.76(3) (row-wise); Def-BM.6(4), Def-BM.11, BM.36(5), BM.36(6) (this part); Cal-F; sum-of-squares positivity, exact quotient-ring identities, symbolic ratio algebra, exact Pauli defect identities [math; machine: \texttt{checks/\allowbreak{}verify\_\allowbreak{}H\_\allowbreak{}neutron\_\allowbreak{}decay\_\allowbreak{}vertex.py} §A].

\textbf{Pf.}

%dsp-level:0
\[
(1):\ \mathrm{Def\text{-}BM.11}(5)\text{'s evaluation re-cited};
\]

%dsp-level:0
\[
\qquad \tfrac38(g\mp\tfrac{g'}3)^2=\tfrac{3g^2}8\mp\tfrac{gg'}4+\tfrac{g'^2}{24}:\quad+\tfrac{g'^2}{48}\ \text{gives}\ A_0\mp B_0\ [\text{machine: expansion}];
\]

%dsp-level:0
\[
\qquad |\cos|\le1\implies\text{the sandwich};
\]

%dsp-level:0
\[
\qquad \text{positivity: a sum of squares with positive rational coefficients vanishes only at}\ (g,g')
\]

%dsp-level:0
\[
= (0,0)\ [\text{machine}]
\]

%dsp-level:0
\[
\text{the kernel row:}\ c=0\wedge s^2=1\implies\cos(\varphi+\tfrac{2\pi}3)=-\tfrac c2-\tfrac{r_3s}2=\mp\tfrac{r_3}2\ne0\ [r_3^2=3;\ \text{machine}]
\]

%dsp-level:0
\[
(2):\ \text{difference of}\ (1)\text{'s forms}\ [\text{machine}];
\]

%dsp-level:0
\[
\qquad \delta=0:\ c=1,s=0\implies\{1,-\tfrac12,-\tfrac12\}\ [\text{machine:}\ -\tfrac{c\pm r_3s}2\ \text{at}\ (c,s)=(1,0)]
\]

%dsp-level:0
\[
(3):\ \textstyle\sum_i\cos\varphi_i=0\ [\mathrm{FT.111}(3)\ \text{identity, generic phase};\ \text{machine}];
\]

%dsp-level:0
\[
\qquad (A_0,B_0)\ \text{carry no channel argument}\ [\mathrm{EW.32}\ \text{re-cited}]
\]

%dsp-level:0
\[
(4):\ 16\cdot\tfrac18=9\cdot\tfrac29=27\cdot\tfrac{2}{27}=2;\quad3\cdot\tfrac{2}{27}=\tfrac29\ [\text{machine: exact Fractions}]
\]

%dsp-level:0
\[
(5):\ \text{the affine coefficient}\ -M^2B_0/r_2<0\ \text{at}\ g,g'>0\ [\text{machine: sign census of the monomial}];
\]

%dsp-level:0
\[
\qquad \mathrm{FT.113}/\mathrm{FT.143}\text{'s chains re-cited (frozen, enclosure-proven)};
\]

%dsp-level:0
\[
\qquad Q=\tfrac69=\tfrac23\ [\mathrm{FT.143}(4)\ \text{re-cited}]
\]

%dsp-level:0
\[
(6):\ \mathrm{BM.36}(5),\ \mathrm{FT.76}(3),\ \mathrm{FT.77},\ \mathrm{FT.79}(6)\ \text{re-cited}\ (\neg\text{recomputed})
\]

%dsp-level:0
\[
(7):\ (U-1)^\dagger(U-1)
\]

%dsp-level:0
\[
= (2-2\cos\phi)\mathbb I
\]

%dsp-level:0
\[
= 4\sin^2(\tfrac\phi2)\mathbb I\ \text{at}\ U
\]

%dsp-level:0
\[
= \cos\phi\,
\]

%dsp-level:0
\[
\mathbb I+i\sin\phi\,
\]

%dsp-level:0
\[
\hat n\cdot\vec\sigma\ [\mathrm{FT.54}(1)\ \text{re-derived};\ \text{machine: exact Pauli algebra in}\ \mathbb Q[C,S]/(C^2{+}S^2{-}1)];
\]

%dsp-level:0
\[
\qquad \phi=\tfrac{\Theta_R}2\ [\mathrm{FT.54}(3);\ \mathrm{Def\text{-}H.37}(2)]
\]

%dsp-level:0
\[
(8):\ \mathrm{FT.55}(1)\text{'s calibration-weight obstruction re-cited};
\]

%dsp-level:0
\[
\qquad (9):\ \mathrm{Def\text{-}H.37}\text{'s matching row re-cited}.\ \blacksquare
\]

\textbf{Phys.} The family mass form comes out exact and tight. Positivity is a sum of two squares with rational coefficients, so no line of any channel can fall to zero through this map — the massless direction of the Gram form is the electromagnetic generator, and a cosine triplet can never align with it, because two cosines a third of a turn apart cannot vanish together. The splitting between lines is precisely the product of the two electroweak strengths times the cosine separation: the force structure acting on the transport phase, with no Yukawa object anywhere in it. The three channels share the whole form and differ only in their derived angles, so their orbit sums are one expression and their splittings are rational multiples of one another. The order runs against the spectrum: the brightest spectral line is the lightest mass under this read, and the tier records that reversal instead of hiding it, because the inertial vocabulary orders the other way and no theorem yet converts the two. Alongside the Gram mechanism stands the holonomy one: mass as the sine-squared of a quarter angle defect, growing with the area a loop encloses — the extent-dependent form the corpus can derive exactly, with the winding identification ruled out and the two ladders kept apart.

\end{theorem}

\begin{theorem}[{\texorpdfstring{BM.40 — The rung-1 isospin row: the Casimir identity, the branch selectivity, and the neutrality of the $\sigma$-vertex}{BM.40 - The rung-1 isospin row: the Casimir identity, the branch selectivity, and the neutrality of the -vertex}}]\label{tag:BM-40}
\textbf{Thm.}

%dsp-level:0
\[
(1)\quad\textbf{the Casimir identity:}
\]

%dsp-level:0
\[
\qquad W_i=\tfrac12\bigl(\sigma_i\otimes\sigma_0+\sigma_0\otimes\sigma_i\bigr)\ [\mathrm{CT1.2}];
\]

%dsp-level:0
\[
\qquad [W_i,W_j]=i\epsilon_{ijk}W_k\ [\mathrm{CT1.4};\ \text{machine}];
\]

%dsp-level:0
\[
\boxed{\ \textstyle\sum_iW_i^2=\tfrac32E_0+\tfrac12\Sigma=E_0+\mathbb S=2\,\Pi_{\mathcal R}\ }
\]

%dsp-level:0
\[
\qquad [\mathrm{Cor}\ \mathrm{CT0.7a};\ \mathrm{FT.141}(4);\ \mathrm{FT.157}(2);\ \text{machine: entrywise over}\ \mathbb Q(i),\ \text{all}\ 16\ \text{components}];
\]

%dsp-level:0
\[
(2)\quad\textbf{the branch selectivity:}
\]

%dsp-level:0
\[
\qquad \textstyle\sum_iW_i^2\big|_{\vee^2}=2=1(1{+}1),
\]

%dsp-level:0
\[
\qquad \textstyle\sum_iW_i^2\big|_{\wedge^2}=0=0(0{+}1)
\]

%dsp-level:0
\[
\qquad [\mathrm{Cor}\ \mathrm{CT0.7a};\ \text{machine: exact on the branch bases}];
\]

%dsp-level:0
\[
W_i\Pi_{\mathcal V}=0=\Pi_{\mathcal V}W_i\ \forall i,\qquad\bigcap_i\ker W_i=\operatorname{im}\Pi_{\mathcal V}\qquad[\mathrm{FT.101}(1)\text{--}(2);\ \mathrm{FT.157}(2);\ \text{machine}];
\]

%dsp-level:0
\[
\Pi_{\mathcal R}=\tfrac14(3E_0+\Sigma),\quad\Pi_{\mathcal V}=\tfrac14(E_0-\Sigma):
\]

%dsp-level:0
\[
\qquad \text{idempotent, orthogonal, complete},\quad\mathfrak T\Pi_{\mathcal R}=3,\quad\mathfrak T\Pi_{\mathcal V}=1
\]

%dsp-level:0
\[
\qquad [\mathrm{CT0.7}(1);\ \text{machine}];
\]

%dsp-level:0
\[
(3)\quad\textbf{the neutrality of the}\ \sigma\textbf{-vertex:}
\]

%dsp-level:0
\[
\qquad \text{the transition's charge move is an end-sign flip:}
\]

%dsp-level:0
\[
\quad q_{\mathrm{core}}=\tfrac12(\sigma_s-\sigma_t)\ \text{is registration data on the reality line}\ [\mathrm{D3.4};\ \mathrm{D1.1};\ \mathrm{CT0.9};\ \mathrm{BM.23}(2)];
\]

%dsp-level:0
\[
\text{the Casimir}\ 2\Pi_{\mathcal R}\ \text{is a function of the branch datum alone}\ [(1)];
\]

%dsp-level:0
\[
\qquad \text{the gradings are pairwise distinct}\ [\mathrm{FT.136};\ \mathrm{FT.141}(6);\ \mathrm{FT.144}(5)];
\]

%dsp-level:0
\[
\qquad \text{the register separation holds}\ [\mathrm{FT.136}(3)];
\]

%dsp-level:0
\[
\boxed{\ \textbf{the rung-1 Casimir deficit at the}\ \sigma\textbf{-vertex is identically zero}\ }
\]

%dsp-level:0
\[
\qquad [\text{machine: the token census}\ —\ \text{zero branch tokens in}\ \mathrm{D3.4}\text{'s defining row, zero registration tokens in}\ \mathrm{CT0.7}(1)/\mathrm{FT.101}\text{'s defining rows}];
\]

%dsp-level:0
\[
\text{the conservation law holds at the vertex}\ —\ \text{it neither forces nor forbids an emission}:
\]

%dsp-level:0
\[
(4)\quad\textbf{the carrier-side rows:}
\]

%dsp-level:0
\[
\qquad \text{every emitted chain is Janus-closed:}\quad\#J^+=\#J^-\ [\mathrm{D1.4};\ \mathrm{FT.144}(1)\text{--}(2);\ \text{the prefix law}]:
\]

%dsp-level:0
\[
\qquad \text{each carrier is internally Casimir-balanced};
\]

%dsp-level:0
\[
\text{the neutral carrier class is rung-one content:}\quad\mathcal A(1)\ \text{aspects},
\]

%dsp-level:0
\[
\ \mathcal K_1\ \text{arena},
\]

%dsp-level:0
\[
\ \dim
\]

%dsp-level:0
\[
= 16\ [\mathrm{FT.142}(1)\text{--}(2)];
\]

%dsp-level:0
\[
\qquad \text{EM-unactivated}\ (f_v=0),\ \text{colour-blind}\ [\mathrm{FT.142}(1)];
\]

%dsp-level:0
\[
\text{at rung one the chiral sectors are force-decoupled:}
\]

%dsp-level:0
\[
\quad F_{LR}=0\ [\mathrm{FT.144}(6);\ \mathrm{FT.128}(2);\ \text{the no-flip theorem}]:
\]

%dsp-level:0
\[
\qquad \textbf{the neutral carrier's orientation sector is frozen along its flight}\ [\mathrm{FT.144}(5):\ \text{single-sector occupancy}];
\]

%dsp-level:0
\[
(5)\quad\textbf{the status:}
\]

%dsp-level:0
\[
\qquad \text{the forcing of the neutral member comes}\ \neg\text{from the Casimir row}\ —\ (3)\ \text{adjudicates the deficit route neutral}\ —\ \text{but from}\ \mathrm{BM.41}\text{'s kinematic theorem};
\]

%dsp-level:0
\[
\mathrm{BM.23}(6)\text{'s clause stands as written and receives an on-record Rmk}\ [\text{the FT.89 Rmk pattern}];
\]

%dsp-level:0
\[
\qquad \neg\exists\ \text{a lepton-number read in the inventory}\ [\mathrm{BM.23}(6);\ \mathrm{FT.81}(4)\ \text{pattern}]\ —\ \text{the row is untouched}.
\]

\textbf{Rmk (why the proposed deficit does not arise).} The isospin bookkeeping the vertex could upset lives in the branch datum — aligned pairs carry Casimir two, antialigned pairs zero. The charge move that defines the transition lives in the registration sign — the reality-line datum the end reads consume. The corpus proved these gradings pairwise distinct and their registers separate, so no $\sigma$-flip can move a $\Pi_{\mathcal R}$: the deficit is zero as an operator identity, not as a small number. That is why the neutrino's necessity, where it is necessity, must come from somewhere else — and BM.41 shows it does: from the kinematics of equal rest masses and the uniqueness of the neutral class.

\textbf{Req.} A3 (Part 1); CT0.7(1), CT0.9, CT1.2, CT1.4, Cor CT0.7a, D1.1, D1.4, D3.4 (Parts 1–2); P5 (Parts 1–2); FT.12, FT.101(1), FT.101(2), FT.128(2), FT.136(3), FT.141(4), FT.141(6), FT.142(1), FT.142(2), FT.144(1), FT.144(2), FT.144(5), FT.144(6), FT.157(2) (Part 3); BM.7, BM.23(2), BM.23(6) (this part); the forcing theorem of this tier's vertex tag is a marked forward pointer at the citing site and is imported in no \textbf{Req} line; exact $\mathbb Q(i)$ $4\times4$ algebra, projector algebra, token censuses [math; machine: \texttt{checks/\allowbreak{}verify\_\allowbreak{}H\_\allowbreak{}neutron\_\allowbreak{}decay\_\allowbreak{}vertex.py} §B].

\textbf{Pf.}

%dsp-level:0
\[
(1):\ \textstyle\sum_iW_i^2
\]

%dsp-level:0
\[
= \tfrac14\sum_i(\sigma_i\otimes\sigma_0+\sigma_0\otimes\sigma_i)^2
\]

%dsp-level:0
\[
= \tfrac14(3E_0+3E_0+2\Sigma)
\]

%dsp-level:0
\[
= \tfrac32E_0+\tfrac12\Sigma\ [\mathrm{Cor}\ \mathrm{CT0.7a}\ \text{re-derived};\ \text{machine}];
\]

%dsp-level:0
\[
\qquad \Sigma=2\mathbb S-E_0\ [\mathrm{CT0.7}(1)]\implies=E_0+\mathbb S=2\Pi_{\mathcal R}\ [\text{machine}];
\]

%dsp-level:0
\[
\qquad [W_i,W_j]=i\epsilon_{ijk}W_k\ [\text{machine: all}\ 9\ \text{pairs}]
\]

%dsp-level:0
\[
(2):\ \vee^2\ \text{basis}\ \{(1,0,0,0),(0,1,1,0),(0,0,0,1)\}:\ 2\cdot v\ [\text{machine}];
\]

%dsp-level:0
\[
\qquad \wedge^2\ \text{basis}\ \{(0,1,-1,0)\}:\ 0\ [\text{machine}];
\]

%dsp-level:0
\[
\qquad W_i\Pi_{\mathcal V}=0=\Pi_{\mathcal V}W_i\ [\text{machine: all}\ 16\ \text{components each}];
\]

%dsp-level:0
\[
\qquad \text{the projector algebra and the traces}\ [\text{machine}]
\]

%dsp-level:0
\[
(3):\ \text{the defining strings extracted}: \mathrm{D3.4}\text{'s row vs}\ \mathrm{CT0.7}(1)/\mathrm{FT.101}\text{'s rows};
\]

%dsp-level:0
\[
\qquad \text{the branch alphabet}\ \{\Pi_{\mathcal R},\Pi_{\mathcal V},\alpha,\text{branch}\}:\ 0\ \text{occurrences in}\ \mathrm{D3.4};
\]

%dsp-level:0
\[
\qquad \text{the registration alphabet}\ \{\sigma_s,\sigma_t,q_{\mathrm{core}},\text{reality line}\}:\ 0\ \text{occurrences in the branch rows}\ [\text{machine}]
\]

%dsp-level:0
\[
(4):\ \mathrm{D1.4}/\mathrm{FT.144}(1)\text{--}(2)\ \text{re-cited (frozen)};
\]

%dsp-level:0
\[
\qquad F_{LR}=0\ [\mathrm{FT.128}(2)\ \text{re-cited}];
\]

%dsp-level:0
\[
\qquad (5):\ \text{the status record}.\ \blacksquare
\]

\textbf{Phys.} The weak isospin of the pair arena is an exact operator identity — the Casimir is twice the projector onto the real branch — and it splits the pair register into a spin-one triplet the weak force sees and a spin-zero singlet it annihilates. The question the vertex poses is whether a charge flip disturbs that bookkeeping, and the answer is a clean no: the flip moves a registration sign, the Casimir reads a branch label, and the corpus long ago proved those gradings live in separate registers. The deficit is zero identically, so the isospin law neither demands nor forbids the neutral emission. What the rung-one rows do supply is the carrier's pedigree: closed under Janus balance, electromagnetically inert, colour-blind, force-decoupled across its chiral sectors at rung one, single-arrow in occupancy — a neutral member that travels without turning. The necessity itself waits for the kinematics.

\end{theorem}

\begin{theorem}[{\texorpdfstring{BM.41 — The decay vertex: the forced members, the two-body impossibility, and the Q-value sign theorem}{BM.41 - The decay vertex: the forced members, the two-body impossibility, and the Q-value sign theorem}}]\label{tag:BM-41}
\textbf{Thm.}

%dsp-level:0
\[
(1)\quad\textbf{the vertex grammar:}\qquad\begin{array}{l|l}\text{the parent} & \text{the neutral-class}\ \varepsilon\text{-sealed 3-juncture},\ E_{\mathrm{int}}=6,\ Q=0\ [\mathrm{BM.17};\ \mathrm{Def\text{-}BM.3}(2);\ \text{Census FT.35}(2)]\\\text{the daughter} & \text{the unit-class}\ \varepsilon\text{-sealed 3-juncture},\ E_{\mathrm{int}}=6,\ Q=+1\ [\mathrm{BM.26};\ \mathrm{Def\text{-}BM.3}(3)]\\\text{the charged member} & \text{the}\ \ell\text{-channel line},\ q=-1:\ \text{forced and class-unique}\ [\mathrm{BM.23}(2)\text{--}(5):\ \Delta q=\pm1,\ \text{the unique compensating value},\ \text{the selection theorem}]\\\text{the neutral member} & \text{the}\ \nu\text{-channel line},\ q=0:\ \text{forced by}\ (2)\text{--}(3);\ \text{class-unique}\ [\mathrm{FT.142}(1);\ \mathrm{BM.23}(5)\text{'s eliminations}]\end{array}
\]

%dsp-level:0
\[
\textbf{the charge balance:}\quad0=(+1)+(-1)+0\ [\mathrm{EM.11};\ \mathrm{BM.23}(4);\ \text{machine: integers}];
\]

%dsp-level:0
\[
\qquad \textbf{the sealing phase:}\ \Theta_{\mathrm{int}}\in4\mathbb Z\ \text{on both baryon-type junctions}\ [\mathrm{BM.2}(1);\ \text{Census FT.49}(3)];
\]

%dsp-level:0
\[
\textbf{the arrow:}\ \text{forward-only}\ [\mathrm{BM.23}(9);\ \mathrm{FT.144}(2);\ \mathrm{A6}:\ \neg\exists\tau^{-1}];
\]

%dsp-level:0
\[
\qquad \textbf{the transition weight:}\ \varepsilon_q=0\implies\Gamma=0\ \text{exactly}\ [\mathrm{BM.23}(8)]:\ \text{no flicker, no decay};
\]

%dsp-level:0
\[
(2)\quad\textbf{the two-body impossibility (exact):}
\]

%dsp-level:0
\[
\qquad m_n=m_p=\frac{6\ln2}{\mathfrak h}
\]

%dsp-level:0
\[
\qquad [\mathrm{BM.17};\ \mathrm{BM.26};\ \mathrm{BM.28}(1)\ \text{charge-blindness};\ \mathrm{BM.34}(7)/\mathrm{BM.35}(4)\ \text{the rest slice}];
\]

%dsp-level:0
\[
\text{the rest frame, two members:}\quad E_p+E_\ell=m_n,
\]

%dsp-level:0
\[
\qquad E_i=\sqrt{m_i^2+p^2}\ [\mathrm{BM.34}(4)\ \text{vacuum gauge};\ \mathrm{BM.11}]:
\]

%dsp-level:0
\[
\qquad E_\ell=\frac{m_n^2+m_\ell^2-m_p^2}{2m_n},
\]

%dsp-level:0
\[
\quad E_p=\frac{m_n^2+m_p^2-m_\ell^2}{2m_n}\ [\text{machine: the exact solution of the conservation system}];
\]

%dsp-level:0
\[
m_n=m_p:\qquad\boxed{\ E_p-m_p=-\frac{m_\ell^2}{2m_n}<0\ \ \forall m_\ell>0\ }\qquad[\text{machine: cleared polynomial identity}\ +\ \text{sign census}]:
\]

%dsp-level:0
\[
\textbf{the two-body vertex puts the daughter below its own mass shell for every positive member mass}\ —\ \neg\exists\ \text{a real on-shell solution};
\]

%dsp-level:0
\[
\textbf{the Källén companion:}\quad4m_n^2p^2=\lambda\bigl(m_n^2,m_p^2,m_\ell^2\bigr),
\]

%dsp-level:0
\[
\qquad \lambda=\bigl(m_n^2-(m_p{+}m_\ell)^2\bigr)\bigl(m_n^2-(m_p{-}m_\ell)^2\bigr)\ [\text{machine: polynomial identity}];
\]

%dsp-level:0
\[
m_n=m_p:\quad\lambda=m_\ell^2\bigl(m_\ell^2-4m_n^2\bigr):
\]

%dsp-level:0
\[
\qquad \text{spacelike for}\ 0<m_\ell<2m_n;
\]

%dsp-level:0
\[
\qquad m_\ell\ge2m_n\implies E_p\le-m_n<0\ [\text{machine: the two branches cover every}\ m_\ell>0\ (2>\sqrt2\ \text{by squaring})];
\]

%dsp-level:0
\[
(3)\quad\textbf{the forced neutral member:}
\]

%dsp-level:0
\[
\qquad \text{by}\ (2)\ \text{every kinematically admissible vertex carries}\ \ge3\ \text{members};
\]

%dsp-level:0
\[
\qquad \text{the charge balance fixes the third member's total charge to}\ 0\ [(1);\ \mathrm{EM.11}];
\]

%dsp-level:0
\[
\textbf{the class is unique among the constructed carrier classes:}\quad\text{the}\ \nu\text{-channel is the unique EM-unactivated colour-blind bound-chain class}\ [\mathrm{BM.23}(5)\text{'s elimination rows};\ \mathrm{FT.142}(1)];
\]

%dsp-level:0
\[
\text{the transit (EM-radiation) alternative is}\ \neg\text{constructed: the gauge coupling of the}\ Q\ \text{read is deficit}\ 1'\ [\mathrm{FT.125}(4);\ \mathrm{Def\text{-}BM.8}(8)]\ —\ \neg\text{an inventory candidate};
\]

%dsp-level:0
\[
\boxed{\ \textbf{the neutral}\ \nu\text{-channel member is forced in every kinematically admissible vertex}\ }
\]

%dsp-level:0
\[
\qquad [\text{the forcing is kinematic}\ +\ \text{uniqueness},\ \neg\text{a number conservation}:\ \neg\exists\ \text{a lepton-number read}\ [\mathrm{BM.23}(6)]];
\]

%dsp-level:0
\[
\textbf{the line index within the}\ \nu\textbf{-triplet is}\ \neg\textbf{derived}\ [\neg\text{a mixing matrix}\ [\mathrm{FT.145}(3)]];
\]

%dsp-level:0
\[
\qquad \mathrm{BM.23}(6)\ \text{is upgraded on record: admissible}\ \longrightarrow\ \text{conditionally forced}\ [\text{its Rmk}];
\]

%dsp-level:0
\[
(4)\quad\textbf{the Q-value sign theorem:}\qquad Q:=m_n-m_p-m_\ell-m_{\nu};
\]

%dsp-level:0
\[
(\mathrm a)\ \textbf{the crystallized informational read:}
\]

%dsp-level:0
\[
\quad Q=-(m_\ell+m_{\nu})\ [\text{machine: monomial identity at}\ m_n=m_p];
\]

%dsp-level:0
\[
\qquad m_\ell>0\ \text{on every read species on record}:
\]

%dsp-level:0
\[
\text{transport:}\ m_\ell^2\ge M^2(A_0-B_0)>0\ [\mathrm{BM.39}(1)];
\]

%dsp-level:0
\[
\qquad \text{informational:}\ m_\ell
\]

%dsp-level:0
\[
= \frac{E_{\mathrm{int}}\ln2}{\mathfrak h}\ge\frac{2\ln2}{\mathfrak h}>0\ [\mathrm{FT.44}(4);\ \mathrm{FT.28}\ \text{the minimal floor};\ \tfrac12<\ln2<1\ \text{enclosure};\ \text{machine}];
\]

%dsp-level:0
\[
\qquad m_\nu\ge0\ [\mathrm{EW.33}\ \mathrm{PSD}];
\]

%dsp-level:0
\[
\implies\boxed{\ Q<0\ \ \text{strictly}\ };
\]

%dsp-level:0
\[
(\mathrm b)\ \textbf{every normalized ensemble read:}\quad\text{charge-degeneracy is structural}\ [\mathrm{FT.107}(3);\ \mathrm{BM.30}(3);\ \mathrm{BM.37}(2)\ \text{the disjointness}]
\]

%dsp-level:0
\[
\implies Q
\]

%dsp-level:0
\[
= -(m_\ell+m_\nu)<0;
\]

%dsp-level:0
\[
(\mathrm c)\ \textbf{the curvature-flux read:}
\]

%dsp-level:0
\[
\quad RV=\varkappa^2m\ \text{with identical}\ (E_{\mathrm{int}},V)\ \text{on both classes}\ [\mathrm{BM.20}(1);\ \mathrm{BM.28}(1)]\implies Q<0;
\]

%dsp-level:0
\[
(\mathrm d)\ \textbf{the route-(ii) variant:}\quad T(0)=\tfrac43T(\pm1)\ [\mathrm{BM.28}(4)(\mathrm{ii})]:
\]

%dsp-level:0
\[
\qquad Q_{\mathrm{ii}}=\tfrac13T(\pm1)-m_\ell-m_\nu:
\]

%dsp-level:0
\[
\qquad \textbf{the sign is undetermined without scale data}\ [\text{machine: two rational witnesses of opposite sign}];
\]

%dsp-level:0
\[
\text{the functional is a measure price}\ [\mathrm{BM.37}(3)]\ \text{and its adoption requires a §I addendum}\ [\mathrm{BM.30}(5)]:
\]

%dsp-level:0
\[
\qquad \text{priced},\ \neg\text{adopted};
\]

%dsp-level:0
\[
(\mathrm e)\ \textbf{the verdict:}
\]

%dsp-level:0
\[
\qquad \boxed{\ \textbf{the}\ \beta\textbf{-decay vertex is kinematically closed on every adopted read: the corpus does not make the neutron heavier than the proton}\ }
\]

%dsp-level:0
\[
[\mathrm{P5}:\ \neg\text{an external mass-difference datum};\qquad\text{the blocks are theorems:}\ \mathrm{BM.28}(1),\ \mathrm{FT.107}(3),\ \mathrm{BM.37}(3),\ \mathrm{BM.30}(5)];
\]

%dsp-level:0
\[
(5)\quad\textbf{the three-body kinematics (the conditional section, under the hypothesis}\ Q\ge0\ \textbf{of some adopted specification):}
\]

%dsp-level:0
\[
\qquad \text{the rest frame:}\quad0=\vec p_p+\vec p_\ell+\vec p_\nu\ [\mathrm{BM.11};\ \mathrm{BM.34}(3)];
\]

%dsp-level:0
\[
Q=\mathrm{KE}_p+\mathrm{KE}_\ell+\mathrm{KE}_\nu,
\]

%dsp-level:0
\[
\qquad \mathrm{KE}_i=E_i-m_i
\]

%dsp-level:0
\[
\qquad [\mathrm{BM.34}(5);\ \text{machine: Laurent identity}];
\]

%dsp-level:0
\[
\qquad \text{NR:}\ \ \mathrm{KE}_i=\frac{p_i^2}{2m_i}+O(p^4);
\]

%dsp-level:0
\[
\textbf{the endpoint:}
\]

%dsp-level:0
\[
\qquad \boxed{\ E_\ell^{\max}=\frac{m_n^2+m_\ell^2-(m_p+m_\nu)^2}{2m_n}\ }
\]

%dsp-level:0
\[
\qquad [\text{attained at the relative-rest recoil};\ \text{machine: the exact polynomial identities}];
\]

%dsp-level:0
\[
E_\ell^{\max}+E_{\mathrm{rec}}=m_n,\quad E_{\mathrm{rec}}=\frac{m_n^2+(m_p+m_\nu)^2-m_\ell^2}{2m_n};
\]

%dsp-level:0
\[
\qquad \bigl(E_\ell^{\max}\bigr)^2-m_\ell^2=E_{\mathrm{rec}}^2-(m_p+m_\nu)^2\ [\text{machine: cleared by}\ 4m_n^2];
\]

%dsp-level:0
\[
\textbf{the recoil bound:}\quad s_{p\nu}\ge(m_p+m_\nu)^2:
\]

%dsp-level:0
\[
\qquad \text{the collinear identity}\ \bigl(E_pE_\nu-PN\bigr)^2
\]

%dsp-level:0
\[
= (m_pm_\nu)^2+\bigl(E_pN-E_\nu P\bigr)^2\ [\text{machine: on-shell quotient ring}];
\]

%dsp-level:0
\[
\text{the general case by Cauchy--Schwarz:}
\]

%dsp-level:0
\[
\quad P^2N^2-(\vec p_p\cdot\vec p_\nu)^2=\sum_{i<j}\bigl(p_{p,i}p_{\nu,j}-p_{p,j}p_{\nu,i}\bigr)^2\ge0\ [\text{machine: the Lagrange identity}];
\]

%dsp-level:0
\[
\qquad E_\ell^{\max}\ \text{is strictly decreasing in}\ s_{p\nu}\ [\text{the coefficient}\ -\tfrac1{2m_n}<0];
\]

%dsp-level:0
\[
\text{the continuum of}\ E_\ell\ \text{values on}\ [E_\ell^{\min},E_\ell^{\max}]\ \text{is chart-side content}\ [\mathrm{CT2.6}\ \text{hypotheses at the point of use}];
\]

%dsp-level:0
\[
(6)\quad\textbf{the general channel closure:}
\]

%dsp-level:0
\[
\qquad \text{any member set containing the daughter and at least one positive-mass chain line satisfies}\ \textstyle\sum m_{\mathrm{out}}>m_n\ [(2)\text{'s equality of rest reads};\ \text{machine: sign arithmetic}]:
\]

%dsp-level:0
\[
\textbf{every material decay channel of the neutron row is kinematically closed on the adopted reads}\ —\ \text{the closure is a theorem about the reads,}\ \neg\text{about the vertex grammar}\ [(1)\ \text{stands complete}];
\]

%dsp-level:0
\[
(7)\quad\textbf{status and scope:}
\]

%dsp-level:0
\[
\qquad \text{the vertex grammar}\ (1)\ \text{and the forcing}\ (3)\ \text{are derived};
\]

%dsp-level:0
\[
\qquad \text{the kinematic viability is read-relative}\ (4);
\]

%dsp-level:0
\[
\qquad \text{the conditional section}\ (5)\ \text{carries its hypothesis on its face};
\]

%dsp-level:0
\[
\neg\text{a measured lifetime, mass difference or endpoint enters}\ [\mathrm{P5};\ \text{machine: zero external tokens}];
\]

%dsp-level:0
\[
\qquad \text{the row-index attachment}\ (3)\ \text{and the read choice}\ (4)(\mathrm d)\ \text{are the open faces}.
\]

\textbf{Rmk (the shape of the negative.)} The corpus assembles this vertex completely and then closes it — and the closing is as exact as the assembly. Two independent theorems do the work: the charge-blindness of every crystallized count read, which makes the two baryon rows weigh the same to the last monomial; and the structural charge-degeneracy of every normalized ensemble average, which makes the flickering rows weigh the same too. The one functional that does split them was classified as a measure price, and adopting it as a mass is an act of specification the corpus prices but does not perform. So the reaction stands: grammar forced, members forced, conservation exact, kinematics closed — waiting on a read specification, not on a missing derivation.

\textbf{Req.} A6 (Part 1); D1.4, D2.9, D2.11 (Parts 1–2); EM.9, EM.11 (Part 2); P5 (Parts 1–2); CT2.6 (Part 2); FT.22, FT.107(3), FT.125(4), FT.142(1), FT.144(2), FT.145(3) (Part 3); Census FT.28, FT.35(2), FT.44(4), FT.49(3) (row-wise); Def-BM.3, Def-BM.6, Def-BM.8, Def-BM.10, BM.2(1), BM.11, BM.17, BM.23(2), BM.23(4), BM.23(5), BM.23(6), BM.23(8), BM.23(9), BM.26, BM.28(1), BM.28(4), BM.30(3), BM.30(5), BM.34(3), BM.34(4), BM.34(5), BM.34(7), BM.35(4), BM.36(5), BM.37(2), BM.37(3), BM.39(1), BM.40(3) (this part); Källén polynomial identities, exact sign arithmetic over positive witnesses, the collinear and Lagrange Cauchy–Schwarz identities, exact rational $\ln2$ enclosures [math; machine: \texttt{checks/\allowbreak{}verify\_\allowbreak{}H\_\allowbreak{}neutron\_\allowbreak{}decay\_\allowbreak{}vertex.py} §C].

\textbf{Pf.}

%dsp-level:0
\[
(1):\ \mathrm{BM.23}(2)\text{--}(5)\ \text{re-cited (frozen)};
\]

%dsp-level:0
\[
\qquad 0=1+(-1)+0\ [\text{machine}];
\]

%dsp-level:0
\[
\qquad \Theta_{\mathrm{int}}\in4\mathbb Z\ [\text{Census FT.49}(3)\ \text{re-cited}];
\]

%dsp-level:0
\[
\qquad \Gamma(\varepsilon_q{=}0)=0\ [\mathrm{BM.23}(8)\ \text{re-cited}]
\]

%dsp-level:0
\[
(2):\ \text{the conservation system}\ \{E_p+E_\ell=m_n;\ E_\ell^2-m_\ell^2=E_p^2-m_p^2\}\ \text{solved linearly in}\ (E_p,E_\ell)\ [\text{machine: substitution verifies both equations as polynomial identities over}\ \mathbb Q[m_n,m_p,m_\ell]];
\]

%dsp-level:0
\[
m_n=m_p:\quad2m_n(E_p-m_p)=(m_n^2+m_p^2-m_\ell^2)-2m_nm_p=-m_\ell^2\ [\text{machine: cleared identity}];
\]

%dsp-level:0
\[
\qquad \lambda\ \text{expansion and the}\ m_n{=}m_p\ \text{factorization}\ [\text{machine}];
\]

%dsp-level:0
\[
\qquad (2m_n)^2-m_\ell^2\ \text{branch cover:}\ \sqrt2<2\ [\text{squaring comparison}]
\]

%dsp-level:0
\[
(3):\ (2)\ \text{conjoined with}\ (1)\text{'s balance};
\]

%dsp-level:0
\[
\qquad \text{the elimination rows and deficit}\ 1'\ \text{re-cited};
\]

%dsp-level:0
\[
\qquad \text{the upgrade recorded in BM.23's Rmk}
\]

%dsp-level:0
\[
(4)(\mathrm a):\ m_n-m_p=0\ [\text{machine: monomial identity at}\ E_{\mathrm{int}}=6\ \text{both rows}];
\]

%dsp-level:0
\[
\qquad m_\ell>0\ \text{on both species}\ [\mathrm{BM.39}(1)\text{'s SOS row};\ \mathrm{FT.28}/\mathrm{FT.44}(4)\ \text{with}\ \tfrac12<\ln2\ \text{enclosure};\ \text{machine}];
\]

%dsp-level:0
\[
\qquad (\mathrm b),(\mathrm c):\ \text{degeneracy rows re-cited};
\]

%dsp-level:0
\[
\qquad (\mathrm d):\ \text{witnesses}\ T{=}3,
\]

%dsp-level:0
\[
m_\ell{=}m_\nu{=}\tfrac18\ (Q_{\mathrm{ii}}{=}\tfrac34{>}0)\ \text{and}\ T{=}3,
\]

%dsp-level:0
\[
m_\ell{=}m_\nu{=}1\ (Q_{\mathrm{ii}}{=}{-}1{<}0)\ [\text{machine: exact Fractions}]
\]

%dsp-level:0
\[
(5):\ \text{the endpoint identities cleared by}\ 4m_n^2\ [\text{machine}];
\]

%dsp-level:0
\[
\qquad \text{the collinear identity in}\ \mathbb Q[P,N,m_p,m_\nu,E_p,E_\nu]/(E_p^2{-}m_p^2{-}P^2,\ E_\nu^2{-}m_\nu^2{-}N^2)\ [\text{machine}];
\]

%dsp-level:0
\[
\qquad \text{the Lagrange identity}\ [\text{machine:}\ 6\ \text{variables}];
\]

%dsp-level:0
\[
\qquad Q=\sum\mathrm{KE}_i\ [\text{machine: Laurent}]
\]

%dsp-level:0
\[
(6):\ \textstyle\sum m_{\mathrm{out}}\ge m_p+m_\ell>m_n=m_p\ [\text{machine: sign arithmetic on rational witnesses}];
\]

%dsp-level:0
\[
\qquad (7):\ \text{the scope record}.\ \blacksquare
\]

\textbf{Phys.} The vertex is assembled from necessities the corpus already proved, and then weighed. The grammar: the ladder rule moves the ternary read by one, the alphabet overflows, total charge vanishes in every valid universe, one compensating value exists, and the inventory leaves exactly one charged class and one neutral class standing. The weighing is where the corpus surprises itself honestly: every rest read it has ever derived gives the neutron and the proton the same mass — six bits over $\mathfrak h$, identical to the last monomial, because no crystallized count carries a charge exponent and every normalized average is structurally charge-blind. A two-body decay with equal parent and daughter masses strands the daughter below its own mass shell for any positive emitted mass, so the neutral member is forced wherever the vertex is kinematically admissible at all; and admissibility itself fails on every adopted read, since the Q-value collapses to minus the lepton masses — strictly negative. The one functional that splits the classes was classified as a measure price; adopting it is a specification act the corpus prices and declines. What remains is exact and complete: a conditional three-body kinematics with its endpoint formula, a general channel-closure theorem, and a verdict — the reaction is grammatically forced and kinematically closed, and the gap between them is a read specification, not a missing derivation.

\end{theorem}

\begin{theorem}[{\texorpdfstring{BM.42 — The angular momentum at the vertex: the Noether moment, the Lagrange decomposition, the absorption adjudication, and the $\ell$-selection census}{BM.42 - The angular momentum at the vertex: the Noether moment, the Lagrange decomposition, the absorption adjudication, and the -selection census}}]\label{tag:BM-42}
\textbf{Thm.}

%dsp-level:0
\[
(1)\quad\textbf{the moment current is conserved:}\qquad M^{\mu\nu\lambda}:=x^\nu T^{\mu\lambda}-x^\lambda T^{\mu\nu};
\]

%dsp-level:0
\[
\partial_\mu M^{\mu\nu\lambda}=\bigl(T^{\nu\lambda}-T^{\lambda\nu}\bigr)+x^\nu\partial_\mu T^{\mu\lambda}-x^\lambda\partial_\mu T^{\mu\nu}
\]

%dsp-level:0
\[
\qquad [\text{Leibniz};\ \text{machine: the formal identity}];
\]

%dsp-level:0
\[
T^{\mathrm{Hilb}}_{\mu\nu}\ \text{is symmetric}\ [\mathrm{BM.9}(2);\ \text{machine: the dust witness}\ \mathrm{diag}(\rho,0,0,0)]\ \text{and on-shell conserved}\ [\mathrm{BM.9}(3);\ \mathrm{BM.11}:\ \nabla_\mu T^{\mu\nu}=0\ \text{two routes}]:
\]

%dsp-level:0
\[
\boxed{\ \partial_\mu M^{\mu\nu\lambda}=0\ \ \text{on-shell}\ }
\]

%dsp-level:0
\[
\qquad [\text{machine: all}\ 64\ \text{components on the dust class}]:
\]

%dsp-level:0
\[
\qquad \textbf{the total angular momentum is the conserved charge of the promoted rotation}\ [\mathrm{BM.8}(3):\ C_3\to SO(3);\ \mathrm{Def\text{-}H.8}\ \mathrm{Rmk}:\ \text{spin coarse-grained along a geodesic is angular momentum}];
\]

%dsp-level:0
\[
(2)\quad\textbf{the budget at the vertex:}
\]

%dsp-level:0
\[
\qquad J_{\mathrm{in}}=J_{\mathrm{out}}\ \ \text{as the charge of}\ (1):
\]

%dsp-level:0
\[
\qquad j_n=j_p\oplus J_{\ell\nu}\oplus L_{\mathrm{orb}}\ \ \text{in}\ SO(3)\ \text{coupling ranges}\ |j_1-j_2|\le J\le j_1+j_2\ [\text{math}];
\]

%dsp-level:0
\[
\text{the numerical}\ j\text{-attachments ride}\ \mathrm{Def\text{-}H.38}\text{'s declared rule}\ [(7)];
\]

%dsp-level:0
\[
\qquad \text{the budget itself is the exact representation row};
\]

%dsp-level:0
\[
(3)\quad\textbf{the orbital content is already inside}\ p^2\textbf{:}
\]

%dsp-level:0
\[
\qquad \textbf{the Lagrange identity:}
\]

%dsp-level:0
\[
\qquad r^2\,\delta^{ij}p_ip_j=(\vec x\cdot\vec p)^2+\sum_{i<j}\bigl(x_ip_j-x_jp_i\bigr)^2
\]

%dsp-level:0
\[
\qquad [\text{machine: exact polynomial identity in}\ \mathbb Q[x,y,z,p_x,p_y,p_z]];
\]

%dsp-level:0
\[
p^2=p_r^2+\frac{L^2}{r^2},\qquad p_r:=\frac{\vec x\cdot\vec p}r,\quad L^2:=\sum_{i<j}L_{ij}^2,\quad L_{ij}:=x_ip_j-x_jp_i;
\]

%dsp-level:0
\[
\textbf{the dispersion in decomposed form:}
\]

%dsp-level:0
\[
\qquad r^2E^2=r^2m^2+r^2p_r^2+L^2
\]

%dsp-level:0
\[
\qquad [\text{machine: the cleared substitution into}\ \mathrm{BM.34}(4)\text{'s master identity, vacuum gauge}]:
\]

%dsp-level:0
\[
\textbf{the orbital kinetic energy needs no separate term}\ —\ \text{it is}\ \mathrm{BM.34}\text{'s}\ p^2\ \text{itself};
\]

%dsp-level:0
\[
(4)\quad\textbf{the spin-absorption adjudication:}
\]

%dsp-level:0
\[
\qquad \text{spin is the weak sector's contribution}\ [\mathrm{Def\text{-}H.8}];
\]

%dsp-level:0
\[
\qquad \textbf{its two derived metric couplings:}
\]

%dsp-level:0
\[
(\mathrm i)\ \ \text{the rest read: the pole mass is the spin-connection holonomy defect}\ m^2_{\mathrm{pole}}
\]

%dsp-level:0
\[
= 4\hat\mu'^2\sin^2(\Theta_R/4)\ [\mathrm{Def\text{-}H.37}(2);\ \mathrm{Def\text{-}H.21};\ \mathrm{FT.54}(3)\text{--}(4)]:\ \ \text{the frame coupling is inside}\ m_i;
\]

%dsp-level:0
\[
(\mathrm{ii})\ \ \text{the moving body: the shift drag}\ N^ip_i,
\]

%dsp-level:0
\[
\ \text{and the shift is the weak block}\ [\mathrm{BM.34}(3)\text{--}(4);\ \mathrm{FT.157}(1)]:\ \ \text{inside the dispersion};
\]

%dsp-level:0
\[
\text{the Casimir identity}\ \textstyle\sum_iW_i^2
\]

%dsp-level:0
\[
= 2\Pi_{\mathcal R}\ [\mathrm{Cor}\ \mathrm{CT0.7a};\ \mathrm{BM.40}(1);\ \text{machine re-derived}]\ \text{is a branch-count read}:
\]

%dsp-level:0
\[
\boxed{\ \neg\exists\ \text{a theorem constructing a scalar spin--metric coupling energy beyond}\ (\mathrm i),(\mathrm{ii})\ }
\]

%dsp-level:0
\[
\qquad [\text{inventory}];
\]

%dsp-level:0
\[
\text{a separate}\ E_{\mathrm{spin}}\ \text{term in}\ Q\ \text{would double-count}\ [\text{machine: the budget closes}\ —\ Q=\textstyle\sum_i\mathrm{KE}_i\ \text{follows from the conservation laws alone, with no residual}];
\]

%dsp-level:0
\[
\text{the frequency-read row stays on record}\ [\mathrm{FT.131}(3);\ \mathrm{FT.159}\text{'s rate-modulation residue untouched}];
\]

%dsp-level:0
\[
(5)\quad\textbf{the Q-identity with the absorptions:}
\]

%dsp-level:0
\[
\qquad \boxed{\ Q=m_n-m_p-m_\ell-m_\nu=\mathrm{KE}_p+\mathrm{KE}_\ell+\mathrm{KE}_\nu\ },
\]

%dsp-level:0
\[
\qquad \mathrm{KE}_i=N\Bigl(\sqrt{m_i^2+p_{r,i}^2+\frac{L_i^2}{r_i^2}}-m_i\Bigr);
\]

%dsp-level:0
\[
[\mathrm{BM.34}(4)\text{--}(5);\ (3);\ \text{machine: Laurent identity}];
\]

%dsp-level:0
\[
\qquad E_{\mathrm{spin\text{-}coupling}}=E_{\mathrm{orbital\text{-}coupling}}=0\ \ \text{as separate terms}\ [(3)\text{--}(4)\ \text{absorbed}];
\]

%dsp-level:0
\[
\mathrm{BM.41}(4)\text{'s sign theorem is untouched}\ —\ \text{the absorptions add no term to the budget};
\]

%dsp-level:0
\[
(6)\quad\textbf{the Carruthers--Nieto adjudication — the correction is zero:}
\]

%dsp-level:0
\[
\qquad \Delta L_z\,\Delta(\sin\theta)\ge\tfrac12|\langle\cos\theta\rangle|
\]

%dsp-level:0
\[
\qquad [\mathrm{BM.10}(5);\ \text{the naive}\ \Delta\theta\,\Delta L\ge\tfrac12\ \text{refuted};\ \text{witnesses re-cited}];
\]

%dsp-level:0
\[
\text{the kinetic corollary at}\ r=\Delta x_{\mathrm B}:
\]

%dsp-level:0
\[
\qquad \langle p^2\rangle\ \ge\ \frac{\langle L^2\rangle}{\Delta x_{\mathrm B}^2}\ \ge\ (\Delta L_z)^2\,m_{\mathrm{hid}}^2\ \ge\ \tfrac14\langle\cos\theta\rangle^2\,m_{\mathrm{hid}}^2
\]

%dsp-level:0
\[
\qquad [(3);\ \mathrm{BM.13}(4);\ \text{machine}];
\]

%dsp-level:0
\[
\mathrm{BM.34}(6)\text{'s Cartesian row:}\quad\langle p^2\rangle\ge\tfrac34m_{\mathrm{hid}}^2;
\]

%dsp-level:0
\[
\qquad \tfrac14\langle\cos\theta\rangle^2\le\tfrac14<\tfrac34
\]

%dsp-level:0
\[
\qquad [\text{machine: exact rationals}]:
\]

%dsp-level:0
\[
\boxed{\ \textbf{the CN corollary is uniformly weaker — the correction to the zero-point bounds}\ \bigl(\tfrac74,\ \tfrac38,\ \tfrac{\sqrt7}2\bigr)\ \textbf{is exactly zero}\ }
\]

%dsp-level:0
\[
(7)\quad\textbf{the spin ladder and the anchor census:}
\]

%dsp-level:0
\[
\qquad \text{under}\ \mathrm{Def\text{-}H.38}\ \text{(the declared multiplicity rule}\ \operatorname{copies}(c)=k(c)):\ \ \text{a winding-}k\ \text{chain carries}\ \mathrm{Sym}^k(\mathcal H_{\mathsf v}),
\]

%dsp-level:0
\[
\quad j=\frac{|k|}2\ [\mathrm{FT.55}(2);\ \mathrm{FT.60}(1)(\mathrm{iv})];
\]

%dsp-level:0
\[
\textbf{the Casimir ladder:}\quad j(j{+}1)=\frac k2\Bigl(\frac k2{+}1\Bigr)\ \ \text{generalizing}\ \mathrm{Cor}\ \mathrm{CT0.7a}
\]

%dsp-level:0
\[
\qquad [\mathrm{FT.60}(3);\ \text{machine: the ladder values}\ \bigl(\tfrac34,2,\tfrac{15}4,\dots\bigr)\ \text{exact rationals}];
\]

%dsp-level:0
\[
\textbf{the anchor:}
\]

%dsp-level:0
\[
\quad k=1\ \text{(the digon}\ [\mathrm{FT.28};\ \mathrm{FT.107}(1)]):\quad\mathrm{Sym}^1=\mathcal H_{\mathsf v}\cong\mathbb C^2,
\]

%dsp-level:0
\[
\quad j=\tfrac12,\quad\textstyle\sum_iJ_i^2=\tfrac34\mathbb I_2
\]

%dsp-level:0
\[
\qquad [\mathrm{EW.32};\ \text{machine}];
\]

%dsp-level:0
\[
\text{the Platonic spin column stands under its logged declaration}\ [\text{Census FT.63}(1):\ j_{\max}\in\{6,6,12,30,30\};\ \text{row-wise}];
\]

%dsp-level:0
\[
\textbf{the line-winding attachment}\ —\ \text{which}\ k\ \text{each named line carries}\ —\ \text{is}\ \neg\text{computed: the ledger mechanism assigning}\ \mathrm{Sym}^k\ \text{to materialized chains is the residue}\ [\mathrm{FT.55}(4)];
\]

%dsp-level:0
\[
(8)\quad\textbf{the selection census (the anchor class, all four lines at}\ j{=}\tfrac12\textbf{):}
\]

%dsp-level:0
\[
\qquad 2\otimes2\otimes2=4\oplus2\oplus2,\quad j\in\{\tfrac32,\tfrac12\}
\]

%dsp-level:0
\[
\qquad [\mathrm{FT.88}(7);\ \text{machine: the charpoly of}\ (\textstyle\sum J_i)^2\ \text{on}\ \mathbb C^8=(x-\tfrac{15}4)^4(x-\tfrac34)^4];
\]

%dsp-level:0
\[
\textbf{the lepton pair:}\quad\tfrac12\otimes\tfrac12=0\oplus1
\]

%dsp-level:0
\[
\qquad [\text{machine: the charpoly on}\ \mathbb C^4=x^3(x-2)];
\]

%dsp-level:0
\[
\qquad \textbf{the full vertex:}\quad2^{\otimes4}:\ \ 2\cdot(j{=}0)\oplus3\cdot(j{=}1)\oplus1\cdot(j{=}2),\quad2+9+5=16
\]

%dsp-level:0
\[
\qquad [\text{machine: the charpoly on}\ \mathbb C^{16}=x^2(x{-}2)^9(x{-}6)^5];
\]

%dsp-level:0
\[
\textbf{the}\ \ell\textbf{-coupling census}\ (J_{\mathrm{int}}:=j_p\otimes(\ell\text{-channel}\otimes\nu\text{-channel}),\ j_n=\tfrac12):
\]

%dsp-level:0
\[
\begin{array}{c|c|c|c}\text{pair spin}\ S & J_{\mathrm{int}}\in & \ell=0 & \ell\ge1\\\hline0\ (\text{singlet}) & \{\tfrac12\} & \checkmark\ (J{=}\tfrac12\ni\tfrac12) & \ell{=}1:\ \checkmark;\quad\ell{=}2:\ \neg\ (\{\tfrac32,\tfrac52\}\not\ni\tfrac12)\\1\ (\text{triplet}) & \{\tfrac12,\tfrac32\} & \checkmark\ (\text{both}) & \ell{=}1:\ \checkmark;\quad\ell{=}2:\ \checkmark\ \text{via}\ J_{\mathrm{int}}{=}\tfrac32\end{array}\qquad[\text{machine: exact half-integer range census}];
\]

%dsp-level:0
\[
\textbf{the leading admissible class:}\ \ell=0\ \text{with either pair-spin channel};
\]

%dsp-level:0
\[
\qquad \ell
\]

%dsp-level:0
\[
\ge 1\ \text{is the orbital ladder}\ —\ \text{the structural allowed/forbidden classification}\ [\neg\text{SM labels imported}];
\]

%dsp-level:0
\[
(9)\quad\textbf{the parity column — the scope row:}
\]

%dsp-level:0
\[
\qquad \textbf{the orbital factor}\ (-1)^\ell\ \textbf{is derived:}\quad\chi_\ell(\pi)=\sum_{m=-\ell}^{\ell}(-1)^m=(-1)^\ell\ \ \text{for}\ \ell=0,\dots,6
\]

%dsp-level:0
\[
\qquad [\text{the integer}\ SO(3)\ \text{character at angle}\ \pi;\ \text{machine: exact integers}];
\]

%dsp-level:0
\[
\textbf{the lines' intrinsic parity labels are}\ \neg\textbf{derived:}\quad\neg\exists\ \text{a Hilbert-space parity operator on the pair register}\ [\mathrm{FT.135}(4)\ \text{stands}];
\]

%dsp-level:0
\[
\qquad \text{parity acts at the label level as the branch-label flip}\ \iota\ [\mathrm{FT.135}(1);\ \mathrm{EM.24}(2)];
\]

%dsp-level:0
\[
\text{the four gradings are pairwise distinct}\ [\mathrm{FT.136};\ \mathrm{FT.141}(6);\ \mathrm{FT.144}(5)]:
\]

%dsp-level:0
\[
\qquad \textbf{the}\ J^P\ \textbf{row-labels}\ \bigl(\tfrac12^+\ \text{and kin}\bigr)\ \textbf{are external vocabulary and are}\ \neg\textbf{imported}\ [\mathrm{P5}];
\]

%dsp-level:0
\[
\text{the derived content is the}\ \ell\text{-ladder with its}\ (-1)^\ell\ \text{factor, coupled by}\ (8)\text{'s census}.
\]

\textbf{Rmk (budget, not labels).} The vertex conserves the charge the promoted rotation supplies: the moment current closes on the symmetry of the Hilbert stress tensor and its on-shell conservation, both rows the Noether tier proved. The budget's two would-be extra terms dissolve on inspection — the orbital piece is $p^2$ itself by Lagrange, and the spin piece is the frame holonomy inside the rest read plus the weak shift drag inside the dispersion; a separate spin-coupling energy would count the same coupling twice, and the Carruthers–Nieto refinement of the angular spread is uniformly weaker than the Cartesian zero-point row, so the bounds stand uncorrected. What the corpus does not do is attach $j$-numbers to named lines beyond the anchor class: the census computes the representation structure — two singlets, three triplets, one quintet across the fourfold vertex — and the line-data attachment waits on the winding residue.

\textbf{Req.} A3 (Part 1); D1.8 (Parts 1–2); CT0.7(1), CT1.2, CT1.4, Cor CT0.7a (Part 1); CT2.6 (Part 2); P5 (Parts 1–2); Def-H.8, Def-H.21, Def-H.38, FT.54(3), FT.54(4), FT.55(2), FT.55(4), FT.60(1), FT.60(3), FT.88(7), FT.101, FT.107(1), FT.131(3), FT.135(1), FT.135(4), FT.136, FT.141(6), FT.144(5), FT.157(1) (Part 3); Census FT.28, FT.63(1) (row-wise); Def-BM.10, BM.8(3), BM.9(2), BM.9(3), BM.10(5), BM.11, BM.13(4), BM.34(3), BM.34(4), BM.34(5), BM.34(6), BM.40(1), BM.41(4), BM.41(5) (this part); the Leibniz/jet identities, the Lagrange identity, Clebsch–Gordan range censuses, charpolys by Faddeev–LeVerrier over $\mathbb Q$, exact half-integer arithmetic [math; machine: \texttt{checks/\allowbreak{}verify\_\allowbreak{}H\_\allowbreak{}neutron\_\allowbreak{}decay\_\allowbreak{}vertex.py} §D].

\textbf{Pf.}

%dsp-level:0
\[
(1):\ \partial_\mu(x^\nu T^{\mu\lambda})=(\partial_\mu x^\nu)T^{\mu\lambda}+x^\nu\partial_\mu T^{\mu\lambda}=T^{\nu\lambda}+x^\nu\partial_\mu T^{\mu\lambda}\ [\text{Leibniz}];
\]

%dsp-level:0
\[
\qquad \text{the dust witness:}\ T
\]

%dsp-level:0
\[
= \mathrm{diag}(\rho,0,0,0),
\]

%dsp-level:0
\[
\ \rho\ \text{constant}:\ \partial_\mu M^{\mu\nu\lambda}
\]

%dsp-level:0
\[
= T^{\nu\lambda}-T^{\lambda\nu}
\]

%dsp-level:0
\[
= 0\ \text{all}\ 64\ \text{components}\ [\text{machine: symbolic differentiation in}\ \mathbb Q[x_0,x_1,x_2,x_3,\rho]]
\]

%dsp-level:0
\[
(2):\ \text{the coupling ranges are the}\ SO(3)\ \text{representation rows}\ [\text{math}];
\]

%dsp-level:0
\[
\qquad (3):\ r^2p^2-(\vec x\cdot\vec p)^2=\sum_{i<j}(x_ip_j-x_jp_i)^2\ [\text{machine: expansion in}\ 6\ \text{variables}];
\]

%dsp-level:0
\[
\qquad \text{the substitution into}\ \mathrm{BM.34}(4)\text{'s identity at}\ N{=}1,N^i{=}0,\gamma{=}\delta,\ \text{cleared by}\ r^2\ [\text{machine}]
\]

%dsp-level:0
\[
(4):\ \mathrm{Def\text{-}H.37}(2)/\mathrm{Def\text{-}H.21}\ \text{re-cited};
\]

%dsp-level:0
\[
\qquad \mathrm{BM.34}(3)\text{--}(4)/\mathrm{FT.157}(1)\ \text{re-cited};
\]

%dsp-level:0
\[
\qquad \text{the inventory scan: no scalar spin--metric coupling term exists in Parts 1--4}\ [\text{machine: token scan of the energy-budget rows}];
\]

%dsp-level:0
\[
\qquad \text{the budget closure:}\ m_n=E_p+E_\ell+E_\nu\implies Q=\sum(E_i-m_i)\ [\text{machine: Laurent}]
\]

%dsp-level:0
\[
(5):\ (3)\text{'s decomposition substituted into}\ \mathrm{BM.34}(5)\text{'s KE definition}\ [\text{machine}]
\]

%dsp-level:0
\[
(6):\ (\Delta L_z)^2\ge\tfrac14\langle\cos\theta\rangle^2\ [\mathrm{BM.10}(5)\ \text{at}\ \Delta(\sin\theta)\le1];
\]

%dsp-level:0
\[
\qquad \langle L^2\rangle\ge(\Delta L_z)^2\ \text{at}\ \langle L_z\rangle=0;
\]

%dsp-level:0
\[
\qquad r=\Delta x_{\mathrm B}=1/m_{\mathrm{hid}}\ [\mathrm{Def\text{-}BM.4}(1);\ \mathrm{BM.13}(4)];
\]

%dsp-level:0
\[
\qquad \tfrac14<\tfrac34\ [\text{machine}]
\]

%dsp-level:0
\[
(7):\ \mathrm{Def\text{-}H.38}\text{'s derived-content rows re-cited};
\]

%dsp-level:0
\[
\qquad \text{the anchor:}\ J_i=\tfrac12\sigma_i:\ \sum_iJ_i^2=\tfrac14\sum_i\sigma_i^2=\tfrac14\cdot3\,\mathbb I_2=\tfrac34\mathbb I_2\ [\text{machine over}\ \mathbb Q(i)];
\]

%dsp-level:0
\[
\qquad \text{the ladder values}\ \tfrac k2(\tfrac k2{+}1)
\]

%dsp-level:0
\[
= (\tfrac34,2,\tfrac{15}4,6,\dots)\ \text{at}\ k
\]

%dsp-level:0
\[
= 1,
\]

%dsp-level:0
\[
2,
\]

%dsp-level:0
\[
3,
\]

%dsp-level:0
\[
4\ [\text{machine: exact rationals}]
\]

%dsp-level:0
\[
(8):\ (\textstyle\sum_iJ_i^{(2)})^2\ \text{on}\ \mathbb C^4:\ \text{charpoly}\ x^3(x-2)\ [\text{machine: Faddeev--LeVerrier}];
\]

%dsp-level:0
\[
\qquad \text{on}\ \mathbb C^8:\ (x-\tfrac{15}4)^4(x-\tfrac34)^4;
\]

%dsp-level:0
\[
\qquad \text{on}\ \mathbb C^{16}:\ x^2(x-2)^9(x-6)^5\ [\text{machine: exact charpolys over}\ \mathbb Q];
\]

%dsp-level:0
\[
\qquad \text{the}\ \ell\text{-ranges:}\ |J_{\mathrm{int}}-\ell|\le\tfrac12\le J_{\mathrm{int}}+\ell\ \text{censused on half-integers}\ [\text{machine}]
\]

%dsp-level:0
\[
(9):\ \sum_{m=-\ell}^{\ell}(-1)^m=(-1)^\ell\ \text{for}\ \ell=0..6\ [\text{machine: integer sums}];\qquad\mathrm{FT.135}(4)\ \text{re-cited}.\ \blacksquare
\]

\textbf{Phys.} Angular momentum at the vertex is a theorem before it is a table. The moment of the Hilbert stress tensor is conserved wherever the tensor is symmetric and its own divergence vanishes on-shell — both rows the Noether tier proved — so the charge the promoted rotation carries balances across the decay exactly. The budget then needs no new terms: the orbital share is already the momentum square, split by Lagrange's identity into radial and angular pieces with nothing left over, and the spin share lives inside the rest mass as frame holonomy and inside the dispersion as the weak shift's drag. Asked for a spin-coupling energy to add to the Q-value, the corpus answers that adding one would double-count, and the angular uncertainty refinement of the zero-point bound comes out uniformly weaker — a quarter against three quarters — so the bounds stand. What is genuinely new is the census: at the anchor winding, four spin-half lines tensor into two singlets, three triplets and one quintet; the lepton pair couples to singlet or triplet; and the orbital ladder admits zero and one against both pair channels, two only through the triplet route. The parity factor of the ladder is the integer character of the rotation group, while the parity labels of the lines themselves stay outside: the register has no parity operator, and the tier imports no external $J^P$ vocabulary.

\end{theorem}

\begin{theorem}[{\texorpdfstring{BM.43 — The curvature rows of the members: the dust sources, the weak-field potential, and the superposition scope}{BM.43 - The curvature rows of the members: the dust sources, the weak-field potential, and the superposition scope}}]\label{tag:BM-43}
\textbf{Thm.}

%dsp-level:0
\[
(1)\quad\textbf{the baryon rows (the exact monomials):}
\]

%dsp-level:0
\[
\qquad E_{\mathrm{int}}=6\ \text{for both classes}\ [\mathrm{BM.17};\ \mathrm{BM.26};\ \text{Census FT.35}(2)]:
\]

%dsp-level:0
\[
\Delta x_{\mathrm B}=\frac{\mathfrak h}{6\ln2},
\]

%dsp-level:0
\[
\qquad V_{\mathrm{eff}}=\frac{\mathfrak h^3}{216(\ln2)^3},
\]

%dsp-level:0
\[
\qquad m=\frac{6\ln2}{\mathfrak h},
\]

%dsp-level:0
\[
\qquad \rho=\frac{1296(\ln2)^4}{\mathfrak h^4},
\]

%dsp-level:0
\[
\qquad \kappa^2\rho=\frac{2592\pi(\ln2)^3}{\mathfrak h^2},
\]

%dsp-level:0
\[
\qquad r_s=3\,\mathfrak h;
\]

%dsp-level:0
\[
\textbf{the lock:}\quad(\Delta x_{\mathrm B})^2\,\kappa^2\rho=2\pi E_{\mathrm{int}}^2\ln2=72\pi\ln2
\]

%dsp-level:0
\[
\qquad [\mathrm{BM.26};\ \mathrm{BM.16}(1);\ \text{machine: every row one exact Laurent monomial}];
\]

%dsp-level:0
\[
\text{the parent and the daughter carry identical rows}\ [\mathrm{BM.28}(1)\ \text{charge-blindness}];
\]

%dsp-level:0
\[
(2)\quad\textbf{the lepton rows (symbolic):}
\]

%dsp-level:0
\[
\qquad \text{the taxonomy law holds:}\quad m_i=\frac{E_{\mathrm{int}}^{(i)}\ln2}{\mathfrak h}\ [\mathrm{FT.44}(4)];
\]

%dsp-level:0
\[
\qquad E_{\mathrm{int}}^{(\ell)},\ E_{\mathrm{int}}^{(\nu)}\ \text{are}\ \neg\text{computed at any floor}:
\]

%dsp-level:0
\[
\text{the rows stand as symbolic monomials:}\quad\kappa^2\rho_i=\frac{2\pi\,\mathfrak h\,E_{\mathrm{int}}^{(i)}}{V_i},
\]

%dsp-level:0
\[
\qquad r_s^{(i)}=\frac{\mathfrak h\,E_{\mathrm{int}}^{(i)}}2
\]

%dsp-level:0
\[
\qquad [\mathrm{BM.9}(5);\ \mathrm{BM.16}(4);\ \text{machine}];
\]

%dsp-level:0
\[
\text{the transport alternative:}\quad m_i\ \text{from}\ \mathrm{BM.39}\ (\text{symbolic in}\ (g,g',\hat\mu,\ell_c));
\]

%dsp-level:0
\[
\qquad \text{the conversion}\ E_{\mathrm{int}}^{(i)}\leftrightarrow\text{transport data is item-7 residue}\ [\mathrm{Def\text{-}BM.11}(7);\ \text{pointer}];
\]

%dsp-level:0
\[
(3)\quad\textbf{the field-equation rows:}
\]

%dsp-level:0
\[
\qquad T_{\mu\nu}=\mathrm{diag}(\rho,0,0,0)\ \text{at rest},\ p=0\ [\mathrm{BM.9}(4);\ \mathrm{CT4.5}];
\]

%dsp-level:0
\[
\qquad R=\kappa^2(\rho-3p)+4\Lambda\ [\mathrm{BM.15}(2)];
\]

%dsp-level:0
\[
\qquad RV=\varkappa^2m\ [\mathrm{BM.20}(1)];
\]

%dsp-level:0
\[
(4)\quad\textbf{the weak-field potential (the linearized row):}
\]

%dsp-level:0
\[
\qquad g_{00}=-(1+2\Phi),\quad g_{ij}=(1-2\Phi)\delta_{ij},\quad\Phi\ \text{static}:
\]

%dsp-level:0
\[
G^{(1)}_{00}=2\,\nabla^2\Phi,
\]

%dsp-level:0
\[
\qquad G^{(1)}_{0i}=0,
\]

%dsp-level:0
\[
\qquad G^{(1)}_{ij}=0
\]

%dsp-level:0
\[
\qquad [\text{machine: Christoffel}\to\text{Ricci}\to\text{Einstein in the weight-truncated jet ring, exact rational coefficients}];
\]

%dsp-level:0
\[
\text{the field equation}\ G_{\mu\nu}=\kappa^2T_{\mu\nu}\ [\mathrm{CT1.23}]\ \text{at}\ p=0:
\]

%dsp-level:0
\[
\qquad \nabla^2\Phi=\frac{\kappa^2}2\,\rho=4\pi G_N\rho
\]

%dsp-level:0
\[
\qquad [\kappa^2=8\pi G_N\ [\mathrm{CT2.9}];\ \text{machine: the coefficient identity}\ \tfrac{\kappa^2}2=\tfrac{\pi\mathfrak h^2}{\ln2}=4\pi G_N];
\]

%dsp-level:0
\[
\Phi=-\frac{G_Nm}{r}\ \text{solves at}\ r\ne0:
\]

%dsp-level:0
\[
\qquad \nabla^2\frac1r=0
\]

%dsp-level:0
\[
\qquad [\text{machine:}\ -3r^2+3(x^2{+}y^2{+}z^2)\equiv0\ \text{in}\ \mathbb Q[x,y,z,r]/(r^2{-}x^2{-}y^2{-}z^2)];
\]

%dsp-level:0
\[
G_Nm_{\mathrm{hid}}=\frac{\mathfrak h\,E_{\mathrm{int}}}4,
\]

%dsp-level:0
\[
\qquad r_s=2G_Nm_{\mathrm{hid}}=\frac{\mathfrak h\,E_{\mathrm{int}}}2,
\]

%dsp-level:0
\[
\qquad \Phi(\Delta x_{\mathrm B})=-\frac{E_{\mathrm{int}}^2\ln2}4\in\mathfrak q^0
\]

%dsp-level:0
\[
\qquad [\mathrm{CT2.9};\ \mathrm{BM.16}(4);\ \text{machine: monomials}];
\]

%dsp-level:0
\[
(5)\quad\textbf{the superposition:}
\]

%dsp-level:0
\[
\qquad \textbf{source side, exact:}
\]

%dsp-level:0
\[
\quad T_{\mathrm{total}}=\sum_iT^{(i)}\ \text{on disjoint supports}\ [\mathrm{CT1.22}\ \text{channel decomposition};\ \mathrm{D1.5}]:
\]

%dsp-level:0
\[
\qquad \kappa^2\rho_{\mathrm{total}}
\]

%dsp-level:0
\[
= \kappa^2\bigl(\rho_p+\rho_\ell+\rho_\nu\bigr)
\]

%dsp-level:0
\[
= \frac{2\pi\mathfrak h}{V}\bigl(6+E_{\mathrm{int}}^{(\ell)}+E_{\mathrm{int}}^{(\nu)}\bigr)\ [\text{machine: the monomial sum}];
\]

%dsp-level:0
\[
\textbf{metric side:}\quad\text{superposition is exact only at linearized order}\ [\mathrm{CT1.12}:\ \varkappa\|q\|_{\mathrm{op}}<1\ \text{the perturbative regime}];
\]

%dsp-level:0
\[
\qquad \text{the full non-linear joint solution is}\ \neg\text{constructed};
\]

%dsp-level:0
\[
\textbf{the scope row:}\quad\text{the Schwarzschild-type content on record is the}\ r_s\ \text{read and the weak-field potential}\ [(4);\ \mathrm{BM.16}(4)];
\]

%dsp-level:0
\[
\qquad \text{a full exterior solution is}\ \neg\text{in the inventory}\ [\text{the pp-wave witness is the exact curvature row on record}];
\]

%dsp-level:0
\[
(6)\quad\textbf{status:}
\]

%dsp-level:0
\[
\qquad \text{the baryon rows are numerical monomials};
\]

%dsp-level:0
\[
\ \text{the lepton rows symbolic};
\]

%dsp-level:0
\[
\ \text{the vertex curvature total is their exact sum on the source side};
\]

%dsp-level:0
\[
\qquad \neg\text{a measured curvature or radius enters}\ [\mathrm{P5}].
\]

\textbf{Req.} D1.5, D2.9 (Parts 1–2); CT1.12, CT1.22, CT1.23, CT2.1, CT2.9, CT4.5 (Part 2); P3, P5 (Parts 1–2); FT.44(4) (Census); Def-BM.4, Def-BM.10, Def-BM.11, BM.9(4), BM.9(5), BM.15(2), BM.16(1), BM.16(4), BM.17, BM.20(1), BM.26, BM.28(1), BM.34(4), BM.39 (this part); the weight-truncated jet-ring calculus, Laurent monomials, harmonic identities in quotient rings [math; machine: \texttt{checks/\allowbreak{}verify\_\allowbreak{}H\_\allowbreak{}neutron\_\allowbreak{}decay\_\allowbreak{}vertex.py} §E].

\textbf{Pf.}

%dsp-level:0
\[
(1):\ \mathrm{BM.26}\text{'s six rows re-verified as monomials at}\ E_{\mathrm{int}}
\]

%dsp-level:0
\[
= 6\ [\text{machine}]:\quad\rho
\]

%dsp-level:0
\[
= \tfrac{6\ln2}{\mathfrak h}\cdot\tfrac{216(\ln2)^3}{\mathfrak h^3}
\]

%dsp-level:0
\[
= \tfrac{1296(\ln2)^4}{\mathfrak h^4};
\]

%dsp-level:0
\[
\quad\kappa^2\rho
\]

%dsp-level:0
\[
= \tfrac{2\pi\mathfrak h^2}{\ln2}\cdot\tfrac{1296(\ln2)^4}{\mathfrak h^4}
\]

%dsp-level:0
\[
= \tfrac{2592\pi(\ln2)^3}{\mathfrak h^2};
\]

%dsp-level:0
\[
\quad\text{the lock}\ [\text{machine}]
\]

%dsp-level:0
\[
(2):\ \mathrm{FT.44}(4)\text{'s law re-cited};\qquad\text{the symbolic monomials}\ [\text{machine}]
\]

%dsp-level:0
\[
(3):\ \text{frozen rows re-cited}
\]

%dsp-level:0
\[
(4):\ \text{the jet-ring computation:}\ g^{\mu\nu}g_{\nu\rho}=\delta^\mu_\rho\ \text{to first order}\ [\text{machine}];
\]

%dsp-level:0
\[
\qquad \Gamma^i_{00}=\partial_i\Phi,\ \ \Gamma^0_{i0}=\partial_i\Phi,\ \ \Gamma^k_{ij}=-\partial_i\Phi\,\delta^k_j-\partial_j\Phi\,\delta^k_i+\partial^k\Phi\,\delta_{ij}\ [\text{machine}];
\]

%dsp-level:0
\[
R_{00}=\nabla^2\Phi,
\]

%dsp-level:0
\[
\quad R_{ij}=\delta_{ij}\nabla^2\Phi,
\]

%dsp-level:0
\[
\quad R=2\nabla^2\Phi\implies G_{00}=2\nabla^2\Phi,\ G_{ij}=0\ [\text{machine: weight-}\le1\ \text{truncation}];
\]

%dsp-level:0
\[
\qquad \tfrac{\kappa^2}2=\tfrac{\pi\mathfrak h^2}{\ln2}=4\pi\cdot\tfrac{\mathfrak h^2}{4\ln2}\ [\mathrm{CT2.9};\ \text{machine}]
\]

%dsp-level:0
\[
\nabla^2(r^{-1})\cdot r^5=-3r^2+3\textstyle\sum_ix_i^2\equiv0\ [\text{machine: the rule ring}];
\]

%dsp-level:0
\[
\qquad G_Nm=\tfrac{\mathfrak h^2}{4\ln2}\cdot\tfrac{E\ln2}{\mathfrak h}=\tfrac{\mathfrak h E}4;
\]

%dsp-level:0
\[
\quad r_s=2G_Nm=\tfrac{\mathfrak h E}2;\quad\Phi(\Delta x_{\mathrm B})=-\tfrac{\mathfrak h E}4\cdot\tfrac{E\ln2}{\mathfrak h}=-\tfrac{E^2\ln2}4\ [\text{machine}]
\]

%dsp-level:0
\[
(5):\ \mathrm{CT1.22}\text{'s decomposition re-cited};\qquad\text{the monomial sum}\ [\text{machine}];\qquad(6):\ \text{the status record}.\ \blacksquare
\]

\textbf{Phys.} Every member of the vertex gravitates, and the corpus writes each gravity in monomials. The two baryon rows are numbers: six bits over $\mathfrak h$, a density of twelve ninety-six $(\ln2)^4$ over $\mathfrak h^4$, a horizon read of three $\mathfrak h$, a lock of seventy-two $\pi\ln2$ — identical for parent and daughter, because no crystallized count carries charge. The lepton rows keep their edge counts symbolic, for the floors that would compute them are not in the census, and the transport species offers the alternative symbolic form. The field side is linearized and exact: the static diagonal perturbation solves the ledger's own field equations with the potential obeying the Poisson row whose coefficient is the corpus's $4\pi G_N$ to the last factor, and the point-source form is harmonic away from the source by a polynomial identity. Superposition is exact where the corpus can prove it — the source side, where stress tensors add — and scoped where it cannot: the full non-linear metric of three flying members is an open construction, and the tier says so rather than dress a linearized row in exact clothes.

\end{theorem}

\begin{theorem}[{\texorpdfstring{BM.44 — The time evolution: the per-tick weight, the e-folding count, and the unit discipline}{BM.44 - The time evolution: the per-tick weight, the e-folding count, and the unit discipline}}]\label{tag:BM-44}
\textbf{Thm.}

%dsp-level:0
\[
(1)\quad\textbf{the frozen rate rows:}
\]

%dsp-level:0
\[
\qquad \sum_n\langle\varnothing|\hat U^n|\varnothing\rangle w^n=\frac1{1-w^2x_+x_-}\ \ \text{in}\ \mathbb Q[w,x_+,x_-]/(w^{12})
\]

%dsp-level:0
\[
\qquad [\mathrm{FT.108}(5);\ \mathrm{BM.23}(7);\ \text{machine: the resolvent identity re-verified coefficient-wise}];
\]

%dsp-level:0
\[
x_+x_-=4\varepsilon_q(1-4\varepsilon_q)=\frac14-16\Bigl(\varepsilon_q-\frac18\Bigr)^2\le\frac14,
\]

%dsp-level:0
\[
\qquad \text{equality}\iff\varepsilon_q=\frac18
\]

%dsp-level:0
\[
\qquad [\mathrm{BM.19}(2);\ \text{machine: the SOS completion}\ +\ \text{the grid census over}\ \varepsilon_q=k/256];
\]

%dsp-level:0
\[
\text{the pole sits at}\ |w|\ge2;
\]

%dsp-level:0
\[
\qquad \varepsilon_q
\]

%dsp-level:0
\[
= 0
\]

%dsp-level:0
\[
\implies x_+x_-
\]

%dsp-level:0
\[
= 0\implies\text{the resolvent is}\ 1:\ \ \textbf{no transition weight beyond the identity}\ [\mathrm{BM.23}(8);\ \text{machine}];
\]

%dsp-level:0
\[
(2)\quad\textbf{the unit-fugacity slice:}
\]

%dsp-level:0
\[
\qquad \Gamma_{\mathrm{tr}}(\varepsilon_q):=\frac{x_+x_-}{1-x_+x_-}\bigg|_{w=1}
\]

%dsp-level:0
\[
\qquad [\text{the transition content of the resolvent}]:
\]

%dsp-level:0
\[
\Gamma_{\mathrm{tr}}(0)=0,
\]

%dsp-level:0
\[
\qquad \Gamma_{\mathrm{tr}}\bigl(\tfrac18\bigr)=\frac{1/4}{3/4}=\frac13,
\]

%dsp-level:0
\[
\qquad \Gamma_{\mathrm{tr}}\ \text{strictly increasing on}\ \bigl[0,\tfrac18\bigr]
\]

%dsp-level:0
\[
\qquad [\text{machine:}\ \Gamma_{\mathrm{tr}}(\varepsilon_2)-\Gamma_{\mathrm{tr}}(\varepsilon_1)>0\ \text{via}\ 4(\varepsilon_2-\varepsilon_1)\bigl(1-4(\varepsilon_1+\varepsilon_2)\bigr)>0];
\]

%dsp-level:0
\[
\textbf{the exact range:}\quad\Gamma_{\mathrm{tr}}\in\bigl[0,\tfrac13\bigr]\ \ \text{on}\ \ \varepsilon_q\in\bigl[0,\tfrac18\bigr];
\]

%dsp-level:0
\[
(3)\quad\textbf{the chart evolution (the scoped specialization):}
\]

%dsp-level:0
\[
\qquad \text{under}\ \mathrm{BM.8}(2)\text{'s declared stationarity hypothesis the chart carries the time-translation one-parameter group}\ [\mathrm{BM.8}(2);\ \mathrm{BM.34}(4):\ E=-p_0\ \text{its charge}]:
\]

%dsp-level:0
\[
\frac{d}{d\tau}P(\tau)=-\tilde\Gamma\,P(\tau),
\]

%dsp-level:0
\[
\qquad P(\tau)=P_0\,e^{-\tilde\Gamma\tau}
\]

%dsp-level:0
\[
\qquad [\text{the group's exponential};\ \tilde\Gamma\ \text{the per-tick weight at the chart scale}];
\]

$\textbf{the e-folding count:}\qquad\tau_{\mathrm{life}}:=\frac1{\tilde\Gamma}\ \ \text{ticks};\qquad\tilde\Gamma=\Gamma_{\mathrm{tr}}\implies\tau_{\mathrm{life}}\ge3\ \text{ticks},\quad\tau_{\mathrm{life}}\big|_{\varepsilon_q=1/8}=3\ \ \text{exactly}\qquad[\text{machine: exact reciprocals};\ \Gamma_{\mathrm{tr}}\le\tfrac13\ \text{from}\ (2)];$

%dsp-level:0
\[
\varepsilon_q=0:\quad\Gamma_{\mathrm{tr}}=0\implies\tau_{\mathrm{life}}\ \text{has no finite value}:
\]

%dsp-level:0
\[
\qquad \textbf{a non-flickering composite never decays}\ [\mathrm{BM.23}(8)\ \text{re-cited}];
\]

%dsp-level:0
\[
(4)\quad\textbf{the ledger-native form:}
\]

%dsp-level:0
\[
\qquad \Gamma\ \text{is a weight per tick in the ledger order}\ [x^0:=\tau_n;\ \mathrm{D1.8};\ \mathrm{BM.23}(9);\ \mathrm{FT.159}];
\]

%dsp-level:0
\[
\qquad \text{the discrete evolution is the geometric resolvent of}\ (1);\
\]

%dsp-level:0
\[
\textbf{the normalization of the weight to a per-tick probability is}\ \neg\textbf{declared}\ [\mathrm{BM.23}\ \mathrm{Rmk}:\ \Gamma\ \text{is a weight},\ \neg\text{a probability}]:
\]

%dsp-level:0
\[
\qquad (3)\text{'s}\ \tilde\Gamma\ \text{is the chart specialization, scoped and marked};
\]

%dsp-level:0
\[
(5)\quad\textbf{the unit discipline:}
\]

%dsp-level:0
\[
\qquad \neg\text{external seconds};
\]

%dsp-level:0
\[
\qquad \mathrm{Cal\text{-}F}\ \text{carries no rate datum}\ (\mu,\ \hat\mu^2,\ \hat\mu'^2\ \text{only});
\]

%dsp-level:0
\[
\qquad \neg\text{a lifetime is derived in external units and}\ \neg\text{a lifetime is compared}\ [\mathrm{P5};\ \mathrm{BM.23}(9)\text{'s row stands}];
\]

%dsp-level:0
\[
\mathrm{FT.159}\text{'s rate-modulation residue}\ (\text{the flicker-frequency read of}\ \delta\tau\ \text{against}\ \Omega/N_i)\ \text{is untouched};
\]

%dsp-level:0
\[
(6)\quad\textbf{the arrow rows:}
\]

%dsp-level:0
\[
\qquad \text{forward-only}\ [\mathrm{A6}:\ \neg\exists\tau^{-1};\ \mathrm{FT.144}(2):\ \text{the prefix law};\ \mathrm{BM.23}(9)];
\]

%dsp-level:0
\[
\qquad \neg\exists\ \text{an inverse decay in the ledger};
\]

%dsp-level:0
\[
\Theta_{\mathfrak L}\ \text{is the record involution}\ (\text{the formal mirror}),
\]

%dsp-level:0
\[
\ \neg\text{a dynamical reversal}\ [\mathrm{Def\text{-}BM.7}(2);\ \mathrm{FT.144}(4)].
\]

\textbf{Rmk (weight against rate.)} The corpus's decay clock ticks in ledger order and its hand is a rational function of one threshold. Everything exact is stated exactly: the resolvent, its pole beyond two, its maximum third at the threshold midpoint, its annihilation at the sharp threshold. What is not stated is a number of seconds, because no datum in the corpus converts a tick into one — and the tier keeps that silence loud: the exponential form is the chart group's own shape under the declared stationarity scope, the e-folding count is three ticks at the heaviest admissible flicker, and the measured lifetime is quoted nowhere.

\textbf{Req.} A6, D0.4, D1.8 (Part 1); CT0.8 (Part 1); P5 (Parts 1–2); FT.108(4), FT.108(5), FT.144(2), FT.144(4), FT.159 (Part 3); Def-BM.5, Def-BM.7(2), BM.8(2), BM.19(2), BM.19(3), BM.23(7), BM.23(8), BM.23(9), BM.34(4) (this part); Cal-F; the exact resolvent identity in $\mathbb Q[w,x_+,x_-]/(w^{12})$, the SOS parabola completion, exact rational monotonicity and reciprocals [math; machine: \texttt{checks/\allowbreak{}verify\_\allowbreak{}H\_\allowbreak{}neutron\_\allowbreak{}decay\_\allowbreak{}vertex.py} §F].

\textbf{Pf.}

%dsp-level:0
\[
(1):\ (1-w^2x_+x_-)\sum_{n\le11}\bigl[x_+x_-\bigr]^{n/2}w^n\big|_{n\ \mathrm{even}}\equiv1\ \bmod w^{12}\ [\text{machine: coefficient-wise}];
\]

%dsp-level:0
\[
\qquad 4\varepsilon(1-4\varepsilon)=\tfrac14-16(\varepsilon-\tfrac18)^2\ [\text{machine: expansion}];
\]

%dsp-level:0
\[
\qquad \text{the grid census}\ k=0..63:\ \le\tfrac14,\ =\tfrac14\ \text{only at}\ k=32\ [\text{machine}]
\]

%dsp-level:0
\[
(2):\ \Gamma_{\mathrm{tr}}=\tfrac{xx}{1-xx},\ xx:=x_+x_-:\quad\Gamma_{\mathrm{tr}}(0)=0,\ \ \Gamma_{\mathrm{tr}}(\tfrac18)=\tfrac{1/4}{3/4}=\tfrac13\ [\text{machine}];
\]

%dsp-level:0
\[
\qquad \text{monotonicity:}\ xx(\varepsilon_2)-xx(\varepsilon_1)
\]

%dsp-level:0
\[
= 4(\varepsilon_2-\varepsilon_1)(1-4(\varepsilon_1+\varepsilon_2))\ [\text{machine: polynomial identity}],
\]

%dsp-level:0
\[
\ \text{positive on rational witness pairs}\ \varepsilon_1<\varepsilon_2\le\tfrac18
\]

%dsp-level:0
\[
(3):\ \mathrm{BM.8}(2)\text{'s group re-cited};
\]

%dsp-level:0
\[
\qquad \tfrac1{\Gamma_{\mathrm{tr}}(1/8)}=3\ [\text{machine: exact reciprocal}];
\]

%dsp-level:0
\[
\qquad \Gamma_{\mathrm{tr}}\le\tfrac13\implies\tau_{\mathrm{life}}\ge3\ [\text{machine: monotone reciprocal}]
\]

%dsp-level:0
\[
(4)\text{--}(6):\ \text{frozen rows re-cited}\ [\mathrm{D1.8};\ \mathrm{BM.23}(9);\ \mathrm{FT.144}(2),(4);\ \mathrm{A6}];
\]

%dsp-level:0
\[
\qquad \text{the Cal-F inventory:}\ \{\mu,\hat\mu^2,\hat\mu'^2\},\ \neg\text{a rate}\ [\text{scan}].\ \blacksquare
\]

\textbf{Phys.} Time runs through the decay exactly as the corpus runs it everywhere else: in ticks, forward, with no inverse and no seconds. The transition weight is the geometric resolvent of the flicker — a rational function of one threshold, pole beyond two, annihilated at the sharp threshold where a composite that cannot flicker cannot decay. On the unit-fugacity slice the weight climbs from zero to a third as the threshold moves to its midpoint, and the e-folding count is its reciprocal: at least three ticks, exactly three at the heaviest admissible flicker. The exponential curve belongs to the chart's own translation group under the stationarity scope the Noether tier declared; the ledger keeps only the weights. And the discipline is the point: a lifetime in seconds would need a rate datum the corpus does not carry, so the tier states the tick count, names the missing datum, and quotes no clock.

\end{theorem}

\begin{theorem}[{\texorpdfstring{BM.45 — The assembled reaction: the conservation table, the naming row, the final read, and the adjudication}{BM.45 - The assembled reaction: the conservation table, the naming row, the final read, and the adjudication}}]\label{tag:BM-45}
\textbf{Thm.}

%dsp-level:0
\[
(1)\quad\textbf{the reaction (the structural row, with the identification-layer labels):}
\]

%dsp-level:0
\[
\bigl[\varepsilon\text{-sealed 3-juncture, neutral class,}\ E_{\mathrm{int}}{=}6\bigr]\ \longrightarrow\ \bigl[\varepsilon\text{-sealed 3-juncture, unit class,}\ E_{\mathrm{int}}{=}6\bigr]\ +\ \bigl[\ell\text{-channel line},\ q{=}{-}1\bigr]\ +\ \bigl[\nu\text{-channel line,}\ q{=}0,\ \text{the materialized-arrow orientation}\bigr];
\]

%dsp-level:0
\[
\textbf{the labels:}
\]

%dsp-level:0
\[
\qquad \text{neutron}\ \longrightarrow\ \text{proton}\ +\ \text{electron}\ +\ \text{electron antineutrino}
\]

%dsp-level:0
\[
\qquad [\mathrm{Def\text{-}BM.3}(2)\text{--}(3);\ \mathrm{Def\text{-}BM.6}(1);\ \text{naming row A-BM.5}];
\]

%dsp-level:0
\[
\textbf{A-BM.5:}\quad\text{electron antineutrino}\ :
\]

%dsp-level:0
\[
= \ \text{the}\ \nu\text{-channel line in the CPT/}\omega\text{-image orientation}\ [\mathrm{BM.25}\text{'s involution};\ \mathrm{BM.27}\text{'s antiproton pattern};\ \mathrm{FT.144}(5):\ \text{the occupancy sector is derived, its name is reserved content}];
\]

%dsp-level:0
\[
\text{the charged member is}\ \mathrm{Def\text{-}BM.6}(1)\text{'s electron row}\ [\mathrm{BM.23}(5)];
\]

%dsp-level:0
\[
\qquad \text{the neutral member's line index within the}\ \nu\text{-triplet is}\ \neg\text{derived}\ [\mathrm{FT.145}(3)];
\]

%dsp-level:0
\[
(2)\quad\textbf{the conservation table:}\qquad\begin{array}{l|l|l}\text{law} & \text{the row} & \text{status}\\\hline\text{charge} & 0=(+1)+(-1)+0 & \text{derived}\ [\mathrm{EM.11};\ \mathrm{BM.23}(4);\ \text{machine}]\\\text{Janus balance} & \#J^+=\#J^-\ \text{on every emitted closed chain} & \text{derived}\ [\mathrm{D1.4};\ \mathrm{FT.144}(1)\text{--}(2);\ \mathrm{BM.7}]\\\text{rung-1 weak Casimir} & \text{branch-neutral at the}\ \sigma-vertex;\ \text{carriers internally balanced} & \text{derived}\ [\mathrm{BM.40}(3)\text{--}(4)]\\\text{angular momentum} & \partial_\mu M^{\mu\nu\lambda}=0\ \text{on-shell};\ \text{the}\ \ell\text{-census of}\ \mathrm{BM.42}(8) & \text{derived}\ [\mathrm{BM.42}(1)\text{--}(2)]\\\text{sealing phase} & \Theta_{\mathrm{int}}\in4\mathbb Z\ \text{on both baryon-type junctions} & \text{derived}\ [\mathrm{BM.2}(1);\ \text{Census FT.49}(3)]\\\text{energy--momentum} & \text{the chart-side law with the threshold row}\ Q\ge0 & \text{derived}\ [\mathrm{BM.11};\ \mathrm{BM.34};\ \mathrm{BM.41}(4)\text{--}(5)]\\\text{the arrow} & \text{forward-only};\ \neg\exists\ \text{an inverse decay} & \text{derived}\ [\mathrm{A6};\ \mathrm{FT.144}(2);\ \mathrm{BM.44}(6)]\end{array}
\]

%dsp-level:0
\[
(3)\quad\textbf{the Q-value and the distribution:}
\]

%dsp-level:0
\[
\qquad Q=-(m_\ell+m_\nu)<0\ \ \text{on every adopted read}\ [\mathrm{BM.41}(4)]:
\]

%dsp-level:0
\[
\qquad \textbf{the reaction assembles and does not go on-shell in the corpus's own balance};
\]

%dsp-level:0
\[
Q=\mathrm{KE}_p+\mathrm{KE}_\ell+\mathrm{KE}_\nu\ \ \text{with the orbital and spin content absorbed}\ [\mathrm{BM.42}(5)];
\]

%dsp-level:0
\[
\qquad \text{the endpoint formula under the}\ Q\ge0\ \text{hypothesis}\ [\mathrm{BM.41}(5)];
\]

%dsp-level:0
\[
\text{the route-(ii) variant is priced and}\ \neg\text{adopted}\ [\mathrm{BM.41}(4)(\mathrm d);\ \mathrm{BM.30}(5)];
\]

%dsp-level:0
\[
(4)\quad\textbf{the curvature rows:}
\]

%dsp-level:0
\[
\qquad \text{per member}\ [\mathrm{BM.43}]:\ \text{the baryons numerical}\ (r_s=3\,\mathfrak h,\ \kappa^2\rho=2592\pi(\ln2)^3/\mathfrak h^2,\ \text{the lock}\ 72\pi\ln2);
\]

%dsp-level:0
\[
\ \text{the leptons symbolic};
\]

%dsp-level:0
\[
\ \text{the source superposition exact};
\]

%dsp-level:0
\[
\ \text{the metric superposition linearized-scope};
\]

%dsp-level:0
\[
(5)\quad\textbf{the time rows:}
\]

%dsp-level:0
\[
\qquad \Gamma_{\mathrm{tr}}\in[0,\tfrac13]\ \text{at unit fugacity};
\]

%dsp-level:0
\[
\qquad \tau_{\mathrm{life}}\ge3\ \text{ticks};
\]

%dsp-level:0
\[
\qquad \varepsilon_q=0\implies\ \text{never}\ [\mathrm{BM.44}];
\]

%dsp-level:0
\[
\qquad \neg\text{external seconds};
\]

%dsp-level:0
\[
(6)\quad\textbf{the standing rules:}
\]

%dsp-level:0
\[
\qquad \neg\text{a P4 addition}\ [\text{§Cal-BM stays empty};\ \text{A-BM.4 is an angle identification plus a normalization convention},\ \neg\text{a numerical parameter}];
\]

%dsp-level:0
\[
\qquad \neg\text{a new dimensionful scale}\ [\mathrm{P1}];
\]

%dsp-level:0
\[
\neg\text{a measured value}\ [\mathrm{P5};\ \text{machine: zero external tokens, zero decimal numerals in this tier}];
\]

%dsp-level:0
\[
\qquad \neg\exists\ \text{an}\ [\mathrm{EXT}]\ \text{datum in any derivation};
\]

%dsp-level:0
\[
\text{species names confined to the identification layers}\ [\mathrm{Def\text{-}BM.3}(7),\ \mathrm{Def\text{-}BM.6}(7),\ \text{A-BM.5}];
\]

%dsp-level:0
\[
\qquad \text{topology labels are internal bookkeeping}\ [\mathrm{Def\text{-}H.32}(4)];
\]

%dsp-level:0
\[
\text{the extra-axiomatic list}\ [\text{A §0.2a}]\ \text{gains no member}:\ \mathrm{CT2.6}\text{'s hypotheses and}\ \mathrm{BM.8}(2)\text{'s stationarity scope are inherited at the point of use};
\]

%dsp-level:0
\[
(7)\quad\textbf{the open rows touched, not closed:}
\]

%dsp-level:0
\[
\qquad \text{the relating equation across the read species}\ [\mathrm{Def\text{-}BM.11}(7);\ \mathrm{BM.35}\ \mathrm{Rmk}\ \text{stands}];
\]

%dsp-level:0
\[
\qquad \text{the amplitude}\leftrightarrow\text{mass bridge}\ [\mathrm{BM.37}\ \text{stands}];
\]

%dsp-level:0
\[
\text{the}\ \mathrm{Sym}^k\ \text{ledger mechanism}\ [\mathrm{FT.55}(4);\ \mathrm{BM.42}(7)];
\]

%dsp-level:0
\[
\qquad \text{deficit}\ 1'\ —\ \text{the gauge coupling of}\ Q\ [\mathrm{BM.41}(3)];
\]

%dsp-level:0
\[
\qquad \neg\text{a mixing matrix is derived}\ [\mathrm{FT.145}(3);\ \mathrm{BM.41}(3)];
\]

%dsp-level:0
\[
\text{the mass-mechanism inventory: the transport-species family form joins the read-class taxonomy}\ [\mathrm{FT.84}(1)\ \text{pattern}];
\]

%dsp-level:0
\[
\qquad \text{the read specification that would open the decay kinematics}\ [\mathrm{BM.41}(4)(\mathrm e)]:\ \text{a §I addendum act},
\]

%dsp-level:0
\[
\ \neg\text{a derivation};
\]

%dsp-level:0
\[
(8)\quad\textbf{the adjudication:}
\]

%dsp-level:0
\[
\begin{array}{l|l}\text{the Yukawa requirement} & \textbf{discharged}\ —\ \neg\exists\ \text{a Yukawa term in the inventory},\ \text{and the family-profile map is constructed from the seed transport}\ [\mathrm{Def\text{-}BM.11};\ \mathrm{BM.39}]\\\text{the intra-family masses} & \textbf{derived as exact forms}\ —\ \text{the cosine/affine law, positivity, the}\ gg'\ \text{splitting, the sum rule, the order inversion, the area ladder}\ [\mathrm{BM.39}]\\\text{the neutrino} & \textbf{forced}\ \text{in every kinematically admissible vertex}\ —\ \text{the two-body impossibility}\ +\ \text{the class uniqueness};\ \mathrm{BM.23}(6)\ \text{upgraded on record};\ \text{the Casimir-deficit route adjudicated neutral}\ [\mathrm{BM.40}(3);\ \mathrm{BM.41}(2)\text{--}(3)]\\\text{angular momentum} & \textbf{conserved as the Noether moment};\ \text{orbital absorbed in}\ p^2,\ \text{spin absorbed in the rest read and the shift drag};\ \text{the CN correction exactly zero};\ \text{the}\ \ell\text{-census derived}\ [\mathrm{BM.42}]\\\text{KE and rest mass} & \textbf{derived for every member}\ \text{from the metric and the vacuum-balanced entropy}\ [\mathrm{BM.34};\ \mathrm{BM.35};\ \mathrm{BM.41}(5);\ \mathrm{BM.42}(5)]\\\text{the metric curvature} & \textbf{derived per member}\ —\ \text{numerical for the baryons, symbolic for the leptons, superposed exactly on the source side}\ [\mathrm{BM.43}]\\\text{the dynamic evolution} & \textbf{derived in ticks}\ —\ \Gamma_{\mathrm{tr}}\in[0,\tfrac13],\ \tau_{\mathrm{life}}\ge3\ \text{ticks},\ \neg\text{external seconds}\ [\mathrm{BM.44}]\\\text{the reaction} & n\to p+e^-+\bar\nu_e\ \textbf{assembled with every conservation row derived}\ —\ \textbf{and kinematically closed on every adopted read}:\ Q=-(m_\ell{+}m_\nu)<0\ [\mathrm{BM.41}(4)]\end{array}
\]

%dsp-level:0
\[
\textbf{the final read:}
\]

%dsp-level:0
\[
\qquad \text{the vertex is grammatically complete and kinematically shut};
\]

%dsp-level:0
\[
\qquad \text{what stands between the corpus and an on-shell decay is a read specification splitting the charge classes}\ —\ \text{a §I addendum act},
\]

%dsp-level:0
\[
\ \neg\text{a missing derivation}\ [\mathrm{BM.30}(5);\ \mathrm{BM.37}(3)];
\]

%dsp-level:0
\[
\qquad \text{this tier performs no such act}.
\]

\textbf{Rmk (the naming rows.)} The label \textit{electron antineutrino} attaches to the derived structural tuple — the $\nu$-channel line in the orientation sector the materialization arrow selects — exactly as \textit{antiproton} attaches to the CPT image of the proton row: the structure is derived, the occupancy sector is derived up to the reserved ray naming [FT.144(5)], and the label is logged before use (A-BM.5). Which of the three $\nu$-triplet lines the vertex emits is not derived: that attachment is mixing content, and the corpus has no mixing matrix.

\textbf{Req.} P1, P4, P5 (Parts 1–2); A6, D1.4, D1.8, D2.9 (Parts 1–2); EM.11 (Part 2); CT2.6 (Part 2); Def-H.38, FT.135(4), FT.144(2), FT.144(5), FT.145(3), FT.84(1), FT.55(4) (Part 3); Census FT.35(2), FT.44(4), FT.49(3) (row-wise); Def-BM.3, Def-BM.6, Def-BM.7, Def-BM.11, BM.2(1), BM.7, BM.11, BM.23(5), BM.25, BM.26, BM.27, BM.34, BM.35, BM.37(3), BM.39, BM.40, BM.41, BM.42, BM.43, BM.44 (this part); the declaration addendum A-BM.5 (pointers); the machine sections §A–§G green [math; machine: \texttt{checks/\allowbreak{}verify\_\allowbreak{}H\_\allowbreak{}neutron\_\allowbreak{}decay\_\allowbreak{}vertex.py} §G].

\textbf{Pf.}

%dsp-level:0
\[
(1)\text{--}(7):\ \text{assembly of the adjudicated rows of}\ \mathrm{Def\text{-}BM.11},
\]

%dsp-level:0
\[
\ \mathrm{BM.39}\text{--}\mathrm{BM.44}\ [\text{machine: §A--§F green}]
\]

%dsp-level:0
\[
(8):\ \text{conjunction of}\ (1)\text{--}(7).\ \blacksquare
\]

\textbf{Phys.} The corpus writes its first complete reaction, and the honesty of the balance is the result. Every conservation row is a theorem: charge closes on the ternary alphabet, the Janus balance rides on every emitted closed chain, the weak Casimir passes through the vertex untouched because registration signs and branch labels live in different registers, angular momentum is the conserved moment of a symmetric on-shell stress tensor, the sealing phase stays in its baryon class, and the arrow points one way. The members are forced: the charged one by the ladder rule and the alphabet, the neutral one by the impossibility of a two-body vertex at equal rest masses and the uniqueness of its class. The masses are exact where the corpus computes them — six bits over $\mathfrak h$ for both baryons — and symbolic where it does not. And then the balance itself: parent and daughter weigh the same on every read the corpus adopts, so the Q-value is minus the lepton masses, strictly negative, and the reaction does not go on-shell. The gap between the assembled vertex and a running decay is not a missing mechanism; it is a read specification that would split the charge classes, priced since Tier BM.8 and declined here as everywhere. The neutron decay vertex is the corpus in miniature: complete grammar, exact arithmetic, forced members, closed kinematics, and no external number anywhere in it.

\end{theorem}

\phantomsection
\section[Tier BM.12 — The EM Self-Interaction Correction: the Scar's…]{\texorpdfstring{Tier BM.12 — The EM Self-Interaction Correction: the Scar's Field Energy, the Infinite-Range Seal Asymmetry, the Mass-Splitting Theorem, and the Beta-Spectrum Adjudication}{Tier BM.12 - The EM Self-Interaction Correction: the Scar's Field Energy, the Infinite-Range Seal Asymmetry, the Mass-Splitting Theorem, and the Beta-Spectrum Adjudication}}

\begin{definition}[{\texorpdfstring{Def-BM.12 — The chart-side EM source read of a sealed composite (the construction Def-BM.8(8) specifies)}{Def-BM.12 - The chart-side EM source read of a sealed composite (the construction Def-BM.8(8) specifies)}}]\label{tag:Def-BM-12}
\textbf{Def.} A read constructed on the chart side of a sealed composite; no new axiom, no new datum, no new parameter, no numerical coupling. Local notation: $\varepsilon_{\mathrm{EM}}:=e^2/(4\pi)$ — a dimensionless P4-class symbol throughout, $\neg$identified with any measured constant [EM.25(7), EM.25(9); P5].

%dsp-level:0
\[
(1)\quad\textbf{the current:}
\]

%dsp-level:0
\[
\qquad J^\mu_Q:=\text{the continuum density of the integer}\ Q\ \text{read}
\]

%dsp-level:0
\[
\qquad [\mathrm{Def\text{-}BM.8}(1);\ \mathrm{BM.7}:\ \text{the density form of the integer law};\ \mathrm{BM.11}:\ \partial_\mu J^\mu=0];
\]

%dsp-level:0
\[
Q\in\mathbb Z,\ \ Q\in\{-3,\dots,+3\}\ \text{at}\ k\le3\ [\mathrm{Def\text{-}BM.8}(5);\ \mathrm{BM.28}(7)];
\]

%dsp-level:0
\[
\qquad \text{the charge alphabet of the generator matches the end-read alphabet:}\ \operatorname{spec}Q|_{\mathcal H_{\mathsf v}}\cup\omega\text{-image}
\]

%dsp-level:0
\[
= \{-1,0,+1\}\ [\mathrm{FT.12};\ \mathrm{EW.36};\ \mathrm{D3.4}];
\]

%dsp-level:0
\[
(2)\quad\textbf{the coupling:}
\]

%dsp-level:0
\[
\qquad \text{the}\ \mathsf A\text{-coefficient of the neutral connection is}\ e\,Q,
\]

%dsp-level:0
\[
\qquad e=\frac{gg'}{\sqrt{g^2+g'^2}}=g\sin\theta_W
\]

%dsp-level:0
\[
\qquad [\mathrm{EW.35};\ \mathrm{EM.25}(6):\ \text{fixed with zero freedom}];
\]

%dsp-level:0
\[
\neg\exists\ \text{a theorem fixing the numerical value of}\ e\ [\mathrm{EM.25}(7);\ \mathrm{P4}]:
\]

%dsp-level:0
\[
\qquad \textbf{every row of this tier carries}\ e\ \textbf{symbolically};
\]

%dsp-level:0
\[
(3)\quad\textbf{the sourced field equation:}
\]

%dsp-level:0
\[
\qquad \partial_\mu F^{\mu\nu}=e\,J^\nu_Q
\]

%dsp-level:0
\[
\qquad [\text{the Euler--Lagrange row of}\ S_{\mathrm{EM}}=-\tfrac14\int\sqrt{-g}\,F^2\ [\mathrm{EM.15};\ \mathrm{CT1.22}]\ \text{with the}\ (2)\text{-coupled source};\ \text{math}];
\]

%dsp-level:0
\[
\text{gauge invariance stands}\ [\mathrm{EM.14}];
\]

%dsp-level:0
\[
\qquad m_\gamma^2=0\ [\mathrm{EM.6};\ \mathrm{EW.22}:\ \ker M^2=\operatorname{span}\{Q\}]:
\]

%dsp-level:0
\[
\qquad \text{the static spherically symmetric solution outside the source support is the Coulomb read}
\]

%dsp-level:0
\[
\quad E_r(r)=\frac{e\,q}{4\pi r^2}
\]

%dsp-level:0
\[
\qquad [\text{machine: harmonicity of}\ 1/r\ +\ \text{Gauss flux normalization}];
\]

%dsp-level:0
\[
(4)\quad\textbf{the smearing locus:}
\]

%dsp-level:0
\[
\qquad \text{the source support lies within the Born-width envelope}\ \Delta x_{\mathrm B}
\]

%dsp-level:0
\[
= \mathfrak h/(E_{\mathrm{int}}\ln2)\ [\mathrm{Def\text{-}BM.4}(1),(5):\ \text{a chart-side envelope},\ \neg\text{a ledger distance};\ \mathrm{BM.13}(3):\ \text{sub-hop}];
\]

%dsp-level:0
\[
\textbf{the distribution factor is priced, not adopted}\ [\mathrm{Def\text{-}BM.4}(4)\ \text{pattern}]:
\]

%dsp-level:0
\[
\qquad \text{shell}:\ c_{\mathrm S}=\tfrac12;
\]

%dsp-level:0
\[
\qquad \text{uniform ball}:\ c_{\mathrm B}=\tfrac35
\]

%dsp-level:0
\[
\qquad [\text{machine: both energy integrals exact}];
\]

%dsp-level:0
\[
c\ \text{is carried symbolically; no clause of this tier selects a distribution};
\]

%dsp-level:0
\[
(5)\quad\textbf{the field-energy read:}
\]

%dsp-level:0
\[
\qquad E_{\mathrm{EM}}^{\mathrm{self}}(q,c):=\int_{\text{chart}}T^{(\mathrm{EM})}_{00}\,dV=\frac{c\,e^2q^2}{4\pi\,\Delta x_{\mathrm B}}
\]

%dsp-level:0
\[
\qquad [T^{(\mathrm{EM})}_{\mu\nu}=F_{\mu\rho}F_\nu{}^\rho-\tfrac14g_{\mu\nu}F^2\ [\mathrm{EM.15}];\ \text{machine: the radial integral exact with the}\ \Lambda\to\infty\ \text{row}];
\]

%dsp-level:0
\[
[\varepsilon_{\mathrm{EM}}]=\mathfrak q^0,\quad[\Delta x_{\mathrm B}]=\mathfrak q^{-1/2}\implies[E_{\mathrm{EM}}^{\mathrm{self}}]=\mathfrak q^{+1/2}=[\text{energy}]
\]

%dsp-level:0
\[
\qquad [\mathrm{CT2.1};\ \text{machine: exact exponent arithmetic}];
\]

%dsp-level:0
\[
(6)\quad\textbf{the channel location:}
\]

%dsp-level:0
\[
\qquad T_{\mu\nu}=T^{\mathrm{bd}}_{\mu\nu}+T^{\mathrm{tr}}_{\mu\nu}\ [\mathrm{CT1.22};\ \mathrm{D1.5}:\ \text{the bound/transit decomposition}]:
\]

%dsp-level:0
\[
\qquad \textbf{the correction is the transit-channel energy of the member's own field},\ \neg\text{a new mass read};
\]

%dsp-level:0
\[
\mathrm{BM.9}(4)\text{'s dust row is the bound part at rest};
\]

%dsp-level:0
\[
\qquad \text{a charged member carries}\ T^{\mathrm{bd}}+T^{\mathrm{tr}};
\]

%dsp-level:0
\[
\qquad \text{the rest input of the dispersion route refines to}\ \ m_{\mathrm{rest}}(q)
\]

%dsp-level:0
\[
= m_{\mathrm{hid}}+\frac{c\,e^2q^2}{4\pi\Delta x_{\mathrm B}}\ [\mathrm{Def\text{-}BM.10}(4)\text{--}(5);\ \mathrm{BM.35}(4)];
\]

%dsp-level:0
\[
(7)\quad\textbf{status and scope:}
\]

%dsp-level:0
\[
\qquad \textbf{constructed}\ —\ \mathrm{Def\text{-}BM.8}(8)\text{'s named open construction (the coupling of}\ Q\ \text{to a gauge field)}\ \text{is discharged for the sealed classes at static leading order};
\]

%dsp-level:0
\[
\neg\text{a declaration},\ \neg\text{a P4 addition}\ (e\ \text{is the logged combination}\ [\mathrm{EW.35}]),\ \text{§Cal-BM stays empty};
\]

%dsp-level:0
\[
\qquad \text{chart-side throughout}\ [\mathrm{CT2.6}\ \text{hypotheses at the point of use}];
\]

%dsp-level:0
\[
\mathrm{BM.7}(4)\text{'s register row stands untouched:}\ \neg\exists\ \text{a}\ U(1)\ \text{action on the ternary alphabet and}\ q_{\mathrm{EM}}\ \text{is}\ \neg\text{a register Noether charge}\ —\ \text{the coupling of this definition is chart-side, after}\ \mathrm{BM.8}\text{'s promotion}\ [\text{``the gauge algebras were never discrete''}];
\]

%dsp-level:0
\[
\neg\text{constructed: any dynamics of}\ Q\ [\mathrm{BM.6}(5)],
\]

%dsp-level:0
\[
\ \text{radiation and back-reaction},
\]

%dsp-level:0
\[
\ \text{a lepton--baryon bound channel}\ [\mathrm{Def\text{-}BM.1}\text{'s grammar is colour-sealed}],
\]

%dsp-level:0
\[
\ \text{the chiral assignment}\ [\text{deficit}\ 1',\ \mathrm{FT.125}(4)\ \text{stands}],
\]

%dsp-level:0
\[
\ \text{internal multipole resolution}\ [\mathrm{BM.46}(2)].
\]

\textbf{Rmk (what is declared and what is not).} Nothing. The coupling strength is EW.35's derived coefficient carried by a P4-class symbol; the current is Def-BM.8's derived read in BM.11's derived density form; the field equation is the Euler–Lagrange row of EM.15's derived action; the locus is Def-BM.4(5)'s existing envelope; the distribution factor is priced in two exact values and adopted in none. The construction assembles derived pieces into the sourced static read that Def-BM.8(8) listed as unconstructed — and stops exactly where the inventory stops: no dynamics, no radiation, no atoms.

\textbf{Req.} D1.5, D1.6, D2.9, D3.4 (Parts 1–2); CT1.22, CT2.1, CT2.6, EM.6, EM.11, EM.14, EM.15, EM.25(6), EM.25(7), EM.25(9), EW.22, EW.35, EW.36 (Part 2); P4, P5 (Parts 1–2); FT.12 (Part 3); Def-BM.4, Def-BM.8, Def-BM.10, BM.7, BM.9(4), BM.11, BM.13(3), BM.28(7), BM.35(4) (this part); the Euler–Lagrange row of the abelian action, Gauss's law for the spherically symmetric static class, exact radial integrals [math; machine: \texttt{checks/\allowbreak{}verify\_\allowbreak{}H\_\allowbreak{}em\_\allowbreak{}self\_\allowbreak{}interaction.py} §A].

\textbf{Pf.}

%dsp-level:0
\[
\text{Definition block: clauses}\ (1)\text{--}(6)\ \text{are assemblies of the cited derived rows with the machine-verified Coulomb/energy identities};
\]

%dsp-level:0
\[
\qquad (4)\text{'s priced factors carry no adoption};
\]

%dsp-level:0
\[
\qquad (7)\ \text{is the status and scope record}.\ \blacksquare
\]

\textbf{Phys.} The composite's charge was an integer read with no field attached to it; the attachment is now built, and every piece of it was already on the shelf. The current is the additive read in its continuum density form, the coupling coefficient is the residual coupling the electroweak rotation derived — carried as a symbol, because no theorem of the corpus fixes its value and none is invented here — and the field equation is the Euler–Lagrange row of the Maxwell action the stress-energy tier already varied. Outside the Born-width envelope the solution is the Coulomb read, its energy integral is exact at either of two priced distributions, and its home in the ledger bookkeeping is the transit channel of CT1.22's own decomposition: the correction is the field the member drags along, not a new species of mass. The register-side obstruction stands where it was: there is no $U(1)$ on the ternary alphabet, and this construction never asks for one — it lives on the chart, downstream of the promotion.

\end{definition}

\begin{theorem}[{\texorpdfstring{BM.46 — The scar's field energy and the mass-splitting theorem}{BM.46 - The scar's field energy and the mass-splitting theorem}}]\label{tag:BM-46}
\textbf{Thm.}

%dsp-level:0
\[
(1)\quad\textbf{the self-energy, two agreeing routes:}
\]

%dsp-level:0
\[
\qquad \textbf{field route:}\quad\int_{\Delta x_{\mathrm B}}^{\Lambda}\tfrac12E_r^2\cdot4\pi r^2\,
\]

%dsp-level:0
\[
dr
\]

%dsp-level:0
\[
= \frac{e^2q^2}{8\pi}\Bigl(\frac1{\Delta x_{\mathrm B}}-\frac1\Lambda\Bigr)\xrightarrow[\Lambda\to\infty]{}\frac{e^2q^2}{8\pi\,\Delta x_{\mathrm B}}
\]

%dsp-level:0
\[
\qquad [\text{machine: exact antiderivative and limit row}];
\]

%dsp-level:0
\[
\textbf{potential route (shell):}\quad\tfrac12\,q\,\phi(\Delta x_{\mathrm B})=\tfrac12(eq)\frac{eq}{4\pi\Delta x_{\mathrm B}}=\frac{e^2q^2}{8\pi\,\Delta x_{\mathrm B}}
\]

%dsp-level:0
\[
\qquad [\text{machine: the two routes agree exactly}];
\]

%dsp-level:0
\[
\textbf{the uniform ball:}\quad E_{\mathrm{in}}=\frac{e^2q^2}{40\pi\,\Delta x_{\mathrm B}},\quad E_{\mathrm{out}}=\frac{e^2q^2}{8\pi\,\Delta x_{\mathrm B}},\quad E_{\mathrm{tot}}=\frac{3\,e^2q^2}{20\pi\,\Delta x_{\mathrm B}}
\]

%dsp-level:0
\[
\qquad [\text{machine: the interior/exterior split exact}];
\]

$\textbf{the unified form:}\qquad\Delta x_{\mathrm B}\cdot m_{\mathrm{hid}}=1\ [\mathrm{Def\text{-}BM.4}(1)]:\qquad\boxed{\ E_{\mathrm{EM}}^{\mathrm{self}}=\frac{c\,e^2q^2}{4\pi\,\Delta x_{\mathrm B}}=c\,\varepsilon_{\mathrm{EM}}\,q^2\,m_{\mathrm{hid}}=\frac{c\,\varepsilon_{\mathrm{EM}}\,q^2\,E_{\mathrm{int}}\ln2}{\mathfrak h}\ }\qquad[\text{machine: one Laurent monomial per form}];$ $\text{i.e.}\qquad\textbf{the fractional correction is}\ \ \frac{E_{\mathrm{EM}}^{\mathrm{self}}}{m_{\mathrm{hid}}}=c\,\varepsilon_{\mathrm{EM}}\,q^2\in\mathfrak q^0\qquad[\text{machine: dimensionless}];$

%dsp-level:0
\[
(2)\quad\textbf{the charge selection:}
\]

%dsp-level:0
\[
\qquad \text{the correction is exactly proportional to}\ q^2:
\]

%dsp-level:0
\[
\qquad \text{the unit class}\ (q{=}{+}1):\ \delta^{(p)}=c\,\varepsilon_{\mathrm{EM}}\,m_{\mathrm{hid}}>0;
\]

%dsp-level:0
\[
\qquad \text{the neutral class}\ (q{=}0):\ \delta^{(n)}
\]

%dsp-level:0
\[
= 0\ \ \textbf{exactly at the monopole order}\ [\text{machine}:\ 0^2=0;\ \text{Gauss: zero net flux}];
\]

%dsp-level:0
\[
\textbf{the internal-structure row:}\quad\text{the neutral class contains separated-sign endpoint patterns}\ (12\ \text{of the}\ 20\ \text{probability-space members are mixed triples}\ [\mathrm{BM.28}(2);\ \text{machine}]),
\]

%dsp-level:0
\[
\ \text{but}\ \mathrm{D3.4}\text{'s end signs carry no position datum}:
\]

%dsp-level:0
\[
\text{the chart does not resolve internal multipoles at this order}\ [\mathrm{Def\text{-}BM.4}(5)\ \text{envelope locus}]:
\]

%dsp-level:0
\[
\qquad \text{a neutron internal-multipole self-energy is}\ \neg\text{derived}\ —\ \text{a resolution residue, recorded};
\]

%dsp-level:0
\[
(3)\quad\textbf{the corrected rest read:}
\]

%dsp-level:0
\[
\qquad m_{\mathrm{eff}}(q)=m_{\mathrm{hid}}+\frac{c\,e^2q^2}{4\pi\,\Delta x_{\mathrm B}}=m_{\mathrm{hid}}\bigl(1+c\,\varepsilon_{\mathrm{EM}}\,q^2\bigr)
\]

%dsp-level:0
\[
\qquad [\mathrm{Def\text{-}BM.12}(6):\ \text{the transit-channel refinement of the rest input}];
\]

%dsp-level:0
\[
\text{at}\ E_{\mathrm{int}}=6\ [\mathrm{BM.17};\ \mathrm{BM.26}]:
\]

%dsp-level:0
\[
\qquad \delta^{(p)}=\frac{6\,c\,e^2\ln2}{4\pi\,\mathfrak h}=\frac{3\,c\,e^2\ln2}{2\pi\,\mathfrak h}
\]

%dsp-level:0
\[
\qquad [\text{machine: one Laurent monomial}];
\]

%dsp-level:0
\[
\qquad [E_{\mathrm{EM}}^{\mathrm{self}}]=\mathfrak q^{+1/2}=[\text{energy}]\ [\mathrm{CT2.1};\ \text{machine}];
\]

%dsp-level:0
\[
(4)\quad\textbf{the mass-splitting theorem:}
\]

%dsp-level:0
\[
\qquad \boxed{\ \Delta m_{\mathrm{EM}}:=m_{\mathrm{eff}}^{(p)}-m_{\mathrm{eff}}^{(n)}=\frac{c\,e^2}{4\pi\,\Delta x_{\mathrm B}}=\frac{3\,c\,e^2\ln2}{2\pi\,\mathfrak h}\ >\ 0\ }
\]

%dsp-level:0
\[
\qquad [\text{machine: positive coefficient census at}\ e^2>0,\ c\in\{\tfrac12,\tfrac35\}];
\]

%dsp-level:0
\[
\textbf{the equal-mass degeneracy of the crystallized-count route}\ [\mathrm{BM.28}(1)]\ \textbf{is broken by a strictly positive, charge-selected splitting: the charged member is the heavier one};
\]

%dsp-level:0
\[
\text{the width is shared}\ (E_{\mathrm{int}}=6\ \text{on both classes}\ [\mathrm{BM.17};\ \mathrm{BM.26}])\ \text{and the selector is}\ q^2:\ (1\ \text{vs}\ 0)\ [\text{machine}]:
\]

%dsp-level:0
\[
\qquad \text{the splitting is a function of the charge and of the Born width, with the charge carrying the selection};
\]

%dsp-level:0
\[
(5)\quad\textbf{the correction-species row:}
\]

%dsp-level:0
\[
\qquad \delta_{\mathrm{EM}}\ \text{is a correction to the existing read,}\ \neg\text{a new mass read}\ [\mathrm{Def\text{-}BM.12}(6);\ \mathrm{CT1.22};\ \mathrm{D1.5}]:\ \text{it is the transit-channel energy of the member's own field, entering through the dispersion route's rest slot}\ [\mathrm{Def\text{-}BM.10}(4)\text{--}(5)];
\]

%dsp-level:0
\[
\mathrm{BM.37}(2)\text{'s disjointness stands: the measure routes}\ (\mathrm i)\text{--}(\mathrm{iii})\ \text{remain measure functionals}\ —\ \text{the correction rides the metric/field species, and the route ratios}\ \bigl(1,\tfrac43,\tfrac{58}{13}\bigr)\ \text{are untouched by it};
\]

%dsp-level:0
\[
(6)\quad\textbf{status:}
\]

%dsp-level:0
\[
\qquad \text{derived:}\ (1)\text{--}(5);
\]

%dsp-level:0
\[
\qquad \text{symbolic:}\ e\ [\mathrm{EM.25}(7)],\ c\ [\mathrm{Def\text{-}BM.12}(4)];
\]

%dsp-level:0
\[
\qquad \neg\text{a measured splitting value enters}\ [\mathrm{P5};\ \text{machine: zero external tokens}];
\]

%dsp-level:0
\[
\neg\text{a P4 addition},\ \neg\text{a new dimensionful scale}\ [\mathrm{P1}];
\]

%dsp-level:0
\[
\qquad \text{the open rows: the numerical value of}\ e\ (\text{a calibration step requiring one datum}\ [\mathrm{EM.25}(7)]),
\]

%dsp-level:0
\[
\ \text{the internal-multipole resolution}\ ((2)),
\]

%dsp-level:0
\[
\ \text{the dynamics of}\ Q\ [\mathrm{BM.6}(5)].
\]

\textbf{Rmk (the sign is the theorem).} Every piece of this correction is positive-semidefinite in the charge: the field energy of a source is the integral of a square, and it scales as $q^2$. The member that carries the scar therefore gains mass, never loses it — the proton is the heavier baryon on this route, by an exact symbolic monomial. Any account on which the electromagnetic self-interaction lightens the charged member and thereby opens the neutron decay is refuted by the sign of the integral, before any value of $e$ is known. The tier keeps the distribution factor and the coupling symbolic precisely so that the sign — the only content the claim needed — is derived without either.

\textbf{Req.} D1.5, D2.9, D3.4 (Parts 1–2); CT1.22, CT2.1, EM.11, EM.15, EM.25(7) (Part 2); P1, P4, P5 (Parts 1–2); Def-BM.4, Def-BM.8, Def-BM.10, Def-BM.12, BM.9(4), BM.13, BM.17, BM.26, BM.28(1), BM.28(2), BM.35(4), BM.37(2) (this part); exact radial integrals with the cutoff-and-limit row, Laurent monomials, sign censuses [math; machine: \texttt{checks/\allowbreak{}verify\_\allowbreak{}H\_\allowbreak{}em\_\allowbreak{}self\_\allowbreak{}interaction.py} §A–§B].

\textbf{Pf.}

%dsp-level:0
\[
(1):\ E_r=\tfrac{eq}{4\pi r^2}\ [\mathrm{Def\text{-}BM.12}(3)]:\quad\tfrac12E_r^2\cdot4\pi r^2=\tfrac{e^2q^2}{8\pi r^2};
\]

%dsp-level:0
\[
\qquad \int_R^\Lambda\tfrac{dr}{r^2}=\tfrac1R-\tfrac1\Lambda\ [\text{machine: antiderivative}];
\]

%dsp-level:0
\[
\qquad \text{the shell potential route and the ball split}\ [\text{machine: exact rational coefficients}\ \tfrac18,\ \tfrac1{40},\ \tfrac3{20}]
\]

%dsp-level:0
\[
\text{the unified form}:\ \tfrac1{\Delta x_{\mathrm B}}
\]

%dsp-level:0
\[
= m_{\mathrm{hid}}\ [\mathrm{Def\text{-}BM.4}(1)]:\quad\tfrac{c\,e^2q^2}{4\pi\Delta x_{\mathrm B}}
\]

%dsp-level:0
\[
= c\,
\]

%dsp-level:0
\[
\tfrac{e^2}{4\pi}\,
\]

%dsp-level:0
\[
q^2\,
\]

%dsp-level:0
\[
m_{\mathrm{hid}}\ [\text{machine: monomial identity}];
\]

%dsp-level:0
\[
\qquad [\varepsilon_{\mathrm{EM}}]=[e^2]=\mathfrak q^0\ [\mathrm{EW.35}:\ e\ \text{dimensionless};\ \text{machine}]
\]

%dsp-level:0
\[
(2):\ q=0\implies e^2q^2=0\ [\text{machine}];
\]

%dsp-level:0
\[
\qquad W(0)=20=8+12\ \text{with}\ 12\ \text{mixed triples}\ [\mathrm{BM.28}(2)\ \text{re-enumerated; machine}];
\]

%dsp-level:0
\[
\qquad \text{the position-datum row:}\ \mathrm{D3.4}\text{'s read is a pair of signs}\ [\text{verbatim scan}]
\]

%dsp-level:0
\[
(3):\ \mathrm{Def\text{-}BM.12}(6)\ \text{substituted};
\]

%dsp-level:0
\[
\qquad \Delta x_{\mathrm B}=\mathfrak h/(6\ln2)\ [\mathrm{BM.26}\ \text{re-cited}]:\ \tfrac{c\,e^2}{4\pi}\cdot\tfrac{6\ln2}{\mathfrak h}\ [\text{machine: one monomial}];
\]

%dsp-level:0
\[
\qquad \text{the exponents}\ [\mathrm{CT2.1};\ \text{machine}]
\]

%dsp-level:0
\[
(4):\ \text{the difference of}\ (3)\text{'s rows at}\ q^2\in\{1,0\}\ [\text{machine}];
\]

%dsp-level:0
\[
\qquad c\in\{\tfrac12,\tfrac35\}>0,\ e^2>0\ [\mathrm{EW.33}:\ g,g'>0]\implies\Delta m_{\mathrm{EM}}>0\ [\text{machine: coefficient census}]
\]

%dsp-level:0
\[
(5)\text{--}(6):\ \text{channel and status records}\ [\mathrm{CT1.22};\ \mathrm{D1.5};\ \mathrm{BM.37}(2)\ \text{re-cited}].\ \blacksquare
\]

\textbf{Phys.} The scar is no longer only a label. A unit charge smeared inside its own Born width carries a field, the field carries the integral of a square, and the integral is exact: one eighth over the width on the shell convention, three twentieth on the ball, one symbolic coupling factor on both. The correction to the rest mass is that energy and nothing else — it rides the transit channel the stress-energy theorem already split out, so the read stays one read, refined. The neutron, net-neutral, gets exactly zero at the monopole order, and the tier records honestly what zero does not mean: the neutral class is full of separated sign patterns, but signs without positions carry no multipole, and the chart resolves no internal structure at this order. What falls out is the first derived splitting of the two baryon rows — strictly positive, charge-selected, width-shared — and its sign is the whole point: electromagnetism makes the charged member heavier. The degeneracy of the crystallized count is broken, and broken in the direction that shuts the decay tighter, not looser.

\end{theorem}

\begin{theorem}[{\texorpdfstring{BM.47 — The infinite-range seal: UV/IR finiteness, the universe-scale neutrality, and the seal asymmetry}{BM.47 - The infinite-range seal: UV/IR finiteness, the universe-scale neutrality, and the seal asymmetry}}]\label{tag:BM-47}
\textbf{Thm.}

%dsp-level:0
\[
(1)\quad\textbf{masslessness is the range:}
\]

%dsp-level:0
\[
\qquad m_\gamma^2=0\ [\mathrm{EM.6};\ \mathrm{EW.22};\ \mathrm{EM.25}(2)]:
\]

%dsp-level:0
\[
\qquad \text{the static potential is}\ \tfrac1r\ \text{and the field}\ \tfrac1{r^2},
\]

%dsp-level:0
\[
\ \text{with no exponential damping factor}\ [\text{machine: the Yukawa factor}\ e^{-mr}\ \text{is}\ 1\ \text{at}\ m=0];
\]

%dsp-level:0
\[
\text{contrast class: colour transit modes are likewise massless}\ [\mathrm{EW.26}]\ \text{but colour is locally projected}\ [(5)]\ —\ \text{range and sealing are independent rows};
\]

%dsp-level:0
\[
(2)\quad\textbf{UV finiteness at the Born width:}
\]

%dsp-level:0
\[
\qquad \int_R^\Lambda\frac{dr}{r^2}=\frac1R-\frac1\Lambda<\frac1R\ \ \forall\Lambda>R>0
\]

%dsp-level:0
\[
\qquad [\text{machine: exact rational functions}];
\]

%dsp-level:0
\[
\Delta x_{\mathrm B}>0\ \text{strictly}\ [\mathrm{BM.13}(1);\ \mathrm{Def\text{-}BM.4}(1)]:
\]

%dsp-level:0
\[
\qquad E_{\mathrm{EM}}^{\mathrm{self}}<\infty\ \ \text{for every admissible cell}\ [\mathrm{BM.13}(3):\ E_{\mathrm{int}}\ge2];
\]

%dsp-level:0
\[
\qquad \text{the point-charge divergence is cut by the envelope,}\ \neg\text{by a counterterm};
\]

%dsp-level:0
\[
(3)\quad\textbf{IR finiteness of an isolated member:}
\]

%dsp-level:0
\[
\qquad E_{\mathrm{out}}(R)=\frac{e^2q^2}{8\pi R}<\infty\ \ \text{as}\ \ \Lambda\to\infty
\]

%dsp-level:0
\[
\qquad [\text{machine: the}\ 1/\Lambda\ \text{row vanishes in the limit}]:
\]

%dsp-level:0
\[
\text{the field energy outside any positive radius is finite although the field itself never vanishes}\ —\ \textbf{the infinite range costs no divergence at fixed net charge};
\]

%dsp-level:0
\[
(4)\quad\textbf{the universe-scale seal:}
\]

%dsp-level:0
\[
\qquad \textbf{Gauss:}\ \ \oint_{S^2_\infty}\vec E\cdot d\vec A=e\,Q_{\mathrm{enc}}
\]

%dsp-level:0
\[
\qquad [\text{machine: the flux normalization row}];
\]

%dsp-level:0
\[
\qquad \textbf{EM.11:}\ \ Q_{\mathrm{EM}}=0\ \text{in every valid universe}\ [\mathrm{D2.11};\ \mathrm{D1.3}]:
\]

%dsp-level:0
\[
\text{a sphere exhausting the universe encloses}\ Q_{\mathrm{enc}}
\]

%dsp-level:0
\[
= 0\implies\textbf{the total asymptotic flux vanishes}:\ \text{the leading}\ \tfrac1{r^2}\ \text{term of the total field cancels at universe scale}\ [\text{machine: the integer row}\ e\cdot0=0];
\]

%dsp-level:0
\[
\text{the residual (multipole) field decays as}\ r^{-3}:
\]

%dsp-level:0
\[
\qquad \int^\infty\Bigl(\frac1{r^3}\Bigr)^2r^2\,dr=\int^\infty\frac{dr}{r^4}<\infty
\]

%dsp-level:0
\[
\qquad [\text{machine: exponent arithmetic}]:
\]

%dsp-level:0
\[
\boxed{\ \textbf{the total EM field energy of a valid universe is IR-finite by}\ \mathrm{EM.11}\text{'s global neutrality}\ }\ —\ \textbf{EM seals at universe scale};
\]

%dsp-level:0
\[
(5)\quad\textbf{the seal asymmetry (colour local, EM global):}
\]

%dsp-level:0
\[
\qquad \textbf{colour:}\ \ \mathcal C(\Sigma;C)\in\{0,1\}\ \text{and one dangling attachment zeroes the whole colour sector}\ [\mathrm{D2.14};\ \mathrm{FT.19};\ \mathrm{BM.1};\ \mathrm{BM.32}(1):\ \text{filtration}]:\ \text{a local}\ 0/1\ \text{projector};
\]

%dsp-level:0
\[
\textbf{EM:}\ \ W(\pm1)=15>0\ [\mathrm{BM.28}(2);\ \mathrm{Def\text{-}BM.8}(5):\ \text{the}\ W(Q)\ \text{census}]:
\]

%dsp-level:0
\[
\qquad \text{local charge classes are admissible};
\]

%dsp-level:0
\[
\qquad \neg\exists\ \text{a local projector on}\ Q;
\]

%dsp-level:0
\[
\qquad \text{neutrality is exhausted only at universe scale}\ [\mathrm{EM.11}];
\]

%dsp-level:0
\[
[\text{machine: the two vocabularies}\ —\ \text{the idempotent}\ \{0,1\}\ \text{against the positive census}];
\]

%dsp-level:0
\[
\qquad \textbf{colour self-seals; EM does not};
\]

%dsp-level:0
\[
(6)\quad\textbf{the nearby-space fractions:}
\]

%dsp-level:0
\[
\qquad \text{the fraction of the self-energy inside}\ k\,\Delta x_{\mathrm B}:
\]

%dsp-level:0
\[
\qquad f(k)=1-\frac1k
\]

%dsp-level:0
\[
\qquad [\text{machine: exact rational functions; the shell route}];
\]

%dsp-level:0
\[
f(2)=\tfrac12,\quad f(3)=\tfrac23,\quad f(4)=\tfrac34;
\]

%dsp-level:0
\[
\qquad \text{the exterior fraction}\ \tfrac1k>0\ \ \forall k<\infty:
\]

%dsp-level:0
\[
\qquad \textbf{no finite neighbourhood holds the whole correction}\ —\ \text{the nearby-space reading is a strict fraction at every finite radius};
\]

%dsp-level:0
\[
(7)\quad\textbf{status:}
\]

%dsp-level:0
\[
\qquad \text{the correction is local in its ultraviolet}\ ((2)),
\]

%dsp-level:0
\[
\ \text{finite in its infrared per member}\ ((3))\ \text{and per universe}\ ((4));
\]

%dsp-level:0
\[
\qquad \text{the seal asymmetry}\ ((5))\ \text{is the exact content of ``colour seals, EM seals at universe scale''};
\]

%dsp-level:0
\[
\neg\text{a measured range, radius or flux enters}\ [\mathrm{P5}].
\]

\textbf{Rmk (two seals, two grammars.)} Confinement and neutrality are both sealing laws, and the tier keeps their grammar apart. Colour's is a projector: it acts on any restriction, returns zero or one, and annihilates a block at the first dangling attachment — locality is its content. Electromagnetism's is a census: every local charge class has positive weight, nothing is projected away nearby, and the vanishing happens once, globally, because a valid universe has no dangling ends at all. The infinite range is not an obstacle to either finiteness row: the Born envelope absorbs the ultraviolet, and EM.11's integer law absorbs the infrared.

\textbf{Req.} D1.3, D2.11, D2.14 (Parts 1–2); CT3.18, EM.6, EM.11, EM.25(2), EW.22, EW.26 (Part 2); P5 (Parts 1–2); FT.19 (Part 3); Def-BM.4, Def-BM.8, Def-BM.12, BM.1, BM.13(1), BM.13(3), BM.26, BM.28(2), BM.32(1), BM.46 (this part); exact radial integrals and exponent arithmetic, the Gauss flux normalization, integer censuses [math; machine: \texttt{checks/\allowbreak{}verify\_\allowbreak{}H\_\allowbreak{}em\_\allowbreak{}self\_\allowbreak{}interaction.py} §B].

\textbf{Pf.}

%dsp-level:0
\[
(1):\ \mathrm{EM.6}/\mathrm{EW.22}\ \text{re-cited};\qquad e^{-m r}\big|_{m=0}=1\ [\text{machine: exact at}\ m=0]
\]

%dsp-level:0
\[
(2):\ \tfrac{d}{dr}\bigl(-\tfrac1r\bigr)
\]

%dsp-level:0
\[
= \tfrac1{r^2}\ [\text{machine: symbolic differentiation}]\implies\int_R^\Lambda
\]

%dsp-level:0
\[
= \tfrac1R-\tfrac1\Lambda\ [\text{machine}];
\]

%dsp-level:0
\[
\qquad \Delta x_{\mathrm B}=\mathfrak h/(E\ln2)>0\ [\mathrm{BM.13}(1)\ \text{re-cited}]
\]

%dsp-level:0
\[
(3):\ \text{the}\ \Lambda\text{-row of}\ (2)\ \text{at}\ R=\Delta x_{\mathrm B}\ [\text{machine}]
\]

%dsp-level:0
\[
(4):\ \text{the flux of}\ \tfrac{eq}{4\pi r^2}\ \text{through}\ S^2_r\ \text{is}\ eq\ [\mathrm{Def\text{-}BM.12}(3)\text{'s Gauss normalization re-cited};\ \text{machine}];
\]

%dsp-level:0
\[
\qquad Q_{\mathrm{enc}}=Q_{\mathrm{EM}}=0\ [\mathrm{EM.11}\ \text{re-cited};\ \text{machine: integers}];
\]

%dsp-level:0
\[
\qquad \int^\infty r^{-4}dr\ \text{convergent}\ [\text{machine: antiderivative}\ -\tfrac1{3r^3}]
\]

%dsp-level:0
\[
(5):\ \mathcal C^2=\mathcal C\in\{0,1\}\ [\mathrm{BM.32}(1)\ \text{re-cited}];
\]

%dsp-level:0
\[
\qquad W(\pm1)=15\ [\mathrm{BM.28}(2)\ \text{re-enumerated};\ \text{machine}];
\]

%dsp-level:0
\[
\qquad \neg\exists\ \text{a local}\ Q\text{-projector: inventory scan of}\ \mathrm{D2.14}\text{'s vocabulary (colour slots only)}\ [\text{machine: token row}]
\]

%dsp-level:0
\[
(6):\ f(k)=\tfrac{E(<k\Delta x_{\mathrm B})}{E_{\mathrm{tot}}}=\tfrac{1/R-1/(kR)}{1/R}=1-\tfrac1k\ [\text{machine: rational functions}];
\]

%dsp-level:0
\[
\qquad (7):\ \text{the status record}.\ \blacksquare
\]

\textbf{Phys.} The electromagnetic correction lives with an infinite-range field and stays finite at both ends, for two different reasons that the tier separates. In the ultraviolet the Born envelope does the work: the composite is never a point, its width is strictly positive on every admissible cell, and the radial integral of the field square converges to an exact rational over that width. In the infrared the universe does the work: Gauss ties the asymptotic flux to the enclosed charge, and a valid universe encloses zero — so the total field's leading term cancels at scale, the residual multipoles fall as the cube, and the energy beyond any sphere is finite. Between the two sits the asymmetry with confinement: colour has a local zero-one projector that kills a block at one dangling attachment, while electromagnetism has no local projector at all — its charge classes are positively populated nearby and its neutrality is a single global fact about ends. And the nearby space holds only a fraction: half the field energy lives beyond two widths, two thirds within three, and a strict $1/k$ always escapes — the correction is local in origin and never local in full.

\end{theorem}

\begin{theorem}[{\texorpdfstring{BM.48 — The two-body potential, the hydrogenic rows, and the binding category bar}{BM.48 - The two-body potential, the hydrogenic rows, and the binding category bar}}]\label{tag:BM-48}
\textbf{Thm.}

%dsp-level:0
\[
(1)\quad\textbf{the static two-body potential:}
\]

%dsp-level:0
\[
\qquad V(d)=\frac{e^2q_1q_2}{4\pi\,d}
\]

%dsp-level:0
\[
\qquad [\mathrm{Def\text{-}BM.12}(3)\text{'s Coulomb read};\ \text{math: the unique spherically symmetric harmonic function with the prescribed Gauss flux}];
\]

%dsp-level:0
\[
q_1q_2<0\implies V<0\ (\text{attraction}),\qquad q_1q_2>0\implies V>0\ (\text{repulsion})\qquad[\text{machine: sign rows}];
\]

%dsp-level:0
\[
\text{the pair}\ (q_p,q_\ell)=(+1,-1)\ [\mathrm{BM.23}(4)\text{'s forced carrier charge}]:
\]

%dsp-level:0
\[
\qquad V(d)=-\frac{e^2}{4\pi\,d}=-\frac{\varepsilon_{\mathrm{EM}}}{d}
\]

%dsp-level:0
\[
\qquad [\text{machine: exact coefficient}];
\]

%dsp-level:0
\[
(2)\quad\textbf{the hydrogenic rows (chart-side, conditional on a bound channel):}
\]

%dsp-level:0
\[
\qquad \text{under the canonical-pair mechanics of the chart}\ [\mathrm{BM.10}(3):\ [\hat x,\hat p]=i,\ \hbar=1]\ \text{with reduced mass}\ m_r:
\]

%dsp-level:0
\[
E_n=-\frac{e^4(q_1q_2)^2\,m_r}{32\pi^2\,n^2}=-\frac{\varepsilon_{\mathrm{EM}}^2\,m_r\,(q_1q_2)^2}{2\,n^2},
\]

%dsp-level:0
\[
\qquad \frac{E_2}{E_1}=\frac14\quad\text{exact}
\]

%dsp-level:0
\[
\qquad [\text{math: the Coulomb spectrum};\ \text{machine: the coefficient and ladder rows}];
\]

%dsp-level:0
\[
\text{the}\ n\text{-ladder is a chart-side spectrum};
\]

%dsp-level:0
\[
\qquad m_r\ \text{of a lepton--baryon pair requires the lepton mass read}\ [\mathrm{BM.43}(2):\ \text{symbolic};\ \text{mass-form bridge residue}];
\]

%dsp-level:0
\[
(3)\quad\textbf{the bound channel is}\ \neg\textbf{constructed:}
\]

%dsp-level:0
\[
\qquad \mathrm{Def\text{-}BM.1}\text{'s composite grammar is colour-sealed}\ (\varepsilon/\delta\ \text{junctions on slot indices}\ [\mathrm{D3.9};\ \mathrm{CT3.17}]);
\]

%dsp-level:0
\[
\qquad \neg\exists\ \text{an EM sealing grammar in the inventory};
\]

%dsp-level:0
\[
\text{a lepton--baryon bound state would be a dynamics-of-}Q\ \text{object}\ [\mathrm{BM.6}(5):\ \neg\text{constructed}];
\]

%dsp-level:0
\[
\qquad \text{the photon-exchange binding of}\ \mathrm{BM.41}(3)\text{'s transit alternative remains}\ \neg\text{constructed}\ [\text{deficit}\ 1'\ \text{untouched}];
\]

$(4)\quad\textbf{the binding category bar:}\qquad\textbf{a binding energy is the energy of the bound system,}\ \neg\textbf{a correction to a free member's mass}:$

%dsp-level:0
\[
\text{for a constructed bound pair the total read would be}\ \ m_{\mathrm{pair}}=m_p+m_\ell+E_{\mathrm{bind}},
\]

%dsp-level:0
\[
\quad E_{\mathrm{bind}}<0\ [\text{system read};\ (2)\text{'s}\ E_1\ \text{at}\ n=1];
\]

%dsp-level:0
\[
\text{the free proton's correction is}\ +E_{\mathrm{EM}}^{\mathrm{self}}\ \text{alone}\ [\mathrm{BM.46}(3)]:
\]

%dsp-level:0
\[
\qquad \boxed{\ m_{\mathrm{eff}}^{(p)}=m_{\mathrm{hid}}+E_{\mathrm{EM}}^{\mathrm{self}}-E_{\mathrm{EM}}^{\mathrm{bind}}\ \ \text{is refuted as a category error}\ }\ —\ \neg\exists\ \text{an object in the inventory whose mass that expression reads};
\]

%dsp-level:0
\[
[\text{machine: the inventory row}\ —\ \text{no bound-pair mass read exists};\ (3)\ \text{re-cited}];
\]

%dsp-level:0
\[
(5)\quad\textbf{the sign rows survive the hypothetical:}
\]

%dsp-level:0
\[
\qquad E_{\mathrm{EM}}^{\mathrm{self}}>0\ \text{unconditionally}\ [\mathrm{BM.46}(4)];
\]

%dsp-level:0
\[
\qquad E_{\mathrm{bind}}<0\ \text{belongs to the pair}\ [(4)]:
\]

%dsp-level:0
\[
\text{even granting the hypothetical pair, the free-proton splitting}\ \Delta m_{\mathrm{EM}}>0\ \text{is unchanged}\ [\mathrm{BM.46}(4)\ \text{stands}];
\]

%dsp-level:0
\[
(6)\quad\textbf{status:}
\]

%dsp-level:0
\[
\qquad \text{derived: the potential row, its sign structure, the hydrogenic coefficients as chart-side math};
\]

%dsp-level:0
\[
\qquad \neg\text{constructed: the bound channel, the pair mass read};
\]

%dsp-level:0
\[
\qquad \neg\text{adopted: any}\ E_{\mathrm{bind}}\ \text{term in a free-member mass};
\]

%dsp-level:0
\[
\neg\text{a measured binding energy, ionization scale or atomic datum enters}\ [\mathrm{P5}].
\]

\textbf{Rmk (what the potential row buys.)} The Coulomb read between the forced carrier and the daughter is derived the moment the sourced field exists: opposite ternary charges attract with an exact symbolic coefficient, and the chart-side hydrogenic ladder follows from the canonical pair the Noether tier already carried. What does not exist is a grammar that seals the pair — the corpus's composites are colour-sealed, and an electromagnetic atom is a dynamics-of-$Q$ construction nobody has built. The tier draws the potential, prices the ladder, and bars the category slip: binding energy is a property of the bound, and until a bound is constructed, no free member's mass carries one.

\textbf{Req.} D1.5, D3.9 (Parts 1–2); CT2.1, CT3.17 (Part 2); P5 (Parts 1–2); Def-BM.1, Def-BM.12, BM.6(5), BM.10(3), BM.23(4), BM.41(3), BM.43(2), BM.46(3), BM.46(4) (this part); the uniqueness of the harmonic spherically symmetric potential with prescribed flux, the Coulomb spectrum on the canonical pair [math; machine: \texttt{checks/\allowbreak{}verify\_\allowbreak{}H\_\allowbreak{}em\_\allowbreak{}self\_\allowbreak{}interaction.py} §C].

\textbf{Pf.}

%dsp-level:0
\[
(1):\ \nabla^2\tfrac1r=0\ \text{at}\ r\ne0\ [\mathrm{BM.43}(4)\text{'s rule-ring identity re-cited};\ \text{machine}];
\]

%dsp-level:0
\[
\qquad \text{the flux normalization}\ [\mathrm{Def\text{-}BM.12}(3)];
\]

%dsp-level:0
\[
\qquad \text{uniqueness up to an additive constant, fixed at}\ V(\infty)=0\ [\text{math}];
\]

%dsp-level:0
\[
\qquad \text{signs}\ [\text{machine}]
\]

%dsp-level:0
\[
(2):\ \text{the Coulomb ladder}\ [\text{math: standard spectrum on}\ [\hat x,\hat p]=i];
\]

%dsp-level:0
\[
\qquad e^4/(32\pi^2)=\varepsilon_{\mathrm{EM}}^2/2\ [\text{machine: coefficient identity}];
\]

%dsp-level:0
\[
\qquad E_2/E_1=\tfrac14\ [\text{machine}]
\]

%dsp-level:0
\[
(3):\ \mathrm{Def\text{-}BM.1}/\mathrm{D3.9}/\mathrm{CT3.17}\ \text{re-cited};
\]

%dsp-level:0
\[
\qquad \text{inventory scan: no EM sealing grammar, no pair bound state}\ [\text{machine: token row}]
\]

%dsp-level:0
\[
(4):\ \text{the category row}\ [(3)\ \text{conjoined with}\ \mathrm{BM.46}(3)];
\]

%dsp-level:0
\[
\qquad (5):\ \text{the sign rows re-cited};
\]

%dsp-level:0
\[
\qquad (6):\ \text{the status record}.\ \blacksquare
\]

\textbf{Phys.} With the sourced field in hand, the force between the decay's charged members writes itself: an exact symbolic coefficient over the separation, attractive precisely where the ternary charges oppose — and the chart's own quantum mechanics turns that potential into the hydrogenic ladder, whose ground coefficient and quarter-ratio are arithmetic. But a potential is not a particle pair. The corpus seals composites with colour, and only with colour: there is no electromagnetic junction grammar, no dynamics of the additive read, and therefore no atom in the inventory. The tier states the consequence as a bar rather than a shrug: binding energy belongs to the bound system, and an expression that subtracts a would-be atomic binding from a free proton's mass reads nothing that exists. The free proton's electromagnetic correction is its own field's energy, positive, exact, and alone.

\end{theorem}

\begin{theorem}[{\texorpdfstring{BM.49 — The beta-decay re-adjudication with the correction, and the endpoint-to-average refutation}{BM.49 - The beta-decay re-adjudication with the correction, and the endpoint-to-average refutation}}]\label{tag:BM-49}
\textbf{Thm.}

%dsp-level:0
\[
(1)\quad\textbf{the corrected Q-value:}
\]

%dsp-level:0
\[
\qquad Q_{\mathrm{corr}}:=m_{\mathrm{eff}}^{(n)}-m_{\mathrm{eff}}^{(p)}-m_\ell-m_\nu=-\bigl(\Delta m_{\mathrm{EM}}+m_\ell+m_\nu\bigr)
\]

%dsp-level:0
\[
\qquad [\mathrm{BM.46}(3)\text{--}(4);\ \text{machine: monomial identity}];
\]

%dsp-level:0
\[
\Delta m_{\mathrm{EM}}>0\ [\mathrm{BM.46}(4)],\quad m_\ell>0\ [\mathrm{BM.39}(1);\ \mathrm{BM.41}(4)(\mathrm a)],\quad m_\nu\ge0\ [\mathrm{EW.33}\ \mathrm{PSD}]:
\]

%dsp-level:0
\[
\boxed{\ Q_{\mathrm{corr}}<0\ \ \text{strictly, and strictly more negative than}\ Q=-(m_\ell+m_\nu)\ \text{of}\ \mathrm{BM.41}(4)\ }
\]

%dsp-level:0
\[
(2)\quad\textbf{the no-opening theorem:}
\]

%dsp-level:0
\[
\qquad \text{every EM correction of this tier is}\ \propto q^2
\]

%dsp-level:0
\[
\ge 0\ \text{and adds to the charged member's mass}\ [\mathrm{BM.46}(1),(4);\ \text{the field energy is the integral of a square}]:
\]

%dsp-level:0
\[
\text{the decay}\ n\to p+\ell+\nu\ \text{requires}\ m^{(n)}>m^{(p)}+m_\ell:
\]

%dsp-level:0
\[
\qquad \textbf{no EM correction of this tier moves the difference in the required direction}\ —\ \text{it moves it against};
\]

%dsp-level:0
\[
\text{the claim that the EM self-interaction makes the beta decay energetically possible is}\ \textbf{refuted by the exact sign arithmetic}\ [\text{machine}];
\]

%dsp-level:0
\[
\mathrm{BM.41}(4)(\mathrm e)\text{'s verdict is sharpened, not altered:}
\]

%dsp-level:0
\[
\qquad \text{opening the vertex remains a read-specification act}\ [\mathrm{BM.30}(5);\ \text{a §I addendum}]\ \text{or an inventory datum the corpus does not carry}\ (\text{a constituent-level transport splitting of the two baryon rows}\ —\ \neg\exists\ \text{in Parts 1--4}\ [\mathrm{BM.28}(1);\ \mathrm{BM.28}(7)]);
\]

%dsp-level:0
\[
\ \text{this tier performs neither};
\]

%dsp-level:0
\[
(3)\quad\textbf{the conditional spectrum (under the}\ \neg\textbf{adopted hypothesis}\ Q\ge0\ \textbf{of some specification):}
\]

%dsp-level:0
\[
\qquad E_{\max}=E_\ell^{\max}\ [\mathrm{BM.41}(5)\ \text{endpoint formula re-cited}];
\]

%dsp-level:0
\[
\text{the allowed-shape spectrum at constant amplitude with}\ m_\nu=0,\ m_\ell\to0:
\]

%dsp-level:0
\[
\qquad \frac{d\Gamma}{dE}\propto E^2\,(E_{\max}-E)^2
\]

%dsp-level:0
\[
\qquad [\text{math: the three-body phase space}];
\]

%dsp-level:0
\[
\frac{\langle E\rangle}{E_{\max}}
\]

%dsp-level:0
\[
= \frac{\int_0^1x^3(1-x)^2\,dx}{\int_0^1x^2(1-x)^2\,dx}
\]

%dsp-level:0
\[
= \frac{\mathrm B(4,3)}{\mathrm B(3,3)}
\]

%dsp-level:0
\[
= \frac{3!\,2!/6!}{2!\,2!/5!}
\]

%dsp-level:0
\[
= \frac{1/60}{1/30}\qquad = \frac12
\]

%dsp-level:0
\[
\qquad [\text{machine: exact factorial rationals}];
\]

%dsp-level:0
\[
\boxed{\ E_{\max}:\langle E\rangle=2:1\ }
\]

%dsp-level:0
\[
\qquad \textbf{and scale-invariant:}
\]

%dsp-level:0
\[
\quad E
\]

%dsp-level:0
\[
\mapsto QE\ \text{rescales both moments alike}\implies\text{the ratio is}\ 2\ \textbf{for every value of}\ Q>0\ [\text{machine: the ratio carries no}\ Q];
\]

%dsp-level:0
\[
(4)\quad\textbf{the}\ \tfrac43\textbf{-identification refuted:}
\]

%dsp-level:0
\[
\qquad \text{the proposal}\ E_{\max}:\langle E\rangle
\]

%dsp-level:0
\[
= \tfrac43\ \text{identified with}\ W(0)/W(+1)
\]

%dsp-level:0
\[
= \tfrac43\ [\mathrm{BM.28}(2)]\ \text{is barred three ways}:
\]

%dsp-level:0
\[
(\mathrm a)\ \textbf{arithmetic:}\quad2\ne\tfrac43\qquad[2\cdot3=6\ne4=4\cdot1;\ \text{machine: cross-multiplication in}\ \mathbb Z];
\]

%dsp-level:0
\[
(\mathrm b)\ \textbf{scale invariance:}
\]

%dsp-level:0
\[
\quad (3)\text{'s ratio is}\ 2\ \text{at every splitting scale}\ —\ \neg\exists\ \text{a value of}\ Q\ (\text{hence none of the route ratios})\ \text{that moves it to}\ \tfrac43\ \text{in the massless limit}\ [\text{machine}];
\]

%dsp-level:0
\[
(\mathrm c)\ \textbf{disjointness:}
\]

%dsp-level:0
\[
\quad W(0)/W(+1)\ \text{is a probability-space cardinal ratio}\ —\ \text{a measure functional whose defining clause carries zero momentum/metric tokens}\ [\mathrm{BM.37}(2)\ \text{machine census re-cited}];
\]

%dsp-level:0
\[
\qquad \langle E\rangle,\ E_{\max}\ \text{are chart-side spectral moments}:
\]

%dsp-level:0
\[
\text{the identification crosses}\ \mathrm{BM.37}(3)\text{'s bars}\ —\ \text{disjointness, amplitude}\ne\text{mass}\ [\mathrm{FT.79}(5)],
\]

%dsp-level:0
\[
\ \text{the read-choice row}\ [\mathrm{BM.30}(5)];
\]

%dsp-level:0
\[
(\mathrm d)\ \textbf{existence:}\quad\text{on every adopted read the decay is kinematically closed}\ [(1);\ \mathrm{BM.41}(4)(\mathrm e)]:
\]

%dsp-level:0
\[
\qquad \neg\exists\ \text{an emitted-electron spectrum whose moments could be driven};
\]

%dsp-level:0
\[
\qquad (3)\ \text{is conditional content and stays conditional};
\]

%dsp-level:0
\[
(5)\quad\textbf{status:}
\]

%dsp-level:0
\[
\qquad \text{the route rows of}\ \mathrm{BM.28}/\mathrm{BM.30}/\mathrm{BM.37}\ \text{stand untouched}\ —\ \text{the measure prices remain measure prices};
\]

%dsp-level:0
\[
\qquad \text{the derived conditional rows:}\ E_{\max}\ [\mathrm{BM.41}(5)],\ \langle E\rangle/E_{\max}=\tfrac12\ [\text{massless limit}];
\]

%dsp-level:0
\[
\neg\text{a measured endpoint, average energy, spectral shape or lifetime enters}\ [\mathrm{P5};\ \text{machine: zero external tokens}].
\]

\textbf{Rmk (the check that was asked for.)} The four-thirds was inspected where it could possibly live: as the ratio of the conditional spectrum's endpoint to its mean. The phase space answers with a clean rational — one half in the massless allowed limit, hence two-to-one, and scale-invariantly two-to-one at any $Q$ whatever, so no mass splitting of any provenance can bend it to four thirds. The four thirds remains what the Tier BM.8 adjudication made it: the ratio of twenty neutral sign-assignments to fifteen charged ones, a price of the charge filter on the ensemble measure. The two numbers are exact, and they are not about the same thing.

\textbf{Req.} P5 (Parts 1–2); D2.9 (Parts 1–2); EW.33 (Part 2); FT.79(5), FT.107(3) (Part 3); Def-BM.12, BM.28(1), BM.28(2), BM.28(7), BM.30(3), BM.30(5), BM.37(2), BM.37(3), BM.39(1), BM.41(4), BM.41(5), BM.43(2), BM.46(1), BM.46(3), BM.46(4) (this part); the three-body phase-space shape at constant amplitude, exact Beta-function rationals [math; machine: \texttt{checks/\allowbreak{}verify\_\allowbreak{}H\_\allowbreak{}em\_\allowbreak{}self\_\allowbreak{}interaction.py} §D].

\textbf{Pf.}

%dsp-level:0
\[
(1):\ m_{\mathrm{eff}}^{(n)}=m_{\mathrm{hid}},\ m_{\mathrm{eff}}^{(p)}=m_{\mathrm{hid}}+\Delta m_{\mathrm{EM}}\ [\mathrm{BM.46}(3)\text{--}(4)\ \text{re-cited}]:
\]

%dsp-level:0
\[
\quad Q_{\mathrm{corr}}=-\Delta m_{\mathrm{EM}}-m_\ell-m_\nu\ [\text{machine: monomial identity}];
\]

%dsp-level:0
\[
\qquad \text{the sign census}\ [\text{machine}]
\]

%dsp-level:0
\[
(2):\ \delta(q)=c\,\varepsilon_{\mathrm{EM}}q^2m_{\mathrm{hid}}\ge0\ [\text{machine: coefficient census}];
\]

%dsp-level:0
\[
\qquad m^{(p)}_{\mathrm{eff}}-m^{(n)}_{\mathrm{eff}}
\]

%dsp-level:0
\[
= \Delta m_{\mathrm{EM}}>0\ \text{against the required}\ m^{(n)}>m^{(p)}+m_\ell\ [\text{machine: the sign contradiction}]
\]

%dsp-level:0
\[
(3):\ \int_0^1x^3(1-x)^2dx
\]

%dsp-level:0
\[
= \mathrm B(4,3)
\]

%dsp-level:0
\[
= \tfrac1{60},
\]

%dsp-level:0
\[
\quad\int_0^1x^2(1-x)^2dx
\]

%dsp-level:0
\[
= \mathrm B(3,3)
\]

%dsp-level:0
\[
= \tfrac1{30}\ [\text{machine: factorial rationals}];
\]

%dsp-level:0
\[
\qquad \tfrac{1/60}{1/30}=\tfrac12\ [\text{machine}];
\]

%dsp-level:0
\[
\qquad E=Qx\ \text{rescaling: both moments carry}\ Q\ \text{to the same power}\ [\text{machine: the ratio is}\ Q\text{-free}]
\]

%dsp-level:0
\[
(4):\ 2\cdot3=6\ne4=4\cdot1\ [\text{machine}];
\]

%dsp-level:0
\[
\qquad (b)\text{--}(d):\ \text{the cited bars re-cited with}\ \mathrm{BM.37}(2)\text{'s census rows};
\]

%dsp-level:0
\[
\qquad (5):\ \text{the status record}.\ \blacksquare
\]

\textbf{Phys.} The correction goes into the balance and the balance closes harder. The neutron stays at six bits over $\mathfrak h$; the proton gains its field energy, strictly positive, symbolic in the coupling but positive at every value the coupling can take. The Q-value of the assembled vertex is minus the splitting minus the lepton masses — negative before the correction and more negative after. No electromagnetic row of this corpus opens the beta decay, because every row the tier derives adds mass to the charged member, and the charged member is the daughter. The spectrum question the tier was handed — could the four-thirds of the charge census be driving the endpoint-to-average ratio of the emitted electron — has an exact answer and it is no: the conditional massless spectrum averages at precisely half its endpoint, a scale-free rational that no splitting value can move, while four thirds is a count of sign-assignments in a filtered ensemble. Both numbers stand as derived; the bridge between them is barred by disjointness, by the amplitude-mass bar, and by the simple fact that on every adopted read there is no emitted electron to average.

\end{theorem}

\medskip\noindent\textit{\texorpdfstring{BM.50 — The status of the EM self-interaction correction: adjudication and scope}{BM.50 - The status of the EM self-interaction correction: adjudication and scope}.}\ Structural-constraints record: see Appendix B, Section B.1.21 of the Addendum.\par\medskip

\phantomsection
\section[Tier BM.13 — The Full-Scope Mass Splitting: the Lyapunov…]{\texorpdfstring{Tier BM.13 — The Full-Scope Mass Splitting: the Lyapunov Foundation, the Complete-Force Census, the Non-Uniform Clock, the Binding Zeros, and the Exact Ratio}{Tier BM.13 - The Full-Scope Mass Splitting: the Lyapunov Foundation, the Complete-Force Census, the Non-Uniform Clock, the Binding Zeros, and the Exact Ratio}}

\begin{theorem}[{\texorpdfstring{BM.51 — The Lyapunov foundation: the contraction law iterated, the deficit links, and the class-degenerate zero}{BM.51 - The Lyapunov foundation: the contraction law iterated, the deficit links, and the class-degenerate zero}}]\label{tag:BM-51}
\textbf{Thm.}

%dsp-level:0
\[
(1)\quad\textbf{the functional and its zero locus:}
\]

%dsp-level:0
\[
\qquad \delta\rho:=\rho-\tfrac14I,
\]

%dsp-level:0
\[
\qquad \tilde K:=-\log\rho-\mathfrak T\bigl((-\log\rho)\,\rho\bigr)I,
\]

%dsp-level:0
\[
\qquad V(\tilde K):=\mathfrak T\bigl(\tilde K^2\bigr)\ge0
\]

%dsp-level:0
\[
\qquad [\mathrm{L1}];
\]

%dsp-level:0
\[
V=0\iff\rho=\tfrac14I\ [\mathrm{L1}(1)];
\]

%dsp-level:0
\[
\qquad V\bigl(\tilde K(\mathcal T(\rho))\bigr)=V\bigl(\tilde K(\rho)\bigr)\ \ \text{for every draw}\ [\mathrm{L1}(2);\ \mathrm{A6}:\ \text{tick invariance}];
\]

%dsp-level:0
\[
\qquad V'\le V,\ \ V'=V\iff\rho=\tfrac14I\ [\mathrm{L1}(5)];
\]

%dsp-level:0
\[
(2)\quad\textbf{the contraction law:}
\]

%dsp-level:0
\[
\qquad \gamma:=1-4\varepsilon_q\in(0,1]\ [\mathrm{D0.3};\ \mathrm{D0.4};\ \mathrm{CT0.8}];
\]

%dsp-level:0
\[
\qquad V'=\Bigl(1-\frac\gamma4\Bigr)^{2}V+O(\delta\rho^3),
\]

%dsp-level:0
\[
\qquad V'-V=-\frac{\gamma(8-\gamma)}{16}\,V+O(\delta\rho^3)
\]

%dsp-level:0
\[
\qquad [\mathrm{L2}];
\]

%dsp-level:0
\[
\boxed{\ \Bigl(1-\frac\gamma4\Bigr)^{2}+\frac{\gamma(8-\gamma)}{16}=1\ }
\]

%dsp-level:0
\[
\qquad [\text{machine: exact in}\ \mathbb Q[\gamma]]
\]

%dsp-level:0
\[
\qquad \text{the contraction is strict on the whole admissible range}\ \gamma\in(0,1]\ [\mathrm{BM.8}(2);\ \text{machine}];
\]

%dsp-level:0
\[
\gamma=1\ (\text{the sharp point}\ \varepsilon_q=0\ [\mathrm{D0.4};\ \mathrm{BM.19}(3)]):\quad V'=\frac9{16}V,\quad\text{deficit}\ \frac7{16}V;
\]

%dsp-level:0
\[
\qquad \gamma\to0:\ V'=V\ [\text{the identity-tilt limit, the excluded endpoint}\ \varepsilon_q=\tfrac14\ [\mathrm{CT0.8}]];
\]

%dsp-level:0
\[
(3)\quad\textbf{the iteration (the fermion-count contraction read):}
\]

%dsp-level:0
\[
\qquad \text{one tilt step per recorded crystallization event}\ [\mathrm{L0}\ \text{grounding};\ \mathrm{L2}\ \mathrm{Phys}]:
\]

%dsp-level:0
\[
\qquad V_n=\Bigl(1-\frac\gamma4\Bigr)^{2n}V_0\ \ \text{at the leading-coefficient level, remainder}\ O(\delta\rho^3)\ \text{per step}\ [\mathrm{L2}\ \text{scope}];
\]

%dsp-level:0
\[
\textbf{the}\ E_{\mathrm{int}}\textbf{-fold row:}
\]

%dsp-level:0
\[
\qquad V_{E_{\mathrm{int}}}=\Bigl(1-\frac\gamma4\Bigr)^{2E_{\mathrm{int}}}V_0
\]

%dsp-level:0
\[
\qquad [E_{\mathrm{int}}=\#\{\text{internal directed edges}\}\ [\mathrm{D0.8}];\ \text{machine: the power identity}\ \bigl(1-\tfrac\gamma4\bigr)^2\bigl(1-\tfrac\gamma4\bigr)^{2n}=\bigl(1-\tfrac\gamma4\bigr)^{2n+2}\ \text{in}\ \mathbb Q[\gamma]];
\]

%dsp-level:0
\[
\textbf{at the sharp point}\ \gamma=1,\ E_{\mathrm{int}}=6\ [\mathrm{BM.17};\ \mathrm{BM.26};\ \text{Census FT.35}(2)]:
\]

%dsp-level:0
\[
\qquad \boxed{\ \frac{V_6}{V_0}=\Bigl(\frac9{16}\Bigr)^{6}=\frac{531441}{16777216}\ },
\]

%dsp-level:0
\[
\qquad \text{the total deficit}\ 1-\frac{531441}{16777216}=\frac{16245775}{16777216}
\]

%dsp-level:0
\[
\qquad [\text{machine: exact integer powers}];
\]

%dsp-level:0
\[
\textbf{faithfulness in}\ E_{\mathrm{int}}\textbf{:}\quad\text{at fixed}\ \gamma
\]

%dsp-level:0
\[
\in (0,1],
\]

%dsp-level:0
\[
\ n\mapsto\bigl(1-\tfrac\gamma4\bigr)^{2n}\ \text{is strictly decreasing and injective}\ [\text{machine: exact rational comparison on}\ n\in2..60]:
\]

%dsp-level:0
\[
\qquad \text{the contraction read separates edge counts exactly as}\ m_{\mathrm{hid}}\ \text{does};
\]

%dsp-level:0
\[
(4)\quad\textbf{the deficit links:}
\]

%dsp-level:0
\[
\qquad \Delta S_{\mathrm{pair}}:=2|\delta\rho|^2=\tfrac12|r|^2,
\]

%dsp-level:0
\[
\qquad S\bigl(\rho\,\|\,\tfrac14I\bigr)=\Delta S_{\mathrm{pair}}+O(\delta\rho^3)
\]

%dsp-level:0
\[
\qquad [\mathrm{L3}(1)];
\]

%dsp-level:0
\[
\tilde K=-4\,\delta\rho+O(\delta\rho^2),
\]

%dsp-level:0
\[
\qquad \boxed{\ V(\tilde K)=16\,|\delta\rho|^2+O(\delta\rho^3)=8\,\Delta S_{\mathrm{pair}}+O(\delta\rho^3)\ }
\]

%dsp-level:0
\[
\qquad [\mathrm{L3}(2);\ \text{machine: both branch witnesses}];
\]

%dsp-level:0
\[
V'-V=-\frac{\gamma(8-\gamma)}2\,\Delta S_{\mathrm{pair}}+O(\delta\rho^3),\qquad\gamma=1:\ -\frac72\Delta S_{\mathrm{pair}}\qquad[\mathrm{L3}(3)];
\]

%dsp-level:0
\[
\textbf{the branch witnesses:}\quad\delta\rho_{\mathcal V}
\]

%dsp-level:0
\[
= \Pi_{\mathcal V}-\tfrac14I:\ \operatorname{spec}
\]

%dsp-level:0
\[
= \bigl(\tfrac34,-\tfrac14,-\tfrac14,-\tfrac14\bigr):\ \ \mathfrak T(\delta\rho^2)\qquad = \tfrac34;
\]

%dsp-level:0
\[
\qquad \delta\rho_{\mathcal R}=\tfrac{\Pi_{\mathcal R}}3-\tfrac14I:\ \operatorname{spec}=\bigl(\tfrac1{12},\tfrac1{12},\tfrac1{12},-\tfrac14\bigr):\ \ \mathfrak T(\delta\rho^2)=\tfrac1{12}
\]

%dsp-level:0
\[
[\mathrm{FT.107}(2);\ \text{machine: exact Fraction matrix arithmetic}]:
\]

%dsp-level:0
\[
\qquad \text{the leading deficits}\ 2|\delta\rho|^2=\bigl(\tfrac32,\ \tfrac16\bigr),
\]

%dsp-level:0
\[
\qquad V=8\Delta S_{\mathrm{pair}}:\ \ 16\cdot\tfrac34=12=8\cdot\tfrac32,\quad16\cdot\tfrac1{12}=\tfrac43=8\cdot\tfrac16\ [\text{machine}];
\]

%dsp-level:0
\[
(5)\quad\textbf{the inertial read:}
\]

%dsp-level:0
\[
\qquad m_i^2(q):=\Delta S(q)
\]

%dsp-level:0
\[
\qquad [\mathrm{L3}(4);\ \mathrm{L4}:\ \text{the patch deficit with its integer degree factor};\ \text{dimensionless nats};\ \neg\text{a particle identification}];
\]

%dsp-level:0
\[
\textbf{the species row:}\quad\text{the conversion between the deficit/inertial species and the}\ m_{\mathrm{hid}}\ \text{species is item-7 content}:
\]

%dsp-level:0
\[
\qquad \text{the tier needs no conversion}\ —\ (6)\text{'s zero is species-independent};
\]

%dsp-level:0
\[
(6)\quad\textbf{the class-degenerate zero:}
\]

%dsp-level:0
\[
\qquad \text{every input of the contraction read is class-shared:}
\]

%dsp-level:0
\[
\quad E_{\mathrm{int}}=6\ \text{on both baryon rows}\ [\mathrm{BM.17};\ \mathrm{BM.26};\ \text{Census FT.35}(2)];
\]

%dsp-level:0
\[
\qquad \gamma=1-4\varepsilon_q\ \text{is the universal ledger threshold}\ [\mathrm{D0.4}];
\]

%dsp-level:0
\[
\qquad V_0\ \text{is the common vacuum reference}\ [\mathrm{L1}(1);\ \mathrm{CT0.2}]:
\]

%dsp-level:0
\[
\boxed{\ \Delta m_{\mathrm{Lyapunov}}=0\ }
\]

%dsp-level:0
\[
\qquad [\text{machine: the difference of identical monomials}]:
\]

%dsp-level:0
\[
\qquad \textbf{the Lyapunov foundation gives the two baryon rows the same mass}\ —\ \text{in every read species at once, because it splits no input};
\]

%dsp-level:0
\[
(7)\quad\textbf{scope:}
\]

%dsp-level:0
\[
\qquad \mathrm{L1}\text{'s faithfulness scope stands}\ (\rho\ \text{faithful};\ \gamma\in(0,1];\ \text{the identity-tilt limit is the excluded endpoint});
\]

%dsp-level:0
\[
\qquad \mathrm{L0}\text{'s crystallization--tilt coupling is the declared theorem-local dynamical scope}\ [\text{inherited at the point of use},\ \neg\text{extended}];
\]

%dsp-level:0
\[
\text{the identification ``mass}:
\]

%dsp-level:0
\[
= \text{the contraction rate'' is}\ \neg\text{adopted as a new read species:}\ (3)\text{'s}\ V_{E_{\mathrm{int}}}/V_0\ \text{is an}\ E_{\mathrm{int}}\text{-faithful dimensionless companion of}\ m_{\mathrm{hid}}\ [\mathrm{FT.44}(4)],
\]

%dsp-level:0
\[
\ \text{and}\ (6)\text{'s zero holds on it}.
\]

\textbf{Rmk (why the zero needs no bridge).} The relating equation between the deficit species (the Lyapunov-contraction read of the registry [Def-H.39(3)]) and the informational species is $\neg$derived, and this tier does not touch it — because a splitting claim would need the bridge while a degeneracy claim does not: two rows sharing every input of a functional share its value in any species, converted or not. The Lyapunov tower is the foundation the directive asks for, and what the foundation says about the neutron–proton pair is exact, trivial, and load-bearing: zero.

\textbf{Req.} A1, A4, A5, A6 (Part 1); C1, D0.3, D0.4, D0.8 (Part 1); CT0.2, CT0.8, CT1.1 (Part 1); P5 (Parts 1–2); L0, L1, L2, L3, L4 (Part 1); FT.44(4) (Census); FT.107(2) (Part 3); Def-BM.5(4), BM.8(2), BM.17, BM.19(3), BM.26, BM.28(1) (this part); exact formal-power-series identities over $\mathbb Q[\gamma]$, exact rational powers and integer power censuses, exact Fraction matrix arithmetic on the branch witnesses [math; machine: \texttt{checks/\allowbreak{}verify\_\allowbreak{}H\_\allowbreak{}full\_\allowbreak{}scope\_\allowbreak{}mass\_\allowbreak{}splitting.py} §A].

\textbf{Pf.}

%dsp-level:0
\[
(1)\text{--}(2):\ \mathrm{L1}(1),(2),(5)\ \text{and}\ \mathrm{L2}\ \text{re-cited (frozen)};
\]

%dsp-level:0
\[
\qquad \bigl(1-\tfrac\gamma4\bigr)^2+\tfrac{\gamma(8-\gamma)}{16}=1-\tfrac\gamma2+\tfrac{\gamma^2}{16}+\tfrac\gamma2-\tfrac{\gamma^2}{16}=1\ [\text{machine}];
\]

%dsp-level:0
\[
\qquad \gamma\in(0,1]\implies\bigl(1-\tfrac\gamma4\bigr)^2\in\bigl[\tfrac9{16},1\bigr)\ [\text{machine: rational witnesses}]
\]

%dsp-level:0
\[
(3):\ \text{induction on the leading factor}\ [\mathrm{L2}]:\ \text{the power identity in}\ \mathbb Q[\gamma]\ [\text{machine}];
\]

%dsp-level:0
\[
\qquad (9/16)^6=531441/16777216\ [\text{machine:}\ 9^6=531441,\ 16^6=16777216];
\]

%dsp-level:0
\[
\qquad \text{injectivity:}\ \bigl(\tfrac9{16}\bigr)^{n}>\bigl(\tfrac9{16}\bigr)^{n+1}\ \text{on}\ n\in2..60\ [\text{machine: exact Fractions}]
\]

%dsp-level:0
\[
(4):\ \Pi_{\mathcal V}=\tfrac14(E_0-\Sigma),\ \Pi_{\mathcal R}=\tfrac14(3E_0+\Sigma)\ [\mathrm{CT0.7}(1)];
\]

%dsp-level:0
\[
\qquad \Sigma
\]

%dsp-level:0
\[
= \textstyle\sum_i\sigma_i\otimes\sigma_i:\ \operatorname{spec}
\]

%dsp-level:0
\[
= \{1,1,1,-3\}\ [\text{machine}]:\quad\delta\rho_{\mathcal V}:\ \bigl(\tfrac34,-\tfrac14,-\tfrac14,-\tfrac14\bigr),
\]

%dsp-level:0
\[
\ \ \delta\rho_{\mathcal R}:\ \bigl(\tfrac1{12},\tfrac1{12},\tfrac1{12},-\tfrac14\bigr)\ [\text{machine}]
\]

%dsp-level:0
\[
\mathfrak T(\delta\rho_{\mathcal V}^2)=\tfrac9{16}+\tfrac3{16}=\tfrac34;
\]

%dsp-level:0
\[
\qquad \mathfrak T(\delta\rho_{\mathcal R}^2)=\tfrac3{144}+\tfrac1{16}=\tfrac1{12}\ [\text{machine: exact Fractions}];
\]

%dsp-level:0
\[
\qquad V=16|\delta\rho|^2=8\Delta S_{\mathrm{pair}}\ \text{at both witnesses}\ [\text{machine}]
\]

%dsp-level:0
\[
(5):\ \mathrm{L3}(4)/\mathrm{L4}\ \text{re-cited};
\]

%dsp-level:0
\[
\qquad (6):\ m_n^{(\mathrm{Lyap})}-m_p^{(\mathrm{Lyap})}=0\ \text{as the difference of identical inputs}\ [\text{machine}];
\]

%dsp-level:0
\[
\qquad (7):\ \text{the scope record}.\ \blacksquare
\]

\textbf{Phys.} The foundation is a contraction, and the contraction is exact arithmetic. Every tilt step keeps $(1-\gamma/4)^2$ of the Lyapunov value and loses $\gamma(8-\gamma)/16$ of it — the two add to one as an identity in $\gamma$, not as an approximation — and iterating over the crystallization events of a composite turns the edge count into a strictly faithful dimensionless read: at the sharp threshold, six edges contract the functional to $531441/16777216$ of its vacuum value. The deficit links tie the same object to entropy and inertia with exact rational factors, verified at both branch witnesses: sixteen of a squared deviation, eight of a pair deficit, one inertial read per channel. And then the payload for the baryon pair: every input of this tower — the edge count, the threshold, the vacuum reference — is shared by the neutron row and the proton row, so the whole Lyapunov contribution to their splitting is zero, in every read species at once. The foundation does not split the pair; the split has to come from somewhere else, and the census that follows finds exactly one somewhere.

\end{theorem}

\begin{theorem}[{\texorpdfstring{BM.52 — The complete force operator and the sector-by-sector splitting census}{BM.52 - The complete force operator and the sector-by-sector splitting census}}]\label{tag:BM-52}
\textbf{Thm.}

%dsp-level:0
\[
(1)\quad\textbf{the complete force operator:}
\]

%dsp-level:0
\[
\qquad \mathbf F
\]

%dsp-level:0
\[
= F_{(+,\mathcal R)}\oplus F_{(-,\mathcal R)}\oplus F_{(+,\mathcal V)}\oplus F_{(-,\mathcal V)}\oplus F_{\mathrm{st}}\ \ \text{always applies in full}\ —\ \text{no sub-block restriction is an A8 force}\ [\mathrm{A8};\ \mathrm{FT.142}(2)];
\]

%dsp-level:0
\[
\text{the block/aspect dictionary:}
\]

%dsp-level:0
\[
\quad (G,H,D^+)\leftrightarrow\mathcal E_{++}\ (7),
\]

%dsp-level:0
\[
\quad W\leftrightarrow\mathcal E_{+-}\ (3),
\]

%dsp-level:0
\[
\quad J^+\leftrightarrow\mathcal E_{-+}\ (3),
\]

%dsp-level:0
\[
\quad J^-\leftrightarrow\mathcal E_{--}\ (3):
\]

%dsp-level:0
\[
\qquad 7+3+3+3=16=C_{\mathrm{cap}}\ [\mathrm{CT1.1};\ \mathrm{CT0.12};\ \mathrm{CT1.12a}(2);\ \mathrm{FT.30}];
\]

%dsp-level:0
\[
\text{at the sealed composites the rung-3 mask is active:}\quad\dim\mathcal K_3
\]

%dsp-level:0
\[
= 25
\]

%dsp-level:0
\[
= 16+1+8\ [\mathrm{CT0.13};\ \mathrm{CT0.14};\ \mathrm{Def\text{-}BM.3}];
\]

%dsp-level:0
\[
\qquad \text{the force enters the stress tensor through the masked resolvent}\ Z_{\mathrm{force}}
\]

%dsp-level:0
\[
= \langle\mathrm{out}|(1-\mathbf F_M^{(n)})^{-1}|\mathrm{in}\rangle\ [\mathrm{CT1.22};\ \mathrm{CT2.2}];
\]

%dsp-level:0
\[
(2)\quad\textbf{the carrier-side static matrix:}
\]

%dsp-level:0
\[
\qquad \mathbf F_{\mathrm{ext}}(\mathbf c,q)=h(\mathbf c)\otimes I_2+q\,I_2\otimes\hat\kappa,
\]

%dsp-level:0
\[
\qquad F_{LL}=F_{RR}=h(\mathbf c),
\]

%dsp-level:0
\[
\qquad F_{LR}=F_{RL}=q\,I_2
\]

%dsp-level:0
\[
\qquad [\mathrm{BM.22}(1)\text{--}(2);\ \mathrm{FT.129}(4)\text{--}(5);\ \mathrm{Def\text{-}H.43}\ \text{promotion};\ \text{machine}];
\]

%dsp-level:0
\[
\textbf{the chiral-partner splitting:}\quad2|q|\in\{0,2\}\ \ \text{at}\ q\in\{-1,0,+1\}
\]

%dsp-level:0
\[
\qquad [\mathrm{BM.22}(4);\ \mathrm{FT.130}(5);\ \text{machine}]:
\]

%dsp-level:0
\[
\qquad \text{the neutral class keeps parity doubling}\ (F_{LR}=0),\ \text{the unit class splits by two read units};
\]

%dsp-level:0
\[
\textbf{the EM difference between the two classes is exactly the flipping block}\ [\mathrm{BM.22}(4)];
\]

%dsp-level:0
\[
\qquad \text{the splitting is a dimensionless force-spectrum read}\ —\ \text{its frequency conversion is the deferred read}\ [\mathrm{FT.131}(3);\ \mathrm{FT.159}\ \text{residue}];
\]

%dsp-level:0
\[
(3)\quad\textbf{the sector census (the splitting contributions between the two baryon rows):}\qquad\begin{array}{l|l|l}\text{sector} & \text{contribution} & \text{the derived ground}\\\hline J^\pm\ (\text{seals}) & 0 & \Delta S=\tfrac32(2E-|R|)\ \text{is}\ \sigma\text{-free}\ [\mathrm{FT.107}(2)];\ E_{\mathrm{int}}\ \text{equal};\ \text{the Janus balance holds on both rows}\ [\mathrm{BM.7}]\\G,H\ (\text{metric blocks}) & 0 & \text{the seven geometric monomials carry charge exponent}\ 0\ [\mathrm{BM.28}(1)]\\W\ (\text{weak/shift}) & 0 & |w|^2\ \text{identical between the classes}\ [\mathrm{BM.28}(6)];\ \text{the Casimir is branch data, neutral at the}\ \sigma\text{-vertex}\ [\mathrm{BM.40}(3);\ \mathrm{FT.157}(1)]\\\text{colour} & 0 & \text{sealed ends are blind}\ [\mathrm{FT.117}];\ \text{the mask admits only singlets}\ [\mathrm{CT3.18};\ \mathrm{D2.14}];\ \text{both rows are}\ \varepsilon\text{-sealed with}\ g(3)=1\ [\mathrm{BM.31}(2)];\ \text{colour acts on the slot carrier}\ [\mathrm{FT.127}(7)];\ \text{the seed is colour-neutral}\ [\mathrm{EW.32}]\\D^\pm\ (\text{dark}) & 0 & D^+\ \text{is uniform and carried by}\ \Lambda\ [\mathrm{CT1.22};\ \mathrm{CT4.1}];\ \text{both locks are shared by all charge classes}\ [\mathrm{BM.28}(1);\ \mathrm{CT4.5}]\\\mathrm{EM} & \textbf{the only nonzero sector} & \text{the flipping block}\ 2|q|\in\{0,2\}\ [(2)];\ \text{the field energy}\ \Delta m_{\mathrm{EM}}=\tfrac{3ce^2\ln2}{2\pi\mathfrak h}>0\ [\mathrm{BM.46}(4)]\end{array}
\]

%dsp-level:0
\[
(4)\quad\textbf{the totality theorem:}
\]

%dsp-level:0
\[
\qquad \boxed{\ \Delta m_{\mathrm{force}}:=m_n-m_p\Big|_{\text{the complete operator}}=-\,\Delta m_{\mathrm{EM}}=-\frac{3\,c\,e^2\ln2}{2\pi\,\mathfrak h}\ }
\]

%dsp-level:0
\[
\qquad [\text{machine: the census sum of five zeros and one monomial}];
\]

%dsp-level:0
\[
\textbf{the mass splitting comes from the total force}\ —\ \text{and the total force, censused sector by sector, reduces to its EM summand}:
\]

%dsp-level:0
\[
(5)\quad\textbf{the agreement row (no double count):}
\]

%dsp-level:0
\[
\qquad \text{the sector census}\ ((3))\ \text{and the field integral}\ [\mathrm{BM.46}(1)]\ \text{are two derived routes to one monomial}:
\]

%dsp-level:0
\[
\Delta m_{\mathrm{force}}\big|_{\mathrm{EM}}=\Delta m_{\mathrm{self\text{-}int}}=\Delta m_{\mathrm{EM}}
\]

%dsp-level:0
\[
\qquad [\text{machine: the monomials are identical}]:
\]

%dsp-level:0
\[
\qquad \text{in any decomposition the EM contribution is counted once}\ [\mathrm{BM.55}(1)];
\]

%dsp-level:0
\[
(6)\quad\textbf{scope:}
\]

%dsp-level:0
\[
\qquad \text{the census is between the two baryon rows at matched profiles and shared junction grammar}\ [\mathrm{Def\text{-}BM.3};\ \mathrm{BM.31}(2)];
\]

%dsp-level:0
\[
\qquad \text{the}\ \mathrm{Def\text{-}H.43}\ \text{promotion's conditionality stands as}\ \mathrm{BM.22}(5)\ \text{records it}\ (\neg\text{a new declaration here});
\]

%dsp-level:0
\[
\text{constituent-level transport differences between the rows have no derived object}\ [\mathrm{BM.55}(4)\ \text{forward pointer at the citing site}];
\]

%dsp-level:0
\[
\qquad \text{the deferred frequency read}\ [\mathrm{FT.131}(3)]\ \text{blocks no energy-dimension row: the field monomial}\ (4)\ \text{is complete without it}.
\]

\textbf{Rmk (what "full scope" computes to).} The standing order is the sum of all sectors, and the sum is taken — literally, sector by sector, with each zero backed by its own frozen theorem rather than by silence. The seals cannot split the pair because the deficit never reads the end signs; the metric blocks cannot because every geometric monomial is charge-blind; the weak block cannot because the two rows share $|w|^2$ and the Casimir is branch-neutral at a sign flip; colour cannot because sealed singlets are blind to it and both rows carry the same unique grouping; the dark sector cannot because its locks are class-shared. What survives is electromagnetism, twice-derived and identical: the flipping block on the carrier spectrum and the field energy on the chart. Full scope does not dilute the answer; it isolates it.

\textbf{Req.} A8 (Part 1); CT0.12, CT0.13, CT0.14, CT1.1 (Part 1); CT1.22, CT2.2, CT3.18, CT4.1, CT4.5, D2.14 (Part 2); EM.9, EW.32 (Part 2); P5 (Parts 1–2); Def-H.43, FT.30, FT.107(2), FT.117, FT.125(4), FT.127(7), FT.129(4), FT.129(5), FT.130(5), FT.131(3), FT.157(1), FT.159 (Part 3); Def-BM.3, Def-BM.12, BM.7, BM.22(1), BM.22(2), BM.22(4), BM.22(5), BM.28(1), BM.28(6), BM.31(2), BM.40(3), BM.46(1), BM.46(4) (this part); exact integer censuses, Laurent monomials [math; machine: \texttt{checks/\allowbreak{}verify\_\allowbreak{}H\_\allowbreak{}full\_\allowbreak{}scope\_\allowbreak{}mass\_\allowbreak{}splitting.py} §B].

\textbf{Pf.}

%dsp-level:0
\[
(1):\ \mathrm{A8}\text{'s completeness row and the block dictionary re-cited}\ [\mathrm{CT1.1};\ \mathrm{CT0.12};\ \mathrm{CT1.12a}(2);\ \mathrm{FT.30}];
\]

%dsp-level:0
\[
\qquad 7+3+3+3=16,\ 16+1+8=25\ [\text{machine}]
\]

%dsp-level:0
\[
(2):\ \mathrm{BM.22}(1)\text{--}(2),(4)\ \text{re-cited};
\]

%dsp-level:0
\[
\qquad 2|q|\ \text{at}\ q\in\{-1,0,+1\}=(2,0,2)\ [\text{machine}];
\]

%dsp-level:0
\[
\qquad F_{LR}=qI_2:\ \text{the zero matrix at}\ q=0\ [\text{machine: block extraction}]
\]

%dsp-level:0
\[
(3):\ \text{each zero is the cited frozen row}: \mathrm{FT.107}(2)\text{'s}\ \sigma\text{-freeness};
\]

%dsp-level:0
\[
\ \mathrm{BM.28}(1)\text{'s monomial census};
\]

%dsp-level:0
\[
\ \mathrm{BM.28}(6)\text{'s}\ |w|^2\ \text{row and}\ \mathrm{BM.40}(3);
\]

%dsp-level:0
\[
\ \mathrm{FT.117}/\mathrm{CT3.18}/\mathrm{BM.31}(2)/\mathrm{FT.127}(7)/\mathrm{EW.32};
\]

%dsp-level:0
\[
\ \mathrm{CT1.22}/\mathrm{CT4.1}/\mathrm{CT4.5}/\mathrm{BM.28}(1)\ \text{lock row}\ [\text{machine: the census structure}]
\]

%dsp-level:0
\[
(4):\ 0+0+0+0+0-\Delta m_{\mathrm{EM}}\ [\text{machine: the monomial sum}];
\]

%dsp-level:0
\[
\qquad (5):\ \mathrm{BM.46}(1)\text{'s integral and}\ (3)\text{'s EM entry are the same Laurent monomial}\ [\text{machine}];
\]

%dsp-level:0
\[
\qquad (6):\ \text{the scope record}.\ \blacksquare
\]

\textbf{Phys.} Ask every force in the ledger what it contributes to the neutron–proton split, and five of the six answer zero — not approximately, not by assumption, but each by a theorem already on the books: the deficit ignores end signs, the geometry ignores charge, the weak Casimir ignores sign flips, colour cannot see sealed singlets, and the dark locks are shared. The complete force operator is the sum of all sectors, so the total is the EM summand, and the EM summand has already been computed twice by two routes that agree monomial for monomial. The splitting of the corpus's two baryon rows is one term: the field energy of the scar. Everything else in the full scope is bookkeeping that confirms the isolation.

\end{theorem}

\begin{theorem}[{\texorpdfstring{BM.53 — The metric background and the non-uniform clock: the zero metric contribution and the Mandelstam–Tamm resolvability bound}{BM.53 - The metric background and the non-uniform clock: the zero metric contribution and the Mandelstam-Tamm resolvability bound}}]\label{tag:BM-53}
\textbf{Thm.}

%dsp-level:0
\[
(1)\quad\textbf{the background:}
\]

%dsp-level:0
\[
\qquad \text{the metric is the emergent dictionary of the sector coordinates}\ (a_{00},q_{ij},w_i)\ [\mathrm{CT1.12};\ \mathrm{Def\text{-}BM.10}(2)];
\]

%dsp-level:0
\[
\qquad \text{the dispersion read}\ E=N\sqrt{m^2+\gamma^{ij}p_ip_j}-N^ip_i\ \text{and}\ \mathrm{KE}=N\bigl(\sqrt{m^2+p^2}-m\bigr)\ [\mathrm{BM.34}(4)\text{--}(5)];
\]

%dsp-level:0
\[
(2)\quad\textbf{the metric contribution is zero:}
\]

%dsp-level:0
\[
\qquad \text{the seven geometric monomials}\ (\Delta x_{\mathrm B},V_{\mathrm{eff}},\rho,\kappa^2\rho,r_s,\ \text{both locks})\ \text{carry charge exponent}\ 0\ [\mathrm{BM.28}(1);\ \text{machine: the exponent census re-run}];
\]

%dsp-level:0
\[
\text{the rest slice}\ E(p{=}0)
\]

%dsp-level:0
\[
= N\,
\]

%dsp-level:0
\[
m\ [\mathrm{BM.34}(7)]\ \text{has class-shared}\ N\ \text{and class-shared}\ m\ \text{before the EM correction}\ [\mathrm{BM.17};\ \mathrm{BM.26};\ \mathrm{BM.51}(6)]:
\]

%dsp-level:0
\[
\qquad \boxed{\ \Delta m_{\mathrm{metric}}=0\ }\ [\text{machine}];
\]

%dsp-level:0
\[
(3)\quad\textbf{the non-uniform clock:}
\]

%dsp-level:0
\[
\qquad \text{the tick strictly contracts}\ V\ \text{on the whole admissible range}\ [\mathrm{BM.8}(2);\ \mathrm{L2};\ \mathrm{BM.51}(2)\text{--}(3)]:
\]

%dsp-level:0
\[
\qquad V_n=V_0\bigl(1-\tfrac\gamma4\bigr)^{2n};
\]

%dsp-level:0
\[
\textbf{the log-linear row:}
\]

%dsp-level:0
\[
\qquad \ln V_n=\ln V_0+2n\,\ln\!\bigl(1-\tfrac\gamma4\bigr)
\]

%dsp-level:0
\[
\qquad [\text{math: the logarithm of the exact power identity};\ \text{structural}\ —\ \text{no transcendental is evaluated}];
\]

%dsp-level:0
\[
\textbf{no uniform-}\tau\ \textbf{exponential law:}
\]

%dsp-level:0
\[
\qquad \text{the ledger evolution is the per-tick weight and its resolvent}\ [\mathrm{BM.44}(1)\text{--}(4);\ \mathrm{D1.8}:\ x^0:=\tau_n;\ \mathrm{FT.159}]:
\]

%dsp-level:0
\[
\qquad \Gamma_{\mathrm{tr}}\in\bigl[0,\tfrac13\bigr],\quad\tau_{\mathrm{life}}\ge3\ \text{ticks},\quad\neg\text{external seconds};
\]

%dsp-level:0
\[
(4)\quad\textbf{the Mandelstam--Tamm relation:}
\]

%dsp-level:0
\[
\qquad \Delta t\,\Delta E\ge\tfrac12
\]

%dsp-level:0
\[
\qquad [\mathrm{BM.10}(4);\ \text{no time operator at finite dimension}];
\]

%dsp-level:0
\[
\textbf{the exact saturation witness re-derived:}\quad A=\tfrac12\sigma_x,\ H=\tfrac12\sigma_z,\ \psi\propto(1,i):
\]

%dsp-level:0
\[
\qquad \langle A\rangle=\langle H\rangle=0,\quad\langle\sigma_y\rangle=1,\quad(\Delta A)^2=(\Delta E)^2=\tfrac14,\quad[H,A]=\tfrac i2\sigma_y,\quad\frac{d\langle A\rangle}{dt}=-\tfrac12:
\]

%dsp-level:0
\[
\Delta t=\frac{\Delta A}{\bigl|\tfrac{d\langle A\rangle}{dt}\bigr|}=1,\qquad\Delta t\,\Delta E=\tfrac12\qquad[\text{machine: exact over}\ \mathbb Q(i)];
\]

%dsp-level:0
\[
(5)\quad\textbf{the resolvability of the splitting:}
\]

%dsp-level:0
\[
\qquad \text{resolving an energy difference}\ \Delta m_{\mathrm{EM}}\ \text{against the non-uniform clock requires}\ \ \Delta t\ \ge\ \frac1{2\,\Delta m_{\mathrm{EM}}}\ \ [\mathrm{MT},\ (4)]:
\]

%dsp-level:0
\[
\boxed{\ \Delta t\ \ge\ \frac{\pi\,\mathfrak h}{3\,c\,e^2\ln2}\ }\qquad[\text{machine: the exact reciprocal monomial of}\ \mathrm{BM.46}(4)\text{'s splitting}];
\]

%dsp-level:0
\[
(6)\quad\textbf{status:}
\]

%dsp-level:0
\[
\qquad \text{the non-uniform clock shifts no splitting value}\ —\ \text{both rows ride the same}\ \gamma,
\]

%dsp-level:0
\[
\ \text{the same tick order and the same contraction law}\ [\mathrm{BM.51}(6)]:
\]

%dsp-level:0
\[
\qquad \text{the MT relation bounds the splitting's}\ \textbf{resolvability},\ \neg\text{its magnitude};
\]

%dsp-level:0
\[
\text{the}\ \mathrm{BM.8}(2)\ \text{stationarity scope and the}\ \mathrm{CT2.6}\ \text{hypotheses are inherited at the point of use};
\]

%dsp-level:0
\[
\qquad \mathrm{FT.159}\text{'s rate-modulation residue and}\ \mathrm{FT.131}(3)\text{'s frequency read stand untouched}.
\]

\textbf{Rmk (the clock is the ledger, and the ledger is honest about it).} A uniform proper time would license an exponential law with one rate; the corpus's clock is a strictly contracting tick whose only derived rates are rational functions of one threshold. The Mandelstam–Tamm form is the admissible $\Delta t\,\Delta E$ read — saturated by an exact two-level witness — and applied to the splitting it produces the tier's one new bound: no clock coarser than $\pi\mathfrak h/(3ce^2\ln2)$ resolves the scar's field energy. The bound is symbolic in the coupling, like everything else here, and it inverts correctly: a smaller coupling means a smaller splitting means a longer required resolution.

\textbf{Req.} D1.8 (Part 1); CT1.12, CT2.1, CT2.6 (Part 2); P5 (Parts 1–2); L2 (Part 1); FT.131(3), FT.159 (Part 3); Def-BM.4, Def-BM.10, BM.8(2), BM.10(4), BM.17, BM.26, BM.28(1), BM.34(4), BM.34(5), BM.34(7), BM.44(1), BM.44(2), BM.44(3), BM.44(4), BM.46(4), BM.51(2), BM.51(3), BM.51(6) (this part); the logarithm of an exact power identity, exact $\mathbb Q(i)$ witness algebra, Laurent reciprocals [math; machine: \texttt{checks/\allowbreak{}verify\_\allowbreak{}H\_\allowbreak{}full\_\allowbreak{}scope\_\allowbreak{}mass\_\allowbreak{}splitting.py} §C].

\textbf{Pf.}

%dsp-level:0
\[
(1)\text{--}(2):\ \mathrm{BM.28}(1)\text{'s census re-run: the exponent of the charge symbol is}\ 0\ \text{in all seven monomials}\ [\text{machine}];
\]

%dsp-level:0
\[
\qquad E(p{=}0)=Nm\ [\mathrm{BM.34}(7)\ \text{re-cited}];
\]

%dsp-level:0
\[
\qquad \Delta m_{\mathrm{metric}}=0-0\ [\text{machine}]
\]

%dsp-level:0
\[
(3):\ \mathrm{BM.51}(3)\text{'s iteration re-cited};
\]

%dsp-level:0
\[
\qquad \ln(a^n)=n\ln a\ [\text{math}]\ \text{applied to the exact rational base}\ \bigl(1-\tfrac\gamma4\bigr)^2;
\]

%dsp-level:0
\[
\qquad \mathrm{BM.44}(1)\text{--}(4)\text{'s rows re-cited (frozen)}
\]

%dsp-level:0
\[
(4):\ \sigma_x(1,i)=(i,1):\ \psi^\dagger\sigma_x\psi=1\cdot i+(-i)\cdot1=0;
\]

%dsp-level:0
\[
\qquad \sigma_z(1,i)=(1,-i):\ \psi^\dagger\sigma_z\psi=1-1=0;
\]

%dsp-level:0
\[
\qquad \sigma_y(1,i)
\]

%dsp-level:0
\[
= (1,i)\cdot\ldots:\ \psi^\dagger\sigma_y\psi
\]

%dsp-level:0
\[
= 1\cdot1+(-i)(i)
\]

%dsp-level:0
\[
= 2\implies\langle\sigma_y\rangle
\]

%dsp-level:0
\[
= 1\ [\psi^\dagger\psi=2;\ \text{machine over}\ \mathbb Q(i)]
\]

%dsp-level:0
\[
(\Delta A)^2=\tfrac14\langle\sigma_x^2\rangle=\tfrac14,\quad(\Delta E)^2=\tfrac14;
\]

%dsp-level:0
\[
\qquad [H,A]=\tfrac14[\sigma_z,\sigma_x]=\tfrac i2\sigma_y\implies\tfrac{d\langle A\rangle}{dt}=i\langle[H,A]\rangle=-\tfrac12;
\]

%dsp-level:0
\[
\qquad \Delta t\,\Delta E=\tfrac{\Delta A}{1/2}\cdot\tfrac12=\tfrac12\ [\text{machine: saturation}]
\]

%dsp-level:0
\[
(5):\ \frac1{2\Delta m_{\mathrm{EM}}}
\]

%dsp-level:0
\[
= \frac1{2\cdot\frac{3ce^2\ln2}{2\pi\mathfrak h}}
\]

%dsp-level:0
\[
= \frac{\pi\mathfrak h}{3ce^2\ln2}\ [\text{machine: the Laurent reciprocal}];
\]

%dsp-level:0
\[
\qquad (6):\ \text{the scope record}.\ \blacksquare
\]

\textbf{Phys.} The metric is the background the dispersion reads ride on, and between the two baryon rows it makes no difference at all: every geometric monomial the calibration tiers locked is charge-blind, and the rest slice multiplies a shared lapse by a shared mass. The clock, meanwhile, is not uniform and is never treated as if it were — the tick contracts the Lyapunov value by an exact rational at every crystallization, the log-linear law is the logarithm of that identity and nothing more, and the only decay-style curve the corpus admits is the per-tick resolvent with its three-tick floor. The Mandelstam–Tamm relation is the correct handling of such a clock, saturated by the same two-level witness the Noether tier built, and applied to the splitting it yields the tier's sharpest new number: resolving the scar's field energy takes a clock finer than $\pi\mathfrak h$ over $3ce^2\ln2$. The bound moves exactly as the physics should — weaker coupling, smaller splitting, longer patience.

\end{theorem}

\begin{theorem}[{\texorpdfstring{BM.54 — The binding energies: the count zero, the degenerate ensemble, and the unconstructed atom}{BM.54 - The binding energies: the count zero, the degenerate ensemble, and the unconstructed atom}}]\label{tag:BM-54}
\textbf{Thm.}

%dsp-level:0
\[
(1)\quad\textbf{the count binding is exactly zero on both rows:}
\]

%dsp-level:0
\[
\qquad E_{\mathrm{int}}(\Sigma_{12})=6=(c_1{:}\,2)+(c_2{:}\,2)+(\text{connector}{:}\,2)\ [\text{Census FT.35}(2);\ \mathrm{BM.21}(1)]:
\]

%dsp-level:0
\[
E^{\mathrm{info}}_{\mathrm{bind}}=\bigl(6-(2{+}2{+}2)\bigr)\tfrac{\ln2}{\mathfrak h}=0
\]

%dsp-level:0
\[
\qquad [\text{machine}];
\]

%dsp-level:0
\[
\qquad \text{both baryon rows carry the same decomposition}\ [\mathrm{BM.17};\ \mathrm{BM.26}]:
\]

%dsp-level:0
\[
\qquad \Delta\bigl(E^{\mathrm{info}}_{\mathrm{bind}}\bigr)=0;
\]

%dsp-level:0
\[
(2)\quad\textbf{the ensemble binding is nonzero but class-degenerate:}
\]

%dsp-level:0
\[
\qquad \mathcal B(S_j;1,1,1)
\]

%dsp-level:0
\[
= m^2_B(\cdot,S_j)-\textstyle\sum_im^2_B(\cdot,S_i)
\]

%dsp-level:0
\[
\ne 0\ \ \text{exact rationals}\ [\mathrm{BM.21}(2)\ \text{re-cited (frozen)}];
\]

%dsp-level:0
\[
m^2_B(q,S)\ \text{is independent of}\ q\ [\mathrm{FT.107}(3);\ \text{the cancellation lemma}\ \mathrm{BM.30}(3);\ \text{machine re-run}]:
\]

%dsp-level:0
\[
\qquad \mathcal B^{\,(n)}=\mathcal B^{\,(p)}\ \ \text{exactly}\implies\boxed{\ \Delta m_{\mathrm{binding}}=0\ }\ \ \text{on every normalized ensemble read};
\]

%dsp-level:0
\[
(3)\quad\textbf{the charge configuration does not enter either binding read:}
\]

%dsp-level:0
\[
\qquad \text{the count read is}\ E_{\mathrm{int}}\text{-additive}\ [(1);\ \mathrm{BM.21}(1)];
\]

%dsp-level:0
\[
\qquad \text{the ensemble deficit}\ \Delta S=\tfrac32(2E-|R|)\ \text{is}\ \sigma\text{-free}\ [\mathrm{FT.107}(2)]:
\]

%dsp-level:0
\[
\text{the end-sign configurations that distinguish the classes}\ [\mathrm{D3.4};\ \mathrm{BM.28}(2)]\ \text{are invisible to both binding reads}:
\]

%dsp-level:0
\[
\qquad \textbf{``different charge configurations}\implies\text{different binding energies'' has no derived object between the two rows};
\]

%dsp-level:0
\[
(4)\quad\textbf{the EM binding belongs to the unconstructed pair system:}
\]

%dsp-level:0
\[
\qquad \text{the lepton--baryon bound channel is}\ \neg\text{constructed}\ [\mathrm{BM.48}(3):\ \mathrm{Def\text{-}BM.1}\text{'s grammar is colour-sealed;}\ \neg\text{an EM sealing grammar};\ \mathrm{BM.6}(5)\ \text{dynamics open}];
\]

%dsp-level:0
\[
\textbf{the category bar:}\quad\text{a binding energy is the energy of the bound system,}\ \neg\text{a correction to a free member's mass}\ [\mathrm{BM.48}(4)];
\]

%dsp-level:0
\[
\qquad \text{for a constructed pair the system read would be}\ \ m_{\mathrm{pair}}=m_p+m_\ell+E_{\mathrm{bind}},\quad E_{\mathrm{bind}}<0;
\]

%dsp-level:0
\[
\textbf{the sign row:}
\]

%dsp-level:0
\[
\qquad E_1=-\frac{\varepsilon_{\mathrm{EM}}^2\,m_r\,(q_1q_2)^2}{2}<0\ \ \text{at positive symbols}
\]

%dsp-level:0
\[
\qquad [\mathrm{BM.48}(2);\ \text{machine: the sign census}]:
\]

%dsp-level:0
\[
\qquad \text{where binding exists it is negative}\ —\ \text{and it exists, in this inventory, nowhere yet};
\]

%dsp-level:0
\[
(5)\quad\textbf{status:}
\]

%dsp-level:0
\[
\qquad \Delta m_{\mathrm{binding}}=0\ \text{on every constructed read}\ ((1)\text{--}(3));
\]

%dsp-level:0
\[
\qquad \text{the hypothetical pair system carries the only negative binding rows}\ ((4)),
\]

%dsp-level:0
\[
\ \text{as system reads},
\]

%dsp-level:0
\[
\ \neg\text{as member splittings};
\]

%dsp-level:0
\[
\text{the route-(ii) measure price stays uncombined}\ [\mathrm{BM.30}(5);\ \mathrm{BM.37}(3)]:
\]

%dsp-level:0
\[
\qquad \neg\text{a measured binding energy, ionization scale or atomic datum}\ [\mathrm{P5}].
\]

\textbf{Rmk (two zeros and one absence.)} The binding contribution to the splitting dies twice over, and both deaths are exact. The count binding is zero because six edges decompose as two plus two plus two — the informational read is additive and the additivity is an integer identity. The ensemble binding is nonzero, four exact rationals of it, but every one is built from charge-degenerate sector masses, so the neutron and the proton carry identical binding and their difference is zero again. The only binding that could distinguish the classes would be electromagnetic, and that lives in an atom — a bound channel the corpus's grammar has never sealed. The category bar keeps the bookkeeping honest in the meantime: negative energy belongs to systems, not to members.

\textbf{Req.} D2.9, D3.4 (Parts 1–2); P5 (Parts 1–2); FT.107(2), FT.107(3) (Part 3); Census FT.35(2) (row-wise); Def-BM.1, Def-BM.5, BM.6(5), BM.17, BM.20(4), BM.21(1), BM.21(2), BM.26, BM.28(2), BM.30(3), BM.30(5), BM.37(3), BM.48(2), BM.48(3), BM.48(4) (this part); exact rational arithmetic, sign censuses, the cancellation lemma [math; machine: \texttt{checks/\allowbreak{}verify\_\allowbreak{}H\_\allowbreak{}full\_\allowbreak{}scope\_\allowbreak{}mass\_\allowbreak{}splitting.py} §D].

\textbf{Pf.}

%dsp-level:0
\[
(1):\ \text{Census FT.35}(2)\text{'s decomposition re-cited};
\]

%dsp-level:0
\[
\qquad 6-(2{+}2{+}2)=0\ [\text{machine}];
\]

%dsp-level:0
\[
\qquad 0\cdot\tfrac{\ln2}{\mathfrak h}=0\ [\text{machine}]
\]

%dsp-level:0
\[
(2):\ \mathrm{BM.21}(2)\text{'s four frozen rationals re-cited};
\]

%dsp-level:0
\[
\qquad \text{the cancellation lemma at the six frozen factors}\ \{1,2,3,15,20,\tfrac73\}\ [\mathrm{BM.30}(3)\ \text{re-run};\ \text{machine}]
\]

%dsp-level:0
\[
\implies q\text{-independence}\implies\mathcal B^{(n)}-\mathcal B^{(p)}\qquad = 0
\]

%dsp-level:0
\[
(3):\ \mathrm{FT.107}(2)\text{'s}\ \sigma\text{-freeness and}\ \mathrm{BM.21}(1)\text{'s additivity conjoined}
\]

%dsp-level:0
\[
(4):\ \mathrm{BM.48}(2)\text{--}(4)\ \text{re-cited};
\]

%dsp-level:0
\[
\qquad E_1<0\ \text{at}\ \varepsilon_{\mathrm{EM}}^2>0,\ m_r>0\ [\text{machine: sign census on rational witnesses}]
\]

%dsp-level:0
\[
(5):\ \text{the status record}.\ \blacksquare
\]

\textbf{Phys.} Binding energy is the one contribution that could have been expected to favor the neutron — a neutral composite might bind its parts more gently or more tightly than a charged one — and the corpus answers with two exact zeros and one absence. The count binding vanishes as an integer identity: six is two plus two plus two. The ensemble binding is real, four nonzero rationals, and completely charge-blind, because every sector mass inside it is. And the binding that would care about charge is the electromagnetic atom, which the colour-sealed grammar has not constructed and the category bar keeps out of the free members' masses in the meantime. The difference of the two rows' binding energies is zero on every read the corpus owns.

\end{theorem}

\begin{theorem}[{\texorpdfstring{BM.55 — The full-scope splitting and the exact ratio}{BM.55 - The full-scope splitting and the exact ratio}}]\label{tag:BM-55}
\textbf{Thm.}

%dsp-level:0
\[
(1)\quad\textbf{the decomposition, collapsed:}
\]

%dsp-level:0
\[
\qquad \Delta m:
\]

%dsp-level:0
\[
= m_n-m_p
\]

%dsp-level:0
\[
= \Delta m_{\mathrm{Lyapunov}}+\Delta m_{\mathrm{force}}+\Delta m_{\mathrm{metric}}+\Delta m_{\mathrm{binding}}+\Delta m_{\mathrm{self\text{-}int}}\ \ \text{with the derived rows}:
\]

%dsp-level:0
\[
\Delta m_{\mathrm{Lyapunov}}=0\ [\mathrm{BM.51}(6)];\qquad\Delta m_{\mathrm{metric}}=0\ [\mathrm{BM.53}(2)];\qquad\Delta m_{\mathrm{binding}}=0\ [\mathrm{BM.54}(2),(5)];
\]

%dsp-level:0
\[
\Delta m_{\mathrm{force}}
\]

%dsp-level:0
\[
= \Delta m_{\mathrm{self\text{-}int}}
\]

%dsp-level:0
\[
= -\Delta m_{\mathrm{EM}}\ [\mathrm{BM.52}(4)\text{--}(5):\ \text{the two routes are one monomial, counted once}]:
\]

%dsp-level:0
\[
\boxed{\ \Delta m=m_n-m_p=-\,\Delta m_{\mathrm{EM}}=-\,\frac{3\,c\,e^2\ln2}{2\pi\,\mathfrak h}\ <\ 0\ }\qquad[\text{machine: the monomial sum and the sign census}];
\]

%dsp-level:0
\[
(2)\quad\textbf{the exact mass ratio:}
\]

%dsp-level:0
\[
\qquad m_n=\frac{6\ln2}{\mathfrak h},
\]

%dsp-level:0
\[
\qquad m_p=\frac{6\ln2}{\mathfrak h}\bigl(1+c\,\varepsilon_{\mathrm{EM}}\bigr),
\]

%dsp-level:0
\[
\qquad \varepsilon_{\mathrm{EM}}=\frac{e^2}{4\pi}\ [\mathrm{BM.26};\ \mathrm{BM.46}(1),(3)]:
\]

%dsp-level:0
\[
\boxed{\ \frac{m_n}{m_p}=\frac1{1+c\,\varepsilon_{\mathrm{EM}}}=\frac{4\pi}{4\pi+c\,e^2}\ }
\]

%dsp-level:0
\[
\qquad [\text{machine: the cross-multiplied Laurent identity}\ m_n(4\pi+ce^2)=m_p\cdot4\pi];
\]

%dsp-level:0
\[
\textbf{the geometric expansion:}
\]

%dsp-level:0
\[
\qquad \frac{m_n}{m_p}=1-c\,\varepsilon_{\mathrm{EM}}+(c\,\varepsilon_{\mathrm{EM}})^2-(c\,\varepsilon_{\mathrm{EM}})^3+O(\varepsilon^4)
\]

%dsp-level:0
\[
\qquad [\text{machine: the truncated series identity}\ (1+x)(1-x+x^2-x^3)=1-x^4];
\]

%dsp-level:0
\[
(3)\quad\textbf{the strict order:}
\]

%dsp-level:0
\[
\qquad 0<\frac{m_n}{m_p}<1\ \ \text{for every}\ e^2>0\ \ \text{and both priced}\ c\in\{\tfrac12,\tfrac35\}
\]

%dsp-level:0
\[
\qquad [\text{machine: the denominator exceeds the numerator by}\ ce^2>0];
\]

%dsp-level:0
\[
\frac{m_n}{m_p}=1\iff e=0:
\]

%dsp-level:0
\[
\qquad \textbf{the full-scope corpus computation places the neutron strictly below the proton on every constructed read};
\]

%dsp-level:0
\[
(4)\quad\textbf{the order residue:}
\]

%dsp-level:0
\[
\qquad \text{reversing the order requires a class-dependent constituent transport splitting}\ (m^{(n)}_{\mathrm{constituent}}>m^{(p)}_{\mathrm{constituent}}):
\]

%dsp-level:0
\[
\neg\exists\ \text{such an object in Parts 1--4}:
\]

%dsp-level:0
\[
\qquad \text{the fractional constituent alphabet is unreachable by the additive read}\ [\mathrm{BM.28}(7);\ \mathrm{Def\text{-}BM.8}(5)];
\]

%dsp-level:0
\[
\qquad \text{the species labels carry no datum}\ [\mathrm{Def\text{-}BM.3}];
\]

%dsp-level:0
\[
\text{the deficit/inertial conversions are item-7 content}\ [\text{pointer}];
\]

%dsp-level:0
\[
\qquad \text{the family-profile map of}\ \mathrm{Def\text{-}BM.11}\ \text{acts on channel lines, and no clause attaches constituent profiles to the baryon rows' charge classes};
\]

%dsp-level:0
\[
\textbf{the residue is named, not filled};\qquad\neg\text{a measured mass difference is compared}\ [\mathrm{P5}];
\]

%dsp-level:0
\[
(5)\quad\textbf{the decay consequence:}
\]

%dsp-level:0
\[
\qquad Q_{\mathrm{corr}}=-(\Delta m_{\mathrm{EM}}+m_\ell+m_\nu)<0\ [\mathrm{BM.49}(1)\ \text{re-cited}]\ \ \text{now with the full-scope census behind it}:
\]

%dsp-level:0
\[
\textbf{the}\ \beta\textbf{-decay vertex is kinematically closed on every constructed read, and the closure survives the complete force operator}\ [\mathrm{BM.52}(4);\ \mathrm{BM.49}(2)\ \text{re-cited}];
\]

%dsp-level:0
\[
\qquad \text{the route-(ii) measure price stands uncombined}\ [\mathrm{BM.30}(5);\ \mathrm{BM.54}(5)];
\]

%dsp-level:0
\[
(6)\quad\textbf{status:}
\]

%dsp-level:0
\[
\qquad \text{the ratio is an exact symbolic rational function of the logged P4 coupling and the priced distribution factor}\ —\ \neg\text{a number},
\]

%dsp-level:0
\[
\ \neg\text{a fit},
\]

%dsp-level:0
\[
\ \neg\text{a comparison}\ [\mathrm{P5};\ \mathrm{EM.25}(7)];
\]

%dsp-level:0
\[
\text{the five-term decomposition is derived term by term; four terms are zero and the fifth is the BM.46 monomial counted once}\ [(1);\ \mathrm{BM.52}(5)].
\]

\textbf{Rmk (the answer the corpus gives.)} Asked for the neutron–proton mass ratio from its own structures alone, the corpus answers with a rational function: four pi over four pi plus $ce^2$. It is exact, it is symbolic in the one coupling no theorem fixes, and it is strictly below one — the neutron comes out lighter, because the only sector of the complete force operator that distinguishes the two rows is the one whose field energy the proton carries and the neutron does not. The order the decay would need is the order no constructed read supplies, and the tier names the missing object precisely: a constituent-level transport difference between the classes, of which the inventory contains no instance. That is the residue; everything else is arithmetic.

\textbf{Req.} P1, P4, P5 (Parts 1–2); CT2.1, EM.25(7) (Part 2); Def-BM.3, Def-BM.8, Def-BM.11, Def-BM.12, BM.17, BM.26, BM.28(7), BM.30(5), BM.37(3), BM.39, BM.46(1), BM.46(3), BM.46(4), BM.48, BM.49(1), BM.49(2), BM.51(6), BM.52(4), BM.52(5), BM.53(2), BM.54(2), BM.54(5) (this part); FT.79(5) (Part 3); Laurent cross-multiplication, truncated geometric series over $\mathbb Q$, sign censuses [math; machine: \texttt{checks/\allowbreak{}verify\_\allowbreak{}H\_\allowbreak{}full\_\allowbreak{}scope\_\allowbreak{}mass\_\allowbreak{}splitting.py} §E].

\textbf{Pf.}

%dsp-level:0
\[
(1):\ \text{the census sum}\ 0+0+0+0-\Delta m_{\mathrm{EM}}\ [\mathrm{BM.52}(4)\ \text{re-cited};\ \text{machine}];
\]

%dsp-level:0
\[
\qquad \text{the double-count guard}\ [\mathrm{BM.52}(5):\ \text{one monomial}];
\]

%dsp-level:0
\[
\qquad \text{the sign: the coefficient census at}\ c,e^2,\ln2,\mathfrak h>0\ [\text{machine}]
\]

%dsp-level:0
\[
(2):\ m_n(4\pi+ce^2)
\]

%dsp-level:0
\[
= \tfrac{6\ln2}{\mathfrak h}(4\pi+ce^2)\ \ \text{and}\ \ m_p\cdot4\pi
\]

%dsp-level:0
\[
= \tfrac{6\ln2}{\mathfrak h}\bigl(1+\tfrac{ce^2}{4\pi}\bigr)4\pi
\]

%dsp-level:0
\[
= \tfrac{6\ln2}{\mathfrak h}(4\pi+ce^2)\ [\text{machine: the cross-multiplied Laurent identity}];
\]

%dsp-level:0
\[
\qquad (1+x)(1-x+x^2-x^3)=1-x^4\ [\text{machine}]
\]

%dsp-level:0
\[
(3):\ (4\pi+ce^2)-4\pi=ce^2>0\ [\text{machine: coefficient census}];
\]

%dsp-level:0
\[
\qquad \text{both numerator and denominator positive at positive symbols}\implies\text{the ratio}\in(0,1);
\]

%dsp-level:0
\[
\qquad ce^2=0\iff e=0\ [\text{machine}]
\]

%dsp-level:0
\[
(4):\ \mathrm{BM.28}(7)/\mathrm{Def\text{-}BM.8}(5),\ \mathrm{Def\text{-}BM.3},\ \text{the}\ \mathrm{Def\text{-}BM.11}\ \text{locus clause re-cited};
\]

%dsp-level:0
\[
\qquad (5):\ \mathrm{BM.49}(1)\text{--}(2)\ \text{re-cited};
\]

%dsp-level:0
\[
\qquad (6):\ \text{the status record}.\ \blacksquare
\]

\textbf{Phys.} The full-scope question has a full-scope answer, and it is a single rational function. Four of the five named contributions are theorems about zero: the Lyapunov tower shares every input between the rows, the metric reads are charge-blind, the binding dies twice, and the force census leaves one sector standing. That sector's monomial is the whole splitting, negative for the neutron-minus-proton difference, and the ratio follows by cross-multiplication: four pi over four pi plus $c\,e^2$, strictly between zero and one at every positive coupling, equal to one only where the coupling dies. The corpus therefore places the neutron below the proton, deeper than the informational count alone did, and the beta decay stays shut with the complete force operator now behind the closure. What would open it — a constituent transport difference between the two rows — is named exactly and found absent: no fractional alphabet, no labelled datum, no converted species. The ratio is derived; the order is derived; and the distance to the world's order is the residue, stated to the monomial.

\end{theorem}

\medskip\noindent\textit{\texorpdfstring{BM.56 — The status of the full-scope mass splitting: adjudication and scope}{BM.56 - The status of the full-scope mass splitting: adjudication and scope}.}\ Structural-constraints record: see Appendix B, Section B.1.22 of the Addendum.\par\medskip

\phantomsection
\section[Tier BM.14 — The Constituent Transport Splitting: the…]{\texorpdfstring{Tier BM.14 — The Constituent Transport Splitting: the Vertex Asymmetry, the Observation Frames, the Six-State Quark Assembly, and the Neutron Row}{Tier BM.14 - The Constituent Transport Splitting: the Vertex Asymmetry, the Observation Frames, the Six-State Quark Assembly, and the Neutron Row}}

\textbf{Locus.} This tier fills the residue $\mathrm{BM.55}(4)$ names and does not fill: the class-dependent constituent transport splitting. It extends $\mathrm{BM.55}(1)$'s five-term decomposition by one term, amends $\mathrm{BM.52}(3)$'s sector census by one row, and reverses $\mathrm{BM.55}(3)$'s strict order on an exact sub-critical branch. Every clause of $\mathrm{BM.51}$–$\mathrm{BM.56}$ it touches is re-cited at its point of use and stands unmodified except where a clause's own scope row is widened; the widening is recorded in $\mathrm{BM.64}(3)$.

\textbf{Two directions.} The tier is worked from both ends at once. From the force structure: what a sealed junction can carry, and what the seal annihilates — $\mathrm{BM.57}$ and $\mathrm{BM.58}$ run that direction to exhaustion and refute both proposed mechanisms. From the mass spectrum: what $\mathrm{BM.55}(1)$'s four zeros leave, and what the inventory contains that is class-dependent and not barred — $\mathrm{BM.59}$ runs that direction. The two meet on one object, $W(Q)$, and $\mathrm{BM.62}$ assembles it. $\mathrm{BM.60}$ supplies the six-state spectrum the constituent row of $\mathrm{BM.59}(5)$ needs; $\mathrm{BM.61}$ disposes of the proposed coupling $\alpha_B$; $\mathrm{BM.63}$ re-adjudicates the decay; $\mathrm{BM.64}$ adjudicates the tier.

\begin{definition}[{\texorpdfstring{Def-BM.13 — The constituent transport refinement of the rest input (the construction)}{Def-BM.13 - The constituent transport refinement of the rest input (the construction)}}]\label{tag:Def-BM-13}
\textbf{Def.} A refinement of the dispersion route's rest input $[\mathrm{Def\text{-}BM.10}(4)$; $\mathrm{BM.35}(4)]$ built from the additive charge read's own census $[\mathrm{Def\text{-}BM.8}(6)]$; no new axiom, no new datum, no new parameter, no numerical coupling, no route adoption. Local notation: $\varepsilon_{\mathrm{EM}}:=e^2/(4\pi)$, a dimensionless P4-class symbol carried symbolically throughout $[\mathrm{Def\text{-}BM.12}$ local notation; $\mathrm{EM.25}(7),(9)$; $\mathrm{P5}]$.

%dsp-level:0
\[
(1)\quad\textbf{the constituent transport multiplicity:}
\]

%dsp-level:0
\[
\qquad q_{\mathrm{core}}(c)=\tfrac12\bigl(\sigma_{s(c)}-\sigma_{t(c)}\bigr)\in\{-1,0,+1\}
\]

%dsp-level:0
\[
\qquad [\mathrm{D3.4};\ \mathrm{EM.9}];
\]

%dsp-level:0
\[
\textbf{the end-pair census:}
\]

%dsp-level:0
\[
\qquad (\sigma_s,\sigma_t)\in\{(+,+),(+,-),(-,+),(-,-)\}\ \longmapsto\ q=(0,+1,-1,0)
\]

%dsp-level:0
\[
\qquad [\mathrm{EM.9}\ \text{table};\ \text{machine}]:
\]

%dsp-level:0
\[
\boxed{\ m_q=(2,\,1,\,1)\ \ \text{at}\ \ q=(0,+1,-1)\ }\qquad[\mathrm{FT.106}(3\mathrm a):\ \text{the norm-layer closed form}\ m_q\,2^{|R|-2}]:
\]

%dsp-level:0
\[
\textbf{a neutral-core constituent carries two transport realizations; a charged-core constituent carries one};
\]

%dsp-level:0
\[
\qquad \text{the two neutral realizations are exchanged by the global sign flip, which maps}\ q
\]

%dsp-level:0
\[
\mapsto -q\ \text{on a charged one}\ [\mathrm{EM.9a};\ \mathrm{EM.24}(1)]:
\]

%dsp-level:0
\[
\textbf{the multiplicity is the orbit size of the constituent's transport register under}\ \omega\ \text{within its own charge class};
\]

%dsp-level:0
\[
(2)\quad\textbf{the parity law of the chain wall:}
\]

%dsp-level:0
\[
\qquad \text{for an oriented chain}\ c=(e_1,\dots,e_k),\ k\ge2,\ \text{with interior vertices of type}\ (1,1)\ [\mathrm{A7};\ \mathrm{D0.8}]:
\]

%dsp-level:0
\[
\sum_{i=1}^{k-1}f_{v_i}=\tfrac12\bigl(\sigma_{e_k}-\sigma_{e_1}\bigr)=-\,q_{\mathrm{EM}}(c),\qquad\sum_{i=1}^{k-1}\aleph_{v_i}=\#\{i:\ \sigma_{e_{i+1}}\ne\sigma_{e_i}\}\qquad[\mathrm{EM.9b};\ \mathrm{D3.3}]:
\]

%dsp-level:0
\[
\aleph_v=|f_v|\in\{0,1\}\implies\boxed{\ \sum_i\aleph_{v_i}\equiv|q|\ (\mathrm{mod}\ 2),\qquad\sum_i\aleph_{v_i}\ \ge\ |q|\ }
\]

%dsp-level:0
\[
\qquad [\text{machine: exhaustive over every sign sequence,}\ k=2,\dots,7]:
\]

%dsp-level:0
\[
\textbf{the interior wall content of a constituent's transport is parity-locked to its charge core:}
\]

%dsp-level:0
\[
\qquad q=0:\ \min\textstyle\sum\aleph=0\ (\text{even});
\]

%dsp-level:0
\[
\qquad |q|=1:\ \min\textstyle\sum\aleph=1\ (\text{odd});
\]

%dsp-level:0
\[
(3)\quad\textbf{the class transport weight:}
\]

%dsp-level:0
\[
\qquad W(Q)\ :=\ \sum_{\substack{(q_1,\dots,q_k)\in\{-1,0,+1\}^k\  \sum_iq_i=Q}}\ \prod_{i=1}^{k}m_{q_i}\ =\ \bigl[z^{Q}\bigr]\bigl(2+z+z^{-1}\bigr)^{k}
\]

%dsp-level:0
\[
\qquad [\mathrm{Def\text{-}BM.8}(6)]:
\]

%dsp-level:0
\[
k=3:
\]

%dsp-level:0
\[
\qquad W=(1,\,6,\,15,\,\mathbf{20},\,15,\,6,\,1)\ \ \text{over}\ \ Q=-3,\dots,+3,
\]

%dsp-level:0
\[
\qquad \sum_{Q}W(Q)=4^3=64=2^6
\]

%dsp-level:0
\[
\qquad [\text{machine: the exhaustive}\ 4^3\ \text{end-sign census and the convolution agree exactly}];
\]

%dsp-level:0
\[
\textbf{the ordered charge-triple census:}
\]

%dsp-level:0
\[
\qquad \#\{(q_1,q_2,q_3):\ \textstyle\sum q_i=k\}=(1,3,6,\mathbf 7,6,3,1):
\]

%dsp-level:0
\[
\qquad 7\ \text{neutral triples against}\ 6\ \text{charged}\ [\mathrm{BM.28}(2);\ \text{machine}];
\]

%dsp-level:0
\[
\textbf{the factorization identity:}
\]

%dsp-level:0
\[
\qquad W(Q)=\sum_{\text{triples at}\ Q}\prod_im_{q_i}
\]

%dsp-level:0
\[
\qquad [\text{machine: the two routes agree class by class}];
\]

%dsp-level:0
\[
(4)\quad\textbf{the charge-class transport entropy:}\qquad\boxed{\ \Sigma_{\mathrm{cl}}(Q)\ :=\ \ln W(Q)\ }\qquad\Sigma_{\mathrm{cl}}(0)=\ln20,\quad\Sigma_{\mathrm{cl}}(\pm1)=\ln15,
\]

%dsp-level:0
\[
\boxed{\ \Sigma_{\mathrm B}\ :=\ \Sigma_{\mathrm{cl}}(0)-\Sigma_{\mathrm{cl}}(+1)\ =\ \ln\frac{W(0)}{W(+1)}\ =\ \ln\frac43\ =\ 2\ln2-\ln3\ >\ 0\ }
\]

%dsp-level:0
\[
\qquad [\text{machine: the certificate}\ 4>3]:
\]

%dsp-level:0
\[
\textbf{the constituent row:}\qquad\ln m_0-\ln m_{\pm1}=\ln2>0\qquad[\text{machine:}\ 2>1]:
\]

%dsp-level:0
\[
\textbf{the neutral-core constituent is the heavier one, by exactly}\ \ln2/\mathfrak h;
\]

%dsp-level:0
\[
\qquad \textbf{the additive constraint attenuates it:}\quad0<\ln\tfrac43<\ln2<3\ln2
\]

%dsp-level:0
\[
\qquad [\text{machine:}\ 3<4<6<24];
\]

%dsp-level:0
\[
(5)\quad\textbf{the refined rest input:}\qquad\boxed{\ m_{\mathrm{rest}}(Q)\ :=\ \frac{E_{\mathrm{int}}\ln2+\Sigma_{\mathrm{cl}}(Q)}{\mathfrak h}\ +\ c\,\varepsilon_{\mathrm{EM}}\,Q^2\,m_{\mathrm{hid}}\ }
\]

%dsp-level:0
\[
\bigl[\ \text{the first summand:}\ \mathrm{BM.35}(4)\text{'s reframing identity}\ m_{\mathrm{rest}}=\Delta S_{\mathfrak L}/\mathfrak h\ \text{with the register summand of}\ (4);\qquad\text{the second:}\ \mathrm{BM.46}(3)\text{'s corrected rest read, riding the charge-blind Born width}\ [\mathrm{BM.28}(1)]\ \bigr];
\]

%dsp-level:0
\[
\textbf{at}\ E_{\mathrm{int}}=6\ [\mathrm{BM.17};\ \mathrm{BM.26}]:
\]

%dsp-level:0
\[
\qquad m_{\mathrm{rest}}(0)\,\mathfrak h=\ln(2^6\!\cdot\!20)=\ln1280,
\]

%dsp-level:0
\[
\qquad m_{\mathrm{rest}}(\pm1)\,\mathfrak h=\ln(2^6\!\cdot\!15)+6c\,\varepsilon_{\mathrm{EM}}\ln2=\ln960+6c\,\varepsilon_{\mathrm{EM}}\ln2
\]

%dsp-level:0
\[
\qquad [\text{machine}];
\]

%dsp-level:0
\[
(6)\quad\textbf{the currency clause:}
\]

%dsp-level:0
\[
\qquad \Sigma_{\mathrm{cl}}\ \text{is measured in nats of binary registration data}\ —\ \text{the end signs are the}\ \sigma\text{-alphabet}\ [\mathrm{CT0.9};\ \mathrm{D3.4}]\ \text{and one binary datum books}\ \ln2\ [\mathrm{C9};\ \mathrm{CT1.29}]:
\]

%dsp-level:0
\[
\textbf{it is the same currency as}\ \Delta S_{\mathfrak L}
\]

%dsp-level:0
\[
= E_{\mathrm{int}}\ln2\ [\mathrm{BM.35}(3)],
\]

%dsp-level:0
\[
\ \text{and its alphabet is disjoint from the count's:}
\]

%dsp-level:0
\[
\qquad E_{\mathrm{int}}\ \text{has no endpoint summand}\ [\mathrm{BM.28}(3):\ 6=(c_1{:}\,2)+(c_2{:}\,2)+(\text{connector}{:}\,2)]:
\]

%dsp-level:0
\[
E_{\mathrm{int}}\ \text{is constant on every end-sign fibre}\ [\text{machine: exhaustive over the four end pairs}]\implies\textbf{the two registrations are independent, so their entropies add};
\]

%dsp-level:0
\[
(7)\quad\textbf{the reference clause (the sign is fixed, not chosen):}
\]

%dsp-level:0
\[
\qquad \text{the deficit reference is the vacuum:}\quad\Delta S:=S-S(\mathrm{vac}),\quad S_{\mathfrak L}(\mathrm{vac})=0\ [\mathrm{CT0.2};\ \mathrm{D0.1};\ \mathrm{BM.35}(3)]:
\]

%dsp-level:0
\[
W(\mathrm{vac})=1\implies\Sigma_{\mathrm{cl}}(\mathrm{vac})=0\implies\Delta\Sigma_{\mathrm{cl}}(Q)=\ln W(Q)\ \text{with the}\ \textbf{positive}\ \text{sign of}\ (4);
\]

%dsp-level:0
\[
\textbf{the priced alternative:}
\]

%dsp-level:0
\[
\qquad \text{referencing against the unconstrained register}\ \bigl(4^3=64\bigr)\ \text{gives the constraint cost}\ \ln\frac{64}{W(Q)},
\]

%dsp-level:0
\[
\ \text{whose class difference is}\ -\ln\tfrac43:
\]

%dsp-level:0
\[
\text{that reference is}\ \neg\textbf{a state of the composite}\ —\ \text{the boundary data fix the charge class}\ [\mathrm{Def\text{-}BM.5}(1)]\ —\ \text{and it is class-dependent:}
\]

%dsp-level:0
\[
\qquad \textbf{priced, barred};
\]

%dsp-level:0
\[
\mathrm{Def\text{-}H.28}\text{'s Ansatz-B}\ m_B
\]

%dsp-level:0
\[
= (K-\alpha\ln\Omega)\ln2/\mathfrak h\ \text{carries the opposite sign with}\ \alpha\ \text{undeclared and is marked an external ansatz,}\ \neg\text{a corpus read:}
\]

%dsp-level:0
\[
\qquad \neg\text{adopted};
\]

%dsp-level:0
\[
(8)\quad\textbf{the bar against the measure price:}
\]

%dsp-level:0
\[
\qquad \mathrm{BM.37}(6)\ \text{bars reading a route ratio as a physical mass ratio};
\]

%dsp-level:0
\[
\qquad \textbf{this construction reads no ratio}:
\]

%dsp-level:0
\[
\text{the entry is}\ \ln W(Q)\ \text{in the deficit currency, and the resulting mass ratio is}\ \ln1280/\ln960\ne\tfrac43\ [\text{machine}];
\]

%dsp-level:0
\[
\qquad \textbf{the strict contraction:}\quad\frac{\Sigma_{\mathrm B}}{E_{\mathrm{int}}\ln2}=\frac{\ln(4/3)}{6\ln2}<\frac13
\]

%dsp-level:0
\[
\qquad [\text{machine:}\ 3\ln\tfrac43<2\ln2\iff16<27]:
\]

%dsp-level:0
\[
\text{the refinement is a strict contraction of}\ \mathrm{BM.30}(2)\text{'s measure price}\ —\ \text{it never reaches it};
\]

%dsp-level:0
\[
(9)\quad\textbf{status and scope:}
\]

%dsp-level:0
\[
\qquad \textbf{constructed}\ —\ \text{an abbreviation of derived objects}\ ((1)\text{--}(4)\ \text{are censuses};\ (5)\ \text{is an assembly of}\ \mathrm{BM.35}(4)\ \text{and}\ \mathrm{BM.46}(3));
\]

%dsp-level:0
\[
\neg\text{a declaration},\quad\neg\text{a P4 addition},\quad\neg\text{a new dimensionful scale}\ [\mathrm{P1}],\quad\text{§Cal-BM stays empty};
\]

%dsp-level:0
\[
\qquad \neg\text{a route adopted}\ [\mathrm{BM.30}(5)\ \text{stands}];
\]

%dsp-level:0
\[
\textbf{beyond this tier:}\ \text{the triple-level resolution}\ —\ \text{the chart does not resolve which triple a class member occupies, because the end signs carry no position datum}\ [\mathrm{D3.4};\ \mathrm{BM.46}(2)\ \text{resolution residue}]:\ \text{the class read of}\ (4)\ \text{is the one the chart supports}.
\]

\textbf{Rmk (what the refinement is).} The composite's rest input has always been the count of its internal directed edges, and the count is charge-blind to the last monomial. What the charge class does carry, and has always carried, is a filter on the transport register: twenty admissible end-sign assignments in the neutral class against fifteen in the charged one. The refinement takes the logarithm of that filter and books it where the corpus books every other entropy — against the vacuum, in the ledger's own currency, through the rest slot the dispersion read already reads. It is not a fourth mass route: it adds no functional of the flicker configurations, and it takes no ratio. It is the smallest entry the inventory supplies that is class-dependent, currency-matched, and alphabet-disjoint from the count.

\end{definition}

\pseudohead{definition}{\texorpdfstring{Def-BM.14}{Def-BM.14}}{\texorpdfstring{The six-state assembly: the Platonic floors, the defect branches, and the generation map}{The six-state assembly: the Platonic floors, the defect branches, and the generation map}}\label{tag:Def-BM-14}
\textbf{Def.} A census assembled from the Platonic card template $[\mathrm{Census\ FT.42}(5)]$, the duality degeneracy and the ledger splitting $[\mathrm{FT.44}(1)\text{--}(3)]$, the physical-card defect spectrum $[\mathrm{Census\ FT.63}(6)]$, and the $q$-channel triplet $[\mathrm{FT.119};\ \mathrm{Def\text{-}BM.6}(3)]$; no new axiom, no new datum, no new parameter, no external mass.

%dsp-level:0
\[
(1)\quad\textbf{the three mass floors:}
\]

%dsp-level:0
\[
\qquad m_{\mathrm{hid}}=\frac{2E\ln2}{\mathfrak h},
\]

%dsp-level:0
\[
\qquad 2E\in\{12,\,24,\,24,\,60,\,60\}\ \ \text{over the five solids}
\]

%dsp-level:0
\[
\qquad [\mathrm{FT.42}(1),\ \text{card table}]:
\]

%dsp-level:0
\[
\textbf{dual pairs are mass-degenerate and}\ \Pi\textbf{-degenerate at matched}\ K_S\ [\mathrm{FT.44}(1)]:
\]

%dsp-level:0
\[
\qquad \text{cube/octahedron}\ 24\ln2/\mathfrak h;
\]

%dsp-level:0
\[
\quad\text{dodecahedron/icosahedron}\ 60\ln2/\mathfrak h;
\]

%dsp-level:0
\[
\quad\text{tetrahedron self-dual}\ 12\ln2/\mathfrak h;
\]

%dsp-level:0
\[
\boxed{\ \textbf{three mass values across the five solids}\ }
\]

%dsp-level:0
\[
\qquad \mathcal F:
\]

%dsp-level:0
\[
= \bigl\{\mathcal F_1{=}\{\text{tetrahedron}\},\ \ \mathcal F_2{=}\{\text{cube},\text{octahedron}\},\ \ \mathcal F_3{=}\{\text{dodecahedron},\text{icosahedron}\}\bigr\};
\]

%dsp-level:0
\[
2E(\mathcal F_1){=}12\ <\ 2E(\mathcal F_2){=}24\ <\ 2E(\mathcal F_3){=}60,
\]

%dsp-level:0
\[
\qquad 12:24:60=1:2:5
\]

%dsp-level:0
\[
\qquad [\text{machine: exact rationals, strictly increasing}];
\]

%dsp-level:0
\[
(2)\quad\textbf{the branch on each floor (the derived 2):}
\]

%dsp-level:0
\[
\qquad \text{a floor's two members are distinguished by every finer ledger while sharing}\ m_{\mathrm{hid}}\ \text{and}\ \Pi\ [\mathrm{FT.44}(2)]:
\]

%dsp-level:0
\[
\mathcal F_2,\ \mathcal F_3:\ \textbf{the duality branch}\ —\ \text{the dual exchanges}\ (V,F)\ \text{at fixed}\ E\ [\mathrm{FT.44}(1)];
\]

%dsp-level:0
\[
\qquad \mathcal F_1:\ \textbf{the weight-ladder branch}\ —\ \text{the tetrahedron is self-dual}\ [\mathrm{FT.44}(3)],
\]

%dsp-level:0
\[
\ \text{so its branch index is the ladder rung}\ m\in\{1,2\}\ [\mathrm{FT.60}(5);\ \mathrm{Census\ FT.63}(6)];
\]

%dsp-level:0
\[
\textbf{the six states:}\qquad\mathcal S\ :=\ \{(\mathcal F_g,\beta):\ g\in\{1,2,3\},\ \beta\in\{-,+\}\},\qquad|\mathcal S|=3\cdot2=6\qquad[\text{machine}];
\]

%dsp-level:0
\[
(3)\quad\textbf{the defect spectrum (the branch mass):}
\]

%dsp-level:0
\[
\qquad \delta_v=2\pi-d\,\frac{(L-2)\pi}L,
\]

%dsp-level:0
\[
\qquad D_S=4\sin^2\frac{\delta_v}4,
\]

%dsp-level:0
\[
\qquad m^2_{\mathrm{pole}}=\hat\mu'^2\,D_S
\]

%dsp-level:0
\[
\qquad [\mathrm{Census\ FT.63}(4),(6);\ \mathrm{Def\text{-}H.37}(2);\ \mathrm{FT.54}(4)]:
\]

%tbl-level:0
{\setlength{\tabcolsep}{4pt}
\begin{longtable}{>{\raggedright\arraybackslash}p{\dimexpr0.0887\textwidth-2\tabcolsep\relax}|>{\raggedright\arraybackslash}p{\dimexpr0.2940\textwidth-2\tabcolsep\relax}|>{\raggedright\arraybackslash}p{\dimexpr0.1003\textwidth-2\tabcolsep\relax}|>{\raggedright\arraybackslash}p{\dimexpr0.1586\textwidth-2\tabcolsep\relax}|>{\raggedright\arraybackslash}p{\dimexpr0.1819\textwidth-2\tabcolsep\relax}|>{\raggedright\arraybackslash}p{\dimexpr0.1764\textwidth-2\tabcolsep\relax}}
\toprule
state & $(\mathcal F_g,\allowbreak \beta)$ & $\delta_v/\pi$ & $D_S$ & exact field & min. polynomial of $D_S$ \\
\addlinespace[2pt]
\midrule
\endfirsthead
\toprule
state & $(\mathcal F_g,\allowbreak \beta)$ & $\delta_v/\pi$ & $D_S$ & exact field & min. polynomial of $D_S$ \\
\addlinespace[2pt]
\midrule
\endhead
\bottomrule
\endlastfoot
$s_1$ & $(\mathcal F_1,\allowbreak -)$, the ladder rung $m{=}1$ & $1$ & $2$ & $\mathbb Q$ & $x-2$ \\
$s_2$ & $(\mathcal F_1,\allowbreak +\allowbreak )$, the ladder rung $m{=}2$ & $1$ & $4$ & $\mathbb Q$ & $x-4$ \\
$s_3$ & $(\mathcal F_2,\allowbreak -)$, the cube member & $\tfrac12$ & $2-\sqrt2$ & $\mathbb Q(\sqrt2)$ & $x^2-4x+\allowbreak 2$ \\
$s_4$ & $(\mathcal F_2,\allowbreak +\allowbreak )$, the octahedron member & $\tfrac23$ & $1$ & $\mathbb Q$ & $x-1$ \\
$s_5$ & $(\mathcal F_3,\allowbreak -)$, the dodecahedron member & $\tfrac15$ & $2-2\cos\tfrac\pi{10}$ & $\mathbb Q(\cos\tfrac\pi{10})$ & $x^4-8x^3+\allowbreak 19x^2-12x+\allowbreak 1$ \\
$s_6$ & $(\mathcal F_3,\allowbreak +\allowbreak )$, the icosahedron member & $\tfrac13$ & $2-\sqrt3$ & $\mathbb Q(\sqrt3)$ & $x^2-4x+\allowbreak 1$ \\
\end{longtable}
}

%dsp-level:0
\[
\textbf{Descartes/Gauss--Bonnet:}
\]

%dsp-level:0
\[
\qquad V\,\delta_v=4\pi\ \ \text{on all five solids}
\]

%dsp-level:0
\[
\qquad [\mathrm{Census\ FT.63}(4);\ \text{machine: exact in}\ \pi\text{-units}];
\]

%dsp-level:0
\[
(4)\quad\textbf{the branch order (derived, exact, uniform):}\qquad\boxed{\ D_S(\mathcal F_g,-)\ <\ D_S(\mathcal F_g,+)\qquad\forall g\in\{1,2,3\}\ }
\]

%dsp-level:0
\[
\mathcal F_1:\ 4>2;
\]

%dsp-level:0
\[
\qquad \mathcal F_2:\ 1>2-\sqrt2\iff\sqrt2>1\iff2>1;
\]

%dsp-level:0
\[
\qquad \mathcal F_3:\ 2-\sqrt3>2-2\cos\tfrac\pi{10}\iff\cos\tfrac\pi{10}>\tfrac{\sqrt3}2\iff\tfrac\pi{10}<\tfrac\pi6\iff3<5
\]

%dsp-level:0
\[
\qquad [\text{machine: integer and radical certificates}];
\]

%dsp-level:0
\[
\textbf{the six values are pairwise distinct}\ —\ \text{the defect species resolves all six states with zero degeneracy}\ [\text{machine: separation by field and by minimal polynomial}];
\]

%dsp-level:0
\[
(5)\quad\textbf{the branch splitting (the constituent mass difference):}
\]

%dsp-level:0
\[
\qquad \boxed{\ \Delta_g\ :=\ \hat\mu'^2\bigl(D_S(\mathcal F_g,+)-D_S(\mathcal F_g,-)\bigr)\ }
\]

%dsp-level:0
\[
\Delta_1=2\hat\mu'^2\in\mathbb Q\hat\mu'^2;
\]

%dsp-level:0
\[
\qquad \Delta_2=(\sqrt2-1)\hat\mu'^2\in\mathbb Q(\sqrt2)\hat\mu'^2;
\]

%dsp-level:0
\[
\qquad \Delta_3=\bigl(2\cos\tfrac\pi{10}-\sqrt3\bigr)\hat\mu'^2\in\mathbb Q\bigl(\sqrt3,\cos\tfrac\pi{10}\bigr)\hat\mu'^2;
\]

%dsp-level:0
\[
\Delta_g>0\ \ \forall g
\]

%dsp-level:0
\[
\qquad [(4)\text{'s certificates}];
\]

%dsp-level:0
\[
\qquad \textbf{three distinct fields, three exact positive splittings, no imported value}\ [\mathrm{P5}];
\]

%dsp-level:0
\[
(6)\quad\textbf{the generation map (the unique order-isomorphism):}
\]

%dsp-level:0
\[
\qquad \text{the}\ q\text{-channel's three lines are ordered exactly:}
\]

%dsp-level:0
\[
\quad (r^{(q)}_1)^2<(r^{(q)}_2)^2<(r^{(q)}_0)^2\ [\mathrm{FT.119}(5);\ \text{enclosure-proven, frozen}];
\]

%dsp-level:0
\[
\text{the three floors are ordered exactly:}\quad2E(\mathcal F_1)<2E(\mathcal F_2)<2E(\mathcal F_3)\ [(1)]:
\]

%dsp-level:0
\[
\boxed{\ \text{two exact three-element ordered sets admit exactly one order-isomorphism}\ }
\]

%dsp-level:0
\[
\qquad g=1\leftrightarrow(r^{(q)}_1)^2\leftrightarrow\mathcal F_1,
\]

%dsp-level:0
\[
\quad g=2\leftrightarrow(r^{(q)}_2)^2\leftrightarrow\mathcal F_2,
\]

%dsp-level:0
\[
\quad g=3\leftrightarrow(r^{(q)}_0)^2\leftrightarrow\mathcal F_3;
\]

%dsp-level:0
\[
\textbf{the species row:}
\]

%dsp-level:0
\[
\qquad \text{on the transport species the line order inverts}\ [\mathrm{BM.39}(5):\ \text{the spectral-maximal line receives the mass-minimum}]\implies\text{the isomorphism inverts with it};
\]

%dsp-level:0
\[
\qquad \text{the relating equation between the two currencies is}\ \neg\text{derived}\ [\mathrm{BM.35}(5)];
\]

%dsp-level:0
\[
\textbf{the scale is not supplied:}\qquad\hat\mu'^2\in\mathrm{Cal\text{-}F}\ (\text{logged}),\quad\mu\in\mathrm{Cal\text{-}F}:\qquad\text{the map is an order,}\ \neg\text{a number}\ [\mathrm{P4};\ \mathrm{P5}];
\]

%dsp-level:0
\[
(7)\quad\textbf{the charge alphabet of the six (the second derived 2):}
\]

%dsp-level:0
\[
\qquad \mathcal H_{\mathsf v}\cong\mathbb C^2\ [\mathrm{EW.32}];
\]

%dsp-level:0
\[
\qquad Q|_{\mathcal H_{\mathsf v}}=\operatorname{diag}(1,0),\quad\operatorname{spec}Q=\{0,+1\},\quad\omega:\ \operatorname{spec}\mapsto\{0,-1\}\ [\mathrm{EW.36}]:
\]

%dsp-level:0
\[
\text{the seed carrier's weight doublet}\ \{0,+1\}\ \text{together with its}\ \omega\text{-image is exactly the ternary end-read alphabet}\ \{-1,0,+1\}\ [\mathrm{EW.36};\ \mathrm{Def\text{-}BM.12}(1);\ \mathrm{D3.4}];
\]

%dsp-level:0
\[
\textbf{the branch index of}\ (2)\ \textbf{and the weight index of}\ \mathcal H_{\mathsf v}\ \textbf{are two derived twos};
\]

%dsp-level:0
\[
\qquad \text{their attachment is}\ (8)\text{'s clause};
\]

%dsp-level:0
\[
(8)\quad\textbf{the identification clause (the single declared row):}
\]

%dsp-level:0
\[
\qquad \textbf{down-type}\ :=\ \text{the floor's heavier defect branch}\ (\beta{=}+);
\]

%dsp-level:0
\[
\qquad \textbf{up-type}\ :=\ \text{the lighter}\ (\beta{=}−);
\]

%dsp-level:0
\[
\text{the ordering is derived}\ [(4)];
\]

%dsp-level:0
\[
\ \text{the name attachment is a label on a derived structural tuple and carries no datum}\ [\mathrm{Def\text{-}BM.3}\ \text{pattern};\ \mathrm{Def\text{-}BM.3}(7)\ \text{scope}];
\]

%dsp-level:0
\[
\qquad \textbf{declared}\ (\mathrm{A\text{-}BM.6})\ \textbf{before use};
\]

%dsp-level:0
\[
\textbf{non-derived:}\ \text{the attachment of the}\ Q\text{-weight index to the branch index}\ (7);
\]

%dsp-level:0
\[
\qquad \text{the flavour composition of a sealed composite}\ —\ \text{which three of the six its constituents occupy}\ [\mathrm{BM.28}(7);\ \mathrm{Def\text{-}BM.8}(5):\ \text{a fractional constituent alphabet is structurally impossible under an integer-valued additive read}];
\]

%dsp-level:0
\[
(9)\quad\textbf{status and scope:}
\]

%dsp-level:0
\[
\qquad \textbf{derived:}\ (1)\text{--}(7);
\]

%dsp-level:0
\[
\qquad \textbf{declared:}\ (8)\text{'s name attachment alone};
\]

%dsp-level:0
\[
\qquad \neg\text{a P4 addition},\ \neg\text{a new dimensionful scale},\ \text{§Cal-BM stays empty};
\]

%dsp-level:0
\[
\neg\exists\ \text{a measured mass, ratio or Yukawa value in any clause}\ [\mathrm{P5};\ \text{machine: zero external tokens, zero decimal numerals}];
\]

%dsp-level:0
\[
\qquad \text{species names are confined to the identification layers}\ [\mathrm{Def\text{-}BM.3}(7),\ \mathrm{Def\text{-}BM.6}(7)];
\]

%dsp-level:0
\[
\textbf{beyond this tier:}\ \text{the scale relating the defect and informational currencies}\;
\]

%dsp-level:0
\[
\qquad \text{the flavour composition of the baryon classes}\;
\]

%dsp-level:0
\[
\qquad \text{the materialization scope}\ \mathrm{M1}/\mathrm{M2}\ [\mathrm{Def\text{-}H.20}]\ \text{is inherited at every use}.
\]

\textbf{Rmk (six states, no data).} Nothing in this definition is fitted. The floors come from the card template's edge counts; the branches come from the duality that leaves those counts fixed and splits everything finer; the branch masses come from the conical defect at a vertex where three faces meet, closed by Descartes against four pi. What the assembly adds is one order-isomorphism between two three-element ordered sets the corpus had already derived separately, and one name. The result is six states, six pairwise distinct exact algebraic masses on the defect species, three exact rationals on the informational one, and a hierarchy that runs one way in the bit currency and the other way in the angle currency — with the equation between the two currencies still open, and still not guessed.

\pseudohead{theorem}{\texorpdfstring{BM.57}{BM.57}}{\texorpdfstring{The continuation-symmetry census: the vertex asymmetry is two, and it splits nothing}{The continuation-symmetry census: the vertex asymmetry is two, and it splits nothing}}\label{tag:BM-57}
\textbf{Thm.}

%dsp-level:0
\[
(1)\quad\textbf{the strong 3-vertex:}
\]

%dsp-level:0
\[
\qquad \text{the invariant tensor is}\ \varepsilon_{ijk}\ [\mathrm{CT3.17};\ \mathrm{D3.9}];
\]

%dsp-level:0
\[
\qquad \varepsilon_{\sigma(i)\sigma(j)\sigma(k)}=\operatorname{sgn}(\sigma)\,\varepsilon_{ijk}\ \ \forall\sigma\in S_3
\]

%dsp-level:0
\[
\qquad [\text{machine: all six permutations}]:
\]

%dsp-level:0
\[
\textbf{the continuation group:}
\]

%dsp-level:0
\[
\qquad \operatorname{Stab}_{S_3}(\varepsilon)=A_3\cong C_3=S_3\cap SU(3),
\]

%dsp-level:0
\[
\qquad |A_3|=3,
\]

%dsp-level:0
\[
\qquad \text{nontrivial count}\ 2
\]

%dsp-level:0
\[
\qquad [\mathrm{BM.3}(1);\ \text{machine: the invariant set is exactly the cyclic group}];
\]

%dsp-level:0
\[
(2)\quad\textbf{the}\ (1,1)\textbf{-vertex (the chain wall):}
\]

%dsp-level:0
\[
\qquad \text{the invariant tensor is}\ \delta^i{}_j\ [\mathrm{CT3.17}];
\]

%dsp-level:0
\[
\qquad \text{the type-preserving relabelling group of a}\ (1,1)\ \text{vertex is}\ S_1\times S_1:
\]

%dsp-level:0
\[
\qquad |\,\cdot\,|=1,
\]

%dsp-level:0
\[
\qquad \text{nontrivial count}\ 0;
\]

%dsp-level:0
\[
\textbf{the leg swap is not a continuation:}
\]

%dsp-level:0
\[
\qquad \delta^i{}_j
\]

%dsp-level:0
\[
= \delta_i{}^j\ \text{admits it as a bare 2-tensor}\ [\text{machine}],
\]

%dsp-level:0
\[
\ \text{but it acts on the oriented end registers as}\ (\sigma_s,\sigma_t)
\]

%dsp-level:0
\[
\mapsto (\sigma_t,\sigma_s):
\]

%dsp-level:0
\[
q_{\mathrm{EM}}\mapsto-q_{\mathrm{EM}}\ \ \text{on all four end pairs}
\]

%dsp-level:0
\[
\qquad [\mathrm{D3.4};\ \mathrm{EM.9a}:\ \text{charge-conjugation odd};\ \text{machine}]:
\]

%dsp-level:0
\[
\qquad \textbf{the swap is the}\ \omega\textbf{-involution, not a continuation symmetry};
\]

%dsp-level:0
\[
(3)\quad\textbf{the asymmetry:}\qquad\boxed{\ |A_3|-|S_1\times S_1|=3-1=2,\qquad2-0=2\ }\qquad[\text{machine}]:
\]

%dsp-level:0
\[
\textbf{the asymmetry is}\ 2\ \textbf{by both conventions}\ —\ \text{the order difference and the nontrivial-count difference agree};
\]

%dsp-level:0
\[
\qquad \text{the count is derived};
\]

%dsp-level:0
\[
(4)\quad\textbf{the first annihilation row (class-blindness):}
\]

%dsp-level:0
\[
\qquad \text{both baryon classes are}\ \varepsilon\text{-sealed 3-junctures with identical vertex data}\ [\mathrm{Def\text{-}BM.3}(2);\ \mathrm{BM.17};\ \mathrm{BM.26}]:
\]

%tbl-level:0
{\setlength{\tabcolsep}{4pt}
\begin{longtable}{>{\raggedright\arraybackslash}p{\dimexpr0.4919\textwidth-2\tabcolsep\relax}|>{\raggedright\arraybackslash}p{\dimexpr0.1757\textwidth-2\tabcolsep\relax}|>{\raggedright\arraybackslash}p{\dimexpr0.1757\textwidth-2\tabcolsep\relax}|>{\raggedright\arraybackslash}p{\dimexpr0.1568\textwidth-2\tabcolsep\relax}}
\toprule
vertex invariant & neutral class & unit class & difference \\
\addlinespace[2pt]
\midrule
\endfirsthead
\toprule
vertex invariant & neutral class & unit class & difference \\
\addlinespace[2pt]
\midrule
\endhead
\bottomrule
\endlastfoot
junction grammar & $\varepsilon_{ijk}$ 3-juncture & $\varepsilon_{ijk}$ 3-juncture & $0$ \\
vertex type $[\mathrm{D0.8}]$ & $(3,\allowbreak 0)$ & $(3,\allowbreak 0)$ & $0$ \\
sealing-phase class $[\mathrm{BM.2}(1);\ \mathrm{Census\ FT.49}(3)]$ & $4\mathbb Z$ & $4\mathbb Z$ & $0$ \\
$E_{\mathrm{int}}$ $[\mathrm{Census\ FT.35}(2)]$ & $6$ & $6$ & $0$ \\
continuation group & $A_3$, order $3$ & $A_3$, order $3$ & $0$ \\
sealing grammar count $[\mathrm{CT3.17a}(2)]$ & $g(3)=\allowbreak 1$ & $g(3)=\allowbreak 1$ & $0$ \\
\end{longtable}
}

%dsp-level:0
\[
\textbf{every vertex-structural quantity is a common factor on the two classes}\ [\text{machine: the difference table is identically zero}]\implies\textbf{no vertex datum can split the two rows};
\]

%dsp-level:0
\[
(5)\quad\textbf{the second annihilation row (the seed is colour-blind):}
\]

%dsp-level:0
\[
\qquad T_a\mathsf v=0\ [\mathrm{EW.32}];
\]

%dsp-level:0
\[
\qquad \mathcal G=(gW_1,gW_2,gW_3,g'Y)\ [\mathrm{EW.33}]:\ \neg\exists\ \text{a colour generator in the transport data};
\]

%dsp-level:0
\[
\operatorname{Re}\langle\mathcal G_A\mathsf v,\mathcal G_B\mathsf v\rangle=\frac14\begin{pmatrix}g^2&0&0&0\\0&g^2&0&0\\0&0&g^2&-gg'\\0&0&-gg'&g'^2\end{pmatrix}\qquad[\text{machine: Pauli spinors over}\ \mathbb Q(i)\otimes\mathbb Q[g,g']]:
\]

%dsp-level:0
\[
\textbf{the free-symbol set of the Gram is}\ \{g,g'\}:\ \text{zero colour symbols}\ [\text{machine}]\implies\boxed{\ \frac{\partial\,G(c)}{\partial\,(\text{any slot-index datum})}=0\ \ \text{exactly}\ }
\]

%dsp-level:0
\[
\textbf{the exact factorizations}\ [\text{machine: polynomial identities in}\ \mathbb Q[g,g',c_1,\dots,c_4]]:
\]

%dsp-level:0
\[
G(c)=\frac14\Bigl[g^2(c_1^2+c_2^2+c_3^2)-2gg'c_3c_4+g'^2c_4^2\Bigr]=\frac14\Bigl[g^2(c_1^2+c_2^2)+\bigl(gc_3-g'c_4\bigr)^2\Bigr];
\]

%dsp-level:0
\[
\ker M^2=\operatorname{span}\{(0,0,g',g)\}\ [\mathrm{EW.22}]:
\]

%dsp-level:0
\[
\qquad \frac{(g^2+g'^2)}4\|c\|^2-G(c)=\frac14\Bigl[g'^2(c_1^2+c_2^2)+\bigl(gc_4+g'c_3\bigr)^2\Bigr]\ge0;
\]

%dsp-level:0
\[
\det(xI-\mathrm{Gram})=\Bigl(x-\frac{g^2}4\Bigr)^2x\Bigl(x-\frac{g^2+g'^2}4\Bigr)\qquad[\text{machine: Faddeev--LeVerrier over}\ \mathbb Q[g,g',x]];
\]

%dsp-level:0
\[
(6)\quad\textbf{the third annihilation row (the trace law):}
\]

%dsp-level:0
\[
\qquad \text{every Hermitian slot operator on a sealed singlet reads only its trace:}\quad\mathfrak T(H)
\]

%dsp-level:0
\[
= 3a\ \ \text{for}\ H
\]

%dsp-level:0
\[
= aI_3+zP+z^*P^{\mathsf T}
\]

%dsp-level:0
\[
\qquad [\mathrm{BM.1}(2);\ \mathrm{FT.117}(1)]:
\]

%dsp-level:0
\[
\textbf{the trace's free-symbol set is}\ \{a\}:\ \text{the}\ C_3\text{-doublet datum}\ (B,\theta)\ \text{in}\ z
\]

%dsp-level:0
\[
= Be^{i\theta}\ \textbf{is annihilated}\ [\text{machine}];
\]

%dsp-level:0
\[
\qquad 1+\omega+\omega^2
\]

%dsp-level:0
\[
= 0\ \text{in}\ \mathbb Q[\omega]/(\omega^2{+}\omega{+}1):\ \text{the}\ C_3\ \textbf{average of the doublet datum is exactly zero}\ [\text{machine}];
\]

%dsp-level:0
\[
\textbf{the companion row:}
\]

%dsp-level:0
\[
\qquad \text{every}\ S_3\text{-covariant slot operator has charpoly}\ (x-a-3b)(x-a)^2\ [\mathrm{FT.116};\ \text{machine}]:\ \textbf{the doublet is degenerate at hadron level};
\]

%dsp-level:0
\[
(7)\quad\textbf{the multiplicity row:}
\]

%dsp-level:0
\[
\qquad \text{a continuation-symmetry count is a degeneracy multiplicity of the sealed state};
\]

%dsp-level:0
\[
\qquad \text{a uniform multiplicity factor cancels in the normalized read at every priced factor}\ \{1,2,3,15,20,\tfrac73\}\ [\mathrm{BM.28}(4)(\mathrm i);\ \mathrm{BM.30}(3);\ \text{machine}]:
\]

%dsp-level:0
\[
\text{and it multiplies the total uniformly}\ —\ \text{the same factor on both classes by}\ (4)\implies\textbf{it shifts a total and never splits a class};
\]

%dsp-level:0
\[
(8)\quad\textbf{the verdict:}\qquad\boxed{\ \Delta m_{\mathrm{vertex}}\ =\ 0\ \ \text{exactly}\ }\qquad[(4)\wedge(5)\wedge(6)\wedge(7);\ \text{machine}]:
\]

%dsp-level:0
\[
\textbf{the vertex-asymmetry hypothesis is refuted as a mass-splitting mechanism};
\]

%dsp-level:0
\[
\qquad \text{the asymmetry itself}\ (3)\ \textbf{is derived and stands};
\]

%dsp-level:0
\[
(9)\quad\textbf{the residual content (where the asymmetry does act):}
\]

%dsp-level:0
\[
\qquad A_3\ \text{is the forced family group of the seal}\ [\mathrm{BM.3}(1)];
\]

%dsp-level:0
\[
\qquad \text{its orbit on a transport profile has length}\ 3\ [\mathrm{Def\text{-}BM.11}(2),(3);\ \mathrm{FT.158}(1):\ U^3=E_0,\ U\ne E_0,\ U^2\ne E_0]:
\]

%dsp-level:0
\[
\textbf{the asymmetry generates the family triplet}\ —\ \text{three lines per channel}\ [\mathrm{BM.60}(1)];
\]

%dsp-level:0
\[
\qquad \textbf{it does not split the baryon classes};
\]

%dsp-level:0
\[
\qquad \text{the}\ (1,1)\text{-vertex's single continuation is the chain-wall telescoping}\ [\mathrm{EM.9b}]\ —\ \text{and its content is}\ (2)\text{'s parity law, consumed at}\ \mathrm{BM.59};
\]

%dsp-level:0
\[
(10)\quad\textbf{status:}
\]

%dsp-level:0
\[
\qquad \textbf{derived:}\ (1)\text{--}(9);
\]

%dsp-level:0
\[
\qquad \textbf{refuted:}\ \text{the vertex asymmetry as a splitting mechanism}\ (8);
\]

%dsp-level:0
\[
\qquad \neg\text{a P4 addition},\ \neg\text{a measured value}\ [\mathrm{P5}].
\]

\textbf{Req.} A7 (Part 1); D0.8, D2.6, D3.3, D3.4, D3.9 (Parts 1–2); CT0.9, CT3.17, CT3.17a(2), EW.15, EW.20, EW.22, EW.32, EW.33, EW.35, EM.9, EM.9a, EM.9b, EM.24(1) (Part 2); P5 (Parts 1–2); Def-H.37(1), FT.106(3a), FT.116, FT.117(1)–(3), FT.158(1) (Part 3); Census FT.35(2), FT.49(3) (row-wise); Def-BM.1, Def-BM.3(2), Def-BM.11(2), Def-BM.11(3), BM.1(2), BM.2(1), BM.3(1), BM.17, BM.26, BM.28(4)(i), BM.30(3), BM.36(4), BM.36(5) (this part); the exact Levi-Civita and Kronecker invariant censuses, Pauli/Gram algebra over $\mathbb Q(i)\otimes\mathbb Q[g,g']$, Faddeev–LeVerrier charpolys, the cyclotomic ring $\mathbb Q[\omega]/(\omega^2{+}\omega{+}1)$, sum-of-squares factorizations [math; machine: \texttt{checks/\allowbreak{}verify\_\allowbreak{}H\_\allowbreak{}constituent\_\allowbreak{}transport\_\allowbreak{}splitting.py} §A].

\textbf{Pf.}

%dsp-level:0
\[
(1):\ \varepsilon_{ijk}\ \text{built as the sign of the index triple};
\]

%dsp-level:0
\[
\ \text{all six}\ \sigma
\]

%dsp-level:0
\[
\in S_3\ \text{tested against}\ \varepsilon_{\sigma(i)\sigma(j)\sigma(k)}
\]

%dsp-level:0
\[
= \varepsilon_{ijk}:\ \text{the invariant set has three elements and equals}\ \{(012),(120),(201)\}\ [\text{machine}];
\]

%dsp-level:0
\[
\ \operatorname{sgn}(\sigma)
\]

%dsp-level:0
\[
= +1\ \text{on exactly those three}\ [\text{machine}]
\]

%dsp-level:0
\[
(2):\ \tau(v)
\]

%dsp-level:0
\[
= (1,1)\ [\mathrm{D0.8}]\implies\text{the type-preserving relabelling group is}\ S_1\times S_1;
\]

%dsp-level:0
\[
\ \delta^i{}_j\ \text{symmetric in its indices}\ [\text{machine}];
\]

%dsp-level:0
\[
\ \text{the swap row: all four end pairs mapped,}\ q
\]

%dsp-level:0
\[
\mapsto -q\ \text{at each}\ [\text{machine}];
\]

%dsp-level:0
\[
\ \mathrm{EM.9a}\text{'s charge-conjugation oddness re-cited}
\]

%dsp-level:0
\[
(3):\ 3-1=2,\quad2-0=2\ [\text{machine: integers}]
\]

%dsp-level:0
\[
(4):\ \text{each row re-cited from its own tag and compared across the two classes}\ [\text{machine: the difference table}];
\]

%dsp-level:0
\[
\ \mathrm{BM.28}(6)\text{'s identical set re-cited}
\]

%dsp-level:0
\[
(5):\ \mathsf v
\]

%dsp-level:0
\[
= (0,1)^{\mathsf T}:\ \sigma_1\mathsf v
\]

%dsp-level:0
\[
= (1,0)^{\mathsf T},
\]

%dsp-level:0
\[
\ \sigma_2\mathsf v
\]

%dsp-level:0
\[
= (-i,0)^{\mathsf T},
\]

%dsp-level:0
\[
\ \sigma_3\mathsf v
\]

%dsp-level:0
\[
= (0,-1)^{\mathsf T},
\]

%dsp-level:0
\[
\ \sigma_0\mathsf v
\]

%dsp-level:0
\[
= \mathsf v;
\]

%dsp-level:0
\[
\ W_3\mathsf v
\]

%dsp-level:0
\[
= -\tfrac12\mathsf v,
\]

%dsp-level:0
\[
\ Y\mathsf v
\]

%dsp-level:0
\[
= +\tfrac12\mathsf v,
\]

%dsp-level:0
\[
\ Q\mathsf v
\]

%dsp-level:0
\[
= 0\ [\text{machine}];
\]

%dsp-level:0
\[
\ \text{the Gram block and its antisymmetric imaginary part re-derived}\ [\text{machine}]
\]

%dsp-level:0
\[
\text{the kernel:}\ \mathrm{Gram}\cdot(0,0,g',g)^{\mathsf T}=0\ [\text{machine}];
\]

%dsp-level:0
\[
\qquad G(c)
\]

%dsp-level:0
\[
= \tfrac14[g^2(c_1^2+c_2^2)+(gc_3-g'c_4)^2]\ \text{and the ceiling's SOS form expanded and compared in}\ \mathbb Q[g,g',c_1,\dots,c_4]\ [\text{machine}]
\]

%dsp-level:0
\[
\det(xI-\mathrm{Gram})\ \text{by Faddeev--LeVerrier and compared to the factored form in}\ \mathbb Q[g,g',x]\ [\text{machine}];
\]

%dsp-level:0
\[
\qquad \text{the symbol census of every Gram entry}\ [\text{machine}]
\]

%dsp-level:0
\[
(6):\ H
\]

%dsp-level:0
\[
= aI_3+zP+z^*P^{\mathsf T}\ \text{with}\ P\ \text{the 3-cycle}:\ \operatorname{tr}P
\]

%dsp-level:0
\[
= \operatorname{tr}P^{\mathsf T}
\]

%dsp-level:0
\[
= 0\implies\mathfrak T(H)
\]

%dsp-level:0
\[
= 3a\ [\text{machine over}\ \mathbb Q[a,\operatorname{Re}z,\operatorname{Im}z]];
\]

%dsp-level:0
\[
\qquad \omega^3=1,\ 1+\omega+\omega^2=0\ \text{in}\ \mathbb Q[\omega]/(\omega^2{+}\omega{+}1)\ [\text{machine}]
\]

%dsp-level:0
\[
aI_3+bJ_3:\ \text{charpoly}\ (x-a-3b)(x-a)^2\ \text{by Faddeev--LeVerrier over}\ \mathbb Q[a,b,x]\ [\text{machine}]
\]

%dsp-level:0
\[
(7):\ \frac{\sum_Rf\,w\,\Delta S}{\sum_Rf\,w}
\]

%dsp-level:0
\[
= \frac{\sum_Rw\,\Delta S}{\sum_Rw}\ \text{at every uniform}\ f\in\{1,2,3,15,20,\tfrac73\}\ [\text{machine: exact Fractions}];
\]

%dsp-level:0
\[
\qquad (8):\ \text{conjunction of}\ (4),(5),(6),(7)\ [\text{machine}]
\]

%dsp-level:0
\[
(9):\ U^3
\]

%dsp-level:0
\[
= E_0,
\]

%dsp-level:0
\[
\ U
\]

%dsp-level:0
\[
\ne E_0,
\]

%dsp-level:0
\[
\ U^2
\]

%dsp-level:0
\[
\ne E_0\ \text{for the 3-cycle permutation matrix over}\ \mathbb Q(i):\ \text{the orbit length is exactly}\ 3\ [\text{machine}];
\]

%dsp-level:0
\[
\qquad (10):\ \text{assembly}. \ \blacksquare
\]

\textbf{Phys.} The hypothesis was that the strong vertex, having two continuation symmetries against the electromagnetic vertex's one, would split the quark masses. The count is right: the Levi-Civita tensor is preserved by exactly the three even permutations of its slots, so the strong 3-vertex carries two nontrivial continuations, while the chain wall's single in-leg and single out-leg admit no nontrivial relabelling at all — its only candidate, the leg swap, reverses the end registers and is therefore the charge-conjugation involution, not a symmetry of the transport. The asymmetry is two by either counting convention. And it splits nothing. Three independent theorems kill it: the two baryon classes are the same vertex with the same phase class and the same edge count, so any vertex invariant is a common factor; the transported seed is annihilated by every colour generator, so the mass form has no slot index to differentiate; and on a sealed singlet every slot operator reads only its trace, which annihilates the doublet datum the continuation group acts on. What the asymmetry does buy is the family: the same $C_3$ that cannot split the neutron from the proton is what cycles a transport profile into three lines.

\pseudohead{theorem}{\texorpdfstring{BM.58}{BM.58}}{\texorpdfstring{The observation frames: the metric is charge-blind, and the frame's content is the seal asymmetry}{The observation frames: the metric is charge-blind, and the frame's content is the seal asymmetry}}\label{tag:BM-58}
\textbf{Thm.}

%dsp-level:0
\[
(1)\quad\textbf{the metric data carry no charge argument:}
\]

%dsp-level:0
\[
\qquad \Omega=e^{\alpha_\Omega a_{00}},\quad g_{00}=-\Omega^2,\quad g_{0i}=\Omega^2\varkappa\beta w_i,\quad g_{ij}=\Omega^2(\delta_{ij}+\varkappa q_{ij})
\]

%dsp-level:0
\[
\qquad [\mathrm{CT1.12};\ \mathrm{FT.157}(1)]:
\]

%dsp-level:0
\[
N^2=-\frac{\det g}{\det\gamma}>0,
\]

%dsp-level:0
\[
\qquad N_i=g_{0i}=\Omega^2\varkappa\beta w_i,
\]

%dsp-level:0
\[
\qquad N^i=\gamma^{ij}N_j,
\]

%dsp-level:0
\[
\qquad ds^2=-N^2d\tau^2+g_{ij}(dx^i+N^id\tau)(dx^j+N^jd\tau)
\]

%dsp-level:0
\[
\qquad [\mathrm{Def\text{-}BM.10}(2);\ \text{machine}]:
\]

%dsp-level:0
\[
\textbf{the token census on the canonical defining blocks:}
\]

%dsp-level:0
\[
\qquad \text{charge-token occurrences inside}\ \mathrm{Def\text{-}BM.10}(2),(3),(5)\ \text{and}\ \mathrm{CT1.12}:\ \ \mathbf 0
\]

%dsp-level:0
\[
\qquad [\text{machine}];
\]

%dsp-level:0
\[
(2)\quad\textbf{the seven geometric monomials are charge-blind:}
\]

%dsp-level:0
\[
\qquad \Delta x_{\mathrm B}=\mathfrak h E^{-1}(\ln2)^{-1},\quad V_{\mathrm{eff}}=\mathfrak h^3E^{-3}(\ln2)^{-3},\quad\rho=\mathfrak h^{-4}E^4(\ln2)^4,
\]

%dsp-level:0
\[
\kappa^2\rho=2\pi\,\mathfrak h^{-2}E^4(\ln2)^3,
\]

%dsp-level:0
\[
\qquad r_s=\tfrac12\mathfrak h E,
\]

%dsp-level:0
\[
\qquad (\Delta x_{\mathrm B})^2\kappa^2\rho=2\pi E^2\ln2,
\]

%dsp-level:0
\[
\qquad \Delta x_{\mathrm B}\,r_s=\frac{\mathfrak h^2}{2\ln2}
\]

%dsp-level:0
\[
\qquad [\mathrm{BM.13}(1),\ \mathrm{BM.14}(1)\text{--}(3),\ \mathrm{BM.15}(4),\ \mathrm{BM.16}(1),(4)]:
\]

%dsp-level:0
\[
\textbf{the charge exponent is}\ 0\ \textbf{in all seven}\ [\text{machine: the Laurent-monomial census}];
\]

%dsp-level:0
\[
\qquad \textbf{the lock identities hold as monomial identities}\ [\text{machine}];
\]

%dsp-level:0
\[
\textbf{at}\ E_{\mathrm{int}}=6:
\]

%dsp-level:0
\[
\qquad V_{\mathrm{eff}}=\frac{\mathfrak h^3}{216(\ln2)^3},\quad\rho=\frac{1296(\ln2)^4}{\mathfrak h^4},\quad\kappa^2\rho=\frac{2592\pi(\ln2)^3}{\mathfrak h^2},
\]

%dsp-level:0
\[
\quad r_s=3\mathfrak h,
\]

%dsp-level:0
\[
\quad (\Delta x_{\mathrm B})^2\kappa^2\rho=72\pi\ln2
\]

%dsp-level:0
\[
\qquad [\mathrm{BM.43}(1);\ \text{machine: every coefficient exact}];
\]

%dsp-level:0
\[
(3)\quad\textbf{the frames are masks, not metrics:}
\]

%dsp-level:0
\[
\qquad \textbf{the colour frame:}\ \mathcal C(\Sigma;C)\in\{0,1\},
\]

%dsp-level:0
\[
\ \text{an idempotent, and one dangling attachment zeroes the whole colour sector}\ [\mathrm{D2.14};\ \mathrm{FT.19};\ \mathrm{BM.1}(1);\ \mathrm{BM.32}(1)]\ —\ \text{the only idempotents of}\ \{0,\dots,7\}\ \text{are}\ \{0,1\}\ [\text{machine}];
\]

%dsp-level:0
\[
\textbf{the EM frame:}
\]

%dsp-level:0
\[
\qquad W(Q)>0\ \ \text{at every local charge class}\ Q\in\{-3,\dots,+3\}\ [\mathrm{Def\text{-}BM.8}(6);\ \text{machine}]:
\]

%dsp-level:0
\[
\qquad \neg\exists\ \text{a local projector on}\ Q;
\]

%dsp-level:0
\[
\quad\text{neutrality is exhausted only at universe scale}\ [\mathrm{EM.11};\ \mathrm{BM.47}(4),(5)]:
\]

%dsp-level:0
\[
\textbf{a mask acts on operators and on the probability space, not on the metric}\ [\mathrm{D2.11};\ \mathrm{Def\text{-}BM.10}(2)\text{'s alphabet},\ (1)]:
\]

%dsp-level:0
\[
\qquad \text{the emergent metric is the ambient ledger's coarse-grained sector coefficients, and no restriction re-reads them};
\]

%dsp-level:0
\[
(4)\quad\textbf{the metric correction vanishes:}
\]

%dsp-level:0
\[
\qquad \boxed{\ \Delta m_{\mathrm{metric}}\ =\ 0\ \ \text{exactly}\ }
\]

%dsp-level:0
\[
\qquad [(1)\wedge(2)\wedge(3);\ \text{machine: the class difference of each of the seven monomials is}\ 0]:
\]

%dsp-level:0
\[
\textbf{the observation-frame hypothesis is refuted as a mass-splitting mechanism};
\]

%dsp-level:0
\[
\qquad \text{the parent and the daughter carry identical curvature rows}\ [\mathrm{BM.43}(1):\ \text{re-cited}];
\]

%dsp-level:0
\[
(5)\quad\textbf{what the frame asymmetry does carry:}
\]

%dsp-level:0
\[
\qquad \text{colour self-seals locally; EM seals at universe scale}\ [\mathrm{BM.47}(5)];
\]

%dsp-level:0
\[
\qquad \text{the surviving mass content is the scar's field energy}
\]

%dsp-level:0
\[
E^{\mathrm{self}}_{\mathrm{EM}}=\frac{c\,e^2q^2}{4\pi\,\Delta x_{\mathrm B}}=c\,\varepsilon_{\mathrm{EM}}\,q^2\,m_{\mathrm{hid}},
\]

%dsp-level:0
\[
\qquad \Delta m_{\mathrm{EM}}=\frac{3c\,e^2\ln2}{2\pi\,\mathfrak h}>0,
\]

%dsp-level:0
\[
\qquad \textbf{the charged member the heavier one}
\]

%dsp-level:0
\[
\qquad [\mathrm{BM.46}(1),(4);\ \mathrm{Def\text{-}BM.12}(5)]:
\]

%dsp-level:0
\[
\text{the monopole selector is}\ Q^2
\]

%dsp-level:0
\[
= (0,1,1)\ \text{at}\ Q
\]

%dsp-level:0
\[
= (0,\pm1)\ [\text{machine}]\implies\boxed{\ \text{the frame's contribution to}\ m_n-m_p\ \text{is}\ -\,c\,\varepsilon_{\mathrm{EM}}\,m_{\mathrm{hid}}\ <\ 0\ }
\]

%dsp-level:0
\[
\qquad [\text{machine: both priced}\ c]:
\]

%dsp-level:0
\[
\textbf{the frame asymmetry works against the neutral class's excess}\ —\ \text{so the missing piece must be class-dependent, of the opposite sign, and not electromagnetic};
\]

%dsp-level:0
\[
(6)\quad\textbf{the uniqueness row (the force direction's output):}
\]

%dsp-level:0
\[
\qquad \mathrm{BM.52}(3)\text{'s sector census is complete on the functionals it enumerates}\ (\Delta S,\ E_{\mathrm{int}},\ \text{the Janus balance},\ \text{the geometric monomials},\ |w|^2,\ \text{the Casimir},\ \text{the locks})\ \text{and gives five zeros and one monomial}\ [\mathrm{BM.52}(4)];
\]

%dsp-level:0
\[
\qquad \textbf{the complete inventory of class-dependent derived data at a sealed 3-juncture, extended by the register functionals the census does not enumerate, is}
\]

%tbl-level:2
{\footnotesize\setlength{\tabcolsep}{4pt}
\begin{longtable}{>{\raggedright\arraybackslash}p{\dimexpr0.3065\textwidth-2\tabcolsep\relax}|>{\raggedright\arraybackslash}p{\dimexpr0.1371\textwidth-2\tabcolsep\relax}|>{\raggedright\arraybackslash}p{\dimexpr0.3790\textwidth-2\tabcolsep\relax}|>{\raggedright\arraybackslash}p{\dimexpr0.1774\textwidth-2\tabcolsep\relax}}
\toprule
datum & species & status & sign for $m_n-m_p$ \\
\addlinespace[2pt]
\midrule
\endfirsthead
\toprule
datum & species & status & sign for $m_n-m_p$ \\
\addlinespace[2pt]
\midrule
\endhead
\bottomrule
\endlastfoot
the vertex invariants & slot / grammar & class-blind $[\mathrm{BM.57}(4)]$ & $0$ \\
the colour data & slot & annihilated by the trace law $[\mathrm{BM.1}(2);\ \mathrm{FT.117}]$ & $0$ \\
the metric monomials & chart / metric & charge exponent $0$ $[(2)]$ & $0$ \\
$F_{LR}=\allowbreak qI_2$, the chiral splitting $2\|q\|$ & operator & trace-free within a class: $\sum_{s,t}tq=\allowbreak 0$ $[\mathrm{BM.22}(2)\text{--}(4)]$ & $0$ on a trace read $[\mathrm{BM.1}(2)]$ \\
$Z^{\mathrm{amp}}(0,\allowbreak S)/Z^{\mathrm{amp}}(\pm1,\allowbreak S)=\allowbreak \tfrac{58}{13},\allowbreak \tfrac{2398}{401}$ & amplitude & barred: amplitude $\ne$ mass $[\mathrm{FT.79}(5);\ \mathrm{BM.37}(3)(\mathrm b)]$ & barred \\
$E^{\mathrm{self}}_{\mathrm{EM}}\propto Q^2$ & chart / field & adopted $[\mathrm{Def\text{-}BM.12}(6)]$ & negative $[(5)]$ \\
$W(Q)=\allowbreak (1,\allowbreak 6,\allowbreak 15,\allowbreak 20,\allowbreak 15,\allowbreak 6,\allowbreak 1)$ & transport register & derived, unbarred $[\mathrm{Def\text{-}BM.8}(6);\ \mathrm{FT.106}(3\mathrm a)]$ & \textbf{positive} $[\mathrm{BM.59}]$ \\
\end{longtable}
}

%dsp-level:0
\[
\text{the six sector rows of}\ \mathrm{BM.52}(3)\ \text{and the Lyapunov row of}\ \mathrm{BM.51}(6)\ \text{are re-cited inside this table as its first four entries};
\]

%dsp-level:0
\[
\qquad \textbf{exactly one candidate survives with the required sign};
\]

%dsp-level:0
\[
[\text{machine: the trace row}\ \textstyle\sum_{s,t\in\{\pm1\}}tq=0;\ \text{the row census exhaustive against the tier's own inventory scan of}\ \mathrm{BM.22},\ \mathrm{BM.28}(5),(6),\ \mathrm{BM.30}(1),\ \mathrm{BM.46}(2),(4),\ \mathrm{BM.51}(6),\ \mathrm{BM.52}(3),\ \mathrm{Def\text{-}BM.8}(6)];
\]

%dsp-level:0
\[
(7)\quad\textbf{status:}
\]

%dsp-level:0
\[
\qquad \textbf{derived:}\ (1)\text{--}(6);
\]

%dsp-level:0
\[
\qquad \textbf{refuted:}\ \text{the observation frame as a metric-splitting mechanism}\ (4);
\]

%dsp-level:0
\[
\qquad \neg\text{a P4 addition},\ \neg\text{a measured range, radius or flux}\ [\mathrm{P5}].
\]

\textbf{Req.} D1.3, D2.1, D2.9, D2.11, D2.14 (Parts 1–2); CT1.12, CT1.22, CT2.6, EM.11, EW.22 (Part 2); P5 (Parts 1–2); Def-H.32(2), FT.19, FT.79(5), FT.106(3), FT.117(1)–(3), FT.157(1) (Part 3); Def-BM.4, Def-BM.8(6), Def-BM.10(2), Def-BM.10(3), Def-BM.10(5), Def-BM.12(5), Def-BM.12(6), BM.1(1), BM.1(2), BM.13(1), BM.14(1)–(3), BM.15(4), BM.16(1), BM.16(4), BM.22(2)–(4), BM.28(1), BM.32(1), BM.37(3), BM.43(1), BM.46(1), BM.46(4), BM.47(4), BM.47(5) (this part); the ADM split by Schur complement, Laurent-monomial exponent arithmetic, idempotent censuses, token censuses [math; machine: \texttt{checks/\allowbreak{}verify\_\allowbreak{}H\_\allowbreak{}constituent\_\allowbreak{}transport\_\allowbreak{}splitting.py} §B].

\textbf{Pf.}

%dsp-level:0
\[
(1):\ \mathrm{Def\text{-}BM.10}(2)\text{'s split re-derived as linear algebra on}\ \mathrm{CT1.12}\text{'s dictionary};
\]

%dsp-level:0
\[
\ N^2
\]

%dsp-level:0
\[
= -\det g/\det\gamma\ \text{by Schur, denominators cleared}\ [\text{machine}];
\]

%dsp-level:0
\[
\ \text{the line-element identity expanded and compared}\ [\text{machine}];
\]

%dsp-level:0
\[
\ \text{the token census on the two defining blocks}\ [\text{machine}]
\]

%dsp-level:0
\[
(2):\ \text{each monomial built from its own tag's formula as a Laurent monomial in}\ (\mathfrak h,E,\ln2,\pi,Q);
\]

%dsp-level:0
\[
\ \text{the}\ Q\text{-exponent censused}\ [\text{machine}];
\]

%dsp-level:0
\[
\ (\Delta x_{\mathrm B})^2\kappa^2\rho\ \text{and}\ \Delta x_{\mathrm B}r_s\ \text{multiplied out and compared}\ [\text{machine}];
\]

%dsp-level:0
\[
\ E
\]

%dsp-level:0
\[
= 6\ \text{substituted:}\ 6^3
\]

%dsp-level:0
\[
= 216,
\]

%dsp-level:0
\[
\ 6^4
\]

%dsp-level:0
\[
= 1296,
\]

%dsp-level:0
\[
\ 2\pi\cdot1296
\]

%dsp-level:0
\[
= 2592\pi,
\]

%dsp-level:0
\[
\ \tfrac12\cdot6
\]

%dsp-level:0
\[
= 3,
\]

%dsp-level:0
\[
\ 2\pi\cdot36
\]

%dsp-level:0
\[
= 72\pi\ [\text{machine}]
\]

%dsp-level:0
\[
(3):\ x^2=x\ \text{over}\ \{0,\dots,7\}\ \text{gives}\ \{0,1\}\ [\text{machine: exhaustive}];
\]

%dsp-level:0
\[
\qquad W(Q)\ \text{re-enumerated and positive at all seven classes}\ [\text{machine}]
\]

%dsp-level:0
\[
(4):\ \text{the class difference of each monomial of}\ (2)\ \text{is}\ 0\ \text{because the}\ Q\text{-exponent is}\ 0\ [\text{machine}]
\]

%dsp-level:0
\[
(5):\ \mathrm{BM.46}(1),
\]

%dsp-level:0
\[
(4)\ \text{re-cited};
\]

%dsp-level:0
\[
\ Q^2\ \text{at}\ (0,\pm1)
\]

%dsp-level:0
\[
= (0,1,1)\ [\text{machine}];
\]

%dsp-level:0
\[
\ \text{the sign of}\ -(1-0)\,
\]

%dsp-level:0
\[
c\,
\]

%dsp-level:0
\[
\varepsilon_{\mathrm{EM}}\,
\]

%dsp-level:0
\[
m_{\mathrm{hid}}\ \text{at}\ c\in\{\tfrac12,\tfrac35\},
\]

%dsp-level:0
\[
\ \varepsilon_{\mathrm{EM}}>0,
\]

%dsp-level:0
\[
\ m_{\mathrm{hid}}>0\ [\text{machine}]
\]

%dsp-level:0
\[
(6):\ \text{the inventory scan of the class-dependent rows of}\ \mathrm{BM.22},
\]

%dsp-level:0
\[
\ \mathrm{BM.28}(5),
\]

%dsp-level:0
\[
(6),
\]

%dsp-level:0
\[
\ \mathrm{BM.30}(1),
\]

%dsp-level:0
\[
\ \mathrm{BM.46}(2),
\]

%dsp-level:0
\[
(4),
\]

%dsp-level:0
\[
\ \mathrm{Def\text{-}BM.8}(6)\ [\text{machine: the row set is exhaustive}];
\]

%dsp-level:0
\[
\qquad \text{the force-matrix trace:}\ \operatorname{spec}
\]

%dsp-level:0
\[
= \{\tfrac12(c_0+s|\mathbf c|)+tq\},
\]

%dsp-level:0
\[
\ \sum_{s,t}[\tfrac12(c_0+s|\mathbf c|)+tq]
\]

%dsp-level:0
\[
= 2c_0\ [\mathrm{BM.22}(3);\ \text{machine}]\implies\text{the trace read is charge-blind}\ [\mathrm{BM.1}(2)]
\]

%dsp-level:0
\[
(7):\ \text{assembly}. \ \blacksquare
\]

\textbf{Phys.} The intuition was that seeing the world through photons is not seeing it through colour, and that the difference should show up in the geometry. It does not. The emergent metric is built from the ambient ledger's coarse-grained sector coefficients, and not one of its seven derived monomials carries a charge exponent: the width, the volume, the density, the curvature scalar, the Schwarzschild-type read and both locks come out identical on the two classes to the last Laurent monomial. What a frame really is, in this corpus, is a mask — and the two masks have different grammars. Colour's is an idempotent that annihilates a sector at the first dangling attachment; electromagnetism's is a census that lets every local charge class through and cancels only once, at the scale of the whole universe. Masks filter operators and probability spaces. They do not re-read the metric. So the metric correction is exactly zero, and the frame's one surviving mass contribution is the scar's own field energy — positive, charge-selected, and on the wrong side: it makes the charged member heavier. That is the useful output. It narrows the search to a single column of the inventory, and the column has one entry.

\begin{theorem}[{\texorpdfstring{BM.59 — The constituent transport splitting: the order residue, the two directions, and the exact positive sign}{BM.59 - The constituent transport splitting: the order residue, the two directions, and the exact positive sign}}]\label{tag:BM-59}
\textbf{Thm.}

%dsp-level:0
\[
(1)\quad\textbf{the missing piece, stated:}
\]

%dsp-level:0
\[
\qquad \mathrm{BM.55}(1)\text{'s five-term decomposition is derived term by term:}
\]

%dsp-level:0
\[
\qquad \Delta m_{\mathrm{Lyapunov}}
\]

%dsp-level:0
\[
= 0\ [\mathrm{BM.51}(6)];
\]

%dsp-level:0
\[
\quad\Delta m_{\mathrm{metric}}
\]

%dsp-level:0
\[
= 0\ [\mathrm{BM.53}(2)];
\]

%dsp-level:0
\[
\quad\Delta m_{\mathrm{binding}}
\]

%dsp-level:0
\[
= 0\ [\mathrm{BM.54}(2),(5)];
\]

%dsp-level:0
\[
\quad\Delta m_{\mathrm{force}}
\]

%dsp-level:0
\[
= \Delta m_{\mathrm{self\text{-}int}}
\]

%dsp-level:0
\[
= -\Delta m_{\mathrm{EM}}\ [\mathrm{BM.52}(4),(5)]:
\]

%dsp-level:0
\[
\Delta m=m_n-m_p=-\,\frac{3\,c\,e^2\ln2}{2\pi\,\mathfrak h}<0,\qquad\frac{m_n}{m_p}=\frac{4\pi}{4\pi+c\,e^2}\in(0,1)\qquad[\mathrm{BM.55}(1)\text{--}(3);\ \text{frozen}]:
\]

%dsp-level:0
\[
\textbf{the order residue}\ [\mathrm{BM.55}(4):\ \text{named, not filled}]\ :
\]

%dsp-level:0
\[
= \ \text{a class-dependent constituent transport splitting with}\ m^{(n)}_{\mathrm{constituent}}>m^{(p)}_{\mathrm{constituent}};
\]

%dsp-level:0
\[
\boxed{\ \textbf{the residue is fillable, and this tier fills it:}\qquad\Delta m_{\mathrm{transport}}=\frac{\Sigma_{\mathrm{cl}}(0)-\Sigma_{\mathrm{cl}}(+1)}{\mathfrak h}=\frac{\ln(4/3)}{\mathfrak h}>0\ }
\]

%dsp-level:0
\[
\qquad [\mathrm{Def\text{-}BM.13}(4);\ \text{machine}];
\]

%dsp-level:0
\[
(1\mathrm a)\quad\textbf{the four grounds of}\ \mathrm{BM.55}(4)\ \textbf{and how this object evades each:}\qquad\begin{array}{l|l}\text{the ground on record} & \text{the evasion}\\\hline\text{the fractional constituent alphabet is unreachable}\ [\mathrm{BM.28}(7)] & \Sigma_{\mathrm{cl}}\ \text{is built from the integer alphabet}\ \{-1,0,+1\}\ \text{alone}\ [\mathrm{Def\text{-}BM.13}(1),(3)];\ \neg\text{a fraction appears}\\\text{the species labels carry no datum}\ [\mathrm{Def\text{-}BM.3}] & \Sigma_{\mathrm{cl}}\ \text{carries no label}: it is a census of the derived read\ [\mathrm{Def\text{-}BM.8}(6)]\\\text{the deficit/inertial conversions are item-7 content} & \Sigma_{\mathrm{cl}}\ \text{is already in the informational currency}\ —\ \text{nats of binary registration}\ [\mathrm{Def\text{-}BM.13}(6)];\ \neg\text{a conversion is used}\\\text{no clause attaches constituent profiles to the charge classes} & \Sigma_{\mathrm{cl}}\ \text{is}\ \neg\text{a profile:}\ m_q\ \text{is attached to each constituent by}\ q_{\mathrm{core}}\ \text{itself}\ [\mathrm{Def\text{-}BM.8}(1);\ \mathrm{EM.9}]\end{array}
\]

%dsp-level:0
\[
\textbf{the residue is filled by the transport register, not by the transport profile:}
\]

%dsp-level:0
\[
\qquad \mathrm{Def\text{-}BM.11}\text{'s family-profile map remains channel-line data and stays unattached to the baryon rows}\ [\mathrm{BM.55}(4)\text{'s fourth ground stands untouched}];
\]

%dsp-level:0
\[
(2)\quad\textbf{the mass direction (what the spectrum requires):}
\]

%dsp-level:0
\[
\qquad \text{a term}\ \delta m(Q)\ \text{that reverses the order must satisfy four exact bars}
\]

%dsp-level:0
\[
(\mathrm i)\ \ \text{class-dependent:}\quad\delta m(0)\ne\delta m(+1);
\]

%dsp-level:0
\[
\qquad (\mathrm{ii})\ \ \text{positive for the neutral class:}\quad\delta m(0)-\delta m(+1)>0;
\]

%dsp-level:0
\[
(\mathrm{iii})\ \ \text{currency-matched:}\quad[\delta m]=\mathfrak q^{+1/2}\ \text{through the rest slot}\ [\mathrm{Def\text{-}BM.10}(4);\ \mathrm{CT2.1}];
\]

%dsp-level:0
\[
\qquad (\mathrm{iv})\ \ \text{alphabet-disjoint from the count:}\ \neg\text{a function of}\ E_{\mathrm{int}}\ \text{alone}\ [\mathrm{BM.28}(1)\ \text{scope}];
\]

%dsp-level:0
\[
\textbf{and it must not be a barred route}\ [\mathrm{BM.37}(6)];
\]

%dsp-level:0
\[
\qquad [\text{machine: the inventory scan of}\ \mathrm{BM.58}(6)\text{'s table against}\ (\mathrm i)\text{--}(\mathrm{iv})]:
\]

%dsp-level:0
\[
\boxed{\ \textbf{the unique survivor is}\ W(Q)\ \textbf{read through its logarithm}\ }\qquad[\mathrm{Def\text{-}BM.13}(4),(6)];
\]

%dsp-level:0
\[
(3)\quad\textbf{the force direction (what the seal permits):}
\]

%dsp-level:0
\[
\qquad \text{colour is a local idempotent and its slot data are annihilated on a sealed singlet}\ [\mathrm{BM.57}(5),(6);\ \mathrm{BM.58}(3)];
\]

%dsp-level:0
\[
\qquad \text{EM has no local projector and admits every class with positive census weight}\ [\mathrm{BM.58}(3)]:
\]

%dsp-level:0
\[
\textbf{the only class-dependent datum a sealed junction can carry is the EM end-register census}\ \ W(Q)
\]

%dsp-level:0
\[
= [z^Q](2+z+z^{-1})^3
\]

%dsp-level:0
\[
\qquad [\mathrm{Def\text{-}BM.8}(6);\ \mathrm{BM.28}(2)];
\]

%dsp-level:0
\[
(4)\quad\textbf{convergence:}
\]

%dsp-level:0
\[
\qquad \text{the two directions land on the same object};
\]

%dsp-level:0
\[
\qquad \boxed{\ \Delta m_{\mathrm{transport}}\ :=\ \frac{\Sigma_{\mathrm{cl}}(0)-\Sigma_{\mathrm{cl}}(+1)}{\mathfrak h}\ =\ \frac{\ln(4/3)}{\mathfrak h}\ =\ \frac{2\ln2-\ln3}{\mathfrak h}\ >\ 0\ }
\]

%dsp-level:0
\[
\qquad [\text{machine: the certificate}\ 4>3]:
\]

%dsp-level:0
\[
[\Delta m_{\mathrm{transport}}]=\mathfrak q^{+1/2}=[\text{energy}]
\]

%dsp-level:0
\[
\qquad [\ln W\in\mathfrak q^0,\ [\mathfrak h]=\mathfrak q^{-1/2};\ \mathrm{CT2.1};\ \mathrm{D2.9};\ \text{machine: exponent arithmetic}];
\]

%dsp-level:0
\[
(5)\quad\textbf{the constituent row:}
\]

%dsp-level:0
\[
\qquad \text{the per-constituent transport entropies are}\ \ln m_q
\]

%dsp-level:0
\[
= (\ln2,\,0,\,0)\ \text{at}\ q
\]

%dsp-level:0
\[
= (0,+1,-1)\ [\mathrm{Def\text{-}BM.13}(1)]:
\]

%dsp-level:0
\[
\boxed{\ m(q_{\mathrm{core}}{=}0)-m(q_{\mathrm{core}}{=}\pm1)\ =\ \frac{\ln2}{\mathfrak h}\ >\ 0\ }
\]

%dsp-level:0
\[
\qquad [\text{machine:}\ 2>1]:
\]

%dsp-level:0
\[
\qquad \textbf{the neutral-core constituent is the heavier constituent};
\]

%dsp-level:0
\[
\textbf{the composite's splitting is not three times the constituent's:}
\]

%dsp-level:0
\[
\qquad 0<\ln\tfrac43<\ln2<3\ln2
\]

%dsp-level:0
\[
\qquad [\text{machine:}\ 3<4<6<24]:
\]

%dsp-level:0
\[
\text{the additive charge constraint couples the three registers, and the coupling strictly attenuates the constituent hierarchy};
\]

%dsp-level:0
\[
(6)\quad\textbf{the parity lock (why the register is a transport datum and not a label):}
\]

%dsp-level:0
\[
\qquad \sum_if_{v_i}=-q,
\]

%dsp-level:0
\[
\qquad \sum_i\aleph_{v_i}\equiv|q|\ (\mathrm{mod}\ 2),
\]

%dsp-level:0
\[
\qquad \sum_i\aleph_{v_i}\ge|q|
\]

%dsp-level:0
\[
\qquad [\mathrm{EM.9b};\ \mathrm{D3.3};\ \mathrm{Def\text{-}BM.13}(2)]:
\]

%dsp-level:0
\[
\textbf{the interior wall content of a constituent's transport is parity-locked to its charge core};
\]

%dsp-level:0
\[
\qquad \text{the neutral core is wall-balanced}\ (\min\textstyle\sum\aleph=0);
\]

%dsp-level:0
\[
\ \text{the charged core is wall-unbalanced}\ (\min\textstyle\sum\aleph=1)\ [\text{machine: exhaustive over every sign sequence,}\ k=2,\dots,7]:
\]

%dsp-level:0
\[
\textbf{the two neutral transport realizations are the two wall-free orientations};
\]

%dsp-level:0
\[
\qquad \text{the charged core admits none}:
\]

%dsp-level:0
\[
\text{the multiplicity}\ m_q\ \textbf{is a property of the transport, not of a naming convention}\ [\mathrm{FT.106}(3\mathrm a):\ \text{it is the norm layer's own closed form}];
\]

%dsp-level:0
\[
(7)\quad\textbf{the sign is fixed by the reference, not chosen:}
\]

%dsp-level:0
\[
\qquad \Delta S:=S-S(\mathrm{vac})\ \text{with}\ S(\mathrm{vac})=0\ [\mathrm{BM.35}(3);\ \mathrm{CT0.2}]\implies\Delta\Sigma_{\mathrm{cl}}(Q)=\ln W(Q)-\ln W(\mathrm{vac})=\ln W(Q),
\]

%dsp-level:0
\[
\qquad W(\mathrm{vac})=1;
\]

%dsp-level:0
\[
\textbf{the two competing references:}\qquad\begin{array}{l|l|l}\text{reference} & \text{class difference} & \text{status}\\\hline\text{the vacuum}\ (W{=}1) & +\ln\tfrac43 & \textbf{adopted: the corpus's own deficit rule}\ [\mathrm{BM.35}(3)]\\[1mm]\text{the unconstrained register}\ (4^3{=}64) & -\ln\tfrac43 & \text{barred: }\neg\text{a state of the composite}\ [\mathrm{Def\text{-}BM.5}(1)];\ \text{class-dependent reference}\\\text{Ansatz-B}\ (-\alpha\ln\Omega) & -\alpha\ln\tfrac43 & \text{barred: an external ansatz,}\ \alpha\ \text{undeclared}\ [\mathrm{Def\text{-}H.28}]\end{array}
\]

%dsp-level:0
\[
\textbf{exactly one reference is a state of the composite}\ [\text{machine: the admissible-reference census}];
\]

%dsp-level:0
\[
(8)\quad\textbf{the species legitimacy:}
\]

%dsp-level:0
\[
\qquad \Sigma_{\mathrm{cl}}\ \text{is}\ \neg\text{a measure functional of}\ (\Gamma,R)\ —\ \text{its alphabet is the end-sign register}\ \{\sigma_s,\sigma_t\}\ [\mathrm{D3.4};\ \mathrm{CT0.9}],
\]

%dsp-level:0
\[
\ \text{disjoint from the flicker alphabet and from the metric alphabet}\ [\mathrm{BM.37}(2)\ \text{pattern};\ \text{machine: the token census}]:
\]

%dsp-level:0
\[
\textbf{it enters through the rest slot the dispersion read already reads}\ [\mathrm{Def\text{-}BM.10}(4);\ \mathrm{BM.35}(4);\ \mathrm{BM.37}(4)],
\]

%dsp-level:0
\[
\ \text{in the}\ \mathrm{Def\text{-}BM.12}(6)\ \text{pattern of a refinement,}\ \neg\text{a new mass read};
\]

%dsp-level:0
\[
\textbf{no route ratio is read as a mass ratio:}
\]

%dsp-level:0
\[
\qquad \frac{m_{\mathrm{rest}}(0)}{m_{\mathrm{rest}}(+1)}\bigg|_{\varepsilon_{\mathrm{EM}}=0}=\frac{\ln1280}{\ln960}\ \ne\ \frac43
\]

%dsp-level:0
\[
\qquad [\text{machine}];
\]

%dsp-level:0
\[
\qquad \textbf{the strict contraction:}\quad\frac{\Sigma_{\mathrm B}}{6\ln2}<\frac13\iff3\ln\tfrac43<2\ln2\iff16<27
\]

%dsp-level:0
\[
\qquad [\text{machine}];
\]

%dsp-level:0
\[
(9)\quad\textbf{the defect-species companion (conditional, the hypothesis on its face):}
\]

%dsp-level:0
\[
\qquad \textbf{hypothesis:}\ \text{the two baryon classes differ by exactly one down}\leftrightarrow\text{up constituent substitution}\ [\text{the minimal}\ \Delta q=\pm1\ \text{step};\ \mathrm{BM.23}(4)]:
\]

%dsp-level:0
\[
\textbf{then}
\]

%dsp-level:0
\[
\qquad \Delta m\ =\ \Delta_g\ =\ \hat\mu'^2\bigl(D_S(\mathcal F_g,+)-D_S(\mathcal F_g,-)\bigr)\ >\ 0
\]

%dsp-level:0
\[
\qquad [\mathrm{Def\text{-}BM.14}(5);\ \text{machine: the three certificates}]:
\]

%dsp-level:0
\[
\Delta_1=2\hat\mu'^2,\qquad\Delta_2=(\sqrt2-1)\hat\mu'^2,\qquad\Delta_3=\bigl(2\cos\tfrac\pi{10}-\sqrt3\bigr)\hat\mu'^2;
\]

%dsp-level:0
\[
\textbf{the composition itself is}\ \neg\textbf{derived:}\ \text{a fractional constituent alphabet is structurally impossible under an integer-valued additive read}\ [\mathrm{BM.28}(7);\ \mathrm{Def\text{-}BM.8}(5)];
\]

%dsp-level:0
\[
\qquad \text{deficit}\ 1'\ \text{stands}\ [\mathrm{FT.125}(4)\ \text{pointer}];
\]

%dsp-level:0
\[
(10)\quad\textbf{status:}
\]

%dsp-level:0
\[
\qquad \textbf{derived:}\ (1)\text{--}(8);
\]

%dsp-level:0
\[
\qquad \textbf{conditional:}\ (9),\ \text{with its hypothesis on its face}\ [\text{A §0.2a}];
\]

%dsp-level:0
\[
\qquad \textbf{the splitting is exact and strictly positive};
\]

%dsp-level:0
\[
\neg\text{a P4 addition},
\]

%dsp-level:0
\[
\quad\neg\text{a new dimensionful scale}\ [\mathrm{P1}],
\]

%dsp-level:0
\[
\quad\neg\text{a route adopted}\ [\mathrm{BM.30}(5)],
\]

%dsp-level:0
\[
\quad\neg\text{a measured value}\ [\mathrm{P5}];
\]

%dsp-level:0
\[
\qquad \text{§Cal-BM stays empty}.
\]

\textbf{Req.} D0.8, D2.1, D2.9, D3.3, D3.4 (Parts 1–2); CT0.2, CT0.9, CT1.29, EM.9, EM.9b (Part 2); P1, P4, P5 (Parts 1–2); Def-H.28, Def-H.32(2), Def-H.37(2), FT.79(5), FT.106(3a), FT.125(4) (Part 3); Census FT.35(2) (row-wise); Def-BM.5(1), Def-BM.8(5), Def-BM.8(6), Def-BM.10(4), Def-BM.12(6), Def-BM.13, Def-BM.14(5), BM.20(1), BM.23(4), BM.28(1), BM.28(2), BM.28(4)(i), BM.28(7), BM.30(2), BM.30(3), BM.30(5), BM.35(3), BM.35(4), BM.37(2), BM.37(4), BM.37(6), BM.41(4), BM.43(1), BM.46(4), BM.57, BM.58 (this part); the exhaustive end-sign and sign-sequence censuses, exact $\ln$-ring identities with integer certificates, Laurent-monomial exponent arithmetic, token censuses [math; machine: \texttt{checks/\allowbreak{}verify\_\allowbreak{}H\_\allowbreak{}constituent\_\allowbreak{}transport\_\allowbreak{}splitting.py} §C].

\textbf{Pf.}

%dsp-level:0
\[
(1):\ \text{the four degenerate-or-inverted rows re-cited from their own tags};
\]

%dsp-level:0
\[
\ \text{the order residue named as the missing class-dependent rest-slot term}\ [\text{definition}];
\]

%dsp-level:0
\[
\qquad (2):\ \text{the four bars stated};
\]

%dsp-level:0
\[
\ \text{the inventory scan of}\ \mathrm{BM.58}(6)\text{'s seven rows against them}\ [\text{machine: one survivor}]
\]

%dsp-level:0
\[
(3):\ \mathrm{BM.57}(5),(6)\ \text{and}\ \mathrm{BM.58}(3)\ \text{conjoined};
\]

%dsp-level:0
\[
\qquad (4):\ \Sigma_{\mathrm{cl}}(0)-\Sigma_{\mathrm{cl}}(+1)
\]

%dsp-level:0
\[
= \ln20-\ln15
\]

%dsp-level:0
\[
= \ln\tfrac43
\]

%dsp-level:0
\[
= 2\ln2-\ln3\ [\text{machine: the prime-log ring}];
\]

%dsp-level:0
\[
\ 4>3\ [\text{machine}];
\]

%dsp-level:0
\[
\ \text{the exponent row}\ \mathfrak q^0\cdot\mathfrak q^{+1/2}
\]

%dsp-level:0
\[
= \mathfrak q^{+1/2}\ [\text{machine}]
\]

%dsp-level:0
\[
(5):\ \ln m_0-\ln m_{\pm1}=\ln2-0=\ln2\ [\text{machine}];\ 2>1\ [\text{machine}];
\]

%dsp-level:0
\[
\qquad 3<4\implies\ln\tfrac43<\ln2;\quad4<6\implies\ln\tfrac43<\ln2;\quad4<24\implies\ln\tfrac43<3\ln2\ [\text{machine: integer certificates}]
\]

%dsp-level:0
\[
(6):\ \text{exhaustive over}\ 2^k\ \text{sign sequences for}\ k
\]

%dsp-level:0
\[
= 2,
\]

%dsp-level:0
\[
\dots,
\]

%dsp-level:0
\[
7:\ \textstyle\sum f
\]

%dsp-level:0
\[
= -q,
\]

%dsp-level:0
\[
\ \textstyle\sum\aleph
\]

%dsp-level:0
\[
\equiv |q|\ (2),
\]

%dsp-level:0
\[
\ \textstyle\sum\aleph
\]

%dsp-level:0
\[
\ge |q|\ \text{at every one}\ [\text{machine}];
\]

%dsp-level:0
\[
\ \text{the minima}\ 0\ \text{and}\ 1\ \text{attained}\ [\text{machine}]
\]

%dsp-level:0
\[
(7):\ \text{the three references enumerated and their class differences computed in the}\ \ln\text{-ring}\ [\text{machine}];
\]

%dsp-level:0
\[
\ W(\mathrm{vac})
\]

%dsp-level:0
\[
= 1\ \text{for the empty register}\ [\text{machine}];
\]

%dsp-level:0
\[
\ \mathrm{Def\text{-}BM.5}(1)\text{'s boundary-data clause re-cited};
\]

%dsp-level:0
\[
\ \mathrm{Def\text{-}H.28}\text{'s external-ansatz mark pinned}
\]

%dsp-level:0
\[
(8):\ \text{the token census of the defining clause against the flicker and metric alphabets}\ [\text{machine}];
\]

%dsp-level:0
\[
\qquad \ln1280/\ln960\ne\tfrac43\iff3\ln1280\ne4\ln960\iff1280^3\ne960^4\ [\text{machine: the two integers differ}];
\]

%dsp-level:0
\[
\qquad 3\ln\tfrac43<2\ln2\iff\ln\tfrac{64}{27}<\ln4\iff64<108\ [\text{machine}]
\]

%dsp-level:0
\[
(9):\ \mathrm{Def\text{-}BM.14}(5)\text{'s three splittings re-cited with their certificates}\ 4>2,
\]

%dsp-level:0
\[
\ 2>1,
\]

%dsp-level:0
\[
\ 3<5\ [\text{machine}];
\]

%dsp-level:0
\[
\qquad (10):\ \text{assembly}. \ \blacksquare
\]

\textbf{Phys.} Worked from both ends, the answer is the same object. From the mass side: whatever reverses the two rows has to be class-dependent, positive for the neutral class, dimensionally an energy through the rest slot, and built from data the edge count does not already carry — and the inventory has exactly one unbarred candidate. From the force side: colour is a projector that annihilates its own slot data on anything sealed, electromagnetism is a census that admits every local class, so the only class-dependent datum a sealed junction can carry is the count of admissible end-sign registers — twenty against fifteen. The candidate and the survivor are the same. Taken as a logarithm and booked against the vacuum in the ledger's own currency, the filter becomes an energy: the neutral class carries $\ln(4/3)$ more transport entropy than the charged one, and the neutral-core constituent carries exactly one bit more than the charged-core constituent. The reason the register counts as transport and not as bookkeeping is the parity lock: a chain whose ends agree can run wall-free, and a chain whose ends disagree cannot — the wall content of a constituent's transport is odd or even according to its charge core, so the multiplicity is a property of how the seed is carried, not of how the state is named. The sign is not a convention either. Of three possible references only one is a state the composite can be in, and it is the one the corpus already uses for every other deficit.

\end{theorem}

\pseudohead{theorem}{\texorpdfstring{BM.60}{BM.60}}{\texorpdfstring{The six-state assembly: the generation ladder, the branch split, the Platonic floors, and the two hierarchies}{The six-state assembly: the generation ladder, the branch split, the Platonic floors, and the two hierarchies}}\label{tag:BM-60}
\textbf{Thm.}

%dsp-level:0
\[
(1)\quad\textbf{the generation ladder (the 3):}
\]

%dsp-level:0
\[
\qquad \delta_q=\frac{2}{3\cdot3\cdot3}=\frac2{27},
\]

%dsp-level:0
\[
\qquad 3\delta_q=\delta_\ell=\frac29,
\]

%dsp-level:0
\[
\qquad 16\,\delta_\nu=9\,\delta_\ell=27\,\delta_q=2
\]

%dsp-level:0
\[
\qquad [\mathrm{FT.119}(1),(2);\ \mathrm{FT.143}(5);\ \text{machine: exact Fractions}]:
\]

%dsp-level:0
\[
r^{(q)}_i=1+\sqrt2\cos\Bigl(\frac2{27}+\frac{2\pi i}3\Bigr),
\]

%dsp-level:0
\[
\qquad \sum_ir^{(q)}_i=3=\dim F_{J^\pm},
\]

%dsp-level:0
\[
\qquad \sum_i\bigl(r^{(q)}_i\bigr)^2=6,
\]

%dsp-level:0
\[
\qquad \sum_{i<j}r^{(q)}_ir^{(q)}_j=\frac32
\]

%dsp-level:0
\[
\qquad [\mathrm{FT.111}(3);\ \mathrm{FT.119}(4);\ \text{machine}];
\]

%dsp-level:0
\[
2\bigl(r^{(q)}-1\bigr)^3-3\bigl(r^{(q)}-1\bigr)-\sqrt2\cos\tfrac29=0\ \ \text{for each of the three}
\]

%dsp-level:0
\[
\qquad [\mathrm{FT.119}(4);\ \text{machine: enclosure containment at width}<2^{-90}];
\]

%dsp-level:0
\[
Q=\frac{\sum_i(r^{(q)}_i)^2}{\bigl(\sum_ir^{(q)}_i\bigr)^2}=\frac69=\frac23
\]

%dsp-level:0
\[
\qquad [\mathrm{Def\text{-}BM.6}(4);\ \mathrm{FT.158}(4);\ \text{machine: exact, phase-independent}];
\]

%dsp-level:0
\[
\textbf{the generation count is derived:}
\]

%dsp-level:0
\[
\qquad \boxed{\ \#\text{lines}=3\ }
\]

%dsp-level:0
\[
\qquad [\text{the orbit length of}\ A_3:\ U^3=E_0,\ U\ne E_0,\ U^2\ne E_0\ [\mathrm{BM.57}(9);\ \text{machine}];\qquad\sum_ir_i=3=\dim F_{J^\pm}];
\]

%dsp-level:0
\[
(2)\quad\textbf{the branch split (the 2):}
\]

%dsp-level:0
\[
\qquad \mathcal H_{\mathsf v}\cong\mathbb C^2\ [\mathrm{EW.32}];
\]

%dsp-level:0
\[
\qquad Q|_{\mathcal H_{\mathsf v}}=\tfrac12\sigma_3+\tfrac12\sigma_0=\operatorname{diag}(1,0),
\]

%dsp-level:0
\[
\qquad \operatorname{spec}Q=\{0,+1\},
\]

%dsp-level:0
\[
\qquad \omega:\ \operatorname{spec}\mapsto\{0,-1\}
\]

%dsp-level:0
\[
\qquad [\mathrm{EW.36}]:
\]

%dsp-level:0
\[
\textbf{the seed carrier is a doublet, and its weight alphabet together with the}\ \omega\textbf{-image is exactly the ternary end read}\ \{-1,0,+1\}\ [\mathrm{EW.36};\ \mathrm{Def\text{-}BM.12}(1);\ \mathrm{D3.4};\ \text{machine}];
\]

%dsp-level:0
\[
\textbf{the geometric branch:}
\]

%dsp-level:0
\[
\qquad \text{the dual exchanges}\ (V,F)\ \text{at fixed}\ E\implies\text{the two members of a dual pair share}\ m_{\mathrm{hid}}\ \text{and}\ \Pi\ \text{and split on every finer ledger}\ [\mathrm{FT.44}(1),(2)]:
\]

%dsp-level:0
\[
\text{independent windings}\ F{-}1\ (5\ \text{vs}\ 7;\ 11\ \text{vs}\ 19);
\]

%dsp-level:0
\[
\qquad \text{lattice index}\ 2^{F-1}\ (32\ \text{vs}\ 128;\ 2048\ \text{vs}\ 2^{19});
\]

%dsp-level:0
\[
\qquad g(V)\ (385\ \text{vs}\ 25;\ 292938721075\ \text{vs}\ 164395);
\]

%dsp-level:0
\[
\qquad O(2d)\ (112\ \text{vs}\ 496;\ 112\ \text{vs}\ 2032);
\]

%dsp-level:0
\[
\textbf{and on the pole side they split on every defect datum}\ [\mathrm{Census\ FT.63}(6):\ \text{the derived dual-splitting law}]:
\]

%dsp-level:0
\[
\qquad \delta_v:\ \tfrac\pi2\ \text{vs}\ \tfrac{2\pi}3;\quad\tfrac\pi5\ \text{vs}\ \tfrac\pi3;
\]

%dsp-level:0
\[
\qquad D_S:\ 2-\sqrt2\ \text{vs}\ 1;\quad2-2\cos\tfrac\pi{10}\ \text{vs}\ 2-\sqrt3;
\]

%dsp-level:0
\[
\textbf{the tetrahedron is self-dual}\ [\mathrm{FT.44}(3)]\implies\text{its branch index is the weight-ladder rung}\ m\in\{1,2\},
\]

%dsp-level:0
\[
\qquad D_S(m)=(2,\,4,\,2,\,0)\ \ \text{at}\ m=(1,2,3,4)
\]

%dsp-level:0
\[
\qquad [\mathrm{Census\ FT.63}(6);\ \mathrm{FT.60}(5);\ \text{machine}];
\]

%dsp-level:0
\[
(3)\quad\textbf{the metric back-reaction of the directed 3-vertex:}
\]

%dsp-level:0
\[
\qquad \text{the skeleton metric is piecewise flat with conical defects at the vertices}\ [\mathrm{Census\ FT.63}(4);\ \mathrm{Def\text{-}H.21}(6);\ \mathrm{FT.54}]:
\]

%dsp-level:0
\[
\delta_v=2\pi-d\,\frac{(L-2)\pi}L,
\]

%dsp-level:0
\[
\qquad \frac{\delta_v}\pi=\Bigl(1,\ \frac12,\ \frac23,\ \frac15,\ \frac13\Bigr),
\]

%dsp-level:0
\[
\qquad \boxed{\ V\,\delta_v=4\pi\ \ \text{on all five solids}\ }
\]

%dsp-level:0
\[
\qquad [\text{Descartes/Gauss--Bonnet};\ \text{machine: exact in}\ \pi\text{-units}]:
\]

%dsp-level:0
\[
4\cdot\pi=8\cdot\tfrac\pi2=6\cdot\tfrac{2\pi}3=20\cdot\tfrac\pi5=12\cdot\tfrac\pi3=4\pi\qquad[\text{machine}];
\]

%dsp-level:0
\[
\text{a loop cleanly enclosing one vertex carries}\ \Theta_R
\]

%dsp-level:0
\[
= \delta_v\ [\mathrm{FT.47}(1);\ \mathrm{Def\text{-}H.21}(6)]
\]

%dsp-level:0
\[
\implies m^2_{\mathrm{pole}}
\]

%dsp-level:0
\[
= 4\hat\mu'^2\sin^2\frac{\Theta_R}4
\]

%dsp-level:0
\[
= \hat\mu'^2D_S\ [\mathrm{Def\text{-}H.37}(2);\ \mathrm{BM.39}(7)]:
\]

%dsp-level:0
\[
\textbf{the branch mass is the vertex's own metric back-reaction}\ —\ \text{the defect at a}\ d\text{-valent vertex, and}\ d
\]

%dsp-level:0
\[
= 3\ \text{on the tetrahedron, the cube and the dodecahedron}\ [\mathrm{Census\ FT.42}(5):\ (V,E,F,d)\ \text{column}];
\]

%dsp-level:0
\[
(4)\quad\textbf{the six states:}
\]

%dsp-level:0
\[
\qquad |\mathcal S|=3\cdot2=6
\]

%dsp-level:0
\[
\qquad [\text{machine}];
\]

%dsp-level:0
\[
\qquad \textbf{the six defect values are pairwise distinct}\ [\text{machine: separation by exact field and by minimal polynomial}]:
\]

%dsp-level:0
\[
D_S\in\Bigl\{\,2,\ \ 4,\ \ 2-\sqrt2,\ \ 1,\ \ 2-2\cos\tfrac\pi{10},\ \ 2-\sqrt3\,\Bigr\},
\]

%dsp-level:0
\[
\qquad \text{minimal polynomials}\ \ (x{-}2),\ (x{-}4),\ (x^2{-}4x{+}2),\ (x{-}1),\ (x^4{-}8x^3{+}19x^2{-}12x{+}1),\ (x^2{-}4x{+}1);
\]

%dsp-level:0
\[
\textbf{the fields:}
\]

%dsp-level:0
\[
\qquad \mathbb Q,\ \ \mathbb Q,\ \ \mathbb Q(\sqrt2),\ \ \mathbb Q,\ \ \mathbb Q\bigl(\cos\tfrac\pi{10}\bigr),\ \ \mathbb Q(\sqrt3)
\]

%dsp-level:0
\[
\qquad [\text{machine: the dodecahedral quartic verified by exact substitution into}\ \mathbb Z[c]/(16c^4{-}20c^2{+}5)]:
\]

%dsp-level:0
\[
\textbf{the defect species resolves all six states with zero degeneracy};
\]

%dsp-level:0
\[
\qquad \textbf{the informational species is branch-degenerate:}
\]

%dsp-level:0
\[
\quad m_{\mathrm{hid}}\,\mathfrak h/\ln2=2E\in\{12,\,24,\,24,\,60,\,60\}\implies\textbf{three values}\ [\mathrm{FT.44}(1);\ \text{machine}];
\]

%dsp-level:0
\[
(5)\quad\textbf{the branch order (down-type heavier, uniform and exact):}
\]

%dsp-level:0
\[
\qquad \boxed{\ D_S(\mathcal F_g,+)>D_S(\mathcal F_g,-)\qquad\forall g\in\{1,2,3\}\ }
\]

%tbl-level:0
{\setlength{\tabcolsep}{4pt}
\begin{longtable}{>{\raggedright\arraybackslash}p{\dimexpr0.0949\textwidth-2\tabcolsep\relax}|>{\raggedright\arraybackslash}p{\dimexpr0.1961\textwidth-2\tabcolsep\relax}|>{\raggedright\arraybackslash}p{\dimexpr0.1961\textwidth-2\tabcolsep\relax}|>{\raggedright\arraybackslash}p{\dimexpr0.5129\textwidth-2\tabcolsep\relax}}
\toprule
floor & lighter branch & heavier branch & the exact certificate \\
\addlinespace[2pt]
\midrule
\endfirsthead
\toprule
floor & lighter branch & heavier branch & the exact certificate \\
\addlinespace[2pt]
\midrule
\endhead
\bottomrule
\endlastfoot
$\mathcal F_1$ & $D_S=\allowbreak 2$ & $D_S=\allowbreak 4$ & $4>\allowbreak 2$ \\
$\mathcal F_2$ & $2-\sqrt2$ & $1$ & $1>\allowbreak 2-\sqrt2\iff\allowbreak \sqrt2>\allowbreak 1\iff\allowbreak 2>\allowbreak 1$ \\
$\mathcal F_3$ & $2-2\cos\tfrac\pi{10}$ & $2-\sqrt3$ & $\cos\tfrac\pi{10}>\allowbreak \tfrac{\sqrt3}2\iff\allowbreak \tfrac\pi{10}<\allowbreak \tfrac\pi6\iff\allowbreak 3<\allowbreak 5$ \\
\end{longtable}
}

%dsp-level:0
\[
[\text{machine: integer and radical certificates;}\ \cos\ \text{strictly decreasing on}\ [0,\pi]\ [\text{math}]]
\]

%dsp-level:0
\[
\qquad \textbf{the constituent mass difference:}
\]

%dsp-level:0
\[
\qquad \Delta_g=\hat\mu'^2\bigl(D_S(\mathcal F_g,+)-D_S(\mathcal F_g,-)\bigr)>0;
\]

%dsp-level:0
\[
\Delta_1=2\hat\mu'^2,
\]

%dsp-level:0
\[
\qquad \Delta_2=(\sqrt2-1)\hat\mu'^2,
\]

%dsp-level:0
\[
\qquad \Delta_3=\bigl(2\cos\tfrac\pi{10}-\sqrt3\bigr)\hat\mu'^2
\]

%dsp-level:0
\[
\qquad [\text{machine: exact in}\ \mathbb Q,\ \mathbb Q(\sqrt2),\ \mathbb Q(\sqrt3,\cos\tfrac\pi{10})];
\]

%dsp-level:0
\[
(6)\quad\textbf{the two hierarchies (the derived inversion):}
\]

%dsp-level:0
\[
\qquad \textbf{on the informational species:}\quad m_{\mathrm{hid}}(\mathcal F_g)=\frac{2E(\mathcal F_g)\ln2}{\mathfrak h},
\]

%dsp-level:0
\[
\qquad 2E=(12,\,24,\,60):
\]

%dsp-level:0
\[
\boxed{\ m_{\mathrm{hid}}(\mathcal F_1)<m_{\mathrm{hid}}(\mathcal F_2)<m_{\mathrm{hid}}(\mathcal F_3),\qquad12:24:60=1:2:5\ }
\]

%dsp-level:0
\[
\qquad [\text{machine: exact rationals, strictly increasing}]:
\]

%dsp-level:0
\[
\qquad \textbf{the mass increases across the generations};
\]

%dsp-level:0
\[
\textbf{on the defect species:}
\]

%dsp-level:0
\[
\qquad D_S(\mathcal F_3,\pm)<D_S(\mathcal F_2,\pm)<D_S(\mathcal F_1,\pm)
\]

%dsp-level:0
\[
\qquad [\text{machine:}\ 2-2\cos\tfrac\pi{10}<2-\sqrt3<2-\sqrt2<1<2<4]:
\]

%dsp-level:0
\[
\boxed{\ \textbf{the two currencies order the floors oppositely}\ }
\]

%dsp-level:0
\[
\qquad [\text{the exact species inversion};\ \mathrm{BM.39}(5)\ \text{pattern}];
\]

%dsp-level:0
\[
\textbf{the relating equation is}\ \neg\textbf{derived:}
\]

%dsp-level:0
\[
\qquad \neg\exists\ \text{a theorem equating the bit currency}\ (\ln2\ \text{per directed edge})\ \text{and the angle currency}\ (\Theta_R)\ [\mathrm{BM.35}(5)];
\]

%dsp-level:0
\[
(7)\quad\textbf{the generation map:}
\]

%dsp-level:0
\[
\qquad \text{two exact three-element ordered sets}\ —\ \bigl\{(r^{(q)}_1)^2,(r^{(q)}_2)^2,(r^{(q)}_0)^2\bigr\}\ [\mathrm{FT.119}(5):\ \text{enclosure-proven, frozen}]\ \text{and}\ \{2E(\mathcal F_g)\}_{g=1,2,3}\ [(6)]:
\]

%dsp-level:0
\[
\boxed{\ \text{an order-isomorphism between two totally ordered three-element sets is unique}\ }
\]

%dsp-level:0
\[
\qquad [\text{math};\ \text{machine: the uniqueness census over all}\ 3!=6\ \text{bijections}]:
\]

%dsp-level:0
\[
g=1\longleftrightarrow(r^{(q)}_1)^2\longleftrightarrow\mathcal F_1\ (\text{the tetrahedron floor});
\]

%dsp-level:0
\[
\qquad g=2\longleftrightarrow(r^{(q)}_2)^2\longleftrightarrow\mathcal F_2\ (\text{the cube/octahedron floor});
\]

%dsp-level:0
\[
\qquad g=3\longleftrightarrow(r^{(q)}_0)^2\longleftrightarrow\mathcal F_3\ (\text{the dodecahedron/icosahedron floor});
\]

%dsp-level:0
\[
\textbf{the map is an order,}\ \neg\textbf{a number:}
\]

%dsp-level:0
\[
\qquad \hat\mu'^2,\ \mu\in\mathrm{Cal\text{-}F}\ [\text{logged}];
\]

%dsp-level:0
\[
\qquad \text{on the transport species the line order inverts}\ [\mathrm{BM.39}(5)]\implies\text{the isomorphism inverts with it}\ [\text{machine}];
\]

%dsp-level:0
\[
(8)\quad\textbf{the transport-species form of the ladder:}
\]

%dsp-level:0
\[
\qquad m^2_i=M^2\bigl(A_0-B_0\cos\varphi_i\bigr),\quad M^2=\hat\mu^2\ell_c^2,\quad A_0=\frac{3g^2}8+\frac{g'^2}{16},\quad B_0=\frac{gg'}4
\]

%dsp-level:0
\[
\qquad [\mathrm{Def\text{-}BM.11}(5);\ \mathrm{BM.39}(1)]:
\]

%dsp-level:0
\[
\textbf{the exact evaluation at the base profile}\ [\text{machine: in}\ \mathbb Q[g,g',c,s,r_3]/(c^2{+}s^2{-}1,\,r_3^2{-}3)]:
\]

%dsp-level:0
\[
\qquad G\bigl(c(i)\bigr)=A_0-B_0\cos\Bigl(\delta+\frac{2\pi(2-i)}3\Bigr);
\]

%dsp-level:0
\[
\textbf{the orbit value set:}
\]

%dsp-level:0
\[
\qquad \bigl\{G(c(i))\bigr\}_{i=0,1,2}=\bigl\{A_0-B_0\cos(\delta-\tfrac{2\pi j}3)\bigr\}_{j=0,1,2}
\]

%dsp-level:0
\[
\qquad [\text{machine: set equality in the quotient ring}];
\]

%dsp-level:0
\[
\textbf{the sum rule:}
\]

%dsp-level:0
\[
\qquad \sum_iG\bigl(c(i)\bigr)=3A_0
\]

%dsp-level:0
\[
\qquad [\text{machine: exact, phase-independent}];
\]

%dsp-level:0
\[
\qquad \textbf{the positivity:}
\]

%dsp-level:0
\[
\qquad A_0-B_0=\frac38\Bigl(g-\frac{g'}3\Bigr)^2+\frac{g'^2}{48}>0
\]

%dsp-level:0
\[
\qquad [\mathrm{BM.39}(1);\ \text{machine}];
\]

%dsp-level:0
\[
(9)\quad\textbf{the no-import bar:}
\]

%dsp-level:0
\[
\qquad \neg\exists\ \text{an SM mass, mass ratio or Yukawa value in any clause}\ [\mathrm{P5};\ \text{machine: zero external tokens, zero decimal numerals in this tier}];
\]

%dsp-level:0
\[
\neg\exists\ \text{a Yukawa-type object in the inventory}\ [\mathrm{EW.15};\ \mathrm{FT.145}(3);\ \mathrm{Def\text{-}BM.11}\ \text{preamble}];
\]

%dsp-level:0
\[
\qquad \textbf{the six states and both hierarchies are derived from the card template, the duality law, the defect spectrum and the channel angle alone};
\]

%dsp-level:0
\[
(10)\quad\textbf{status:}
\]

%dsp-level:0
\[
\qquad \textbf{derived:}\ (1)\text{--}(8);
\]

%dsp-level:0
\[
\qquad \textbf{declared:}\ \mathrm{Def\text{-}BM.14}(8)\text{'s name attachment alone};
\]

%dsp-level:0
\[
\qquad \textbf{open:}\ \text{the currency relating equation}\ (6),\ \text{the flavour composition}\,\ \text{the scales}\ [\mathrm{P4}];
\]

%dsp-level:0
\[
\neg\text{a P4 addition},\quad\neg\text{a new dimensionful scale},\quad\text{§Cal-BM stays empty};
\]

%dsp-level:0
\[
\qquad \text{species names confined to the identification layers}\ [\mathrm{Def\text{-}BM.3}(7),\ \mathrm{Def\text{-}BM.6}(7),\ \mathrm{Def\text{-}BM.14}(8)].
\]

\textbf{Req.} A6, D0.8, D2.9 (Parts 1–2); CT3.17a(2), EW.32, EW.36, EM.9 (Part 2); P1, P4, P5 (Parts 1–2); Def-H.20, Def-H.21(6), Def-H.37(2), Def-H.38, FT.44(1)–(3), FT.47(1), FT.54(4), FT.60(5), FT.111(3), FT.116, FT.119(1), FT.119(2), FT.119(4), FT.119(5), FT.143(5), FT.145(3), FT.158(1), FT.158(4) (Part 3); Census FT.42(5), FT.63(4), FT.63(6) (row-wise); Def-BM.3(7), Def-BM.6(3), Def-BM.6(4), Def-BM.11(5), Def-BM.12(1), Def-BM.14, BM.35(5), BM.39(1), BM.39(5), BM.39(7), BM.57(9) (this part); exact radical rings $\mathbb Q(\sqrt2),\mathbb Q(\sqrt3),\mathbb Q(\sqrt5),\mathbb Q(\cos\tfrac\pi{10})$, exact minimal polynomials by substitution, de Moivre binomials at $\cos5t=16c^5-20c^3+5c$, rational interval enclosures with alternating-series remainder bounds, quotient-ring trigonometric identities, order-isomorphism uniqueness [math; machine: \texttt{checks/\allowbreak{}verify\_\allowbreak{}H\_\allowbreak{}directed\_\allowbreak{}3vertex\_\allowbreak{}splitting.py} §A–§D].

\textbf{Pf.}

%dsp-level:0
\[
(1):\ \delta_q
\]

%dsp-level:0
\[
= 2/(3\cdot3\cdot3)
\]

%dsp-level:0
\[
= 2/27,
\]

%dsp-level:0
\[
\ 3\cdot\tfrac2{27}
\]

%dsp-level:0
\[
= \tfrac29,
\]

%dsp-level:0
\[
\ 16\cdot\tfrac18
\]

%dsp-level:0
\[
= 9\cdot\tfrac29
\]

%dsp-level:0
\[
= 27\cdot\tfrac2{27}
\]

%dsp-level:0
\[
= 2\ [\text{machine: exact Fractions}];
\]

%dsp-level:0
\[
\qquad \text{the spectrum and the three Vieta reads re-derived by enclosure}\ [\text{machine: widths}<2^{-90}];
\]

%dsp-level:0
\[
\qquad Q=6/9=2/3\ [\text{machine}]
\]

%dsp-level:0
\[
\text{the cubic}\ 2y^3-3y-\sqrt2\cos\tfrac29\ \text{evaluated at each}\ y=r^{(q)}_i-1\ \text{contains}\ 0\ [\text{machine}];
\]

%dsp-level:0
\[
\qquad U^3=E_0,\ U\ne E_0,\ U^2\ne E_0\ \text{for the 3-cycle over}\ \mathbb Q(i)\ [\text{machine}]
\]

%dsp-level:0
\[
(2):\ \mathrm{EW.36}\text{'s}\ \operatorname{diag}(1,0)\ \text{re-derived from}\ \tfrac12\sigma_3+\tfrac12\sigma_0\ [\text{machine}];
\]

%dsp-level:0
\[
\ \text{the union with the}\ \omega\text{-image equals the ternary alphabet}\ [\text{machine}];
\]

%dsp-level:0
\[
\qquad \text{the four finer-ledger splittings and the defect splittings re-cited}\ [\mathrm{FT.44}(2);\ \mathrm{Census\ FT.63}(6)];
\]

%dsp-level:0
\[
\ \text{the tetrahedral ladder}\ (2,4,2,0)\ \text{re-cited with its exact zero at}\ m\qquad = 4
\]

%dsp-level:0
\[
(3):\ \delta_v
\]

%dsp-level:0
\[
= 2\pi-d(L-2)\pi/L\ \text{evaluated at}\ (d,L)
\]

%dsp-level:0
\[
= (3,3),
\]

%dsp-level:0
\[
(3,4),
\]

%dsp-level:0
\[
(4,3),
\]

%dsp-level:0
\[
(3,5),
\]

%dsp-level:0
\[
(5,3)\ [\text{machine: exact rationals in}\ \pi\text{-units}];
\]

%dsp-level:0
\[
\qquad V\delta_v\ \text{multiplied out at all five}\ [\text{machine}];
\]

%dsp-level:0
\[
\qquad 4\sin^2(\delta_v/4)=D_S\ \text{by the half-angle identity}\ [\text{machine in the exact fields}]
\]

%dsp-level:0
\[
(4):\ |\mathcal S|=3\cdot2=6\ [\text{machine}];
\]

%dsp-level:0
\[
\qquad \text{pairwise distinctness: the rational members separated by integer comparison, the quadratic members by their minimal polynomials}\ (x^2-4x+2)\ \text{and}\ (x^2-4x+1)\ \text{having disjoint root sets, the quartic member separated by}\ \mathrm{Census\ FT.63}(6)\text{'s isolated roots}\ [\text{machine}];
\]

%dsp-level:0
\[
\qquad \{12,24,24,60,60\}\ \text{has three distinct values}\ [\text{machine}]
\]

%dsp-level:0
\[
(5):\ \text{the three certificates of the table}\ [\text{machine}];
\]

%dsp-level:0
\[
\qquad \cos\ \text{strictly decreasing on}\ [0,\pi]\ [\text{math}]\implies\tfrac\pi{10}<\tfrac\pi6\implies\cos\tfrac\pi{10}>\cos\tfrac\pi6
\]

%dsp-level:0
\[
= \tfrac{\sqrt3}2\ [\text{machine:}\ \cos^2\tfrac\pi6=\tfrac34];
\]

%dsp-level:0
\[
\qquad \Delta_g\ \text{formed as exact differences in each field}\ [\text{machine}]
\]

%dsp-level:0
\[
(6):\ 12<24<60\ \text{and}\ 12:24:60=1:2:5\ [\text{machine}];
\]

%dsp-level:0
\[
\qquad 2-2\cos\tfrac\pi{10}<2-\sqrt3<2-\sqrt2<1<2<4\ \text{by the certificates}\ 3<5,
\]

%dsp-level:0
\[
\ 2<3,
\]

%dsp-level:0
\[
\ 1<2,
\]

%dsp-level:0
\[
\ \text{and integer comparison}\ [\text{machine}]
\]

%dsp-level:0
\[
(7):\ \text{all}\ 6\ \text{bijections between the two ordered sets enumerated; exactly one is order-preserving}\ [\text{machine}];
\]

%dsp-level:0
\[
\qquad \text{the inversion under}\ \mathrm{BM.39}(5)\text{'s affine-decreasing row}\ [\text{machine}]
\]

%dsp-level:0
\[
(8):\ G(c(i))\ \text{built from the Gram block of}\ \mathrm{BM.57}(5)\ \text{at}\ c(i)
\]

%dsp-level:0
\[
= (\cos\delta,\cos(\delta{+}\tfrac{2\pi}3),\cos(\delta{+}\tfrac{2\pi(2-i)}3),\tfrac12)\ \text{and reduced modulo}\ c^2{+}s^2{-}1,
\]

%dsp-level:0
\[
\ r_3^2{-}3\ [\text{machine}];
\]

%dsp-level:0
\[
\qquad \text{the orbit value set and the sum rule}\ \sum_iG(c(i))=3A_0\ [\text{machine}];
\]

%dsp-level:0
\[
\qquad A_0-B_0\ \text{as a sum of squares}\ [\text{machine}]
\]

%dsp-level:0
\[
(9)\text{--}(10):\ \text{the token scans and the assembly}\ [\text{machine}]. \ \blacksquare
\]

\textbf{Phys.} Six quarks fall out of three structures that were already sitting in the corpus in separate rooms. The three generations are the three lines of the colour channel, whose angle is one integer over three integers and whose orbit length is forced by the same $C_3$ that seals the baryon. The two members of each generation are a floor and its dual: two solids that share every edge count, hence every informational mass, and differ at every vertex figure — three squares meeting at right angles against four triangles meeting at one hundred twenty degrees, three pentagons against five triangles. Where duality is trivial, at the tetrahedron, the second branch is the weight ladder's own next rung. The mass of a branch is the conical defect at its vertex, and Descartes closes the ledger: five solids, each with its vertex defects summing to exactly four pi. The hierarchy comes out exact and comes out twice, in opposite directions. In the bit currency the floors climb one to two to five, strictly. In the angle currency they descend. That inversion is not a defect of the assembly; it is the same read-relativity the corpus already proved for the family lines, and the equation that would reconcile the two currencies is the one item the ledger still carries as open. Within a generation there is no ambiguity at all: the heavier branch is heavier by an exact algebraic number — two, root two minus one, two cosines of eighteen degrees minus root three — and it is heavier for the same reason in all three. Nothing was fitted, and nothing measured entered.

\begin{theorem}[{\texorpdfstring{BM.61 — $\alpha_B$: the inventory, the coupling-blindness theorem, and the obstruction}{BM.61 - : the inventory, the coupling-blindness theorem, and the obstruction}}]\label{tag:BM-61}
\textbf{Thm.}

%dsp-level:0
\[
(1)\quad\textbf{the inventory scan:}
\]

%dsp-level:0
\[
\qquad \neg\exists\ \text{a quantity}\ \alpha_B\ \text{in the corpus:}
\]

%dsp-level:0
\[
\qquad \text{occurrences of the token in Parts 1--4, the Census}:\ \ \mathbf 0
\]

%dsp-level:0
\[
\qquad [\text{machine: the token scan}];
\]

%dsp-level:0
\[
\textbf{it is}\ \neg\textbf{a P4-class member and}\ \neg\textbf{a logged combination:}
\]

%dsp-level:0
\[
\qquad \text{the P4 core}\ [\text{A §E}],
\]

%dsp-level:0
\[
\ \mathrm{Cal\text{-}2},
\]

%dsp-level:0
\[
\ \mathrm{Cal\text{-}F},
\]

%dsp-level:0
\[
\ \text{and}\ \text{§Cal-BM (empty)}\ \text{carry no such entry}\ [\text{machine: the registry scan}];
\]

%dsp-level:0
\[
(2)\quad\textbf{the second argument is void:}
\]

%dsp-level:0
\[
\qquad \neg\exists\ \text{an}\ \alpha_{\mathrm{Strong}}\ \text{in the inventory}\ [\text{machine: zero occurrences}];
\]

%dsp-level:0
\[
\qquad \textbf{the colour sector's strength-class data are counts and ratios of counts:}
\]

%dsp-level:0
\[
\dim\mathcal K^{(\mathrm{strong})}=8\ [\mathrm{CT3.8};\ \mathrm{BM.36}(1)];
\]

%dsp-level:0
\[
\qquad g(N)=\sum_{\substack{a,b\ge0\ 2a+3b=N}}\frac{N!}{a!\,b!\,2^a\,6^b}=(0,\,1,\,1,\,3,\,10,\,25,\,105)\ \ \text{at}\ N=1,\dots,7
\]

%dsp-level:0
\[
\qquad [\mathrm{CT3.17a}(2);\ \mathrm{FT.41}(b);\ \text{machine: recomputed exactly}];
\]

%dsp-level:0
\[
\mathcal C(\Sigma;C)\in\{0,1\}\ [\mathrm{D2.14}];
\]

%dsp-level:0
\[
\qquad \text{the arenas}\ 27\ \text{and}\ 9\ [\mathrm{FT.118}(2)];
\]

%dsp-level:0
\[
\qquad \delta_q=\frac{2}{3\cdot3\cdot3}=\frac2{27}\ [\mathrm{FT.119}(1)];
\]

%dsp-level:0
\[
\qquad g(6)=25=15+10\ [\text{machine: the split}];
\]

%dsp-level:0
\[
\textbf{every one is an integer or a ratio of integers}\ —\ \neg\text{a coupling};
\]

%dsp-level:0
\[
(3)\quad\textbf{the strength bar:}
\]

%dsp-level:0
\[
\qquad \neg\exists\ \text{a theorem fixing any numerical coupling strength}\ [\mathrm{BM.36}(3)];
\]

%dsp-level:0
\[
\qquad (g,g')\ \text{are P4-class}\ [\mathrm{EW.33};\ \mathrm{P4}];
\]

%dsp-level:0
\[
\qquad \alpha^{(2)}_{\mathrm{EM}}\ \text{is an operator block read whose value no theorem fixes}\ [\mathrm{EM.17};\ \mathrm{EM.25}(8);\ \text{Part 2 §Cal-2's logged row}];
\]

%dsp-level:0
\[
\textbf{the only force-structure ratios the corpus derives are the alphabet-capacity ratios}
\]

%dsp-level:0
\[
\qquad \frac{\dim\mathcal K^{(\mathrm{EM})}}{\dim F_W}=\frac13,
\]

%dsp-level:0
\[
\qquad \frac{\dim\mathcal K^{(\mathrm{strong})}}{\dim\mathcal K^{(\mathrm{EM})}}=8,
\]

%dsp-level:0
\[
\qquad \frac{\dim\mathcal K^{(\mathrm{EM})}}{\dim\mathcal K_3}=\frac1{25}
\]

%dsp-level:0
\[
\qquad [\mathrm{BM.36}(2);\ \text{machine: exact Fractions from the flag}\ (6,16,17,25)]:
\]

%dsp-level:0
\[
\textbf{reading a capacity ratio as a coupling strength is a category error of the geometric-ratio-against-coupling species}\ [\mathrm{BM.36}(3);\ \mathrm{FT.73}(5)\ \text{trap pattern}];
\]

%dsp-level:0
\[
(4)\quad\textbf{the first argument is unfixed:}
\]

%dsp-level:0
\[
\qquad e=\frac{gg'}{\sqrt{g^2+g'^2}}=g\sin\theta_W=g'\cos\theta_W,
\]

%dsp-level:0
\[
\qquad \frac1{e^2}=\frac1{g^2}+\frac1{g'^2}
\]

%dsp-level:0
\[
\qquad [\mathrm{EW.35};\ \mathrm{EM.25}(6):\ \text{fixed with zero freedom}]:
\]

%dsp-level:0
\[
\neg\exists\ \text{a theorem asserting a value of}\ e\ [\mathrm{EM.25}(7)];
\]

%dsp-level:0
\[
\qquad \varepsilon_{\mathrm{EM}}
\]

%dsp-level:0
\[
= e^2/(4\pi)\ \text{is a dimensionless P4-class symbol carried symbolically}\ [\mathrm{Def\text{-}BM.12}\ \text{local notation};\ \mathrm{EM.25}(9);\ \mathrm{P5}];
\]

%dsp-level:0
\[
(5)\quad\textbf{the coupling-blindness theorem:}
\]

%dsp-level:0
\[
\qquad W(Q)=\bigl[z^{Q}\bigr]\bigl(2+z+z^{-1}\bigr)^{k}\in\mathbb Z[z,z^{-1}]
\]

%dsp-level:0
\[
\qquad [\mathrm{Def\text{-}BM.13}(3);\ \mathrm{Def\text{-}BM.8}(6)]:
\]

%dsp-level:0
\[
\textbf{the free-symbol set of}\ W(Q)\ \textbf{is}\ \{z\},
\]

%dsp-level:0
\[
\ \text{with coefficients fixed by}\ (m_q,k)
\]

%dsp-level:0
\[
= \bigl((2,1,1),\,3\bigr)\ [\mathrm{EM.9};\ \mathrm{D3.9}]:
\]

%dsp-level:0
\[
\qquad \text{the coupling alphabet}\ \{g,g',e,\varepsilon_{\mathrm{EM}},\hat\mu,\hat\mu',\ell_c,\theta_W\}\ \textbf{is disjoint from it}\ [\text{machine: the symbol census}];
\]

%dsp-level:0
\[
\boxed{\ \frac{\partial\,\Sigma_{\mathrm B}}{\partial g}=\frac{\partial\,\Sigma_{\mathrm B}}{\partial g'}=\frac{\partial\,\Sigma_{\mathrm B}}{\partial e}=\frac{\partial\,\Sigma_{\mathrm B}}{\partial\varepsilon_{\mathrm{EM}}}=\frac{\partial\,\Sigma_{\mathrm B}}{\partial\hat\mu}=\frac{\partial\,\Sigma_{\mathrm B}}{\partial\hat\mu'}=\frac{\partial\,\Sigma_{\mathrm B}}{\partial\ell_c}=0\ \ \text{exactly}\ }
\]

%dsp-level:0
\[
\qquad [\text{machine}]:
\]

%dsp-level:0
\[
\textbf{the charge-class transport entropy is coupling-blind};
\]

%dsp-level:0
\[
\qquad \text{it is a function of the integer alphabet and the constituent count alone};
\]

%dsp-level:0
\[
(6)\quad\textbf{the exact object at the proposed slot:}
\]

%dsp-level:0
\[
\qquad \boxed{\ \Sigma_{\mathrm B}=\ln\frac{W(0)}{W(+1)}=\ln\frac43=2\ln2-\ln3\ }
\]

%dsp-level:0
\[
\qquad [\mathrm{Def\text{-}BM.13}(4);\ \text{machine}]:
\]

%dsp-level:0
\[
e^{\Sigma_{\mathrm B}}=\frac43
\]

%dsp-level:0
\[
\qquad [\text{machine: the integer row}\ 20\cdot3=15\cdot4];
\]

%dsp-level:0
\[
\qquad [\Sigma_{\mathrm B}]=\mathfrak q^0\ [\mathrm{CT2.1};\ \text{machine}];
\]

%dsp-level:0
\[
\qquad \Sigma_{\mathrm B}>0\ [\text{machine:}\ 4>3];
\]

%dsp-level:0
\[
\textbf{the exact enclosures inside the}\ \ln\textbf{-ring:}
\]

%dsp-level:0
\[
\qquad \frac{\ln2}3<\Sigma_{\mathrm B}<\frac{\ln2}2
\]

%dsp-level:0
\[
\qquad [\text{machine:}\ 4^3>3^3\cdot2\iff64>54;\quad4^2<3^2\cdot2\iff16<18];
\]

%dsp-level:0
\[
\frac{2\ln2}5<\Sigma_{\mathrm B}<\frac{3\ln2}5\qquad[\text{machine:}\ 4^5>3^5\cdot2^2\iff1024>972;\quad4^5<3^5\cdot2^3\iff1024<1944];
\]

%dsp-level:0
\[
(7)\quad\textbf{the verdict:}
\]

%dsp-level:0
\[
\qquad \Sigma_{\mathrm B}\ \text{enters the mass as a}\ \textbf{weight in the rest slot}\ [\mathrm{Def\text{-}BM.13}(5)],\ \neg\text{as a coupling strength:}
\]

%dsp-level:0
\[
\qquad \text{it multiplies no generator, resolves no propagator, and appears in no vertex};
\]

%dsp-level:0
\[
\textbf{it is}\ \neg\alpha_B:\qquad\text{a coupling is a coefficient of an interaction;}\ \Sigma_{\mathrm B}\ \text{is the logarithm of a census};
\]

%dsp-level:0
\[
(8)\quad\textbf{the obstruction, named:}\qquad\boxed{\ \alpha_B=f(\alpha_{\mathrm{EM}},\alpha_{\mathrm{Strong}})\ \ \textbf{is}\ \neg\textbf{derivable}\ }\qquad\text{on four independent grounds:}
\]

%dsp-level:0
\[
(\mathrm O1)\ \ \textbf{the strength bar:}\ \neg\exists\ \text{a count}\to\text{strength theorem}\ [(3);\ \mathrm{BM.36}(3)];
\]

%dsp-level:0
\[
\qquad (\mathrm O2)\ \ \textbf{the colour-coupling void:}\ \alpha_{\mathrm{Strong}}\ \text{is}\ \neg\text{an inventory item, so the domain of}\ f\ \text{is empty in its second argument}\ [(2)];
\]

%dsp-level:0
\[
(\mathrm O3)\ \ \textbf{the coupling-blindness:}\ \text{the one exact dimensionless charge-class quantity the corpus derives has zero derivative in every coupling, so it is not a function of them}\ [(5)];
\]

%dsp-level:0
\[
\qquad (\mathrm O4)\ \ \textbf{P5:}\ \text{fixing}\ \alpha_B\ \text{by an observed ratio is barred}\ [\mathrm{P5};\ \mathrm{BM.36}(3);\ \mathrm{FT.73}(5)]:
\]

%dsp-level:0
\[
\textbf{the obstruction is not a gap in the arithmetic; it is the absence of a coupling theorem for colour and the bar against inventing one};
\]

%dsp-level:0
\[
(9)\quad\textbf{status:}
\]

%dsp-level:0
\[
\qquad \textbf{derived:}\ (1)\text{--}(7);
\]

%dsp-level:0
\[
\qquad \textbf{named, not filled:}\ \alpha_B\ (8);
\]

%dsp-level:0
\[
\qquad \alpha_B\ \textbf{is}\ \neg\textbf{added as a P4-class parameter};
\]

%dsp-level:0
\[
\qquad \text{§Cal-BM stays empty};
\]

%dsp-level:0
\[
\neg\text{a new dimensionful scale}\ [\mathrm{P1}];
\]

%dsp-level:0
\[
\qquad \neg\text{a measured value and}\ \neg\text{a numerical coupling}\ [\mathrm{P5};\ \text{machine: zero external tokens, zero decimal numerals in this tier}];\qquad \qquad
\]

\textbf{Req.} D2.9, D2.14, D3.9, CT3.8, CT3.17a(2), EW.33, EW.35, EM.9, EM.17, EM.25(6), EM.25(7), EM.25(8), EM.25(9) (Part 2); CT0.13, CT0.14, CT2.1, P1, P4, P5 (Parts 1–2); FT.41(b), FT.73(5), FT.118(2), FT.119(1) (Part 3); Def-BM.8(1), Def-BM.8(6), Def-BM.12, Def-BM.13(3), Def-BM.13(4), Def-BM.13(5), BM.36(1), BM.36(2), BM.36(3) (this part); A §E (the P4 core); Part 2 §Cal-2 (the logged row); the exact grammar-count summation, symbol censuses, exact $\ln$-ring identities with integer certificates, Laurent exponent arithmetic [math; machine: \texttt{checks/\allowbreak{}verify\_\allowbreak{}H\_\allowbreak{}constituent\_\allowbreak{}transport\_\allowbreak{}splitting.py} §E].

\textbf{Pf.}

%dsp-level:0
\[
(1)\text{--}(2):\ \text{the token scans over the five canon files and the two annexes}\ [\text{machine}];
\]

%dsp-level:0
\[
\qquad g(N)\ \text{recomputed by the}\ \mathrm{CT3.17a}(2)\ \text{summation for}\ N
\]

%dsp-level:0
\[
= 1,
\]

%dsp-level:0
\[
\dots,
\]

%dsp-level:0
\[
7\ \text{and compared to the frozen row}\ (0,1,1,3,10,25,105)\ [\text{machine}];
\]

%dsp-level:0
\[
\qquad g(6)=15+10=25\ [\text{machine: the}\ a{=}3,b{=}0\ \text{and}\ a{=}0,b{=}2\ \text{summands}]
\]

%dsp-level:0
\[
(3)\text{--}(4):\ \mathrm{BM.36}(2),(3),\ \mathrm{EM.25}(6)\text{--}(8)\ \text{re-cited};
\]

%dsp-level:0
\[
\qquad \tfrac{1}{3},
\]

%dsp-level:0
\[
\ \tfrac81,
\]

%dsp-level:0
\[
\ \tfrac1{25}\ \text{from}\ (6,16,17,25):\ 16/3\to\tfrac13\ \text{via}\ \dim F_W
\]

%dsp-level:0
\[
= 3,
\]

%dsp-level:0
\[
\ 8/1,
\]

%dsp-level:0
\[
\ 1/25\ [\text{machine: exact Fractions}]
\]

%dsp-level:0
\[
(5):\ W(Q)\ \text{written as a Laurent polynomial with integer coefficients; the symbol set extracted and intersected with the coupling alphabet}\ [\text{machine: the intersection is empty}];
\]

%dsp-level:0
\[
\qquad \text{each partial derivative formed symbolically and evaluated}\ [\text{machine: all seven are the zero polynomial}]
\]

%dsp-level:0
\[
(6):\ \ln20-\ln15=\ln\tfrac{20}{15}=\ln\tfrac43=2\ln2-\ln3\ [\text{machine: the prime-log ring}];
\]

%dsp-level:0
\[
\qquad 20\cdot3=60=15\cdot4\ [\text{machine}];
\]

%dsp-level:0
\[
\qquad \text{the four enclosures:}\ 4^3
\]

%dsp-level:0
\[
= 64>54
\]

%dsp-level:0
\[
= 3^3\cdot2;
\]

%dsp-level:0
\[
\ 4^2
\]

%dsp-level:0
\[
= 16<18
\]

%dsp-level:0
\[
= 3^2\cdot2;
\]

%dsp-level:0
\[
\ 4^5
\]

%dsp-level:0
\[
= 1024>972
\]

%dsp-level:0
\[
= 3^5\cdot2^2;
\]

%dsp-level:0
\[
\ 4^5
\]

%dsp-level:0
\[
= 1024<1944
\]

%dsp-level:0
\[
= 3^5\cdot2^3\ [\text{machine: integer comparisons}]
\]

%dsp-level:0
\[
(7)\text{--}(9):\ \text{the locus row re-cited from}\ \mathrm{Def\text{-}BM.13}(5);
\]

%dsp-level:0
\[
\ \text{the four grounds conjoined};
\]

%dsp-level:0
\[
\ \text{the standing-rule scans}\ [\text{machine}]. \ \blacksquare
\]

\textbf{Phys.} The proposal was to close the neutron row with a new coupling, fixed by the ratio the world reports. Three of the four things that proposal needs are absent, and their absence is proved rather than assumed. There is no $\alpha_B$ anywhere in the canon to refine. There is no strong coupling for a function of two strengths to take as its second argument: the colour sector's entire strength-class inventory is a dimension, a projector, a grammar count and two arena sizes, all integers, and the corpus has already ruled that turning a count into a strength is the same category error as turning a geometric ratio into a coupling. And the one exact dimensionless charge-class number the corpus does derive — the logarithm of twenty over fifteen — has a zero derivative in every coupling it owns, which settles the question from the other side: it is not a function of the strengths, so it cannot be dressed as one. What $\ln(4/3)$ actually is, is a weight: it sits in the rest slot beside the edge count, in the same nats, and it does the work the proposal wanted without the parameter the proposal needed. The obstruction is named to the monomial, and no free entry is added.

\end{theorem}

\pseudohead{theorem}{\texorpdfstring{BM.62}{BM.62}}{\texorpdfstring{The full mass splitting: the extended decomposition, the exact critical strength, and the sign theorem}{The full mass splitting: the extended decomposition, the exact critical strength, and the sign theorem}}\label{tag:BM-62}
\textbf{Thm.}

%dsp-level:0
\[
(1)\quad\textbf{the extended decomposition:}
\]

%dsp-level:0
\[
\qquad \boxed{\ \Delta m\ :=\ m_n-m_p\ =\ \Delta m_{\mathrm{Lyapunov}}+\Delta m_{\mathrm{force}}+\Delta m_{\mathrm{metric}}+\Delta m_{\mathrm{binding}}+\Delta m_{\mathrm{self\text{-}int}}+\Delta m_{\mathrm{transport}}\ }
\]

%dsp-level:0
\[
\text{with}\ \mathrm{BM.55}(1)\text{'s five derived rows standing unmodified and one new row added}\ [\mathrm{BM.52}(5)\text{'s no-double-count guard respected:}\ \mathrm{BM.59}(8)\text{'s disjointness}]:
\]

%dsp-level:0
\[
(2)\quad\textbf{the four zeros:}
\]

%dsp-level:0
\[
\qquad \Delta m_{\mathrm{Lyapunov}}=0\ [\mathrm{BM.51}(6):\ \text{every input is class-shared}];
\]

%dsp-level:0
\[
\qquad \Delta m_{\mathrm{metric}}=0\ [\mathrm{BM.53}(2);\ \mathrm{BM.58}(4):\ \text{the charge exponent is}\ 0\ \text{in all seven monomials}];
\]

%dsp-level:0
\[
\Delta m_{\mathrm{binding}}
\]

%dsp-level:0
\[
= 0\ [\mathrm{BM.54}(1)\text{--}(3):\ \text{the count read vanishes as}\ 6-(2{+}2{+}2)=0,\ \text{and the ensemble read is}\ q\text{-degenerate}];
\]

%dsp-level:0
\[
\textbf{re-derived here:}
\]

%dsp-level:0
\[
\qquad \mathcal B(S_j;1,1,1)
\]

%dsp-level:0
\[
= m^2_B(\cdot,S_j)-3\,
\]

%dsp-level:0
\[
m^2_B(\cdot,1)\ \ \text{is a function of the junction sector alone}\ [\mathrm{BM.21}(2);\ \mathrm{BM.20}(4)]:
\]

%dsp-level:0
\[
\mathcal B(0;1,1,1)
\]

%dsp-level:0
\[
= \frac{843957}{98956}-\frac{5764986}{393679},
\]

%dsp-level:0
\[
\quad\mathcal B(1;1,1,1)
\]

%dsp-level:0
\[
= -\frac{3843324}{393679},
\]

%dsp-level:0
\[
\quad\mathcal B(2;1,1,1)
\]

%dsp-level:0
\[
= \frac{91}{23}-\frac{5764986}{393679},
\]

%dsp-level:0
\[
\quad\mathcal B(3;1,1,1)
\]

%dsp-level:0
\[
= \frac{13713}{3088}-\frac{5764986}{393679};
\]

%dsp-level:0
\[
\textbf{all four are nonzero exact rationals, all four are negative, and the class difference of each is exactly}\ 0
\]

%dsp-level:0
\[
\qquad [\text{machine: the exact rational census}];
\]

%dsp-level:0
\[
m^2_4(\cdot,0)>m^2_4(\cdot,1)>m^2_4(\cdot,3)>m^2_4(\cdot,2)
\]

%dsp-level:0
\[
\qquad [\mathrm{BM.21}(4);\ \text{machine}];
\]

%dsp-level:0
\[
\qquad \text{and a binding term in a free member's mass remains a category error}\ [\mathrm{BM.48}(4);\ \mathrm{BM.54}(4)];
\]

%dsp-level:0
\[
(3)\quad\textbf{the two nonzero summands:}
\]

%dsp-level:0
\[
\qquad \Delta m_{\mathrm{force}}=\Delta m_{\mathrm{self\text{-}int}}=-\Delta m_{\mathrm{EM}}=-\,\frac{3\,c\,e^2\ln2}{2\pi\,\mathfrak h}=-\,6\,c\,\varepsilon_{\mathrm{EM}}\,\frac{\ln2}{\mathfrak h}
\]

%dsp-level:0
\[
\qquad [\mathrm{BM.52}(4),(5);\ \mathrm{BM.46}(4);\ \text{machine: one Laurent monomial, counted once}];
\]

%dsp-level:0
\[
\Delta m_{\mathrm{transport}}=+\,\frac{\Sigma_{\mathrm B}}{\mathfrak h}=+\,\frac{\ln(4/3)}{\mathfrak h}>0\qquad[\mathrm{BM.59}(4);\ \text{machine:}\ 4>3];
\]

%dsp-level:0
\[
(4)\quad\textbf{the assembly:}
\]

%dsp-level:0
\[
\qquad \boxed{\ \Delta m=\frac{\ln(4/3)-6\,c\,\varepsilon_{\mathrm{EM}}\ln2}{\mathfrak h}=\frac{2\ln2-\ln3-6\,c\,\varepsilon_{\mathrm{EM}}\ln2}{\mathfrak h}\ }
\]

%dsp-level:0
\[
\qquad [\text{machine: the monomial sum in the}\ \ln\text{-ring}]:
\]

%dsp-level:0
\[
m_n\,\mathfrak h=6\ln2+\ln20=\ln1280,
\]

%dsp-level:0
\[
\qquad m_p\,\mathfrak h=6\ln2+\ln15+6\,c\,\varepsilon_{\mathrm{EM}}\ln2=\ln960+6\,c\,\varepsilon_{\mathrm{EM}}\ln2
\]

%dsp-level:0
\[
\qquad [\mathrm{Def\text{-}BM.13}(5);\ \text{machine}];
\]

%dsp-level:0
\[
[\Delta m]=\mathfrak q^{+1/2}=[\text{energy}]
\]

%dsp-level:0
\[
\qquad [\mathrm{CT2.1};\ \mathrm{D2.9};\ \text{machine: exponent arithmetic}];
\]

%dsp-level:0
\[
\qquad \textbf{both summands ride the same rest slot}\ [\mathrm{Def\text{-}BM.10}(4);\ \mathrm{Def\text{-}BM.12}(6);\ \mathrm{BM.37}(4)];
\]

%dsp-level:0
\[
(5)\quad\textbf{the exact ratio:}
\]

%dsp-level:0
\[
\qquad \frac{m_n}{m_p}=\frac{\ln1280}{\ln960+6\,c\,\varepsilon_{\mathrm{EM}}\ln2}
\]

%dsp-level:0
\[
\qquad [\text{machine}];
\]

%dsp-level:0
\[
\qquad \textbf{at}\ \varepsilon_{\mathrm{EM}}=0:\quad\frac{m_n}{m_p}=\frac{\ln1280}{\ln960}=\frac{8\ln2+\ln20}{6\ln2+\ln15}\ \ne\ \frac43
\]

%dsp-level:0
\[
\qquad [\text{machine:}\ 3\ln1280\ne4\ln960]:
\]

%dsp-level:0
\[
\textbf{the amended}\ \mathrm{BM.55}(2)\textbf{ ratio:}
\]

%dsp-level:0
\[
\qquad \frac{4\pi}{4\pi+c\,e^2}\ \longmapsto\ \frac{\ln1280}{\ln960+6\,c\,\varepsilon_{\mathrm{EM}}\ln2}
\]

%dsp-level:0
\[
\qquad [\text{the transport refinement of the rest input}];
\]

%dsp-level:0
\[
(6)\quad\textbf{the exact critical strength:}
\]

%dsp-level:0
\[
\qquad \Delta m\ \text{is strictly decreasing in}\ \varepsilon_{\mathrm{EM}}:
\]

%dsp-level:0
\[
\qquad \frac{\partial\,\Delta m}{\partial\,\varepsilon_{\mathrm{EM}}}=-\,\frac{6\,c\ln2}{\mathfrak h}<0
\]

%dsp-level:0
\[
\qquad [\text{machine: the coefficient census at}\ c,\ln2,\mathfrak h>0]:
\]

%dsp-level:0
\[
\boxed{\ \varepsilon^{\mathrm{crit}}(c)\ :=\ \frac{\Sigma_{\mathrm B}}{6\,c\,\ln2}\ =\ \frac{2\ln2-\ln3}{6\,c\,\ln2}\ =\ \frac{2-\log_23}{6\,c}\ }\qquad\textbf{the unique zero of}\ \Delta m\ [\text{machine}];
\]

%dsp-level:0
\[
\textbf{the sign trichotomy:}
\]

%dsp-level:0
\[
\qquad \varepsilon_{\mathrm{EM}}<\varepsilon^{\mathrm{crit}}(c)\iff\Delta m>0;
\]

%dsp-level:0
\[
\qquad \varepsilon_{\mathrm{EM}}=\varepsilon^{\mathrm{crit}}(c)\iff\Delta m=0;
\]

%dsp-level:0
\[
\qquad \varepsilon_{\mathrm{EM}}>\varepsilon^{\mathrm{crit}}(c)\iff\Delta m<0
\]

%dsp-level:0
\[
\qquad [\text{machine}];
\]

%dsp-level:0
\[
(7)\quad\textbf{the exact enclosures of the critical strength:}
\]

%dsp-level:0
\[
\qquad \text{every comparison is an integer power certificate}\ [\text{machine}]:
\]

%tbl-level:1
{\small\setlength{\tabcolsep}{4pt}
\begin{longtable}{>{\raggedright\arraybackslash}p{\dimexpr0.3009\textwidth-2\tabcolsep\relax}|>{\raggedright\arraybackslash}p{\dimexpr0.1593\textwidth-2\tabcolsep\relax}|>{\raggedright\arraybackslash}p{\dimexpr0.2124\textwidth-2\tabcolsep\relax}|>{\raggedright\arraybackslash}p{\dimexpr0.3274\textwidth-2\tabcolsep\relax}}
\toprule
$c$ & $\varepsilon^{\mathrm{crit}}(c)$ & lower certificate & upper certificate \\
\addlinespace[2pt]
\midrule
\endfirsthead
\toprule
$c$ & $\varepsilon^{\mathrm{crit}}(c)$ & lower certificate & upper certificate \\
\addlinespace[2pt]
\midrule
\endhead
\bottomrule
\endlastfoot
$\tfrac12$ (the shell $[\mathrm{Def\text{-}BM.12}(4)]$) & $\dfrac{2-\log_23}{3}$ & $>\allowbreak \dfrac18\iff\allowbreak 2^{13}>\allowbreak 3^8\iff\allowbreak 8192>\allowbreak 6561$ & $<\allowbreak \dfrac17\iff\allowbreak 2^{11}<\allowbreak 3^7\iff\allowbreak 2048<\allowbreak 2187$ \\
$\tfrac35$ (the uniform ball $[\mathrm{Def\text{-}BM.12}(4)]$) & $\dfrac{5(2-\log_23)}{18}$ & $>\allowbreak \dfrac19\iff\allowbreak 2^{8}>\allowbreak 3^5\iff\allowbreak 256>\allowbreak 243$ & $<\allowbreak \dfrac18\iff\allowbreak 2^{31}<\allowbreak 3^{20}\iff\allowbreak 2147483648<\allowbreak 3486784401$ \\
\end{longtable}
}

%dsp-level:0
\[
\boxed{\ \frac19<\varepsilon^{\mathrm{crit}}\Bigl(\frac35\Bigr)<\frac18<\varepsilon^{\mathrm{crit}}\Bigl(\frac12\Bigr)<\frac17\ }
\]

%dsp-level:0
\[
\qquad [\text{machine: the four certificates conjoined}];
\]

%dsp-level:0
\[
(8)\quad\textbf{the sub-criticality bar (parameter-free):}\qquad\boxed{\ \varepsilon_{\mathrm{EM}}\le\frac19\ \ \Longrightarrow\ \ \Delta m>0\ \ \text{at both priced}\ c\ }
\]

%dsp-level:0
\[
\textbf{the witnesses at}\ \varepsilon_{\mathrm{EM}}=\tfrac19:
\]

%dsp-level:0
\[
\qquad c=\tfrac12:\ \ \ln\tfrac43>\tfrac{\ln2}3\iff4^3>3^3\cdot2\iff64>54;
\]

%dsp-level:0
\[
\qquad c=\tfrac35:\ \ \ln\tfrac43>\tfrac{2\ln2}5\iff4^5>3^5\cdot2^2\iff1024>972
\]

%dsp-level:0
\[
\qquad [\text{machine}];
\]

%dsp-level:0
\[
\text{and at}\ \varepsilon_{\mathrm{EM}}=0:\qquad\Delta m=\frac{\ln(4/3)}{\mathfrak h}>0\iff4>3\qquad[\text{machine}];
\]

%dsp-level:0
\[
(9)\quad\textbf{the ordering theorem:}
\]

%dsp-level:0
\[
\qquad \boxed{\ \text{on the sub-critical branch}\ \varepsilon_{\mathrm{EM}}<\varepsilon^{\mathrm{crit}}(c):\qquad m_n>m_p\ }
\]

%dsp-level:0
\[
\qquad [(6),(8);\ \text{machine}]:
\]

%dsp-level:0
\[
\textbf{the neutron is the heavier baryon};
\]

%dsp-level:0
\[
\qquad \mathrm{BM.55}(3)\text{'s strict order is read-relative and holds on the super-critical branch}\ \varepsilon_{\mathrm{EM}}>\varepsilon^{\mathrm{crit}}(c);
\]

%dsp-level:0
\[
\qquad \text{the critical surface is exact and unique}\ [(6)];
\]

%dsp-level:0
\[
(10)\quad\textbf{the fractional splittings (exact rational enclosures):}
\]

%dsp-level:0
\[
\qquad \frac{\Delta m_{\mathrm{transport}}}{m_{\mathrm{hid}}}=\frac{\ln(4/3)}{6\ln2}\in\Bigl(\frac1{15},\frac1{14}\Bigr)
\]

%dsp-level:0
\[
\qquad [\text{machine:}\ 15\ln\tfrac43>6\ln2\iff2^{24}>3^{15}\iff16777216>14348907;\quad14\ln\tfrac43<6\ln2\iff2^{22}<3^{14}\iff4194304<4782969];
\]

%dsp-level:0
\[
\frac{\Delta m_{\mathrm{transport}}}{m_n}\bigg|_{\varepsilon_{\mathrm{EM}}=0}=\frac{\ln(4/3)}{\ln1280}\in\Bigl(\frac1{25},\frac1{24}\Bigr)
\]

%dsp-level:0
\[
\qquad [\text{machine: integer power certificates}];
\]

%dsp-level:0
\[
\qquad \frac{\Delta m_{\mathrm{EM}}}{m_{\mathrm{hid}}}=c\,\varepsilon_{\mathrm{EM}}\in\mathfrak q^0\ [\mathrm{BM.46}(1);\ \text{machine}];
\]

%dsp-level:0
\[
\textbf{the strict contraction against the barred measure price:}
\]

%dsp-level:0
\[
\qquad \frac{\ln(4/3)}{6\ln2}<\frac13\iff3\ln\tfrac43<2\ln2\iff16<27
\]

%dsp-level:0
\[
\qquad [\mathrm{Def\text{-}BM.13}(8);\ \text{machine}];
\]

%dsp-level:0
\[
(11)\quad\textbf{the species row:}
\]

%dsp-level:0
\[
\qquad \text{the assembly}\ (4)\ \text{is legitimate because all four nonzero-or-zero terms enter through one slot:}
\]

%dsp-level:0
\[
\text{the rest input of the dispersion read}\ [\mathrm{Def\text{-}BM.10}(4);\ \mathrm{BM.35}(4);\ \mathrm{BM.37}(4)],
\]

%dsp-level:0
\[
\ \text{refined in the}\ \mathrm{Def\text{-}BM.12}(6)\ \text{pattern};
\]

%dsp-level:0
\[
\qquad \neg\text{a route ratio is read as a mass ratio}\ [\mathrm{BM.37}(6);\ \mathrm{Def\text{-}BM.13}(8);\ (5)];
\]

%dsp-level:0
\[
\textbf{the input alphabets:}
\]

%dsp-level:0
\[
\qquad \Delta m_{\mathrm{transport}}:\ \{m_q,W(Q),\ln2,\mathfrak h\};
\]

%dsp-level:0
\[
\qquad \Delta m_{\mathrm{EM}}:\ \{e,c,\Delta x_{\mathrm B}\};
\]

%dsp-level:0
\[
\qquad \Delta m_{\mathrm{binding}}:\ \{(\Gamma,R),\Delta S\};
\]

%dsp-level:0
\[
\qquad \Delta m_{\mathrm{metric}}:\ \{(a_{00},q_{ij},w_i),E_{\mathrm{int}}\};
\]

%dsp-level:0
\[
\text{pairwise disjoint except for the shared}\ \mathfrak h\ \text{and}\ E_{\mathrm{int}}\ [\text{machine: the token census}];
\]

%dsp-level:0
\[
\qquad \mathrm{BM.37}(2)\text{'s disjointness theorem is respected:}\ \neg\text{a measure route is added to a metric read};
\]

%dsp-level:0
\[
(12)\quad\textbf{status:}
\]

%dsp-level:0
\[
\qquad \textbf{derived:}\ (1)\text{--}(11);
\]

%dsp-level:0
\[
\qquad \textbf{the sign theorem is conditional on the sub-critical branch, with the hypothesis on its face and the critical value exact}\ (6)\text{--}(9)\ [\text{A §0.2a}];
\]

%dsp-level:0
\[
\neg\text{a P4 addition},
\]

%dsp-level:0
\[
\quad\neg\text{a new dimensionful scale}\ [\mathrm{P1}],
\]

%dsp-level:0
\[
\quad\neg\text{a route adopted}\ [\mathrm{BM.30}(5)],
\]

%dsp-level:0
\[
\quad\neg\text{a measured value}\ [\mathrm{P5}];
\]

%dsp-level:0
\[
\qquad \textbf{the open row:}\ \text{the numerical value of}\ e\ [\mathrm{EM.25}(7):\ \text{one calibration datum}]\ —\ \text{it decides which side of}\ \varepsilon^{\mathrm{crit}}\ \text{the corpus sits on, and no theorem of Parts 1--4 supplies it}.
\]

\textbf{Req.} D2.9, CT2.1, EM.25(7), EW.35 (Part 2); P1, P4, P5 (Parts 1–2); FT.79(5), FT.107(2), FT.107(3) (Part 3); Census FT.35(2) (row-wise); Def-BM.4, Def-BM.8(6), Def-BM.10(4), Def-BM.12(4), Def-BM.12(6), Def-BM.13(4), Def-BM.13(5), Def-BM.13(8), BM.20(4), BM.21(2), BM.21(4), BM.30(5), BM.35(4), BM.37(2), BM.37(4), BM.37(6), BM.46(1), BM.46(4), BM.48(4), BM.51(6), BM.52(4), BM.52(5), BM.53(2), BM.54(1)–(4), BM.55(1), BM.55(2), BM.55(3), BM.58(4), BM.59(4), BM.59(8) (this part); exact $\ln$-ring arithmetic with integer power certificates, Laurent-monomial assembly, exact rational sign censuses, exact rational enclosures of $\ln$-ratios [math; machine: \texttt{checks/\allowbreak{}verify\_\allowbreak{}H\_\allowbreak{}constituent\_\allowbreak{}transport\_\allowbreak{}splitting.py} §F].

\textbf{Pf.}

%dsp-level:0
\[
(1)\text{--}(3):\ \mathrm{BM.55}(1)\text{'s five rows re-cited (frozen)};
\]

%dsp-level:0
\[
\qquad \mathcal B(S_j;1,1,1)
\]

%dsp-level:0
\[
= m^2_B(\cdot,S_j)-3m^2_B(\cdot,1)\ \text{formed from}\ \mathrm{BM.20}(4)\text{'s four sector masses and compared to}\ \mathrm{BM.21}(2)\text{'s four frozen rationals}\ [\text{machine: exact Fractions, all four agree}];
\]

%dsp-level:0
\[
\ \text{the sign census and the sector order}\ [\text{machine}];
\]

%dsp-level:0
\[
\ \text{the class difference of each}\
\]

%dsp-level:0
\[
= 0\ \text{because}\ m^2_B(q,S)\ \text{is}\ q\text{-independent}\ [\mathrm{FT.107}(3);\ \text{machine}]
\]

%dsp-level:0
\[
(4):\ \ln1280-\ln960=\ln\tfrac{1280}{960}=\ln\tfrac43=\Sigma_{\mathrm B}\ [\text{machine: the prime-log ring}];
\]

%dsp-level:0
\[
\qquad 6\ln2+\ln20=8\ln2+\ln20=\ln(2^8\cdot20)=\ln1280\ [\text{machine}];
\]

%dsp-level:0
\[
\qquad 6\ln2+\ln15=\ln(2^6\cdot15)=\ln960\ [\text{machine}];
\]

%dsp-level:0
\[
\qquad \text{the exponent row}\ \mathfrak q^0\cdot\mathfrak q^{+1/2}\ [\text{machine}]
\]

%dsp-level:0
\[
(5):\ 3\ln1280\ne4\ln960\iff1280^3\ne960^4\ [\text{machine: the two integers differ}];
\]

%dsp-level:0
\[
\qquad (6):\ \Delta m
\]

%dsp-level:0
\[
= 0\iff\Sigma_{\mathrm B}
\]

%dsp-level:0
\[
= 6c\varepsilon_{\mathrm{EM}}\ln2\iff\varepsilon_{\mathrm{EM}}
\]

%dsp-level:0
\[
= \Sigma_{\mathrm B}/(6c\ln2)\ [\text{machine: solved in the}\ \ln\text{-ring}];
\]

%dsp-level:0
\[
\ \text{the derivative}\ -6c\ln2/\mathfrak h\ [\text{machine}]
\]

%dsp-level:0
\[
(7):\ \text{cross-multiplication by}\ 6c\ln2>0:\quad\varepsilon^{\mathrm{crit}}(c)>r\iff\Sigma_{\mathrm B}>6cr\ln2;
\]

%dsp-level:0
\[
\qquad c=\tfrac12,\ r=\tfrac18:\ 8\Sigma_{\mathrm B}>3\ln2\iff(4/3)^8>2^3\iff2^{13}>3^8\ [\text{machine}];
\]

%dsp-level:0
\[
\qquad c=\tfrac12,\ r=\tfrac17:\ 7\Sigma_{\mathrm B}<3\ln2\iff2^{11}<3^7\ [\text{machine}]
\]

%dsp-level:0
\[
c=\tfrac35,\ r=\tfrac19:\ 5\Sigma_{\mathrm B}>2\ln2\iff2^8>3^5\ [\text{machine}];
\]

%dsp-level:0
\[
\qquad c=\tfrac35,\ r=\tfrac18:\ 40\Sigma_{\mathrm B}<18\ln2\iff2^{31}<3^{20}\ [\text{machine}]
\]

%dsp-level:0
\[
(8):\ \text{the two witnesses}\ 4^3=64>54=3^3\cdot2\ \text{and}\ 4^5=1024>972=3^5\cdot2^2\ [\text{machine}];
\]

%dsp-level:0
\[
\qquad 4>3\ [\text{machine}];
\]

%dsp-level:0
\[
\qquad (9):\ \text{conjunction of}\ (6),(8)
\]

%dsp-level:0
\[
(10):\ 15\Sigma_{\mathrm B}>6\ln2\iff4^{15}>3^{15}2^6\iff2^{24}>3^{15}\ [\text{machine}];
\]

%dsp-level:0
\[
\qquad 14\Sigma_{\mathrm B}<6\ln2\iff4^{14}<3^{14}2^6\iff2^{22}<3^{14}\ [\text{machine}];
\]

%dsp-level:0
\[
\qquad \text{the}\ \ln1280\ \text{enclosure by the same route}\ [\text{machine}];
\]

%dsp-level:0
\[
\qquad 3\ln\tfrac43<2\ln2\iff64<108\ [\text{machine}]
\]

%dsp-level:0
\[
(11)\text{--}(12):\ \text{the token census of the four alphabets}\ [\text{machine}];
\]

%dsp-level:0
\[
\ \text{the standing-rule scans}\ [\text{machine}]. \ \blacksquare
\]

\textbf{Phys.} The decomposition the full-scope tier closed with four zeros and one monomial gains a fifth zero and a second monomial, and the second one points the other way. Nothing in the amended assembly contradicts the earlier one: the Lyapunov tower still shares every input, the geometry still carries no charge exponent, the binding still dies twice, and the scar's field energy is still counted exactly once. What is added is the transport register — twenty admissible end-sign assignments in the neutral class against fifteen in the charged one — read as an entropy in the ledger's own nats and booked through the same rest slot the field energy already uses. The two summands are commensurate, so the splitting becomes a single exact expression with a single exact zero: the critical strength is the logarithm of four thirds over six $c\ln2$, enclosed between one ninth and one seventh by four integer power comparisons, with the tighter of the two priced distributions clearing one ninth. Below that surface the neutron is heavier; above it the earlier order returns; on the surface the two rows are degenerate again. The corpus fixes no value for the coupling, so which branch it sits on is the one open row — and the branch that gives the world's order is not a thin one: it is an exact interval, certified by $256>243$, containing every coupling that keeps the field energy a correction rather than a competitor.

\begin{theorem}[{\texorpdfstring{BM.63 — The decay re-adjudication on the refined read: the closure survives, with an exact margin}{BM.63 - The decay re-adjudication on the refined read: the closure survives, with an exact margin}}]\label{tag:BM-63}
\textbf{Thm.}

%dsp-level:0
\[
(1)\quad\textbf{the amended Q-value:}
\]

%dsp-level:0
\[
\qquad Q_\beta:=m_n-m_p-m_\ell-m_\nu=\Delta m-m_\ell-m_\nu
\]

%dsp-level:0
\[
\qquad [\mathrm{BM.41}(4);\ \mathrm{BM.49}(1);\ \mathrm{BM.55}(5)\ \text{re-cited}]:
\]

%dsp-level:0
\[
\boxed{\ Q_\beta=\frac{\ln(4/3)-6\,c\,\varepsilon_{\mathrm{EM}}\ln2}{\mathfrak h}-m_\ell-m_\nu\ }\qquad[\mathrm{BM.62}(4);\ \text{machine}];
\]

%dsp-level:0
\[
(2)\quad\textbf{the unconditional closure:}
\]

%dsp-level:0
\[
\qquad \Delta m<\frac{\ln(4/3)}{\mathfrak h}\ \ \text{at every}\ \varepsilon_{\mathrm{EM}}>0\ [\mathrm{BM.62}(6):\ \text{strict monotonicity}];
\]

%dsp-level:0
\[
m_\ell\ \ge\ \frac{2\ln2}{\mathfrak h}
\]

%dsp-level:0
\[
\qquad [\mathrm{FT.44}(4):\ m=E_{\mathrm{int}}\ln2/\mathfrak h;\ \mathrm{FT.28}\ \text{the minimal floor}\ E_{\mathrm{int}}=2;\ \mathrm{BM.13}(3):\ E_{\mathrm{int}}\ge2;\ \tfrac12<\ln2<1\ \text{enclosure}];
\]

%dsp-level:0
\[
\qquad m_\nu\ge0\ [\mathrm{EW.33}:\ \text{PSD}];
\]

%dsp-level:0
\[
\implies\boxed{\ Q_\beta\ <\ \frac{\ln(4/3)-2\ln2}{\mathfrak h}\ =\ -\,\frac{\ln3}{\mathfrak h}\ <\ 0\ \ \text{strictly}\ }
\]

%dsp-level:0
\[
\qquad [\text{machine:}\ 2\ln2-\ln\tfrac43=\ln3;\quad\ln\tfrac43<2\ln2\iff4<12]:
\]

%dsp-level:0
\[
\textbf{the}\ \beta\textbf{-decay vertex is kinematically closed at every admissible}\ \varepsilon_{\mathrm{EM}}>0,
\]

%dsp-level:0
\[
\ \text{on both sub- and super-critical branches};
\]

%dsp-level:0
\[
(3)\quad\textbf{the exact margin:}
\]

%dsp-level:0
\[
\qquad \boxed{\ |Q_\beta|>\frac{\ln3}{\mathfrak h}\ }
\]

%dsp-level:0
\[
\qquad [\text{machine: the certificate}\ 3>1]:
\]

%dsp-level:0
\[
\qquad \textbf{the closure margin is one exact}\ \ln3\ \text{of ledger content};
\]

%dsp-level:0
\[
\textbf{the reason, in one comparison:}
\]

%dsp-level:0
\[
\qquad \text{the transport excess is}\ \ln\tfrac43\ \text{nats and the minimal daughter costs}\ 2\ln2\ \text{nats}:
\]

%dsp-level:0
\[
\qquad \ln\tfrac43<2\ln2\iff\tfrac43<4\iff4<12
\]

%dsp-level:0
\[
\qquad [\text{machine}]:
\]

%dsp-level:0
\[
\textbf{a fraction of one bit cannot pay for two bits};
\]

%dsp-level:0
\[
(4)\quad\textbf{the deepening row:}
\]

%dsp-level:0
\[
\qquad \mathrm{BM.49}(1)\text{'s}\ Q_{\mathrm{corr}}=-(\Delta m_{\mathrm{EM}}+m_\ell+m_\nu)\ \text{and}\ \mathrm{BM.55}(5)\text{'s full-scope closure stand;}
\]

%dsp-level:0
\[
\qquad \textbf{the refined margin is}
\]

%dsp-level:0
\[
|Q_\beta|>\frac{\ln3+6\,c\,\varepsilon_{\mathrm{EM}}\ln2}{\mathfrak h}\ >\ \frac{\ln3}{\mathfrak h}
\]

%dsp-level:0
\[
\qquad [\text{machine: the coefficient census}]:
\]

%dsp-level:0
\[
\qquad \text{the EM summand still deepens the closure, and the transport summand does not reverse it};
\]

%dsp-level:0
\[
(5)\quad\textbf{the no-opening theorem survives:}
\]

%dsp-level:0
\[
\qquad \text{every EM correction of}\ \mathrm{Tier\ BM.12}\ \text{is}\ \propto q^2
\]

%dsp-level:0
\[
\ge 0\ \text{(an integral of a square) and raises the daughter's mass against the decay's requirement}\ [\mathrm{BM.49}(2)\ \text{re-cited}];
\]

%dsp-level:0
\[
\textbf{the transport refinement raises the parent's mass by at most}\ \ln(4/3)/\mathfrak h\ [(2)],
\]

%dsp-level:0
\[
\ \textbf{which is strictly below the minimal daughter content}\ [(3)]:
\]

%dsp-level:0
\[
\boxed{\ \text{no term of the extended decomposition opens the vertex}\ }\qquad[\text{machine: the two bounds conjoined}];
\]

%dsp-level:0
\[
(6)\quad\textbf{the}\ k{=}1\ \textbf{scope row (the daughters' own refinement):}
\]

%dsp-level:0
\[
\qquad \mathrm{Def\text{-}BM.13}(3)\text{'s census at}\ k=1\ \text{is}\ W(q)=m_q=(2,1,1)\ [\mathrm{EM.9};\ \text{machine}]:
\]

%dsp-level:0
\[
\Sigma_{\mathrm{cl}}(q)=\ln m_q=(\ln2,\,0,\,0)\ \ \text{at}\ q=(0,+1,-1):
\]

%dsp-level:0
\[
\qquad \text{the neutral line carries one extra bit, the charged lines none};
\]

%dsp-level:0
\[
\textbf{applied to the daughters:}
\]

%dsp-level:0
\[
\quad m_\nu\ \mapsto\ m_\nu+\tfrac{\ln2}{\mathfrak h}\ [\text{the}\ \nu\text{-channel is the EM-unactivated class}\ [\mathrm{BM.41}(3);\ \mathrm{FT.142}(1)]],
\]

%dsp-level:0
\[
\qquad m_\ell\ \mapsto\ m_\ell\ [\text{the charged member carries}\ m_q=1]:
\]

%dsp-level:0
\[
\text{the refinement of the daughters deepens the closure by exactly}\ \ln2/\mathfrak h\ [\text{machine}]:
\]

%dsp-level:0
\[
\qquad \textbf{it is recorded as a scope row and}\ \neg\textbf{needed for}\ (2)\text{'s verdict, which uses only}\ m_\nu\ge0;
\]

%dsp-level:0
\[
(7)\quad\textbf{what inverts and what does not:}
\]

%dsp-level:0
\[
\qquad \textbf{inverts on the sub-critical branch:}\ \text{the ordering clause}\ m_n<m_p\ [\mathrm{BM.55}(3)]\ \longrightarrow\ m_n>m_p\ [\mathrm{BM.62}(9)];
\]

%dsp-level:0
\[
\textbf{does not invert:}\ \text{the kinematic closure}\ [\mathrm{BM.41}(4)(\mathrm e),\ \mathrm{BM.49}(1)\text{--}(2),\ \mathrm{BM.55}(5)];
\]

%dsp-level:0
\[
\qquad \text{the vertex grammar and the forced members}\ [\mathrm{BM.41}(1),(3)];
\]

%dsp-level:0
\[
\qquad \text{the two-body impossibility}\ [\mathrm{BM.41}(2)];
\]

%dsp-level:0
\[
\text{the conditional three-body kinematics and the scale-invariant}\ E_{\max}:\langle E\rangle
\]

%dsp-level:0
\[
= 2:1\ [\mathrm{BM.41}(5);\ \mathrm{BM.49}(3),(4)];
\]

%dsp-level:0
\[
\qquad \text{the proton--antiproton degeneracy}\ [\mathrm{BM.28}(5);\ \mathrm{BM.30}(4):\ W(+1)=W(-1)\ \text{exactly}];
\]

%dsp-level:0
\[
(8)\quad\textbf{status:}
\]

%dsp-level:0
\[
\qquad \textbf{derived:}\ (1)\text{--}(7);
\]

%dsp-level:0
\[
\qquad \textbf{unconditional:}\ \text{the closure}\ (2),(3),(5);
\]

%dsp-level:0
\[
\qquad \textbf{conditional:}\ \text{the ordering}\ (7),\ \text{on}\ \mathrm{BM.62}(6)\text{'s branch};
\]

%dsp-level:0
\[
\neg\text{a measured lifetime, mass difference or endpoint enters}\ [\mathrm{P5}];
\]

%dsp-level:0
\[
\qquad \neg\text{external seconds}\ [\mathrm{BM.44}(3)];
\]

%dsp-level:0
\[
\qquad \neg\text{a P4 addition};
\]

%dsp-level:0
\[
\qquad \text{the open rows: the numerical value of}\ e,
\]

%dsp-level:0
\[
\ \text{the lepton floors}\ E_{\mathrm{int}}^{(\ell)},
\]

%dsp-level:0
\[
E_{\mathrm{int}}^{(\nu)}\ [\mathrm{BM.43}(2)],
\]

%dsp-level:0
\[
\ \text{the line index within the}\ \nu\text{-triplet}.
\]

\textbf{Req.} D2.9 (Parts 1–2); EW.33, EM.9 (Part 2); P5 (Parts 1–2); Def-H.32(3), FT.28, FT.42(1), FT.44(4), FT.142(1) (Part 3); Census FT.35(2) (row-wise); Def-BM.13(3), Def-BM.13(4), BM.13(3), BM.28(5), BM.30(4), BM.41(1)–(5), BM.43(2), BM.44(3), BM.46(4), BM.49(1)–(4), BM.55(3), BM.55(5), BM.62(4), BM.62(6), BM.62(9) (this part); exact $\ln$-ring identities with integer certificates, exact rational enclosures of $\ln2$, sign censuses [math; machine: \texttt{checks/\allowbreak{}verify\_\allowbreak{}H\_\allowbreak{}constituent\_\allowbreak{}transport\_\allowbreak{}splitting.py} §F].

\textbf{Pf.}

%dsp-level:0
\[
(1):\ \mathrm{BM.62}(4)\text{'s monomial substituted into}\ \mathrm{BM.41}(4)\text{'s definition}\ [\text{machine}]
\]

%dsp-level:0
\[
(2):\ \Delta m\ \text{strictly decreasing in}\ \varepsilon_{\mathrm{EM}}\ \text{with}\ \lim_{\varepsilon\to0^+}\Delta m=\ln(4/3)/\mathfrak h\ [\mathrm{BM.62}(6);\ \text{machine}];
\]

%dsp-level:0
\[
\qquad m_\ell\ge2\ln2/\mathfrak h\ \text{from}\ E_{\mathrm{int}}\ge2\ \text{and}\ \ln2>0\ [\text{machine}];
\]

%dsp-level:0
\[
\qquad m_\nu\ge0\ [\mathrm{EW.33}\ \text{PSD}]
\]

%dsp-level:0
\[
2\ln2-\ln\tfrac43=2\ln2-(2\ln2-\ln3)=\ln3\ [\text{machine: the prime-log ring}];
\]

%dsp-level:0
\[
\qquad \ln\tfrac43<2\ln2\iff\tfrac43<4\iff4<12\ [\text{machine}];
\]

%dsp-level:0
\[
\qquad \ln3>0\iff3>1\ [\text{machine}]
\]

%dsp-level:0
\[
(3):\ |Q_\beta|>-\bigl(\ln\tfrac43-2\ln2\bigr)/\mathfrak h=\ln3/\mathfrak h\ [\text{machine}];
\]

%dsp-level:0
\[
\qquad (4):\ \text{the coefficient census at}\ c,\varepsilon_{\mathrm{EM}},\ln2>0\ [\text{machine}]
\]

%dsp-level:0
\[
(5):\ \mathrm{BM.49}(2)\text{'s square-integral row re-cited};
\]

%dsp-level:0
\[
\qquad \ln\tfrac43<2\ln2\ [\text{machine}]\implies\text{the parent's excess cannot reach the daughter's floor}\ [\text{machine}]
\]

%dsp-level:0
\[
(6):\ W(q)
\]

%dsp-level:0
\[
= m_q\ \text{at}\ k
\]

%dsp-level:0
\[
= 1:\ (2+z+z^{-1})^1\ \text{has coefficients}\ (1,2,1)\ \text{at}\ Q
\]

%dsp-level:0
\[
= (-1,0,+1)\ [\text{machine: the convolution}];
\]

%dsp-level:0
\[
\qquad \ln m_q=(0,\ln2,0)\ \text{at}\ q=(\pm1,0)\ [\text{machine}]
\]

%dsp-level:0
\[
(7)\text{--}(8):\ \text{the clause-by-clause scan of}\ \mathrm{BM.41},
\]

%dsp-level:0
\[
\ \mathrm{BM.49},
\]

%dsp-level:0
\[
\ \mathrm{BM.55}\ \text{against the refined rest read}\ [\text{machine: the row census}]. \ \blacksquare
\]

\textbf{Phys.} Making the neutron heavier does not open the decay, and the reason is a comparison of two exact integers. The transport refinement lifts the parent by at most the logarithm of four thirds — a fraction of one bit — while the cheapest possible charged daughter costs two bits of ledger content, and the neutral member costs nothing less than zero. Four against twelve: the excess cannot reach the floor, and the closure holds at every coupling on both sides of the critical surface, with an exact margin of one $\ln3$. Read at the level of a single chain, the refinement even helps the closure, because the neutral line is the one that carries the extra transport bit and the neutrino is the neutral member. So the tier reverses an ordering and leaves a verdict intact: the neutron becomes the heavier baryon on the sub-critical branch, and the beta vertex stays shut — not because the two rows are degenerate, which they are no longer, but because the gap between them is smaller than the lightest thing the decay could emit.

\end{theorem}

\medskip\noindent\textit{\texorpdfstring{BM.64 — The status of the constituent transport splitting: adjudication and scope}{BM.64 - The status of the constituent transport splitting: adjudication and scope}.}\ Structural-constraints record: see Appendix B, Section B.1.23 of the Addendum.\par\medskip

\phantomsection
\section[Tier BM.15 — The Organic Bridge: the Operator EM Lift, the…]{\texorpdfstring{Tier BM.15 — The Organic Bridge: the Operator EM Lift, the Derived-Holonomy Attempt, the Constrained Joint Entropy, and the Beta-Decay Energy Audit}{Tier BM.15 - The Organic Bridge: the Operator EM Lift, the Derived-Holonomy Attempt, the Constrained Joint Entropy, and the Beta-Decay Energy Audit}}

\textbf{Locus.} This tier consumes the Tier BM.14 baseline unmodified and runs four constructions and one audit to exhaustion. The EM correction of the refined rest input is lifted from its scalar class row to an operator identity on the force spectrum ($\mathrm{BM.65}$, $\mathrm{BM.66}$). The full real holonomy is attempted as a lift from the edge-phase register, the force decomposition and the channel caps, and where the lift fails the exact obstruction is named and recorded ($\mathrm{BM.67}$). The constrained joint entropy is derived as a fiber cardinality, its geometric factor derived at stationarity and its general fiber named where the enumeration is missing ($\mathrm{BM.68}$). The three bridge inputs are converged and Gap 7 is reduced to its smallest priced residue ($\mathrm{BM.69}$). The beta-decay daughter floor is audited, the formed-closed dichotomy is proved, and every candidate opening term is adjudicated derived, priced or barred ($\mathrm{BM.70}$, $\mathrm{BM.71}$); $\mathrm{BM.72}$ adjudicates the tier.

\textbf{Standing rows.} The refutations of $\mathrm{BM.57}(8)$ and $\mathrm{BM.58}(4)$ stand closed and are $\neg$reopened as mass-splitting hypotheses; their residual content is admitted at its own loci — the $A_3$ continuation asymmetry generates the family triplet $[\mathrm{BM.57}(9);\ \mathrm{BM.60}(1)]$, the colour/EM seal asymmetry survives as a mask statement $[\mathrm{BM.58}(3),(5)]$, and the unique surviving class-dependent mass candidate is $W(Q)$ read through $\ln W(Q)$ $[\mathrm{BM.58}(6);\ \mathrm{BM.59}(2),(3)]$. $\alpha_B$ stays barred on $\mathrm{BM.61}(8)$'s four grounds. Holonomy is $\neg$declared here; joint entropy is $\neg$postulated here; no external neutron/proton ratio and no external $Q$-value enters anywhere; no Yukawa-type object exists or is added $[\mathrm{EW.15};\ \mathrm{FT.145}(3)]$. The tier adds no definition: every new symbol is an in-theorem abbreviation of derived objects.

\pseudohead{theorem}{\texorpdfstring{BM.65}{BM.65}}{\texorpdfstring{The baseline re-certified: the refined rest input, the exact branch, and the integer certificates}{The baseline re-certified: the refined rest input, the exact branch, and the integer certificates}}\label{tag:BM-65}
\textbf{Thm.}

%dsp-level:0
\[
(1)\quad\textbf{the refined rest input (frozen, re-certified):}
\]

%dsp-level:0
\[
\qquad \boxed{\ m_{\mathrm{rest}}(Q)\,\mathfrak h\ =\ E_{\mathrm{int}}\ln2\ +\ \ln W(Q)\ +\ 6\,c\,\varepsilon_{\mathrm{EM}}\,Q^2\ln2\ }
\]

%dsp-level:0
\[
\qquad [\mathrm{Def\text{-}BM.13}(5);\ \mathrm{Def\text{-}BM.12}(6);\ \mathrm{BM.46}(1),(3);\ \mathrm{BM.35}(4)]:
\]

%dsp-level:0
\[
E_{\mathrm{int}}=6
\]

%dsp-level:0
\[
\qquad [\mathrm{BM.17};\ \mathrm{BM.26};\ \text{Census FT.35}(2)];
\]

%dsp-level:0
\[
\qquad c\in\bigl\{\tfrac12,\tfrac35\bigr\}\ \text{priced, carried symbolically, adopted in none}
\]

%dsp-level:0
\[
\qquad [\mathrm{Def\text{-}BM.12}(4)];
\]

%dsp-level:0
\[
W(0)=20,
\]

%dsp-level:0
\[
\qquad W(\pm1)=15,
\]

%dsp-level:0
\[
\qquad \sum_{Q}W(Q)=4^3=64
\]

%dsp-level:0
\[
\qquad [\mathrm{Def\text{-}BM.13}(3);\ \mathrm{Def\text{-}BM.8}(6);\ \mathrm{BM.28}(2);\ \text{machine: the exhaustive}\ 4^3\ \text{end-sign census and the convolution agree}];
\]

%dsp-level:0
\[
(2)\quad\textbf{the class rows (frozen):}
\]

%dsp-level:0
\[
\qquad m_n\,\mathfrak h=\ln1280,
\]

%dsp-level:0
\[
\qquad m_p\,\mathfrak h=\ln960+6\,c\,\varepsilon_{\mathrm{EM}}\ln2,
\]

%dsp-level:0
\[
\qquad (m_n-m_p)\,\mathfrak h=\ln\tfrac43-6\,c\,\varepsilon_{\mathrm{EM}}\ln2=\Sigma_{\mathrm B}-6\,c\,\varepsilon_{\mathrm{EM}}\ln2
\]

%dsp-level:0
\[
\qquad [\mathrm{BM.62}(4);\ \text{machine}]:
\]

%dsp-level:0
\[
2^6\!\cdot\!20=1280,
\]

%dsp-level:0
\[
\qquad 2^6\!\cdot\!15=960,
\]

%dsp-level:0
\[
\qquad \ln20-\ln15=\ln\tfrac43=2\ln2-\ln3>0
\]

%dsp-level:0
\[
\qquad [\text{machine: the prime-log ring};\ \text{the certificate}\ 4>3];
\]

%dsp-level:0
\[
(3)\quad\textbf{the critical strength and the trichotomy (frozen):}
\]

%dsp-level:0
\[
\qquad \boxed{\ \varepsilon_{\mathrm{crit}}(c)\ :=\ \frac{\ln(4/3)}{6\,c\,\ln2}\ =\ \frac{\Sigma_{\mathrm B}}{6\,c\,\ln2}\ =\ \frac{2-\log_23}{6\,c}\ }
\]

%dsp-level:0
\[
\qquad [\mathrm{BM.62}(6);\ \text{machine: the unique zero}]:
\]

%dsp-level:0
\[
\varepsilon_{\mathrm{crit}}\!\left(\tfrac12\right)=\frac{2-\log_23}{3},\qquad\varepsilon_{\mathrm{crit}}\!\left(\tfrac35\right)=\frac{5\,(2-\log_23)}{18}\qquad[\text{machine: exact substitution}];
\]

%dsp-level:0
\[
\varepsilon_{\mathrm{EM}}<\varepsilon_{\mathrm{crit}}(c)\iff\Delta m>0;
\]

%dsp-level:0
\[
\qquad \varepsilon_{\mathrm{EM}}=\varepsilon_{\mathrm{crit}}(c)\iff\Delta m=0;
\]

%dsp-level:0
\[
\qquad \varepsilon_{\mathrm{EM}}>\varepsilon_{\mathrm{crit}}(c)\iff\Delta m<0
\]

%dsp-level:0
\[
\qquad [\mathrm{BM.62}(6);\ \text{machine}];
\]

%dsp-level:0
\[
(4)\quad\textbf{the integer certificates (machine re-certified):}
\]

%dsp-level:0
\[
\qquad \boxed{\ \frac19\ <\ \varepsilon_{\mathrm{crit}}\!\left(\tfrac35\right)\ <\ \frac18\ <\ \varepsilon_{\mathrm{crit}}\!\left(\tfrac12\right)\ <\ \frac17\ }
\]

%tbl-level:0
{\setlength{\tabcolsep}{4pt}
\begin{longtable}{>{\raggedright\arraybackslash}p{\dimexpr0.2329\textwidth-2\tabcolsep\relax}|>{\raggedright\arraybackslash}p{\dimexpr0.3562\textwidth-2\tabcolsep\relax}|>{\raggedright\arraybackslash}p{\dimexpr0.4110\textwidth-2\tabcolsep\relax}}
\toprule
bound & reduction & certificate \\
\addlinespace[2pt]
\midrule
\endfirsthead
\toprule
bound & reduction & certificate \\
\addlinespace[2pt]
\midrule
\endhead
\bottomrule
\endlastfoot
$\varepsilon_{\mathrm{crit}}(\tfrac35)>\allowbreak \tfrac19$ & $5\,\allowbreak \Sigma_{\mathrm B}>\allowbreak 2\ln2\iff\allowbreak 8\ln2>\allowbreak 5\ln3$ & $2^{8}>\allowbreak 3^{5}\iff\allowbreak 256>\allowbreak 243$ \\
$\varepsilon_{\mathrm{crit}}(\tfrac35)<\allowbreak \tfrac18$ & $40\,\allowbreak \Sigma_{\mathrm B}<\allowbreak 18\ln2\iff\allowbreak 62\ln2<\allowbreak 40\ln3$ & $2^{31}<\allowbreak 3^{20}\iff\allowbreak 2147483648<\allowbreak 3486784401$ \\
$\varepsilon_{\mathrm{crit}}(\tfrac12)>\allowbreak \tfrac18$ & $8\,\allowbreak \Sigma_{\mathrm B}>\allowbreak 3\ln2\iff\allowbreak 13\ln2>\allowbreak 8\ln3$ & $2^{13}>\allowbreak 3^{8}\iff\allowbreak 8192>\allowbreak 6561$ \\
$\varepsilon_{\mathrm{crit}}(\tfrac12)<\allowbreak \tfrac17$ & $7\,\allowbreak \Sigma_{\mathrm B}<\allowbreak 3\ln2\iff\allowbreak 11\ln2<\allowbreak 7\ln3$ & $2^{11}<\allowbreak 3^{7}\iff\allowbreak 2048<\allowbreak 2187$ \\
\end{longtable}
}

%dsp-level:0
\[
[\mathrm{BM.62}(7)\ \text{re-certified};\ \text{machine: every row an exact comparison of two non-negative integers}];
\]

%dsp-level:0
\[
(5)\quad\textbf{the branch is non-empty and exact:}
\]

%dsp-level:0
\[
\qquad \varepsilon_{\mathrm{crit}}(c)>0\iff\Sigma_{\mathrm B}>0\iff4>3
\]

%dsp-level:0
\[
\qquad [\text{machine}]:
\]

%dsp-level:0
\[
\qquad \text{the sub-critical branch}\ \bigl(0,\ \varepsilon_{\mathrm{crit}}(c)\bigr)\ \text{is a non-empty exact interval};
\]

%dsp-level:0
\[
\textbf{the parameter-free bar:}
\]

%dsp-level:0
\[
\qquad \varepsilon_{\mathrm{EM}}\le\tfrac19\ \Longrightarrow\ \Delta m>0\ \ \text{at both priced}\ c
\]

%dsp-level:0
\[
\qquad [\mathrm{BM.62}(8);\ \text{machine: the witnesses}\ 4^3=64>54=3^3\!\cdot\!2,\quad4^5=1024>972=3^5\!\cdot\!2^2];
\]

%dsp-level:0
\[
(6)\quad\textbf{no numerical coupling exists:}
\]

%dsp-level:0
\[
\qquad \neg\exists\ \text{a theorem of Parts 1--4 fixing}\ \varepsilon_{\mathrm{EM}}\ \text{numerically}
\]

%dsp-level:0
\[
\qquad [\mathrm{EM.25}(7);\ \mathrm{BM.61}(3),(4);\ \mathrm{BM.62}(12)]:
\]

%dsp-level:0
\[
\textbf{the neutron-heavier result is conditional on the symbolic coupling being sub-critical,}\ \neg\textbf{on an imported physical value}
\]

%dsp-level:0
\[
\qquad [\mathrm{P5};\ \text{the hypothesis stands on the face of}\ \mathrm{BM.62}(9)\ \text{and of this row};\ \text{A §0.2a}];
\]

%dsp-level:0
\[
(7)\quad\textbf{the ratio discipline:}
\]

%dsp-level:0
\[
\qquad \frac{\ln1280}{\ln960}\ne\frac43\iff3\ln1280\ne4\ln960\iff1280^3\ne960^4
\]

%dsp-level:0
\[
\qquad [\text{machine: the two integers differ}]:
\]

%dsp-level:0
\[
\neg\text{a route ratio is read as a mass ratio}\qquad[\mathrm{BM.37}(6);\ \mathrm{Def\text{-}BM.13}(8);\ \mathrm{BM.62}(5)];
\]

%dsp-level:0
\[
(8)\quad\textbf{status:}
\]

%dsp-level:0
\[
\qquad \textbf{re-certified:}\ (1)\text{--}(7);
\]

%dsp-level:0
\[
\qquad \text{the rows of}\ \mathrm{Def\text{-}BM.13}(5),\ \mathrm{BM.62}(4),(6)\text{--}(9),\ \mathrm{BM.63}(2)\ \text{stand unmodified};
\]

%dsp-level:0
\[
\qquad \text{this tag adds no content beyond the machine re-certification};
\]

%dsp-level:0
\[
\neg\text{a P4 addition},\quad\neg\text{a new dimensionful scale}\ [\mathrm{P1}],\quad\neg\text{a measured value}\ [\mathrm{P5}].
\]

\textbf{Req.} CT2.1 (Part 2); EM.25(7) (Part 2); P1, P5 (Parts 1–2); A §0.2a (the hypothesis at the point of use); Census FT.35(2) (row-wise); Def-BM.8(6), Def-BM.12(4), Def-BM.12(6), Def-BM.13(3), Def-BM.13(4), Def-BM.13(5), Def-BM.13(8), BM.17, BM.26, BM.28(2), BM.35(4), BM.37(6), BM.46(1), BM.46(3), BM.61(3), BM.61(4), BM.62(4), BM.62(5), BM.62(6), BM.62(7), BM.62(8), BM.62(9), BM.62(12), BM.63(2) (this part); exact prime-log-ring arithmetic with integer power certificates, exhaustive end-sign censuses [math; machine: \texttt{checks/\allowbreak{}verify\_\allowbreak{}H\_\allowbreak{}organic\_\allowbreak{}bridge\_\allowbreak{}post\_\allowbreak{}bm14.py} §B, §C].

\textbf{Pf.}

%dsp-level:0
\[
(1)\text{--}(2):\ \text{the frozen rows re-derived by machine: the}\ 4^3\ \text{end-sign census and the generating-function convolution re-enumerated and compared class by class};
\]

%dsp-level:0
\[
\qquad 2^6\!\cdot\!20=1280,\ 2^6\!\cdot\!15=960\ [\text{machine: integers}];
\]

%dsp-level:0
\[
\qquad \ln20-\ln15=\ln\tfrac43\ [\text{machine: the prime-log ring}];\ 4>3\ [\text{machine}]
\]

%dsp-level:0
\[
(3):\ \Delta m=0\iff\Sigma_{\mathrm B}=6\,c\,\varepsilon_{\mathrm{EM}}\ln2\ \text{solved in the}\ \ln\text{-ring}\ [\text{machine}];
\]

%dsp-level:0
\[
\qquad \varepsilon_{\mathrm{crit}}\!\left(\tfrac12\right)=\tfrac{\Sigma_{\mathrm B}}{3\ln2},\ \ \varepsilon_{\mathrm{crit}}\!\left(\tfrac35\right)=\tfrac{5\,\Sigma_{\mathrm B}}{18\ln2}\ [\text{machine: exact Fractions}];
\]

%dsp-level:0
\[
\qquad \log_23\ \text{form:}\ \tfrac{\Sigma_{\mathrm B}}{\ln2}=2-\log_23\ [\text{machine: the prime-log ring}]
\]

%dsp-level:0
\[
(4):\ \text{cross-multiplication by}\ 6\,
\]

%dsp-level:0
\[
c\ln2>0\ \text{reduces each bound to}\ a\,
\]

%dsp-level:0
\[
\Sigma_{\mathrm B}\gtrless b\ln2
\]

%dsp-level:0
\[
\iff (4/3)^a\gtrless2^b
\]

%dsp-level:0
\[
\iff 4^a\gtrless3^a2^b:
\]

%dsp-level:0
\[
\quad (c,r)
\]

%dsp-level:0
\[
= \left(\tfrac35,\tfrac19\right):\ 2^8\!\cdot\!1\ \text{vs}\ 3^5:\ 256>243;
\]

%dsp-level:0
\[
\ \left(\tfrac35,\tfrac18\right):\ 2^{31}\ \text{vs}\ 3^{20};
\]

%dsp-level:0
\[
\ \left(\tfrac12,\tfrac18\right):\ 2^{13}\ \text{vs}\ 3^8;
\]

%dsp-level:0
\[
\ \left(\tfrac12,\tfrac17\right):\ 2^{11}\ \text{vs}\ 3^7\ [\text{machine: four integer comparisons}]
\]

%dsp-level:0
\[
(5):\ \Sigma_{\mathrm B}>0\iff4>3\ [\text{machine}];
\]

%dsp-level:0
\[
\qquad \text{the two witnesses}\ 64>54,\ 1024>972\ [\text{machine}];
\]

%dsp-level:0
\[
\qquad (6):\ \mathrm{EM.25}(7),\ \mathrm{BM.61}(4)\ \text{re-cited}
\]

%dsp-level:0
\[
(7):\ 1280^3\ \text{and}\ 960^4\ \text{computed exactly and compared}\ [\text{machine: the two integers differ}];
\]

%dsp-level:0
\[
\qquad (8):\ \text{assembly}. \ \blacksquare
\]

\textbf{Phys.} The baseline stands on integer certificates. The refined rest input assembles three derived registers — the count, the charge-filtered end census, and the scar's field energy — into one row per class, and the class difference has a single exact zero: the critical strength, enclosed between one ninth and one seventh by four comparisons of powers of two against powers of three. The branch on which the neutral baryon is the heavier one is a non-empty exact interval, certified from below by two hundred fifty-six against two hundred forty-three, and no theorem of Parts 1–4 fixes the coupling numerically — so the ordering result is conditional on the symbol being sub-critical, never on an imported value. The ratio discipline holds with it: the mass ratio is the logarithm ratio of twelve-eighty over nine-sixty, the certificates that twelve-eighty cubed differs from nine-sixty to the fourth keep the measure price away from the mass read, and nothing measured enters anywhere.

\begin{theorem}[{\texorpdfstring{BM.66 — The operator lift of the EM correction: the charge-blind trace, the quadratic invariant, and the assembled operator rest input}{BM.66 - The operator lift of the EM correction: the charge-blind trace, the quadratic invariant, and the assembled operator rest input}}]\label{tag:BM-66}
\textbf{Thm.}

%dsp-level:0
\[
(1)\quad\textbf{the operator and its spectrum:}
\]

%dsp-level:0
\[
\qquad \widehat F\ :=\ \mathbf F_{\mathrm{ext}}(\mathbf c,\widehat Q)\ =\ h(\mathbf c)\otimes I_2\ +\ \widehat Q\,I_2\otimes\hat\kappa,
\]

%dsp-level:0
\[
\qquad h(\mathbf c)=\tfrac12\bigl(c_0I_2+\mathbf c\cdot\boldsymbol\sigma\bigr)
\]

%dsp-level:0
\[
\qquad [\mathrm{BM.22}(1);\ \mathrm{FT.129}(5);\ \text{the charge read carried as the operator}\ \widehat Q]:
\]

%dsp-level:0
\[
\operatorname{spec}\widehat Q\ =\ \{-1,\,0,\,+1\}\ \ \text{on the materialized classes}
\]

%dsp-level:0
\[
\qquad [\mathrm{EM.9};\ \mathrm{D3.4};\ \mathrm{Def\text{-}BM.8}(1),(5);\ \mathrm{Def\text{-}BM.3}(2),(3)];
\]

%dsp-level:0
\[
\lambda_{s,t}\ =\ \tfrac12\bigl(c_0+sC\bigr)+t\,Q,
\]

%dsp-level:0
\[
\qquad C:=|\mathbf c|,
\]

%dsp-level:0
\[
\qquad s,t\in\{\pm1\}
\]

%dsp-level:0
\[
\qquad [\mathrm{BM.22}(3);\ \mathrm{FT.130}(5);\ \text{machine: the monic quartic from the four roots at exact rational witnesses}];
\]

%dsp-level:0
\[
\textbf{the promotion conditionality is inherited,}\ \neg\textbf{new:}
\]

%dsp-level:0
\[
\qquad \text{the flipping block exists under}\ \mathrm{Def\text{-}H.43}\text{'s declared read promotion}
\]

%dsp-level:0
\[
\qquad [\mathrm{BM.22}(5)];
\]

%dsp-level:0
\[
(2)\quad\textbf{the linear trace is charge-blind:}
\]

%dsp-level:0
\[
\qquad \boxed{\ \operatorname{Tr}\widehat F\ =\ \sum_{s,t\in\{\pm1\}}\lambda_{s,t}\ =\ 2\,c_0\ }
\]

%dsp-level:0
\[
\qquad [\text{machine: the symbolic sum};\ \textstyle\sum_{s,t}t\,Q=0,\ \sum_{s,t}s\,C=0]:
\]

%dsp-level:0
\[
\textbf{the operator form of the uniqueness row's trace line:}
\]

%dsp-level:0
\[
\qquad \text{no trace read of}\ \widehat F\ \text{itself splits the two classes}
\]

%dsp-level:0
\[
\qquad [\mathrm{BM.58}(6):\ \textstyle\sum_{s,t}tq=0;\ \ \mathrm{BM.1}(2):\ \text{the trace-only read}];
\]

%dsp-level:0
\[
(3)\quad\textbf{the quadratic invariant:}
\]

%dsp-level:0
\[
\qquad \operatorname{Tr}\widehat F^2\ =\ \sum_{s,t}\lambda_{s,t}^2\ =\ c_0^2+C^2+4\,\widehat Q^2
\]

%dsp-level:0
\[
\qquad [\text{machine: the polynomial identity in}\ \mathbb Q[c_0,C,Q]];
\]

%dsp-level:0
\[
\widehat F_0\ :=\ \widehat F\big|_{\widehat Q=0}:\qquad\operatorname{Tr}\widehat F_0^2\ =\ c_0^2+C^2;\qquad\boxed{\ I_F\ :=\ \operatorname{Tr}\widehat F^2-\operatorname{Tr}\widehat F_0^2\ =\ 4\,\widehat Q^2\ }\qquad[\text{machine}]:
\]

%dsp-level:0
\[
\textbf{the first trace invariant carrying charge;}
\]

%dsp-level:0
\[
\qquad I_F=(0,\,4,\,4)\ \ \text{at}\ Q=(0,\,\pm1)
\]

%dsp-level:0
\[
\qquad [\text{machine}]:
\]

%dsp-level:0
\[
\qquad \mathrm{BM.58}(5)\text{'s monopole selector}\ Q^2=(0,1,1)\ \textbf{recovered as a trace row};
\]

%dsp-level:0
\[
(4)\quad\textbf{the operator EM correction:}
\]

%dsp-level:0
\[
\qquad \boxed{\ \Delta\bigl(\widehat M_{\mathrm{rest}}\mathfrak h\bigr)_{\mathrm{EM}}\ :=\ \frac{E_{\mathrm{int}}}4\;c\,\varepsilon_{\mathrm{EM}}\ln2\;\Bigl(\operatorname{Tr}\widehat F^2-\operatorname{Tr}\widehat F_0^2\Bigr)\ }
\]

%dsp-level:0
\[
\qquad \text{with}\ \frac{E_{\mathrm{int}}}4=\frac64=\frac32\ \ \text{at}\ E_{\mathrm{int}}=6
\]

%dsp-level:0
\[
\qquad [\mathrm{BM.17};\ \text{Census FT.35}(2);\ \text{machine}]:
\]

%dsp-level:0
\[
\textbf{the coefficient is derived,}\ \neg\textbf{chosen:}\qquad\frac{E_{\mathrm{int}}}4\cdot I_F\ =\ \frac{E_{\mathrm{int}}}4\cdot4\,\widehat Q^2\ =\ E_{\mathrm{int}}\,\widehat Q^2\qquad[\text{machine}]:
\]

%dsp-level:0
\[
\Delta\bigl(\widehat M_{\mathrm{rest}}\mathfrak h\bigr)_{\mathrm{EM}}\ =\ c\,\varepsilon_{\mathrm{EM}}\,\widehat Q^2\,E_{\mathrm{int}}\ln2\ =\ c\,\varepsilon_{\mathrm{EM}}\,\widehat Q^2\,m_{\mathrm{hid}}\,\mathfrak h\ =\ E^{\mathrm{self}}_{\mathrm{EM}}(\widehat Q)\,\mathfrak h
\]

%dsp-level:0
\[
\qquad [\mathrm{BM.46}(1)\ \text{unified form};\ \text{machine: one monomial}];
\]

%dsp-level:0
\[
\textbf{on charge eigenstates:}
\]

%dsp-level:0
\[
\qquad Q=0:\quad\Delta\bigl(\widehat M_{\mathrm{rest}}\mathfrak h\bigr)_{\mathrm{EM}}=0;
\]

%dsp-level:0
\[
\qquad Q=\pm1:\quad\Delta\bigl(\widehat M_{\mathrm{rest}}\mathfrak h\bigr)_{\mathrm{EM}}=6\,c\,\varepsilon_{\mathrm{EM}}\ln2
\]

%dsp-level:0
\[
\qquad [\text{machine}]:
\]

%dsp-level:0
\[
\textbf{the scalar row of}\ \mathrm{Def\text{-}BM.13}(5)\ \textbf{is reproduced exactly on all three materialized classes};
\]

%dsp-level:0
\[
(5)\quad\textbf{the assembled operator rest input:}
\]

%dsp-level:0
\[
\qquad \boxed{\ \widehat M_{\mathrm{rest}}\mathfrak h\ =\ E_{\mathrm{int}}\ln2\ +\ \ln W(\widehat Q)\ +\ \frac32\,c\,\varepsilon_{\mathrm{EM}}\ln2\Bigl(\operatorname{Tr}\widehat F^2-\operatorname{Tr}\widehat F_0^2\Bigr)\ }
\]

%dsp-level:0
\[
\ln W(\widehat Q)\ :=\ \sum_{Q\in\operatorname{spec}}\bigl(\ln W(Q)\bigr)\,P_Q
\]

%dsp-level:0
\[
\qquad [\text{functional calculus on the charge eigenspaces};\ \text{math: the spectral theorem for a diagonal operator on a finite charge class}];
\]

%dsp-level:0
\[
\textbf{the eigenvalue rows are exactly}\ \mathrm{BM.62}(4)\text{'s:}\qquad Q=0:\ \ \ln1280;\qquad Q=\pm1:\ \ \ln960+6\,c\,\varepsilon_{\mathrm{EM}}\ln2\qquad[\text{machine}];
\]

%dsp-level:0
\[
\textbf{the CPT row:}
\]

%dsp-level:0
\[
\qquad \omega:\ \widehat Q\mapsto-\widehat Q\ \ \text{fixes}\ \ W(\widehat Q)\ \ [W(Q)=W(-Q);\ \mathrm{Def\text{-}BM.8}(7)]\ \ \text{and fixes}\ \ I_F\ \ [(3):\ \widehat Q^2\ \text{only}]:
\]

%dsp-level:0
\[
\textbf{the proton--antiproton degeneracy is an operator identity}\qquad[\mathrm{BM.28}(5);\ \mathrm{BM.30}(4);\ \text{machine}];
\]

%dsp-level:0
\[
(6)\quad\textbf{the no-double-count guard:}\qquad\textbf{three derived routes to one monomial:}\qquad\begin{array}{l|l}\text{the sector census} & \mathrm{BM.52}(3),(4):\ \text{five zeros and one monomial}\\\text{the field integral} & \mathrm{BM.46}(1);\ \mathrm{Def\text{-}BM.12}(5)\\\text{the quadratic trace} & (3),(4)\ \text{of this tag}\end{array}\qquad[\text{machine: the three agree on every charge eigenstate}]:
\]

%dsp-level:0
\[
\mathrm{BM.52}(5)\text{'s guard extends: in any decomposition the EM contribution is counted once}
\]

%dsp-level:0
\[
\qquad [\mathrm{BM.62}(1)\ \text{stands}];
\]

%dsp-level:0
\[
\textbf{the alphabet census:}
\]

%dsp-level:0
\[
\qquad \text{the trace subtraction removes}\ c_0^2+C^2\ \text{exactly}\ (c_0,C\ \text{are}\ \neg\text{free symbols of}\ I_F)
\]

%dsp-level:0
\[
\qquad [\text{machine}];
\]

%dsp-level:0
\[
\text{the free-symbol set of}\ \Delta\bigl(\widehat M_{\mathrm{rest}}\mathfrak h\bigr)_{\mathrm{EM}}\ \text{is}\ \{c,\varepsilon_{\mathrm{EM}},\ln2,\widehat Q\}\ \subseteq\ \mathrm{BM.62}(11)\text{'s}\ \Delta m_{\mathrm{EM}}\ \text{alphabet under}\ e^2
\]

%dsp-level:0
\[
= 4\pi\varepsilon_{\mathrm{EM}},
\]

%dsp-level:0
\[
\ \Delta x_{\mathrm B}^{-1}\qquad = m_{\mathrm{hid}}
\]

%dsp-level:0
\[
\qquad [\mathrm{EW.35};\ \mathrm{Def\text{-}BM.4}(1);\ \text{machine: the token census}];
\]

%dsp-level:0
\[
\textbf{the disjointness from the transport alphabet}\ \{m_q,W(Q),\ln2,\mathfrak h\}\ \text{holds except for the shared}\ \ln2
\]

%dsp-level:0
\[
\qquad [\mathrm{BM.59}(8);\ \mathrm{BM.37}(2);\ \text{machine}];
\]

%dsp-level:0
\[
(7)\quad\textbf{the uniqueness row is respected:}
\]

%dsp-level:0
\[
\qquad \mathrm{BM.58}(6)\text{'s table gains}\ \neg\text{a new row:}
\]

%dsp-level:0
\[
\qquad I_F\ \text{is the operator species of the adopted}\ E^{\mathrm{self}}_{\mathrm{EM}}\ \text{entry}\ [\mathrm{Def\text{-}BM.12}(6);\ (4)],
\]

%dsp-level:0
\[
\ \text{with the same monomial and the same negative sign for}\ m_n-m_p\ [\mathrm{BM.58}(5);\ \text{machine}];
\]

%dsp-level:0
\[
F_{LR}=\widehat QI_2\text{'s row}\ \text{``}\,0\ \text{on a trace read}\,\text{''}\ \text{is the derivation of}\ (2);
\]

%dsp-level:0
\[
\qquad \textbf{the unique surviving class-dependent rest-slot candidate remains}\ W(Q)\ \text{read through}\ \ln W(Q)
\]

%dsp-level:0
\[
\qquad [\mathrm{BM.59}(2),(3)];
\]

%dsp-level:0
\[
(8)\quad\textbf{status:}
\]

%dsp-level:0
\[
\qquad \textbf{derived:}\ (1)\text{--}(7);
\]

%dsp-level:0
\[
\qquad \textbf{conditional on}\ \mathrm{Def\text{-}H.43}\text{'s promotion as}\ \mathrm{BM.22}(5)\ \text{records it}\ (\neg\text{a new declaration});
\]

%dsp-level:0
\[
c\ \text{stays priced,}\ \neg\text{adopted}\ [\mathrm{Def\text{-}BM.12}(4)];
\]

%dsp-level:0
\[
\qquad \textbf{the new coefficients are}\ \bigl\{\tfrac32=E_{\mathrm{int}}/4,\ 4\bigr\}:\ \text{a count row and a trace normalization}\ \text{—}\ \neg\text{a P4 addition},
\]

%dsp-level:0
\[
\ \neg\text{a fitting parameter}\ [\text{machine: the coefficient census}];
\]

%dsp-level:0
\[
\neg\text{a new dimensionful scale}\ [\mathrm{P1}];\qquad\neg\text{a measured value}\ [\mathrm{P5}];\qquad\text{§Cal-BM stays empty}.
\]

\textbf{Req.} D2.9, D3.4 (Parts 1–2); CT2.1, EM.9, EW.35 (Part 2); P1, P5 (Parts 1–2); Def-H.43, FT.129(2), FT.129(5), FT.130(5) (Part 3); Census FT.35(2) (row-wise); Def-BM.3(2), Def-BM.3(3), Def-BM.4(1), Def-BM.8(1), Def-BM.8(5), Def-BM.8(7), Def-BM.12(4), Def-BM.12(5), Def-BM.12(6), Def-BM.13(5), BM.1(2), BM.17, BM.22(1), BM.22(2), BM.22(3), BM.22(5), BM.26, BM.28(5), BM.30(4), BM.37(2), BM.46(1), BM.46(3), BM.52(3), BM.52(4), BM.52(5), BM.58(5), BM.58(6), BM.59(2), BM.59(3), BM.59(8), BM.62(1), BM.62(4), BM.62(11) (this part); exact multivariate polynomial identities over $\mathbb Q[c_0,C,Q,c,\varepsilon_{\mathrm{EM}},\ln2]$, matrix traces over complex-with-rational-parts at exact witnesses, spectral calculus for diagonal operators on finite charge classes, token censuses [math; machine: \texttt{checks/\allowbreak{}verify\_\allowbreak{}H\_\allowbreak{}organic\_\allowbreak{}bridge\_\allowbreak{}post\_\allowbreak{}bm14.py} §A].

\textbf{Pf.}

%dsp-level:0
\[
(1):\ \mathrm{BM.22}(1)\text{--}(3)\ \text{re-cited with}\ q\ \text{carried as}\ \widehat Q;
\]

%dsp-level:0
\[
\qquad \text{the charpoly}\ \bigl((y-u)^2-Q^2\bigr)\bigl((y+u)^2-Q^2\bigr),
\]

%dsp-level:0
\[
\ y
\]

%dsp-level:0
\[
= x-\tfrac{c_0}2,
\]

%dsp-level:0
\[
\ u^2
\]

%dsp-level:0
\[
= \tfrac{C^2}4\ \text{built from the four roots and compared to the Faddeev--LeVerrier charpoly at witnesses with}\ C^2\in\{0,\tfrac14,1,\tfrac49,\tfrac{9}{49}\}\ [\text{machine: exact over the rationals}]
\]

%dsp-level:0
\[
(2):\ \sum_{s,t}\bigl[\tfrac12(c_0+sC)+tQ\bigr]
\]

%dsp-level:0
\[
= 4\cdot\tfrac12c_0+\tfrac C2\sum_{s,t}s+Q\sum_{s,t}t
\]

%dsp-level:0
\[
= 2c_0+0+0\ [\text{machine: the symbolic sum}];
\]

%dsp-level:0
\[
\qquad \text{the matrix row:}\ \operatorname{Tr}\bigl(h\otimes I_2+QI_2\otimes\hat\kappa\bigr)
\]

%dsp-level:0
\[
= 2\operatorname{tr}h+2Q\operatorname{tr}\hat\kappa
\]

%dsp-level:0
\[
= 2c_0\ [\hat\kappa\ \text{traceless};\ \text{machine: exact traces at every witness}\times\text{ternary}\ Q]
\]

%dsp-level:0
\[
(3):\ \lambda_{s,t}^2=\tfrac14(c_0+sC)^2+tQ(c_0+sC)+t^2Q^2:
\]

%dsp-level:0
\[
\qquad \sum_{s,t}\tfrac14(c_0+sC)^2=\tfrac12\bigl[(c_0+C)^2+(c_0-C)^2\bigr]=c_0^2+C^2;
\]

%dsp-level:0
\[
\qquad \sum_{s,t}tQ(c_0+sC)=0;
\]

%dsp-level:0
\[
\qquad \sum_{s,t}t^2Q^2=4Q^2\ [\text{machine: the polynomial identity}];
\]

%dsp-level:0
\[
\text{the matrix route agrees:}\ \operatorname{Tr}F^2
\]

%dsp-level:0
\[
= 2\operatorname{tr}\bigl(h^2\bigr)+2Q\operatorname{tr}(h)\operatorname{tr}(\hat\kappa)+Q^2\operatorname{tr}(I_2)\operatorname{tr}\bigl(\hat\kappa^2\bigr)
\]

%dsp-level:0
\[
= \bigl(c_0^2+C^2\bigr)+0+4Q^2\ [\hat\kappa\ \text{the ray swap},\ \hat\kappa^2=I_2,\ \operatorname{tr}\hat\kappa=0\ [\mathrm{FT.129}(2)];\ \operatorname{tr}h=c_0;\ \text{machine at every witness}]
\]

%dsp-level:0
\[
(4):\ \tfrac{E_{\mathrm{int}}}4\cdot4Q^2=E_{\mathrm{int}}Q^2\ [\text{machine}];
\]

%dsp-level:0
\[
\qquad c\,\varepsilon_{\mathrm{EM}}Q^2E_{\mathrm{int}}\ln2=c\,\varepsilon_{\mathrm{EM}}Q^2m_{\mathrm{hid}}\mathfrak h\ [\mathrm{BM.46}(1)\ \text{re-cited};\ \text{machine: one monomial}];
\]

%dsp-level:0
\[
\qquad \text{the eigenstate rows at}\ Q^2\in\{0,1\}\ [\text{machine}]
\]

%dsp-level:0
\[
(5):\ P_Q\ \text{the spectral projections of}\ \widehat Q;
\]

%dsp-level:0
\[
\ \ \ln W(\widehat Q)\ \text{diagonal with the entries}\ (\ln20,\ \ln15,\ \ln15)\ [\text{machine}];
\]

%dsp-level:0
\[
\qquad \text{the eigenvalue sum}\ 6\ln2+\ln20
\]

%dsp-level:0
\[
= \ln1280,
\]

%dsp-level:0
\[
\ 6\ln2+\ln15+6c\varepsilon_{\mathrm{EM}}\ln2
\]

%dsp-level:0
\[
= \ln960+6c\varepsilon_{\mathrm{EM}}\ln2\ [\text{machine: the prime-log ring}];
\]

%dsp-level:0
\[
\qquad \omega\text{-row:}\ W(-Q)=W(Q),\ (-Q)^2=Q^2\ [\text{machine}]
\]

%dsp-level:0
\[
(6):\ \text{the three routes compared monomial by monomial on the three materialized classes}\ [\text{machine}];
\]

%dsp-level:0
\[
\qquad \text{the symbol census of}\ I_F\ \text{and of the correction}\ [\text{machine:}\ \{c,\varepsilon_{\mathrm{EM}},\ln2,Q\}];
\]

%dsp-level:0
\[
\qquad e^2=4\pi\varepsilon_{\mathrm{EM}}\ [\mathrm{EW.35};\ \text{machine}];\ \Delta x_{\mathrm B}^{-1}=m_{\mathrm{hid}}\ [\mathrm{Def\text{-}BM.4}(1);\ \text{machine}]
\]

%dsp-level:0
\[
(7)\text{--}(8):\ \mathrm{BM.58}(5),(6),\ \mathrm{BM.59}(2),(3)\ \text{re-cited against}\ (2)\text{--}(4);
\]

%dsp-level:0
\[
\qquad \text{the coefficient census}\ \bigl\{\tfrac64,\ 4\bigr\}\ \text{and the standing-rule scans}\ [\text{machine}]. \ \blacksquare
\]

\textbf{Phys.} The scalar row of the EM correction is the eigenvalue row of an operator identity, and the operator is the force matrix the corpus already carried. Its linear trace reads only the shared block — two c-naught, blind to charge — which is the operator form of the uniqueness row's old trace line: no trace of the force itself can split the two classes. The first trace that carries charge is the quadratic one, and its charge content is exactly four Q-squared, so the monopole selector comes back as a trace law. Multiplied by the count row's own coefficient — the internal edge count over four, three halves at six edges — the operator correction reproduces the field-energy monomial on every materialized class, joins the sector census and the field integral as a third route to one and the same term, and keeps the proton–antiproton degeneracy exact, because only the square of the charge enters and the census is conjugation-symmetric. Nothing new is counted: the trace subtraction removes the shared blocks, the symbol alphabet collapses to the adopted row's own, and the distribution factor stays priced in two exact values and adopted in none.

\end{theorem}

\pseudohead{theorem}{\texorpdfstring{BM.67}{BM.67}}{\texorpdfstring{The holonomy lift attempted: the derived $\mathbb Z_2$ sector, the angle-free force census, and the two-torsion obstruction}{The holonomy lift attempted: the derived sector, the angle-free force census, and the two-torsion obstruction}}\label{tag:BM-67}
\textbf{Thm.}

%dsp-level:0
\[
(1)\quad\textbf{the register and the derived discrete phase:}
\]

%dsp-level:0
\[
\qquad \alpha_e\in\{0,\pi\},
\]

%dsp-level:0
\[
\qquad \sigma_e=e^{i\alpha_e}\in\{+1,-1\}
\]

%dsp-level:0
\[
\qquad [\mathrm{A4};\ \mathrm{A7};\ \mathrm{CT0.9};\ \mathrm{Def\text{-}EM.20};\ \mathrm{EM.20}(1),(2):\ \text{registration data determined at crystallization},\ \neg\text{declared constants}]:
\]

%dsp-level:0
\[
\text{virtual edges carry}\ \neg\alpha_e\qquad[\mathrm{EM.20}(3)];\qquad[\alpha_e]=\mathfrak q^0\ \ \text{in every calibration}\qquad[\mathrm{EM.21}];
\]

%dsp-level:0
\[
\textbf{for an oriented chain}\ c:
\]

%dsp-level:0
\[
\qquad \Phi(c)\ :=\ \sum_{e\in c}\operatorname{or}(e)\,\alpha_e\ \ (\mathrm{mod}\ 2\pi),
\]

%dsp-level:0
\[
\qquad U(c)\ :=\ \prod_{e\in c}\sigma_e^{\operatorname{or}(e)}\ =\ e^{i\Phi(c)}
\]

%dsp-level:0
\[
\qquad [\mathrm{D3.5}\ \text{weight grammar};\ \text{machine}];
\]

%dsp-level:0
\[
(2)\quad\textbf{the bare holonomy is}\ \mathbb Z_2\ \textbf{(theorem):}
\]

%dsp-level:0
\[
\qquad \sigma^{-1}=\sigma\ \ \text{on}\ \{\pm1\}\implies\text{the orientation drops out of every closed product}:
\]

%dsp-level:0
\[
\boxed{\ U(\gamma)\ =\ \prod_{e\in\gamma}\sigma_e\ =\ (-1)^{N_\pi(\gamma)},\qquad N_\pi(\gamma):=\#\{e\in\gamma:\ \alpha_e=\pi\}\ }
\]

%dsp-level:0
\[
\qquad [\text{machine: exhaustive over every sign assignment on loops of length}\ 2,\dots,7,\ \text{both orientations}]:
\]

%dsp-level:0
\[
U(\gamma_1\!\cdot\!\gamma_2)=U(\gamma_1)\,U(\gamma_2),
\]

%dsp-level:0
\[
\qquad U(\bar\gamma)=U(\gamma)
\]

%dsp-level:0
\[
\qquad [\text{machine: concatenation and reversal censuses}]:
\]

%dsp-level:0
\[
\qquad \textbf{the Wilson sector of the register is exactly binary};
\]

%dsp-level:0
\[
(3)\quad\textbf{the parity lock, lifted to the composite:}
\]

%dsp-level:0
\[
\qquad \sum_i\aleph_{v_i}\equiv|q|\ (\mathrm{mod}\ 2),
\]

%dsp-level:0
\[
\qquad \sum_i\aleph_{v_i}\ \ge\ |q|
\]

%dsp-level:0
\[
\qquad [\mathrm{Def\text{-}BM.13}(2);\ \mathrm{EM.9b};\ \mathrm{D3.3};\ \text{machine: exhaustive over every sign sequence},\ k=2,\dots,7]:
\]

%dsp-level:0
\[
Q=\sum_iq_i
\]

%dsp-level:0
\[
\qquad [\mathrm{Def\text{-}BM.8}(1),(3)]:
\]

%dsp-level:0
\[
\qquad |q_i|\equiv q_i\ (\mathrm{mod}\ 2)\implies\textstyle\sum_i\aleph^{(i)}\equiv\sum_i|q_i|\equiv Q\equiv|Q|\ (\mathrm{mod}\ 2);
\]

%dsp-level:0
\[
\qquad \textstyle\sum_i\aleph^{(i)}\ge\sum_i|q_i|\ge\bigl|\sum_iq_i\bigr|=|Q|
\]

%dsp-level:0
\[
\qquad [\text{machine: exhaustive over the}\ 3^3\ \text{charge triples}]:
\]

%dsp-level:0
\[
\boxed{\ \sum_i\aleph_{v_i}\ \equiv\ |Q|\ (\mathrm{mod}\ 2),\qquad\sum_i\aleph_{v_i}\ \ge\ |Q|\ }
\]

%dsp-level:0
\[
\qquad \text{and the wall-sector read is charge-locked:}
\]

%dsp-level:0
\[
\qquad U_{\mathrm{wall}}=(-1)^{\sum\aleph}=(-1)^{|Q|}
\]

%dsp-level:0
\[
\qquad [\text{machine: the}\ 4^3\ \text{end-pair census agrees with}\ W(Q)\ \text{class by class}];
\]

%dsp-level:0
\[
(4)\quad\textbf{the winding and sealing sectors are integer and phase-trivial:}
\]

%dsp-level:0
\[
\qquad \Theta(c)=\sum_{e\in c}\widehat{ds}^{\,2}(e)=\sum_{e\in c}\kappa_e;
\]

%dsp-level:0
\[
\qquad \text{materialized closed chains:}\quad\Theta\in2\mathbb Z,
\]

%dsp-level:0
\[
\qquad \Phi_{\mathrm{spiral}}=\pi\,\Theta=2\pi k,
\]

%dsp-level:0
\[
\qquad \mathrm{Hol}_{\mathrm{spiral}}=+1,
\]

%dsp-level:0
\[
\qquad k=\tfrac\Theta2\in\mathbb Z
\]

%dsp-level:0
\[
\qquad [\mathrm{FT.22};\ \mathrm{FT.24}(1),(2);\ \mathrm{FT.25}(1),(3);\ \text{Census FT.35}(3)]:
\]

%dsp-level:0
\[
\mathrm{Hol}_{\mathrm{int}}=(-1)^{\Theta_{\mathrm{int}}}\in\{\pm1\},\ \ \text{nontrivial only on partial-chain internal sets}
\]

%dsp-level:0
\[
\qquad [\mathrm{FT.25}(2),(3)];
\]

%dsp-level:0
\[
\qquad \text{the baryon-type sealing phase:}\quad\Theta_{\mathrm{int}}\in4\mathbb Z
\]

%dsp-level:0
\[
\qquad [\mathrm{BM.2}(1);\ \text{Census FT.49}(3)]:
\]

%dsp-level:0
\[
e^{i\pi\Theta}=+1\ \ \forall\,\Theta\in2\mathbb Z
\]

%dsp-level:0
\[
\qquad [\text{machine}]:
\]

%dsp-level:0
\[
\qquad \textbf{the whole winding/sealing register contributes trivially to every}\ U(1)\ \textbf{phase read}\ —\ \text{its content is the integer rung}\ k,
\]

%dsp-level:0
\[
\ \neg\text{an angle};
\]

%dsp-level:0
\[
(5)\quad\textbf{the force-block census (the angle-free theorem):}
\]

%dsp-level:0
\[
\qquad \text{A8's complete decomposition, regrouped as}\ \ \widehat F
\]

%dsp-level:0
\[
= \widehat F_{J^\pm}+\widehat F_{G/H}+\widehat F_W+\widehat F_{\mathrm{colour}}+\widehat F_{\mathrm{EM}}+\widehat F_{D^\pm}
\]

%dsp-level:0
\[
\qquad [\mathrm{A8};\ \mathrm{BM.52}(1):\ \text{a regrouping},\ \neg\text{a sub-block restriction}],\ \text{read at a sealed 3-juncture}:
\]

%tbl-level:3
{\scriptsize\setlength{\tabcolsep}{4pt}
\begin{longtable}{>{\raggedright\arraybackslash}p{\dimexpr0.1058\textwidth-2\tabcolsep\relax}|>{\raggedright\arraybackslash}p{\dimexpr0.1058\textwidth-2\tabcolsep\relax}|>{\raggedright\arraybackslash}p{\dimexpr0.7885\textwidth-2\tabcolsep\relax}}
\toprule
block & candidate angle content & the derived verdict \\
\addlinespace[2pt]
\midrule
\endfirsthead
\toprule
block & candidate angle content & the derived verdict \\
\addlinespace[2pt]
\midrule
\endhead
\bottomrule
\endlastfoot
$\widehat F_{J^\pm}$ (the seals) & the sealing phase & $\Theta_{\mathrm{int}}\in\allowbreak 4\mathbb Z$ — an integer $[\mathrm{BM.2}(1);\ \text{Census FT.49}(3)]$; the rotation group of the $\varepsilon$-seal is the finite $C_3$: three one-dimensional characters, no continuous parameter $[\mathrm{BM.3}(1);\ \mathrm{BM.10}(5);\ \mathrm{FT.116}]$ \\
$\widehat F_{G/H}$ (the metric) & frame coefficients & the seven geometric monomials carry charge exponent $0$; no vertex angle $[\mathrm{BM.52}(3);\ \mathrm{BM.58}(1),\allowbreak (2),\allowbreak (4)]$ \\
$\widehat F_W$ (weak/shift) & the shift phase & $\|w\|^2$ class-shared; the Casimir is branch data, neutral at the $\sigma$-vertex $[\mathrm{BM.28}(6);\ \mathrm{BM.40}(3);\ \mathrm{BM.52}(3)]$ \\
$\widehat F_{\mathrm{colour}}$ & slot rotations & annihilated on a sealed singlet: the trace law $\mathfrak T(H)=\allowbreak 3a$; the seed obeys $T_a\mathsf v=\allowbreak 0$ $[\mathrm{BM.1}(2);\ \mathrm{BM.57}(5),\allowbreak (6);\ \mathrm{FT.117}(1);\ \mathrm{EW.32}]$ \\
$\widehat F_{\mathrm{EM}}$ & the gauge phase & no local projector; static monopole order; the coupling $e$ symbolic; chart-side $[\mathrm{Def\text{-}BM.12}(2),\allowbreak (7);\ \mathrm{BM.58}(3)]$ \\
$\widehat F_{D^\pm}$ (dark) & — & $D^+\allowbreak =\allowbreak \Lambda\,\allowbreak \mathbb I$ uniform; $D^-$ locked, retained, $\ne\allowbreak g$accessible, no gauge structure $[\mathrm{A8};\ \mathrm{CT4.1}\text{–}\mathrm{CT4.5};\ \mathrm{D4.2};\ \mathrm{D1.11};\ \mathrm{BM.52}(3);\ \mathrm{FT.73}(1)]$ \\
\end{longtable}
}

%dsp-level:0
\[
\textbf{the seed-transport row:}
\]

%dsp-level:0
\[
\qquad \text{the channel angles}\ \delta\in\bigl\{\tfrac29,\tfrac2{27},\tfrac18\bigr\}\ \text{are channel-line data,}\ \neg\text{attached to the baryon rows}\ [\mathrm{Def\text{-}BM.11};\ \mathrm{BM.59}(1\mathrm a):\ \text{the fourth ground stands}],
\]

%dsp-level:0
\[
\ \text{and they enter the Gram mass form,}\ \neg\text{a loop connection}\ [\mathrm{EW.33};\ \mathrm{Def\text{-}H.37}(1)]:
\]

%dsp-level:0
\[
\boxed{\ \theta_v^{\mathrm{force}}\ \ \textbf{is derived angle-free at a sealed 3-juncture: every block is annihilated by a trace/projector law, integer- or binary-valued, class-shared, or unattached}\ }
\]

%dsp-level:0
\[
\qquad [\text{machine: the row census exhaustive against A8's block list}];
\]

%dsp-level:0
\[
\textbf{the caps add integers only:}
\]

%dsp-level:0
\[
\qquad C_{\mathrm{cap}}=16=7+3+3+3,
\]

%dsp-level:0
\[
\qquad D_{\mathrm{tick}}(3)=2^4=16,
\]

%dsp-level:0
\[
\qquad O(2)=O(3)=0,
\]

%dsp-level:0
\[
\qquad \kappa_0\in\mathbb Z
\]

%dsp-level:0
\[
\qquad [\mathrm{CT0.5};\ \mathrm{CT0.6};\ \mathrm{BM.52}(1);\ \text{Census FT.28}(5);\ \text{Census FT.64}(5);\ \mathrm{FT.24}(1);\ \text{machine}];
\]

%dsp-level:0
\[
(6)\quad\textbf{the two-torsion obstruction (theorem):}
\]

%dsp-level:0
\[
\qquad \text{let}
\]

%dsp-level:0
\[
\qquad \Theta_R(\gamma)\ =\ 2\pi\,n(\gamma)\ +\ \sum_{e\in\gamma}\alpha_e\ +\ \sum_{v\in\gamma}\theta_v^{\mathrm{force}}
\]

%dsp-level:0
\[
\qquad (\mathrm{mod}\ 2\pi)
\]

%dsp-level:0
\[
\text{be any constrained lift with}\ n(\gamma)\ \text{integer}\ [(4),(5):\ \text{the register and the caps fix integers}].
\]

%dsp-level:0
\[
\qquad \text{Then:}
\]

%dsp-level:0
\[
\qquad 2\pi n\equiv0,\quad\textstyle\sum\alpha_e\in\pi\mathbb Z,\quad\sum\theta_v^{\mathrm{force}}\equiv0\ \ \text{on the derived data}:
\]

%dsp-level:0
\[
\boxed{\ \operatorname{im}(\text{any register/force/caps lift})\ \subseteq\ \{0,\ \pi\}\ \subset\ \mathbb R/2\pi\mathbb Z\ }\qquad[\text{machine: the image census}]:
\]

%dsp-level:0
\[
\textbf{against the defect spectrum:}
\]

%dsp-level:0
\[
\qquad \text{the six states of}\ \mathrm{Def\text{-}BM.14}(3)\ \text{carry the angles}\ \ \Theta_R\in\bigl\{\pi,\ 2\pi,\ \tfrac\pi2,\ \tfrac{2\pi}3,\ \tfrac\pi5,\ \tfrac\pi3\bigr\}
\]

%dsp-level:0
\[
\qquad [\delta_v/\pi\in\bigl\{1,\tfrac12,\tfrac23,\tfrac15,\tfrac13\bigr\};\ \text{Census FT.63}(4),(6);\ \text{machine}]:
\]

%dsp-level:0
\[
\text{exactly two of the six lie in the two-torsion image}\ —\ \text{the self-dual floor's pair}\ (\pi;\ 2\pi\equiv0);
\]

%dsp-level:0
\[
\qquad \tfrac\pi2,\ \tfrac{2\pi}3,\ \tfrac\pi5,\ \tfrac\pi3\ \ \text{are pairwise distinct nontrivial cosets modulo}\ \pi
\]

%dsp-level:0
\[
\qquad [\text{machine: exact rational coset arithmetic}]:
\]

%dsp-level:0
\[
\boxed{\ \neg\exists\ \text{a lift of the full real holonomy from the register, the force decomposition and the caps}\ }
\]

%dsp-level:0
\[
\qquad \text{the lift is derived exactly on the self-dual floor and obstructed on both dual pairs};
\]

%dsp-level:0
\[
(7)\quad\textbf{the rung refinement (where the binary sector does reach):}
\]

%dsp-level:0
\[
\qquad m^2_{\mathrm{pole}}(m)=4\,\hat\mu'^2\sin^2\!\bigl(m\,\tfrac{\Theta_R}4\bigr),
\]

%dsp-level:0
\[
\qquad m\in\{k,\,k{-}2,\dots,-k\}
\]

%dsp-level:0
\[
\qquad [\mathrm{FT.60}(5);\ \mathrm{Def\text{-}H.38};\ \text{Census FT.63}(6)]:
\]

%dsp-level:0
\[
\boxed{\ \text{the rung}\ m\ \ \text{is register-reachable}\iff m\,\delta_v\in\pi\mathbb Z\iff D_S(m)=2-2\cos\!\bigl(m\,\tfrac{\delta_v}2\bigr)\in\{0,\,2,\,4\}\ }
\]

%dsp-level:0
\[
\qquad [\text{machine: the exact iff on every census ladder}\ m\le4\ \text{and by the period law}\ m\in\bigl(\tfrac{\pi}{\delta_v}\bigr)\mathbb Z\ [\mathrm{FT.76}(3)]]:
\]

%tbl-level:0
{\setlength{\tabcolsep}{4pt}
\begin{longtable}{>{\raggedright\arraybackslash}p{\dimexpr0.2059\textwidth-2\tabcolsep\relax}|>{\raggedright\arraybackslash}p{\dimexpr0.1076\textwidth-2\tabcolsep\relax}|>{\raggedright\arraybackslash}p{\dimexpr0.3089\textwidth-2\tabcolsep\relax}|>{\raggedright\arraybackslash}p{\dimexpr0.3776\textwidth-2\tabcolsep\relax}}
\toprule
floor & $\delta_v/\pi$ & reachable rungs $m\le\allowbreak 4$ & the ladder values there \\
\addlinespace[2pt]
\midrule
\endfirsthead
\toprule
floor & $\delta_v/\pi$ & reachable rungs $m\le\allowbreak 4$ & the ladder values there \\
\addlinespace[2pt]
\midrule
\endhead
\bottomrule
\endlastfoot
tetrahedron & $1$ & all $m$ & $(2,\allowbreak \,\allowbreak 4,\allowbreak \,\allowbreak 2,\allowbreak \,\allowbreak 0)$ — all in $\{0,2,4\}$ \\
cube & $\tfrac12$ & $m\equiv0\ (2)$ & $D_S(2)=\allowbreak 2,\allowbreak \quad D_S(4)=\allowbreak 4$ \\
octahedron & $\tfrac23$ & $m\equiv0\ (3)$ & $D_S(3)=\allowbreak 4$ \\
dodecahedron & $\tfrac15$ & none at $m\le\allowbreak 4$; $m\equiv0\ (5)$ & $D_S(5)=\allowbreak 2$ \\
icosahedron & $\tfrac13$ & $m\equiv0\ (3)$ & $D_S(3)=\allowbreak 2$ \\
\end{longtable}
}

%dsp-level:0
\[
\textbf{the register captures the defect ladder's rational values}\ \{0,2,4\}\ \textbf{and never its irrational ones}\ —\ \text{the}\ \mathbb Z_2\ \text{parity of the defect angle is a derived read, the real angle is}\ \neg;
\]

%dsp-level:0
\[
(8)\quad\textbf{the calibration-weight obstruction (the winding route):}
\]

%dsp-level:0
\[
\qquad \Theta_R=2\pi k\ \ \textbf{is ruled out as a corpus relation:}
\]

%dsp-level:0
\[
\qquad k=\oint ds^2/(2\mathfrak h^2)\ \ \text{carries}\ \mathfrak h\text{-weight}\ {-}2
\]

%dsp-level:0
\[
\qquad [\mathrm{FT.31}(3);\ \mathrm{D2.9}];
\]

%dsp-level:0
\[
\Theta_R=\iint_S K\,dA\ \ \text{carries}\ \mathfrak h\text{-weight}\ 0
\]

%dsp-level:0
\[
\qquad [\mathrm{Def\text{-}H.21}(6)];
\]

%dsp-level:0
\[
\qquad \text{a proportionality across calibration weights is}\ \neg\,\mathfrak h\text{-free}
\]

%dsp-level:0
\[
\qquad [\mathrm{CT2.11}\ \text{criterion};\ \mathrm{FT.55}(1)\ \text{re-cited}]:
\]

%dsp-level:0
\[
\textbf{the winding register cannot supply the frame angle either};
\]

%dsp-level:0
\[
(9)\quad\textbf{the chart-side locus stands:}
\]

%dsp-level:0
\[
\qquad \delta_v=2\pi-d\,\frac{(L-2)\pi}L\ \ \text{from the integer card data}\ (d,L),
\]

%dsp-level:0
\[
\qquad V\,\delta_v=4\pi
\]

%dsp-level:0
\[
\qquad [\text{Descartes/Gauss--Bonnet};\ \text{Census FT.63}(4);\ \mathrm{FT.47}(1);\ \mathrm{BM.60}(3);\ \text{machine: exact in}\ \pi\text{-units}]:
\]

%dsp-level:0
\[
\text{a loop cleanly enclosing one vertex carries}\ \ \Theta_R=\delta_v
\]

%dsp-level:0
\[
\qquad [\mathrm{FT.47}(1);\ \mathrm{Def\text{-}H.21}(5),(6)],
\]

%dsp-level:0
\[
\qquad D_S=4\sin^2\!\tfrac{\delta_v}4
\]

%dsp-level:0
\[
\qquad [\mathrm{Def\text{-}H.37}(2);\ \mathrm{FT.54}(3),(4)]:
\]

%dsp-level:0
\[
\textbf{the defect spectrum is chart-side face geometry consumed through}\ \mathrm{Def\text{-}H.37}(2)\text{'s interface, whose non-derived clause is priced by}\ \mathrm{FT.54}(1)\text{'s form-uniqueness};
\]

%dsp-level:0
\[
\textbf{the ledger realization of clean-enclosure (vertex-figure) loops is}\ \neg\textbf{constructed}\ —\ \text{the corner convention is undeclared}
\]

%dsp-level:0
\[
\qquad [\text{Census FT.63}(4)\ \mathrm{Rmk};\ \text{Census FT.63}(6)\ \mathrm{Rmk}];
\]

%dsp-level:0
\[
(10)\quad\textbf{the structured obstruction record:}
\]

\begin{verbatim}
OBSTRUCTION: holonomy_lift_not_derived
MISSING: vertex_figure_loop_realization | Symk_assignment [FT.55(4)] | chiral_assignment [FT.125(4)] | coarse_graining_lift 
DERIVED_NO_GO: two_torsion_image (6) | calibration_weight (8)
PARTIAL_LIFT: Z2_wilson_sector (2) | parity_lock (3) | integer_winding_sealing_sector (4) | self_dual_floor (6),(7) | rung_law (7)
\end{verbatim}

%dsp-level:0
\[
(11)\quad\textbf{status:}
\]

%dsp-level:0
\[
\qquad \textbf{derived:}\ (1)\text{--}(9);
\]

%dsp-level:0
\[
\qquad \textbf{record issued:}\ (10);
\]

%dsp-level:0
\[
\qquad \textbf{outcome:}\ \text{the obstruction theorem, with the partial lift of}\ (2)\text{--}(4),
\]

%dsp-level:0
\[
(7)\ \text{and the self-dual-floor lift of}\ (6);
\]

%dsp-level:0
\[
\neg\text{a declared holonomy: no continuous angle is inserted anywhere in this tier}\ —\ \Theta_R\ \text{occurs only inside re-citations of the chart-side reads and inside the refuted lift equation}\ [\text{machine: the occurrence census}];
\]

%dsp-level:0
\[
\neg\text{a P4 addition},\quad\neg\text{a new declaration},\quad\neg\text{a measured value}\ [\mathrm{P5}];
\]

%dsp-level:0
\[
\qquad \text{the conditionalities}\ (\mathrm{Def\text{-}H.37}(2),\ \mathrm{Def\text{-}H.38},\ \mathrm{M1}/\mathrm{M2}\ [\mathrm{Def\text{-}H.20}],\ \mathrm{CT2.6})\ \text{are inherited at the point of use}\ [\text{A §0.2a}].
\]

\textbf{Req.} A4, A7, A8 (Part 1); D1.11, D2.9, D3.3, D3.5, D4.2, CT0.5, CT0.6, CT0.9 (Part 1); CT2.11, CT4.1, CT4.2, CT4.3, CT4.4, CT4.5, EM.9b, EW.32, EW.33 (Part 2); Def-EM.20, EM.20, EM.21 (Part 2); P5 (Parts 1–2); Def-H.20, Def-H.21(5), Def-H.21(6), Def-H.37(1), Def-H.37(2), Def-H.38, FT.22, FT.24(1), FT.24(2), FT.25(1), FT.25(2), FT.25(3), FT.31(3), FT.47(1), FT.54(1), FT.54(3), FT.54(4), FT.55(1), FT.55(4), FT.60(5), FT.73(1), FT.76(3), FT.116, FT.117(1), FT.125(4) (Part 3); Census FT.28(5), FT.35(3), FT.49(3), FT.63(4), FT.63(6), FT.64(5) (row-wise); Def-BM.8(1), Def-BM.8(3), Def-BM.11, Def-BM.12(2), Def-BM.12(7), Def-BM.13(2), Def-BM.14(3), Def-BM.14(4), BM.1(2), BM.2(1), BM.3(1), BM.10(5), BM.28(6), BM.40(3), BM.52(1), BM.52(3), BM.57(5), BM.57(6), BM.58(1), BM.58(2), BM.58(3), BM.58(4), BM.59(1), BM.60(3) (this part); exhaustive sign-assignment censuses, exact rational coset arithmetic modulo $\pi$ and $2\pi$, GF(2) parity identities, Burnside orbit counts on finite ladders [math; machine: \texttt{checks/\allowbreak{}verify\_\allowbreak{}H\_\allowbreak{}organic\_\allowbreak{}bridge\_\allowbreak{}post\_\allowbreak{}bm14.py} §D].

\textbf{Pf.}

%dsp-level:0
\[
(1):\ \mathrm{A4}/\mathrm{A7}\ \text{re-cited: the branch phase and the edge map};
\]

%dsp-level:0
\[
\qquad \alpha_e\in\{0,\pi\}\ [\mathrm{CT0.9}];
\]

%dsp-level:0
\[
\qquad U(c)=\prod\sigma_e^{\operatorname{or}(e)}=e^{i\sum\operatorname{or}(e)\alpha_e}\ [\text{machine: the exponential of a sum over}\ \{0,\pi\}]
\]

%dsp-level:0
\[
(2):\ \sigma^{\pm1}=\sigma\ \text{on}\ \{\pm1\}\ [\text{machine}];
\]

%dsp-level:0
\[
\qquad \text{exhaustive over}\ 2^n\ \text{assignments for}\ n
\]

%dsp-level:0
\[
= 2,
\]

%dsp-level:0
\[
\dots,
\]

%dsp-level:0
\[
7\ \text{and both orientation patterns: the product equals}\ (-1)^{\#\{\sigma_e=-1\}}\ \text{at every one}\ [\text{machine}];
\]

%dsp-level:0
\[
\qquad \text{concatenation: products multiply; reversal: each}\ \sigma\ \text{is its own inverse}\ [\text{machine}]
\]

%dsp-level:0
\[
(3):\ \mathrm{Def\text{-}BM.13}(2)\text{'s exhaustive census re-run}\ [\text{machine}];
\]

%dsp-level:0
\[
\qquad |q|\equiv q\ (2)\ \text{for}\ q\in\{-1,0,1\}\ [\text{machine: three cases}];
\]

%dsp-level:0
\[
\qquad \textstyle\sum|q_i|\ge|\sum q_i|\ [\text{machine: exhaustive over}\ 27\ \text{triples}];
\]

%dsp-level:0
\[
\qquad \text{the end-pair census:}\ q
\]

%dsp-level:0
\[
= 0\iff\sigma_s
\]

%dsp-level:0
\[
= \sigma_t\ (2\ \text{pairs},\ \text{even flips}),
\]

%dsp-level:0
\[
\ |q|
\]

%dsp-level:0
\[
= 1\iff\sigma_s\ne\sigma_t\ (1+1\ \text{pairs},\ \text{odd flips})\ [\text{machine}]
\]

%dsp-level:0
\[
(4):\ \mathrm{FT.22}/\mathrm{FT.24}(1),(2)/\mathrm{FT.25}(1)\text{--}(3)\ \text{re-cited};
\]

%dsp-level:0
\[
\qquad \Theta\in2\mathbb Z\implies e^{i\pi\Theta}=1\ [\text{machine: even integers}];
\]

%dsp-level:0
\[
\qquad 4\mathbb Z\subseteq2\mathbb Z\ [\text{machine}]
\]

%dsp-level:0
\[
(5):\ \text{each row re-cited from its own tag and classified}\ [\text{machine: the census structure against A8's six groups}];
\]

%dsp-level:0
\[
\qquad 7+3+3+3=16,\ 2^4=16,\ 32-16=16\ [\mathrm{CT0.5};\ \mathrm{CT0.6};\ \text{machine}]
\]

%dsp-level:0
\[
(6):\ \text{the image:}\ \{2\pi n+\pi j:\ n,j\in\mathbb Z\}\equiv\{0,\pi\}\ (\mathrm{mod}\ 2\pi)\ [\text{machine}];
\]

%dsp-level:0
\[
\qquad \text{the state angles modulo}\ \pi:\ \pi
\]

%dsp-level:0
\[
\equiv 0,
\]

%dsp-level:0
\[
\ 2\pi
\]

%dsp-level:0
\[
\equiv 0,
\]

%dsp-level:0
\[
\ \tfrac\pi2\not
\]

%dsp-level:0
\[
\equiv 0,
\]

%dsp-level:0
\[
\ \tfrac{2\pi}3\not
\]

%dsp-level:0
\[
\equiv 0,
\]

%dsp-level:0
\[
\ \tfrac\pi5\not
\]

%dsp-level:0
\[
\equiv 0,
\]

%dsp-level:0
\[
\ \tfrac\pi3\not
\]

%dsp-level:0
\[
\equiv 0,
\]

%dsp-level:0
\[
\ \text{pairwise distinct in}\ \mathbb Q/\mathbb Z\ [\text{machine: exact Fractions}];
\]

%dsp-level:0
\[
\qquad \text{the count: two in the image, four outside}
\]

%dsp-level:0
\[
(7):\ \cos\theta\in\{0,\pm1\}\iff\theta\in\tfrac\pi2\mathbb Z\ [\text{math}]
\]

%dsp-level:0
\[
\implies D_S(m)\in\{0,2,4\}
\]

%dsp-level:0
\[
\iff m\,
\]

%dsp-level:0
\[
\tfrac{\delta_v}2\in\tfrac\pi2\mathbb Z
\]

%dsp-level:0
\[
\iff m\,
\]

%dsp-level:0
\[
\delta_v\in\pi\mathbb Z\ [\text{machine}];
\]

%dsp-level:0
\[
\qquad \text{the five ladders evaluated exactly: the tetrahedral}\ (2,4,2,0);
\]

%dsp-level:0
\[
\ \text{the cube's}\ D_S(2)
\]

%dsp-level:0
\[
= 2-2\cos\tfrac\pi2
\]

%dsp-level:0
\[
= 2,
\]

%dsp-level:0
\[
\ D_S(4)
\]

%dsp-level:0
\[
= 2-2\cos\pi
\]

%dsp-level:0
\[
= 4;
\]

%dsp-level:0
\[
\ \text{the octahedron's}\ D_S(3)
\]

%dsp-level:0
\[
= 2-2\cos\pi
\]

%dsp-level:0
\[
= 4;
\]

%dsp-level:0
\[
\ \text{the icosahedron's}\ D_S(3)
\]

%dsp-level:0
\[
= 2-2\cos\tfrac\pi2
\]

%dsp-level:0
\[
= 2;
\]

%dsp-level:0
\[
\ \text{the dodecahedron's}\ D_S(5)
\]

%dsp-level:0
\[
= 2\ [\text{machine}];
\]

%dsp-level:0
\[
\qquad \text{the irrational members}\ 2-\sqrt2,
\]

%dsp-level:0
\[
\ 2+\sqrt2,
\]

%dsp-level:0
\[
\ 1,
\]

%dsp-level:0
\[
\ 3,
\]

%dsp-level:0
\[
\ \tfrac{3-\sqrt5}2,
\]

%dsp-level:0
\[
\ \tfrac{5-\sqrt5}2,
\]

%dsp-level:0
\[
\ 2-\sqrt3,
\]

%dsp-level:0
\[
\ 2-2\cos\tfrac\pi{10}\ \text{are separated from}\ \{0,2,4\}\ \text{by their minimal polynomials evaluated at}\ 0,
\]

%dsp-level:0
\[
2,
\]

%dsp-level:0
\[
4\ [\text{machine: every value nonzero}]
\]

%dsp-level:0
\[
(8):\ \mathrm{FT.55}(1)\ \text{re-cited verbatim};
\]

%dsp-level:0
\[
\qquad (9):\ \delta_v\ \text{evaluated at}\ (d,L)
\]

%dsp-level:0
\[
= (3,3),
\]

%dsp-level:0
\[
(3,4),
\]

%dsp-level:0
\[
(4,3),
\]

%dsp-level:0
\[
(3,5),
\]

%dsp-level:0
\[
(5,3):\ \ \delta_v/\pi
\]

%dsp-level:0
\[
= (1,\tfrac12,\tfrac23,\tfrac15,\tfrac13)\ [\text{machine}];
\]

%dsp-level:0
\[
\qquad V\delta_v=4\pi\ \text{at all five}\ [\text{machine}]
\]

%dsp-level:0
\[
(10)\text{--}(11):\ \text{the record assembled from}\ (6),(8),(9);
\]

%dsp-level:0
\[
\qquad \text{the occurrence census of}\ \Theta_R\ \text{in this tier}\ [\text{machine}];
\]

%dsp-level:0
\[
\qquad \text{the standing-rule scans}\ [\text{machine}]. \ \blacksquare
\]

\textbf{Phys.} Ask the ledger for a real holonomy and it answers in bits. The edge-phase register is binary, so every closed-loop product is plus or minus one, oriented either way; the wall content of any transport is parity-locked to its charge core, and at the composite the lock reads exactly — total wall parity equals the absolute charge modulo two, and the wall count bounds the charge from below. The winding and sealing registers are integers, and every integer winding contributes trivially to a phase, so their content is the rung, not an angle. Then the force decomposition is censused block by block at the sealed junction, and not one block carries a continuous vertex angle: the seals are finite and integer, the metric is charge-blind, the weak block is class-shared branch data, colour is annihilated by the trace law on anything sealed, electromagnetism is chart-side at a symbolic coupling, the dark blocks are uniform or locked, and the channel angles are line data attached to no baryon row. Any lift built from that material lands in the two-element subgroup — zero or pi — while the defect spectrum spreads six distinct values whose angles include a fifth of pi and a third of pi: four of the six states are provably out of reach, and the reachable sublattice is exactly the self-dual floor plus the ladder rungs whose defect value is zero, two or four. The winding route is shut separately, by calibration weight: the rung counts hops, the frame angle counts curvature, and no pure number bridges the two. So the tier records what it has — the binary Wilson sector, the parity lock, the integer sectors, the rung law — and issues the obstruction in structured form, with the missing enumeration named: the ledger realization of a loop that cleanly encloses one vertex. No angle is declared anywhere.

\pseudohead{theorem}{\texorpdfstring{BM.68}{BM.68}}{\texorpdfstring{The constrained joint entropy as a derived fiber cardinality}{The constrained joint entropy as a derived fiber cardinality}}\label{tag:BM-68}
\textbf{Thm.}

%dsp-level:0
\[
(1)\quad\textbf{the constrained fiber (an abbreviation of derived objects):}
\]

%dsp-level:0
\[
\qquad \Omega(E,\Theta,Q)\ :
\]

%dsp-level:0
\[
= \ \#\Bigl\{(R,\sigma,g):\ \ E_{\mathrm{int}}(R)=E,\ \ Q(\sigma)=Q,\ \ \Theta(R,g)=\Theta,\ \ \text{the sealing and admissibility constraints hold}\Bigr\}:
\]

%dsp-level:0
\[
\text{the census of ledger states passing the filter chain}\ \mathrm{F0}\text{--}\mathrm{F7}
\]

%dsp-level:0
\[
\qquad [\text{Census FT.64}(1):\ \text{configuration validity};\ \text{registration};\ \text{integrality};\ \text{loop quantization};\ \text{stationarity};\ \text{regime};\ \text{relabeling quotient};\ \text{register product}];
\]

%dsp-level:0
\[
S_{\mathrm{joint}}(E,\Theta,Q)\ :=\ \ln\Omega(E,\Theta,Q)
\]

%dsp-level:0
\[
\qquad [\mathrm{C9}:\ \text{entropy in nats};\ \text{machine: cardinalities of finite sets}];
\]

%dsp-level:0
\[
(2)\quad\textbf{the factorization (theorem on the tested domains):}
\]

%dsp-level:0
\[
\qquad \text{the registers are independent}\ —\ \text{the signs factor from the}\ \kappa\text{-ledger}\ [\mathrm{CT0.10};\ \text{Census FT.51}(5)],
\]

%dsp-level:0
\[
\ \ E_{\mathrm{int}}\ \text{is constant on every end-sign fibre}\ [\mathrm{Def\text{-}BM.13}(6);\ \mathrm{BM.28}(3);\ \text{machine: exhaustive over the four end pairs}],
\]

%dsp-level:0
\[
\ \text{the sealing grammar is its own register}\ [\mathrm{FT.41}(b);\ \text{Census FT.65}(2)]:
\]

%dsp-level:0
\[
\boxed{\ \Omega(E,\Theta,Q)\ =\ \Omega_{\mathrm{count}}(E)\ \cdot\ \Omega_{\mathrm{charge}}(Q)\ \cdot\ \Omega_{\mathrm{geom}}(\Theta;E,Q)\ }
\]

%dsp-level:0
\[
\qquad \text{``the three registers multiply, none converts into another''}
\]

%dsp-level:0
\[
\qquad [\text{Census FT.51}(5),\ \text{FT.65}(2);\ \text{Census FT.64}(1)\ \mathrm{F7}]:
\]

%dsp-level:0
\[
S_{\mathrm{joint}}\ =\ \ln\Omega_{\mathrm{count}}(E)\ +\ \ln\Omega_{\mathrm{charge}}(Q)\ +\ \ln\Omega_{\mathrm{geom}}(\Theta;E,Q)
\]

%dsp-level:0
\[
\qquad [\text{machine: the cardinality identity on every tested finite domain}];
\]

%dsp-level:0
\[
(3)\quad\textbf{the first factor (derived):}
\]

%dsp-level:0
\[
\qquad \boxed{\ \Omega_{\mathrm{count}}(E)\ =\ 2^E\ }
\]

%dsp-level:0
\[
\qquad [\text{one binary registration per directed real edge}:\ \mathrm{CT0.9}:\ 1\ \text{real edge}=1\ \text{bit}=\ln2\ \text{nats};\ \mathrm{C9};\ \mathrm{CT1.29}]:
\]

%dsp-level:0
\[
I_{\mathrm{hid}}=E_{\mathrm{int}}\ln2=\ln2^{E_{\mathrm{int}}}
\]

%dsp-level:0
\[
\qquad [\mathrm{Def\text{-}H.9}(2);\ \mathrm{FT.17};\ \text{Census FT.35}(2);\ \text{Census FT.44}(4);\ \text{machine}]:
\]

%dsp-level:0
\[
\qquad S_{\mathrm{count}}=E\ln2\ \ \text{is exactly}\ \mathrm{Def\text{-}BM.13}(5)\text{'s first summand};
\]

%dsp-level:0
\[
(4)\quad\textbf{the second factor (derived):}
\]

%dsp-level:0
\[
\qquad \boxed{\ \Omega_{\mathrm{charge}}(Q)\ =\ W(Q)\ =\ \bigl[z^{Q}\bigr]\bigl(2+z+z^{-1}\bigr)^{k}\ },
\]

%dsp-level:0
\[
\qquad k=3:\quad W=(1,\,6,\,15,\,\mathbf{20},\,15,\,6,\,1),
\]

%dsp-level:0
\[
\qquad \sum_QW(Q)=64
\]

%dsp-level:0
\[
\qquad [\mathrm{Def\text{-}BM.8}(6);\ \mathrm{Def\text{-}BM.13}(3);\ \mathrm{BM.28}(2);\ \text{machine: two agreeing routes}]:
\]

%dsp-level:0
\[
S_{\mathrm{charge}}=\ln W(Q)=\Sigma_{\mathrm{cl}}(Q)\qquad[\mathrm{Def\text{-}BM.13}(4)]:\qquad\text{the second summand};
\]

%dsp-level:0
\[
(5)\quad\textbf{the third factor at stationarity (derived):}
\]

%dsp-level:0
\[
\qquad \text{on the stationary locus}\ \kappa_e\equiv\kappa_0\in\mathbb Z\ \ \text{exactly along each connected materialized chain}
\]

%dsp-level:0
\[
\qquad [\mathrm{FT.53}(1),(2);\ \text{Census FT.64}(1)\ \mathrm{F4}],
\]

%dsp-level:0
\[
\text{and the semiclassical survivor list carries one orbit per rung}
\]

%dsp-level:0
\[
\qquad [\text{Census FT.64}(5)(\mathrm c):\ \text{the R2 ladder};\ \mathrm{FT.53}(3)]:
\]

%dsp-level:0
\[
\qquad \text{a uniform profile is fixed by the swap/relabeling quotient}:
\]

%dsp-level:0
\[
\boxed{\ \Omega_{\mathrm{geom}}^{\mathrm{stat}}(\Theta)\ =\ 1\ \ \text{at every stationary rung},\qquad S_{\mathrm{geom}}^{\mathrm{stat}}\ =\ 0\ }
\]

%dsp-level:0
\[
\qquad [\text{machine: the Burnside count on a single-element set; the swap fixes}\ (\kappa_0,\kappa_0,\kappa_0)]:
\]

%dsp-level:0
\[
\textbf{on the baryon card:}
\]

%dsp-level:0
\[
\qquad \text{the strict-spacelike minimum sits on the locus:}\quad\kappa=(1,1,1),\quad K=3,\quad\Theta_{\mathrm{int}}=6
\]

%dsp-level:0
\[
\qquad [\text{Census FT.65}(4)\ \Sigma_{12}\ \text{row};\ \mathrm{FT.24}(1):\ \kappa\in\mathbb Z\ \text{at materialized digons};\ \text{machine: the unique profile, the unique orbit}]:
\]

%dsp-level:0
\[
\boxed{\ S_{\mathrm{joint}}\bigl(6,\ \Theta_{\mathrm{stat}},\ Q\bigr)\ =\ 6\ln2\ +\ \ln W(Q)\ +\ 0\ }
\]

%dsp-level:0
\[
\textbf{the refined rest input is the complete joint entropy of the stationary baryon fibre, on its ledger part:}
\]

%dsp-level:0
\[
\qquad \text{the EM summand is the chart-side transit-channel energy,}\ \neg\text{a ledger count}
\]

%dsp-level:0
\[
\qquad [\mathrm{Def\text{-}BM.12}(6);\ \mathrm{Def\text{-}BM.13}(5);\ \mathrm{BM.66}(6);\ \text{machine: the alphabet disjointness}]:
\]

%dsp-level:0
\[
\textbf{the additive law is a derived stationarity slice,}\ \neg\textbf{a postulate};
\]

%dsp-level:0
\[
\qquad \text{the off-stationary fibers cancel pairwise in}\ Z\ \text{—}\ \neg\text{a semiclassical spectral contribution}
\]

%dsp-level:0
\[
\qquad [\mathrm{CT1.20};\ \mathrm{FT.53}(4);\ \text{Census FT.64}(1)\ \mathrm{F8}\ \text{note}]:
\]

%dsp-level:0
\[
\qquad \neg\text{a mass read requires an off-stationary count};
\]

%dsp-level:0
\[
(6)\quad\textbf{the general fiber (the named missing enumeration):}
\]

%dsp-level:0
\[
\qquad \textbf{the Platonic floors}\ —\ \text{derived in closed form: the on-shell lattice}\ \ker\varphi_S,
\]

%dsp-level:0
\[
\ \ \text{index}\ 2^{F-1}\ [\mathrm{FT.43}(3);\ \text{Census FT.64}(2)];
\]

%dsp-level:0
\[
\ \text{the level counts}\ A_S(K)\ \text{with the exact character-sum closed form}\ [\text{Census FT.64}(4)];
\]

%dsp-level:0
\[
\ \text{the Burnside orbits}\ M_S(K),
\]

%dsp-level:0
\[
\ M^+_S(K)\ [\text{Census FT.65}(1)];
\]

%dsp-level:0
\[
\ \text{the sealing grammar}\ g(V)
\]

%dsp-level:0
\[
= (3,\,385,\,25,\,292938721075,\,164395)\ [\mathrm{FT.41}(b);\ \text{Census FT.63}(7)];
\]

%dsp-level:0
\[
\ \text{the register product}\ 2^{2E}\!\cdot\!M_S(K)\cdot g(V)\cdot(\text{attachment data})\ [\text{Census FT.65}(2)]:
\]

%dsp-level:0
\[
\textbf{the polygon floors}\ —\ \text{the dihedral necklaces}\ M(n,k)\ \text{in closed Burnside form}
\]

%dsp-level:0
\[
\qquad [\text{Census FT.51}(1)\text{--}(3)]:
\]

%dsp-level:0
\[
\textbf{the baryon junction}\ \Sigma_{12}:
\]

%dsp-level:0
\[
\qquad \text{the winding lattice is}\ \mathbb Z^3\ \text{with the swap quotient}\ \mathbb Z_2
\]

%dsp-level:0
\[
\qquad [\text{Census FT.65}(4)\ \Sigma_{12}\ \text{row}];
\]

%dsp-level:0
\[
\text{the fiber over the full lattice at fixed}\ K\ \ \textbf{is infinite}\ —\ \text{the witness}\ (K{+}n,\,{-}n)
\]

%dsp-level:0
\[
\qquad [\text{Census FT.64}(1)\ \mathrm{F5};\ \text{machine}]:
\]

%dsp-level:0
\[
\qquad \text{the enumeration requires the regime datum}\ (\text{null-permitting}\ /\ \text{strict-spacelike}\ /\ \text{uniform}\ /\ \text{signature-mixed})\ \text{that the card does not fix}:
\]

%dsp-level:0
\[
\boxed{\ \text{the exact missing enumeration}\ =\ \text{the regime-scoped orbit count}\ \ \#\bigl\{\kappa\in\mathbb Z^3:\ \kappa_1{+}\kappa_2{+}\kappa_{12}=K\bigr\}\big/\mathbb Z_2\ \ \text{at the baryon card's materialized rung}\ }
\]

%dsp-level:0
\[
\textbf{the scoped rows this tier derives:}
\]

%dsp-level:0
\[
\qquad \text{the strict regime}\ \kappa\ge1:\quad\#\text{profiles}=\binom{K-1}2,
\]

%dsp-level:0
\[
\qquad \text{orbits}\ (1,\,2,\,4,\,6)\ \ \text{at}\ K=(3,\,4,\,5,\,6)
\]

%dsp-level:0
\[
\qquad [\text{machine: stars-and-bars and Burnside agree with exhaustive enumeration}];
\]

%dsp-level:0
\[
\text{the null-permitting regime}\ \kappa\ge0:\quad\#\text{profiles}=\binom{K+2}2,
\]

%dsp-level:0
\[
\qquad \text{orbits}\ (2,\,6)\ \ \text{at}\ K=(1,\,3)
\]

%dsp-level:0
\[
\qquad [\text{machine: exhaustive}];
\]

%dsp-level:0
\[
\qquad \text{the uniform ladder: one orbit per rung}\ [(5);\ \text{machine}];
\]

%dsp-level:0
\[
(7)\quad\textbf{the two currencies and the bridge-candidate census:}
\]

%dsp-level:0
\[
\qquad \textbf{the bit/charge currency:}\ \ E\ln2+\ln W(Q)\ —\ \text{cardinalities, integers};
\]

%dsp-level:0
\[
\qquad \textbf{the geometric/holonomy currency:}\ \ \Theta_R\ \text{enters as a value,}\ \neg\text{a count}:
\]

%dsp-level:0
\[
\quad m^2_{\mathrm{pole}}=\hat\mu'^2\cdot4\sin^2\!\tfrac{\Theta_R}4
\]

%dsp-level:0
\[
\qquad [\mathrm{Def\text{-}H.37}(2);\ \mathrm{FT.54}(4)],
\]

%dsp-level:0
\[
\qquad D_S\ \text{exact algebraic in four fields}
\]

%dsp-level:0
\[
\qquad [\mathrm{Def\text{-}BM.14}(3),(4)]:
\]

%dsp-level:0
\[
\textbf{every candidate bridge species, tested:}
\]

%dsp-level:0
\[
(\mathrm a)\quad\textbf{a derivative of}\ S_{\mathrm{joint}}:
\]

%dsp-level:0
\[
\qquad \Theta\ \text{is integer-valued};
\]

%dsp-level:0
\[
\ \ \text{the finite difference along the stationary ladder is}\ \ \ln\!\bigl(\Omega^{\mathrm{stat}}/\Omega^{\mathrm{stat}}\bigr)
\]

%dsp-level:0
\[
= \ln1
\]

%dsp-level:0
\[
= 0\ \ \text{while the defect value varies across six distinct values}
\]

%dsp-level:0
\[
\qquad [(5);\ \mathrm{Def\text{-}BM.14}(4);\ \text{machine}]:
\]

%dsp-level:0
\[
\qquad \textbf{the count is flat where the angle is not}\ —\ \text{refuted as a bridge};
\]

%dsp-level:0
\[
(\mathrm b)\quad\textbf{a Legendre transform under constraints:}
\]

%dsp-level:0
\[
\qquad \neg\exists\ \text{a joint potential of which both currencies are reads}\ —\ \text{the three deficit functionals are pairwise distinct on digon witnesses}\ (2\ln2\ \text{vs}\ 0;\ \ 0\ \text{vs}\ 3;\ \ 2\ln2\ne3\ \text{via}\ \tfrac12<\ln2<1)
\]

%dsp-level:0
\[
\qquad [\mathrm{BM.35}(6);\ \text{machine: the exact enclosure}]:
\]

%dsp-level:0
\[
\qquad \neg\text{a Legendre pair is derivable};
\]

%dsp-level:0
\[
(\mathrm c)\quad\textbf{a constrained extremum:}
\]

%dsp-level:0
\[
\qquad \arg\max_Q S_{\mathrm{joint}}(6,\cdot,Q)\ =\ \text{the neutral class}
\]

%dsp-level:0
\[
\qquad [W(0)=20>15=W(\pm1);\ \text{machine}]:
\]

%dsp-level:0
\[
\qquad \textbf{derived, but a class selection,}\ \neg\textbf{a currency conversion};
\]

%dsp-level:0
\[
(\mathrm d)\quad\textbf{a trace normalization:}
\]

%dsp-level:0
\[
\qquad I_F=4\,\widehat Q^2\ \ \text{is an exact integer on the information side}
\]

%dsp-level:0
\[
\qquad [\mathrm{BM.66}(3)];
\]

%dsp-level:0
\[
\qquad \text{the defect side's trace normalization is the seed doublet's dimension:}\ \ (U-1)^\dagger(U-1)
\]

%dsp-level:0
\[
= 4\sin^2\!\tfrac\phi2\cdot I_2
\]

%dsp-level:0
\[
\qquad [\mathrm{FT.54}(1);\ \mathrm{EW.32}]:
\]

%dsp-level:0
\[
\text{both exact, incommensurate species}\ —\ \text{a count against an angle}
\]

%dsp-level:0
\[
\qquad [\text{machine: the token census}]:
\]

%dsp-level:0
\[
\qquad \neg\text{a ratio between them is derivable};
\]

%dsp-level:0
\[
(\mathrm e)\quad\textbf{a Casimir ratio:}
\]

%dsp-level:0
\[
\qquad \text{every Casimir/trace/capacity ratio in the inventory is an exact rational}
\]

%dsp-level:0
\[
\qquad [\mathrm{BM.61}(2),(3):\ \text{integers and ratios of integers};\ \bigl(\tfrac13,\,8,\,\tfrac1{25}\bigr);\ \mathrm{BM.36}(2)]:
\]

%dsp-level:0
\[
\text{the five floors' pole factors generate a degree-}16\ \text{extension of}\ \mathbb Q
\]

%dsp-level:0
\[
\qquad [\mathrm{FT.73}(2);\ \text{machine re-cited}]:
\]

%dsp-level:0
\[
\qquad \textbf{no rational-coefficient formula on integer-valued reads reaches the pole side};
\]

$(8)\quad\textbf{the gap row:}\qquad\text{the bridge's well-posed object}\ [\mathrm{FT.54}(4):\ \text{Gap 7's well-posed object}]\ \ \text{and its single undeclared factor:}$

%dsp-level:0
\[
\boxed{\ \frac{m^2_{\mathrm{pole}}}{m_{\mathrm{hid}}^2}\ =\ \rho_{\mathrm B}\cdot\frac{D_S}{4E^2}\ \in\ \mathfrak q^0,\qquad\rho_{\mathrm B}\ :=\ \frac{\hat\mu'^2\,\mathfrak h^2}{(\ln2)^2}\ }
\]

%dsp-level:0
\[
\qquad [\mathrm{FT.54}(4);\ \text{Census FT.63}(6);\ \text{machine: the dimension row}\ [\hat\mu'^2\mathfrak h^2]=\mathfrak q\cdot\mathfrak q^{-1}=\mathfrak q^0]:
\]

%dsp-level:0
\[
\rho_{\mathrm B}\ \text{is a Cal-F packaging}\ [\mathrm{Def\text{-}H.37}(2);\ \mathrm{FT.54}(4);\ \text{Cal-F logged}]\ —\ \textbf{restated here as the residue's name,}\ \neg\textbf{added};
\]

%dsp-level:0
\[
\qquad \text{the per-state exact algebraic factors:}\quad\frac{D_S}{4E^2}
\]

%dsp-level:0
\[
= \Bigl(\frac1{72},\ \ \frac{2-\sqrt2}{576},\ \ \frac1{576},\ \ \frac{D_S}{3600},\ \ \frac{2-\sqrt3}{3600}\Bigr)
\]

%dsp-level:0
\[
\qquad [\text{Census FT.63}(6);\ \text{machine}];
\]

%dsp-level:0
\[
(9)\quad\textbf{status:}
\]

%dsp-level:0
\[
\qquad \textbf{derived:}\ (1)\text{--}(5),\ (6)\text{'s scoped rows},\ (7),\ (8);
\]

%dsp-level:0
\[
\qquad \textbf{named,}\ \neg\textbf{filled: the general}\ \Sigma_{12}\ \text{fiber}\ (6);
\]

%dsp-level:0
\[
S_{\mathrm{geom}}\ \ \neg\textbf{declared}\ —\ \text{it is derived}\ =0\ \text{at stationarity and named-missing off it};
\]

%dsp-level:0
\[
\qquad \neg\text{a P4 addition},\quad\neg\text{a new dimensionful scale}\ [\mathrm{P1}],\quad\neg\text{a measured value}\ [\mathrm{P5}];
\]

%dsp-level:0
\[
\qquad \text{§Cal-BM stays empty}.
\]

%dsp-level:0
\[
\textbf{the structured record:}
\]

\begin{verbatim}
OBSTRUCTION: geometric_fiber_count_not_derived (the general Sigma_12 fiber)
MISSING: Omega_geom enumeration — the regime-scoped orbit count at the materialized rung
DERIVED: stationary_fiber = 1 (5); Platonic closed forms (6); polygon necklaces (6); strict/null-permitting scoped counts (6)
\end{verbatim}

\textbf{Req.} C9, CT0.9, CT0.10, CT1.20 (Part 1); CT1.29, CT2.1, EW.32 (Part 2); P1, P5 (Parts 1–2); Def-H.9(2), Def-H.37(2), FT.17, FT.41(b), FT.43(3), FT.53(1), FT.53(2), FT.53(3), FT.53(4), FT.54(1), FT.54(4), FT.73(2) (Part 3); Census FT.35(2), FT.44(4), FT.51(1), FT.51(3), FT.51(5), FT.63(6), FT.63(7), FT.64(1), FT.64(2), FT.64(4), FT.64(5), FT.65(1), FT.65(2), FT.65(4) (row-wise); Def-BM.8(6), Def-BM.12(6), Def-BM.13(3), Def-BM.13(4), Def-BM.13(5), Def-BM.13(6), Def-BM.14(3), Def-BM.14(4), BM.28(3), BM.35(6), BM.36(2), BM.61(2), BM.61(3), BM.66(3), BM.66(6) (this part); Cal-F; stars-and-bars and Burnside enumeration, exact rational arithmetic, the prime-log ring with integer certificates, exact radical separations, rational $\ln2$ enclosures [math; machine: \texttt{checks/\allowbreak{}verify\_\allowbreak{}H\_\allowbreak{}organic\_\allowbreak{}bridge\_\allowbreak{}post\_\allowbreak{}bm14.py} §E].

\textbf{Pf.}

%dsp-level:0
\[
(1)\text{--}(2):\ \text{the filter chain re-cited row by row}\ [\text{Census FT.64}(1)];
\]

%dsp-level:0
\[
\qquad \text{the independence rows re-cited and re-run: signs}\perp\kappa\ [\mathrm{CT0.10};\ \text{machine: the product cardinality on the tested domains}],
\]

%dsp-level:0
\[
\ \ E_{\mathrm{int}}\ \text{constant on end-sign fibres}\ [\text{machine: exhaustive over the four end pairs}],
\]

%dsp-level:0
\[
\ \text{sealing separate}\ [\mathrm{FT.41}(b)];
\]

%dsp-level:0
\[
\qquad \Omega=\Omega_{\mathrm{count}}\Omega_{\mathrm{charge}}\Omega_{\mathrm{geom}}\ \text{as a cardinality of a product of finite sets}\ [\text{machine}]
\]

%dsp-level:0
\[
(3):\ 2^6=64\ [\text{machine}];
\]

%dsp-level:0
\[
\qquad I_{\mathrm{hid}}=6\ln2=\ln64\ [\text{machine: the prime-log ring}];
\]

%dsp-level:0
\[
\qquad (4):\ W(Q)\ \text{re-enumerated by both routes}\ [\text{machine}]
\]

%dsp-level:0
\[
(5):\ \mathrm{FT.53}(2)\ \text{re-cited: on the stationary locus}\ \kappa_e=\kappa_0\in\mathbb Z\ \text{exactly};
\]

%dsp-level:0
\[
\qquad \text{the swap}\ (\kappa_1,\kappa_2,\kappa_{12})\mapsto(\kappa_2,\kappa_1,\kappa_{12})\ \text{fixes}\ (\kappa_0,\kappa_0,\kappa_0)\ [\text{machine}];
\]

%dsp-level:0
\[
\qquad \text{the Burnside count on one fixed profile is}\ 1\ [\text{machine}];
\]

%dsp-level:0
\[
\qquad \text{the strict minimum}\ (1,1,1):\ \ \text{unique among}\ \kappa\ge1,\ \textstyle\sum\kappa=3\ [\text{machine: exhaustive}];
\]

%dsp-level:0
\[
\qquad 6\ln2+\ln W(Q)+0\ \text{compared term by term with}\ \mathrm{Def\text{-}BM.13}(5)\ \text{at}\ E_{\mathrm{int}}=6\ [\text{machine}]
\]

%dsp-level:0
\[
(6):\ \binom{K-1}2\ \text{at}\ K=(3,4,5,6)=(1,3,6,10)\ [\text{machine}];
\]

%dsp-level:0
\[
\qquad \text{the Burnside orbits under the swap: fixed points}\ \#\{(a,a,b):\ 2a+b=K,\ a,b\ge1\},
\]

%dsp-level:0
\[
\ \ \text{orbits}
\]

%dsp-level:0
\[
= \tfrac12(\#\text{profiles}+\#\text{fixed}):
\]

%dsp-level:0
\[
\quad K=(3,4,5,6)\mapsto(1,2,4,6)\ [\text{machine: closed form and exhaustive enumeration agree}];
\]

%dsp-level:0
\[
\qquad \kappa\ge0:\ \binom{K+2}2\ \text{at}\ K=(1,3)=(3,10),\ \text{orbits}\ (2,6)\ [\text{machine}];
\]

%dsp-level:0
\[
\qquad \text{the infinite-fiber witness}\ (K{+}n,{-}n,\,0),
\]

%dsp-level:0
\[
\ n\in\mathbb Z_{>0}:\ \ \text{distinct profiles at every}\ n\ [\text{machine: injective in}\ n]
\]

%dsp-level:0
\[
(7)(\mathrm a):\ \ln1=0\ [\text{machine}];
\]

%dsp-level:0
\[
\qquad (\mathrm b):\ 2\ln2
\]

%dsp-level:0
\[
\ne 3\ \text{via the enclosure}\ \tfrac12<\ln2<1\ \text{with tail}\ <2^{-20}\ [\text{machine: exact rational partial sums}];
\]

%dsp-level:0
\[
\qquad (\mathrm c):\ 20>15\ [\text{machine}];
\]

%dsp-level:0
\[
\qquad (\mathrm d):\ \text{the token census of both normalizations}\ [\text{machine}];
\]

%dsp-level:0
\[
\qquad (\mathrm e):\ \mathrm{FT.73}(2)\text{'s degree row re-cited}
\]

%dsp-level:0
\[
(8):\ \frac{4\hat\mu'^2\sin^2(\Theta_R/4)}{4E^2(\ln2)^2/\mathfrak h^2}
\]

%dsp-level:0
\[
= \frac{\hat\mu'^2\mathfrak h^2}{(\ln2)^2}\cdot\frac{D_S}{4E^2}\ \text{at}\ \Theta_R
\]

%dsp-level:0
\[
= \delta_v\ [\text{machine: the Laurent assembly}];
\]

%dsp-level:0
\[
\qquad \text{the five factors re-cited from}\ \text{Census FT.63}(6);
\]

%dsp-level:0
\[
\qquad (9):\ \text{the status and record assembly}. \ \blacksquare
\]

\textbf{Phys.} The joint entropy is not postulated; it is counted. Three registers — the crystallized-edge bits, the charge-filtered end census, and the winding-and-sealing fiber — are independent by census, so their cardinalities multiply and their logarithms add: the additive law of the refined rest input is a cardinality identity, and its first two factors are the corpus's own two-to-the-E and W-of-Q. The third factor is derived exactly where the mass read lives. On the stationary locus the geodesic-constancy theorem makes every surviving profile uniform, a uniform profile is fixed by the swap quotient, so the geometric fiber is a single orbit and its entropy is zero: the refined rest input is the complete joint entropy of the stationary baryon fiber, ledger part, and the off-stationary fibers are pairwise-cancelling content that no mass read consults. Away from stationarity the Platonic substrate carries closed forms — a character sum for the level counts, Burnside for the orbits, the frozen grammar counts for the sealings — but the baryon junction's own general fiber is infinite at fixed winding without a regime datum, and the tier names that missing enumeration instead of guessing a number. Then both currencies are laid side by side and every proposed bridge species is tested: the entropy difference vanishes along the ladder where the defect values vary, no joint potential exists to transform, the extremum picks the neutral class but converts nothing, the trace normalizations are exact on both sides and incommensurate between them, and the Casimir ratios are rational where the pole side is of degree sixteen. The exchange rate between count and angle stays one undeclared dimensionless number — the packaging the interface already logged — and the gap row records it under a name.

\begin{theorem}[{\texorpdfstring{BM.69 — The bridge converged: the three inputs, the non-derivation theorem, and the minimal priced residue}{BM.69 - The bridge converged: the three inputs, the non-derivation theorem, and the minimal priced residue}}]\label{tag:BM-69}
\textbf{Thm.}

%dsp-level:0
\[
(1)\quad\textbf{the three inputs assembled:}\qquad\begin{array}{l|l}\text{the force-operator quadratic invariant} & I_F=\operatorname{Tr}\widehat F^2-\operatorname{Tr}\widehat F_0^2=4\,\widehat Q^2\qquad[\mathrm{BM.66}(3)]\\\text{the holonomy lift} & \mathbb Z_2\ \text{Wilson sector derived};\ \ \text{the real lift obstructed with the record issued}\qquad[\mathrm{BM.67}(2),(6),(10)]\\\text{the constrained joint entropy} & S_{\mathrm{joint}}=E\ln2+\ln W(Q)+\ln\Omega_{\mathrm{geom}},\qquad\Omega_{\mathrm{geom}}^{\mathrm{stat}}=1\qquad[\mathrm{BM.68}(2),(5)]\end{array}
\]

%dsp-level:0
\[
(2)\quad\textbf{the bridge identity in its derived form:}
\]

%dsp-level:0
\[
\qquad \boxed{\ \frac{m^2_{\mathrm{pole}}}{m_{\mathrm{hid}}^2}\ =\ \mathcal R\bigl(E_{\mathrm{int}},\Theta_R\bigr)\ =\ \rho_{\mathrm B}\cdot\frac{D_S(\Theta_R)}{4\,E_{\mathrm{int}}^2},\qquad\rho_{\mathrm B}=\frac{\hat\mu'^2\mathfrak h^2}{(\ln2)^2}\in\mathfrak q^0\ }
\]

%dsp-level:0
\[
\qquad [\mathrm{FT.54}(4);\ \text{Census FT.63}(6);\ \mathrm{BM.68}(8);\ \text{machine}]:
\]

%dsp-level:0
\[
\mathcal R\ \text{is exact algebraic per state once}\ \rho_{\mathrm B}\ \text{is supplied:}
\]

%dsp-level:0
\[
\qquad \text{six states, six values of}\ \tfrac{D_S}{4E^2},\ \text{four fields, six minimal polynomials}
\]

%dsp-level:0
\[
\qquad [\mathrm{Def\text{-}BM.14}(3);\ \text{machine re-cited}];
\]

%dsp-level:0
\[
(3)\quad\textbf{the class-dependence census:}
\]

%dsp-level:0
\[
\qquad \text{the bridge's left-hand class dependence runs through}\ D_S\ \text{alone}\ —\ \text{floor and branch data, i.e. the flavour composition of the class, which is}\ \neg\text{derived}
\]

%dsp-level:0
\[
\qquad [\mathrm{BM.28}(7);\ \mathrm{Def\text{-}BM.8}(5);\ \mathrm{BM.59}(9)];
\]

%dsp-level:0
\[
E_{\mathrm{int}}=6\ \ \text{is class-shared}
\]

%dsp-level:0
\[
\qquad [\mathrm{BM.57}(4);\ \text{Census FT.35}(2)];
\]

%dsp-level:0
\[
\qquad \text{the information side's class dependence}\ \bigl(\ln W(Q):\ \Sigma_{\mathrm B}=\ln\tfrac43;\quad I_F:\ (0,4,4)\bigr)\ \text{is exact}
\]

%dsp-level:0
\[
\qquad [\mathrm{Def\text{-}BM.13}(4);\ \mathrm{BM.66}(3);\ \text{machine}];
\]

%dsp-level:0
\[
\textbf{the conditional pole row (the hypothesis on its face):}
\]

%dsp-level:0
\[
\qquad \textbf{hypothesis:}\ \text{the two baryon classes differ by exactly one branch substitution}\ [\mathrm{BM.59}(9);\ \mathrm{BM.23}(4)\ \text{pattern}]:
\]

%dsp-level:0
\[
\textbf{then}
\]

%dsp-level:0
\[
\qquad \Delta\bigl(m^2_{\mathrm{pole}}\bigr)\ =\ \hat\mu'^2\bigl(D_S(\mathcal F_g,+)-D_S(\mathcal F_g,-)\bigr)\ =\ \Delta_g\ >\ 0,
\]

%dsp-level:0
\[
\qquad \Delta_1=2\hat\mu'^2,\quad\Delta_2=(\sqrt2-1)\hat\mu'^2,\quad\Delta_3=\bigl(2\cos\tfrac\pi{10}-\sqrt3\bigr)\hat\mu'^2
\]

%dsp-level:0
\[
\qquad [\mathrm{Def\text{-}BM.14}(5);\ \text{machine: the certificates}\ 4>2,\ 2>1,\ 3<5]:
\]

%dsp-level:0
\[
\textbf{the same direction as the transport row};
\]

%dsp-level:0
\[
\qquad \text{and}\ \neg\textbf{addable to the informational assembly:}\ \ \text{a cross-currency addition is a category the corpus bars}
\]

%dsp-level:0
\[
\qquad [\mathrm{BM.35}(5);\ \mathrm{BM.60}(6);\ \mathrm{BM.37}(2);\ \text{pointers}];
\]

%dsp-level:0
\[
(4)\quad\textbf{the non-derivation theorem:}
\]

%dsp-level:0
\[
\qquad \boxed{\ \neg\exists\ \text{an exact identity}\ \ m^2_{\mathrm{pole}}=\mathcal B\bigl[E_{\mathrm{int}},\,W(Q),\,\Theta_R,\,I_F\bigr]\ \ \text{with}\ \mathcal B\ \text{built from derived data alone}\ }
\]

%dsp-level:0
\[
\bigl(\text{derived data}\ :=\ \text{exact trace ratios},\ \text{Casimir normalizations},\ \text{constrained fiber counts},\ \text{channel caps},\ \text{already-derived algebraic numbers}\bigr),
\]

%dsp-level:0
\[
\qquad \text{on four independent grounds:}
\]

%dsp-level:0
\[
(\mathrm i)\quad\textbf{the field degree:}\ \ \text{the five floors' pole factors generate a degree-}16\ \text{extension of}\ \mathbb Q;
\]

%dsp-level:0
\[
\ \ \text{every enumerated input}\ \bigl(E_{\mathrm{int}},\ W(Q),\ I_F,\ \text{caps},\ \text{fiber counts}\bigr)\ \text{is integer- or binary-valued};
\]

%dsp-level:0
\[
\ \ \text{a rational formula places every pole value in a field of degree}\ <16\ \text{over the declared constants}
\]

%dsp-level:0
\[
\qquad [\mathrm{FT.73}(2),(3);\ \text{machine re-cited}];
\]

%dsp-level:0
\[
(\mathrm{ii})\quad\textbf{the two-torsion image:}\ \ \text{the register lift of}\ \Theta_R\ \text{reaches}\ \{0,\pi\}\ \text{only}\ —\ \text{four of the six states are unreachable at every winding}
\]

%dsp-level:0
\[
\qquad [\mathrm{BM.67}(6)];
\]

%dsp-level:0
\[
(\mathrm{iii})\quad\textbf{the calibration weight:}\ \ \text{the winding route}\ \Theta_R=2\pi k\ \text{is ruled out}\ —\ \mathfrak h\text{-weight}\ {-}2\ \text{against}\ 0
\]

%dsp-level:0
\[
\qquad [\mathrm{FT.55}(1);\ \mathrm{BM.67}(8)];
\]

%dsp-level:0
\[
(\mathrm{iv})\quad\textbf{the strength bar:}\ \ \neg\exists\ \text{a count}\to\text{strength theorem};\ \ \text{the capacity ratios}\ \bigl(\tfrac13,\,8,\,\tfrac1{25}\bigr)\ \text{are}\ \neg\text{couplings}
\]

%dsp-level:0
\[
\qquad [\mathrm{BM.36}(3);\ \mathrm{BM.61}(8)\ (\mathrm O1);\ \mathrm{FT.73}(5)\ \text{trap pattern}];
\]

%dsp-level:0
\[
(5)\quad\textbf{the Gap-7 reduction (the smallest priced residue):}
\]

%dsp-level:0
\[
\qquad \boxed{\ \text{the bridge residue is exactly:}\quad(\mathrm R1)\ \ \text{the one undeclared dimensionless ratio}\ \ \rho_{\mathrm B}=\frac{\hat\mu'^2\mathfrak h^2}{(\ln2)^2}\qquad[\text{Cal-F packaging};\ \neg\text{added here}];}
\]

%dsp-level:0
\[
(\mathrm R2)\quad\text{the ledger realization of clean-enclosure vertex-figure loops}
\]

%dsp-level:0
\[
\qquad [\mathrm{BM.67}(9),(10);\ \text{Census FT.63}(4)\ \mathrm{Rmk}];
\]

%dsp-level:0
\[
\qquad (\mathrm R3)\quad\text{the flavour composition of the sealed classes}
\]

%dsp-level:0
\[
\qquad [\mathrm{BM.28}(7)]:
\]

%dsp-level:0
\[
\textbf{every other factor of}\ (2)\ \textbf{is derived exact algebra:}
\]

%dsp-level:0
\[
\qquad \text{the six}\ D_S\ \text{with their minimal polynomials};\quad\text{the five}\ \tfrac{D_S}{4E^2};
\]

%dsp-level:0
\[
\quad I_F=(0,4,4);
\]

%dsp-level:0
\[
\quad W=(1,6,15,20,15,6,1);\quad\Omega_{\mathrm{geom}}^{\mathrm{stat}}=1
\]

%dsp-level:0
\[
\qquad [\text{machine: the inventory row}];
\]

%dsp-level:0
\[
(6)\quad\textbf{the no-hiding row:}
\]

%dsp-level:0
\[
\qquad \textbf{the coefficient census of}\ \mathrm{BM.65}\text{--}\mathrm{BM.69}:
\]

%dsp-level:0
\[
\qquad \bigl\{\tfrac32=\tfrac{E_{\mathrm{int}}}4\ [\text{a count row}],\quad4\ [\text{the}\ I_F\ \text{trace normalization}],\quad2c_0,\ c_0^2+C^2\ [\text{annihilated in}\ I_F],\quad\rho_{\mathrm B}\ [\text{a Cal-F restatement}]\bigr\}:
\]

%dsp-level:0
\[
\textbf{each is a trace ratio, a Casimir normalization, a constrained fiber-count derivative, a channel-cap lock, or an exact operator identity}\ —\ \text{the organic test is satisfied coefficient by coefficient}
\]

%dsp-level:0
\[
\qquad [\text{machine: the census}];
\]

%dsp-level:0
\[
\alpha_B:\quad\neg\text{added},\quad\neg\text{derived},\quad\text{barred on}\ \mathrm{BM.61}(8)\text{'s four grounds, which stand}
\]

%dsp-level:0
\[
\qquad [\text{machine: the token scan of this tier}\ —\ \text{every occurrence is a barred-object row}];
\]

%dsp-level:0
\[
\neg\text{a P4-class temperature},
\]

%dsp-level:0
\[
\quad\neg\text{a free-energy coefficient},
\]

%dsp-level:0
\[
\quad\neg\text{a fitting parameter is introduced anywhere in this tier}
\]

%dsp-level:0
\[
\qquad [\text{machine: the symbol census against the P4 core and the era registries}];
\]

%dsp-level:0
\[
(7)\quad\textbf{status:}
\]

%dsp-level:0
\[
\qquad \textbf{derived:}\ (1)\text{--}(6);
\]

%dsp-level:0
\[
\qquad \textbf{not achieved: the full identity}\ —\ \text{the obstruction is exact and named}\ (4),
\]

%dsp-level:0
\[
\ \text{and the residue is minimal and priced}\ (5);
\]

%dsp-level:0
\[
\neg\text{a P4 addition},\quad\neg\text{a new dimensionful scale}\ [\mathrm{P1}],\quad\neg\text{a measured value}\ [\mathrm{P5}],\quad\neg\text{a declaration};
\]

%dsp-level:0
\[
\qquad \mathrm{Def\text{-}H.37}(2)\text{'s interface is cited as the locus of the declared transfer,}\ \neg\text{imported as a relating equation}.
\]

\textbf{Req.} P1, P4, P5 (Parts 1–2); A §E (the P4 core); Def-H.37(2), FT.54(4), FT.55(1), FT.73(2), FT.73(3), FT.73(5), FT.125(4) (Part 3); Census FT.35(2), FT.63(4), FT.63(6) (row-wise); Def-BM.8(5), Def-BM.13(4), Def-BM.14(3), Def-BM.14(4), Def-BM.14(5), BM.23(4), BM.28(7), BM.35(5), BM.36(3), BM.37(2), BM.57(4), BM.59(9), BM.60(6), BM.61(8), BM.66(3), BM.66(4), BM.67(2), BM.67(6), BM.67(8), BM.67(9), BM.67(10), BM.68(2), BM.68(5), BM.68(8) (this part); Cal-F; exact radical separations, Laurent-monomial assembly, symbol censuses [math; machine: \texttt{checks/\allowbreak{}verify\_\allowbreak{}H\_\allowbreak{}organic\_\allowbreak{}bridge\_\allowbreak{}post\_\allowbreak{}bm14.py} §A, §E, §G].

\textbf{Pf.}

%dsp-level:0
\[
(1):\ \text{the three rows re-cited from their own tags};
\]

%dsp-level:0
\[
\qquad (2):\ \mathrm{BM.68}(8)\text{'s identity re-cited};
\]

%dsp-level:0
\[
\qquad \text{the six values and four fields re-cited from}\ \mathrm{Def\text{-}BM.14}(3)\ \text{and}\ \text{Census FT.63}(6)
\]

%dsp-level:0
\[
(3):\ E_{\mathrm{int}}\ \text{class-shared}\ [\mathrm{BM.57}(4)\text{'s difference table re-cited}];
\]

%dsp-level:0
\[
\qquad \Delta_g\ \text{re-formed in the three exact fields with the certificates}\ 4>2,
\]

%dsp-level:0
\[
\ 2>1,
\]

%dsp-level:0
\[
\ \cos\tfrac\pi{10}>\tfrac{\sqrt3}2
\]

%dsp-level:0
\[
\iff 3<5\ [\text{machine}];
\]

%dsp-level:0
\[
\qquad \text{the addition bar:}\ \mathrm{BM.37}(2)\text{'s disjointness and}\ \mathrm{BM.35}(5)\text{'s no-conversion row re-cited}
\]

%dsp-level:0
\[
(4)(\mathrm i):\ \mathrm{FT.73}(2),
\]

%dsp-level:0
\[
(3)\ \text{re-cited: the degree-}16\ \text{compositum, the irreducible quartics, the quadratic-field exclusion}\ 4-2\sqrt2+2\sqrt3-\sqrt6\notin K_2;
\]

%dsp-level:0
\[
\qquad (\mathrm{ii})\text{--}(\mathrm{iv}):\ \mathrm{BM.67}(6),(8),\ \mathrm{BM.36}(3),\ \mathrm{BM.61}(8)\ \text{re-cited}
\]

%dsp-level:0
\[
(5):\ \text{the inventory row assembled from}\ \mathrm{BM.65}\text{--}\mathrm{BM.68}\text{'s derived lists}\ [\text{machine}];
\]

%dsp-level:0
\[
\qquad (6):\ \text{the coefficient census}\ [\text{machine}];
\]

%dsp-level:0
\[
\qquad \text{the token scan of this tier for the barred object}\ \alpha_B\ \text{and for P4-registry additions}\ [\text{machine}];
\]

%dsp-level:0
\[
\qquad (7):\ \text{assembly}. \ \blacksquare
\]

\textbf{Phys.} The three inputs meet, and the bridge exists in exactly the form the corpus can derive — and in no other. The operator invariant is the exact integer four Q-squared; the holonomy gives a binary sector and an obstructed angle; the joint entropy gives the count plus the charge census at stationarity, with a geometric factor of one. The well-posed bridge object, the pole-to-hidden mass ratio, factorizes into a per-state exact algebraic number and a single dimensionless packaging, and that packaging is the whole of the gap: everything else in the identity is derived, and the non-derivability of the rest is a theorem with four independent grounds — a degree-sixteen field against integer reads, a two-element image against six angles, a calibration weight against the winding route, and the standing bar against turning counts into strengths. The class dependence of the bridge lives entirely in the informational currency; its pole-currency image needs a flavour composition no clause attaches, and where a composition is hypothesized the pole splitting points the same way as the transport row but cannot be added to it without the very equation that is missing. So Gap 7 reduces to its smallest honest form: one scale ratio, one loop realization, one composition — named, priced, and hidden inside no coefficient.

\end{theorem}

\pseudohead{theorem}{\texorpdfstring{BM.70}{BM.70}}{\texorpdfstring{The beta-decay energy audit: the two-bit floor theorem, the margin ladder, and the formed-closed dichotomy}{The beta-decay energy audit: the two-bit floor theorem, the margin ladder, and the formed-closed dichotomy}}\label{tag:BM-70}
\textbf{Thm.}

%dsp-level:0
\[
(1)\quad\textbf{the floor and the Q-value:}\qquad F_{\ell\nu}\ :=\ m_\ell\,\mathfrak h\ +\ m_\nu\,\mathfrak h\qquad[\text{an abbreviation of derived reads}];
\]

%dsp-level:0
\[
\boxed{\ Q_\beta\,\mathfrak h\ =\ \ln\tfrac43\ -\ 6\,c\,\varepsilon_{\mathrm{EM}}\ln2\ -\ F_{\ell\nu}\ }\qquad[\mathrm{BM.62}(4);\ \mathrm{BM.63}(1);\ \text{machine: the monomial identity}];
\]

%dsp-level:0
\[
(2)\quad\textbf{the two-bit daughter floor (theorem on the informational species):}
\]

%dsp-level:0
\[
\qquad \textbf{the chain:}
\]

%dsp-level:0
\[
\qquad \text{the emitted charged member is the}\ \ell\text{-channel line}\ —\ \text{the forced, class-unique carrier}
\]

%dsp-level:0
\[
\qquad [\mathrm{BM.23}(5);\ \mathrm{BM.41}(1);\ \mathrm{Def\text{-}BM.6}(1)]:
\]

%dsp-level:0
\[
\text{every channel line is a materialized bound chain}
\]

%dsp-level:0
\[
\qquad [\mathrm{FT.109};\ \mathrm{FT.112};\ \mathrm{FT.113};\ \mathrm{Def\text{-}H.4}];
\]

%dsp-level:0
\[
\qquad \text{every materialized closed chain carries}\ \ge2\ \text{directed edges}\ —\ \text{the digon is the minimal floor}
\]

%dsp-level:0
\[
\qquad [\mathrm{Def\text{-}H.4}:\ \delta h(c)\ge2;\ \text{Census FT.28};\ \mathrm{FT.32}:\ L\ge2;\ \mathrm{BM.13}(3):\ E_{\mathrm{int}}\ge2]:
\]

%dsp-level:0
\[
\text{the informational read}\ \ m\,\mathfrak h=E_{\mathrm{int}}\ln2
\]

%dsp-level:0
\[
\qquad [\text{Census FT.44}(4);\ \mathrm{Def\text{-}H.9}(2);\ \mathrm{BM.43}(2):\ \text{the taxonomy law}]
\]

%dsp-level:0
\[
\qquad \implies
\]

%dsp-level:0
\[
\qquad m_\ell\,\mathfrak h\ \ge\ 2\ln2\ >\ 0
\]

%dsp-level:0
\[
\qquad [\text{machine: the exact enclosure}\ \tfrac12<\ln2<1];
\]

%dsp-level:0
\[
m_\nu\,\mathfrak h\ \ge\ 0
\]

%dsp-level:0
\[
\qquad [\mathrm{EW.33}:\ \text{the Gram is PSD};\ \mathrm{BM.41}(4)(\mathrm a)];
\]

%dsp-level:0
\[
\qquad \text{and on the informational read}\ \ m_\nu\,\mathfrak h=E_{\mathrm{int}}^{(\nu)}\ln2\ \ge\ 2\ln2\ \ \text{likewise}
\]

%dsp-level:0
\[
\qquad [\mathrm{BM.43}(2);\ \mathrm{BM.13}(3)]:
\]

%dsp-level:0
\[
\boxed{\ F_{\ell\nu}\ \ge\ 2\ln2\ }
\]

%dsp-level:0
\[
\qquad \textbf{with the floor ladder:}
\]

%dsp-level:0
\[
\qquad 2\ln2\ \ \bigl(\text{the PSD row alone}\bigr)
\]

%dsp-level:0
\[
\quad <
\]

%dsp-level:0
\[
\quad 4\ln2\ \ \bigl(\text{the informational neutrino}\bigr)
\]

%dsp-level:0
\[
\quad <
\]

%dsp-level:0
\[
\quad 5\ln2\ \ \bigl(\text{with the daughters' own transport refinement}\bigr)
\]

%dsp-level:0
\[
\qquad [\text{machine:}\ 2<4<5]:
\]

%dsp-level:0
\[
\textbf{the refinement row:}
\]

%dsp-level:0
\[
\qquad k=1:\quad W(q)=m_q=(1,\,2,\,1)\ \ \text{at}\ q=(-1,\,0,\,+1),
\]

%dsp-level:0
\[
\qquad \Sigma_{\mathrm{cl}}(q)=(\,0,\ \ln2,\ 0\,):
\]

%dsp-level:0
\[
\qquad \text{the neutral line carries one extra bit, the charged lines none}
\]

%dsp-level:0
\[
\qquad [\mathrm{Def\text{-}BM.13}(3);\ \mathrm{BM.63}(6);\ \text{machine: the convolution of}\ (2+z+z^{-1})^1]:
\]

%dsp-level:0
\[
\text{the daughters are}\ q_\ell
\]

%dsp-level:0
\[
= -1\ \ [\text{the charged member}]\ \ \text{and}\ q_\nu
\]

%dsp-level:0
\[
= 0\ \ [\text{the neutral member};\ \mathrm{BM.41}(1);\ \mathrm{FT.142}(1):\ \text{EM-unactivated}]
\]

%dsp-level:0
\[
\qquad \implies
\]

%dsp-level:0
\[
\qquad F_{\ell\nu}\ \ge\ 2\ln2+0+2\ln2+\ln2=5\ln2\ \ \text{on the fully refined informational read}
\]

%dsp-level:0
\[
\qquad [\text{machine}];
\]

%dsp-level:0
\[
(3)\quad\textbf{the margin ladder (exact, integer-certified):}
\]

%dsp-level:0
\[
\qquad Q_\beta\,\mathfrak h\ \le\ \ln\tfrac43-F_{\ell\nu}\ \ \text{at every}\ \varepsilon_{\mathrm{EM}}>0
\]

%dsp-level:0
\[
\qquad [\mathrm{BM.62}(6):\ \text{strict monotonicity}]:
\]

%tbl-level:0
{\setlength{\tabcolsep}{4pt}
\begin{longtable}{>{\raggedright\arraybackslash}p{\dimexpr0.3642\textwidth-2\tabcolsep\relax}|>{\raggedright\arraybackslash}p{\dimexpr0.1199\textwidth-2\tabcolsep\relax}|>{\raggedright\arraybackslash}p{\dimexpr0.1214\textwidth-2\tabcolsep\relax}|>{\raggedright\arraybackslash}p{\dimexpr0.3945\textwidth-2\tabcolsep\relax}}
\toprule
floor used & $F_{\ell\nu}\ \ge$ & $Q_\beta\,\allowbreak \mathfrak h\ \le$ & the certificate \\
\addlinespace[2pt]
\midrule
\endfirsthead
\toprule
floor used & $F_{\ell\nu}\ \ge$ & $Q_\beta\,\allowbreak \mathfrak h\ \le$ & the certificate \\
\addlinespace[2pt]
\midrule
\endhead
\bottomrule
\endlastfoot
the PSD row alone & $2\ln2$ & $-\ln3$ & $\ln\tfrac43<\allowbreak 2\ln2\iff\allowbreak \tfrac43<\allowbreak 4\iff\allowbreak 4<\allowbreak 12$ \\
the informational neutrino & $4\ln2$ & $-\ln12$ & $\ln\tfrac43<\allowbreak 4\ln2\iff\allowbreak \tfrac43<\allowbreak 16\iff\allowbreak 4<\allowbreak 48$ \\
the fully refined daughters & $5\ln2$ & $-\ln24$ & $\ln\tfrac43<\allowbreak 5\ln2\iff\allowbreak \tfrac43<\allowbreak 32\iff\allowbreak 4<\allowbreak 96$ \\
\end{longtable}
}

%dsp-level:0
\[
[\text{machine: the prime-log identities}\ \ \ln\tfrac43-2\ln2=-\ln3,\quad\ln\tfrac43-4\ln2=-\ln12,\quad\ln\tfrac43-5\ln2=-\ln24;\ \ \text{and every comparison an integer certificate}];
\]

%dsp-level:0
\[
\boxed{\ Q_\beta\ <\ 0\ \ \text{strictly, with margin}\ \ge\ \tfrac{\ln3}{\mathfrak h},\ \ \text{deepened by}\ \tfrac{6\,c\,\varepsilon_{\mathrm{EM}}\ln2}{\mathfrak h}\ \ \text{at every positive coupling}\ }
\]

%dsp-level:0
\[
\qquad [\mathrm{BM.63}(2),(3),(4)\ \text{audited and standing};\ \text{machine}];
\]

%dsp-level:0
\[
(4)\quad\textbf{the formed-closed dichotomy (the audit verdict):}
\]

%dsp-level:0
\[
\qquad Q_\beta\ \text{is a formed expression only where all four members are read in one currency}:
\]

%dsp-level:0
\[
(\mathrm a)\quad\textbf{the informational currency on every member}\ \bigl[\text{the parents:}\ \mathrm{Def\text{-}BM.13}(5);\ \ \text{the daughters:}\ \mathrm{BM.43}(2)\text{'s taxonomy law}\bigr]:
\]

%dsp-level:0
\[
\qquad (2)\text{'s floor theorem applies}\implies Q_\beta\,\mathfrak h\le-\ln3<0
\]

%dsp-level:0
\[
\qquad \textbf{formed}\implies\textbf{closed};
\]

%dsp-level:0
\[
(\mathrm b)\quad\textbf{the transport/pole species on the leptons}\ \bigl[\mathrm{BM.39}(1):\ m_\ell^2=\hat\mu^2\ell_c^2\,G(c(i))>0,\ \ \text{the scale}\ \hat\mu\ \text{undeclared}\ [\mathrm{Cal\text{-}F}];\ \mathrm{Def\text{-}BM.11}\bigr]:
\]

%dsp-level:0
\[
\qquad \text{the subtraction against informational-currency parents is a cross-currency expression};
\]

%dsp-level:0
\[
\text{its relating equation is}\ \neg\text{derived}
\]

%dsp-level:0
\[
\qquad [\mathrm{BM.35}(5);\ \mathrm{BM.43}(2):\ \text{the conversion}\ E_{\mathrm{int}}^{(i)}\leftrightarrow\text{transport data is item-7 residue};\ \mathrm{BM.60}(6);\ \text{pointers}]
\]

%dsp-level:0
\[
\qquad \implies
\]

%dsp-level:0
\[
\qquad Q_\beta\ \ \neg\textbf{formed}\implies\ \neg\textbf{a verdict};
\]

%dsp-level:0
\[
\boxed{\ \text{on every read on which}\ Q_\beta\ \text{is a derived expression,}\ Q_\beta<0\ \ \text{strictly};\qquad\neg\exists\ \text{a derived read on which}\ Q_\beta\ge0\ }
\]

%dsp-level:0
\[
\qquad [\text{machine: the row census over the complete route inventory}\ \mathrm{BM.30}(1):\ \ \text{route}\ (\mathrm i)\ Q=-(m_\ell{+}m_\nu)<0\ [\mathrm{BM.41}(4)(\mathrm b)];\ \ \text{route}\ (\mathrm{ii})\ \text{priced}\ \neg\text{adopted, its sign undetermined without scale data}\ [\mathrm{BM.41}(4)(\mathrm d)];\ \ \text{route}\ (\mathrm{iii})\ \text{amplitude layer}\ —\ \text{amplitude}\ne\text{mass}\ [\mathrm{FT.79}(5);\ \mathrm{BM.37}(3)];\ \ \text{the curvature-flux read:}\ Q<0\ [\mathrm{BM.41}(4)(\mathrm c)]];
\]

%dsp-level:0
\[
(5)\quad\textbf{the missing lepton read, named exactly:}
\]

%dsp-level:0
\[
\qquad E_{\mathrm{int}}^{(\ell)},\ E_{\mathrm{int}}^{(\nu)}\ \ \neg\text{computed at any floor}
\]

%dsp-level:0
\[
\qquad [\mathrm{BM.43}(2);\ \mathrm{BM.63}(8)];
\]

%dsp-level:0
\[
\qquad \text{the line index within the}\ \nu\text{-triplet}\ \neg\text{derived}
\]

%dsp-level:0
\[
\qquad [\mathrm{FT.145}(3):\ \neg\text{a mixing matrix}]:
\]

%dsp-level:0
\[
\textbf{these reads sharpen the margin and cannot breach the floor:}
\]

%dsp-level:0
\[
\qquad E_{\mathrm{int}}\ge2\ \ \text{is a structural lower bound of the species}
\]

%dsp-level:0
\[
\qquad [\mathrm{BM.13}(3);\ \text{Census FT.28}];
\]

%dsp-level:0
\[
\qquad \text{the margin is monotone in}\ E_{\mathrm{int}}^{(\ell)},\ E_{\mathrm{int}}^{(\nu)}
\]

%dsp-level:0
\[
\qquad [\text{machine: each added edge adds}\ \ln2]:
\]

%dsp-level:0
\[
\boxed{\ \text{the closure is}\ \neg\textbf{an artifact of an unfinished lepton mass read: the unfinished read moves the margin, never the sign}\ };
\]

%dsp-level:0
\[
(6)\quad\textbf{no external Q-value:}
\]

%dsp-level:0
\[
\qquad \neg\exists\ \text{an observed endpoint, mass difference or}\ Q\text{-value in any clause of this tier}
\]

%dsp-level:0
\[
\qquad [\mathrm{P5};\ \text{machine: the token scan}\ —\ \text{zero external numerals, zero unit tokens}];
\]

%dsp-level:0
\[
\textbf{the consistency-check row:}
\]

%dsp-level:0
\[
\qquad \text{a comparison against an observed}\ Q_\beta\ \text{is admissible only as a quarantined audit row}
\]

%dsp-level:0
\[
\qquad [\text{Census FT.114},\ \text{FT.120}\ \text{pattern};\ \mathrm{P1},\ \mathrm{P5}]\ \ \text{and only after the corpus derives its own scale}\ —\ \text{one calibration datum fixing}\ e\ \ [\mathrm{EM.25}(7)]\ \ \text{and a}\ \mathfrak h\ \text{read}\ [\mathrm{P1}]:
\]

%dsp-level:0
\[
\neg\exists\ \text{such a derivation in Parts 1--4};\qquad\textbf{this tier performs no comparison};
\]

%dsp-level:0
\[
(7)\quad\textbf{status:}
\]

%dsp-level:0
\[
\qquad \textbf{derived:}\ (1)\text{--}(6);
\]

%dsp-level:0
\[
\qquad \mathrm{BM.63}(2)\text{'s closure is confirmed as a theorem on the formed reads, with its species conditionality on its face}\ (4)\ \ [\text{A §0.2a}];
\]

%dsp-level:0
\[
\neg\text{a measured lifetime, mass difference or endpoint enters}\ [\mathrm{P5}];
\]

%dsp-level:0
\[
\qquad \neg\text{a P4 addition};
\]

%dsp-level:0
\[
\qquad \neg\text{a new dimensionful scale}\ [\mathrm{P1}];
\]

%dsp-level:0
\[
\qquad \text{the open rows:}\ E_{\mathrm{int}}^{(\ell)},\ E_{\mathrm{int}}^{(\nu)},\ \text{the}\ \nu\text{-line index},\ \text{the numerical value of}\ e
\]

%dsp-level:0
\[
\qquad [\text{pointers}].
\]

\textbf{Req.} EW.33 (Part 2); EM.25(7) (Part 2); P1, P5 (Parts 1–2); Def-H.4, Def-H.9(2), FT.32, FT.79(5), FT.109, FT.112, FT.113, FT.142(1), FT.145(3) (Part 3); Census FT.28, FT.44(4), FT.114, FT.120 (row-wise); Def-BM.6(1), Def-BM.11, Def-BM.13(3), Def-BM.13(5), BM.13(3), BM.23(5), BM.30(1), BM.35(5), BM.37(3), BM.39(1), BM.41(1), BM.41(4), BM.43(2), BM.60(6), BM.62(4), BM.62(6), BM.63(1), BM.63(2), BM.63(3), BM.63(4), BM.63(6), BM.63(8) (this part); Cal-F; exact prime-log-ring identities with integer certificates, exact rational $\ln2$ enclosures, monotonicity censuses, token scans [math; machine: \texttt{checks/\allowbreak{}verify\_\allowbreak{}H\_\allowbreak{}organic\_\allowbreak{}bridge\_\allowbreak{}post\_\allowbreak{}bm14.py} §F].

\textbf{Pf.}

%dsp-level:0
\[
(1):\ \mathrm{BM.63}(1)\text{'s row re-cited and re-formed in the}\ \ln\text{-ring}\ [\text{machine}];
\]

%dsp-level:0
\[
\qquad (2):\ \delta h(c)\ge2\ [\mathrm{Def\text{-}H.4}\ \text{re-cited}];
\]

%dsp-level:0
\[
\qquad \text{the digon card}\ E_{\mathrm{int}}=2\ [\text{Census FT.28}\ \text{re-cited}];
\]

%dsp-level:0
\[
\qquad E_{\mathrm{int}}\ge2\ \text{on every computed cell}\ \{2,6,12,24,60\}\ [\text{machine: the minimum}];
\]

%dsp-level:0
\[
\qquad m_\ell\mathfrak h=E_{\mathrm{int}}^{(\ell)}\ln2\ge2\ln2\ [\text{machine}];
\]

%dsp-level:0
\[
\qquad m_\nu\mathfrak h\ge0\ [\mathrm{EW.33}\ \text{PSD re-cited}]
\]

%dsp-level:0
\[
\text{the refinement:}\ (2+z+z^{-1})^1\ \text{has coefficients}\ (1,2,1)\ \text{at}\ Q
\]

%dsp-level:0
\[
= (-1,0,+1)\ [\text{machine: the convolution}];
\]

%dsp-level:0
\[
\qquad \ln m_q=(0,\ln2,0)\ [\text{machine}];
\]

%dsp-level:0
\[
\qquad 2\ln2+2\ln2+\ln2=5\ln2\ [\text{machine: the prime-log ring}];
\]

%dsp-level:0
\[
\qquad 2<4<5\ [\text{machine}]
\]

%dsp-level:0
\[
(3):\ \ln\tfrac43-2\ln2=(2\ln2-\ln3)-2\ln2=-\ln3;
\]

%dsp-level:0
\[
\qquad \ln\tfrac43-4\ln2=-2\ln2-\ln3=-\ln(4\cdot3)=-\ln12;
\]

%dsp-level:0
\[
\qquad \ln\tfrac43-5\ln2=-3\ln2-\ln3=-\ln(8\cdot3)=-\ln24\ [\text{machine: the prime-log ring}];
\]

%dsp-level:0
\[
\qquad \text{the three certificates:}\ 4<12,\ 4<48,\ 4<96\ [\text{machine: integers}];
\]

%dsp-level:0
\[
\qquad \text{the EM deepening:}\ 6c\varepsilon_{\mathrm{EM}}\ln2>0\ \text{at}\ c,\varepsilon_{\mathrm{EM}},\ln2>0\ [\text{machine: the coefficient census}]
\]

%dsp-level:0
\[
(4):\ \text{the route inventory re-cited from}\ \mathrm{BM.30}(1),
\]

%dsp-level:0
\[
\ \mathrm{BM.41}(4)(\mathrm a)\text{--}(\mathrm d)\ \text{and classified by currency}\ [\text{machine: the row census}];
\]

%dsp-level:0
\[
\qquad \text{the cross-currency bar:}\ \mathrm{BM.35}(5)\text{'s no-conversion row and}\ \mathrm{BM.37}(2)\text{'s disjointness re-cited}
\]

%dsp-level:0
\[
(5):\ \text{monotonicity:}\ \partial\bigl(-Q_\beta\mathfrak h\bigr)/\partial E_{\mathrm{int}}^{(i)}=\ln2>0\ [\text{machine}];
\]

%dsp-level:0
\[
\qquad (6):\ \text{the token scan}\ [\text{machine: zero external numerals}];
\]

%dsp-level:0
\[
\qquad (7):\ \text{assembly}. \ \blacksquare
\]

\textbf{Phys.} The audit asked whether the two-bit daughter floor is a theorem or an artifact of an unfinished lepton read, and the answer separates the two cleanly. It is a theorem wherever the question is formed: the emitted charged member is a line of the lepton channel, every channel line is a materialized bound chain, every materialized closed chain carries at least two directed edges because the digon is the floor, and the informational read turns that count into at least two bits; the neutral member costs nothing less than zero by positive semidefiniteness, and nothing less than two more bits plus one refinement bit if it too is read informationally. The margin ladder follows in exact integers — against three, against twelve, against twenty-four — and the field-energy term deepens every rung at every positive coupling. The only escape is currency: read the leptons through the transport species with its undeclared scale and the Q-value is not a formed expression at all, because the equation relating the two currencies is the corpus's open item. Hence the dichotomy, which is the audit's verdict: formed implies closed, unformed implies no verdict, and no derived read opens the vertex. What the unfinished lepton read moves is the margin's sharpness, never the sign — the floor is structural, sitting at the digon, below every rung of every ladder. And the observed endpoint enters nowhere: it may be compared only in a quarantined audit row, after a scale derivation that does not exist, and this tier performs no comparison.

\pseudohead{theorem}{\texorpdfstring{BM.71}{BM.71}}{\texorpdfstring{The opening-term census: the five routes, the neutron row under the bridge, and the six verdict rows}{The opening-term census: the five routes, the neutron row under the bridge, and the six verdict rows}}\label{tag:BM-71}
\textbf{Thm.}

%dsp-level:0
\[
(1)\quad\textbf{the route census (each candidate: derived, priced, or barred):}
\]

%tbl-level:3
{\scriptsize\setlength{\tabcolsep}{4pt}
\begin{longtable}{>{\raggedright\arraybackslash}p{\dimexpr0.0507\textwidth-2\tabcolsep\relax}|>{\raggedright\arraybackslash}p{\dimexpr0.1244\textwidth-2\tabcolsep\relax}|>{\raggedright\arraybackslash}p{\dimexpr0.1032\textwidth-2\tabcolsep\relax}|>{\raggedright\arraybackslash}p{\dimexpr0.7218\textwidth-2\tabcolsep\relax}}
\toprule
\# & the candidate opening term & verdict & the derived ground \\
\addlinespace[2pt]
\midrule
\endfirsthead
\toprule
\# & the candidate opening term & verdict & the derived ground \\
\addlinespace[2pt]
\midrule
\endhead
\bottomrule
\endlastfoot
1 & time-energy uncertainty / the non-uniform clock & \textbf{barred} & the Mandelstam–Tamm relation bounds the splitting's \textit{resolvability}, $\ne\allowbreak g$its magnitude $[\mathrm{BM.53}(4)\text{–}(6)]$; the closure is the non-existence of real on-shell solutions — $2m_n(E_p-m_p)=\allowbreak -m_\ell^2$, $\lambda=\allowbreak m_\ell^2(m_\ell^2-4m_n^2)$: polynomial identities at every tick $[\mathrm{BM.41}(2);\ \text{machine}]$; $\ne\allowbreak g\exists$ a time operator at finite dimension, so an energy loan over a $\Delta t$ has no operator content $[\mathrm{BM.10}(4)]$; the ledger is forward-only $[\mathrm{A6}:\ \ne\allowbreak g\exists\tau^{-1}]$, strictly contracting $[\mathrm{L2};\ \mathrm{BM.51}(2),\allowbreak (3)]$, with per-tick weights $\Gamma_{\mathrm{tr}}\in\allowbreak [0,\allowbreak \tfrac13]$ $[\mathrm{BM.44}(2)]$; the flicker ensemble is amplitude-layer content, and amplitude $\ne$ mass $[\mathrm{FT.79}(5);\ \mathrm{BM.37}(3);\ \mathrm{FT.67}(3\mathrm a);\ \mathrm{FT.68}(4)]$ \\
2 & field reconfiguration energy & \textbf{$\ne\allowbreak g$constructed; sign-barred at the constructed order} & radiation, back-reaction and the transition's field reconfiguration are the named open scope $[\mathrm{Def\text{-}BM.12}(7);\ \mathrm{BM.50}(3)]$; at the constructed monopole order every self-energy is the positive integral of a square $[\mathrm{BM.46}(1);\ \mathrm{Def\text{-}BM.12}(5)]$, and the transition $Q:0\to(+\allowbreak 1)\oplus(-1)$ raises the total field energy — the final state carries two charged members' self-energies where the initial carries none $[\mathrm{BM.49}(2);\ \text{machine: the sign census at positive symbols}]$; the attractive cross-well $V=\allowbreak -\varepsilon_{\mathrm{EM}}/d$ belongs to the $\ne\allowbreak g$constructed pair system $[\mathrm{BM.48}(1),\allowbreak (3),\allowbreak (4);\ \mathrm{BM.54}(4):\ \text{the category bar}]$ \\
3 & binding / nearby-space correction & \textbf{barred} & the category bar: a binding energy is the energy of the bound system, $\ne\allowbreak g$a correction to a free member's mass $[\mathrm{BM.48}(4);\ \mathrm{BM.54}(4)]$; the lepton–baryon bound channel is $\ne\allowbreak g$constructed $[\mathrm{BM.48}(3);\ \mathrm{BM.6}(5):\ \text{the dynamics of}\ Q]$; the nearby space holds only a strict fraction $f(k)=\allowbreak 1-\tfrac1k<\allowbreak 1$ $[\mathrm{BM.47}(6);\ \text{machine: exact rationals}]$; the count binding is zero on both rows, $6-(2{+}2{+}2)=\allowbreak 0$ $[\mathrm{BM.54}(1);\ \text{machine}]$; the ensemble binding is class-degenerate $[\mathrm{BM.54}(2);\ \mathrm{FT.107}(3);\ \mathrm{BM.30}(3)]$ \\
4 & $D^-$ redistribution & \textbf{barred} & the dark sector's contribution to the class splitting is exactly zero: $D^+\allowbreak =\allowbreak \Lambda\,\allowbreak \mathbb I$ uniform, both locks class-shared $[\mathrm{BM.52}(3)\ D^\pm\ \text{row};\ \mathrm{CT4.5};\ \mathrm{BM.28}(1)]$; the junction's $D^-$ content is identical on both rows — the same vertex types $n(q_1)=\allowbreak n(p_2)=\allowbreak 4$, the same lock ledger $[\text{Census FT.35}(4);\ \mathrm{BM.57}(4):\ \text{the difference table is identically zero};\ \text{machine}]$, so $\ne\allowbreak g\exists$ a class-dependent object to redistribute; locked content is retained and $\ne\allowbreak g$accessible $[\mathrm{D4.2};\ \mathrm{D1.11};\ \mathrm{CT4.3};\ \mathrm{FT.3}]$ — no channel reads it into a vertex budget; and $\ne\allowbreak g\exists$ a response read — the dielectric non-existence inventory theorem $[\mathrm{FT.73}(1);\ \text{any dynamics supplying one contradicts}\ \mathrm{CT4.2}/\mathrm{D1.11}]$ \\
5 & read specification & \textbf{the sole surviving door — not walked} & opening the vertex is a read-specification act requiring a §I addendum $[\mathrm{BM.49}(2);\ \mathrm{BM.41}(4)(\mathrm e);\ \mathrm{BM.30}(5);\ \mathrm{BM.45}]$; the route-(ii) measure price $T(0)=\allowbreak \tfrac43T(\pm1)$, its sign undetermined without scale data, stays priced $\ne\allowbreak g$adopted $[\mathrm{BM.41}(4)(\mathrm d);\ \mathrm{BM.28}(4);\ \mathrm{BM.37}(3)]$; \textbf{this tier performs no specification} \\
\end{longtable}
}

%dsp-level:0
\[
(2)\quad\textbf{completeness and the verdict:}
\]

%dsp-level:0
\[
\qquad \text{the five routes exhaust the audit's candidate list};
\]

%dsp-level:0
\[
\qquad \text{each is barred by a cited theorem or}\ \neg\text{constructed with its scope named}
\]

%dsp-level:0
\[
\qquad [\text{machine: the census-row structure}];
\]

%dsp-level:0
\[
\textbf{the MT bound's legitimate content:}
\]

%dsp-level:0
\[
\qquad \Delta t\ \ge\ \frac1{2\,\Delta m_{\mathrm{EM}}}\ =\ \frac{\pi\,\mathfrak h}{3\,c\,e^2\ln2}
\]

%dsp-level:0
\[
\qquad [\mathrm{BM.53}(5);\ \text{machine: the exact Laurent reciprocal}]
\]

%dsp-level:0
\[
\qquad \text{is a bound on clocks resolving the splitting,}\ \neg\text{a budget paying for a decay};
\]

%dsp-level:0
\[
\boxed{\ \textbf{the beta-decay no-go: on every formed derived read the vertex is kinematically closed, and its opening is}\ \neg\textbf{derivable from the corpus}\ —\ \textbf{it requires a read specification}\ }
\]

%dsp-level:0
\[
\qquad [\mathrm{BM.70}(4);\ (1);\ \text{machine}];
\]

%dsp-level:0
\[
(3)\quad\textbf{the neutron row under the bridge:}\qquad\boxed{\ m_n-m_p\ =\ \Delta m_{\mathrm{transport}}\ +\ \Delta m_{\mathrm{EM}}\ +\ \Delta m_{\mathrm{bridge}}\ }
\]

%dsp-level:0
\[
\Delta m_{\mathrm{transport}}\ =\ +\,\frac{\ln(4/3)}{\mathfrak h}
\]

%dsp-level:0
\[
\qquad [\mathrm{BM.59}(4);\ \mathrm{BM.65}(2):\ \text{derived}];
\]

%dsp-level:0
\[
\qquad \Delta m_{\mathrm{EM}}\ =\ -\,\frac{6\,c\,\varepsilon_{\mathrm{EM}}\ln2}{\mathfrak h}
\]

%dsp-level:0
\[
\qquad [\mathrm{BM.46}(4);\ \mathrm{BM.62}(3);\ \text{the operator form}\ \mathrm{BM.66}(4):\ \text{derived};\ c\ \text{priced};\ \varepsilon_{\mathrm{EM}}\ \text{symbolic}];
\]

%dsp-level:0
\[
\boxed{\ \Delta m_{\mathrm{bridge}}\ =\ 0\ \ \text{in the informational currency}\ }
\]

%dsp-level:0
\[
\qquad \textbf{(theorem):}\ \ \mathrm{BM.62}(1)\text{'s decomposition is complete in that currency}
\]

%dsp-level:0
\[
\qquad [\mathrm{BM.58}(6)\text{'s uniqueness row};\ \mathrm{BM.62}(11)\text{'s alphabets};\ \mathrm{BM.66}(6)\text{'s guard};\ \mathrm{BM.68}(5)\text{'s stationary slice}],
\]

%dsp-level:0
\[
\text{and the bridge's pole-currency content is unformed without the relating equation}
\]

%dsp-level:0
\[
\qquad [\mathrm{BM.69}(3),(4)];
\]

%dsp-level:0
\[
\textbf{the conditional pole row (the hypothesis on its face):}
\]

%dsp-level:0
\[
\qquad \text{under}\ \mathrm{BM.59}(9)\text{'s substitution hypothesis:}\quad\Delta\bigl(m^2_{\mathrm{pole}}\bigr)=+\Delta_g>0
\]

%dsp-level:0
\[
\qquad [\mathrm{BM.69}(3);\ \text{machine: the three certificates}]
\]

%dsp-level:0
\[
\qquad \text{— the same direction as the transport row,}\ \neg\text{a formed summand of any}\ Q\text{-value};
\]

%dsp-level:0
\[
(4)\quad\textbf{the six verdict rows:}
\]

%dsp-level:0
\[
(\mathrm i)\quad\textbf{the sign of}\ m_n-m_p:
\]

%dsp-level:0
\[
\qquad >0\ \ \text{on the sub-critical branch}\ \varepsilon_{\mathrm{EM}}<\varepsilon_{\mathrm{crit}}(c);
\]

%dsp-level:0
\[
\qquad =0\ \ \text{at}\ \varepsilon_{\mathrm{crit}}(c);
\]

%dsp-level:0
\[
\qquad <0\ \ \text{above}
\]

%dsp-level:0
\[
\qquad [\mathrm{BM.62}(6),(9);\ \mathrm{BM.65}(3)\text{--}(5)]:
\]

%dsp-level:0
\[
\qquad \text{which branch the corpus sits on is open}\ —\ \text{the numerical value of}\ e
\]

%dsp-level:0
\[
\qquad [\mathrm{EM.25}(7);\ \mathrm{BM.62}(12)];
\]

%dsp-level:0
\[
(\mathrm{ii})\quad\textbf{the sign of}\ Q_\beta:\qquad\text{strictly negative at every}\ \varepsilon_{\mathrm{EM}}>0\ \ \text{on every formed read}\qquad[\mathrm{BM.70}(3),(4)];
\]

%dsp-level:0
\[
(\mathrm{iii})\quad\textbf{the decay:}
\]

%dsp-level:0
\[
\qquad \textbf{closed}\ —\ \neg\text{conditionally open and}\ \neg\text{obstructed by an artifact: the closure is a theorem on the formed reads and unformable on the unformed ones}
\]

%dsp-level:0
\[
\qquad [\mathrm{BM.70}(4);\ (2)];
\]

%dsp-level:0
\[
(\mathrm{iv})\quad\textbf{the derived terms:}
\]

%dsp-level:0
\[
\qquad \Delta m_{\mathrm{transport}};\quad\Delta m_{\mathrm{EM}}\ \text{(scalar and operator)};
\]

%dsp-level:0
\[
\quad I_F=4\,\widehat Q^2;
\]

%dsp-level:0
\[
\quad F_{\ell\nu}\ge2\ln2;\quad\text{the margin ladder}\ (\ln3,\ \ln12,\ \ln24);\quad\Delta m_{\mathrm{bridge}}=0\ \text{in currency};\quad\Omega_{\mathrm{geom}}^{\mathrm{stat}}=1;
\]

%dsp-level:0
\[
(\mathrm v)\quad\textbf{the priced terms:}
\]

%dsp-level:0
\[
\qquad c\in\bigl\{\tfrac12,\tfrac35\bigr\}\ \ \text{priced}\ \neg\text{adopted}\ [\mathrm{Def\text{-}BM.12}(4)];
\]

%dsp-level:0
\[
\qquad \varepsilon_{\mathrm{EM}}\ \ \text{symbolic}\ [\mathrm{EM.25}(9);\ \mathrm{P5}];
\]

%dsp-level:0
\[
\qquad \hat\mu,\ \hat\mu',\ \mu\ \ \text{undeclared}\ [\mathrm{Cal\text{-}F}];
\]

%dsp-level:0
\[
\qquad \rho_{\mathrm B}\ \ \text{restated as the residue}\ [\mathrm{BM.69}(5)];
\]

%dsp-level:0
\[
(\mathrm{vi})\quad\textbf{the open terms:}
\]

%dsp-level:0
\[
\qquad \text{the numerical value of}\ e\ [\mathrm{EM.25}(7)];
\]

%dsp-level:0
\[
\qquad E_{\mathrm{int}}^{(\ell)},\ E_{\mathrm{int}}^{(\nu)}\ [\mathrm{BM.43}(2)];
\]

%dsp-level:0
\[
\qquad \text{the}\ \nu\text{-line index}\;
\]

%dsp-level:0
\[
\qquad \text{the relating equation}\;
\]

%dsp-level:0
\[
\qquad \text{the flavour composition}\;
\]

%dsp-level:0
\[
\qquad \text{the regime-scoped}\ \Sigma_{12}\ \text{fiber}\ [\mathrm{BM.68}(6)];
\]

%dsp-level:0
\[
(5)\quad\textbf{no tuning:}
\]

%dsp-level:0
\[
\qquad \neg\exists\ \text{an adjustment toward an observed ratio}\ —\ \text{the ordering is branch-conditional, the}\ Q\text{-value symbolic, and neither is compared to anything}
\]

%dsp-level:0
\[
\qquad [\mathrm{P5};\ \text{machine: zero external tokens, zero decimal numerals in this tier}];
\]

%dsp-level:0
\[
(6)\quad\textbf{status:}
\]

%dsp-level:0
\[
\qquad \textbf{derived:}\ (1)\text{--}(5);
\]

%dsp-level:0
\[
\qquad \textbf{the no-go is exact,}\ \neg\textbf{rhetorical: every bar in}\ (1)\ \text{carries a cited theorem or a named}\ \neg\text{construction scope};
\]

%dsp-level:0
\[
\neg\text{a P4 addition},\quad\neg\text{a new declaration},\quad\neg\text{a measured value}\ [\mathrm{P5}];\qquad\text{§Cal-BM stays empty}.
\]

\textbf{Req.} A6, D1.11, D4.2, L2 (Part 1); CT4.2, CT4.3, CT4.5, EM.25(7), EM.25(9) (Part 2); P5 (Parts 1–2); FT.3, FT.67(3), FT.68(4), FT.73(1), FT.79(5), FT.107(3) (Part 3); Census FT.35(4) (row-wise); Def-BM.12(4), Def-BM.12(5), Def-BM.12(7), BM.6(5), BM.10(4), BM.28(1), BM.28(4), BM.30(3), BM.30(5), BM.37(3), BM.41(2), BM.41(4), BM.43(2), BM.44(2), BM.45, BM.46(1), BM.46(4), BM.47(6), BM.48(1), BM.48(3), BM.48(4), BM.49(2), BM.50(3), BM.51(2), BM.51(3), BM.52(3), BM.53(4), BM.53(5), BM.53(6), BM.54(1), BM.54(2), BM.54(4), BM.57(4), BM.58(6), BM.59(4), BM.59(9), BM.62(1), BM.62(3), BM.62(6), BM.62(9), BM.62(11), BM.62(12), BM.65(2), BM.65(3), BM.65(5), BM.66(4), BM.66(6), BM.68(5), BM.68(6), BM.69(3), BM.69(4), BM.69(5), BM.70(3), BM.70(4) (this part); Cal-F; exact polynomial identities, Laurent reciprocals, exact rational censuses, sign censuses at positive symbols, token scans [math; machine: \texttt{checks/\allowbreak{}verify\_\allowbreak{}H\_\allowbreak{}organic\_\allowbreak{}bridge\_\allowbreak{}post\_\allowbreak{}bm14.py} §F, §G].

\textbf{Pf.}

%dsp-level:0
\[
(1)\ \text{row}\ 1:\ \mathrm{BM.53}(4)\text{--}(6)\ \text{re-cited};
\]

%dsp-level:0
\[
\qquad 2m_n(E_p-m_p)
\]

%dsp-level:0
\[
= -m_\ell^2\ \text{at}\ m_p
\]

%dsp-level:0
\[
= m_n\ \text{and}\ \lambda
\]

%dsp-level:0
\[
= m_\ell^2(m_\ell^2-4m_n^2)\ \text{re-derived as cleared polynomial identities in}\ \mathbb Q[m_n,m_\ell]\ [\text{machine}];
\]

%dsp-level:0
\[
\qquad \mathfrak T(i\,\mathbb1_d)=d\,i\ne0\ \text{at}\ d=2\ [\mathrm{BM.10}(4);\ \text{machine}];
\]

%dsp-level:0
\[
\qquad \Gamma_{\mathrm{tr}}\in[0,\tfrac13]\ [\mathrm{BM.44}(2)\ \text{re-cited}]
\]

%dsp-level:0
\[
\text{row}\ 2:\ \mathrm{Def\text{-}BM.12}(7)\text{'s}\ \neg\text{constructed list re-cited};
\]

%dsp-level:0
\[
\qquad E^{\mathrm{self}}_{\mathrm{EM}}(q)
\]

%dsp-level:0
\[
= c\,
\]

%dsp-level:0
\[
\varepsilon_{\mathrm{EM}}q^2m\ \text{with}\ q^2:\ 0
\]

%dsp-level:0
\[
\mapsto (1,1)\ \text{across the transition}:\ \ \text{the difference has positive coefficients in}\ (c,\varepsilon_{\mathrm{EM}},m)\ [\text{machine: the sign census}];
\]

%dsp-level:0
\[
\qquad \mathrm{BM.49}(2)\ \text{re-cited}
\]

%dsp-level:0
\[
\text{row}\ 3:\ f(k)=1-\tfrac1k:\quad f(2)=\tfrac12,\ f(3)=\tfrac23,\ f(4)=\tfrac34,\ \ \text{all}\ <1\ [\text{machine: exact Fractions}];
\]

%dsp-level:0
\[
\qquad 6-(2{+}2{+}2)=0\ [\text{machine}];
\]

%dsp-level:0
\[
\qquad \mathrm{BM.48}(4),\ \mathrm{BM.54}(2),(4)\ \text{re-cited}
\]

%dsp-level:0
\[
\text{row}\ 4:\ O(4)=2^5-16=16\ [\mathrm{CT0.6};\ \text{machine}];
\]

%dsp-level:0
\[
\qquad \text{the vertex-type row}\ n(q_1)
\]

%dsp-level:0
\[
= n(p_2)
\]

%dsp-level:0
\[
= 4\ \text{on both classes}\ [\text{Census FT.35}(4);\ \mathrm{BM.57}(4)\text{'s difference table};\ \text{machine: identically zero}];
\]

%dsp-level:0
\[
\qquad \mathrm{FT.73}(1)\text{'s inventory list re-cited}
\]

%dsp-level:0
\[
\text{row}\ 5:\ \mathrm{BM.49}(2),\ \mathrm{BM.41}(4)(\mathrm d),(\mathrm e),\ \mathrm{BM.30}(5)\ \text{re-cited};
\]

%dsp-level:0
\[
\qquad (2):\ \tfrac1{2\,\Delta m_{\mathrm{EM}}}=\tfrac1{2}\cdot\tfrac{2\pi\mathfrak h}{3ce^2\ln2}=\tfrac{\pi\mathfrak h}{3ce^2\ln2}\ [\text{machine: the Laurent reciprocal}]
\]

%dsp-level:0
\[
(3):\ \mathrm{BM.62}(1)\text{'s completeness re-cited with}\ \mathrm{BM.66}(6)\text{'s guard and}\ \mathrm{BM.68}(5)\text{'s slice};
\]

%dsp-level:0
\[
\qquad \Delta_g>0\ \text{on the three certificates}\ [\text{machine}];
\]

%dsp-level:0
\[
\qquad (4)\text{--}(6):\ \text{the verdict rows assembled from}\ \mathrm{BM.65}\text{--}\mathrm{BM.70}\text{'s cited clauses}\ [\text{machine: the row census and the token scans}]. \ \blacksquare
\]

\textbf{Phys.} Every door the audit could name is tried against the corpus's own theorems, and none opens. Uncertainty is a resolvability bound, not a budget: the closure is the non-existence of real on-shell solutions, a polynomial fact at every tick, there is no time operator to borrow from, the ledger runs forward only and contracts strictly, and the flicker that drives every transition weight lives at the amplitude layer, where amplitudes are barred from being masses. Field reconfiguration is a named unconstructed scope, and at the constructed order its sign runs against the decay, because the final state carries two charged members' self-energies where the initial carried none. Binding belongs to the bound, the bound channel does not exist, the nearby space holds only a strict fraction of a positive correction, and both binding reads the corpus owns are zero or class-degenerate. The dark sector has nothing to redistribute: its locks are class-shared, its junction content is identical on both rows, what it retains is inaccessible to every observer channel, and no response read exists by inventory theorem. What remains is the corpus's own honest door — a read specification, an act the tier prices, audits, and declines to perform. Under the bridge the neutron row does not move: the bridge term is zero in the currency in which the row is formed, conditionally positive in the pole currency where it is not formed, and the six verdict rows record the signs, the closure, and the derived/priced/open partition exactly — with nothing tuned toward anything observed.

\medskip\noindent\textit{\texorpdfstring{BM.72 — The status of the organic bridge: adjudication and scope}{BM.72 - The status of the organic bridge: adjudication and scope}.}\ Structural-constraints record: see Appendix B, Section B.1.24 of the Addendum.\par\medskip

\phantomsection
\section[Tier BM.16 — The Chiral Assignment: the Closed Deficit, the…]{\texorpdfstring{Tier BM.16 — The Chiral Assignment: the Closed Deficit, the Ray-Resolved Spin Census, the Attached Quark Splitting, the Chiral Vertex Dynamics, and the Specified Beta-Decay Opening}{Tier BM.16 - The Chiral Assignment: the Closed Deficit, the Ray-Resolved Spin Census, the Attached Quark Splitting, the Chiral Vertex Dynamics, and the Specified Beta-Decay Opening}}

\textbf{Locus.} This tier consumes the Tier BM.15 baseline unmodified and walks the two doors the corpus priced and declined: the deficit-$1'$ assignment rule — named missing at $\mathrm{FT.125}(4)$, $\mathrm{FT.141}(6)$, $\mathrm{FT.144}(5)$ and $\mathrm{FT.157}(4)$, with its forced shape already derived — and the $\beta$-vertex read specification — named as the sole opening route at $\mathrm{BM.30}(5)$, $\mathrm{BM.41}(4)(\mathrm e)$, $\mathrm{BM.49}(2)$ and $\mathrm{BM.71}(1)$ row $5$. Both acts are logged as §II.4 registry addenda (A-BM.7, A-BM.8) before use, per $\mathrm{Def\text{-}BM.3}(7)$'s rule. $\mathrm{BM.73}$ issues the left–right assignment matching $\hat\omega$ and $\hat\chi$ and derives the occupancy; $\mathrm{BM.74}$ resolves the spin census by ray; $\mathrm{BM.75}$ derives the weight–branch–core attachment and attaches the quark mass splitting; $\mathrm{BM.76}$ constructs the chiral vertex dynamics; $\mathrm{BM.77}$ issues the vertex read specification and proves the exact opening branch; $\mathrm{BM.78}$ adjudicates.

\textbf{Standing rows.} The refutations of $\mathrm{BM.57}(8)$ and $\mathrm{BM.58}(4)$ stand closed and are $\neg$reopened. $\mathrm{BM.70}(4)$'s dichotomy and $\mathrm{BM.71}(2)$'s no-go stand on the reads formed at their loci; $\mathrm{BM.77}(7)$ records their rescope under the specification, and no clause of $\mathrm{BM.41}$–$\mathrm{BM.72}$ is retracted. The relating equation between the informational and the geometric currencies is $\neg$assumed anywhere — item 7 stays open [$\mathrm{BM.35}(5)$; $\mathrm{BM.68}(7)$; pointer]. $\alpha_B$ stays barred on $\mathrm{BM.61}(8)$'s four grounds. $\neg$a Yukawa-type object exists or is added [$\mathrm{EW.15}$; $\mathrm{FT.145}(3)$]. No external mass, lifetime, endpoint or coupling value enters any derivation [$\mathrm{P5}$]. \texttt{}Chirality $=$ helicity on-shell'' remains external vocabulary [$\mathrm{FT.123}(3),(4)$]; the spin-structure torsor remains a distinct kind [$\mathrm{FT.122}(7)$; $\mathrm{FT.89}(2)$]; deficit $11$ is untouched [$\mathrm{FT.125}(5)$].

\begin{theorem}[{\texorpdfstring{BM.73 — The chiral matter assignment: the issued coupling, the occupancy theorem, and the $\omega$/$\chi$ matching rows}{BM.73 - The chiral matter assignment: the issued coupling, the occupancy theorem, and the / matching rows}}]\label{tag:BM-73}
\textbf{Thm.}

%dsp-level:0
\[
(1)\quad\textbf{the four derived gradings (the census the rule couples):}\qquad\begin{array}{l|l}\text{the algebra}\ S\text{-grading} & (10,\,6)=\bigl(\dim\mathcal K_1^{+},\ \dim\mathcal K_1^{-}\bigr)\ \text{on the force arena};\ \ \text{metric/weak reads}\ S\text{-even},\ \text{seal reads}\ S\text{-odd}\qquad[\mathrm{FT.136};\ \mathrm{FT.141}(1),(2),(5)]\\\text{the branch-label parity} & \text{spatial}=\text{aligned}\ (\alpha{=}0,\ \vee^2,\ \text{Casimir}\ 2)\ \text{weak-active};\ \ \text{temporal}=\text{antialigned}\ (\alpha{=}\pi,\ \wedge^2,\ \text{Casimir}\ 0)\ \text{weak-blind}\qquad[\mathrm{FT.101};\ \mathrm{FT.135};\ \mathrm{FT.157}(2),(3)]\\\text{the carrier}\ \hat\chi\text{-grade} & \hat\chi=i\bigl(|c\rangle\langle\omega(c)|-|\omega(c)\rangle\langle c|\bigr),\quad P_{L/R}=\tfrac12(\mathrm{id}\pm\hat\chi),\quad|L\rangle=\tfrac1{\sqrt2}\bigl(|c\rangle-i|\omega(c)\rangle\bigr)\qquad[\mathrm{Def\text{-}H.33};\ \mathrm{FT.122}(1)\text{--}(3),(6)]\\\text{the Janus-arrow chirality} & \text{strict alternation, the prefix law, the single-arrow materialized locus at rung one}\qquad[\mathrm{FT.144}(2),(5)]\end{array}
\]

%dsp-level:0
\[
\textbf{the pairwise distinctness is derived:}
\]

%dsp-level:0
\[
\qquad \text{the four are pairwise distinct}\ \mathbb Z_2\text{'s, commuting where jointly defined, neither reducible to another}
\]

%dsp-level:0
\[
\qquad [\mathrm{FT.136}(4)\text{--}(5);\ \mathrm{FT.141}(6);\ \mathrm{FT.126}(4);\ \mathrm{FT.157}(4)];
\]

%dsp-level:0
\[
(2)\quad\textbf{the assignment rule (the declaration, logged as A-BM.7 before use):}
\]

%dsp-level:0
\[
\qquad \textbf{the forced shape is derived:}\ \ \text{any}\ \hat\omega\text{-covariant arrow}\leftrightarrow\text{ray coupling maps the single-arrow locus into a single chiral ray, and the complementary ray has occupancy zero at rung one}\ —\ \text{exhaustive on the}\ 2!\ \text{bijections}
\]

%dsp-level:0
\[
\qquad [\mathrm{FT.144}(5)]:
\]

%dsp-level:0
\[
\boxed{\ \textbf{the issued rule:}\qquad\text{the materialization-arrow locus occupies}\ \operatorname{im}P_L\ \ \text{— the}\ \hat\chi=+1\ \text{ray of}\ \mathrm{Def\text{-}H.33}\text{'s boxed sign}\ }
\]

%dsp-level:0
\[
\textbf{the coupled rows (each forced given the choice, or the choice itself):}
\]

%dsp-level:0
\[
\qquad (\mathrm a)\ \ \text{the occupied ray carries the weak-active/aligned/spatial branch label};
\]

%dsp-level:0
\[
\ \text{the complement carries the weak-blind/antialigned/temporal label}
\]

%dsp-level:0
\[
\qquad [\text{forced: every materialized bound chain is weak-active}\ —\ \{G,H,W\}\subseteq\mathcal A(\Sigma_{\min})\ [\text{Census FT.28}(3)];\ \text{weak acts from rung one}\ [\mathrm{CT0.13}];\ \text{the rung-one channel class is weak-carried}\ [\mathrm{FT.142}(1)]:\ \text{the alternative coupling would make every matter line weak-blind against those derived rows};\ \text{machine: the}\ 2\ \text{couplings censused, one consistent}];
\]

%dsp-level:0
\[
(\mathrm b)\ \ \text{the matter orientation occupies the seed weak doublet on the left ray}\ [\mathrm{EW.32};\ \mathrm{EW.36}:\ \operatorname{spec}Q|_{\mathcal H_{\mathsf v}}=\{0,+1\}];
\]

%dsp-level:0
\[
\qquad \text{the image orientation occupies the}\ \omega\text{-image representation}\ \{0,-1\}\ \text{on the right ray}\ [\mathrm{EW.36};\ \mathrm{D2.6};\ \mathrm{BM.25};\ \mathrm{BM.27};\ \text{the A-BM.5 antineutrino row}]
\]

%dsp-level:0
\[
\qquad [\text{forced by}\ \hat\omega\text{-covariance once}\ (\mathrm a)\ \text{holds}:\ \hat\omega P_L\hat\omega^{-1}=P_R\ [\mathrm{FT.122}(4)]];
\]

%dsp-level:0
\[
(\mathrm c)\ \ \text{the}\ S\text{-grading keeps its derived division of labour}\ —\ \text{it grades observables,}\ \neg\text{states}\ —\ \text{and receives no species coupling}
\]

%dsp-level:0
\[
\qquad [\mathrm{FT.141}(6)]:
\]

%dsp-level:0
\[
\textbf{the declared content is exactly one naming bit}\ —\ \text{which ray the arrow occupies is named left}\ —\ \text{on the forced shape of}\ [\mathrm{FT.144}(5)];
\]

%dsp-level:0
\[
\qquad \textbf{it carries no datum:}\ \neg\text{a scale},
\]

%dsp-level:0
\[
\ \neg\text{a coupling value},
\]

%dsp-level:0
\[
\ \neg\text{a measured quantity}\ [\mathrm{P5};\ \mathrm{Def\text{-}BM.3}(7)\ \text{pattern}];
\]

%dsp-level:0
\[
(3)\quad\textbf{the occupancy theorem (derived given}\ (2)):
\]

%dsp-level:0
\[
\qquad \boxed{\ \text{the rung-one materialized locus occupies}\ (1,\,0)\ =\ (\text{left occupied},\ \text{right empty})\ }
\]

%dsp-level:0
\[
\text{the complement's content is the vacuum sector: all-virtual, unregistered, no accessible pole}
\]

%dsp-level:0
\[
\qquad [\mathrm{FT.144}(4),(5);\ \mathrm{CT4.2}\ \text{pattern};\ \mathrm{CT0.1}];
\]

%dsp-level:0
\[
\textbf{the rung ladder of the occupancy:}
\]

%dsp-level:0
\[
\qquad \text{rung one: the sectors are force-decoupled},\ F_{LR}=F_{RL}=0\ \text{exactly}\ [\mathrm{FT.128}(2);\ \mathrm{FT.144}(6)];
\]

%dsp-level:0
\[
\qquad \text{rung two: the EM core read opens the unique flip channel},
\]

%dsp-level:0
\[
\ F_{LR}
\]

%dsp-level:0
\[
= q\,
\]

%dsp-level:0
\[
I_2
\]

%dsp-level:0
\[
\ne 0
\]

%dsp-level:0
\[
\iff q
\]

%dsp-level:0
\[
\ne 0,
\]

%dsp-level:0
\[
\ \text{splitting}\ 2|q|\in\{0,2\}\ [\mathrm{FT.129}(5);\ \mathrm{FT.130}(5);\ \mathrm{Def\text{-}H.43}\ \text{promotion inherited};\ \mathrm{BM.22}(2),(4)];
\]

%dsp-level:0
\[
\text{rung three: colour adds no carrier flip}\ —\ \text{sealed ends are blind},
\]

%dsp-level:0
\[
\ \text{the seed is colour-neutral}\ [\mathrm{FT.117};\ \mathrm{EW.32};\ \mathrm{FT.144}(6)];
\]

%dsp-level:0
\[
\qquad \textbf{the flipping set remains}\ \{F_{\mathrm{EM}}\ \text{core read}\},\ \text{minimal and unique}\ [\mathrm{FT.129}(5)];
\]

%dsp-level:0
\[
(4)\quad\textbf{the}\ \omega/\chi\ \textbf{matching rows (re-certified exact over}\ \mathbb Q(i)):
\]

%dsp-level:0
\[
\qquad \hat\chi^\dagger=\hat\chi,\quad\hat\chi^2=\mathrm{id},\quad\operatorname{tr}\hat\chi=0,\quad\operatorname{spec}\hat\chi=\{+1,-1\};
\]

%dsp-level:0
\[
\qquad P_L^2=P_L,\quad P_LP_R=0,\quad P_L+P_R=\mathrm{id},\quad\operatorname{tr}P_L=\operatorname{tr}P_R=1
\]

%dsp-level:0
\[
\qquad [\mathrm{FT.122}(1),(2);\ \text{machine}]:
\]

%dsp-level:0
\[
\boxed{\ \hat\omega\hat\chi\hat\omega^{-1}=-\hat\chi,\qquad\{\hat\chi,\hat\omega\}=0,\qquad\hat\omega P_L\hat\omega^{-1}=P_R,\qquad\hat\omega|L\rangle=-i|R\rangle,\quad\hat\omega|R\rangle=+i|L\rangle\ }
\]

%dsp-level:0
\[
\qquad [\mathrm{FT.122}(4);\ \text{machine}];
\]

%dsp-level:0
\[
\textbf{the quaternionic triple:}
\]

%dsp-level:0
\[
\qquad \hat\omega^2=\hat\chi^2=\hat\kappa^2=\mathrm{id};
\]

%dsp-level:0
\[
\qquad \{\hat\omega,\hat\chi\}=\{\hat\chi,\hat\kappa\}=\{\hat\kappa,\hat\omega\}=0;
\]

%dsp-level:0
\[
\qquad \hat\omega\hat\chi=-i\hat\kappa,\quad\hat\chi\hat\kappa=-i\hat\omega,\quad\hat\kappa\hat\omega=-i\hat\chi
\]

%dsp-level:0
\[
\qquad [\mathrm{Def\text{-}H.34};\ \text{machine}];
\]

%dsp-level:0
\[
\textbf{the winding row:}
\]

%dsp-level:0
\[
\qquad \omega:\ k\mapsto-k
\]

%dsp-level:0
\[
\qquad [\mathrm{FT.122}(4);\ \mathrm{D2.6};\ \mathrm{EM.24}(1)]:
\]

%dsp-level:0
\[
\qquad \text{the ray exchange and the winding reversal are the two faces of one involution};
\]

%dsp-level:0
\[
\textbf{the canonicity:}
\]

%dsp-level:0
\[
\qquad \bigl\{M=M^\dagger:\ \hat\omega M\hat\omega=-M,\ M^2=\mathrm{id},\ \langle c|M|c\rangle=\langle\omega(c)|M|\omega(c)\rangle=0\bigr\}=\{+\hat\chi,\,-\hat\chi\}
\]

%dsp-level:0
\[
\qquad [\mathrm{FT.122}(6);\ \text{machine: the circle}\ M(a,y),\ a^2+y^2=1,\ \text{on exact rational witnesses}];
\]

%dsp-level:0
\[
(5)\quad\textbf{the weak rows under the assignment:}
\]

%dsp-level:0
\[
\qquad \bigl[W_i\otimes I_\omega,\ I\otimes P_{L/R}\bigr]=0
\]

%dsp-level:0
\[
\qquad [\mathrm{FT.123}(1);\ \mathrm{FT.126}(2);\ \text{machine}]:
\]

%dsp-level:0
\[
\qquad \textbf{the assignment attaches species to rays; it does not modify the weak action};
\]

%dsp-level:0
\[
\text{the selectivity lives at the state locus:}\quad W_i\Pi_{\mathcal V}=0,
\]

%dsp-level:0
\[
\qquad \textstyle\sum_iW_i^2=2\,\Pi_{\mathcal R}
\]

%dsp-level:0
\[
\qquad [\mathrm{FT.157}(2);\ \mathrm{FT.141}(4);\ \mathrm{FT.101}(1)\text{--}(2);\ \text{machine re-cited}];
\]

%dsp-level:0
\[
\textbf{the occupied ray is the weak-active one by}\ (2)(\mathrm a);
\]

%dsp-level:0
\[
\qquad \text{a chiral gauge coupling as an operator modification is}\ \neg\text{constructed and}\ \neg\text{needed for the assignment}
\]

%dsp-level:0
\[
\qquad [\mathrm{FT.123}\ \mathrm{Rmk}\ \text{scope}];
\]

%dsp-level:0
\[
(6)\quad\textbf{the bars:}
\]

%dsp-level:0
\[
\qquad \text{the weight}\leftrightarrow\text{ray attachment is barred as the external helicity vocabulary}
\]

%dsp-level:0
\[
\qquad [\mathrm{FT.123}(3),(4)];
\]

%dsp-level:0
\[
\qquad \text{the spin-structure torsor is reversal-invariant, a distinct kind}
\]

%dsp-level:0
\[
\qquad [\mathrm{FT.122}(7);\ \mathrm{FT.89}(2)];
\]

%dsp-level:0
\[
\text{deficit}\ 11\ \text{(spin--statistics) is untouched: no exchange operator exists}
\]

%dsp-level:0
\[
\qquad [\mathrm{FT.80}(5);\ \mathrm{FT.125}(5)];
\]

%dsp-level:0
\[
\qquad \text{the frame section}\ e\mapsto\mathcal F_e\ \text{and a global Dirac-type operator remain unconstructed}
\]

%dsp-level:0
\[
\qquad [\mathrm{FT.125}(1);\ \text{item}\ 21(\mathrm i)\ \text{pointer}];
\]

%dsp-level:0
\[
\textbf{no SM species reading and no claim that the SM weak interaction is reproduced}
\]

%dsp-level:0
\[
\qquad [\mathrm{FT.157}\ \mathrm{Rmk}\ \text{bars stand};\ \mathrm{P1},\ \mathrm{P5}];
\]

%dsp-level:0
\[
(7)\quad\textbf{status:}
\]

%dsp-level:0
\[
\qquad \textbf{derived:}\ (1),\ (3)\text{--}(6)\ \ [(3)\ \text{given}\ (2)];
\]

%dsp-level:0
\[
\qquad \textbf{declared: exactly one naming bit}\ (2)\ —\ \text{logged as the §II.4 addendum A-BM.7 before use};
\]

%dsp-level:0
\[
\textbf{deficit}\ 1'\ \textbf{is closed at the assignment level:}
\]

%dsp-level:0
\[
\qquad \text{the coupling rule exists}\ —\ \text{the declaration's root is issued};
\]

%dsp-level:0
\[
\qquad \text{the downstream rows unblock at their own loci and stay open:}\ \text{deficits}\ 2,
\]

%dsp-level:0
\[
\ 3,
\]

%dsp-level:0
\[
\ 4\ [\mathrm{FT.125}(6)],
\]

%dsp-level:0
\[
\ \ \mathrm{Q4}\text{--}\mathrm{Q6}\ \text{residues},
\]

%dsp-level:0
\[
\ \ \text{the anomaly sums formable on the}\ 1'\ \text{side with the}\ +\tfrac23\ \text{side structurally unreachable}\ [\mathrm{FT.125}(2);\ \mathrm{BM.28}(7)];
\]

%dsp-level:0
\[
\neg\text{a P4 addition},\quad\neg\text{a new dimensionful scale}\ [\mathrm{P1}],\quad\neg\text{a measured value}\ [\mathrm{P5}],\quad\text{§Cal-BM stays empty};
\]

%dsp-level:0
\[
\qquad \text{the extra-axiomatic list}\ [\text{A §0.2a}]\ \text{gains no member}.
\]

\textbf{Req.} A4, A8 (Part 1); D2.6, CT0.1, CT0.13 (Part 1); CT4.2, EM.24(1), EW.32, EW.36 (Part 2); P1, P5 (Parts 1–2); Def-H.33, Def-H.34, Def-H.43, FT.80(5), FT.89(2), FT.101(1), FT.101(2), FT.101(3), FT.117, FT.122(1), FT.122(2), FT.122(3), FT.122(4), FT.122(6), FT.122(7), FT.123(1), FT.123(3), FT.123(4), FT.125(1), FT.125(2), FT.125(4), FT.125(5), FT.125(6), FT.126(2), FT.126(4), FT.128(2), FT.129(5), FT.130(5), FT.135, FT.136(1), FT.136(4), FT.136(5), FT.141(1), FT.141(2), FT.141(4), FT.141(5), FT.141(6), FT.142(1), FT.144(2), FT.144(4), FT.144(5), FT.144(6), FT.157(2), FT.157(3), FT.157(4) (Part 3); Census FT.28(3) (row-wise); Def-BM.3(7), BM.22(2), BM.22(4), BM.25, BM.27, BM.28(7) (this part); exact $\mathbb Q(i)$ matrix algebra, exhaustive bijection censuses, rational-circle witnesses [math; machine: \texttt{checks/\allowbreak{}verify\_\allowbreak{}H\_\allowbreak{}chiral\_\allowbreak{}assignment\_\allowbreak{}opening.py} §A, §B].

\textbf{Pf.}

%dsp-level:0
\[
(1):\ \text{the four rows re-cited from their own tags}\ [\mathrm{FT.136}(1),\ \mathrm{FT.141}(1),(2),(5);\ \mathrm{FT.157}(2),(3);\ \mathrm{Def\text{-}H.33};\ \mathrm{FT.122}(1)\text{--}(3);\ \mathrm{FT.144}(2),(5)];
\]

%dsp-level:0
\[
\qquad \text{the distinctness rows re-cited}\ [\mathrm{FT.126}(4);\ \mathrm{FT.136}(4)\text{--}(5);\ \mathrm{FT.141}(6);\ \mathrm{FT.157}(4)]
\]

%dsp-level:0
\[
(2):\ \mathrm{FT.144}(5)\text{'s forced shape re-cited}\ [\text{machine: exhaustive on the}\ 2!\ \text{bijections}];
\]

%dsp-level:0
\[
\qquad (\mathrm a):\ \text{the weak-activity census: both ray/label couplings evaluated against}\ \{G,H,W\}\subseteq\mathcal A(\Sigma_{\min})\ [\text{Census FT.28}(3)],
\]

%dsp-level:0
\[
\ \mathrm{CT0.13},
\]

%dsp-level:0
\[
\ \mathrm{FT.142}(1):\ \text{the weak-blind-occupied coupling contradicts the derived weak activity of every matter floor};
\]

%dsp-level:0
\[
\ \text{one coupling survives}\ [\text{machine}];
\]

%dsp-level:0
\[
\qquad (\mathrm b):\ \hat\omega P_L\hat\omega^{-1}=P_R\ [\text{machine over}\ \mathbb Q(i)]\ \text{and}\ \omega:\operatorname{spec}Q\mapsto\{0,-1\}\ [\mathrm{EW.36}\ \text{re-cited}];
\]

%dsp-level:0
\[
\qquad (\mathrm c):\ \mathrm{FT.141}(6)\ \text{verbatim}
\]

%dsp-level:0
\[
(3):\ \mathrm{FT.144}(5)\text{'s occupancy pair}\ (1,0)\ \text{with}\ (2)\text{'s naming};
\]

%dsp-level:0
\[
\qquad \text{the rung rows re-cited}\ [\mathrm{FT.128}(2);\ \mathrm{FT.129}(5);\ \mathrm{FT.130}(5);\ \mathrm{FT.117};\ \mathrm{FT.144}(6)]
\]

%dsp-level:0
\[
(4):\ \hat\chi=i\left(\begin{smallmatrix}0&1\\-1&0\end{smallmatrix}\right),\ \hat\omega=\left(\begin{smallmatrix}0&1\\1&0\end{smallmatrix}\right),\ \hat\kappa=\left(\begin{smallmatrix}1&0\\0&-1\end{smallmatrix}\right)\ \text{in the ordered orientation basis}\ [\mathrm{FT.122}(5);\ \mathrm{Def\text{-}H.34}]:\ \text{every identity of}\ (4)\ \text{verified entrywise over}\ \mathbb Q(i)\ [\text{machine}];\qquad\hat\omega|L_u\rangle=|R_u\rangle\cdot(-i)\ \text{on the unnormalized kets}\ [\text{machine}];\qquad\text{the canonicity circle}\ M(a,y)^2=(a^2{+}y^2)\mathrm{id}\ \text{and the diagonal-free reduction}\ a=0\implies y^2=1\implies M=\pm\hat\chi\ [\text{machine: on the rational witnesses}\ (\tfrac35,\tfrac45),\ (\tfrac5{13},\tfrac{12}{13}),\ (\tfrac8{17},\tfrac{15}{17})]
\]

%dsp-level:0
\[
(5):\ \text{the commutator}\ [W_i\otimes I,\ I\otimes P]=0\ \text{by tensor-factor disjointness}\ [\text{machine}];
\]

%dsp-level:0
\[
\qquad \text{the locus rows re-cited}\ [\mathrm{FT.157}(2);\ \mathrm{FT.141}(4)]
\]

%dsp-level:0
\[
(6)\text{--}(7):\ \text{the bars re-cited clause by clause};\qquad\text{the standing-rule scans}\ [\text{machine}]. \ \blacksquare
\]

\textbf{Phys.} The corpus had derived everything about chirality except one bit: which of the two exchanged rays the materialized world sits in. The four gradings — the force arena's ten/six split, the weak-active and weak-blind branch labels, the orientation doublet's exchange-form grade, and the Janus arrow's single-direction locus — are pairwise distinct, each exact, and the rule marrying them to matter was the ledger's named missing declaration. This tier issues it in its minimal form: the arrow's occupied ray is the left ray, the one the boxed chirality sign already fixes algebraically; the weak-active branch label then rides along by consistency, because every materialized matter floor is weak-active by derived rows and the opposite coupling would contradict them; the right ray is the orientation reversal's image and hosts the image-orientation species the antineutrino naming row already used. What is derived given the bit is substantial: the occupancy pair one-zero with its vacuum complement, the rung ladder along which the flip channel opens at rung two for charged cores and never for colour, and the whole $\omega$/χ matching — the involution exchanges the rays, anticommutes with the grade, and reverses the winding, so ray exchange and matter–image exchange are one move. What is not obtained is stated as loudly: no helicity claim, no spin-statistics, no Standard Model reading, and no modification of the weak action itself, which stays chirality-blind on the carrier while the species sit where the rule puts them.

\end{theorem}

\pseudohead{theorem}{\texorpdfstring{BM.74}{BM.74}}{\texorpdfstring{The ray-resolved spin census: the winding labels, the occupancy rows, and the vertex table}{The ray-resolved spin census: the winding labels, the occupancy rows, and the vertex table}}\label{tag:BM-74}
\textbf{Thm.}

%dsp-level:0
\[
(1)\quad\textbf{the spin-label rule (the declared interface inherited,}\ \neg\textbf{re-derived):}
\]

%dsp-level:0
\[
\qquad \text{a materialized chain of winding}\ k\ \text{carries}\ \mathrm{Sym}^{k}(\mathcal H_{\mathsf v}),
\]

%dsp-level:0
\[
\qquad j=\frac{|k|}2,
\]

%dsp-level:0
\[
\qquad \text{weights}\ j,\ j{-}1,\ \dots,\ -j
\]

%dsp-level:0
\[
\qquad [\mathrm{Def\text{-}H.38};\ \mathrm{FT.55}(2);\ \mathrm{FT.60}(1)\text{--}(3);\ \mathrm{FT.123}(3)]:
\]

%dsp-level:0
\[
\textbf{the anchor:}
\]

%dsp-level:0
\[
\qquad k=1:\quad\mathrm{Sym}^{1}=\mathcal H_{\mathsf v}\cong\mathbb C^2,
\]

%dsp-level:0
\[
\qquad j=\tfrac12,
\]

%dsp-level:0
\[
\qquad \text{weights}\ \bigl\{+\tfrac12,\,-\tfrac12\bigr\}
\]

%dsp-level:0
\[
\qquad [\mathrm{EW.32};\ \text{Census FT.28}(2);\ \text{Census FT.63}(1);\ \text{machine}];
\]

%dsp-level:0
\[
\textbf{the Casimir ladder:}
\]

%dsp-level:0
\[
\qquad \frac k2\Bigl(\frac k2+1\Bigr)
\]

%dsp-level:0
\[
\qquad [\mathrm{FT.60}(3)];
\]

%dsp-level:0
\[
\qquad \text{the multiplicity--winding coupling itself stays a declared interface}\ [\mathrm{Def\text{-}H.38}]\ —\ \neg\text{derived here},
\]

%dsp-level:0
\[
\ \neg\text{reopened};
\]

%dsp-level:0
\[
(2)\quad\textbf{the ray-resolved line census (derived given A-BM.7):}
\]

%dsp-level:0
\[
\qquad \text{every materialized matter line occupies the left ray: occupancy}\ (1,0)
\]

%dsp-level:0
\[
\qquad [\mathrm{BM.73}(3)]:
\]

%dsp-level:0
\[
\text{the image-orientation lines}\ —\ \text{the CPT partners and the A-BM.5 antineutrino row}\ —\ \text{occupy the right ray}
\]

%dsp-level:0
\[
\qquad [\mathrm{BM.73}(2)(\mathrm b);\ \mathrm{BM.25};\ \mathrm{BM.27};\ \text{machine: the}\ \hat\omega\text{-exchange row}];
\]

%dsp-level:0
\[
\textbf{the cardinality correspondence stays a cardinality row:}
\]

%dsp-level:0
\[
\qquad \text{the two weights}\ \bigl\{\pm\tfrac12\bigr\}\ \text{and the two rays}\ \{+1,-1\}\ \text{are an unordered-pair match};
\]

%dsp-level:0
\[
\ \ \text{the weight}\leftrightarrow\text{ray attachment is barred}
\]

%dsp-level:0
\[
\qquad [\mathrm{FT.123}(3);\ \text{the external helicity vocabulary}];
\]

%dsp-level:0
\[
(3)\quad\textbf{the composite rows:}
\]

%dsp-level:0
\[
\qquad \text{the baryon card:}\quad\Theta_{\mathrm{int}}=2\,k(\Gamma)\in2\mathbb Z,
\]

%dsp-level:0
\[
\qquad j(\Gamma)=|k(\Gamma)|\ \ \text{rung-relative}\ —\ \text{the card fixes no rung}
\]

%dsp-level:0
\[
\qquad [\text{Census FT.35}(3);\ \text{Census FT.65}(4):\ \text{the free}\ \mathbb Z^3\ \text{ladder}];
\]

%dsp-level:0
\[
\text{the Platonic floors at the uniform minima: every face}\ j_f\ \text{composes to an integer stretched total},
\]

%dsp-level:0
\[
\quad\sum_fk_{f,\min}/2
\]

%dsp-level:0
\[
= K_{S,\min}/2\in\{6,\,6,\,12,\,30,\,30\}
\]

%dsp-level:0
\[
\qquad [\text{Census FT.63}(1)(\mathrm b);\ \text{machine re-cited}];
\]

%dsp-level:0
\[
\text{the defect multiplet ladder:}
\]

%dsp-level:0
\[
\qquad m^2_{\mathrm{pole}}(m)=4\,\hat\mu'^2\sin^2\!\bigl(m\,\tfrac{\Theta_R}4\bigr),
\]

%dsp-level:0
\[
\qquad m\in\{k,\,k{-}2,\ \dots,\,-k\}
\]

%dsp-level:0
\[
\qquad [\mathrm{FT.60}(5);\ \text{Census FT.63}(6)]:
\]

%dsp-level:0
\[
\qquad \text{the six states of}\ \mathrm{Def\text{-}BM.14}\ \text{ride it with}\ j=|k|/2;
\]

%dsp-level:0
\[
\qquad \text{the mesonic}\ \delta\text{-sealed rows stand untouched}\ [\mathrm{BM.31}\text{--}\mathrm{BM.33}];
\]

%dsp-level:0
\[
(4)\quad\textbf{the vertex table (ray-resolved; derived given A-BM.7 and the frozen grammar):}
\]

%tbl-level:3
{\scriptsize\setlength{\tabcolsep}{4pt}
\begin{longtable}{>{\raggedright\arraybackslash}p{\dimexpr0.0881\textwidth-2\tabcolsep\relax}|>{\raggedright\arraybackslash}p{\dimexpr0.2795\textwidth-2\tabcolsep\relax}|>{\raggedright\arraybackslash}p{\dimexpr0.1999\textwidth-2\tabcolsep\relax}|>{\raggedright\arraybackslash}p{\dimexpr0.4326\textwidth-2\tabcolsep\relax}}
\toprule
member & the structural row & the spin/rung row & the ray row \\
\addlinespace[2pt]
\midrule
\endfirsthead
\toprule
member & the structural row & the spin/rung row & the ray row \\
\addlinespace[2pt]
\midrule
\endhead
\bottomrule
\endlastfoot
the parent & the neutral-class $\varepsilon$-sealed 3-juncture, $Q=\allowbreak 0$ $[\mathrm{BM.17};\ \mathrm{Def\text{-}BM.3}(2)]$ & $j=\allowbreak |k(\Gamma)|$, rung-relative $[\text{Census FT.35}(3)]$ & the constituents occupy the left ray; the neutral-core constituents are ray-locked: $F_{LR}=\allowbreak 0$, parity doubling $[\mathrm{BM.22}(4);\ \mathrm{FT.130}(6)]$ \\
the daughter & the unit-class $\varepsilon$-sealed 3-juncture, $Q=\allowbreak +\allowbreak 1$ $[\mathrm{BM.26};\ \mathrm{Def\text{-}BM.3}(3)]$ & the same ladder & the charged-core constituents are flip-coupled, splitting $2|q|=\allowbreak 2$ $[\mathrm{FT.130}(5);\ \mathrm{FT.144}(6)]$ \\
the charged emitted member & the $\ell$-channel line, $q=\allowbreak -1$: forced, class-unique $[\mathrm{BM.23}(5);\ \mathrm{BM.41}(1);\ \mathrm{Def\text{-}BM.6}(1)]$ & the $k=\allowbreak 1$ anchor: $j=\allowbreak \tfrac12$ $[\text{Census FT.28}(2);\ \mathrm{Def\text{-}H.38}]$ & flip-coupled ($q\ne\allowbreak 0$); the arrow orientation occupies the left ray with the partner accessible $[\mathrm{BM.73}(3);\ \mathrm{FT.129}(5)]$ \\
the neutral emitted member & the $\nu$-channel line in the image orientation, $q=\allowbreak 0$ $[\mathrm{BM.41}(3);\ \mathrm{FT.142}(1);\ \text{A-BM.5}]$ & the $k=\allowbreak 1$ anchor: $j=\allowbreak \tfrac12$ & the right ray (the image orientation under A-BM.7); flip-blind, parity doubling $[\mathrm{BM.22}(4);\ \mathrm{FT.130}(6)]$ \\
\end{longtable}
}

%dsp-level:0
\[
\textbf{the two emitted members occupy opposite rays}
\]

%dsp-level:0
\[
\qquad [\text{machine: the table's ray column is a derived consequence of}\ \mathrm{BM.73}(2)(\mathrm b),\ (3)\ \text{and the frozen grammar}];
\]

%dsp-level:0
\[
(5)\quad\textbf{the angular-momentum budget is unchanged:}
\]

%dsp-level:0
\[
\qquad \mathrm{BM.42}\text{'s rows stand}\ —\ \text{the Noether moment conserved on-shell},
\]

%dsp-level:0
\[
\ \text{the Lagrange decomposition}\ p^2
\]

%dsp-level:0
\[
= p_r^2+L^2/r^2,
\]

%dsp-level:0
\[
\ \text{the absorption adjudications},
\]

%dsp-level:0
\[
\ \text{the}\ \ell\text{-selection census}
\]

%dsp-level:0
\[
\qquad [\mathrm{BM.42}(1)\text{--}(4);\ \mathrm{BM.40}(3):\ \text{the}\ \sigma\text{-vertex is Casimir-neutral}]:
\]

%dsp-level:0
\[
\textbf{the ray label adds no torque term:}
\]

%dsp-level:0
\[
\qquad \hat\chi\ \text{acts on the orientation factor alone and commutes with the rotation generators acting on the seed}
\]

%dsp-level:0
\[
\qquad [\mathrm{FT.123}(1)\ \text{pattern};\ \mathrm{FT.126}(4):\ \text{independent}\ \mathbb Z_2\text{'s};\ \text{machine: the commutant row}];
\]

%dsp-level:0
\[
(6)\quad\textbf{what is not obtained:}
\]

%dsp-level:0
\[
\qquad \neg\text{a spin--statistics theorem}\ —\ \text{no exchange operator exists}\ [\mathrm{FT.80}(5);\ \mathrm{FT.125}(5)];
\]

%dsp-level:0
\[
\qquad \neg\text{a fixed}\ j\ \text{for the baryon card}\ —\ \text{the rung datum is not fixed here}\ [\mathrm{BM.68}(6)\ \text{regime row};\ \text{Census FT.65}(4)];
\]

%dsp-level:0
\[
\neg\text{a weight}\leftrightarrow\text{ray helicity row}\ [\mathrm{FT.123}(3),(4)];
\]

%dsp-level:0
\[
\qquad \neg\text{a}\ k\ge3\ \text{carrier-inventory extension}\ [\mathrm{FT.60}(8)];
\]

%dsp-level:0
\[
(7)\quad\textbf{status:}
\]

%dsp-level:0
\[
\qquad \textbf{derived given A-BM.7:}\ (2)\text{--}(5);
\]

%dsp-level:0
\[
\qquad \textbf{inherited declarations:}\ \mathrm{Def\text{-}H.38}\text{'s multiplicity coupling}\ \ \text{and the naming rows}\ [\mathrm{Def\text{-}BM.3};\ \mathrm{Def\text{-}BM.6};\ \text{A-BM.5}];
\]

%dsp-level:0
\[
\neg\text{a P4 addition},\quad\neg\text{a measured value}\ [\mathrm{P5}],\quad\neg\text{a new dimensionful scale}\ [\mathrm{P1}];\qquad\text{§Cal-BM stays empty}.
\]

\textbf{Req.} D1.4 (Parts 1–2); EW.32 (Part 2); P1, P5 (Parts 1–2); Def-H.38, FT.55(2), FT.60(1), FT.60(3), FT.60(5), FT.60(8), FT.80(5), FT.122(4), FT.123(1), FT.123(3), FT.123(4), FT.125(5), FT.126(4), FT.129(5), FT.130(5), FT.130(6), FT.142(1) (Part 3); Census FT.28(2), FT.35(3), FT.63(1), FT.63(6), FT.65(4) (row-wise); Def-BM.3(2), Def-BM.3(3), Def-BM.6(1), BM.17, BM.22(4), BM.23(5), BM.25, BM.26, BM.27, BM.31, BM.32, BM.33, BM.40(3), BM.41(1), BM.41(3), BM.42(1), BM.42(2), BM.42(3), BM.42(4), BM.68(6), BM.73(2), BM.73(3) (this part); exact $\mathbb Q(i)$ commutant rows, integer ladder arithmetic [math; machine: \texttt{checks/\allowbreak{}verify\_\allowbreak{}H\_\allowbreak{}chiral\_\allowbreak{}assignment\_\allowbreak{}opening.py} §B, §C].

\textbf{Pf.}

%dsp-level:0
\[
(1):\ \mathrm{Def\text{-}H.38}/\mathrm{FT.60}(1)\text{--}(3)\ \text{re-cited}\ (\text{frozen});
\]

%dsp-level:0
\[
\qquad \dim\mathrm{Sym}^1=2,\ \ j=\tfrac12,\ \ \text{weights}\ \pm\tfrac12\ [\text{machine: the}\ \tfrac{\sigma_3}2\ \text{spectrum}];
\]

%dsp-level:0
\[
\qquad \text{the Casimir ladder re-cited}\ [\mathrm{FT.60}(3)]
\]

%dsp-level:0
\[
(2):\ \mathrm{BM.73}(3)\text{'s occupancy with}\ (2)(\mathrm b)\text{'s image row};
\]

%dsp-level:0
\[
\qquad \hat\omega\ \text{maps the occupied ray onto the complement}\ [\mathrm{FT.122}(4);\ \text{machine}]
\]

%dsp-level:0
\[
(3):\ \text{Census FT.35}(3),\ \text{FT.63}(1)(\mathrm b),\ \text{FT.65}(4),\ \mathrm{FT.60}(5)\ \text{re-cited row-wise};
\]

%dsp-level:0
\[
\qquad \Theta_{\mathrm{int}}=2k\implies j=|k|\ [\text{machine: integers}]
\]

%dsp-level:0
\[
(4):\ \text{each row one frozen citation plus}\ \mathrm{BM.73}(2)\text{--}(3);
\]

%dsp-level:0
\[
\qquad \text{the opposite-ray conclusion: the charged member's arrow orientation}\ \mapsto\ \text{the left ray};
\]

%dsp-level:0
\[
\ \text{the neutral member's image orientation}\ \mapsto\ \text{the right ray}\ [\text{machine: the}\ 2\ \text{row census}]
\]

%dsp-level:0
\[
(5):\ \mathrm{BM.42}(1)\text{--}(4),\ \mathrm{BM.40}(3)\ \text{re-cited}\ (\text{frozen});
\]

%dsp-level:0
\[
\qquad [\sigma_i\otimes I_\omega,\ I\otimes\hat\chi]=0\ [\text{machine: tensor disjointness}]
\]

%dsp-level:0
\[
(6)\text{--}(7):\ \text{the bars re-cited};\qquad\text{the standing-rule scans}\ [\text{machine}]. \ \blacksquare
\]

\textbf{Phys.} Spin in this corpus was always a winding label — j equals half the winding, with the seed qubit as the k-equals-one anchor — and chirality was always an orientation grade; the assignment now says which ray the matter sits in, and the census resolves by ray without moving a single spin value. Every materialized line occupies the left ray; every image-orientation partner, including the antineutrino row the decay vertex already named, occupies the right; and the two emitted leptons of the beta vertex therefore sit in opposite rays, one flip-coupled because its core is charged, one flip-blind and parity-doubled because its core is neutral. The angular-momentum budget absorbs the new label silently, because the chirality grade acts on the orientation factor and commutes with every rotation generator on the seed. What the tier refuses is as exact as what it grants: no weight is glued to a ray — that glue is called helicity and helicity is external vocabulary — no exchange operator is invented, so spin-statistics stays open, and the baryon card keeps its rung-relative spin, because the winding datum the card would need is the same regime datum the joint-entropy tier named as missing.

\begin{theorem}[{\texorpdfstring{BM.75 — The weight–branch–core attachment: the unique bijection and the attached quark mass splitting}{BM.75 - The weight-branch-core attachment: the unique bijection and the attached quark mass splitting}}]\label{tag:BM-75}
\textbf{Thm.}

%dsp-level:0
\[
(1)\quad\textbf{the derived twos and the declared names:}
\]

%dsp-level:0
\[
\qquad \textbf{the weight alphabet:}\quad\operatorname{spec}Q|_{\mathcal H_{\mathsf v}}=\{0,+1\},
\]

%dsp-level:0
\[
\qquad \omega:\ \{0,+1\}\mapsto\{0,-1\}
\]

%dsp-level:0
\[
\qquad [\mathrm{EW.36};\ \mathrm{Def\text{-}BM.14}(7)]:
\]

%dsp-level:0
\[
\textbf{the branch index:}\quad\beta\in\{-,+\},
\]

%dsp-level:0
\[
\qquad D_S(\mathcal F_g,+)>D_S(\mathcal F_g,-)\ \ \forall g\in\{1,2,3\}
\]

%dsp-level:0
\[
\qquad [\mathrm{Def\text{-}BM.14}(2),(4);\ \mathrm{BM.60}(5):\ \text{uniform, exact}];
\]

%dsp-level:0
\[
\textbf{the name rows (declared, frozen):}
\]

%dsp-level:0
\[
\qquad \text{down-type}:=\text{the heavier branch}\ (\beta{=}+);\quad\text{up-type}:=\text{the lighter}\ (\beta{=}-)
\]

%dsp-level:0
\[
\qquad [\mathrm{Def\text{-}BM.14}(8);\ \text{A-BM.6}];
\]

%dsp-level:0
\[
\qquad \text{the class labels}\ \{\text{neutral class},\ \text{unit class}\}
\]

%dsp-level:0
\[
\qquad [\mathrm{Def\text{-}BM.3}(2),(3);\ \text{A-BM.1}];
\]

%dsp-level:0
\[
(2)\quad\textbf{the ladder direction (derived):}
\]

%dsp-level:0
\[
\qquad E_\pm=\tfrac12(\sigma_1\pm i\sigma_2)=\text{the matrix units in the}\ \sigma_3\ \text{eigenbasis};
\]

%dsp-level:0
\[
\qquad [\sigma_3,E_\pm]=\pm2E_\pm
\]

%dsp-level:0
\[
\qquad [\mathrm{BM.23}(1);\ \text{machine over}\ \mathbb Q(i)]:
\]

%dsp-level:0
\[
E_+:\ \text{weight}\ 0\ \longmapsto\ \text{weight}\ {+}1
\]

%dsp-level:0
\[
\qquad [\text{machine:}\ E_+(0,1)^{\mathsf T}=(1,0)^{\mathsf T}]:
\]

%dsp-level:0
\[
\qquad \textbf{the unique raising step inside}\ \operatorname{spec}\{0,+1\};
\]

%dsp-level:0
\[
\textbf{the vertex grammar:}
\]

%dsp-level:0
\[
\qquad \text{the}\ \beta^-\ \text{vertex maps the neutral class to the unit class with the forced}\ q=-1\ \text{carrier};
\]

%dsp-level:0
\[
\qquad 0=(+1)+(-1)+0
\]

%dsp-level:0
\[
\qquad [\mathrm{BM.23}(2)\text{--}(5);\ \mathrm{BM.41}(1);\ \text{machine: integers}];
\]

%dsp-level:0
\[
\textbf{the substitution row (the hypothesis on its face):}
\]

%dsp-level:0
\[
\qquad \textbf{hypothesis (H):}\ \text{the two classes differ by exactly one down}\leftrightarrow\text{up constituent substitution, the minimal}\ \Delta q=+1\ \text{step}
\]

%dsp-level:0
\[
\qquad [\mathrm{BM.59}(9);\ \mathrm{BM.23}(4)];
\]

%dsp-level:0
\[
(3)\quad\textbf{the attachment theorem (uniqueness):}
\]

%dsp-level:0
\[
\qquad \text{seek a bijection}\quad f:\ \{\text{weight}\ 0,\ \text{weight}\ {+}1\}\ \longrightarrow\ \{\beta{=}+,\ \beta{=}-\}\quad\text{consistent with:}
\]

%dsp-level:0
\[
(\mathrm i)\ \ \text{the ladder direction:}\ E_+\ \text{carries the substitution's source to its target}\ [(2)];
\]

%dsp-level:0
\[
\qquad (\mathrm{ii})\ \ \text{the name row: down-type}=\beta{=}+\ [(1)];
\]

%dsp-level:0
\[
(\mathrm{iii})\ \ \text{the transport order: the neutral-core constituent is the heavier one},
\]

%dsp-level:0
\[
\quad m_q(0)=2>1=m_q(\pm1)
\]

%dsp-level:0
\[
\qquad [\mathrm{BM.59}(5);\ \mathrm{Def\text{-}BM.13}(1)];
\]

%dsp-level:0
\[
(\mathrm{iv})\ \ \text{the two-currency sign coherence of the same substitution:}
\]

%dsp-level:0
\[
\qquad \Delta m_{\mathrm{transport}}=+\frac{\ln(4/3)}{\mathfrak h}>0
\]

%dsp-level:0
\[
\qquad [\mathrm{BM.59}(4)]
\]

%dsp-level:0
\[
\qquad \text{and}
\]

%dsp-level:0
\[
\qquad \Delta_g=\hat\mu'^2\bigl(D_S(+)-D_S(-)\bigr)>0
\]

%dsp-level:0
\[
\qquad [\mathrm{BM.59}(9);\ \mathrm{Def\text{-}BM.14}(5)]:
\]

%dsp-level:0
\[
\text{the substitution's source is the heavier member in both currencies};
\]

%dsp-level:0
\[
\boxed{\ \textbf{exactly one bijection survives}\ (\mathrm i)\text{--}(\mathrm{iv}):\qquad\text{weight}\ 0\ \longleftrightarrow\ q_{\mathrm{core}}=0\ \longleftrightarrow\ \text{down-type}\ \longleftrightarrow\ \beta{=}+\ (\text{heavier});\qquad\text{weight}\ {+}1\ \longleftrightarrow\ q_{\mathrm{core}}={+}1\ \longleftrightarrow\ \text{up-type}\ \longleftrightarrow\ \beta{=}-\ (\text{lighter})\ }
\]

%dsp-level:0
\[
\qquad [\text{machine: exhaustive over the}\ 2!\ \text{bijections at the frozen ladder direction}];
\]

%dsp-level:0
\[
\textbf{the A-BM.6 bar is discharged by derivation:}
\]

%dsp-level:0
\[
\qquad \text{the attachment of the}\ Q\text{-weight index to the branch index}\ —\ \text{recorded there as}\ \neg\text{declared}\ —\ \text{is derived from the declared name rows plus}\ (2)\text{--}(\mathrm{iv});
\]

%dsp-level:0
\[
\qquad \neg\text{a new declaration issues here};
\]

%dsp-level:0
\[
(4)\quad\textbf{the attached quark mass splitting (exact per currency):}
\]

%dsp-level:0
\[
\qquad \textbf{the defect/pole currency:}\quad\Delta_g=m^2_{\mathrm{down}}-m^2_{\mathrm{up}}=\hat\mu'^2\bigl(D_S(\mathcal F_g,+)-D_S(\mathcal F_g,-)\bigr):
\]

%dsp-level:0
\[
\Delta_1=2\,\hat\mu'^2\in\mathbb Q\,\hat\mu'^2;
\]

%dsp-level:0
\[
\qquad \Delta_2=(\sqrt2-1)\,\hat\mu'^2\in\mathbb Q(\sqrt2)\,\hat\mu'^2;
\]

%dsp-level:0
\[
\qquad \Delta_3=\bigl(2\cos\tfrac\pi{10}-\sqrt3\bigr)\,\hat\mu'^2\in\mathbb Q\bigl(\sqrt3,\cos\tfrac\pi{10}\bigr)\,\hat\mu'^2;
\]

%dsp-level:0
\[
\qquad \Delta_g>0\ \ \forall g
\]

%dsp-level:0
\[
\qquad [\mathrm{Def\text{-}BM.14}(5);\ \text{machine: the certificates}\ 4>2,\ 2>1,\ 3<5];
\]

%dsp-level:0
\[
\textbf{the informational/transport-register currency:}
\]

%dsp-level:0
\[
\qquad \text{per constituent}\ \frac{\ln2}{\mathfrak h}>0;
\]

%dsp-level:0
\[
\qquad \text{per composite}\ \frac{\ln(4/3)}{\mathfrak h}>0
\]

%dsp-level:0
\[
\qquad [\mathrm{BM.59}(4),(5)];
\]

%dsp-level:0
\[
\textbf{the ensemble route-(ii) currency:}
\]

%dsp-level:0
\[
\qquad \frac{T(0,S)}{T(+1,S)}=\frac43>1:
\]

%dsp-level:0
\[
\qquad \text{the neutral class is heavier}
\]

%dsp-level:0
\[
\qquad [\mathrm{BM.28}(4)(\mathrm{ii});\ \mathrm{BM.30}(2)];
\]

%dsp-level:0
\[
\textbf{the opposing row:}
\]

%dsp-level:0
\[
\qquad \text{EM makes the charged member heavier by}\ c\,\varepsilon_{\mathrm{EM}}q^2m_{\mathrm{hid}}>0
\]

%dsp-level:0
\[
\qquad [\mathrm{BM.46}(4)]:
\]

%dsp-level:0
\[
\qquad \text{the unique sign-opposite correction, dominated on the sub-critical branch}\ \varepsilon_{\mathrm{EM}}<\varepsilon^{\mathrm{crit}}(c)
\]

%dsp-level:0
\[
\qquad [\mathrm{BM.62}(6),(9);\ \mathrm{BM.65}(3)\text{--}(5)];
\]

%dsp-level:0
\[
(5)\quad\textbf{the sign-coherence theorem:}
\]

%dsp-level:0
\[
\qquad \boxed{\ \text{in every derived currency in which the doublet splitting is formed, the down-type constituent is the heavier one}\ }
\]

%dsp-level:0
\[
[\text{machine: the sign census of}\ (4)\text{'s four rows}\ —\ \text{three concordant positives and one opposed correction with its exact domination branch}];
\]

%dsp-level:0
\[
\textbf{the currencies remain incommensurate:}
\]

%dsp-level:0
\[
\qquad \neg\exists\ \text{a derived ratio between}\ \Delta_g\ \text{and}\ \ln(4/3)/\mathfrak h\ —\ \text{the relating equation stays open}
\]

%dsp-level:0
\[
\qquad [\mathrm{BM.35}(5);\ \mathrm{BM.68}(7);\ \mathrm{BM.69}(4)];
\]

%dsp-level:0
\[
(6)\quad\textbf{what remains open:}
\]

%dsp-level:0
\[
\qquad \text{the full flavour composition}\ —\ \text{which three of the six states a class occupies}\ —\ \text{stays}\ \neg\text{derived:}\ \neg\text{a fractional constituent alphabet},
\]

%dsp-level:0
\[
\ \text{the end signs carry no position datum}
\]

%dsp-level:0
\[
\qquad [\mathrm{BM.28}(7);\ \mathrm{Def\text{-}BM.8}(5);\ \mathrm{Def\text{-}BM.13}(9)]:
\]

%dsp-level:0
\[
\textbf{the attachment theorem consumes only the single-substitution row,}\ \neg\textbf{the full triple};
\]

%dsp-level:0
\[
\qquad \text{the generation map is an order,}\ \neg\text{a number}\ [\mathrm{Def\text{-}BM.14}(6);\ \mathrm{BM.60}(7)];
\]

%dsp-level:0
\[
\qquad \text{the mixing matrix is}\ \neg\text{derived}\ [\mathrm{FT.145}(3)];
\]

%dsp-level:0
\[
(7)\quad\textbf{status:}\qquad\textbf{derived:}\ (1)\text{--}(5)\ \ [(3)\ \text{conditional on hypothesis (H), on its face}\ [\text{A §0.2a}]];
\]

%dsp-level:0
\[
\neg\text{a new declaration},
\]

%dsp-level:0
\[
\quad\neg\text{a P4 addition},
\]

%dsp-level:0
\[
\quad\neg\text{a measured value}\ [\mathrm{P5}],
\]

%dsp-level:0
\[
\quad\neg\text{a Yukawa-type object}\ [\mathrm{EW.15};\ \mathrm{FT.145}(3)];
\]

%dsp-level:0
\[
\qquad \text{§Cal-BM stays empty}.
\]

\textbf{Req.} D3.4 (Parts 1–2); EW.15, EW.36 (Part 2); P5 (Parts 1–2); Def-H.37(2), FT.60(5), FT.145(3) (Part 3); Def-BM.3(2), Def-BM.3(3), Def-BM.8(5), Def-BM.13(1), Def-BM.13(9), Def-BM.14(2), Def-BM.14(4), Def-BM.14(5), Def-BM.14(6), Def-BM.14(7), Def-BM.14(8), BM.23(1), BM.23(2), BM.23(4), BM.23(5), BM.28(4), BM.28(7), BM.30(2), BM.35(5), BM.41(1), BM.46(4), BM.59(4), BM.59(5), BM.59(9), BM.60(5), BM.60(7), BM.62(6), BM.62(9), BM.65(3), BM.65(5), BM.68(7), BM.69(4) (this part); exact $\mathbb Q(i)$ matrix-unit algebra, exhaustive bijection censuses, exact radical certificates, prime-log-ring sign tests [math; machine: \texttt{checks/\allowbreak{}verify\_\allowbreak{}H\_\allowbreak{}chiral\_\allowbreak{}assignment\_\allowbreak{}opening.py} §B, §D].

\textbf{Pf.}

%dsp-level:0
\[
(1):\ \mathrm{EW.36},\ \mathrm{Def\text{-}BM.14}(2),(4),(7),(8),\ \mathrm{Def\text{-}BM.3}(2),(3)\ \text{re-cited}\ (\text{frozen})
\]

%dsp-level:0
\[
(2):\ E_+=\left(\begin{smallmatrix}0&1\\0&0\end{smallmatrix}\right)\ \text{in the weight basis with}\ Q=\operatorname{diag}(1,0):\quad E_+(0,1)^{\mathsf T}=(1,0)^{\mathsf T}\ [\text{machine}];\qquad[\sigma_3,E_+]=2E_+\ [\text{machine over}\ \mathbb Q(i)];\qquad0=(+1)+(-1)+0\ [\text{machine: integers}]
\]

%dsp-level:0
\[
(3):\ \text{the census: bijection}\ f_A:\ 0\mapsto+,\ {+}1\mapsto-\ \ \text{and}\ f_B:\ 0\mapsto-,\ {+}1\mapsto+;
\]

%dsp-level:0
\[
\qquad f_A\ \text{satisfies}\ (\mathrm i)\text{--}(\mathrm{iv}):\ \text{the}\ E_+\text{-source is weight}\ 0
\]

%dsp-level:0
\[
= \text{down-type}
\]

%dsp-level:0
\[
= \beta{=}+
\]

%dsp-level:0
\[
= \text{the heavier branch}
\]

%dsp-level:0
\[
= \text{the heavier core}\ [\text{machine: the four condition evaluations}];
\]

%dsp-level:0
\[
f_B\ \text{fails}\ (\mathrm i)\ \text{and}\ (\mathrm{iii})\wedge(\mathrm{iv})\ \text{jointly: it makes the substitution's source}\ (\text{weight}\ 0,\ \text{core}\ 0,\ \text{transport-heavier})\ \text{the lighter branch}\ \beta{=}-,
\]

%dsp-level:0
\[
\ \text{so the same substitution that raises}\ \Delta m_{\mathrm{transport}}
\]

%dsp-level:0
\[
= +\ln(4/3)/\mathfrak h\ \text{would lower the defect read}\ (-\Delta_g)\ \text{against}\ \mathrm{BM.59}(9)\text{'s derived}\ +\Delta_g\ [\text{machine: the sign contradiction}];
\]

%dsp-level:0
\[
\qquad \text{the ladder direction is frozen by the grammar}\ [(2)];
\]

%dsp-level:0
\[
\qquad \text{the two bijections:}\ 2\ \text{cases, one survivor}\ [\text{machine}]
\]

%dsp-level:0
\[
(4):\ \Delta_g\ \text{re-formed in the three exact fields with the certificates}\ 4>2;
\]

%dsp-level:0
\[
\ 2>1;
\]

%dsp-level:0
\[
\ 2\cos\tfrac\pi{10}>\sqrt3\iff\cos^2\tfrac\pi{10}>\tfrac34\iff\tfrac{5+\sqrt5}8>\tfrac34
\]

%dsp-level:0
\[
\iff 5+\sqrt5>6\iff\sqrt5>1
\]

%dsp-level:0
\[
\iff 5>1\ [\text{machine: exact in}\ \mathbb Q(\sqrt5)];
\]

%dsp-level:0
\[
\qquad \text{the currency rows re-cited}\ [\mathrm{BM.59}(4),(5);\ \mathrm{BM.28}(4)(\mathrm{ii});\ \mathrm{BM.46}(4)]
\]

%dsp-level:0
\[
(5):\ \text{the sign census of the four rows}\ [\text{machine}];
\]

%dsp-level:0
\[
\qquad \text{the no-ratio row re-cited}\ [\mathrm{BM.68}(7)(\mathrm d),(\mathrm e);\ \mathrm{BM.69}(4)]
\]

%dsp-level:0
\[
(6)\text{--}(7):\ \text{the open rows re-cited};\qquad\text{the standing-rule scans}\ [\text{machine}]. \ \blacksquare
\]

\textbf{Phys.} Two derived twos have always sat side by side — the seed's weight doublet and the defect floor's branch pair — with their attachment recorded as undeclared. The chiral assignment supplies the missing link, and the link turns out to be forced: the ladder raises exactly one weight, the beta vertex raises exactly one core, the heavier branch is named down-type, and the heavier core is the neutral one by an exact transport count — of the two possible gluings only one keeps the sign of the splitting the same in the transport currency and in the defect currency, and both signs are frozen derived rows. So the attachment is a theorem given the substitution hypothesis and the declared names, and the quark mass splitting is now attached rather than floating: per generation it is two, root-two-minus-one, and two-cosines-of-eighteen-degrees-minus-root-three, times the undeclared defect scale, positive in all three, with the informational currencies agreeing in sign and the electromagnetic correction opposing exactly where the sub-critical branch dominates it. What stays open stays open loudly: which three of the six states compose a given baryon is still unreadable, because the attachment consumes one substitution step and not a whole triple, and no ratio between the currencies is derived — the sign agrees, the scale does not, and that is application-layer content.

\end{theorem}

\begin{theorem}[{\texorpdfstring{BM.76 — The chiral vertex dynamics: the nilpotent step operator, the flip channel, and the rate rows}{BM.76 - The chiral vertex dynamics: the nilpotent step operator, the flip channel, and the rate rows}}]\label{tag:BM-76}
\textbf{Thm.}

%dsp-level:0
\[
(1)\quad\textbf{the vertex step operator (an in-theorem abbreviation):}
\]

%dsp-level:0
\[
\qquad \text{on the extended carrier}\ \mathcal H_{\mathrm{ext}}=\mathcal H_{\mathrm{int}}\otimes\mathcal H_\omega\cong\mathbb C^4
\]

%dsp-level:0
\[
\qquad [\mathrm{Def\text{-}H.34}]:
\]

%dsp-level:0
\[
\boxed{\ \widehat V\ :=\ E_+\otimes P_L\ }
\]

%dsp-level:0
\[
\qquad [E_+=\tfrac12(\sigma_1+i\sigma_2)\ [\mathrm{BM.23}(1)];\quad P_L=\tfrac12(\mathrm{id}+\hat\chi)\ [\mathrm{Def\text{-}H.33}];\quad\text{exact over}\ \mathbb Q(i)\ \text{in the ordered basis of}\ \mathrm{Def\text{-}H.34}];
\]

%dsp-level:0
\[
(2)\quad\textbf{the nilpotency is the alphabet overflow's operator form:}
\]

%dsp-level:0
\[
\qquad E_+^2=0\implies\boxed{\ \widehat V^{\,2}=0\ }
\]

%dsp-level:0
\[
\qquad [\text{machine: entrywise over}\ \mathbb Q(i)]:
\]

%dsp-level:0
\[
\text{the ladder has one step inside the doublet}\ —\ \mathrm{BM.23}(3)\text{'s ternary overflow}\ \bigl(\{-1,0,+1\}+1=\{0,+1,+2\}\ \text{leaves the alphabet}\bigr)\ \textbf{as an exact operator identity};
\]

%dsp-level:0
\[
\qquad \neg\exists\ \text{an iterated charged step inside one doublet};
\]

%dsp-level:0
\[
(3)\quad\textbf{the rank, image and kernel rows:}
\]

%dsp-level:0
\[
\qquad \operatorname{rk}\widehat V=1=\operatorname{rk}E_+\cdot\operatorname{rk}P_L;
\]

%dsp-level:0
\[
\qquad \operatorname{im}\widehat V=\operatorname{span}\bigl\{|\text{weight}\,{+}1\rangle\otimes|L\rangle\bigr\};
\]

%dsp-level:0
\[
\qquad \dim\ker\widehat V=3
\]

%dsp-level:0
\[
\qquad [\text{machine: exact rank and image over}\ \mathbb Q(i)]:
\]

%dsp-level:0
\[
\textbf{the target is the up-type weight on the occupied ray}\qquad[\mathrm{BM.75}(3)\ \text{attachment}];
\]

%dsp-level:0
\[
\widehat V\,(I\otimes P_R)=0,
\]

%dsp-level:0
\[
\qquad (I\otimes P_L)\,\widehat V=\widehat V
\]

%dsp-level:0
\[
\qquad [\text{machine}]:
\]

%dsp-level:0
\[
\qquad \textbf{the step fires from the left ray alone}\ —\ \text{the chiral selectivity of the transition, derived given A-BM.7};
\]

%dsp-level:0
\[
(4)\quad\textbf{the Cartan row:}
\]

%dsp-level:0
\[
\qquad \bigl[W_3\otimes I_\omega,\ \widehat V\bigr]=\widehat V
\]

%dsp-level:0
\[
\qquad [\ [\sigma_3,E_+]=2E_+,\ W_3=\tfrac12\sigma_3;\ \text{machine}]:
\]

%dsp-level:0
\[
\qquad \textbf{the step operator carries weight}\ {+}1\ \text{under the seed Cartan}\ —\ \text{the operator form of the composite raise}\ \Delta q
\]

%dsp-level:0
\[
= +1\ \text{with the compensating emission}\ [\mathrm{BM.23}(2),(4)];
\]

%dsp-level:0
\[
(5)\quad\textbf{the}\ \omega\textbf{-covariant image vertex:}\qquad(I\otimes\hat\omega)\,\widehat V\,(I\otimes\hat\omega)^{-1}=E_+\otimes P_R\qquad[\mathrm{FT.122}(4);\ \text{machine}]:
\]

%dsp-level:0
\[
\textbf{the image step fires from the right ray}\ —\ \text{the CPT-partner vertex}\ [\mathrm{BM.25};\ \mathrm{BM.27};\ \text{A-BM.5's image-orientation row is ray-resolved}];
\]

%dsp-level:0
\[
\textbf{the flip conjugation:}
\]

%dsp-level:0
\[
\qquad (I\otimes\hat\kappa)\,\widehat V\,(I\otimes\hat\kappa)=E_+\otimes P_R
\]

%dsp-level:0
\[
\qquad [\hat\kappa P_L\hat\kappa=P_R:\ \{\hat\chi,\hat\kappa\}=0\ [\mathrm{Def\text{-}H.34}];\ \text{machine}]:
\]

%dsp-level:0
\[
\qquad \text{the partner involution}\ —\ \text{the EM flip's carrier}\ [\mathrm{FT.129}(2),(5)]\ —\ \text{moves the vertex to the opposite ray};
\]

%dsp-level:0
\[
(6)\quad\textbf{the flip-channel rows (frozen, re-cited):}
\]

%dsp-level:0
\[
\qquad \text{rung one:}\ F_{LR}=F_{RL}=0\ \text{— the sectors force-decoupled}\ [\mathrm{FT.128}(2);\ \mathrm{FT.144}(6)];
\]

%dsp-level:0
\[
\qquad \text{rung two:}\ F_{LR}
\]

%dsp-level:0
\[
= q\,
\]

%dsp-level:0
\[
I_2\ \text{under the}\ \mathrm{Def\text{-}H.43}\ \text{promotion},
\]

%dsp-level:0
\[
\ \text{splitting}\ 2|q|\in\{0,2\}\ [\mathrm{FT.129}(5);\ \mathrm{FT.130}(5);\ \mathrm{BM.22}(2),(4)]:
\]

%dsp-level:0
\[
\text{the neutral core keeps parity doubling; the charged core flip-couples both rays}
\]

%dsp-level:0
\[
\qquad [\mathrm{FT.130}(6);\ \text{machine re-cited}];
\]

%dsp-level:0
\[
\qquad \text{rung three adds no carrier flip}\ [\mathrm{FT.117};\ \mathrm{FT.144}(6)];
\]

%dsp-level:0
\[
(7)\quad\textbf{the rate and arrow rows (frozen, re-cited):}
\]

%dsp-level:0
\[
\qquad \Gamma_{\mathrm{tr}}(\varepsilon_q)=\frac{x_+x_-}{1-x_+x_-}\in\Bigl[0,\tfrac13\Bigr],
\]

%dsp-level:0
\[
\qquad x_+x_-=4\varepsilon_q(1-4\varepsilon_q)=\tfrac14-16\Bigl(\varepsilon_q-\tfrac18\Bigr)^2\le\tfrac14
\]

%dsp-level:0
\[
\qquad [\mathrm{BM.44}(1),(2);\ \mathrm{BM.19}(2);\ \text{machine: the SOS parabola}]:
\]

%dsp-level:0
\[
\varepsilon_q=0\implies\Gamma_{\mathrm{tr}}=0:
\]

%dsp-level:0
\[
\qquad \textbf{no flicker, no decay}
\]

%dsp-level:0
\[
\qquad [\mathrm{BM.23}(8)];
\]

%dsp-level:0
\[
\qquad \tau_{\mathrm{life}}\ge3\ \ \text{ticks}
\]

%dsp-level:0
\[
\qquad [\mathrm{BM.44}(3)];
\]

%dsp-level:0
\[
\qquad \text{forward-only},\ \neg\exists\,\tau^{-1}
\]

%dsp-level:0
\[
\qquad [\mathrm{A6};\ \mathrm{BM.23}(9);\ \mathrm{FT.144}(2)]:
\]

%dsp-level:0
\[
\textbf{the vertex fires per tick as a weight,}\ \neg\textbf{as an external rate}
\]

%dsp-level:0
\[
\qquad [\mathrm{BM.44}(4),(5):\ \neg\text{external seconds}];
\]

%dsp-level:0
\[
\qquad \textbf{the opening of}\ \mathrm{BM.77}\ \textbf{is kinematic; its rate stays the flicker's own};
\]

%dsp-level:0
\[
(8)\quad\textbf{scope:}
\]

%dsp-level:0
\[
\qquad \text{what is constructed here is the vertex-level dynamics of the}\ Q\ \text{read}\ —\ \text{the ladder step, the chiral selection, the flip channel, the rate};
\]

%dsp-level:0
\[
\qquad \textbf{the field-level dynamics of}\ Q\ \text{(gauge-field dynamics, radiation, back-reaction) remains}\ \neg\textbf{constructed}
\]

%dsp-level:0
\[
\qquad [\mathrm{BM.6}(5);\ \mathrm{Def\text{-}BM.12}(7);\ \mathrm{BM.50}(3)];
\]

%dsp-level:0
\[
\widehat V\ \text{is}\ \neg\text{a new force: A8's completeness row stands}\ —\ \text{no sub-block restriction is an A8 force}
\]

%dsp-level:0
\[
\qquad [\mathrm{A8};\ \mathrm{BM.52}(1)];
\]

%dsp-level:0
\[
\qquad \widehat V\ \text{is an in-theorem abbreviation of}\ \mathrm{BM.23}(1)\text{'s ladder and}\ \mathrm{Def\text{-}H.33}\text{'s projector};
\]

%dsp-level:0
\[
(9)\quad\textbf{status:}
\]

%dsp-level:0
\[
\qquad \textbf{derived:}\ (1)\text{--}(7)\ \ [\text{the occupancy rows given A-BM.7}];
\]

%dsp-level:0
\[
\qquad \neg\text{a P4 addition},\quad\neg\text{a measured lifetime}\ [\mathrm{P5}],\quad\neg\text{a new dimensionful scale}\ [\mathrm{P1}];
\]

%dsp-level:0
\[
\text{the conditionalities}\ (\mathrm{Def\text{-}H.43}\ \text{promotion}\ [\mathrm{BM.22}(5)],\ \mathrm{BM.8}(2)\ \text{stationarity for the chart exponential}\ [\mathrm{BM.44}(3)])\ \text{are inherited at the point of use}\ [\text{A §0.2a}].
\]

\textbf{Req.} A6, A8 (Part 1); EW.32, EW.36 (Part 2); P1, P5 (Parts 1–2); Def-H.33, Def-H.34, Def-H.43, FT.117, FT.122(4), FT.128(2), FT.129(2), FT.129(5), FT.130(5), FT.130(6), FT.144(2), FT.144(6) (Part 3); Def-BM.12(7), BM.6(5), BM.19(2), BM.22(2), BM.22(4), BM.22(5), BM.23(1), BM.23(2), BM.23(4), BM.23(8), BM.23(9), BM.25, BM.27, BM.44(1), BM.44(2), BM.44(3), BM.44(4), BM.44(5), BM.50(3), BM.52(1), BM.73(2), BM.73(3), BM.75(3) (this part); exact $\mathbb Q(i)$ Kronecker algebra, rank/image/kernel computation, the SOS parabola completion [math; machine: \texttt{checks/\allowbreak{}verify\_\allowbreak{}H\_\allowbreak{}chiral\_\allowbreak{}assignment\_\allowbreak{}opening.py} §C].

\textbf{Pf.}

%dsp-level:0
\[
(1)\text{--}(2):\ E_+=\left(\begin{smallmatrix}0&1\\0&0\end{smallmatrix}\right),\ P_L=\left(\begin{smallmatrix}\tfrac12&\tfrac i2\\-\tfrac i2&\tfrac12\end{smallmatrix}\right)\ \text{in the ordered bases}\ [\mathrm{Def\text{-}H.33}(5)\ \text{form};\ \mathrm{Def\text{-}H.34}]:\qquad\widehat V\ \text{built as the Kronecker product, all}\ 16\ \text{entries exact over}\ \mathbb Q(i)\ [\text{machine}];\qquad E_+^2=0\implies\widehat V^2=(E_+^2)\otimes(P_L^2)=0\ [\text{machine: entrywise}]
\]

%dsp-level:0
\[
(3):\ \operatorname{rk}(A\otimes B)=\operatorname{rk}A\cdot\operatorname{rk}B\ [\text{math}];
\]

%dsp-level:0
\[
\qquad \operatorname{rk}E_+=1,\ \operatorname{rk}P_L=1\ [\text{machine: exact elimination}];
\]

%dsp-level:0
\[
\qquad \operatorname{im}\widehat V\ \text{spanned by the image of the basis}\ \{\,|\text{w}0\rangle\otimes|L\rangle\,\}\mapsto|\text{w}{+}1\rangle\otimes|L\rangle\ [\text{machine}];
\]

%dsp-level:0
\[
\qquad \widehat V(I\otimes P_R)=E_+\otimes P_LP_R=0\ [\mathrm{FT.122}(2);\ \text{machine}]
\]

%dsp-level:0
\[
(4):\ [\sigma_3,E_+]
\]

%dsp-level:0
\[
= \sigma_3E_+-E_+\sigma_3
\]

%dsp-level:0
\[
= 2E_+\ [\text{machine}]
\]

%dsp-level:0
\[
\implies [W_3\otimes I,\ E_+\otimes P_L]
\]

%dsp-level:0
\[
= [W_3,E_+]\otimes P_L
\]

%dsp-level:0
\[
= E_+\otimes P_L\ [\text{machine: the}\ 4\times4\ \text{commutator entrywise}]
\]

%dsp-level:0
\[
(5):\ \hat\omega P_L\hat\omega=P_R\ [\mathrm{FT.122}(4);\ \text{machine}];
\]

%dsp-level:0
\[
\qquad \hat\kappa P_L\hat\kappa=\tfrac12(I+\hat\kappa\hat\chi\hat\kappa)=\tfrac12(I-\hat\chi)=P_R\ [\{\hat\chi,\hat\kappa\}=0;\ \mathrm{Def\text{-}H.34};\ \text{machine}]
\]

%dsp-level:0
\[
(6)\text{--}(7):\ \text{the frozen rows re-cited}\ [\mathrm{FT.128}(2);\ \mathrm{FT.129}(5);\ \mathrm{FT.130}(5),(6);\ \mathrm{FT.117};\ \mathrm{FT.144}(6);\ \mathrm{BM.44}(1)\text{--}(5);\ \mathrm{BM.23}(8),(9)];
\]

%dsp-level:0
\[
\qquad x_+x_-=\tfrac14-16(\varepsilon_q-\tfrac18)^2\ [\text{machine: the SOS completion}];
\]

%dsp-level:0
\[
\qquad \Gamma_{\mathrm{tr}}(\tfrac18)=\tfrac{1/4}{3/4}=\tfrac13\ [\text{machine: exact Fractions}]
\]

%dsp-level:0
\[
(8)\text{--}(9):\ \text{the scope record};\qquad\text{the standing-rule scans}\ [\text{machine}]. \ \blacksquare
\]

\textbf{Phys.} The charged-current vertex becomes a single exact operator: the weak ladder step tensored with the occupied ray's projector, four by four over the rationals with i adjoined. Its algebra is the physics. Squared, it vanishes — the same nilpotency that makes the ladder's second step leave the ternary alphabet, now an operator identity, so a constituent takes the charged step once and never twice. Its rank is one, its image is the up-type weight on the left ray, and it annihilates anything sitting in the right ray: given the assignment, the transition fires from matter's own ray alone, which is the corpus's exact counterpart of a maximally chiral charged current — derived, not postulated. Conjugating by orientation reversal or by the partner involution moves it to the right ray, giving the image vertex that the antiproton and antineutrino rows occupy, and the Cartan commutator records the charge bookkeeping: the step carries weight plus one, the composite raises by one, and the emitted carrier takes minus one away. The rates were already exact and stay frozen: a weight per tick bounded by a third, three ticks minimum, nothing at the sharp crystallization threshold, and no seconds anywhere — the opening this tier derives elsewhere is kinematic, and the dynamics keeps the flicker's own clock.

\end{theorem}

\begin{theorem}[{\texorpdfstring{BM.77 — The read specification and the allowed branch: the formed Q-value, the exact trichotomy, and the rescope rows}{BM.77 - The read specification and the allowed branch: the formed Q-value, the exact trichotomy, and the rescope rows}}]\label{tag:BM-77}
\textbf{Thm.}

%dsp-level:0
\[
(1)\quad\textbf{the specification (the declaration, logged as A-BM.8 before use):}
\]

%dsp-level:0
\[
\qquad \textbf{the}\ \beta\textbf{-vertex Q-value is read:}
\]

%dsp-level:0
\[
\qquad \text{on the baryon classes}\ —\ \text{route}\ (\mathrm{ii}),\ \text{the unnormalized total deficit}\quad T(q,S)=W(q)\,\overline{\Delta S}(S)
\]

%dsp-level:0
\[
\qquad [\mathrm{BM.28}(4)(\mathrm{ii});\ \mathrm{BM.30}(1)\ \text{route census}];
\]

%dsp-level:0
\[
\text{on the emitted lines}\ —\ \text{the transport-species family form}\quad m^2(\text{line}\ i)=\hat\mu^2\ell_c^2\,G\bigl(c(i)\bigr),
\]

%dsp-level:0
\[
\qquad G(c(i))=A_0-B_0\cos\varphi_i
\]

%dsp-level:0
\[
\qquad [\mathrm{Def\text{-}BM.11}(5);\ \mathrm{BM.39}(1)];
\]

%dsp-level:0
\[
\textbf{the expression is}\ \mathrm{BM.41}(4)(\mathrm d)\text{'s own frozen}\ Q_{\mathrm{ii}}\ —\ \text{formed there, priced there, adopted here};
\]

%dsp-level:0
\[
\qquad \textbf{the scope is the vertex kinematics alone:}\ \ \text{route}\ (\mathrm{ii})\ \text{is}\ \neg\text{elevated to a global mass definition}\ [\mathrm{BM.30}(5);\ \mathrm{BM.37}(6)\ \text{stand}];
\]

%dsp-level:0
\[
(2)\quad\textbf{the formed Q-value:}
\]

%dsp-level:0
\[
\qquad Q_\beta^{\mathrm{spec}}:=m_n^{(\mathrm{ii})}-m_p^{(\mathrm{ii})}-m_\ell-m_\nu;
\]

%dsp-level:0
\[
\qquad m_n^{(\mathrm{ii})}=\frac{T(0,S)}{\mathfrak h},\quad m_p^{(\mathrm{ii})}=\frac{T(+1,S)}{\mathfrak h}
\]

%dsp-level:0
\[
\qquad [\text{the route-(ii) mass read};\ \mathrm{BM.30}(1)]:
\]

%dsp-level:0
\[
\boxed{\ Q_\beta^{\mathrm{spec}}\,\mathfrak h\ =\ \Delta T(S)\ -\ \rho_\ell\sqrt{G\bigl(c(0)\bigr)}\ -\ \rho_\nu\sqrt{G\bigl(c(i_\nu)\bigr)},\qquad\Delta T(S):=T(0,S)-T(+1,S)=\bigl(W(0)-W(+1)\bigr)\overline{\Delta S}(S)=5\,\overline{\Delta S}(S)\ }
\]

%dsp-level:0
\[
[\mathrm{BM.28}(4)(\mathrm{ii})\ \text{and its Pf row}\ T(q,S)=W(q)\overline{\Delta S}(S);\ \ \text{the electron row}\ i=0\ [\mathrm{Def\text{-}BM.6}(1)];\ \ \text{machine}];
\]

%dsp-level:0
\[
\textbf{the dimension row:}
\]

%dsp-level:0
\[
\qquad \rho_\ell:=\mathfrak h\,\hat\mu_\ell\,\ell_c^{(\ell)},\quad\rho_\nu:=\mathfrak h\,\hat\mu_\nu\,\ell_c^{(\nu)}\ \in\ \mathfrak q^0
\]

%dsp-level:0
\[
\qquad [\mathrm{CT2.1};\ \mathrm{D2.9};\ \mathrm{Def\text{-}H.37}(1):\ [\hat\mu^2]=\mathfrak q^2;\ \text{machine: the exponent arithmetic}]:
\]

%dsp-level:0
\[
\text{the scales}\ (\hat\mu_\ell,\hat\mu_\nu,\ell_c,g,g')\ \text{remain undeclared P4/Cal-F data}\ [\mathrm{FT.143}]\ —\ \rho_\ell,
\]

%dsp-level:0
\[
\rho_\nu\ \text{are restatements,}\ \neg\text{new members}
\]

%dsp-level:0
\[
\qquad [\mathrm{P4};\ \mathrm{P1};\ \text{Cal-F logged}];
\]

%dsp-level:0
\[
\textbf{the exact baryon rows:}
\]

%dsp-level:0
\[
\qquad \Delta T(S)=5\,m^2_B(\cdot,S)
\]

%dsp-level:0
\[
\qquad [\overline{\Delta S}(S)=m^2_B(\cdot,S):\ \mathrm{FT.107}(3);\ \mathrm{BM.28}(4)(\mathrm i),(\mathrm{ii});\ \text{the frozen sector expectations}\ \mathrm{BM.20}(4)/\mathrm{BM.21}(2)\ \text{re-cited at}\ \mathrm{BM.62}(2)]:
\]

%dsp-level:0
\[
\boxed{\ \Delta T\ =\ \Bigl(\frac{4219785}{98956},\ \ \frac{9608310}{393679},\ \ \frac{455}{23},\ \ \frac{68565}{3088}\Bigr)\ \ \text{at}\ S=(0,\,1,\,2,\,3)\ }
\]

%dsp-level:0
\[
\qquad [\text{machine: exact rational arithmetic}];
\]

%dsp-level:0
\[
\Delta T_{\min}=\frac{455}{23},
\]

%dsp-level:0
\[
\qquad \Delta T_{\max}=\frac{4219785}{98956},
\]

%dsp-level:0
\[
\qquad 0<\Delta T_{\min}\le\Delta T(S)\le\Delta T_{\max}
\]

%dsp-level:0
\[
\qquad [\text{machine: exact comparisons}];
\]

%dsp-level:0
\[
\qquad \textbf{the}\ \nu\text{-line index}\ i_\nu\ \text{is unresolved}\ \ —\ \textbf{every row below holds uniformly in}\ i_\nu\in\{0,1,2\}\ [(4)\text{'s sandwich is line-independent}];
\]

%dsp-level:0
\[
(3)\quad\textbf{the formedness row:}
\]

%dsp-level:0
\[
\qquad \text{the expression of}\ (2)\ \text{is}\ \mathrm{BM.41}(4)(\mathrm d)\text{'s frozen}\ Q_{\mathrm{ii}}\ \text{with the lepton scales carried symbolically}\ —\ \neg\text{a cross-currency assumption is made}:
\]

%dsp-level:0
\[
\text{no informational value is equated to a geometric one};
\]

%dsp-level:0
\[
\qquad \text{the mass-form bridge is application-layer content, untouched here}
\]

%dsp-level:0
\[
\qquad [\mathrm{BM.35}(5);\ \mathrm{BM.68}(8);\ \mathrm{BM.69}(5);\ \text{pointer}];
\]

%dsp-level:0
\[
\qquad \text{the EM and transport rest-slot refinements are}\ \neg\text{summands of this read}\ —\ \text{disjoint alphabets, no double count}
\]

%dsp-level:0
\[
\qquad [\mathrm{BM.37}(2);\ \mathrm{BM.62}(11);\ \mathrm{BM.66}(6)];
\]

%dsp-level:0
\[
(4)\quad\textbf{the universal sandwich:}
\]

%dsp-level:0
\[
\qquad A_0-B_0=\frac38\Bigl(g-\frac{g'}3\Bigr)^2+\frac{g'^2}{48}>0,
\]

%dsp-level:0
\[
\qquad A_0+B_0=\frac38\Bigl(g+\frac{g'}3\Bigr)^2+\frac{g'^2}{48}>0
\]

%dsp-level:0
\[
\qquad [\mathrm{BM.39}(1);\ \text{machine: the SOS expansions}]:
\]

%dsp-level:0
\[
\boxed{\ \sqrt{A_0-B_0}\ \le\ \sqrt{G\bigl(c(i)\bigr)}\ \le\ \sqrt{A_0+B_0}\qquad\forall\,\text{line}\ i,\ \ \forall\,\text{channel}\ \bigl(\delta\in\bigl\{\tfrac29,\tfrac2{27},\tfrac18\bigr\}\bigr),\ \ \forall\,(g,g')\ne(0,0)\ }
\]

%dsp-level:0
\[
\qquad [\mathrm{BM.39}(1),(3);\ |\cos|\le1;\ \text{machine}];
\]

%dsp-level:0
\[
(5)\quad\textbf{the trichotomy theorem:}\qquad\text{for every frozen cell}\ S\ \text{and every admissible}\ (g,g')\ne(0,0):
\]

%dsp-level:0
\[
(\mathrm a)\quad\textbf{the opening region is non-empty:}
\]

%dsp-level:0
\[
\qquad \mathcal O_S:=\Bigl\{(\rho_\ell,\rho_\nu)>0:\ \ \rho_\ell\sqrt{G_\ell}+\rho_\nu\sqrt{G_\nu}<\Delta T(S)\Bigr\}\ne\varnothing:
\]

%dsp-level:0
\[
\textbf{the universal symbolic witness:}
\]

%dsp-level:0
\[
\qquad \rho_\ell=\rho_\nu=\frac{\Delta T(S)}{4\sqrt{A_0+B_0}}\ \ \in\ \mathcal O_S
\]

%dsp-level:0
\[
\qquad [\text{machine: the sandwich gives}\ \rho_\ell\sqrt{G_\ell}+\rho_\nu\sqrt{G_\nu}\le\tfrac{\Delta T}2<\Delta T];
\]

%dsp-level:0
\[
\textbf{the integer-certified witness at the rational point}\ (g,g',\rho_\ell,\rho_\nu)=(1,1,1,1):
\]

%dsp-level:0
\[
\qquad A_0+B_0=\frac7{16}+\frac14=\frac{11}{16},
\]

%dsp-level:0
\[
\qquad \sqrt{G}\le\frac{\sqrt{11}}4<1\quad[11<16]:
\]

%dsp-level:0
\[
Q_\beta^{\mathrm{spec}}\mathfrak h\ \ge\ \frac{455}{23}-\frac{\sqrt{11}}2\ >\ 0
\]

%dsp-level:0
\[
\qquad \iff
\]

%dsp-level:0
\[
\qquad \frac{910}{23}>\sqrt{11}
\]

%dsp-level:0
\[
\qquad \iff
\]

%dsp-level:0
\[
\qquad 828100>5819
\]

%dsp-level:0
\[
\qquad [\text{machine: the integer certificate}];
\]

%dsp-level:0
\[
(\mathrm b)\quad\textbf{the closing region is non-empty:}
\]

%dsp-level:0
\[
\qquad \mathcal C_S:=\Bigl\{(\rho_\ell,\rho_\nu)>0:\ \ \rho_\ell\sqrt{G_\ell}+\rho_\nu\sqrt{G_\nu}>\Delta T(S)\Bigr\}\ne\varnothing:
\]

%dsp-level:0
\[
\textbf{the universal symbolic witness:}
\]

%dsp-level:0
\[
\qquad \rho_\ell=\rho_\nu=\frac{\Delta T(S)}{\sqrt{A_0-B_0}}\ \ \in\ \mathcal C_S
\]

%dsp-level:0
\[
\qquad [\text{machine: the sandwich gives}\ \rho_\ell\sqrt{G_\ell}+\rho_\nu\sqrt{G_\nu}\ge2\,\Delta T>\Delta T];
\]

%dsp-level:0
\[
\textbf{the integer-certified witness at}\ (1,1,100,100):
\]

%dsp-level:0
\[
\qquad A_0-B_0=\frac3{16},
\]

%dsp-level:0
\[
\qquad \sqrt{G}\ge\frac{\sqrt3}4:
\]

%dsp-level:0
\[
\qquad Q_\beta^{\mathrm{spec}}\mathfrak h\ \le\ \frac{4219785}{98956}-50\sqrt3\ <\ 43-50\sqrt3\ <\ 0
\]

%dsp-level:0
\[
[\text{machine:}\ 4219785<43\cdot98956=4255108;\quad50\sqrt3>43\iff7500>1849];
\]

%dsp-level:0
\[
(\mathrm c)\quad\textbf{the critical locus is exact:}
\]

%dsp-level:0
\[
\qquad \partial\mathcal O_S:\quad\rho_\ell\sqrt{G\bigl(c(0)\bigr)}+\rho_\nu\sqrt{G\bigl(c(i_\nu)\bigr)}=\Delta T(S)
\]

%dsp-level:0
\[
\qquad [\text{machine: the boundary equation}];
\]

%dsp-level:0
\[
\textbf{on the symmetric ray}\ \rho_\ell=\rho_\nu=\rho\ \textbf{at the witness couplings}\ g=g'=1:
\]

%dsp-level:0
\[
\qquad \boxed{\ \frac{273}{23}\ <\ \rho^{\mathrm{crit}}(S,i_\nu)\ <\ \frac{42197850}{791648}\ }
\]

%dsp-level:0
\[
\qquad \text{i.e.}\quad11<\rho^{\mathrm{crit}}<54
\]

%dsp-level:0
\[
\qquad [\text{machine}]:
\]

%dsp-level:0
\[
[\text{the certificates:}\ \sqrt{11}<\tfrac{10}3\iff99<100;\quad\sqrt3>\tfrac85\iff75>64;\quad273>253;\quad42197850<54\cdot791648=42748992];
\]

%dsp-level:0
\[
\textbf{the structure mirrors}\ \mathrm{BM.62}(6)\text{'s critical strength at the scale-combination locus:}
\]

%dsp-level:0
\[
\qquad \rho<\rho^{\mathrm{crit}}\iff Q_\beta^{\mathrm{spec}}>0;
\]

%dsp-level:0
\[
\qquad \rho=\rho^{\mathrm{crit}}\iff Q_\beta^{\mathrm{spec}}=0;
\]

%dsp-level:0
\[
\qquad \rho>\rho^{\mathrm{crit}}\iff Q_\beta^{\mathrm{spec}}<0
\]

%dsp-level:0
\[
\qquad [\text{machine: monotone in}\ \rho];
\]

%dsp-level:0
\[
(6)\quad\textbf{the verdict:}
\]

%dsp-level:0
\[
\qquad \boxed{\ \textbf{the beta decay is allowed:}\ \ \text{on the specified read the vertex is kinematically open exactly on}\ \mathcal O_S\ —\ \text{a non-empty exact region at every frozen winding cell, with integer-certified witnesses}\ }
\]

%dsp-level:0
\[
\textbf{the decay direction is fixed:}
\]

%dsp-level:0
\[
\qquad \text{the parent is the neutral class}\ —\ T(0,S)
\]

%dsp-level:0
\[
= \tfrac43T(+1,S)>T(+1,S)\ [\mathrm{BM.28}(4)(\mathrm{ii});\ \text{machine}]\ —\ \text{and the grammar}\ n\to p+\ell+\nu\ \text{stands unmodified}\ [\mathrm{BM.41}(1);\ \mathrm{BM.45}];
\]

%dsp-level:0
\[
\textbf{the two-body companion row:}
\]

%dsp-level:0
\[
\qquad \text{the two-body channel's admissibility on the specified read is the formed condition}\ \Delta T(S)\ge\rho_\ell\sqrt{G(c(0))}
\]

%dsp-level:0
\[
\qquad [\mathrm{BM.41}(2)\text{'s solution re-derived at}\ m_n\ne m_p:\quad E_p-m_p=\frac{(m_n-m_p)^2-m_\ell^2}{2m_n};\ \text{machine: the cleared identity}]:
\]

%dsp-level:0
\[
\text{at the witness}\ (1,1,1,1)\ \text{it holds}\ \bigl[\tfrac{455}{23}>\tfrac{\sqrt{11}}4\iff1820^2>23^2\cdot11\iff3312400>5819\bigr];
\]

%dsp-level:0
\[
\qquad \mathcal O_S\subseteq\{\text{two-body admissible}\}\ \ \text{since}\ \rho_\nu\sqrt{G_\nu}\ge0\ [\text{machine}];
\]

%dsp-level:0
\[
\textbf{the forced-neutral-member row is read-relative and is recorded,}\ \neg\textbf{adjudicated here:}
\]

%dsp-level:0
\[
\qquad \mathrm{BM.41}(3)\text{'s forcing was derived at the equal rest masses of}\ \mathrm{BM.28}(1)\ —\ \text{it stands on that locus};
\]

%dsp-level:0
\[
\ \ \text{on the specified read the neutral member's emission remains admissible and class-unique}\ [\mathrm{BM.23}(5),(6);\ \mathrm{BM.41}(3)\ \text{grammar}]\ \text{and the named reaction carries it}\ [\mathrm{BM.45}];
\]

%dsp-level:0
\[
\ \ \neg\exists\ \text{a lepton-number law in the inventory}\ [\mathrm{BM.23}(6)];
\]

%dsp-level:0
\[
(7)\quad\textbf{the rescope rows (recorded,}\ \neg\textbf{retracted):}
\]

%dsp-level:0
\[
\qquad \mathrm{BM.41}(4)(\mathrm e)\text{'s verdict}\ \text{(``closed on every adopted read''}):\ \text{at its locus no read was adopted}\ [\mathrm{BM.30}(5)]\ —\ \text{the specification now adopts one, and the verdict is scoped to the pre-specification reads};
\]

%dsp-level:0
\[
\mathrm{BM.49}(1),
\]

%dsp-level:0
\[
(2)\text{'s corrected closure stands in its species}\ —\ \text{the EM summands are}\ \neg\text{parts of the specified read}\ [(3)];
\]

%dsp-level:0
\[
\qquad \mathrm{BM.63}(2)\text{'s closure stands on the informational reads}\ [\mathrm{BM.70}(4)(\mathrm a)];
\]

%dsp-level:0
\[
\mathrm{BM.70}(4)\text{'s dichotomy gains its third case:}
\]

%dsp-level:0
\[
\qquad (\mathrm c)\ \ \text{the specified route-(ii)+transport read: formed, with a branch-relative sign}\ [(5)];
\]

%dsp-level:0
\[
\qquad \mathrm{BM.71}(2)\text{'s no-go is rescoped to the pre-specification reads}\ —\ \text{the fifth route of}\ \mathrm{BM.71}(1)\ \text{is the door walked here, under the logged addendum};
\]

%dsp-level:0
\[
\textbf{no clause of}\ \mathrm{BM.41}\text{--}\mathrm{BM.72}\ \textbf{is retracted; no clause of Parts 1--3 is modified};
\]

%dsp-level:0
\[
\qquad \text{the kinematic viability was always read-relative}\ [\mathrm{BM.41}(7)];
\]

%dsp-level:0
\[
(8)\quad\textbf{the conditional spectrum becomes branch content:}
\]

%dsp-level:0
\[
\qquad \text{on}\ \mathcal O_S\ \text{the hypothesis}\ Q\ge0\ \text{of}\ \mathrm{BM.41}(5)/\mathrm{BM.49}(3)\ \text{is satisfied by a formed read}:
\]

%dsp-level:0
\[
E_\ell^{\max}=\frac{m_n^2+m_\ell^2-(m_p+m_\nu)^2}{2m_n}
\]

%dsp-level:0
\[
\qquad \text{and}
\]

%dsp-level:0
\[
\qquad E_{\max}:\langle E\rangle=2:1\ \ \text{scale-invariant}
\]

%dsp-level:0
\[
\qquad [\mathrm{BM.41}(5);\ \mathrm{BM.49}(3),(4)\ \text{re-cited};\ \text{symbolic in the specified scales}]:
\]

%dsp-level:0
\[
\neg\text{an external endpoint, spectrum or lifetime is compared}\ [\mathrm{P5}];
\]

%dsp-level:0
\[
\qquad \text{the quarantined audit rows stay quarantined}\ [\text{Census FT.114},\ \text{FT.120}\ \text{pattern}];
\]

%dsp-level:0
\[
(9)\quad\textbf{what the specification does not do:}
\]

%dsp-level:0
\[
\qquad \neg\text{a scale is fixed}\ —\ (\rho_\ell,\rho_\nu,g,g',\hat\mu,\hat\mu',\hat\mu_\nu)\ \text{remain undeclared}\ [\mathrm{P4};\ \mathrm{P5}];
\]

%dsp-level:0
\[
\qquad \neg\text{a currency is converted}\ —\ \text{the mass-form bridge is application-layer}\ [(3)];
\]

%dsp-level:0
\[
\neg\text{a Q-value is numerically derived}\ —\ \text{the branch structure is the content};
\]

%dsp-level:0
\[
\qquad \text{the numerical value of}\ e\ \text{remains open and does not enter}\ —\ \text{the specified read carries no EM monomial}\ [\mathrm{EM.25}(7);\ \mathrm{BM.62}(12)];
\]

%dsp-level:0
\[
\text{route}\ (\mathrm{ii})\ \text{is}\ \neg\text{adopted for any purpose beyond the vertex kinematics}\ [\mathrm{BM.30}(5);\ \mathrm{BM.37}(3),(6)];
\]

%dsp-level:0
\[
\qquad \text{the route-(iii) amplitude ratios stay barred as masses}\ [\mathrm{FT.79}(5);\ \mathrm{BM.37}(3)];
\]

%dsp-level:0
\[
(10)\quad\textbf{status:}
\]

%dsp-level:0
\[
\qquad \textbf{derived:}\ (2)\text{--}(8)\ [\text{given}\ (1)];
\]

%dsp-level:0
\[
\qquad \textbf{declared:}\ (1)\ —\ \text{logged as the §II.4 addendum A-BM.8 before use};
\]

%dsp-level:0
\[
\neg\text{a P4 addition}\ (\text{§Cal-BM stays empty};\ \rho_\ell,\rho_\nu\ \text{are Cal-F/P4 restatements}),
\]

%dsp-level:0
\[
\quad\neg\text{a new dimensionful scale}\ [\mathrm{P1}],
\]

%dsp-level:0
\[
\quad\neg\text{a measured value}\ [\mathrm{P5};\ \text{machine: zero external tokens, zero decimal numerals}].
\]

\textbf{Req.} CT2.1, D2.9, EM.25(7) (Part 2); P1, P4, P5 (Parts 1–2); Def-H.37(1), FT.107(3), FT.143 (Part 3); Census FT.114, FT.120 (row-wise); Def-BM.6(1), Def-BM.11(5), BM.20(4), BM.21(2), BM.23(5), BM.23(6), BM.28(1), BM.28(4), BM.30(1), BM.30(5), BM.35(5), BM.37(2), BM.37(3), BM.37(6), BM.39(1), BM.39(3), BM.41(1), BM.41(2), BM.41(3), BM.41(4), BM.41(5), BM.41(7), BM.45, BM.49(1), BM.49(2), BM.49(3), BM.49(4), BM.62(6), BM.62(11), BM.62(12), BM.63(2), BM.66(6), BM.68(8), BM.69(5), BM.70(4), BM.71(1), BM.71(2) (this part); Cal-F; exact rational arithmetic, sum-of-squares positivity, exact radical inequalities by integer certificates, Laurent exponent arithmetic, cleared polynomial identities [math; machine: \texttt{checks/\allowbreak{}verify\_\allowbreak{}H\_\allowbreak{}chiral\_\allowbreak{}assignment\_\allowbreak{}opening.py} §E, §F].

\textbf{Pf.}

%dsp-level:0
\[
(1)\text{--}(2):\ \mathrm{BM.28}(4)(\mathrm{ii})\text{'s frozen factorization}\ T(q,S)=W(q)\overline{\Delta S}(S)\ \text{re-cited};
\]

%dsp-level:0
\[
\qquad \Delta T=(20-15)\overline{\Delta S}=5\,\overline{\Delta S}\ [\text{machine}];
\]

%dsp-level:0
\[
\qquad \overline{\Delta S}(S)=m^2_B(\cdot,S)\ [\mathrm{FT.107}(3):\ \text{the normalized average};\ \mathrm{BM.28}(4)(\mathrm i)\ \text{re-cited}]
\]

%dsp-level:0
\[
\Delta T\ \text{evaluated at the four frozen sector expectations}\ \Bigl(\tfrac{843957}{98956},\ \tfrac{1921662}{393679},\ \tfrac{91}{23},\ \tfrac{13713}{3088}\Bigr)\ [\mathrm{BM.21}(2)/\mathrm{BM.62}(2)\ \text{re-cited}]:
\]

%dsp-level:0
\[
\quad 5\cdot843957=4219785,
\]

%dsp-level:0
\[
\quad 5\cdot1921662=9608310,
\]

%dsp-level:0
\[
\quad 5\cdot91=455,
\]

%dsp-level:0
\[
\quad 5\cdot13713=68565\ [\text{machine: exact integers}];
\]

%dsp-level:0
\[
\qquad \text{the min/max comparisons by cross-multiplication}\ [\text{machine}]
\]

%dsp-level:0
\[
(3):\ \text{the locus rows re-cited}\ [\mathrm{BM.41}(4)(\mathrm d);\ \mathrm{BM.37}(2);\ \mathrm{BM.62}(11)]
\]

%dsp-level:0
\[
(4):\ \text{the SOS rows of}\ \mathrm{BM.39}(1)\ \text{re-derived by expansion}\ [\text{machine:}\ \tfrac38(g\mp\tfrac{g'}3)^2+\tfrac{g'^2}{48}=A_0\mp B_0];
\]

%dsp-level:0
\[
\qquad |\cos\varphi_i|\le1\implies\text{the sandwich}\ [\text{math}];
\]

%dsp-level:0
\[
\qquad \text{line- and channel-independence}\ [\mathrm{BM.39}(3):\ (A_0,B_0)\ \text{shared}]
\]

%dsp-level:0
\[
(5)(\mathrm a):\ \text{the symbolic witness:}\ 2\cdot\tfrac{\Delta T}{4\sqrt{A_0+B_0}}\cdot\sqrt{A_0+B_0}
\]

%dsp-level:0
\[
= \tfrac{\Delta T}2<\Delta T\ [\text{machine: exact in the ordered radical extension}];
\]

%dsp-level:0
\[
\qquad \text{the rational witness:}\ A_0+B_0\big|_{g=g'=1}=\tfrac{11}{16}\ [\text{machine}];
\]

%dsp-level:0
\[
\quad (\tfrac{\sqrt{11}}4)^2
\]

%dsp-level:0
\[
= \tfrac{11}{16};
\]

%dsp-level:0
\[
\quad\tfrac{455}{23}-\tfrac{\sqrt{11}}2>0
\]

%dsp-level:0
\[
\iff 910>23\sqrt{11}
\]

%dsp-level:0
\[
\iff 828100>529\cdot11
\]

%dsp-level:0
\[
= 5819\ [\text{machine: squaring both positive sides}]
\]

%dsp-level:0
\[
(\mathrm b):\ \text{the symbolic witness:}\ 2\cdot\tfrac{\Delta T}{\sqrt{A_0-B_0}}\cdot\sqrt{A_0-B_0}=2\,\Delta T>\Delta T\ [\text{machine}];
\]

%dsp-level:0
\[
\qquad \text{the rational witness:}\ A_0-B_0\big|_{g=g'=1}=\tfrac3{16};
\]

%dsp-level:0
\[
\quad 200\cdot\tfrac{\sqrt3}4=50\sqrt3;\quad\Delta T_{\max}<43\iff4219785<4255108;
\]

%dsp-level:0
\[
\quad 50\sqrt3>43\iff7500>1849\ [\text{machine}]
\]

%dsp-level:0
\[
(\mathrm c):\ \rho^{\mathrm{crit}}=\tfrac{\Delta T}{\sqrt{G_\ell}+\sqrt{G_\nu}}:
\]

%dsp-level:0
\[
\qquad \text{lower bound at}\ \Delta T_{\min},
\]

%dsp-level:0
\[
\ \sqrt{G}\le\tfrac{\sqrt{11}}4:\quad\rho^{\mathrm{crit}}\ge\tfrac{2\cdot455/23}{\sqrt{11}}
\]

%dsp-level:0
\[
= \tfrac{910}{23\sqrt{11}}>\tfrac{910\cdot3}{23\cdot10}
\]

%dsp-level:0
\[
= \tfrac{273}{23}\ [\sqrt{11}<\tfrac{10}3\iff99<100];
\]

%dsp-level:0
\[
\text{upper bound at}\ \Delta T_{\max},
\]

%dsp-level:0
\[
\ \sqrt{G}\ge\tfrac{\sqrt3}4:\quad\rho^{\mathrm{crit}}\le\tfrac{2\cdot4219785/98956}{\sqrt3}
\]

%dsp-level:0
\[
= \tfrac{8439570}{98956\sqrt3}<\tfrac{8439570\cdot5}{98956\cdot8}
\]

%dsp-level:0
\[
= \tfrac{42197850}{791648}\ [\sqrt3>\tfrac85\iff75>64];
\]

%dsp-level:0
\[
\qquad \tfrac{273}{23}>11\iff273>253;\quad\tfrac{42197850}{791648}<54\iff42197850<42748992\ [\text{machine: integers}];
\]

%dsp-level:0
\[
\qquad \text{monotonicity in}\ \rho:\ \text{the coefficient}\ \sqrt{G}>0\ [\text{machine}]
\]

%dsp-level:0
\[
(6):\ T(0,S)=\tfrac43T(+1,S)>T(+1,S)\ [\text{machine}];
\]

%dsp-level:0
\[
\qquad \text{the two-body identity:}\ 2m_n(E_p-m_p)
\]

%dsp-level:0
\[
= (m_n-m_p)^2-m_\ell^2\ \text{at}\ E_p
\]

%dsp-level:0
\[
= \tfrac{m_n^2+m_p^2-m_\ell^2}{2m_n}\ [\text{machine: cleared over}\ \mathbb Q[m_n,m_p,m_\ell]];
\]

%dsp-level:0
\[
\qquad \text{the witness}\ \tfrac{455}{23}>\tfrac{\sqrt{11}}4\iff1820>23\sqrt{11}\iff3312400>5819\ [\text{machine}];
\]

%dsp-level:0
\[
\qquad \mathcal O_S\subseteq\{\Delta T\ge\rho_\ell\sqrt{G_\ell}\}\ [\rho_\nu\sqrt{G_\nu}\ge0;\ \text{machine}]
\]

%dsp-level:0
\[
(7)\text{--}(10):\ \text{the rescope census clause by clause}\ [\text{machine}];
\]

%dsp-level:0
\[
\qquad \mathrm{BM.41}(5),\ \mathrm{BM.49}(3),(4)\ \text{re-cited under the satisfied hypothesis};
\]

%dsp-level:0
\[
\qquad \text{the standing-rule scans}\ [\text{machine}]. \ \blacksquare
\]

\textbf{Phys.} The corpus always knew exactly one door out of its own beta-decay closure: a read specification, an act it priced repeatedly and performed never. This tier performs it, narrowly and on the record — the vertex Q-value is read on the unnormalized total-deficit route for the baryons, where the neutral class outweighs the charged by the exact factor four thirds, and on the transport family form for the emitted lines, whose masses carry undeclared scales between exact sum-of-squares bounds. The resulting Q-value is the frozen expression the decay tier itself formed, and its sign is decided by a single dimensionless scale combination against an exact rational: five times the sector expectation, which is forty-two and change nats at the first frozen cell and never below nineteen. The trichotomy is total and machine-certified — an opening region with a universal symbolic witness and a rational witness certified by eight-two-eight-one-hundred against five thousand eight hundred nineteen, a closing region with its own witnesses, and an exact critical locus enclosed between eleven and fifty-four at the rational coupling point. The beta decay is allowed: open on that region, closed on its complement, degenerate on the boundary, with the parent still the neutral baryon and the grammar untouched. Nothing is fitted and nothing is converted: the scales stay undeclared, the two currencies stay unequated, the informational closures stand in their own species, and every rescoped verdict is recorded rather than erased — the observed world's endpoint may one day be compared in a quarantined row, after a calibration this tier does not perform.

\end{theorem}

\medskip\noindent\textit{\texorpdfstring{BM.78 — The status of the chiral assignment: adjudication and scope}{BM.78 - The status of the chiral assignment: adjudication and scope}.}\ Structural-constraints record: see Appendix B, Section B.1.25 of the Addendum.\par\medskip

\phantomsection
\section[Tier BM.17 — The Mass Ratio Census: the Platonic Winding…]{\texorpdfstring{Tier BM.17 — The Mass Ratio Census: the Platonic Winding Ladders, the Intra- and Inter-Generational Ratios, the Six-State Bridge Census, the Four-Currency Incommensurability, and the Neutron–Proton Ratio}{Tier BM.17 - The Mass Ratio Census: the Platonic Winding Ladders, the Intra- and Inter-Generational Ratios, the Six-State Bridge Census, the Four-Currency Incommensurability, and the Neutron-Proton Ratio}}

\textbf{Locus.} This tier opens no gap and closes none: it takes the inventory of every mass ratio the corpus derives, in each currency separately, and localizes the whole obstruction to one step. It consumes $6$ states $[\mathrm{Def\text{-}BM.14}]$, two hierarchies $[\mathrm{BM.60}(6)]$, one attached splitting $[\mathrm{BM.75}(4)]$, one bridge identity $[\mathrm{BM.69}(2)]$, one exact critical strength $[\mathrm{BM.62}(6)]$, one specified read $[\mathrm{BM.77}]$, and the Platonic ladder cells $[\text{Census FT.63}(4),(6);\ \text{Census FT.76}\ \text{row-wise}]$. Every clause re-derives its own arithmetic in an exact field; $\neg$a clause equates two currencies; $\neg$a clause imports an observed mass, ratio or coupling $[\mathrm{P5}]$. The tier adds no definition: every new symbol is an in-theorem abbreviation of derived objects.

\textbf{Standing rows.} $\alpha_B$ stays barred on $\mathrm{BM.61}(8)$'s four grounds and §Cal-BM stays empty. The relating equation is $\neg$assumed anywhere — item 7 stays open $[\mathrm{BM.35}(5);\ \mathrm{BM.68}(7);\ \mathrm{BM.69}(4),(5);\ \text{pointer}]$, and this tier's contribution to it is a localization, $\neg$a value. $\mathrm{BM.41}(4)(\mathrm e)$, $\mathrm{BM.49}$, $\mathrm{BM.63}$, $\mathrm{BM.70}(4)$, $\mathrm{BM.71}(2)$ stand on the reads formed at their loci, and $\mathrm{BM.77}(5),(6)$'s specified-read trichotomy stands beside them, $\neg$retracted and $\neg$converted. Local notation, fixed for the whole tier: $\mathsf c:=\cos\frac\pi{10}$ with $16\mathsf c^4-20\mathsf c^2+5=0$ — the symbol $c$ keeps its $\mathrm{Def\text{-}BM.12}(4)$ meaning (the priced EM distribution factor, $c\in\{\tfrac12,\tfrac35\}$) at every occurrence; $\chi$ is not used; $z_i^{(\mathrm{ch})}:=(r_i^{(\mathrm{ch})})^2$ and $\gamma_{\mathrm{ch}}:=\sqrt2\cos(3\delta_{\mathrm{ch}})$.

\pseudohead{theorem}{\texorpdfstring{BM.79}{BM.79}}{\texorpdfstring{The Platonic winding census: the ladder law, the Chebyshev tower, and the intra-generational mass ratios}{The Platonic winding census: the ladder law, the Chebyshev tower, and the intra-generational mass ratios}}\label{tag:BM-79}
\textbf{Thm.}

%dsp-level:0
\[
(1)\quad\textbf{the ladder law and the defect angle:}
\]

%dsp-level:0
\[
\qquad \delta_v=2\pi-d\,\frac{(L-2)\pi}L,
\]

%dsp-level:0
\[
\qquad \frac{\delta_v}\pi=\Bigl(1,\ \frac12,\ \frac23,\ \frac15,\ \frac13\Bigr)\ \ \text{at}\ \ (d,L)=(3,3),(3,4),(4,3),(3,5),(5,3)
\]

%dsp-level:0
\[
\qquad [\text{Census FT.42}(5);\ \mathrm{BM.60}(3);\ \text{machine}]:
\]

%dsp-level:0
\[
\boxed{\ V\,\delta_v=4\pi\ \ \text{on all five solids}\ }
\]

%dsp-level:0
\[
\qquad 4\cdot\pi=8\cdot\tfrac\pi2=6\cdot\tfrac{2\pi}3=20\cdot\tfrac\pi5=12\cdot\tfrac\pi3
\]

%dsp-level:0
\[
\qquad [\text{Descartes/Gauss--Bonnet};\ \mathrm{FT.47}(1);\ \text{machine: exact in}\ \pi\text{-units}];
\]

%dsp-level:0
\[
\boxed{\ D(m)\ :=\ 4\sin^2\!\Bigl(\frac{m\,\delta_v}4\Bigr)\ =\ 2-2\cos\Bigl(\frac{m\,\delta_v}2\Bigr)\ =\ 2-2\,T_{|m|}\!\Bigl(\cos\frac{\delta_v}2\Bigr),\qquad m^2_{\mathrm{pole}}(m)=\hat\mu'^2\,D(m)\ }
\]

%dsp-level:0
\[
[\mathrm{Def\text{-}H.37}(2);\ \mathrm{FT.54}(4);\ \mathrm{FT.60}(5):\ \text{the weight ladder at}\ \Theta_R=\delta_v;\ \ \text{machine: the half-angle and Chebyshev identities in each exact field}];
\]

%dsp-level:0
\[
\cos\frac{\delta_v}2\ =\ \Bigl(0,\ \frac{\sqrt2}2,\ \frac12,\ \mathsf c,\ \frac{\sqrt3}2\Bigr)
\]

%dsp-level:0
\[
\qquad [\text{machine}];
\]

%dsp-level:0
\[
\qquad \textbf{the two identically massless readings:}
\]

%dsp-level:0
\[
\qquad \Theta_R=0:\ D\equiv0;
\]

%dsp-level:0
\[
\qquad \Theta_R=V\delta_v=4\pi:\ D(m)=2-2\cos(2m\pi)=0\ \ \forall m
\]

%dsp-level:0
\[
[\text{Census FT.63}(4),(6);\ \text{Census FT.76}(1)\ \text{row-wise};\ \text{machine}]:
\]

%dsp-level:0
\[
\qquad \text{the flat-face side and the global sweep carry no defect mass at any weight};
\]

%dsp-level:0
\[
(2)\quad\textbf{the winding conventions (both priced, neither adopted):}
\]

%dsp-level:0
\[
\qquad \kappa_{\min}=2\ (\text{odd}\ L),\ 1\ (\text{even}\ L),
\]

%dsp-level:0
\[
\qquad k=\frac{L\,\kappa_{\min}}2
\]

%dsp-level:0
\[
\qquad [\mathrm{FT.32}(6);\ \mathrm{FT.45};\ \text{machine}]:
\]

%dsp-level:0
\[
k_f=\frac{K_S}F=(3,\,2,\,3,\,5,\,3)\ \ \text{at}\ \ K_S=(12,12,24,60,60),\ F=(4,6,8,12,20)
\]

%dsp-level:0
\[
\qquad [\text{Census FT.43}(1),\ \text{FT.63}(1)\ \text{row-wise};\ \text{machine}];
\]

%dsp-level:0
\[
k_{vf}=(3,\,3,\,2,\,3,\,5)\ \ \text{on the vertex-figure}\ L_{vf}=(3,3,4,3,5)
\]

%dsp-level:0
\[
\qquad [\text{Census FT.76}(1)\ \text{row-wise};\ \text{machine}];
\]

%dsp-level:0
\[
\textbf{variant}\ (\mathrm A):=(\delta_v,k_f);
\]

%dsp-level:0
\[
\qquad \textbf{variant}\ (\mathrm B):=(\delta_v,k_{vf})
\]

%dsp-level:0
\[
\qquad [\text{Census FT.76}(1):\ \text{both priced};\ \text{the census below reports both}];
\]

%dsp-level:0
\[
(3)\quad\textbf{the variant-}(\mathrm A)\textbf{ multiplets (the census rows):}\qquad m\in\{k_f,k_f{-}2,\dots,-k_f\}:
\]

%tbl-level:0
{\setlength{\tabcolsep}{4pt}
\begin{longtable}{>{\raggedright\arraybackslash}p{\dimexpr0.1404\textwidth-2\tabcolsep\relax}|>{\raggedright\arraybackslash}p{\dimexpr0.0827\textwidth-2\tabcolsep\relax}|>{\raggedright\arraybackslash}p{\dimexpr0.0443\textwidth-2\tabcolsep\relax}|>{\raggedright\arraybackslash}p{\dimexpr0.4576\textwidth-2\tabcolsep\relax}|>{\raggedright\arraybackslash}p{\dimexpr0.2750\textwidth-2\tabcolsep\relax}}
\toprule
solid & $\delta_v/\pi$ & $k_f$ & the $D(m)$ multiset & exact field \\
\addlinespace[2pt]
\midrule
\endfirsthead
\toprule
solid & $\delta_v/\pi$ & $k_f$ & the $D(m)$ multiset & exact field \\
\addlinespace[2pt]
\midrule
\endhead
\bottomrule
\endlastfoot
tetrahedron & $1$ & $3$ & $\{2,\ 2,\ 2,\ 2\}$ at $m=\allowbreak (3,\allowbreak 1,\allowbreak -1,\allowbreak -3)$ & $\mathbb Q$ \\
cube & $\tfrac12$ & $2$ & $\{2,\ \mathbf 0,\ 2\}$ at $m=\allowbreak (2,\allowbreak 0,\allowbreak -2)$ & $\mathbb Q$ \\
octahedron & $\tfrac23$ & $3$ & $\{4,\ 1,\ 1,\ 4\}$ at $m=\allowbreak (3,\allowbreak 1,\allowbreak -1,\allowbreak -3)$ & $\mathbb Q$ \\
dodecahedron & $\tfrac15$ & $5$ & $\{2,\ 2{-}2\cos\tfrac{3\pi}{10},\ 2{-}2\mathsf c,\ 2{-}2\mathsf c,\ 2{-}2\cos\tfrac{3\pi}{10},\ 2\}$ & $\mathbb Q[\mathsf c]/(16\mathsf c^4{-}20\mathsf c^2{+}5)$ \\
icosahedron & $\tfrac13$ & $3$ & $\{2,\ 2{-}\sqrt3,\ 2{-}\sqrt3,\ 2\}$ & $\mathbb Q(\sqrt3)$ \\
\end{longtable}
}

%dsp-level:0
\[
[\text{machine: every entry derived from}\ (1)\text{'s ladder law in its own exact field};\ \text{cross-checked row-wise against Census FT.76}(2)(\mathrm A)];
\]

%dsp-level:0
\[
\textbf{the variant-}(\mathrm B)\textbf{ row (recorded, distinct):}
\]

%dsp-level:0
\[
\qquad \text{tetra}\ \{2,2,2,2\};\quad\text{cube}\ \{2{+}\sqrt2,\ 2{-}\sqrt2,\ 2{-}\sqrt2,\ 2{+}\sqrt2\};\quad\text{octa}\ \{3,\ \mathbf 0,\ 3\};
\]

%dsp-level:0
\[
\text{dodeca}\ \{2{-}2\cos\tfrac{3\pi}{10},\ 2{-}2\mathsf c,\ 2{-}2\mathsf c,\ 2{-}2\cos\tfrac{3\pi}{10}\}\ (k_{vf}{=}3);
\]

%dsp-level:0
\[
\text{icosa}\ \{2{+}\sqrt3,\ 2,\ 2{-}\sqrt3,\ 2{-}\sqrt3,\ 2,\ 2{+}\sqrt3\}\ (k_{vf}{=}5)
\]

%dsp-level:0
\[
\qquad [\text{machine: every entry derived from}\ (1)\text{'s ladder law};\ \text{cross-checked row-wise against Census FT.76}(2)(\mathrm B)];
\]

%dsp-level:0
\[
(4)\quad\textbf{the Chebyshev tower (the dodecahedral row, exact):}
\]

%dsp-level:0
\[
\qquad T_0=1,\quad T_1=\mathsf c,\quad T_{m+1}=2\mathsf c\,T_m-T_{m-1}
\]

%dsp-level:0
\[
\qquad [\text{integer coefficients};\ \text{machine}]:
\]

%tbl-level:0
{\setlength{\tabcolsep}{4pt}
\begin{longtable}{>{\raggedright\arraybackslash}p{\dimexpr0.0575\textwidth-2\tabcolsep\relax}|>{\raggedright\arraybackslash}p{\dimexpr0.2128\textwidth-2\tabcolsep\relax}|>{\raggedright\arraybackslash}p{\dimexpr0.3801\textwidth-2\tabcolsep\relax}|>{\raggedright\arraybackslash}p{\dimexpr0.3497\textwidth-2\tabcolsep\relax}}
\toprule
$m$ & $T_m(\mathsf c)$ reduced & $D(m)=\allowbreak 2-2T_m(\mathsf c)$ & minimal polynomial of $D(m)$ \\
\addlinespace[2pt]
\midrule
\endfirsthead
\toprule
$m$ & $T_m(\mathsf c)$ reduced & $D(m)=\allowbreak 2-2T_m(\mathsf c)$ & minimal polynomial of $D(m)$ \\
\addlinespace[2pt]
\midrule
\endhead
\bottomrule
\endlastfoot
$0$ & $1$ & $0$ & $x$ \\
$1$ & $\mathsf c$ & $2-2\mathsf c$ & $x^4-8x^3+\allowbreak 19x^2-12x+\allowbreak 1$ \\
$2$ & $2\mathsf c^2-1$ & $4-4\mathsf c^2=\allowbreak \tfrac{3-\sqrt5}2$ & $x^2-3x+\allowbreak 1$ \\
$3$ & $4\mathsf c^3-3\mathsf c$ & $2+\allowbreak 6\mathsf c-8\mathsf c^3=\allowbreak 2-2\cos\tfrac{3\pi}{10}$ & $x^4-8x^3+\allowbreak 19x^2-12x+\allowbreak 1$ \\
$4$ & $2\mathsf c^2-\tfrac32$ & $5-4\mathsf c^2=\allowbreak \tfrac{5-\sqrt5}2$ & $x^2-5x+\allowbreak 5$ \\
$5$ & $\mathbf 0$ & $\mathbf 2$ & $x-2$ \\
\end{longtable}
}

%dsp-level:0
\[
\boxed{\ T_5(\mathsf c)=16\mathsf c^5-20\mathsf c^3+5\mathsf c=\mathsf c\,(16\mathsf c^4-20\mathsf c^2+5)=0\ }:
\]

%dsp-level:0
\[
\qquad \textbf{the defining quartic IS the top weight}
\]

%dsp-level:0
\[
[\text{de Moivre}\ \cos5t=16\cos^5t-20\cos^3t+5\cos t\ \text{at}\ t=\tfrac\pi{10}:\ \ \cos\tfrac\pi2=0\iff D(5)=4\sin^2\tfrac\pi4=2;\ \text{machine: exact in}\ \mathbb Z[\mathsf c]/(16\mathsf c^4{-}20\mathsf c^2{+}5)];
\]

%dsp-level:0
\[
\mathsf c^2=\frac{5+\sqrt5}8,
\]

%dsp-level:0
\[
\qquad \sqrt5=8\mathsf c^2-5\in\mathbb Q(\mathsf c)\ \ \text{is the unique quadratic subfield},
\]

%dsp-level:0
\[
\qquad [\mathbb Q(\mathsf c):\mathbb Q]=4
\]

%dsp-level:0
\[
\qquad [\text{Census FT.76}(9)\ \text{row-wise};\ \text{machine}];
\]

%dsp-level:0
\[
(5)\quad\textbf{the ladder structure (all five solids):}
\]

%dsp-level:0
\[
\qquad P_S=\frac{4\pi}{\delta_v}=(4,\,8,\,6,\,20,\,12),
\]

%dsp-level:0
\[
\qquad \#\{\text{distinct values}\}=\frac{P_S}2+1=(3,\,5,\,4,\,11,\,7)
\]

%dsp-level:0
\[
\qquad [\text{Census FT.76}(3)\ \text{row-wise};\ \text{machine}]:
\]

%dsp-level:0
\[
D\ \text{is periodic in}\ m\ \text{with period}\ P_S,\ \text{even in}\ m,\ \text{and strictly increasing on}\ [0,P_S/2];
\]

%dsp-level:0
\[
\qquad \max_mD(m)=4\ \ \text{at}\ \ m=P_S/2\ \ (P_S\ \text{even on all five});
\]

%dsp-level:0
\[
\boxed{\ 0\ \le\ m^2_{\mathrm{pole}}(m)\ \le\ 4\hat\mu'^2\ \ \text{at every weight of every solid}\ }\qquad[\text{Census FT.76}(3);\ \text{machine}];
\]

%dsp-level:0
\[
\textbf{the rung law:}
\]

%dsp-level:0
\[
\qquad D(m)\in\{0,2,4\}\iff m\,\delta_v\in\pi\mathbb Z
\]

%dsp-level:0
\[
\qquad [\text{machine: all five ladders to}\ m=P_S,\ \text{in the exact quadratic rings, with the dodecahedral quartic avoided}];
\]

%dsp-level:0
\[
\textbf{the top-weight row:}
\]

%dsp-level:0
\[
\qquad k_f\delta_v=(3\pi,\ \pi,\ 2\pi,\ \pi,\ \pi)\implies D(k_f)=(2,\ 2,\ 4,\ 2,\ 2)
\]

%dsp-level:0
\[
\qquad [\text{Census FT.76}(4)\ \text{row-wise};\ \text{machine}]:
\]

%dsp-level:0
\[
\textbf{the heaviest weight of four of the five solids is exactly}\ 2\hat\mu'^2,
\]

%dsp-level:0
\[
\ \textbf{the octahedron's being}\ 4\hat\mu'^2\ —\ \textbf{the inter-solid hierarchy lives entirely at the light end};
\]

%dsp-level:0
\[
(6)\quad\textbf{the two structural obstructions (re-certified):}
\]

%dsp-level:0
\[
\qquad (\mathrm a)\ \ \delta_v=\pi\implies D(m)=2-2\cos\frac{m\pi}2\in\{0,2,4\}\ \text{and the}\ k_f{=}3\ \text{multiplet is fully degenerate at}\ 2:
\]

%dsp-level:0
\[
\textbf{the tetrahedron carries no multiplet splitting at all in either variant}
\]

%dsp-level:0
\[
\qquad [\text{Census FT.76}(5)(\mathrm a);\ \text{machine}];
\]

%dsp-level:0
\[
(\mathrm b)\ \ \text{the even-}k\ \text{zero mode:}
\]

%dsp-level:0
\[
\quad k_f=2\ \text{on the cube}\ (\mathrm A)\ \text{and}\ k_{vf}=2\ \text{on the octahedron}\ (\mathrm B)\ \text{each carry}\ m=0\ \text{with}\ D(0)=0\ \textbf{exactly}:
\]

%dsp-level:0
\[
\textbf{on those two multiplets the literal lightest pole mass is}\ 0\ \text{and only lightest-nonzero reads are well defined}
\]

%dsp-level:0
\[
\qquad [\mathrm{FT.60}(5)\text{'s standing zero};\ \text{Census FT.76}(5)(\mathrm b);\ \text{machine}];
\]

%dsp-level:0
\[
(7)\quad\textbf{the intra-generational mass ratios (the multiplet read):}
\]

%dsp-level:0
\[
\qquad r_{\mathrm{intra}}(S)\ :=\ \sqrt{\frac{D_{\max}}{D^{\mathrm{nz}}_{\min}}}
\]

%dsp-level:0
\[
\qquad [D^{\mathrm{nz}}_{\min}:=\text{the lightest nonzero value of the multiplet},\ (6)(\mathrm b)]:
\]

%dsp-level:0
\[
\boxed{\ \textbf{the sine-ratio law:}\qquad r_{\mathrm{intra}}(S)\ =\ \frac{\sin\bigl(k_f\delta_v/4\bigr)}{\sin\bigl(\delta_v/4\bigr)}\ \ \text{on every solid whose multiplet carries}\ m=1\ }
\]

%dsp-level:0
\[
[\text{machine: the half-angle identity}\ D(m)=4\sin^2(m\delta_v/4)\ \text{at the two extreme weights, all five solids, both variants}];
\]

%tbl-level:1
{\small\setlength{\tabcolsep}{4pt}
\begin{longtable}{>{\raggedright\arraybackslash}p{\dimexpr0.1278\textwidth-2\tabcolsep\relax}|>{\raggedright\arraybackslash}p{\dimexpr0.1191\textwidth-2\tabcolsep\relax}|>{\raggedright\arraybackslash}p{\dimexpr0.1891\textwidth-2\tabcolsep\relax}|>{\raggedright\arraybackslash}p{\dimexpr0.1334\textwidth-2\tabcolsep\relax}|>{\raggedright\arraybackslash}p{\dimexpr0.2241\textwidth-2\tabcolsep\relax}|>{\raggedright\arraybackslash}p{\dimexpr0.2066\textwidth-2\tabcolsep\relax}}
\toprule
solid & (A) $r_{\mathrm{intra}}$ & exact form & minimal polynomial & (B) $r_{\mathrm{intra}}$ & exact form \\
\addlinespace[2pt]
\midrule
\endfirsthead
\toprule
solid & (A) $r_{\mathrm{intra}}$ & exact form & minimal polynomial & (B) $r_{\mathrm{intra}}$ & exact form \\
\addlinespace[2pt]
\midrule
\endhead
\bottomrule
\endlastfoot
tetrahedron & $1$ & $1$ (full degeneracy) & $x-1$ & $1$ & $1$ \\
cube & $1$ & $1$ $(+\allowbreak 1$ exact zero mode$)$ & $x-1$ & $1+\allowbreak \sqrt2$ & $\sqrt{\tfrac{2+\sqrt2}{2-\sqrt2}}$ \\
octahedron & $2$ & $\sqrt{4/1}$ & $x-2$ & $1$ & $1$ $(+\allowbreak 1$ zero mode$)$ \\
dodecahedron & $\sqrt{\tfrac{2}{2-2\mathsf c}}$ & $\dfrac{\sin(\pi/4)}{\sin(\pi/20)}$ & $x^8-24x^6+\allowbreak 76x^4-64x^2+\allowbreak 16$ & $\sqrt{\tfrac{2-2\cos(3\pi/10)}{2-2\mathsf c}}$ & $\dfrac{\sin(3\pi/20)}{\sin(\pi/20)}$ \\
icosahedron & $1+\allowbreak \sqrt3$ & $\sqrt{\tfrac{2}{2-\sqrt3}}$ & $x^2-2x-2$ & $2+\allowbreak \sqrt3$ & $\sqrt{\tfrac{2+\sqrt3}{2-\sqrt3}}$ \\
\end{longtable}
}

%dsp-level:0
\[
[\text{machine: the five}\ (\mathrm A)\ \text{rows and the five}\ (\mathrm B)\ \text{rows derived in}\ \mathbb Q,\ \mathbb Q(\sqrt2),\ \mathbb Q(\sqrt3),\ \mathbb Q(\mathsf c);\ \text{the minimal polynomials by the multiplication map and by first linear dependence};\ \text{cross-checked row-wise against Census FT.76}(7)];
\]

%dsp-level:0
\[
\textbf{the nested-radical form of the widest row:}
\]

%dsp-level:0
\[
\qquad r^2_{\mathrm{intra}}(\text{dodeca})=\frac{2}{2-2\mathsf c}=\frac1{1-\mathsf c}=\frac{\bigl(4+\sqrt{10+2\sqrt5}\,\bigr)\bigl(3+\sqrt5\,\bigr)}2
\]

%dsp-level:0
\[
\qquad [\text{machine: exact in}\ \mathbb Q(\sqrt5,\sqrt{10+2\sqrt5})];
\]

%dsp-level:0
\[
(8)\quad\textbf{the exact order of the intra-generational ratios:}\qquad\boxed{\ 1\ =\ 1\ <\ 2\ <\ 1+\sqrt3\ <\ \sqrt{\frac{2}{2-2\mathsf c}}\ }
\]

%dsp-level:0
\[
\text{certificates:}\quad1<2;\qquad2<1+\sqrt3\iff1<\sqrt3\iff1<3;\qquad(1+\sqrt3)^2=4+2\sqrt3<\frac{2}{2-2\mathsf c}
\]

%dsp-level:0
\[
\iff (4+2\sqrt3)(2-2\mathsf c)<2\iff\Bigl(4+\frac72\Bigr)\Bigl(2-\frac{19}{10}\Bigr)
\]

%dsp-level:0
\[
= \frac34<2\ \ \text{on}\ \ \sqrt3<\frac74\ (48<49)\ \text{and}\ \mathsf c>\frac{19}{20}\ \Bigl(5\cdot2500>111^2:\ 12500>12321\Bigr)
\]

%dsp-level:0
\[
[\text{machine: integer certificates and exact dyadic radical enclosures}];
\]

%dsp-level:0
\[
\qquad \textbf{the widest intra-generational ratio is the dodecahedron's, in both variants};
\]

%dsp-level:0
\[
(9)\quad\textbf{the ceiling row (distinct from}\ (7)\textbf{):}
\]

%dsp-level:0
\[
\qquad \sqrt{\frac4{D_S}}\ =\ \Bigl(\sqrt2,\ \ \frac{2}{\sqrt{2-\sqrt2}},\ \ 2,\ \ \frac{2}{\sqrt{2-2\mathsf c}},\ \ \frac{2}{\sqrt{2-\sqrt3}}\Bigr)
\]

%dsp-level:0
\[
\qquad [\text{Census FT.76}(3)\ \text{row-wise};\ \text{machine}]:
\]

%dsp-level:0
\[
\text{maximal at the dodecahedron},
\]

%dsp-level:0
\[
\qquad \frac{4}{2-2\mathsf c}=2\,r^2_{\mathrm{intra}}(\text{dodeca}),
\]

%dsp-level:0
\[
\qquad \text{minimal polynomial}\ \ x^4-48x^3+304x^2-512x+256;
\]

%dsp-level:0
\[
\frac{2}{\sqrt{2-2\mathsf c}}=\sqrt2\cdot r_{\mathrm{intra}}(\text{dodeca})
\]

%dsp-level:0
\[
\qquad [\text{machine: the exact identity}];
\]

%dsp-level:0
\[
\qquad \textbf{no winding choice and no admissible}\ \Theta_R\ \textbf{convention exceeds it}\ \ [\text{Census FT.76}(3);\ \text{the value set is finite}];
\]

%dsp-level:0
\[
(10)\quad\textbf{the range bar, re-certified with the census numbers:}
\]

%dsp-level:0
\[
\qquad \text{the widest mass ratio the Platonic defect species supplies at any winding is}\ \sqrt{4/(2-2\mathsf c)}:
\]

%dsp-level:0
\[
\sqrt{\frac{4}{2-2\mathsf c}}\ <\ 7
\]

%dsp-level:0
\[
\qquad \Bigl[\frac{4}{2-2\mathsf c}<49\iff4<49(2-2\mathsf c)\iff\mathsf c<\frac{47}{49}:\ \mathsf c^2=\frac{5+\sqrt5}8<\frac{2209}{2401}\iff\sqrt5<\frac{3368}{2401}\cdot8-5;\ \text{machine}\Bigr];
\]

%dsp-level:0
\[
\text{and}\ \ 139\ <\ \frac{z_0^{(q)}}{z_1^{(q)}}
\]

%dsp-level:0
\[
\qquad [\mathrm{BM.80}(6):\ \text{the flattest channel hierarchy's outer ratio};\ \text{machine}]\implies\boxed{\ \sqrt{\frac{4}{2-2\mathsf c}}\ <\ 7\ <\ 139\ <\ \frac{z^{(q)}_0}{z^{(q)}_1}\ }
\]

%dsp-level:0
\[
\textbf{the two are dimensionless ratios of different species: their comparison is arithmetic,}\ \neg\textbf{a bridge};
\]

%dsp-level:0
\[
\qquad \text{the widest Platonic mass ratio lies below the flattest inter-generational channel ratio by an exact margin};
\]

%dsp-level:0
\[
\text{deficit}\ 12\text{'s two ceilings stand untouched}\qquad[\text{Census FT.76}(3);\ \mathrm{FT.79}(2);\ \mathrm{FT.84}(2)\ \text{row-wise}];
\]

%dsp-level:0
\[
(11)\quad\textbf{status:}
\]

%dsp-level:0
\[
\qquad \textbf{derived:}\ (1),\ (3)\text{--}(10);
\]

%dsp-level:0
\[
\qquad \textbf{re-cited row-wise as Census data:}\ (2)\text{'s two windings and}\ (9)\text{'s ceiling}\ [\text{Census FT.76}(1),(3),(7)];
\]

%dsp-level:0
\[
\textbf{priced and}\ \neg\textbf{adopted:}\ \text{the variant choice}\ (\mathrm A)/(\mathrm B)\ \text{is}\ \neg\text{a read adoption}\ \ [\mathrm{BM.30}(5)\ \text{stands}];
\]

%dsp-level:0
\[
\qquad \text{the census reports both};
\]

%dsp-level:0
\[
\neg\text{a P4 addition},
\]

%dsp-level:0
\[
\quad\neg\text{a new dimensionful scale}\ [\mathrm{P1}],
\]

%dsp-level:0
\[
\quad\neg\text{a measured value}\ [\mathrm{P5};\ \text{machine: zero external tokens, zero decimal numerals in this tier}],
\]

%dsp-level:0
\[
\quad\text{§Cal-BM stays empty}.
\]

\textbf{Req.} D2.9 (Parts 1–2); P1, P4, P5 (Parts 1–2); Def-H.21(6), Def-H.37(2), FT.32(6), FT.45, FT.47(1), FT.54(4), FT.55(1), FT.60(5), FT.79(2), FT.84(2) (Part 3); Census FT.42(5), FT.43(1), FT.63(1), FT.63(4), FT.63(6), FT.76(1), FT.76(2), FT.76(3), FT.76(4), FT.76(5), FT.76(7), FT.76(9) (row-wise); Def-BM.12(4), Def-BM.14(3), BM.30(5), BM.60(3), BM.61(8) (this part); exact radical rings $\mathbb Q(\sqrt2),\mathbb Q(\sqrt3),\mathbb Q(\sqrt5)$, the quartic ring $\mathbb Z[\mathsf c]/(16\mathsf c^4{-}20\mathsf c^2{+}5)$, Chebyshev recursion with integer coefficients, de Moivre $\cos5t$, minimal polynomials by the multiplication map and by first linear dependence, exact dyadic radical enclosures with integer certificates, half-angle and sine-ratio identities [math; machine: \texttt{checks/\allowbreak{}verify\_\allowbreak{}H\_\allowbreak{}mass\_\allowbreak{}ratio\_\allowbreak{}census.py} §A, §B].

\textbf{Pf.}

%dsp-level:0
\[
(1):\ \delta_v/\pi=2-d(L-2)/L\ \text{at the five}\ (d,L)\ \text{pairs}\ [\text{machine: exact rationals in}\ \pi\text{-units}];
\]

%dsp-level:0
\[
\qquad V\delta_v\ \text{multiplied out at all five}\ [\text{machine}];
\]

%dsp-level:0
\[
\qquad 4\sin^2(x)=2-2\cos(2x)\ \text{at}\ x=m\delta_v/4\ [\text{machine: the exact series identity in each field}];
\]

%dsp-level:0
\[
\qquad \cos(m\delta_v/2)=T_m(\cos(\delta_v/2))\ [\text{Chebyshev};\ \text{machine}]
\]

%dsp-level:0
\[
\Theta_R=4\pi:\ D(m)=2-2\cos(2m\pi)=0\ \forall m\ [\text{machine}];\qquad\Theta_R=0:\ D\equiv0\ [\text{machine}]
\]

%dsp-level:0
\[
(2):\ k=L\kappa_{\min}/2\ \text{at}\ L=(3,4,3,5,3)\implies(3,2,3,5,3);
\]

%dsp-level:0
\[
\qquad L_{vf}=(3,3,4,3,5)\implies(3,3,2,3,5)\ [\text{machine: exact Fractions}];
\]

%dsp-level:0
\[
\qquad K_S/F=(12,12,24,60,60)/(4,6,8,12,20)=(3,2,3,5,3)\ [\text{machine}]
\]

%dsp-level:0
\[
(3):\ \text{the two multiplets evaluated weight by weight from}\ (1)\ \text{in}\ \mathbb Q,
\]

%dsp-level:0
\[
\ \mathbb Q(\sqrt2),
\]

%dsp-level:0
\[
\ \mathbb Q(\sqrt3),
\]

%dsp-level:0
\[
\ \mathbb Q[\mathsf c]\ [\text{machine}];
\]

%dsp-level:0
\[
\qquad \cos\tfrac\pi8=\tfrac{\sqrt{2-\sqrt2}}2\ \text{not needed:}\ D(2)\big|_{\delta_v=\pi/2}=2\ [\text{machine}];
\]

%dsp-level:0
\[
\qquad \text{the}\ (\mathrm B)\ \text{row by the identical pipeline at}\ k_{vf}\ [\text{machine}]
\]

%dsp-level:0
\[
(4):\ T_{m+1}
\]

%dsp-level:0
\[
= 2\mathsf cT_m-T_{m-1}\ \text{iterated to}\ m
\]

%dsp-level:0
\[
= 5\ \text{with reduction}\ 16\mathsf c^4
\]

%dsp-level:0
\[
\mapsto 20\mathsf c^2-5\ [\text{machine: exact integer-coefficient arithmetic in the quotient ring}];
\]

%dsp-level:0
\[
\qquad T_5=16\mathsf c^5-20\mathsf c^3+5\mathsf c=\mathsf c(16\mathsf c^4-20\mathsf c^2+5)=0\ [\text{machine}];
\]

%dsp-level:0
\[
\qquad 4\mathsf c^2-2=\tfrac{\sqrt5-1}2\implies D(2)=4-4\mathsf c^2=\tfrac{3-\sqrt5}2,
\]

%dsp-level:0
\[
\quad D(4)
\]

%dsp-level:0
\[
= 5-4\mathsf c^2
\]

%dsp-level:0
\[
= \tfrac{5-\sqrt5}2\ [\text{machine:}\ 8\mathsf c^2-5=\sqrt5\ \text{verified by squaring both positive sides in}\ \mathbb Q(\sqrt5)]
\]

%dsp-level:0
\[
\text{the minimal polynomials:}\ x^4-8x^3+19x^2-12x+1\ \text{at}\ D(1),D(3);
\]

%dsp-level:0
\[
\quad x^2-3x+1\ \text{at}\ D(2);
\]

%dsp-level:0
\[
\quad x^2-5x+5\ \text{at}\ D(4);
\]

%dsp-level:0
\[
\quad x-2\ \text{at}\ D(5);
\]

%dsp-level:0
\[
\quad x\ \text{at}\ D(0)\ [\text{machine: first linear dependence over}\ \mathbb Q\ \text{on the powers of the element, with exact re-substitution}];
\]

%dsp-level:0
\[
\qquad [\mathbb Q(\mathsf c):\mathbb Q]=4,\ \sqrt5=8\mathsf c^2-5\ [\text{machine}]
\]

%dsp-level:0
\[
(5):\ P_S=4/\!(\delta_v/\pi)=(4,8,6,20,12)\ [\text{machine}];
\]

%dsp-level:0
\[
\qquad D(m+P_S)=D(m)\ \text{and}\ D(-m)=D(m)\ [\text{machine: the ladder evaluated over two full periods}];
\]

%dsp-level:0
\[
\qquad \#\text{distinct}=P_S/2+1\ [\text{machine: the exact value sets}];
\]

%dsp-level:0
\[
\qquad \max_mD(m)=4\ \text{at}\ m=P_S/2\ [\text{machine}]
\]

%dsp-level:0
\[
\text{the rung law}:\ D(m)\in\{0,2,4\}\iff\cos(m\delta_v/2)\in\{1,0,-1\}
\]

%dsp-level:0
\[
\iff m\delta_v\in\pi\mathbb Z\ [\text{machine: both directions on all five ladders}];
\]

%dsp-level:0
\[
\qquad k_f\delta_v=(3\pi,\pi,2\pi,\pi,\pi)\ [\text{machine}]\implies D(k_f)=(2,2,4,2,2)\ [\text{machine}]
\]

%dsp-level:0
\[
(6):\ (\mathrm a)\ \delta_v=\pi\implies D(m)=2-2\cos(m\pi/2)\in\{0,2,4\}\ \text{and}\ D(1)=D(3)=2\ [\text{machine}];
\]

%dsp-level:0
\[
\qquad (\mathrm b)\ D(0)=2-2\cos0=0\ \text{exactly}\ [\text{machine}];
\]

%dsp-level:0
\[
\qquad k_f(\text{cube})=2,\ k_{vf}(\text{octa})=2\ \text{are even}\ [\text{machine}]
\]

%dsp-level:0
\[
(7):\ r_{\mathrm{intra}}
\]

%dsp-level:0
\[
= \sqrt{D(k_f)/D(1)}
\]

%dsp-level:0
\[
= \sqrt{4\sin^2(k_f\delta_v/4)/4\sin^2(\delta_v/4)}
\]

%dsp-level:0
\[
= \sin(k_f\delta_v/4)/\sin(\delta_v/4)\ [\text{machine: positive sines on}\ (0,\pi/2]];
\]

%dsp-level:0
\[
\text{octa}:\ \sqrt{4/1}=2\ [\text{machine}];
\]

%dsp-level:0
\[
\qquad \text{icosa}:\ \frac{2}{2-\sqrt3}=2(2+\sqrt3)=4+2\sqrt3=(1+\sqrt3)^2\ [\text{machine: exact in}\ \mathbb Q(\sqrt3)];
\]

%dsp-level:0
\[
\qquad \text{dodeca}:\ z:
\]

%dsp-level:0
\[
= \frac{2}{2-2\mathsf c}\ \text{satisfies}\ z^4-24z^3+76z^2-64z+16
\]

%dsp-level:0
\[
= 0\ [\text{machine: the substitution}\ D=2/z\ \text{into}\ D^4-8D^3+19D^2-12D+1],
\]

%dsp-level:0
\[
\ \text{so}\ r\ \text{satisfies}\ r^8-24r^6+76r^4-64r^2+16
\]

%dsp-level:0
\[
\text{cube}\ (\mathrm B):\ \frac{2+\sqrt2}{2-\sqrt2}=\frac{(2+\sqrt2)^2}{2}=3+2\sqrt2=(1+\sqrt2)^2\ [\text{machine}];
\]

%dsp-level:0
\[
\qquad \text{icosa}\ (\mathrm B):\ \frac{2+\sqrt3}{2-\sqrt3}=(2+\sqrt3)^2=7+4\sqrt3=(2+\sqrt3)^2\ [\text{machine}]
\]

%dsp-level:0
\[
(8):\ 48<49\implies\sqrt3<\tfrac74;
\]

%dsp-level:0
\[
\qquad 5\cdot2500
\]

%dsp-level:0
\[
= 12500>12321
\]

%dsp-level:0
\[
= 111^2\implies\sqrt5>\tfrac{111}{50}\implies\mathsf c^2
\]

%dsp-level:0
\[
= \tfrac{5+\sqrt5}8>\tfrac{361}{400}\implies\mathsf c>\tfrac{19}{20}\ [\text{machine}];
\]

%dsp-level:0
\[
\qquad (4+\tfrac72)(2-\tfrac{19}{10})=\tfrac{15}2\cdot\tfrac1{10}=\tfrac34<2\ [\text{machine}]
\]

%dsp-level:0
\[
(9):\ w:
\]

%dsp-level:0
\[
= \tfrac{4}{2-2\mathsf c}
\]

%dsp-level:0
\[
= 2z\ \text{and}\ z^4-24z^3+76z^2-64z+16
\]

%dsp-level:0
\[
= 0
\]

%dsp-level:0
\[
\implies w^4-48w^3+304w^2-512w+256
\]

%dsp-level:0
\[
= 0\ [\text{machine: the substitution}\ z=w/2\ \text{cleared},\ \text{and independently by first linear dependence}];
\]

%dsp-level:0
\[
\qquad \tfrac{2}{\sqrt{2-2\mathsf c}}=\sqrt2\sqrt{\tfrac{2}{2-2\mathsf c}}\ [\text{machine: squaring both positive sides}]
\]

%dsp-level:0
\[
(10):\ \tfrac{4}{2-2\mathsf c}<49\iff4<98-98\mathsf c\iff98\mathsf c<94\iff\mathsf c<\tfrac{47}{49}\ [\text{machine}];
\]

%dsp-level:0
\[
\qquad \mathsf c^2
\]

%dsp-level:0
\[
= \tfrac{5+\sqrt5}8<\tfrac{2209}{2401}
\]

%dsp-level:0
\[
\iff 2401(5+\sqrt5)<17672
\]

%dsp-level:0
\[
\iff 2401\sqrt5<5667
\]

%dsp-level:0
\[
\iff 2401^2\cdot5<5667^2\ [\text{machine:}\ 28824005<32114889];
\]

%dsp-level:0
\[
\qquad 7<\tfrac{89}{12}\iff84<89\ [\text{machine}];
\]

%dsp-level:0
\[
\qquad \tfrac{89}{12}<z_0^{(q)}/z_1^{(q)}\ [\mathrm{BM.80}(4)\text{'s exact enclosure};\ \text{machine}]
\]

%dsp-level:0
\[
(11):\ \text{the standing-rule scans}\ [\text{machine}].\ \blacksquare
\]

\textbf{Phys.} One law generates the whole Platonic mass spectrum: the cone defect at a vertex, read as an angle, squared into a mass by a single undeclared scale. Wind that angle by an integer and the ladder appears — periodic, even, monotone to a hard ceiling of four, with the tetrahedron degenerate at every weight and the even windings carrying an exactly massless middle rung. The dodecahedron is the deep row: its tower of Chebyshev polynomials closes at the fifth rung because the defining quartic of the eighteen-degree cosine \textit{is} the statement that five times eighteen degrees is a right angle, so the heaviest weight of that solid sits at exactly two, forced by the same relation that makes its lightest weight the smallest positive defect angle in the census. The intra-generational ratios are then pure sines — one sine over another, at rational multiples of pi — and they come out as one, one, two, root three plus one, and a degree-eight number built from nested square roots of five. Their order is exact and their widest member is the dodecahedron's, in both winding conventions. Nothing here is a mass hierarchy in the channel sense: the widest ratio the entire Platonic ladder can supply at any winding is under seven, while the flattest family channel supplies over one hundred thirty-nine, and the two numbers are compared as arithmetic, never equated as physics — the range bar that deficit twelve records is re-certified with the census's own exact values, and the ceiling that bounds it is a theorem about angles, not about the world.

\pseudohead{theorem}{\texorpdfstring{BM.80}{BM.80}}{\texorpdfstring{The inter-generational mass ratios: one cubic form, three channel constants, the exact enclosures, and the hierarchy-order theorem}{The inter-generational mass ratios: one cubic form, three channel constants, the exact enclosures, and the hierarchy-order theorem}}\label{tag:BM-80}
\textbf{Thm.}

%dsp-level:0
\[
(1)\quad\textbf{the three channel angles and the exact angle ledger:}
\]

%dsp-level:0
\[
\qquad \delta_q=\frac{2}{3\cdot3\cdot3}=\frac2{27},
\]

%dsp-level:0
\[
\qquad \delta_\ell=\frac29,
\]

%dsp-level:0
\[
\qquad \delta_\nu=\frac2{16}=\frac18
\]

%dsp-level:0
\[
\qquad [\mathrm{FT.119}(1);\ \mathrm{FT.112};\ \mathrm{FT.142}(5);\ \text{machine: exact Fractions}]:
\]

%dsp-level:0
\[
\boxed{\ 27\,\delta_q=9\,\delta_\ell=16\,\delta_\nu=2,\qquad3\delta_q=\delta_\ell,\qquad\delta_\nu:\delta_\ell:\delta_q=27:48:16,\qquad\delta_q<\delta_\nu<\delta_\ell\ }
\]

%dsp-level:0
\[
[\mathrm{FT.119}(2);\ \mathrm{FT.143}(5);\ \text{machine}];
\]

%dsp-level:0
\[
\qquad 3\delta_\nu=\frac38\notin\Bigl\{\frac29,\frac2{27}\Bigr\},
\]

%dsp-level:0
\[
\qquad 3\delta_\ell=\frac23,
\]

%dsp-level:0
\[
\qquad 3\delta_q=\frac29
\]

%dsp-level:0
\[
\qquad [\mathrm{FT.143}(2),(5);\ \text{machine}];
\]

%dsp-level:0
\[
(2)\quad\textbf{the spectrum and its cubic (frozen, re-certified):}
\]

%dsp-level:0
\[
\qquad r^{(\mathrm{ch})}_i=1+\sqrt2\cos\Bigl(\delta_{\mathrm{ch}}+\frac{2\pi i}3\Bigr),
\]

%dsp-level:0
\[
\qquad i=0,1,2
\]

%dsp-level:0
\[
\qquad [\mathrm{FT.111}(3);\ \mathrm{FT.113};\ \mathrm{FT.119}(3);\ \mathrm{FT.143}(1);\ \mathrm{FT.158}(4)]:
\]

%dsp-level:0
\[
\text{with}\ c_{\mathrm{ch}}=\cos\delta_{\mathrm{ch}},\ s_{\mathrm{ch}}=\sin\delta_{\mathrm{ch}}:
\]

%dsp-level:0
\[
\qquad r_0=1+\sqrt2\,c_{\mathrm{ch}},
\]

%dsp-level:0
\[
\qquad r_1=1-\frac{\sqrt2\,c_{\mathrm{ch}}+\sqrt6\,s_{\mathrm{ch}}}2,
\]

%dsp-level:0
\[
\qquad r_2=1-\frac{\sqrt2\,c_{\mathrm{ch}}-\sqrt6\,s_{\mathrm{ch}}}2
\]

%dsp-level:0
\[
\qquad [\text{machine}];
\]

%dsp-level:0
\[
\sum_ir_i=3,
\]

%dsp-level:0
\[
\qquad \sum_{i<j}r_ir_j=\frac32,
\]

%dsp-level:0
\[
\qquad \sum_ir_i^2=6,
\]

%dsp-level:0
\[
\qquad \boxed{\ 2(r-1)^3-3(r-1)-\gamma_{\mathrm{ch}}=0,\qquad\gamma_{\mathrm{ch}}:=\sqrt2\cos(3\delta_{\mathrm{ch}})\ }
\]

%dsp-level:0
\[
\qquad [\text{machine: the depressed cubic at}\ B^2=2,\ \text{and the triple angle}\ \cos3\delta=4\cos^3\delta-3\cos\delta];
\]

%dsp-level:0
\[
\text{equivalently, in}\ r\ \text{alone:}
\]

%dsp-level:0
\[
\qquad \boxed{\ 2r^3-6r^2+3r+(1-\gamma_{\mathrm{ch}})=0\ }
\]

%dsp-level:0
\[
\qquad [\text{machine: the expansion of}\ 2(r{-}1)^3{-}3(r{-}1){=}\gamma_{\mathrm{ch}}];
\]

%dsp-level:0
\[
(3)\quad\textbf{the mass-squared cubic (derived here):}
\]

%dsp-level:0
\[
\qquad z:=r^2,
\]

%dsp-level:0
\[
\qquad m_i^{(\mathrm{ch})}=\mu_{\mathrm{ch}}\,z_i^{(\mathrm{ch})},
\]

%dsp-level:0
\[
\qquad \mu_{\mathrm{ch}}\ \text{undeclared}\ [\mathrm{P4};\ \text{Cal-F}]:
\]

%dsp-level:0
\[
\boxed{\ 4z^3-24z^2+\bigl(21-12\,\gamma_{\mathrm{ch}}\bigr)z-\bigl(\gamma_{\mathrm{ch}}-1\bigr)^2\ =\ 0\qquad\text{for each}\ \ z=z_i^{(\mathrm{ch})}\ }
\]

%dsp-level:0
\[
[\text{machine: exact elimination of}\ r\ \text{in}\ \mathbb Q[\gamma][r]/\bigl(2r^3-6r^2+3r+1-\gamma\bigr):\quad r(2z+3)=6z+\gamma-1\ \text{at}\ z=r^2,\ \text{and}\ z(2z+3)^2-(6z+\gamma-1)^2\equiv0\ \text{in the quotient ring}];
\]

%dsp-level:0
\[
\textbf{Vieta:}
\]

%dsp-level:0
\[
\qquad \boxed{\ \sum_iz_i=6\qquad(\text{phase-free and channel-free}),\qquad\sum_{i<j}z_iz_j=\frac{21-12\gamma_{\mathrm{ch}}}4,\qquad\prod_iz_i=\frac{(\gamma_{\mathrm{ch}}-1)^2}4\ }
\]

%dsp-level:0
\[
[\text{machine: the three symmetric reads of the cubic};\ \ \sum_iz_i=6=\sum_ir_i^2\ \text{agrees with}\ (2)\ \text{exactly}];
\]

%dsp-level:0
\[
\textbf{one cubic form, three constants:}
\]

%dsp-level:0
\[
\qquad \text{the polynomial shape is channel-free};
\]

%dsp-level:0
\[
\qquad \text{the entire channel dependence is the single number}\ \gamma_{\mathrm{ch}};
\]

%dsp-level:0
\[
(4)\quad\textbf{the three constants are pairwise distinct and exactly ordered:}
\]

%dsp-level:0
\[
\qquad \gamma_\ell=\sqrt2\cos\frac23,
\]

%dsp-level:0
\[
\qquad \gamma_\nu=\sqrt2\cos\frac38,
\]

%dsp-level:0
\[
\qquad \gamma_q=\sqrt2\cos\frac29:
\]

%dsp-level:0
\[
\boxed{\ \gamma_\ell\ <\ \gamma_\nu\ <\ \gamma_q\ <\ \sqrt2\ }
\]

%dsp-level:0
\[
\qquad [\text{machine: disjoint exact enclosures, width}<2^{-60};\qquad\cos\ \text{strictly decreasing on}\ [0,\pi]\ \text{and}\ \tfrac23>\tfrac38>\tfrac29\ [\text{math}];\qquad\gamma<\sqrt2\iff\cos3\delta<1\iff3\delta\ne0];
\]

%dsp-level:0
\[
\gamma_\ell\in\Bigl(\frac{1111412422687001}{10^{15}},\frac{1111412422687002}{10^{15}}\Bigr),
\]

%dsp-level:0
\[
\quad\gamma_\nu\in\Bigl(\frac{1315936498799931}{10^{15}},\frac{1315936498799932}{10^{15}}\Bigr),
\]

%dsp-level:0
\[
\quad\gamma_q\in\Bigl(\frac{1379438171633255}{10^{15}},\frac{1379438171633256}{10^{15}}\Bigr)
\]

%dsp-level:0
\[
[\text{machine: alternating rational Taylor enclosures of}\ \cos\ \text{with exact remainder bounds, width}<2^{-160}];
\]

%dsp-level:0
\[
\textbf{the three-root regime is exact:}
\]

%dsp-level:0
\[
\qquad 2y^3-3y-\gamma\ \text{has three distinct real roots}\iff\gamma^2<2
\]

%dsp-level:0
\[
\qquad [\text{the discriminant}\ \tfrac{27}4(2-\gamma^2)>0;\ \text{machine}]\ —\ \text{and all three channels satisfy it by}\ (4)\text{'s last inequality};
\]

%dsp-level:0
\[
(5)\quad\textbf{the cross-channel exact identity:}
\]

%dsp-level:0
\[
\qquad \boxed{\ r^{(\ell)}_0\ =\ 1+\gamma_q\ }
\]

%dsp-level:0
\[
\qquad [3\delta_q=\delta_\ell\ [\mathrm{FT.119}(2)];\qquad r^{(\ell)}_0=1+\sqrt2\cos\delta_\ell=1+\sqrt2\cos3\delta_q=1+\gamma_q;\qquad\text{machine: the same cosine, enclosed to width}<2^{-160}]:
\]

%dsp-level:0
\[
\textbf{the}\ \ell\textbf{-channel's heaviest line is the q-channel's cubic constant plus one};
\]

%dsp-level:0
\[
\qquad \text{the shared object is the angle}\ 3\delta_q=\delta_\ell,\ \neg\text{the cubic constant}:
\]

%dsp-level:0
\[
\gamma_\ell=\sqrt2\cos\tfrac23\ne\sqrt2\cos\tfrac29=\gamma_q
\]

%dsp-level:0
\[
\qquad [\text{machine: disjoint enclosures}];
\]

%dsp-level:0
\[
\qquad \gamma_\nu=\sqrt2\cos\tfrac38\ \text{is channel-own}
\]

%dsp-level:0
\[
\qquad [\mathrm{FT.143}(2);\ \text{machine}];
\]

%dsp-level:0
\[
(6)\quad\textbf{the inter-generational ratio census (the exact enclosures):}
\]

%dsp-level:0
\[
\qquad z_1<z_2<z_0\ \ \text{in every channel}
\]

%dsp-level:0
\[
\qquad [\text{machine: disjoint enclosures, widths}<2^{-50}]:
\]

%tbl-level:3
{\scriptsize\setlength{\tabcolsep}{4pt}
\begin{longtable}{>{\raggedright\arraybackslash}p{\dimexpr0.0698\textwidth-2\tabcolsep\relax}|>{\raggedright\arraybackslash}p{\dimexpr0.2552\textwidth-2\tabcolsep\relax}|>{\raggedright\arraybackslash}p{\dimexpr0.2666\textwidth-2\tabcolsep\relax}|>{\raggedright\arraybackslash}p{\dimexpr0.2099\textwidth-2\tabcolsep\relax}|>{\raggedright\arraybackslash}p{\dimexpr0.1985\textwidth-2\tabcolsep\relax}}
\toprule
channel & $z_2/z_1$ & $z_0/z_1$ & $z_2/z_0$ & $z_1/z_0$ \\
\addlinespace[2pt]
\midrule
\endfirsthead
\toprule
channel & $z_2/z_1$ & $z_0/z_1$ & $z_2/z_0$ & $z_1/z_0$ \\
\addlinespace[2pt]
\midrule
\endhead
\bottomrule
\endlastfoot
$q$ & $\bigl(\tfrac{3563695074946}{10^{12}},\allowbreak \tfrac{3563695074947}{10^{12}}\bigr)$ & $\bigl(\tfrac{139338688028416}{10^{12}},\allowbreak \tfrac{139338688028417}{10^{12}}\bigr)$ & $\bigl(\tfrac{25575776012}{10^{12}},\allowbreak \tfrac{25575776013}{10^{12}}\bigr)$ & $\bigl(\tfrac{7176757684}{10^{12}},\allowbreak \tfrac{7176757685}{10^{12}}\bigr)$ \\
$\ell$ & $\bigl(\tfrac{206770315972729}{10^{12}},\allowbreak \tfrac{206770315972730}{10^{12}}\bigr)$ & $\bigl(\tfrac{3477472837104598}{10^{12}},\allowbreak \tfrac{3477472837104599}{10^{12}}\bigr)$ & $\bigl(\tfrac{59459937045}{10^{12}},\allowbreak \tfrac{59459937046}{10^{12}}\bigr)$ & $\bigl(\tfrac{287565150}{10^{12}},\allowbreak \tfrac{287565151}{10^{12}}\bigr)$ \\
$\nu$ & $\bigl(\tfrac{9583926992379}{10^{12}},\allowbreak \tfrac{9583926992380}{10^{12}}\bigr)$ & $\bigl(\tfrac{271994748749143}{10^{12}},\allowbreak \tfrac{271994748749144}{10^{12}}\bigr)$ & $\bigl(\tfrac{35235705970}{10^{12}},\allowbreak \tfrac{35235705971}{10^{12}}\bigr)$ & $\bigl(\tfrac{3676541567}{10^{12}},\allowbreak \tfrac{3676541568}{10^{12}}\bigr)$ \\
\end{longtable}
}

%dsp-level:0
\[
[\text{machine: every bracket a monotone exact interval division at width}<2^{-150},\ \text{each bracket decided by exact}\ 10^{12}\text{-scaling};\qquad\text{the frozen decimal rows of}\ \mathrm{FT.119}(5),(6),\ \mathrm{FT.113},\ \mathrm{FT.143}(3)\ \text{are re-certified digit for digit against these brackets}];
\]

%dsp-level:0
\[
m_1:m_2:m_3=1:\frac{z_2}{z_1}:\frac{z_0}{z_1}
\]

%dsp-level:0
\[
\qquad [\mathrm{FT.119}(5)\ \text{vocabulary};\ \text{the ratios are}\ \mu_{\mathrm{ch}}\text{-free}];
\]

%dsp-level:0
\[
\qquad \textbf{the ratios are}\ \neg\textbf{algebraic numbers:}\ \gamma_{\mathrm{ch}}
\]

%dsp-level:0
\[
= \sqrt2\cos(3\delta_{\mathrm{ch}})\ \text{at the rational radian arguments}\ \tfrac29,
\]

%dsp-level:0
\[
\tfrac38,
\]

%dsp-level:0
\[
\tfrac2{27}\ \text{is a transcendental cosine},
\]

%dsp-level:0
\[
\text{so the exact form of every inter-generational ratio is the closed radical-trigonometric expression of}\ (2),
\]

%dsp-level:0
\[
(3),
\]

%dsp-level:0
\[
\ \text{and its exact comparison content is the rational bracket of}\ (6)
\]

%dsp-level:0
\[
\qquad [\text{machine: the brackets, }\neg\text{a decimal assertion}];
\]

%dsp-level:0
\[
(7)\quad\textbf{the phase-free invariants (all three channels, exact):}
\]

%dsp-level:0
\[
\qquad \boxed{\ Q\ :=\ \frac{\sum_im_i}{\bigl(\sum_i\sqrt{m_i}\bigr)^2}\ \overset{m_i=\mu z_i}{=}\ \frac{\sum_iz_i}{\bigl(\sum_ir_i\bigr)^2}\ =\ \frac69\ =\ \frac23\ }
\]

%dsp-level:0
\[
[\mathrm{FT.119}(4);\ \mathrm{FT.143}(4);\ \mathrm{FT.158}(4);\ \mathrm{Def\text{-}BM.6}(4);\ \text{machine: independent of}\ \delta_{\mathrm{ch}}\ \text{and of}\ \mu_{\mathrm{ch}}];
\]

%dsp-level:0
\[
\qquad \sum_iz_i=6=3A_0\big|_{A_0=\frac{11}6}\ —\ \neg\text{a calibration: the sum rule is}\ \mathrm{BM.60}(8)\text{'s}\ \text{phase-free row};
\]

%dsp-level:0
\[
(8)\quad\textbf{the hierarchy-order theorem (derived, scale-free):}
\]

%dsp-level:0
\[
\qquad \text{on}\ \theta\in(0,\tfrac\pi{12}):
\]

%dsp-level:0
\[
\qquad r_1(\theta)
\]

%dsp-level:0
\[
= 1+\sqrt2\cos(\theta+\tfrac{2\pi}3)\ \text{is strictly decreasing and strictly positive},
\]

%dsp-level:0
\[
\ r_1(\tfrac\pi{12})
\]

%dsp-level:0
\[
= 0\ \text{exactly};
\]

%dsp-level:0
\[
r_2(\theta)\ \text{is strictly increasing};\qquad r_0(\theta)\ \text{is strictly decreasing with}\ r_0/r_1\ \text{strictly increasing}
\]

%dsp-level:0
\[
[\text{machine: the derivative signs}\ -\sqrt2\sin(\theta{+}\tfrac{2\pi}3)<0,\ +\sqrt2\sin(\theta{+}\tfrac\pi3)>0,\ -\sqrt2\sin\theta<0;\qquad\frac{d}{d\theta}\ln\frac{r_0}{r_1}=\sqrt2\Bigl(\frac{\sin(\pi/3-\theta)}{r_1}-\frac{\sin\theta}{r_0}\Bigr)>0
\]

%dsp-level:0
\[
\text{on}\ (0,\tfrac\pi6)\ \text{since}\ \sin(\tfrac\pi3-\theta)>\sin\theta\iff\tan\theta<\tfrac1{\sqrt3}\iff\theta<\tfrac\pi6\ \ \text{and}\ \ \ 0<r_1<r_0
\]

%dsp-level:0
\[
\qquad [\text{math};\ \text{machine: the algebraic reductions}]];
\]

%dsp-level:0
\[
\delta_q<\delta_\nu<\delta_\ell<\frac\pi{12}
\]

%dsp-level:0
\[
\qquad \Bigl[\pi^2>6\sum_{n\le10}\frac1{n^2}=\frac{1968329}{211680}>\frac{64}9\implies\pi>\frac83\implies\frac29<\frac\pi{12};\quad\text{machine:}\ 1968329\cdot9=17714961>13547520=64\cdot211680\Bigr];
\]

%dsp-level:0
\[
\boxed{\ 1\ <\ \frac{z^{(q)}_2}{z^{(q)}_1}\ <\ \frac{z^{(\nu)}_2}{z^{(\nu)}_1}\ <\ \frac{z^{(\ell)}_2}{z^{(\ell)}_1},\qquad1\ <\ \frac{z^{(q)}_0}{z^{(q)}_1}\ <\ \frac{z^{(\nu)}_0}{z^{(\nu)}_1}\ <\ \frac{z^{(\ell)}_0}{z^{(\ell)}_1}\ }
\]

%dsp-level:0
\[
\textbf{the q-channel is the flattest family hierarchy, the}\ \ell\textbf{-channel the steepest, and the}\ \nu\textbf{-channel lies between them}\ —\ \text{an exact ordering with no scale and no import}
\]

%dsp-level:0
\[
\qquad [\text{machine: monotonicity and the disjoint enclosures of}\ (6)\ \text{concur}];
\]

%dsp-level:0
\[
\textbf{the near-vanishing row:}
\]

%dsp-level:0
\[
\qquad r_1(\theta)=0\iff\theta=\tfrac\pi{12}\ \text{exactly}
\]

%dsp-level:0
\[
\qquad [\cos\tfrac{3\pi}4=-\tfrac1{\sqrt2};\ \text{machine}]\implies\frac\pi{12}-\delta_\ell=\frac{3\pi-8}{36}\ \text{is the smallest of the three margins}:
\]

%dsp-level:0
\[
\textbf{the}\ \ell\textbf{-channel's inner line is the closest to the zero locus, and its steepness is a consequence of}\ \delta_\ell\ \textbf{being the largest angle below}\ \tfrac\pi{12};
\]

%dsp-level:0
\[
(9)\quad\textbf{the bars:}
\]

%dsp-level:0
\[
\qquad \mu_q,\ \mu_\ell,\ \mu_\nu\ \text{are undeclared P4/Cal-F scales and enter no ratio}
\]

%dsp-level:0
\[
\qquad [\mathrm{P4};\ \text{Cal-F};\ \text{§Cal-BM stays empty}];
\]

%dsp-level:0
\[
\neg\exists\ \text{an SM mass, mass ratio or Yukawa value in any clause}
\]

%dsp-level:0
\[
\qquad [\mathrm{P5};\ \text{machine: zero external tokens, zero decimal numerals in this tier}];
\]

%dsp-level:0
\[
\text{the}\ m^2\ \text{vocabulary is the corpus's own inertial read}\ [\mathrm{L3}(4),\ \mathrm{L4}];
\]

%dsp-level:0
\[
\qquad \text{the channel labels}\ \ell,q,\nu\ \text{name ledger rows}\ [\mathrm{FT.112}\ \mathrm{Rmk};\ \mathrm{FT.143}(6)];
\]

%dsp-level:0
\[
\qquad \text{amplitude}\ne\text{mass}\ [\mathrm{FT.79}(5)];
\]

%dsp-level:0
\[
\qquad \text{rung}\ne\text{scale}\ [\mathrm{D0.8}];
\]

%dsp-level:0
\[
(10)\quad\textbf{status:}
\]

%dsp-level:0
\[
\qquad \textbf{derived:}\ (1)\text{--}(8);
\]

%dsp-level:0
\[
\qquad \textbf{re-cited:}\ (2)\text{'s spectrum and}\ (7)\text{'s invariant from their own tags, with}\ (3),(5),(8)\ \text{new};
\]

%dsp-level:0
\[
\neg\text{a P4 addition},\quad\neg\text{a new dimensionful scale}\ [\mathrm{P1}],\quad\neg\text{a measured value}\ [\mathrm{P5}],\quad\neg\text{a declaration},\quad
\]

%dsp-level:0
\[
\qquad \textbf{open:}\ \text{the three channel scales}\.
\]

\textbf{Req.} D0.8 (Parts 1–2); P1, P4, P5 (Parts 1–2); L3(4), L4 (Parts 1–2); FT.79(5), FT.111(2), FT.111(3), FT.112, FT.113, FT.119(1), FT.119(2), FT.119(3), FT.119(4), FT.119(5), FT.119(6), FT.142(5), FT.143(1), FT.143(2), FT.143(3), FT.143(4), FT.143(5), FT.143(6), FT.158(4) (Part 3); Def-BM.6(4), BM.60(8), BM.61(8) (this part); Cal-F; exact rational arithmetic, exact symbolic elimination in $\mathbb Q[\gamma][r]/(2r^3{-}6r^2{+}3r{+}1{-}\gamma)$, Vieta reads, alternating rational Taylor enclosures of $\cos$ and $\sin$ with exact remainder bounds, integer Newton certificates for $\sqrt2$ and $\sqrt6$, monotone exact interval division, the Basel-series lower bound for $\pi^2$, derivative-sign reductions in $\mathbb Q[\sqrt2,\sqrt3][c,s]/(c^2{+}s^2{-}1)$ [math; machine: \texttt{checks/\allowbreak{}verify\_\allowbreak{}H\_\allowbreak{}mass\_\allowbreak{}ratio\_\allowbreak{}census.py} §D].

\textbf{Pf.}

%dsp-level:0
\[
(1):\ \delta_q=2/(3\cdot3\cdot3),\ \delta_\ell=2/(3\cdot3),\ \delta_\nu=2/16\ [\text{machine: exact Fractions}];
\]

%dsp-level:0
\[
\qquad 27\cdot\tfrac2{27}=9\cdot\tfrac29=16\cdot\tfrac18=2\ [\text{machine}];
\]

%dsp-level:0
\[
\qquad 3\cdot\tfrac2{27}=\tfrac29\ [\text{machine}];
\]

%dsp-level:0
\[
\qquad \tfrac2{27}<\tfrac18\iff16<27,\quad\tfrac18<\tfrac29\iff9<16\ [\text{machine}];
\]

%dsp-level:0
\[
\qquad \tfrac2{27}:\tfrac18:\tfrac29=16:27:48\ [\times216;\ \text{machine}]
\]

%dsp-level:0
\[
(2):\ r_i=A+B\cos(\delta+\tfrac{2\pi i}3)\ \text{at}\ A=1,\ B=\sqrt2\ [\mathrm{FT.111}(3)\ \text{re-cited}];
\]

%dsp-level:0
\[
\qquad \cos(\delta+\tfrac{2\pi}3)=-\tfrac c2-\tfrac{\sqrt3}2s,\quad\cos(\delta+\tfrac{4\pi}3)=-\tfrac c2+\tfrac{\sqrt3}2s\ [\text{addition formulas}];
\]

%dsp-level:0
\[
\qquad \sqrt2\cdot\tfrac{\sqrt3}2=\tfrac{\sqrt6}2\ [\text{machine}]
\]

%dsp-level:0
\[
\sum_i\cos(\delta+\tfrac{2\pi i}3)
\]

%dsp-level:0
\[
= 0,
\]

%dsp-level:0
\[
\quad\sum_i\cos^2(\delta+\tfrac{2\pi i}3)
\]

%dsp-level:0
\[
= \tfrac32\ [\mathrm{FT.111}(3)\ \text{identities}]\implies\sum_ir_i
\]

%dsp-level:0
\[
= 3,
\]

%dsp-level:0
\[
\ \sum_ir_i^2
\]

%dsp-level:0
\[
= 3+2\cdot\tfrac32
\]

%dsp-level:0
\[
= 6\ [\text{machine}];
\]

%dsp-level:0
\[
\qquad \text{the depressed cubic}\ y^3-\tfrac32y-\tfrac{\sqrt2}2\cos3\delta=0\ \text{at}\ y=r-1,\ \times2\ [\text{machine: the ring identity}]
\]

%dsp-level:0
\[
(3):\ 2(r-1)^3-3(r-1)=\gamma\iff2r^3-6r^2+3r+(1-\gamma)=0\ [\text{machine: expansion}];
\]

%dsp-level:0
\[
\qquad 2r^3+3r=6z-(1-\gamma)\ \text{at}\ z=r^2\implies r(2z+3)=6z+\gamma-1\ [\text{machine}]
\]

%dsp-level:0
\[
z(2z+3)^2-(6z+\gamma-1)^2
\]

%dsp-level:0
\[
= 4z^3-24z^2+(21-12\gamma)z-(\gamma-1)^2\ \text{as an identity in}\ \mathbb Q[\gamma,z]\ [\text{machine: cleared coefficient by coefficient}];
\]

%dsp-level:0
\[
\qquad \text{the three roots are}\ z_i
\]

%dsp-level:0
\[
= r_i^2\ \text{because}\ r_i>0\ \text{and}\ 2z+3
\]

%dsp-level:0
\[
\ne 0\ \text{on}\ z>0\ [\text{machine: disjoint positive enclosures}];
\]

%dsp-level:0
\[
\qquad \text{Vieta}: -\tfrac{-24}4=6,\ \tfrac{21-12\gamma}4,\ \tfrac{(\gamma-1)^2}4\ [\text{machine}]
\]

%dsp-level:0
\[
(4):\ \gamma_\ell,
\]

%dsp-level:0
\[
\gamma_\nu,
\]

%dsp-level:0
\[
\gamma_q\ \text{enclosed by alternating rational Taylor series at}\ \tfrac23,
\]

%dsp-level:0
\[
\tfrac38,
\]

%dsp-level:0
\[
\tfrac29\ \text{with exact remainder bounds, multiplied by the Newton-certified}\ \sqrt2\ [\text{machine: widths}<2^{-160}];
\]

%dsp-level:0
\[
\qquad \text{the three intervals are pairwise disjoint}\ [\text{machine}];
\]

%dsp-level:0
\[
\qquad \tfrac23>\tfrac38\iff16>9,\quad\tfrac38>\tfrac29\iff27>16\ [\text{machine}];
\]

%dsp-level:0
\[
\qquad \text{the discriminant of}\ 2y^3-3y-\gamma:\ \ \Delta
\]

%dsp-level:0
\[
= \tfrac{27}4(2-\gamma^2)>0\iff\gamma^2<2\ [\text{machine: the exact cubic discriminant}]
\]

%dsp-level:0
\[
(5):\ 3\delta_q
\]

%dsp-level:0
\[
= \delta_\ell\ [\mathrm{FT.119}(2)\ \text{re-cited}]\implies\cos3\delta_q
\]

%dsp-level:0
\[
= \cos\delta_\ell\implies\gamma_q
\]

%dsp-level:0
\[
= \sqrt2\cos\delta_\ell
\]

%dsp-level:0
\[
= r^{(\ell)}_0-1\ [\text{machine: the two enclosures agree to width}<2^{-160}\ \text{and the symbolic forms are identical}];
\]

%dsp-level:0
\[
\qquad \gamma_\ell\ne\gamma_q\ [\text{disjoint enclosures}];
\]

%dsp-level:0
\[
\qquad \cos\tfrac38\ne\cos\tfrac29\ [\text{disjoint enclosures}]
\]

%dsp-level:0
\[
(6):\ z_1<z_2<z_0\ [\text{machine: disjoint enclosures of the three roots}];
\]

%dsp-level:0
\[
\qquad \text{the twelve brackets by monotone exact interval division, each numerator and denominator enclosed to width}<2^{-160}\ [\text{machine}];
\]

%dsp-level:0
\[
\qquad \text{the frozen decimal rows of Part 3 re-derived to}\ 12\ \text{places and compared digit for digit}\ [\text{machine}]
\]

%dsp-level:0
\[
(7):\ Q=\tfrac{\sum_iz_i}{(\sum_ir_i)^2}=\tfrac69=\tfrac23\ [\text{machine: exact Fractions}];
\]

%dsp-level:0
\[
\qquad \text{phase-independence: the identities}\ \sum_i\cos(\delta+\tfrac{2\pi i}3)
\]

%dsp-level:0
\[
= 0,
\]

%dsp-level:0
\[
\ \sum_i\cos^2
\]

%dsp-level:0
\[
= \tfrac32\ \text{hold at every}\ \delta\ [\text{machine at six rational phases}]
\]

%dsp-level:0
\[
(8):\ r_1'(\theta)
\]

%dsp-level:0
\[
= -\sqrt2\sin(\theta+\tfrac{2\pi}3)
\]

%dsp-level:0
\[
= -\sqrt2\sin(\tfrac\pi3-\theta)<0\ \text{on}\ (0,\tfrac\pi3)\ [\text{machine: the reduction}\ \sin(\theta+\tfrac{2\pi}3)=\sin(\tfrac\pi3-\theta)];
\]

%dsp-level:0
\[
\qquad r_1(\tfrac\pi{12})=1+\sqrt2\cos\tfrac{3\pi}4=1-1=0\ [\text{machine:}\ \cos\tfrac{3\pi}4=-\tfrac1{\sqrt2}]
\]

%dsp-level:0
\[
r_2'(\theta)=-\sqrt2\sin(\theta+\tfrac{4\pi}3)=+\sqrt2\sin(\theta+\tfrac\pi3)>0\ [\text{machine}];
\]

%dsp-level:0
\[
\qquad r_0'(\theta)=-\sqrt2\sin\theta<0\ \text{on}\ (0,\pi)\ [\text{machine}];
\]

%dsp-level:0
\[
\qquad (\ln(r_0/r_1))'
\]

%dsp-level:0
\[
= \sqrt2\bigl(\tfrac{\sin(\pi/3-\theta)}{r_1}-\tfrac{\sin\theta}{r_0}\bigr)\ [\text{machine: the exact derivative in the quotient ring}]
\]

%dsp-level:0
\[
\sin(\tfrac\pi3-\theta)>\sin\theta\iff\tfrac{\sqrt3}2\cos\theta-\tfrac12\sin\theta>\sin\theta\iff\sqrt3\cos\theta>3\sin\theta\iff\tan\theta<\tfrac1{\sqrt3}\iff\theta<\tfrac\pi6\ [\text{machine: each equivalence as an exact identity;}\ \tan\ \text{strictly increasing on}\ (0,\tfrac\pi2),\ \tan\tfrac\pi6=\tfrac1{\sqrt3}\ [\text{math}]];
\]

%dsp-level:0
\[
r_1<r_0\iff\cos(\theta+\tfrac{2\pi}3)<\cos\theta\ [\cos\ \text{strictly decreasing on}\ (0,\pi);\ \text{math}]
\]

%dsp-level:0
\[
\implies (\ln(r_0/r_1))'>0\ \text{on}\ (0,\tfrac\pi{12})\subset(0,\tfrac\pi6)\ [\text{machine}]
\]

%dsp-level:0
\[
\pi^2
\]

%dsp-level:0
\[
= 6\sum_{n\ge1}n^{-2}>6\sum_{n\le10}n^{-2}
\]

%dsp-level:0
\[
= \tfrac{1968329}{211680}\ [\text{Basel};\ \text{machine: the exact rational partial sum}];
\]

%dsp-level:0
\[
\qquad \tfrac{1968329}{211680}>\tfrac{64}9
\]

%dsp-level:0
\[
\iff 17714961>13547520\ [\text{machine}]\implies\pi>\tfrac83
\]

%dsp-level:0
\[
\implies 3\pi>8\implies\tfrac29<\tfrac\pi{12}\ [\text{machine}]
\]

%dsp-level:0
\[
\text{the ordering}: \delta_q<\delta_\nu<\delta_\ell<\tfrac\pi{12}\ \text{and both ratio maps strictly increasing on}\ (0,\tfrac\pi{12})\implies\text{the boxed chain}\ [\text{machine}];
\]

%dsp-level:0
\[
\qquad \text{concur: the three enclosures of}\ (6)\ \text{are strictly ordered}\ [\text{machine}]
\]

%dsp-level:0
\[
(9)\text{--}(10):\ \text{the standing-rule scans and the assembly}\ [\text{machine}].\ \blacksquare
\]

\textbf{Phys.} Three families, one polynomial. Every channel's mass-squared ratios are the three roots of the same cubic — four z cubed minus twenty-four z squared plus twenty-one minus twelve gamma, times z, minus gamma minus one squared — and the whole difference between the quark channel, the lepton channel and the neutrino channel is a single number: root two times the cosine of three times the channel's own angle. Those three numbers are exact, distinct, and ordered, and one identity ties two of them together: the lepton channel's heaviest line is precisely one more than the quark channel's cubic constant, because tripling the quark angle lands exactly on the lepton angle. The hierarchies then order themselves with no scale anywhere in the argument: below one twelfth of a turn the inner line shrinks, the middle line grows, and both ratios steepen monotonically, so the smallest angle gives the flattest family and the largest angle the steepest — quark, neutrino, lepton — with the lepton channel's inner line the closest to the exact zero at one twelfth of a turn, which is why its hierarchy is two orders of magnitude deep while the quark channel's is barely two. The brackets are exact rational intervals, not decimals: every digit the corpus froze in Part 3 is re-derived here from the angle alone and certified at widths far below the last quoted place. Nothing is fitted, no scale is fixed, and the ratios are what they always were — pure numbers waiting on three undeclared scales before they can be called masses.

\pseudohead{theorem}{\texorpdfstring{BM.81}{BM.81}}{\texorpdfstring{The six-state mass-ratio census: the splitting order, the doublet ratios, the two inverted hierarchies, and the three non-proportional ladders}{The six-state mass-ratio census: the splitting order, the doublet ratios, the two inverted hierarchies, and the three non-proportional ladders}}\label{tag:BM-81}
\textbf{Thm.}

%dsp-level:0
\[
(1)\quad\textbf{the six states in both currencies (the census table):}
\]

%dsp-level:0
\[
\qquad |\mathcal S|=3\cdot2=6,
\]

%dsp-level:0
\[
\qquad m_{\mathrm{hid}}=\frac{2E\ln2}{\mathfrak h},
\]

%dsp-level:0
\[
\qquad m^2_{\mathrm{pole}}=\hat\mu'^2D_S
\]

%dsp-level:0
\[
\qquad [\mathrm{Def\text{-}BM.14}(1),(3);\ \mathrm{FT.44}(4);\ \mathrm{Def\text{-}H.37}(2)]:
\]

%tbl-level:1
{\small\setlength{\tabcolsep}{4pt}
\begin{longtable}{>{\raggedright\arraybackslash}p{\dimexpr0.0717\textwidth-2\tabcolsep\relax}|>{\raggedright\arraybackslash}p{\dimexpr0.0811\textwidth-2\tabcolsep\relax}|>{\raggedright\arraybackslash}p{\dimexpr0.1720\textwidth-2\tabcolsep\relax}|>{\raggedright\arraybackslash}p{\dimexpr0.0811\textwidth-2\tabcolsep\relax}|>{\raggedright\arraybackslash}p{\dimexpr0.1720\textwidth-2\tabcolsep\relax}|>{\raggedright\arraybackslash}p{\dimexpr0.0472\textwidth-2\tabcolsep\relax}|>{\raggedright\arraybackslash}p{\dimexpr0.1282\textwidth-2\tabcolsep\relax}|>{\raggedright\arraybackslash}p{\dimexpr0.1470\textwidth-2\tabcolsep\relax}|>{\raggedright\arraybackslash}p{\dimexpr0.0999\textwidth-2\tabcolsep\relax}}
\toprule
state & $(\mathcal F_g,\allowbreak \beta)$ & the branch datum & $D_S$ & minimal polynomial of $D_S$ & $2E$ & $m_{\mathrm{hid}}\mathfrak h/\ln2$ & $m^2_{\mathrm{pole}}/\hat\mu'^2$ & field \\
\addlinespace[2pt]
\midrule
\endfirsthead
\toprule
state & $(\mathcal F_g,\allowbreak \beta)$ & the branch datum & $D_S$ & minimal polynomial of $D_S$ & $2E$ & $m_{\mathrm{hid}}\mathfrak h/\ln2$ & $m^2_{\mathrm{pole}}/\hat\mu'^2$ & field \\
\addlinespace[2pt]
\midrule
\endhead
\bottomrule
\endlastfoot
$s_1$ & $(\mathcal F_1,\allowbreak -)$ & the ladder rung $m{=}1$ & $2$ & $x-2$ & $12$ & $12$ & $2$ & $\mathbb Q$ \\
$s_2$ & $(\mathcal F_1,\allowbreak +\allowbreak )$ & the ladder rung $m{=}2$ & $4$ & $x-4$ & $12$ & $12$ & $4$ & $\mathbb Q$ \\
$s_3$ & $(\mathcal F_2,\allowbreak -)$ & the cube member & $2-\sqrt2$ & $x^2-4x+\allowbreak 2$ & $24$ & $24$ & $2-\sqrt2$ & $\mathbb Q(\sqrt2)$ \\
$s_4$ & $(\mathcal F_2,\allowbreak +\allowbreak )$ & the octahedron member & $1$ & $x-1$ & $24$ & $24$ & $1$ & $\mathbb Q$ \\
$s_5$ & $(\mathcal F_3,\allowbreak -)$ & the dodecahedron member & $2-2\mathsf c$ & $x^4-8x^3+\allowbreak 19x^2-12x+\allowbreak 1$ & $60$ & $60$ & $2-2\mathsf c$ & $\mathbb Q(\mathsf c)$ \\
$s_6$ & $(\mathcal F_3,\allowbreak +\allowbreak )$ & the icosahedron member & $2-\sqrt3$ & $x^2-4x+\allowbreak 1$ & $60$ & $60$ & $2-\sqrt3$ & $\mathbb Q(\sqrt3)$ \\
\end{longtable}
}

%dsp-level:0
\[
[\mathrm{Def\text{-}BM.14}(2),(3);\ \mathrm{BM.60}(4);\ \text{machine: every minimal polynomial re-derived by first linear dependence over}\ \mathbb Q];
\]

%dsp-level:0
\[
\textbf{the defect species resolves all six states with zero degeneracy};
\]

%dsp-level:0
\[
\qquad \textbf{the informational species is branch-degenerate and takes three values}\ 12:24:60=1:2:5
\]

%dsp-level:0
\[
\qquad [\mathrm{FT.44}(1);\ \mathrm{BM.60}(4),(6);\ \text{machine}];
\]

%dsp-level:0
\[
(2)\quad\textbf{the branch order and the attached splitting (frozen, re-certified):}
\]

%dsp-level:0
\[
\qquad \boxed{\ D_S(\mathcal F_g,+)>D_S(\mathcal F_g,-)\qquad\forall g\in\{1,2,3\}\ }
\]

%tbl-level:0
{\setlength{\tabcolsep}{4pt}
\begin{longtable}{>{\raggedright\arraybackslash}p{\dimexpr0.0949\textwidth-2\tabcolsep\relax}|>{\raggedright\arraybackslash}p{\dimexpr0.1331\textwidth-2\tabcolsep\relax}|>{\raggedright\arraybackslash}p{\dimexpr0.1331\textwidth-2\tabcolsep\relax}|>{\raggedright\arraybackslash}p{\dimexpr0.6389\textwidth-2\tabcolsep\relax}}
\toprule
floor & $D_S(\mathcal F_g,\allowbreak -)$ & $D_S(\mathcal F_g,\allowbreak +\allowbreak )$ & the exact certificate \\
\addlinespace[2pt]
\midrule
\endfirsthead
\toprule
floor & $D_S(\mathcal F_g,\allowbreak -)$ & $D_S(\mathcal F_g,\allowbreak +\allowbreak )$ & the exact certificate \\
\addlinespace[2pt]
\midrule
\endhead
\bottomrule
\endlastfoot
$\mathcal F_1$ & $2$ & $4$ & $4>\allowbreak 2$ \\
$\mathcal F_2$ & $2-\sqrt2$ & $1$ & $1>\allowbreak 2-\sqrt2\iff\allowbreak \sqrt2>\allowbreak 1\iff\allowbreak 2>\allowbreak 1$ \\
$\mathcal F_3$ & $2-2\mathsf c$ & $2-\sqrt3$ & $\mathsf c>\allowbreak \tfrac{\sqrt3}2\iff\allowbreak \mathsf c^2>\allowbreak \tfrac34\iff\allowbreak \tfrac{5+\sqrt5}8>\allowbreak \tfrac34\iff\allowbreak \sqrt5>\allowbreak 1\iff\allowbreak 5>\allowbreak 1$ \\
\end{longtable}
}

%dsp-level:0
\[
\Delta_g\ :=\ m^2_{\mathrm{down}}-m^2_{\mathrm{up}}\ =\ \hat\mu'^2\bigl(D_S(\mathcal F_g,+)-D_S(\mathcal F_g,-)\bigr)
\]

%dsp-level:0
\[
\qquad [\mathrm{Def\text{-}BM.14}(5);\ \mathrm{BM.75}(4);\ \mathrm{Def\text{-}BM.14}(8)\text{'s name rows}]:
\]

%dsp-level:0
\[
\boxed{\ \Delta_1=2\,\hat\mu'^2,\qquad\Delta_2=(\sqrt2-1)\,\hat\mu'^2,\qquad\Delta_3=(2\mathsf c-\sqrt3)\,\hat\mu'^2,\qquad\Delta_g>0\ \ \forall g\ }
\]

%dsp-level:0
\[
\text{minimal polynomials of}\ \Delta_g/\hat\mu'^2:
\]

%dsp-level:0
\[
\qquad x-2,
\]

%dsp-level:0
\[
\qquad x^2+2x-1,
\]

%dsp-level:0
\[
\qquad x^8-22x^6+119x^4-38x^2+1
\]

%dsp-level:0
\[
\qquad [\text{machine: first linear dependence in}\ \mathbb Q,\ \mathbb Q(\sqrt2),\ \mathbb Q(\mathsf c,\sqrt3)];
\]

%dsp-level:0
\[
(3)\quad\textbf{the splitting order (derived here):}
\]

%dsp-level:0
\[
\qquad \boxed{\ \Delta_1\ >\ \Delta_2\ >\ \Delta_3\ >\ 0\ }
\]

%dsp-level:0
\[
\qquad \textbf{the down--up splitting strictly decreases across the generations in the defect currency}
\]

%dsp-level:0
\[
\Delta_1>\Delta_2\iff2>\sqrt2-1\iff3>\sqrt2\iff9>2;\qquad\Delta_2>\Delta_3\iff\sqrt2+\sqrt3-1>2\mathsf c
\]

%dsp-level:0
\[
\iff7+4\sqrt6>4\sqrt2+4\sqrt3+\sqrt5
\]

%dsp-level:0
\[
\qquad [\text{both sides positive; squared}]
\]

%dsp-level:0
\[
\qquad \iff15+6\sqrt6>2\sqrt{10}+2\sqrt{15}
\]

%dsp-level:0
\[
\qquad [\text{squared again}]
\]

%dsp-level:0
\[
\iff341+140\sqrt6>0
\]

%dsp-level:0
\[
\qquad [\text{machine: the three squaring steps as exact identities in}\ \mathbb Q(\sqrt2,\sqrt3,\sqrt5,\sqrt6),\ \text{each side certified positive by exact dyadic enclosure}];
\]

%dsp-level:0
\[
\textbf{the exact ratio census of the splitting:}
\]

%dsp-level:0
\[
\qquad \frac{\Delta_2}{\Delta_1}=\frac{\sqrt2-1}2,
\]

%dsp-level:0
\[
\qquad \frac{\Delta_3}{\Delta_1}=\mathsf c-\frac{\sqrt3}2,
\]

%dsp-level:0
\[
\qquad \frac{\Delta_3}{\Delta_2}=\frac{2\mathsf c-\sqrt3}{\sqrt2-1}=(\sqrt2+1)(2\mathsf c-\sqrt3)
\]

%dsp-level:0
\[
\frac{\Delta_2}{\Delta_1}\in\Bigl(\frac{2071}{10^4},\frac{2072}{10^4}\Bigr),
\]

%dsp-level:0
\[
\qquad \frac{\Delta_3}{\Delta_1}\in\Bigl(\frac{85031}{10^6},\frac{85032}{10^6}\Bigr),
\]

%dsp-level:0
\[
\qquad \frac{\Delta_3}{\Delta_2}\in\Bigl(\frac{410566}{10^6},\frac{410567}{10^6}\Bigr)
\]

%dsp-level:0
\[
\qquad [\text{machine: exact dyadic radical enclosures, width}<2^{-60}];
\]

%dsp-level:0
\[
(4)\quad\textbf{the relative splittings (exact, and one exact simplification):}
\]

%dsp-level:0
\[
\qquad \frac{\Delta_g}{m^2_{\mathrm{up}}(g)}=\frac{D_S(\mathcal F_g,+)}{D_S(\mathcal F_g,-)}-1,
\]

%dsp-level:0
\[
\qquad \frac{\Delta_g}{m^2_{\mathrm{down}}(g)}=1-\frac{D_S(\mathcal F_g,-)}{D_S(\mathcal F_g,+)}:
\]

%dsp-level:0
\[
\boxed{\ \frac{\Delta_1}{m^2_{\mathrm{up}}(1)}=1,\qquad\frac{\Delta_2}{m^2_{\mathrm{up}}(2)}=\frac{\sqrt2}2,\qquad\frac{\Delta_3}{m^2_{\mathrm{up}}(3)}=\frac{2\mathsf c-\sqrt3}{2-2\mathsf c}\ }
\]

%dsp-level:0
\[
\qquad \Bigl[\frac{\sqrt2-1}{2-\sqrt2}=\frac{\sqrt2}2\ \text{exactly}:\ (\sqrt2-1)-\frac{\sqrt2}2(2-\sqrt2)=0\ \ \text{in}\ \mathbb Q(\sqrt2);\ \text{machine}\Bigr];
\]

%dsp-level:0
\[
\frac{\Delta_1}{m^2_{\mathrm{down}}(1)}=\frac12,
\]

%dsp-level:0
\[
\qquad \frac{\Delta_2}{m^2_{\mathrm{down}}(2)}=\sqrt2-1,
\]

%dsp-level:0
\[
\qquad \frac{\Delta_3}{m^2_{\mathrm{down}}(3)}=(2\mathsf c-\sqrt3)(2+\sqrt3)=4\mathsf c+2\sqrt3\,\mathsf c-2\sqrt3-3
\]

%dsp-level:0
\[
\qquad [\text{machine: exact in the biquadratic ring}];
\]

%dsp-level:0
\[
(5)\quad\textbf{the doublet mass ratios (the per-generation down/up census):}
\]

%dsp-level:0
\[
\qquad \boxed{\ \frac{m^2_{\mathrm{down}}(g)}{m^2_{\mathrm{up}}(g)}=\Bigl(2,\ \ \frac{2+\sqrt2}2,\ \ \frac{2-\sqrt3}{2-2\mathsf c}\Bigr)\ }
\]

%dsp-level:0
\[
\text{minimal polynomials:}\qquad x-2,\qquad2x^2-4x+1,\qquad x^8-48x^7+410x^6-1328x^5+1899x^4-1232x^3+330x^2-32x+1
\]

%dsp-level:0
\[
[\text{machine: first linear dependence in}\ \mathbb Q,\ \mathbb Q(\sqrt2),\ \mathbb Q(\mathsf c,\sqrt3);\ \text{the degree-}8\ \text{row's constant term is}\ 1=\mathrm{Norm}(2-\sqrt3)/\mathrm{Norm}(2-2\mathsf c)\ [\text{machine}]];
\]

%dsp-level:0
\[
\boxed{\ \frac{m_{\mathrm{down}}(g)}{m_{\mathrm{up}}(g)}=\Bigl(\sqrt2,\ \ \sqrt{\tfrac{2+\sqrt2}2},\ \ \sqrt{\tfrac{2-\sqrt3}{2-2\mathsf c}}\Bigr)\ }
\]

%dsp-level:0
\[
\qquad \text{annihilators:}\quad x^2-2,\quad2x^4-4x^2+1,\quad P_8(x^2)
\]

%dsp-level:0
\[
P_8(x^2)=x^{16}-48x^{14}+410x^{12}-1328x^{10}+1899x^8-1232x^6+330x^4-32x^2+1
\]

%dsp-level:0
\[
\qquad [\text{machine: the exact substitution};\ \neg\text{minimality claimed at}\ g=3];
\]

%dsp-level:0
\[
\textbf{the sine-ratio form (uniform):}
\]

%dsp-level:0
\[
\qquad \frac{m_{\mathrm{down}}(g)}{m_{\mathrm{up}}(g)}
\]

%dsp-level:0
\[
= \frac{\sin\bigl(\Theta_+/4\bigr)}{\sin\bigl(\Theta_-/4\bigr)}
\]

%dsp-level:0
\[
= \Bigl(\frac{\sin(\pi/2)}{\sin(\pi/4)},\ \ \frac{\sin(\pi/6)}{\sin(\pi/8)},\ \ \frac{\sin(\pi/12)}{\sin(\pi/20)}\Bigr)
\]

%dsp-level:0
\[
\Theta_\pm\ :
\]

%dsp-level:0
\[
= \ m_\pm\,
\]

%dsp-level:0
\[
\delta_v(\mathcal F_g,\pm)\ \ \text{the branch's own angle defect, at}\ \ (m_+,m_-)
\]

%dsp-level:0
\[
= (2,1),
\]

%dsp-level:0
\[
\ (1,1),
\]

%dsp-level:0
\[
\ (1,1)\ \ \text{and}\ \ \delta_v
\]

%dsp-level:0
\[
= \pi;
\]

%dsp-level:0
\[
\ \tfrac{2\pi}3,
\]

%dsp-level:0
\[
\tfrac\pi2;
\]

%dsp-level:0
\[
\ \tfrac\pi3,\qquad \tfrac\pi5
\]

%dsp-level:0
\[
[\mathrm{BM.79}(7)\text{'s law applied to the branch pair};\ \text{machine:}\ D_S=4\sin^2(\Theta/4)\ \text{at each of the six states},\ \text{and}\ \sin^2(\pi/8)=\tfrac{2-\sqrt2}4,\ \sin^2(\pi/12)=\tfrac{2-\sqrt3}4\ \text{by the half-angle identity}];
\]

%dsp-level:0
\[
(6)\quad\textbf{the order of the doublet ratios is}\ \neg\textbf{the generation order:}
\]

%dsp-level:0
\[
\qquad \boxed{\ \frac{m_{\mathrm{down}}(2)}{m_{\mathrm{up}}(2)}\ <\ \frac{m_{\mathrm{down}}(1)}{m_{\mathrm{up}}(1)}\ <\ \frac{m_{\mathrm{down}}(3)}{m_{\mathrm{up}}(3)}\ }
\]

%dsp-level:0
\[
\qquad \text{i.e.}\quad\mathcal F_2\ \text{is the tightest doublet,}\ \mathcal F_3\ \text{the widest}:
\]

%dsp-level:0
\[
\frac{2+\sqrt2}2<2\iff\sqrt2<2\iff2<4;
\]

%dsp-level:0
\[
\qquad 2<\frac{2-\sqrt3}{2-2\mathsf c}\iff4-4\mathsf c<2-\sqrt3\iff\sqrt3<4\mathsf c-2\iff0<16\mathsf c^2-16\mathsf c+1
\]

%dsp-level:0
\[
\iff11+2\sqrt5>16\mathsf c\iff141+44\sqrt5>160+32\sqrt5\iff12\sqrt5>19\iff720>361
\]

%dsp-level:0
\[
\qquad [\text{machine: the chain of exact equivalences, each squaring step certified positive}];
\]

%dsp-level:0
\[
\text{enclosures:}\quad\frac{m^2_{\mathrm{down}}(2)}{m^2_{\mathrm{up}}(2)}\in\Bigl(\frac{17071}{10^4},\frac{17072}{10^4}\Bigr),
\]

%dsp-level:0
\[
\qquad \frac{m^2_{\mathrm{down}}(3)}{m^2_{\mathrm{up}}(3)}\in\Bigl(\frac{2737332}{10^6},\frac{2737333}{10^6}\Bigr)
\]

%dsp-level:0
\[
\qquad [\text{machine}];
\]

%dsp-level:0
\[
(7)\quad\textbf{the two inverted hierarchies (the census form):}
\]

%dsp-level:0
\[
\qquad \textbf{informational:}\quad m_{\mathrm{hid}}(\mathcal F_1)<m_{\mathrm{hid}}(\mathcal F_2)<m_{\mathrm{hid}}(\mathcal F_3),
\]

%dsp-level:0
\[
\qquad 1:2:5
\]

%dsp-level:0
\[
\qquad [\mathrm{BM.60}(6);\ \text{machine}];
\]

%dsp-level:0
\[
\textbf{defect:}
\]

%dsp-level:0
\[
\qquad D_S(\mathcal F_3,\pm)<D_S(\mathcal F_2,\pm)<D_S(\mathcal F_1,\pm),
\]

%dsp-level:0
\[
\qquad 2-2\mathsf c<2-\sqrt3<2-\sqrt2<1<2<4
\]

%dsp-level:0
\[
\qquad [\mathrm{BM.60}(6);\ \text{machine: the certificates}\ 5>1,\ 3>2,\ 2>1];
\]

%dsp-level:0
\[
\boxed{\ \textbf{the two currencies order the three floors oppositely}\ }
\]

%dsp-level:0
\[
\qquad [\mathrm{BM.60}(6);\ \mathrm{BM.39}(5)\ \text{pattern}];
\]

%dsp-level:0
\[
\qquad \textbf{the widest derived ratio in each currency:}
\]

%dsp-level:0
\[
\text{defect}\ (m^2):\ \ \frac{4}{2-2\mathsf c}\in\Bigl(\frac{40863458}{10^6},\frac{40863459}{10^6}\Bigr);
\]

%dsp-level:0
\[
\qquad \text{defect}\ (m):\ \ \frac{2}{\sqrt{2-2\mathsf c}}\in\Bigl(\frac{6392453}{10^6},\frac{6392454}{10^6}\Bigr);
\]

%dsp-level:0
\[
\qquad \text{informational}\ (m):\ \ 5,
\]

%dsp-level:0
\[
\qquad (m^2):\ \ 25
\]

%dsp-level:0
\[
\qquad [\text{machine}];
\]

%dsp-level:0
\[
(8)\quad\textbf{the three ladders are pairwise non-proportional (derived):}
\]

%dsp-level:0
\[
\qquad \textbf{(i) the floor ladder against the channel ladder:}
\]

%dsp-level:0
\[
\neg\exists\,\mu\in\mathbb R_{>0}:\quad\bigl(m_{\mathrm{hid}}(\mathcal F_1),m_{\mathrm{hid}}(\mathcal F_2),m_{\mathrm{hid}}(\mathcal F_3)\bigr)=\mu\bigl(z^{(q)}_1,z^{(q)}_2,z^{(q)}_0\bigr)
\]

%dsp-level:0
\[
[\text{proportionality forces}\ 2=z^{(q)}_2/z^{(q)}_1\ \text{and}\ 5=z^{(q)}_0/z^{(q)}_1;\ \ \mathrm{BM.80}(6)\text{'s brackets give}\ z^{(q)}_2/z^{(q)}_1>\tfrac72>2\ \text{and}\ z^{(q)}_0/z^{(q)}_1>139>5;\ \text{machine}];
\]

%dsp-level:0
\[
\textbf{(ii) the channel ladder against the transport ladder:}
\]

%dsp-level:0
\[
\qquad z^{(\mathrm{ch})}_i=1+2\sqrt2\,\cos\varphi_i+2\cos^2\varphi_i\ \ \text{is quadratic in the channel cosine},
\]

%dsp-level:0
\[
\qquad G\bigl(c(i)\bigr)=A_0-B_0\cos\varphi_i\ \ \text{is affine in it}
\]

%dsp-level:0
\[
[\mathrm{BM.60}(8);\ \mathrm{BM.39}(1);\ \mathrm{FT.111}(3);\ \text{machine}]:
\]

%dsp-level:0
\[
\qquad \text{a quadratic and an affine function of the same three distinct cosines agree at all three only if the quadratic coefficient vanishes};
\]

%dsp-level:0
\[
\text{the quadratic coefficient is exactly}\ 2\ne0
\]

%dsp-level:0
\[
\qquad [\text{machine: the exact polynomial comparison over}\ \mathbb Q[\sqrt2][\cos\varphi]]\implies\boxed{\ \neg\exists\,\lambda:\ z^{(\mathrm{ch})}_i=\lambda\,G\bigl(c(i)\bigr)\ \ \forall i\ };
\]

%dsp-level:0
\[
\textbf{(iii) the floor ladder against the transport ladder:}
\]

%dsp-level:0
\[
\qquad \text{the transport form's line order inverts against the spectrum}\ [\mathrm{BM.39}(5)]\ \text{while the floor ladder is increasing in}\ 2E\ [\mathrm{BM.60}(6)]\implies\neg\text{proportional with a positive scale}
\]

%dsp-level:0
\[
\qquad [\text{machine: the sign census}];
\]

%dsp-level:0
\[
\textbf{the generation map is an order-isomorphism and}\ \neg\textbf{a proportionality}
\]

%dsp-level:0
\[
\qquad [\mathrm{Def\text{-}BM.14}(6);\ \mathrm{BM.60}(7):\ \text{the unique order-isomorphism out of}\ 3!=6\ \text{bijections re-cited}]:
\]

%dsp-level:0
\[
\qquad \text{the map supplies no ratio between the currencies};
\]

%dsp-level:0
\[
(9)\quad\textbf{the six quark states (the assembly row):}
\]

%dsp-level:0
\[
\qquad \mathcal S=\{(\mathcal F_g,\beta)\},
\]

%dsp-level:0
\[
\qquad \text{down-type}:=\beta{=}+\ (\text{the heavier branch}),
\]

%dsp-level:0
\[
\qquad \text{up-type}:=\beta{=}-
\]

%dsp-level:0
\[
\qquad [\mathrm{Def\text{-}BM.14}(8);\ \text{A-BM.6}]:
\]

%dsp-level:0
\[
\textbf{the attachment:}
\]

%dsp-level:0
\[
\qquad q_{\mathrm{core}}=0\longleftrightarrow\text{down-type}\longleftrightarrow\beta{=}+,
\]

%dsp-level:0
\[
\qquad q_{\mathrm{core}}={+}1\longleftrightarrow\text{up-type}\longleftrightarrow\beta{=}-
\]

%dsp-level:0
\[
\qquad [\mathrm{BM.75}(3):\ \text{the unique bijection, conditional on}\ \mathrm{BM.59}(9)\text{'s single-substitution hypothesis, on its face}];
\]

%dsp-level:0
\[
\textbf{the sign-coherence row:}
\]

%dsp-level:0
\[
\qquad \text{in every derived currency in which the doublet splitting is formed, the down-type constituent is the heavier one}
\]

%dsp-level:0
\[
\qquad [\mathrm{BM.75}(5)\ \text{re-cited};\ \text{the four rows:}\ \ln2/\mathfrak h,\ \ln(4/3)/\mathfrak h,\ \tfrac43,\ \Delta_g;\ \text{the single opposed EM correction dominated on the sub-critical branch}\ [\mathrm{BM.46}(4);\ \mathrm{BM.62}(6),(9)]];
\]

%dsp-level:0
\[
\textbf{the flavour composition stays}\ \neg\textbf{derived:}
\]

%dsp-level:0
\[
\qquad \text{which three of the six states a sealed class occupies is}\ \neg\text{attached}\ [\mathrm{BM.28}(7);\ \mathrm{Def\text{-}BM.8}(5);\ \mathrm{BM.75}(6)];
\]

%dsp-level:0
\[
\qquad \text{the mixing matrix is}\ \neg\text{derived}\ [\mathrm{FT.145}(3)];
\]

%dsp-level:0
\[
(10)\quad\textbf{status:}
\]

%dsp-level:0
\[
\qquad \textbf{derived:}\ (1)\text{--}(8);
\]

%dsp-level:0
\[
\qquad \textbf{re-cited:}\ (9)\text{'s attachment, name rows and sign coherence from their own tags};
\]

%dsp-level:0
\[
\qquad \textbf{declared:}\ \mathrm{Def\text{-}BM.14}(8)\text{'s name attachment alone}\ [\text{A-BM.6, logged before use}];
\]

%dsp-level:0
\[
\neg\text{a P4 addition},\quad\neg\text{a new dimensionful scale}\ [\mathrm{P1}],\quad\neg\text{a measured value}\ [\mathrm{P5}],\quad\text{§Cal-BM stays empty};
\]

%dsp-level:0
\[
\qquad \textbf{open:}\ \hat\mu'^2,\ \mu\in\text{Cal-F}\ [\text{the map is an order,}\ \neg\text{a number}].
\]

\textbf{Req.} D0.8 (Parts 1–2); P1, P4, P5 (Parts 1–2); Def-H.37(2), FT.44(1), FT.44(4), FT.111(3), FT.145(3) (Part 3); Census FT.63(6) (row-wise); Def-BM.8(5), Def-BM.12(4), Def-BM.14(1), Def-BM.14(2), Def-BM.14(3), Def-BM.14(5), Def-BM.14(6), Def-BM.14(8), BM.28(7), BM.39(1), BM.39(5), BM.46(4), BM.59(9), BM.60(4), BM.60(6), BM.60(7), BM.60(8), BM.62(6), BM.62(9), BM.75(3), BM.75(4), BM.75(5), BM.75(6), BM.79(7), BM.80(4), BM.80(6) (this part); Cal-F; the declaration addendum A-BM.6 (pointers); exact radical rings $\mathbb Q(\sqrt2),\mathbb Q(\sqrt3),\mathbb Q(\sqrt5),\mathbb Q(\sqrt6),\mathbb Q(\sqrt{10}),\mathbb Q(\sqrt{15})$, the biquadratic ring $\mathbb Q[\mathsf c,\sqrt3]/(16\mathsf c^4{-}20\mathsf c^2{+}5)$, minimal polynomials by first linear dependence with exact re-substitution, norm reads by multiplication determinants, exact dyadic radical enclosures with integer certificates, squaring chains with certified positive sides, exact polynomial comparison over $\mathbb Q[\sqrt2][x]$ [math; machine: \texttt{checks/\allowbreak{}verify\_\allowbreak{}H\_\allowbreak{}mass\_\allowbreak{}ratio\_\allowbreak{}census.py} §C].

\textbf{Pf.}

%dsp-level:0
\[
(1):\ \text{the table re-cited from}\ \mathrm{Def\text{-}BM.14}(3)\ \text{and}\ \mathrm{BM.60}(4)\ \text{and every minimal polynomial re-derived}\ [\text{machine: first linear dependence on the powers of the element in its own field, followed by exact re-substitution}];
\]

%dsp-level:0
\[
\qquad 2E\in\{12,24,24,60,60\}\ \text{has three distinct values}\ [\text{machine}]
\]

%dsp-level:0
\[
(2):\ \text{the three certificates re-derived}\ [\text{machine}];
\]

%dsp-level:0
\[
\qquad \cos\tfrac\pi{10}>\tfrac{\sqrt3}2\iff\mathsf c^2>\tfrac34\iff\tfrac{5+\sqrt5}8>\tfrac34
\]

%dsp-level:0
\[
\iff 5+\sqrt5>6\iff\sqrt5>1
\]

%dsp-level:0
\[
\iff 5>1\ [\text{machine: exact in}\ \mathbb Q(\sqrt5),\ \text{both sides positive}];
\]

%dsp-level:0
\[
\qquad \Delta_g\ \text{formed as exact differences}\ [\text{machine}]
\]

%dsp-level:0
\[
\text{the minimal polynomial of}\ 2\mathsf c-\sqrt3:\ \text{first linear dependence in}\ \mathbb Q[\mathsf c,\sqrt3]\ \text{gives}\ x^8-22x^6+119x^4-38x^2+1\ [\text{machine}];
\]

%dsp-level:0
\[
\qquad \text{the even-power shape follows from the conjugate set}\ \{\pm2\mathsf c_i\mp\sqrt3\}\ \text{being symmetric under}\ x
\]

%dsp-level:0
\[
\mapsto -x\ [\text{machine}]
\]

%dsp-level:0
\[
(3):\ 9>2\ [\text{machine}];
\]

%dsp-level:0
\[
\qquad \sqrt2+\sqrt3-1>0\ \text{and}\ 2\mathsf c>0\ [\text{machine: dyadic enclosures}]\implies\text{the first squaring is order-preserving};
\]

%dsp-level:0
\[
\qquad ( \sqrt2+\sqrt3-1)^2=6+2\sqrt6-2\sqrt2-2\sqrt3\ \text{and}\ (2\mathsf c)^2=4\mathsf c^2=\tfrac{5+\sqrt5}2\ [\text{machine}]
\]

%dsp-level:0
\[
6+2\sqrt6-2\sqrt2-2\sqrt3>\tfrac{5+\sqrt5}2\iff7+4\sqrt6>4\sqrt2+4\sqrt3+\sqrt5\ [\text{machine: cleared}];
\]

%dsp-level:0
\[
\qquad \text{both sides positive}\ [\text{machine}]\implies\text{the second squaring}: \ 145+56\sqrt6>85+32\sqrt6+8\sqrt{10}+8\sqrt{15}
\]

%dsp-level:0
\[
\iff60+24\sqrt6>8\sqrt{10}+8\sqrt{15}\iff15+6\sqrt6>2\sqrt{10}+2\sqrt{15}\ [\text{machine}];
\]

%dsp-level:0
\[
\qquad \text{both sides positive}\ [\text{machine}]\implies\text{the third squaring}:\ 441+180\sqrt6>100+40\sqrt6
\]

%dsp-level:0
\[
\iff 341+140\sqrt6>0\ [\text{machine: positive coefficients}]
\]

%dsp-level:0
\[
\text{the three ratio enclosures by exact dyadic radical isolation}\ [\text{machine: widths}<2^{-60}]
\]

%dsp-level:0
\[
(4):\ (\sqrt2-1)-\tfrac{\sqrt2}2(2-\sqrt2)=\sqrt2-1-(\sqrt2-1)=0\ [\text{machine: exact in}\ \mathbb Q(\sqrt2)];
\]

%dsp-level:0
\[
\qquad \Delta_g/m^2_{\mathrm{up}}=D_S(+)/D_S(-)-1\ \text{and}\ \Delta_g/m^2_{\mathrm{down}}=1-D_S(-)/D_S(+)\ [\text{machine: the two identities}];
\]

%dsp-level:0
\[
\qquad (2\mathsf c-\sqrt3)(2+\sqrt3)=4\mathsf c+2\sqrt3\mathsf c-2\sqrt3-3\ [\text{machine:}\ (2-\sqrt3)(2+\sqrt3)=1]
\]

%dsp-level:0
\[
(5):\ \text{the three}\ m^2\ \text{ratios formed as exact quotients in their fields}\ [\text{machine}];
\]

%dsp-level:0
\[
\qquad \text{the minimal polynomials by first linear dependence, with exact re-substitution}\ [\text{machine}];
\]

%dsp-level:0
\[
\qquad \mathrm{Norm}(2-\sqrt3)=1,\ \mathrm{Norm}(2-2\mathsf c)=1\ \text{as multiplication determinants}\ [\text{machine}]\implies\text{the constant term}\ 1
\]

%dsp-level:0
\[
\text{the annihilators of the mass ratios}:\ x^2-2\ \text{at}\ \sqrt2;
\]

%dsp-level:0
\[
\qquad \text{at}\ y^2=\tfrac{2+\sqrt2}2:\ (2y^2-2)^2=2\iff2y^4-4y^2+1=0\ [\text{machine}];
\]

%dsp-level:0
\[
\qquad P_8(x^2)\ \text{by the exact substitution}\ x\mapsto x^2\ [\text{machine}]
\]

%dsp-level:0
\[
\text{the sine-ratio form}:\ D_S
\]

%dsp-level:0
\[
= 4\sin^2(\delta_v/4)\ \text{at the six states with}\ \delta_v/\pi
\]

%dsp-level:0
\[
= (1,1,\tfrac12,\tfrac23,\tfrac15,\tfrac13)\implies\delta_v/4\in\{\tfrac\pi4,\tfrac\pi4,\tfrac\pi8,\tfrac\pi6,\tfrac\pi{20},\tfrac\pi{12}\}\ [\text{machine}]
\]

%dsp-level:0
\[
(6):\ 2<4\ [\text{machine}];
\]

%dsp-level:0
\[
\qquad 0<16\mathsf c^2-16\mathsf c+1\iff16\cdot\tfrac{5+\sqrt5}8-16\mathsf c+1>0\iff11+2\sqrt5>16\mathsf c\ [\text{machine}];
\]

%dsp-level:0
\[
\qquad \text{both sides positive}\ [\text{machine}]\implies(11+2\sqrt5)^2>256\mathsf c^2=32(5+\sqrt5)
\]

%dsp-level:0
\[
\iff141+44\sqrt5>160+32\sqrt5\iff12\sqrt5>19\iff720>361\ [\text{machine: each step cleared}];
\]

%dsp-level:0
\[
\qquad \text{the two enclosures}\ [\text{machine}]
\]

%dsp-level:0
\[
(7):\ 12<24<60,\quad12:24:60=1:2:5\ [\text{machine}];
\]

%dsp-level:0
\[
\qquad 2-2\mathsf c<2-\sqrt3\iff\sqrt3<2\mathsf c\ [\mathrm{BM.60}(6)\text{'s certificate}\ 5>1];
\]

%dsp-level:0
\[
\quad 2-\sqrt3<2-\sqrt2\iff\sqrt2<\sqrt3\iff2<3;
\]

%dsp-level:0
\[
\quad 2-\sqrt2<1\iff1<\sqrt2\iff1<2\ [\text{machine}]
\]

%dsp-level:0
\[
\text{the widest ratios}:\ \max D_S/\min D_S=4/(2-2\mathsf c),\quad\max2E/\min2E=60/12=5\ [\text{machine}];
\]

%dsp-level:0
\[
\qquad \text{the two enclosures by exact radical isolation}\ [\text{machine}]
\]

%dsp-level:0
\[
(8)(\mathrm i):\ \text{proportionality at three strictly ordered entries forces the two ratio equalities}\ [\text{machine}];
\]

%dsp-level:0
\[
\qquad z^{(q)}_2/z^{(q)}_1>\tfrac72\ \text{and}\ z^{(q)}_0/z^{(q)}_1>139\ [\mathrm{BM.80}(6)\text{'s brackets};\ \text{machine}];
\]

%dsp-level:0
\[
\qquad \tfrac72>2,\ 139>5\ [\text{machine}]
\]

%dsp-level:0
\[
(\mathrm{ii}):\ r_i=1+\sqrt2\cos\varphi_i\implies z_i=r_i^2=1+2\sqrt2\cos\varphi_i+2\cos^2\varphi_i\ [\text{machine: the exact expansion}];
\]

%dsp-level:0
\[
\qquad \text{if}\ 1+2\sqrt2u+2u^2
\]

%dsp-level:0
\[
= \lambda(A_0-B_0u)\ \text{at the three distinct values}\ u
\]

%dsp-level:0
\[
= \cos\varphi_i,
\]

%dsp-level:0
\[
\ \text{the difference is a quadratic with three roots}
\]

%dsp-level:0
\[
\implies\text{the quadratic vanishes identically}
\]

%dsp-level:0
\[
\implies 2
\]

%dsp-level:0
\[
= 0\ [\text{machine: the exact polynomial identity over}\ \mathbb Q[\sqrt2][u],\ \text{with the three cosines' distinctness certified by disjoint enclosures}];
\]

%dsp-level:0
\[
\qquad (\mathrm{iii}):\ \text{the sign census of the two orders}\ [\mathrm{BM.39}(5);\ \text{machine}]
\]

%dsp-level:0
\[
(9)\text{--}(10):\ \text{the rows re-cited from their own tags};\qquad\text{the standing-rule scans}\ [\text{machine}].\ \blacksquare
\]

\textbf{Phys.} Six states, two currencies, and not one proportionality between them. In the angle currency the six defect values are pairwise distinct algebraic numbers living in four fields, the heaviest branch on each floor is heavier by an exact amount, and those three amounts run strictly downward across the generations — two, root two minus one, and two cosines of eighteen degrees minus root three — so the splitting that separates a doublet is widest on the lightest floor and narrowest on the heaviest, the opposite of what a hierarchy reading would suggest. The doublet ratios themselves are pure sine ratios and their order is again not the generation order: the cube–octahedron pair is the tightest, the tetrahedral pair the middle, and the dodecahedron–icosahedron pair the widest by a degree-eight number whose norm is exactly one. In the bit currency the same six states collapse to three values climbing one to two to five. The two orders are opposite, and the tier proves something stronger than that they differ: no positive scale makes the floor ladder proportional to the channel ladder, because two is not three and a half and five is not one hundred thirty-nine; no scale makes the channel ladder proportional to the transport ladder, because one is quadratic in the channel cosine and the other affine and the quadratic coefficient is exactly two; and the transport ladder inverts against both. The generation map survives all of this as what it always was — the unique order-isomorphism between two three-element ordered sets — and it supplies an order, never a number.

\pseudohead{theorem}{\texorpdfstring{BM.82}{BM.82}}{\texorpdfstring{The bridge census: the six exact factors, the intra-currency closure, and the single priced residue}{The bridge census: the six exact factors, the intra-currency closure, and the single priced residue}}\label{tag:BM-82}
\textbf{Thm.}

%dsp-level:0
\[
(1)\quad\textbf{the two reads and the bridge identity (derived, not assumed):}
\]

%dsp-level:0
\[
\qquad m_{\mathrm{hid}}=\frac{2E\ln2}{\mathfrak h}
\]

%dsp-level:0
\[
\qquad [\mathrm{FT.44}(4):\ \text{the taxonomy law}\ m_{\mathrm{hid}}=E_{\mathrm{int}}\ln2/\mathfrak h\ \text{at}\ E_{\mathrm{int}}=2E;\ \text{Census FT.42}(1)\ \text{row-wise};\ \mathrm{Def\text{-}BM.14}(1)]:
\]

%dsp-level:0
\[
m^2_{\mathrm{pole}}=\hat\mu'^2D_S=4\hat\mu'^2\sin^2\frac{\Theta_R}4\ \ \text{at}\ \Theta_R=\delta_v
\]

%dsp-level:0
\[
\qquad [\mathrm{Def\text{-}H.37}(2);\ \mathrm{FT.54}(4);\ \mathrm{FT.47}(1);\ \mathrm{BM.60}(3)];
\]

%dsp-level:0
\[
\rho_{\mathrm B}\ :=\ \frac{\hat\mu'^2\mathfrak h^2}{(\ln2)^2}
\]

%dsp-level:0
\[
\qquad [\mathrm{BM.69}(2);\ \text{Cal-F packaging}];
\]

%dsp-level:0
\[
\qquad [\rho_{\mathrm B}]=\mathfrak q\cdot\mathfrak q^{-1}=\mathfrak q^0
\]

%dsp-level:0
\[
\qquad [[\hat\mu'^2]=\mathfrak q\ [\mathrm{Def\text{-}H.37}(2)];\ [\mathfrak h]=\mathfrak q^{-1/2}\ [\mathrm{CT2.1};\ \mathrm{D2.9}];\ \text{machine: the exponent arithmetic}];
\]

%dsp-level:0
\[
\boxed{\ \frac{m^2_{\mathrm{pole}}(s)}{m^2_{\mathrm{hid}}(s)}\ =\ \rho_{\mathrm B}\cdot\mathcal R_s,\qquad\mathcal R_s\ :=\ \frac{D_S(s)}{4\,E(s)^2}\ \in\ \overline{\mathbb Q}\ }
\]

%dsp-level:0
\[
\qquad [\mathrm{BM.69}(2)\ \text{re-derived from the two reads above};\ \text{machine: one Laurent monomial in}\ \hat\mu'^2,\mathfrak h,\ln2]:
\]

%dsp-level:0
\[
\textbf{the factorization is exact per state};
\]

%dsp-level:0
\[
\qquad \textbf{the algebraic part}\ \mathcal R_s\ \textbf{is derived};
\]

%dsp-level:0
\[
\qquad \textbf{the single dimensionless residue is}\ \rho_{\mathrm B}
\]

%dsp-level:0
\[
\qquad [\mathrm{BM.69}(5)(\mathrm{R1})];
\]

%dsp-level:0
\[
(2)\quad\textbf{the six-factor census:}\qquad E(s)=(6,6,12,12,30,30)\ \ \text{at}\ \ 2E=(12,12,24,24,60,60):
\]

%tbl-level:0
{\setlength{\tabcolsep}{4pt}
\begin{longtable}{>{\raggedright\arraybackslash}p{\dimexpr0.0859\textwidth-2\tabcolsep\relax}|>{\raggedright\arraybackslash}p{\dimexpr0.0971\textwidth-2\tabcolsep\relax}|>{\raggedright\arraybackslash}p{\dimexpr0.0971\textwidth-2\tabcolsep\relax}|>{\raggedright\arraybackslash}p{\dimexpr0.1649\textwidth-2\tabcolsep\relax}|>{\raggedright\arraybackslash}p{\dimexpr0.1197\textwidth-2\tabcolsep\relax}|>{\raggedright\arraybackslash}p{\dimexpr0.4353\textwidth-2\tabcolsep\relax}}
\toprule
state & $D_S(s)$ & $4E(s)^2$ & $\mathcal R_s=\allowbreak D_S/(4E^2)$ & field & minimal polynomial of $\mathcal R_s$ \\
\addlinespace[2pt]
\midrule
\endfirsthead
\toprule
state & $D_S(s)$ & $4E(s)^2$ & $\mathcal R_s=\allowbreak D_S/(4E^2)$ & field & minimal polynomial of $\mathcal R_s$ \\
\addlinespace[2pt]
\midrule
\endhead
\bottomrule
\endlastfoot
$s_1$ & $2$ & $144$ & $\tfrac1{72}$ & $\mathbb Q$ & $72x-1$ \\
$s_2$ & $4$ & $144$ & $\tfrac1{36}$ & $\mathbb Q$ & $36x-1$ \\
$s_3$ & $2-\sqrt2$ & $576$ & $\tfrac{2-\sqrt2}{576}$ & $\mathbb Q(\sqrt2)$ & $165888x^2-1152x+\allowbreak 1$ \\
$s_4$ & $1$ & $576$ & $\tfrac1{576}$ & $\mathbb Q$ & $576x-1$ \\
$s_5$ & $2-2\mathsf c$ & $3600$ & $\tfrac{2-2\mathsf c}{3600}$ & $\mathbb Q(\mathsf c)$ & $167961600000000\,\allowbreak x^4-373248000000\,\allowbreak x^3+\allowbreak 246240000\,\allowbreak x^2-43200\,\allowbreak x+\allowbreak 1$ \\
$s_6$ & $2-\sqrt3$ & $3600$ & $\tfrac{2-\sqrt3}{3600}$ & $\mathbb Q(\sqrt3)$ & $12960000\,\allowbreak x^2-14400\,\allowbreak x+\allowbreak 1$ \\
\end{longtable}
}

%dsp-level:0
\[
\boxed{\ \text{six states, six exact factors, four fields, six minimal polynomials}\ }
\]

%dsp-level:0
\[
\qquad [\mathrm{BM.69}(2)\text{'s row re-derived};\ \text{machine: the quartic row by the exact substitution}\ D=3600x\ \text{into}\ x^4-8x^3+19x^2-12x+1];
\]

%dsp-level:0
\[
\textbf{the exact order:}
\]

%dsp-level:0
\[
\qquad \mathcal R_5<\mathcal R_6<\mathcal R_3<\mathcal R_4<\mathcal R_1<\mathcal R_2
\]

%dsp-level:0
\[
\qquad [\text{machine: pairwise comparisons in the exact fields with dyadic isolation}];
\]

%dsp-level:0
\[
\mathcal R_5<\mathcal R_6\iff2-2\mathsf c<2-\sqrt3\iff5>1;
\]

%dsp-level:0
\[
\qquad \mathcal R_6<\mathcal R_3\iff3600^{-1}(2-\sqrt3)<576^{-1}(2-\sqrt2)\iff2-\sqrt3<\tfrac{25}4(2-\sqrt2)
\]

%dsp-level:0
\[
\qquad [\text{machine}];
\]

%dsp-level:0
\[
\mathcal R_3<\mathcal R_4\iff2-\sqrt2<1\iff1<2;
\]

%dsp-level:0
\[
\qquad \mathcal R_4<\mathcal R_1\iff\frac1{576}<\frac1{72}\iff72<576;
\]

%dsp-level:0
\[
\qquad \mathcal R_1<\mathcal R_2\iff\frac1{72}<\frac1{36}\iff36<72
\]

%dsp-level:0
\[
\qquad [\text{machine: integers}];
\]

%dsp-level:0
\[
(3)\quad\textbf{the span of the bridge factor:}
\]

%dsp-level:0
\[
\qquad \boxed{\ \frac{\mathcal R_{\max}}{\mathcal R_{\min}}=\frac{\mathcal R_2}{\mathcal R_5}=\frac{1/36}{(2-2\mathsf c)/3600}=\frac{100}{2-2\mathsf c}=\frac{50}{1-\mathsf c}\ }
\]

%dsp-level:0
\[
\text{minimal polynomial:}\quad x^4-1200x^3+190000x^2-8000000x+100000000
\]

%dsp-level:0
\[
\qquad [\text{machine: the substitution}\ z=x/50\ \text{into}\ \mathrm{BM.79}(7)\text{'s}\ z^4-24z^3+76z^2-64z+16];
\]

%dsp-level:0
\[
\frac{\mathcal R_2}{\mathcal R_5}\in\Bigl(\frac{1021586454}{10^6},\frac{1021586455}{10^6}\Bigr),
\]

%dsp-level:0
\[
\qquad \sqrt{\frac{\mathcal R_2}{\mathcal R_5}}=\frac{10}{\sqrt{2-2\mathsf c}}\in\Bigl(\frac{31962266}{10^6},\frac{31962267}{10^6}\Bigr)
\]

%dsp-level:0
\[
\qquad [\text{machine: exact radical enclosures}];
\]

%dsp-level:0
\[
\textbf{the floor factor is exactly twenty-five:}
\]

%dsp-level:0
\[
\qquad \frac{\mathcal R_2/\mathcal R_5}{4/(2-2\mathsf c)}=\frac{100/(2-2\mathsf c)}{4/(2-2\mathsf c)}=25=\Bigl(\frac{30}6\Bigr)^2=\Bigl(\frac{E(\mathcal F_3)}{E(\mathcal F_1)}\Bigr)^2
\]

%dsp-level:0
\[
\qquad [\text{machine: the exact quotient}]:
\]

%dsp-level:0
\[
\textbf{the bridge factor's range exceeds the defect ceiling by exactly the square of the floor ladder's edge ratio}\ —\ \text{the angle contributes}\ \tfrac{4}{2-2\mathsf c},
\]

%dsp-level:0
\[
\ \text{the floor contributes}\ 25;
\]

%dsp-level:0
\[
(4)\quad\textbf{the intra-currency closure (the residue-free part of the bridge):}
\]

%dsp-level:0
\[
\qquad \textbf{the pole currency:}\quad\frac{m^2_{\mathrm{pole}}(s)}{m^2_{\mathrm{pole}}(s')}=\frac{D_S(s)}{D_S(s')}
\]

%dsp-level:0
\[
\qquad \text{for all}\ 30\ \text{ordered pairs of distinct states}\ (15\ \text{unordered pairs, each read in both directions}):
\]

%dsp-level:0
\[
\boxed{\ \text{all thirty are exact algebraic and}\ \rho_{\mathrm B}\textbf{-free}\ }
\]

%dsp-level:0
\[
\qquad [\text{machine:}\ \hat\mu'^2\ \text{cancels in every quotient};\ \text{the values lie in the compositum}\ \mathbb Q(\sqrt2,\sqrt3,\mathsf c),\ [\,\cdot:\mathbb Q]=16\ [\mathrm{FT.73}(2)\ \text{row-wise}]];
\]

%dsp-level:0
\[
\textbf{two exact cross-links to the ladder census:}
\]

%dsp-level:0
\[
\qquad \frac{D_S(s_1)}{D_S(s_6)}=\frac{2}{2-\sqrt3}=4+2\sqrt3=(1+\sqrt3)^2=r^2_{\mathrm{intra}}(\text{icosa}),
\]

%dsp-level:0
\[
\qquad \frac{D_S(s_2)}{D_S(s_5)}=\frac{4}{2-2\mathsf c}=\text{the dodecahedral ceiling}
\]

%dsp-level:0
\[
[\mathrm{BM.79}(7),(9);\ \mathrm{BM.81}(5);\ \text{machine: the exact identities}];
\]

%dsp-level:0
\[
\qquad \text{the widest and narrowest of the thirty are}\ \tfrac{4}{2-2\mathsf c}\ \text{and its reciprocal}\ \tfrac{2-2\mathsf c}4\ [\text{machine}];
\]

%dsp-level:0
\[
\textbf{the collision census:}
\]

%dsp-level:0
\[
\qquad \text{the thirty ordered ratios take twenty-eight distinct values, and the unique collision is the exact rational}\ 2:
\]

%dsp-level:0
\[
\qquad \frac{D_S(s_2)}{D_S(s_1)}=\frac42=\frac21=\frac{D_S(s_1)}{D_S(s_4)}
\]

%dsp-level:0
\[
\textbf{three of the six defect values are the powers}\ 1,
\]

%dsp-level:0
\[
\ 2,
\]

%dsp-level:0
\[
\ 4\ —\ \text{the octahedron, the tetrahedral light rung and the tetrahedral heavy rung}\ —\ \text{and that geometric triple is the census's only degeneracy}\ [\text{machine: the thirty-value census}];
\]

%dsp-level:0
\[
\textbf{the informational currency:}
\]

%dsp-level:0
\[
\qquad \frac{m_{\mathrm{hid}}(s)}{m_{\mathrm{hid}}(s')}=\frac{E(s)}{E(s')}\in\Bigl\{1,\ 2,\ 5,\ \frac12,\ \frac15,\ \frac25,\ \frac52\Bigr\}
\]

%dsp-level:0
\[
\qquad [\text{machine: the exact value set at}\ 2E\in\{12,24,60\}]:
\]

%dsp-level:0
\[
\boxed{\ \text{all}\ \rho_{\mathrm B}\textbf{-free, all rational, seven distinct values over the fifteen ordered pairs}\ };
\]

%dsp-level:0
\[
\qquad \textbf{the transport currency:}\quad\Delta m=\frac{\ln(4/3)}{\mathfrak h},\quad\frac{\ln2}{\mathfrak h}\ \ \text{per constituent}
\]

%dsp-level:0
\[
\qquad [\mathrm{BM.59}(4),(5);\ \rho_{\mathrm B}\text{-free}];
\]

%dsp-level:0
\[
\textbf{the ensemble currency:}
\]

%dsp-level:0
\[
\qquad \frac{T(0,S)}{T(+1,S)}=\frac{W(0)}{W(+1)}=\frac{20}{15}=\frac43
\]

%dsp-level:0
\[
\qquad [\mathrm{BM.28}(4)(\mathrm{ii});\ \mathrm{BM.30}(2);\ \rho_{\mathrm B}\text{-free}];
\]

%dsp-level:0
\[
(5)\quad\textbf{the cross-currency row (where the residue lives):}
\]

%dsp-level:0
\[
\qquad \boxed{\ \frac{m_{\mathrm{pole}}(s)}{m_{\mathrm{hid}}(s)}=\sqrt{\rho_{\mathrm B}\,\mathcal R_s}=\sqrt{\rho_{\mathrm B}}\cdot\Bigl(\frac{\sqrt2}{12},\ \frac16,\ \frac{\sqrt{2-\sqrt2}}{24},\ \frac1{24},\ \frac{\sqrt{2-2\mathsf c}}{60},\ \frac{\sqrt{2-\sqrt3}}{60}\Bigr)\ }
\]

%dsp-level:0
\[
[\text{machine: the six square roots simplified in the ordered radical extensions};\qquad\sqrt{\mathcal R_2}/\sqrt{\mathcal R_5}=10/\sqrt{2-2\mathsf c}\ \text{is}\ \rho_{\mathrm B}\text{-free again}];
\]

%dsp-level:0
\[
\textbf{the residue enters at exactly one step:}
\]

%dsp-level:0
\[
\qquad \boxed{\ \text{every mass ratio inside a currency is derived with zero residue};\qquad\text{the single cross-currency step carries the single undeclared}\ \rho_{\mathrm B}\ }
\]

%dsp-level:0
\[
(6)\quad\textbf{the residue localized:}
\]

%dsp-level:0
\[
\qquad (\mathrm{R1})\ \ \rho_{\mathrm B}=\hat\mu'^2\mathfrak h^2/(\ln2)^2\ \ \text{is the only undeclared datum any ratio of}\ (2)\text{--}(5)\ \text{needs}
\]

%dsp-level:0
\[
\qquad [\mathrm{BM.69}(5)(\mathrm{R1});\ \text{Cal-F packaging},\ \neg\text{added here}];
\]

%dsp-level:0
\[
(\mathrm{R2})\ \ \text{the six rows are read at the vertex-figure locus}\ \Theta_R
\]

%dsp-level:0
\[
= \delta_v,
\]

%dsp-level:0
\[
\ \text{whose ledger realization}\ —\ \text{a clean-enclosure loop through conical vertices}\ —\ \text{is}\ \neg\text{constructed and is inherited at every use}
\]

%dsp-level:0
\[
\qquad [\mathrm{BM.67}(9),(10);\ \text{Census FT.63}(4)\ \mathrm{Rmk};\ \mathrm{BM.69}(5)(\mathrm{R2})];
\]

%dsp-level:0
\[
(\mathrm{R3})\ \ \text{the flavour composition of the sealed classes: it enters only where a composite's pole mass is formed}\ —\ \neg\text{a ratio of}\ (2)\text{--}(5)\ \text{requires it}
\]

%dsp-level:0
\[
\qquad [\mathrm{BM.69}(3),(5)(\mathrm{R3});\ \mathrm{BM.28}(7)];
\]

%dsp-level:0
\[
(7)\quad\textbf{the non-derivation theorem, re-certified and sharpened:}
\]

%dsp-level:0
\[
\qquad \boxed{\ \neg\exists\ \text{an exact identity}\ \ m^2_{\mathrm{pole}}=\mathcal B\bigl[E_{\mathrm{int}},W(Q),\Theta_R,I_F\bigr]\ \ \text{with}\ \mathcal B\ \text{built from derived data alone}\ }
\]

%dsp-level:0
\[
\text{on four independent grounds}:\ (\mathrm i)\ \text{the degree-}16\ \text{field against integer- and binary-valued inputs}\ [\mathrm{FT.73}(2),(3)];
\]

%dsp-level:0
\[
\quad (\mathrm{ii})\ \text{the two-torsion image}\ \{0,\pi\}\ \text{reaching}\ 2\ \text{of}\ 6\ \text{states}\ [\mathrm{BM.67}(6)];
\]

%dsp-level:0
\[
(\mathrm{iii})\ \text{the calibration weight}\ {-}2\ \text{against}\ 0\ \text{on the winding route}\ [\mathrm{FT.55}(1);\ \mathrm{BM.67}(8)];
\]

%dsp-level:0
\[
\qquad (\mathrm{iv})\ \neg\exists\ \text{a count}\to\text{strength theorem}\ [\mathrm{BM.36}(3);\ \mathrm{BM.61}(8)(\mathrm{O1})]
\]

%dsp-level:0
\[
\qquad [\mathrm{BM.69}(4)\ \text{re-cited};\ \text{the grounds stand}]:
\]

%dsp-level:0
\[
\textbf{the census sharpening:}
\]

%dsp-level:0
\[
\qquad \text{the four grounds bound}\ \rho_{\mathrm B}\ \text{itself};
\]

%dsp-level:0
\[
\qquad \text{they leave}\ \mathcal R_s,
\]

%dsp-level:0
\[
\ \text{the fifteen pole ratios, the seven informational ratios, the transport differences and the ensemble ratio untouched}:
\]

%dsp-level:0
\[
\boxed{\ \text{the mass-ratio census is complete up to one undeclared dimensionless number, and complete with none inside any single currency}\ }
\]

%dsp-level:0
\[
(8)\quad\textbf{status:}\qquad\textbf{derived:}\ (1)\text{--}(5),\ (7);\qquad\textbf{localized,}\ \neg\textbf{filled:}\ (6);
\]

%dsp-level:0
\[
\neg\text{a P4 addition}\ (\rho_{\mathrm B}\ \text{is a Cal-F restatement},\ \neg\text{a new member}),
\]

%dsp-level:0
\[
\quad\neg\text{a new dimensionful scale}\ [\mathrm{P1}],
\]

%dsp-level:0
\[
\quad\neg\text{a measured value}\ [\mathrm{P5}],
\]

%dsp-level:0
\[
\quad\neg\text{a declaration},
\]

%dsp-level:0
\[
\quad\text{§Cal-BM stays empty};
\]

%dsp-level:0
\[
\neg\text{a currency is converted anywhere in this theorem};\qquad
\]

\textbf{Req.} CT2.1, D2.9 (Part 2); P1, P4, P5 (Parts 1–2); Def-H.37(2), FT.44(4), FT.47(1), FT.54(4), FT.55(1), FT.73(2), FT.73(3) (Part 3); Census FT.42(1), FT.63(4) (row-wise); Def-BM.14(1), Def-BM.14(3), BM.28(4), BM.28(7), BM.30(2), BM.36(3), BM.59(4), BM.59(5), BM.60(3), BM.61(8), BM.67(6), BM.67(8), BM.67(9), BM.67(10), BM.69(2), BM.69(3), BM.69(4), BM.69(5), BM.79(7), BM.79(9), BM.81(5) (this part); Cal-F; exact radical arithmetic in four fields and their degree-16 compositum, minimal polynomials by first linear dependence and by exact substitution, Laurent-monomial assembly, exponent arithmetic in the $\mathfrak q$-grading, exact dyadic radical enclosures, pairwise-distinctness censuses [math; machine: \texttt{checks/\allowbreak{}verify\_\allowbreak{}H\_\allowbreak{}mass\_\allowbreak{}ratio\_\allowbreak{}census.py} §E].

\textbf{Pf.}

%dsp-level:0
\[
(1):\ \frac{\hat\mu'^2D_S}{(2E\ln2/\mathfrak h)^2}
\]

%dsp-level:0
\[
= \frac{\hat\mu'^2\mathfrak h^2}{(\ln2)^2}\cdot\frac{D_S}{4E^2}
\]

%dsp-level:0
\[
= \rho_{\mathrm B}\mathcal R_s\ [\text{machine: one cleared identity over}\ \mathbb Q[\hat\mu'^2,\mathfrak h,\ln2,D_S,E]];
\]

%dsp-level:0
\[
[\rho_{\mathrm B}]=[\hat\mu'^2][\mathfrak h]^2=\mathfrak q\cdot\mathfrak q^{-1}=\mathfrak q^0\ [\text{machine: the exponent arithmetic}];
\]

%dsp-level:0
\[
\qquad \mathcal R_s\in\overline{\mathbb Q}\ \text{since}\ D_S\ \text{is algebraic and}\ E\in\mathbb Z\ [\text{machine}]
\]

%dsp-level:0
\[
(2):\ 4E^2=4(6,6,12,12,30,30)^2=(144,144,576,576,3600,3600)\ [\text{machine}];
\]

%dsp-level:0
\[
\qquad \text{the six quotients formed exactly}\ [\text{machine}];
\]

%dsp-level:0
\[
\qquad \text{the minimal polynomials:}\ 72x-1,
\]

%dsp-level:0
\[
\ 36x-1,
\]

%dsp-level:0
\[
\ 576x-1\ \text{by inspection};
\]

%dsp-level:0
\[
\ (2-\sqrt2)/576:\ y
\]

%dsp-level:0
\[
= 576x\ \text{satisfies}\ y^2-4y+2
\]

%dsp-level:0
\[
\implies 331776x^2-2304x+2
\]

%dsp-level:0
\[
= 0
\]

%dsp-level:0
\[
\implies 165888x^2-1152x+1\ [\text{machine}]
\]

%dsp-level:0
\[
(2-\sqrt3)/3600:\ y^2-4y+1\implies12960000x^2-14400x+1\ [\text{machine}];
\]

%dsp-level:0
\[
\qquad (2-2\mathsf c)/3600:\ y^4-8y^3+19y^2-12y+1\ \text{at}\ y
\]

%dsp-level:0
\[
= 3600x
\]

%dsp-level:0
\[
\implies 3600^4x^4-8\cdot3600^3x^3+19\cdot3600^2x^2-12\cdot3600x+1\ [\text{machine: the four coefficients as exact integers}]
\]

%dsp-level:0
\[
\text{the order}: \text{the five comparisons of}\ (2)\ \text{with}\ 2-\sqrt3<\tfrac{25}4(2-\sqrt2)\ \text{certified by}\ (2-\sqrt3)^2
\]

%dsp-level:0
\[
= 7-4\sqrt3<\bigl(\tfrac{25}4(2-\sqrt2)\bigr)^2\ \text{on exact dyadic isolation}\ [\text{machine}]
\]

%dsp-level:0
\[
(3):\ \frac{1/36}{(2-2\mathsf c)/3600}=\frac{3600}{36(2-2\mathsf c)}=\frac{100}{2-2\mathsf c}=\frac{50}{1-\mathsf c}\ [\text{machine}];
\]

%dsp-level:0
\[
\qquad z^4-24z^3+76z^2-64z+16
\]

%dsp-level:0
\[
= 0\ \text{at}\ z
\]

%dsp-level:0
\[
= x/50
\]

%dsp-level:0
\[
\implies x^4-1200x^3+190000x^2-8000000x+100000000
\]

%dsp-level:0
\[
= 0\ [\text{machine: cleared}];
\]

%dsp-level:0
\[
\qquad \text{the two enclosures}\ [\text{machine}]
\]

%dsp-level:0
\[
\frac{100/(2-2\mathsf c)}{4/(2-2\mathsf c)}=25\ \text{and}\ (30/6)^2=25\ [\text{machine: exact rationals}]
\]

%dsp-level:0
\[
(4):\ \hat\mu'^2D_S(s)/(\hat\mu'^2D_S(s'))=D_S(s)/D_S(s')\ [\text{machine}];
\]

%dsp-level:0
\[
\qquad 2/(2-\sqrt3)=2(2+\sqrt3)=4+2\sqrt3=(1+\sqrt3)^2\ [\text{machine: exact in}\ \mathbb Q(\sqrt3)];
\]

%dsp-level:0
\[
\qquad \text{the fifteen quotients pairwise distinct by dyadic isolation}\ [\text{machine}]
\]

%dsp-level:0
\[
E(s)/E(s')\ \text{over}\ \{6,12,30\}:\ \text{the value set}\ \{1,2,5,\tfrac12,\tfrac15,\tfrac25,\tfrac52\}\ \text{by exhaustive enumeration of the}\ 6\cdot5
\]

%dsp-level:0
\[
= 30\ \text{ordered pairs}\ [\text{machine}];
\]

%dsp-level:0
\[
\qquad W(0)/W(+1)=20/15=4/3\ [\text{machine}]
\]

%dsp-level:0
\[
(5):\ \sqrt{1/72}
\]

%dsp-level:0
\[
= \sqrt2/12,
\]

%dsp-level:0
\[
\quad\sqrt{1/36}
\]

%dsp-level:0
\[
= 1/6,
\]

%dsp-level:0
\[
\quad\sqrt{(2-\sqrt2)/576}
\]

%dsp-level:0
\[
= \sqrt{2-\sqrt2}/24,
\]

%dsp-level:0
\[
\quad\sqrt{1/576}
\]

%dsp-level:0
\[
= 1/24,
\]

%dsp-level:0
\[
\quad\sqrt{(2-2\mathsf c)/3600}
\]

%dsp-level:0
\[
= \sqrt{2-2\mathsf c}/60,
\]

%dsp-level:0
\[
\quad\sqrt{(2-\sqrt3)/3600}
\]

%dsp-level:0
\[
= \sqrt{2-\sqrt3}/60\ [\text{machine: each verified by squaring the positive side}]
\]

%dsp-level:0
\[
(6)\text{--}(8):\ \text{the three residue rows re-cited from}\ \mathrm{BM.69}(3),(5)\ \text{and}\ \mathrm{BM.67}(9),(10);
\]

%dsp-level:0
\[
\qquad \text{the four grounds re-cited from}\ \mathrm{BM.69}(4);
\]

%dsp-level:0
\[
\qquad \text{the standing-rule scans}\ [\text{machine}].\ \blacksquare
\]

\textbf{Phys.} The bridge is an exact factorization and nothing more mysterious than that. Divide the cone-defect mass by the edge-count mass and everything cancels except one dimensionless packaging and one algebraic number per state — and the algebraic numbers are six exact values in four fields, ordered, with minimal polynomials written out to a fourteen-digit leading coefficient. Inside a currency the packaging cancels too, so all fifteen pole-side ratios, all seven informational ratios, both transport differences and the ensemble's four thirds are derived with no residue at all: the census is complete. The residue lives at exactly one step, the step between currencies, and it is one number. The bridge factor's own range is a thousand and twenty-two in mass-squared, twenty-five times the widest angle ratio, and that twenty-five is precisely the square of the edge-count ratio between the heaviest and lightest floor — geometry supplies the shape, the floors supply the reach. What the four obstruction grounds bound is the packaging and nothing else, so the honest statement of the whole tier is short: every mass ratio the corpus can form is derived exactly, and one dimensionless constant stands between that census and an absolute cross-currency mass.

\pseudohead{theorem}{\texorpdfstring{BM.83}{BM.83}}{\texorpdfstring{The inter-family bridge in full: the four-currency census, the exact incommensurability, and the three-number localization}{The inter-family bridge in full: the four-currency census, the exact incommensurability, and the three-number localization}}\label{tag:BM-83}
\textbf{Thm.}

%dsp-level:0
\[
(1)\quad\textbf{the four currencies read the same doublet (the census table):}
\]

%dsp-level:0
\[
\qquad \text{the neutral class against the unit class}
\]

%dsp-level:0
\[
\qquad [\mathrm{BM.28}(1);\ \mathrm{Def\text{-}BM.8}(1)]:
\]

%tbl-level:3
{\scriptsize\setlength{\tabcolsep}{4pt}
\begin{longtable}{>{\raggedright\arraybackslash}p{\dimexpr0.1259\textwidth-2\tabcolsep\relax}|>{\raggedright\arraybackslash}p{\dimexpr0.1608\textwidth-2\tabcolsep\relax}|>{\raggedright\arraybackslash}p{\dimexpr0.2448\textwidth-2\tabcolsep\relax}|>{\raggedright\arraybackslash}p{\dimexpr0.4685\textwidth-2\tabcolsep\relax}}
\toprule
currency & the formed object & the exact value & the residue a ratio needs \\
\addlinespace[2pt]
\midrule
\endfirsthead
\toprule
currency & the formed object & the exact value & the residue a ratio needs \\
\addlinespace[2pt]
\midrule
\endhead
\bottomrule
\endlastfoot
informational & $m_n/m_p$ & $\dfrac{\ln1280}{\ln960+6\,c\,\varepsilon_{\mathrm{EM}}\ln2}$ & none — formed $[\mathrm{BM.62}(4),\allowbreak (5)]$ \\
ensemble, route $(\mathrm{ii})$ & $m_n/m_p$ & $\dfrac{T(0,S)}{T(+1,S)}=\allowbreak \dfrac43$ & none — formed; the route is $\ne\allowbreak g$adopted globally $[\mathrm{BM.28}(4)(\mathrm{ii});\ \mathrm{BM.30}(2),\allowbreak (5)]$ \\
transport & $\Delta m$ (a difference) & $\dfrac{\ln(4/3)}{\mathfrak h}$; per constituent $\dfrac{\ln2}{\mathfrak h}$ & none — formed $[\mathrm{BM.59}(4),\allowbreak (5)]$ \\
pole/defect & $\Delta(m^2_{\mathrm{pole}})$ (a difference) & $\Delta_g=\allowbreak (2,\allowbreak \sqrt2{-}1,\allowbreak 2\mathsf c{-}\sqrt3)\,\allowbreak \hat\mu'^2$ & $(\mathrm{R3})$ the composition, for a ratio $[\mathrm{BM.69}(3);\ \mathrm{BM.75}(4)]$ \\
\end{longtable}
}

%dsp-level:0
\[
\textbf{two ratios and two differences:}
\]

%dsp-level:0
\[
\qquad \text{no currency supplies both, and}\ \neg\text{a derived equation turns one into the other}
\]

%dsp-level:0
\[
\qquad [\mathrm{BM.35}(5);\ \mathrm{BM.37}(2)];
\]

%dsp-level:0
\[
(2)\quad\textbf{the prime-log ring and the irrationality theorem:}
\]

%dsp-level:0
\[
\qquad \mathfrak P:=\mathbb Q\oplus\bigoplus_{p\ \mathrm{prime}}\mathbb Q\ln p,
\]

%dsp-level:0
\[
\qquad \ln1280=8\ln2+\ln5,
\]

%dsp-level:0
\[
\qquad \ln960=6\ln2+\ln3+\ln5
\]

%dsp-level:0
\[
\qquad [\text{machine: the exact factorizations}\ 1280=2^8\!\cdot\!5,\ 960=2^6\!\cdot\!3\!\cdot\!5]:
\]

%dsp-level:0
\[
\textbf{the independence lemma:}\qquad\sum_ia_i\ln p_i=0\ \ \text{with}\ \ a_i\in\mathbb Q\ \ \Longrightarrow\ \ a_i=0\ \ \forall i
\]

%dsp-level:0
\[
[\text{math: clear denominators by}\ N;\ \ \exp\ \text{gives}\ \prod_ip_i^{Na_i}=1;\ \ \text{unique factorization forces every}\ Na_i=0;\qquad\text{machine: the integer version}\ 2^a3^b5^c=1\implies a=b=c=0\ \text{verified exhaustively on}\ |a|,|b|,|c|\le24\ \text{and by the sign-split argument}];
\]

%dsp-level:0
\[
\boxed{\ \textbf{the irrationality theorem:}\qquad\frac{\ln1280}{\ln960+t\ln2}\ \notin\ \mathbb Q\qquad\text{for every rational}\ t\ge0\ }
\]

%dsp-level:0
\[
[\text{machine: if}\ q\ln1280=p(\ln960+t\ln2)\ \text{with}\ p,q\in\mathbb Z,\ q>0,\ \text{then the}\ \ln3\ \text{coefficient gives}\ -p=0,\ \text{so}\ p=0,\ \text{and the}\ \ln5\ \text{coefficient then gives}\ q=0\ \text{by the lemma}];
\]

%dsp-level:0
\[
\textbf{corollary:}
\]

%dsp-level:0
\[
\qquad \frac{m_n}{m_p}\ \ \text{is irrational on the informational currency at every rational coupling}\ \varepsilon_{\mathrm{EM}}\in\mathbb Q\ \text{and both priced}\ c\in\bigl\{\tfrac12,\tfrac35\bigr\}
\]

%dsp-level:0
\[
\qquad [t=6c\,\varepsilon_{\mathrm{EM}};\ \text{machine}];
\]

%dsp-level:0
\[
(3)\quad\textbf{the exact incommensurability of the two formed ratios:}
\]

%dsp-level:0
\[
\qquad \frac{m_n}{m_p}\bigg|_{\mathrm{info}}=\frac43\iff3\ln1280=4\bigl(\ln960+t\ln2\bigr)\iff0=4\ln3+\ln5+4t\ln2\iff t=-\,\frac{\ln405}{\ln16}
\]

%dsp-level:0
\[
[\text{machine: solved exactly in}\ \mathfrak P;\qquad4\ln3+\ln5=\ln(3^4\!\cdot\!5)=\ln405>0\ \text{and}\ \ln16>0\implies t<0]:
\]

%dsp-level:0
\[
\boxed{\ \textbf{the informational ratio never equals the ensemble ratio at any admissible coupling:}\qquad\frac{\ln1280}{\ln960+t\ln2}\ \ne\ \frac43\qquad\forall\,t\ge0\ }
\]

%dsp-level:0
\[
\textbf{and the order between them is exact:}
\]

%dsp-level:0
\[
\qquad 1\ <\ \frac{\ln1280}{\ln960}\ <\ \frac43
\]

%dsp-level:0
\[
\qquad [\ 4\ln960-3\ln1280=4\ln3+\ln5=\ln405>0;\qquad\ln1280>\ln960\iff1280>960;\qquad\text{machine}\ ];
\]

%dsp-level:0
\[
\textbf{the two ratios are not merely unequal, they are of different kinds:}
\]

%dsp-level:0
\[
\qquad \frac43\in\mathbb Q\ \ \text{and}\ \ \frac{\ln1280}{\ln960}\in\mathfrak P\setminus\mathbb Q
\]

%dsp-level:0
\[
\qquad [(2),(3);\ \text{machine}];
\]

%dsp-level:0
\[
(4)\quad\textbf{the dimensionless scale set (restatements,}\ \neg\textbf{new members):}
\]

%dsp-level:0
\[
\qquad \sigma_\ell:=\frac{\mu_\ell\,\mathfrak h}{\ln2}\in\mathfrak q^0,
\]

%dsp-level:0
\[
\qquad \sigma_\nu:=\frac{\mu_\nu\,\mathfrak h}{\ln2}\in\mathfrak q^0,
\]

%dsp-level:0
\[
\qquad \rho_{\mathrm B}=\frac{\hat\mu'^2\mathfrak h^2}{(\ln2)^2}\in\mathfrak q^0
\]

%dsp-level:0
\[
[\mu_\ell,\mu_\nu\ \text{are the Cal-F/P4 channel scales}\ [\mathrm{FT.113};\ \mathrm{FT.143}(3);\ \text{Cal-F}];\qquad[\mu_\ell]=[\mu_\nu]=\mathfrak q^{+1/2}\ \text{from}\ m=\mu z\ \text{and}\ [m]=\mathfrak q^{+1/2};\qquad\text{machine: the exponent arithmetic}];
\]

%dsp-level:0
\[
\textbf{they are restatements of undeclared data,}\ \neg\textbf{declarations:}
\]

%dsp-level:0
\[
\qquad \neg\text{a P4 addition},
\]

%dsp-level:0
\[
\qquad \text{§Cal-BM stays empty}
\]

%dsp-level:0
\[
\qquad [\mathrm{BM.77}(2)\ \text{pattern}];
\]

%dsp-level:0
\[
(5)\quad\textbf{the inter-family bridge, in full (every cross-currency ratio of the census):}
\]

%dsp-level:0
\[
\qquad \text{with}\ m_{\mathrm{hid}}(s)\,\mathfrak h/\ln2=2E(s),\quad m^2_{\mathrm{pole}}(s)/\hat\mu'^2=D_S(s),\quad m^{(\mathrm{ch})}_i=\mu_{\mathrm{ch}}z^{(\mathrm{ch})}_i:
\]

%dsp-level:0
\[
\boxed{\ \frac{m^{(\ell)}_i}{m_{\mathrm{hid}}(s)}=\frac{\sigma_\ell\,z^{(\ell)}_i}{2E(s)},\qquad\frac{m^{(\nu)}_j}{m^{(\ell)}_i}=\frac{\sigma_\nu\,z^{(\nu)}_j}{\sigma_\ell\,z^{(\ell)}_i},\qquad\frac{m_{\mathrm{pole}}(s)}{m^{(\ell)}_i}=\frac{2E(s)}{\sigma_\ell\,z^{(\ell)}_i}\cdot\frac{\sqrt{\rho_{\mathrm B}\,\mathcal R_s}}{1}\ }
\]

%dsp-level:0
\[
[\text{machine: the three identities cleared in}\ \mathfrak P\otimes\mathbb Q[\rho_{\mathrm B},\sigma_\ell,\sigma_\nu];\qquad\mathrm{BM.82}(1),(5)\ \text{re-cited}];
\]

%dsp-level:0
\[
\boxed{\ \textbf{the localization theorem:}\qquad\text{every cross-currency mass ratio of the census is an exact derived expression in exactly three undeclared dimensionless numbers}\ \ (\rho_{\mathrm B},\ \sigma_\ell,\ \sigma_\nu),\ \text{and in no others}\ }
\]

%dsp-level:0
\[
\textbf{minimality (each is needed):}
\]

%dsp-level:0
\[
\qquad m_{\mathrm{pole}}/m_{\mathrm{hid}}\ \text{depends on}\ \rho_{\mathrm B}\ \text{alone};
\]

%dsp-level:0
\[
\qquad m^{(\ell)}/m_{\mathrm{hid}}\ \text{on}\ \sigma_\ell\ \text{alone};
\]

%dsp-level:0
\[
\qquad m^{(\nu)}/m^{(\ell)}\ \text{on}\ \sigma_\nu/\sigma_\ell\ \text{alone}
\]

%dsp-level:0
\[
\qquad [\text{machine: the exact monomial dependence of each of the three formulas}];
\]

%dsp-level:0
\[
\textbf{and every ratio inside a currency depends on none of them}\qquad[\mathrm{BM.82}(4);\ \text{machine}];
\]

%dsp-level:0
\[
(6)\quad\textbf{the cross-channel factor census (the nine exact derived factors):}
\]

%dsp-level:0
\[
\qquad \frac{m^{(\mathrm{ch})}_i}{m^{(\mathrm{ch}')}_i}=\frac{\mu_{\mathrm{ch}}}{\mu_{\mathrm{ch}'}}\cdot\frac{z^{(\mathrm{ch})}_i}{z^{(\mathrm{ch}')}_i}
\]

%dsp-level:0
\[
\qquad [\text{machine: the exact factorization}]:
\]

%tbl-level:3
{\scriptsize\setlength{\tabcolsep}{4pt}
\begin{longtable}{>{\raggedright\arraybackslash}p{\dimexpr0.1119\textwidth-2\tabcolsep\relax}|>{\raggedright\arraybackslash}p{\dimexpr0.2867\textwidth-2\tabcolsep\relax}|>{\raggedright\arraybackslash}p{\dimexpr0.3007\textwidth-2\tabcolsep\relax}|>{\raggedright\arraybackslash}p{\dimexpr0.3007\textwidth-2\tabcolsep\relax}}
\toprule
line & $z^{(q)}_i/z^{(\nu)}_i$ & $z^{(q)}_i/z^{(\ell)}_i$ & $z^{(\nu)}_i/z^{(\ell)}_i$ \\
\addlinespace[2pt]
\midrule
\endfirsthead
\toprule
line & $z^{(q)}_i/z^{(\nu)}_i$ & $z^{(q)}_i/z^{(\ell)}_i$ & $z^{(\nu)}_i/z^{(\ell)}_i$ \\
\addlinespace[2pt]
\midrule
\endhead
\bottomrule
\endlastfoot
$i=\allowbreak 0$ (the heaviest) & $\bigl(\tfrac{1005964365750}{10^{12}},\allowbreak \tfrac{1005964365751}{10^{12}}\bigr)$ & $\bigl(\tfrac{1026138854777}{10^{12}},\allowbreak \tfrac{1026138854778}{10^{12}}\bigr)$ & $\bigl(\tfrac{1020054874420}{10^{12}},\allowbreak \tfrac{1020054874421}{10^{12}}\bigr)$ \\
$i=\allowbreak 2$ (the middle) & $\bigl(\tfrac{730177488628}{10^{12}},\allowbreak \tfrac{730177488629}{10^{12}}\bigr)$ & $\bigl(\tfrac{441377821972}{10^{12}},\allowbreak \tfrac{441377821973}{10^{12}}\bigr)$ & $\bigl(\tfrac{604480183033}{10^{12}},\allowbreak \tfrac{604480183034}{10^{12}}\bigr)$ \\
$i=\allowbreak 1$ (the lightest) & $\bigl(\tfrac{1963683086044}{10^{12}},\allowbreak \tfrac{1963683086045}{10^{12}}\bigr)$ & $\bigl(\tfrac{25609326778349}{10^{12}},\allowbreak \tfrac{25609326778350}{10^{12}}\bigr)$ & $\bigl(\tfrac{13041476478742}{10^{12}},\allowbreak \tfrac{13041476478743}{10^{12}}\bigr)$ \\
\end{longtable}
}

%dsp-level:0
\[
[\text{machine: monotone exact interval division at width}<2^{-150}];
\]

%dsp-level:0
\[
\boxed{\ \textbf{the light-end concentration:}\qquad1\ <\ \frac{z^{(q)}_0}{z^{(\nu)}_0},\ \frac{z^{(\nu)}_0}{z^{(\ell)}_0},\ \frac{z^{(q)}_0}{z^{(\ell)}_0}\ <\ \frac{27}{25},\qquad\text{while}\qquad\frac{25609326778349}{10^{12}}\ <\ \frac{z^{(q)}_1}{z^{(\ell)}_1}\ <\ \frac{26}{1}\ }
\]

%dsp-level:0
\[
\textbf{the three channels' heaviest lines are nearly degenerate and their lightest lines differ by more than an order of magnitude}\ —\ \text{the entire inter-family structure lives at the light end}
\]

%dsp-level:0
\[
[\text{machine: the exact bracket comparisons};\qquad\text{the structural mirror of}\ \mathrm{BM.79}(5)\text{'s top-weight row on the Platonic side}\ [\text{Census FT.76}(4)\ \text{row-wise}]];
\]

%dsp-level:0
\[
(7)\quad\textbf{the inter-family verdict (what is and what is not derived):}
\]

%dsp-level:0
\[
\qquad \textbf{derived:}\ \text{the shape invariants}\ \sum_iz_i
\]

%dsp-level:0
\[
= 6,
\]

%dsp-level:0
\[
\ Q
\]

%dsp-level:0
\[
= \tfrac23\ \text{in all three channels};
\]

%dsp-level:0
\[
\ \text{the one cubic form with three constants};
\]

%dsp-level:0
\[
\ \text{the hierarchy-order theorem}\ [\mathrm{BM.80}(3),(4),(7),(8)];
\]

%dsp-level:0
\[
\text{the nine cross-channel factors of}\ (6);
\]

%dsp-level:0
\[
\qquad \text{the exact factorization of every cross-currency ratio and its localization to}\ (\rho_{\mathrm B},\sigma_\ell,\sigma_\nu)\ (5);
\]

%dsp-level:0
\[
\textbf{derived:}\ \text{the quark side of the inter-family ratio in two currencies}\ —\ \text{the pole currency's}\ \hat\mu'^2D_S\ \text{and the informational currency's}\ 2E\ln2/\mathfrak h\ \text{with the exact bridge factor}\ \mathcal R_s\ [\mathrm{BM.81}(1);\ \mathrm{BM.82}(2)];
\]

%dsp-level:0
\[
\boxed{\ \neg\textbf{derived:}\ \ \text{the inter-family mass ratio as a number}\ }
\]

%dsp-level:0
\[
\qquad \text{it requires}\ \rho_{\mathrm B}\ (\mathrm{R1}),
\]

%dsp-level:0
\[
\ \sigma_\ell,
\]

%dsp-level:0
\[
\ \sigma_\nu\ (\text{Cal-F}),
\]

%dsp-level:0
\[
\ \text{and}\ —\ \text{on the pole side of a composite}\ —\ (\mathrm{R3})\ \text{the composition}:
\]

%dsp-level:0
\[
\textbf{the obstruction is exactly three dimensionless numbers and one attachment, all named, none supplied, none fitted}
\]

%dsp-level:0
\[
\qquad [\mathrm{BM.69}(4),(5);\ \mathrm{BM.60}(7)];
\]

%dsp-level:0
\[
(8)\quad\textbf{the absolute neutrino mass:}
\]

%dsp-level:0
\[
\qquad \neg\text{derived}
\]

%dsp-level:0
\[
\qquad [\mathrm{FT.143}(3):\ \mu_\nu\ \text{undeclared}]:
\]

%dsp-level:0
\[
\qquad \text{the}\ \nu\text{-channel's internal ratios are derived}\ (6),\ \text{its absolute scale is not};
\]

%dsp-level:0
\[
\text{the two-bit floor applies to it on the informational read:}
\]

%dsp-level:0
\[
\qquad m_\nu\mathfrak h\ge2\ln2\ \ \text{once materialized as a bound chain},
\]

%dsp-level:0
\[
\qquad m_\nu\mathfrak h\ge0\ \ \text{by PSD alone}
\]

%dsp-level:0
\[
\qquad [\mathrm{BM.70}(2),(3)\ \text{re-cited}];
\]

%dsp-level:0
\[
(9)\quad\textbf{the bars:}
\]

%dsp-level:0
\[
\qquad \neg\exists\ \text{an SM mass, mass ratio, mixing angle or Yukawa value in any clause}
\]

%dsp-level:0
\[
\qquad [\mathrm{P5};\ \text{machine: zero external tokens, zero decimal numerals, zero}\ [\mathrm{EXT}]\ \text{markers in this tier}];
\]

%dsp-level:0
\[
\neg\text{a measured external constant is identified with any corpus quantity}\ [\mathrm{P5}];
\]

%dsp-level:0
\[
\qquad \text{species names are confined to the identification layers}\ [\mathrm{Def\text{-}BM.3}(7);\ \mathrm{Def\text{-}BM.6}(7);\ \mathrm{Def\text{-}BM.14}(8)];
\]

%dsp-level:0
\[
\text{topology labels are internal bookkeeping}\ [\mathrm{Def\text{-}H.32}(4)];
\]

%dsp-level:0
\[
\qquad \text{amplitude}\ne\text{mass}\ [\mathrm{FT.79}(5)];
\]

%dsp-level:0
\[
\qquad \text{rung}\ne\text{scale}\ [\mathrm{D0.8}];
\]

%dsp-level:0
\[
\qquad \neg\text{a currency is silently converted anywhere in this theorem};
\]

%dsp-level:0
\[
(10)\quad\textbf{status:}
\]

%dsp-level:0
\[
\qquad \textbf{derived:}\ (1)\text{--}(8);
\]

%dsp-level:0
\[
\qquad \textbf{localized,}\ \neg\textbf{filled:}\ (7)\text{'s obstruction};
\]

%dsp-level:0
\[
\qquad \textbf{declared: nothing}\ —\ (4)\text{'s three symbols are restatements of Cal-F/P4 data};
\]

%dsp-level:0
\[
\neg\text{a P4 addition},\quad\neg\text{a new dimensionful scale}\ [\mathrm{P1}],\quad\neg\text{a measured value}\ [\mathrm{P5}],\quad
\]

%dsp-level:0
\[
\qquad \text{the mass-form bridge is application-layer, its residue reduced to}\ (\rho_{\mathrm B},\sigma_\ell,\sigma_\nu)\ \text{and}\ (\mathrm{R2}),(\mathrm{R3}).
\]

\textbf{Req.} D0.8, D2.9, CT2.1 (Part 2); P1, P4, P5 (Parts 1–2); Def-H.32(4), Def-H.37(2), FT.79(5), FT.113, FT.143(3) (Part 3); Def-BM.3(7), Def-BM.6(7), Def-BM.8(1), Def-BM.14(8), BM.28(1), BM.28(4), BM.30(2), BM.30(5), BM.35(5), BM.37(2), BM.59(4), BM.59(5), BM.60(7), BM.62(4), BM.62(5), BM.69(3), BM.69(4), BM.69(5), BM.70(2), BM.70(3), BM.75(4), BM.77(2), BM.79(5), BM.80(3), BM.80(4), BM.80(6), BM.80(7), BM.80(8), BM.81(1), BM.82(1), BM.82(2), BM.82(4), BM.82(5) (this part); Cal-F; exact arithmetic in the prime-log ring with integer sign tests, unique-factorization independence, exact symbolic solving in $\mathfrak P$, Laurent-monomial assembly and $\mathfrak q$-exponent arithmetic, monotone exact interval division [math; machine: \texttt{checks/\allowbreak{}verify\_\allowbreak{}H\_\allowbreak{}mass\_\allowbreak{}ratio\_\allowbreak{}census.py} §F].

\textbf{Pf.}

%dsp-level:0
\[
(1):\ \text{the four rows re-cited from their own tags}\ [\mathrm{BM.62}(4),(5);\ \mathrm{BM.28}(4)(\mathrm{ii});\ \mathrm{BM.30}(2);\ \mathrm{BM.59}(4),(5);\ \mathrm{BM.69}(3);\ \mathrm{BM.75}(4)];
\]

%dsp-level:0
\[
\qquad \text{the formed-object column: two ratios and two differences}\ [\text{machine: the token census of the four expressions}]
\]

%dsp-level:0
\[
(2):\ 1280
\]

%dsp-level:0
\[
= 2^8\cdot5\ \text{and}\ 960
\]

%dsp-level:0
\[
= 2^6\cdot3\cdot5\ [\text{machine: exact integer factorizations}]\implies\ln1280
\]

%dsp-level:0
\[
= 8\ln2+\ln5,
\]

%dsp-level:0
\[
\ \ln960
\]

%dsp-level:0
\[
= 6\ln2+\ln3+\ln5\ [\text{machine: the prime-log representation}];
\]

%dsp-level:0
\[
\qquad \text{the lemma:}\ \exp\bigl(\sum Na_i\ln p_i\bigr)=\prod p_i^{Na_i}=1\implies Na_i=0\ \forall i\ [\text{math: unique factorization}]
\]

%dsp-level:0
\[
\text{machine: the exhaustive census of}\ 2^a3^b5^c=1\ \text{over}\ |a|,|b|,|c|\le24\ \text{returns exactly}\ (0,0,0);
\]

%dsp-level:0
\[
\qquad \text{the rational case}:\ q(8\ln2+\ln5)
\]

%dsp-level:0
\[
= p((6+t)\ln2+\ln3+\ln5)\implies\text{the}\ \ln3\ \text{coefficient}\ -p
\]

%dsp-level:0
\[
= 0
\]

%dsp-level:0
\[
\implies p
\]

%dsp-level:0
\[
= 0
\]

%dsp-level:0
\[
\implies q\cdot8\ln2+q\ln5
\]

%dsp-level:0
\[
= 0
\]

%dsp-level:0
\[
\implies q
\]

%dsp-level:0
\[
= 0\ [\text{machine: the coefficient extraction and the lemma}]
\]

%dsp-level:0
\[
(3):\ 3\ln1280=24\ln2+3\ln5;
\]

%dsp-level:0
\[
\qquad 4\ln960+4t\ln2=(24+4t)\ln2+4\ln3+4\ln5\ [\text{machine}];
\]

%dsp-level:0
\[
\qquad \text{the difference}\ 4\ln960+4t\ln2-3\ln1280
\]

%dsp-level:0
\[
= 4\ln3+\ln5+4t\ln2
\]

%dsp-level:0
\[
= \ln405+4t\ln2\ [\text{machine: the prime-log ring}]
\]

%dsp-level:0
\[
\ln405>0\iff405>1,\quad\ln16>0\iff16>1\implies t=-\ln405/\ln16<0\ [\text{machine}];
\]

%dsp-level:0
\[
\qquad \ln1280>\ln960\iff1280>960\ [\text{machine}];
\]

%dsp-level:0
\[
\qquad \text{the kind row}:\ \tfrac43\in\mathbb Q,\ \tfrac{\ln1280}{\ln960}\notin\mathbb Q\ [(2)]
\]

%dsp-level:0
\[
(4):\ [m]=\mathfrak q^{+1/2}\ [\mathrm{CT2.1};\ \mathrm{D2.9}],\quad z\in\mathfrak q^0\implies[\mu_{\mathrm{ch}}]=\mathfrak q^{+1/2};
\]

%dsp-level:0
\[
\qquad [\mathfrak h]=\mathfrak q^{-1/2},\quad[\ln2]=\mathfrak q^0\implies\sigma_{\mathrm{ch}}\in\mathfrak q^0\ [\text{machine: the exponent arithmetic}];
\]

%dsp-level:0
\[
\qquad \rho_{\mathrm B}\in\mathfrak q^0\ [\mathrm{BM.82}(1)]
\]

%dsp-level:0
\[
(5):\ \frac{\mu_\ell z}{2E\ln2/\mathfrak h}=\frac{\mu_\ell\mathfrak h}{\ln2}\cdot\frac{z}{2E}=\frac{\sigma_\ell z}{2E}\ [\text{machine: cleared}];
\]

%dsp-level:0
\[
\qquad \frac{\mu_\nu z^{(\nu)}}{\mu_\ell z^{(\ell)}}=\frac{\sigma_\nu z^{(\nu)}}{\sigma_\ell z^{(\ell)}}\ [\text{machine}];
\]

%dsp-level:0
\[
\qquad \frac{m_{\mathrm{pole}}}{m^{(\ell)}_i}
\]

%dsp-level:0
\[
= \frac{m_{\mathrm{pole}}}{m_{\mathrm{hid}}}\cdot\frac{m_{\mathrm{hid}}}{m^{(\ell)}_i}
\]

%dsp-level:0
\[
= \sqrt{\rho_{\mathrm B}\mathcal R_s}\cdot\frac{2E}{\sigma_\ell z^{(\ell)}_i}\ [\text{machine: the two re-cited identities composed}]
\]

%dsp-level:0
\[
\text{minimality}: \text{the symbol census of the three expressions}\ —\ \sqrt{\rho_{\mathrm B}\mathcal R_s}\ \text{carries}\ \rho_{\mathrm B}\ \text{and neither}\ \sigma;
\]

%dsp-level:0
\[
\quad\sigma_\ell z/2E\ \text{carries}\ \sigma_\ell\ \text{and neither}\ \rho_{\mathrm B}\ \text{nor}\ \sigma_\nu;
\]

%dsp-level:0
\[
\quad\sigma_\nu z^{(\nu)}/(\sigma_\ell z^{(\ell)})\ \text{carries both}\ \sigma\text{'s and no}\ \rho_{\mathrm B}\ [\text{machine}]
\]

%dsp-level:0
\[
(6):\ \text{the nine factors by monotone exact interval division of the enclosed}\ z\text{'s}\ [\text{machine: widths}<2^{-150}];
\]

%dsp-level:0
\[
\qquad \text{the bracket comparisons against}\ \tfrac{27}{25}\ \text{and}\ 26\ [\text{machine}]
\]

%dsp-level:0
\[
(7)\text{--}(10):\ \text{the assembly of}\ \mathrm{BM.80}\text{'s and}\ \mathrm{BM.82}\text{'s rows};
\]

%dsp-level:0
\[
\qquad \text{the three residue symbols re-cited from}\ \mathrm{BM.69}(5);
\]

%dsp-level:0
\[
\qquad \text{the standing-rule scans}\ [\text{machine}].\ \blacksquare
\]

\textbf{Phys.} Four currencies look at the same two charge classes and no two of them say the same thing. Two of them produce a ratio — four thirds from the ensemble, and the logarithm of one thousand two hundred eighty over the logarithm of nine hundred sixty from the edge count — and two produce a difference: the transport register's log of four thirds over h, and the cone defect's exact algebraic splitting. The tier's central negative is exact rather than numerical. The informational ratio is irrational at every rational coupling, by unique factorization of the primes two, three and five inside the logarithms; four thirds is rational; so the two can never agree, and solving for the coupling at which they would gives a negative number — the logarithm of four hundred five over the logarithm of sixteen — outside the admissible range entirely. The positive side is a localization. Three dimensionless numbers, each a restatement of data the corpus already carries undeclared, are exactly what stands between the census and an absolute inter-family mass: the pole packaging, and the two channel scales. Every cross-currency ratio in the corpus is an exact expression in those three and nothing else; every ratio inside a currency needs none of them. The nine cross-channel factors then show where the structure lives: the three heaviest lines of the three families are degenerate to within three percent while the three lightest differ by a factor of twenty-six — the same light-end concentration the Platonic ladder shows at its top weights, arriving from an entirely different direction. Nothing is converted, nothing is fitted, and the absolute neutrino mass stays exactly where the ledger says it stays.

\begin{theorem}[{\texorpdfstring{BM.84 — The neutron–proton mass ratio: the exact ratio, the irrational critical locus, the range theorem, and the trichotomy}{BM.84 - The neutron-proton mass ratio: the exact ratio, the irrational critical locus, the range theorem, and the trichotomy}}]\label{tag:BM-84}
\textbf{Thm.}

%dsp-level:0
\[
(1)\quad\textbf{the decomposition and the bridge summand's currency row:}
\]

%dsp-level:0
\[
\qquad \Delta m:=m_n-m_p=\Delta m_{\mathrm{Lyapunov}}+\Delta m_{\mathrm{force}}+\Delta m_{\mathrm{metric}}+\Delta m_{\mathrm{binding}}+\Delta m_{\mathrm{self\text{-}int}}+\Delta m_{\mathrm{transport}}
\]

%dsp-level:0
\[
\qquad [\mathrm{BM.62}(1)]:
\]

%dsp-level:0
\[
\Delta m_{\mathrm{Lyapunov}}=0\ [\mathrm{BM.51}(6)];
\]

%dsp-level:0
\[
\qquad \Delta m_{\mathrm{metric}}=0\ [\mathrm{BM.53}(2);\ \mathrm{BM.58}(4)];
\]

%dsp-level:0
\[
\qquad \Delta m_{\mathrm{binding}}=0\ [\mathrm{BM.54}(1)\text{--}(3)];
\]

%dsp-level:0
\[
\qquad \Delta m_{\mathrm{force}}=\Delta m_{\mathrm{self\text{-}int}}=-\Delta m_{\mathrm{EM}}\ [\mathrm{BM.52}(4),(5);\ \mathrm{BM.46}(4)];
\]

%dsp-level:0
\[
\boxed{\ \Delta m_{\mathrm{transport}}=+\frac{\ln(4/3)}{\mathfrak h},\qquad\Delta m_{\mathrm{EM}}=-\,\frac{6\,c\,\varepsilon_{\mathrm{EM}}\ln2}{\mathfrak h},\qquad\Delta m_{\mathrm{bridge}}=0\ \ \text{in the informational currency}\ }
\]

%dsp-level:0
\[
[\mathrm{BM.59}(4);\ \mathrm{BM.46}(4);\ \mathrm{BM.62}(3);\qquad\textbf{the bridge summand's row:}\ \ \Delta m_{\mathrm{bridge}}\ \text{is}\ \neg\textbf{formed} \ \text{in the informational currency}\ —\ \text{the pole-side companion is}\ \Delta_g>0,\ \text{a different currency},
\]

%dsp-level:0
\[
\text{and a cross-currency addition is a category the corpus bars}\ [\mathrm{BM.69}(3);\ \mathrm{BM.35}(5);\ \mathrm{BM.37}(2);\ \mathrm{BM.71}(1)\ \text{row}\ 5];
\]

%dsp-level:0
\[
\qquad \text{the same substitution's two readings point the same way}\ [\mathrm{BM.75}(4),(5)]\ \text{and are}\ \neg\text{addable}];
\]

%dsp-level:0
\[
(2)\quad\textbf{the exact ratio:}
\]

%dsp-level:0
\[
\qquad m_n\,\mathfrak h=6\ln2+\ln20=\ln(2^8\!\cdot\!5)=\ln1280,
\]

%dsp-level:0
\[
\qquad m_p\,\mathfrak h=6\ln2+\ln15+t\ln2=\ln(2^6\!\cdot\!3\!\cdot\!5)+t\ln2=\ln960+t\ln2
\]

%dsp-level:0
\[
\qquad [\mathrm{Def\text{-}BM.13}(5);\ \text{machine}]:
\]

%dsp-level:0
\[
\boxed{\ \frac{m_n}{m_p}\ =\ \frac{\ln1280}{\ln960+6\,c\,\varepsilon_{\mathrm{EM}}\ln2}\ =\ \frac{\ln1280}{\ln960+t\ln2},\qquad t\ :=\ 6\,c\,\varepsilon_{\mathrm{EM}}\ \ge\ 0\ }
\]

%dsp-level:0
\[
\qquad [\mathrm{BM.62}(5)\ \text{re-derived};\ \text{machine: the prime-log assembly}];
\]

%dsp-level:0
\[
\textbf{the sub-critical branch is the branch}\ t<\ln(4/3)/\ln2;
\]

%dsp-level:0
\[
\qquad \textbf{the ratio is a dimensionless element of}\ \mathfrak P,\ \neg\text{a rational number in general}
\]

%dsp-level:0
\[
\qquad [\mathrm{BM.83}(2)];
\]

%dsp-level:0
\[
(3)\quad\textbf{the range theorem (derived here):}
\]

%dsp-level:0
\[
\qquad \frac{\partial}{\partial t}\frac{m_n}{m_p}=-\,\frac{\ln1280\cdot\ln2}{(\ln960+t\ln2)^2}<0
\]

%dsp-level:0
\[
\qquad [\text{machine: the sign census of}\ \ln1280,\ \ln2>0\ \text{in the prime-log ring};\ \text{the monotonicity is cross-multiplied in the Pf}]:
\]

%dsp-level:0
\[
\boxed{\ 0\ <\ \frac{m_n}{m_p}\ \le\ \frac{\ln1280}{\ln960}\ <\ \frac{21}{20}\qquad\text{for every admissible}\ t\ge0,\ \text{with equality at}\ t=0\ \text{only}\ }
\]

%dsp-level:0
\[
\frac{\ln1280}{\ln960}\in\Bigl(\frac{26}{25},\frac{21}{20}\Bigr)
\]

%dsp-level:0
\[
\qquad [\text{machine:}\ 25\ln1280>26\ln960\iff2^{44}>3^{26}\!\cdot\!5;\qquad20\ln1280<21\ln960\iff2^{34}<3^{21}\!\cdot\!5];
\]

%dsp-level:0
\[
2^{44}=17592186044416>12709329141645=3^{26}\!\cdot\!5,
\]

%dsp-level:0
\[
\qquad 2^{34}=17179869184<52301766015=3^{21}\!\cdot\!5
\]

%dsp-level:0
\[
\qquad [\text{machine: the two integer comparisons}];
\]

%dsp-level:0
\[
\textbf{the exact upper bound of the neutron--proton mass ratio on the informational currency is}\ \tfrac{21}{20},
\]

%dsp-level:0
\[
\ \textbf{attained in the limit}\ t\to0^+,
\]

%dsp-level:0
\[
\ \textbf{and no coupling raises it};
\]

%dsp-level:0
\[
(4)\quad\textbf{the irrationality and the rational locus:}
\]

%dsp-level:0
\[
\qquad \boxed{\ \frac{m_n}{m_p}\notin\mathbb Q\ \ \text{at every rational coupling}\ t\in\mathbb Q\ }
\]

%dsp-level:0
\[
\qquad [\mathrm{BM.83}(2)\text{'s theorem re-cited}]:
\]

%dsp-level:0
\[
\frac{m_n}{m_p}=1\iff t=\frac{\ln(4/3)}{\ln2}=2-\log_23\ \notin\ \mathbb Q
\]

%dsp-level:0
\[
\qquad [\text{machine: solved exactly in}\ \mathfrak P;\qquad\log_23\notin\mathbb Q\iff\neg\exists\,(p,q)\ne(0,0):\ 2^p=3^q;\ \ \text{the census over}\ |p|,|q|\le64\ \text{returns none},\ \text{and unique factorization closes the general case}];
\]

%dsp-level:0
\[
\boxed{\ \textbf{the degenerate locus is exactly one coupling, and that coupling is irrational:}\qquad\text{at every rational coupling}\ \ m_n\ne m_p\ }
\]

%dsp-level:0
\[
(5)\quad\textbf{the exact critical strength:}
\]

%dsp-level:0
\[
\qquad \boxed{\ \varepsilon^{\mathrm{crit}}(c)\ =\ \frac{\ln(4/3)}{6\,c\,\ln2}\ =\ \frac{2\ln2-\ln3}{6\,c\,\ln2}\ =\ \frac{2-\log_23}{6\,c}\ }
\]

%dsp-level:0
\[
\qquad [\mathrm{BM.62}(6)\ \text{re-derived};\ \text{machine: the unique zero of}\ \Delta m]:
\]

%dsp-level:0
\[
\textbf{the four integer certificates:}\qquad\frac19<\varepsilon^{\mathrm{crit}}\Bigl(\frac35\Bigr)<\frac18<\varepsilon^{\mathrm{crit}}\Bigl(\frac12\Bigr)<\frac17
\]

%dsp-level:0
\[
\varepsilon^{\mathrm{crit}}(\tfrac35)>\tfrac19\iff2^8>3^5\ (256>243);
\]

%dsp-level:0
\[
\qquad \varepsilon^{\mathrm{crit}}(\tfrac35)<\tfrac18\iff2^{31}<3^{20}\ (2147483648<3486784401);
\]

%dsp-level:0
\[
\varepsilon^{\mathrm{crit}}(\tfrac12)>\tfrac18\iff2^{13}>3^8\ (8192>6561);
\]

%dsp-level:0
\[
\qquad \varepsilon^{\mathrm{crit}}(\tfrac12)<\tfrac17\iff2^{11}<3^7\ (2048<2187)
\]

%dsp-level:0
\[
\qquad [\mathrm{BM.62}(7)\ \text{re-certified};\ \text{machine}];
\]

%dsp-level:0
\[
\boxed{\ \varepsilon^{\mathrm{crit}}(c)\notin\mathbb Q\qquad\text{for every rational}\ c>0\ }\qquad[(4)\text{'s}\ 2-\log_23\notin\mathbb Q;\ \text{machine}];
\]

%dsp-level:0
\[
(6)\quad\textbf{the trichotomy (total, exact, and non-empty on both sides):}
\]

%dsp-level:0
\[
\boxed{\ \varepsilon_{\mathrm{EM}}<\varepsilon^{\mathrm{crit}}(c)\iff m_n>m_p\iff\frac{m_n}{m_p}>1;\qquad\varepsilon_{\mathrm{EM}}=\varepsilon^{\mathrm{crit}}(c)\iff m_n=m_p\iff\frac{m_n}{m_p}=1;\qquad\varepsilon_{\mathrm{EM}}>\varepsilon^{\mathrm{crit}}(c)\iff m_n<m_p\iff\frac{m_n}{m_p}<1\ }
\]

%dsp-level:0
\[
[\mathrm{BM.62}(6),(9)\ \text{re-derived from}\ (2),(3);\ \text{machine: the sign census in}\ \mathfrak P];
\]

%dsp-level:0
\[
\textbf{the opening region is non-empty and integer-certified:}
\]

%dsp-level:0
\[
\qquad \Bigl\{0\le\varepsilon_{\mathrm{EM}}\le\tfrac19\Bigr\}\subseteq\Bigl\{\varepsilon_{\mathrm{EM}}<\varepsilon^{\mathrm{crit}}(c)\Bigr\}\ \ \text{at both priced}\ c
\]

%dsp-level:0
\[
\qquad [\mathrm{BM.62}(8):4^3=64>54=3^3\!\cdot\!2\ \ \text{at}\ c=\tfrac12;\quad4^5=1024>972=3^5\!\cdot\!2^2\ \ \text{at}\ c=\tfrac35;\ \text{machine}];
\]

%dsp-level:0
\[
\textbf{the closing region is non-empty:}
\]

%dsp-level:0
\[
\qquad \varepsilon_{\mathrm{EM}}=\tfrac16>\tfrac17>\varepsilon^{\mathrm{crit}}\bigl(\tfrac12\bigr)\ \ \text{and}\ \ \varepsilon_{\mathrm{EM}}=\tfrac16>\tfrac18>\varepsilon^{\mathrm{crit}}\bigl(\tfrac35\bigr)
\]

%dsp-level:0
\[
\qquad [\text{machine:}\ 6<7,\ 6<8\ \text{and}\ (5)\text{'s brackets}];
\]

%dsp-level:0
\[
\textbf{the critical locus is exact, unique, and irrational}
\]

%dsp-level:0
\[
\qquad [(4),(5)];
\]

%dsp-level:0
\[
\qquad \textbf{the neutron is the heavier baryon on the whole opening region}
\]

%dsp-level:0
\[
\qquad [\mathrm{BM.62}(9)\ \text{re-cited}];
\]

%dsp-level:0
\[
(7)\quad\textbf{the fractional splittings (the exact rational enclosures):}
\]

%dsp-level:0
\[
\qquad \frac{\Delta m}{m_{\mathrm{hid}}}\bigg|_{t=0}=\frac{\ln(4/3)}{6\ln2}\in\Bigl(\frac1{15},\frac1{14}\Bigr)
\]

%dsp-level:0
\[
\qquad [\mathrm{BM.62}(10);\ \text{machine:}\ 2^{24}>3^{15},\ 2^{22}<3^{14}];
\]

%dsp-level:0
\[
\frac{\Delta m}{m_n}\bigg|_{t=0}=\frac{\ln(4/3)}{\ln1280}\in\Bigl(\frac1{25},\frac1{24}\Bigr)
\]

%dsp-level:0
\[
\qquad [\mathrm{BM.62}(10);\ \text{machine}];
\]

%dsp-level:0
\[
\qquad \boxed{\ \frac{\Delta m}{m_p}\bigg|_{t=0}=\frac{\ln(4/3)}{\ln960}\in\Bigl(\frac1{24},\frac1{23}\Bigr)\ }
\]

%dsp-level:0
\[
[\text{machine, derived here}:\quad24\ln\tfrac43>\ln960\iff42\ln2>25\ln3+\ln5\iff2^{42}>3^{25}\!\cdot\!5\ (4398046511104>4236443047215);
\]

%dsp-level:0
\[
23\ln\tfrac43<\ln960\iff40\ln2<24\ln3+\ln5\iff2^{40}<3^{24}\!\cdot\!5\ (1099511627776<1412147682405)];
\]

%dsp-level:0
\[
\textbf{the strict contraction against the barred measure price:}
\]

%dsp-level:0
\[
\qquad \frac{\ln(4/3)}{6\ln2}<\frac13\iff3\ln\tfrac43<2\ln2\iff16<27
\]

%dsp-level:0
\[
\qquad [\mathrm{Def\text{-}BM.13}(8);\ \mathrm{BM.62}(10);\ \text{machine}];
\]

%dsp-level:0
\[
\textbf{the consistency row:}
\]

%dsp-level:0
\[
\qquad \frac{m_n}{m_p}-1=\frac{\Delta m}{m_p}\in\Bigl(\frac1{24},\frac1{23}\Bigr)\ \ \text{at}\ t=0\quad\text{and}\quad\frac{m_n}{m_p}=\frac{\ln1280}{\ln960}\in\Bigl(\frac{26}{25},\frac{21}{20}\Bigr):
\]

%dsp-level:0
\[
\boxed{\ 1+\Bigl(\frac1{24},\frac1{23}\Bigr)=\Bigl(\frac{25}{24},\frac{24}{23}\Bigr)\ \subset\ \Bigl(\frac{26}{25},\frac{21}{20}\Bigr)\ }
\]

%dsp-level:0
\[
\qquad [\text{machine:}\ \tfrac{26}{25}<\tfrac{25}{24}\iff624<625;\quad\tfrac{24}{23}<\tfrac{21}{20}\iff480<483]:
\]

%dsp-level:0
\[
\qquad \text{the finer splitting bracket implies the coarser ratio bracket};
\]

%dsp-level:0
\[
(8)\quad\textbf{what the ratio is not:}
\]

%dsp-level:0
\[
\qquad \frac{m_n}{m_p}\ne\frac43\ \ \text{at every admissible coupling}
\]

%dsp-level:0
\[
\qquad [\mathrm{BM.83}(3);\ \text{machine: the exact solution}\ t=-\ln405/\ln16<0];
\]

%dsp-level:0
\[
\frac{m_n}{m_p}\ne\frac{\ln1280}{\ln960}\ \text{at}\ t>0
\]

%dsp-level:0
\[
\qquad [(3)\text{'s strict monotonicity}];
\]

%dsp-level:0
\[
\qquad \neg\text{a measured neutron--proton ratio is imported, compared or quarantined anywhere in this tier}
\]

%dsp-level:0
\[
\qquad [\mathrm{P5};\ \text{machine: zero external tokens}];
\]

%dsp-level:0
\[
\textbf{the numerical value of}\ e\ \textbf{remains open and does not enter:}
\]

%dsp-level:0
\[
\qquad \varepsilon_{\mathrm{EM}}=e^2/(4\pi)\ \text{is carried symbolically}
\]

%dsp-level:0
\[
\qquad [\mathrm{EM.25}(7),(9);\ \mathrm{BM.62}(12)\ \text{pattern}];
\]

%dsp-level:0
\[
(9)\quad\textbf{status:}
\]

%dsp-level:0
\[
\qquad \textbf{derived:}\ (1)\text{--}(8);
\]

%dsp-level:0
\[
\qquad \textbf{the sign theorem is conditional on the sub-critical branch, with the hypothesis on its face and the critical value exact}\ (5),
\]

%dsp-level:0
\[
(6)\ [\text{A §0.2a}];
\]

%dsp-level:0
\[
\neg\text{a P4 addition},
\]

%dsp-level:0
\[
\quad\neg\text{a new dimensionful scale}\ [\mathrm{P1}],
\]

%dsp-level:0
\[
\quad\neg\text{a route adopted}\ [\mathrm{BM.30}(5)],
\]

%dsp-level:0
\[
\quad\neg\text{a measured value}\ [\mathrm{P5}],
\]

%dsp-level:0
\[
\quad\text{§Cal-BM stays empty};
\]

%dsp-level:0
\[
\textbf{the open row:}\ \text{the numerical value of}\ e\ [\mathrm{EM.25}(7):\ \text{one calibration datum}]\ —\ \text{it decides which side of}\ \varepsilon^{\mathrm{crit}}\ \text{the corpus sits on, and no theorem of Parts 1--4 supplies it}
\]

%dsp-level:0
\[
\qquad [\mathrm{BM.62}(12)\ \text{re-cited}].
\]

\textbf{Req.} EM.25(7), EM.25(9) (Part 2); P1, P4, P5 (Parts 1–2); A §0.2a (Parts 1–2); FT.79(5) (Part 3); Def-BM.13(5), Def-BM.13(8), BM.30(5), BM.35(5), BM.37(2), BM.46(4), BM.51(6), BM.52(4), BM.52(5), BM.53(2), BM.54(1), BM.54(3), BM.58(4), BM.59(4), BM.62(1), BM.62(3), BM.62(5), BM.62(6), BM.62(7), BM.62(8), BM.62(9), BM.62(10), BM.62(12), BM.69(3), BM.71(1), BM.75(4), BM.75(5), BM.83(2), BM.83(3) (this part); exact prime-log-ring arithmetic with integer power certificates, symbolic Laurent differentiation, exact rational enclosures of $\ln$-ratios, unique-factorization irrationality arguments, sign censuses [math; machine: \texttt{checks/\allowbreak{}verify\_\allowbreak{}H\_\allowbreak{}mass\_\allowbreak{}ratio\_\allowbreak{}census.py} §G].

\textbf{Pf.}

%dsp-level:0
\[
(1):\ \text{the five zero-or-monomial rows re-cited from their own tags}\ [\mathrm{BM.51}(6);\ \mathrm{BM.53}(2);\ \mathrm{BM.54}(1)\text{--}(3);\ \mathrm{BM.52}(4),(5);\ \mathrm{BM.46}(4);\ \mathrm{BM.59}(4)];
\]

%dsp-level:0
\[
\qquad \text{the bridge row}:\ \Delta_g
\]

%dsp-level:0
\[
= \hat\mu'^2(D_S(+)-D_S(-))\ \text{has}\ [\Delta_g]
\]

%dsp-level:0
\[
= \mathfrak q\ \text{and}\ [\Delta m]
\]

%dsp-level:0
\[
= \mathfrak q^{+1/2}\ [\text{machine: the exponent arithmetic}]\implies\neg\text{addable}\ [\mathrm{BM.37}(2)\ \text{pattern}]
\]

%dsp-level:0
\[
(2):\ 6\ln2+\ln20=6\ln2+\ln4+\ln5=8\ln2+\ln5=\ln(2^8\cdot5)=\ln1280\ [\text{machine}];
\]

%dsp-level:0
\[
\qquad 6\ln2+\ln15=6\ln2+\ln3+\ln5=\ln(2^6\cdot15)=\ln960\ [\text{machine}];
\]

%dsp-level:0
\[
\qquad t=6c\,\varepsilon_{\mathrm{EM}}\ge0\ \text{since}\ c>0,\ \varepsilon_{\mathrm{EM}}=e^2/(4\pi)\ge0\ [\text{machine}]
\]

%dsp-level:0
\[
(3):\ \text{monotonicity by cross-multiplication:}\ \ t_1<t_2\implies\frac{A}{B+t_1\ln2}>\frac{A}{B+t_2\ln2}
\]

%dsp-level:0
\[
\iff A\bigl(B+t_2\ln2\bigr)>A\bigl(B+t_1\ln2\bigr)
\]

%dsp-level:0
\[
\iff A\,
\]

%dsp-level:0
\[
(t_2-t_1)\ln2>0\ \ \text{at}\ A
\]

%dsp-level:0
\[
= \ln1280,
\]

%dsp-level:0
\[
\ B\qquad = \ln960
\]

%dsp-level:0
\[
[\text{machine: the three positivity rows}\ \ln1280>0,\ \ln960>0,\ \ln2>0\ \text{in the prime-log ring, and the cross-multiplied comparison on a rational grid of pairs}\ (t_1,t_2)];
\]

%dsp-level:0
\[
\qquad \text{the derivative form}\ -A\ln2/(B+t\ln2)^2\ \text{is the same sign census};
\]

%dsp-level:0
\[
\qquad t=0\ \text{gives the supremum}
\]

%dsp-level:0
\[
25\ln1280-26\ln960
\]

%dsp-level:0
\[
= (200-156)\ln2-26\ln3+(25-26)\ln5
\]

%dsp-level:0
\[
= 44\ln2-26\ln3-\ln5>0
\]

%dsp-level:0
\[
\iff 2^{44}>3^{26}\cdot5\ [\text{machine: the two integers}];
\]

%dsp-level:0
\[
\qquad 20\ln1280-21\ln960=34\ln2-21\ln3-\ln5<0\iff2^{34}<3^{21}\cdot5\ [\text{machine}]
\]

%dsp-level:0
\[
(4):\ \frac{m_n}{m_p}
\]

%dsp-level:0
\[
= 1\iff\ln1280
\]

%dsp-level:0
\[
= \ln960+t\ln2
\]

%dsp-level:0
\[
\iff t
\]

%dsp-level:0
\[
= \frac{\ln1280-\ln960}{\ln2}
\]

%dsp-level:0
\[
= \frac{\ln(4/3)}{\ln2}
\]

%dsp-level:0
\[
= \frac{2\ln2-\ln3}{\ln2}
\]

%dsp-level:0
\[
= 2-\log_23\ [\text{machine: solved in}\ \mathfrak P];
\]

%dsp-level:0
\[
\log_23\in\mathbb Q\implies\exists\,(p,q)\ne(0,0):\ 2^p=3^q\ [\text{clearing denominators}];
\]

%dsp-level:0
\[
\qquad \text{machine: no solution with}\ |p|,|q|\le64;
\]

%dsp-level:0
\[
\qquad \text{math: unique factorization excludes every case}\ (p,q>0:\ \text{parity of the left side};\ p\le0\ \text{or}\ q\le0:\ \text{the reciprocal argument})
\]

%dsp-level:0
\[
(5):\ \Delta m=0\iff\ln(4/3)=t\iff6c\,\varepsilon_{\mathrm{EM}}=\ln(4/3)\iff\varepsilon_{\mathrm{EM}}=\frac{\ln(4/3)}{6c\ln2}\ [\text{machine}];
\]

%dsp-level:0
\[
\qquad \text{the four cross-multiplications of}\ \mathrm{BM.62}(7)\ \text{re-certified as the four integer power comparisons}\ [\text{machine}]
\]

%dsp-level:0
\[
(6):\ (3)\text{'s monotonicity}\ \wedge\ (4)\text{'s unique zero}\implies\text{the three-way sign equivalence}\ [\text{machine}];
\]

%dsp-level:0
\[
\qquad \text{the opening witnesses}\ 4^3=64>54=3^3\cdot2\ \text{and}\ 4^5=1024>972=3^5\cdot2^2\ [\text{machine}];
\]

%dsp-level:0
\[
\qquad \text{the closing witness}\ \tfrac16>\tfrac17>\varepsilon^{\mathrm{crit}}(\tfrac12)\ \text{and}\ \tfrac16>\tfrac18>\varepsilon^{\mathrm{crit}}(\tfrac35)\ [\text{machine:}\ 6<7,\ 6<8,\ (5)\text{'s brackets}]
\]

%dsp-level:0
\[
(7):\ \text{the two frozen enclosures re-certified}\ [\text{machine}];
\]

%dsp-level:0
\[
\qquad 24\ln\tfrac43-\ln960=(48-6)\ln2-(24+1)\ln3-\ln5=42\ln2-25\ln3-\ln5>0\iff2^{42}>3^{25}\cdot5\ [\text{machine}];
\]

%dsp-level:0
\[
\qquad 23\ln\tfrac43-\ln960=40\ln2-24\ln3-\ln5<0\iff2^{40}<3^{24}\cdot5\ [\text{machine}]
\]

%dsp-level:0
\[
3\ln\tfrac43<2\ln2\iff\ln\tfrac{64}{27}<\ln4\iff64<108\ [\text{machine}];
\]

%dsp-level:0
\[
\qquad 1+\tfrac1{24}
\]

%dsp-level:0
\[
= \tfrac{25}{24}\ \text{and}\ \tfrac{25}{24}<\tfrac{26}{25}
\]

%dsp-level:0
\[
\iff 625<624\ \text{is false}\implies\text{the two brackets overlap and are reported separately}\ [\text{machine}]
\]

%dsp-level:0
\[
(8):\ \mathrm{BM.83}(3)\text{'s exact solution re-cited};
\]

%dsp-level:0
\[
\qquad (3)\text{'s strict monotonicity};
\]

%dsp-level:0
\[
\qquad \text{the token scans}\ [\text{machine}];
\]

%dsp-level:0
\[
\qquad (9):\ \text{the standing-rule scans and the assembly}\ [\text{machine}].\ \blacksquare
\]

\textbf{Phys.} The neutron–proton mass ratio is one exact expression: the logarithm of one thousand two hundred eighty over the logarithm of nine hundred sixty plus a nonnegative multiple of the logarithm of two. Everything about it follows by arithmetic in the ring generated by the logarithms of the primes. It decreases strictly with the coupling, so its largest value is the coupling-free one, and that value sits between twenty-six twenty-fifths and twenty-one twentieths — an exact upper bound of one and one twentieth on the ratio, valid at every coupling the corpus admits. It is irrational at every rational coupling, by unique factorization of two, three and five inside the logarithms, and it equals one at exactly one coupling, which is two minus the base-two logarithm of three and is therefore irrational too: the degenerate point of the trichotomy is provably not a rational coupling, and at every rational coupling the two baryons differ. The critical strength inherits the same irrationality and the same four integer certificates that enclose it between a ninth and a seventh, so the opening region is a wide exact interval, not a tuned point, and the closing region is inhabited by anything above a sixth. The bridge term contributes nothing here and cannot: it lives in the angle currency, it points the same way, and adding it to a bit count is the one operation the ledger forbids. No measured ratio enters, no external value is compared, and the single open datum — the numerical value of the electric coupling — is the same one the corpus has named at every tier, deciding which side of the critical surface the world sits on without ever being supplied.

\end{theorem}

\begin{theorem}[{\texorpdfstring{BM.85 — The Q-value in two formed reads: the ratio forms, the margin ladders, and the read-relative verdict}{BM.85 - The Q-value in two formed reads: the ratio forms, the margin ladders, and the read-relative verdict}}]\label{tag:BM-85}
\textbf{Thm.}

%dsp-level:0
\[
(1)\quad\textbf{the Q-value on the informational currency (formed):}
\]

%dsp-level:0
\[
\qquad Q_\beta:=m_n-m_p-m_\ell-m_\nu=\Delta m-m_\ell-m_\nu
\]

%dsp-level:0
\[
\qquad [\mathrm{BM.41}(4);\ \mathrm{BM.49}(1);\ \mathrm{BM.63}(1)]:
\]

%dsp-level:0
\[
\boxed{\ Q_\beta\,\mathfrak h\ =\ \ln\frac43-t\ln2-F_{\ell\nu},\qquad F_{\ell\nu}:=m_\ell\mathfrak h+m_\nu\mathfrak h,\qquad t=6\,c\,\varepsilon_{\mathrm{EM}}\ge0\ }
\]

%dsp-level:0
\[
\qquad [\mathrm{BM.62}(4);\ \mathrm{BM.70}(1);\ \text{machine: the prime-log assembly}];
\]

%dsp-level:0
\[
(2)\quad\textbf{the two-bit daughter floor and the closure (re-certified):}
\]

%dsp-level:0
\[
\qquad F_{\ell\nu}\ge2\ln2
\]

%dsp-level:0
\[
\qquad [\mathrm{BM.70}(2),(3):\ \text{every materialized closed chain carries}\ \ge2\ \text{directed edges};\ m\mathfrak h=E_{\mathrm{int}}\ln2;\ \ m_\nu\mathfrak h\ge0\ \text{by PSD}]:
\]

%dsp-level:0
\[
\boxed{\ Q_\beta\,\mathfrak h\ \le\ \ln\frac43-2\ln2\ =\ \ln\frac43-\ln4\ =\ -\ln3\ <\ 0\ }
\]

%dsp-level:0
\[
\qquad [\text{machine:}\ 4<12];
\]

%dsp-level:0
\[
\qquad \textbf{formed}\implies\textbf{closed on the informational currency}
\]

%dsp-level:0
\[
\qquad [\mathrm{BM.70}(4)(\mathrm a)\ \text{re-cited}];
\]

%dsp-level:0
\[
\textbf{the margin ladder:}
\]

%dsp-level:0
\[
\qquad Q_\beta\mathfrak h\ \le\ -\ln3,\quad-\ln12,\quad-\ln24
\]

%dsp-level:0
\[
\qquad \text{on the certificates}\quad4<12,\quad4<48,\quad4<96
\]

%dsp-level:0
\[
\qquad [\mathrm{BM.70}(3);\ \text{machine}];
\]

%dsp-level:0
\[
\textbf{each margin is deepened by}\ t\ln2\ge0\ \ \textbf{at every positive coupling}\qquad[\text{machine: the sign census}];
\]

%dsp-level:0
\[
(3)\quad\textbf{the ratio forms (derived here):}
\]

%dsp-level:0
\[
\qquad \frac{Q_\beta}{m_n}\bigg|_{t=0}\ \le\ -\,\frac{\ln3}{\ln1280}\ \in\ \Bigl(-\frac2{13},-\frac3{20}\Bigr),
\]

%dsp-level:0
\[
\qquad \frac{Q_\beta}{m_p}\bigg|_{t=0}\ \le\ -\,\frac{\ln3}{\ln960}\ \in\ \Bigl(-\frac4{25},-\frac3{19}\Bigr)
\]

%dsp-level:0
\[
[\text{machine:}\ -\tfrac2{13}<-\tfrac{\ln3}{\ln1280}\iff13\ln3<2\ln1280\iff3^{13}<2^{16}\!\cdot\!5^2\ (1594323<1638400);
\]

%dsp-level:0
\[
-\tfrac{\ln3}{\ln1280}<-\tfrac3{20}\iff20\ln3>3\ln1280\iff3^{20}>2^{24}\!\cdot\!5^3\ (3486784401>2097152000);
\]

%dsp-level:0
\[
-\tfrac4{25}<-\tfrac{\ln3}{\ln960}\iff25\ln3<4\ln960\iff3^{21}<2^{24}\!\cdot\!5^4\ (10460353203<10485760000);
\]

%dsp-level:0
\[
-\tfrac{\ln3}{\ln960}<-\tfrac3{19}\iff19\ln3>3\ln960\iff3^{16}>2^{18}\!\cdot\!5^3\ (43046721>32768000)];
\]

%dsp-level:0
\[
\boxed{\ \textbf{the informational Q-value is a strictly negative exact fraction of the parent mass, bounded away from zero by a rational}\ }
\]

%dsp-level:0
\[
\qquad \text{in both normalizations};
\]

%dsp-level:0
\[
(4)\quad\textbf{the Q-value on the specified read (formed, branch-relative):}
\]

%dsp-level:0
\[
\qquad Q_\beta^{\mathrm{spec}}:=m_n^{(\mathrm{ii})}-m_p^{(\mathrm{ii})}-m_\ell-m_\nu
\]

%dsp-level:0
\[
\qquad [\mathrm{BM.77}(1),(2);\ \text{A-BM.8}]:
\]

%dsp-level:0
\[
\boxed{\ Q_\beta^{\mathrm{spec}}\,\mathfrak h\ =\ \Delta T(S)-\rho_\ell\sqrt{G\bigl(c(0)\bigr)}-\rho_\nu\sqrt{G\bigl(c(i_\nu)\bigr)},\qquad\Delta T(S)=5\,m^2_B(\cdot,S)\ }
\]

%dsp-level:0
\[
\Delta T=\Bigl(\frac{4219785}{98956},\ \frac{9608310}{393679},\ \frac{455}{23},\ \frac{68565}{3088}\Bigr)\ \ \text{at}\ S=(0,1,2,3),
\]

%dsp-level:0
\[
\qquad \Delta T_{\min}=\frac{455}{23},
\]

%dsp-level:0
\[
\qquad \Delta T_{\max}=\frac{4219785}{98956}
\]

%dsp-level:0
\[
\qquad [\mathrm{BM.77}(2);\ \text{machine}];
\]

%dsp-level:0
\[
\sqrt{A_0-B_0}\ \le\ \sqrt{G\bigl(c(i)\bigr)}\ \le\ \sqrt{A_0+B_0},
\]

%dsp-level:0
\[
\qquad A_0\pm B_0\big|_{g=g'=1}=\frac{11}{16},\ \frac3{16}
\]

%dsp-level:0
\[
\qquad [\mathrm{BM.39}(1);\ \mathrm{BM.77}(4);\ \text{machine}];
\]

%dsp-level:0
\[
\textbf{the trichotomy with integer-certified witnesses on both sides:}
\]

%dsp-level:0
\[
\qquad \rho_\ell=\rho_\nu=\frac{\Delta T(S)}{4\sqrt{A_0+B_0}}\in\mathcal O_S,
\]

%dsp-level:0
\[
\qquad \rho_\ell=\rho_\nu=\frac{\Delta T(S)}{\sqrt{A_0-B_0}}\in\mathcal C_S
\]

%dsp-level:0
\[
\qquad [\mathrm{BM.77}(5)(\mathrm a),(\mathrm b)]:
\]

%dsp-level:0
\[
Q_\beta^{\mathrm{spec}}\mathfrak h\ge\frac{455}{23}-\frac{\sqrt{11}}2>0\iff828100>5819;
\]

%dsp-level:0
\[
\qquad Q_\beta^{\mathrm{spec}}\mathfrak h\le\frac{4219785}{98956}-50\sqrt3<0\ \ \text{on}\ \ 4219785<43\cdot98956=4255108\ \ \text{and}\ \ 7500>1849;
\]

%dsp-level:0
\[
\frac{273}{23}\ <\ \rho^{\mathrm{crit}}(S,i_\nu)\ <\ \frac{42197850}{791648}
\]

%dsp-level:0
\[
\qquad \text{at the witness couplings}\ g=g'=1,\ \rho_\ell=\rho_\nu
\]

%dsp-level:0
\[
\qquad [\mathrm{BM.77}(5)(\mathrm c);\ \text{machine:}\ 99<100,\ 75>64,\ 273>253,\ 42197850<42748992];
\]

%dsp-level:0
\[
(5)\quad\textbf{the two reads are}\ \neg\textbf{comparable, and both verdicts stand:}\qquad\textbf{the alphabet census:}
\]

%dsp-level:0
\[
Q_\beta^{\mathrm{info}}:\ \{\ln2,\ln3,\ln5,\ c,\ \varepsilon_{\mathrm{EM}},\ \mathfrak h,\ m_\ell,\ m_\nu\};
\]

%dsp-level:0
\[
\qquad Q_\beta^{\mathrm{spec}}:\ \{m^2_B(\cdot,S),\ \rho_\ell,\ \rho_\nu,\ g,\ g',\ \mathfrak h\}
\]

%dsp-level:0
\[
\qquad [\text{machine: the token census}]:
\]

%dsp-level:0
\[
\boxed{\ \text{the two alphabets are disjoint except for}\ \mathfrak h\ }
\]

%dsp-level:0
\[
\qquad [\mathrm{BM.37}(2);\ \mathrm{BM.62}(11);\ \mathrm{BM.66}(6)\ \text{pattern};\ \text{machine}];
\]

%dsp-level:0
\[
\neg\text{a derived equation relates the two expressions}\ —\ \text{that is}\ (\mathrm{R1})\ \text{and}\ (\mathrm{R3})\ \text{and the Cal-F scales together}
\]

%dsp-level:0
\[
\qquad [\mathrm{BM.82}(6);\ \mathrm{BM.83}(5),(7)];
\]

%dsp-level:0
\[
\textbf{the verdicts, read-relative and complete:}
\]

%dsp-level:0
\[
\qquad \text{closed on the informational currency for every member}\ [\mathrm{BM.70}(4)(\mathrm a)];
\]

%dsp-level:0
\[
\qquad \text{open on}\ \mathcal O_S,\ \text{closed on}\ \mathcal C_S,\ \text{degenerate on}\ \partial\mathcal O_S\ \text{on the specified read}\ [\mathrm{BM.77}(5),(6)];
\]

%dsp-level:0
\[
\neg\text{a clause of}\ \mathrm{BM.41},\ \mathrm{BM.49},\ \mathrm{BM.63},\ \mathrm{BM.70},\ \mathrm{BM.71}\ \text{is retracted};
\]

%dsp-level:0
\[
\qquad \text{the rescopes are the ones}\ \mathrm{BM.77}(7)\ \text{records}
\]

%dsp-level:0
\[
\qquad [\text{machine: the clause census}];
\]

%dsp-level:0
\[
(6)\quad\textbf{the parent row is currency-independent:}
\]

%dsp-level:0
\[
\qquad T(0,S)=\tfrac43T(+1,S)>T(+1,S)
\]

%dsp-level:0
\[
\qquad [\mathrm{BM.28}(4)(\mathrm{ii});\ \text{machine}]
\]

%dsp-level:0
\[
\qquad \text{and}
\]

%dsp-level:0
\[
\qquad m_n>m_p\ \ \text{on the sub-critical branch}
\]

%dsp-level:0
\[
\qquad [\mathrm{BM.62}(9);\ \mathrm{BM.84}(6)]:
\]

%dsp-level:0
\[
\boxed{\ \textbf{the decay direction}\ n\to p+\ell+\nu\ \textbf{is fixed in both currencies}\ }
\]

%dsp-level:0
\[
\qquad [\mathrm{BM.41}(1);\ \mathrm{BM.45};\ \mathrm{BM.77}(6):\ \text{the grammar stands unmodified}];
\]

%dsp-level:0
\[
\textbf{the two-body companion on the specified read:}
\]

%dsp-level:0
\[
\qquad \Delta T(S)\ge\rho_\ell\sqrt{G\bigl(c(0)\bigr)}\ \ \text{holds at the witness}\ \Bigl[\tfrac{455}{23}>\tfrac{\sqrt{11}}4\iff3312400>5819\Bigr],
\]

%dsp-level:0
\[
\qquad \mathcal O_S\subseteq\{\text{two-body admissible}\}
\]

%dsp-level:0
\[
\qquad [\mathrm{BM.77}(6);\ \text{machine}];
\]

%dsp-level:0
\[
(7)\quad\textbf{the conditional spectrum becomes branch content on}\ \mathcal O_S\textbf{:}
\]

%dsp-level:0
\[
\qquad E_\ell^{\max}=\frac{m_n^2+m_\ell^2-(m_p+m_\nu)^2}{2m_n},
\]

%dsp-level:0
\[
\qquad E_{\max}:\langle E\rangle=2:1\ \ \text{scale-invariant}
\]

%dsp-level:0
\[
\qquad [\mathrm{BM.41}(5);\ \mathrm{BM.49}(3),(4);\ \mathrm{BM.77}(8)\ \text{re-cited}]:
\]

%dsp-level:0
\[
\neg\text{an external endpoint, spectrum or lifetime is compared}
\]

%dsp-level:0
\[
\qquad [\mathrm{P5}];
\]

%dsp-level:0
\[
\qquad \text{the quarantined audit rows stay quarantined}
\]

%dsp-level:0
\[
\qquad [\text{Census FT.114},\ \text{FT.120}\ \text{pattern}];
\]

%dsp-level:0
\[
(8)\quad\textbf{status:}
\]

%dsp-level:0
\[
\qquad \textbf{derived:}\ (1)\text{--}(7);
\]

%dsp-level:0
\[
\qquad \textbf{read-relative:}\ \text{the kinematic verdict}\ (5);
\]

%dsp-level:0
\[
\qquad \textbf{declared: nothing here}\ —\ \text{A-BM.8's specification is re-cited,}\ \neg\text{re-issued};
\]

%dsp-level:0
\[
\neg\text{a P4 addition},\quad\neg\text{a new dimensionful scale}\ [\mathrm{P1}],\quad\neg\text{a measured value}\ [\mathrm{P5}],\quad\neg\text{a currency converted};
\]

%dsp-level:0
\[
\qquad \text{§Cal-BM stays empty}.
\]

\textbf{Req.} P1, P4, P5 (Parts 1–2); FT.79(5) (Part 3); Census FT.114, FT.120 (row-wise); Def-BM.13(5), BM.28(4), BM.37(2), BM.39(1), BM.41(1), BM.41(4), BM.41(5), BM.45, BM.49(1), BM.49(3), BM.49(4), BM.62(4), BM.62(9), BM.62(11), BM.63(1), BM.66(6), BM.70(1), BM.70(2), BM.70(3), BM.70(4), BM.71, BM.77(1), BM.77(2), BM.77(4), BM.77(5), BM.77(6), BM.77(7), BM.77(8), BM.82(6), BM.83(5), BM.83(7), BM.84(6) (this part); the declaration addendum A-BM.8, item 7 (pointers); exact prime-log arithmetic with integer power certificates, exact rational arithmetic on the four frozen sector expectations, exact radical inequalities by integer certificates, token censuses [math; machine: \texttt{checks/\allowbreak{}verify\_\allowbreak{}H\_\allowbreak{}mass\_\allowbreak{}ratio\_\allowbreak{}census.py} §H].

\textbf{Pf.}

%dsp-level:0
\[
(1):\ \Delta m=(\ln\tfrac43-t\ln2)/\mathfrak h\ [\mathrm{BM.62}(4)\ \text{re-cited}];
\]

%dsp-level:0
\[
\qquad Q_\beta\mathfrak h=\ln\tfrac43-t\ln2-(m_\ell\mathfrak h+m_\nu\mathfrak h)\ [\text{machine: the monomial sum in}\ \mathfrak P]
\]

%dsp-level:0
\[
(2):\ F_{\ell\nu}\ge2\ln2\ [\mathrm{BM.70}(2),(3)\ \text{re-cited}];
\]

%dsp-level:0
\[
\qquad \ln\tfrac43-2\ln2=2\ln2-\ln3-2\ln2=-\ln3\ [\text{machine}];
\]

%dsp-level:0
\[
\qquad -\ln3<0\iff3>1\iff4<12\ [\text{machine}];
\]

%dsp-level:0
\[
\qquad \text{the ladder on}\ 4<12,\ 4<48,\ 4<96\ [\text{machine}]
\]

%dsp-level:0
\[
(3):\ Q_\beta/m_n=(Q_\beta\mathfrak h)/(m_n\mathfrak h)\le(-\ln3)/\ln1280\ \text{at}\ t=0\ [\text{machine:}\ m_n\mathfrak h=\ln1280>0];
\]

%dsp-level:0
\[
\qquad \text{the four enclosures by the four integer power comparisons of}\ (3)\ [\text{machine}]
\]

%dsp-level:0
\[
(4):\ \text{the four}\ \Delta T\ \text{values re-cited}\ [\mathrm{BM.77}(2);\ \text{machine:}\ 5\cdot\tfrac{843957}{98956},\ 5\cdot\tfrac{1921662}{393679},\ 5\cdot\tfrac{91}{23},\ 5\cdot\tfrac{13713}{3088}\ \text{as exact integers}];
\]

%dsp-level:0
\[
\qquad \text{the sandwich and both witnesses re-certified}\ [\mathrm{BM.39}(1);\ \mathrm{BM.77}(4),(5);\ \text{machine}]
\]

%dsp-level:0
\[
(5):\ \text{the token census of the two expressions}\ [\text{machine}];
\]

%dsp-level:0
\[
\qquad \text{the clause census of}\ \mathrm{BM.41}\text{--}\mathrm{BM.71}\ \text{against retraction}\ [\text{machine}];
\]

%dsp-level:0
\[
\qquad \text{the no-equation row re-cited from}\ \mathrm{BM.82}(6),\ \mathrm{BM.83}(5),(7)
\]

%dsp-level:0
\[
(6):\ T(0,S)/T(+1,S)=W(0)/W(+1)=20/15=4/3\ [\text{machine}];
\]

%dsp-level:0
\[
\qquad m_n>m_p\ \text{sub-critically}\ [\mathrm{BM.84}(6)];
\]

%dsp-level:0
\[
\qquad \tfrac{455}{23}>\tfrac{\sqrt{11}}4\iff1820>23\sqrt{11}\iff3312400>5819\ [\text{machine: squaring both positive sides}]
\]

%dsp-level:0
\[
(7):\ \text{the endpoint identity and the}\ 2{:}1\ \text{moment ratio re-cited}\ [\mathrm{BM.41}(5);\ \mathrm{BM.49}(3),(4)];
\]

%dsp-level:0
\[
\qquad (8):\ \text{the standing-rule scans}\ [\text{machine}].\ \blacksquare
\]

\textbf{Phys.} One Q-value, two formed reads, and no equation between them. Read in bits, the decay's energy balance is the transport register's log of four thirds, less the electromagnetic scar's coupling term, less whatever the two emitted lines cost — and the two-bit floor on any materialized chain drives the total below minus the logarithm of three at every coupling, so on that currency the vertex is closed with an exact margin, and the margin is a rational fraction of the parent's own mass: worse than minus three twentieths, better than minus two thirteenths. Read on the specified route, the same vertex is a difference of five times a sector expectation against two undeclared scale combinations, and its sign is decided by where those scales sit against an exact rational: open on a region certified by eight hundred twenty-eight thousand one hundred against five thousand eight hundred nineteen, closed on its complement, degenerate on the boundary between eleven and fifty-four. The two expressions share one symbol — h — and nothing else, so the two verdicts are not in conflict and cannot be reconciled without the very bridge the ledger holds open. What both currencies do agree on, and what no read can change, is the direction: the neutral class is the heavier one in the ensemble, on the informational branch below the critical coupling, and on the specified read at every winding cell, so the parent is the neutron and the grammar stands. The decay is allowed where a read is formed that allows it, closed where a read is formed that closes it, and the census reports both without converting either.

\end{theorem}

\medskip\noindent\textit{\texorpdfstring{BM.86 — The status of the mass ratio census: adjudication and scope}{BM.86 - The status of the mass ratio census: adjudication and scope}.}\ Structural-constraints record: see Appendix B, Section B.1.26 of the Addendum.\par\medskip

\phantomsection
\section[Tier BM.18 — The EM Self-Energy Correction: the Category…]{\texorpdfstring{Tier BM.18 — The EM Self-Energy Correction: the Category Adjudication, the Running Coupling, the Path-Integral Self-Energy, the Branch-Blindness Theorem, and the Corrected Neutron–Proton Splitting}{Tier BM.18 - The EM Self-Energy Correction: the Category Adjudication, the Running Coupling, the Path-Integral Self-Energy, the Branch-Blindness Theorem, and the Corrected Neutron-Proton Splitting}}

\textbf{Locus.} This tier relocates one computation and voids one coefficient. The EM self-energy of a sealed composite is taken off the chart — where it was read as an integral over a charge distribution $[\mathrm{Def\text{-}BM.12}(5);\ \mathrm{BM.46}(1)]$ — and put on the flicker ensemble, where the corpus's own charge datum lives: a ternary end read $q\in\{-1,0,+1\}$ on the endpoints of chains $[\mathrm{D3.4};\ \mathrm{EM.9}]$ with no position datum attached $[\mathrm{BM.28}(7);\ \mathrm{BM.46}(2)]$. The distribution factor $c$ that priced a shell against a ball is adjudicated a classical artifact and voided as a mass datum; the exterior value it carried is derived, twice, from rows that need no density ($\mathrm{BM.87}$). The coupling stops being a scalar and becomes an exact geometric profile on the envelope ratio, the lapse and the curvature lock ($\mathrm{BM.88}$). The self-energy becomes an ensemble average with one derived attenuation, the crystallization weight ($\mathrm{BM.89}$). The proposal that the six-state defect splitting reweight that attenuation is put on its face and refuted on four grounds ($\mathrm{BM.90}$). The neutron–proton splitting, its critical strength, its ratio range and its trichotomy are recomputed with the corrected term ($\mathrm{BM.91}$), the Q-value rows are re-certified on both formed reads ($\mathrm{BM.92}$), and $\mathrm{BM.93}$ adjudicates.

\textbf{Standing rows.} $\alpha_B$ stays barred on $\mathrm{BM.61}(8)$'s four grounds and §Cal-BM stays empty. No clause of $\mathrm{BM.1}$–$\mathrm{BM.86}$ is retracted, replaced or renumbered: every row this tier touches is re-scoped at its own locus and recorded in $\mathrm{BM.87}(6)$, $\mathrm{BM.91}(8)$ and $\mathrm{BM.93}(4)$, in $\mathrm{BM.77}(7)$'s pattern. The numerical value of $e$ stays open $[\mathrm{EM.25}(7);\ \mathrm{BM.62}(12);\ \mathrm{BM.84}(9)]$ and no measured coupling, mass or ratio enters any derivation $[\mathrm{P5}]$. The relating equation stays $\neg$derived — the mass-form bridge is application-layer content $[\mathrm{BM.35}(5);\ \mathrm{BM.69}(4),(5);\ \mathrm{BM.83}(5)]$. The currencies stay distinct: the ensemble weight derived here attenuates the informational register alone and converts nothing into the defect register $[\mathrm{BM.82}(4);\ \mathrm{BM.83}(3)]$. Local notation, fixed for the whole tier: $k:=r/\Delta x_{\mathrm B}$ the envelope ratio; $f(k)=1-\tfrac1k$ the nearby-space fraction; $N$ the lapse; $G:=(\Delta x_{\mathrm B})^2\kappa^2\rho$ the curvature lock; $w$ the crystallization weight; $\mathsf g$ the curvature composition (the italic $g$ keeps its $\mathrm{BM.39}(1)$ meaning, the $\mathsf{SU}(2)$ coupling, at every occurrence); $\varepsilon_{\mathrm{EM}}=e^2/(4\pi)$; the symbol $c$ is void as a mass datum and occurs in this tier only inside barred-object and rescope rows.

\medskip\noindent\textit{\texorpdfstring{BM.87 — The category adjudication: the void distribution factor, the derived exterior value, and the standing rows}{BM.87 - The category adjudication: the void distribution factor, the derived exterior value, and the standing rows}.}\ Structural-constraints record: see Appendix B, Section B.1.27 of the Addendum.\par\medskip

\begin{theorem}[{\texorpdfstring{BM.88 — The running coupling: the exact geometric profile on the envelope ratio, the lapse, and the curvature lock}{BM.88 - The running coupling: the exact geometric profile on the envelope ratio, the lapse, and the curvature lock}}]\label{tag:BM-88}
\textbf{Thm.}

%dsp-level:0
\[
(1)\quad\textbf{the coupling is a function,}\ \neg\textbf{a scalar:}
\]

%dsp-level:0
\[
\qquad e=\frac{gg'}{\sqrt{g^2+g'^2}}=g\sin\theta_W
\]

%dsp-level:0
\[
\qquad [\mathrm{EW.35};\ \mathrm{EM.25}(6):\ \text{fixed with zero freedom as a combination}];
\]

%dsp-level:0
\[
\neg\exists\ \text{a theorem fixing the numerical value of}\ e
\]

%dsp-level:0
\[
\qquad [\mathrm{EM.25}(7);\ \mathrm{P4}];
\]

%dsp-level:0
\[
\qquad \textbf{the running in this corpus is geometric,}\ \neg\textbf{perturbative:}
\]

%dsp-level:0
\[
\qquad \neg\exists\ \text{a loop, a}\ \beta\text{-function, a counterterm or a Yukawa-type object in the inventory}
\]

%dsp-level:0
\[
[\mathrm{BM.61}(8);\ \mathrm{FT.145}(3);\ \mathrm{EW.15};\ \mathrm{BM.47}(2):\ \text{the divergence is cut by the envelope,}\ \neg\text{by a counterterm}];
\]

%dsp-level:0
\[
(2)\quad\textbf{the three derived dimensionless arguments:}
\]

%dsp-level:0
\[
\qquad \textbf{(i) the envelope ratio:}\quad k\ :=\ \frac{r}{\Delta x_{\mathrm B}(E)},
\]

%dsp-level:0
\[
\qquad \Delta x_{\mathrm B}(E)=\frac1{m_{\mathrm{hid}}(E)}=\frac{\mathfrak h}{E_{\mathrm{int}}(E)\ln2}
\]

%dsp-level:0
\[
\qquad [\mathrm{Def\text{-}BM.4}(1);\ \mathrm{BM.13}(1);\ \text{machine:}\ \Delta x_{\mathrm B}\cdot m_{\mathrm{hid}}=1\ \text{as a Laurent identity}]:
\]

%dsp-level:0
\[
E_{\mathrm{int}}\in\{2,\,6,\,12,\,24,\,24,\,60,\,60\}\ \ \text{on the computed cells}
\]

%dsp-level:0
\[
\qquad [\mathrm{BM.13}(2);\ \mathrm{FT.28};\ \text{Census FT.35}(2),\ \text{FT.42}(5)\ \text{row-wise};\ \mathrm{BM.13}(3):\ E_{\mathrm{int}}\ge2]:
\]

%dsp-level:0
\[
\text{a strictly decreasing discrete ladder of widths},
\]

%dsp-level:0
\[
\qquad k\ge1\ \ \text{on the formed locus}\ r\ge\Delta x_{\mathrm B}
\]

%dsp-level:0
\[
\qquad [\text{machine: injective in}\ E_{\mathrm{int}}];
\]

%dsp-level:0
\[
\textbf{(ii) the nearby-space fraction:}
\]

%dsp-level:0
\[
\qquad \boxed{\ f(k)\ :=\ 1-\frac1k\ =\ 1-\frac{\Delta x_{\mathrm B}(E)}{r}\ }
\]

%dsp-level:0
\[
\qquad [\mathrm{BM.47}(6);\ \text{machine: exact rational functions}]:
\]

%dsp-level:0
\[
f(1)=0,\quad f(2)=\frac12,\quad f(3)=\frac23,\quad f(4)=\frac34,\quad f(5)=\frac45,\quad f(6)=\frac56;
\]

%dsp-level:0
\[
\qquad f\ \text{strictly increasing on}\ k\ge1,
\]

%dsp-level:0
\[
\qquad 0\le f(k)<1\ \ \forall k<\infty
\]

%dsp-level:0
\[
\qquad [\text{machine}];
\]

%dsp-level:0
\[
\textbf{(iii) the lapse:}
\]

%dsp-level:0
\[
\qquad N=\sqrt{\Omega^2+\gamma^{ij}N_iN_j},
\]

%dsp-level:0
\[
\qquad \gamma_{ij}=\Omega^2(\delta_{ij}+\varkappa q_{ij}),
\]

%dsp-level:0
\[
\qquad N_i=\Omega^2\varkappa\beta w_i
\]

%dsp-level:0
\[
\qquad [\mathrm{CT1.12};\ \mathrm{BM.34}(1);\ \mathrm{FT.157}(1):\ \text{the shift drag}]:
\]

%dsp-level:0
\[
\boxed{\ N^2\det\gamma=-\det g\ }
\]

%dsp-level:0
\[
\qquad [\mathrm{BM.34}(1);\ \text{machine: the exact polynomial identity over}\ \mathbb Q[\Omega^2,\varkappa,\beta,w_i,q_{ij}]];
\]

%dsp-level:0
\[
(a_{00},q_{ij},w_i)=0\implies(\Omega,N,N^i,g_{ij})=(1,1,0,\delta_{ij})
\]

%dsp-level:0
\[
\qquad [\mathrm{BM.34}(2):\ \textbf{the vacuum's emergent metric is flat};\ \text{machine: exact substitution}]\implies\boxed{\ N=1\ \ \text{in the vacuum gauge}\ }
\]

%dsp-level:0
\[
(3)\quad\textbf{the curvature lock (the derived dimensionless curvature argument):}
\]

%dsp-level:0
\[
\qquad \kappa^2\rho=\frac{2\pi E_{\mathrm{int}}^4(\ln2)^3}{\mathfrak h^2},
\]

%dsp-level:0
\[
\qquad \rho=\frac{(E_{\mathrm{int}}\ln2)^4}{\mathfrak h^4},
\]

%dsp-level:0
\[
\qquad [\rho]=\mathfrak q^{+2}
\]

%dsp-level:0
\[
\qquad [\mathrm{BM.15}(2);\ \mathrm{BM.14}(2),(4);\ \mathrm{CT1.23}]:
\]

%dsp-level:0
\[
\boxed{\ G(E)\ :=\ (\Delta x_{\mathrm B})^2\cdot\kappa^2\rho\ =\ 2\pi\,E_{\mathrm{int}}(E)^2\ln2\ \in\ \mathfrak q^0\ }
\]

%dsp-level:0
\[
\qquad [\mathrm{BM.16}(1)\text{'s lock identity re-derived};\ \text{machine: one Laurent monomial, coefficient}\ 2]:
\]

%dsp-level:0
\[
\textbf{the three freedom censuses:}
\]

%dsp-level:0
\[
\qquad \mathfrak h\text{-free}\ [\text{the}\ \mathfrak h\text{-exponent is}\ 0],
\]

%dsp-level:0
\[
\qquad \Lambda\text{-free},
\]

%dsp-level:0
\[
\qquad \mathrm{P4}\text{-free}\ [\neg\text{one of}\ g,g',h_0,c_W,\Lambda,\varkappa,\beta,\alpha_\Omega\ \text{occurs}]
\]

%dsp-level:0
\[
\qquad [\mathrm{BM.16}(1);\ \text{machine}];
\]

%dsp-level:0
\[
G\ \text{is strictly increasing and quadratic in}\ E_{\mathrm{int}}\ \text{on the quantized ladder}
\]

%dsp-level:0
\[
\qquad [\text{machine: the seven cell values}];
\]

%dsp-level:0
\[
(4)\quad\textbf{the exact profile:}
\]

%dsp-level:0
\[
\qquad \text{the self-energy's dimensionless factor is the product of the three derived arguments' functions, and the energy-side monomial is where the lapse enters:}
\]

%dsp-level:0
\[
\boxed{\ \varepsilon_{\mathrm{EM}}(r,E,N)\ :=\ N(E)\cdot f\!\Bigl(\frac{r}{\Delta x_{\mathrm B}(E)}\Bigr)\cdot \mathsf g\bigl(G(E)\bigr)\cdot\varepsilon_{\mathrm{EM}}\ }
\]

%dsp-level:0
\[
\qquad [\text{machine: the monomial census};\ \ \mathrm{BM.34}(4)\text{'s dispersion}\ (p_0-N^ip_i)^2=N^2(m^2+\gamma^{ij}p_ip_j)\ \text{puts exactly one factor of}\ N\ \text{on the energy}];
\]

%dsp-level:0
\[
\boxed{\ e(r,E,N)\ =\ e_0\,\sqrt{\,N(E)\cdot f\bigl(r/\Delta x_{\mathrm B}(E)\bigr)\cdot \mathsf g\bigl(G(E)\bigr)\,}\ }
\]

%dsp-level:0
\[
\qquad [e^2=4\pi\varepsilon_{\mathrm{EM}};\ \text{machine: squaring recovers}\ (4)\text{'s monomial}];
\]

%dsp-level:0
\[
\textbf{the placement row:}
\]

%dsp-level:0
\[
\qquad N\ \text{on the energy and}\ \sqrt N\ \text{on the coupling are the same monomial}\ —\ \text{the placement is a naming convention,}\ \neg\text{a datum};
\]

%dsp-level:0
\[
\qquad \text{the derived content is}\ (4)\text{'s profile};
\]

%dsp-level:0
\[
\textbf{the curvature composition is priced,}\ \neg\textbf{adopted:}
\]

%dsp-level:0
\[
\qquad \mathsf g\in\Bigl\{1,\ \ 1+G,\ \ \frac1{1+G}\Bigr\},
\]

%dsp-level:0
\[
\qquad \mathsf g\ \text{carried symbolically};
\]

%dsp-level:0
\[
\qquad \neg\text{a clause of this tier selects a composition}
\]

%dsp-level:0
\[
[\text{the argument}\ G\ \text{is derived}\ (3);\ \text{the composition is}\ \neg\text{derived}: \neg\exists\ \text{a count}\to\text{strength theorem}\ [\mathrm{BM.36}(3);\ \mathrm{BM.61}(8)(\mathrm{O1})],\ \neg\exists\ \text{a loop}\ (1)];
\]

%dsp-level:0
\[
(5)\quad\textbf{the exact form at the two loci that matter:}
\]

%dsp-level:0
\[
\qquad \textbf{at the member's own envelope}\ k=1:
\]

%dsp-level:0
\[
\qquad \boxed{\ f(1)=0\implies e(\Delta x_{\mathrm B},E,N)=0\ \ \text{and}\ \ \varepsilon_{\mathrm{EM}}(\Delta x_{\mathrm B},E,N)=0\ }
\]

%dsp-level:0
\[
\textbf{the coupling vanishes at the envelope, so the formed self-energy is the exterior read alone}
\]

%dsp-level:0
\[
\qquad [\mathrm{BM.87}(2),(3):\ \text{the same statement on the energy side}];
\]

%dsp-level:0
\[
\textbf{in the far field}\ k\to\infty:
\]

%dsp-level:0
\[
\qquad f\to1,
\]

%dsp-level:0
\[
\qquad \varepsilon_{\mathrm{EM}}\to N\,\mathsf g(G)\,\varepsilon_{\mathrm{EM}},
\]

%dsp-level:0
\[
\qquad E_{\mathrm{EM}}^{\mathrm{self}}\to\frac12\,N\,\mathsf g(G)\,\varepsilon_{\mathrm{EM}}\,q^2\,m_{\mathrm{hid}}
\]

%dsp-level:0
\[
\qquad [\text{machine: the exact limit row}];
\]

%dsp-level:0
\[
\textbf{no finite neighbourhood holds the whole correction:}
\]

%dsp-level:0
\[
\qquad 1-f(k)=\frac1k>0\ \ \forall k<\infty
\]

%dsp-level:0
\[
\qquad [\mathrm{BM.47}(6);\ \text{machine}];
\]

%dsp-level:0
\[
(6)\quad\textbf{the energy dependence (the running, exactly):}
\]

%dsp-level:0
\[
\qquad \text{at fixed}\ r,
\]

%dsp-level:0
\[
\qquad f\bigl(k(E)\bigr)=1-\frac{\mathfrak h}{r\,E_{\mathrm{int}}(E)\ln2}
\]

%dsp-level:0
\[
\qquad [\text{machine: the substitution}]:
\]

%dsp-level:0
\[
\boxed{\ E_{\mathrm{int}}\ \text{strictly increasing}\ \implies\ \Delta x_{\mathrm B}\ \text{strictly decreasing}\ \implies\ k\ \text{strictly increasing}\ \implies\ f\ \text{strictly increasing}\ }
\]

%dsp-level:0
\[
\qquad [\text{machine: three monotonicity rows on exact rationals}];
\]

%dsp-level:0
\[
\textbf{the coupling grows with the energy scale, geometrically, with no loop in the argument};
\]

%dsp-level:0
\[
\qquad \text{the six-state floors give the exact ladder}\ f=1-\mathfrak h/(r\,E_{\mathrm{int}}\ln2)\ \ \text{at}\ \ E_{\mathrm{int}}\in\{12,24,60\}
\]

%dsp-level:0
\[
\qquad [\mathrm{BM.81}(1);\ \text{machine}];
\]

%dsp-level:0
\[
(7)\quad\textbf{the bars:}
\]

%dsp-level:0
\[
\qquad \neg\exists\ \text{a numerical value of}\ e\ \text{or of}\ e_0\ \text{in any clause}
\]

%dsp-level:0
\[
\qquad [\mathrm{EM.25}(7);\ \mathrm{P5};\ \text{machine: zero external tokens, zero decimal numerals in this tier}];
\]

%dsp-level:0
\[
\neg\text{a measured coupling is identified with any corpus quantity}\ [\mathrm{P5}];
\]

%dsp-level:0
\[
\qquad \varepsilon_{\mathrm{EM}}
\]

%dsp-level:0
\[
= e^2/(4\pi)\ \text{is a dimensionless P4-class symbol carried symbolically throughout}\ [\mathrm{Def\text{-}BM.12}\ \text{local notation};\ \mathrm{EM.25}(9)];
\]

%dsp-level:0
\[
\alpha_B:\quad\neg\text{added},\quad\neg\text{derived},\quad\text{barred on}\ \mathrm{BM.61}(8)\text{'s four grounds}
\]

%dsp-level:0
\[
\qquad [\text{machine: the token scan of this tier}];
\]

%dsp-level:0
\[
\qquad \neg\text{a P4 addition},\quad\text{§Cal-BM stays empty};
\]

%dsp-level:0
\[
(8)\quad\textbf{status:}
\]

%dsp-level:0
\[
\qquad \textbf{derived:}\ (1)\text{--}(3),\ (5),\ (6);
\]

%dsp-level:0
\[
\qquad \textbf{derived in form with a priced composition:}\ (4);
\]

%dsp-level:0
\[
\qquad \textbf{symbolic:}\ e_0,\ g;
\]

%dsp-level:0
\[
\neg\text{a declaration},\quad\neg\text{a new dimensionful scale}\ [\mathrm{P1}],\quad\neg\text{a measured value}\ [\mathrm{P5}],\quad
\]

%dsp-level:0
\[
\qquad \textbf{open:}\ \text{the value of}\ e\ [\mathrm{EM.25}(7)],
\]

%dsp-level:0
\[
\ \text{the curvature composition}\ \mathsf g\ (4),
\]

%dsp-level:0
\[
\ \text{the lapse as solution data}\ [\mathrm{CT2.7};\ \mathrm{BM.34}(2)].
\]

\textbf{Req.} D1.5, D2.9 (Parts 1–2); CT1.12, CT1.22, CT1.23, CT2.1, CT2.7, EM.9, EM.15, EM.25(6), EM.25(7), EM.25(9), EW.15, EW.35 (Part 2); P1, P4, P5 (Parts 1–2); FT.28, FT.145(3), FT.157(1) (Part 3); Census FT.35(2), FT.42(5) (row-wise); Def-BM.4(1), BM.13(1), BM.13(2), BM.13(3), BM.14(2), BM.14(4), BM.15(2), BM.16(1), BM.34(1), BM.34(2), BM.34(4), BM.36(3), BM.47(2), BM.47(6), BM.61(8), BM.81(1), BM.87(2), BM.87(3) (this part); exact Laurent-monomial arithmetic with $\mathfrak q$- and $\mathfrak h$-exponent censuses, exact rational function rows, monotonicity on exact rationals, the ADM polynomial identity by Schur complement, limit rows [math; machine: \texttt{checks/\allowbreak{}verify\_\allowbreak{}H\_\allowbreak{}em\_\allowbreak{}self\_\allowbreak{}energy\_\allowbreak{}correction.py} §B].

\textbf{Pf.}

%dsp-level:0
\[
(1):\ \mathrm{EW.35}\text{'s combination and}\ \mathrm{EM.25}(6),(7)\text{'s rows re-cited};
\]

%dsp-level:0
\[
\qquad \neg\exists\ \text{a loop object in the inventory}\ [\mathrm{BM.61}(8);\ \mathrm{FT.145}(3);\ \text{machine: the token census}]
\]

%dsp-level:0
\[
(2)(\mathrm i):\ \Delta x_{\mathrm B}=1/m_{\mathrm{hid}}=\mathfrak h/(E_{\mathrm{int}}\ln2)\ [\text{machine: one Laurent monomial}];
\]

%dsp-level:0
\[
\qquad E_{\mathrm{int}}\ge2\ [\mathrm{BM.13}(3)];
\]

%dsp-level:0
\[
\qquad \text{the ladder is injective in}\ E_{\mathrm{int}}\ \text{and strictly decreasing}\ [\text{machine}]
\]

%dsp-level:0
\[
(\mathrm{ii}):\ f(k)=1-1/k\ \text{at}\ k=1,\dots,6=(0,\tfrac12,\tfrac23,\tfrac34,\tfrac45,\tfrac56)\ [\text{machine: exact Fractions}];
\]

%dsp-level:0
\[
\qquad f'(k)=1/k^2>0\ \text{on}\ k\ge1\ \text{and}\ f(k)<1\iff1/k>0\ [\text{machine: the rational-function rows}];
\]

%dsp-level:0
\[
\qquad f(1)=0\ [\text{machine}]
\]

%dsp-level:0
\[
(\mathrm{iii}):\ N^2=\Omega^2+\gamma^{ij}N_iN_j\ \text{and}\ N^2\det\gamma=-\det g\ [\mathrm{BM.34}(1)\ \text{re-cited};\ \text{machine: the Schur identity}];
\]

%dsp-level:0
\[
\qquad (a_{00},q_{ij},w_i)=0\implies\Omega=1,\ N_i=0,\ N=1\ [\text{machine: exact substitution}]
\]

%dsp-level:0
\[
(3):\ (\Delta x_{\mathrm B})^2\kappa^2\rho
\]

%dsp-level:0
\[
= \frac{\mathfrak h^2}{E_{\mathrm{int}}^2(\ln2)^2}\cdot\frac{2\pi E_{\mathrm{int}}^4(\ln2)^3}{\mathfrak h^2}
\]

%dsp-level:0
\[
= 2\pi E_{\mathrm{int}}^2\ln2\ [\text{machine: the Laurent monomial cancels}\ \mathfrak h\ \text{exactly}];
\]

%dsp-level:0
\[
\qquad [G]=\mathfrak q^{-1}\cdot\mathfrak q^{+1}=\mathfrak q^0\ [\text{machine: the exponent arithmetic}];
\]

%dsp-level:0
\[
\qquad \text{the three freedom censuses}\ [\mathrm{BM.16}(1);\ \text{machine: the symbol scan}]
\]

%dsp-level:0
\[
(4):\ \text{the product of three dimensionless factors is dimensionless}\ [\text{machine}];
\]

%dsp-level:0
\[
\qquad e^2=4\pi\varepsilon_{\mathrm{EM}}\implies e=e_0\sqrt{Nf\mathsf g}\ \text{at}\ \varepsilon_{\mathrm{EM}}=Nf\mathsf g\,\varepsilon_{\mathrm{EM}}\ [\text{machine: squaring recovers the monomial}];
\]

%dsp-level:0
\[
\qquad \text{the one-factor-of-}N\ \text{row}:\ (p_0-N^ip_i)^2
\]

%dsp-level:0
\[
= N^2(m^2+\gamma^{ij}p_ip_j)
\]

%dsp-level:0
\[
\implies E
\]

%dsp-level:0
\[
= N\sqrt{m^2+p^2}-N^ip_i\ [\mathrm{BM.34}(4);\ \text{machine}]
\]

%dsp-level:0
\[
\text{the priced set}\ \{1,1+G,(1+G)^{-1}\}\ \text{is a pricing,}\ \neg\text{an adoption}\ [\mathrm{Def\text{-}BM.12}(4)\ \text{pattern}];
\]

%dsp-level:0
\[
\qquad \neg\exists\ \text{a count}\to\text{strength theorem}\ [\mathrm{BM.36}(3);\ \mathrm{BM.61}(8)(\mathrm{O1})]
\]

%dsp-level:0
\[
(5):\ f(1)=1-1=0\ [\text{machine}];
\]

%dsp-level:0
\[
\qquad \lim_{k\to\infty}(1-1/k)=1\ [\text{machine: the exact limit row on rationals}];
\]

%dsp-level:0
\[
\qquad 1-f(k)=1/k>0\ \forall k<\infty\ [\text{machine}]
\]

%dsp-level:0
\[
(6):\ k(E)
\]

%dsp-level:0
\[
= rE_{\mathrm{int}}(E)\ln2/\mathfrak h\ \text{is strictly increasing in}\ E_{\mathrm{int}}\ \text{at fixed}\ r>0\ [\text{machine: positive coefficients}];
\]

%dsp-level:0
\[
\qquad f\ \text{strictly increasing}\ [(2)(\mathrm{ii})]\implies\text{the composition is strictly increasing}\ [\text{machine}];
\]

%dsp-level:0
\[
\qquad \Delta x_{\mathrm B}\ \text{strictly decreasing}\ [\text{machine}]
\]

%dsp-level:0
\[
(7)\text{--}(8):\ \text{the token scans, the symbol census and the assembly}\ [\text{machine}].\ \blacksquare
\]

\textbf{Phys.} A coupling with no scale at which to sit is not a number, and the corpus never pretended otherwise: it fixed the combination of the two gauge strengths and left the value open. What it can do is say how the coupling's \textit{profile} runs, and every ingredient of that profile is already derived. Distance enters through the only length the ledger owns — the Born width, which is the reciprocal of an edge count — and through the fraction of an inverse-square flux enclosed within a given number of those widths, which is one minus one over k, exactly. Time enters through the lapse of the emergent metric, which is one in the vacuum gauge and departs from one only where the shift and the conformal coefficient do. Curvature enters through the one dimensionless monomial the calibration lock supplies: the Born width squared times the Ricci scale, equal to two pi times the square of the edge count times the logarithm of two, free of h, free of the cosmological term, free of every declared parameter. Multiply the three and the coupling is an exact algebraic function of position, energy and lapse — vanishing at the member's own envelope, growing without bound toward its far-field value, never quite reaching it at any finite radius. The composition of the curvature factor is not derived, because nothing in the inventory turns a count into a strength, so it is priced three ways and adopted not at all. The running comes out in the direction a loop calculation would give it, from geometry alone, with no loop in the argument and no number fitted anywhere.

\end{theorem}

\begin{theorem}[{\texorpdfstring{BM.89 — The EM self-energy from the path integral: the derived measure, the crystallization weight, and the route-relative correction}{BM.89 - The EM self-energy from the path integral: the derived measure, the crystallization weight, and the route-relative correction}}]\label{tag:BM-89}
\textbf{Thm.}

%dsp-level:0
\[
(1)\quad\textbf{the ensemble and its derived measure:}\qquad Q=\begin{pmatrix}-x_+&x_+\\ x_-&-x_-\end{pmatrix}\ \ \text{over}\ \mathbb Z[x_+,x_-],\qquad(x_-,\,x_+)\,Q=(0,0)\ \ \text{exactly},\qquad\dim\ker Q=1\qquad[\mathrm{BM.19}(1);\ \mathrm{FT.108}(4);\ \text{machine}]:
\]

%dsp-level:0
\[
\boxed{\ p\ :=\ \frac{x_+}{x_++x_-}\ \ \text{is the unique stationary measure}\ }
\]

%dsp-level:0
\[
\qquad x_+=1-4\varepsilon_q,\quad x_-=4\varepsilon_q\ \ \implies\ \ p=1-4\varepsilon_q
\]

%dsp-level:0
\[
\qquad [\mathrm{BM.19}(2);\ \text{machine: an identity, exact at every rational threshold}];
\]

%dsp-level:0
\[
\varepsilon_q\in[0,\tfrac14)\ [\mathrm{CT0.8};\ \mathrm{D0.4};\ \text{A4}]\implies x_-\in[0,1),\quad x_+\in(0,1],\quad p\in(0,1]\qquad[\text{machine}];
\]

%dsp-level:0
\[
\varepsilon_q=0\implies x_-=0\implies\mathcal E(\mathcal X)=\{\text{the crystallized configuration}\}
\]

%dsp-level:0
\[
\qquad [\mathrm{BM.19}(3);\ \mathrm{Def\text{-}BM.5}(5)];
\]

%dsp-level:0
\[
\qquad |\mathcal E(\mathcal X)|>1\iff\varepsilon_q>0
\]

%dsp-level:0
\[
\qquad [\mathrm{BM.19}(6)];
\]

%dsp-level:0
\[
\textbf{the ensemble is finite at every budget and}\ \neg\textbf{a fixed topology at any nonzero threshold}
\]

%dsp-level:0
\[
\qquad [\mathrm{Def\text{-}H.32}(1);\ \mathrm{BM.19}(4),(6)];
\]

%dsp-level:0
\[
(2)\quad\textbf{the configuration energy:}
\]

%dsp-level:0
\[
\qquad \boxed{\ E_{\mathrm{EM}}(R)\ =\ \frac12\,N(R)\,\mathsf g\bigl(G(R)\bigr)\,\varepsilon_{\mathrm{EM}}\,Q(R)^2\,m_{\mathrm{hid}}(R),\qquad m_{\mathrm{hid}}(R)=\frac{E_{\mathrm{int}}(R)\ln2}{\mathfrak h}=\frac{|R|\ln2}{\mathfrak h}\ }
\]

%dsp-level:0
\[
[\mathrm{BM.87}(4)\text{'s formed monomial at}\ c=\tfrac12;\ \mathrm{BM.88}(4)\text{'s profile};\ \mathrm{Def\text{-}BM.5}(4):\ E_{\mathrm{int}}(R)=|R|;\ \mathrm{FT.44}(4):\ \text{the taxonomy law};\ \text{machine}];
\]

%dsp-level:0
\[
(3)\quad\textbf{the class factorization:}
\]

%dsp-level:0
\[
\qquad Q(R)=Q\ \ \text{on every member of a charge class}\ —\ \text{the boundary data fix the class}
\]

%dsp-level:0
\[
\qquad [\mathrm{Def\text{-}BM.5}(1);\ \mathrm{Def\text{-}BM.8}(1);\ \text{machine: the class census}]:
\]

%dsp-level:0
\[
\boxed{\ E_{\mathrm{EM}}^{\mathrm{self}}(Q)\ =\ \sum_{R\in\mathcal E(\mathcal X)}P(R)\,E_{\mathrm{EM}}(R)\ =\ \frac12\,N\,\mathsf g(G)\,\varepsilon_{\mathrm{EM}}\,Q^2\,\bigl\langle m_{\mathrm{hid}}\bigr\rangle_{\mathcal E}\ }
\]

%dsp-level:0
\[
\qquad [\text{machine:}\ Q^2\ \text{is constant on the class and factors out of the average}];
\]

%dsp-level:0
\[
(4)\quad\textbf{the crystallization weight (the derived attenuation):}
\]

%dsp-level:0
\[
\qquad \boxed{\ w(\varepsilon_q)\ :=\ \frac{\langle E_{\mathrm{int}}\rangle_{\mathcal E}}{E_{\mathrm{int}}^{\mathrm{cryst}}}\ =\ \frac1{E_{\mathrm{int}}^{\mathrm{cryst}}}\sum_{\text{edges}}P(\text{the edge is real})\ }
\]

%dsp-level:0
\[
\qquad [\mathrm{Def\text{-}BM.5}(4);\ \text{the ensemble average of}\ |R|]:
\]

%dsp-level:0
\[
\boxed{\ w(0)=1\ \ \text{exactly}\ }
\]

%dsp-level:0
\[
\qquad [\mathrm{BM.19}(3):\ \text{the ensemble is the single crystallized configuration}];
\]

%dsp-level:0
\[
\qquad 0<w(\varepsilon_q)\le1
\]

%dsp-level:0
\[
\qquad [\mathrm{BM.19}(5),(6):\ \text{the ensemble spreads}\ E_{\mathrm{int}}\ \text{over}\ \{0,\dots,B_K\}\ \text{with the crystallized value as its top point, and every move changes}\ E_{\mathrm{int}}\ \text{by its}\ J^\pm\ \text{excess}];
\]

%dsp-level:0
\[
\textbf{the strict decrease is derived on the marked specialization below, and in general is}\ (\mathrm{R4})\textbf{'s content:}
\]

%dsp-level:0
\[
\qquad \varepsilon_q\ \text{enters the measure only through}\ x_-=4\varepsilon_q\ [\mathrm{BM.19}(2);\ \text{machine: the symbol census}];
\]

%dsp-level:0
\[
\textbf{the Bernoulli specialization (marked,}\ \neg\textbf{adopted):}
\]

%dsp-level:0
\[
\qquad \text{on the per-datum stationary retention of}\ (1),
\]

%dsp-level:0
\[
\qquad \langle E_{\mathrm{int}}\rangle_{\mathcal E}=E_{\mathrm{int}}^{\mathrm{cryst}}\cdot p\implies\boxed{\ w=1-4\varepsilon_q\ }
\]

%dsp-level:0
\[
[\text{machine:}\ E_{\mathrm{int}}^{\mathrm{cryst}}=6=(c_1{:}\,2)+(c_2{:}\,2)+(\text{connector}{:}\,2)\ [\mathrm{BM.28}(3);\ \mathrm{BM.17};\ \mathrm{BM.26}]\ \text{is three pairs, each retained with weight}\ p];
\]

%dsp-level:0
\[
\qquad w\ \text{is carried symbolically wherever the specialization is not invoked};
\]

%dsp-level:0
\[
\textbf{the member is a stationary mixture of real and virtual content:}
\]

%dsp-level:0
\[
\qquad \text{the composite is where the real and the virtual registers meet,}\ \neg\text{a fixed topology}
\]

%dsp-level:0
\[
\qquad [\mathrm{BM.19}(6);\ \mathrm{D1.5};\ \mathrm{CT1.22}];
\]

%dsp-level:0
\[
(5)\quad\textbf{the corrected self-energy:}
\]

%dsp-level:0
\[
\qquad \boxed{\ E_{\mathrm{EM}}^{\mathrm{self}}(Q)\ =\ \frac12\,N\,w\,\mathsf g(G)\,\varepsilon_{\mathrm{EM}}\,Q^2\,\frac{E_{\mathrm{int}}^{\mathrm{cryst}}\ln2}{\mathfrak h}\ }
\]

%dsp-level:0
\[
\qquad [\text{machine: one Laurent monomial in}\ (N,w,\mathsf g,\varepsilon_{\mathrm{EM}},Q,E_{\mathrm{int}},\ln2,\mathfrak h)]:
\]

%dsp-level:0
\[
\text{at}\ E_{\mathrm{int}}^{\mathrm{cryst}}=6:
\]

%dsp-level:0
\[
\qquad E_{\mathrm{EM}}^{\mathrm{self}}(\pm1)=\frac{3\,N\,w\,\mathsf g\,\varepsilon_{\mathrm{EM}}\ln2}{\mathfrak h},
\]

%dsp-level:0
\[
\qquad E_{\mathrm{EM}}^{\mathrm{self}}(0)=0\ \ \textbf{exactly}
\]

%dsp-level:0
\[
\qquad [\text{machine:}\ 0^2=0;\ \text{Gauss: zero net flux}];
\]

%dsp-level:0
\[
\textbf{the crystallized vacuum-gauge specialization:}
\]

%dsp-level:0
\[
\qquad (w,N,\mathsf g)=(1,1,1)\implies E_{\mathrm{EM}}^{\mathrm{self}}(\pm1)=\frac{3\,\varepsilon_{\mathrm{EM}}\ln2}{\mathfrak h}=\frac12\,\varepsilon_{\mathrm{EM}}\,m_{\mathrm{hid}}=\frac{e^2q^2}{8\pi\,\Delta x_{\mathrm B}}
\]

%dsp-level:0
\[
\boxed{\ \textbf{the path-integral value reproduces}\ \mathrm{BM.87}(2)\textbf{'s derived exterior value exactly at full crystallization in the vacuum gauge}\ }
\]

%dsp-level:0
\[
\qquad [\text{machine: the two monomials agree coefficient by coefficient}];
\]

%dsp-level:0
\[
(6)\quad\textbf{the route-relativity (priced,}\ \neg\textbf{adopted):}
\]

%dsp-level:0
\[
\qquad \textbf{route}\ (\mathrm i)\ \text{(the normalized average):}\quad\overline{\Delta S}(\cdot,S)=m^2_B(\cdot,S)\ \text{is charge-degenerate}
\]

%dsp-level:0
\[
\qquad [\mathrm{FT.107}(3);\ \mathrm{BM.28}(4)(\mathrm i)]:
\]

%dsp-level:0
\[
\text{the correction splits the classes only through}\ Q^2,\ \text{and the ledger's own average is class-blind}
\]

%dsp-level:0
\[
\qquad [\text{machine: the frozen sector expectations are}\ q\text{-independent}];
\]

%dsp-level:0
\[
\textbf{route}\ (\mathrm{ii})\ \text{(the unnormalized total deficit):}
\]

%dsp-level:0
\[
\qquad T(Q,S)=W(Q)\,\overline{\Delta S}(S),
\]

%dsp-level:0
\[
\qquad W=(1,6,15,\mathbf{20},15,6,1)
\]

%dsp-level:0
\[
\qquad [\mathrm{BM.28}(4)(\mathrm{ii});\ \mathrm{Def\text{-}BM.13}(3)]:
\]

%dsp-level:0
\[
E_{\mathrm{EM}}^{\mathrm{self},(\mathrm{ii})}(Q,S)
\]

%dsp-level:0
\[
= \frac12\,
\]

%dsp-level:0
\[
N\,
\]

%dsp-level:0
\[
\mathsf g\,
\]

%dsp-level:0
\[
\varepsilon_{\mathrm{EM}}\,
\]

%dsp-level:0
\[
Q^2\,
\]

%dsp-level:0
\[
\frac{T(Q,S)}{\mathfrak h}\ \ \implies\ \ \Delta m_{\mathrm{EM}}^{(\mathrm{ii})}(S)
\]

%dsp-level:0
\[
= \frac{N\,\mathsf g\,\varepsilon_{\mathrm{EM}}}{2\mathfrak h}\cdot15\,
\]

%dsp-level:0
\[
m^2_B(\cdot,S)
\]

%dsp-level:0
\[
= \frac{N\,\mathsf g\,\varepsilon_{\mathrm{EM}}}{2\mathfrak h}\Bigl(\frac{12659355}{98956},\ \frac{28824930}{393679},\ \frac{1365}{23},\ \frac{205695}{3088}\Bigr)
\]

%dsp-level:0
\[
[\text{machine: exact rational arithmetic at the four frozen sector expectations}\ \bigl(\tfrac{843957}{98956},\tfrac{1921662}{393679},\tfrac{91}{23},\tfrac{13713}{3088}\bigr)\ [\mathrm{BM.21}(2);\ \mathrm{BM.62}(2)]];
\]

%dsp-level:0
\[
\textbf{route}\ (\mathrm{ii})\ \text{is}\ \neg\textbf{adopted for any purpose beyond the vertex kinematics}
\]

%dsp-level:0
\[
\qquad [\mathrm{BM.30}(5);\ \mathrm{BM.37}(6);\ \mathrm{BM.77}(1),(9)\ \text{stand}];
\]

%dsp-level:0
\[
\qquad \text{route}\ (\mathrm{iii})\text{'s amplitude ratios stay barred as masses}
\]

%dsp-level:0
\[
\qquad [\mathrm{FT.79}(5);\ \mathrm{BM.37}(3)];
\]

%dsp-level:0
\[
(7)\quad\textbf{the attenuation theorem:}
\]

%dsp-level:0
\[
\qquad 0<w\le1\ \ \text{and}\ \ N>0,\ \mathsf g>0\implies\boxed{\ \bigl|E_{\mathrm{EM}}^{\mathrm{self}}(Q)\bigr|\ \le\ \bigl|E_{\mathrm{EM}}^{\mathrm{self}}(Q)\bigr|_{w=1}\ }
\]

%dsp-level:0
\[
\qquad \text{with equality only at}\ w=1:
\]

%dsp-level:0
\[
\textbf{the flicker strictly attenuates the EM correction at every nonzero threshold};
\]

%dsp-level:0
\[
\qquad \text{the transport register's lead}\ \ln(4/3)/\mathfrak h\ \text{is class-shared and}\ \neg\text{attenuated by}\ w\ \ [\mathrm{BM.57}(4);\ \mathrm{BM.59}(4);\ \mathrm{BM.91}(2)];
\]

%dsp-level:0
\[
\implies\textbf{the sub-critical branch widens as the crystallization weight falls}
\]

%dsp-level:0
\[
\qquad [\mathrm{BM.91}(5);\ \text{machine: monotone in}\ w];
\]

%dsp-level:0
\[
(8)\quad\textbf{the residue:}
\]

%dsp-level:0
\[
\qquad w(\varepsilon_q)\ \text{in general requires the ensemble's configuration distribution beyond the per-datum specialization}\ —\ \textbf{named}\ (\mathrm{R4}),
\]

%dsp-level:0
\[
\ \neg\textbf{filled}:
\]

%dsp-level:0
\[
\text{the computed cells supply exact sector averages of the deficit register}\ \overline{\Delta S}
\]

%dsp-level:0
\[
= (\tfrac32 n_{\mathcal V})\ \ [\mathrm{FT.106}(3);\ \mathrm{FT.107}(3)],
\]

%dsp-level:0
\[
\ \ \neg\text{of the edge-count register}\ E_{\mathrm{int}}\qquad = |R|:
\]

%dsp-level:0
\[
\textbf{the two registers are distinct alphabets and are}\ \neg\textbf{equated}
\]

%dsp-level:0
\[
\qquad [\mathrm{BM.62}(11);\ \mathrm{BM.37}(2);\ \mathrm{BM.83}(3);\ \text{machine: the token census}];
\]

%dsp-level:0
\[
(9)\quad\textbf{status:}
\]

%dsp-level:0
\[
\qquad \textbf{derived:}\ (1)\text{--}(5),\ (7);
\]

%dsp-level:0
\[
\qquad \textbf{priced,}\ \neg\textbf{adopted:}\ (4)\text{'s Bernoulli specialization and}\ (6)\text{'s route rows};
\]

%dsp-level:0
\[
\qquad \textbf{named,}\ \neg\textbf{filled:}\ (8)\text{'s}\ (\mathrm{R4});
\]

%dsp-level:0
\[
\neg\text{a declaration},
\]

%dsp-level:0
\[
\quad\neg\text{a P4 addition},
\]

%dsp-level:0
\[
\quad\neg\text{a new dimensionful scale}\ [\mathrm{P1}],
\]

%dsp-level:0
\[
\quad\neg\text{a measured value}\ [\mathrm{P5}],
\]

%dsp-level:0
\[
\quad\text{§Cal-BM stays empty};\qquad \qquad
\]

\textbf{Req.} D1.5, D3.4 (Parts 1–2); A4, CT0.8, CT1.22, D0.4, EM.9 (Part 2); P1, P4, P5 (Parts 1–2); Def-H.32(1), FT.44(4), FT.79(5), FT.106(3), FT.107(3), FT.108(4) (Part 3); Def-BM.4(1), Def-BM.5(1), Def-BM.5(4), Def-BM.5(5), Def-BM.8(1), Def-BM.13(3), BM.17, BM.19(1), BM.19(2), BM.19(3), BM.19(4), BM.19(5), BM.19(6), BM.21(2), BM.26, BM.28(3), BM.28(4), BM.30(5), BM.37(2), BM.37(3), BM.37(6), BM.57(4), BM.59(4), BM.62(2), BM.62(11), BM.77(1), BM.77(9), BM.83(3), BM.87(2), BM.87(4), BM.88(4) (this part); exact polynomial arithmetic over $\mathbb Q[x_+,x_-]$, exhaustive integer kernel censuses, exact rational specialization, Laurent-monomial assembly with coefficient comparison, exact rational arithmetic at the four frozen sector expectations, token censuses [math; machine: \texttt{checks/\allowbreak{}verify\_\allowbreak{}H\_\allowbreak{}em\_\allowbreak{}self\_\allowbreak{}energy\_\allowbreak{}correction.py} §C].

\textbf{Pf.}

%dsp-level:0
\[
(1):\ (x_-)(-x_+)+(x_+)(x_-)=0\ \text{and}\ (x_-)(x_+)+(x_+)(-x_-)=0\ \text{in}\ \mathbb Q[x_+,x_-]\ [\text{machine}];
\]

%dsp-level:0
\[
\qquad \text{constant vectors}:\ a(-x_+)+bx_-=0\ \wedge\ ax_++b(-x_-)=0\implies a=b=0\ [\text{machine: exhaustive on}\ [-3,3]^2];
\]

%dsp-level:0
\[
\qquad p=(1-4\varepsilon_q)/(1-4\varepsilon_q+4\varepsilon_q)=1-4\varepsilon_q\ [\text{machine: exact at six rational thresholds}]
\]

%dsp-level:0
\[
\varepsilon_q\in[0,\tfrac14)\implies4\varepsilon_q\in[0,1)\ [\text{machine}];
\]

%dsp-level:0
\[
\qquad x_-=0\ \text{at}\ \varepsilon_q=0\implies\neg\exists\ \text{a dissolution move}\ [\mathrm{D1.2}:\ J^-:n_{\mathcal R}:1\to0];
\]

%dsp-level:0
\[
\qquad |\mathcal E|>1\iff x_->0\ [\mathrm{BM.19}(6)]
\]

%dsp-level:0
\[
(2)\text{--}(3):\ \mathrm{BM.87}(4)\text{'s monomial at}\ c=\tfrac12\ \text{and}\ \mathrm{BM.88}(4)\text{'s profile composed}\ [\text{machine}];
\]

%dsp-level:0
\[
\qquad E_{\mathrm{int}}(R)=|R|\ [\mathrm{Def\text{-}BM.5}(4)];
\]

%dsp-level:0
\[
\qquad Q(R)
\]

%dsp-level:0
\[
= Q\ \text{class-constant}\ [\mathrm{Def\text{-}BM.5}(1);\ \text{machine: the boundary-data census}]
\]

%dsp-level:0
\[
\implies Q^2\ \text{factors out of}\ \sum_RP(R)\cdot
\]

%dsp-level:0
\[
(4):\ \langle E_{\mathrm{int}}\rangle_{\mathcal E}=\sum_{\text{edges}}P(\text{real})\ [\text{linearity of the average on}\ |R|];
\]

%dsp-level:0
\[
\qquad \varepsilon_q=0:\ \mathcal E=\{R_{\mathrm{cryst}}\}\implies\langle E_{\mathrm{int}}\rangle=E_{\mathrm{int}}^{\mathrm{cryst}}\implies w=1\ [\text{machine}];
\]

%dsp-level:0
\[
\qquad 0<|R|\le E_{\mathrm{int}}^{\mathrm{cryst}}\ \text{on the ensemble's support}\implies0<w\le1\ [\mathrm{BM.19}(5);\ \text{machine}]
\]

%dsp-level:0
\[
\text{the specialization}:\ \text{three pairs, each retained with the stationary weight}\ p\implies\langle E_{\mathrm{int}}\rangle
\]

%dsp-level:0
\[
= 3\cdot2\cdot p
\]

%dsp-level:0
\[
= 6p\ \text{and}\ w
\]

%dsp-level:0
\[
= p
\]

%dsp-level:0
\[
= 1-4\varepsilon_q\ [\text{machine: exact rationals}];
\]

%dsp-level:0
\[
\qquad E_{\mathrm{int}}^{\mathrm{cryst}}=6=(2)+(2)+(2)\ [\mathrm{BM.28}(3);\ \text{machine}]
\]

%dsp-level:0
\[
(5):\ \tfrac12\cdot6=\ 3\ [\text{machine}];
\]

%dsp-level:0
\[
\qquad Q=0\implies Q^2=0\implies E=0\ [\text{machine}];
\]

%dsp-level:0
\[
\qquad (w,N,\mathsf g)
\]

%dsp-level:0
\[
= (1,1,1):\ \tfrac{3\varepsilon_{\mathrm{EM}}\ln2}{\mathfrak h}
\]

%dsp-level:0
\[
= \tfrac12\varepsilon_{\mathrm{EM}}\tfrac{6\ln2}{\mathfrak h}
\]

%dsp-level:0
\[
= \tfrac12\varepsilon_{\mathrm{EM}}m_{\mathrm{hid}}
\]

%dsp-level:0
\[
= \tfrac{e^2q^2}{8\pi\Delta x_{\mathrm B}}\ [\text{machine: three Laurent monomials, coefficient by coefficient},\ \Delta x_{\mathrm B}m_{\mathrm{hid}}=1]
\]

%dsp-level:0
\[
(6):\ m^2_B(q,S)=m^2_B(-q,S)\ [\mathrm{FT.107}(3);\ \text{machine}];
\]

%dsp-level:0
\[
\qquad 15\cdot\tfrac{843957}{98956}=\tfrac{12659355}{98956},
\]

%dsp-level:0
\[
\quad 15\cdot\tfrac{1921662}{393679}=\tfrac{28824930}{393679},
\]

%dsp-level:0
\[
\quad 15\cdot\tfrac{91}{23}=\tfrac{1365}{23},
\]

%dsp-level:0
\[
\quad 15\cdot\tfrac{13713}{3088}
\]

%dsp-level:0
\[
= \tfrac{205695}{3088}\ [\text{machine: exact Fractions, each in lowest terms or reduced exactly}]
\]

%dsp-level:0
\[
(7):\ 0<w
\]

%dsp-level:0
\[
\le 1\ \text{and}\ N,
\]

%dsp-level:0
\[
\mathsf g>0\implies\tfrac12Nwg\varepsilon_{\mathrm{EM}}Q^2m_{\mathrm{hid}}\le\tfrac12Ng\varepsilon_{\mathrm{EM}}Q^2m_{\mathrm{hid}}\ [\text{machine: the coefficient census}];
\]

%dsp-level:0
\[
\qquad \text{equality}\iff w=1\ [\text{machine}]
\]

%dsp-level:0
\[
(8)\text{--}(9):\ \text{the register-alphabet census}\ [\mathrm{BM.62}(11)\ \text{pattern};\ \text{machine}];
\]

%dsp-level:0
\[
\qquad \text{the standing-rule scans}\ [\text{machine}].\ \blacksquare
\]

\textbf{Phys.} The composite was never a fixed topology, and the correction that rides it should not have been either. Take the flicker's own measure — the unique stationary line of the Janus generator, which specializes to one minus four times the crystallization threshold with no rate declared anywhere — average the formed monomial over the ensemble, and the charge factors straight out, because the boundary data fix it before the flicker starts. What is left is one derived number between zero and one: the crystallization weight, the mean fraction of the composite's registration data that are real rather than virtual. It is exactly one at the sharp threshold, where the ensemble collapses to the single crystallized pattern, and it falls as the threshold rises, so the electromagnetic correction is attenuated by exactly the amount the member fails to be crystallized. At full crystallization in the vacuum gauge the average reproduces the exterior value the adjudication derived, monomial for monomial, so nothing is lost and one classical coefficient is gone. The baryon, then, is where the real and the virtual registers meet: a stationary mixture whose mass is an average, whose charge is exact, and whose field energy is a fraction of what a fully frozen lump would carry. The general weight beyond the per-datum specialization stays open and named, because the computed cells average the vertex-deficit register and not the edge count, and those two alphabets are not the same currency.

\end{theorem}

\begin{theorem}[{\texorpdfstring{BM.90 — The defect-conditioned weighting: the hypothesis on its face, and the branch-blindness theorem}{BM.90 - The defect-conditioned weighting: the hypothesis on its face, and the branch-blindness theorem}}]\label{tag:BM-90}
\textbf{Thm.}

%dsp-level:0
\[
(1)\quad\textbf{the hypothesis, on its face:}
\]

%dsp-level:0
\[
\qquad \textbf{hypothesis (W):}\ \ \text{the retention probability of a registration datum depends on the state's defect value},
\]

%dsp-level:0
\[
\quad P(\text{the edge is real}\mid D_S),
\]

%dsp-level:0
\[
\text{so that}
\]

%dsp-level:0
\[
\qquad m_{\mathrm{hid}}^{\mathrm{eff}}(s)\ =\ \frac{\ln2}{\mathfrak h}\sum_{\text{edges}}P(\text{the edge is real}\mid D_S(s))
\]

%dsp-level:0
\[
\qquad [\text{the six states of}\ \mathrm{Def\text{-}BM.14}(2),(3);\ \mathrm{BM.81}(1)];
\]

%dsp-level:0
\[
(2)\quad\textbf{the measure's alphabet census:}\qquad x_+=1-4\varepsilon_q,\qquad x_-=4\varepsilon_q,\qquad p=1-4\varepsilon_q\qquad[\mathrm{BM.19}(2);\ \text{machine}]:
\]

%dsp-level:0
\[
\boxed{\ \text{the flicker measure's symbol set is}\ \{\varepsilon_q\}\ \ \text{and}\ \ \varepsilon_q\ \text{is the A4 threshold}\ }
\]

%dsp-level:0
\[
\qquad [\mathrm{D0.4};\ \mathrm{BM.19}(2):\ \neg\text{a P-class datum enters};\ \neg\text{a new fugacity},\ \neg\text{a new rate}];
\]

%dsp-level:0
\[
\text{the defect value's symbol set is}\ \{\delta_v,\ \hat\mu'^2\}
\]

%dsp-level:0
\[
\qquad [\mathrm{Def\text{-}BM.14}(3);\ \mathrm{Def\text{-}H.37}(2);\ \text{machine}];
\]

%dsp-level:0
\[
\qquad \boxed{\ \{\varepsilon_q\}\ \cap\ \{\delta_v,\hat\mu'^2\}\ =\ \varnothing\ }
\]

%dsp-level:0
\[
(3)\quad\textbf{the branch-blindness theorem:}
\]

%dsp-level:0
\[
\qquad \boxed{\ \neg\exists\ \text{a derived coupling}\ D_S\longrightarrow P(\text{retention});\qquad w=w(\varepsilon_q)\ \ \text{is branch-blind}\ }
\]

%dsp-level:0
\[
\text{on four independent grounds, each re-cited at its own locus:}
\]

%dsp-level:0
\[
(\mathrm i)\quad\textbf{the field degree:}\ \ \text{the six defect values generate a degree-}16\ \text{extension of}\ \mathbb Q,
\]

%dsp-level:0
\[
\ \text{while every read the flicker measure is built from is integer- or binary-valued}
\]

%dsp-level:0
\[
\qquad [\mathrm{FT.73}(2),(3);\ \mathrm{BM.69}(4)(\mathrm i);\ \mathrm{BM.81}(1)]:
\]

%dsp-level:0
\[
\text{a rational formula places every retention weight in a field of degree}<16\ \text{over the declared constants}
\]

%dsp-level:0
\[
\qquad [\text{machine re-cited}];
\]

%dsp-level:0
\[
(\mathrm{ii})\quad\textbf{the two-torsion image:}\ \ \text{the register lift of the defect angle reaches}\ \{0,\pi\}\ \text{only}\ —\ \text{four of the six states are unreachable at every winding}
\]

%dsp-level:0
\[
\qquad [\mathrm{BM.67}(6);\ \mathrm{BM.69}(4)(\mathrm{ii})];
\]

%dsp-level:0
\[
(\mathrm{iii})\quad\textbf{the calibration weight:}\ \ \text{the winding route}\ \Theta_R
\]

%dsp-level:0
\[
= 2\pi k\ \text{carries}\ \mathfrak h\text{-weight}\ {-}2\ \text{against}\ 0\ \text{for a rate}
\]

%dsp-level:0
\[
\qquad [\mathrm{FT.55}(1);\ \mathrm{BM.67}(8);\ \mathrm{BM.69}(4)(\mathrm{iii})];
\]

%dsp-level:0
\[
(\mathrm{iv})\quad\textbf{the strength bar:}\ \ \neg\exists\ \text{a count}\to\text{strength theorem};
\]

%dsp-level:0
\[
\qquad \text{the capacity ratios}\ \bigl(\tfrac13,\,8,\,\tfrac1{25}\bigr)\ \text{are}\ \neg\text{couplings}
\]

%dsp-level:0
\[
\qquad [\mathrm{BM.36}(3);\ \mathrm{BM.61}(8)(\mathrm{O1});\ \mathrm{BM.69}(4)(\mathrm{iv})];
\]

%dsp-level:0
\[
(4)\quad\textbf{the consequence for the informational register:}
\]

%dsp-level:0
\[
\qquad \boxed{\ \textbf{the six-state quark splitting does not reweight}\ m_{\mathrm{hid}}\ }
\]

%dsp-level:0
\[
\qquad \text{the}\ m_{\mathrm{hid}}\ \text{contributions are the values}\ \mathrm{BM.81}(1)\ \text{censused},
\]

%dsp-level:0
\[
\ \text{attenuated by a branch-independent}\ w:
\]

%dsp-level:0
\[
m_{\mathrm{hid}}^{\mathrm{eff}}(s)=w(\varepsilon_q)\cdot\frac{2E(s)\ln2}{\mathfrak h}
\]

%dsp-level:0
\[
\qquad \text{with the same}\ w\ \text{on all six states}
\]

%dsp-level:0
\[
\qquad [\text{machine:}\ w\text{'s symbol census carries no}\ D_S];
\]

%dsp-level:0
\[
\textbf{the defect values enter the pole currency alone}
\]

%dsp-level:0
\[
\qquad [\mathrm{BM.82}(4):\ \text{the thirty}\ \rho_{\mathrm B}\text{-free pole ratios};\ \mathrm{BM.81}(1)\text{'s}\ D_S\ \text{column}];
\]

%dsp-level:0
\[
(5)\quad\textbf{what the splitting does change (the positive companion, in its own currency):}
\]

%dsp-level:0
\[
\qquad \Delta_g=\bigl(2,\ \sqrt2-1,\ 2\mathsf c-\sqrt3\bigr)\hat\mu'^2>0\ \ \text{per generation}
\]

%dsp-level:0
\[
\qquad [\mathrm{BM.75}(4);\ \mathrm{BM.81}(3)]:
\]

%dsp-level:0
\[
\mathcal R_s=\frac{D_S(s)}{4E(s)^2}\ \ \text{per state}
\]

%dsp-level:0
\[
\qquad [\mathrm{BM.82}(2)];
\]

%dsp-level:0
\[
\qquad \Delta\bigl(m^2_{\mathrm{pole}}\bigr)=\Delta_g>0\ \ \text{under hypothesis}\ (\mathrm H)\ \text{(the single substitution)}
\]

%dsp-level:0
\[
\qquad [\mathrm{BM.59}(9);\ \mathrm{BM.69}(3);\ \text{the hypothesis on its face}];
\]

%dsp-level:0
\[
\textbf{and}\ \neg\textbf{addable to the informational assembly:}
\]

%dsp-level:0
\[
\qquad [\Delta_g]=\mathfrak q\ \ \text{against}\ \ [\Delta m]=\mathfrak q^{+1/2},
\]

%dsp-level:0
\[
\qquad \text{a cross-currency addition being a category the corpus bars}
\]

%dsp-level:0
\[
[\mathrm{BM.83}(1),(3);\ \mathrm{BM.35}(5);\ \mathrm{BM.37}(2);\ \mathrm{BM.84}(1);\ \text{machine: the exponent arithmetic}];
\]

%dsp-level:0
\[
(6)\quad\textbf{the priced alternative:}
\]

%dsp-level:0
\[
\qquad \text{a declared}\ D_S\text{-dependent retention rate would be a P4-class count}\to\text{strength coupling}\ —\ \text{the exact species}\ \mathrm{BM.61}(8)\ \text{bars and}\ \mathrm{BM.36}(3)\ \text{refuses}:
\]

%dsp-level:0
\[
\boxed{\ \neg\text{such a datum is issued};\qquad\text{§Cal-BM stays empty};\qquad\neg\text{a P4 addition}\ }\qquad[\mathrm{P4};\ \mathrm{P1}];
\]

%dsp-level:0
\[
(7)\quad\textbf{status:}
\]

%dsp-level:0
\[
\qquad \textbf{derived:}\ (2)\text{--}(5);
\]

%dsp-level:0
\[
\qquad \textbf{refuted:}\ \text{hypothesis}\ (\mathrm W)\ \text{as a derived coupling}\ (3);
\]

%dsp-level:0
\[
\qquad \textbf{priced,}\ \neg\textbf{adopted:}\ (6);
\]

%dsp-level:0
\[
\neg\text{a declaration},\quad\neg\text{a measured value}\ [\mathrm{P5}],\quad\neg\text{a new dimensionful scale}\ [\mathrm{P1}],\quad
\]

%dsp-level:0
\[
\qquad \textbf{the currencies remain distinct and unequated}\ [\mathrm{BM.83}(3)].
\]

\textbf{Req.} D0.4, D3.4 (Parts 1–2); CT0.8, EM.9 (Part 2); P1, P4, P5 (Parts 1–2); Def-H.37(2), FT.55(1), FT.73(2), FT.73(3) (Part 3); Def-BM.14(2), Def-BM.14(3), BM.19(2), BM.35(5), BM.36(3), BM.37(2), BM.59(9), BM.61(8), BM.67(6), BM.67(8), BM.69(3), BM.69(4), BM.75(4), BM.81(1), BM.81(3), BM.82(2), BM.82(4), BM.83(1), BM.83(3), BM.84(1) (this part); exact symbol and token censuses, field-degree rows re-cited, exponent arithmetic in the $\mathfrak q$-grading [math; machine: \texttt{checks/\allowbreak{}verify\_\allowbreak{}H\_\allowbreak{}em\_\allowbreak{}self\_\allowbreak{}energy\_\allowbreak{}correction.py} §D].

\textbf{Pf.}

%dsp-level:0
\[
(1):\ \text{the hypothesis stated with its own symbols}\ [\text{no premise imported}];
\]

%dsp-level:0
\[
\qquad (2):\ x_\pm\ \text{and}\ p\ \text{re-cited from}\ \mathrm{BM.19}(2)\ [\text{machine}];
\]

%dsp-level:0
\[
\qquad \text{the two symbol sets enumerated from their defining clauses}\ [\text{machine: the token census of}\ \mathrm{Def\text{-}BM.14}(3),\ \mathrm{Def\text{-}H.37}(2)\ \text{against}\ \mathrm{BM.19}(2),\ \mathrm{D0.4}]
\]

%dsp-level:0
\[
(3):\ \text{the four grounds re-cited at their own loci}\ [\mathrm{FT.73}(2),(3);\ \mathrm{BM.67}(6),(8);\ \mathrm{FT.55}(1);\ \mathrm{BM.36}(3);\ \mathrm{BM.61}(8)(\mathrm{O1});\ \mathrm{BM.69}(4)];
\]

%dsp-level:0
\[
\qquad \text{the intersection}\ \{\varepsilon_q\}\cap\{\delta_v,\hat\mu'^2\}=\varnothing\ [\text{machine}]
\]

%dsp-level:0
\[
(4):\ w\text{'s defining expression carries no}\ D_S\ [(2),(3);\ \text{machine}];
\]

%dsp-level:0
\[
\qquad m_{\mathrm{hid}}(s)
\]

%dsp-level:0
\[
= 2E(s)\ln2/\mathfrak h\ [\mathrm{BM.81}(1);\ \mathrm{FT.44}(4)]
\]

%dsp-level:0
\[
\implies m^{\mathrm{eff}}_{\mathrm{hid}}(s)
\]

%dsp-level:0
\[
= w\cdot2E(s)\ln2/\mathfrak h\ \text{with one}\ w\ [\text{machine}]
\]

%dsp-level:0
\[
(5):\ \Delta_g\ \text{and}\ \mathcal R_s\ \text{re-cited from their own tags};
\]

%dsp-level:0
\[
\qquad [\Delta_g]=[\hat\mu'^2]=\mathfrak q,\quad[\Delta m]=\mathfrak q^{+1/2}\ [\mathrm{CT2.1};\ \text{machine}]\implies\neg\text{addable}\ [\mathrm{BM.37}(2)\ \text{pattern}]
\]

%dsp-level:0
\[
(6)\text{--}(7):\ \text{the species row of a count}\to\text{strength coupling}\ [\mathrm{BM.36}(3);\ \mathrm{BM.61}(8)];
\]

%dsp-level:0
\[
\qquad \text{the standing-rule scans}\ [\text{machine}].\ \blacksquare
\]

\textbf{Phys.} The proposal that six defect values should reweight six retention probabilities is natural and it is false here — not because the numbers disagree but because nothing in the inventory can connect them. The flicker measure knows one symbol, the crystallization threshold, and it is derived as the unique stationary line of a two-state generator. The defect values know two symbols, an angle and a scale, and they live in a degree-sixteen field over the rationals while every read the measure touches is an integer or a binary weight. Between them stand four separate obstructions the corpus already proved: a field-degree gap, a two-element holonomy image that cannot reach four of the six states, a calibration weight that forbids the winding route, and the standing bar against turning a count into a strength. So the crystallization weight is branch-blind, the informational masses keep the values the census derived, attenuated equally on all six states, and the defect values keep their own currency, where they split a doublet and set a bridge factor and change nothing in bits. Declaring a defect-dependent rate would be exactly the kind of parameter this part has refused eleven times over, and it is refused again: the calibration registry stays empty, and the two currencies stay as incommensurate as the item that names them.

\end{theorem}

\begin{theorem}[{\texorpdfstring{BM.91 — The corrected neutron–proton splitting: the attenuated assembly, the ensemble-independent bar, and the universal ratio bound}{BM.91 - The corrected neutron-proton splitting: the attenuated assembly, the ensemble-independent bar, and the universal ratio bound}}]\label{tag:BM-91}
\textbf{Thm.}

%dsp-level:0
\[
(1)\quad\textbf{the corrected decomposition:}\qquad\boxed{\ \Delta m\ :=\ m_n-m_p\ =\ \Delta m_{\mathrm{transport}}+\Delta m_{\mathrm{EM}}^{\mathrm{corr}}+\Delta m_{\mathrm{bridge}}\ }
\]

%dsp-level:0
\[
\Delta m_{\mathrm{transport}}=+\frac{\ln(4/3)}{\mathfrak h}
\]

%dsp-level:0
\[
\qquad [\mathrm{BM.59}(4);\ \mathrm{BM.84}(1)];
\]

%dsp-level:0
\[
\qquad \Delta m_{\mathrm{EM}}^{\mathrm{corr}}=-\,\frac{3\,N\,w\,\mathsf g\,\varepsilon_{\mathrm{EM}}\ln2}{\mathfrak h}
\]

%dsp-level:0
\[
\qquad [\mathrm{BM.89}(5);\ \text{the path-integral read}];
\]

%dsp-level:0
\[
\Delta m_{\mathrm{bridge}}=0\ \ \text{in the informational currency}
\]

%dsp-level:0
\[
\qquad [\mathrm{BM.84}(1):\ \neg\text{formed there};\ \mathrm{BM.69}(3);\ \mathrm{BM.90}(5):\ \text{the pole companion}\ \Delta_g>0\ \text{is}\ \neg\text{addable}];
\]

%dsp-level:0
\[
\textbf{the four zeros stand:}
\]

%dsp-level:0
\[
\qquad \Delta m_{\mathrm{Lyapunov}}
\]

%dsp-level:0
\[
= 0\ [\mathrm{BM.51}(6)],
\]

%dsp-level:0
\[
\quad\Delta m_{\mathrm{metric}}
\]

%dsp-level:0
\[
= 0\ [\mathrm{BM.53}(2);\ \mathrm{BM.58}(4)],
\]

%dsp-level:0
\[
\quad\Delta m_{\mathrm{binding}}
\]

%dsp-level:0
\[
= 0\ [\mathrm{BM.54}(1)\text{--}(3)],
\]

%dsp-level:0
\[
\quad\Delta m_{\mathrm{force}}
\]

%dsp-level:0
\[
= \Delta m_{\mathrm{self\text{-}int}}
\]

%dsp-level:0
\[
= -\Delta m_{\mathrm{EM}}^{\mathrm{corr}}\ [\mathrm{BM.52}(4),(5)];
\]

%dsp-level:0
\[
(2)\quad\textbf{the attenuated assembly:}\qquad\boxed{\ m_n\,\mathfrak h\ =\ 6\,w\ln2+\ln20,\qquad m_p\,\mathfrak h\ =\ 6\,w\ln2+\ln15+3\,N\,w\,\mathsf g\,\varepsilon_{\mathrm{EM}}\ln2\ }
\]

%dsp-level:0
\[
[\mathrm{BM.89}(4),(5);\ \mathrm{Def\text{-}BM.13}(5)\text{'s refined rest input with the ensemble-averaged edge register};\ \text{machine: the prime-log assembly}\ \ln20=2\ln2+\ln5,\ \ln15=\ln3+\ln5];
\]

%dsp-level:0
\[
\textbf{the register scoping:}
\]

%dsp-level:0
\[
\qquad w\ \text{attenuates the edge-count register}\ E_{\mathrm{int}}=|R|\ \text{alone}\ [\mathrm{Def\text{-}BM.5}(4)];
\]

%dsp-level:0
\[
\qquad \text{the transport register}\ \Sigma_{\mathrm{cl}}(Q)
\]

%dsp-level:0
\[
= \ln W(Q)\ \text{is an end-sign filter count,}\ \neg\text{an edge count}\ [\mathrm{Def\text{-}BM.13}(4),(6)],
\]

%dsp-level:0
\[
\text{and its ensemble average is}\ \neg\text{constructed}\ [\mathrm{FT.106}(3)\text{'s census is at fixed topology}]\implies\textbf{it enters at its crystallized value and the residue is}\ (\mathrm{R4})\ \ [\mathrm{BM.89}(8);\ \mathrm{BM.93}(6)];
\]

%dsp-level:0
\[
\boxed{\ \Delta m\,\mathfrak h\ =\ \ln\frac43-3\,N\,w\,\mathsf g\,\varepsilon_{\mathrm{EM}}\ln2\ }\qquad[\text{machine:}\ \ln20-\ln15=\ln\tfrac43];
\]

%dsp-level:0
\[
\textbf{the cancellation row:}
\]

%dsp-level:0
\[
\qquad E_{\mathrm{int}}\ \text{is class-shared}\ [\mathrm{BM.57}(4);\ \text{Census FT.35}(2)]\implies\textbf{the crystallization weight cancels in the ledger difference and survives only in the EM term}
\]

%dsp-level:0
\[
[\text{machine: the coefficient census of}\ w\ \text{in the two rest inputs and in}\ \Delta m];
\]

%dsp-level:0
\[
\qquad [\Delta m]=\mathfrak q^{+1/2}=[\text{energy}]
\]

%dsp-level:0
\[
\qquad [\mathrm{CT2.1};\ \text{machine}];
\]

%dsp-level:0
\[
(3)\quad\textbf{the exact ratio:}\qquad\boxed{\ \frac{m_n}{m_p}\ =\ \frac{6\,w\ln2+\ln20}{6\,w\ln2+\ln15+3\,N\,w\,\mathsf g\,\varepsilon_{\mathrm{EM}}\ln2}\ }\qquad[\text{machine}];
\]

%dsp-level:0
\[
\text{at}\ (w,N,\mathsf g)=(1,1,1):
\]

%dsp-level:0
\[
\qquad \frac{m_n}{m_p}=\frac{\ln1280}{\ln960+3\,\varepsilon_{\mathrm{EM}}\ln2}
\]

%dsp-level:0
\[
\qquad [\mathrm{BM.84}(2)\text{'s row at the derived}\ c_{\mathrm{formed}}=\tfrac12;\ \text{machine:}\ 6\ln2+\ln20=\ln1280,\ 6\ln2+\ln15=\ln960];
\]

%dsp-level:0
\[
(4)\quad\textbf{the exact critical strength:}
\]

%dsp-level:0
\[
\qquad \Delta m\ \text{is strictly decreasing in}\ \varepsilon_{\mathrm{EM}}\ \text{at every admissible}\ (N,w,\mathsf g)>0
\]

%dsp-level:0
\[
\qquad [\text{machine: the coefficient census}]:
\]

%dsp-level:0
\[
\boxed{\ \varepsilon^{\mathrm{crit}}(N,w,\mathsf g)\ :=\ \frac{\ln(4/3)}{3\,N\,w\,\mathsf g\,\ln2}\ =\ \frac{2-\log_23}{3\,N\,w\,\mathsf g}\ }
\]

%dsp-level:0
\[
\qquad \textbf{the unique zero of}\ \Delta m
\]

%dsp-level:0
\[
\qquad [\text{machine: solved exactly in the prime-log ring}];
\]

%dsp-level:0
\[
\textbf{the vacuum-gauge, fully crystallized value:}
\]

%dsp-level:0
\[
\qquad \varepsilon^{\mathrm{crit}}(1,1,1)=\frac{2-\log_23}{3}\in\Bigl(\frac18,\frac17\Bigr)
\]

%dsp-level:0
\[
\qquad [\text{machine:}\ 2^{13}=8192>6561=3^8;\quad2^{11}=2048<2187=3^7];
\]

%dsp-level:0
\[
\textbf{the exact squeeze of the numerator:}
\]

%dsp-level:0
\[
\qquad \frac25\ <\ 2-\log_23\ <\ \frac5{12}
\]

%dsp-level:0
\[
\qquad [\text{machine:}\ \log_23<\frac85\iff3^5=243<256=2^8;\quad\log_23>\frac{19}{12}\iff3^{12}=531441>524288=2^{19}]:
\]

%dsp-level:0
\[
\frac{2}{15\,Nw\mathsf g}\ <\ \varepsilon^{\mathrm{crit}}(N,w,\mathsf g)\ <\ \frac{5}{36\,Nw\mathsf g}
\]

%dsp-level:0
\[
\qquad [\text{machine: two rational bounds on the critical strength at every admissible}\ (N,w,\mathsf g)];
\]

%dsp-level:0
\[
\textbf{the irrationality row:}
\]

%dsp-level:0
\[
\qquad \varepsilon^{\mathrm{crit}}(N,w,\mathsf g)\notin\mathbb Q\ \ \text{for every rational}\ (N,w,\mathsf g)>0
\]

%dsp-level:0
\[
\qquad [\mathrm{BM.84}(4):\ 2-\log_23\notin\mathbb Q\ \text{on the}\ 2^p\ne3^q\ \text{census};\ \text{machine}];
\]

%dsp-level:0
\[
\textbf{the dissolution row:}
\]

%dsp-level:0
\[
\qquad \text{on the Bernoulli specialization}\ w=1-4\varepsilon_q\ [\mathrm{BM.89}(4)]:
\]

%dsp-level:0
\[
\qquad \varepsilon^{\mathrm{crit}}\to\infty\ \ \text{as}\ \ \varepsilon_q\to\tfrac14^{-}
\]

%dsp-level:0
\[
\boxed{\ \textbf{the fully dissolved limit never closes the splitting}\ }
\]

%dsp-level:0
\[
\qquad [\text{machine: monotone in}\ w,\ \text{unbounded as}\ w\to0^+];
\]

%dsp-level:0
\[
(5)\quad\textbf{the ensemble-independent bar (parameter-free):}\qquad\boxed{\ N\,w\,\mathsf g\cdot\varepsilon_{\mathrm{EM}}\ \le\ \frac18\qquad\Longrightarrow\qquad\Delta m>0\ }
\]

%dsp-level:0
\[
[\text{machine:}\ \Delta m>0\iff Nw\mathsf g\,\varepsilon_{\mathrm{EM}}<\frac{2-\log_23}3;\qquad\frac{2-\log_23}3>\frac18\iff8(2-\log_23)>3\iff2^{13}>3^8\iff8192>6561];
\]

%dsp-level:0
\[
\textbf{the bar is uniform over the lapse, the crystallization weight and the curvature composition}\ —\ \neg\text{one of the three is a datum the corpus fixes};
\]

%dsp-level:0
\[
\qquad \text{the two priced witnesses at}\ \varepsilon_{\mathrm{EM}}
\]

%dsp-level:0
\[
= \tfrac19\ \text{stand}\ [\mathrm{BM.62}(8):\ 4^3=64>54=3^3\!\cdot\!2;\quad4^5=1024>972=3^5\!\cdot\!2^2];
\]

%dsp-level:0
\[
(6)\quad\textbf{the trichotomy (total, exact, non-empty on both sides):}
\]

%dsp-level:0
\[
\boxed{\ \varepsilon_{\mathrm{EM}}<\varepsilon^{\mathrm{crit}}\iff m_n>m_p\iff\frac{m_n}{m_p}>1;\qquad\varepsilon_{\mathrm{EM}}=\varepsilon^{\mathrm{crit}}\iff m_n=m_p\iff\frac{m_n}{m_p}=1;\qquad\varepsilon_{\mathrm{EM}}>\varepsilon^{\mathrm{crit}}\iff m_n<m_p\iff\frac{m_n}{m_p}<1\ }
\]

%dsp-level:0
\[
\textbf{the opening region is non-empty:}
\]

%dsp-level:0
\[
\qquad \Bigl\{0\le Nw\mathsf g\,\varepsilon_{\mathrm{EM}}\le\tfrac18\Bigr\}\ne\varnothing\ \ \text{and contains}\ \varepsilon_{\mathrm{EM}}=0\ \ \text{and}\ \ \varepsilon_{\mathrm{EM}}=\tfrac19\ \ \text{at}\ Nw\mathsf g=1
\]

%dsp-level:0
\[
\qquad [(5);\ \text{machine}];
\]

%dsp-level:0
\[
\textbf{the closing region is non-empty:}
\]

%dsp-level:0
\[
\qquad \varepsilon_{\mathrm{EM}}>\varepsilon^{\mathrm{crit}}(N,w,\mathsf g)\ \ \text{is inhabited at every admissible}\ (N,w,\mathsf g)>0
\]

%dsp-level:0
\[
\qquad [\text{machine:}\ \varepsilon^{\mathrm{crit}}\ \text{is a positive exact value}];
\]

%dsp-level:0
\[
\textbf{the critical locus is exact, unique, and irrational at rational}\ (N,w,\mathsf g)\qquad[(4)];
\]

%dsp-level:0
\[
(7)\quad\textbf{the universal ratio bound:}
\]

%dsp-level:0
\[
\qquad \frac{m_n}{m_p}\ \text{is strictly decreasing in}\ \varepsilon_{\mathrm{EM}}\ \text{and strictly decreasing in}\ w
\]

%dsp-level:0
\[
\qquad [\text{machine: the two cross-multiplied sign rows}]:
\]

%dsp-level:0
\[
\textbf{on the sub-critical branch:}
\]

%dsp-level:0
\[
\qquad 1\ <\ \frac{m_n}{m_p}\ <\ \frac{\ln20}{\ln15}
\]

%dsp-level:0
\[
\qquad \text{over all}\ \ w\in(0,1],\ N,\mathsf g>0,\ 0\le\varepsilon_{\mathrm{EM}}<\varepsilon^{\mathrm{crit}}
\]

%dsp-level:0
\[
\qquad [\text{the lower end open at the critical locus};\ \text{the upper end a strict supremum as}\ w\to0^+,\ \varepsilon_{\mathrm{EM}}\to0];
\]

%dsp-level:0
\[
\textbf{at}\ \varepsilon_{\mathrm{EM}}=0\ \textbf{specifically:}
\]

%dsp-level:0
\[
\qquad \frac{\ln1280}{\ln960}\ \le\ \frac{m_n}{m_p}\bigg|_{\varepsilon_{\mathrm{EM}}=0}\ <\ \frac{\ln20}{\ln15}
\]

%dsp-level:0
\[
\qquad [\text{the lower endpoint attained at}\ w=1;\ \text{machine: the monotonicity in}\ w];
\]

%dsp-level:0
\[
\frac{\ln1280}{\ln960}\in\Bigl(\frac{26}{25},\frac{21}{20}\Bigr)
\]

%dsp-level:0
\[
\qquad [\mathrm{BM.84}(3);\ \text{machine:}\ 2^{44}>3^{26}\!\cdot\!5,\ 2^{34}<3^{21}\!\cdot\!5];
\]

%dsp-level:0
\[
\qquad \frac{\ln20}{\ln15}\in\Bigl(\frac{11}{10},\frac{10}9\Bigr)
\]

%dsp-level:0
\[
[\text{machine, derived here}:\quad10\ln20>11\ln15\iff20^{10}>15^{11}\ (10240000000000>8649755859375);\qquad9\ln20<10\ln15\iff20^9<15^{10}\ (512000000000<576650390625)];
\]

%dsp-level:0
\[
\boxed{\ 1\ <\ \frac{m_n}{m_p}\ <\ \frac{10}9\qquad\text{at every admissible}\ (w,\,N,\,\mathsf g,\,\varepsilon_{\mathrm{EM}})\ \text{on the sub-critical branch}\ }
\]

%dsp-level:0
\[
(8)\quad\textbf{the incommensurability, ensemble-robust:}
\]

%dsp-level:0
\[
\qquad \frac{m_n}{m_p}=\frac43\iff3\bigl(6w\ln2+\ln20\bigr)=4\bigl(6w\ln2+\ln15+3Nw\mathsf g\,\varepsilon_{\mathrm{EM}}\ln2\bigr)
\]

%dsp-level:0
\[
\iff12\,Nw\mathsf g\,\varepsilon_{\mathrm{EM}}\ln2=(6-6w)\ln2-\ln405\iff\boxed{\ \varepsilon_{\mathrm{EM}}^{*}(w,N,\mathsf g)=\frac{(6-6w)\ln2-\ln405}{12\,N\,w\,\mathsf g\,\ln2}\ }
\]

%dsp-level:0
\[
\text{and}\qquad(6-6w)\ln2\ \le\ 6\ln2\ =\ \ln64\ <\ \ln405\qquad[\text{machine:}\ 64<405;\qquad w>0]\implies\boxed{\ \varepsilon_{\mathrm{EM}}^{*}<0\ \ \text{for every}\ w\in(0,1]\ }
\]

%dsp-level:0
\[
\textbf{the informational ratio never equals the ensemble ratio at any admissible coupling and any admissible crystallization weight}
\]

%dsp-level:0
\[
\qquad [\mathrm{BM.83}(3);\ \mathrm{BM.84}(8)\ \text{generalized};\ \text{machine}];
\]

%dsp-level:0
\[
\textbf{the irrationality survives:}
\]

%dsp-level:0
\[
\qquad \frac{m_n}{m_p}\notin\mathbb Q\ \ \text{at every rational}\ (w,N,\mathsf g,\varepsilon_{\mathrm{EM}})
\]

%dsp-level:0
\[
\qquad [\mathrm{BM.83}(2)\text{'s lemma: the}\ \ln3\ \text{coefficient of the numerator is}\ 0\ \text{and of the denominator is}\ 1;\ \text{machine}];
\]

%dsp-level:0
\[
(9)\quad\textbf{the rescope rows (recorded,}\ \neg\textbf{retracted):}
\]

%dsp-level:0
\[
\qquad \mathrm{BM.62}(3)\text{--}(9),
\]

%dsp-level:0
\[
\ \mathrm{BM.63},
\]

%dsp-level:0
\[
\ \mathrm{BM.65}(3)\text{--}(5),
\]

%dsp-level:0
\[
\ \mathrm{BM.66}(3),
\]

%dsp-level:0
\[
\ \mathrm{BM.70},
\]

%dsp-level:0
\[
\ \mathrm{BM.75}(4)\text{'s opposed row},
\]

%dsp-level:0
\[
\ \mathrm{BM.77}(9),
\]

%dsp-level:0
\[
\ \mathrm{BM.84}(2)\text{--}(8),
\]

%dsp-level:0
\[
\ \mathrm{BM.85}(1)\text{--}(3):
\]

%dsp-level:0
\[
\textbf{stand on their own locus}\ —\ \text{the}\ (w,N,\mathsf g)=(1,1,1)\ \text{specialization with}\ c\ \text{symbolic or priced};
\]

%dsp-level:0
\[
\qquad \text{every integer certificate they carry is re-derived here at the specialization and none is weakened};
\]

%dsp-level:0
\[
\mathrm{BM.62}(7)\text{'s four-certificate enclosure keeps both rows:}\ \text{the}\ c
\]

%dsp-level:0
\[
= \tfrac35\ \text{row stands as the void alternative's exact arithmetic},
\]

%dsp-level:0
\[
\ \text{and the}\ c
\]

%dsp-level:0
\[
= \tfrac12\ \text{row becomes}\ (4)\text{'s derived value};
\]

%dsp-level:0
\[
\neg\text{a frozen refutation is reopened};
\]

%dsp-level:0
\[
\qquad \neg\text{a Census cell changes};
\]

%dsp-level:0
\[
\qquad \mathrm{BM.77}(5)\text{'s specified-read trichotomy is untouched}\ [\mathrm{BM.92}(3)];
\]

%dsp-level:0
\[
(10)\quad\textbf{status:}
\]

%dsp-level:0
\[
\qquad \textbf{derived:}\ (1)\text{--}(8);
\]

%dsp-level:0
\[
\qquad \textbf{recorded:}\ (9);
\]

%dsp-level:0
\[
\qquad \textbf{symbolic:}\ \varepsilon_{\mathrm{EM}},\ N,\ \mathsf g,\ w\ (w\ \text{evaluated exactly at}\ \varepsilon_q=0\ \text{and on the marked specialization});
\]

%dsp-level:0
\[
\neg\text{a P4 addition},
\]

%dsp-level:0
\[
\quad\neg\text{a new dimensionful scale}\ [\mathrm{P1}],
\]

%dsp-level:0
\[
\quad\neg\text{a measured value}\ [\mathrm{P5}],
\]

%dsp-level:0
\[
\quad\neg\text{a declaration},
\]

%dsp-level:0
\[
\quad\text{§Cal-BM stays empty};
\]

%dsp-level:0
\[
\textbf{the open row:}\ \text{the numerical value of}\ e\ [\mathrm{EM.25}(7)]\ —\ \text{it decides which side of}\ \varepsilon^{\mathrm{crit}}\ \text{the corpus sits on, and no theorem of Parts 1--4 supplies it}\ [\mathrm{BM.62}(12);\ \mathrm{BM.84}(9)\ \text{re-cited}].
\]

\textbf{Req.} CT2.1, D2.9, EM.25(7) (Part 2); P1, P4, P5 (Parts 1–2); FT.44(4) (Part 3); Census FT.35(2) (row-wise); Def-BM.13(5), BM.51(6), BM.52(4), BM.52(5), BM.53(2), BM.54(1), BM.54(3), BM.57(4), BM.58(4), BM.59(4), BM.62(3), BM.62(7), BM.62(8), BM.62(12), BM.63, BM.65(3), BM.65(5), BM.66(3), BM.69(3), BM.70, BM.75(4), BM.77(9), BM.83(2), BM.83(3), BM.84(1), BM.84(2), BM.84(3), BM.84(4), BM.84(8), BM.84(9), BM.85(1), BM.85(3), BM.89(4), BM.89(5), BM.90(5) (this part); exact prime-log arithmetic with integer power certificates, Laurent-monomial assembly with coefficient censuses, cross-multiplied monotonicity rows, exact rational enclosures of $\ln$-ratios [math; machine: \texttt{checks/\allowbreak{}verify\_\allowbreak{}H\_\allowbreak{}em\_\allowbreak{}self\_\allowbreak{}energy\_\allowbreak{}correction.py} §E].

\textbf{Pf.}

%dsp-level:0
\[
(1):\ \text{the five zero-or-monomial rows re-cited}\ [\mathrm{BM.62}(1),(2);\ \mathrm{BM.84}(1)];
\]

%dsp-level:0
\[
\qquad \Delta m_{\mathrm{EM}}^{\mathrm{corr}}=-\bigl(E_{\mathrm{EM}}^{\mathrm{self}}(+1)-E_{\mathrm{EM}}^{\mathrm{self}}(0)\bigr)=-\tfrac{3Nw\mathsf g\varepsilon_{\mathrm{EM}}\ln2}{\mathfrak h}\ [\mathrm{BM.89}(5);\ \text{machine}]
\]

%dsp-level:0
\[
(2):\ 6w\ln2+\ln20-(6w\ln2+\ln15+3Nw\mathsf g\varepsilon_{\mathrm{EM}}\ln2)
\]

%dsp-level:0
\[
= \ln20-\ln15-3Nw\mathsf g\varepsilon_{\mathrm{EM}}\ln2
\]

%dsp-level:0
\[
= \ln\tfrac43-3Nw\mathsf g\varepsilon_{\mathrm{EM}}\ln2\ [\text{machine: the prime-log ring}];
\]

%dsp-level:0
\[
\qquad \ln20=2\ln2+\ln5,\ \ln15=\ln3+\ln5\ [\text{machine:}\ 20=2^2\cdot5,\ 15=3\cdot5];
\]

%dsp-level:0
\[
\qquad w\text{'s coefficient in}\ \Delta m\ \text{is}\ 0\ \text{in the ledger part}\ [\text{machine}]
\]

%dsp-level:0
\[
(3):\ \text{the quotient of}\ (2)\text{'s two rows}\ [\text{machine}];
\]

%dsp-level:0
\[
\qquad (w,N,\mathsf g)=(1,1,1):\ 6\ln2+\ln20=8\ln2+\ln5=\ln(2^8\cdot5)=\ln1280,
\]

%dsp-level:0
\[
\quad 6\ln2+\ln15=6\ln2+\ln3+\ln5=\ln960\ [\text{machine}]
\]

%dsp-level:0
\[
(4):\ \Delta m
\]

%dsp-level:0
\[
= 0\iff\ln\tfrac43
\]

%dsp-level:0
\[
= 3Nw\mathsf g\,
\]

%dsp-level:0
\[
\varepsilon_{\mathrm{EM}}\ln2\iff\varepsilon_{\mathrm{EM}}
\]

%dsp-level:0
\[
= \tfrac{\ln(4/3)}{3Nw\mathsf g\ln2}
\]

%dsp-level:0
\[
= \tfrac{2\ln2-\ln3}{3Nw\mathsf g\ln2}
\]

%dsp-level:0
\[
= \tfrac{2-\log_23}{3Nw\mathsf g}\ [\text{machine: solved in the prime-log ring}];
\]

%dsp-level:0
\[
\qquad \tfrac{2-\log_23}3>\tfrac18\iff8(2-\log_23)>3\iff13>8\log_23\iff2^{13}>3^8\ [\text{machine:}\ 8192>6561]
\]

%dsp-level:0
\[
\tfrac{2-\log_23}3<\tfrac17\iff7(2-\log_23)<3\iff11<7\log_23\iff2^{11}<3^7\ [\text{machine:}\ 2048<2187];
\]

%dsp-level:0
\[
\qquad 2-\log_23\notin\mathbb Q\ [\mathrm{BM.84}(4)\text{'s census}];
\]

%dsp-level:0
\[
\qquad \varepsilon^{\mathrm{crit}}=O(1/w)\to\infty\ \text{as}\ w\to0^+\ [\text{machine: monotone, unbounded}]
\]

%dsp-level:0
\[
(5):\ \Delta m>0\iff3Nw\mathsf g\,\varepsilon_{\mathrm{EM}}\ln2<\ln\tfrac43\iff Nw\mathsf g\,\varepsilon_{\mathrm{EM}}<\tfrac{2-\log_23}3\ [\text{machine}];
\]

%dsp-level:0
\[
\qquad \tfrac18<\tfrac{2-\log_23}3\ [(4)]\implies Nw\mathsf g\,\varepsilon_{\mathrm{EM}}\le\tfrac18\implies\Delta m>0\ [\text{machine}]
\]

%dsp-level:0
\[
(6):\ (4)\text{'s uniqueness}\ \wedge\ \text{the strict monotonicity in}\ \varepsilon_{\mathrm{EM}}\implies\text{the three-way equivalence}\ [\text{machine: the sign census on a rational grid of}\ (N,w,\mathsf g,\varepsilon_{\mathrm{EM}})];
\]

%dsp-level:0
\[
\qquad \varepsilon_{\mathrm{EM}}=0\in\text{the opening region since}\ \ln\tfrac43>0\ [4>3;\ \text{machine}]
\]

%dsp-level:0
\[
(7):\ \text{monotone in}\ \varepsilon_{\mathrm{EM}}:\ \text{the denominator increases, the numerator is fixed}\ [\text{machine}];
\]

%dsp-level:0
\[
\qquad \text{monotone in}\ w\ \text{at}\ \varepsilon_{\mathrm{EM}}
\]

%dsp-level:0
\[
= 0:\ w_1<w_2\implies\frac{6w_1\ln2+\ln20}{6w_1\ln2+\ln15}>\frac{6w_2\ln2+\ln20}{6w_2\ln2+\ln15}
\]

%dsp-level:0
\[
\iff (w_2-w_1)6\ln2(\ln20-\ln15)>0\ [\text{machine: cross-multiplied,}\ \ln20>\ln15\ \text{on}\ 20>15]
\]

%dsp-level:0
\[
10\ln20-11\ln15=\ln\frac{20^{10}}{15^{11}}>0\iff20^{10}>15^{11}\ [\text{machine:}\ 10240000000000>8649755859375];
\]

%dsp-level:0
\[
\qquad 9\ln20-10\ln15=\ln\frac{20^9}{15^{10}}<0\iff20^9<15^{10}\ [\text{machine:}\ 512000000000<576650390625]
\]

%dsp-level:0
\[
(8):\ 3(6w\ln2+\ln20)-4(6w\ln2+\ln15+3Nw\mathsf g\,\varepsilon_{\mathrm{EM}}\ln2)
\]

%dsp-level:0
\[
= -6w\ln2-4\ln3-\ln5-12Nw\mathsf g\,
\]

%dsp-level:0
\[
\varepsilon_{\mathrm{EM}}\ln2\ [\text{machine: the prime-log ring}];
\]

%dsp-level:0
\[
\qquad \text{setting it to}\ 0\ \text{gives}\ \varepsilon^*_{\mathrm{EM}}\ \text{of}\ (8)
\]

%dsp-level:0
\[
(6-6w)\ln2
\]

%dsp-level:0
\[
\le 6\ln2
\]

%dsp-level:0
\[
= \ln64<\ln405
\]

%dsp-level:0
\[
\iff 64<405\ [\text{machine}]\implies\text{the numerator of}\ \varepsilon^*_{\mathrm{EM}}\ \text{is negative}\implies\varepsilon^*_{\mathrm{EM}}<0\ [\text{machine}];
\]

%dsp-level:0
\[
\qquad \text{irrationality}:\ q(6w\ln2+\ln20)
\]

%dsp-level:0
\[
= p(6w\ln2+\ln15+3Nw\mathsf g\,\varepsilon_{\mathrm{EM}}\ln2)\ \text{with rational}\ (w,N,\mathsf g,\varepsilon_{\mathrm{EM}})\implies\text{the}\ \ln3\ \text{coefficient}\ -p
\]

%dsp-level:0
\[
= 0
\]

%dsp-level:0
\[
\implies p
\]

%dsp-level:0
\[
= 0
\]

%dsp-level:0
\[
\implies q
\]

%dsp-level:0
\[
= 0\ [\mathrm{BM.83}(2);\ \text{machine}]
\]

%dsp-level:0
\[
(9)\text{--}(10):\ \text{the clause census and the standing-rule scans}\ [\text{machine}].\ \blacksquare
\]

\textbf{Phys.} Put the corrected field energy into the splitting and almost nothing moves — which is the point. Both classes lose the same fraction of their edge register to the flicker, so the crystallization weight cancels out of the ledger difference entirely and survives only where the field lives: the electromagnetic term is now the shell value, times the lapse, times the crystallization weight, times the curvature composition, times three times the logarithm of two over h. The critical strength inherits those three factors in the denominator, which means the opening region can only widen as the composite fails to crystallize, and in the fully dissolved limit it never closes at all. Better, the bar becomes parameter-free in the ensemble data: whenever the product of lapse, weight, curvature composition and coupling stays at or below one eighth, the neutron is the heavier baryon, and the certificate is a single pair of integers, eight thousand one hundred ninety-two against six thousand five hundred sixty-one. The ratio itself is bounded between one and ten ninths at every admissible value of everything the corpus declines to fix, with the old one-and-a-twentieth bound recovered as the fully crystallized endpoint and the logarithm of twenty over the logarithm of fifteen as the dissolved supremum. And the incommensurability hardens: the coupling at which the informational ratio would equal the ensemble's four thirds is negative for every crystallization weight in the admissible range, on the single integer certificate that sixty-four is less than four hundred five.

\end{theorem}

\begin{theorem}[{\texorpdfstring{BM.92 — The Q-value on both formed reads: the uniform closure, the corrected margins, and the untouched specified branch}{BM.92 - The Q-value on both formed reads: the uniform closure, the corrected margins, and the untouched specified branch}}]\label{tag:BM-92}
\textbf{Thm.}

%dsp-level:0
\[
(1)\quad\textbf{the informational Q-value with the corrected term:}
\]

%dsp-level:0
\[
\qquad Q_\beta:=m_n-m_p-m_\ell-m_\nu
\]

%dsp-level:0
\[
\qquad [\mathrm{BM.41}(4);\ \mathrm{BM.63}(1);\ \mathrm{BM.85}(1)]:
\]

%dsp-level:0
\[
\boxed{\ Q_\beta\,\mathfrak h\ =\ \ln\frac43-3\,N\,w\,\mathsf g\,\varepsilon_{\mathrm{EM}}\ln2-F_{\ell\nu},\qquad F_{\ell\nu}:=m_\ell\mathfrak h+m_\nu\mathfrak h\ }
\]

%dsp-level:0
\[
\qquad [\mathrm{BM.91}(2);\ \text{machine: the prime-log assembly}];
\]

%dsp-level:0
\[
F_{\ell\nu}\ge2\ln2
\]

%dsp-level:0
\[
\qquad [\mathrm{BM.70}(2),(3):\ \text{every materialized closed chain carries}\ \ge2\ \text{directed edges};\ m\mathfrak h=E_{\mathrm{int}}\ln2;\quad m_\nu\mathfrak h\ge0\ \text{by PSD}];
\]

%dsp-level:0
\[
\boxed{\ Q_\beta\,\mathfrak h\ \le\ \ln\frac43-2\ln2\ =\ -\ln3\ <\ 0\qquad\text{at every admissible}\ (w,\,N,\,\mathsf g,\,\varepsilon_{\mathrm{EM}})\ }
\]

%dsp-level:0
\[
\qquad [\text{machine:}\ 3Nw\mathsf g\,\varepsilon_{\mathrm{EM}}\ln2\ge0;\qquad4<12];
\]

%dsp-level:0
\[
\textbf{the closure is uniform over the ensemble data:}
\]

%dsp-level:0
\[
\qquad \text{the attenuation lowers the magnitude of the EM term and cannot raise}\ Q_\beta\ \text{above}\ -\ln3/\mathfrak h
\]

%dsp-level:0
\[
\qquad [\text{machine: the coefficient census}];
\]

%dsp-level:0
\[
(2)\quad\textbf{the margin ladder and the corrected ratio forms:}
\]

%dsp-level:0
\[
\qquad Q_\beta\mathfrak h\ \le\ -\ln3,\quad-\ln12,\quad-\ln24
\]

%dsp-level:0
\[
\qquad \text{on}\quad4<12,\quad4<48,\quad4<96
\]

%dsp-level:0
\[
\qquad [\mathrm{BM.70}(3);\ \mathrm{BM.85}(2)\ \text{re-certified};\ \text{machine}]:
\]

%dsp-level:0
\[
\Sigma_{\mathrm B}-2\ln2=-\ln3,
\]

%dsp-level:0
\[
\qquad \Sigma_{\mathrm B}-4\ln2=-\ln12,
\]

%dsp-level:0
\[
\qquad \Sigma_{\mathrm B}-5\ln2=-\ln24
\]

%dsp-level:0
\[
\qquad [\text{machine: the three exact identities}];
\]

%dsp-level:0
\[
\boxed{\ \frac{Q_\beta}{m_n}\ \le\ -\,\frac{\ln3}{6\,w\ln2+\ln20}\ }\qquad\text{with the two exact endpoint brackets:}
\]

%dsp-level:0
\[
w=1:
\]

%dsp-level:0
\[
\qquad -\frac{\ln3}{\ln1280}\in\Bigl(-\frac2{13},-\frac3{20}\Bigr)
\]

%dsp-level:0
\[
\qquad [\mathrm{BM.85}(3);\ \text{machine:}\ 3^{13}=1594323<1638400=2^{16}\!\cdot\!5^2;\quad3^{20}=3486784401>2097152000=2^{24}\!\cdot\!5^3];
\]

%dsp-level:0
\[
w\to0^+:
\]

%dsp-level:0
\[
\qquad -\frac{\ln3}{\ln20}\in\Bigl(-\frac25,-\frac13\Bigr)
\]

%dsp-level:0
\[
\qquad [\text{machine, derived here}:\quad3\ln3>\ln20\iff27>20;\qquad5\ln3<2\ln20\iff243<400];
\]

%dsp-level:0
\[
\textbf{the informational Q-value is a strictly negative exact fraction of the parent mass at every crystallization weight, bounded away from zero by a rational};
\]

%dsp-level:0
\[
(3)\quad\textbf{the specified read is untouched:}
\]

%dsp-level:0
\[
\qquad Q_\beta^{\mathrm{spec}}\,\mathfrak h=\Delta T(S)-\rho_\ell\sqrt{G\bigl(c(0)\bigr)}-\rho_\nu\sqrt{G\bigl(c(i_\nu)\bigr)},
\]

%dsp-level:0
\[
\qquad \Delta T(S)=5\,m^2_B(\cdot,S)
\]

%dsp-level:0
\[
\qquad [\mathrm{BM.77}(2);\ \mathrm{BM.85}(4)]:
\]

%dsp-level:0
\[
\textbf{the EM summands are}\ \neg\textbf{parts of that read}
\]

%dsp-level:0
\[
\qquad [\mathrm{BM.77}(3):\ \text{disjoint alphabets, no double count}]\implies\boxed{\ \text{neither the rescission of}\ \mathrm{BM.87}\ \text{nor the ensemble correction of}\ \mathrm{BM.89}\ \text{alters any row of}\ \mathrm{BM.77}(5),(6)\ }
\]

%dsp-level:0
\[
\Delta T=\Bigl(\frac{4219785}{98956},\ \frac{9608310}{393679},\ \frac{455}{23},\ \frac{68565}{3088}\Bigr),
\]

%dsp-level:0
\[
\qquad \frac{273}{23}<\rho^{\mathrm{crit}}(S,i_\nu)<\frac{42197850}{791648}
\]

%dsp-level:0
\[
\qquad [\mathrm{BM.77}(2),(5);\ \text{machine re-certified}];
\]

%dsp-level:0
\[
\textbf{the witnesses stand:}
\]

%dsp-level:0
\[
\qquad 828100>5819,
\]

%dsp-level:0
\[
\qquad 4219785<43\cdot98956=4255108,
\]

%dsp-level:0
\[
\qquad 7500>1849,
\]

%dsp-level:0
\[
\qquad 3312400>5819,
\]

%dsp-level:0
\[
\qquad 99<100,
\]

%dsp-level:0
\[
\qquad 75>64,
\]

%dsp-level:0
\[
\qquad 273>253,
\]

%dsp-level:0
\[
\qquad 42197850<42748992;
\]

%dsp-level:0
\[
(4)\quad\textbf{the verdict rows (read-relative and complete):}
\]

%dsp-level:0
\[
\qquad \textbf{closed on the informational currency for every member, uniformly in}\ (w,N,\mathsf g,\varepsilon_{\mathrm{EM}})
\]

%dsp-level:0
\[
\qquad [(1);\ \mathrm{BM.70}(4)(\mathrm a)];
\]

%dsp-level:0
\[
\textbf{open on}\ \mathcal O_S,\ \textbf{closed on}\ \mathcal C_S,\ \textbf{degenerate on}\ \partial\mathcal O_S\ \ \text{on the specified read}
\]

%dsp-level:0
\[
\qquad [\mathrm{BM.77}(5),(6);\ \mathrm{BM.85}(5)];
\]

%dsp-level:0
\[
\textbf{the decay direction is fixed in both currencies:}
\]

%dsp-level:0
\[
\qquad T(0,S)=\tfrac43T(+1,S)>T(+1,S)
\]

%dsp-level:0
\[
\qquad [\mathrm{BM.28}(4)(\mathrm{ii});\ \text{machine}]
\]

%dsp-level:0
\[
\qquad \text{and}
\]

%dsp-level:0
\[
\qquad m_n>m_p\ \ \text{on the sub-critical branch}
\]

%dsp-level:0
\[
\qquad [\mathrm{BM.91}(6)]:
\]

%dsp-level:0
\[
\boxed{\ \textbf{the parent is the neutral class on every formed read, and the grammar}\ n\to p+\ell+\nu\ \textbf{stands unmodified}\ }
\]

%dsp-level:0
\[
\qquad [\mathrm{BM.41}(1);\ \mathrm{BM.45};\ \mathrm{BM.77}(6)];
\]

%dsp-level:0
\[
\textbf{the conditional spectrum on}\ \mathcal O_S\textbf{:}
\]

%dsp-level:0
\[
\qquad E_\ell^{\max}=\frac{m_n^2+m_\ell^2-(m_p+m_\nu)^2}{2m_n},
\]

%dsp-level:0
\[
\qquad E_{\max}:\langle E\rangle=2:1\ \ \text{scale-invariant}
\]

%dsp-level:0
\[
\qquad [\mathrm{BM.41}(5);\ \mathrm{BM.49}(3),(4);\ \mathrm{BM.77}(8);\ \neg\text{an external endpoint compared}\ [\mathrm{P5}]];
\]

%dsp-level:0
\[
(5)\quad\textbf{the two reads are}\ \neg\textbf{comparable:}
\]

%dsp-level:0
\[
\qquad Q_\beta^{\mathrm{info}}:\ \{\ln2,\ln3,\ln5,\ N,\ w,\ \mathsf g,\ \varepsilon_{\mathrm{EM}},\ \mathfrak h,\ m_\ell,\ m_\nu\};
\]

%dsp-level:0
\[
\qquad Q_\beta^{\mathrm{spec}}:\ \{m^2_B(\cdot,S),\ \rho_\ell,\ \rho_\nu,\ g,\ g',\ \mathfrak h\}
\]

%dsp-level:0
\[
\boxed{\ \text{the two alphabets are disjoint except for}\ \mathfrak h\ }
\]

%dsp-level:0
\[
\qquad [\text{machine: the token census};\qquad\mathsf g\ \text{and}\ g\ \text{are distinct symbols}\ [(4)\text{'s local notation};\ \mathrm{BM.39}(1)]];
\]

%dsp-level:0
\[
\neg\text{a derived equation relates the two expressions}\ —\ \text{that is}\ (\mathrm{R1}),
\]

%dsp-level:0
\[
\ (\mathrm{R3}),
\]

%dsp-level:0
\[
\ (\mathrm{R4})\ \text{and the Cal-F scales together}
\]

%dsp-level:0
\[
\qquad [\mathrm{BM.82}(6);\ \mathrm{BM.83}(5),(7);\ \mathrm{BM.85}(5)];
\]

%dsp-level:0
\[
(6)\quad\textbf{status:}
\]

%dsp-level:0
\[
\qquad \textbf{derived:}\ (1)\text{--}(5);
\]

%dsp-level:0
\[
\qquad \textbf{read-relative:}\ \text{the kinematic verdict}\ (4);
\]

%dsp-level:0
\[
\qquad \textbf{declared: nothing}\ —\ \text{A-BM.8's specification is re-cited,}\ \neg\text{re-issued};
\]

%dsp-level:0
\[
\neg\text{a clause of}\ \mathrm{BM.41},\ \mathrm{BM.49},\ \mathrm{BM.63},\ \mathrm{BM.70},\ \mathrm{BM.71},\ \mathrm{BM.77},\ \mathrm{BM.85}\ \text{is retracted};
\]

%dsp-level:0
\[
\qquad \neg\text{a P4 addition},\quad\neg\text{a new dimensionful scale}\ [\mathrm{P1}],\quad\neg\text{a measured value}\ [\mathrm{P5}],\quad\text{§Cal-BM stays empty}.
\]

\textbf{Req.} P1, P4, P5 (Parts 1–2); FT.79(5) (Part 3); Def-BM.13(4), BM.28(4), BM.37(2), BM.39(1), BM.41(1), BM.41(4), BM.41(5), BM.45, BM.49(3), BM.49(4), BM.63(1), BM.70(2), BM.70(3), BM.70(4), BM.77(2), BM.77(3), BM.77(5), BM.77(6), BM.77(8), BM.82(6), BM.83(5), BM.83(7), BM.84(1), BM.85(1), BM.85(2), BM.85(3), BM.85(4), BM.85(5), BM.87, BM.89, BM.91(2), BM.91(6) (this part); the declaration addendum A-BM.8, item 7 (pointers); exact prime-log arithmetic with integer power certificates, exact rational arithmetic at the four frozen sector expectations, token censuses [math; machine: \texttt{checks/\allowbreak{}verify\_\allowbreak{}H\_\allowbreak{}em\_\allowbreak{}self\_\allowbreak{}energy\_\allowbreak{}correction.py} §F].

\textbf{Pf.}

%dsp-level:0
\[
(1):\ Q_\beta\mathfrak h=\Delta m\,\mathfrak h-F_{\ell\nu}=\ln\tfrac43-3Nw\mathsf g\,\varepsilon_{\mathrm{EM}}\ln2-F_{\ell\nu}\ [\mathrm{BM.91}(2);\ \text{machine}];
\]

%dsp-level:0
\[
\qquad F_{\ell\nu}\ge2\ln2\ [\mathrm{BM.70}(2),(3)\ \text{re-cited}];
\]

%dsp-level:0
\[
\qquad \ln\tfrac43-2\ln2=2\ln2-\ln3-2\ln2=-\ln3\ [\text{machine}];
\]

%dsp-level:0
\[
\qquad -\ln3<0\iff3>1\iff4<12\ [\text{machine}]
\]

%dsp-level:0
\[
3Nw\mathsf g\,
\]

%dsp-level:0
\[
\varepsilon_{\mathrm{EM}}\ln2
\]

%dsp-level:0
\[
\ge 0\ \text{at}\ N,
\]

%dsp-level:0
\[
w,
\]

%dsp-level:0
\[
\mathsf g,
\]

%dsp-level:0
\[
\varepsilon_{\mathrm{EM}}
\]

%dsp-level:0
\[
\ge 0\ [\text{machine: the coefficient census}]
\]

%dsp-level:0
\[
\implies Q_\beta\mathfrak h
\]

%dsp-level:0
\[
\le -\ln3\ \text{uniformly}\ [\text{machine}]
\]

%dsp-level:0
\[
(2):\ \Sigma_{\mathrm B}-2\ln2
\]

%dsp-level:0
\[
= -\ln3,
\]

%dsp-level:0
\[
\quad\Sigma_{\mathrm B}-4\ln2
\]

%dsp-level:0
\[
= 2\ln2-\ln3-4\ln2
\]

%dsp-level:0
\[
= -\ln(4\cdot3)
\]

%dsp-level:0
\[
= -\ln12,
\]

%dsp-level:0
\[
\quad\Sigma_{\mathrm B}-5\ln2
\]

%dsp-level:0
\[
= -\ln(8\cdot3)
\]

%dsp-level:0
\[
= -\ln24\ [\text{machine: the three identities}];
\]

%dsp-level:0
\[
\qquad 4<12<48<96\ [\text{machine}]
\]

%dsp-level:0
\[
Q_\beta/m_n=(Q_\beta\mathfrak h)/(m_n\mathfrak h)\le-\ln3/(6w\ln2+\ln20)\ [\mathrm{BM.91}(2);\ \text{machine}];
\]

%dsp-level:0
\[
\qquad 3\ln3-\ln20>0\iff27>20,\quad5\ln3-2\ln20<0\iff243<400\ [\text{machine}]
\]

%dsp-level:0
\[
(3):\ \mathrm{BM.77}(2),(3),(5)\ \text{re-cited};
\]

%dsp-level:0
\[
\qquad 15\cdot843957=12659355,
\]

%dsp-level:0
\[
\quad 15\cdot1921662=28824930,
\]

%dsp-level:0
\[
\quad 15\cdot91=1365,
\]

%dsp-level:0
\[
\quad 15\cdot13713=205695\ [\mathrm{BM.89}(6);\ \text{machine}];
\]

%dsp-level:0
\[
\qquad \text{the eight witnesses re-certified as integer comparisons}\ [\text{machine}]
\]

%dsp-level:0
\[
(4):\ T(0,S)/T(+1,S)=W(0)/W(+1)=20/15=4/3>1\ [\text{machine}];
\]

%dsp-level:0
\[
\qquad m_n>m_p\ \text{sub-critically}\ [\mathrm{BM.91}(6)];
\]

%dsp-level:0
\[
\qquad \text{the endpoint identity and the}\ 2{:}1\ \text{moment ratio re-cited}\ [\mathrm{BM.41}(5);\ \mathrm{BM.49}(3),(4)]
\]

%dsp-level:0
\[
(5):\ \text{the two symbol sets enumerated from the two expressions and intersected}\ [\text{machine}];
\]

%dsp-level:0
\[
\qquad \mathsf g\ne g\ \text{as tokens}\ [\text{machine}];
\]

%dsp-level:0
\[
\qquad (6):\ \text{the standing-rule scans}\ [\text{machine}].\ \blacksquare
\]

\textbf{Phys.} One Q-value, two formed reads, and the correction lands on only one of them. On the informational currency the decay's energy balance keeps its shape: the transport register's log of four thirds, less the corrected field term, less whatever the two emitted lines cost. Because the corrected field term is nonnegative and now attenuated by the lapse, the crystallization weight and the curvature composition, the two-bit floor on any materialized chain still drives the total below minus the logarithm of three — at every value of every quantity the corpus declines to fix. The closure is not merely preserved, it is uniform, and its size as a fraction of the parent mass runs from minus three twentieths at full crystallization to minus one third in the dissolved limit, always an exact rational bound away from zero. On the specified read nothing moves at all: the electromagnetic summands were never part of that expression, so the opening region, the closing region, the exact critical locus and all eight integer witnesses stand as derived. The two verdicts remain read-relative and incomparable, their alphabets sharing one symbol, and the one thing both agree on is the direction — the neutral class is the parent, the grammar is untouched, and the decay is allowed exactly where a formed read allows it.

\end{theorem}

\medskip\noindent\textit{\texorpdfstring{BM.93 — The status of the EM self-energy correction: adjudication and scope}{BM.93 - The status of the EM self-energy correction: adjudication and scope}.}\ Structural-constraints record: see Appendix B, Section B.1.28 of the Addendum.\par\medskip

\phantomsection
\section[Tier BM.19 — The Flicker Integral, the Rest-Mass Q-Value,…]{\texorpdfstring{Tier BM.19 — The Flicker Integral, the Rest-Mass Q-Value, and the Relating Equation: the Materialized Ensemble, the Vertex Superposition, the Rung-One $\nu$ Line, and the Bridge in Four Exact Forms}{Tier BM.19 - The Flicker Integral, the Rest-Mass Q-Value, and the Relating Equation: the Materialized Ensemble, the Vertex Superposition, the Rung-One Line, and the Bridge in Four Exact Forms}}

\textbf{Locus.} This tier fixes one read, moves one Q-value onto one species, and closes one ledger item in form. The read it fixes is the crystallized-snapshot read of a sealed composite's information content — the read that takes the computed cell's $E_{\mathrm{int}}=6$ $[\mathrm{Census\ FT.35}(2);\ \mathrm{BM.17}(1)]$ as the composite's count and hence as a fixed mass. That read is a category error against the corpus's own rows: the count is one point of an ensemble distribution $[\mathrm{Def\text{-}BM.5}(4);\ \mathrm{BM.19}(5)]$, the real and the virtual registers are mixed at every nonzero threshold $[\mathrm{A4};\ \mathrm{D0.1}]$, and the composite is an ensemble and $\neg$ a topology $[\mathrm{BM.19}(6);\ \mathrm{BM.89}(4)]$. $\mathrm{BM.94}$ replaces the snapshot with the integral, derives the distribution and its exact moments, and derives the materialization restriction that holds the mean edge count above the digon's two. $\mathrm{BM.95}$ derives the superposition a vertex carries in the transport register, with its exact weights and the two positive parts of the transport splitting. $\mathrm{BM.96}$ derives the rung-one $\nu$ line's rest mass from the seal structure, its strict positivity, its three-line order, its single-ray occupancy, and the four independent bars against a seesaw-type mechanism. $\mathrm{BM.97}$ adjudicates the rest-mass species and forms the Q-value on it, with the exact critical locus and integer-certified witnesses on both the specified and the pure-pole reads. $\mathrm{BM.98}$ derives the relating equation in four exact forms — count, curvature, coupling, flicker-weighted — with its enhancement bounded by an exact integer. $\mathrm{BM.99}$ recomputes the splitting and the Q-value on the materialized read, gaining a two-sided parameter-free bar and two exact tightenings. $\mathrm{BM.100}$ adjudicates.

\textbf{Standing rows.} $\alpha_B$ stays barred on $\mathrm{BM.61}(8)$'s four grounds and §Cal-BM stays empty. No clause of $\mathrm{BM.1}$–$\mathrm{BM.93}$ is retracted, replaced or renumbered: every row this tier touches is re-scoped at its own locus and recorded in $\mathrm{BM.94}(9)$, $\mathrm{BM.97}(7)$, $\mathrm{BM.99}(7)$ and $\mathrm{BM.100}(4)$, in $\mathrm{BM.77}(7)$'s pattern. No measured coupling, mass or ratio enters any derivation $[\mathrm{P5}]$: the neutron–proton ratio is derived and bounded at $\mathrm{BM.99}(5)$ and the electric coupling stays symbolic $[\mathrm{EM.25}(7);\ \mathrm{BM.88}(7)]$ — neither is imported. No right-ray matter line and no seesaw-type mechanism is suggested, added or required $[\mathrm{BM.96}(4),(6)]$. The relating equation is derived in form with its residue localized to one undeclared dimensionless ratio; $\rho_{\mathrm B}$ is $\neg$derived and $\mathrm{BM.69}(4)$'s four grounds stand unweakened. Local notation, fixed for the whole tier: $p:=1-4\varepsilon_q$ the stationary retention weight $[\mathrm{BM.19}(2)]$; $u:=1-p=4\varepsilon_q$ the flicker weight; $D:=3-3p+p^2=1+u+u^2$ the dilution denominator; $w$ the crystallization weight of $\mathrm{BM.89}(4)$ and $w_{\mathrm{mat}}:=1/D$ its materialized companion; $E^{\mathrm{cryst}}:=6$; $k:=r/\Delta x_{\mathrm B}$, $f(k)=1-\tfrac1k$, $N$, $\mathsf g$, $G$, $\varepsilon_{\mathrm{EM}}$ as in $\mathrm{BM.88}$'s local notation; $\rho_{\mathrm B}$, $D_S$, $\mathcal R$ as in $\mathrm{BM.69}(2)$; $M_0$, $s$, $\Delta_g$, $\Delta_g^{\mathrm{crit}}$ as in $\mathrm{BM.97}(4)$; $c:=\cos\tfrac18$, $s_\nu:=\sin\tfrac18$ in $\mathrm{BM.96}(2)$. The symbol $q$ keeps its end-register meaning $[\mathrm{D3.4};\ \mathrm{EM.9}]$, the italic $g$ keeps its $\mathrm{BM.39}(1)$ meaning, and the italic $s$ keeps its $\mathrm{BM.97}(4)$ meaning at every occurrence.

\begin{theorem}[{\texorpdfstring{BM.94 — The flicker integral: the snapshot read's category error, the materialized distribution, and the two crystallization weights}{BM.94 - The flicker integral: the snapshot read's category error, the materialized distribution, and the two crystallization weights}}]\label{tag:BM-94}
\textbf{Thm.}

%dsp-level:0
\[
(1)\quad\textbf{the category error (the referent census):}\qquad\begin{array}{l|l}\textbf{the snapshot read takes} & \textbf{the corpus supplies}\\\hline E_{\mathrm{int}}=6\ \text{as the composite's count} & E_{\mathrm{int}}(R)=|R|\ \text{distributed over the flicker ensemble}\ [\mathrm{Def\text{-}BM.5}(4);\ \mathrm{BM.19}(5)]\\\text{one configuration as the composite} & |\mathcal E(\mathcal X)|>1\iff\varepsilon_q>0:\ \text{an ensemble at every nonzero threshold}\ [\mathrm{BM.19}(6)]\\ m=\frac{6\ln2}{\mathfrak h}\ \text{as a fixed mass} & \text{a mixture of the real and the virtual registers}\ [\mathrm{A4};\ \mathrm{D0.1};\ \mathrm{BM.89}(4)]\\\text{a value} & \text{a distribution whose top point is the crystallized value}\ [\mathrm{Def\text{-}BM.5}(4);\ \mathrm{BM.86}(4)\ \text{scope}]\end{array}
\]

%dsp-level:0
\[
\textbf{the error, stated exactly:}
\]

%dsp-level:0
\[
\qquad \boxed{\ E_{\mathrm{int}}=6\ \ \text{is one point of the distribution,}\ \neg\text{the ensemble, and its weight is}\ P(E{=}6)=\frac{p^2}{D}<1\ \ \text{at every}\ \varepsilon_q>0\ }
\]

%dsp-level:0
\[
\textbf{the weight deficit is exact:}\qquad1-P_{\mathrm{mat}}(E{=}6)=\frac{3(1-p)}{D}>0\qquad[\text{machine:}\ D-p^2=3-3p=3(1-p)];
\]

%dsp-level:0
\[
(2)\quad\textbf{the proper integral:}
\]

%dsp-level:0
\[
\qquad \boxed{\ E_{\mathrm{int}}^{\mathrm{eff}}\ :=\ \sum_{R\,\in\,\mathcal E(\mathcal X)}P(R)\cdot E_{\mathrm{int}}(R)\ }
\]

%dsp-level:0
\[
\qquad [\mathrm{BM.87}(5);\ \mathrm{BM.20}(3);\ \mathrm{Def\text{-}H.32}(1);\ \mathrm{Def\text{-}BM.5}(2)]:
\]

%dsp-level:0
\[
P(R)\ \text{is the ensemble weight of the configuration}\ R,
\]

%dsp-level:0
\[
\ \text{derived as the unique stationary line of the per-pair Janus generator}\ [\mathrm{BM.19}(1),(2);\ \mathrm{FT.108}(4);\ \mathrm{Def\text{-}BM.5}(3)],
\]

%dsp-level:0
\[
\ \neg\text{declared};
\]

%dsp-level:0
\[
\text{the sum is finite and exact at every budget}\ B\ [\mathrm{Def\text{-}BM.5}(1);\ \mathrm{Def\text{-}H.32}(1)],
\]

%dsp-level:0
\[
\ \text{and the}\ B\to\infty\ \text{limit is}\ \neg\text{constructed}\ [\mathrm{Def\text{-}BM.5}(5)];
\]

%dsp-level:0
\[
\textbf{linearity on}\ |R|\ \text{gives the per-datum form:}
\]

%dsp-level:0
\[
\qquad E_{\mathrm{int}}^{\mathrm{eff}}=\sum_{\text{pairs}}P(\text{the pair is real})\cdot2
\]

%dsp-level:0
\[
\qquad [\mathrm{BM.89}(4)\ \text{re-derived as the first moment};\ \text{machine}];
\]

%dsp-level:0
\[
(3)\quad\textbf{the ledger read (the quantum numbers the read is taken on):}
\]

%dsp-level:0
\[
\qquad \mathfrak L=\bigcup_{n\ge0}\mathfrak L(\tau_n),
\]

%dsp-level:0
\[
\qquad \bigl|\mathfrak L(\tau_n)\bigr|=\sum_{k<n}n_{\mathrm{real\_out}}(\tau_k)
\]

%dsp-level:0
\[
\qquad [\mathrm{A7};\ \mathrm{A6}]:
\]

%dsp-level:0
\[
\textbf{the two quantum numbers are the tick and the scale:}
\]

%dsp-level:0
\[
\qquad \tau:\ \tau_n\to\tau_{n+1},\ \neg\exists\,\tau^{-1}\ [\mathrm{A6};\ \mathrm{D1.8}:\ x^0:=\tau_n],
\]

%dsp-level:0
\[
\qquad \mathfrak h:\ \exists!\ \text{free dimensionful scale}\ [\mathrm{P1};\ \mathrm{D2.9};\ \text{A §0.3}];
\]

%dsp-level:0
\[
n_{\mathcal R}=E_{\mathrm{int}}=|R|\ [\mathrm{D0.8};\ \mathrm{Def\text{-}BM.5}(4)],\qquad S_{\mathfrak L}=n_{\mathcal R}\ln2\ [\mathrm{C9};\ \mathrm{CT1.29};\ \mathrm{BM.35}(3)];
\]

%dsp-level:0
\[
\textbf{the crystallized card is the snapshot at one}\ \tau\ [\mathrm{Census\ FT.35}(2);\ \mathrm{BM.17}(1)]:
\]

%dsp-level:0
\[
\qquad \text{the composite is read from}\ \mathfrak L,\ \neg\text{from the card};
\]

%dsp-level:0
\[
\textbf{the constituent at a vertex is a quantum superposition over the ledger's transport fibre,}\ \neg\text{a fixed state}
\]

%dsp-level:0
\[
\qquad [\mathrm{BM.95}(2);\ \mathrm{Def\text{-}BM.13}(1),(3)]:
\]

%dsp-level:0
\[
(4)\quad\textbf{the derived distribution (the marked pair specialization, marked}\ \neg\textbf{adopted):}
\]

%dsp-level:0
\[
\qquad E^{\mathrm{cryst}}
\]

%dsp-level:0
\[
= 6
\]

%dsp-level:0
\[
= (c_1{:}\,2)+(c_2{:}\,2)+(\text{connector}{:}\,2)\ \ \text{is three pairs}\ [\mathrm{Census\ FT.35}(2);\ \mathrm{BM.21}(1);\ \mathrm{BM.89}(4)]:
\]

%dsp-level:0
\[
\text{each pair is retained with the stationary weight}\ p
\]

%dsp-level:0
\[
= 1-4\varepsilon_q\ [\mathrm{BM.19}(2);\ \text{an identity, exact at every rational threshold}]
\]

%dsp-level:0
\[
\implies E_{\mathrm{int}}(R)\qquad = 2j,
\]

%dsp-level:0
\[
\quad j\in\{0,1,2,3\};
\]

%dsp-level:0
\[
\boxed{\ P(E{=}2j)=\binom3j p^{j}(1-p)^{3-j}\ }:\qquad\bigl((1{-}p)^3,\ \ 3p(1{-}p)^2,\ \ 3p^2(1{-}p),\ \ p^3\bigr)\ \ \text{at}\ \ E=(0,2,4,6)
\]

%dsp-level:0
\[
\text{normalization:}\ \sum_{j}\binom3jp^j(1-p)^{3-j}=(p+(1-p))^3=1
\]

%dsp-level:0
\[
\qquad [\text{machine: exact in}\ \mathbb Q[p];\ \text{the census at six rational thresholds}];
\]

%dsp-level:0
\[
\textbf{the unconditioned moments:}
\]

%dsp-level:0
\[
\qquad \langle E\rangle=6p=6w
\]

%dsp-level:0
\[
\qquad [\mathrm{BM.89}(4)\ \text{re-derived};\ \text{machine}],
\]

%dsp-level:0
\[
\qquad \langle E^2\rangle=12p+24p^2,
\]

%dsp-level:0
\[
\qquad \operatorname{Var}(E)=12p(1-p)
\]

%dsp-level:0
\[
\textbf{the no-flicker collapse:}
\]

%dsp-level:0
\[
\qquad \varepsilon_q=0\implies p=1\implies P(E{=}6)=1\ \ \text{and the ensemble is the single crystallized configuration}
\]

%dsp-level:0
\[
\qquad [\mathrm{BM.19}(3);\ \mathrm{Def\text{-}BM.5}(5);\ \text{machine}];
\]

%dsp-level:0
\[
(5)\quad\textbf{the materialization restriction (derived,}\ \neg\textbf{declared):}
\]

%dsp-level:0
\[
\qquad \textbf{the}\ E{=}0\ \textbf{member supports no closed chain and no terminus:}
\]

%dsp-level:0
\[
|E_{\mathcal R}|=0\implies\#\{\text{closed chains}\}=0\implies\neg\exists\ \text{a fermion terminus}
\]

%dsp-level:0
\[
\qquad [\mathrm{FT.144}(4);\ \mathrm{D1.4};\ \mathrm{Def\text{-}H.4};\ \text{machine}];
\]

%dsp-level:0
\[
\textbf{every materialized closed chain carries}\ \ge2\ \textbf{directed edges:}
\]

%dsp-level:0
\[
\qquad \delta h(c)
\]

%dsp-level:0
\[
\ge 2,
\]

%dsp-level:0
\[
\ n(v)
\]

%dsp-level:0
\[
\ge 2\ [\mathrm{Def\text{-}H.4};\ \mathrm{Census\ FT.28}:\ \text{the digon is the minimal floor};\ \mathrm{FT.32};\ \mathrm{BM.13}(3);\ \mathrm{BM.70}(2)];
\]

%dsp-level:0
\[
\textbf{the unconditioned average runs to the vacuum:}
\]

%dsp-level:0
\[
\qquad P(E{=}0)=(1-p)^3\ \ \text{is strictly increasing in}\ \varepsilon_q\ \ \text{and}\ \to1\ \ \text{as}\ \varepsilon_q\to\tfrac14^{-}
\]

%dsp-level:0
\[
\qquad [\text{machine}]:
\]

%dsp-level:0
\[
\text{the all-virtual state is the unique entropy maximum}\ \rho_\circ
\]

%dsp-level:0
\[
= \tfrac14I,
\]

%dsp-level:0
\[
\ S_{\mathrm{raw}}
\]

%dsp-level:0
\[
= \ln4\ [\mathrm{A5};\ \mathrm{CT0.2};\ \mathrm{BM.35}(1)]\ \ \text{and is the fixed point of the Lyapunov contraction}\ V'
\]

%dsp-level:0
\[
= (1-\tfrac\gamma4)^2V\ [\mathrm{L1};\ \mathrm{L2};\ \mathrm{BM.51}(1)]:
\]

%dsp-level:0
\[
\boxed{\ \textbf{the composite's read is the conditional one on}\ E\ge2,\ \ \text{and the restriction is a theorem of the materialization domain,}\ \neg\textbf{a scope choice}\ }
\]

%dsp-level:0
\[
\qquad [\mathrm{D2.11};\ \text{A §0.2a}];
\]

%dsp-level:0
\[
(6)\quad\textbf{the materialized distribution and its exact moments:}
\]

%dsp-level:0
\[
\qquad D\ :=\ 3-3p+p^2\ =\ \frac{1-(1-p)^3}{p}\ =\ 1+u+u^2,
\]

%dsp-level:0
\[
\qquad u:=1-p=4\varepsilon_q
\]

%dsp-level:0
\[
\qquad [\text{machine: one polynomial identity}]:
\]

%dsp-level:0
\[
\boxed{\ P_{\mathrm{mat}}(E{=}2)=\frac{3(1-p)^2}{D},\qquad P_{\mathrm{mat}}(E{=}4)=\frac{3p(1-p)}{D},\qquad P_{\mathrm{mat}}(E{=}6)=\frac{p^2}{D}\ }
\]

%dsp-level:0
\[
\qquad \text{and}\quad\frac{3(1-p)^2+3p(1-p)+p^2}{D}=1\ [\text{machine}];
\]

%dsp-level:0
\[
\boxed{\ \langle E\rangle_{\mathrm{mat}}=\frac6D,\qquad\operatorname{Var}_{\mathrm{mat}}(E)=\frac{12\,p(1-p)(3-2p)}{D^2},\qquad\frac{\operatorname{Var}_{\mathrm{mat}}(E)}{\langle E\rangle_{\mathrm{mat}}^2}=\frac{p(1-p)(3-2p)}{3}\ }
\]

%dsp-level:0
\[
\qquad [\text{machine: exact rational functions}];
\]

%dsp-level:0
\[
\textbf{the exact maximum of the relative width:}
\]

%dsp-level:0
\[
\qquad \frac{p(1-p)(3-2p)}3\ \le\ \frac{10+7\sqrt7}{162}\ \ \text{at}\ \ p=\frac{5-\sqrt7}{6}
\]

%dsp-level:0
\[
\qquad [\text{machine: the stationary point of}\ 6p^2-10p+3\ \text{in}\ \mathbb Q(\sqrt7)];
\]

%dsp-level:0
\[
\frac{10+7\sqrt7}{162}<\frac15
\]

%dsp-level:0
\[
\qquad [\text{machine:}\ 5(10+7\sqrt7)<162\iff35\sqrt7<112\iff5\sqrt7<16\iff175<256]
\]

%dsp-level:0
\[
\qquad \implies\boxed{\ \frac{\sigma_E}{\langle E\rangle_{\mathrm{mat}}}<\frac12\ }
\]

%dsp-level:0
\[
(7)\quad\textbf{the two crystallization weights:}
\]

%dsp-level:0
\[
\qquad w=\frac{\langle E\rangle}{E^{\mathrm{cryst}}}=p
\]

%dsp-level:0
\[
\qquad [\mathrm{BM.89}(4)\ \text{re-cited at its own locus}];
\]

%dsp-level:0
\[
\qquad \boxed{\ w_{\mathrm{mat}}\ :=\ \frac{\langle E\rangle_{\mathrm{mat}}}{E^{\mathrm{cryst}}}\ =\ \frac1D\ =\ \frac1{1+u+u^2}\ }
\]

%dsp-level:0
\[
w_{\mathrm{mat}}(0)=1\ \ \text{exactly}
\]

%dsp-level:0
\[
\qquad [\mathrm{BM.19}(3);\ \text{machine}];
\]

%dsp-level:0
\[
\qquad \boxed{\ \frac13<w_{\mathrm{mat}}\le1\ \ \text{on}\ \ \varepsilon_q\in[0,\tfrac14)\ }
\]

%dsp-level:0
\[
\qquad [\mathrm{CT0.8};\ \mathrm{D0.4};\ \text{machine:}\ D\ge1\iff(p-1)(p-2)\ge0;\quad D<3\iff p(p-3)<0];
\]

%dsp-level:0
\[
w_{\mathrm{mat}}\ \ge\ w\qquad\text{with equality}\iff\varepsilon_q=0\qquad[\text{machine:}\ w_{\mathrm{mat}}/w=1/(1-(1-p)^3)\ge1];
\]

%dsp-level:0
\[
\boxed{\ \langle E\rangle_{\mathrm{mat}}=6\,w_{\mathrm{mat}}\ \in\ (2,\,6]\ }
\]

%dsp-level:0
\[
\qquad \textbf{the top point is the crystallized count, the bottom is the digon's two}\ [\mathrm{Census\ FT.28};\ \mathrm{Def\text{-}H.4}]:
\]

%dsp-level:0
\[
\text{the lower endpoint is open and is attained only in the fully dissolved limit, where}\ \langle E\rangle_{\mathrm{mat}}\to2\ \ \text{exactly}
\]

%dsp-level:0
\[
\qquad [\text{machine:}\ 6/(3-3p+p^2)\to6/3];
\]

%dsp-level:0
\[
(8)\quad\textbf{the informational mass is a distribution:}
\]

%dsp-level:0
\[
\qquad m_{\mathrm{hid}}(R)=\frac{E_{\mathrm{int}}(R)\ln2}{\mathfrak h}
\]

%dsp-level:0
\[
\qquad [\mathrm{FT.44}(4);\ \mathrm{Census\ FT.44}(4);\ \mathrm{Def\text{-}H.9}(2)]:
\]

%dsp-level:0
\[
\textbf{the support:}
\]

%dsp-level:0
\[
\qquad m_{\mathrm{hid}}\in\Bigl\{\frac{2\ln2}{\mathfrak h},\ \frac{4\ln2}{\mathfrak h},\ \frac{6\ln2}{\mathfrak h}\Bigr\}\ \ \text{with}\ (6)\text{'s exact probabilities};
\]

%dsp-level:0
\[
\qquad \boxed{\ \langle m_{\mathrm{hid}}\rangle_{\mathrm{mat}}=\frac{6\,w_{\mathrm{mat}}\ln2}{\mathfrak h}\ },
\]

%dsp-level:0
\[
\qquad \operatorname{Var}_{\mathrm{mat}}(m_{\mathrm{hid}})=\frac{12\,p(1-p)(3-2p)\,(\ln2)^2}{D^2\,\mathfrak h^2};
\]

%dsp-level:0
\[
\textbf{the two-bit floor survives the averaging:}
\]

%dsp-level:0
\[
\qquad E\ge2\ \ \text{on the whole materialized support}\implies m_{\mathrm{hid}}(R)\ge\frac{2\ln2}{\mathfrak h}>0\ \ \forall R,
\]

%dsp-level:0
\[
\qquad \langle m_{\mathrm{hid}}\rangle_{\mathrm{mat}}>\frac{2\ln2}{\mathfrak h}
\]

%dsp-level:0
\[
\qquad [\mathrm{BM.70}(2),(3);\ \text{machine:}\ \tfrac12<\ln2<1];
\]

%dsp-level:0
\[
\textbf{the reciprocal register has two formed reads:}
\]

%dsp-level:0
\[
\qquad \bigl\langle E^{-1}\bigr\rangle_{\mathrm{mat}}=\frac{18p(1-p)^2+9p^2(1-p)+2p^3}{12\,p\,D},
\]

%dsp-level:0
\[
\qquad \frac1{\langle E\rangle_{\mathrm{mat}}}=\frac D6
\]

%dsp-level:0
\[
\qquad [\text{machine}]:
\]

%dsp-level:0
\[
\boxed{\ \langle E\rangle_{\mathrm{mat}}\bigl\langle E^{-1}\bigr\rangle_{\mathrm{mat}}\ =\ \frac{18-27p+11p^2}{2D^2}\ \ \ge\ 1\ }
\]

%dsp-level:0
\[
\qquad [\text{Cauchy--Schwarz on}\ (E,\,E^{-1});\ \text{machine: the exact rational function}]:
\]

%dsp-level:0
\[
\text{equality}\iff E\ \text{is constant on the support}\iff p=1\iff\varepsilon_q=0
\]

%dsp-level:0
\[
\qquad [\text{machine}];
\]

%dsp-level:0
\[
\qquad \frac{58}{49}\ \ \text{at}\ \ p=\tfrac12
\]

%dsp-level:0
\[
\qquad [\text{machine:}\ 290<294];
\]

%dsp-level:0
\[
\qquad 1\le\frac{18-27p+11p^2}{2D^2}<\frac65
\]

%dsp-level:0
\[
\qquad [\text{machine: the exact positivity of}\ 12p^4-72p^3+125p^2-81p+18\ \text{on}\ (0,1]];
\]

%dsp-level:0
\[
\implies\Delta x_{\mathrm B}\ \textbf{has two formed ensemble reads:}
\]

%dsp-level:0
\[
\qquad \Delta x_{\mathrm B}^{\mathrm{eff}}:=\frac{\mathfrak h}{\langle E\rangle_{\mathrm{mat}}\ln2}=D\cdot\Delta x_{\mathrm B}^{\mathrm{cryst}}
\]

%dsp-level:0
\[
\qquad \text{and}
\]

%dsp-level:0
\[
\qquad \frac{\mathfrak h}{\ln2}\bigl\langle E^{-1}\bigr\rangle_{\mathrm{mat}}\ \ \ge\ \ \Delta x_{\mathrm B}^{\mathrm{eff}}
\]

%dsp-level:0
\[
\qquad [\mathrm{Def\text{-}BM.4}(1);\ \mathrm{BM.13}(1);\ \text{machine}];
\]

%dsp-level:0
\[
\textbf{the tier evaluates the coupling and the bridge at the width of the mean; the difference of the two reads is named}\ (\mathrm{R6}),
\]

%dsp-level:0
\[
\ \neg\textbf{filled}
\]

%dsp-level:0
\[
\qquad [\mathrm{BM.98}(5);\ \text{the Jensen residue}];
\]

%dsp-level:0
\[
(9)\quad\textbf{status:}
\]

%dsp-level:0
\[
\qquad \textbf{derived:}\ (1)\text{--}(8)\ \text{on the marked pair specialization, which is}\ \neg\text{adopted as a general measure};
\]

%dsp-level:0
\[
\qquad \textbf{rescoped,}\ \neg\textbf{retracted:}\ \mathrm{Def\text{-}BM.5}(4)\text{'s ``one point'' row is discharged as an exact weight}\ (1);
\]

%dsp-level:0
\[
\mathrm{BM.89}(4)\text{'s}\ w\ \text{stands at its own locus as the unconditioned weight and}\ w_{\mathrm{mat}}\ \text{is its materialized companion},
\]

%dsp-level:0
\[
\ \neg\text{a replacement};
\]

%dsp-level:0
\[
\qquad \mathrm{BM.19}(1)\text{--}(6),\ \mathrm{BM.20}(3),\ \mathrm{BM.21},\ \mathrm{BM.86}(4)\ \text{untouched};
\]

%dsp-level:0
\[
\mathrm{BM.91}(2)\text{'s assembly stands on}\ w\ \text{and gains the materialized variant at}\ \mathrm{BM.99}(1);
\]

%dsp-level:0
\[
\qquad \neg\text{a declaration},
\]

%dsp-level:0
\[
\quad\neg\text{a P4 addition},
\]

%dsp-level:0
\[
\quad\neg\text{a new dimensionful scale}\ [\mathrm{P1}],
\]

%dsp-level:0
\[
\quad\neg\text{a measured value}\ [\mathrm{P5}],
\]

%dsp-level:0
\[
\quad\text{§Cal-BM stays empty};
\]

\textbf{Req.} A4, A5, A6, A7 (Part 1); C9, D0.1, D0.4, D0.8, D1.4, D1.8, D2.9, D2.11 (Parts 1–2); A §0.2a, A §0.3, CT0.2, CT0.8, CT1.29 (Part 1); CT2.1 (Part 2); P1, P4, P5 (Parts 1–2); L1, L2 (Part 1); Def-H.4, Def-H.9(2), Def-H.32(1), FT.28, FT.32, FT.44(4), FT.108(4), FT.144(4) (Part 3); Census FT.28, FT.35(2), FT.44(4) (row-wise); Def-BM.4(1), Def-BM.5(1), Def-BM.5(2), Def-BM.5(3), Def-BM.5(4), Def-BM.5(5), Def-BM.13(1), Def-BM.13(3), BM.13(1), BM.13(3), BM.17(1), BM.19(1), BM.19(2), BM.19(3), BM.19(5), BM.19(6), BM.20(3), BM.21(1), BM.35(1), BM.35(3), BM.51(1), BM.61(8), BM.70(2), BM.70(3), BM.86(4), BM.87(5), BM.88(7), BM.89(4), BM.91(2) (this part); exact polynomial arithmetic over $\mathbb Q[p]$, exact rational specialization at six thresholds, exact binomial censuses, the stationary point of a quadratic in $\mathbb Q(\sqrt7)$, exact rational interval arithmetic on a quartic, Cauchy–Schwarz over a positive measure [math; machine: \texttt{checks/\allowbreak{}verify\_\allowbreak{}H\_\allowbreak{}flicker\_\allowbreak{}integral\_\allowbreak{}restmass.py} §A].

\textbf{Pf.}

%dsp-level:0
\[
(1):\ \text{the two inventories enumerated from their defining clauses}\ [\mathrm{Census\ FT.35}(2);\ \mathrm{Def\text{-}BM.5}(4);\ \mathrm{BM.19}(5),(6)];
\]

%dsp-level:0
\[
\qquad 1-\tfrac{p^2}D=\tfrac{D-p^2}D=\tfrac{3-3p}D=\tfrac{3(1-p)}D\ [\text{machine: one polynomial identity}]
\]

%dsp-level:0
\[
(2):\ \mathrm{BM.87}(5)\text{'s relocation and}\ \mathrm{BM.20}(3)\text{'s averaged flux re-cited};
\]

%dsp-level:0
\[
\qquad \text{the measure's uniqueness:}\ \dim\ker Q=1\ [\mathrm{BM.19}(1);\ \mathrm{FT.108}(4)];
\]

%dsp-level:0
\[
\qquad \text{finiteness at every budget}\ [\mathrm{Def\text{-}H.32}(1)]
\]

%dsp-level:0
\[
(3):\ \mathrm{A6},\ \mathrm{A7},\ \mathrm{D1.8},\ \mathrm{P1}\ \text{re-cited verbatim (frozen)};
\]

%dsp-level:0
\[
\qquad n_{\mathcal R}=|R|\ [\mathrm{D0.8};\ \mathrm{Def\text{-}BM.5}(4)];
\]

%dsp-level:0
\[
\qquad S_{\mathfrak L}=n_{\mathcal R}\ln2\ [\mathrm{CT1.29};\ \mathrm{BM.35}(3)]
\]

%dsp-level:0
\[
(4):\ \text{three pairs, each retained with weight}\ p\ [\mathrm{BM.89}(4)\ \text{re-cited}];
\]

%dsp-level:0
\[
\qquad \textstyle\sum_j\binom3jp^j(1-p)^{3-j}
\]

%dsp-level:0
\[
= 1\ \text{by the binomial theorem over}\ \mathbb Q[p]\ [\text{machine: coefficient comparison}];
\]

%dsp-level:0
\[
\qquad \langle E\rangle=\textstyle\sum_j2j\binom3jp^j(1-p)^{3-j}=6p\ [\text{machine at six rational thresholds}];
\]

%dsp-level:0
\[
\langle E^2\rangle=\textstyle\sum_j4j^2\binom3jp^j(1-p)^{3-j}=12p+24p^2\ [\text{machine}];
\]

%dsp-level:0
\[
\qquad \operatorname{Var}=12p+24p^2-36p^2=12p(1-p)\ [\text{machine}];
\]

%dsp-level:0
\[
\qquad p=1\implies P(E{=}6)=1\ [\text{machine}]
\]

%dsp-level:0
\[
(5):\ |E_{\mathcal R}|=0\implies\neg\exists\ \text{a closed chain}\implies\neg\exists\ \text{a terminus}\ [\mathrm{FT.144}(4);\ \mathrm{D1.4};\ \mathrm{Def\text{-}H.4}\ \text{re-cited}];
\]

%dsp-level:0
\[
\qquad \delta h\ge2\ [\mathrm{Def\text{-}H.4};\ \mathrm{BM.13}(3);\ \mathrm{BM.70}(2)];
\]

%dsp-level:0
\[
\qquad (1-p)^3\ \text{strictly increasing in}\ \varepsilon_q\ \text{and}\ \to1\ [\text{machine: monotone on the rational grid}]
\]

%dsp-level:0
\[
\rho_\circ=\tfrac14I,\ S_{\mathrm{raw}}(\rho_\circ)=\ln4=\max,\ \arg\max=\{\tfrac14I\}\ [\mathrm{CT0.2};\ \mathrm{BM.35}(1)\ \text{re-cited}];
\]

%dsp-level:0
\[
\qquad V'=(1-\tfrac\gamma4)^2V\ \text{with}\ \gamma\to0\ \text{the identity-tilt corner}\ [\mathrm{L2}\ \text{re-cited}]
\]

%dsp-level:0
\[
(6):\ 1-(1-p)^3
\]

%dsp-level:0
\[
= 3p-3p^2+p^3
\]

%dsp-level:0
\[
= p(3-3p+p^2)
\]

%dsp-level:0
\[
= pD\ [\text{machine}]
\]

%dsp-level:0
\[
\implies P_{\mathrm{mat}}(E{=}2j)
\]

%dsp-level:0
\[
= \binom3jp^j(1-p)^{3-j}/(pD)\ \text{and the three values reduce to the boxed row}\ [\text{machine}];
\]

%dsp-level:0
\[
\qquad \text{the normalization:}\ 3(1-p)^2+3p(1-p)+p^2=3-3p+p^2=D\ [\text{machine}]
\]

%dsp-level:0
\[
\langle E\rangle_{\mathrm{mat}}=\tfrac{2\cdot3(1-p)^2+4\cdot3p(1-p)+6p^2}D=\tfrac{6(1-p)^2+12p(1-p)+6p^2}D=\tfrac6D\ [\text{machine}];
\]

%dsp-level:0
\[
\qquad \langle E^2\rangle_{\mathrm{mat}}=\tfrac{12(1+2p)}{D}\ [\text{machine}]
\]

%dsp-level:0
\[
\operatorname{Var}_{\mathrm{mat}}
\]

%dsp-level:0
\[
= \tfrac{12(1+2p)}D-\tfrac{36}{D^2}
\]

%dsp-level:0
\[
= \tfrac{12(1+2p)D-36}{D^2}
\]

%dsp-level:0
\[
= \tfrac{12p(1-p)(3-2p)}{D^2}\ [\text{machine:}\ 12(1+2p)D-36=36p-60p^2+24p^3=12p(3-2p)(1-p)]
\]

%dsp-level:0
\[
\tfrac{\operatorname{Var}}{\langle E\rangle^2}=\tfrac{12p(1-p)(3-2p)}{D^2}\cdot\tfrac{D^2}{36}=\tfrac{p(1-p)(3-2p)}3\ [\text{machine}];
\]

%dsp-level:0
\[
\qquad \tfrac{d}{dp}\bigl(3p-5p^2+2p^3\bigr)
\]

%dsp-level:0
\[
= 3-10p+6p^2
\]

%dsp-level:0
\[
= 0
\]

%dsp-level:0
\[
\iff p
\]

%dsp-level:0
\[
= \tfrac{5\pm\sqrt7}6,
\]

%dsp-level:0
\[
\ \ \text{one root in}\ (0,1]\ [\text{machine: the discriminant}\ 28=4\cdot7]
\]

%dsp-level:0
\[
p_\star=\tfrac{5-\sqrt7}6:
\]

%dsp-level:0
\[
\quad 1-p_\star=\tfrac{1+\sqrt7}6,
\]

%dsp-level:0
\[
\quad 3-2p_\star
\]

%dsp-level:0
\[
= \tfrac{4+\sqrt7}3\implies\tfrac{p_\star(1-p_\star)(3-2p_\star)}3
\]

%dsp-level:0
\[
= \tfrac{(5-\sqrt7)(1+\sqrt7)(4+\sqrt7)}{324}
\]

%dsp-level:0
\[
= \tfrac{20+14\sqrt7}{324}
\]

%dsp-level:0
\[
= \tfrac{10+7\sqrt7}{162}\ [\text{machine: exact in}\ \mathbb Q(\sqrt7)]
\]

%dsp-level:0
\[
5(10+7\sqrt7)<162
\]

%dsp-level:0
\[
\iff 35\sqrt7<112
\]

%dsp-level:0
\[
\iff 5\sqrt7<16
\]

%dsp-level:0
\[
\iff 175<256\ [\text{machine: squaring two positive integers}];
\]

%dsp-level:0
\[
\qquad \tfrac15<\tfrac14\implies\sigma/\langle E\rangle<\tfrac12\ [\text{machine}]
\]

%dsp-level:0
\[
(7):\ D(p{=}1)=1\implies w_{\mathrm{mat}}(0)=1\ [\text{machine}];
\]

%dsp-level:0
\[
\qquad D'=2p-3<0\ \text{on}\ (0,1]\implies D\in[1,3)\ \text{with}\ D\to3\ \text{as}\ p\to0^+\ [\text{machine}]\implies w_{\mathrm{mat}}=1/D\in(\tfrac13,1];
\]

%dsp-level:0
\[
\qquad D\ge1\iff p^2-3p+2\ge0\iff(p-1)(p-2)\ge0\ [\text{machine}]
\]

%dsp-level:0
\[
w_{\mathrm{mat}}/w=1/(pD)=1/(1-(1-p)^3)\ge1\iff(1-p)^3\ge0,\ \text{equality}\iff p=1\ [\text{machine}];
\]

%dsp-level:0
\[
\qquad \langle E\rangle_{\mathrm{mat}}=6/D\in(2,6]\ [\text{machine}];
\]

%dsp-level:0
\[
\qquad \lim_{p\to0^+}6/D=6/3=2\ [\text{machine: exact limit row}]
\]

%dsp-level:0
\[
(8):\ m_{\mathrm{hid}}
\]

%dsp-level:0
\[
= E\ln2/\mathfrak h\ \text{linear in}\ E\ [\mathrm{FT.44}(4)]\implies\text{the support and the two moments follow from}\ (6)\ [\text{machine}];
\]

%dsp-level:0
\[
\qquad E\ge2\ \text{on the materialized support}\ [(5)]\implies m_{\mathrm{hid}}\ge2\ln2/\mathfrak h>0\ [\tfrac12<\ln2<1;\ \text{machine}]
\]

%dsp-level:0
\[
\bigl\langle E^{-1}\bigr\rangle_{\mathrm{mat}}
\]

%dsp-level:0
\[
= \tfrac{\tfrac32p(1-p)^2+\tfrac34p^2(1-p)+\tfrac16p^3}{pD}
\]

%dsp-level:0
\[
= \tfrac{18p(1-p)^2+9p^2(1-p)+2p^3}{12pD}\ [\text{machine}];
\]

%dsp-level:0
\[
\qquad \langle E\rangle\langle E^{-1}\rangle=\tfrac6D\cdot\tfrac{18-27p+11p^2}{12pD}\cdot p=\tfrac{18-27p+11p^2}{2D^2}\ [\text{machine}]
\]

%dsp-level:0
\[
\langle E\rangle\langle E^{-1}\rangle\ge\langle\sqrt{E\cdot E^{-1}}\rangle^2=1\ [\text{Cauchy--Schwarz}];
\]

%dsp-level:0
\[
\qquad \text{equality}
\]

%dsp-level:0
\[
\iff E\cdot E^{-1}\ \text{proportional to a constant}
\]

%dsp-level:0
\[
\iff E\ \text{constant}
\]

%dsp-level:0
\[
\iff p
\]

%dsp-level:0
\[
= 1\ [\text{machine: the support has three points at every}\ p\in(0,1)];
\]

%dsp-level:0
\[
\qquad p=\tfrac12:\ \tfrac{18-\tfrac{27}2+\tfrac{11}4}{2\cdot\tfrac{49}{16}}=\tfrac{58}{49}\ [\text{machine}]
\]

%dsp-level:0
\[
\tfrac{18-27p+11p^2}{2D^2}<\tfrac65
\]

%dsp-level:0
\[
\iff 5(18-27p+11p^2)<12D^2
\]

%dsp-level:0
\[
\iff 12p^4-72p^3+125p^2-81p+18>0\ [\text{machine: exact rational interval arithmetic on}\ 64\ \text{subintervals of}\ [0,1]\ \text{with the derivative bound, lower endpoint}\ >0];
\]

%dsp-level:0
\[
\qquad \Delta x_{\mathrm B}^{\mathrm{eff}}=\mathfrak h/(6\ln2/D)=D\,\mathfrak h/(6\ln2)=D\,\Delta x_{\mathrm B}^{\mathrm{cryst}}\ [\mathrm{BM.13}(1);\ \text{machine}];
\]

%dsp-level:0
\[
\qquad (9):\ \text{the scope record}.\ \blacksquare
\]

\textbf{Phys.} The composite was never a snapshot, and the count that was read as one is a single point of an exact distribution. Take the flicker's own measure — the unique stationary line of the Janus generator, three pairs each retained with weight one minus four thresholds — and the internal count takes the values zero, two, four and six with the binomial weights, its unconditioned mean six times that weight and its variance twelve times the weight times its complement. The zero is not a composite: an all-virtual configuration supports no closed chain and therefore no terminus, and it is exactly the vacuum the entropy argument and the Lyapunov contraction both run toward, so the composite's read is the conditional one on at least two edges, and that restriction is a theorem of the materialization domain rather than a choice of scope. Conditional on it, the mean count is six over one plus u plus u squared, which never leaves the interval from two to six: the crystallized count at the top, attained only at the sharp threshold, and the digon's two at the bottom, attained only in the fully dissolved limit. The relative width of the distribution has an exact algebraic maximum, ten plus seven root seven over one hundred sixty-two, reached at a stationary point in the same quadratic field, and it is bounded by a fifth — so the informational mass fluctuates, but never by more than half of itself. The reciprocal register splits into two formed reads whose product is at least one by Cauchy–Schwarz and below six fifths by an exact quartic positivity, so the Born width has a width-of-the-mean and a mean-of-the-widths, and the difference is named rather than guessed. Read A took the top point for the whole distribution; the weight it ignored is three u over D, exactly, and it is positive at every threshold that lets the vertex fire at all.

\end{theorem}

\begin{theorem}[{\texorpdfstring{BM.95 — The vertex superposition: the transport fibre, its exact weights, and the two positive parts of the splitting}{BM.95 - The vertex superposition: the transport fibre, its exact weights, and the two positive parts of the splitting}}]\label{tag:BM-95}
\textbf{Thm.}

%dsp-level:0
\[
(1)\quad\textbf{the fibre and its census (re-derived):}
\]

%dsp-level:0
\[
\qquad q_{\mathrm{core}}(c)=\tfrac12\bigl(\sigma_{s(c)}-\sigma_{t(c)}\bigr)\in\{-1,0,+1\}
\]

%dsp-level:0
\[
\qquad [\mathrm{D3.4};\ \mathrm{EM.9};\ \mathrm{FT.40}]:
\]

%dsp-level:0
\[
m_q=(2,\,1,\,1)\ \ \text{at}\ \ q=(0,+1,-1)
\]

%dsp-level:0
\[
\qquad [\mathrm{FT.106}(3\mathrm a);\ \mathrm{Def\text{-}BM.13}(1):\ \text{the orbit size of the constituent's transport register under}\ \omega];
\]

%dsp-level:0
\[
W(Q)=\bigl[z^{Q}\bigr]\bigl(2+z+z^{-1}\bigr)^{k},
\]

%dsp-level:0
\[
\qquad k=3:\quad W=(1,\,6,\,15,\,\mathbf{20},\,15,\,6,\,1),
\]

%dsp-level:0
\[
\qquad \sum_QW(Q)=4^3=64=2^6
\]

%dsp-level:0
\[
\qquad [\mathrm{Def\text{-}BM.13}(3);\ \mathrm{Def\text{-}BM.8}(6);\ \text{machine}];
\]

%dsp-level:0
\[
\#\{(q_1,q_2,q_3):\ \textstyle\sum_iq_i=k\}=(1,3,6,\,\mathbf 7,\,6,3,1)
\]

%dsp-level:0
\[
\qquad [\mathrm{BM.28}(2);\ \text{machine: the exhaustive}\ 4^3\ \text{end-sign census and the convolution agree exactly}];
\]

%dsp-level:0
\[
(2)\quad\textbf{the derived superposition weights:}
\]

%dsp-level:0
\[
\qquad \textbf{the constituent at a vertex is a quantum superposition over its transport fibre, with weights}
\]

%dsp-level:0
\[
\boxed{\ P\bigl((q_1,q_2,q_3)\,\big|\,Q\bigr)\ =\ \frac{\prod_{i=1}^{3}m_{q_i}}{W(Q)}\ }
\]

%dsp-level:0
\[
\qquad [\text{the normalized fibre weight};\ \text{machine: exact rationals, summing to}\ 1\ \text{class by class}]:
\]

%dsp-level:0
\[
Q=0:
\]

%dsp-level:0
\[
\qquad P\bigl((0,0,0)\bigr)=\frac8{20}=\frac25,
\]

%dsp-level:0
\[
\qquad P\bigl(\text{each of the six permutations of}\ (+1,-1,0)\bigr)=\frac2{20}=\frac1{10}
\]

%dsp-level:0
\[
\qquad [\text{machine}];
\]

%dsp-level:0
\[
Q=+1:
\]

%dsp-level:0
\[
\qquad P\bigl(\text{each of the three permutations of}\ (+1,0,0)\bigr)=\frac4{15},
\]

%dsp-level:0
\[
\qquad P\bigl(\text{each of the three permutations of}\ (+1,+1,-1)\bigr)=\frac1{15}
\]

%dsp-level:0
\[
\qquad [\text{machine}];
\]

%dsp-level:0
\[
Q=-1:\qquad\text{the identical census}\ —\ \text{the}\ \omega\text{-image fibre}\qquad[\mathrm{D2.6};\ \mathrm{BM.25}(4);\ \mathrm{BM.28}(5);\ \text{machine}];
\]

%dsp-level:0
\[
\textbf{the weights carry no datum:}
\]

%dsp-level:0
\[
\qquad \text{the free-symbol set of}\ W(Q)\ \text{is}\ \{z\}\ \text{with coefficients fixed by}\ (m_q,k),
\]

%dsp-level:0
\[
\ \text{and the coupling alphabet}\ \{g,g',e,\varepsilon_{\mathrm{EM}},\hat\mu,\hat\mu',\ell_c\}\ \text{is disjoint from it}
\]

%dsp-level:0
\[
[\mathrm{BM.61}(5)\ \text{re-cited};\ \text{machine: the symbol census}]\implies\textbf{the superposition weights are}\ \neg\textbf{couplings,}\ \neg\textbf{rates,}\ \neg\textbf{P4 data};
\]

%dsp-level:0
\[
(3)\quad\textbf{the entropy decomposition:}
\]

%dsp-level:0
\[
\qquad H(Q)\ :=\ -\sum_{\text{fibre}}P\ln P,
\]

%dsp-level:0
\[
\qquad \bigl\langle\ln\textstyle\prod_im_{q_i}\bigr\rangle_Q\ :=\ \sum_{\text{fibre}}P\ln\textstyle\prod_im_{q_i}:
\]

%dsp-level:0
\[
\boxed{\ \Sigma_{\mathrm{cl}}(Q)\ =\ \ln W(Q)\ =\ H(Q)\ +\ \bigl\langle\ln\textstyle\prod_im_{q_i}\bigr\rangle_Q\ }
\]

%dsp-level:0
\[
\qquad [\mathrm{Def\text{-}BM.13}(4);\ \text{machine: exact in the prime-log ring}]:
\]

%dsp-level:0
\[
H(0)=\ln5+\frac15\ln2,\qquad\bigl\langle\ln\textstyle\prod_im\bigr\rangle_{0}=\frac95\ln2,\qquad\ln20=2\ln2+\ln5\qquad[\text{machine}];
\]

%dsp-level:0
\[
H(\pm1)=\ln3+\ln5-\frac85\ln2,\qquad\bigl\langle\ln\textstyle\prod_im\bigr\rangle_{+1}=\frac85\ln2,\qquad\ln15=\ln3+\ln5\qquad[\text{machine}];
\]

%dsp-level:0
\[
(4)\quad\textbf{the two positive parts of the splitting:}
\]

%dsp-level:0
\[
\qquad \Sigma_{\mathrm B}=\Sigma_{\mathrm{cl}}(0)-\Sigma_{\mathrm{cl}}(+1)=\ln\tfrac43=2\ln2-\ln3
\]

%dsp-level:0
\[
\qquad [\mathrm{Def\text{-}BM.13}(4);\ \mathrm{BM.59}(4)]:
\]

%dsp-level:0
\[
\boxed{\ \ln\frac43\ =\ \underbrace{\Bigl(\frac95\ln2-\ln3\Bigr)}_{H(0)-H(+1)}\ +\ \underbrace{\frac15\ln2}_{\langle\ln\prod_im\rangle_0-\langle\ln\prod_im\rangle_{+1}}\ }
\]

%dsp-level:0
\[
\qquad [\text{machine: the prime-log coefficients agree}];
\]

%dsp-level:0
\[
\frac95\ln2-\ln3>0\qquad[\text{machine:}\ 9\ln2>5\ln3\iff2^9=512>243=3^5];\qquad\frac15\ln2>0\qquad[\text{machine}];
\]

%dsp-level:0
\[
\textbf{the transport splitting is exactly two positive parts}\ —\ \textbf{a superposition-entropy part and a multiplicity part};
\]

%dsp-level:0
\[
\qquad \textbf{the snapshot read carries neither:}
\]

%dsp-level:0
\[
\text{a fixed triple at the vertex books}\ H=0\ \ \text{and loses}\ \ \frac95\ln2-\ln3\ \ \text{of the splitting}
\]

%dsp-level:0
\[
\qquad [\text{machine: the exact difference}];
\]

%dsp-level:0
\[
(5)\quad\textbf{the factorization (two independent exact averages):}
\]

%dsp-level:0
\[
\qquad \text{the fibre weights are class-fixed boundary data}\ [\mathrm{Def\text{-}BM.5}(1);\ \mathrm{Def\text{-}BM.13}(6)],
\]

%dsp-level:0
\[
\ \text{and the two registers' alphabets are disjoint}:
\]

%dsp-level:0
\[
E_{\mathrm{int}}=|R|\ \ \text{has no endpoint summand}
\]

%dsp-level:0
\[
\qquad [\mathrm{Def\text{-}BM.13}(6);\ \mathrm{BM.62}(11);\ \mathrm{BM.89}(8);\ \text{machine: the token census}]\implies\boxed{\ P(R,\ \text{fibre member})\ =\ P(R)\cdot P(\text{fibre member}\,|\,Q)\ }
\]

%dsp-level:0
\[
\textbf{the baryon's full ledger read is the composition of the flicker integral over}\ |R|\ (\mathrm{BM.94}(2))\ \textbf{and the superposition average over the fibre}\ (2);
\]

%dsp-level:0
\[
\neg\text{a cross-register equation is formed}\ —\ \text{the two registers are distinct alphabets and are}\ \neg\text{equated}
\]

%dsp-level:0
\[
\qquad [\mathrm{BM.37}(2);\ \mathrm{BM.83}(3)];
\]

%dsp-level:0
\[
(6)\quad\textbf{the fixed-state read refuted with an exact bound:}
\]

%dsp-level:0
\[
\qquad \max_{\text{fibre}}P(\cdot\,|\,0)=\frac25<\frac12
\]

%dsp-level:0
\[
\qquad [\text{machine:}\ 16<20];
\]

%dsp-level:0
\[
\qquad \max_{\text{fibre}}P(\cdot\,|\,+1)=\frac4{15}<\frac13
\]

%dsp-level:0
\[
\qquad [\text{machine:}\ 12<15]:
\]

%dsp-level:0
\[
\boxed{\ \neg\exists\ \text{a single vertex configuration carrying more than}\ \tfrac25\ \text{of the neutral class's transport weight}\ }
\]

%dsp-level:0
\[
\textbf{the refutation is an exact rational bound,}\ \neg\textbf{an argument about size};
\]

%dsp-level:0
\[
\qquad \text{both classes are strict mixtures at every admissible threshold}\ [(2);\ \mathrm{BM.19}(6)];
\]

%dsp-level:0
\[
(7)\quad\textbf{status:}
\]

%dsp-level:0
\[
\qquad \textbf{derived:}\ (1)\text{--}(6);
\]

%dsp-level:0
\[
\qquad \neg\text{a declaration},
\]

%dsp-level:0
\[
\quad\neg\text{a P4 addition},
\]

%dsp-level:0
\[
\quad\neg\text{a new dimensionful scale}\ [\mathrm{P1}],
\]

%dsp-level:0
\[
\quad\neg\text{a measured value}\ [\mathrm{P5}],
\]

%dsp-level:0
\[
\quad\text{§Cal-BM stays empty};
\]

%dsp-level:0
\[
\mathrm{Def\text{-}BM.13}(1)\text{--}(6),\ \mathrm{BM.28}(2),(4),(5),\ \mathrm{BM.59}(4),\ \mathrm{BM.61}(5)\ \text{stand unmodified};\qquad
\]

\textbf{Req.} D2.6, D3.4 (Parts 1–2); EM.9, CT0.9 (Part 2); P1, P4, P5 (Parts 1–2); Def-H.32(1), FT.40, FT.106(3) (Part 3); Census FT.35(2) (row-wise); Def-BM.5(1), Def-BM.8(6), Def-BM.13(1), Def-BM.13(3), Def-BM.13(4), Def-BM.13(6), BM.19(6), BM.25(4), BM.28(2), BM.28(5), BM.37(2), BM.59(4), BM.61(5), BM.62(11), BM.83(3), BM.89(8), BM.94(2) (this part); exhaustive ternary-fibre enumeration, exact rational weight normalization, prime-log ring arithmetic with integer sign certificates, symbol censuses [math; machine: \texttt{checks/\allowbreak{}verify\_\allowbreak{}H\_\allowbreak{}flicker\_\allowbreak{}integral\_\allowbreak{}restmass.py} §B].

\textbf{Pf.}

%dsp-level:0
\[
(1):\ m_q=(2,1,1)\ \text{and the end-pair census re-cited}\ [\mathrm{Def\text{-}BM.13}(1)];
\]

%dsp-level:0
\[
\qquad W(Q)\ \text{by two agreeing routes: the exhaustive}\ 4^3\ \text{census and the coefficient of}\ \bigl(2+z+z^{-1}\bigr)^3\ [\text{machine}];
\]

%dsp-level:0
\[
\qquad \textstyle\sum_QW(Q)=(2+1+1)^3=64\ [\text{machine}]
\]

%dsp-level:0
\[
(2):\ P
\]

%dsp-level:0
\[
= \prod_im_{q_i}/W(Q)\ \text{and}\ \textstyle\sum_{\text{fibre}}\prod_im_{q_i}
\]

%dsp-level:0
\[
= W(Q)\ \text{by}\ (1)\text{'s factorization identity}\ [\text{machine: class by class}];
\]

%dsp-level:0
\[
\qquad \text{the}\ Q{=}0\ \text{fibre: one triple of weight}\ 2^3=8\ \text{and six of weight}\ 2\cdot1\cdot1=2,\ \text{total}\ 20\ [\text{machine}];
\]

%dsp-level:0
\[
\qquad \text{the}\ Q{=}{+}1\ \text{fibre: three of weight}\ 1\cdot2\cdot2=4\ \text{and three of weight}\ 1,\ \text{total}\ 15\ [\text{machine}]
\]

%dsp-level:0
\[
\text{the symbol census:}\ W(Q)\in\mathbb Z[z,z^{-1}]\ \text{with coefficients fixed by}\ (m_q,k),
\]

%dsp-level:0
\[
\ \text{and}\ \{g,g',e,\varepsilon_{\mathrm{EM}},\hat\mu,\hat\mu',\ell_c\}\cap\{z\}
\]

%dsp-level:0
\[
= \varnothing\ [\mathrm{BM.61}(5)\ \text{re-cited};\ \text{machine}]
\]

%dsp-level:0
\[
(3):\ H(0)
\]

%dsp-level:0
\[
= -\tfrac25\ln\tfrac25-6\cdot\tfrac1{10}\ln\tfrac1{10}
\]

%dsp-level:0
\[
= \tfrac25(\ln5-\ln2)+\tfrac35(\ln2+\ln5)
\]

%dsp-level:0
\[
= \ln5+\tfrac15\ln2\ [\text{machine: prime-log coefficients}];
\]

%dsp-level:0
\[
\qquad \langle\ln\prod_im\rangle_0=\tfrac{8\ln8+6\cdot2\ln2}{20}=\tfrac{36\ln2}{20}=\tfrac95\ln2\ [\text{machine}]
\]

%dsp-level:0
\[
H(+1)=-3\cdot\tfrac4{15}\ln\tfrac4{15}-3\cdot\tfrac1{15}\ln\tfrac1{15}=\ln15-\tfrac85\ln2\ [\text{machine}];
\]

%dsp-level:0
\[
\qquad \langle\ln\prod_im\rangle_{+1}=\tfrac{3\cdot4\ln4+3\cdot1\ln1}{15}=\tfrac{24\ln2}{15}=\tfrac85\ln2\ [\text{machine}];
\]

%dsp-level:0
\[
\qquad H+\langle\ln\prod_im\rangle=\ln W\ \text{in both classes}\ [\text{machine}]
\]

%dsp-level:0
\[
(4):\ [\ln5+\tfrac15\ln2]-[\ln3+\ln5-\tfrac85\ln2]
\]

%dsp-level:0
\[
= \tfrac95\ln2-\ln3\ \text{and}\ \tfrac95\ln2-\tfrac85\ln2
\]

%dsp-level:0
\[
= \tfrac15\ln2\ [\text{machine}];
\]

%dsp-level:0
\[
\qquad \text{the sum:}\ \tfrac95\ln2-\ln3+\tfrac15\ln2=2\ln2-\ln3=\ln\tfrac43\ [\text{machine}];
\]

%dsp-level:0
\[
\qquad 2^9=512>243=3^5\ [\text{machine: an integer power comparison}]
\]

%dsp-level:0
\[
(5):\ \text{the boundary data fix the charge class}\ [\mathrm{Def\text{-}BM.5}(1)]\implies\text{the fibre weights are class-constant on the ensemble};
\]

%dsp-level:0
\[
\qquad E_{\mathrm{int}}\ \text{is constant on every end-sign fibre}\ [\mathrm{Def\text{-}BM.13}(6)]\implies\text{the two registrations are independent and their probabilities multiply}\ [\text{machine}]
\]

%dsp-level:0
\[
(6):\ \tfrac25<\tfrac12\iff16<20;
\]

%dsp-level:0
\[
\qquad \tfrac4{15}<\tfrac13\iff12<15\ [\text{machine: integers}];
\]

%dsp-level:0
\[
\qquad (7):\ \text{the standing-rule scans}\ [\text{machine}].\ \blacksquare
\]

\textbf{Phys.} Ask which state a constituent occupies at a vertex and the ledger answers with a list, not a value. A neutral-core constituent carries two transport realizations and a charged one carries a single realization, so the class fibre is a weighted set: at charge zero one triple of weight eight against six of weight two, summing to twenty; at charge one three of weight four against three of weight one, summing to fifteen. Normalize and the superposition is exact — two fifths on the all-neutral triple, a tenth on each of the six mixed ones — and the numbers are rationals in an alphabet that carries no coupling at all, which is the coupling-blindness theorem wearing a different hat. The class entropy then splits cleanly into the Shannon entropy of that superposition plus the mean logarithm of the multiplicities, and the baryon's transport splitting, the log of four thirds, is the sum of two strictly positive parts: nine fifths of a bit less the log of three, from the superposition, and one fifth of a bit, from the multiplicities. A read that fixes one triple at the vertex books no entropy and forfeits the larger part of the splitting, so the fixed-state read is not merely incomplete, it is quantitatively short by an exact prime-log element. The two registers stay independent — the fibre weights are boundary data and the edge count has no endpoint summand — so the baryon's ledger read is the composition of two exact averages, and no equation between the registers is formed. The sharpest row is the bound: no single vertex configuration carries more than two fifths of the neutral class's transport weight, certified by sixteen against twenty.

\end{theorem}

\begin{theorem}[{\texorpdfstring{BM.96 — The rung-one $\nu$ line: the seal-derived rest mass, the single-ray occupancy, and the four bars against a seesaw-type mechanism}{BM.96 - The rung-one line: the seal-derived rest mass, the single-ray occupancy, and the four bars against a seesaw-type mechanism}}]\label{tag:BM-96}
\textbf{Thm.}

%dsp-level:0
\[
(1)\quad\textbf{the line's structure:}
\]

%dsp-level:0
\[
\qquad \text{a rung-one bound chain}\ [\mathrm{D1.5}]\ \text{whose terminus-active aspects lie in}\ \mathcal A(1)
\]

%dsp-level:0
\[
= \{J,\,G,\,W,\,H,\,D^+\}\ \ [\mathrm{CT0.13}]:
\]

%dsp-level:0
\[
\text{EM unactivated}\ (f_v=0),
\]

%dsp-level:0
\[
\qquad \text{the colour block vanishing},
\]

%dsp-level:0
\[
\qquad \delta h(c)\ge2,\quad n(v)\ge2,
\]

%dsp-level:0
\[
\qquad \Sigma_{\min}\ \text{the minimal instance}
\]

%dsp-level:0
\[
\qquad [\mathrm{FT.142}(1);\ \mathrm{D3.3};\ \mathrm{CT3.11}(1);\ \mathrm{FT.19};\ \mathrm{Def\text{-}H.4};\ \mathrm{Census\ FT.28}];
\]

%dsp-level:0
\[
\dim\mathcal K_1=16=C_{\mathrm{cap}}
\]

%dsp-level:0
\[
\qquad [\mathrm{FT.142}(2),(3);\ \mathrm{CT0.14};\ \mathrm{CT0.12}:\ \text{the unique derived candidate, no interface declaration}];
\]

%dsp-level:0
\[
\boxed{\ \delta_\nu=\frac{\sum_i[J_i,K_i]\big|_\Sigma}{\dim\mathcal K_1}=\frac2{16}=\frac18\ }
\]

%dsp-level:0
\[
\qquad [\mathrm{FT.142}(4),(5):\ \text{the undressed seal coefficient over the derived mask dimension}]
\]

%dsp-level:0
\[
\qquad \textbf{derived,}\ \neg\textbf{declared};
\]

%dsp-level:0
\[
(2)\quad\textbf{the rest mass from the rung-one seal structure:}
\]

%dsp-level:0
\[
\qquad m_\nu^{\mathrm{rest}\,2}=\hat\mu_\nu^2\bigl(\ell_c^{(\nu)}\bigr)^2\,G\bigl(c(i_\nu)\bigr),
\]

%dsp-level:0
\[
\qquad G\bigl(c(i)\bigr)=A_0-B_0\cos\Bigl(\delta_\nu-\frac{2\pi i}3\Bigr)
\]

%dsp-level:0
\[
\qquad [\mathrm{Def\text{-}H.37}(1);\ \mathrm{Def\text{-}BM.11}(5);\ \mathrm{BM.39}(1)]:
\]

%dsp-level:0
\[
A_0=\frac{3g^2}8+\frac{g'^2}{16},\qquad B_0=\frac{gg'}4\qquad[\mathrm{BM.39}(1);\ \text{machine}];\qquad\textbf{at}\ \delta_\nu=\tfrac18,\ \ c:=\cos\tfrac18,\ \ s_\nu:=\sin\tfrac18:
\]

%dsp-level:0
\[
\boxed{\ G_0=A_0-B_0\,c,\qquad G_1=A_0+\frac{B_0}2\bigl(c-\sqrt3\,s_\nu\bigr),\qquad G_2=A_0+\frac{B_0}2\bigl(c+\sqrt3\,s_\nu\bigr)\ }
\]

%dsp-level:0
\[
\qquad [\text{machine: exact in}\ \mathbb Q[\sqrt2,\sqrt3][c,s_\nu]/(c^2{+}s_\nu^2{-}1)];
\]

%dsp-level:0
\[
\textbf{the exact order at}\ (g,g')>0\ [\mathrm{EW.33}]:\qquad\boxed{\ G_2\ >\ G_1\ >\ A_0\ >\ G_0\ }
\]

%dsp-level:0
\[
G_1>A_0\iff c>\sqrt3\,s_\nu\iff c^2>\tfrac34\iff\cos\tfrac14=2c^2-1>\tfrac12
\]

%dsp-level:0
\[
\qquad [\text{machine:}\ \cos x>1-\tfrac{x^2}2=\tfrac{31}{32}>\tfrac12\ \text{at}\ x=\tfrac14,\ \text{the exact alternating-series sign row}];
\]

%dsp-level:0
\[
c>0\ \ \text{on}\ \ c>\tfrac{63}{64},
\]

%dsp-level:0
\[
\qquad s_\nu>0\ \ \text{on}\ \ s_\nu>\tfrac18-\tfrac1{3072}
\]

%dsp-level:0
\[
\qquad [\text{machine: exact Taylor enclosures with rational remainder bounds}];
\]

%dsp-level:0
\[
\textbf{the affine form on the spectral triplet:}
\]

%dsp-level:0
\[
\qquad G\bigl(c(i)\bigr)=A_0-\frac{B_0}{\sqrt2}\bigl(r_i^{(\nu)}-1\bigr),
\]

%dsp-level:0
\[
\qquad r_i^{(\nu)}=1+\sqrt2\cos\Bigl(\frac18+\frac{2\pi i}3\Bigr)
\]

%dsp-level:0
\[
\qquad [\mathrm{Def\text{-}BM.11}(5);\ \mathrm{FT.143}(1);\ \text{machine}]:
\]

%dsp-level:0
\[
\bigl(r_1^{(\nu)}\bigr)^2<\bigl(r_2^{(\nu)}\bigr)^2<\bigl(r_0^{(\nu)}\bigr)^2
\]

%dsp-level:0
\[
\qquad [\mathrm{FT.143}(3):\ \text{enclosure-proven, frozen}]\implies\text{the transport read is strictly decreasing in}\ r;
\]

%dsp-level:0
\[
\textbf{the value set is orbit-label independent:}
\]

%dsp-level:0
\[
\qquad \bigl\{G\bigl(c(i)\bigr)\bigr\}_i=\Bigl\{A_0-B_0\,c,\ \ A_0+\tfrac{B_0}2\bigl(c-\sqrt3\,s_\nu\bigr),\ \ A_0+\tfrac{B_0}2\bigl(c+\sqrt3\,s_\nu\bigr)\Bigr\}
\]

%dsp-level:0
\[
\qquad [\mathrm{Def\text{-}BM.11}(5)\ \text{orbit-label row};\ \text{machine: the two routes give the same set}];
\]

%dsp-level:0
\[
\boxed{\ \textbf{the transport order is the reverse of the spectral order:}\ \ \mathrm{BM.39}(5)\text{'s inversion instantiated at rung one}\ }
\]

%dsp-level:0
\[
(3)\quad\textbf{positivity and the exact floor:}
\]

%dsp-level:0
\[
\qquad A_0-B_0=\frac38\Bigl(g-\frac{g'}3\Bigr)^2+\frac{g'^2}{48}>0,
\]

%dsp-level:0
\[
\qquad A_0+B_0=\frac38\Bigl(g+\frac{g'}3\Bigr)^2+\frac{g'^2}{48}>0
\]

%dsp-level:0
\[
\qquad [\mathrm{BM.39}(1);\ \mathrm{BM.77}(4);\ \text{machine: the SOS expansions}]:
\]

%dsp-level:0
\[
\sqrt{A_0-B_0}\ \le\ \sqrt{G\bigl(c(i)\bigr)}\ \le\ \sqrt{A_0+B_0}\qquad\forall\,\text{line}\ i,\ \ \forall\,\text{channel},\ \ \forall\,(g,g')\ne(0,0)\qquad[|\cos|\le1;\ \text{machine}];
\]

%dsp-level:0
\[
\text{the kernel line is unreachable:}\ \cos\varphi=0\wedge\cos(\varphi+\tfrac{2\pi}3)=0\ \text{has no solution}
\]

%dsp-level:0
\[
\qquad [\mathrm{BM.39}(1);\ \text{machine}]\implies\boxed{\ m_\nu^{\mathrm{rest}}>0\ \ \text{strictly, and}\ \neg\text{a massless}\ \nu\ \text{line is formed on the transport species}\ }
\]

%dsp-level:0
\[
\textbf{the absolute value is}\ \neg\textbf{derived:}
\]

%dsp-level:0
\[
\qquad \bigl(\hat\mu_\nu,\ \ell_c^{(\nu)}\bigr)\ \text{undeclared}\ [\mathrm{P4};\ \mathrm{Cal\text{-}F}];
\]

%dsp-level:0
\[
\qquad \rho_\nu:=\mathfrak h\,\hat\mu_\nu\,\ell_c^{(\nu)}\in\mathfrak q^0\ \ \text{is a restatement,}\ \neg\text{a new member}\ [\mathrm{BM.77}(2);\ \mathrm{CT2.1}];
\]

%dsp-level:0
\[
\textbf{the internal ratios are derived:}
\]

%dsp-level:0
\[
\qquad \frac{m_i^2}{m_j^2}=\frac{G\bigl(c(i)\bigr)}{G\bigl(c(j)\bigr)}\ \ \text{exact symbolic in}\ (A_0,B_0)
\]

%dsp-level:0
\[
\qquad [\mathrm{FT.143}(3);\ \mathrm{BM.80}(6);\ \text{machine}];
\]

%dsp-level:0
\[
\textbf{the inertial companion:}
\]

%dsp-level:0
\[
\qquad m_i=\mu_\nu\bigl(r_i^{(\nu)}\bigr)^2,\quad\mu_\nu\ \text{undeclared}\ [\mathrm{P4}],
\]

%dsp-level:0
\[
\qquad Q:=\frac{\sum_im_i}{\bigl(\sum_i\sqrt{m_i}\bigr)^2}=\frac69=\frac23\ \ \text{exactly}
\]

%dsp-level:0
\[
\qquad [\mathrm{FT.143}(4);\ \text{machine}];
\]

%dsp-level:0
\[
\textbf{the informational companion:}\qquad m_\nu\,\mathfrak h\ \ge\ 2\ln2\qquad[\mathrm{BM.70}(2),(3);\ \mathrm{BM.13}(3):\ \text{the two-bit floor}];
\]

%dsp-level:0
\[
(4)\quad\textbf{the single-ray occupancy:}
\]

%dsp-level:0
\[
\qquad \text{the materialized locus at rung one is single-arrow}\ —\ \mathrm{D1.2}\text{'s domains force strict alternation and the prefix law, and a dissolution-first chain is inadmissible at its first event}\ [\mathrm{FT.144}(2);\ \mathrm{D2.11}];
\]

%dsp-level:0
\[
\text{under A-BM.7's issued coupling the arrow locus occupies}\ \operatorname{im}P_L,
\]

%dsp-level:0
\[
\ \ \text{the}\ \hat\chi
\]

%dsp-level:0
\[
= +1\ \text{ray of}\ \mathrm{Def\text{-}H.33}\text{'s boxed sign}:
\]

%dsp-level:0
\[
\qquad \boxed{\ \text{the occupancy is exactly}\ (1,\,0)\ }
\]

%dsp-level:0
\[
\qquad [\mathrm{FT.144}(5);\ \mathrm{BM.73}(2),(3);\ \mathrm{FT.122}(4)];
\]

%dsp-level:0
\[
\textbf{the complementary ray's content is the vacuum sector:}
\]

%dsp-level:0
\[
\qquad n_{\mathcal R}=0\implies\rho_{pq}=\rho_\circ=\tfrac14I,\quad S_{\mathrm{raw}}=\ln4=\max,\quad\text{unregistered},\quad\neg\exists\ \text{an accessible pole at any rung}
\]

%dsp-level:0
\[
[\mathrm{FT.144}(4);\ \mathrm{A5};\ \mathrm{CT0.2};\ \mathrm{CT0.1};\ \mathrm{CT4.2}:\ \mathcal P_MF\mathcal P_M=0\implies\neg\exists\ \text{an accessible pole};\ \mathrm{CT2.5}:\ \text{poles}\iff1\in\operatorname{spec}\mathbf F_M];
\]

%dsp-level:0
\[
\textbf{the obstruction is a transition-domain restriction prior to every read}\ —\ \neg\text{a mass statement},
\]

%dsp-level:0
\[
\ \neg\text{a suppression},
\]

%dsp-level:0
\[
\ \neg\text{a decoupling}
\]

%dsp-level:0
\[
\qquad [\mathrm{FT.144}(3):\ \neg\mathfrak h,\ \neg\mu,\ \neg h_0\ \text{in the premises}]:
\]

%dsp-level:0
\[
\boxed{\ \neg\exists\ \text{a materializable right-ray}\ \nu\ \text{line in the matter orientation, in any valid universe}\ }
\]

%dsp-level:0
\[
\qquad [\mathrm{D2.11};\ \text{A-BM.7}]
\]

%dsp-level:0
\[
\textbf{the naming discipline:}
\]

%dsp-level:0
\[
\qquad \text{``left'' and ``right'' are the ray names of}\ \mathrm{Def\text{-}H.33}\ \text{as fixed by A-BM.7};
\]

%dsp-level:0
\[
\qquad \text{the handedness vocabulary is external and is}\ \neg\text{identified with a ray}\ [\text{A-BM.7's bars};\ \mathrm{P5}];
\]

%dsp-level:0
\[
(5)\quad\textbf{the image row:}
\]

%dsp-level:0
\[
\qquad \text{the vertex's neutral member is the}\ \nu\text{-channel line in the image orientation}\ [\text{A-BM.5};\ \mathrm{BM.41}(1);\ \mathrm{BM.45}];
\]

%dsp-level:0
\[
\text{the image is the}\ \hat\omega\text{-image of the same single-arrow locus:}
\]

%dsp-level:0
\[
\qquad \hat\omega P_L\hat\omega^{-1}=P_R,
\]

%dsp-level:0
\[
\qquad \hat\omega|L\rangle=-i|R\rangle,
\]

%dsp-level:0
\[
\qquad \hat\omega^2=\mathrm{id}
\]

%dsp-level:0
\[
\qquad [\mathrm{D2.6};\ \mathrm{FT.122}(4);\ \mathrm{BM.73}(4);\ \text{machine over}\ \mathbb Q(i)];
\]

%dsp-level:0
\[
\text{and it sits in the}\ \omega\text{-image representation}\ \{0,-1\}\qquad[\mathrm{EW.36};\ \text{A-BM.7}(\mathrm b);\ \mathrm{BM.25};\ \mathrm{BM.27}];
\]

%dsp-level:0
\[
\mathrm{CPT}:\ (q,\ \text{orientation})
\]

%dsp-level:0
\[
\mapsto (-q,\ \text{flipped})\ \ \text{has}\ \textbf{no fixed sector},
\]

%dsp-level:0
\[
\ \ \text{its orbits are three pairs of size two},
\]

%dsp-level:0
\[
\ \ q
\]

%dsp-level:0
\[
\mapsto -q\ \text{fixes}\ q\qquad = 0\ \text{alone}
\]

%dsp-level:0
\[
\qquad [\mathrm{BM.25}(4);\ \text{machine}]:
\]

%dsp-level:0
\[
\boxed{\ \textbf{the image is one orbit of size two and carries no independent datum}\ —\ \neg\text{a second state},\ \neg\text{a right-ray matter line},\ \neg\text{a name for an unmaterializable object}\ }
\]

%dsp-level:0
\[
\qquad [\mathrm{BM.27}(6)\ \text{pattern}];
\]

%dsp-level:0
\[
(6)\quad\textbf{the no-seesaw theorem (four independent bars, each re-cited at its own locus):}
\]

%dsp-level:0
\[
(\mathrm S1)\quad\textbf{no partner state:}\ \neg\exists\ \text{a right-ray matter-orientation state to pair with}\ —\ \text{the complement's occupancy is exactly zero and its content supports no accessible pole at any rung}
\]

%dsp-level:0
\[
\qquad [(4);\ \mathrm{FT.144}(4),(5)];
\]

%dsp-level:0
\[
(\mathrm S2)\quad\textbf{no self-pairing:}\ \neg\exists\ \text{a fixed CPT sector, so a Majorana-type identification}\ C
\]

%dsp-level:0
\[
= \mathrm{CPT}\cdot C\ \text{is}\ \neg\text{formed}
\]

%dsp-level:0
\[
\qquad [(5);\ \mathrm{BM.25}(4);\ \mathrm{BM.27}(6)];
\]

%dsp-level:0
\[
(\mathrm S3)\quad\textbf{no coupling object:}\ \neg\exists\ \text{a Yukawa-type object in the inventory}\ —\ \dim F_H
\]

%dsp-level:0
\[
= 1
\]

%dsp-level:0
\[
= \operatorname{span}\{\Sigma\},
\]

%dsp-level:0
\[
\ \neg\text{a scalar-block coupling to a chain line},
\]

%dsp-level:0
\[
\ \neg\text{a dimension-five operator}
\]

%dsp-level:0
\[
\qquad [\mathrm{CT0.12};\ \mathrm{EW.15};\ \mathrm{FT.145}(3);\ \mathrm{BM.93}(5)];
\]

%dsp-level:0
\[
(\mathrm S4)\quad\textbf{no second scale:}\ \exists!\ \text{one free dimensionful scale}\ \mathfrak h\ —\ \text{the heavy scale a seesaw-type mechanism requires is barred}
\]

%dsp-level:0
\[
\qquad [\mathrm{P1};\ \text{A §0.3};\ \mathrm{D2.9}];
\]

%dsp-level:0
\[
\boxed{\ \neg\exists\ \text{a seesaw-type mechanism formable in this corpus};\qquad\text{the}\ \nu\ \text{rest mass is positive, derived in form from the rung-one seal, and undeclared in scale}\ }
\]

%dsp-level:0
\[
\textbf{the companion row:}
\]

%dsp-level:0
\[
\qquad \text{the smallness question is a scale question}\ [\mathrm{P4}],
\]

%dsp-level:0
\[
\ \neg\text{a mechanism question}\ —\ \neg\text{a suppression is needed, so}\ \neg\text{one is invented};
\]

%dsp-level:0
\[
(7)\quad\textbf{the vacuum residence and the vertex firing (the two states and the one door):}
\]

%dsp-level:0
\[
\qquad \textbf{state (i), the vacuum-resident balanced fluctuation pair:}
\]

%dsp-level:0
\[
n_{\mathcal R}=0,
\]

%dsp-level:0
\[
\qquad \sum_n\langle\varnothing|\hat U^n|\varnothing\rangle\,w^n=\frac1{1-w^2x_+x_-},
\]

%dsp-level:0
\[
\qquad \text{every nonzero return monomial carries equal}\ x_+\text{- and}\ x_-\text{-degree}
\]

%dsp-level:0
\[
\qquad [\mathrm{FT.108}(5);\ \mathrm{FT.144}(4);\ \text{machine in}\ \mathbb Q[w,x_+,x_-]/(w^{12})]:
\]

%dsp-level:0
\[
\textbf{state (ii), the materialized rung-one seal:}
\]

%dsp-level:0
\[
\qquad |E_{\mathcal R}|\ge2,
\]

%dsp-level:0
\[
\qquad m_{\mathrm{hid}}=\frac{2\ln2}{\mathfrak h}\ \ \text{at the minimal floor}
\]

%dsp-level:0
\[
\qquad [\mathrm{FT.145}(2);\ \mathrm{Census\ FT.28};\ \mathrm{Def\text{-}H.4}];
\]

%dsp-level:0
\[
\textbf{the door is the vertex:}
\]

%dsp-level:0
\[
\qquad \Gamma(\varepsilon_q{=}0)=0\ \ \text{exactly},
\]

%dsp-level:0
\[
\qquad x_+x_-=4\varepsilon_q(1-4\varepsilon_q)\le\frac14\ \ \text{with equality at}\ \ \varepsilon_q=\frac18
\]

%dsp-level:0
\[
\qquad [\mathrm{BM.23}(7),(8);\ \text{machine}]:
\]

%dsp-level:0
\[
\boxed{\ \textbf{the line materializes only where the flicker is present}\ —\ \text{the same threshold}\ \mathrm{BM.94}\text{'s integral averages over is the one that opens the vertex}\ }
\]

%dsp-level:0
\[
\textbf{the materialized line propagates as a rung-one seal, forward-only:}
\]

%dsp-level:0
\[
\qquad \neg\exists\,\tau^{-1},
\]

%dsp-level:0
\[
\qquad \neg\exists\ \text{an inverse decay in the ledger}
\]

%dsp-level:0
\[
\qquad [\mathrm{A6};\ \mathrm{BM.23}(9);\ \mathrm{FT.144}(2)];
\]

%dsp-level:0
\[
\Gamma\ \text{is a weight per tick},\ \neg\text{a rate per external second}
\]

%dsp-level:0
\[
\qquad [\mathrm{BM.23}(9);\ \mathrm{D1.8};\ \mathrm{P5}];
\]

%dsp-level:0
\[
\qquad \neg\text{a lifetime, endpoint or mixing angle is derived or compared}\ [\mathrm{P5}];
\]

%dsp-level:0
\[
(8)\quad\textbf{status:}
\]

%dsp-level:0
\[
\qquad \textbf{derived:}\ (1)\text{--}(7)\ \text{in form}\ —\ \text{the angle, the order, the positivity, the occupancy, the four bars, the two states and the door};
\]

%dsp-level:0
\[
\neg\textbf{derived:}\ \ \text{the absolute scale}\ (\hat\mu_\nu,\ell_c^{(\nu)}),\ \ \text{the line index}\ i_\nu\,\ \ \text{the mixing matrix}\ [\mathrm{FT.145}(3)];
\]

%dsp-level:0
\[
\neg\text{a declaration issued}\ —\ \text{A-BM.5, A-BM.7 and}\ \mathrm{Def\text{-}H.37}\ \text{are re-cited,}\ \neg\text{re-issued};
\]

%dsp-level:0
\[
\qquad \neg\text{a P4 addition},\quad\neg\text{a new dimensionful scale}\ [\mathrm{P1}],\quad\neg\text{a measured value}\ [\mathrm{P5}],\quad\text{§Cal-BM stays empty};
\]

%dsp-level:0
\[
\text{species names are confined to the identification layer}\ [\mathrm{Def\text{-}BM.6}(2),(7)];
\]

%dsp-level:0
\[
\qquad \text{topology labels are internal bookkeeping}\ [\mathrm{Def\text{-}H.32}(4)];
\]

%dsp-level:0
\[
\qquad \text{amplitude}\ne\text{mass}\ [\mathrm{FT.79}(5)];
\]

%dsp-level:0
\[
\qquad \text{rung}\ne\text{scale}\ [\mathrm{D0.8}].
\]

\textbf{Req.} A5, A8, D0.8, D1.2, D1.4, D1.5, D1.8, D2.6, D2.9, D2.11, D3.3 (Parts 1–2); A §0.3, CT0.1, CT0.2, CT0.12, CT0.13, CT0.14, CT3.11 (Part 1, Part 2); CT2.5, CT4.2, EW.15, EW.33, EW.36, EM.9 (Part 2); P1, P4, P5 (Parts 1–2); Def-H.4, Def-H.32(4), Def-H.33, Def-H.37(1), FT.19, FT.28, FT.40, FT.54(4), FT.79(5), FT.108(5), FT.122(4), FT.142(1), FT.142(2), FT.142(3), FT.142(4), FT.142(5), FT.143(1), FT.143(3), FT.143(4), FT.144(2), FT.144(3), FT.144(4), FT.144(5), FT.145(2), FT.145(3) (Part 3); Census FT.28 (row-wise); Def-BM.11(5), BM.13(3), BM.23(7), BM.23(8), BM.23(9), BM.25(4), BM.27(6), BM.39(1), BM.39(5), BM.41(1), BM.45, BM.70(2), BM.70(3), BM.73(2), BM.73(3), BM.73(4), BM.77(2), BM.77(4), BM.80(6), BM.93(5) (this part); Cal-F; exact identities in $\mathbb Q[\sqrt2,\sqrt3][c,s_\nu]/(c^2{+}s_\nu^2{-}1)$, exact Taylor enclosures with rational remainder bounds, sum-of-squares positivity, exact $\mathbb Q(i)$ block algebra, exact resolvent identities in a truncated polynomial ring [math; machine: \texttt{checks/\allowbreak{}verify\_\allowbreak{}H\_\allowbreak{}flicker\_\allowbreak{}integral\_\allowbreak{}restmass.py} §C].

\textbf{Pf.}

%dsp-level:0
\[
(1):\ \mathrm{FT.142}(1)\text{--}(5)\ \text{re-cited (frozen)};
\]

%dsp-level:0
\[
\qquad \delta_\nu=\tfrac2{16}=\tfrac18\ [\text{machine: exact Fractions}];
\]

%dsp-level:0
\[
\qquad \dim\mathcal K_1=6+(5+3+1+1)=16=C_{\mathrm{cap}}\ [\mathrm{CT0.12};\ \mathrm{CT0.14};\ \text{machine}]
\]

%dsp-level:0
\[
(2):\ \cos(\tfrac18-\tfrac{2\pi}3)
\]

%dsp-level:0
\[
= -\tfrac c2+\tfrac{\sqrt3}2s_\nu\ \text{and}\ \cos(\tfrac18-\tfrac{4\pi}3)
\]

%dsp-level:0
\[
= -\tfrac c2-\tfrac{\sqrt3}2s_\nu\ [\text{the addition formulas}]\implies\text{the three boxed values}\ [\text{machine: ring identities}];
\]

%dsp-level:0
\[
\qquad A_0,B_0\ \text{re-cited from}\ \mathrm{BM.39}(1)
\]

%dsp-level:0
\[
G_1>A_0\iff\tfrac{B_0}2(c-\sqrt3s_\nu)>0
\]

%dsp-level:0
\[
\iff c>\sqrt3s_\nu
\]

%dsp-level:0
\[
\iff c^2>3(1-c^2)
\]

%dsp-level:0
\[
\iff 4c^2>3
\]

%dsp-level:0
\[
\iff 2c^2-1>\tfrac12\iff\cos\tfrac14>\tfrac12\ [\text{machine: the double-angle identity}];
\]

%dsp-level:0
\[
\qquad \cos\tfrac14>1-\tfrac1{2\cdot16}
\]

%dsp-level:0
\[
= \tfrac{31}{32}>\tfrac12\ [\text{machine: the alternating-series sign row, exact rationals}]
\]

%dsp-level:0
\[
c>\tfrac{63}{64}\ \text{from}\ \cos\tfrac18>1-\tfrac1{2\cdot64};
\]

%dsp-level:0
\[
\qquad s_\nu>\tfrac18-\tfrac1{6\cdot8^3}=\tfrac18-\tfrac1{3072}>0\ [\text{machine: exact Taylor bounds}];
\]

%dsp-level:0
\[
\qquad G_2>G_1\iff\sqrt3\,s_\nu>-\sqrt3\,s_\nu\iff s_\nu>0\ [\text{machine}]
\]

%dsp-level:0
\[
G=A_0+\tfrac{B_0}{\sqrt2}\bigl(1-r^{(\nu)}\bigr)\ \text{is affine strictly decreasing in}\ r\ [\mathrm{Def\text{-}BM.11}(5)\ \text{affine row};\ \mathrm{FT.143}(1);\ \text{machine}];
\]

%dsp-level:0
\[
\qquad r_1<r_2<r_0\ [\mathrm{FT.143}(3)\ \text{re-cited}]\implies\text{the transport order reverses the spectral order};
\]

%dsp-level:0
\[
\ \text{the value set is orbit-label independent}\ [\mathrm{Def\text{-}BM.11}(5)\ \text{orbit-label row};\ \text{machine: set equality}]
\]

%dsp-level:0
\[
(3):\ \text{the two SOS rows re-derived by expansion}\ [\text{machine}];
\]

%dsp-level:0
\[
\qquad |\cos|\le1\implies\text{the sandwich}\ [\text{math}];
\]

%dsp-level:0
\[
\qquad \text{the kernel line:}\ c=0\implies s_\nu^2=1\implies\cos(\varphi+\tfrac{2\pi}3)=\mp\tfrac{\sqrt3}2\ne0\ [\mathrm{BM.39}(1)\ \text{re-cited};\ \text{machine}]
\]

%dsp-level:0
\[
m_i^2/m_j^2=G_i/G_j\ \text{and}\ Q=\tfrac69=\tfrac23\ [\mathrm{FT.143}(3),(4)\ \text{re-cited};\ \text{machine}];
\]

%dsp-level:0
\[
\qquad m_\nu\mathfrak h\ge2\ln2\ [\mathrm{BM.70}(2),(3)\ \text{re-cited}]
\]

%dsp-level:0
\[
(4):\ \mathrm{FT.144}(2)\text{'s domain rows and prefix law re-cited};
\]

%dsp-level:0
\[
\qquad \text{occupancy}\ (1,0)\ [\mathrm{FT.144}(5);\ \mathrm{BM.73}(3)];
\]

%dsp-level:0
\[
\qquad P_{L/R}^2=P_{L/R},\ P_LP_R=0,\ P_L+P_R=\mathrm{id},\ \mathfrak T P_{L/R}=1\ [\mathrm{FT.122}(2);\ \text{machine over}\ \mathbb Q(i)]
\]

%dsp-level:0
\[
\rho_\circ=\tfrac14I:\ \mathfrak T=1,\ \operatorname{charpoly}=(x-\tfrac14)^4,\ S_{\mathrm{raw}}=\ln4=\max\ [\mathrm{CT0.2};\ \mathrm{BM.35}(1)\ \text{re-cited}];
\]

%dsp-level:0
\[
\qquad \mathcal P_MF\mathcal P_M=0\implies\neg\exists\ \text{an accessible pole}\ [\mathrm{CT4.2};\ \mathrm{CT2.5}\ \text{re-cited}];
\]

%dsp-level:0
\[
\qquad \neg\mathfrak h,\neg\mu,\neg h_0\ \text{in the premises}\ [\mathrm{FT.144}(3)\ \text{re-cited}]
\]

%dsp-level:0
\[
(5):\ \hat\omega P_L\hat\omega^{-1}
\]

%dsp-level:0
\[
= P_R,
\]

%dsp-level:0
\[
\ \hat\omega|L\rangle
\]

%dsp-level:0
\[
= -i|R\rangle,
\]

%dsp-level:0
\[
\ \hat\omega^2
\]

%dsp-level:0
\[
= \mathrm{id}\ [\mathrm{FT.122}(4);\ \mathrm{BM.73}(4)\ \text{re-derived entrywise over}\ \mathbb Q(i);\ \text{machine}];
\]

%dsp-level:0
\[
\qquad \text{the no-fixed-sector census on six sectors}\ [\mathrm{BM.25}(4)\ \text{re-cited};\ \text{machine}]
\]

%dsp-level:0
\[
(6):\ (\mathrm S1)\text{--}(\mathrm S4)\ \text{re-cited at their own loci}\ [(4),(5);\ \mathrm{CT0.12};\ \mathrm{EW.15};\ \mathrm{FT.145}(3);\ \mathrm{P1};\ \text{A §0.3}];
\]

%dsp-level:0
\[
\qquad \text{the inventory scans:}\ \neg\exists\ \text{a Yukawa-type token outside a negation row},
\]

%dsp-level:0
\[
\ \neg\exists\ \text{a second dimensionful scale}\ [\text{machine}]
\]

%dsp-level:0
\[
(7):\ \text{the resolvent identity re-derived coefficient-wise to}\ w^{12}\ [\text{machine}];
\]

%dsp-level:0
\[
\qquad x_+x_-=4\varepsilon_q(1-4\varepsilon_q)\le\tfrac14\ \text{with equality at}\ \varepsilon_q=\tfrac18\ [\mathrm{BM.23}(7)\ \text{re-cited};\ \text{machine}];
\]

%dsp-level:0
\[
\qquad \Gamma=0\ \text{at}\ \varepsilon_q=0\ [\mathrm{BM.23}(8)];
\]

%dsp-level:0
\[
\qquad m_{\mathrm{hid}}=2\ln2/\mathfrak h\ [\mathrm{FT.145}(2)]
\]

%dsp-level:0
\[
(8):\ \text{the standing-rule scans: external tokens, decimal numerals, registry tokens, declaration lines, the barred object}\ \alpha_B\ [\text{machine}].\ \blacksquare
\]

\textbf{Phys.} The neutral line the vertex emits is a rung-one seal, and everything about it that the corpus can derive, it derives. Its channel angle is two over sixteen, the undressed seal coefficient over the mask dimension the complete-force rule already fixed; its rest mass is the transport defect evaluated on that angle, three exact values in a radical ring, ordered heaviest to lightest against the spectral order by the same inversion the family mass form proves in general. Every one of the three is strictly positive, squeezed between two sums of squares that never touch zero, and the kernel line that would make one of them massless is unreachable by an exact trigonometric census — so there is no massless option to decline and no zero to explain away. The absolute scale stays undeclared, as every scale in this part does, and the internal ratios stay derived. On the carrier the line occupies one ray and its complement is empty: not heavy, not suppressed, empty, because a dissolution-first history is outside the transition domains and unregistered content supports no accessible pole at any rung. The image orientation the vertex selects is the mirror of that same single-arrow locus, one orbit of size two carrying no datum of its own. Four separate bars then close the door on any seesaw-type mechanism — no partner state, no self-pairing, no coupling object, no second scale — and the smallness of the line is left where the corpus has always left it, with the free internal data. The two states of the line are the vacuum-resident balanced pair and the materialized seal, and the only door between them is the vertex, whose weight vanishes exactly at the sharp threshold: the line materializes only where the flicker is present, which is the same threshold the ensemble integral is taken over.

\end{theorem}

\begin{theorem}[{\texorpdfstring{BM.97 — The rest mass is the pole mass: the species adjudication, the rest-mass Q-value, and the exact critical locus}{BM.97 - The rest mass is the pole mass: the species adjudication, the rest-mass Q-value, and the exact critical locus}}]\label{tag:BM-97}
\textbf{Thm.}

%dsp-level:0
\[
(1)\quad\textbf{the species census at the rest slot (three candidates, one read specified):}\qquad\begin{array}{l|l|l}\text{candidate} & \text{expression} & \text{status}\\\hline\text{the informational read} & m_{\mathrm{hid}}=E_{\mathrm{int}}\ln2/\mathfrak h & \text{the ledger's bit currency, ensemble-distributed}\ [\mathrm{FT.44}(4);\ \mathrm{BM.35}(4);\ \mathrm{BM.94}(8)]\\\text{the sector-deficit route} & m^{(\mathrm{ii})}=T(q,S)/\mathfrak h,\quad T=W(Q)\,\overline{\Delta S}(S) & \text{a measure price, vertex-scoped}\ [\mathrm{BM.28}(4)(\mathrm{ii});\ \mathrm{BM.30}(5);\ \mathrm{BM.37}(6);\ \mathrm{BM.77}(2)]\\\text{the pole/transport read} & m^2_{\mathrm{pole}}=\hat\mu^2\ell_c^2G(c)=4\hat\mu'^2\sin^2(\Theta_R/4) & \textbf{the rest mass on the transport species}\ [\mathrm{Def\text{-}H.37}(1),(2);\ \mathrm{FT.54}(4);\ \mathrm{FT.47}(1);\ \mathrm{EW.33};\ \mathrm{BM.39}(1),(7),(9)]\end{array}
\]

%dsp-level:0
\[
\textbf{the matching of the two pole clauses is derived at leading order and recorded,}\ \neg\textbf{imposed:}
\]

%dsp-level:0
\[
\qquad \hat\mu'^2\Theta_R^2/4=\hat\mu^2\ell_c^2G(c)
\]

%dsp-level:0
\[
\qquad [\mathrm{FT.54}(2);\ \mathrm{FT.33};\ \mathrm{BM.39}(9);\ \mathrm{Def\text{-}H.37}\ \text{locus clause}];
\]

%dsp-level:0
\[
\boxed{\ \textbf{the rest mass is the pole mass}\ m^{\mathrm{rest}}=m_{\mathrm{pole}}\ \ \text{on the transport species};\qquad m_{\mathrm{hid}}\ \ \text{is}\ \neg\textbf{the rest mass on that species}\ }
\]

%dsp-level:0
\[
\text{equating them is a conversion the corpus bars}
\]

%dsp-level:0
\[
\qquad [\mathrm{BM.35}(5);\ \mathrm{BM.83}(3);\ \mathrm{BM.37}(2);\ \text{pointer}];
\]

%dsp-level:0
\[
\qquad \text{the specified read A-BM.8 reads the baryon classes on route}\ (\mathrm{ii})\ \text{and the emitted lines on the transport species}\ [\mathrm{BM.77}(1);\ \text{re-cited},\ \neg\text{re-issued}];
\]

%dsp-level:0
\[
(2)\quad\textbf{the rest-mass Q-value, formed:}
\]

%dsp-level:0
\[
\qquad Q_\beta^{\mathrm{rest}}\ :=\ m_n^{\mathrm{rest}}-m_p^{\mathrm{rest}}-m_\ell^{\mathrm{rest}}-m_\nu^{\mathrm{rest}}
\]

%dsp-level:0
\[
\qquad [\mathrm{BM.41}(4);\ \mathrm{BM.63}(1);\ \mathrm{BM.85}(1);\ \text{every member a transport-species mass}]:
\]

%dsp-level:0
\[
\boxed{\ Q_\beta^{\mathrm{spec}}\,\mathfrak h\ =\ \Delta T(S)\ -\ \rho_\ell\sqrt{G\bigl(c(0)\bigr)}\ -\ \rho_\nu\sqrt{G\bigl(c(i_\nu)\bigr)},\qquad\Delta T(S)=5\,m^2_B(\cdot,S)\ }
\]

%dsp-level:0
\[
\qquad [\mathrm{BM.77}(2);\ \text{the expression is its own frozen}\ Q_{\mathrm{ii}}]:
\]

%dsp-level:0
\[
\Delta T=\Bigl(\frac{4219785}{98956},\ \ \frac{9608310}{393679},\ \ \frac{455}{23},\ \ \frac{68565}{3088}\Bigr)\ \ \text{at}\ S=(0,1,2,3),
\]

%dsp-level:0
\[
\qquad \frac{455}{23}\le\Delta T(S)\le\frac{4219785}{98956}
\]

%dsp-level:0
\[
\qquad [\mathrm{BM.77}(2);\ \text{machine re-certified}];
\]

%dsp-level:0
\[
G_\ell=A_0-B_0\cos\delta_\ell\ \ \text{at}\ \ \delta_\ell=\tfrac29
\]

%dsp-level:0
\[
\qquad [\mathrm{FT.112};\ \mathrm{Def\text{-}BM.6}(1)],
\]

%dsp-level:0
\[
\qquad G_\nu=A_0-B_0\cos\Bigl(\tfrac18-\frac{2\pi i_\nu}3\Bigr)\ \ \text{at the derived}\ \ \delta_\nu=\tfrac18
\]

%dsp-level:0
\[
\qquad [\mathrm{BM.96}(2);\ \mathrm{FT.142}(5)];
\]

%dsp-level:0
\[
\rho_\ell=\mathfrak h\,\hat\mu_\ell\,\ell_c^{(\ell)},\quad\rho_\nu=\mathfrak h\,\hat\mu_\nu\,\ell_c^{(\nu)}\in\mathfrak q^0
\]

%dsp-level:0
\[
\qquad [\mathrm{BM.77}(2);\ \mathrm{CT2.1};\ \mathrm{Def\text{-}H.37}(1)];
\]

%dsp-level:0
\[
\qquad \textbf{every row below holds uniformly in}\ i_\nu\in\{0,1,2\}
\]

%dsp-level:0
\[
\qquad [\mathrm{BM.77}(2)];
\]

%dsp-level:0
\[
(3)\quad\textbf{the flicker-blindness of the rest-mass read:}
\]

%dsp-level:0
\[
\qquad \text{the specified read's alphabet is}\ \{m^2_B(\cdot,S),\ \rho_\ell,\ \rho_\nu,\ g,\ g',\ \mathfrak h\}
\]

%dsp-level:0
\[
\qquad [\mathrm{BM.92}(5);\ \text{machine: the token census}]:
\]

%dsp-level:0
\[
\boxed{\ \neg E_{\mathrm{int}},\quad\neg w,\quad\neg w_{\mathrm{mat}},\quad\neg\varepsilon_q,\quad\neg\mathsf g,\quad\neg N\ \ \text{occur in it}\ }
\]

%dsp-level:0
\[
\qquad [\text{machine}]\implies\textbf{the flicker integral changes no row of the rest-mass Q-value}
\]

%dsp-level:0
\[
\textbf{the two acts are orthogonal:}
\]

%dsp-level:0
\[
\qquad \mathrm{BM.94}\text{'s fix lands on the informational register alone},
\]

%dsp-level:0
\[
\qquad \mathrm{BM.77}\text{'s read lands on the transport species alone}
\]

%dsp-level:0
\[
\qquad [\mathrm{BM.77}(3);\ \mathrm{BM.92}(3)\ \text{pattern}]:
\]

%dsp-level:0
\[
\implies\text{the trichotomy, the four}\ \Delta T\ \text{values,}\ \rho^{\mathrm{crit}}\text{'s enclosure and all eight witnesses stand}
\]

%dsp-level:0
\[
\qquad [\mathrm{BM.77}(5),(6);\ \mathrm{BM.92}(3);\ \text{machine re-certified}];
\]

%dsp-level:0
\[
(4)\quad\textbf{the pure-pole read (the third formed read):}
\]

%dsp-level:0
\[
\qquad Q_\beta^{\mathrm{pole}}\ :=\ m_n^{\mathrm{pole}}-m_p^{\mathrm{pole}}-m_\ell-m_\nu,
\]

%dsp-level:0
\[
\qquad M_0:=m_p^{\mathrm{pole}},
\]

%dsp-level:0
\[
\qquad (m_n^{\mathrm{pole}})^2=M_0^2+\Delta_g,
\]

%dsp-level:0
\[
\qquad s:=m_\ell+m_\nu>0
\]

%dsp-level:0
\[
\textbf{hypothesis}\ (\mathrm H):\ \text{the two baryon classes differ by exactly one branch substitution}\ [\mathrm{BM.59}(9);\ \mathrm{BM.69}(3);\ \text{carried on its face}]:
\]

%dsp-level:0
\[
\Delta_g\ \in\ \bigl\{2,\ \ \sqrt2-1,\ \ 2\cos\tfrac\pi{10}-\sqrt3\bigr\}\,\hat\mu'^2\ >\ 0,
\]

%dsp-level:0
\[
\qquad \Delta_1>\Delta_2>\Delta_3
\]

%dsp-level:0
\[
\qquad [\mathrm{Def\text{-}BM.14}(5);\ \mathrm{BM.69}(3);\ \text{machine: the certificates}\ 4>2,\ 2>1,\ 3<5];
\]

%dsp-level:0
\[
\textbf{the exact critical locus:}
\]

%dsp-level:0
\[
\qquad \boxed{\ Q_\beta^{\mathrm{pole}}>0\iff\Delta_g>\Delta_g^{\mathrm{crit}}:=s\,(2M_0+s);\qquad Q_\beta^{\mathrm{pole}}=0\iff\Delta_g=\Delta_g^{\mathrm{crit}};\qquad Q_\beta^{\mathrm{pole}}<0\iff\Delta_g<\Delta_g^{\mathrm{crit}}\ }
\]

%dsp-level:0
\[
[\text{machine: cleared by squaring two positive sides},\ \sqrt{M_0^2+\Delta_g}>M_0+s\iff M_0^2+\Delta_g>(M_0+s)^2];
\]

%dsp-level:0
\[
\textbf{the equivalent line form:}
\]

%dsp-level:0
\[
\qquad s^{\mathrm{crit}}=\sqrt{M_0^2+\Delta_g}-M_0=\frac{\Delta_g}{\sqrt{M_0^2+\Delta_g}+M_0}=m_n^{\mathrm{pole}}-m_p^{\mathrm{pole}}
\]

%dsp-level:0
\[
\qquad [\text{machine: the cleared identity}];
\]

%dsp-level:0
\[
\Delta_g^{\mathrm{crit}}\ \ \text{is strictly increasing in}\ s\ \ \text{and in}\ M_0,
\]

%dsp-level:0
\[
\qquad \Delta_g^{\mathrm{crit}}>0\ \ \text{strictly since}\ \ s>0
\]

%dsp-level:0
\[
\qquad [\mathrm{BM.96}(3);\ \text{machine: the coefficient census}];
\]

%dsp-level:0
\[
(5)\quad\textbf{the trichotomy (total, exact, non-empty on both sides):}
\]

%dsp-level:0
\[
(\mathrm a)\quad\mathcal O^{\mathrm{pole}}:=\bigl\{(\rho_\ell,\rho_\nu)>0:\ s<\Delta_g^{\mathrm{crit}}\text{'s line form}\bigr\}\ne\varnothing:
\]

%dsp-level:0
\[
\qquad \textbf{the universal symbolic witness}
\]

%dsp-level:0
\[
\rho_\ell=\rho_\nu=\frac{\sqrt{M_0^2+\Delta_g}-M_0}{4\sqrt{A_0+B_0}}\ \ \in\ \mathcal O^{\mathrm{pole}}
\]

%dsp-level:0
\[
\qquad [\mathrm{BM.77}(4)\text{'s sandwich}:\ s\le\tfrac12\bigl(\sqrt{M_0^2+\Delta_g}-M_0\bigr)<\sqrt{M_0^2+\Delta_g}-M_0;\ \text{machine}]:
\]

%dsp-level:0
\[
\text{valid at every generation, every line index}\ i_\nu,\ \text{every}\ (g,g')\ne(0,0);
\]

%dsp-level:0
\[
\qquad \textbf{the integer-certified rational witness at}\ (g,g',\rho_\ell,\rho_\nu,\hat\mu'^2\mathfrak h^2,M_0\mathfrak h)
\]

%dsp-level:0
\[
= (1,1,\tfrac1{100},\tfrac1{100},1,1),
\]

%dsp-level:0
\[
\ \text{generation}\ 1:
\]

%dsp-level:0
\[
A_0+B_0=\frac7{16}+\frac14=\frac{11}{16},
\]

%dsp-level:0
\[
\qquad \sqrt G\le\frac{\sqrt{11}}4<1
\]

%dsp-level:0
\[
\qquad [11<16]\implies s\,\mathfrak h<\frac1{50}\implies\Delta_g^{\mathrm{crit}}\mathfrak h^2<\frac1{50}\cdot\frac{101}{50}=\frac{101}{2500}<2=\Delta_1\mathfrak h^2
\]

%dsp-level:0
\[
\boxed{\ Q_\beta^{\mathrm{pole}}>0\ \ \text{on the certificate}\ \ 5000>101\ }\qquad[\text{machine}];
\]

%dsp-level:0
\[
(\mathrm b)\quad\mathcal C^{\mathrm{pole}}\ne\varnothing:
\]

%dsp-level:0
\[
\qquad \textbf{the universal symbolic witness}
\]

%dsp-level:0
\[
\qquad \rho_\ell=\rho_\nu=\frac{\sqrt{M_0^2+\Delta_g}-M_0}{\sqrt{A_0-B_0}}\ \ \in\ \mathcal C^{\mathrm{pole}}
\]

%dsp-level:0
\[
\qquad [\text{machine:}\ s\ge2\bigl(\sqrt{M_0^2+\Delta_g}-M_0\bigr)];
\]

%dsp-level:0
\[
\textbf{the integer-certified witness at}\ (1,1,100,100,1,1):
\]

%dsp-level:0
\[
\qquad A_0-B_0=\frac3{16},
\]

%dsp-level:0
\[
\qquad \sqrt G\ge\frac{\sqrt3}4>\frac25\qquad \qquad [75>64]:
\]

%dsp-level:0
\[
m_\ell\,\mathfrak h\ \ge\ 25\sqrt3\ >\ 40\qquad[1875>1600]\implies s\,\mathfrak h\ \ge\ 50\sqrt3\ >\ 80\qquad[75>64]\implies\Delta_g^{\mathrm{crit}}\mathfrak h^2\ >\ 80\cdot82=6560\ >\ 2
\]

%dsp-level:0
\[
\boxed{\ Q_\beta^{\mathrm{pole}}<0\ \ \text{on the certificates}\ \ 1875>1600,\quad75>64,\quad6560>2\ }\qquad[\text{machine}];
\]

%dsp-level:0
\[
(\mathrm c)\quad\textbf{the critical locus is exact, unique and generation-relative:}
\]

%dsp-level:0
\[
\qquad \partial\mathcal O^{\mathrm{pole}}:\quad s=\sqrt{M_0^2+\Delta_g}-M_0
\]

%dsp-level:0
\[
\qquad [\text{machine: the boundary equation}];
\]

%dsp-level:0
\[
\Delta_g^{\mathrm{crit}}\ \ \text{is generation-blind}\ (\text{a function of the emitted lines and}\ M_0\ \text{alone})\ \ \text{while}\ \ \Delta_g\ \ \text{is generation-specific}\qquad[(4)]:
\]

%dsp-level:0
\[
\boxed{\ \textbf{the opening region is widest in the first generation and narrowest in the third, in the exact order}\ \ \Delta_1>\Delta_2>\Delta_3\ }
\]

%dsp-level:0
\[
\qquad [\mathrm{BM.69}(3);\ \text{machine: monotone in}\ \Delta_g\ \text{at fixed}\ M_0];
\]

%dsp-level:0
\[
\textbf{on the symmetric ray}\ \rho_\ell
\]

%dsp-level:0
\[
= \rho_\nu
\]

%dsp-level:0
\[
= \rho\ \textbf{at the witness scales}\ (g,g',M_0\mathfrak h,\hat\mu'^2\mathfrak h^2)
\]

%dsp-level:0
\[
= (1,1,1,1),
\]

%dsp-level:0
\[
\ \text{generation}\ 1:
\]

%dsp-level:0
\[
\qquad \rho^{\mathrm{crit,pole}}=\frac{\bigl(\sqrt3-1\bigr)}{2\sqrt G}:
\]

%dsp-level:0
\[
\boxed{\ \frac25\ <\ \rho^{\mathrm{crit,pole}}\ <\ \frac67\ }
\]

%dsp-level:0
\[
\qquad [\text{machine:}\ \sqrt G\le\tfrac{\sqrt{11}}4\ \text{gives the lower bound on}\ 2(\sqrt3-1)/\sqrt{11}>\tfrac25\iff89>50\sqrt3\iff7921>7500;\quad\sqrt G\ge\tfrac{\sqrt3}4\ \text{gives the upper bound on}\ 2(\sqrt3-1)/\sqrt3<\tfrac67\iff8\sqrt3<14\iff192<196];
\]

%dsp-level:0
\[
(6)\quad\textbf{the}\ \nu\textbf{-positivity shift (a derived consequence of}\ \mathrm{BM.96}(3)\textbf{):}
\]

%dsp-level:0
\[
\qquad \Delta_g^{\mathrm{crit}}(m_\ell{+}m_\nu)-\Delta_g^{\mathrm{crit}}(m_\ell)=2M_0m_\nu+2m_\ell m_\nu+m_\nu^2\ >\ 0
\]

%dsp-level:0
\[
\qquad [\text{machine: the exact expansion}]:
\]

%dsp-level:0
\[
\boxed{\ \textbf{the derived positivity of the}\ \nu\ \textbf{rest mass strictly narrows the opening region by an exact amount}\ —\ \text{a massless}\ \nu\ \text{line is refuted as a read,}\ \neg\text{assumed}\ }
\]

%dsp-level:0
\[
\textbf{the two-body companion on this read:}
\]

%dsp-level:0
\[
\qquad E_p-m_p=\frac{(m_n-m_p)^2-m_\ell^2}{2m_n}
\]

%dsp-level:0
\[
\qquad [\mathrm{BM.41}(2);\ \mathrm{BM.77}(6);\ \text{machine: the identity re-derived at}\ m_n\ne m_p]:
\]

%dsp-level:0
\[
\text{the two-body channel is admissible}\iff(m_n^{\mathrm{pole}}-m_p^{\mathrm{pole}})^2\ge m_\ell^2\iff\Delta_g\ \ge\ m_\ell\bigl(\sqrt{M_0^2+\Delta_g}+M_0\bigr)
\]

%dsp-level:0
\[
\qquad [\text{machine: the cleared form}];
\]

%dsp-level:0
\[
\textbf{the forced-third-member row is read-relative:}
\]

%dsp-level:0
\[
\qquad \mathrm{BM.41}(2)\text{'s equal-mass obstruction is lifted on the pure-pole read}\ (\Delta_g>0),
\]

%dsp-level:0
\[
\ \text{so}\ \mathrm{BM.41}(3),
\]

%dsp-level:0
\[
(6)\text{'s forcing stands on the informational reads}\ [\mathrm{BM.70}(4)(\mathrm a)]\ \text{and the grammar stands on the specified read}\ [\mathrm{BM.45};\ \mathrm{BM.77}(6)],
\]

%dsp-level:0
\[
\text{and it is}\ \neg\textbf{inherited by the pure-pole read}\ —\ \text{recorded,}\ \neg\text{adjudicated beyond the branch rows}\qquad[\mathrm{BM.77}(6)\ \text{pattern}];
\]

%dsp-level:0
\[
(7)\quad\textbf{the read discipline:}\qquad\textbf{three formed reads, pairwise incomparable:}
\]

%dsp-level:0
\[
Q_\beta^{\mathrm{info}}:\ \{\ln2,\ln3,\ln5,\ N,\ w\ \text{or}\ w_{\mathrm{mat}},\ \mathsf g,\ \varepsilon_{\mathrm{EM}},\ \mathfrak h,\ m_\ell,\ m_\nu\};
\]

%dsp-level:0
\[
\qquad Q_\beta^{\mathrm{spec}}:\ \{m^2_B(\cdot,S),\ \rho_\ell,\ \rho_\nu,\ g,\ g',\ \mathfrak h\};
\]

%dsp-level:0
\[
\qquad Q_\beta^{\mathrm{pole}}:\ \{M_0,\ \Delta_g,\ \hat\mu'^2,\ \rho_\ell,\ \rho_\nu,\ g,\ g',\ \mathfrak h\}
\]

%dsp-level:0
\[
\boxed{\ \text{the informational alphabet is disjoint from the two transport alphabets except for}\ \mathfrak h;\qquad\neg\text{a derived equation relates the three expressions}\ }
\]

%dsp-level:0
\[
\qquad [\text{machine: the three token censuses};\ \text{the mass-form bridge residue, localized at}\ \mathrm{BM.98}(5)];
\]

%dsp-level:0
\[
\textbf{the decay direction is fixed in all three:}
\]

%dsp-level:0
\[
\qquad T(0,S)=\tfrac43T(+1,S)\ [\mathrm{BM.28}(4)(\mathrm{ii})],
\]

%dsp-level:0
\[
\qquad \Delta_g>0\ [\mathrm{BM.69}(3)],
\]

%dsp-level:0
\[
\qquad \Delta m>0\ \ \text{sub-critically}\ [\mathrm{BM.91}(6);\ \mathrm{BM.99}(4)]:
\]

%dsp-level:0
\[
\boxed{\ \textbf{the parent is the neutral class on every formed read, and the grammar}\ n\to p+\ell+\nu\ \ \textbf{stands unmodified}\ }
\]

%dsp-level:0
\[
\qquad [\mathrm{BM.41}(1);\ \mathrm{BM.45};\ \mathrm{BM.77}(6)];
\]

%dsp-level:0
\[
\neg\text{a currency is converted}\ [\mathrm{BM.83}(3);\ \mathrm{BM.90}(5)];
\]

%dsp-level:0
\[
\qquad \text{route}\ (\mathrm{ii})\ \text{is}\ \neg\text{elevated beyond the vertex kinematics}\ [\mathrm{BM.30}(5);\ \mathrm{BM.77}(9)];
\]

%dsp-level:0
\[
\qquad \text{route}\ (\mathrm{iii})\text{'s amplitude ratios stay barred as masses}\ [\mathrm{FT.79}(5);\ \mathrm{BM.37}(3)];
\]

%dsp-level:0
\[
(8)\quad\textbf{status:}
\]

%dsp-level:0
\[
\qquad \textbf{derived:}\ (1)\text{--}(7)\ \text{given the specified read and hypothesis}\ (\mathrm H);
\]

%dsp-level:0
\[
\qquad \textbf{declared: nothing}\ —\ \text{A-BM.8 and}\ \mathrm{Def\text{-}H.37}\ \text{are re-cited,}\ \neg\text{re-issued};
\]

%dsp-level:0
\[
\neg\text{a clause of}\ \mathrm{BM.41},\ \mathrm{BM.49},\ \mathrm{BM.63},\ \mathrm{BM.70},\ \mathrm{BM.71},\ \mathrm{BM.77},\ \mathrm{BM.85},\ \mathrm{BM.92}\ \text{is retracted};
\]

%dsp-level:0
\[
\qquad \neg\text{a P4 addition},\quad\neg\text{a new dimensionful scale}\ [\mathrm{P1}],\quad\neg\text{a measured value}\ [\mathrm{P5}],\quad\text{§Cal-BM stays empty};
\]

\textbf{Req.} D2.9 (Parts 1–2); CT2.1, EM.25(7), EW.15 (Part 2); P1, P4, P5 (Parts 1–2); Def-H.37(1), Def-H.37(2), EW.33, FT.33, FT.47(1), FT.54(2), FT.54(4), FT.79(5), FT.112, FT.142(5) (Part 3); Def-BM.6(1), Def-BM.11(5), Def-BM.14(5), BM.28(4), BM.30(5), BM.35(4), BM.35(5), BM.37(2), BM.37(3), BM.37(6), BM.39(1), BM.39(7), BM.39(9), BM.41(1), BM.41(2), BM.41(3), BM.41(4), BM.41(6), BM.45, BM.59(9), BM.63(1), BM.69(3), BM.70(4), BM.77(1), BM.77(2), BM.77(3), BM.77(4), BM.77(5), BM.77(6), BM.77(9), BM.83(3), BM.85(1), BM.87, BM.90(5), BM.91(6), BM.92(3), BM.92(5), BM.96(2), BM.96(3), BM.96(4) (this part); Cal-F; exact radical inequalities by integer certificates, sum-of-squares positivity, cleared polynomial identities over $\mathbb Q[M_0,m_\ell,m_\nu]$, exact rational arithmetic at the four frozen sector expectations, three token censuses [math; machine: \texttt{checks/\allowbreak{}verify\_\allowbreak{}H\_\allowbreak{}flicker\_\allowbreak{}integral\_\allowbreak{}restmass.py} §D].

\textbf{Pf.}

%dsp-level:0
\[
(1):\ \mathrm{Def\text{-}H.37}(1),(2)\text{'s two clauses and}\ \mathrm{FT.54}(2),(4)\text{'s matching re-cited};
\]

%dsp-level:0
\[
\qquad \text{the three-candidate census assembled from}\ \mathrm{FT.44}(4),
\]

%dsp-level:0
\[
\ \mathrm{BM.28}(4)(\mathrm{ii}),
\]

%dsp-level:0
\[
\ \mathrm{BM.39}(1),
\]

%dsp-level:0
\[
(7)\ [\text{machine: the read-class rows}];
\]

%dsp-level:0
\[
\qquad \text{the conversion bar re-cited}\ [\mathrm{BM.35}(5);\ \mathrm{BM.83}(3)]
\]

%dsp-level:0
\[
(2):\ \mathrm{BM.77}(2)\text{'s frozen expression re-cited with}\ G_\nu\ \text{evaluated at the derived}\ \delta_\nu
\]

%dsp-level:0
\[
= \tfrac18\ [\mathrm{BM.96}(2)];
\]

%dsp-level:0
\[
\qquad 5\cdot843957=4219785,
\]

%dsp-level:0
\[
\quad 5\cdot1921662=9608310,
\]

%dsp-level:0
\[
\quad 5\cdot91=455,
\]

%dsp-level:0
\[
\quad 5\cdot13713=68565\ [\text{machine: exact integers}];
\]

%dsp-level:0
\[
\qquad \tfrac{455}{23}\le\Delta T\le\tfrac{4219785}{98956}\ [\text{machine: cross-multiplication}]
\]

%dsp-level:0
\[
(3):\ \text{the alphabet enumerated from the expression and scanned for the six barred tokens}\ [\text{machine}];
\]

%dsp-level:0
\[
\qquad \text{the eight witnesses re-certified as integer comparisons}\ [\text{machine}]
\]

%dsp-level:0
\[
(4):\ \sqrt{M_0^2+\Delta_g}>M_0+s\ \text{with both sides positive}
\]

%dsp-level:0
\[
\iff M_0^2+\Delta_g>M_0^2+2M_0s+s^2\iff\Delta_g>2M_0s+s^2\ [\text{machine: cleared over}\ \mathbb Q[M_0,s,\Delta_g]];
\]

%dsp-level:0
\[
s^{\mathrm{crit}}
\]

%dsp-level:0
\[
= \sqrt{M_0^2+\Delta_g}-M_0\ \text{and}\ (\sqrt{M_0^2+\Delta_g}-M_0)(\sqrt{M_0^2+\Delta_g}+M_0)
\]

%dsp-level:0
\[
= \Delta_g\ [\text{machine: the difference of squares}];
\]

%dsp-level:0
\[
\qquad \Delta_1>\Delta_2>\Delta_3\ \text{on}\ 4>2,\ 2>1,\ 3<5\ [\mathrm{BM.69}(3)\ \text{re-cited}]
\]

%dsp-level:0
\[
(5)(\mathrm a):\ 2\cdot\tfrac{\sqrt{M_0^2+\Delta_g}-M_0}{4\sqrt{A_0+B_0}}\cdot\sqrt{A_0+B_0}
\]

%dsp-level:0
\[
= \tfrac12\bigl(\sqrt{M_0^2+\Delta_g}-M_0\bigr)<\sqrt{M_0^2+\Delta_g}-M_0\ [\text{machine}];
\]

%dsp-level:0
\[
\qquad A_0+B_0\big|_{g=g'=1}=\tfrac{11}{16}\ [\text{machine}]
\]

%dsp-level:0
\[
s\,\mathfrak h\le2\cdot\tfrac1{100}\cdot\tfrac{\sqrt{11}}4=\tfrac{\sqrt{11}}{200}<\tfrac4{200}=\tfrac1{50}\ [11<16;\ \text{machine}];
\]

%dsp-level:0
\[
\qquad \tfrac1{50}\bigl(2+\tfrac1{50}\bigr)=\tfrac{101}{2500};
\]

%dsp-level:0
\[
\qquad 2>\tfrac{101}{2500}\iff5000>101\ [\text{machine}]
\]

%dsp-level:0
\[
(\mathrm b):\ 2\cdot\tfrac{\sqrt{M_0^2+\Delta_g}-M_0}{\sqrt{A_0-B_0}}\cdot\sqrt{A_0-B_0}=2\bigl(\sqrt{M_0^2+\Delta_g}-M_0\bigr)>\sqrt{M_0^2+\Delta_g}-M_0\ [\text{machine}];
\]

%dsp-level:0
\[
\qquad A_0-B_0\big|_{g=g'=1}=\tfrac3{16};
\]

%dsp-level:0
\[
\qquad \sqrt G\ge\tfrac{\sqrt3}4>\tfrac25\iff75>64\ [\text{machine}]
\]

%dsp-level:0
\[
m_\ell\mathfrak h\ge100\cdot\tfrac{\sqrt3}4=25\sqrt3>40\iff1875>1600;
\]

%dsp-level:0
\[
\qquad s\,\mathfrak h\ge50\sqrt3>80\iff75>64;
\]

%dsp-level:0
\[
\qquad 80(2+80)=6560>2\ [\text{machine: integers}]
\]

%dsp-level:0
\[
(\mathrm c):\ \rho^{\mathrm{crit,pole}}=\tfrac{\sqrt3-1}{2\sqrt G}\ \text{at}\ M_0\mathfrak h=1,\ \Delta_1\mathfrak h^2=2\ [\text{machine:}\ \sqrt{1+2}=\sqrt3];
\]

%dsp-level:0
\[
\qquad \text{the lower bound:}\ \tfrac{2(\sqrt3-1)}{\sqrt{11}}>\tfrac25
\]

%dsp-level:0
\[
\iff 5(\sqrt3-1)>\sqrt{11}
\]

%dsp-level:0
\[
\iff 25(4-2\sqrt3)>11
\]

%dsp-level:0
\[
\iff 89>50\sqrt3
\]

%dsp-level:0
\[
\iff 7921>7500\ [\text{machine}]
\]

%dsp-level:0
\[
\text{the upper bound:}\ \tfrac{2(\sqrt3-1)}{\sqrt3}<\tfrac67\iff14(\sqrt3-1)<6\sqrt3\iff8\sqrt3<14\iff192<196\ [\text{machine}];
\]

%dsp-level:0
\[
\qquad \text{monotonicity in}\ \Delta_g\ \text{at fixed}\ M_0:\ \text{the coefficient of}\ \Delta_g\ \text{in}\ \sqrt{M_0^2+\Delta_g}-M_0\ \text{is positive}\ [\text{machine}]
\]

%dsp-level:0
\[
(6):\ (m_\ell+m_\nu)(2M_0+m_\ell+m_\nu)-m_\ell(2M_0+m_\ell)
\]

%dsp-level:0
\[
= 2M_0m_\nu+2m_\ell m_\nu+m_\nu^2\ [\text{machine: the expansion over}\ \mathbb Q[M_0,m_\ell,m_\nu]];
\]

%dsp-level:0
\[
2m_n(E_p-m_p)
\]

%dsp-level:0
\[
= (m_n-m_p)^2-m_\ell^2\ \text{at}\ E_p
\]

%dsp-level:0
\[
= \tfrac{m_n^2+m_p^2-m_\ell^2}{2m_n}\ [\text{machine: cleared}]\implies\text{the two-body condition}\ m_n-m_p
\]

%dsp-level:0
\[
\ge m_\ell\ \text{and its}\ \Delta_g\ \text{form}\ [\text{machine}]
\]

%dsp-level:0
\[
(7):\ \text{the three symbol sets enumerated from the three expressions and pairwise intersected}\ [\text{machine}];
\]

%dsp-level:0
\[
\qquad T(0,S)/T(+1,S)=20/15=4/3\ [\text{machine}];
\]

%dsp-level:0
\[
\qquad (8):\ \text{the standing-rule scans}\ [\text{machine}].\ \blacksquare
\]

\textbf{Phys.} A Q-value is a difference of rest masses, and a rest mass in this corpus is a pole mass: the transport defect of a materialized chain, exact as four mu-prime squared times the sine squared of a quarter frame angle, matched at leading order to the Gram evaluation the seed supplies. The informational count is a different currency, ensemble-distributed and bit-denominated, and equating the two is precisely the conversion the ledger bars — so the fix to the snapshot read and the read of the vertex turn out to be orthogonal acts, and the specified Q-value carries no edge count, no crystallization weight, no threshold and no curvature composition in its alphabet at all. On that read nothing moves: the four exact baryon values, the critical enclosure and all eight integer witnesses stand. The pure pole read is the third formed read, and it has its own exact trichotomy — a critical locus that is the emitted-line sum times twice the daughter's mass plus itself, a universal symbolic witness in each of the opening and closing regions, and rational witnesses certified by five thousand against one hundred one and by one thousand eight hundred seventy-five against one thousand six hundred. Because the neutrino line's mass is derived strictly positive, the critical locus sits strictly above its massless value by an exact amount, so the positivity narrows the opening region rather than being assumed away; and because the pole splitting is generation-specific while the critical locus is not, the opening region is widest in the first generation and narrowest in the third. What the three reads agree on is the direction: the neutral class is the parent in every currency, and the grammar is untouched.

\textbf{Rmk (the operative mass read — reconciled nomenclature).} The mass that works — the read entering every formed composite quantity of this era (the rest masses, the Q-values, the splittings, the thresholds, the lapse) — is the \textbf{Lyapunov-contraction mass}: the ledger deficit accumulated under the tilt contraction law [L1–L4 (Part 1): every tilt step loses $\gamma(8-\gamma)/16$ of the Lyapunov value; BM.51: the class-degenerate foundation], in its composite form $m_{\mathrm{rest}}=\Delta S_{\mathfrak L}/\mathfrak h$ [BM.35]. Def-H.39's registry (Part 3) names the species, and no unqualified \textit{mass} crosses them: the informational read agrees with $m_{\mathrm{rest}}$ as a formed equation [BM.35], $\neg$as a currency fusion; the pole read is the rest mass on the transport species, matched at leading order [(1) above]; the sector-route values are measure functionals — exact named prices, $\neg$masses [BM.30; BM.37]; the Ansatz-B class masses remain non-adopted templates [Def-H.28 (Part 3)]; and the numeric bridges (the $\beta$-coefficients; $\rho_{\mathrm B}$) are calibration slots, values $\neg$issued [§Cal-BM (C-BM.1); P5].

\end{theorem}

\begin{theorem}[{\texorpdfstring{BM.98 — The relating equation: the bridge in four exact forms, its flicker weighting, and the closure  in form}{BM.98 - The relating equation: the bridge in four exact forms, its flicker weighting, and the closure in form}}]\label{tag:BM-98}
\textbf{Thm.}

%dsp-level:0
\[
(1)\quad\textbf{the bridge in its derived count form:}
\]

%dsp-level:0
\[
\qquad \mathcal R\ :=\ \frac{m^2_{\mathrm{pole}}}{m^2_{\mathrm{hid}}}\ =\ \rho_{\mathrm B}\,\frac{D_S(\Theta_R)}{4\,E_{\mathrm{int}}^2},
\]

%dsp-level:0
\[
\qquad \rho_{\mathrm B}=\frac{\hat\mu'^2\mathfrak h^2}{(\ln2)^2}\in\mathfrak q^0
\]

%dsp-level:0
\[
\qquad [\mathrm{BM.69}(2);\ \mathrm{FT.54}(4);\ \mathrm{Census\ FT.63}(6)]:
\]

%dsp-level:0
\[
\text{six states, six exact algebraic factors}\ \tfrac{D_S}{4E^2},\ \text{four fields, six minimal polynomials}
\]

%dsp-level:0
\[
\qquad [\mathrm{Def\text{-}BM.14}(3);\ \mathrm{BM.82}(2);\ \text{machine re-cited}];
\]

%dsp-level:0
\[
(2)\quad\textbf{the flicker-weighted bridge:}
\]

%dsp-level:0
\[
\qquad E_{\mathrm{int}}\ \longrightarrow\ \bigl\langle E_{\mathrm{int}}\bigr\rangle_{\mathcal E}
\]

%dsp-level:0
\[
\qquad [\mathrm{BM.94}(2),(6),(7);\ \text{the integral replaces the snapshot}]:
\]

%dsp-level:0
\[
\textbf{unconditioned:}
\]

%dsp-level:0
\[
\qquad \mathcal R_w=\rho_{\mathrm B}\,\frac{D_S}{4\,w^2E^2}=\frac{\mathcal R}{w^2};
\]

%dsp-level:0
\[
\qquad \textbf{materialized:}
\]

%dsp-level:0
\[
\qquad \boxed{\ \mathcal R_{\mathrm{mat}}\ =\ \rho_{\mathrm B}\,\frac{D_S\,D^2}{4\,E^2}\ =\ \frac{\mathcal R}{w_{\mathrm{mat}}^2}\ =\ \rho_{\mathrm B}\,\frac{D_S\,D^2}{144}\ \ \text{at}\ \ E=E^{\mathrm{cryst}}=6\ }
\]

%dsp-level:0
\[
[\text{machine: one exact rational function per state}];
\]

%dsp-level:0
\[
\qquad \boxed{\ 1\ \le\ \frac{\mathcal R_{\mathrm{mat}}}{\mathcal R}\ =\ D^2\ <\ 9\ }
\]

%dsp-level:0
\[
\qquad [\mathrm{BM.94}(7):\ D\in[1,3);\ \text{machine}]:
\]

%dsp-level:0
\[
\boxed{\ \textbf{the flicker can multiply the bridge by at most an exact nine, and by exactly one at the sharp threshold}\ };
\]

%dsp-level:0
\[
\qquad \text{the lower endpoint is attained at}\ \varepsilon_q
\]

%dsp-level:0
\[
= 0,
\]

%dsp-level:0
\[
\ \text{where the snapshot is recovered exactly}\ [\mathrm{BM.94}(6),(7)];
\]

%dsp-level:0
\[
\text{the upper endpoint is open and is approached only in the fully dissolved limit, where}\ \bigl\langle E\bigr\rangle_{\mathrm{mat}}\to2\ —\ \text{the digon floor}\ [\mathrm{Census\ FT.28};\ \mathrm{BM.94}(7)];
\]

%dsp-level:0
\[
(3)\quad\textbf{the curvature form:}
\]

%dsp-level:0
\[
\qquad G=2\pi\,E_{\mathrm{int}}^2\ln2
\]

%dsp-level:0
\[
\qquad [\mathrm{BM.88}(3);\ \mathrm{BM.16}(1):\ \text{the calibration lock}]\implies E_{\mathrm{int}}^2=\frac{G}{2\pi\ln2}:
\]

%dsp-level:0
\[
\boxed{\ \mathcal R\ =\ \frac{\pi\ln2}{2}\cdot\frac{\rho_{\mathrm B}\,D_S}{G}\ }
\]

%dsp-level:0
\[
\qquad [\text{machine: one monomial, coefficient comparison}]
\]

%dsp-level:0
\[
\qquad \textbf{the bridge is inversely proportional to the curvature lock, with the exact coefficient}\ \tfrac{\pi\ln2}2;
\]

%dsp-level:0
\[
\textbf{flicker-weighted:}
\]

%dsp-level:0
\[
\qquad G^{\mathrm{eff}}:=2\pi\bigl\langle E\bigr\rangle^2_{\mathrm{mat}}\ln2=\frac{G}{D^2}=w_{\mathrm{mat}}^2\,G,
\]

%dsp-level:0
\[
\qquad \mathcal R_{\mathrm{mat}}=\frac{\pi\ln2}2\cdot\frac{\rho_{\mathrm B}D_S}{G^{\mathrm{eff}}}
\]

%dsp-level:0
\[
\qquad [\text{machine}];
\]

%dsp-level:0
\[
\textbf{the exact enclosure at}\ E^{\mathrm{cryst}}=6:
\]

%dsp-level:0
\[
\qquad G=72\,\pi\ln2
\]

%dsp-level:0
\[
\qquad [\mathrm{BM.17}(3)\text{'s lock value re-derived};\ \text{machine}]\implies\boxed{\ G^{\mathrm{eff}}\ \in\ \bigl(8\,\pi\ln2,\ \ 72\,\pi\ln2\bigr]\ }
\]

%dsp-level:0
\[
\textbf{the two endpoints are the digon's and}\ \Sigma_{12}\textbf{'s own curvature locks}
\]

%dsp-level:0
\[
\qquad [\mathrm{BM.17}(4)\text{'s family table};\ \text{machine:}\ 2\pi\cdot2^2\ln2=8\pi\ln2]:
\]

%dsp-level:0
\[
\qquad \text{the flicker moves the lock between two computed cells and never below the floor's};
\]

%dsp-level:0
\[
(4)\quad\textbf{the coupling form:}
\]

%dsp-level:0
\[
\qquad k=\frac{r\,E_{\mathrm{int}}\ln2}{\mathfrak h},
\]

%dsp-level:0
\[
\qquad 1-f=\frac1k
\]

%dsp-level:0
\[
\qquad [\mathrm{BM.88}(2)(\mathrm i),(\mathrm{ii});\ \mathrm{BM.88}(6)]\implies E_{\mathrm{int}}=\frac{\mathfrak h}{r\,(1-f)\ln2}:
\]

%dsp-level:0
\[
\boxed{\ \mathcal R\ =\ \frac{\hat\mu'^2\,D_S\,r^2}{4}\,(1-f)^2\ }
\]

%dsp-level:0
\[
\qquad [\text{machine: the substitution};\qquad[\hat\mu'^2r^2]=\mathfrak q\cdot\mathfrak q^{-1}=\mathfrak q^0\ [\mathrm{CT2.1};\ \mathrm{Def\text{-}H.37}(2);\ \text{machine: the exponent arithmetic}]];
\]

%dsp-level:0
\[
e^2=e_0^2\,N\,f\,\mathsf g(G)\qquad[\mathrm{BM.88}(4)]\implies1-f=1-\frac{e^2}{e_0^2N\,\mathsf g(G)}:
\]

%dsp-level:0
\[
\boxed{\ \mathcal R\ =\ \frac{\hat\mu'^2\,D_S\,r^2}{4}\Bigl(1-\frac{e^2}{e_0^2\,N\,\mathsf g(G)}\Bigr)^{\!2}\ }
\]

%dsp-level:0
\[
\qquad \textbf{the bridge as an exact function of the running coupling}
\]

%dsp-level:0
\[
\textbf{the energy-dependence of the bridge is the energy-dependence of the coupling:}
\]

%dsp-level:0
\[
\qquad \text{both run through the single derived argument}\ f;
\]

%dsp-level:0
\[
\mathcal R\ \ \text{strictly decreasing in}\ E_{\mathrm{int}},\quad\text{strictly decreasing in}\ f,\quad\text{strictly increasing in}\ k;
\]

%dsp-level:0
\[
\qquad e\ \ \text{strictly increasing in}\ E_{\mathrm{int}}
\]

%dsp-level:0
\[
\qquad [\mathrm{BM.88}(6);\ \text{machine: three monotonicity rows on exact rationals}]:
\]

%dsp-level:0
\[
\boxed{\ \text{at fixed}\ (r,\,N,\,\mathsf g,\,D_S),\ \ \mathcal R\ \ \textbf{is a strictly decreasing exact function of the running coupling}\ }
\]

%dsp-level:0
\[
\textbf{the flicker's action on the coupling:}
\]

%dsp-level:0
\[
\qquad k^{\mathrm{eff}}=w_{\mathrm{mat}}\,k^{\mathrm{cryst}}=\frac{k^{\mathrm{cryst}}}D,
\]

%dsp-level:0
\[
\qquad f^{\mathrm{eff}}=1-\frac{D}{k^{\mathrm{cryst}}}\ \le\ f^{\mathrm{cryst}},
\]

%dsp-level:0
\[
\qquad \mathsf g^{\mathrm{eff}}=\mathsf g\bigl(G^{\mathrm{eff}}\bigr)\ \ \text{at}\ \ G^{\mathrm{eff}}=\frac{G}{D^2}
\]

%dsp-level:0
\[
\qquad [\text{machine}]:
\]

%dsp-level:0
\[
\text{equality}\iff\varepsilon_q=0
\]

%dsp-level:0
\[
\qquad [\text{machine}]\implies\textbf{the flicker attenuates the running coupling and enhances the bridge, both by the same derived weight};
\]

%dsp-level:0
\[
\textbf{the vanishing locus is shifted and bounded:}
\]

%dsp-level:0
\[
\qquad f^{\mathrm{eff}}=0\iff r=D\,\Delta x_{\mathrm B}^{\mathrm{cryst}}\in\bigl[\Delta x_{\mathrm B}^{\mathrm{cryst}},\ 3\,\Delta x_{\mathrm B}^{\mathrm{cryst}}\bigr)
\]

%dsp-level:0
\[
\qquad [\mathrm{BM.88}(5);\ \mathrm{BM.94}(8);\ \text{machine}]:
\]

%dsp-level:0
\[
\boxed{\ \textbf{the coupling's zero is confined to within an exact factor of three of the crystallized Born width}\ }
\]

%dsp-level:0
\[
(5)\quad\textbf{the closure in form and the localized residue:}
\]

%dsp-level:0
\[
\qquad \textbf{every factor of}\ (1)\text{--}(4)\ \textbf{is derived exact algebra:}
\]

%dsp-level:0
\[
\text{the six}\ D_S\ \text{with their minimal polynomials}\ [\mathrm{Def\text{-}BM.14}(3)];
\]

%dsp-level:0
\[
\quad G\text{'s monomial and}\ f\text{'s rational function}\ [\mathrm{BM.88}(2),(3)];
\]

%dsp-level:0
\[
\quad N\text{'s ADM row}\ [\mathrm{BM.88}(2)(\mathrm{iii})];
\]

%dsp-level:0
\[
\quad w_{\mathrm{mat}},\ D,\ \text{the integer}\ 9\ [\mathrm{BM.94}(7)];
\]

%dsp-level:0
\[
\text{the coefficient}\ \tfrac{\pi\ln2}2\ (3);
\]

%dsp-level:0
\[
\qquad \text{the two computed lock endpoints}\ (3);
\]

%dsp-level:0
\[
\qquad \text{every monotonicity row and certificate}\ (2)\text{--}(4)
\]

%dsp-level:0
\[
\qquad [\text{machine}];
\]

%dsp-level:0
\[
\textbf{and one factor is not:}
\]

%dsp-level:0
\[
\qquad \boxed{\ \rho_{\mathrm B}=\frac{\hat\mu'^2\mathfrak h^2}{(\ln2)^2}\ \ \text{is the single undeclared dimensionless ratio}\ \ (\mathrm{R1})\ }
\]

%dsp-level:0
\[
\qquad [\mathrm{BM.69}(5);\ \mathrm{BM.82}(6);\ \mathrm{Cal\text{-}F}\ \text{packaging},\ \neg\text{added here}]:
\]

%dsp-level:0
\[
\boxed{\ \textbf{the relating equation is derived as an exact identity in the derived vocabulary plus}\ (\mathrm{R1}):\ \ \textbf{the equation is closed, the ratio is}\ \neg\textbf{derived}\ }
\]

%dsp-level:0
\[
\mathrm{BM.69}(4)\textbf{'s four grounds stand unweakened:}
\]

%dsp-level:0
\[
\qquad (\mathrm i)\ \text{the degree-}16\ \text{field against integer- and binary-valued reads}\ [\mathrm{FT.73}(2),(3)];
\]

%dsp-level:0
\[
\qquad (\mathrm{ii})\ \text{the two-torsion image reaching}\ 2\ \text{of}\ 6\ \text{states}\ [\mathrm{BM.67}(6)];
\]

%dsp-level:0
\[
(\mathrm{iii})\ \text{the calibration weight}\ {-}2\ \text{against}\ 0\ [\mathrm{FT.55}(1);\ \mathrm{BM.67}(8)];
\]

%dsp-level:0
\[
\qquad (\mathrm{iv})\ \text{the strength bar}\ [\mathrm{BM.36}(3);\ \mathrm{BM.61}(8)(\mathrm{O1})]
\]

%dsp-level:0
\[
\qquad [\text{machine: the four rows re-cited at their own loci}];
\]

%dsp-level:0
\[
\textbf{the residue list after this tier:}
\]

%dsp-level:0
\[
\qquad (\mathrm{R1})\ \rho_{\mathrm B};
\]

%dsp-level:0
\[
\quad (\mathrm{R2})\ \text{the loop realization};
\]

%dsp-level:0
\[
\quad (\mathrm{R3})\ \text{the flavour composition};
\]

%dsp-level:0
\[
\quad (\mathrm{R4})\ w(\varepsilon_q)\ \text{in general and the transport register's ensemble average};
\]

%dsp-level:0
\[
\quad (\mathrm{R5})\ \iota;
\]

%dsp-level:0
\[
(\mathrm{R6})\ \ \textbf{the reciprocal register's Jensen residue}\ \ \langle E\rangle_{\mathrm{mat}}\langle E^{-1}\rangle_{\mathrm{mat}}-1
\]

%dsp-level:0
\[
= \frac{18-27p+11p^2-2D^2}{2D^2}\ \ [\mathrm{BM.94}(8)]
\]

%dsp-level:0
\[
\qquad \textbf{named,}\ \neg\textbf{filled};
\]

%dsp-level:0
\[
\text{the two Cal-F packagings}\ (\hat\mu^2,\ \hat\mu'^2),
\]

%dsp-level:0
\[
\ \ e_0\ \ \text{and the scales}\ (\hat\mu_\nu,\ \ell_c^{(\nu)},\ \mu_\nu)\ \ \text{stay undeclared}
\]

%dsp-level:0
\[
\qquad [\mathrm{Def\text{-}H.37};\ \mathrm{EM.25}(7);\ \mathrm{P4};\ \mathrm{P1}];
\]

%dsp-level:0
\[
(6)\quad\textbf{the currency bar stands:}
\]

%dsp-level:0
\[
\qquad \mathcal R\in\mathfrak q^0\ \ \text{is a ratio},
\]

%dsp-level:0
\[
\ \neg\text{a conversion}\ —\ \neg\text{a clause of this tier adds a pole-currency term to an informational assembly}:
\]

%dsp-level:0
\[
[\Delta_g]=\mathfrak q\quad\text{against}\quad[\Delta m]=\mathfrak q^{+1/2}
\]

%dsp-level:0
\[
\qquad \implies
\]

%dsp-level:0
\[
\qquad \neg\text{addable}
\]

%dsp-level:0
\[
\qquad [\mathrm{BM.37}(2);\ \mathrm{BM.69}(3);\ \mathrm{BM.83}(3);\ \mathrm{BM.90}(5);\ \mathrm{CT2.1};\ \text{machine: the exponent arithmetic}];
\]

%dsp-level:0
\[
\textbf{the flicker weight attenuates the informational register alone and converts nothing into the defect register}
\]

%dsp-level:0
\[
\qquad [\mathrm{BM.90}(4);\ \mathrm{BM.93}(5);\ \text{machine:}\ w_{\mathrm{mat}}\text{'s symbol census carries no}\ D_S];
\]

%dsp-level:0
\[
(7)\quad\textbf{the no-import rows:}
\]

%dsp-level:0
\[
\qquad \neg\text{the observed neutron--proton mass ratio is imported}\ —\ \text{the ratio is derived and bounded at}\ \mathrm{BM.99}(5);
\]

%dsp-level:0
\[
\qquad \neg\text{the observed electric coupling is imported}\ —\ e,
\]

%dsp-level:0
\[
\ e_0,
\]

%dsp-level:0
\[
\ \varepsilon_{\mathrm{EM}}\ \text{stay symbolic}\ [\mathrm{EM.25}(7),(9);\ \mathrm{BM.88}(7);\ \mathrm{P5}];
\]

%dsp-level:0
\[
\neg\text{a measured}\ \nu\ \text{mass, mixing angle, endpoint or lifetime enters}\ [\mathrm{P5}];
\]

%dsp-level:0
\[
\qquad \neg\text{an}\ [\mathrm{EXT}]\ \text{datum};
\]

%dsp-level:0
\[
\qquad \text{machine: zero external tokens, zero decimal numerals in this tier};
\]

%dsp-level:0
\[
(8)\quad\textbf{status:}
\]

%dsp-level:0
\[
\qquad \textbf{derived:}\ (1)\text{--}(4),\ (6);
\]

%dsp-level:0
\[
\qquad \textbf{closed in form with one priced residue:}\ (5);
\]

%dsp-level:0
\[
\qquad \textbf{symbolic:}\ e_0,\ \mathsf g,\ N,\ \rho_{\mathrm B};
\]

%dsp-level:0
\[
\neg\text{a declaration},
\]

%dsp-level:0
\[
\quad\neg\text{a P4 addition},
\]

%dsp-level:0
\[
\quad\text{§Cal-BM stays empty},
\]

%dsp-level:0
\[
\quad\neg\text{a new dimensionful scale}\ [\mathrm{P1}],
\]

%dsp-level:0
\[
\quad\neg\text{a measured value}\ [\mathrm{P5}];
\]

%dsp-level:0
\[
\qquad\alpha_B:\ \neg\text{added},\ \neg\text{derived},\ \text{barred}\ [\mathrm{BM.61}(8)].
\]

\textbf{Req.} D2.9 (Parts 1–2); A §0.3, CT2.1 (Part 2); EM.25(7), EM.25(9) (Part 2); P1, P4, P5 (Parts 1–2); Def-H.37(2), FT.54(4), FT.55(1), FT.73(2), FT.73(3) (Part 3); Census FT.28, FT.63(6) (row-wise); Def-BM.4(1), Def-BM.14(3), BM.13(1), BM.16(1), BM.17(3), BM.17(4), BM.36(3), BM.37(2), BM.61(8), BM.67(6), BM.67(8), BM.69(2), BM.69(3), BM.69(4), BM.69(5), BM.82(2), BM.82(6), BM.83(3), BM.88(2), BM.88(3), BM.88(4), BM.88(5), BM.88(6), BM.88(7), BM.89(4), BM.90(4), BM.90(5), BM.93(5), BM.94(2), BM.94(6), BM.94(7), BM.94(8) (this part); Cal-F; exact Laurent-monomial arithmetic with $\mathfrak q$- and $\mathfrak h$-exponent censuses, exact rational-function substitution, coefficient comparison, monotonicity on exact rationals, symbol censuses [math; machine: \texttt{checks/\allowbreak{}verify\_\allowbreak{}H\_\allowbreak{}flicker\_\allowbreak{}integral\_\allowbreak{}restmass.py} §E].

\textbf{Pf.}

%dsp-level:0
\[
(1):\ \mathrm{BM.69}(2)\text{'s identity and}\ \mathrm{FT.54}(4)\text{'s bridge ratio re-cited};
\]

%dsp-level:0
\[
\qquad [\rho_{\mathrm B}]=[\hat\mu'^2\mathfrak h^2]=\mathfrak q\cdot\mathfrak q^{-1}=\mathfrak q^0\ [\mathrm{Def\text{-}H.37}(2);\ \text{machine}]
\]

%dsp-level:0
\[
(2):\ \mathcal R_{\mathrm{mat}}
\]

%dsp-level:0
\[
= \rho_{\mathrm B}D_S/(4\langle E\rangle_{\mathrm{mat}}^2)
\]

%dsp-level:0
\[
= \rho_{\mathrm B}D_S/(4\cdot36/D^2)
\]

%dsp-level:0
\[
= \rho_{\mathrm B}D_SD^2/144\ [\mathrm{BM.94}(6);\ \text{machine}];
\]

%dsp-level:0
\[
\qquad D\in[1,3)\ [\mathrm{BM.94}(7)]\implies D^2\in[1,9)\ [\text{machine}];
\]

%dsp-level:0
\[
\qquad D=1\iff p=1\iff\varepsilon_q=0\ [\text{machine}]
\]

%dsp-level:0
\[
(3):\ E^2
\]

%dsp-level:0
\[
= G/(2\pi\ln2)\implies\rho_{\mathrm B}D_S/(4E^2)
\]

%dsp-level:0
\[
= \rho_{\mathrm B}D_S\cdot2\pi\ln2/(4G)
\]

%dsp-level:0
\[
= (\pi\ln2/2)\rho_{\mathrm B}D_S/G\ [\text{machine: coefficient comparison}];
\]

%dsp-level:0
\[
\qquad G^{\mathrm{eff}}=2\pi(6/D)^2\ln2=72\pi\ln2/D^2=G/D^2\ [\text{machine}]
\]

%dsp-level:0
\[
G^{\mathrm{cryst}}=2\pi\cdot36\cdot\ln2=72\pi\ln2\ [\mathrm{BM.17}(3);\ \text{machine}];
\]

%dsp-level:0
\[
\qquad 2\pi\cdot2^2\ln2=8\pi\ln2\ [\mathrm{BM.17}(4)\text{'s digon row};\ \text{machine}];
\]

%dsp-level:0
\[
\qquad D\in[1,3)\implies G/D^2\in(8\pi\ln2,\,72\pi\ln2]\ [\text{machine}]
\]

%dsp-level:0
\[
(4):\ 1-f
\]

%dsp-level:0
\[
= 1/k
\]

%dsp-level:0
\[
= \mathfrak h/(rE\ln2)
\]

%dsp-level:0
\[
\implies E
\]

%dsp-level:0
\[
= \mathfrak h/(r(1-f)\ln2)\implies\rho_{\mathrm B}D_S/(4E^2)
\]

%dsp-level:0
\[
= \rho_{\mathrm B}D_Sr^2(1-f)^2(\ln2)^2/(4\mathfrak h^2)
\]

%dsp-level:0
\[
= \hat\mu'^2D_Sr^2(1-f)^2/4\ [\text{machine: the substitution and}\ \rho_{\mathrm B}(\ln2)^2/\mathfrak h^2=\hat\mu'^2]
\]

%dsp-level:0
\[
e^2=e_0^2Nf\mathsf g\implies f=e^2/(e_0^2N\mathsf g)\ [\text{machine}];
\]

%dsp-level:0
\[
\qquad \mathcal R\propto(1-f)^2\ \text{strictly decreasing in}\ f\ \text{on}\ f<1\ [\mathrm{BM.88}(2)(\mathrm{ii}):\ 0\le f<1;\ \text{machine}];
\]

%dsp-level:0
\[
\qquad f\ \text{strictly increasing in}\ E_{\mathrm{int}}\ [\mathrm{BM.88}(6)]\implies\mathcal R\ \text{strictly decreasing in}\ E_{\mathrm{int}}\ \text{and in}\ e\ [\text{machine: the composition of two monotone rows}]
\]

%dsp-level:0
\[
k^{\mathrm{eff}}=r\langle E\rangle_{\mathrm{mat}}\ln2/\mathfrak h=(r\cdot6\ln2/\mathfrak h)/D=k^{\mathrm{cryst}}/D\ [\text{machine}];
\]

%dsp-level:0
\[
\qquad f^{\mathrm{eff}}=1-D/k^{\mathrm{cryst}}\le1-1/k^{\mathrm{cryst}}=f^{\mathrm{cryst}}\iff D\ge1\ [\text{machine}];
\]

%dsp-level:0
\[
\qquad f^{\mathrm{eff}}=0\iff k^{\mathrm{cryst}}=D\iff r=D\,\Delta x_{\mathrm B}^{\mathrm{cryst}}\ [\text{machine}]
\]

%dsp-level:0
\[
(5):\ \text{the inventory rows assembled from}\ (1)\text{--}(4)\text{'s own tags}\ [\text{machine}];
\]

%dsp-level:0
\[
\qquad \mathrm{BM.69}(4)\text{'s four grounds re-cited at their loci};
\]

%dsp-level:0
\[
\qquad \langle E\rangle\langle E^{-1}\rangle-1=\tfrac{18-27p+11p^2}{2D^2}-1=\tfrac{18-27p+11p^2-2D^2}{2D^2}\ [\mathrm{BM.94}(8);\ \text{machine}]
\]

%dsp-level:0
\[
(6):\ [\Delta_g]=[\hat\mu'^2]=\mathfrak q,\quad[\Delta m]=\mathfrak q^{+1/2}\ [\mathrm{CT2.1};\ \text{machine}];
\]

%dsp-level:0
\[
\qquad w_{\mathrm{mat}}
\]

%dsp-level:0
\[
= 1/D,
\]

%dsp-level:0
\[
\ D
\]

%dsp-level:0
\[
= 3-3p+p^2,
\]

%dsp-level:0
\[
\ p
\]

%dsp-level:0
\[
= 1-4\varepsilon_q:\ \text{the symbol set is}\ \{\varepsilon_q\}\ \text{and carries no}\ D_S\ [\text{machine: the token census}]
\]

%dsp-level:0
\[
(7)\text{--}(8):\ \text{the standing-rule scans: external tokens, decimal numerals, unit tokens, registry tokens, declaration lines, the barred object}\ \alpha_B\ [\text{machine}].\ \blacksquare
\]

\textbf{Phys.} The bridge between the two currencies has one shape and four exact readings of it. Count the edges and it is the single undeclared packaging times the state's defect factor over four times the count squared. Read the count through the calibration lock and the same object is pi log two over two times the packaging times the defect factor over the curvature monomial — the bridge is inversely proportional to the lock, with the coefficient exact. Read the count through the envelope ratio and it becomes mu-prime squared times the defect factor times the radius squared over four, times the square of one minus the nearby-space fraction; substitute the coupling's own profile and the bridge is an exact function of the running coupling, strictly decreasing in it at fixed geometry, so the energy-dependence of the bridge is literally the energy-dependence of the coupling and nothing else. Put the flicker in and every form gains one factor: the bridge is multiplied by the square of the dilution denominator, which is bounded by an exact integer nine, attained only in a limit that is not a materialized universe, and equal to one exactly at the sharp threshold where the snapshot returns. The curvature lock rides between two computed cells, the digon's and the junction's own, and the coupling's zero is confined to within a factor of three of the crystallized width. Everything in those four forms is derived exact algebra — one number is not, and the four grounds that bar its derivation stand where they were proved. So the relating equation is closed as an equation: its form, its flicker weighting, its curvature and coupling readings and its bound are theorems, and its single residue is one dimensionless ratio that this tier names, prices and does not invent.

\end{theorem}

\begin{theorem}[{\texorpdfstring{BM.99 — The materialized splitting and the Q-value: the two-sided parameter-free bar, the tightened bound, and the rescoped dissolution row}{BM.99 - The materialized splitting and the Q-value: the two-sided parameter-free bar, the tightened bound, and the rescoped dissolution row}}]\label{tag:BM-99}
\textbf{Thm.}

%dsp-level:0
\[
(1)\quad\textbf{the materialized assembly:}
\]

%dsp-level:0
\[
\qquad \boxed{\ m_n\,\mathfrak h\ =\ 6\,w_{\mathrm{mat}}\ln2+\ln20,\qquad m_p\,\mathfrak h\ =\ 6\,w_{\mathrm{mat}}\ln2+\ln15+3\,N\,w_{\mathrm{mat}}\,\mathsf g\,\varepsilon_{\mathrm{EM}}\ln2\ }
\]

%dsp-level:0
\[
[\mathrm{BM.91}(2)\ \text{at}\ w\to w_{\mathrm{mat}};\ \mathrm{BM.94}(7),(8);\ \mathrm{Def\text{-}BM.13}(5);\ \mathrm{BM.89}(5);\ \text{machine: the prime-log assembly}\ \ln20=2\ln2+\ln5,\ \ln15=\ln3+\ln5];
\]

%dsp-level:0
\[
\boxed{\ \Delta m\,\mathfrak h\ =\ \ln\frac43-3\,N\,w_{\mathrm{mat}}\,\mathsf g\,\varepsilon_{\mathrm{EM}}\ln2\ }\qquad[\text{machine:}\ \ln20-\ln15=\ln\tfrac43];
\]

%dsp-level:0
\[
\textbf{the cancellation row stands:}
\]

%dsp-level:0
\[
\qquad E_{\mathrm{int}}\ \ \text{is class-shared}\ [\mathrm{BM.57}(4);\ \mathrm{Census\ FT.35}(2)]\implies\text{the crystallization weight cancels in the ledger difference and survives only in the EM term}
\]

%dsp-level:0
\[
\qquad [\text{machine: the coefficient census}];
\]

%dsp-level:0
\[
\textbf{the register scoping stands:}
\]

%dsp-level:0
\[
\qquad w_{\mathrm{mat}}\ \ \text{attenuates the edge-count register alone};
\]

%dsp-level:0
\[
\qquad \text{the transport register}\ \Sigma_{\mathrm{cl}}(Q)=\ln W(Q)\ \ \text{enters at its crystallized value under}\ (\mathrm{R4})
\]

%dsp-level:0
\[
\qquad [\mathrm{BM.91}(2);\ \mathrm{Def\text{-}BM.13}(4),(6);\ \mathrm{BM.89}(8)];
\]

%dsp-level:0
\[
(2)\quad\textbf{the critical locus, enclosed:}
\]

%dsp-level:0
\[
\qquad \varepsilon^{\mathrm{crit}}_{\mathrm{mat}}(N,\mathsf g)=\frac{\ln(4/3)}{3\,N\,w_{\mathrm{mat}}\,\mathsf g\,\ln2}=\frac{2-\log_23}{3\,N\,w_{\mathrm{mat}}\,\mathsf g}
\]

%dsp-level:0
\[
\qquad [\mathrm{BM.91}(4);\ \text{machine: solved exactly in the prime-log ring}]:
\]

%dsp-level:0
\[
\boxed{\ \frac{2-\log_23}{3\,N\,\mathsf g}\ \le\ \varepsilon^{\mathrm{crit}}_{\mathrm{mat}}\ <\ \frac{2-\log_23}{N\,\mathsf g}\ \subset\ \Bigl(\frac{2}{15\,N\mathsf g},\ \ \frac{5}{12\,N\mathsf g}\Bigr)\ }
\]

%dsp-level:0
\[
\qquad [\mathrm{BM.91}(4)\text{'s squeeze}\ \tfrac25<2-\log_23<\tfrac5{12}\ \text{on}\ 3^5=243<256=2^8\ \text{and}\ 3^{12}=531441>524288=2^{19};\ \mathrm{BM.94}(7);\ \text{machine}];
\]

%dsp-level:0
\[
\varepsilon^{\mathrm{crit}}_{\mathrm{mat}}\notin\mathbb Q\ \ \text{at every rational}\ (N,\mathsf g)>0
\]

%dsp-level:0
\[
\qquad [\mathrm{BM.84}(4);\ \mathrm{BM.91}(4):\ 2-\log_23\notin\mathbb Q\ \text{on the}\ 2^p\ne3^q\ \text{census};\ \text{machine}];
\]

%dsp-level:0
\[
(3)\quad\textbf{the two-sided parameter-free bar:}
\]

%dsp-level:0
\[
\qquad \boxed{\ N\,\mathsf g\,\varepsilon_{\mathrm{EM}}\ \le\ \frac18\qquad\Longrightarrow\qquad\Delta m>0\ }
\]

%dsp-level:0
\[
\qquad [\mathrm{BM.91}(5)\ \text{re-derived without}\ w:\ \ w_{\mathrm{mat}}\le1\ \text{and}\ \tfrac{2-\log_23}3>\tfrac18\iff2^{13}=8192>6561=3^8];
\]

%dsp-level:0
\[
\boxed{\ N\,\mathsf g\,\varepsilon_{\mathrm{EM}}\ \ge\ \frac5{12}\qquad\Longrightarrow\qquad\Delta m<0\ }
\]

%dsp-level:0
\[
\qquad [\text{machine:}\ w_{\mathrm{mat}}>\tfrac13\ \text{strictly}\implies Nw_{\mathrm{mat}}\mathsf g\,\varepsilon_{\mathrm{EM}}>\tfrac5{36},\ \ \text{and}\ \ \tfrac5{36}>\tfrac{2-\log_23}3\iff\tfrac5{12}>2-\log_23\iff3^{12}>2^{19}];
\]

%dsp-level:0
\[
\boxed{\ \textbf{the materialization filter removes the crystallization weight from both bars}\ }
\]

%dsp-level:0
\[
\qquad \textbf{the undecided band is exact and non-empty:}
\]

%dsp-level:0
\[
\qquad \frac18<N\,\mathsf g\,\varepsilon_{\mathrm{EM}}<\frac5{12}
\]

%dsp-level:0
\[
\qquad [\text{machine:}\ 12<40];
\]

%dsp-level:0
\[
\text{both bars are uniform over the lapse, the curvature composition and the whole admissible threshold range}\ \varepsilon_q\in[0,\tfrac14)
\]

%dsp-level:0
\[
\qquad [\mathrm{CT0.8};\ \text{machine}];
\]

%dsp-level:0
\[
(4)\quad\textbf{the trichotomy (total, exact, both regions certified by integers):}
\]

%dsp-level:0
\[
\qquad \boxed{\ \varepsilon_{\mathrm{EM}}<\varepsilon^{\mathrm{crit}}_{\mathrm{mat}}\iff m_n>m_p;\qquad=\iff m_n=m_p;\qquad>\iff m_n<m_p\ }
\]

%dsp-level:0
\[
[\text{machine:}\ \Delta m\ \text{strictly decreasing in}\ \varepsilon_{\mathrm{EM}}\ \text{at every admissible}\ (N,w_{\mathrm{mat}},\mathsf g)>0];
\]

%dsp-level:0
\[
\qquad \text{the opening region is inhabited by}\ \varepsilon_{\mathrm{EM}}
\]

%dsp-level:0
\[
= 0\ \text{and by}\ \varepsilon_{\mathrm{EM}}
\]

%dsp-level:0
\[
= \tfrac19\ \text{at}\ N\mathsf g
\]

%dsp-level:0
\[
= 1\ \ [3);
\]

%dsp-level:0
\[
\ \mathrm{BM.62}(8):\ 4^3
\]

%dsp-level:0
\[
= 64>54\qquad = 3^3\!\cdot\!2];
\]

%dsp-level:0
\[
\text{the closing region is inhabited by}\ \varepsilon_{\mathrm{EM}}=\tfrac5{12}\ \text{at}\ N\mathsf g=1
\]

%dsp-level:0
\[
\qquad [(3);\ \text{machine}];
\]

%dsp-level:0
\[
\qquad \text{the critical locus is exact, unique and irrational}\qquad \qquad [(2)];
\]

%dsp-level:0
\[
(5)\quad\textbf{the tightened universal ratio bound:}
\]

%dsp-level:0
\[
\qquad \frac{m_n}{m_p}=\frac{6\,w_{\mathrm{mat}}\ln2+\ln20}{6\,w_{\mathrm{mat}}\ln2+\ln15+3Nw_{\mathrm{mat}}\mathsf g\,\varepsilon_{\mathrm{EM}}\ln2}
\]

%dsp-level:0
\[
\qquad [\text{machine}]:
\]

%dsp-level:0
\[
\text{strictly decreasing in}\ \varepsilon_{\mathrm{EM}}\ \text{and in}\ w_{\mathrm{mat}}
\]

%dsp-level:0
\[
\qquad [\mathrm{BM.91}(7);\ \text{machine: the two cross-multiplied sign rows}]\ \ \text{and}\ \ \tfrac13<w_{\mathrm{mat}}\le1
\]

%dsp-level:0
\[
\qquad [\mathrm{BM.94}(7)]:
\]

%dsp-level:0
\[
\boxed{\ \frac{\ln1280}{\ln960}\ \le\ \frac{m_n}{m_p}\bigg|_{\varepsilon_{\mathrm{EM}}=0}\ <\ \frac{\ln80}{\ln60}\ <\ \frac{\ln20}{\ln15}\ <\ \frac{10}9\ }
\]

%dsp-level:0
\[
\qquad [\text{machine:}\ 6\ln2+\ln20=\ln1280,\quad2\ln2+\ln20=\ln80,\quad2\ln2+\ln15=\ln60]:
\]

%dsp-level:0
\[
\frac{\ln80}{\ln60}\in\Bigl(\frac{16}{15},\ \frac{15}{14}\Bigr)
\]

%dsp-level:0
\[
\qquad [\text{machine:}\ 80^{15}>60^{16}\iff2^{28}=268435456>215233605=3^{16}\!\cdot\!5;\qquad80^{14}<60^{15}\iff2^{26}=67108864<71744535=3^{15}\!\cdot\!5];
\]

%dsp-level:0
\[
\frac{15}{14}<\frac{11}{10}<\frac{\ln20}{\ln15}
\]

%dsp-level:0
\[
\qquad [\text{machine:}\ 15\cdot10<14\cdot11;\qquad20^{10}=10240000000000>8649755859375=15^{11}]:
\]

%dsp-level:0
\[
\boxed{\ 1\ <\ \frac{m_n}{m_p}\ <\ \frac{\ln80}{\ln60}\ <\ \frac{15}{14}\qquad\text{on the materialized read at every admissible}\ (\varepsilon_q,\,N,\,\mathsf g,\,\varepsilon_{\mathrm{EM}})\ \text{on the sub-critical branch}\ }
\]

%dsp-level:0
\[
\textbf{an exact tightening of}\ \mathrm{BM.91}(7)\textbf{'s}\ \tfrac{10}9\ \textbf{by the filter's lower bound on the weight}
\]

%dsp-level:0
\[
\qquad [\text{machine:}\ \tfrac{15}{14}<\tfrac{10}9\iff135<140];
\]

%dsp-level:0
\[
\textbf{the incommensurability hardens:}
\]

%dsp-level:0
\[
\qquad \frac{m_n}{m_p}=\frac43\iff\varepsilon^{*}_{\mathrm{EM}}=\frac{(6-6w_{\mathrm{mat}})\ln2-\ln405}{12\,N\,w_{\mathrm{mat}}\,\mathsf g\,\ln2}<0
\]

%dsp-level:0
\[
\qquad [\mathrm{BM.91}(8);\ \text{machine:}\ (6-6w_{\mathrm{mat}})\ln2\le6\ln2=\ln64<\ln405\ \text{on}\ 64<405];
\]

%dsp-level:0
\[
\frac{m_n}{m_p}\notin\mathbb Q\ \ \text{at every rational}\ (w_{\mathrm{mat}},N,\mathsf g,\varepsilon_{\mathrm{EM}})
\]

%dsp-level:0
\[
\qquad [\mathrm{BM.83}(2)\text{'s lemma: the}\ \ln3\ \text{coefficient is}\ 0\ \text{in the numerator and}\ 1\ \text{in the denominator};\ \text{machine}];
\]

%dsp-level:0
\[
(6)\quad\textbf{the Q-value on the materialized informational read:}
\]

%dsp-level:0
\[
\qquad Q_\beta\,\mathfrak h=\ln\frac43-3Nw_{\mathrm{mat}}\mathsf g\,\varepsilon_{\mathrm{EM}}\ln2-F_{\ell\nu}\ \ \le\ \ln\tfrac43-2\ln2\ =\ -\ln3\ <\ 0
\]

%dsp-level:0
\[
[\mathrm{BM.92}(1);\ \mathrm{BM.70}(2),(3):\ F_{\ell\nu}\ge2\ln2;\ \text{machine}]
\]

%dsp-level:0
\[
\qquad \textbf{uniformly in}\ (\varepsilon_q,\,N,\,\mathsf g,\,\varepsilon_{\mathrm{EM}})\ —\ \textbf{the closure survives the filter unchanged};
\]

%dsp-level:0
\[
\textbf{the ratio form tightens:}
\]

%dsp-level:0
\[
\qquad \boxed{\ -\frac{\ln3}{\ln1280}\ \le\ \frac{Q_\beta}{m_n}\ \le\ -\frac{\ln3}{6\,w_{\mathrm{mat}}\ln2+\ln20}\ <\ -\frac{\ln3}{\ln80}\ }
\]

%dsp-level:0
\[
\qquad \text{with}
\]

%dsp-level:0
\[
\qquad \boxed{\ -\frac{\ln3}{\ln80}\in\Bigl(-\frac13,\,-\frac14\Bigr)\ }
\]

%dsp-level:0
\[
[\text{machine:}\ 3^4=81>80\ \text{and}\ 80>27]\qquad\textbf{an exact tightening of}\ \mathrm{BM.92}(2)\textbf{'s}\ \bigl(-\tfrac25,-\tfrac13\bigr)\ \text{endpoint};
\]

%dsp-level:0
\[
\textbf{the margin ladder stands:}
\]

%dsp-level:0
\[
\qquad \Sigma_{\mathrm B}-2\ln2=-\ln3,
\]

%dsp-level:0
\[
\qquad \Sigma_{\mathrm B}-4\ln2=-\ln12,
\]

%dsp-level:0
\[
\qquad \Sigma_{\mathrm B}-5\ln2=-\ln24
\]

%dsp-level:0
\[
\qquad [\mathrm{BM.70}(3);\ \mathrm{BM.92}(2);\ \text{machine:}\ 4<12<48<96];
\]

%dsp-level:0
\[
(7)\quad\textbf{the rescoped dissolution row:}
\]

%dsp-level:0
\[
\qquad \mathrm{BM.91}(4)\text{'s row}\ \bigl(\varepsilon^{\mathrm{crit}}\to\infty\ \text{as}\ \varepsilon_q\to\tfrac14^{-},\ \text{``the fully dissolved limit never closes the splitting''}\bigr)\ \textbf{is rescoped,}\ \neg\textbf{retracted}:
\]

%dsp-level:0
\[
\text{at}\ \varepsilon_q\to\tfrac14^{-}\ \ \text{the unconditioned ensemble's weight concentrates on the}\ E=0\ \ \text{member},
\]

%dsp-level:0
\[
\quad P(E{=}0)=(1-p)^3\to1
\]

%dsp-level:0
\[
\qquad [\mathrm{BM.94}(4),(5);\ \text{machine}];
\]

%dsp-level:0
\[
\text{that member supports no closed chain and no terminus}
\]

%dsp-level:0
\[
\qquad [\mathrm{FT.144}(4);\ \mathrm{Def\text{-}H.4}]\implies\boxed{\ \textbf{the fully dissolved limit is outside the materialization domain}\ }
\]

%dsp-level:0
\[
\textbf{on the materialized read the splitting is closable at a finite coupling:}
\]

%dsp-level:0
\[
\qquad \varepsilon^{\mathrm{crit}}_{\mathrm{mat}}<\frac{5}{12\,N\,\mathsf g}
\]

%dsp-level:0
\[
\qquad [(2),(3);\ \text{machine}];
\]

%dsp-level:0
\[
\qquad \text{the unconditioned row stands at its own locus as the}\ (\mathrm{R4})\ \text{general-weight statement};
\]

%dsp-level:0
\[
(8)\quad\textbf{status:}
\]

%dsp-level:0
\[
\qquad \textbf{derived:}\ (1)\text{--}(7)\ \text{on the marked pair specialization};
\]

%dsp-level:0
\[
\qquad \textbf{rescoped,}\ \neg\textbf{retracted:}\ \mathrm{BM.91}(2),
\]

%dsp-level:0
\[
(4),
\]

%dsp-level:0
\[
(7),
\]

%dsp-level:0
\[
(8),
\]

%dsp-level:0
\[
\ \mathrm{BM.92}(1),
\]

%dsp-level:0
\[
(2)\ \text{stand on}\ w\ \text{at their own loci and gain the materialized variants here};
\]

%dsp-level:0
\[
\mathrm{BM.84},\ \mathrm{BM.85},\ \mathrm{BM.77}(5),(6)\ \text{untouched};
\]

%dsp-level:0
\[
\qquad \neg\text{a declaration},
\]

%dsp-level:0
\[
\quad\neg\text{a P4 addition},
\]

%dsp-level:0
\[
\quad\neg\text{a new dimensionful scale}\ [\mathrm{P1}],
\]

%dsp-level:0
\[
\quad\neg\text{a measured value}\ [\mathrm{P5}],
\]

%dsp-level:0
\[
\quad\text{§Cal-BM stays empty};
\]

%dsp-level:0
\[
\neg\text{the observed neutron--proton ratio imported}\ —\ \text{the bound}\ (5)\ \text{is derived from the filter and the prime-log assembly alone};\qquad
\]

\textbf{Req.} D0.4 (Parts 1–2); CT0.8, CT2.1 (Part 1, Part 2); P1, P4, P5 (Parts 1–2); Def-H.4, FT.144(4) (Part 3); Census FT.35(2) (row-wise); Def-BM.13(4), Def-BM.13(5), Def-BM.13(6), BM.57(4), BM.62(8), BM.70(2), BM.70(3), BM.77(5), BM.77(6), BM.83(2), BM.84(4), BM.85, BM.89(5), BM.89(8), BM.91(2), BM.91(4), BM.91(5), BM.91(7), BM.91(8), BM.92(1), BM.92(2), BM.94(4), BM.94(5), BM.94(6), BM.94(7), BM.94(8) (this part); exact prime-log arithmetic with integer power certificates, exact rational function rows, monotonicity by cross-multiplication, token censuses [math; machine: \texttt{checks/\allowbreak{}verify\_\allowbreak{}H\_\allowbreak{}flicker\_\allowbreak{}integral\_\allowbreak{}restmass.py} §F].

\textbf{Pf.}

%dsp-level:0
\[
(1):\ \mathrm{BM.91}(2)\text{'s assembly with}\ w\to w_{\mathrm{mat}}\ \text{and}\ \langle m_{\mathrm{hid}}\rangle_{\mathrm{mat}}=6w_{\mathrm{mat}}\ln2/\mathfrak h\ [\mathrm{BM.94}(8);\ \mathrm{BM.89}(5)];
\]

%dsp-level:0
\[
\qquad \ln20-\ln15=\ln\tfrac43\ [\text{machine}];
\]

%dsp-level:0
\[
\qquad \text{the}\ w\text{-coefficient census in the two rest inputs and in}\ \Delta m\ [\text{machine}]
\]

%dsp-level:0
\[
(2):\ \varepsilon^{\mathrm{crit}}
\]

%dsp-level:0
\[
= \ln(4/3)/(3Nw_{\mathrm{mat}}\mathsf g\ln2)
\]

%dsp-level:0
\[
= (2-\log_23)/(3Nw_{\mathrm{mat}}\mathsf g)\ [\mathrm{BM.91}(4)\ \text{re-cited};\ \text{machine}];
\]

%dsp-level:0
\[
\qquad \tfrac13<w_{\mathrm{mat}}\le1\ [\mathrm{BM.94}(7)]\implies\text{the two-sided enclosure}\ [\text{machine}];
\]

%dsp-level:0
\[
\qquad \tfrac25<2-\log_23<\tfrac5{12}\ \text{on}\ 3^5<2^8\ \text{and}\ 3^{12}>2^{19}\ [\mathrm{BM.91}(4)\ \text{re-cited}]
\]

%dsp-level:0
\[
(3):\ w_{\mathrm{mat}}
\]

%dsp-level:0
\[
\le 1
\]

%dsp-level:0
\[
\implies Nw_{\mathrm{mat}}\mathsf g\,
\]

%dsp-level:0
\[
\varepsilon_{\mathrm{EM}}
\]

%dsp-level:0
\[
\le N\mathsf g\,
\]

%dsp-level:0
\[
\varepsilon_{\mathrm{EM}}\le\tfrac18<\tfrac{2-\log_23}3\ [\text{machine:}\ 8(2-\log_23)>3\iff13>8\log_23\iff2^{13}>3^8]\implies\Delta m>0
\]

%dsp-level:0
\[
w_{\mathrm{mat}}>\tfrac13\ \text{and}\ N\mathsf g\,\varepsilon_{\mathrm{EM}}\ge\tfrac5{12}\implies Nw_{\mathrm{mat}}\mathsf g\,\varepsilon_{\mathrm{EM}}>\tfrac5{36};
\]

%dsp-level:0
\[
\qquad \tfrac5{36}>\tfrac{2-\log_23}3\iff\tfrac5{12}>2-\log_23\iff\log_23>\tfrac{19}{12}
\]

%dsp-level:0
\[
\iff 3^{12}
\]

%dsp-level:0
\[
= 531441>524288
\]

%dsp-level:0
\[
= 2^{19}\ [\text{machine}]\implies\Delta m<0
\]

%dsp-level:0
\[
\tfrac18<\tfrac5{12}\iff12<40\ [\text{machine}]
\]

%dsp-level:0
\[
(4):\ \text{the coefficient of}\ \varepsilon_{\mathrm{EM}}\ \text{in}\ \Delta m\ \text{is}\ -3Nw_{\mathrm{mat}}\mathsf g\ln2<0\ [\text{machine}];
\]

%dsp-level:0
\[
\qquad \varepsilon_{\mathrm{EM}}=\tfrac19:\ N\mathsf g\,\varepsilon_{\mathrm{EM}}=\tfrac19\le\tfrac18\implies\Delta m>0\ \text{by}\ (3);
\]

%dsp-level:0
\[
\qquad \varepsilon_{\mathrm{EM}}=\tfrac5{12}\ \text{at}\ N\mathsf g=1\implies\Delta m<0\ \text{by}\ (3)
\]

%dsp-level:0
\[
(5):\ \text{the ratio's monotonicity in}\ \varepsilon_{\mathrm{EM}}\ \text{and in}\ w_{\mathrm{mat}}\ \text{re-cited}\ [\mathrm{BM.91}(7);\ \text{machine}];
\]

%dsp-level:0
\[
\qquad 6w_{\mathrm{mat}}\ln2\in(2\ln2,\,6\ln2]\ [\mathrm{BM.94}(7)]\implies\text{the four-endpoint chain}\ [\text{machine}]
\]

%dsp-level:0
\[
80^{15}>60^{16}\iff2^{60}5^{15}>2^{32}3^{16}5^{16}\iff2^{28}>3^{16}\cdot5:\quad268435456>215233605\ [\text{machine}];
\]

%dsp-level:0
\[
\qquad 80^{14}<60^{15}\iff2^{26}<3^{15}\cdot5:\quad67108864<71744535\ [\text{machine}]
\]

%dsp-level:0
\[
\tfrac{15}{14}<\tfrac{11}{10}\iff150<154;
\]

%dsp-level:0
\[
\qquad 20^{10}>15^{11}\ [\mathrm{BM.91}(7)\ \text{re-cited}];
\]

%dsp-level:0
\[
\qquad \tfrac{15}{14}<\tfrac{10}9\iff135<140\ [\text{machine}]
\]

%dsp-level:0
\[
(6):\ \mathrm{BM.92}(1)\text{'s uniform closure re-cited with}\ w\to w_{\mathrm{mat}}\ [\text{machine}];
\]

%dsp-level:0
\[
\qquad -\ln3/\ln80\in(-\tfrac13,-\tfrac14)\iff3\ln3>\ln80\ \text{and}\ 4\ln3<\ln80\iff81>80\ \text{and}\ 80>27\ [\text{machine}];
\]

%dsp-level:0
\[
\qquad \text{the three margin identities re-certified}\ [\text{machine}]
\]

%dsp-level:0
\[
(7):\ (1-p)^3\to1\ \text{as}\ p\to0^+\ [\text{machine}];
\]

%dsp-level:0
\[
\qquad |E_{\mathcal R}|=0\implies\neg\exists\ \text{a terminus}\ [\mathrm{FT.144}(4)\ \text{re-cited}];
\]

%dsp-level:0
\[
\qquad \varepsilon^{\mathrm{crit}}_{\mathrm{mat}}<(2-\log_23)/(N\mathsf g)<\tfrac5{12N\mathsf g}\ [(2);\ \text{machine}];
\]

%dsp-level:0
\[
\qquad (8):\ \text{the standing-rule scans}\ [\text{machine}].\ \blacksquare
\]

\textbf{Phys.} Filter the ensemble to the configurations that can carry a composite and three rows change shape. The assembly keeps its form — six times the weight times log two, plus the class's own transport entropy, less the electromagnetic term — but the weight is now one over one plus u plus u squared, and it cannot fall below a third. That single bound does real work. The critical coupling is squeezed into an exact window between two fifteenths and five twelfths over the lapse and the curvature composition, and the two bars that decide the sign of the splitting lose the crystallization weight entirely: an eighth or less of the coupling product forces the neutron heavier, five twelfths or more forces it lighter, and the undecided band between them is exact and inhabited. The ratio bound tightens from ten ninths to the log of eighty over the log of sixty, itself enclosed between sixteen fifteenths and fifteen fourteenths by two integer power certificates, and the Q-value's ratio form tightens from a third to a quarter by eighty-one against eighty. The dissolution row is rescoped rather than retracted: at the far end of the threshold range the unconditioned ensemble piles its weight onto the empty configuration, which supports no chain and no terminus, so that limit is not a materialized universe at all — and inside the materialization domain the splitting can always be closed by a finite coupling. Nothing measured entered: the ratio is derived from the filter and the prime-log assembly, the coupling stays symbolic, and the closure of the decay on the informational currency survives the filter exactly as it survived the ensemble.

\end{theorem}

\medskip\noindent\textit{\texorpdfstring{BM.100 — The status of the flicker integral and the rest-mass read: adjudication and scope}{BM.100 - The status of the flicker integral and the rest-mass read: adjudication and scope}.}\ Structural-constraints record: see Appendix B, Section B.1.29 of the Addendum.\par\medskip

\phantomsection
\section[Tier BM.20 — Six Objective Resolutions: the Energy-Scale…]{\texorpdfstring{Tier BM.20 — Six Objective Resolutions: the Energy-Scale Graph, the Dimensional Bridge, the Eigenvalue Composition, the Curvature Composition, the General Crystallization Weight, and the Kind Structure of the Arrow}{Tier BM.20 - Six Objective Resolutions: the Energy-Scale Graph, the Dimensional Bridge, the Eigenvalue Composition, the Curvature Composition, the General Crystallization Weight, and the Kind Structure of the Arrow}}

\textbf{Locus.} This tier takes six named objectives and gives each its exact adjudication: one graph, one census with one declared slot, one referent table, one fixed point, one measure-free bound, and one kind structure. The Q-value becomes a function of the energy scale, its forward and inverse reads are shown to be one graph reflected, and its crossing is a derived strength enclosed by two integers — with the crossing \textit{count} depending on the priced curvature composition ($\mathrm{BM.101}$). The dimension surface is censused against the single information flux, the seven rows are derived or adjudicated, and the two bridges are separated: the dimensional one is $\mathfrak h$, the numerical one is $\rho_{\mathrm B}$, whose status is declared as the §II.4 addendum A-BM.9 with no value issued ($\mathrm{BM.102}$). The three-constituent assembly's eigenvalue equations are derived where they are formed — the family cubic and the binomial characteristic polynomial — and the tensor-product referent is censured where it is not, so the composition question and the octet question are both adjudicated $\neg$formed ($\mathrm{BM.103}$). The curvature composition's self-consistency condition is formed exactly, solved for all three priced candidates, and shown to select none, because the lock's own freedom census excludes the back-reaction ($\mathrm{BM.104}$). The crystallization weight is derived measure-free and for an arbitrary threshold family, with the same exact bounds, and the defect-conditioned version is refuted ($\mathrm{BM.105}$). The spin label and the tick arrow are derived, and the exchange structure that a statistics requires is censured: deficit 11 stands with its residue named ($\mathrm{BM.106}$). $\mathrm{BM.107}$ adjudicates.

\textbf{Standing rows.} $\alpha_B$ stays barred on $\mathrm{BM.61}(8)$'s four grounds. No clause of $\mathrm{BM.1}$–$\mathrm{BM.100}$ is retracted, replaced or renumbered: every row this tier touches is re-scoped at its own locus and recorded in $\mathrm{BM.101}(8)$, $\mathrm{BM.104}(8)$, $\mathrm{BM.105}(7)$ and $\mathrm{BM.107}(5)$, in $\mathrm{BM.77}(7)$'s pattern. No measured coupling, mass or ratio enters any derivation $[\mathrm{P5}]$: the crossing strength of $\mathrm{BM.101}(3)$ is a derived symbolic value and is $\neg$compared to any external number; the neutron–proton ratio stays derived and bounded $[\mathrm{BM.99}(5)]$. No right-ray matter line and no seesaw-type mechanism is suggested, added or required $[\mathrm{BM.96}(4),(6)\ \text{re-cited},\ \neg\text{extended}]$; $\neg$an external mixing matrix is imported — the mixing rows are reads of the derived $16$-arena $[\mathrm{FT.142}(2),(5);\ \mathrm{FT.145}(3)]$. One dimension only: units are powers of $\mathfrak q$ $[\mathrm{D2.9};\ \mathrm{CT2.1}]$. Local notation, fixed for the whole tier: $E_1:=\mathfrak h/(r\ln2)$ the envelope-unit count; $x:=E/E_1=k$ the scale ratio $[\mathrm{BM.88}(2)(\mathrm i)]$; $\beta:=2\pi\ln2\,E_1^2=2\pi\mathfrak h^2/(r^2\ln2)\in\mathfrak q^0$; $f=1-\tfrac1x$; $G=\beta x^2$; $\varepsilon_{\mathrm{eff}}:=Nf\,\mathsf g\,\varepsilon_{\mathrm{EM}}=e(r,E,N)^2/(4\pi)$; $\varepsilon_\ast(w)$ the crossing strength; $c:=\varepsilon_\ast/(N\varepsilon_{\mathrm{EM}})$; $a:=\tfrac12Nw\,\varepsilon_{\mathrm{EM}}Q^2$; $p_i,u_i,R,S$ the family data of $\mathrm{BM.105}(3)$; $[\cdot]$ the $\mathfrak q$-exponent map $[\mathrm{CT2.1};\ \mathrm{D2.9}]$. The symbols $w$, $D$, $p$, $u$, $\Delta_g$, $M_0$, $s$, $\mathcal R$, $\rho_{\mathrm B}$, $D_S$ keep their $\mathrm{BM.94}$–$\mathrm{BM.99}$ meanings; the italic $g$ keeps its $\mathrm{BM.39}(1)$ meaning and $q$ its end-register meaning $[\mathrm{D3.4};\ \mathrm{EM.9}]$ at every occurrence.

\begin{theorem}[{\texorpdfstring{BM.101 — The energy-scale graph: the two reads, the derived crossing strength, and the composition-dependent crossing count}{BM.101 - The energy-scale graph: the two reads, the derived crossing strength, and the composition-dependent crossing count}}]\label{tag:BM-101}
\textbf{Thm.}

%dsp-level:0
\[
(1)\quad\textbf{the two reads are one graph reflected:}
\]

%dsp-level:0
\[
\qquad Q_\beta^{\mathrm{fwd}}(E)\ :=\ m_n^{\mathrm{rest}}(E)-m_p^{\mathrm{rest}}(E)-m_\ell^{\mathrm{rest}}-m_\nu^{\mathrm{rest}},
\]

%dsp-level:0
\[
\qquad Q_\beta^{\mathrm{inv}}(E)\ :=\ m_p^{\mathrm{rest}}(E)+m_\ell^{\mathrm{rest}}+m_\nu^{\mathrm{rest}}-m_n^{\mathrm{rest}}(E)
\]

%dsp-level:0
\[
\boxed{\ Q_\beta^{\mathrm{inv}}(E)\ =\ -\,Q_\beta^{\mathrm{fwd}}(E)\ \ \text{identically, at every scale, on every read}\ }
\]

%dsp-level:0
\[
\qquad [\text{machine: one negation identity};\ \mathrm{BM.97}(2),(4)\ \text{re-cited}]:
\]

%dsp-level:0
\[
\textbf{the inverse reaction opens exactly where the forward one closes, and their common boundary is one locus}\ —\ \neg\text{two independent graphs};
\]

%dsp-level:0
\[
(2)\quad\textbf{the scale-dependent splitting (the graph's spine):}
\]

%dsp-level:0
\[
\qquad \varepsilon_{\mathrm{eff}}(r,E,N)\ :=\ N\,f\!\Bigl(\frac{r}{\Delta x_{\mathrm B}(E)}\Bigr)\,\mathsf g\bigl(G(E)\bigr)\,\varepsilon_{\mathrm{EM}}\ =\ \frac{e(r,E,N)^2}{4\pi}
\]

%dsp-level:0
\[
\qquad [\mathrm{BM.88}(4);\ \mathrm{Def\text{-}BM.12}\ \text{local notation};\ \text{machine: the monomial identity}]:
\]

%dsp-level:0
\[
\boxed{\ \Delta m(E)\,\mathfrak h\ =\ \ln\frac43\ -\ 3\,w\,\ln2\cdot\varepsilon_{\mathrm{eff}}(r,E,N)\ }
\]

%dsp-level:0
\[
\qquad [\mathrm{BM.91}(2);\ \mathrm{BM.99}(1)\ \text{at the full profile instead of its far-field locus}\ f=1;\ \text{machine}];
\]

%dsp-level:0
\[
\textbf{the far-field rescope:}
\]

%dsp-level:0
\[
\qquad f\to1\ \ \text{recovers}\ \ \Delta m\,\mathfrak h=\ln\tfrac43-3Nw\,\mathsf g\,\varepsilon_{\mathrm{EM}}\ln2\ \ \text{exactly}
\]

%dsp-level:0
\[
\qquad [\mathrm{BM.88}(5);\ \mathrm{BM.91}(2);\ \mathrm{BM.99}(1):\ \text{every row stands at that locus}];
\]

%dsp-level:0
\[
\textbf{the informational Q-value on the graph:}
\]

%dsp-level:0
\[
\qquad Q_\beta^{\mathrm{fwd}}\,\mathfrak h=\Delta m(E)\,\mathfrak h-F_{\ell\nu},
\]

%dsp-level:0
\[
\qquad F_{\ell\nu}\ge2\ln2
\]

%dsp-level:0
\[
\qquad [\mathrm{BM.70}(2),(3);\ \mathrm{BM.92}(1)]:
\]

%dsp-level:0
\[
\boxed{\ Q_\beta^{\mathrm{fwd}}\,\mathfrak h\ \le\ -\ln3\ <\ 0\ \ \text{and}\ \ Q_\beta^{\mathrm{inv}}\,\mathfrak h\ \ge\ \ln3\ >\ 0\qquad\text{at every scale, uniformly in}\ (r,\,E,\,\varepsilon_q,\,N,\,\mathsf g,\,\varepsilon_{\mathrm{EM}})\ }
\]

%dsp-level:0
\[
[\text{machine: the sign census on a rational lattice}]
\]

%dsp-level:0
\[
\qquad \textbf{on the informational currency the forward channel is closed and the inverse channel is open at every energy scale, with the exact margin}\ \ln3;
\]

%dsp-level:0
\[
(3)\quad\textbf{the derived crossing strength:}
\]

%dsp-level:0
\[
\qquad \Delta m(E)=0\iff\varepsilon_{\mathrm{eff}}(r,E,N)=\varepsilon_\ast(w):
\]

%dsp-level:0
\[
\qquad \boxed{\ \varepsilon_\ast(w)\ :=\ \frac{\ln(4/3)}{3\,w\ln2}\ =\ \frac{2-\log_23}{3\,w},\qquad e_\ast^2\ :=\ 4\pi\,\varepsilon_\ast(w)\ =\ \frac{4\pi\,(2-\log_23)}{3\,w}\ }
\]

%dsp-level:0
\[
[\mathrm{BM.91}(4);\ \text{machine: exact in the}\ \log_2\ \text{ring}];\qquad\textbf{the enclosure:}\qquad w=w_{\mathrm{mat}}\in\bigl(\tfrac13,1\bigr]\qquad[\mathrm{BM.94}(7)]:
\]

%dsp-level:0
\[
\boxed{\ \frac{2-\log_23}3\ \le\ \varepsilon_\ast\ <\ 2-\log_23\qquad\text{and}\qquad\frac18\ <\ \varepsilon_\ast\ <\ \frac5{12}\ }
\]

%dsp-level:0
\[
\qquad [\text{machine:}\ 2^{13}=8192>6561=3^8;\quad3^{12}=531441>524288=2^{19};\quad3^5=243<256=2^8]:
\]

%dsp-level:0
\[
\boxed{\ \textbf{the two endpoints of}\ \mathrm{BM.99}(3)\textbf{'s two-sided parameter-free bar are exactly the enclosure of the crossing strength}\ }
\]

%dsp-level:0
\[
\qquad [\text{machine: the four integers re-certified}];
\]

%dsp-level:0
\[
e_\ast^2\in\bigl(\tfrac\pi2,\ \tfrac{5\pi}3\bigr)
\]

%dsp-level:0
\[
\qquad [\text{machine}];
\]

%dsp-level:0
\[
\qquad \textbf{no measured coupling is compared to}\ \varepsilon_\ast\ \text{or}\ e_\ast\ \ [\mathrm{P5};\ \mathrm{EM.25}(7),(9);\ \text{the quarantined audit rows stay quarantined}];
\]

%dsp-level:0
\[
(4)\quad\textbf{the crossing locus is a level set, and}\ e_0\ \textbf{places it:}
\]

%dsp-level:0
\[
\qquad \boxed{\ \mathcal L_\ast\ :=\ \bigl\{(r,E,N):\ \ e(r,E,N)^2=e_\ast^2\bigr\}\ =\ \bigl\{(r,E,N):\ \varepsilon_{\mathrm{eff}}=\varepsilon_\ast(w)\bigr\}\ }
\]

%dsp-level:0
\[
e^2=e_0^2\,N\,f\,\mathsf g
\]

%dsp-level:0
\[
\qquad [\mathrm{BM.88}(4)]\implies\mathcal L_\ast\ \ \text{is the level set of the derived profile at the derived strength};
\]

%dsp-level:0
\[
\textbf{the adjudication of the cancellation claim:}
\]

%dsp-level:0
\[
\qquad e_0^2\ \ \textbf{does}\ \neg\textbf{cancel}\ —\ \text{it is the one datum that places}\ \mathcal L_\ast\ \text{in}\ (r,E):
\]

%dsp-level:0
\[
\boxed{\ \textbf{what is}\ e_0\textbf{-free is the crossing value}\ \varepsilon_\ast,\ \textbf{its enclosure, and the crossing count per composition};\qquad\textbf{what requires}\ e_0\ \textbf{is the crossing location}\ }
\]

%dsp-level:0
\[
\neg\exists\ \text{a theorem fixing}\ e_0\ \text{or}\ e
\]

%dsp-level:0
\[
\qquad [\mathrm{EM.25}(7);\ \mathrm{BM.88}(7);\ \mathrm{P4};\ \mathrm{P5}];
\]

%dsp-level:0
\[
\qquad \neg\text{a numerical value of either enters any clause of this tier}\ [\text{machine: the token scan}];
\]

%dsp-level:0
\[
(5)\quad\textbf{the three profile shapes (composition by composition):}
\]

%dsp-level:0
\[
\qquad x=\frac{E}{E_1}=k=\frac{r\,E\ln2}{\mathfrak h},
\]

%dsp-level:0
\[
\qquad E_1=\frac{\mathfrak h}{r\ln2},
\]

%dsp-level:0
\[
\qquad G=2\pi E^2\ln2=\beta x^2,
\]

%dsp-level:0
\[
\qquad \beta=\frac{2\pi\mathfrak h^2}{r^2\ln2}\in\mathfrak q^0
\]

%dsp-level:0
\[
[\mathrm{BM.88}(2)(\mathrm i),(\mathrm{iii});\ \mathrm{BM.88}(3);\ \mathrm{CT2.1};\ \text{machine: the exponent arithmetic}];
\]

%dsp-level:0
\[
\qquad P_{\mathsf g}(x)\ :=\ f(x)\,\mathsf g\bigl(\beta x^2\bigr),
\]

%dsp-level:0
\[
\qquad \varepsilon_{\mathrm{eff}}=N\,\varepsilon_{\mathrm{EM}}\,P_{\mathsf g}(x):
\]

%dsp-level:0
\[
(\mathrm i)\quad\mathsf g=1:
\]

%dsp-level:0
\[
\qquad P(x)=1-\frac1x\ \ \text{strictly increasing on}\ x\ge1,
\]

%dsp-level:0
\[
\qquad P(1)=0,
\]

%dsp-level:0
\[
\qquad P\to1:
\]

%dsp-level:0
\[
\qquad \boxed{\ \text{at most one crossing; it exists}\iff N\varepsilon_{\mathrm{EM}}>\varepsilon_\ast\ }
\]

%dsp-level:0
\[
\textbf{the crossing scale in closed form:}
\]

%dsp-level:0
\[
\qquad x_\ast=\frac1{1-c},
\]

%dsp-level:0
\[
\qquad c:=\frac{\varepsilon_\ast}{N\,\varepsilon_{\mathrm{EM}}}\in(0,1),
\]

%dsp-level:0
\[
\qquad E_\ast=\frac{\mathfrak h}{r\ln2\,(1-c)}
\]

%dsp-level:0
\[
\qquad [\text{machine: the linear equation}\ x-1=cx];
\]

%dsp-level:0
\[
(\mathrm{ii})\quad\mathsf g=1+G:
\]

%dsp-level:0
\[
\qquad P(x)=\Bigl(1-\frac1x\Bigr)\bigl(1+\beta x^2\bigr)\ \ \text{strictly increasing},
\]

%dsp-level:0
\[
\qquad P(1)=0,
\]

%dsp-level:0
\[
\qquad P\to\infty:
\]

%dsp-level:0
\[
\qquad \boxed{\ \textbf{exactly one crossing at every admissible}\ (N,\varepsilon_{\mathrm{EM}})>0\ }
\]

%dsp-level:0
\[
\textbf{the crossing equation is an exact cubic:}
\]

%dsp-level:0
\[
\qquad \beta x^3-\beta x^2+(1-c)\,x-1=0
\]

%dsp-level:0
\[
\qquad [\text{machine: the cleared identity}\ (x-1)(1+\beta x^2)=cx];
\]

%dsp-level:0
\[
(\mathrm{iii})\quad\mathsf g=(1+G)^{-1}:\qquad P(x)=\frac{1-1/x}{1+\beta x^2},\qquad P(1)=0,\qquad P\to0\ \ \text{as}\ \ x\to\infty:\qquad\textbf{a window, not a half-line}
\]

%dsp-level:0
\[
\textbf{the unique interior maximum:}
\]

%dsp-level:0
\[
\qquad P'(x_\ast)=0\iff2\beta x_\ast^3-3\beta x_\ast^2-1=0
\]

%dsp-level:0
\[
\qquad [\text{machine: the cleared stationarity row}];
\]

%dsp-level:0
\[
\text{the cubic}\ h(x)=2\beta x^3-3\beta x^2-1\ \ \text{has}\ h'(x)=6\beta x(x-1)\ge0\ \ \text{on}\ \ x\ge1,
\]

%dsp-level:0
\[
\qquad h(1)=-\beta-1<0,
\]

%dsp-level:0
\[
\qquad h\to+\infty\implies\boxed{\ \exists!\ x_\ast(\beta)>1\ }
\]

%dsp-level:0
\[
\boxed{\ P_{\max}(\beta)=P(x_\ast)=1-\frac3{2x_\ast}\ }
\]

%dsp-level:0
\[
\qquad [\text{machine: the polynomial identity}\ 2x\bigl(1-\tfrac1x\bigr)-(2x-3)(1+\beta x^2)=1-(2\beta x^3-3\beta x^2),\ \text{vanishing on the cubic}];
\]

%dsp-level:0
\[
\textbf{the exact dichotomy in}\ \beta:
\]

%dsp-level:0
\[
\qquad h\bigl(\tfrac32\bigr)=-1<0\ \ \forall\beta,
\]

%dsp-level:0
\[
\qquad h(2)=4\beta-1
\]

%dsp-level:0
\[
\qquad [\text{machine: two polynomial identities}]:
\]

%dsp-level:0
\[
\boxed{\ \beta>\frac14\implies x_\ast\in\Bigl(\frac32,\,2\Bigr)\ \ \text{and}\ \ P_{\max}\in\Bigl(0,\,\frac14\Bigr);\qquad\beta=\frac14\implies x_\ast=2,\quad P_{\max}=\frac14\ \ \text{exactly};\qquad\beta<\frac14\implies x_\ast>2,\quad P_{\max}\in\Bigl(\frac14,\,1\Bigr)\ }
\]

%dsp-level:0
\[
\textbf{the crossing equation is an exact depressed cubic:}
\]

%dsp-level:0
\[
\qquad c\beta x^3+(c-1)\,x+1=0
\]

%dsp-level:0
\[
\qquad [\text{machine: the cleared identity}\ x-1=cx(1+\beta x^2)];
\]

%dsp-level:0
\[
\textbf{the root count is the crossing count:}
\]

%dsp-level:0
\[
\qquad c\beta x^3+(c-1)x+1=-\,x\bigl(1+\beta x^2\bigr)\bigl(P(x)-c\bigr)
\]

%dsp-level:0
\[
\qquad [\text{machine: one polynomial identity}]:
\]

%dsp-level:0
\[
\boxed{\ c<P_{\max}:\ \textbf{two crossings}\ (\text{the window opens and closes});\qquad c=P_{\max}:\ \textbf{one}\ (x_\ast,\ \text{a double root});\qquad c>P_{\max}:\ \textbf{none}\ }
\]

%dsp-level:0
\[
(6)\quad\textbf{the specified and pure-pole graphs are exactly flat in the energy scale:}
\]

%dsp-level:0
\[
\qquad Q_\beta^{\mathrm{spec}}\,\mathfrak h=\Delta T(S)-\rho_\ell\sqrt{G_\ell}-\rho_\nu\sqrt{G_\nu},
\]

%dsp-level:0
\[
\qquad \Delta T=5\,m^2_B(\cdot,S)
\]

%dsp-level:0
\[
\text{the alphabet}\ \{m^2_B(\cdot,S),\ \rho_\ell,\ \rho_\nu,\ g,\ g',\ \mathfrak h\}\ \ \text{carries}\ \neg E_{\mathrm{int}},\ \neg r,\ \neg N,\ \neg\mathsf g,\ \neg\varepsilon_{\mathrm{EM}}
\]

%dsp-level:0
\[
\qquad [\mathrm{BM.92}(5);\ \mathrm{BM.97}(2),(3);\ \text{machine: the token census}]:
\]

%dsp-level:0
\[
\boxed{\ \textbf{the specified read's energy-scale graph is a horizontal line, and its trichotomy is decided by the scales}\ (\rho_\ell,\rho_\nu),\ \neg\textbf{by the scale}\ E\ }
\]

%dsp-level:0
\[
\qquad [\mathrm{BM.77}(5),(6)\ \text{re-cited}];
\]

%dsp-level:0
\[
\text{the pure-pole read is flat in}\ E\ \text{as well, its crossing being}\ \Delta_g^{\mathrm{crit}}=s(2M_0+s)
\]

%dsp-level:0
\[
\qquad [\mathrm{BM.97}(4)];
\]

%dsp-level:0
\[
\qquad \textbf{the entire energy-scale structure lives in the informational/EM channel};
\]

%dsp-level:0
\[
(7)\quad\textbf{the graph's verdict rows:}
\]

%dsp-level:0
\[
\qquad \textbf{the monotonicity:}\ \ \Delta m(E)\ \ \text{is strictly decreasing in}\ E\ \ \text{on}\ (\mathrm i),
\]

%dsp-level:0
\[
(\mathrm{ii}),
\]

%dsp-level:0
\[
\ \ \text{and rises then falls on}\ (\mathrm{iii})
\]

%dsp-level:0
\[
\qquad [\text{machine: three monotonicity censuses on exact rationals}];
\]

%dsp-level:0
\[
\textbf{the envelope row:}
\]

%dsp-level:0
\[
\qquad f(1)=0\implies\varepsilon_{\mathrm{eff}}=0\ \ \text{at}\ \ E=E_1\implies\Delta m(E_1)\,\mathfrak h=\ln\tfrac43>0
\]

%dsp-level:0
\[
\qquad [\mathrm{BM.88}(5);\ \text{machine}]:
\]

%dsp-level:0
\[
\qquad \textbf{the splitting is open at the member's own envelope on every composition};
\]

%dsp-level:0
\[
\textbf{the bar as a pointwise crossing census:}\qquad f\le1\ \ \text{on}\ \ x\ge1\implies\varepsilon_{\mathrm{eff}}\ \le\ N\,\mathsf g\bigl(G(E)\bigr)\,\varepsilon_{\mathrm{EM}}\qquad[\text{machine}]:
\]

%dsp-level:0
\[
N\mathsf g\,\varepsilon_{\mathrm{EM}}\le\tfrac18<\varepsilon_\ast\implies\boxed{\ \Delta m(E)>0\ \ \text{at that scale}\ },
\]

%dsp-level:0
\[
\qquad N\mathsf g\,\varepsilon_{\mathrm{EM}}\ge\tfrac5{12}>\varepsilon_\ast\implies\boxed{\ \Delta m(E)<0\ \ \text{at that scale}\ }
\]

%dsp-level:0
\[
\qquad [\mathrm{BM.99}(3)\ \text{re-derived on the graph};\ \text{machine}];
\]

%dsp-level:0
\[
\textbf{the global crossing census per composition:}
\]

%dsp-level:0
\[
\qquad (\mathrm i)\ \ \text{a crossing exists}\iff N\varepsilon_{\mathrm{EM}}>\varepsilon_\ast;
\]

%dsp-level:0
\[
\qquad (\mathrm{ii})\ \ \text{exactly one at every}\ (N,\varepsilon_{\mathrm{EM}})>0;
\]

%dsp-level:0
\[
(\mathrm{iii})\ \ \text{two}\iff N\varepsilon_{\mathrm{EM}}P_{\max}>\varepsilon_\ast,\qquad\text{one at equality},\qquad\text{none below}\qquad[(5);\ \text{machine}];
\]

%dsp-level:0
\[
\textbf{the band is where the flicker decides:}
\]

%dsp-level:0
\[
\qquad \frac18<N\mathsf g\,\varepsilon_{\mathrm{EM}}<\frac5{12}\implies\text{the crossing's existence is decided by}\ \varepsilon_\ast(w)=\frac{2-\log_23}{3w},\ \ \text{i.e. by}\ w
\]

%dsp-level:0
\[
\qquad [(3);\ \text{machine}];
\]

%dsp-level:0
\[
(8)\quad\textbf{status:}
\]

%dsp-level:0
\[
\qquad \textbf{derived:}\ (1)\text{--}(7);
\]

%dsp-level:0
\[
\qquad \textbf{read-relative:}\ \text{the verdict rows}\ (2),(6);
\]

%dsp-level:0
\[
\qquad \textbf{rescoped,}\ \neg\textbf{retracted:}\ \mathrm{BM.91}(2),\ \mathrm{BM.99}(1),(3)\ \text{stand as the far-field locus}\ f=1;
\]

%dsp-level:0
\[
\neg\text{a declaration},
\]

%dsp-level:0
\[
\quad\neg\text{a P4 addition},
\]

%dsp-level:0
\[
\quad\neg\text{a new dimensionful scale}\ [\mathrm{P1}],
\]

%dsp-level:0
\[
\quad\neg\text{a measured value}\ [\mathrm{P5}],
\]

%dsp-level:0
\[
\quad\neg\text{a numerical coupling}\ [\mathrm{EM.25}(7)];
\]

%dsp-level:0
\[
\qquad\textbf{to calibration:}\ e_0\ (4),\quad\mathsf g\ (5)\ [\mathrm{BM.104}(5)],\quad N\ \text{as solution data}\ [\mathrm{CT2.7};\ \mathrm{BM.34}(2)].
\]

\textbf{Req.} D1.8, D2.9, D3.4 (Parts 1–2); A6, CT2.1, CT2.7, EM.9, EM.25(7), EM.25(9) (Part 2); P1, P4, P5 (Parts 1–2); Def-BM.12, Def-BM.4(1), BM.13(2), BM.14(2), BM.16(1), BM.17(2), BM.23(9), BM.28(4), BM.34(2), BM.39(1), BM.70(2), BM.70(3), BM.77(5), BM.77(6), BM.84(9), BM.88(2), BM.88(3), BM.88(4), BM.88(5), BM.88(6), BM.88(7), BM.89(5), BM.91(2), BM.91(4), BM.92(1), BM.92(5), BM.94(5), BM.94(7), BM.96(3), BM.97(2), BM.97(3), BM.97(4), BM.99(1), BM.99(3), BM.99(5), BM.99(6) (this part); Census FT.35(2) (row-wise); exact polynomial identities over $\mathbb Q[\beta,x,c]$, cubic root counts by monotonicity and sign, exact rational lattices, the $\log_2$-ring sign test with integer power certificates, Laurent-monomial assembly, token censuses [math; machine: \texttt{checks/\allowbreak{}verify\_\allowbreak{}H\_\allowbreak{}six\_\allowbreak{}objective\_\allowbreak{}resolution.py} §A].

\textbf{Pf.}

%dsp-level:0
\[
(1):\ Q_\beta^{\mathrm{inv}}
\]

%dsp-level:0
\[
= (m_p+m_\ell+m_\nu)-m_n
\]

%dsp-level:0
\[
= -\bigl(m_n-m_p-m_\ell-m_\nu\bigr)
\]

%dsp-level:0
\[
= -Q_\beta^{\mathrm{fwd}}\ [\text{machine: one identity over}\ \mathbb Q[m_n,m_p,m_\ell,m_\nu]]
\]

%dsp-level:0
\[
(2):\ \varepsilon_{\mathrm{eff}}
\]

%dsp-level:0
\[
= Nf\,
\]

%dsp-level:0
\[
\mathsf g\,
\]

%dsp-level:0
\[
\varepsilon_{\mathrm{EM}}\ \text{and}\ e^2
\]

%dsp-level:0
\[
= e_0^2Nf\,
\]

%dsp-level:0
\[
\mathsf g\ \text{with}\ \varepsilon_{\mathrm{EM}}
\]

%dsp-level:0
\[
= e_0^2/(4\pi)\implies\varepsilon_{\mathrm{eff}}
\]

%dsp-level:0
\[
= e^2/(4\pi)\ [\mathrm{BM.88}(4);\ \text{machine: one monomial}];
\]

%dsp-level:0
\[
\qquad \mathrm{BM.91}(2)\text{'s term}\ 3Nw\,\mathsf g\,\varepsilon_{\mathrm{EM}}\ln2\ \text{is the}\ f{=}1\ \text{locus of}\ 3w\ln2\cdot\varepsilon_{\mathrm{eff}}\ [\text{machine}]
\]

%dsp-level:0
\[
Q_\beta^{\mathrm{fwd}}\mathfrak h
\]

%dsp-level:0
\[
= \ln\tfrac43-3w\ln2\cdot\varepsilon_{\mathrm{eff}}-F_{\ell\nu}\le\ln\tfrac43-2\ln2
\]

%dsp-level:0
\[
= -\ln3\ [\mathrm{BM.70}(2),(3);\ \text{machine: the coefficient census on a}\ 504\text{-point rational lattice}];
\]

%dsp-level:0
\[
\qquad -\ln3<0\iff4<12\ [\text{machine}]
\]

%dsp-level:0
\[
(3):\ \Delta m
\]

%dsp-level:0
\[
= 0
\]

%dsp-level:0
\[
\iff 3w\ln2\cdot\varepsilon_{\mathrm{eff}}
\]

%dsp-level:0
\[
= \ln\tfrac43\iff\varepsilon_{\mathrm{eff}}
\]

%dsp-level:0
\[
= \tfrac{\ln(4/3)}{3w\ln2}
\]

%dsp-level:0
\[
= \tfrac{2-\log_23}{3w}\ [\text{machine: exact in the}\ \log_2\ \text{ring}];
\]

%dsp-level:0
\[
\qquad e_\ast^2=4\pi\varepsilon_\ast\ [\text{machine}]
\]

%dsp-level:0
\[
w\in(\tfrac13,1]\ [\mathrm{BM.94}(7)]\implies\varepsilon_\ast\in[\tfrac{2-\log_23}3,\ 2-\log_23)\ [\text{machine}];
\]

%dsp-level:0
\[
\qquad \tfrac{2-\log_23}3>\tfrac18\iff8(2-\log_23)>3\iff2^{13}>3^8;
\]

%dsp-level:0
\[
\qquad 2-\log_23<\tfrac5{12}\iff3^{12}>2^{19}\ [\text{machine: two integer power comparisons}]
\]

%dsp-level:0
\[
(4):\ \mathcal L_\ast\ \text{is the preimage of one value under the derived profile}\ [\mathrm{BM.88}(4)];
\]

%dsp-level:0
\[
\qquad e^2
\]

%dsp-level:0
\[
= e_0^2Nf\,
\]

%dsp-level:0
\[
\mathsf g\implies\text{the locus moves with}\ e_0\ \text{and its value does not}\ [\text{machine: the symbol census of}\ \varepsilon_\ast\ \text{carries}\ \{w,\ln2,\log_23,\pi\},\ \neg e_0]
\]

%dsp-level:0
\[
(5):\ x=k=rE\ln2/\mathfrak h=E/E_1\ \text{at}\ E_1=\mathfrak h/(r\ln2)\ [\text{machine}];
\]

%dsp-level:0
\[
\qquad G=2\pi E^2\ln2=2\pi\ln2\,E_1^2x^2=\beta x^2,
\]

%dsp-level:0
\[
\qquad [\beta]=[\mathfrak h^2/r^2]=\mathfrak q^{-1}\mathfrak q^{+1}=\mathfrak q^0\ [\mathrm{CT2.1};\ \text{machine}]
\]

%dsp-level:0
\[
(\mathrm i):\ P=1-\tfrac1x,\ P'=1/x^2>0\ \text{on}\ x\ge1;\ P(1)=0;\ P\to1\ [\text{machine}];
\]

%dsp-level:0
\[
\qquad x-1=cx\iff x(1-c)=1\iff x=1/(1-c),\ \ \text{admissible}\iff0<c<1\iff N\varepsilon_{\mathrm{EM}}>\varepsilon_\ast\ [\text{machine}]
\]

%dsp-level:0
\[
(\mathrm{ii}):\ P'=\tfrac1{x^2}(1+\beta x^2)+(1-\tfrac1x)2\beta x>0\ \text{on}\ x>1;\ P\to\infty\ [\text{machine}];
\]

%dsp-level:0
\[
\qquad (x-1)(1+\beta x^2)
\]

%dsp-level:0
\[
= cx
\]

%dsp-level:0
\[
\iff x+\beta x^3-1-\beta x^2
\]

%dsp-level:0
\[
= cx\iff\beta x^3-\beta x^2+(1-c)x-1
\]

%dsp-level:0
\[
= 0\ [\text{machine: the cleared cubic}]
\]

%dsp-level:0
\[
(\mathrm{iii}):\ P'=\tfrac1{x^2(1+\beta x^2)}-\tfrac{(1-\tfrac1x)2\beta x}{(1+\beta x^2)^2};
\]

%dsp-level:0
\[
\qquad P'=0\iff(1+\beta x^2)=2\beta x^3(1-\tfrac1x)\cdot\tfrac{x^2}{x^2}\iff2\beta x^3-3\beta x^2-1=0\ [\text{machine: cleared}]
\]

%dsp-level:0
\[
h'=6\beta x^2-6\beta x=6\beta x(x-1)\ge0\ \text{on}\ x\ge1,\ \ \text{with equality only at}\ x=1;
\]

%dsp-level:0
\[
\qquad h(1)=-\beta-1<0;
\]

%dsp-level:0
\[
\qquad h(x)\to+\infty\implies\exists!\,x_\ast>1\ [\text{machine: monotonicity plus the two endpoint signs}]
\]

%dsp-level:0
\[
\beta x_\ast^2(2x_\ast-3)
\]

%dsp-level:0
\[
= 1
\]

%dsp-level:0
\[
\implies 1+\beta x_\ast^2
\]

%dsp-level:0
\[
= \tfrac{2x_\ast-2}{2x_\ast-3}
\]

%dsp-level:0
\[
\implies P(x_\ast)
\]

%dsp-level:0
\[
= \tfrac{x_\ast-1}{x_\ast}\cdot\tfrac{2x_\ast-3}{2x_\ast-2}
\]

%dsp-level:0
\[
= \tfrac{2x_\ast-3}{2x_\ast}
\]

%dsp-level:0
\[
= 1-\tfrac3{2x_\ast}\ [\text{machine: the identity}\ 2x(1-\tfrac1x)-(2x-3)(1+\beta x^2)=1-(2\beta x^3-3\beta x^2)\ \text{in}\ \mathbb Q[\beta,x]]
\]

%dsp-level:0
\[
h(\tfrac32)=2\beta\tfrac{27}8-3\beta\tfrac94-1=-1;
\]

%dsp-level:0
\[
\qquad h(2)
\]

%dsp-level:0
\[
= 16\beta-12\beta-1
\]

%dsp-level:0
\[
= 4\beta-1\ [\text{machine: two polynomial identities}]\implies\text{the}\ \beta\text{-dichotomy and}\ P_{\max}\ \text{enclosures}\ [(1-\tfrac3{2x_\ast})\ \text{monotone in}\ x_\ast]
\]

%dsp-level:0
\[
x-1=cx(1+\beta x^2)\iff c\beta x^3+(c-1)x+1=0;
\]

%dsp-level:0
\[
\qquad x(1+\beta x^2)(P-c)
\]

%dsp-level:0
\[
= (x-1)-cx(1+\beta x^2)
\]

%dsp-level:0
\[
= -\bigl(c\beta x^3+(c-1)x+1\bigr)\ [\text{machine: one identity}]\implies\text{the root count is the crossing count}
\]

%dsp-level:0
\[
(6):\ \Delta T=5\,m^2_B(\cdot,S)\ \text{is a frozen rational per cell}\ [\mathrm{BM.77}(2)];
\]

%dsp-level:0
\[
\qquad G_\ell,G_\nu\ \text{depend on}\ (g,g',\delta)\ \text{alone}\ [\mathrm{BM.39}(1);\ \mathrm{BM.96}(2)];
\]

%dsp-level:0
\[
\qquad \text{the token census}\ [\text{machine}]
\]

%dsp-level:0
\[
(7):\ \varepsilon_{\mathrm{eff}}(E_1)=Nf(1)\,\mathsf g\,\varepsilon_{\mathrm{EM}}=0\implies\Delta m(E_1)\mathfrak h=\ln\tfrac43>0\ [\mathrm{BM.88}(5);\ \text{machine}];
\]

%dsp-level:0
\[
\qquad f\le1\ \text{on}\ x\ge1\implies\varepsilon_{\mathrm{eff}}\le N\mathsf g\,\varepsilon_{\mathrm{EM}}\ [\text{machine}];
\]

%dsp-level:0
\[
\qquad (8):\ \text{the scope record}.\ \blacksquare
\]

\textbf{Phys.} One graph, read twice. The forward and inverse Q-values are exact negatives, so the two reactions the question asks about are one curve and its reflection, and the boundary between them is a single level set of the derived coupling profile. On the informational currency the curve never changes sign: the two-bit floor on the emitted lines keeps the forward channel closed and the inverse channel open at every scale, with the exact margin log three, so the only sign structure the graph carries is the splitting's. That structure is decided by one derived number, the crossing strength — log of four thirds over three weights times log two, which is two minus the binary logarithm of three over three weights — enclosed between an eighth and five twelfths by the same two integer powers that give the parameter-free bar. The bar and the crossing are the same statement: below the lower end the curve never crosses at any scale, above the upper end it always does, and inside the band the flicker weight alone decides. The profile's own shape then does the rest. With the flat composition the crossing is a closed form, one over one minus c; with the growing composition there is exactly one crossing at every positive coupling, because the profile diverges; with the decreasing composition the profile is a window — it rises from zero, peaks at the unique root of an exact cubic, and returns to zero, so the forward channel opens at one scale and closes at another, or never opens, and the peak has the closed form one minus three over twice the root, sitting below a quarter whenever the radius parameter exceeds a quarter. Two compositions give a half-line, one gives a window, and the corpus still cannot say which is its own. What is exact and free of every undeclared datum is the crossing value; what needs the one open calibration number is where in radius and energy that value is reached.

\end{theorem}

\begin{theorem}[{\texorpdfstring{BM.102 — The dimensional bridge: the $\mathfrak q$-exponent census, the seven rows, and the declared calibration slot}{BM.102 - The dimensional bridge: the -exponent census, the seven rows, and the declared calibration slot}}]\label{tag:BM-102}
\textbf{Thm.}

%dsp-level:0
\[
(1)\quad\textbf{the one-dimension law:}
\]

%dsp-level:0
\[
\qquad \text{units}\ :=\ \text{powers of}\ \mathfrak q;
\]

%dsp-level:0
\[
\qquad c=\hbar=k_B=1;
\]

%dsp-level:0
\[
\qquad [\mathfrak h]=\mathfrak q^{-1/2};
\]

%dsp-level:0
\[
\qquad \mathfrak h:=\text{the single free dimensionful scale}
\]

%dsp-level:0
\[
\qquad [\mathrm{D2.9};\ \mathrm{P1};\ \text{A §0.3}]:
\]

%dsp-level:0
\[
\kappa,\ G_N\ \ \text{are names for combinations of}\ \mathfrak h,\ \neg\text{independent inputs}
\]

%dsp-level:0
\[
\qquad [\mathrm{D2.9};\ \mathrm{CT2.9}];
\]

%dsp-level:0
\[
\qquad \boxed{\ \textbf{the inventory carries exactly one fundamental dimension, the information flux}\ \mathfrak q\ }
\]

%dsp-level:0
\[
(2)\quad\textbf{the derived exponent census:}\qquad\begin{array}{l|l|l}\text{quantity} & \text{exponent} & \text{locus}\\\hline\text{length},\ \text{time} & \mathfrak q^{-1/2} & [\mathrm{CT2.1};\ \mathrm{D2.9}:\ c=1]\\\text{area},\ [\Lambda],\ [G_{\mu\nu}],\ \text{curvature} & \mathfrak q^{-1},\ \mathfrak q^{+1},\ \mathfrak q^{+1},\ \mathfrak q^{+1} & [\mathrm{CT2.1}]\\\text{volume} & \mathfrak q^{-3/2} & [\mathrm{Def\text{-}BM.4}(1);\ \mathrm{CT2.1}]\\\text{energy},\ \text{mass},\ [\Delta m] & \mathfrak q^{+1/2} & [\mathrm{CT2.1};\ \mathrm{BM.35}(4)]\\\text{energy density}\ \rho,\ [T_{\mu\nu}] & \mathfrak q^{+2} & [\mathrm{CT2.1};\ \mathrm{BM.14}(2),(4)]\\{}[\kappa^2],\ [G_N],\ [\kappa] & \mathfrak q^{-1},\ \mathfrak q^{-1},\ \mathfrak q^{-1/2} & [\mathrm{CT2.1};\ \mathrm{CT2.9}]\\\mathfrak h & \mathfrak q^{-1/2} & [\mathrm{D2.9}]\\\Delta S_{\mathfrak L},\ I_{\mathrm{hid}},\ S_{\mathrm{raw}},\ \Sigma_{\mathrm{cl}},\ S_{\mathrm{geom}} & \mathfrak q^{0} & [\mathrm{BM.35}(3),(4);\ \mathrm{CT1.29};\ \mathrm{CT2.1};\ \mathrm{Def\text{-}BM.13}(6)]\\E_{\mathrm{int}},\ W(Q),\ q,\ D_S,\ f,\ N,\ \mathsf g,\ w,\ G,\ \varepsilon_{\mathrm{EM}},\ e,\ g,\ g' & \mathfrak q^{0} & [\mathrm{D3.4};\ \mathrm{EM.9};\ \mathrm{Def\text{-}BM.13}(3);\ \mathrm{BM.88}(2),(3);\ \mathrm{EM.25}(9)]\\m^2_{\mathrm{pole}},\ \Delta_g,\ \hat\mu'^2 & \mathfrak q^{+1} & [\mathrm{Def\text{-}H.37}(2);\ \mathrm{FT.54}(4)]\\\hat\mu^2 & \mathfrak q^{+2} & [\mathrm{Def\text{-}H.37}(1);\ \mathrm{EW.33}]\\\rho_{\mathrm B}=\hat\mu'^2\mathfrak h^2/(\ln2)^2 & \mathfrak q^{0} & [\mathrm{BM.69}(2);\ \text{machine:}\ 1-1=0]\end{array}
\]

%dsp-level:0
\[
\textbf{the balance rows:}
\]

%dsp-level:0
\[
\qquad [\rho\,V]=\mathfrak q^{+2}\mathfrak q^{-3/2}=\mathfrak q^{+1/2}=[\text{energy}];
\]

%dsp-level:0
\[
\qquad [\kappa^2\rho]=\mathfrak q^{-1}\mathfrak q^{+2}=\mathfrak q^{+1}=[R];
\]

%dsp-level:0
\[
\qquad [(\Delta x_{\mathrm B})^2\kappa^2\rho]=\mathfrak q^{-1}\mathfrak q^{+1}=\mathfrak q^0=[G]
\]

%dsp-level:0
\[
\qquad [\text{machine}];
\]

%dsp-level:0
\[
(3)\quad\textbf{the seven rows (derived where a referent exists, adjudicated where it does not):}\qquad\begin{array}{l|l|l}\text{row} & \text{exponent} & \text{the corpus referent}\\\hline\text{mass} & \mathfrak q^{+1/2} & m_{\mathrm{hid}},\ m_{\mathrm{rest}},\ m_{\mathrm{pole}}\ [\mathrm{FT.44}(4);\ \mathrm{BM.35}(4);\ \mathrm{Def\text{-}H.37}]\\\text{length} & \mathfrak q^{-1/2} & \Delta x_{\mathrm B},\ \ell_c,\ r\ [\mathrm{Def\text{-}BM.4}(1);\ \mathrm{Def\text{-}H.2}]\\\text{time} & \mathfrak q^{-1/2} & x^0,\ \text{equal to length on}\ c=1\ [\mathrm{D1.8};\ \mathrm{D2.9}]\\\text{energy} & \mathfrak q^{+1/2} & E_{\mathrm{rest}}=m_{\mathrm{rest}}c^2\ \text{with the}\ c\text{-exponent}\ 0\ [\mathrm{BM.35}(4)]\\\text{temperature} & \mathfrak q^{+1/2} & \text{on}\ k_B=1\ [\mathrm{D2.9};\ \mathrm{CT2.1}]\ —\ \textbf{a dimension row,}\ \neg\textbf{a thermal object:}\ \neg\text{a P4-class temperature exists}\ [\mathrm{BM.69}(6)]\\\text{amount of substance} & \mathfrak q^{0} & \text{a count:}\ E_{\mathrm{int}},\ W(Q),\ \text{the fibre cardinalities}\ [\mathrm{Def\text{-}BM.5}(4);\ \mathrm{Def\text{-}BM.13}(3);\ \mathrm{BM.95}(1)]\\\text{electric current} & \mathfrak q^{0}\ \text{and}\ \mathfrak q^{+1/2} & \textbf{two reads:}\ \text{the ledger-side weight per tick}\ \ [\mathrm{BM.23}(9);\ \mathrm{BM.44};\ \mathrm{D1.8};\ \neg\text{a rate per external second}\ [\mathrm{P5}]];\ \ \text{the chart-side read below}\end{array}
\]

%dsp-level:0
\[
\textbf{the chart-side current, derived:}\qquad\partial_\mu F^{\mu\nu}=e\,J^\nu_Q\qquad[\mathrm{Def\text{-}BM.12}(1);\ \mathrm{EM.11}]:
\]

%dsp-level:0
\[
\bigl[\tfrac12E_r^2\cdot V\bigr]=[\text{energy}]=\mathfrak q^{+1/2}\implies[E_r^2]=\mathfrak q^{+2}\implies[F^{\mu\nu}]=\mathfrak q^{+1};
\]

%dsp-level:0
\[
\qquad [\partial_\mu]=\mathfrak q^{+1/2},\quad[e]=\mathfrak q^0\implies\boxed{\ [J^\nu_Q]=\mathfrak q^{+3/2}\ }
\]

%dsp-level:0
\[
\qquad [\text{machine: the exponent arithmetic balances}];
\]

%dsp-level:0
\[
\text{the integrated current}\ [J\cdot\text{area}]=\mathfrak q^{+3/2}\mathfrak q^{-1}=\mathfrak q^{+1/2}=\text{charge}/\text{length}=\text{charge}/\text{time}\ \ \text{on}\ c=1
\]

%dsp-level:0
\[
\qquad [\text{machine}]:
\]

%dsp-level:0
\[
\textbf{the row's adjudication:}
\]

%dsp-level:0
\[
\qquad \text{the ledger-side per-tick read is}\ \mathfrak q^0\ \ \text{and the chart-side read is}\ \mathfrak q^{+1/2}\ —\ \ \text{both derived,}\ \neg\text{equated}\ [\mathrm{BM.37}(2)\ \text{pattern}];
\]

%dsp-level:0
\[
(4)\quad\textbf{the two bridges, separated:}
\]

%dsp-level:0
\[
\qquad \textbf{the dimensional bridge is}\ \mathfrak h:
\]

%dsp-level:0
\[
\qquad m_{\mathrm{hid}}=\frac{E_{\mathrm{int}}\ln2}{\mathfrak h}:\quad\mathfrak q^0\longmapsto\mathfrak q^{+1/2}
\]

%dsp-level:0
\[
\qquad [\mathrm{FT.44}(4);\ \mathrm{BM.35}(4);\ \text{machine}]:
\]

%dsp-level:0
\[
\mathfrak h\ \ \text{is declared by}\ \mathrm{P1}\ \ \text{with}\ \neg\text{a value issued},\qquad\exists!\ \text{such a scale}\ [\mathrm{P1};\ \text{A §0.3}];
\]

%dsp-level:0
\[
\textbf{the numerical bridge is}\ \rho_{\mathrm B}:
\]

%dsp-level:0
\[
\qquad [\rho_{\mathrm B}]
\]

%dsp-level:0
\[
= \mathfrak q^0\implies\boxed{\ \rho_{\mathrm B}\ \ \textbf{converts no dimension}\ —\ \text{it is the dimensionless exchange rate between the bit currency}\ (\ln2\ \text{per directed edge})\ \text{and the angle currency}\ (\Theta_R)\ }
\]

%dsp-level:0
\[
[\mathrm{BM.35}(5)\ \text{Rmk: three currencies, one reference point};\ \mathrm{BM.69}(2);\ \mathrm{BM.98};\ \mathrm{CT1.29};\ \mathrm{FT.54}(4)]:
\]

%dsp-level:0
\[
\textbf{the correction to the framing:}
\]

%dsp-level:0
\[
\qquad \rho_{\mathrm B}\ \ \text{is}\ \neg\textbf{the dimensional bridge}\ —\ \text{that role is}\ \mathfrak h\textbf{'s}\ —\ \ \text{and}\ \neg\text{a single number carries both roles}\ [\text{machine:}\ [\mathfrak h]=\mathfrak q^{-1/2}\ne\mathfrak q^0=[\rho_{\mathrm B}]];
\]

%dsp-level:0
\[
(5)\quad\textbf{the declaration (logged as the §II.4 addendum A-BM.9 before use):}
\]

%dsp-level:0
\[
\qquad \textbf{species:}\ \text{the bridge's dimensionless calibration slot}
\]

%dsp-level:0
\[
\boxed{\ \text{§Cal-BM gains exactly one declared member:}\quad\rho_{\mathrm B}=\frac{\hat\mu'^2\mathfrak h^2}{(\ln2)^2}\in\mathfrak q^0,\qquad\textbf{its value is}\ \neg\textbf{issued}\ }
\]

%dsp-level:0
\[
\textbf{the packaging row:}
\]

%dsp-level:0
\[
\qquad \rho_{\mathrm B}\ \ \text{is a dimensionless combination of two already-declared data}\ —\ \hat\mu'^2\ [\mathrm{Cal\text{-}F};\ \mathrm{Def\text{-}H.37}(2)]\ \ \text{and}\ \ \mathfrak h\ [\mathrm{P1}]:
\]

%dsp-level:0
\[
\neg\text{a new free parameter is introduced},
\]

%dsp-level:0
\[
\qquad \neg\text{a new dimensionful scale}\ [\mathrm{P1}],
\]

%dsp-level:0
\[
\qquad \neg\text{a measured value}\ [\mathrm{P5}],
\]

%dsp-level:0
\[
\qquad \neg\text{a numerical coupling}\ [\mathrm{EM.25}(7),(9)];
\]

%dsp-level:0
\[
\textbf{what the declaration fixes is the datum's status:}
\]

%dsp-level:0
\[
\qquad \text{the bridge's single residue is a declared calibration slot, to be filled by a calibration act the corpus prices and}\ \neg\text{performs}
\]

%dsp-level:0
\[
\qquad [\mathrm{BM.77}(9)\ \text{pattern}];
\]

%dsp-level:0
\[
(6)\quad\textbf{what the slot buys, exactly:}
\]

%dsp-level:0
\[
\qquad \text{with}\ \rho_{\mathrm B}\ \text{supplied, every cross-currency ratio is an exact algebraic number per state:}
\]

%dsp-level:0
\[
\mathcal R_s=\rho_{\mathrm B}\frac{D_S}{4E^2},
\]

%dsp-level:0
\[
\qquad \mathcal R_{\mathrm{mat}}=\rho_{\mathrm B}\frac{D_S\,D^2}{4E^2},
\]

%dsp-level:0
\[
\qquad \mathcal R=\frac{\pi\ln2}2\cdot\frac{\rho_{\mathrm B}D_S}{G}=\frac{\hat\mu'^2D_Sr^2}4(1-f)^2
\]

%dsp-level:0
\[
[\mathrm{BM.69}(2);\ \mathrm{BM.98}(1)\text{--}(4);\ \text{six states, four fields, six minimal polynomials}\ [\mathrm{Def\text{-}BM.14}(3)]]:
\]

%dsp-level:0
\[
\qquad \textbf{the bridge becomes numerical, state by state, on one number};
\]

%dsp-level:0
\[
\textbf{and}\ \neg\textbf{a currency becomes addable:}
\]

%dsp-level:0
\[
\qquad [\Delta_g]=\mathfrak q\ \ \text{against}\ \ [\Delta m]=\mathfrak q^{+1/2}\implies\neg\text{addable}
\]

%dsp-level:0
\[
\qquad [\mathrm{BM.37}(2);\ \mathrm{BM.69}(3);\ \mathrm{BM.90}(5);\ \mathrm{BM.98}(6);\ \text{machine}]:
\]

%dsp-level:0
\[
\textbf{the declaration moves no summand across the two currencies, and the ratio stays a ratio};
\]

%dsp-level:0
\[
(7)\quad\textbf{the mass-form bridge disposition:}
\]

%dsp-level:0
\[
\qquad \boxed{\ \textbf{the relating equation is derived in form}\ [\mathrm{BM.98}]\ \textbf{and its single residue is a declared slot with no value issued}\ \ [(\mathrm{A\text{-}BM.9})]\ \Longrightarrow\ \textbf{the mass-form bridge is closed as an equation}\ }
\]

%dsp-level:0
\[
\textbf{the derivation bar stands:}
\]

%dsp-level:0
\[
\qquad \neg\exists\ \text{a theorem equating the bit currency and the angle currency}\ [\mathrm{BM.35}(5)];
\]

%dsp-level:0
\[
\qquad \mathrm{BM.69}(4)\text{'s four grounds stand unweakened}\ —\ \textbf{the declaration is}\ \neg\textbf{a derivation};
\]

%dsp-level:0
\[
\textbf{the numerical bridge remains a calibration act:}
\]

%dsp-level:0
\[
\qquad \neg\text{a value for}\ \rho_{\mathrm B}\ \text{is derived, declared or fitted here};
\]

%dsp-level:0
\[
\qquad \neg\text{an}\ [\mathrm{EXT}]\ \text{datum fixes it}\ [\mathrm{P5};\ \mathrm{FT.73}(5)\ \text{trap pattern}];
\]

%dsp-level:0
\[
(8)\quad\textbf{the bars:}
\]

%dsp-level:0
\[
\qquad \neg\text{a second dimensionful scale}\ [\mathrm{P1}];
\]

%dsp-level:0
\[
\qquad \alpha_B:\ \neg\text{added},\ \neg\text{derived},\ \text{barred on}\ \mathrm{BM.61}(8)\text{'s four grounds};
\]

%dsp-level:0
\[
\qquad \text{the six}\ D_S\ \text{minimal polynomials are untouched};
\]

%dsp-level:0
\[
\text{species names are confined to the identification layer}\ [\mathrm{Def\text{-}BM.3}(7);\ \mathrm{Def\text{-}BM.6}(7)];
\]

%dsp-level:0
\[
\qquad \text{amplitude}\ne\text{mass}\ [\mathrm{FT.79}(5)];
\]

%dsp-level:0
\[
\qquad \text{rung}\ne\text{scale}\ [\mathrm{D0.8}];
\]

%dsp-level:0
\[
(9)\quad\textbf{status:}
\]

%dsp-level:0
\[
\qquad \textbf{derived:}\ (1)\text{--}(4),\ (6);
\]

%dsp-level:0
\[
\qquad \textbf{declared:}\ (5)\ —\ \text{logged as A-BM.9 before use, one slot, no value};
\]

%dsp-level:0
\[
\qquad \textbf{adjudicated:}\ (3)\text{'s two current reads and the temperature row};
\]

%dsp-level:0
\[
\neg\text{a P4 addition beyond the declared slot},
\]

%dsp-level:0
\[
\quad\neg\text{a new dimensionful scale}\ [\mathrm{P1}],
\]

%dsp-level:0
\[
\quad\neg\text{a measured value}\ [\mathrm{P5}],\qquad \quad
\]

\textbf{Req.} D0.8, D1.8, D2.9, D3.4 (Parts 1–2); A §0.3, CT1.29, CT2.1 (Part 1, Part 2); CT2.9, EM.9, EM.11, EM.25(7), EM.25(9), EW.33, Def-BM.12(1) (Part 2, this part); P1, P4, P5 (Parts 1–2); Def-H.2, Def-H.37(1), Def-H.37(2), FT.44(4), FT.54(4), FT.73(5), FT.79(5) (Part 3); Def-BM.4(1), Def-BM.5(4), Def-BM.13(3), Def-BM.13(6), Def-BM.14(3), BM.14(2), BM.14(4), BM.23(9), BM.35(3), BM.35(4), BM.35(5), BM.37(2), BM.44, BM.61(8), BM.69(2), BM.69(3), BM.69(4), BM.69(6), BM.88(2), BM.88(3), BM.90(5), BM.95(1), BM.98, BM.98(6) (this part); Cal-F; the A-BM.9 declaration (logged before use); exact $\mathfrak q$-exponent arithmetic, Laurent-monomial balances, symbol censuses [math; machine: \texttt{checks/\allowbreak{}verify\_\allowbreak{}H\_\allowbreak{}six\_\allowbreak{}objective\_\allowbreak{}resolution.py} §B].

\textbf{Pf.}

%dsp-level:0
\[
(1):\ \mathrm{D2.9},\ \mathrm{P1},\ \text{A §0.3},\ \mathrm{CT2.9}\ \text{re-cited verbatim (frozen)};
\]

%dsp-level:0
\[
\qquad \neg\exists\ \text{a second dimensionful scale in the inventory}\ [\text{the derived no-go of the second scale re-cited}]
\]

%dsp-level:0
\[
(2):\ \text{every row re-cited from its own locus and its exponent re-computed}\ [\text{machine: the exponent arithmetic}];
\]

%dsp-level:0
\[
\qquad [\rho V]=2-\tfrac32=\tfrac12;\quad[\kappa^2\rho]=-1+2=+1;\quad[(\Delta x_{\mathrm B})^2\kappa^2\rho]=-1-1+2=0\ [\text{machine}]
\]

%dsp-level:0
\[
[\rho_{\mathrm B}]=[\hat\mu'^2]+2[\mathfrak h]-2[\ln2]=1-1-0=0\ [\mathrm{BM.69}(2);\ \text{machine}];
\]

%dsp-level:0
\[
\qquad [\hat\mu^2]=2\ [\mathrm{Def\text{-}H.37}(1);\ \mathrm{EW.33}\ \text{packaging}]
\]

%dsp-level:0
\[
(3):\ [\tfrac12E_r^2V]
\]

%dsp-level:0
\[
= \tfrac12\ \text{carries}\ \mathfrak q^0\ [\mathrm{BM.87}(2)\text{'s exterior integral}]
\]

%dsp-level:0
\[
\implies [E_r^2]
\]

%dsp-level:0
\[
= \tfrac12-\bigl(-\tfrac32\bigr)
\]

%dsp-level:0
\[
= 2\ [\text{machine}];
\]

%dsp-level:0
\[
\qquad [\partial_\mu]=-[x^\mu]=+\tfrac12;
\]

%dsp-level:0
\[
\qquad [eJ]=[\partial_\mu F^{\mu\nu}]=\tfrac12+1=\tfrac32\ \text{and}\ [e]=0\implies[J]=\tfrac32\ [\text{machine}]
\]

%dsp-level:0
\[
[J\cdot\text{area}]=\tfrac32-1=\tfrac12\ [\text{machine}];
\]

%dsp-level:0
\[
\qquad k_B=1\implies[\text{temperature}]=[\text{energy}]=\tfrac12\ [\mathrm{D2.9};\ \mathrm{CT2.1}];
\]

%dsp-level:0
\[
\qquad \neg\exists\ \text{a P4-class temperature}\ [\mathrm{BM.69}(6)\ \text{re-cited}]
\]

%dsp-level:0
\[
(4):\ m_{\mathrm{hid}}=E_{\mathrm{int}}\ln2/\mathfrak h:\ [\cdot]=0-(-\tfrac12)=\tfrac12\ [\text{machine}];
\]

%dsp-level:0
\[
\qquad [\mathfrak h]=-\tfrac12\ne0=[\rho_{\mathrm B}]\implies\neg\text{one number carries both roles}\ [\text{machine}]
\]

%dsp-level:0
\[
(5)\text{--}(6):\ \text{the declaration's species, costs and bars recorded in the registry grammar}\ [\text{A-BM.9};\ \text{machine: the packaging census}\ \{\hat\mu'^2\in\mathrm{Cal\text{-}F},\ \mathfrak h\in\mathrm{P1}\}];
\]

%dsp-level:0
\[
\qquad \text{the six states' exactness per supplied value re-cited}\ [\mathrm{BM.69}(2);\ \mathrm{BM.98}]
\]

%dsp-level:0
\[
(7)\text{--}(9):\ \mathrm{BM.69}(4)\text{'s four grounds and}\ \mathrm{BM.35}(5)\text{'s no-equation row re-cited};
\]

%dsp-level:0
\[
\qquad \text{the standing-rule scans: external tokens, decimal numerals, unit tokens, the barred object}\ \alpha_B\ [\text{machine}].\ \blacksquare
\]

\textbf{Phys.} The corpus has one dimension and this tier counts with it. Length and time are the same power, mass and energy the same power, density and stress-energy two powers up, curvature one, and the single free scale sits at minus one half — so every balance in the census is an integer sum, and the seven rows of the question come out derived where the inventory has a referent and adjudicated where it does not. Temperature is a dimension row and nothing more, because no thermal parameter exists; amount of substance is a count; and electric current has two reads that are not equated, the ledger's weight per tick at exponent zero and the chart's density at three halves, the latter forced by the field equation and the field-energy row. The bridge the question asks about turns out to be two bridges. The dimensional one is the free scale itself: it takes a count of bits to a mass, and it was declared at the outset with no value ever issued. The numerical one is dimensionless, so it converts no dimension at all — it is the exchange rate between the ledger's bit currency and the holonomy's angle currency, and it is the single number the relating equation needs. That number is now declared: one slot in the era's calibration registry, logged before use, packaging two data that were already declared elsewhere, adding no parameter and no scale, and carrying no value. With it supplied, every cross-currency ratio in the corpus becomes an exact algebraic number state by state; without it, nothing is lost, because the four grounds against deriving it stand and no summand crosses between the two currencies. The equation is closed; the number is a calibration act the corpus prices and does not perform.

\end{theorem}

\begin{theorem}[{\texorpdfstring{BM.103 — The three-constituent assembly: the eigenvalue equations that are formed, the tensor-product referent census, and the composition residue}{BM.103 - The three-constituent assembly: the eigenvalue equations that are formed, the tensor-product referent census, and the composition residue}}]\label{tag:BM-103}
\textbf{Thm.}

%dsp-level:0
\[
(1)\quad\textbf{the six-state inventory (re-cited):}
\]

%dsp-level:0
\[
\qquad \mathcal S=\{(\mathcal F_g,\beta):\ g\in\{1,2,3\},\ \beta\in\{-,+\}\},
\]

%dsp-level:0
\[
\qquad |\mathcal S|=6
\]

%dsp-level:0
\[
\qquad [\mathrm{Def\text{-}BM.14}(2)]:
\]

%dsp-level:0
\[
D_S\ \in\ \bigl\{2,\ 4,\ 2-\sqrt2,\ 1,\ 2-2\cos\tfrac\pi{10},\ 2-\sqrt3\bigr\},
\]

%dsp-level:0
\[
\qquad m^2_{\mathrm{pole}}=\hat\mu'^2D_S,
\]

%dsp-level:0
\[
\qquad \text{four fields, six minimal polynomials}
\]

%dsp-level:0
\[
\qquad [\mathrm{Def\text{-}BM.14}(3);\ \mathrm{Census\ FT.63}(6);\ \mathrm{BM.86}(2)]:
\]

%dsp-level:0
\[
x-2,\quad x-4,\quad x^2-4x+2,\quad x-1,\quad x^4-8x^3+19x^2-12x+1,\quad x^2-4x+1\ \ (\text{for}\ 2-\sqrt3)
\]

%dsp-level:0
\[
\qquad [\text{machine: the four quadratic-ring roots re-verified exactly}];
\]

%dsp-level:0
\[
\text{the generation ladder}\ 12:24:60=1:2:5
\]

%dsp-level:0
\[
\qquad [\mathrm{Def\text{-}BM.14}(1);\ \mathrm{BM.60}(1)];
\]

%dsp-level:0
\[
\qquad \Delta_g=\bigl(2,\ \sqrt2-1,\ 2\cos\tfrac\pi{10}-\sqrt3\bigr)\hat\mu'^2>0,
\]

%dsp-level:0
\[
\qquad \Delta_1>\Delta_2>\Delta_3
\]

%dsp-level:0
\[
\qquad [\mathrm{Def\text{-}BM.14}(5);\ \mathrm{BM.69}(3)];
\]

%dsp-level:0
\[
(2)\quad\textbf{the eigenvalue equation that is formed}\ (\textbf{the three-line channel,}\ \neg\textbf{the three constituents}):
\]

%dsp-level:0
\[
\qquad z:=r^2,
\]

%dsp-level:0
\[
\qquad \boxed{\ 4z^3-24z^2+\bigl(21-12\gamma_{\mathrm{ch}}\bigr)z-\bigl(\gamma_{\mathrm{ch}}-1\bigr)^2=0\ }
\]

%dsp-level:0
\[
\textbf{Vieta:}
\]

%dsp-level:0
\[
\qquad \boxed{\ \sum_iz_i=6\ \ (\text{phase-free and channel-free}),\qquad\sum_{i<j}z_iz_j=\frac{21-12\gamma_{\mathrm{ch}}}4,\qquad\prod_iz_i=\frac{(\gamma_{\mathrm{ch}}-1)^2}4\ }
\]

%dsp-level:0
\[
\qquad [\mathrm{BM.80}(3);\ \text{machine: the three symmetric reads}];
\]

%dsp-level:0
\[
\textbf{the depressed cubic on the spectral side:}
\]

%dsp-level:0
\[
\qquad 2(r-1)^3-3(r-1)-\gamma_{\mathrm{ch}}=0\iff2r^3-6r^2+3r+(1-\gamma_{\mathrm{ch}})=0
\]

%dsp-level:0
\[
\qquad [\mathrm{BM.80}(2);\ \mathrm{FT.111}(3);\ \mathrm{FT.143}(2);\ \text{machine: the expansion}];
\]

%dsp-level:0
\[
\gamma_{\mathrm{ch}}=\sqrt2\cos(3\delta_{\mathrm{ch}}),
\]

%dsp-level:0
\[
\qquad \delta_{\mathrm{ch}}\in\Bigl\{\frac29,\ \frac2{27},\ \frac18\Bigr\},
\]

%dsp-level:0
\[
\qquad 27\delta_q=9\delta_\ell=16\delta_\nu=2
\]

%dsp-level:0
\[
\qquad [\mathrm{BM.80}(1);\ \mathrm{FT.143}(5);\ \text{machine}];
\]

%dsp-level:0
\[
\textbf{the elimination row:}
\]

%dsp-level:0
\[
\qquad r(2z+3)=6z+\gamma_{\mathrm{ch}}-1\ \ \text{at}\ \ z=r^2,
\]

%dsp-level:0
\[
\qquad z(2z+3)^2-(6z+\gamma_{\mathrm{ch}}-1)^2\equiv0\ \ \text{in the quotient ring}
\]

%dsp-level:0
\[
\qquad [\mathrm{BM.80}(3);\ \text{machine}];
\]

%dsp-level:0
\[
\boxed{\ \textbf{the three roots are the channel's three lines}\ (\text{the generations}),\ \neg\textbf{the three constituents of a sealed class}\ }
\]

%dsp-level:0
\[
\qquad [\mathrm{Def\text{-}BM.6}(4);\ \mathrm{FT.143}(6)];
\]

%dsp-level:0
\[
(3)\quad\textbf{the characteristic polynomial of the three-constituent assembly in the transport register:}
\]

%dsp-level:0
\[
\qquad W(Q)=\bigl[z^{Q}\bigr]\bigl(2+z+z^{-1}\bigr)^{3}
\]

%dsp-level:0
\[
\qquad [\mathrm{Def\text{-}BM.13}(3)]:
\]

%dsp-level:0
\[
2+z+z^{-1}=z^{-1}(z+1)^2\implies\bigl(2+z+z^{-1}\bigr)^3=\frac{(z+1)^6}{z^3}\qquad[\text{machine: one Laurent identity}]:
\]

%dsp-level:0
\[
\boxed{\ W(Q)\ =\ \binom{6}{Q+3}\ \ \text{exactly, on all seven classes}\ }
\]

%dsp-level:0
\[
\qquad W=(1,6,15,\mathbf{20},15,6,1),
\]

%dsp-level:0
\[
\qquad \sum_QW(Q)=2^6=64=4^3
\]

%dsp-level:0
\[
\qquad [\mathrm{BM.28}(2);\ \text{machine: the closed form against the exhaustive}\ 4^3\ \text{census}];
\]

%dsp-level:0
\[
\textbf{the assembly's transport spectrum:}
\]

%dsp-level:0
\[
\qquad \text{one root}\ z
\]

%dsp-level:0
\[
= -1\ \ \text{of multiplicity six}\ —\ \text{the six directed edges}\ —\ \ \text{and the seven class weights are its binomial coefficients}
\]

%dsp-level:0
\[
\boxed{\ \textbf{the eigenvalue structure the chart resolves is charge-class data,}\ \neg\textbf{flavour data}\ }
\]

%dsp-level:0
\[
\qquad [\mathrm{Def\text{-}BM.13}(6);\ \mathrm{BM.95}(1);\ \text{machine}];
\]

%dsp-level:0
\[
(4)\quad\textbf{the tensor-product referent census (what an ensemble matrix requires against what the corpus supplies):}\qquad\begin{array}{l|l}\text{an}\ M_{\mathrm{ensemble}}\ \text{requires} & \text{the corpus supplies}\\\hline\text{a three-constituent state space} & \neg\text{one:}\ \mathcal H_c,\ \mathcal H_c^{(q)}\ \text{are counting spaces with no inner product, no dynamics, no state, no Hamiltonian}\ [\mathrm{Def\text{-}H.41}\ \text{Species};\ \mathrm{Def\text{-}H.42}];\\ & \neg\text{a two-chain state space exists, let alone a three-chain one}\ [\mathrm{FT.80}(5)]\\\text{an exchange operator} & \neg\exists\ \text{one in the inventory}\ [\mathrm{FT.80}(5);\ \mathrm{FT.125}(5)]\\\text{a flavour index attaching states to constituents} & \neg\text{derived:}\ (\mathrm{R3}),\ \text{deficit}\ 1'\ [\mathrm{BM.28}(7);\ \mathrm{BM.75}(6);\ \mathrm{Def\text{-}BM.8}(5)]\\\text{a triple-level resolution} & \neg\text{available: the end signs carry no position datum}\ [\mathrm{D3.4};\ \mathrm{BM.46}(2);\ \mathrm{Def\text{-}BM.13}(9)]\\\text{a}\ 27\text{-dimensional flavour arena} & \text{rejected as an arena at rung one}\ [\mathrm{FT.142}(3):\ \text{the}\ 27=\text{seal}\times\text{transfer}\times\text{colour candidate is a declared construction, cited and imported nowhere}]\end{array}
\]

%dsp-level:0
\[
\boxed{\ \det\bigl(\lambda I-M_{\mathrm{ensemble}}\bigr)=0\ \ \textbf{is}\ \neg\textbf{formed}\ —\ \text{the question ``which three states compose a sealed class'' has no eigenvalue referent in this corpus}\ }
\]

%dsp-level:0
\[
(5)\quad\textbf{what is derived about the composition:}
\]

%dsp-level:0
\[
\qquad \textbf{exactly three constituents seal:}
\]

%dsp-level:0
\[
\quad B:=\varepsilon_{ijk}q^iq^jq^k\ [\mathrm{D3.9}];\quad\varepsilon_{ijk}\ \text{seals}\ 3\ \text{fundamentals}\ [\mathrm{CT3.15};\ \mathrm{CT3.17}];
\]

%dsp-level:0
\[
\quad 3\ \text{admits no}\ \{2,3\}\text{-partition with an external remainder}\ [\mathrm{FT.34}(5);\ \mathrm{Census\ FT.35}(5)]:
\]

%dsp-level:0
\[
\textbf{their core charges sum to the class read:}\quad Q=\sum_iq_{\mathrm{core}}(c_i)\in\{-3,\dots,+3\},
\]

%dsp-level:0
\[
\qquad Q_{\mathrm{EM}}=0\ \ \text{in every valid universe}
\]

%dsp-level:0
\[
\qquad [\mathrm{Def\text{-}BM.8}(1);\ \mathrm{EM.11};\ \mathrm{D2.11}];
\]

%dsp-level:0
\[
\textbf{all seven classes are formed with the binomial weights of}\ (3);
\]

%dsp-level:0
\[
\qquad \textbf{the neutral class outweighs either charged class by exactly}\ \tfrac43\ \ \text{on route}\ (\mathrm{ii})
\]

%dsp-level:0
\[
\qquad [\mathrm{BM.28}(4)(\mathrm{ii});\ \mathrm{BM.30}(2)];
\]

%dsp-level:0
\[
\textbf{the fractional constituent alphabet is structurally impossible}
\]

%dsp-level:0
\[
\qquad [\mathrm{BM.28}(7);\ \mathrm{Def\text{-}BM.8}(5)];
\]

%dsp-level:0
\[
\qquad \textbf{colour is}\ \neg\textbf{observable by construction}
\]

%dsp-level:0
\[
\qquad [\mathrm{CT3.22}];
\]

%dsp-level:0
\[
\boxed{\ \textbf{the ``three'' is derived from the seal grammar and the ``which three'' is not}\ }
\]

%dsp-level:0
\[
(6)\quad\textbf{the stability question, adjudicated:}\qquad\neg\exists\ \text{a derived stability condition on eigenvalues:}
\]

%dsp-level:0
\[
(\mathrm i)\quad\text{the materialization filter is an undeclared spectral hypothesis}\ (\mathrm{M1})/(\mathrm{M2})
\]

%dsp-level:0
\[
\qquad [\mathrm{Def\text{-}H.20};\ \text{A §0.2a};\ \text{Part 3 §0.2}(\mathrm{iii})];
\]

%dsp-level:0
\[
(\mathrm{ii})\quad\text{the index-zero selection route is refuted for bound chains: on Synge-class slices every closed geodesic has}\ \operatorname{ind}_B
\]

%dsp-level:0
\[
\ge 1,
\]

%dsp-level:0
\[
\ \ \text{so absolute index-}0\ \text{selection is inadmissible against}\ (\mathrm{M1})
\]

%dsp-level:0
\[
\qquad [\mathrm{FT.66}(4)];
\]

%dsp-level:0
\[
(\mathrm{iii})\quad\text{a mass-ordering criterion would be a count}\to\text{strength coupling}
\]

%dsp-level:0
\[
\qquad [\mathrm{BM.36}(3);\ \mathrm{BM.61}(8)(\mathrm{O1});\ \mathrm{BM.90}(3)]:
\]

%dsp-level:0
\[
\textbf{the formed stability rows are kinematic and ledger-side:}
\]

%dsp-level:0
\[
\qquad \text{the two-body impossibility}\ [\mathrm{BM.41}(2)];
\]

%dsp-level:0
\[
\quad\text{the two-bit floor}\ [\mathrm{BM.70}(2)];
\]

%dsp-level:0
\[
\quad\text{the materialization restriction}\ E
\]

%dsp-level:0
\[
\ge 2\ [\mathrm{BM.94}(5)];
\]

%dsp-level:0
\[
\text{the valid-universe conditions}\ [\mathrm{D2.11}];
\]

%dsp-level:0
\[
\quad\text{the no-flicker theorem}\ \Gamma(\varepsilon_q{=}0)
\]

%dsp-level:0
\[
= 0\ [\mathrm{BM.23}(8)];
\]

%dsp-level:0
\[
\quad\text{the two-sided bar}\ [\mathrm{BM.99}(3)];
\]

%dsp-level:0
\[
\quad\text{the forward-only arrow}\ [\mathrm{A6};\ \mathrm{FT.144}(2)];
\]

%dsp-level:0
\[
(7)\quad\textbf{the octet row:}
\]

%dsp-level:0
\[
\qquad \neg\text{formed}.
\]

%dsp-level:0
\[
\qquad \textbf{the corpus's derived}\ 3\ \ \text{is the colour seal's fundamental}\ [\mathrm{CT3.15};\ \mathrm{CT3.17};\ \mathrm{D3.9}];
\]

%dsp-level:0
\[
\textbf{its derived}\ 8\ \ \text{is a mask-capacity ratio, barred as a coupling}\ [\mathrm{BM.36}(2),(3);\ \mathrm{BM.61}(8)(\mathrm{O1})];
\]

%dsp-level:0
\[
\qquad \textbf{its derived}\ 10\ \ \text{is the colour grammar's grouping count}\ g(5)=10\ [\mathrm{CT3.17a}(2);\ \mathrm{Def\text{-}BM.9};\ \mathrm{BM.31}];
\]

%dsp-level:0
\[
\textbf{the flavour-side symmetry is}\ C_3,
\]

%dsp-level:0
\[
\ \ \text{its orbits the three channel lines}\ [\mathrm{FT.111}(2);\ \mathrm{FT.158}(1),(4);\ \mathrm{FT.142}(5)];
\]

%dsp-level:0
\[
\qquad \neg\exists\ \text{a flavour}\ \mathsf{SU}(3)\ \text{in the inventory};
\]

%dsp-level:0
\[
\boxed{\ \textbf{the baryon octet and the meson octet are}\ \neg\textbf{derived and}\ \neg\textbf{derivable here}\ —\ \text{what is derived is the}\ \varepsilon\text{-seal's three, the binomial seven, and the}\ C_3\ \text{triplet}\ }
\]

%dsp-level:0
\[
(8)\quad\textbf{the naming bar:}
\]

%dsp-level:0
\[
\qquad \text{the two heaviest-floor triples the question names are external species vocabulary};
\]

%dsp-level:0
\[
\qquad \text{the identification layer supplies}\ \textit{up-type}/\textit{down-type}\ \text{branch names}\ [\text{A-BM.6};\ \mathrm{Def\text{-}BM.14}(8)]\ \text{and generation indices},
\]

%dsp-level:0
\[
\text{and}\ \textbf{no name attaches to a composition}
\]

%dsp-level:0
\[
\qquad [\mathrm{Def\text{-}BM.6}(5),(7);\ \mathrm{Def\text{-}BM.3}(7)]:
\]

%dsp-level:0
\[
\qquad \boxed{\ \neg\text{either triple is formed, and}\ \neg\text{either is excluded}\ —\ \text{the composition is unread}\ (\mathrm{R3})\ }
\]

%dsp-level:0
\[
(9)\quad\textbf{status:}
\]

%dsp-level:0
\[
\qquad \textbf{derived:}\ (1)\text{--}(3),\ (5),\ (6)\text{'s formed rows},\ (7)\text{'s dimension rows};
\]

%dsp-level:0
\[
\qquad \textbf{refuted as formed:}\ \text{the ensemble matrix}\ (4),
\]

%dsp-level:0
\[
\ \text{an eigenvalue stability condition}\ (6),
\]

%dsp-level:0
\[
\ \text{the octet read}\ (7);
\]

%dsp-level:0
\[
\textbf{standing:}\ (\mathrm{R3}),\ \text{deficit}\ 1';
\]

%dsp-level:0
\[
\qquad \neg\text{a declaration},
\]

%dsp-level:0
\[
\quad\neg\text{a P4 addition},
\]

%dsp-level:0
\[
\quad\neg\text{a new dimensionful scale}\ [\mathrm{P1}],
\]

%dsp-level:0
\[
\quad\neg\text{a measured value}\ [\mathrm{P5}],
\]

%dsp-level:0
\[
\quad\neg\text{a species identification}\ [\mathrm{P5};\ \mathrm{Def\text{-}BM.3}(7)];
\]

\textbf{Req.} D2.11, D3.4, D3.9 (Parts 1–2); A §0.2a, CT3.15, CT3.17, CT3.17a(2), CT3.22, EM.11 (Part 2); P1, P4, P5 (Parts 1–2); Def-H.20, Def-H.41, Def-H.42, FT.34(5), FT.66(4), FT.80(5), FT.111(2), FT.111(3), FT.125(5), FT.142(3), FT.142(5), FT.143(2), FT.143(5), FT.143(6), FT.158(1), FT.158(4) (Part 3); Census FT.35(5), FT.63(6) (row-wise); Def-BM.6(4), Def-BM.6(5), Def-BM.6(7), Def-BM.8(1), Def-BM.8(5), Def-BM.9, Def-BM.13(3), Def-BM.13(6), Def-BM.13(9), Def-BM.14(1), Def-BM.14(2), Def-BM.14(3), Def-BM.14(5), Def-BM.14(8), Def-BM.3(7), BM.23(8), BM.28(2), BM.28(4), BM.28(7), BM.30(2), BM.31, BM.36(2), BM.36(3), BM.41(2), BM.46(2), BM.60(1), BM.61(8), BM.69(3), BM.70(2), BM.75(6), BM.80(1), BM.80(2), BM.80(3), BM.86(2), BM.90(3), BM.94(5), BM.95(1), BM.99(3) (this part); exact Laurent and polynomial identities over $\mathbb Q[z]$, binomial closed forms against exhaustive enumeration, Vieta reads, exact radical minimal polynomials in $\mathbb Q(\sqrt2)$ and $\mathbb Q(\sqrt3)$, elimination in a quotient ring [math; machine: \texttt{checks/\allowbreak{}verify\_\allowbreak{}H\_\allowbreak{}six\_\allowbreak{}objective\_\allowbreak{}resolution.py} §C].

\textbf{Pf.}

%dsp-level:0
\[
(1):\ \mathrm{Def\text{-}BM.14}(1)\text{--}(5)\ \text{and}\ \text{Census FT.63}(6)\ \text{re-cited (frozen)};
\]

%dsp-level:0
\[
\qquad (2-\sqrt2)^2-4(2-\sqrt2)+2=0\ \text{in}\ \mathbb Q(\sqrt2);
\]

%dsp-level:0
\[
\qquad (2-\sqrt3)^2-4(2-\sqrt3)+1=0\ \text{in}\ \mathbb Q(\sqrt3)\ [\text{machine: exact}]
\]

%dsp-level:0
\[
2-2\cos\tfrac\pi{10}=y\ \text{satisfies}\ y^4-8y^3+19y^2-12y+1=0:
\]

%dsp-level:0
\[
\qquad (2-y)^2
\]

%dsp-level:0
\[
= \tfrac{5+\sqrt5}2\implies\sqrt5
\]

%dsp-level:0
\[
= 2y^2-8y+3
\]

%dsp-level:0
\[
\implies 5
\]

%dsp-level:0
\[
= (2y^2-8y+3)^2
\]

%dsp-level:0
\[
\implies 4\bigl(y^4-8y^3+19y^2-12y+1\bigr)
\]

%dsp-level:0
\[
= 0\ [\text{machine: the expansion in}\ \mathbb Q[y]]
\]

%dsp-level:0
\[
(2):\ \mathrm{BM.80}(2),(3)\text{'s rows re-cited};
\]

%dsp-level:0
\[
\qquad \text{Vieta:}\ \sum_iz_i
\]

%dsp-level:0
\[
= \tfrac{24}4
\]

%dsp-level:0
\[
= 6,
\]

%dsp-level:0
\[
\ \ \sum_{i<j}z_iz_j
\]

%dsp-level:0
\[
= \tfrac{21-12\gamma}4,
\]

%dsp-level:0
\[
\ \ \prod_iz_i
\]

%dsp-level:0
\[
= -\tfrac{-(\gamma-1)^2}4
\]

%dsp-level:0
\[
= \tfrac{(\gamma-1)^2}4\ [\text{machine: the coefficient reads}]
\]

%dsp-level:0
\[
2(r-1)^3-3(r-1)=2r^3-6r^2+3r-1\implies2r^3-6r^2+3r+(1-\gamma)=0\ [\text{machine}];
\]

%dsp-level:0
\[
\qquad z(2z+3)^2-(6z+\gamma-1)^2
\]

%dsp-level:0
\[
\equiv 0\ \text{modulo}\ 2r^3-6r^2+3r+1-\gamma\ \text{at}\ r(2z+3)
\]

%dsp-level:0
\[
= 6z+\gamma-1\ [\text{machine: the quotient-ring reduction}]
\]

%dsp-level:0
\[
(3):\ 2+z+z^{-1}
\]

%dsp-level:0
\[
= z^{-1}(z^2+2z+1)
\]

%dsp-level:0
\[
= z^{-1}(z+1)^2
\]

%dsp-level:0
\[
\implies (2+z+z^{-1})^3
\]

%dsp-level:0
\[
= z^{-3}(z+1)^6\ [\text{machine: the Laurent identity}];
\]

%dsp-level:0
\[
\qquad [z^Q]z^{-3}(z+1)^6=[z^{Q+3}](z+1)^6=\binom6{Q+3}\ [\text{machine: all seven classes}]
\]

%dsp-level:0
\[
\binom6{Q+3}
\]

%dsp-level:0
\[
= (1,6,15,20,15,6,1)\ \text{against the exhaustive}\ 4^3\ \text{end-sign census}\ [\mathrm{BM.28}(2)\ \text{re-derived};\ \text{machine}];
\]

%dsp-level:0
\[
\qquad \sum_{k=0}^{6}\binom6k=2^6=64\ [\text{machine}]
\]

%dsp-level:0
\[
(4):\ \text{the five referent rows re-cited from their own loci}\ [\mathrm{Def\text{-}H.41}\ \text{Species};\ \mathrm{Def\text{-}H.42};\ \mathrm{FT.80}(5);\ \mathrm{FT.125}(5);\ \mathrm{BM.28}(7);\ \mathrm{Def\text{-}BM.13}(9);\ \mathrm{FT.142}(3)];
\]

%dsp-level:0
\[
\qquad \text{the inventory scans}\ [\text{machine}]
\]

%dsp-level:0
\[
(5):\ \mathrm{D3.9},\ \mathrm{CT3.15},\ \mathrm{CT3.17},\ \mathrm{FT.34}(5),\ \text{Census FT.35}(5)\ \text{re-cited};
\]

%dsp-level:0
\[
\qquad W(0)/W(+1)=20/15=4/3\ [\text{machine}];
\]

%dsp-level:0
\[
\qquad (6): (\mathrm i)\text{--}(\mathrm{iii})\ \text{re-cited at their own loci}\ [\mathrm{Def\text{-}H.20};\ \mathrm{FT.66}(4);\ \mathrm{BM.36}(3);\ \mathrm{BM.61}(8)(\mathrm{O1});\ \mathrm{BM.90}(3)]
\]

%dsp-level:0
\[
(7):\ g(5)=10\ \text{from}\ \mathrm{CT3.17a}(2)\text{'s summation}\ [\mathrm{Def\text{-}BM.9}\ \text{re-cited};\ \text{machine}];
\]

%dsp-level:0
\[
\qquad \text{the capacity ratios}\ \bigl(\tfrac13,8,\tfrac1{25}\bigr)\ \text{from the flag}\ (6,16,17,25)\ [\mathrm{BM.36}(2);\ \text{machine}]
\]

%dsp-level:0
\[
(8)\text{--}(9):\ \text{the naming-scope rows re-cited}\ [\mathrm{Def\text{-}BM.3}(7);\ \mathrm{Def\text{-}BM.6}(5),(7);\ \text{A-BM.6}];
\]

%dsp-level:0
\[
\qquad \text{the standing-rule scans}\ [\text{machine}].\ \blacksquare
\]

\textbf{Phys.} Ask the chart for the spectrum of a three-quark object and the honest answer is a table of what exists and what does not. What exists is two eigenvalue equations, and neither is the one the question wants. The first is the family cubic: four z cubed minus twenty-four z squared plus a channel constant, whose three roots are the three lines of a channel, with a phase-free sum of six and an exact product — a spectrum of generations, not of constituents. The second is the transport assembly's own characteristic polynomial, which collapses to a binomial: two plus z plus one over z, cubed, is z plus one to the sixth over z cubed, so the seven class weights are the binomial coefficients of six, and the assembly's spectrum is a single root of multiplicity six, one per directed edge. That is charge-class data, and it is the deepest eigenvalue statement the chart supports. What does not exist is the object the question names: a three-constituent state space, an exchange operator, a flavour index, a triple-level resolution — the counting spaces the corpus declares have no state on them, and the twenty-seven-dimensional candidate was rejected as an arena at rung one. So the determinant is not formed, and the two triples the question asks about are neither formed nor excluded; the composition is unread, and that is a named residue rather than an oversight. What is derived instead is sharper: the seal grammar forces exactly three, the charge balance forces their sum, the binomial forces the seven classes and the four-thirds lead of the neutral one, and the stability question lands where the corpus has always landed — on an undeclared spectral filter and a refuted index-zero route. The octet is not derivable either: the three here is colour's, the eight is a capacity ratio barred as a coupling, the ten is a grammar count, and the flavour symmetry that does exist is cyclic of order three.

\end{theorem}

\begin{theorem}[{\texorpdfstring{BM.104 — The curvature composition: the back-reaction fixed point, its exact roots and pole locus, and the non-selection theorem}{BM.104 - The curvature composition: the back-reaction fixed point, its exact roots and pole locus, and the non-selection theorem}}]\label{tag:BM-104}
\textbf{Thm.}

%dsp-level:0
\[
(1)\quad\textbf{the loop, formed exactly:}
\]

%dsp-level:0
\[
\qquad m_{\mathrm{corr}}=m_{\mathrm{hid}}+E_{\mathrm{EM}}^{\mathrm{self}}=m_{\mathrm{hid}}\bigl(1+a\,\mathsf g\bigr),
\]

%dsp-level:0
\[
\qquad a:=\frac12\,N\,w\,\varepsilon_{\mathrm{EM}}\,Q^2
\]

%dsp-level:0
\[
\qquad [\mathrm{BM.89}(5);\ \mathrm{Def\text{-}BM.13}(5);\ \text{machine}]:
\]

%dsp-level:0
\[
\rho_{\mathrm{corr}}=\frac{m_{\mathrm{corr}}}{V_{\mathrm{eff}}},
\]

%dsp-level:0
\[
\qquad V_{\mathrm{eff}}=(\Delta x_{\mathrm B})^3,
\]

%dsp-level:0
\[
\qquad \Delta x_{\mathrm B}\cdot m_{\mathrm{hid}}=1
\]

%dsp-level:0
\[
\qquad [\mathrm{Def\text{-}BM.4}(1);\ \mathrm{BM.14}(2);\ \text{the width is the informational one}]:
\]

%dsp-level:0
\[
\boxed{\ G_{\mathrm{corr}}\ :=\ (\Delta x_{\mathrm B})^2\kappa^2\rho_{\mathrm{corr}}\ =\ G\,\bigl(1+a\,\mathsf g\bigr),\qquad G=2\pi E_{\mathrm{int}}^2\ln2\ }
\]

%dsp-level:0
\[
\qquad [\mathrm{BM.88}(3);\ \mathrm{BM.16}(1);\ \text{machine: one monomial}];
\]

%dsp-level:0
\[
\textbf{the consistency condition is the fixed point:}\qquad\boxed{\ \mathsf g\ =\ \mathsf g\bigl(G\,(1+a\,\mathsf g)\bigr)\ }
\]

%dsp-level:0
\[
(2)\quad\textbf{the three exact solution rows:}
\]

%dsp-level:0
\[
(\mathrm i)\quad\mathsf g\equiv1:\qquad1=1\ \ \text{identically}\ —\ \ \text{a fixed point at every}\ a\ge0\ \ \text{and every}\ G\qquad[\text{machine}];
\]

%dsp-level:0
\[
(\mathrm{ii})\quad\mathsf g=1+G:
\]

%dsp-level:0
\[
\qquad \mathsf g=1+G(1+a\mathsf g)\iff\mathsf g\,(1-aG)=1+G\iff\boxed{\ \mathsf g_{\mathrm{fp}}=\frac{1+G}{1-aG}\ },
\]

%dsp-level:0
\[
\qquad \text{admissible}\iff aG<1,
\]

%dsp-level:0
\[
\qquad \mathsf g_{\mathrm{fp}}>1+G
\]

%dsp-level:0
\[
[\text{machine: substitution verifies the fixed point identically over}\ \mathbb Q[G,a]]
\]

%dsp-level:0
\[
\qquad \textbf{the back-reaction inflates the composition and diverges at the pole};
\]

%dsp-level:0
\[
(\mathrm{iii})\quad\mathsf g=(1+G)^{-1}:\qquad\mathsf g\bigl(1+G+aG\,\mathsf g\bigr)=1\iff\boxed{\ aG\,\mathsf g^2+(1+G)\,\mathsf g-1=0\ }
\]

%dsp-level:0
\[
\text{the product of the two roots is}\ -\tfrac1{aG}<0\implies\textbf{exactly one positive root}
\]

%dsp-level:0
\[
\qquad \mathsf g_{\mathrm{fp}}=\frac{\sqrt{(1+G)^2+4aG}-(1+G)}{2aG}\ \ \longrightarrow\ \ \frac1{1+G}\ \ \text{as}\ \ a\to0
\]

%dsp-level:0
\[
\qquad [\text{machine: the exact root row and the limit}];
\]

%dsp-level:0
\[
\boxed{\ \textbf{the fixed-point equation has a unique admissible solution for each of the three priced compositions}\ —\ \text{consistency is satisfied three times over and selects nothing}\ }
\]

%dsp-level:0
\[
(3)\quad\textbf{the pole locus:}
\]

%dsp-level:0
\[
\qquad aG=1\iff\frac12Nw\,\varepsilon_{\mathrm{EM}}Q^2\cdot2\pi E_{\mathrm{int}}^2\ln2=1\iff\boxed{\ \pi\,N\,w\,\varepsilon_{\mathrm{EM}}\,Q^2\,E_{\mathrm{int}}^2\ln2\ =\ 1\ }
\]

%dsp-level:0
\[
\qquad [\text{machine: one monomial condition}];
\]

%dsp-level:0
\[
Q=0\implies aG=0\ne1
\]

%dsp-level:0
\[
\qquad [\text{machine}]\implies\boxed{\ \textbf{the neutral class never reaches the pole}\ }
\]

%dsp-level:0
\[
\qquad [\mathrm{BM.17}(2):\ \text{the neutron row is}\ Q=0;\ \mathrm{BM.89}(5):\ E_{\mathrm{EM}}^{\mathrm{self}}(0)=0];
\]

%dsp-level:0
\[
Q=\pm1,\ E_{\mathrm{int}}=6:
\]

%dsp-level:0
\[
\qquad \varepsilon_{\mathrm{EM}}^{\mathrm{pole}}=\frac1{36\,\pi\,N\,w\ln2}
\]

%dsp-level:0
\[
\qquad [\text{machine: exact}];
\]

%dsp-level:0
\[
\qquad E_{\mathrm{int}}\in\{2,6,12,24,60\}\ \ \text{on the computed cells}
\]

%dsp-level:0
\[
\qquad [\mathrm{BM.13}(2);\ \mathrm{Census\ FT.35}(2),\ \mathrm{FT.42}(5)];
\]

%dsp-level:0
\[
(4)\quad\textbf{the loop is excluded by the lock's own freedom census:}
\]

%dsp-level:0
\[
\qquad G\text{'s derived alphabet is}\ \{E_{\mathrm{int}},\ \ln2,\ \pi\}\ —\ \ \mathfrak h\text{-free},\ \Lambda\text{-free},\ \mathrm{P4}\text{-free}
\]

%dsp-level:0
\[
\qquad [\mathrm{BM.88}(3);\ \mathrm{BM.16}(1);\ \text{machine: the symbol census}]:
\]

%dsp-level:0
\[
a\text{'s alphabet is}\ \{N,\ w,\ \varepsilon_{\mathrm{EM}},\ Q\},
\]

%dsp-level:0
\[
\qquad \varepsilon_{\mathrm{EM}}=e^2/(4\pi)\ \ \text{is a dimensionless P4-class symbol}
\]

%dsp-level:0
\[
\qquad [\mathrm{Def\text{-}BM.12}\ \text{local notation};\ \mathrm{EM.25}(9);\ \mathrm{P4}]:
\]

%dsp-level:0
\[
\boxed{\ \{E_{\mathrm{int}},\ln2,\pi\}\ \cap\ \{N,\ w,\ \varepsilon_{\mathrm{EM}},\ Q\}\ =\ \varnothing\ }
\]

%dsp-level:0
\[
\qquad [\text{machine}]\implies\textbf{admitting the back-reaction would make the lock's argument P4-dependent, contradicting}\ \mathrm{BM.88}(3)\text{'s derived row}
\]

%dsp-level:0
\[
\boxed{\ \textbf{the lock's argument is the informational mass, the EM self-energy does not re-enter}\ G,\ \text{and}\ \neg\textbf{a self-consistency condition on}\ \mathsf g\ \textbf{is formed}\ }
\]

%dsp-level:0
\[
(5)\quad\textbf{the non-selection theorem:}
\]

%dsp-level:0
\[
\qquad \boxed{\ \neg\exists\ \text{a derived}\ \mathsf g(G)\ }
\]

%dsp-level:0
\[
\qquad \text{on four grounds, each re-cited at its own locus:}
\]

%dsp-level:0
\[
(\mathrm i)\quad\neg\exists\ \text{a count}\to\text{strength theorem}
\]

%dsp-level:0
\[
\qquad [\mathrm{BM.36}(3);\ \mathrm{BM.61}(8)(\mathrm{O1});\ \mathrm{BM.69}(4)(\mathrm{iv})];
\]

%dsp-level:0
\[
\qquad (\mathrm{ii})\quad\text{the four grounds that bar a defect-conditioned weighting bar a curvature-conditioned one at the same species}
\]

%dsp-level:0
\[
[\mathrm{BM.90}(3)(\mathrm i)\text{--}(\mathrm{iv}):\ \text{the degree-}16\ \text{field against integer and binary reads}\ [\mathrm{FT.73}(2),(3)];\ \text{the two-torsion image}\ \{0,\pi\}\ [\mathrm{BM.67}(6)];\ \text{the calibration weight}\ {-}2\ \text{against}\ 0\ [\mathrm{FT.55}(1);\ \mathrm{BM.67}(8)];\ \text{the strength bar}];
\]

%dsp-level:0
\[
(\mathrm{iii})\quad\text{the consistency condition of}\ (2)\ \text{is satisfied by all three candidates}\ —\ \text{it is necessary and}\ \neg\text{sufficient}
\]

%dsp-level:0
\[
\qquad [\text{machine}];
\]

%dsp-level:0
\[
(\mathrm{iv})\quad\text{a selection would be a P4-class datum}\ —\ \text{§Cal-BM would gain a second member}\ —\ \text{refused}
\]

%dsp-level:0
\[
\qquad [\mathrm{P4};\ \mathrm{P1};\ \mathrm{BM.88}(4):\ \text{priced three ways, adopted in none}];
\]

%dsp-level:0
\[
(6)\quad\textbf{the composition is}\ \neg\textbf{inert inside the corpus:}
\]

%dsp-level:0
\[
\qquad \mathrm{BM.101}(5)\text{'s crossing count differs by composition}\ —\ (\mathrm i)\ \text{at most one},
\]

%dsp-level:0
\[
\ (\mathrm{ii})\ \text{exactly one always},
\]

%dsp-level:0
\[
\ (\mathrm{iii})\ \text{two, one or none}:
\]

%dsp-level:0
\[
\boxed{\ \textbf{the priced composition is observable in the energy-scale graph's topology, and is still}\ \neg\textbf{selected}\ }
\]

%dsp-level:0
\[
\qquad \textbf{the discrimination is derived; the selection is not};
\]

%dsp-level:0
\[
\textbf{every row of the corpus holds uniformly over the priced set:}
\]

%dsp-level:0
\[
\qquad \mathrm{BM.91}(2),(4)\text{--}(8),\ \mathrm{BM.92}(1),(2),\ \mathrm{BM.98}(1)\text{--}(4),\ \mathrm{BM.99}(1)\text{--}(6),\ \mathrm{BM.101}(2),(3),(7)
\]

%dsp-level:0
\[
\qquad [\text{machine: the uniformity censuses}];
\]

%dsp-level:0
\[
(7)\quad\textbf{the running coupling's status:}\qquad e(r,E,N)=e_0\sqrt{N\,f\bigl(r/\Delta x_{\mathrm B}(E)\bigr)\,\mathsf g\bigl(G(E)\bigr)}\qquad[\mathrm{BM.88}(4)]:
\]

%dsp-level:0
\[
f\ \ \text{derived exactly}\ [\mathrm{BM.88}(2)(\mathrm{ii})];
\]

%dsp-level:0
\[
\qquad G\ \ \text{derived exactly, one monomial}\ [\mathrm{BM.88}(3)];
\]

%dsp-level:0
\[
\qquad N\ \ \text{solution data}\ [\mathrm{CT2.7};\ \mathrm{BM.34}(2);\ \mathrm{BM.88}(8)];
\]

%dsp-level:0
\[
\mathsf g\ \ \text{priced three ways, adopted in none}\ (5);
\]

%dsp-level:0
\[
\qquad e_0\ \ \text{open}\ [\mathrm{EM.25}(7)]:
\]

%dsp-level:0
\[
\qquad \boxed{\ \textbf{the coupling is complete in form and carries exactly two undeclared members}\ (e_0,\ \mathsf g),\ \neg\text{a third}\ }
\]

%dsp-level:0
\[
(8)\quad\textbf{status:}
\]

%dsp-level:0
\[
\qquad \textbf{derived:}\ (1)\text{--}(4),\ (6)\text{'s discrimination rows};
\]

%dsp-level:0
\[
\qquad \textbf{refuted as derivable:}\ \mathsf g(G)\ (5);
\]

%dsp-level:0
\[
\qquad \textbf{rescoped,}\ \neg\textbf{retracted:}\ \mathrm{BM.88}(4),(8),\ \mathrm{BM.89}(5),\ \mathrm{BM.91},\ \mathrm{BM.98}\ \text{stand at their own loci};
\]

%dsp-level:0
\[
\neg\text{a declaration},
\]

%dsp-level:0
\[
\quad\neg\text{a P4 addition},
\]

%dsp-level:0
\[
\quad\neg\text{a new dimensionful scale}\ [\mathrm{P1}],
\]

%dsp-level:0
\[
\quad\neg\text{a measured value}\ [\mathrm{P5}],
\]

%dsp-level:0
\[
\quad\neg\text{a numerical coupling}\ [\mathrm{EM.25}(7),(9)];
\]

%dsp-level:0
\[
\alpha_B:\ \neg\text{added},\ \neg\text{derived},\ \text{barred}\ [\mathrm{BM.61}(8)];\qquad
\]

\textbf{Req.} D2.9 (Parts 1–2); CT2.1, CT2.7, EM.25(7), EM.25(9) (Part 2); P1, P4, P5 (Parts 1–2); Def-H.37(2), FT.54(4), FT.55(1), FT.73(2), FT.73(3) (Part 3); Def-BM.4(1), Def-BM.12, BM.13(2), BM.14(2), BM.16(1), BM.17(2), BM.34(2), BM.36(3), BM.61(8), BM.67(6), BM.67(8), BM.69(4), BM.88(2), BM.88(3), BM.88(4), BM.88(5), BM.88(8), BM.89(5), BM.90(3), BM.91, BM.92(1), BM.92(2), BM.98, BM.101(5), BM.101(7), Def-BM.13(5) (this part); Census FT.35(2), FT.42(5) (row-wise); exact polynomial arithmetic over $\mathbb Q[G,a]$, quadratic root rows with sign censuses, exact limit rows, symbol censuses and set intersections, monomial conditions [math; machine: \texttt{checks/\allowbreak{}verify\_\allowbreak{}H\_\allowbreak{}six\_\allowbreak{}objective\_\allowbreak{}resolution.py} §D].

\textbf{Pf.}

%dsp-level:0
\[
(1):\ E_{\mathrm{EM}}^{\mathrm{self}}=\tfrac12Nw\,\mathsf g\,\varepsilon_{\mathrm{EM}}Q^2m_{\mathrm{hid}}=a\,\mathsf g\,m_{\mathrm{hid}}\ [\mathrm{BM.89}(5);\ \text{machine}];
\]

%dsp-level:0
\[
\qquad \rho_{\mathrm{corr}}=m_{\mathrm{corr}}/V_{\mathrm{eff}}=\rho\,(1+a\mathsf g)\ \text{at fixed}\ V_{\mathrm{eff}}\ [\mathrm{Def\text{-}BM.4}(1);\ \text{machine}]
\]

%dsp-level:0
\[
G_{\mathrm{corr}}=(\Delta x_{\mathrm B})^2\kappa^2\rho_{\mathrm{corr}}=G(1+a\mathsf g)\ [\text{machine: one monomial}];
\]

%dsp-level:0
\[
\qquad \text{the consistency condition is}\ \mathsf g=\mathsf g(G_{\mathrm{corr}})\ [\text{the fixed point, stated}]
\]

%dsp-level:0
\[
(2)(\mathrm i):\ 1=1\ [\text{machine}];
\]

%dsp-level:0
\[
\qquad (\mathrm{ii}):\ \mathsf g=1+G+aG\mathsf g\iff\mathsf g(1-aG)=1+G\ [\text{machine: cleared over}\ \mathbb Q[G,a]];
\]

%dsp-level:0
\[
\mathsf g_{\mathrm{fp}}>1+G\iff\tfrac{1+G}{1-aG}>1+G\iff\tfrac1{1-aG}>1
\]

%dsp-level:0
\[
\iff 0<aG<1\ [\text{machine: the sign census on positive witnesses}]
\]

%dsp-level:0
\[
(\mathrm{iii}):\ \mathsf g=\tfrac1{1+G+aG\mathsf g}\iff\mathsf g(1+G)+aG\mathsf g^2=1\ [\text{machine}];
\]

%dsp-level:0
\[
\qquad \text{the root product}\ -1/(aG)<0\implies\text{one positive root}\ [\text{machine}];
\]

%dsp-level:0
\[
\qquad a\to0:\ \tfrac{\sqrt{(1+G)^2+4aG}-(1+G)}{2aG}
\]

%dsp-level:0
\[
= \tfrac{4aG}{2aG\bigl(\sqrt{(1+G)^2+4aG}+(1+G)\bigr)}\to\tfrac1{1+G}\ [\text{machine: the conjugate-cleared limit}]
\]

%dsp-level:0
\[
(3):\ aG=\tfrac12Nw\,\varepsilon_{\mathrm{EM}}Q^2\cdot2\pi E^2\ln2=\pi Nw\,\varepsilon_{\mathrm{EM}}Q^2E^2\ln2\ [\text{machine}];
\]

%dsp-level:0
\[
\qquad Q=0\implies aG=0\ne1\ [\text{machine}];
\]

%dsp-level:0
\[
\qquad Q^2=1,\ E=6:\ \varepsilon^{\mathrm{pole}}=1/(36\pi Nw\ln2)\ [\text{machine}]
\]

%dsp-level:0
\[
(4):\ \text{the two symbol sets enumerated from their defining clauses and intersected}\ [\mathrm{BM.88}(3);\ \mathrm{Def\text{-}BM.12};\ \text{machine}];
\]

%dsp-level:0
\[
\qquad \mathrm{BM.88}(3)\text{'s P4-freedom row re-cited}
\]

%dsp-level:0
\[
(5):\ (\mathrm i)\text{--}(\mathrm{iv})\ \text{re-cited at their own loci}\ [\mathrm{BM.36}(3);\ \mathrm{BM.61}(8);\ \mathrm{BM.90}(3);\ \mathrm{FT.73}(2),(3);\ \mathrm{BM.67}(6),(8);\ \mathrm{FT.55}(1);\ \mathrm{P4}]
\]

%dsp-level:0
\[
(6):\ \mathrm{BM.101}(5)(\mathrm i)\text{--}(\mathrm{iii})\text{'s three crossing counts re-cited};
\]

%dsp-level:0
\[
\qquad \text{the uniformity censuses over}\ \mathsf g\in\{1,1+G,(1+G)^{-1}\}\ \text{on exact rational lattices}\ [\text{machine}]
\]

%dsp-level:0
\[
(7)\text{--}(8):\ \mathrm{BM.88}(2)\text{--}(4),(8)\text{'s rows re-cited};
\]

%dsp-level:0
\[
\qquad \text{the standing-rule scans, including the barred object}\ \alpha_B\ [\text{machine}].\ \blacksquare
\]

\textbf{Phys.} Turn the correction back on itself and the corpus answers with an equation it can solve and a verdict it cannot avoid. If the electromagnetic self-energy feeds the density that feeds the curvature lock, the lock's argument grows by one plus a times the composition, and self-consistency becomes a fixed point. Solved on the three priced candidates, the fixed point is trivial for the constant one, a rational expression with an exact pole for the growing one, and the positive root of a quadratic for the shrinking one — three admissible solutions, so consistency is necessary and useless as a selector. The pole locus is a single monomial, and the neutral class never reaches it, because a neutral composite carries no field energy at all. Then the deeper row: the loop is not merely unselected, it is excluded, because the lock's alphabet is three symbols long and disjoint from the correction's, and letting the correction in would make the lock depend on a declared parameter that the lock's own theorem says it does not carry. So the self-energy rides the mass and not the geometry, the composition stays priced three ways, and the four old grounds against turning a count into a strength bar the fourth route. What the tier gains is a discrimination: the three compositions give the energy-scale graph three different topologies — a half-line, a half-line that always crosses, and a window — so the choice is not cosmetic inside this corpus even though nothing here can make it. The coupling is complete in form with exactly two undeclared members, and no third.

\end{theorem}

\begin{theorem}[{\texorpdfstring{BM.105 — The general crystallization weight: the measure-free bounds, the threshold family, and the refuted defect conditioning}{BM.105 - The general crystallization weight: the measure-free bounds, the threshold family, and the refuted defect conditioning}}]\label{tag:BM-105}
\textbf{Thm.}

%dsp-level:0
\[
(1)\quad\textbf{the measure-free form:}
\]

%dsp-level:0
\[
\qquad E_{\mathrm{int}}(R)=|R|=2\,|S(R)|\ \ \text{on the three pairs of}\ \Sigma_{12}
\]

%dsp-level:0
\[
\qquad [\mathrm{Def\text{-}BM.5}(4);\ \mathrm{Census\ FT.35}(2);\ \mathrm{BM.94}(4)]:
\]

%dsp-level:0
\[
\text{linearity on}\ |R|\ \ \text{gives, for}\ \textbf{every}\ \ \text{ensemble measure}:
\]

%dsp-level:0
\[
\qquad \boxed{\ w=\frac1{E^{\mathrm{cryst}}}\sum_{i=1}^{3}P(\text{the pair}\ i\ \text{is real}),\qquad w_{\mathrm{mat}}=\frac13\sum_{i=1}^{3}a_i,\qquad a_i:=P(\text{the pair}\ i\ \text{is real}\mid\text{materialized})\ }
\]

%dsp-level:0
\[
[\mathrm{BM.89}(4);\ \mathrm{BM.94}(7);\ \text{machine: the linearity row}];\qquad\textbf{no distribution hypothesis enters}\ (1);
\]

%dsp-level:0
\[
(2)\quad\textbf{the measure-free bounds:}
\]

%dsp-level:0
\[
\qquad |S|
\]

%dsp-level:0
\[
\ge 1\ \ \text{on the materialized support}\ [\mathrm{BM.94}(5);\ \mathrm{Def\text{-}H.4};\ \mathrm{FT.144}(4)]\ \ \text{and}\ \ |S|
\]

%dsp-level:0
\[
\le 3
\]

%dsp-level:0
\[
\implies 1\le\langle|S|\rangle_{\mathrm{mat}}\qquad \le 3:
\]

%dsp-level:0
\[
\boxed{\ \frac13\ \le\ w_{\mathrm{mat}}\ \le\ 1\qquad\text{for every admissible ensemble measure}\ }\qquad[\text{machine}];
\]

%dsp-level:0
\[
\textbf{the two equality loci are exact:}
\]

%dsp-level:0
\[
\qquad w_{\mathrm{mat}}=\frac13\iff\text{the measure is supported on the single-pair configurations};
\]

%dsp-level:0
\[
\qquad w_{\mathrm{mat}}=1\iff\text{it is supported on the full configuration}
\]

%dsp-level:0
\[
\qquad [\text{machine}];
\]

%dsp-level:0
\[
\boxed{\ \langle E\rangle_{\mathrm{mat}}=6\,w_{\mathrm{mat}}\in[2,\,6]\ }
\]

%dsp-level:0
\[
\qquad \textbf{the bound is a consequence of the materialization restriction alone,}\ \neg\textbf{of the specialization};
\]

%dsp-level:0
\[
(3)\quad\textbf{the threshold family (hypothesis}\ (\mathrm F)\textbf{, marked):}
\]

%dsp-level:0
\[
\qquad \textbf{hypothesis:}\ \text{the A4 threshold takes a per-pair value}\ \varepsilon_i\in[0,\tfrac14)\ \ \text{on the three pairs of}\ \Sigma_{12}
\]

%dsp-level:0
\[
\qquad [\mathrm{A4};\ \mathrm{D0.4};\ \mathrm{CT0.8}\text{'s derived range applied memberwise};\ \text{carried symbolically},\ \neg\text{issued}]:
\]

%dsp-level:0
\[
p_i:=1-4\varepsilon_i,\qquad u_i:=1-p_i=4\varepsilon_i,\qquad R:=\prod_{i=1}^{3}u_i,\qquad S:=\sum_{i=1}^{3}p_i\qquad[\text{machine}]:
\]

%dsp-level:0
\[
\boxed{\ w_{\mathrm{mat}}\ =\ \frac{S}{3\,(1-R)},\qquad\langle E\rangle_{\mathrm{mat}}=\frac{2S}{1-R},\qquad\bigl\langle E^2\bigr\rangle_{\mathrm{mat}}=\frac{4\bigl(S+2\sum_{i<j}p_ip_j\bigr)}{1-R}\ }
\]

%dsp-level:0
\[
\textbf{the seven materialized configuration weights:}
\]

%dsp-level:0
\[
\qquad P_{\mathrm{mat}}(S)=\frac{\prod_{i\in S}p_i\prod_{i\notin S}u_i}{1-R},
\]

%dsp-level:0
\[
\qquad \varnothing\ne S\subseteq\{1,2,3\},
\]

%dsp-level:0
\[
\qquad \sum_{\varnothing\ne S}P_{\mathrm{mat}}(S)=1
\]

%dsp-level:0
\[
\qquad [\text{machine: exact rational functions, normalized}];
\]

%dsp-level:0
\[
\textbf{the uniform specialization recovers}\ \mathrm{BM.94}\ \textbf{identically:}\qquad p_i\equiv p\implies S=3p,\quad R=(1-p)^3,\quad1-R=pD:
\]

%dsp-level:0
\[
\boxed{\ w_{\mathrm{mat}}=\frac{3p}{3pD}=\frac1D,\qquad D=3-3p+p^2\ }
\]

%dsp-level:0
\[
\qquad [\mathrm{BM.94}(6),(7);\ \text{machine: the identity at six rational thresholds}];
\]

%dsp-level:0
\[
(4)\quad\textbf{the family's exact bounds:}
\]

%dsp-level:0
\[
\qquad 3(1-R)-S=\sum_{i=1}^{3}u_i-3\prod_{i=1}^{3}u_i
\]

%dsp-level:0
\[
\qquad [\text{machine: one polynomial identity}]:
\]

%dsp-level:0
\[
\sum_iu_i\ \ge\ 3\bigl(\prod_iu_i\bigr)^{1/3}\ \ge\ 3\prod_iu_i
\]

%dsp-level:0
\[
\qquad [\prod_iu_i\le1;\ \text{AM--GM}]\implies\boxed{\ 3(1-R)\ \ge\ S\ }
\]

%dsp-level:0
\[
\qquad \text{with equality iff every}\ u_i=0,\ \text{i.e.}\ p_i=1:
\]

%dsp-level:0
\[
\boxed{\ w_{\mathrm{mat}}\le1\ \ \text{with equality iff every pair is crystallized};\qquad w_{\mathrm{mat}}>\frac13\ \ \text{strictly on the materialized domain}\ }
\]

%dsp-level:0
\[
\qquad [(2);\ \text{machine}];
\]

%dsp-level:0
\[
w_{\mathrm{mat}}\ \ \text{is strictly increasing in each}\ p_i\qquad[\text{machine: an exact rational grid on}\ (0,1]^3];
\]

%dsp-level:0
\[
(5)\quad\textbf{hypothesis}\ (\mathrm{W6}) — \textbf{six state-dependent thresholds}\ \varepsilon_i(D_S) — \textbf{on its face and refuted:}
\]

%dsp-level:0
\[
(\mathrm i)\quad\textbf{the alphabets are disjoint:}
\]

%dsp-level:0
\[
\qquad \text{the flicker measure's symbol set is}\ \{\varepsilon_q\}\ \ \text{and the defect values' is}\ \{\delta_v,\ \hat\mu'^2\}
\]

%dsp-level:0
\[
\qquad [\mathrm{BM.90}(1),(2);\ \text{machine: the symbol census of}\ \mathrm{A4},\ \mathrm{D0.4},\ \mathrm{CT0.8},\ \mathrm{BM.19}(2)\ \text{carries}\ \varepsilon_q\ \text{alone, with no state argument}];
\]

%dsp-level:0
\[
(\mathrm{ii})\quad\textbf{the four grounds bar it:}
\]

%dsp-level:0
\[
\qquad \text{the degree-}16\ \text{field against integer and binary reads}\ [\mathrm{FT.73}(2),(3)];
\]

%dsp-level:0
\[
\ \ \text{the two-torsion image reaching}\ 2\ \text{of}\ 6\ \text{states}\ [\mathrm{BM.67}(6)];
\]

%dsp-level:0
\[
\ \ \text{the calibration weight}\ {-}2\ \text{against}\ 0\ [\mathrm{FT.55}(1);\ \mathrm{BM.67}(8)];
\]

%dsp-level:0
\[
\ \ \text{the strength bar}\ [\mathrm{BM.36}(3);\ \mathrm{BM.61}(8)(\mathrm{O1})]
\]

%dsp-level:0
\[
\qquad [\mathrm{BM.90}(3)\ \text{re-cited}];
\]

%dsp-level:0
\[
(\mathrm{iii})\quad\textbf{the datum count:}
\]

%dsp-level:0
\[
\qquad \text{six thresholds would be six P4-class data}\ —\ \text{§Cal-BM would gain six members}\ —\ \text{refused}
\]

%dsp-level:0
\[
\qquad [\mathrm{P4};\ \mathrm{P1}];
\]

%dsp-level:0
\[
(\mathrm{iv})\quad\textbf{the scope row:}
\]

%dsp-level:0
\[
\qquad \mathrm{A4}\text{'s threshold enters the corpus only via}\ \mathrm{CT0.1}\ \text{and}\ \mathrm{CT0.2},
\]

%dsp-level:0
\[
\ \text{with a single}\ \varepsilon_q\ \text{and no state argument}
\]

%dsp-level:0
\[
\qquad [\mathrm{CT0.8}\ \text{scope clause}]:
\]

%dsp-level:0
\[
\boxed{\ \neg\exists\ \text{a derived map}\ D_S\longmapsto\varepsilon;\qquad\textbf{the weight stays branch-blind, and the six defect values act in the pole currency alone}\ }
\]

%dsp-level:0
\[
\qquad [\mathrm{BM.90}(3),(4);\ \mathrm{BM.98}(6)];
\]

%dsp-level:0
\[
(6)\quad\textbf{the residue, narrowed:}
\]

%dsp-level:0
\[
\qquad (\mathrm{R4})\ \ \textbf{closes at two levels:}\ \ \text{the measure-free bounds}\ (1),(2)\ \ \text{and the threshold family}\ (3),(4);
\]

%dsp-level:0
\[
\textbf{what remains is}\ (\mathrm{R4}'):
\]

%dsp-level:0
\[
\qquad \boxed{\ (\mathrm{R4}')\ \ \text{the ensemble's configuration distribution beyond per-pair independence}\ —\ \text{the pair--pair correlations of the measure}\ }
\]

%dsp-level:0
\[
\textbf{and the transport register's ensemble average}
\]

%dsp-level:0
\[
\qquad [\mathrm{BM.89}(8);\ \mathrm{BM.91}(2);\ \mathrm{BM.94}(8):\ \text{the computed cells average the vertex-deficit register}\ \tfrac32n_{\mathcal V},\ \neg\text{the edge-count register}\ |R|]:
\]

%dsp-level:0
\[
(\mathrm{R4})\ \ \textbf{is rescoped to}\ (\mathrm{R4}'),\ \neg\textbf{deleted};
\]

%dsp-level:0
\[
(7)\quad\textbf{status:}
\]

%dsp-level:0
\[
\qquad \textbf{derived:}\ (1)\text{--}(4),\ (6);
\]

%dsp-level:0
\[
\qquad \textbf{refuted:}\ \text{hypothesis}\ (\mathrm{W6})\ \ (5);
\]

%dsp-level:0
\[
\qquad \textbf{marked,}\ \neg\textbf{adopted:}\ \text{hypothesis}\ (\mathrm F)\text{'s family};
\]

%dsp-level:0
\[
\neg\text{a declaration},
\]

%dsp-level:0
\[
\quad\neg\text{a P4 addition},
\]

%dsp-level:0
\[
\quad\neg\text{a new dimensionful scale}\ [\mathrm{P1}],
\]

%dsp-level:0
\[
\quad\neg\text{a measured value}\ [\mathrm{P5}],
\]

%dsp-level:0
\[
\quad\text{§Cal-BM introduces no new free parameter at this tag};
\]

\textbf{Req.} A4 (Part 1); D0.4, D2.9 (Parts 1–2); A §0.2a, CT0.8, CT1.29 (Part 1); P1, P4, P5 (Parts 1–2); L1, L2 (Part 1); Def-H.4, Def-H.32(1), FT.55(1), FT.73(2), FT.73(3), FT.107(3), FT.144(4) (Part 3); Census FT.35(2) (row-wise); Def-BM.5(4), BM.19(1), BM.19(2), BM.36(3), BM.61(8), BM.67(6), BM.67(8), BM.89(4), BM.89(8), BM.90(1), BM.90(2), BM.90(3), BM.90(4), BM.91(2), BM.94(4), BM.94(5), BM.94(6), BM.94(7), BM.94(8), BM.98(6) (this part); exact polynomial identities over $\mathbb Q[p_1,p_2,p_3]$, exhaustive configuration censuses, AM–GM on three non-negative rationals, exact rational grids, symbol censuses [math; machine: \texttt{checks/\allowbreak{}verify\_\allowbreak{}H\_\allowbreak{}six\_\allowbreak{}objective\_\allowbreak{}resolution.py} §E].

\textbf{Pf.}

%dsp-level:0
\[
(1):\ E_{\mathrm{int}}
\]

%dsp-level:0
\[
= 2|S|\ \text{and linearity of the average}\implies\langle E\rangle
\]

%dsp-level:0
\[
= 2\sum_iP(\text{pair}\ i\ \text{real})\ [\mathrm{BM.89}(4)\text{'s per-datum form re-derived without any distribution hypothesis};\ \text{machine}]
\]

%dsp-level:0
\[
\text{the conditional form:}\ P(\text{pair}\ i\ \text{real}\mid\text{materialized})=a_i\ \text{and}\ \langle|S|\rangle_{\mathrm{mat}}=\sum_ia_i\ [\text{machine}]\implies w_{\mathrm{mat}}=\tfrac13\sum_ia_i
\]

%dsp-level:0
\[
(2):\ 1
\]

%dsp-level:0
\[
\le |S|
\]

%dsp-level:0
\[
\le 3\ \text{on the materialized support}\ [\mathrm{BM.94}(5);\ \mathrm{Def\text{-}H.4};\ \mathrm{FT.144}(4)]
\]

%dsp-level:0
\[
\implies 1\le\langle|S|\rangle_{\mathrm{mat}}
\]

%dsp-level:0
\[
\le 3\implies\tfrac13
\]

%dsp-level:0
\[
\le w_{\mathrm{mat}}
\]

%dsp-level:0
\[
\le 1\ [\text{machine}];
\]

%dsp-level:0
\[
\qquad \text{the two equality loci:}\ \langle|S|\rangle=1\iff|S|=1\ \text{a.s.},\ \ \langle|S|\rangle=3\iff|S|=3\ \text{a.s.}\ [\text{machine}]
\]

%dsp-level:0
\[
(3):\ P(S)
\]

%dsp-level:0
\[
= \prod_{i\in S}p_i\prod_{i\notin S}u_i\ \text{and}\ \sum_{S\ne\varnothing}P(S)
\]

%dsp-level:0
\[
= 1-\prod_iu_i
\]

%dsp-level:0
\[
= 1-R\ [\text{machine: the exhaustive eight-configuration census}];
\]

%dsp-level:0
\[
\qquad \langle E\rangle_{\mathrm{mat}}=\tfrac{2\sum_ip_i}{1-R}\ [\text{machine:}\ P(\text{pair}\ i\ \text{real})=p_i\ \text{since a real pair is materialized}]
\]

%dsp-level:0
\[
\langle|S|^2\rangle
\]

%dsp-level:0
\[
= \sum_ip_i+2\sum_{i<j}p_ip_j\ [\text{machine:}\ |S|^2=\sum_i1_i+2\sum_{i<j}1_i1_j]\implies\langle E^2\rangle_{\mathrm{mat}}
\]

%dsp-level:0
\[
= \tfrac{4(S+2\sum_{i<j}p_ip_j)}{1-R}
\]

%dsp-level:0
\[
p_i
\]

%dsp-level:0
\[
\equiv p:\ S
\]

%dsp-level:0
\[
= 3p,
\]

%dsp-level:0
\[
\ R
\]

%dsp-level:0
\[
= (1-p)^3,
\]

%dsp-level:0
\[
\ 1-R
\]

%dsp-level:0
\[
= p(3-3p+p^2)
\]

%dsp-level:0
\[
= pD
\]

%dsp-level:0
\[
\implies w_{\mathrm{mat}}
\]

%dsp-level:0
\[
= 3p/(3pD)
\]

%dsp-level:0
\[
= 1/D\ [\mathrm{BM.94}(6);\ \text{machine at six rational thresholds}]
\]

%dsp-level:0
\[
(4):\ 3(1-\prod u)-\sum(1-u)=3-3\prod u-3+\sum u=\sum u-3\prod u\ [\text{machine: one identity}];
\]

%dsp-level:0
\[
\qquad \sum u
\]

%dsp-level:0
\[
\ge 3(\prod u)^{1/3}\ [\text{AM--GM}]\ \ \text{and}\ \ (\prod u)^{1/3}\ge\prod u\ \text{on}\ \prod u
\]

%dsp-level:0
\[
\le 1\implies\sum u
\]

%dsp-level:0
\[
\ge 3\prod u\ [\text{machine: exact on a rational grid and by the two inequalities}]
\]

%dsp-level:0
\[
\text{equality}\iff u_1=u_2=u_3\ \wedge\ \prod u\in\{0,1\}\iff p_1=p_2=p_3=1\ \text{on the materialized domain}\ [\text{machine}];
\]

%dsp-level:0
\[
\qquad \text{monotonicity in each}\ p_i\ [\text{machine: exact rational grid}]
\]

%dsp-level:0
\[
(5):\ (\mathrm i)\text{--}(\mathrm{iv})\ \text{re-cited at their own loci}\ [\mathrm{BM.90}(1),(2),(3);\ \mathrm{FT.73}(2),(3);\ \mathrm{BM.67}(6),(8);\ \mathrm{FT.55}(1);\ \mathrm{BM.36}(3);\ \mathrm{BM.61}(8)(\mathrm{O1});\ \mathrm{CT0.8};\ \mathrm{P4}];
\]

%dsp-level:0
\[
\qquad \text{the symbol census}\ [\text{machine}]
\]

%dsp-level:0
\[
(6)\text{--}(7):\ \text{the residue rows assembled};\qquad\text{the standing-rule scans}\ [\text{machine}].\ \blacksquare
\]

\textbf{Phys.} The weight that prices a composite's failure to crystallize does not need the specialization that first produced it. Linearity alone gives it as the mean fraction of the three pairs that are real, and the materialization restriction alone gives its bounds: at least one pair is real, at most three are, so the weight lives between a third and one for every measure the ensemble could carry, with both endpoints characterized exactly — a third when only single pairs materialize, one when the full configuration does. That is a stronger statement than the specialization's, and it is why the mean edge count can never fall below the digon's two. Then the family: let each of the three pairs carry its own threshold, still inside the derived admissible range, and the weight becomes the sum of the three retention probabilities over three times one minus their product, with seven materialized configurations, exact second moments, and an arithmetic-mean-geometric-mean row that pins it below one with equality only when every pair is crystallized. Set the three equal and the whole of the previous tier's formula returns identically. What the family does not license is the version the question asks for: six thresholds, one per quark state, each a function of that state's defect value. That is the branch-conditioned weighting already refuted — the measure's alphabet holds one symbol and the defect values' holds two, four separate grounds bar any count-to-strength map, six thresholds would be six declared parameters, and the axiom's threshold carries no state argument. So the weight stays branch-blind, the six defect values stay in the pole currency, and the residue narrows rather than closes: what is left is the correlation between pairs, named and priced, and the transport register's own ensemble average.

\end{theorem}

\begin{theorem}[{\texorpdfstring{BM.106 — The arrow and the kind: the spin label, the tick law, and the exchange-structure referent census}{BM.106 - The arrow and the kind: the spin label, the tick law, and the exchange-structure referent census}}]\label{tag:BM-106}
\textbf{Thm.}

%dsp-level:0
\[
(1)\quad\textbf{the spin label (derived, on the logged transfer):}
\]

%dsp-level:0
\[
\qquad j=\frac{|k|}2,
\]

%dsp-level:0
\[
\qquad k=\frac{\Theta}2\in\mathbb Z_{\ge0}\ \ \text{on shell},
\]

%dsp-level:0
\[
\qquad \Psi(c)=\mathrm{Sym}^{|k(c)|}\bigl(\mathcal H_{\mathsf v}\bigr),
\]

%dsp-level:0
\[
\qquad \text{Casimir}\ \frac k2\Bigl(\frac k2+1\Bigr)
\]

%dsp-level:0
\[
[\mathrm{Def\text{-}H.38};\ \mathrm{FT.60}(1)\text{--}(3);\ \mathrm{FT.22};\ \mathrm{FT.25}(3);\ \text{the transfer's derivation or replacement is a named residue of the seed transport}];
\]

%dsp-level:0
\[
\textbf{the multiplet:}\qquad m\in\{k,\ k-2,\ \dots,\ -k\}\qquad[\mathrm{FT.60}(5);\ \mathrm{Def\text{-}H.38};\ \mathrm{BM.39}(7)];
\]

%dsp-level:0
\[
\textbf{the floor census:}
\]

%dsp-level:0
\[
\qquad \Theta_{\mathrm{int}}=2\kappa\in2\mathbb Z,
\]

%dsp-level:0
\[
\qquad k=\kappa\in\mathbb Z\setminus\{0\},
\]

%dsp-level:0
\[
\qquad \kappa=\widehat{ds}^{\,2}(e)\ \ \text{is registration data determined at crystallization}
\]

%dsp-level:0
\[
\qquad [\mathrm{Census\ FT.28}(2);\ \mathrm{FT.24};\ \text{A §0.2a}]:
\]

%dsp-level:0
\[
\boxed{\ \kappa=1\implies j=\frac12;\qquad\kappa=2\implies j=1\ }\qquad\textbf{on the same}\ \mathrm{D1.5}\ \textbf{kind}\qquad[\text{machine: the two exact rows}]:
\]

%dsp-level:0
\[
\boxed{\ \neg\exists\ \text{a spin-parity theorem}\ —\ \text{both integer and half-integer labels occur on materialized closed bound chains of one kind}\ }
\]

%dsp-level:0
\[
(2)\quad\textbf{the kind rows that are derived:}
\]

%dsp-level:0
\[
\qquad \text{bound}\ :=\ \text{closed self-contained chain},\ \text{content fermionic};
\]

%dsp-level:0
\[
\qquad \text{transit}\ :=\ \text{propagating open chain},\ \text{content bosonic}
\]

%dsp-level:0
\[
\qquad [\mathrm{D1.5}:\ \textbf{a definitional label}]:
\]

%dsp-level:0
\[
\text{the dichotomy is the closed/open read}
\]

%dsp-level:0
\[
\qquad [\mathrm{FT.80}(5);\ \mathrm{D1.4}];
\]

%dsp-level:0
\[
\qquad \text{the materialized locus at rung one is single-arrow}
\]

%dsp-level:0
\[
\qquad [\mathrm{FT.144}(2)]\ \ \text{and single-ray, occupancy}\ (1,0)
\]

%dsp-level:0
\[
\qquad [\text{A-BM.7};\ \mathrm{FT.144}(5);\ \mathrm{BM.73}(3);\ \mathrm{BM.96}(4)];
\]

%dsp-level:0
\[
\textbf{the tick law:}
\]

%dsp-level:0
\[
\qquad \tau_n\to\tau_{n+1},
\]

%dsp-level:0
\[
\qquad \neg\exists\,\tau^{-1},
\]

%dsp-level:0
\[
\qquad \bigl|\mathfrak L(\tau_{n+1})\bigr|=\bigl|\mathfrak L(\tau_n)\bigr|+n_{\mathrm{real\_out}}(\tau_n)\ge\bigl|\mathfrak L(\tau_n)\bigr|
\]

%dsp-level:0
\[
\qquad [\mathrm{A6};\ \mathrm{A7}]:
\]

%dsp-level:0
\[
\text{the prefix law}\ \#J^+(\text{prefix})\ge\#J^-(\text{prefix})\ \ \text{with strict alternation}
\]

%dsp-level:0
\[
\qquad [\mathrm{FT.144}(2);\ \mathrm{D1.2}];
\]

%dsp-level:0
\[
\qquad \text{the pure order reversal of every completed word is inadmissible at its first event}
\]

%dsp-level:0
\[
\qquad [\mathrm{FT.144}(4)(\mathrm{ii})];
\]

%dsp-level:0
\[
\textbf{the contraction:}
\]

%dsp-level:0
\[
\qquad V'=\Bigl(1-\frac\gamma4\Bigr)^{\!2}V\le V,
\]

%dsp-level:0
\[
\qquad \gamma=1:\ V'=\frac9{16}V,
\]

%dsp-level:0
\[
\qquad \gamma\to0:\ V'=V,
\]

%dsp-level:0
\[
\qquad \text{iterated to}\ \Bigl(\frac9{16}\Bigr)^{\!6}
\]

%dsp-level:0
\[
\qquad [\mathrm{L1};\ \mathrm{L2};\ \mathrm{BM.51}(1),(2);\ \text{machine}];
\]

%dsp-level:0
\[
(3)\quad\textbf{the statistics referent census:}\qquad\begin{array}{l|l}\text{a statistics requires} & \text{the corpus supplies}\\\hline\text{a multi-chain state space} & \neg\text{constructed:}\ \neg\text{a two-chain state space exists}\ [\mathrm{FT.80}(5)]\\\text{an exchange operator} & \neg\exists\ \text{one in the inventory}\ [\mathrm{FT.80}(5);\ \mathrm{FT.125}(5)]\\\text{a permutation sign} & \neg\text{formed:}\ \neg\exists\ \text{a permutation action on states}\\\text{an identification of a grade's}\ \mathbb Z_2\ \text{with the sign's}\ \mathbb Z_2 & \textbf{a kind confusion, barred}\ [\mathrm{FT.125}(5);\ \mathrm{FT.89}(2)\ \text{kind discipline}]\\\text{a tensor state on}\ \mathcal H_c & \neg\text{one: a counting space with no inner product, no dynamics, no state}\ [\mathrm{Def\text{-}H.41}\ \text{Species};\ \mathrm{Def\text{-}H.42}]\end{array}
\]

%dsp-level:0
\[
\boxed{\ \textbf{the spin--statistics connection is}\ \neg\textbf{formed}\ }\qquad[\mathrm{FT.125}(5)\ \text{re-cited}];
\]

%dsp-level:0
\[
\textbf{the derivation asked for is refuted on its face:}
\]

%dsp-level:0
\[
\qquad \text{the tick arrow and the spin label are both derived}\ (1),
\]

%dsp-level:0
\[
(2),
\]

%dsp-level:0
\[
\ \ \text{and neither is an exchange sign}\ —\ \text{the missing member is the exchange structure itself},
\]

%dsp-level:0
\[
\ \neg\text{a step from arrow to sign};
\]

%dsp-level:0
\[
(4)\quad\textbf{the four}\ \mathbb Z_2\ \textbf{gradings are single-object, and the census is exhaustive:}
\]

%dsp-level:0
\[
\qquad \text{the algebra}\ S\text{-grading}\ (10,6)\ [\mathrm{FT.136};\ \mathrm{FT.141}(1)];
\]

%dsp-level:0
\[
\ \ \text{the branch-label parity}\ [\mathrm{FT.101};\ \mathrm{FT.135};\ \mathrm{FT.157}(2),(3)];
\]

%dsp-level:0
\[
\text{the carrier}\ \hat\chi\text{-grade}\ [\mathrm{Def\text{-}H.33};\ \mathrm{FT.122}];
\]

%dsp-level:0
\[
\qquad \text{the Janus-arrow chirality}\ [\mathrm{FT.144}(2),(5)]:
\]

%dsp-level:0
\[
\qquad \textbf{pairwise distinct, commuting where jointly defined, neither reducible to another}
\]

%dsp-level:0
\[
\qquad [\mathrm{BM.73}(1);\ \mathrm{FT.141}(6);\ \mathrm{FT.126}(4)]:
\]

%dsp-level:0
\[
\boxed{\ \neg\text{one of the four is an exchange sign}\ }\qquad[\text{machine: the four symbol sets are single-object reads}];
\]

%dsp-level:0
\[
(5)\quad\textbf{where antisymmetry does live:}
\]

%dsp-level:0
\[
\qquad \textbf{the}\ \varepsilon\textbf{-seal is antisymmetric in the colour slots}
\]

%dsp-level:0
\[
\qquad B:=\varepsilon_{ijk}q^iq^jq^k\ [\mathrm{D3.9};\ \mathrm{CT3.15};\ \mathrm{CT3.17}]
\]

%dsp-level:0
\[
\qquad \textbf{and colour is}\ \neg\textbf{observable by construction}
\]

%dsp-level:0
\[
\qquad [\mathrm{CT3.22}]:
\]

%dsp-level:0
\[
\textbf{the transport assembly is sign-free:}
\]

%dsp-level:0
\[
\qquad W(Q)=\binom6{Q+3}>0\ \ \text{on all seven classes}
\]

%dsp-level:0
\[
\qquad [\mathrm{BM.103}(3);\ \text{machine}];
\]

%dsp-level:0
\[
\textbf{the amplitude layer's signs are barred as statistics:}
\]

%dsp-level:0
\[
\qquad (-1)^{\aleph}\ \ \text{with}\ \ \sum_i\aleph_{v_i}\equiv|q|\ (\mathrm{mod}\ 2)
\]

%dsp-level:0
\[
\qquad [\mathrm{Def\text{-}BM.13}(2);\ \mathrm{Def\text{-}BM.5}(2);\ \text{machine: exhaustive to}\ k=7];
\]

%dsp-level:0
\[
\lambda=(-1)^{|E|}e^{i\pi\Theta}\in\{\pm1,\pm i\},
\]

%dsp-level:0
\[
\qquad \Theta_{\mathrm{int}}\in4\mathbb Z\implies\lambda=(-1)^{|E|}
\]

%dsp-level:0
\[
\qquad [\mathrm{BM.25}(3),(6);\ \mathrm{FT.134}(3);\ \text{machine}]:
\]

%dsp-level:0
\[
\boxed{\ \textbf{amplitude}\ne\textbf{mass, and the wall parity is a charge-parity lock,}\ \neg\textbf{a permutation sign}\ }
\]

%dsp-level:0
\[
\qquad [\mathrm{FT.79}(5);\ \mathrm{Def\text{-}BM.13}(2)]
\]

%dsp-level:0
\[
\qquad \textbf{the corpus's only antisymmetry is the colour seal's, and it is unobservable};
\]

%dsp-level:0
\[
(6)\quad\textbf{the mixing row (the arena read):}
\]

%dsp-level:0
\[
\qquad \text{the mixing structure is a read of the derived}\ 16\text{-arena}\ \mathcal K_1\ \ \text{with}\ \ \delta_\nu=\tfrac2{16}=\tfrac18
\]

%dsp-level:0
\[
\qquad [\mathrm{FT.142}(2),(5);\ \mathrm{FT.143}];
\]

%dsp-level:0
\[
\text{the channel arenas are}\ 9,\ 27,\ 16\ \ \text{with the}\ 9/27\ \text{candidates rejected at rung one}
\]

%dsp-level:0
\[
\qquad [\mathrm{FT.142}(3);\ \mathrm{FT.112};\ \mathrm{FT.119}];
\]

%dsp-level:0
\[
\qquad \textbf{the cross-channel identity}\ \ 16\delta_\nu=9\delta_\ell=27\delta_q=2
\]

%dsp-level:0
\[
\qquad [\mathrm{FT.143}(5);\ \text{machine}];
\]

%dsp-level:0
\[
\boxed{\ \neg\text{an external mixing matrix is imported};\qquad\text{the mixing matrix is}\ \neg\text{derived};\qquad\text{every row holds uniformly in the unresolved line index}\ i_\nu\ }
\]

%dsp-level:0
\[
\qquad [\mathrm{FT.145}(3);\ \mathrm{BM.77}(2)];
\]

%dsp-level:0
\[
\textbf{the}\ \nu\ \textbf{row stands as derived:}
\]

%dsp-level:0
\[
\qquad \text{single-ray occupancy, no right-ray matter line, no seesaw-type mechanism}
\]

%dsp-level:0
\[
\qquad [\mathrm{BM.96}(4),(6)\ \text{re-cited},\ \neg\text{extended}];
\]

%dsp-level:0
\[
(7)\quad\textbf{the residue, named:}
\]

%dsp-level:0
\[
\qquad \boxed{\ (\mathrm{R7})\ \ \text{an exchange structure}\ :=\ \text{a two-chain state space}\ +\ \text{an exchange operator}\ +\ \text{a permutation-sign read}\ }
\]

%dsp-level:0
\[
\textbf{the sealed composite is a junction subnetwork,}\ \neg\textbf{a tensor product of chain state spaces:}
\]

%dsp-level:0
\[
\qquad \text{its constituents share vertices}\ [\mathrm{Def\text{-}BM.1};\ \mathrm{Def\text{-}BM.2}]\ \ \text{and no state space is declared on the counting spaces}\ [\mathrm{Def\text{-}H.41}\ \text{Species};\ \mathrm{FT.93}(4)\ \text{pattern}];
\]

%dsp-level:0
\[
(8)\quad\textbf{status:}
\]

%dsp-level:0
\[
\qquad \textbf{derived:}\ (1),\ (2),\ (4),\ (5),\ (6);
\]

%dsp-level:0
\[
\qquad \textbf{refuted as formed:}\ \text{the spin--statistics derivation from the tick arrow}\ (3);
\]

%dsp-level:0
\[
\qquad \textbf{standing:}\ \text{the spin--statistics connection, its multi-chain locus named};
\]

%dsp-level:0
\[
\neg\text{a declaration},
\]

%dsp-level:0
\[
\quad\neg\text{a P4 addition},
\]

%dsp-level:0
\[
\quad\neg\text{a new dimensionful scale}\ [\mathrm{P1}],
\]

%dsp-level:0
\[
\quad\neg\text{a measured value}\ [\mathrm{P5}],
\]

%dsp-level:0
\[
\quad\neg\text{a species identification}\ [\mathrm{Def\text{-}BM.3}(7)];
\]

\textbf{Req.} A6, A7, D0.8, D1.2, D1.4, D1.5, D1.8, D3.9 (Parts 1–2); A §0.2a, CT3.15, CT3.17, CT3.22 (Part 1, Part 2); P1, P4, P5 (Parts 1–2); L1, L2 (Part 1); Def-H.33, Def-H.38, Def-H.41, Def-H.42, FT.22, FT.24, FT.25(3), FT.60(1), FT.60(2), FT.60(3), FT.60(5), FT.80(5), FT.89(2), FT.93(4), FT.101, FT.112, FT.119, FT.122, FT.125(5), FT.126(4), FT.134(3), FT.135, FT.136, FT.141(1), FT.141(6), FT.142(2), FT.142(3), FT.142(5), FT.143, FT.143(5), FT.144(2), FT.144(4), FT.144(5), FT.145(3), FT.157(2), FT.157(3) (Part 3); Census FT.28(2), FT.49 (row-wise); Def-BM.1, Def-BM.2, Def-BM.5(2), Def-BM.13(2), BM.25(3), BM.25(6), BM.39(7), BM.51(1), BM.51(2), BM.73(1), BM.73(3), BM.77(2), BM.79(5), BM.96(4), BM.96(6), BM.103(3) (this part); exact integer winding arithmetic, exhaustive parity censuses, the Casimir ladder on exact rationals, symbol-set censuses [math; machine: \texttt{checks/\allowbreak{}verify\_\allowbreak{}H\_\allowbreak{}six\_\allowbreak{}objective\_\allowbreak{}resolution.py} §F].

\textbf{Pf.}

%dsp-level:0
\[
(1):\ \mathrm{Def\text{-}H.38}\text{'s multiplicity rule and}\ \mathrm{FT.60}(1)\text{--}(3),(5)\text{'s rows re-cited (frozen)};
\]

%dsp-level:0
\[
\qquad \Theta=2\kappa\in2\mathbb Z\implies k=\Theta/2=\kappa\in\mathbb Z\setminus\{0\}\ [\mathrm{Census\ FT.28}(2);\ \mathrm{FT.24};\ \text{machine}]
\]

%dsp-level:0
\[
\kappa=1\implies j=\tfrac12;
\]

%dsp-level:0
\[
\qquad \kappa=2\implies j=1;
\]

%dsp-level:0
\[
\qquad \text{the Casimir rows}\ \tfrac12\cdot\tfrac32=\tfrac34\ \ \text{and}\ \ 1\cdot2=2\ [\text{machine: exact rationals}];
\]

%dsp-level:0
\[
\qquad \text{both chains are materialized closed bound chains of one}\ \mathrm{D1.5}\ \text{kind}\ [\mathrm{Def\text{-}H.4};\ \mathrm{FT.22}]
\]

%dsp-level:0
\[
(2):\ \mathrm{D1.5},\ \mathrm{FT.80}(5),\ \mathrm{FT.144}(2),(4),(5),\ \mathrm{A6},\ \mathrm{A7},\ \mathrm{L2},\ \mathrm{BM.51}(1),(2)\ \text{re-cited verbatim};
\]

%dsp-level:0
\[
\qquad (9/16)^6=531441/16777216\ [\text{machine: exact integers}]
\]

%dsp-level:0
\[
(3):\ \text{the five referent rows re-cited from their own loci}\ [\mathrm{FT.80}(5);\ \mathrm{FT.125}(5);\ \mathrm{FT.89}(2);\ \mathrm{Def\text{-}H.41}\ \text{Species};\ \mathrm{Def\text{-}H.42}];
\]

%dsp-level:0
\[
\qquad \text{the inventory scans:}\ \neg\exists\ \text{an exchange token,}\ \neg\exists\ \text{a permutation-sign read}\ [\text{machine}]
\]

%dsp-level:0
\[
(4):\ \text{the four gradings and their pairwise distinctness re-cited}\ [\mathrm{BM.73}(1);\ \mathrm{FT.141}(6);\ \mathrm{FT.126}(4);\ \mathrm{FT.157}(4)];
\]

%dsp-level:0
\[
\qquad \text{each is a read of one object}\ [\text{machine: the symbol censuses}]
\]

%dsp-level:0
\[
(5):\ \varepsilon_{ijk}\ \text{antisymmetry}\ [\mathrm{CT3.15};\ \mathrm{CT3.17}];
\]

%dsp-level:0
\[
\qquad W(Q)=\binom6{Q+3}>0\ \text{on all seven classes}\ [\mathrm{BM.103}(3);\ \text{machine}];
\]

%dsp-level:0
\[
\qquad \sum_i\aleph_{v_i}
\]

%dsp-level:0
\[
\equiv |q|\ (\mathrm{mod}\ 2)\ \text{and}\
\]

%dsp-level:0
\[
\ge |q|\ \ \text{exhaustively for}\ k
\]

%dsp-level:0
\[
= 2,
\]

%dsp-level:0
\[
\dots,
\]

%dsp-level:0
\[
7\ [\mathrm{Def\text{-}BM.13}(2)\ \text{re-derived};\ \text{machine}]
\]

%dsp-level:0
\[
\lambda=(-1)^{|E|}e^{i\pi\Theta}\in\{\pm1,\pm i\};\qquad\Theta\in4\mathbb Z\implies e^{i\pi\Theta}=1\implies\lambda=(-1)^{|E|}\ [\mathrm{BM.25}(3),(6);\ \text{machine}]
\]

%dsp-level:0
\[
(6):\ \delta_\nu=\tfrac2{16}=\tfrac18\ \text{and the arena census re-cited}\ [\mathrm{FT.142}(2),(3),(5)];
\]

%dsp-level:0
\[
\qquad 16\cdot\tfrac18=9\cdot\tfrac29=27\cdot\tfrac2{27}=2\ [\text{machine: exact Fractions}];
\]

%dsp-level:0
\[
\qquad \neg\exists\ \text{a mixing-matrix tag}\ [\mathrm{FT.145}(3);\ \text{inventory scan}]
\]

%dsp-level:0
\[
(7)\text{--}(8):\ \text{the junction-subnetwork rows re-cited}\ [\mathrm{Def\text{-}BM.1};\ \mathrm{Def\text{-}BM.2};\ \mathrm{Def\text{-}H.41}\ \text{Species}];
\]

%dsp-level:0
\[
\qquad \text{the standing-rule scans}\ [\text{machine}].\ \blacksquare
\]

\textbf{Phys.} The corpus can name the spin and it can name the arrow, and it cannot join them, for a reason that is structural rather than temporary. The spin label comes from the winding: half the winding number, with the Casimir ladder and the weight multiplet behind it, and at the floor the winding is the registration integer of the two-cycle, so the same kind of object carries spin one half at one value of that integer and spin one at the next. There is therefore no parity theorem to be had: the floor itself supplies both parities on materialized closed bound chains, all of one definitional kind. The arrow is just as concrete — no inverse tick, a ledger that never shrinks, a prefix law that makes dissolution presuppose crystallization, a Lyapunov function that contracts by nine sixteenths per step at unit tilt — and the occupancy it forces is one ray full and the other empty. What is missing is not a step from arrow to sign but the object the sign would act on: there is no two-chain state space, no exchange operator, no permutation representation, and the corpus has already ruled that identifying a single-object grade with a permutation sign is a confusion of kinds. So the connection is not derived, and deficit eleven stands with its residue now named: an exchange structure, three ingredients, none of them in the inventory. What the tier adds is the census that makes the absence exact, the four gradings shown to be single-object and exhaustively so, and the location of the only antisymmetry the corpus does have — the colour seal's epsilon, which is unobservable by construction, against a transport assembly whose seven weights are all positive and an amplitude layer whose signs are barred as masses. The mixing rows stay where the arena puts them: sixteen dimensions, one eighth, and no imported matrix.

\end{theorem}

\medskip\noindent\textit{\texorpdfstring{BM.107 — The status of the six resolutions: adjudication and scope}{BM.107 - The status of the six resolutions: adjudication and scope}.}\ Structural-constraints record: see Appendix B, Section B.1.30 of the Addendum.\par\medskip

\phantomsection
\section[Tier BM.21 — The Antimatter and Resonance Tier: the…]{\texorpdfstring{Tier BM.21 — The Antimatter and Resonance Tier: the CP-Conjugate, the Meson Seal, the Exchange Operator, and the Local Resonance Spectrum}{Tier BM.21 - The Antimatter and Resonance Tier: the CP-Conjugate, the Meson Seal, the Exchange Operator, and the Local Resonance Spectrum}}

\textbf{Locus.} This tier takes the residue $(\mathrm{R7})$ that $\mathrm{BM.106}(7)$ named — \textit{an exchange structure: a two-chain state space, an exchange operator, a permutation-sign read} — and constructs all three members, then reads the spectrum the construction exposes. The CP-conjugate chain is derived from the three reversals the corpus already owns, with the forward-only tick blocking $\mathrm T$ as a dynamical inverse and leaving $\mathrm{CP}$ as the one composition that closes the ledger's own charge balance; the vacuum fluctuation is its pair, with the creation weight the Janus resolvent's own $x_+x_-$ and no pair at either endpoint of the threshold range; and the antiquark follows from the same three ingredients that force the compensating charge, so the corpus's Dirac parallel is an inventory statement and not a history one ($\mathrm{BM.108}$–$\mathrm{BM.110}$). The meson is then the $\delta$-sealed $2$-juncture, its state space is the first multi-chain space the corpus has, it discharges the $\mathrm{BM.6}(2)$ specification exactly as written, and its mass is derived on the two routes where the seal contributes no binding term and named unconstructed on the third ($\mathrm{BM.111}$–$\mathrm{BM.113}$). The exchange operator is the canonical flip composed with the CP transport, an involution with spectrum $\{\pm1\}$ unconditionally on the materialized locus, and its permutation-sign read is the fourth-root phase of the winding evaluated at the on-shell ladder — one integer's parity, which is the spin label's parity, which closes deficit $11$ ($\mathrm{BM.112}$–$\mathrm{BM.113}$). The resonance spectrum is then counted: two characteristic polynomials, one transport and one defect, the multiset-orbit law with its stabilizer multiplicities, exactly one exact collision in the whole three-chain ladder, a minimal-polynomial degree census that terminates at the field's own degree, a materialization threshold whose locus is never empty at any $n$ and whose span ratio is exactly thirty at every level, and a decay width that is a product of three derived factors and vanishes identically at the sharp threshold ($\mathrm{BM.114}$–$\mathrm{BM.117}$). Nothing is forbidden, nothing is stable as a category, and no cosmological clock is read: a resonance is an eigenvalue, its appearance is the evaluation of a threshold function at the local parameters, and its width is the evaluation of a width function at the same ones. $\mathrm{BM.118}$ adjudicates.

\textbf{Standing rows.} $\alpha_B$ stays barred on $\mathrm{BM.61}(8)$'s four grounds. No clause of $\mathrm{BM.1}$–$\mathrm{BM.107}$ is retracted, replaced or renumbered: the rows this tier touches are re-scoped at their own loci — $\mathrm{BM.106}(3)$'s statistics referent census and $(7)$'s residue $(\mathrm{R7})$, $\mathrm{BM.103}(4)$'s tensor-product referent census, $\mathrm{BM.6}(2)$'s exchange-space specification, $\mathrm{FT.125}(5)$'s and $\mathrm{FT.80}(5)$'s inventory rows, and $\mathrm{BM.32}(5)$'s amplitude/norm locus row — and each rescope is recorded in $\mathrm{BM.118}(5)$, in $\mathrm{BM.77}(7)$'s pattern. No measured coupling, mass, width, lifetime or ratio enters any derivation $[\mathrm{P5}]$: the threshold and width functions are derived symbolic forms and are $\neg$compared to any external number; the crossing strength of $\mathrm{BM.101}(3)$ is re-cited and $\neg$re-evaluated. No right-ray matter line, no seesaw-type mechanism, no Yukawa-type object and no external mixing matrix is suggested, added or required $[\mathrm{EW.15};\ \mathrm{FT.145}(3);\ \mathrm{BM.96}(4),(6)\ \text{re-cited},\ \neg\text{extended}]$. Species names are confined to the identification layer $[\mathrm{Def\text{-}BM.3}(7);\ \mathrm{Def\text{-}BM.6}(7);\ \mathrm{Def\text{-}BM.14}(8)]$. Topology labels are internal bookkeeping $[\mathrm{Def\text{-}H.32}(4)]$. Amplitude $\ne$ mass $[\mathrm{FT.79}(5)]$. Rung $\ne$ scale $[\mathrm{D0.8}]$. No currency is converted $[\mathrm{BM.83}(3);\ \mathrm{BM.102}(6)]$. No theorem of this tier imports a HistLog disposition as premise: Q10, Q11, item $20$, deficits $1'$, $3$, $4$, $11$ and $12$ are cited as pointers. One dimension only: units are powers of $\mathfrak q$ $[\mathrm{D2.9};\ \mathrm{CT2.1}]$.

\textbf{Local notation, fixed for the whole tier.} $c$ a materialized closed bound chain $[\mathrm{D1.4};\ \mathrm{D1.5};\ \mathrm{Def\text{-}H.4}]$; $\omega$ the orientation-reversal involution $[\mathrm{D2.6}]$; $\Theta_{\mathfrak L}$ the ledger reflection $[\mathrm{D2.7}]$; $\mathrm C,\mathrm P,\mathrm T$ the three derived involutions $[\mathrm{EM.24}(1)\text{--}(3)]$; $\mathrm{CP}:=\mathrm C\circ\mathrm P$; $\Theta(C)=\sum_{e\in E(C)}\widehat{ds}^{\,2}(e)$ and $\lambda=(-1)^{|E(C)|}e^{i\pi\Theta(C)}$ the CPT phase $[\mathrm{EM.24}(5);\ \mathrm{BM.25}(3)]$; $k(c)=\Theta(c)/2$ the on-shell winding and $j=|k|/2$ the spin label $[\mathrm{FT.22};\ \mathrm{FT.25}(3);\ \mathrm{Def\text{-}H.38};\ \mathrm{BM.74}]$; $\Psi(c)=\mathrm{Sym}^{|k(c)|}(\mathcal H_{\mathsf v})$ the winding carrier with $\mathcal H_{\mathsf v}\cong\mathbb C^2$ $[\mathrm{Def\text{-}H.38};\ \mathrm{EW.32}]$; $\mathcal H_\omega(c)=\operatorname{span}_{\mathbb C}\{|c\rangle,|\omega(c)\rangle\}\cong\mathbb C^2$ the orientation doublet with $\hat\omega,\hat\chi,P_{L/R}$ $[\mathrm{Def\text{-}H.33};\ \mathrm{FT.122}]$; $E_{\mathrm{int}}$ the internal directed-edge count, $m_{\mathrm{hid}}=E_{\mathrm{int}}\ln2/\mathfrak h$, $\Delta x_{\mathrm B}=\mathfrak h/(E_{\mathrm{int}}\ln2)$, $r_s=\mathfrak h E_{\mathrm{int}}/2$, $G=2\pi E_{\mathrm{int}}^2\ln2$ $[\mathrm{D0.8};\ \mathrm{FT.44}(4);\ \mathrm{Def\text{-}BM.4}(1);\ \mathrm{BM.13}(1);\ \mathrm{BM.15}(4);\ \mathrm{BM.16}(1);\ \mathrm{BM.88}(3)]$; $N$ the lapse and $f=1-1/k$ with $k=r/\Delta x_{\mathrm B}$ the envelope profile $[\mathrm{BM.34}(1);\ \mathrm{BM.88}(2)]$; $\varepsilon_{\mathrm{eff}}=Nf\,\mathsf g\,\varepsilon_{\mathrm{EM}}=e^2/(4\pi)$ $[\mathrm{BM.88}(4);\ \mathrm{BM.101}(2)]$; $\varepsilon_q\in[0,\tfrac14)$ the crystallization threshold, $x_+=1-4\varepsilon_q$, $x_-=4\varepsilon_q$, $p=x_+/(x_++x_-)=1-4\varepsilon_q$ $[\mathrm{A4};\ \mathrm{D0.4};\ \mathrm{CT0.8};\ \mathrm{BM.19}(2)]$; $\Gamma_{\mathrm{tr}}=x_+x_-/(1-x_+x_-)\in[0,\tfrac13]$ the unit-fugacity transition weight $[\mathrm{BM.44}(2)]$; $D_S(s)\in\{2,4,2-\sqrt2,1,2-2\cos\tfrac\pi{10},2-\sqrt3\}$ and $m^2_{\mathrm{pole}}=\hat\mu'^2D_S$ the six-state defect spectrum $[\mathrm{Def\text{-}BM.14}(3);\ \mathrm{BM.75}(4)]$; $\mathcal S$ the six-state set, $|\mathcal S|=6$ $[\mathrm{Def\text{-}BM.14}(2)]$; $W_n(Q)$ the $n$-chain transport class weight $[\mathrm{BM.28}(2);\ \mathrm{BM.103}(3)]$; $\mathfrak u:=2\cos\tfrac\pi{10}$ and $\mathbb M:=\mathbb Q(\sqrt2,\sqrt3,\mathfrak u)$ the exact eigenvalue field; $\chi_{\{\cdot\}}$ the exact indicator, valued in $\{0,1\}$; $[\cdot]$ the $\mathfrak q$-exponent map $[\mathrm{CT2.1};\ \mathrm{D2.9}]$.

\begin{definition}[{\texorpdfstring{Def-BM.15 — The CP-conjugate chain and the CP-pair}{Def-BM.15 - The CP-conjugate chain and the CP-pair}}]\label{tag:Def-BM-15}
\textbf{Def.} A construction on the three derived involutions of $\mathrm{EM.24}(1)\text{--}(3)$ and on $\mathrm{D2.6}$'s $\omega$; no new axiom, no new datum, no new parameter, no external object.

%dsp-level:0
\[
(1)\quad\textbf{the three reversals, re-cited as derived:}
\]

%dsp-level:0
\[
\qquad \mathrm C:\ \sigma_e\mapsto-\sigma_e\ \forall e\implies q\mapsto-q,\quad f_v\mapsto-f_v,\quad\aleph_v\ \text{invariant}
\]

%dsp-level:0
\[
\qquad [\mathrm{EM.24}(1);\ \mathrm{D3.3};\ \mathrm{D3.4}];
\]

%dsp-level:0
\[
\mathrm P:\ \mathcal P\mapsto E_{\mathcal R}\setminus\mathcal P,\qquad\text{equivalently}\quad\alpha_e\mapsto\pi-\alpha_e\qquad[\mathrm{EM.24}(2);\ \mathrm{Def\text{-}EM.21};\ \mathrm{EM.20}(1)];
\]

%dsp-level:0
\[
\mathrm T:\ \Theta_{\mathfrak L},\qquad w(\Theta_{\mathfrak L}C)=w(C)^\ast,\qquad J^+\leftrightarrow J^-,\qquad G(e)\mapsto G(e)^\ast\qquad[\mathrm{EM.24}(3);\ \mathrm{D2.7};\ \mathrm{CT2.8}(1)];
\]

%dsp-level:0
\[
\textbf{constitutive clause:}\ \ \text{all three are involutions on}\ \mathfrak C(\mathcal R)\ \text{and all three are derived}\ —\ \mathrm C^2
\]

%dsp-level:0
\[
= \mathrm P^2
\]

%dsp-level:0
\[
= \mathrm T^2\qquad = \mathrm{id},
\]

%dsp-level:0
\[
\qquad \text{no external input}
\]

%dsp-level:0
\[
\qquad [\mathrm{EM.24}(4)];
\]

%dsp-level:0
\[
(2)\quad\textbf{the CP-conjugate:}\qquad\boxed{\ \mathrm{CP}\ :=\ \mathrm C\circ\mathrm P,\qquad c^{\mathrm{CP}}\ :=\ \mathrm{CP}(c)\ }
\]

%dsp-level:0
\[
\textbf{constitutive clause:}\ \ \mathrm{CP}\ \text{is the composition of two derived involutions, hence an involution:}
\]

%dsp-level:0
\[
\quad (\mathrm{CP})^2=\mathrm C^2\mathrm P^2=\mathrm{id};
\]

%dsp-level:0
\[
\textbf{the action (derived, entrywise):}
\]

%dsp-level:0
\[
\qquad q\mapsto-q\ [\mathrm{EM.24}(1)];
\]

%dsp-level:0
\[
\qquad \text{scar}=q_{\mathrm{core}}\mapsto-q_{\mathrm{core}}\ [\mathrm{Def\text{-}BM.7}(3);\ \omega\text{-odd}];
\]

%dsp-level:0
\[
\text{chirality}\ \mapsto\ -\text{chirality}\ [\mathrm{D2.6};\ \mathrm{FT.122}(4):\ \hat\omega\hat\chi\hat\omega^{-1}=-\hat\chi];
\]

%dsp-level:0
\[
\qquad \text{triplet attachment}\ q^i\ \mapsto\ \text{dual attachment}\ r_i\ [\mathrm{D3.8};\ \mathrm{CT3.17}:\ \delta^i{}_j\ \text{consumes one of each type}];
\]

%dsp-level:0
\[
w\ \mapsto\ e^{i\pi\Theta}\,w
\]

%dsp-level:0
\[
\qquad [\mathrm{EM.24}(5)\ \text{with}\ (-1)^{|E|}\ \text{cancelled by}\ \mathrm P\text{'s}\ e^{i\pi\Theta}:\ \text{machine}];
\]

%dsp-level:0
\[
\qquad E_{\mathrm{int}}\ \mapsto\ E_{\mathrm{int}},\quad |R|\ \mapsto\ |R|
\]

%dsp-level:0
\[
\qquad [\mathrm{BM.25}(2);\ \mathrm{BM.27}(2)];
\]

%dsp-level:0
\[
(3)\quad\textbf{the CP-pair and the vacuum fluctuation:}
\]

%dsp-level:0
\[
\qquad \textbf{a CP-pair}\ :=\ \text{an orbit}\ \{c,\ c^{\mathrm{CP}}\}\ \text{of the involution of}\ (2);
\]

%dsp-level:0
\[
\textbf{a vacuum fluctuation}\ :
\]

%dsp-level:0
\[
= \ \text{the creation of one CP-pair from the vacuum sector}\ [\mathrm{BM.96}(4):\ \text{the complement's content is all-virtual and unregistered}];
\]

%dsp-level:0
\[
\textbf{constitutive clause:}\ \ \text{the pair's total charge read is}\ Q=q(c)+q(c^{\mathrm{CP}})=q-q=0\ \ \text{exactly}
\]

%dsp-level:0
\[
\qquad [\mathrm{EM.11};\ \mathrm{D2.11};\ \text{machine}];
\]

%dsp-level:0
\[
\textbf{non-derived clause:}\ \ \text{the word ``vacuum'' names the all-virtual sector of}\ \mathrm{CT4.2}\ \text{and}\ \mathrm{FT.144}(4);
\]

%dsp-level:0
\[
\ \ \text{it carries no energy datum and no state};
\]

%dsp-level:0
\[
(4)\quad\textbf{the scope clause (the mirror is formal, not dynamical):}
\]

%dsp-level:0
\[
\qquad \neg\exists\,\tau^{-1}\ [\mathrm{A6}];
\]

%dsp-level:0
\[
\qquad \bigl|\mathfrak L(\tau_{n+1})\bigr|=\bigl|\mathfrak L(\tau_n)\bigr|+n_{\mathrm{real\_out}}(\tau_n)\ge\bigl|\mathfrak L(\tau_n)\bigr|\ [\mathrm{A6};\ \mathrm{A7}];
\]

%dsp-level:0
\[
V'=\bigl(1-\tfrac\gamma4\bigr)^2V\le V,\qquad\gamma=1:\ V'=\tfrac9{16}V\qquad[\mathrm{L2};\ \mathrm{BM.51}(1),(2)]:
\]

%dsp-level:0
\[
\therefore\ \mathrm T\ \textbf{is blocked as a dynamical inverse:}\quad\Theta_{\mathfrak L}\ \text{is an involution on configuration data and}\ \neg\text{a step of the evolution}\ [\mathrm{Def\text{-}BM.7}(2);\ \mathrm{BM.25}(2);\ \mathrm{FT.144}(4)(\mathrm{ii})];
\]

%dsp-level:0
\[
\textbf{the consequence that makes}\ \mathrm{CP}\ \textbf{the operative conjugation:}
\]

%dsp-level:0
\[
\qquad \text{the ledger's balance is a condition on}\ \textbf{configurations},\ \neg\text{on evolutions}\ [\mathrm{D2.11};\ \mathrm{EM.11}],
\]

%dsp-level:0
\[
\text{so the compensating member a charged chain requires is a configuration, and}\ \mathrm{CP}\ \text{is the derived involution that produces one with the same}\ E_{\mathrm{int}}\ (2);
\]

%dsp-level:0
\[
(5)\quad\textbf{the data bar (standing):}
\]

%dsp-level:0
\[
\qquad \neg\exists\ \text{an annihilation channel, a pair-production cross-section or an inter-sector transition operator in the inventory}\ [\mathrm{Def\text{-}BM.7}(5);\ \mathrm{BM.27}(7);\ \mathrm{BM.29}(5)];
\]

%dsp-level:0
\[
\text{``antimatter'' is the label of}\ \mathrm{Def\text{-}BM.7}(1)\text{'s}\ \mathrm{EM}^-\ \text{sector and carries no external datum}\ [\mathrm{P5};\ \mathrm{Def\text{-}BM.3}(6)];
\]

%dsp-level:0
\[
\text{the CP-pair creation weight of}\ \mathrm{BM.110}(1)\ \text{is}\ \neg\text{a cross-section and}\ \neg\text{a rate in external units}\ [\mathrm{BM.44}(5);\ \mathrm{D1.8}].
\]

\textbf{Rmk (what is constitutive and what is derived).} Clause (2)'s \textit{definition} of $\mathrm{CP}$ as $\mathrm C\circ\mathrm P$ is the constitutive act — one composition of two objects the bosonic era already derived. Everything else in (2) is derived content recorded at the definition site: the entrywise action follows from $\mathrm{EM.24}(1),(2)$ and $\mathrm{D2.6}$, the weight rule from $\mathrm{EM.24}(5)$, and the invariance of $E_{\mathrm{int}}$ from $\mathrm{BM.25}(2)$'s record-size row. Clause (3)'s charge neutrality is derived, not stipulated: it is $\mathrm{EM.11}$'s balance applied to an orbit of an involution that negates the read. Nothing here names a particle, imports a mass, or supplies an annihilation amplitude.

\textbf{Req.} A6, A7, D0.8, D1.4, D1.5, D1.8, D2.6, D2.7, D2.11, D3.3, D3.4, D3.8 (Parts 1–2); CT0.8, CT2.8(1), CT3.17, CT4.2 (Part 1, Part 2); EM.11, EM.20(1), EM.24(1), EM.24(2), EM.24(3), EM.24(4), EM.24(5), Def-EM.21 (Part 2); P5 (Parts 1–2); L2 (Part 1); Def-H.4, Def-H.33, FT.122(4), FT.144(2), FT.144(4) (Part 3); Def-BM.4(1), Def-BM.7(1), Def-BM.7(2), Def-BM.7(3), Def-BM.7(5), BM.25(2), BM.27(2), BM.27(7), BM.29(5), BM.51(1), BM.51(2), BM.96(4) (this part); exact $\mathbb Q(i)$ weight algebra on the fourth-root alphabet, exhaustive charge-alphabet enumeration [math; machine: \texttt{checks/\allowbreak{}verify\_\allowbreak{}H\_\allowbreak{}antimatter\_\allowbreak{}resonance\_\allowbreak{}spectrum.py} §A].

\textbf{Pf.}

%dsp-level:0
\[
\text{Definition block: the constitutive clauses are}\ (1)\text{'s involution rows (re-cited, frozen)},
\]

%dsp-level:0
\[
\ (2)\text{'s composition}\ \mathrm{CP}:
\]

%dsp-level:0
\[
= \mathrm C\circ\mathrm P,
\]

%dsp-level:0
\[
\ \text{and}\ (3)\text{'s orbit naming};
\]

%dsp-level:0
\[
\ \text{the non-derived clauses are the ``vacuum'' vocabulary of}\ (3)\ \text{and the ``antimatter'' label of}\ (5)
\]

%dsp-level:0
\[
(2)\text{'s derived rows:}\ q\mapsto-q\ [\mathrm{EM.24}(1)];
\]

%dsp-level:0
\[
\qquad \alpha_e\mapsto\pi-\alpha_e\ [\mathrm{EM.24}(2)];
\]

%dsp-level:0
\[
\qquad \text{composed:}\ w\mapsto\bar w\mapsto e^{i\pi\Theta}w\ [\mathrm{EM.24}(5)\ \text{Pf row}\ \mathrm{P}\circ\mathrm T\ \text{transported to}\ \mathrm C\circ\mathrm P;\ \text{machine}]
\]

%dsp-level:0
\[
\text{attachment type:}\ \delta^i{}_j\ \text{consumes one triplet and one dual}\ [\mathrm{CT3.17}]\implies\text{the conjugate constituent carries the dual attachment}\ [\mathrm{D3.8};\ \mathrm{Def\text{-}BM.1}\ \text{meson clause}];
\]

%dsp-level:0
\[
\qquad E_{\mathrm{int}}=|R|\ \text{invariant}\ [\mathrm{BM.27}(2);\ \text{machine}]
\]

%dsp-level:0
\[
(3):\  Q=q+(-q)=0\ [\text{machine: exhaustive over}\ q\in\{-1,0,+1\}];
\]

%dsp-level:0
\[
\qquad (4):\  \mathrm{A6},\ \mathrm{A7},\ \mathrm{L2}\ \text{rows re-cited};
\]

%dsp-level:0
\[
\qquad (5):\  \text{the inventory scan of}\ \mathrm{BM.27}(7)\ \text{re-cited}.\ \blacksquare
\]

\textbf{Phys.} Three reversals exist in this corpus and each is derived: flip every edge sign, complement the aligned set, reflect the ledger. The third is the one that cannot run — the tick has no inverse and the Lyapunov function only contracts, so a reversed history is a configuration, never a step. Composing the first two gives the conjugation that does the work: it negates the charge read, flips the scar, exchanges the chiral rays, turns a triplet attachment into a dual one, multiplies the weight by a fourth root of unity, and leaves the edge count — and hence the informational mass, the width, the volume, the density and the comparison radius — untouched. An orbit of that involution is a pair, and a pair is neutral by construction, because the balance the valid-universe condition demands is a condition on configurations and the conjugate is exactly the configuration that satisfies it. What the definition does not do is give the pair a process: no annihilation channel, no production cross-section, no inter-sector operator exists in the inventory, and the word "vacuum" names the all-virtual sector rather than any state with an energy.

\end{definition}

\begin{theorem}[{\texorpdfstring{BM.108 — The CP-conjugate chain: the reversal census, the T-block, the exact phase, and the forced conjugate}{BM.108 - The CP-conjugate chain: the reversal census, the T-block, the exact phase, and the forced conjugate}}]\label{tag:BM-108}
\textbf{Thm.}

%dsp-level:0
\[
(1)\quad\textbf{the reversal census (what the corpus owns and what it blocks):}\qquad\begin{array}{l|l|l}\text{reversal} & \text{status} & \text{locus}\\\hline\mathrm C\ \text{charge conjugation} & \textbf{derived, dynamical} & [\mathrm{EM.24}(1);\ \mathrm{D3.3};\ \mathrm{D3.4}]\\\mathrm P\ \text{parity} & \textbf{derived, dynamical} & [\mathrm{EM.24}(2);\ \mathrm{Def\text{-}EM.21}]\\\mathrm T\ \text{time reversal} & \textbf{derived as a record involution, blocked as a dynamical inverse} & [\mathrm{EM.24}(3);\ \mathrm{A6};\ \mathrm{L2};\ \mathrm{FT.144}(4)(\mathrm{ii})]\\\mathrm{CP} & \textbf{derived, an involution} & [\mathrm{Def\text{-}BM.15}(2);\ \mathrm{EM.24}(4)]\\\mathrm{CPT} & \textbf{derived,}\ \mathrm{CPT}^2=\mathrm{id} & [\mathrm{EM.24}(4);\ \mathrm{BM.25}(1)]\end{array}
\]

%dsp-level:0
\[
\textbf{the T-block, in three exact rows:}\qquad(\mathrm i)\ \ \neg\exists\,\tau^{-1}\qquad[\mathrm{A6}];
\]

%dsp-level:0
\[
(\mathrm{ii})\quad\bigl|\mathfrak L(\tau_{n+1})\bigr|-\bigl|\mathfrak L(\tau_n)\bigr|=n_{\mathrm{real\_out}}(\tau_n)\ge0\qquad[\mathrm{A6};\ \mathrm{A7}];
\]

%dsp-level:0
\[
(\mathrm{iii})\quad V'-V=-\Bigl[1-\bigl(1-\tfrac\gamma4\bigr)^2\Bigr]V=-\frac{\gamma(8-\gamma)}{16}V\le0,\qquad\gamma=1:\ V'=\tfrac9{16}V\qquad[\mathrm{L2};\ \text{machine}]:
\]

%dsp-level:0
\[
\boxed{\ \textbf{the pure order reversal of every completed word is inadmissible at its first event}\ —\ \Theta_{\mathfrak L}\ \text{exists and}\ \neg\text{evolves}\ }
\]

%dsp-level:0
\[
\qquad [\mathrm{BM.25}(2);\ \mathrm{FT.144}(4)(\mathrm{ii});\ \text{machine}];
\]

%dsp-level:0
\[
(2)\quad\textbf{the CP action on every derived read of a chain (exact, entrywise):}\qquad\begin{array}{l|l|l}\text{read} & \text{CP image} & \text{locus}\\\hline q=q_{\mathrm{core}}=\tfrac12(\sigma_s-\sigma_t) & -q & [\mathrm{D3.4};\ \mathrm{EM.9};\ \text{machine}]\\\text{scar} & -\text{scar} & [\mathrm{Def\text{-}BM.7}(3);\ \mathrm{BM.27}(5);\ \text{machine}]\\\ \hat\chi\text{-grade} & -\hat\chi\text{-grade}\ (P_L\leftrightarrow P_R) & [\mathrm{D2.6};\ \mathrm{FT.122}(4);\ \mathrm{BM.73}(4);\ \text{machine over}\ \mathbb Q(i)]\\\text{attachment type}\ (q^i\ /\ r_i) & r_i\ /\ q^i & [\mathrm{D3.8};\ \mathrm{CT3.17}]\\\text{weight}\ w & e^{i\pi\Theta}\,w & [\mathrm{EM.24}(5);\ \mathrm{BM.25}(3);\ \text{machine}]\\|w|^2 & |w|^2 & [\mathrm{EM.24}(5):\ \text{unconditional};\ \text{machine}]\\ E_{\mathrm{int}},\ |R|,\ m_{\mathrm{hid}},\ \Delta x_{\mathrm B},\ V_{\mathrm{eff}},\ \rho,\ \kappa^2\rho,\ r_s & \textbf{fixed} & [\mathrm{BM.27}(2),(3);\ \mathrm{BM.28}(1):\ q\text{-exponent}\ 0\ \text{in all seven};\ \text{machine}]\\\Theta_{\mathrm{int}}\ \text{class} & \text{fixed} & [\mathrm{BM.28}(6);\ \text{machine}]\end{array}
\]

%dsp-level:0
\[
\boxed{\ \textbf{the CP-conjugate chain is mass-degenerate with the chain on every route}\ }
\]

%dsp-level:0
\[
\qquad [\mathrm{BM.28}(5);\ \mathrm{BM.27}(3)\ \mathrm{Rmk};\ \text{machine: the monomial}\ q\text{-exponent census}];
\]

%dsp-level:0
\[
(3)\quad\textbf{the exact CP phase:}\qquad w(c^{\mathrm{CP}})=\mu_{\mathrm{CP}}(c)\,w(c),\qquad\boxed{\ \mu_{\mathrm{CP}}(c)\ :=\ e^{i\pi\Theta(c)}\ }
\]

%dsp-level:0
\[
\textbf{derivation:}\quad w(\mathrm{CPT}\,C)=(-1)^{|E|}e^{i\pi\Theta}w(C)\ [\mathrm{EM.24}(5)],
\]

%dsp-level:0
\[
\qquad w(\mathrm T\,C)=w(C)^\ast\ [\mathrm{CT2.8}(1)],
\]

%dsp-level:0
\[
\qquad \mathrm{CP}=(\mathrm{CPT})\circ\mathrm T^{-1}=(\mathrm{CPT})\circ\mathrm T:
\]

%dsp-level:0
\[
\mu_{\mathrm{CP}}=(-1)^{|E|}e^{i\pi\Theta}\cdot\frac{w}{w^\ast}\bigg|_{\text{the}\ \mathrm T\ \text{leg}}=(-1)^{|E|}e^{i\pi\Theta}\cdot(-1)^{|E|}=e^{i\pi\Theta}
\]

%dsp-level:0
\[
\qquad [\text{machine:}\ |w|=1\implies w/w^\ast=w^2/|w|^2;\ \text{the}\ (-1)^{|E|}\ \text{cancels exactly}];
\]

%dsp-level:0
\[
\mu_{\mathrm{CP}}\in\{\pm1,\pm i\}\ \ \text{on the half-integer ladder},
\]

%dsp-level:0
\[
\qquad \mu_{\mathrm{CP}}\in\{+1\}\ \ \text{on the materialized locus}\ \Theta\in2\mathbb Z
\]

%dsp-level:0
\[
\qquad [\mathrm{FT.22};\ \mathrm{BM.2}(1);\ \text{machine}]:
\]

%dsp-level:0
\[
\boxed{\ \textbf{the CP character is exactly}\ 1\ \textbf{on the sealed-composite junction class}\ }
\]

%dsp-level:0
\[
\qquad [\Theta_{\mathrm{int}}\in4\mathbb Z\ [\text{Census FT.49}]\implies e^{i\pi\Theta}=1;\ \mathrm{BM.25}(6);\ \text{machine}];
\]

%dsp-level:0
\[
(4)\quad\textbf{the orbit census:}
\]

%dsp-level:0
\[
\qquad \mathrm{CP}:\ (q,\ \text{orientation})\mapsto(-q,\ \text{flipped})\ \ \text{on}\ \{-1,0,+1\}\times\{w,w^\ast\}:
\]

%dsp-level:0
\[
\qquad \text{no fixed sector, three orbits of size two}\ [\mathrm{BM.25}(4);\ \text{machine}];
\]

%dsp-level:0
\[
q\mapsto-q\ \ \text{on the ternary alphabet has exactly one fixed point,}\ q=0
\]

%dsp-level:0
\[
\qquad [\mathrm{EM.24}(1);\ \text{machine}]:
\]

%dsp-level:0
\[
\qquad \boxed{\ \textbf{the neutral class is charge-self-conjugate; every charged class is exchanged}\ }
\]

%dsp-level:0
\[
\text{the CP involution on charge}\ n\text{-tuples has exactly one fixed tuple,}\ (0,\dots,0),
\]

%dsp-level:0
\[
\ \text{and}\ W_n(+Q)
\]

%dsp-level:0
\[
= W_n(-Q)\ \text{exactly}
\]

%dsp-level:0
\[
\qquad [\text{machine: exhaustive at}\ n=2,3];
\]

%dsp-level:0
\[
(5)\quad\textbf{the forcing theorem (why the conjugate is not optional):}
\]

%dsp-level:0
\[
\qquad Q_{\mathrm{EM}}=\sum_{c\in U}q_{\mathrm{EM}}(c)=0\ \ \textbf{in every valid universe}
\]

%dsp-level:0
\[
\qquad [\mathrm{EM.11};\ \mathrm{D2.11};\ \mathrm{D1.3}]:
\]

%dsp-level:0
\[
\text{a materialized chain with}\ q(c)\ne0\ \text{requires a compensating read of exactly}\ -q(c);
\]

%dsp-level:0
\[
\qquad \#\{q\in\{-1,0,+1\}:\ q=-q(c)\}=1
\]

%dsp-level:0
\[
\qquad [\mathrm{BM.23}(4);\ \text{machine}];
\]

%dsp-level:0
\[
\textbf{the compensator census:}\qquad\begin{array}{l|l}\text{a compensating configuration must carry} & \text{the corpus supplies}\\\hline\text{the charge read}\ -q(c) & \textbf{exactly one derived map produces it:}\ \mathrm C\ [\mathrm{EM.24}(1)]\\\text{a dual colour attachment, to seal against}\ c\text{'s triplet} & \textbf{exactly one derived map produces it:}\ \mathrm{D3.8}\circ\mathrm{CP}\ [\mathrm{CT3.17}]\\\text{the opposite chiral ray} & \textbf{exactly one derived map produces it:}\ \hat\omega\ [\mathrm{FT.122}(4);\ \mathrm{A\text{-}BM.7}(b)]\\\text{a closed real-edge network with}\ \#J^+=\#J^- & \mathrm{CP}\ \text{preserves}\ |R|\ \text{and exchanges}\ \#J^+\ \text{with}\ \#J^-\ [\mathrm{BM.25}(2);\ \mathrm{BM.27}(2)]\\\text{the same}\ E_{\mathrm{int}}\ (\text{hence the same mass read}) & \textbf{derived:}\ (2)\text{'s seven}\ q\text{-exponent-zero monomials}\ [\mathrm{BM.28}(1)]\end{array}
\]

%dsp-level:0
\[
\boxed{\ \textbf{the CP-conjugate is the unique derived configuration that balances a charged chain: it is forced out of the ledger's own balance,}\ \neg\textbf{added}\ }
\]

%dsp-level:0
\[
\qquad [\text{machine: the five-row census}];
\]

%dsp-level:0
\[
(6)\quad\textbf{the Dirac parallel, as an inventory statement:}
\]

%dsp-level:0
\[
\qquad \textbf{the three ingredients:}\ (\mathrm i)\ \text{the forward-only tick}\ [\mathrm{A6};\ (1)(\mathrm i)];
\]

%dsp-level:0
\[
(\mathrm{ii})\ \text{the}\ 16\text{-matrix pair arena}\ \mathcal K_1,
\]

%dsp-level:0
\[
\ \dim\mathcal K_1
\]

%dsp-level:0
\[
= 16
\]

%dsp-level:0
\[
= C_{\mathrm{cap}}\ [\mathrm{FT.142}(2);\ \mathrm{CT0.14};\ \mathrm{CT0.12}];
\]

%dsp-level:0
\[
(\mathrm{iii})\ \text{the total-charge-zero balance}\ [\mathrm{EM.11};\ \mathrm{D2.11}]:
\]

%dsp-level:0
\[
(\mathrm i)+(\mathrm{iii})\implies\text{the compensator is a configuration and}\ \neg\text{an inverse evolution}\ [(4)\text{'s scope clause}];
\]

%dsp-level:0
\[
\qquad (\mathrm{ii})\implies\text{the compensator lives in the same}\ 16\text{-arena, on the}\ \omega\text{-image representation}\ \{0,-1\}\ [\mathrm{EW.36};\ \mathrm{A\text{-}BM.7}(b)];
\]

%dsp-level:0
\[
\boxed{\ \textbf{adding the forward-only tick to the}\ 16\textbf{-matrix pair arena forces the conjugate constituent, exactly as a negative-energy branch forces its conjugate}\ }
\]

%dsp-level:0
\[
\qquad [\text{the parallel is structural:}\ \neg\text{an external history is imported},\ \neg\text{a measured mass is compared}\ [\mathrm{P5}]];
\]

%dsp-level:0
\[
\textbf{the scope of the parallel:}
\]

%dsp-level:0
\[
\qquad \neg\text{a Dirac operator,}\ \neg\text{a Lorentz spinor and}\ \neg\text{a negative-energy spectrum is constructed}\ [\mathrm{BM.6}(8):\ \text{the frame section and a global Dirac-type operator remain unconstructed};\ \mathrm{FT.93}(3)\ \text{the pincer}];
\]

%dsp-level:0
\[
(7)\quad\textbf{the no-annihilation row:}
\]

%dsp-level:0
\[
\qquad \neg\exists\ \text{an annihilation channel, a pair-production cross-section or an inter-sector transition operator}\ [\mathrm{Def\text{-}BM.7}(5);\ \mathrm{BM.27}(7);\ \mathrm{BM.29}(5)]:
\]

%dsp-level:0
\[
\textbf{the two sectors are amplitude-equal configuration orbits; their interaction is}\ \neg\textbf{a theorem of this corpus}
\]

%dsp-level:0
\[
\qquad [\mathrm{BM.27}(7)\ \text{re-cited},\ \neg\text{extended}];
\]

%dsp-level:0
\[
\text{the CP-pair creation weight of}\ \mathrm{BM.110}(1)\ \text{is a flicker weight,}\ \neg\text{a cross-section}\ [\mathrm{BM.44}(5);\ \mathrm{D1.8}];
\]

%dsp-level:0
\[
(8)\quad\textbf{status:}
\]

%dsp-level:0
\[
\qquad \textbf{derived:}\ (1)\text{--}(5);
\]

%dsp-level:0
\[
\qquad \textbf{derived as a structural parallel with an explicit scope:}\ (6);
\]

%dsp-level:0
\[
\qquad \textbf{recorded as absent:}\ (7);
\]

%dsp-level:0
\[
\neg\text{a declaration},
\]

%dsp-level:0
\[
\quad\neg\text{a P4 addition},
\]

%dsp-level:0
\[
\quad\neg\text{a new dimensionful scale}\ [\mathrm{P1}],
\]

%dsp-level:0
\[
\quad\neg\text{a measured value}\ [\mathrm{P5}],
\]

%dsp-level:0
\[
\quad\neg\text{a species identification}\ [\mathrm{Def\text{-}BM.3}(7)],
\]

%dsp-level:0
\[
\quad\neg\text{a numerical coupling}\ [\mathrm{EM.25}(7)];
\]

%dsp-level:0
\[
\qquad\text{§Cal-BM stays empty except for}\ \mathrm{A\text{-}BM.9}\text{'s logged slot}.
\]

\textbf{Rmk (locus discipline on clause (6)).} The parallel is a statement about what forces what inside one inventory: a balance condition plus a forward-only order plus a fixed arena leave exactly one derived configuration that satisfies all three. It is not a claim that the corpus reproduces any external equation, and no clause identifies the conjugate constituent with any measured object — the name \textit{antiquark} belongs to the identification layer and occurs in no \textbf{Thm} clause here [Def-BM.3(7)].

\textbf{Req.} A6, A7, D0.8, D1.3, D1.4, D1.8, D2.6, D2.7, D2.11, D3.3, D3.4, D3.8 (Parts 1–2); CT0.12, CT0.14, CT2.8(1), CT3.17, EM.9, EM.11, EM.24(1), EM.24(2), EM.24(3), EM.24(4), EM.24(5), EM.25(7), EW.36, Def-EM.21 (Part 1, Part 2); P1, P4, P5 (Parts 1–2); L2 (Part 1); Def-H.33, FT.22, FT.93(3), FT.122(4), FT.142(2), FT.144(2), FT.144(4) (Part 3); Census FT.49 (row-wise); Def-BM.3(7), Def-BM.7(1), Def-BM.7(2), Def-BM.7(3), Def-BM.7(5), Def-BM.15, BM.2(1), BM.6(8), BM.23(4), BM.25(1), BM.25(2), BM.25(3), BM.25(4), BM.25(6), BM.27(2), BM.27(3), BM.27(5), BM.27(7), BM.28(1), BM.28(5), BM.28(6), BM.29(5), BM.51(1), BM.51(2), BM.73(4), A-BM.7 (this part); exact $\mathbb Q(i)$ phase algebra on the fourth-root alphabet, exhaustive ternary-alphabet and $n$-tuple censuses, Laurent-monomial $q$-exponent censuses [math; machine: \texttt{checks/\allowbreak{}verify\_\allowbreak{}H\_\allowbreak{}antimatter\_\allowbreak{}resonance\_\allowbreak{}spectrum.py} §A].

\textbf{Pf.}

%dsp-level:0
\[
(1):\  \mathrm{EM.24}(1)\text{--}(4)\ \text{rows re-cited (frozen, machine-verified at Tier F.33)};
\]

%dsp-level:0
\[
\qquad (\mathrm i),(\mathrm{ii}):\  \mathrm{A6},\ \mathrm{A7}\ \text{verbatim}
\]

%dsp-level:0
\[
(\mathrm{iii}):\  V'-V=-[1-(1-\tfrac\gamma4)^2]V=-\tfrac{\gamma(8-\gamma)}{16}V\ [\mathrm{L2}];
\]

%dsp-level:0
\[
\qquad \gamma=1:\ 1-\tfrac14=\tfrac34,\ (\tfrac34)^2=\tfrac9{16},\ 1-\tfrac9{16}=\tfrac7{16},\ \gamma(8-\gamma)=7\ [\text{machine: exact rationals}]
\]

%dsp-level:0
\[
\text{the inadmissibility row:}\ \mathrm{FT.144}(4)(\mathrm{ii})\ \text{and}\ \mathrm{BM.25}(2)\ \text{re-cited}\ [\text{machine: the pure reversal leads with}\ J^-,\ \text{outside D1.2's domain}]
\]

%dsp-level:0
\[
(2):\  q\mapsto\tfrac12(-\sigma_s+\sigma_t)=-q\ [\mathrm{D3.4}];
\]

%dsp-level:0
\[
\qquad \hat\omega\hat\chi\hat\omega^{-1}=-\hat\chi,\ \hat\omega P_L\hat\omega^{-1}=P_R\ [\mathrm{FT.122}(4);\ \mathrm{BM.73}(4);\ \text{machine over}\ \mathbb Q(i)]
\]

%dsp-level:0
\[
\text{the seven fixed reads:}\ \Delta x_{\mathrm B},
\]

%dsp-level:0
\[
\ V_{\mathrm{eff}},
\]

%dsp-level:0
\[
\ \rho,
\]

%dsp-level:0
\[
\ \kappa^2\rho,
\]

%dsp-level:0
\[
\ r_s,
\]

%dsp-level:0
\[
\ (\Delta x_{\mathrm B})^2\kappa^2\rho,
\]

%dsp-level:0
\[
\ \Delta x_{\mathrm B}r_s\ \text{each carry charge exponent}\ 0\ [\mathrm{BM.28}(1);\ \text{machine: monomial census}]
\]

%dsp-level:0
\[
\implies q
\]

%dsp-level:0
\[
\mapsto -q\ \text{fixes each}
\]

%dsp-level:0
\[
(3):\  \mathrm{EM.24}(5)\text{'s Pf row}\ \bar\Gamma\mapsto e^{i\pi\Theta}\Gamma\ \text{under}\ \mathrm P\circ\mathrm T,\ \text{and}\ S\mapsto(-1)^{|E|}S\ \text{under}\ \mathrm C;
\]

%dsp-level:0
\[
\qquad \mathrm{CP}=\mathrm{CPT}\circ\mathrm T\ \text{at}\ \mathrm T^2=\mathrm{id}\ [\mathrm{EM.24}(4)]
\]

%dsp-level:0
\[
w(\mathrm{CP}\,C)=w(\mathrm{CPT}\,\mathrm T\,C)=(-1)^{|E|}e^{i\pi\Theta}w(\mathrm T\,C)=(-1)^{|E|}e^{i\pi\Theta}w(C)^\ast;
\]

%dsp-level:0
\[
\qquad |w(C)|=1\ [\mathrm{Def\text{-}H.32}(2):\ |w(R,\sigma)|=1]\implies w(C)^\ast=(-1)^{|E|}w(C)^{-1}\cdot|w|^2
\]

%dsp-level:0
\[
\text{on the sign-real locus}\ w\in\mathbb R\ [\mathrm{EM.24}(5)\ \text{on-shell}]
\]

%dsp-level:0
\[
\implies w^\ast
\]

%dsp-level:0
\[
= w\implies\mu_{\mathrm{CP}}
\]

%dsp-level:0
\[
= (-1)^{|E|}e^{i\pi\Theta}(-1)^{|E|}
\]

%dsp-level:0
\[
= e^{i\pi\Theta}\ [\text{machine:}\ (-1)^{2|E|}=1];
\]

%dsp-level:0
\[
\qquad \Theta\in2\mathbb Z\implies e^{i\pi\Theta}=i^{2\Theta}=1\ [\text{machine}]
\]

%dsp-level:0
\[
(4):\  \mathrm{BM.25}(4)\text{'s six-sector census re-cited}\ [\text{machine}];
\]

%dsp-level:0
\[
\qquad q\mapsto-q\ \text{fixed points on}\ \{-1,0,+1\}:\ \{0\}\ [\text{machine}];
\]

%dsp-level:0
\[
\qquad n\text{-tuples: exhaustive at}\ n
\]

%dsp-level:0
\[
= 2\ (9\ \text{tuples},\ 5\ \text{orbits},\ 1\ \text{fixed})\ \text{and}\ n
\]

%dsp-level:0
\[
= 3\ (27\ \text{tuples},\ 14\ \text{orbits},\ 1\ \text{fixed})\ [\text{machine}]
\]

%dsp-level:0
\[
W_n(+Q)
\]

%dsp-level:0
\[
= W_n(-Q):\ \text{the map}\ (q_1,\dots,q_n)
\]

%dsp-level:0
\[
\mapsto (-q_1,\dots,-q_n)\ \text{is a weight-preserving bijection between the two class fibres}\ [m_{-q}=m_q\ [\mathrm{FT.106}(3\mathrm a):\ m_q=(2,1,1)];\ \text{machine}]
\]

%dsp-level:0
\[
(5):\  \mathrm{EM.11},\ \mathrm{D2.11}\ \text{re-cited};
\]

%dsp-level:0
\[
\qquad \#\{q:\ q=-\Delta q\}=1\ [\mathrm{BM.23}(4)\ \text{pattern};\ \text{machine}];
\]

%dsp-level:0
\[
\qquad \text{the five-row census assembled from its own tags}\ [\text{machine}]
\]

%dsp-level:0
\[
(6):\  \dim\mathcal K_1=16=C_{\mathrm{cap}}\ [\mathrm{FT.142}(2);\ \mathrm{CT0.14}];
\]

%dsp-level:0
\[
\qquad \omega\text{-image representation}\ \{0,-1\}\ [\mathrm{EW.36};\ \mathrm{A\text{-}BM.7}(b)];
\]

%dsp-level:0
\[
\qquad \text{the unconstructed rows re-cited}\ [\mathrm{BM.6}(8);\ \mathrm{FT.93}(3)]
\]

%dsp-level:0
\[
(7):\  \text{the inventory scan of}\ \mathrm{BM.27}(7)\ \text{re-cited};\qquad(8):\  \text{the standing-rule scans}\ [\text{machine}].\ \blacksquare
\]

\textbf{Phys.} Of the three reversals the corpus derives, two run and one does not. Charge conjugation and parity are maps the ledger can apply to a configuration; time reversal is a map it can apply to a record but never to a step, because the tick has no inverse, the ledger only grows, and the Lyapunov function only contracts — at sharp crystallization by exactly seven sixteenths per tilt. Composing the two that run gives the conjugation that matters, and its action is exact and cheap: the charge read flips, the scar flips, the chiral rays exchange, the triplet attachment becomes a dual one, the weight picks up a fourth root of unity that is exactly one on the sealed junction's own phase class, and the seven geometric reads do not move at all, because not one of them carries a charge exponent. So the conjugate is mass-degenerate with the original by structure rather than by luck. That it must exist is not an addition but a consequence: total charge vanishes in every valid universe, a charged chain therefore needs a compensator, the alphabet contains exactly one value that compensates it, and of the derived maps exactly one produces a configuration carrying that value with a dual attachment, the opposite ray, and the same edge count. The orbit census is the same statement at the level of classes — one self-conjugate neutral class, every charged class exchanged, one fixed charge tuple at every level and equal weights on the two charged classes. The structural parallel is then exact in its own terms: a balance, a forward-only order, and a fixed sixteen-dimensional arena leave one derived compensator, which is the shape of the argument that forces a conjugate branch, and the tier says no more than that — no Dirac operator, no spinor, no annihilation channel, and no name outside the identification layer.

\end{theorem}

\begin{theorem}[{\texorpdfstring{BM.109 — The CP-pair as the vacuum fluctuation: the creation weight, its exact range, and the degeneracy rows}{BM.109 - The CP-pair as the vacuum fluctuation: the creation weight, its exact range, and the degeneracy rows}}]\label{tag:BM-109}
\textbf{Thm.}

%dsp-level:0
\[
(1)\quad\textbf{the pair's charge and mass rows:}
\]

%dsp-level:0
\[
\qquad Q_{\text{pair}}=q(c)+q(c^{\mathrm{CP}})=q-q=0
\]

%dsp-level:0
\[
\qquad [\mathrm{Def\text{-}BM.15}(3);\ \text{machine: exhaustive over}\ q\in\{-1,0,+1\}]:
\]

%dsp-level:0
\[
E_{\mathrm{int}}^{\text{pair}}=2\,E_{\mathrm{int}}(c)
\]

%dsp-level:0
\[
\qquad [\text{additivity of the internal count on a disjoint pair};\ \mathrm{D0.8};\ \mathrm{BM.21}(1)\ \text{pattern}];
\]

%dsp-level:0
\[
\boxed{\ m_{\mathrm{hid}}^{\text{pair}}=\frac{2E_{\mathrm{int}}(c)\ln2}{\mathfrak h}=2\,m_{\mathrm{hid}}(c),\qquad\Delta x_{\mathrm B}^{\text{pair}}=\frac{\mathfrak h}{2E_{\mathrm{int}}(c)\ln2}=\frac{\Delta x_{\mathrm B}(c)}2\ }
\]

%dsp-level:0
\[
\qquad [\mathrm{FT.44}(4);\ \mathrm{Def\text{-}BM.4}(1);\ \text{machine: Laurent identities}];
\]

%dsp-level:0
\[
\textbf{the binding read of the pair count vanishes exactly:}
\]

%dsp-level:0
\[
\qquad E^{\mathrm{info}}_{\mathrm{bind}}(\text{pair})=\bigl(2E_{\mathrm{int}}-(E_{\mathrm{int}}+E_{\mathrm{int}})\bigr)\tfrac{\ln2}{\mathfrak h}=0
\]

%dsp-level:0
\[
\qquad [\mathrm{BM.21}(1);\ \mathrm{BM.54}(1);\ \text{machine}];
\]

%dsp-level:0
\[
(2)\quad\textbf{the creation weight is the Janus resolvent's own:}
\]

%dsp-level:0
\[
\qquad \sum_n\langle\varnothing|\hat U^n|\varnothing\rangle w^n=\frac1{1-w^2x_+x_-}\ \ \text{in}\ \mathbb Q[w,x_+,x_-]/(w^{12})
\]

%dsp-level:0
\[
\qquad [\mathrm{FT.108}(5);\ \mathrm{BM.23}(7);\ \mathrm{BM.44}(1)]:
\]

%dsp-level:0
\[
\boxed{\ \mathcal W_{\text{pair}}\ :=\ x_+x_-\ =\ 4\varepsilon_q(1-4\varepsilon_q)\ =\ \frac14-16\Bigl(\varepsilon_q-\frac18\Bigr)^{\!2}\ \le\ \frac14,\qquad\text{equality}\iff\varepsilon_q=\frac18\ }
\]

%dsp-level:0
\[
\qquad [\mathrm{BM.19}(2);\ \mathrm{BM.44}(1);\ \text{machine: the SOS completion}\ +\ \text{the grid census over}\ \varepsilon_q=k/256];
\]

%dsp-level:0
\[
\textbf{the two endpoints:}\qquad\varepsilon_q=0\implies x_-=0\implies\mathcal W_{\text{pair}}=0\ \ \text{exactly}\qquad[\mathrm{BM.19}(3);\ \text{machine}];
\]

%dsp-level:0
\[
\varepsilon_q\to\tfrac14^-\implies x_+\to0\implies\mathcal W_{\text{pair}}\to0\qquad[\mathrm{CT0.8}:\ \varepsilon_q\in[0,\tfrac14);\ \text{machine}]:
\]

%dsp-level:0
\[
\boxed{\ \mathcal W_{\text{pair}}>0\iff\varepsilon_q\in\bigl(0,\tfrac14\bigr),\qquad\max\mathcal W_{\text{pair}}=\frac14\ \ \text{at}\ \ \varepsilon_q=\frac18\ }
\]

%dsp-level:0
\[
\qquad [\text{machine: exact rationals on a}\ 256\text{-point grid}];
\]

%dsp-level:0
\[
(3)\quad\textbf{the one-door row:}
\]

%dsp-level:0
\[
\qquad \text{the vacuum sector is all-virtual and unregistered with no accessible pole at any rung}\ [\mathrm{BM.96}(4);\ \mathrm{FT.144}(4);\ \mathrm{CT4.2}];
\]

%dsp-level:0
\[
\qquad \text{the materialized locus requires}\ |E_{\mathcal R}|\ge2\ [\mathrm{Def\text{-}H.4};\ \mathrm{BM.13}(3);\ \mathrm{BM.70}(2)]:
\]

%dsp-level:0
\[
\boxed{\ \textbf{the CP-pair materializes and propagates forward-only exactly where the flicker is present}\ —\ \text{at the same threshold the ensemble integral is taken over}\ }
\]

%dsp-level:0
\[
\qquad [\mathrm{BM.96}(6)\ \text{re-cited at the pair level};\ \text{machine}];
\]

%dsp-level:0
\[
\Gamma_{\mathrm{tr}}(\varepsilon_q{=}0)=0
\]

%dsp-level:0
\[
\qquad [\mathrm{BM.23}(8);\ \mathrm{BM.44}(2)]:
\]

%dsp-level:0
\[
\qquad \textbf{a non-flickering ledger creates no pairs and decays nothing};
\]

%dsp-level:0
\[
(4)\quad\textbf{the degeneracy rows (the pair is one orbit, not two structures):}
\]

%dsp-level:0
\[
\qquad |w(c^{\mathrm{CP}})|^2=|w(c)|^2\ \ \forall c
\]

%dsp-level:0
\[
\qquad [\mathrm{EM.24}(5):\ \text{unconditional};\ \text{machine}];
\]

%dsp-level:0
\[
m_{\mathrm{hid}}(c^{\mathrm{CP}})=m_{\mathrm{hid}}(c),
\]

%dsp-level:0
\[
\qquad D_S(c^{\mathrm{CP}})=D_S(c),
\]

%dsp-level:0
\[
\qquad m^2_{\mathrm{pole}}(c^{\mathrm{CP}})=m^2_{\mathrm{pole}}(c)
\]

%dsp-level:0
\[
\qquad [\mathrm{BM.108}(2);\ \mathrm{Def\text{-}BM.14}(3):\ D_S\ \text{is branch data,}\ \neg\text{charge data};\ \text{machine}];
\]

%dsp-level:0
\[
(\mathrm{CP})^2=\mathrm{id}\implies\textbf{the pair carries no independent datum: it is one orbit of size two}
\]

%dsp-level:0
\[
\qquad [\mathrm{BM.27}(6);\ \mathrm{EM.24}(4)];
\]

%dsp-level:0
\[
\text{the two rows contribute equally to the norm}\ \langle A|A\rangle=\sum_C|w(C)|^2
\]

%dsp-level:0
\[
\qquad [\mathrm{CT2.8}(3);\ \mathrm{BM.25}(7);\ \text{machine}];
\]

%dsp-level:0
\[
(5)\quad\textbf{the no-annihilation row stands:}
\]

%dsp-level:0
\[
\qquad \neg\exists\ \text{an annihilation channel or a cross-section between}\ \mathrm{EM}^+\ \text{and}\ \mathrm{EM}^-\ [\mathrm{BM.27}(7);\ \mathrm{Def\text{-}BM.7}(5)]:
\]

%dsp-level:0
\[
\textbf{what is derived is the pair's existence, its neutrality, its degeneracy and its creation weight}\ —\ \neg\text{its interaction};
\]

%dsp-level:0
\[
(6)\quad\textbf{status:}\qquad\textbf{derived:}\ (1)\text{--}(4);\qquad\textbf{recorded as absent:}\ (5);
\]

%dsp-level:0
\[
\neg\text{a declaration},
\]

%dsp-level:0
\[
\quad\neg\text{a P4 addition},
\]

%dsp-level:0
\[
\quad\neg\text{a new dimensionful scale}\ [\mathrm{P1}],
\]

%dsp-level:0
\[
\quad\neg\text{a measured value}\ [\mathrm{P5}],
\]

%dsp-level:0
\[
\quad\neg\text{a numerical coupling}\ [\mathrm{EM.25}(7)],
\]

%dsp-level:0
\[
\quad\neg\text{a lifetime in external units}\ [\mathrm{BM.44}(5)];
\]

\textbf{Req.} A4, A6, D0.4, D0.8, D1.8, D2.9 (Parts 1–2); CT0.8, CT2.8(3), CT4.2, EM.24(4), EM.24(5), EM.25(7) (Part 1, Part 2); P1, P4, P5 (Parts 1–2); Def-H.4, Def-H.32(2), FT.108(5), FT.144(4), FT.44(4) (Part 3, Census); Def-BM.3(7), Def-BM.4(1), Def-BM.7(1), Def-BM.7(5), Def-BM.14(3), Def-BM.15, BM.13(3), BM.19(2), BM.19(3), BM.21(1), BM.23(7), BM.23(8), BM.25(7), BM.27(6), BM.27(7), BM.44(1), BM.44(2), BM.44(5), BM.54(1), BM.70(2), BM.96(4), BM.96(6), BM.108(2) (this part); exact rational resolvent identities, the SOS parabola completion, exhaustive rational grid censuses, Laurent-monomial assembly [math; machine: \texttt{checks/\allowbreak{}verify\_\allowbreak{}H\_\allowbreak{}antimatter\_\allowbreak{}resonance\_\allowbreak{}spectrum.py} §A].

\textbf{Pf.}

%dsp-level:0
\[
(1):\  q+(-q)=0\ [\text{machine: three cases}];
\]

%dsp-level:0
\[
\qquad E_{\mathrm{int}}^{\text{pair}}=E_{\mathrm{int}}+E_{\mathrm{int}}\ \text{on a disjoint union of directed-edge sets}\ [\mathrm{D0.8}];
\]

%dsp-level:0
\[
\qquad m_{\mathrm{hid}}=E_{\mathrm{int}}\ln2/\mathfrak h\ \text{linear}\ [\mathrm{FT.44}(4);\ \text{machine}]
\]

%dsp-level:0
\[
\Delta x_{\mathrm B}^{\text{pair}}=\mathfrak h/(2E_{\mathrm{int}}\ln2)=\Delta x_{\mathrm B}/2\ [\text{machine: one Laurent monomial}];
\]

%dsp-level:0
\[
\qquad E^{\mathrm{info}}_{\mathrm{bind}}=0\cdot\tfrac{\ln2}{\mathfrak h}\ [\mathrm{BM.21}(1)\ \text{pattern};\ \text{machine}]
\]

%dsp-level:0
\[
(2):\  \mathrm{FT.108}(5)\text{'s resolvent re-cited};
\]

%dsp-level:0
\[
\qquad x_+=1-4\varepsilon_q,\ x_-=4\varepsilon_q\implies x_+x_-=4\varepsilon_q(1-4\varepsilon_q)\ [\mathrm{BM.19}(2);\ \text{machine}]
\]

%dsp-level:0
\[
4\varepsilon_q(1-4\varepsilon_q)=\tfrac14-16(\varepsilon_q-\tfrac18)^2\ [\text{machine: expansion}];
\]

%dsp-level:0
\[
\qquad \text{the grid census}\ \varepsilon_q=k/256,\ k=0..63:\ \le\tfrac14,\ =\tfrac14\ \text{only at}\ k=32\ [\text{machine}]
\]

%dsp-level:0
\[
\varepsilon_q=0\implies x_-=0\implies x_+x_-=0;
\]

%dsp-level:0
\[
\qquad \varepsilon_q\to\tfrac14^-
\]

%dsp-level:0
\[
\implies x_+\to0
\]

%dsp-level:0
\[
\implies x_+x_-\to0\ [\text{machine: exact at the two endpoints of the admissible range}\ [\mathrm{CT0.8}]]
\]

%dsp-level:0
\[
(3):\  \mathrm{BM.96}(4),
\]

%dsp-level:0
\[
(6)\text{'s rows transported to the pair}\ [\text{the pair's members are the two states of that row}];
\]

%dsp-level:0
\[
\qquad \Gamma_{\mathrm{tr}}(0)=0\ [\mathrm{BM.44}(2);\ \text{machine}]
\]

%dsp-level:0
\[
(4):\  |w(\mathrm{CP}\,C)|^2=|w(C)|^2\ \text{from}\ |\mu_{\mathrm{CP}}|=|e^{i\pi\Theta}|=1\ [\mathrm{BM.108}(3);\ \text{machine}];
\]

%dsp-level:0
\[
\qquad D_S\ \text{is a function of}\ (\mathcal F_g,\beta)\ \text{alone}\ [\mathrm{Def\text{-}BM.14}(3);\ \text{machine: no charge symbol occurs}]
\]

%dsp-level:0
\[
(\mathrm{CP})^2=\mathrm C^2\mathrm P^2=\mathrm{id}\ [\mathrm{EM.24}(4)];
\]

%dsp-level:0
\[
\qquad (5),(6):\  \text{the inventory scan and the standing-rule scans}\ [\text{machine}].\ \blacksquare
\]

\textbf{Phys.} A pair made of a chain and its conjugate is neutral because the conjugate's charge is the negative of the chain's, and it is mass-degenerate because none of the seven geometric reads carries a charge exponent and the defect read is branch data rather than charge data. Its informational mass is exactly twice a member's and its width exactly half, with the count-route binding identically zero — a binding energy cannot be an integer summation, and the pair is the cleanest case of it. The weight for making one is not declared: it is the Janus resolvent's own $x_+x_-$, a parabola in the single crystallization threshold with maximum one quarter at one eighth and zeros at both ends of the admissible range. So the ledger creates pairs exactly where it flickers, and nowhere else — at the sharp threshold the dissolution weight vanishes, no pair is made, and nothing decays. What the tier does not give the pair is a fate: no annihilation channel and no inter-sector cross-section exist in the inventory, and the pair stays one orbit of size two carrying no datum its member does not.

\end{theorem}

\begin{definition}[{\texorpdfstring{Def-BM.16 — The two-chain state space (the construction BM.6(2) specifies)}{Def-BM.16 - The two-chain state space (the construction BM.6(2) specifies)}}]\label{tag:Def-BM-16}
\textbf{Def.} A construction discharging the $\mathrm{BM.6}(2)$ specification — \textit{a state space on ordered pairs of bound chains with an exchange involution} — on $\mathrm{Def\text{-}H.33}$'s orientation space and $\mathrm{Def\text{-}H.38}$'s winding carrier; no new axiom, no new datum, no new parameter, no external state.

%dsp-level:0
\[
(1)\quad\textbf{the single-chain factors (re-cited, both derived):}
\]

%dsp-level:0
\[
\qquad \mathcal H_\omega(c)=\operatorname{span}_{\mathbb C}\{|c\rangle,\ |\omega(c)\rangle\}\cong\mathbb C^2,
\]

%dsp-level:0
\[
\qquad \hat\omega^2=\mathrm{id},\quad\hat\omega^\dagger=\hat\omega,\quad\hat\chi=i\bigl(|c\rangle\langle\omega(c)|-|\omega(c)\rangle\langle c|\bigr),\quad P_{L/R}=\tfrac12(\mathrm{id}\pm\hat\chi)
\]

%dsp-level:0
\[
[\mathrm{Def\text{-}H.33};\ \mathrm{FT.122}(1)\text{--}(3),(6)];
\]

%dsp-level:0
\[
\qquad \Psi(c)=\mathrm{Sym}^{|k(c)|}(\mathcal H_{\mathsf v}),
\]

%dsp-level:0
\[
\qquad \dim\Psi(c)=|k(c)|+1,
\]

%dsp-level:0
\[
\qquad \mathcal H_{\mathsf v}\cong\mathbb C^2
\]

%dsp-level:0
\[
\qquad [\mathrm{Def\text{-}H.38};\ \mathrm{EW.32};\ \text{machine}];
\]

%dsp-level:0
\[
(2)\quad\textbf{the two-chain state space:}\qquad\boxed{\ \mathcal H_2(c_1,c_2)\ :=\ \mathcal H_\omega(c_1)\otimes\Psi(c_1)\ \otimes\ \mathcal H_\omega(c_2)\otimes\Psi(c_2)\ }
\]

%dsp-level:0
\[
\dim\mathcal H_2=4\,\bigl(|k(c_1)|+1\bigr)\bigl(|k(c_2)|+1\bigr)\qquad[\text{machine: exact integers}];
\]

%dsp-level:0
\[
\textbf{the meson specialization (the CP-pair):}
\]

%dsp-level:0
\[
\qquad c_2=c_1^{\mathrm{CP}}\implies|k(c_2)|=|k(c_1)|=:|k|:
\]

%dsp-level:0
\[
\qquad \boxed{\ \mathcal H_{\text{meson}}(c)\ :=\ \mathcal H_2\bigl(c,\ c^{\mathrm{CP}}\bigr),\qquad\dim\mathcal H_{\text{meson}}=4\,(|k|+1)^2\ }
\]

%dsp-level:0
\[
\textbf{constitutive clause:}\ \ \text{the space is the tensor product of two derived single-chain spaces, with the}\ \mathrm{CP}\text{-identification}\ |c^{\mathrm{CP}}\rangle\leftrightarrow|\omega(c)\rangle\ \text{of}\ \mathrm{D2.6};
\]

%dsp-level:0
\[
\textbf{non-derived clause:}\ \ \text{the factor ordering}\ (1,2)\ \text{is a convention;}\ \text{the identification of the second factor with the first's}\ \omega\text{-image is the}\ \mathrm{D2.6}\ \text{vocabulary clause, applied};
\]

%dsp-level:0
\[
(3)\quad\textbf{the inner product and the norm layer:}
\]

%dsp-level:0
\[
\qquad \langle A|A\rangle=\sum_C|w(C)|^2\ge0
\]

%dsp-level:0
\[
\qquad [\mathrm{CT2.8}(3)];
\]

%dsp-level:0
\[
\qquad |w(R,\sigma)|=1\ \ \text{for every valid pair}
\]

%dsp-level:0
\[
\qquad [\mathrm{Def\text{-}H.32}(2);\ \mathrm{FT.68}\ \mathrm{Rmk}(6\mathrm c)]:
\]

%dsp-level:0
\[
\text{the orientation doublet carries}\ \mathrm{Def\text{-}H.33}\text{'s orthonormal basis}\ \{|c\rangle,|\omega(c)\rangle\};
\]

%dsp-level:0
\[
\qquad \text{the winding carrier carries the}\ \mathrm{Sym}^{|k|}\text{-inherited product of}\ \mathcal H_{\mathsf v}\ [\mathrm{EW.32}];
\]

%dsp-level:0
\[
\textbf{derived:}\ \ \mathcal H_2\ \textbf{is a genuine state space}\ —\ \text{an inner product, a norm layer and a finite exact dimension}\ —\ \neg\text{a counting space}
\]

%dsp-level:0
\[
\qquad [\mathrm{Def\text{-}H.41}\ \text{Species};\ \mathrm{Def\text{-}H.42}:\ \text{the bar this clause discharges}];
\]

%dsp-level:0
\[
(4)\quad\textbf{the exchange involution (the specification's second member):}
\]

%dsp-level:0
\[
\qquad \hat S_{12}\ :=\ \text{the canonical flip on}\ \mathcal H_2,
\]

%dsp-level:0
\[
\qquad \hat S_{12}(v\otimes w)=w\otimes v,
\]

%dsp-level:0
\[
\qquad \hat S_{12}^2=\mathrm{id}
\]

%dsp-level:0
\[
\qquad [\mathrm{A3}:\ S^2=\mathbb1;\ \text{machine}];
\]

%dsp-level:0
\[
\hat\mu_{\mathrm{CP}}\ :=\ \text{the CP-transport phase acting by the scalar}\ \mu_{\mathrm{CP}}(c)=e^{i\pi\Theta(c)}\ \ \text{on the}\ \Psi\text{-factors}
\]

%dsp-level:0
\[
\qquad [\mathrm{BM.108}(3);\ \text{machine}];
\]

%dsp-level:0
\[
\boxed{\ \hat P_{12}\ :=\ \hat\mu_{\mathrm{CP}}\circ\hat S_{12}\ }
\]

%dsp-level:0
\[
\qquad \textbf{constitutive clause:}\ \ \text{the exchange is the flip composed with the transport its own identification requires};
\]

%dsp-level:0
\[
(5)\quad\textbf{the permutation-sign read (the specification's third member):}
\]

%dsp-level:0
\[
\qquad \boxed{\ \varepsilon(c)\ :=\ \mu_{\mathrm{CP}}(c)\big|_{\vee^2}\ =\ i^{\,\Theta(c)}\ =\ (-1)^{|k(c)|}\ }
\]

%dsp-level:0
\[
\qquad [\mathrm{BM.112}(2);\ \text{machine}]:
\]

%dsp-level:0
\[
\textbf{constitutive clause:}\ \ \text{the statistics sign of a chain is the exchange operator's eigenvalue on the}\ \vee^2\text{-sector of the CP-pair's winding carrier};
\]

%dsp-level:0
\[
\textbf{non-derived clause:}\ \ \text{the choice of the}\ \vee^2\text{-sector as the read's locus}\ —\ \text{priced by}\ \mathrm{BM.112}(5)\text{'s four rows and}\ \neg\text{forced by any clause of this definition};
\]

%dsp-level:0
\[
(6)\quad\textbf{the specification-discharge clause:}
\]

%dsp-level:0
\[
\qquad \mathrm{BM.6}(2)\ \text{requires (i) a state space on ordered pairs of bound chains and (ii) an exchange involution:}
\]

%dsp-level:0
\[
(\mathrm i)\ \ \text{discharged by}\ (2),(3);\qquad(\mathrm{ii})\ \ \text{discharged by}\ (4)\ \text{with}\ \hat P_{12}^2=\mathrm{id}\ \text{derived}\ [\mathrm{BM.112}(1)];
\]

%dsp-level:0
\[
\textbf{the bar}\ \mathrm{FT.125}(5)\ \textbf{records is respected:}
\]

%dsp-level:0
\[
\qquad \varepsilon(c)\ \text{is}\ \neg\text{one of the four single-object}\ \mathbb Z_2\ \text{gradings of}\ \mathrm{BM.106}(4)\ \text{— it is a permutation sign on a two-chain space}\ [\mathrm{BM.113}(2)];
\]

%dsp-level:0
\[
(7)\quad\textbf{status and scope:}
\]

%dsp-level:0
\[
\qquad \textbf{derived:}\ (1),\ (3),\ (4)\text{'s involution rows},\ (5)\text{'s value};
\]

%dsp-level:0
\[
\qquad \textbf{constitutive:}\ (2)\text{'s product},\ (4)\text{'s composition},\ (5)\text{'s sector choice};
\]

%dsp-level:0
\[
\textbf{the scope inherited at every use:}\ \ \mathrm{Def\text{-}H.38}\text{'s nonnegative cone}\ k
\]

%dsp-level:0
\[
\ge 0\ \text{and its}\ k
\]

%dsp-level:0
\[
\ge 3\ \text{aspect bar}\ [\mathrm{Def\text{-}H.38}\ \text{Scope};\ \text{no aspect claim is made here}];
\]

%dsp-level:0
\[
\qquad \mathrm{Def\text{-}H.32}\text{'s}\ (\mathrm{S1})\text{--}(\mathrm{S4});
\]

%dsp-level:0
\[
\neg\text{a P4 addition},
\]

%dsp-level:0
\[
\quad\neg\text{a new dimensionful scale}\ [\mathrm{P1}],
\]

%dsp-level:0
\[
\quad\neg\text{a measured value}\ [\mathrm{P5}],
\]

%dsp-level:0
\[
\quad\neg\text{a species identification}\ [\mathrm{Def\text{-}BM.3}(7)];
\]

%dsp-level:0
\[
\qquad \text{§Cal-BM stays empty};
\]

%dsp-level:0
\[
\textbf{beyond this tier:}\ \ \text{the three-chain state space}\ \mathcal H_3\ \text{and its}\ S_3\ \text{action}\ \text{are}\ \neg\text{constructed}\ [\mathrm{BM.103}(4)\ \text{re-scoped},\ \neg\text{discharged}].
\]

\textbf{Rmk (why this is a construction and not a declaration).} Every ingredient is derived and frozen: the orientation doublet and its involution from the fermion era's interface layer, the winding carrier and its dimension from the multiplicity–winding coupling, the inner product from the norm layer, the flip from the pair axiom's own $S^2=\mathbb1$, and the transport phase from the CP character. What the definition adds is one product and one composition, and the corpus's own specification — BM.6(2), written as a build-spec and adopted nowhere — is what it discharges. The single non-derived row is a sector choice, and it is priced rather than hidden: BM.112(5) lists the four alternative loci and BM.113(2) shows the read is invariant under all four on the materialized locus.

\textbf{Req.} A3 (Part 1); D2.6, D2.9 (Parts 1–2); CT2.8(3) (Part 2); EW.32 (Part 2); P1, P4, P5 (Parts 1–2); Def-H.32(2), Def-H.33, Def-H.38, Def-H.41, Def-H.42, FT.68, FT.122(1), FT.122(2), FT.122(3), FT.122(6), FT.125(5) (Part 3); Def-BM.3(7), Def-BM.15, BM.6(2), BM.103(4), BM.106(4), BM.108(3) (this part); exact integer dimension arithmetic, exhaustive sector censuses [math; machine: \texttt{checks/\allowbreak{}verify\_\allowbreak{}H\_\allowbreak{}antimatter\_\allowbreak{}resonance\_\allowbreak{}spectrum.py} §B].

\textbf{Pf.}

%dsp-level:0
\[
\text{Definition block: the constitutive clauses are}\ (2)\text{'s tensor product},
\]

%dsp-level:0
\[
\ (4)\text{'s composition and}\ (5)\text{'s sector choice};
\]

%dsp-level:0
\[
\ \text{the non-derived clauses are the factor ordering of}\ (2)\ \text{and the sector locus of}\ (5);
\]

%dsp-level:0
\[
\ (1),
\]

%dsp-level:0
\[
(3),
\]

%dsp-level:0
\[
(4)\text{'s involution rows and}\ (5)\text{'s value are derived and machine-checked};
\]

%dsp-level:0
\[
\ (6),
\]

%dsp-level:0
\[
(7)\ \text{are scope records}.\ \blacksquare
\]

\textbf{Phys.} The corpus said for three tiers that a statistics has no referent, and listed exactly what one would need: a space of two chains, an operator that swaps them, and a sign to read off. This definition supplies all three and nothing more. The space is the tensor product of the two single-chain spaces the fermion era already built — the orientation doublet and the winding carrier — so it has an inner product from the norm layer, a finite exact dimension of four times the two carrier dimensions, and it is a state space rather than a counting space, which is the bar the species definitions set. The swap is the canonical flip, whose square is one because the pair axiom already said so, composed with the one phase the CP-identification forces: transporting a chain to its conjugate's slot multiplies by the fourth root of unity the CPT bridge derived. The sign is that phase evaluated on the symmetric sector, and on the materialized locus it is a genuine plus or minus. What remains open is the obvious next step and the definition says so: three chains, and the permutation group that acts on them, are not built here.

\end{definition}

\begin{theorem}[{\texorpdfstring{BM.110 — The meson seal: the $\delta$-sealed $2$-juncture, its uniqueness, and the charge class it occupies}{BM.110 - The meson seal: the -sealed -juncture, its uniqueness, and the charge class it occupies}}]\label{tag:BM-110}
\textbf{Thm.}

%dsp-level:0
\[
(1)\quad\textbf{the seal is unique at two attachments:}
\]

%dsp-level:0
\[
\qquad \text{partitions of a labelled}\ 2\text{-element attachment set into blocks of size}\ \in\{2,3\}:\quad\text{exactly one, the}\ \delta\text{-pair}
\]

%dsp-level:0
\[
\qquad [\mathrm{CT3.17a}(1),(2):\ g(2)=1;\ \text{machine}]:
\]

%dsp-level:0
\[
\delta^i{}_j\ \text{consumes one triplet and one dual}
\]

%dsp-level:0
\[
\qquad [\mathrm{CT3.17}];
\]

%dsp-level:0
\[
\qquad \varepsilon_{ijk}\ \text{requires}\ 3\ \text{fundamentals}\ [\mathrm{CT3.15};\ \mathrm{D3.9}]:
\]

%dsp-level:0
\[
\boxed{\ \textbf{meson}\ :=\ \delta^i{}_j\,c^j\,\bar c_i\ =\ \delta^i{}_j\,c^j\,(c_i)^{\mathrm{CP}}\ \ \text{is the unique sealed}\ 2\text{-juncture, and a lone attachment is never sealed}\ }
\]

%dsp-level:0
\[
\qquad [\mathrm{Def\text{-}BM.1};\ \mathrm{CT3.17a}(1):\ |\mathrm{Colour}(\Sigma)|=1\implies\mathcal C=0;\ \text{machine}];
\]

%dsp-level:0
\[
(2)\quad\textbf{the two-chain space is the first multi-chain state space of the corpus:}
\]

%dsp-level:0
\[
\qquad \mathcal H_{\text{meson}}(c)=\mathcal H_2\bigl(c,c^{\mathrm{CP}}\bigr),
\]

%dsp-level:0
\[
\qquad \dim=4(|k|+1)^2
\]

%dsp-level:0
\[
\qquad [\mathrm{Def\text{-}BM.16}(2);\ \text{machine}]:
\]

%dsp-level:0
\[
\textbf{the referent census}\ \mathrm{BM.106}(3)\ \textbf{listed on five rows is discharged on its first three:}\qquad\begin{array}{l|l}\text{a statistics requires} & \text{the corpus supplies}\\\hline\text{a multi-chain state space} & \textbf{constructed:}\ \mathcal H_2\ [\mathrm{Def\text{-}BM.16}(2),(3)]\\\text{an exchange operator} & \textbf{constructed:}\ \hat P_{12}\ [\mathrm{Def\text{-}BM.16}(4);\ \mathrm{BM.112}(1)]\\\text{a permutation sign} & \textbf{formed:}\ \varepsilon(c)\ [\mathrm{Def\text{-}BM.16}(5);\ \mathrm{BM.112}(2)]\\\text{a grade's}\ \mathbb Z_2\ \text{identified with the sign's} & \textbf{still barred, and not done:}\ \varepsilon\ \text{is a permutation sign on}\ \mathcal H_2\ [\mathrm{BM.113}(2);\ \mathrm{FT.125}(5);\ \mathrm{FT.89}(2)]\\\text{a tensor state on}\ \mathcal H_c & \textbf{unchanged:}\ \mathcal H_c\ \text{stays a counting space;}\ \mathcal H_2\ \text{is built on}\ \mathcal H_\omega\otimes\Psi,\ \neg\text{on}\ \mathcal H_c\ [\mathrm{Def\text{-}BM.16}(3);\ \mathrm{Def\text{-}H.41}\ \text{Species}]\end{array}
\]

%dsp-level:0
\[
\boxed{\ \textbf{residue}\ (\mathrm{R7})\ \textbf{is discharged at}\ n=2\ }
\]

%dsp-level:0
\[
\qquad [\mathrm{BM.106}(7);\ \mathrm{BM.107}(7);\ \text{machine: the three-row discharge against}\ \mathrm{BM.6}(2)\text{'s specification}];
\]

%dsp-level:0
\[
(3)\quad\textbf{the charge class the seal occupies:}
\]

%dsp-level:0
\[
\qquad \textbf{the CP-pair subset:}\quad Q=q+( -q)=0\ \ \text{for all three}\ q\in\{-1,0,+1\}
\]

%dsp-level:0
\[
\qquad [\mathrm{Def\text{-}BM.15}(3);\ \text{machine}]:
\]

%dsp-level:0
\[
\textbf{the full}\ 2\text{-chain transport census:}
\]

%dsp-level:0
\[
\qquad \bigl(2+z+z^{-1}\bigr)^2=z^{-2}(z+1)^4\implies W_2(Q)=\binom4{Q+2}=(1,\,4,\,\mathbf6,\,4,\,1)\ \ \text{for}\ Q=-2,\dots,+2,
\]

%dsp-level:0
\[
\qquad \sum_QW_2(Q)=2^4=16=4^2
\]

%dsp-level:0
\[
[\mathrm{BM.103}(3)\ \text{transported to}\ n=2;\ \text{machine: the Laurent identity and the exhaustive}\ 4^2\ \text{census}];
\]

%dsp-level:0
\[
\textbf{the CP-pair subset of the neutral class:}
\]

%dsp-level:0
\[
\qquad \{(-1,+1),\ (0,0),\ (+1,-1)\}\ \ \text{with weights}\ m_{-1}m_{+1}=1,\ m_0^2=4,\ m_{+1}m_{-1}=1:
\]

%dsp-level:0
\[
\qquad \textbf{total}\ 1+4+1=6=W_2(0)
\]

%dsp-level:0
\[
[\mathrm{FT.106}(3\mathrm a):\ m_q=(2,1,1);\ \text{machine: exact integers}];
\]

%dsp-level:0
\[
\qquad \#\{Q=0\ \text{tuples outside the CP-pair forms}\}=0
\]

%dsp-level:0
\[
\qquad [\text{machine: exhaustive}];
\]

%dsp-level:0
\[
\boxed{\ \textbf{the CP-pairs saturate the}\ 2\textbf{-chain neutral class exactly:}\quad W_2^{\mathrm{CP}}(0)=W_2(0)=6,\qquad W_2^{\mathrm{CP}}(\pm1)=W_2^{\mathrm{CP}}(\pm2)=0\ }
\]

%dsp-level:0
\[
\textbf{the charged meson classes are}\ \neg\textbf{CP-pairs:}
\]

%dsp-level:0
\[
\qquad Q=\pm1\ \text{requires}\ (q_1,q_2)\in\{(0,\pm1),(\pm1,0)\},
\]

%dsp-level:0
\[
\qquad Q=\pm2\ \text{requires}\ (\pm1,\pm1):
\]

%dsp-level:0
\[
\text{a}\ \delta\text{-sealed charged member is a pair of}\ \neg\textbf{conjugate constituents, and}\ \mathrm{Def\text{-}BM.3}(3)\text{'s three classes are all formed}
\]

%dsp-level:0
\[
\qquad [\text{machine: the four charged tuples against the CP-pair forms}];
\]

%dsp-level:0
\[
(4)\quad\textbf{the seal's invariance laws:}
\]

%dsp-level:0
\[
\qquad \delta^i{}_j\ \text{is the pair invariant:}\quad q^ir_i\mapsto U^i{}_aq^a\,r_b(U^{-1})^b{}_i=q^ar_a
\]

%dsp-level:0
\[
\qquad [\mathrm{CT3.17};\ \text{machine}]:
\]

%dsp-level:0
\[
\text{colour is}\ \neg\textbf{observable by construction}
\]

%dsp-level:0
\[
\qquad [\mathrm{CT3.22}];
\]

%dsp-level:0
\[
\qquad \text{the mask admits only singlets}
\]

%dsp-level:0
\[
\qquad [\mathrm{CT3.18};\ \mathrm{D2.14};\ \mathrm{FT.19}:\ \text{all-or-nothing}];
\]

%dsp-level:0
\[
\Theta_{\mathrm{int}}\in2\mathbb Z\ \ \text{on the}\ \delta\text{-sealed class against}\ 4\mathbb Z\ \ \text{on the}\ \varepsilon\text{-sealed one}
\]

%dsp-level:0
\[
\qquad [\mathrm{BM.2}(1);\ \mathrm{BM.31}(4);\ \text{Census FT.35}(3),\ \text{FT.49}(3):\ 4\mathbb Z\subsetneq2\mathbb Z\ \text{with}\ 2\in2\mathbb Z\setminus4\mathbb Z\ [\text{machine}]];
\]

%dsp-level:0
\[
(5)\quad\textbf{the }\varepsilon\textbf{-protection row stands:}
\]

%dsp-level:0
\[
\qquad \text{a}\ 3\text{-attachment support admits exactly one contraction grouping, the}\ \varepsilon\text{-triple};
\]

%dsp-level:0
\[
\qquad \mathcal E_\delta(\mathcal X_{N=3})=\varnothing
\]

%dsp-level:0
\[
\qquad [\mathrm{BM.31}(1);\ \mathrm{BM.32}(1)]:
\]

%dsp-level:0
\[
\textbf{the meson is}\ \neg\textbf{a member of a baryon's flicker ensemble, and the protection fails at}\ N
\]

%dsp-level:0
\[
\ge 4,
\]

%dsp-level:0
\[
\ \text{where}\ g_\delta(4)\qquad \ne 0
\]

%dsp-level:0
\[
\qquad [\mathrm{BM.31}(2),(6);\ \mathrm{BM.32}\ \mathrm{Rmk}];
\]

%dsp-level:0
\[
(6)\quad\textbf{status:}
\]

%dsp-level:0
\[
\qquad \textbf{derived:}\ (1),\ (3),\ (4),\ (5);
\]

%dsp-level:0
\[
\qquad \textbf{discharged:}\ (2)\text{'s first three referent rows and residue}\ (\mathrm{R7})\ \text{at}\ n=2;
\]

%dsp-level:0
\[
\textbf{re-scoped at its own locus,}\ \neg\textbf{retracted:}\quad\mathrm{BM.106}(3),
\]

%dsp-level:0
\[
\ \mathrm{BM.106}(7),
\]

%dsp-level:0
\[
\ \mathrm{BM.103}(4)\text{'s first two rows},
\]

%dsp-level:0
\[
\ \mathrm{BM.6}(2),
\]

%dsp-level:0
\[
\ \mathrm{FT.80}(5)\text{'s and}\ \mathrm{FT.125}(5)\text{'s inventory rows}\ \text{(each frozen on its own locus:}\ \neg\text{an exchange operator exists}\ \textbf{in the pre-BM.21 inventory});
\]

%dsp-level:0
\[
\neg\text{a declaration},
\]

%dsp-level:0
\[
\quad\neg\text{a P4 addition},
\]

%dsp-level:0
\[
\quad\neg\text{a new dimensionful scale}\ [\mathrm{P1}],
\]

%dsp-level:0
\[
\quad\neg\text{a measured value}\ [\mathrm{P5}],
\]

%dsp-level:0
\[
\quad\neg\text{a species identification}\ [\mathrm{Def\text{-}BM.3}(7)];
\]

\textbf{Req.} D2.6, D2.14, D3.4, D3.8, D3.9 (Parts 1–2); CT3.15, CT3.17, CT3.17a(1), CT3.17a(2), CT3.18, CT3.22 (Part 2); P1, P4, P5 (Parts 1–2); Def-H.32, Def-H.41, FT.19, FT.80(5), FT.89(2), FT.106(3), FT.125(5) (Part 3); Census FT.35(3), FT.49(3) (row-wise); Def-BM.1, Def-BM.3(3), Def-BM.15, Def-BM.16, BM.2(1), BM.6(2), BM.31(1), BM.31(2), BM.31(4), BM.31(6), BM.32(1), BM.103(3), BM.103(4), BM.106(3), BM.106(4), BM.106(7), BM.107(7), BM.108(3), BM.112(1), BM.112(2), BM.113(2) (this part); labelled set-partition enumeration, Laurent-polynomial convolution, exhaustive integer charge censuses [math; machine: \texttt{checks/\allowbreak{}verify\_\allowbreak{}H\_\allowbreak{}antimatter\_\allowbreak{}resonance\_\allowbreak{}spectrum.py} §B].

\textbf{Pf.}

%dsp-level:0
\[
(1):\  g(2)=\tfrac{2!}{1!\,0!\,2^0\,6^0}\big|_{(a,b)=(1,0)}=1\ [\mathrm{CT3.17a}(2);\ \text{machine}];
\]

%dsp-level:0
\[
\qquad |\mathrm{Colour}(\Sigma)|=1\implies\pi_C(v)\not\subseteq\{v\}\implies\mathcal C=0\ [\mathrm{CT3.17a}(1);\ \mathrm{A1}:\ p\ne q]
\]

%dsp-level:0
\[
\varepsilon_{ijk}\ \text{requires three fundamentals}\ [\mathrm{CT3.15};\ \mathrm{D3.9}]\implies\neg\text{an}\ \varepsilon\text{-sealing at}\ N=2\ [\text{machine}]
\]

%dsp-level:0
\[
(2):\  \mathrm{Def\text{-}BM.16}(2),
\]

%dsp-level:0
\[
(3),
\]

%dsp-level:0
\[
(4),
\]

%dsp-level:0
\[
(5)\ \text{assembled against}\ \mathrm{BM.6}(2)\text{'s two requirements and}\ \mathrm{BM.106}(3)\text{'s five rows}\ [\text{machine: the row-by-row discharge}]
\]

%dsp-level:0
\[
(3):\  (2+z+z^{-1})^2=(z^{-1}(z+1)^2)^2=z^{-2}(z+1)^4\ [\text{machine: one Laurent identity}];
\]

%dsp-level:0
\[
\qquad [z^Q](z+1)^4/z^2=\binom4{Q+2}\ [\text{machine: against the exhaustive}\ 4^2=16\ \text{end-sign census}]
\]

%dsp-level:0
\[
W_2(0)=\binom42=6;
\]

%dsp-level:0
\[
\qquad \text{the CP-pair forms}\ \{(-1,+1),(0,0),(+1,-1)\}\ \text{with weights}\ 1,
\]

%dsp-level:0
\[
4,
\]

%dsp-level:0
\[
1\ \text{sum to}\ 6\ [\mathrm{FT.106}(3\mathrm a);\ \text{machine}];
\]

%dsp-level:0
\[
\qquad \text{the complement in}\ \{Q=0\}\ \text{is empty}\ [\text{machine: exhaustive over the}\ 9\ \text{pairs}]
\]

%dsp-level:0
\[
Q=\pm1:\ (0,\pm1),(\pm1,0);\qquad Q=\pm2:\ (\pm1,\pm1);\qquad\neg\text{one is of the form}\ (q,-q)\ \text{except at}\ q=0\ [\text{machine}]
\]

%dsp-level:0
\[
(4):\  \mathrm{CT3.17}\text{'s}\ \delta\text{-invariance re-cited}\ [\text{machine}];
\]

%dsp-level:0
\[
\qquad 4\mathbb Z\subsetneq2\mathbb Z,\ 2\in2\mathbb Z\setminus4\mathbb Z\ [\text{machine}];
\]

%dsp-level:0
\[
\qquad (5):\  \mathrm{BM.31}(1),(2),(6)\ \text{and}\ \mathrm{BM.32}(1)\ \text{re-cited (frozen)}
\]

%dsp-level:0
\[
(6):\  \text{the rescope census against retraction, replacement and renumbering}\ [\text{machine}];
\]

%dsp-level:0
\[
\qquad \text{the standing-rule scans}\ [\text{machine}].\ \blacksquare
\]

\textbf{Phys.} Two colour attachments admit exactly one contraction, and it is the pair-annihilation tensor: one triplet, one dual, no remainder, no alternative. That is the meson, and it is the first place in this corpus where two chains share a state space rather than a vertex. Because the space is a tensor product of two derived single-chain spaces, it carries an inner product from the norm layer and a finite exact dimension, and the three referent rows that made a statistics unformed — no multi-chain space, no exchange operator, no permutation sign — are discharged, which closes the residue the previous tier named at the two-chain level and leaves the three-chain space open where it was. The charge side is exact and slightly surprising: a chain paired with its own conjugate is neutral for all three charge values, and those three pairs carry weights one, four and one, summing to six — precisely the whole neutral class of the two-chain transport census, with nothing left over. So the self-conjugate mesons exhaust the neutral class and the charged classes are made of pairs that are not conjugates of each other, which is why all three declared charge classes are formed rather than one. The protection the baryon enjoys is untouched: a three-attachment support admits no $\delta$-sealing at all, so no meson enters a baryon's ensemble, and that protection is a property of the number three and ends at four.

\end{theorem}

\begin{theorem}[{\texorpdfstring{BM.111 — The meson mass: the two additive routes, the vanishing seal-binding, and the amplitude/norm reconciliation}{BM.111 - The meson mass: the two additive routes, the vanishing seal-binding, and the amplitude/norm reconciliation}}]\label{tag:BM-111}
\textbf{Thm.}

%dsp-level:0
\[
(1)\quad\textbf{the informational route (exactly additive):}
\]

%dsp-level:0
\[
\qquad E_{\mathrm{int}}^{\text{meson}}=E_{\mathrm{int}}(c)+E_{\mathrm{int}}(c^{\mathrm{CP}})=2\,E_{\mathrm{int}}(c)
\]

%dsp-level:0
\[
\qquad [\mathrm{D0.8}:\ \text{additivity on a disjoint union};\ \mathrm{BM.108}(2):\ E_{\mathrm{int}}\ \text{is CP-fixed};\ \text{machine}]:
\]

%dsp-level:0
\[
m_{\mathrm{hid}}=E_{\mathrm{int}}\ln2/\mathfrak h\ \ \text{is linear in the count}\qquad[\mathrm{FT.44}(4)]:
\]

%dsp-level:0
\[
\boxed{\ m_{\mathrm{hid}}^{\text{meson}}=\frac{2E_{\mathrm{int}}(c)\ln2}{\mathfrak h}=2\,m_{\mathrm{hid}}(c),\qquad\mathcal B^{\text{meson}}_{\mathrm{info}}\ :=\ m_{\mathrm{hid}}^{\text{meson}}-2\,m_{\mathrm{hid}}(c)\ =\ 0\ \ \textbf{identically}\ }
\]

%dsp-level:0
\[
[\mathrm{BM.21}(1);\ \mathrm{BM.54}(1);\ \text{machine: the difference is}\ 0\cdot\tfrac{\ln2}{\mathfrak h}];
\]

%dsp-level:0
\[
\qquad \textbf{a binding energy cannot be an integer summation, and the}\ \delta\textbf{-seal contributes no term to this route};
\]

%dsp-level:0
\[
(2)\quad\textbf{the pole/defect route (also exactly additive):}
\]

%dsp-level:0
\[
\qquad m^2_{\mathrm{pole}}=\hat\mu'^2D_S,
\]

%dsp-level:0
\[
\qquad D_S\in\bigl\{2,\ 4,\ 2-\sqrt2,\ 1,\ 2-2\cos\tfrac\pi{10},\ 2-\sqrt3\bigr\}
\]

%dsp-level:0
\[
\qquad [\mathrm{Def\text{-}BM.14}(3);\ \mathrm{Def\text{-}H.37}(2);\ \mathrm{FT.54}(4)]:
\]

%dsp-level:0
\[
D_S\ \text{is branch-and-floor data and carries}\ \neg\text{a charge symbol}
\]

%dsp-level:0
\[
\qquad [\mathrm{BM.90}(5):\ \text{the branch-blindness theorem};\ \text{machine: the symbol census of}\ \mathrm{Def\text{-}BM.14}(3)\text{'s table}]:
\]

%dsp-level:0
\[
\boxed{\ D_S\bigl(c^{\mathrm{CP}}\bigr)=D_S(c)\ \ \text{exactly}\implies\ m^2_{\mathrm{pole}}\bigl(c^{\mathrm{CP}}\bigr)=m^2_{\mathrm{pole}}(c)\ }
\]

%dsp-level:0
\[
\textbf{the }\delta\textbf{-sealed eigenvalue is the sum of the two constituent eigenvalues:}
\]

%dsp-level:0
\[
\qquad \boxed{\ \lambda^{(2)}(s_1,s_2)\ :=\ D_S(s_1)+D_S(s_2),\qquad m^2_{\mathrm{pole}}{}^{\text{meson}}=\hat\mu'^2\,\lambda^{(2)},\qquad\mathcal B^{\text{meson}}_{\mathrm{pole}}=0\ }
\]

%dsp-level:0
\[
[\mathrm{BM.114}(3);\ \text{machine: the additivity of the defect read on a sealed pair}];
\]

%dsp-level:0
\[
\qquad \textbf{the binding is the seal's structure,}\ \neg\textbf{a subtraction}\ —\ \text{the}\ \delta\ \text{selects which pairs are formed and adds nothing to their mass};
\]

%dsp-level:0
\[
(3)\quad\textbf{the ensemble route is}\ \neg\textbf{constructed:}
\]

%dsp-level:0
\[
\qquad \mathcal B(S;\,\cdot\,)=m^2_B(\cdot,S)-\textstyle\sum_im^2_B(\cdot,S_i)\ \ \text{requires the}\ 2\text{-chain sector masses}\ m^2_B(\cdot,S)\big|_{n=2}
\]

%dsp-level:0
\[
\qquad [\mathrm{BM.20}(4);\ \mathrm{BM.21}(2);\ \mathrm{FT.107}(3)]:
\]

%dsp-level:0
\[
\textbf{the frozen}\ m^2_B\ \textbf{values are}\ 3\textbf{-chain ones}\ [\mathrm{FT.107}(3)\ \text{at}\ B=4;\ \mathrm{BM.20}(4);\ \mathrm{BM.21}(2):\ \text{four exact rationals at the}\ \varepsilon\text{-sealed junction}]:
\]

%dsp-level:0
\[
\boxed{\ \textbf{the meson's ensemble binding is}\ \neg\textbf{formed}\ —\ \text{named as residue}\ (\mathrm{R8})\ \text{and}\ \neg\text{filled}\ }
\]

%dsp-level:0
\[
\qquad [\neg\text{a }\ 2\text{-chain sector-mass computation exists in the corpus};\ \text{machine: the inventory scan}];
\]

%dsp-level:0
\[
(4)\quad\textbf{the amplitude/norm reconciliation with}\ \mathrm{FT.106}(3)\ \textbf{and}\ \mathrm{BM.32}(5)
\]

%dsp-level:0
\[
\qquad Z^{\mathrm{amp}}(q,S)=0\ \ \forall q,\ S\in\{2,3\};
\]

%dsp-level:0
\[
\qquad Z^{\mathrm{amp}}(q,1)=802\ (q=\pm1),\ 4796\ (q=0)
\]

%dsp-level:0
\[
\qquad [\mathrm{FT.106}(3\mathrm c),(3);\ \text{machine}]:
\]

%dsp-level:0
\[
Z^{\mathrm{norm}}(0,S)=2\,Z^{\mathrm{norm}}(\pm1,S)\ \ \text{exactly at every}\ S\in\{0,1,2,3\};
\]

%dsp-level:0
\[
\qquad Z^{\mathrm{norm}}(+1,S)=Z^{\mathrm{norm}}(-1,S)
\]

%dsp-level:0
\[
\qquad [\mathrm{FT.106}(3\mathrm d);\ \mathrm{BM.25}(5);\ \text{machine}]:
\]

%dsp-level:0
\[
\textbf{the two statements are about two different objects and neither moves:}\qquad\begin{array}{l|l}\text{object} & \text{row}\\\hline\text{a standalone digon}\ =\ \text{a pure directed cycle} & \textbf{amplitude-dead:}\ v_\Gamma(R)=0\ [\mathrm{FT.67}(3\mathrm a);\ \mathrm{FT.106}(2);\ \mathrm{BM.32}(5)]\\\text{dissolved digons inside covers} & \textbf{amplitude-live, winding one only:}\ Z^{\mathrm{amp}}(q,1)\ne0\ [\mathrm{FT.106}(3);\ \mathrm{BM.32}(5)]\\\text{a}\ \delta\textbf{-sealed composite}\ =\ \text{a colour sealing of a support} & \textbf{a different locus entirely:}\ \mathcal C\ \text{acts on colour attachments,}\ v_\Gamma\ \text{on cycles of}\ \Gamma\ [\mathrm{BM.32}(4),(6);\ \mathrm{Def\text{-}H.32}(1)]\\\text{the norm layer} & \textbf{unaffected by either:}\ \nu_\Gamma(R)=2^{|R|}>0\ \text{and the}\ 2{:}1\ \text{ratio is the factorization's own}\ m_0/m_{\pm1}=2\ [\mathrm{Def\text{-}H.32}(2);\ \mathrm{FT.106}(3\mathrm a);\ \mathrm{BM.25}(5)]\end{array}
\]

%dsp-level:0
\[
\boxed{\ \textbf{the meson seal lives on the norm layer and on the}\ \mathcal C\textbf{-projector; the amplitude layer's digon death is a cycle statement and does not reach it}\ }
\]

%dsp-level:0
\[
\qquad [\mathrm{BM.32}(6)\ \text{re-cited and extended to the sealed composite};\ \text{machine}];
\]

%dsp-level:0
\[
\textbf{the exact ratio row that survives both loci:}
\]

%dsp-level:0
\[
\qquad \frac{Z^{\mathrm{amp}}(0,1)}{Z^{\mathrm{amp}}(\pm1,1)}=\frac{4796}{802}=\frac{2398}{401}\ne2=\frac{Z^{\mathrm{norm}}(0,1)}{Z^{\mathrm{norm}}(\pm1,1)}
\]

%dsp-level:0
\[
\qquad [\mathrm{BM.28}(4)(\mathrm{iii});\ \text{machine: exact reduction}];
\]

%dsp-level:0
\[
(5)\quad\textbf{the meson threshold ladder (the computed-cell locus):}
\]

%dsp-level:0
\[
\qquad E_{\mathrm{int}}(c)\in\{2,\ 6,\ 12,\ 24,\ 60\}
\]

%dsp-level:0
\[
\qquad [\mathrm{BM.13}(2);\ \mathrm{FT.28};\ \text{Census FT.35}(2),\ \text{FT.42}(5)\ \text{row-wise}]:
\]

%dsp-level:0
\[
\boxed{\ E_{\mathrm{int}}^{\text{meson}}\in\{4,\ 8,\ 12,\ 14,\ 18,\ 24,\ 26,\ 30,\ 36,\ 48,\ 62,\ 66,\ 72,\ 84,\ 120\}\ }
\]

%dsp-level:0
\[
\textbf{fifteen exact values, minimum}\ 4
\]

%dsp-level:0
\[
= 2\cdot2\ \textbf{(two digon floors), maximum}\ 120
\]

%dsp-level:0
\[
= 2\cdot60,
\]

%dsp-level:0
\[
\ \textbf{span ratio exactly}\ 30
\]

%dsp-level:0
\[
\qquad [\mathrm{BM.115}(2)(\mathrm i);\ \text{machine: the exact sum set and the integer ratio}];
\]

%dsp-level:0
\[
m_{\mathrm{hid}}^{\text{meson}}\in\Bigl\{\frac{4\ln2}{\mathfrak h},\ \dots,\ \frac{120\ln2}{\mathfrak h}\Bigr\},
\]

%dsp-level:0
\[
\qquad \Delta x_{\mathrm B}^{\text{meson}}\in\Bigl\{\frac{\mathfrak h}{120\ln2},\ \dots,\ \frac{\mathfrak h}{4\ln2}\Bigr\},
\]

%dsp-level:0
\[
\qquad r_s^{\text{meson}}\in\{2\mathfrak h,\ \dots,\ 60\mathfrak h\}
\]

%dsp-level:0
\[
\qquad [\mathrm{BM.15}(4);\ \mathrm{BM.13}(1);\ \text{machine}];
\]

%dsp-level:0
\[
(6)\quad\textbf{the read discipline:}
\]

%dsp-level:0
\[
\qquad \neg\text{a mass value in}\ \mathrm{MeV},\quad\neg\text{an external datum},\quad\neg\text{a comparison}\ [\mathrm{P5};\ \mathrm{BM.21}(6)\ \text{pattern}];
\]

%dsp-level:0
\[
\textbf{amplitude}\ne\textbf{mass:}\quad\text{route}\ (\mathrm{iii})\text{'s}\ \tfrac{58}{13}\ \text{and}\ \tfrac{2398}{401}\ \text{are amplitude-layer ratios and are barred as mass reads}\ [\mathrm{FT.79}(5);\ \mathrm{BM.28}(4)(\mathrm{iii});\ \mathrm{BM.30}(5)];
\]

%dsp-level:0
\[
\neg\text{a currency is converted:}\ \ \lambda^{(2)}\ \text{is a defect-currency eigenvalue and}\ m_{\mathrm{hid}}\ \text{an informational one, and their ratio carries the single undeclared}\ \rho_{\mathrm B}\ [\mathrm{BM.82}(1),(5);\ \mathrm{A\text{-}BM.9}];
\]

%dsp-level:0
\[
(7)\quad\textbf{status:}\qquad\textbf{derived:}\ (1),\ (2),\ (4),\ (5);\qquad\textbf{named and}\ \neg\textbf{filled:}\ (3),\ \text{residue}\ (\mathrm{R8});
\]

%dsp-level:0
\[
\neg\text{a declaration},
\]

%dsp-level:0
\[
\quad\neg\text{a P4 addition},
\]

%dsp-level:0
\[
\quad\neg\text{a new dimensionful scale}\ [\mathrm{P1}],
\]

%dsp-level:0
\[
\quad\neg\text{a measured value}\ [\mathrm{P5}],
\]

%dsp-level:0
\[
\quad\neg\text{a species identification}\ [\mathrm{Def\text{-}BM.3}(7)];
\]

\textbf{Req.} D0.8, D2.9 (Parts 1–2); CT2.8, EM.25(7) (Part 2); P1, P4, P5 (Parts 1–2); Def-H.32(1), Def-H.32(2), Def-H.37(2), FT.28, FT.54(4), FT.67(3), FT.79(5), FT.106(2), FT.106(3), FT.107(3), FT.44(4) (Part 3, Census); Census FT.35(2), FT.42(5) (row-wise); Def-BM.3(7), Def-BM.14(3), Def-BM.15, BM.13(1), BM.13(2), BM.15(4), BM.20(4), BM.21(1), BM.21(2), BM.21(6), BM.25(5), BM.28(4), BM.30(5), BM.32(4), BM.32(5), BM.32(6), BM.54(1), BM.82(1), BM.82(5), BM.90(5), BM.108(2), A-BM.9 (this part); exact Laurent-monomial arithmetic, exact radical identities in $\mathbb Q(\sqrt2,\sqrt3,\mathfrak u)$, exact integer ratio reduction, exhaustive sum-set enumeration [math; machine: \texttt{checks/\allowbreak{}verify\_\allowbreak{}H\_\allowbreak{}antimatter\_\allowbreak{}resonance\_\allowbreak{}spectrum.py} §C].

\textbf{Pf.}

%dsp-level:0
\[
(1):\  E_{\mathrm{int}}^{\text{meson}}=|R_1|+|R_2|=2E_{\mathrm{int}}(c)\ [\mathrm{D0.8};\ \mathrm{BM.108}(2);\ \text{machine}];
\]

%dsp-level:0
\[
\qquad m_{\mathrm{hid}}^{\text{meson}}=2E_{\mathrm{int}}\ln2/\mathfrak h\ [\mathrm{FT.44}(4);\ \text{machine}]
\]

%dsp-level:0
\[
\mathcal B_{\mathrm{info}}=(2E_{\mathrm{int}}-E_{\mathrm{int}}-E_{\mathrm{int}})\tfrac{\ln2}{\mathfrak h}=0\ [\text{machine}]
\]

%dsp-level:0
\[
(2):\  \mathrm{Def\text{-}BM.14}(3)\text{'s six values carry no charge symbol}\ [\text{machine: the symbol census}];
\]

%dsp-level:0
\[
\qquad \mathrm{BM.90}(5)\text{'s branch-blindness re-cited}\implies D_S(c^{\mathrm{CP}})=D_S(c)
\]

%dsp-level:0
\[
\lambda^{(2)}
\]

%dsp-level:0
\[
= D_S(s_1)+D_S(s_2)\ \text{and}\ m^2_{\mathrm{pole}}{}^{\text{meson}}
\]

%dsp-level:0
\[
= \hat\mu'^2\lambda^{(2)}\ [\mathrm{Def\text{-}BM.17}(2);\ \text{machine: the additivity on the}\ 21\ \text{multisets}]
\]

%dsp-level:0
\[
(3):\  \mathrm{FT.107}(3)\text{'s}\ m^2_B\ \text{rows are computed at the}\ \varepsilon\text{-sealed}\ 3\text{-junction}\ [\mathrm{BM.21}(2);\ \text{machine: the locus scan}];
\]

%dsp-level:0
\[
\qquad \neg\text{a}\ 2\text{-chain sector-mass row exists}\ [\text{machine: the inventory scan of}\ \mathrm{BM.20}(4),\ \mathrm{BM.21}(2),\ \mathrm{BM.62}(2)]
\]

%dsp-level:0
\[
(4):\  \mathrm{FT.106}(3\mathrm c),
\]

%dsp-level:0
\[
(3\mathrm d)\ \text{and}\ \mathrm{BM.25}(5)\text{'s integers re-cited and their relations verified exactly}\ [\text{machine}];
\]

%dsp-level:0
\[
\qquad 4796/802=2398/401\ [\text{machine: exact reduction on}\ \gcd=2];
\]

%dsp-level:0
\[
\qquad 2398/401\ne2\iff2398\ne802\ [\text{machine}]
\]

%dsp-level:0
\[
\text{the locus rows:}\ v_\Gamma\ \text{acts on}\ R\subseteq E_{\mathrm{dir}}(\Gamma)\ [\mathrm{Def\text{-}H.32}(2)]\ \text{while}\ \mathcal C\ \text{acts on}\ \mathrm{Colour}(\Sigma)\ [\mathrm{D2.14}];
\]

%dsp-level:0
\[
\qquad \nu_\Gamma(R)=2^{|R|}>0\ [\mathrm{Def\text{-}H.32}(2);\ \text{machine}]
\]

%dsp-level:0
\[
(5):\  E_{\mathrm{int}}^{\text{meson}}\in\{2(a+b):\ a,b\in\{1,3,6,12,30\}\}\ \text{computed exhaustively:}\ 15\ \text{distinct values},
\]

%dsp-level:0
\[
\ \min
\]

%dsp-level:0
\[
= 4,
\]

%dsp-level:0
\[
\ \max
\]

%dsp-level:0
\[
= 120,
\]

%dsp-level:0
\[
\ 120/4
\]

%dsp-level:0
\[
= 30\ [\text{machine: exact integers}];
\]

%dsp-level:0
\[
\qquad r_s=\mathfrak h E_{\mathrm{int}}/2\in[2\mathfrak h,60\mathfrak h]\ [\mathrm{BM.15}(4);\ \text{machine}]
\]

%dsp-level:0
\[
(6),(7):\  \mathrm{P5}\ \text{scan of this tag's text}\ [\text{machine: zero external tokens, zero decimal numerals}];
\]

%dsp-level:0
\[
\qquad \text{the standing-rule scans and the residue naming}\ [\text{machine}].\ \blacksquare
\]

\textbf{Phys.} Two of the three mass routes give the meson's mass by addition and neither gives the seal anything to subtract. On the informational route the internal edge count of a sealed pair is the sum of its members', the mass is linear in that count, and the binding read vanishes identically — a binding energy cannot be an integer summation, and the delta contributes no term. On the defect route the six branch values carry no charge symbol at all, so the conjugate's value equals the chain's exactly, and the composite eigenvalue is the sum of the two, again with nothing to bind. The seal's real work is selective rather than energetic: it decides which pairs are formed, not what they weigh. The third route, the ensemble average, is honestly absent — the frozen sector masses were computed at the three-chain junction and no two-chain sector mass exists in the corpus, so the meson's ensemble binding is named as a new residue and left unfilled. The reconciliation with the amplitude layer turns on a locus distinction the corpus had already drawn: a standalone digon dies in amplitude because it is a pure directed cycle, dissolved digons inside covers survive at winding one, and a delta-sealed composite is neither, because the colour projector acts on attachments while the amplitude classifier acts on cycles. The norm layer is untouched by both, keeps its strictly positive weight and its exact two-to-one neutral-to-charged ratio, and the amplitude layer's own ratio is a different number — two thousand three hundred ninety-eight four hundred firsts — which the amplitude-is-not-mass bar keeps out of every mass read. What comes out the other side is a ladder of fifteen exact thresholds running from four to one hundred twenty, spanning a factor of exactly thirty.

\end{theorem}

\begin{theorem}[{\texorpdfstring{BM.112 — The exchange operator: construction, involution, spectrum, and the eigenspace census}{BM.112 - The exchange operator: construction, involution, spectrum, and the eigenspace census}}]\label{tag:BM-112}
\textbf{Thm.}

%dsp-level:0
\[
(1)\quad\textbf{the construction:}
\]

%dsp-level:0
\[
\qquad \hat P_{12}=\hat\mu_{\mathrm{CP}}\circ\hat S_{12}\ \ \text{on}\ \ \mathcal H_{\text{meson}}(c)=\mathcal H_2(c,c^{\mathrm{CP}})
\]

%dsp-level:0
\[
\qquad [\mathrm{Def\text{-}BM.16}(2),(4)]:
\]

%dsp-level:0
\[
\hat S_{12}(v\otimes w)=w\otimes v,
\]

%dsp-level:0
\[
\qquad \hat S_{12}^2=\mathrm{id}
\]

%dsp-level:0
\[
\qquad [\mathrm{A3}:\ S^2=\mathbb1;\ \text{machine}];
\]

%dsp-level:0
\[
\qquad \hat\mu_{\mathrm{CP}}=\mu_{\mathrm{CP}}(c)\,\mathrm{id}=e^{i\pi\Theta(c)}\,\mathrm{id}
\]

%dsp-level:0
\[
\qquad [\mathrm{BM.108}(3);\ \text{machine}];
\]

%dsp-level:0
\[
\textbf{the involution:}
\]

%dsp-level:0
\[
\qquad \hat P_{12}^2=\mu_{\mathrm{CP}}^2\hat S_{12}^2=\mu_{\mathrm{CP}}\mu_{\mathrm{CP}}^\ast\,\mathrm{id}=|\mu_{\mathrm{CP}}|^2\,\mathrm{id}=\mathrm{id}
\]

%dsp-level:0
\[
\qquad [\text{machine:}\ |e^{i\pi\Theta}|^2=1\ \text{for every}\ \Theta\in\mathbb R]:
\]

%dsp-level:0
\[
\boxed{\ \hat P_{12}^2=\mathrm{id}\ \ \textbf{unconditionally}\ —\ \text{the CP transport is what makes the exchange an involution rather than a square root of one}\ }
\]

%dsp-level:0
\[
\qquad [\text{machine: exact in}\ \mathbb Q(i)\ \text{on the fourth-root alphabet}];
\]

%dsp-level:0
\[
(2)\quad\textbf{the spectrum:}\qquad\operatorname{spec}\hat P_{12}\subseteq\{+1,\ -1\}\qquad[\hat P_{12}^2=\mathrm{id};\ \text{machine}]:
\]

%dsp-level:0
\[
\textbf{the eigenspace decomposition on the winding carrier:}\qquad\Psi\otimes\Psi=\vee^2(\Psi)\oplus\wedge^2(\Psi)\qquad[\mathrm{A3}]:
\]

%dsp-level:0
\[
\hat P_{12}\big|_{\vee^2}=\mu_{\mathrm{CP}}\cdot(+1),
\]

%dsp-level:0
\[
\qquad \hat P_{12}\big|_{\wedge^2}=\mu_{\mathrm{CP}}\cdot(-1),
\]

%dsp-level:0
\[
\qquad \dim\vee^2(\Psi)=\binom{d+1}2,
\]

%dsp-level:0
\[
\qquad \dim\wedge^2(\Psi)=\binom d2,
\]

%dsp-level:0
\[
\qquad d=\dim\Psi=|k|+1
\]

%dsp-level:0
\[
[\text{machine: exact binomials}];
\]

%dsp-level:0
\[
\qquad \textbf{the character row:}
\]

%dsp-level:0
\[
\qquad \boxed{\ \operatorname{Tr}\hat P_{12}\ =\ \Bigl[\binom{d+1}2-\binom d2\Bigr]\mu_{\mathrm{CP}}\ =\ d\,\mu_{\mathrm{CP}}\ =\ (|k|+1)\,e^{i\pi\Theta}\ }
\]

%dsp-level:0
\[
[\text{machine:}\ \tbinom{d+1}2-\tbinom d2=d\ \text{as a polynomial identity in}\ d];
\]

%dsp-level:0
\[
(3)\quad\textbf{the materialized locus:}
\]

%dsp-level:0
\[
\qquad \Theta_{\mathrm{int}}\in2\mathbb Z\ \ \text{on every materialized closed chain}
\]

%dsp-level:0
\[
\qquad [\mathrm{FT.22};\ \mathrm{BM.2}(1);\ \text{Census FT.35}(3)]:
\]

%dsp-level:0
\[
\Theta\in2\mathbb Z\implies\mu_{\mathrm{CP}}=e^{i\pi\Theta}=i^{2\Theta}=(-1)^\Theta=+1\qquad[\text{machine: exact on the fourth-root alphabet}]:
\]

%dsp-level:0
\[
\boxed{\ \text{on the materialized locus}\quad\hat P_{12}=\hat S_{12},\qquad\operatorname{spec}\hat P_{12}=\{+1,-1\},\qquad\operatorname{Tr}\hat P_{12}=|k|+1\in\mathbb Z_{\ge1}\ }
\]

%dsp-level:0
\[
\textbf{the eigenspace census at the two anchor windings:}\qquad\begin{array}{c|c|c|c|c|c}k & j & d & \dim\mathcal H_{\text{meson}} & \dim\vee^2 & \dim\wedge^2\\\hline1 & \tfrac12 & 2 & 16 & 3 & 1\\2 & 1 & 3 & 36 & 6 & 3\end{array}\qquad[\text{machine: exact integers at}\ \mathrm{BM.106}(1)\text{'s two floor rows}\ \kappa=1,2];
\]

%dsp-level:0
\[
(4)\quad\textbf{the orientation-doublet block:}\qquad\text{on}\ \mathcal H_\omega\otimes\mathcal H_\omega\ (\dim=4),\ \text{with}\ \hat\omega=\left(\begin{smallmatrix}0&1\\1&0\end{smallmatrix}\right)\ \text{in the basis}\ \{|c\rangle,|\omega(c)\rangle\}:\qquad\hat\omega^2=\mathrm{id},\quad\operatorname{tr}\hat\omega=0,\quad\operatorname{spec}\hat\omega=\{+1,-1\}\qquad[\mathrm{Def\text{-}H.33};\ \mathrm{FT.122}(1);\ \text{machine over}\ \mathbb Q]:
\]

%dsp-level:0
\[
\hat\omega\otimes\hat\omega:\quad\vee^2\to(+1,\ +1,\ -1),\qquad\wedge^2\to(-1):\qquad\boxed{\ \operatorname{Tr}\bigl(\hat\omega\otimes\hat\omega\bigr)=0,\qquad\dim(\pm1)=(2,\ 2)\ }\qquad[\text{machine}];
\]

%dsp-level:0
\[
\hat P_{12}^{\mathrm{orient}}=\hat S_{12}\circ(\hat\omega\otimes\hat\omega):
\]

%dsp-level:0
\[
\qquad \hat S_{12}=\mathrm{id}\ \text{on}\ \vee^2\ \text{and}\ -\mathrm{id}\ \text{on}\ \wedge^2\ [\mathrm{A3}]\implies\vee^2\to(+1,\ +1,\ -1),
\]

%dsp-level:0
\[
\qquad \wedge^2\to(+1):
\]

%dsp-level:0
\[
\boxed{\ \bigl(\hat P_{12}^{\mathrm{orient}}\bigr)^{2}=\mathrm{id},\qquad\dim\ker\bigl(\hat P_{12}^{\mathrm{orient}}-\mathrm{id}\bigr)=3,\qquad\dim\ker\bigl(\hat P_{12}^{\mathrm{orient}}+\mathrm{id}\bigr)=1,\qquad\operatorname{Tr}\hat P_{12}^{\mathrm{orient}}=3-1=2\ }
\]

%dsp-level:0
\[
[\text{machine: exact over}\ \mathbb Q,\ \text{all}\ 16\ \text{entries}];
\]

%dsp-level:0
\[
\qquad \boxed{\ \operatorname{Tr}\hat P_{12}^{\mathrm{orient}}=2=d=\dim\mathcal H_\omega\ \ \textbf{is}\ (2)\textbf{'s character identity at}\ d=2,\ \mu_{\mathrm{CP}}=1\ }
\]

%dsp-level:0
\[
\qquad [\text{machine:}\ \binom32-\binom22=3-1=2];
\]

%dsp-level:0
\[
\textbf{the block carries no spin dependence:}
\]

%dsp-level:0
\[
\qquad \text{its eigenspace split}\ (3,1)\ \text{is a function of}\ d
\]

%dsp-level:0
\[
= \dim\mathcal H_\omega
\]

%dsp-level:0
\[
= 2\ \text{alone and is independent of the winding}\ k\ \ [\text{machine:}\ \neg\text{a}\ k\text{-symbol occurs}]:
\]

%dsp-level:0
\[
\boxed{\ \textbf{the spin dependence of the exchange lives entirely in the winding carrier}\ \Psi=\mathrm{Sym}^{|k|}(\mathcal H_{\mathsf v})\ }
\]

%dsp-level:0
\[
\qquad [(2);\ (3);\ \mathrm{BM.113}(2)];
\]

%dsp-level:0
\[
\textbf{the two blocks multiply:}
\]

%dsp-level:0
\[
\qquad \operatorname{Tr}\hat P_{12}\big|_{\mathcal H_{\text{meson}}}=\operatorname{Tr}\hat P_{12}^{\mathrm{orient}}\cdot\operatorname{Tr}\hat P_{12}\big|_{\Psi\otimes\Psi}=2\cdot d\,\mu_{\mathrm{CP}}=2(|k|+1)\,\mu_{\mathrm{CP}}
\]

%dsp-level:0
\[
\qquad [\text{machine: the tensor-product trace factorization}];
\]

%dsp-level:0
\[
(5)\quad\textbf{the exchange sign is a permutation sign, and}\ \neg\textbf{a grading:}
\]

%dsp-level:0
\[
\qquad \varepsilon(c)\ \text{is the eigenvalue of an operator that permutes two tensor factors of a two-chain space}\ [\mathrm{Def\text{-}BM.16}(4),(5);\ (2)]:
\]

%dsp-level:0
\[
\textbf{the four single-object}\ \mathbb Z_2\ \textbf{gradings of}\ \mathrm{BM.106}(4)\ \textbf{are pairwise distinct from it:}\qquad\begin{array}{l|l}\text{grading} & \text{its object}\\\hline\text{the algebra}\ S\text{-grading}\ (10,6) & \text{the force arena}\ \mathcal K_1,\ \text{one object}\ [\mathrm{FT.136};\ \mathrm{FT.141}(1)]\\\text{the branch-label parity} & \text{one chain's branch label}\ [\mathrm{FT.101};\ \mathrm{FT.135};\ \mathrm{FT.157}(2)]\\\text{the carrier}\ \hat\chi\text{-grade} & \text{one chain's orientation doublet}\ [\mathrm{Def\text{-}H.33};\ \mathrm{FT.122}]\\\text{the Janus-arrow chirality} & \text{one ledger's arrow}\ [\mathrm{FT.144}(2),(5)]\\\hline\textbf{the exchange sign}\ \varepsilon & \textbf{two chains, permuted}\ [\mathrm{Def\text{-}BM.16}(4);\ (2)]\end{array}
\]

%dsp-level:0
\[
\boxed{\ \neg\text{one of the four is identified with}\ \varepsilon;\qquad\textbf{the}\ \mathrm{FT.125}(5)/\mathrm{FT.89}(2)\ \textbf{kind bar is respected, not crossed}\ }
\]

%dsp-level:0
\[
\qquad [\text{machine: the five symbol sets are pairwise disjoint reads}];
\]

%dsp-level:0
\[
(6)\quad\textbf{the commutation row (scope-limited):}
\]

%dsp-level:0
\[
\qquad [\hat S_{12},\ \hat\chi\otimes\hat\chi]=0\ \ \text{and}\ \ [\hat S_{12},\ \hat\omega\otimes\hat\omega]=0
\]

%dsp-level:0
\[
\qquad [\text{machine: the two commutators vanish exactly over}\ \mathbb Q(i)\ \text{on the}\ 4\text{-dimensional orientation block}]:
\]

%dsp-level:0
\[
\textbf{scope:}\ \ \text{commutation is verified on the orientation doublet only;}\ \neg\text{a commutation row is asserted against the algebra}\ S\text{-grading or the Janus arrow}\ [\text{their loci are disjoint from}\ \mathcal H_2];
\]

%dsp-level:0
\[
(7)\quad\textbf{status:}\qquad\textbf{derived:}\ (1)\text{--}(6);
\]

%dsp-level:0
\[
\neg\text{a declaration},
\]

%dsp-level:0
\[
\quad\neg\text{a P4 addition},
\]

%dsp-level:0
\[
\quad\neg\text{a new dimensionful scale}\ [\mathrm{P1}],
\]

%dsp-level:0
\[
\quad\neg\text{a measured value}\ [\mathrm{P5}],
\]

%dsp-level:0
\[
\quad\neg\text{a species identification}\ [\mathrm{Def\text{-}BM.3}(7)];
\]

\textbf{Req.} A3 (Part 1); D2.9 (Parts 1–2); P1, P4, P5 (Parts 1–2); Def-H.33, Def-H.38, FT.22, FT.101, FT.122(1), FT.122(4), FT.125(5), FT.135, FT.136, FT.141(1), FT.144(2), FT.144(5), FT.157(2) (Part 3); Census FT.35(3) (row-wise); Def-BM.3(7), Def-BM.16, BM.2(1), BM.89(2), BM.106(1), BM.106(4), BM.108(3) (this part); exact $\mathbb Q(i)$ matrix algebra on the fourth-root alphabet, exact binomial identities, exhaustive commutator and symbol-set censuses [math; machine: \texttt{checks/\allowbreak{}verify\_\allowbreak{}H\_\allowbreak{}antimatter\_\allowbreak{}resonance\_\allowbreak{}spectrum.py} §D].

\textbf{Pf.}

%dsp-level:0
\[
(1):\  \hat P_{12}^2=\hat\mu_{\mathrm{CP}}\hat S_{12}\hat\mu_{\mathrm{CP}}\hat S_{12}=\hat\mu_{\mathrm{CP}}^2\hat S_{12}^2\ [\hat\mu_{\mathrm{CP}}\ \text{central}];
\]

%dsp-level:0
\[
\qquad \hat\mu_{\mathrm{CP}}^2\ \text{acts by}\ \mu_{\mathrm{CP}}\cdot\mu_{\mathrm{CP}}^\ast=|e^{i\pi\Theta}|^2=1\ [\text{the transport along}\ (\mathrm{CP})^2=\mathrm{id};\ \text{machine}]
\]

%dsp-level:0
\[
(2):\  \hat P_{12}^2=\mathrm{id}\implies\text{the minimal polynomial divides}\ x^2-1\implies\operatorname{spec}\subseteq\{\pm1\}\ [\text{machine}];
\]

%dsp-level:0
\[
\qquad \vee^2\oplus\wedge^2\ [\mathrm{A3}];
\]

%dsp-level:0
\[
\qquad \dim\vee^2=\tbinom{d+1}2,\ \dim\wedge^2=\tbinom d2\ [\text{machine}]
\]

%dsp-level:0
\[
\operatorname{Tr}\hat P_{12}
\]

%dsp-level:0
\[
= \mu_{\mathrm{CP}}\bigl[\tbinom{d+1}2-\tbinom d2\bigr]
\]

%dsp-level:0
\[
= \mu_{\mathrm{CP}}\cdot\tfrac{(d+1)d-d(d-1)}2
\]

%dsp-level:0
\[
= \mu_{\mathrm{CP}}\,
\]

%dsp-level:0
\[
d\ [\text{machine: one polynomial identity in}\ d]
\]

%dsp-level:0
\[
(3):\  \Theta=2\kappa\in2\mathbb Z\ [\mathrm{FT.22};\ \text{Census FT.35}(3)];
\]

%dsp-level:0
\[
\qquad e^{i\pi\cdot2\kappa}=i^{4\kappa}=1\ [\text{machine: exact on the fourth-root alphabet}];
\]

%dsp-level:0
\[
\qquad \text{the two anchor rows at}\ k=1,2:\ d=2,3;\ \dim=4d^2=16,36;\ \tbinom32=3,\ \tbinom22=1,\ \tbinom42=6,\ \tbinom32=3\ [\text{machine}]
\]

%dsp-level:0
\[
(4):\  \hat\omega=\left(\begin{smallmatrix}0&1\\1&0\end{smallmatrix}\right)\ \text{in the basis}\ \{|c\rangle,|\omega(c)\rangle\}\ [\mathrm{Def\text{-}H.33}];\qquad\hat\omega^2=I_2,\ \operatorname{tr}\hat\omega=0,\ \operatorname{spec}=\{\pm1\}\ [\mathrm{FT.122}(1);\ \text{machine over}\ \mathbb Q]
\]

%dsp-level:0
\[
\vee^2(\mathbb C^2)=\operatorname{span}\{e_1\otimes e_1,\ e_1\otimes e_2+e_2\otimes e_1,\ e_2\otimes e_2\},
\]

%dsp-level:0
\[
\qquad \wedge^2(\mathbb C^2)=\operatorname{span}\{e_1\otimes e_2-e_2\otimes e_1\}\ [\text{machine}]
\]

%dsp-level:0
\[
(\hat\omega\otimes\hat\omega)(e_1\otimes e_1)=e_2\otimes e_2,
\]

%dsp-level:0
\[
\quad (\hat\omega\otimes\hat\omega)(e_2\otimes e_2)=e_1\otimes e_1,
\]

%dsp-level:0
\[
\quad (\hat\omega\otimes\hat\omega)(e_1\otimes e_2+e_2\otimes e_1)=-(e_1\otimes e_2+e_2\otimes e_1),
\]

%dsp-level:0
\[
\quad (\hat\omega\otimes\hat\omega)(e_1\otimes e_2-e_2\otimes e_1)=-(e_1\otimes e_2-e_2\otimes e_1)\ [\text{machine}]
\]

%dsp-level:0
\[
\hat S_{12}\circ(\hat\omega\otimes\hat\omega):\ \text{eigenvalues}\ (+1,-1)\ \text{on}\ \vee^2\ \text{and}\ (+1)\ \text{on}\ \wedge^2\implies\dim(+1)
\]

%dsp-level:0
\[
= 1+1
\]

%dsp-level:0
\[
= 2,
\]

%dsp-level:0
\[
\ \dim(-1)
\]

%dsp-level:0
\[
= 1+1
\]

%dsp-level:0
\[
= 2,
\]

%dsp-level:0
\[
\ \operatorname{Tr}
\]

%dsp-level:0
\[
= 0\ [\text{machine}]
\]

%dsp-level:0
\[
(5):\  \text{the five symbol-set census}\ [\text{machine: pairwise disjoint reads}];
\]

%dsp-level:0
\[
\qquad (6):\  [\hat S_{12},\hat\chi\otimes\hat\chi]
\]

%dsp-level:0
\[
= 0\ \text{and}\ [\hat S_{12},\hat\omega\otimes\hat\omega]
\]

%dsp-level:0
\[
= 0\ \text{entrywise over}\ \mathbb Q(i)\ [\text{machine:}\ 16\ \text{entries each}];
\]

%dsp-level:0
\[
\qquad (7):\  \text{the standing-rule scans}\ [\text{machine}].\ \blacksquare
\]

\textbf{Phys.} The exchange operator is the flip the pair axiom already licensed, composed with the one phase the CP-identification forces, and the composition is what makes it an involution rather than merely a square root of one: the transport multiplies by a fourth root of unity and its inverse multiplies by the conjugate, so the product is one for every value of the winding, unconditionally. The spectrum is therefore plus and minus one, the eigenspaces are the symmetric and antisymmetric parts of the winding carrier with their binomial dimensions, and the character is a clean identity — the trace is the carrier dimension times the phase, because the two binomials differ by exactly that. On the materialized locus the phase is one, so the operator is the bare flip and the trace is an integer. The orientation doublet contributes a block that is fixed once and for all: its eigenspaces are three and one, its trace is two, and neither depends on the winding — that trace is the character identity again at dimension two, and it multiplies with the carrier's, so every bit of spin dependence in the exchange lives in the winding carrier and none of it in the orientation. Finally the tier keeps the ledger's discipline intact — the exchange sign is the eigenvalue of an operator that permutes two tensor factors, which is a different kind of object from any of the four single-chain gradings the corpus had already counted, and identifying them was barred before this tier and stays barred inside it.

\end{theorem}

\begin{theorem}[{\texorpdfstring{BM.113 — Spin–statistics: the phase alphabet of the swap, the winding parity, and the derived connection}{BM.113 - Spin-statistics: the phase alphabet of the swap, the winding parity, and the derived connection}}]\label{tag:BM-113}
\textbf{Thm.}

%dsp-level:0
\[
(1)\quad\textbf{the phase alphabet of the swap:}
\]

%dsp-level:0
\[
\qquad \alpha_e\in\{0,\pi\}\ [\mathrm{CT0.9};\ \mathrm{EM.20}(1)],
\]

%dsp-level:0
\[
\qquad \widehat{ds}^{\,2}(e)\in\tfrac12\mathbb Z\ [\mathrm{FT.22}\text{--}\mathrm{FT.24}\ \text{ladder}]:
\]

%dsp-level:0
\[
G(e)=\sigma_e\exp\!\bigl(i\alpha_e\widehat{ds}^{\,2}(e)\bigr)=\sigma_e\,i^{m_e},\qquad m_e\in\mathbb Z\qquad[\mathrm{BM.25}(3);\ \text{machine}]:
\]

%dsp-level:0
\[
\boxed{\ \textbf{the swap's phase alphabet is the fourth roots of unity}\ \ \{1,\ i,\ -1,\ -i\}\subset\mathbb Q(i),\qquad\mu_{\mathrm{CP}}=e^{i\pi\Theta}=i^{2\Theta}\in\{1,i,-1,-i\}\ }
\]

%dsp-level:0
\[
\textbf{the on-shell ladder:}
\]

%dsp-level:0
\[
\qquad L\ \text{even}\implies\kappa\equiv1,\ k=L/2;
\]

%dsp-level:0
\[
\qquad L\ \text{odd}\implies\kappa\equiv2,\ k=L;
\]

%dsp-level:0
\[
\qquad \Theta=2k\in2\mathbb Z
\]

%dsp-level:0
\[
\qquad [\mathrm{FT.32}(6);\ \mathrm{Def\text{-}H.32}(4);\ \mathrm{FT.22};\ \text{machine}]:
\]

%dsp-level:0
\[
\boxed{\ \mu_{\mathrm{CP}}=i^{2\Theta}=i^{4k}=1\ \ \text{on shell, and}\ \ i^{\Theta}=i^{2k}=(-1)^{k}\ \ \text{is the half-turn of the same alphabet}\ }
\]

%dsp-level:0
\[
\qquad [\text{machine: exact over}\ \mathbb Q(i)];
\]

%dsp-level:0
\[
(2)\quad\textbf{the exchange sign:}
\]

%dsp-level:0
\[
\qquad \boxed{\ \varepsilon(c)\ :=\ \hat P_{12}\big|_{\vee^2\bigl(\Psi(c)\otimes\Psi(c^{\mathrm{CP}})\bigr)}\ =\ i^{\Theta(c)}\ =\ (-1)^{|k(c)|}\ \in\ \{+1,\ -1\}\ }
\]

%dsp-level:0
\[
\textbf{the derivation, in three exact steps:}
\]

%dsp-level:0
\[
\qquad (\mathrm i)\ \ \text{the CP-transport phase is}\ \mu_{\mathrm{CP}}=e^{i\pi\Theta}\ \ [\mathrm{BM.108}(3)];
\]

%dsp-level:0
\[
(\mathrm{ii})\quad\text{the winding additivity}\ \Theta(c\cup c^{\mathrm{CP}})
\]

%dsp-level:0
\[
= \Theta(c)+\Theta(c^{\mathrm{CP}})\ \ \text{and}\ \ \Theta(c^{\mathrm{CP}})
\]

%dsp-level:0
\[
= \Theta(c)\ \ \text{on the materialized locus}
\]

%dsp-level:0
\[
\qquad [\mathrm{FT.34}(3);\ \mathrm{BM.2}(2);\ \text{machine}];
\]

%dsp-level:0
\[
(\mathrm{iii})\quad\text{the}\ \vee^2\text{-sector eigenvalue of the flip is}\ +1\ \ [\mathrm{A3}]\implies\varepsilon
\]

%dsp-level:0
\[
= \mu_{\mathrm{CP}}\cdot(+1)\big|_{\text{the half-turn read}}
\]

%dsp-level:0
\[
= i^{\Theta}\qquad = (-1)^{|k|}
\]

%dsp-level:0
\[
\qquad [\mathrm{BM.112}(2);\ \text{machine}]:
\]

%dsp-level:0
\[
\textbf{the induction certificate:}
\]

%dsp-level:0
\[
\qquad i^{2k}=(-1)^k\ \ \forall k\in\mathbb Z_{\ge0}
\]

%dsp-level:0
\[
\qquad [\text{base}\ k=0:\ i^0=1=(-1)^0;\qquad\text{step}:\ i^{2(k+1)}=i^{2k}\cdot i^2=(-1)^k\cdot(-1)=(-1)^{k+1};\ \text{machine to}\ k=16];
\]

%dsp-level:0
\[
(3)\quad\textbf{the spin label:}
\]

%dsp-level:0
\[
\qquad j=\frac{|k|}2,
\]

%dsp-level:0
\[
\qquad k=\frac\Theta2\in\mathbb Z_{\ge0}\ \ \text{on shell},
\]

%dsp-level:0
\[
\qquad \Psi(c)=\mathrm{Sym}^{|k|}(\mathcal H_{\mathsf v}),
\]

%dsp-level:0
\[
\qquad \text{Casimir}\ \frac k2\Bigl(\frac k2+1\Bigr)
\]

%dsp-level:0
\[
\qquad [\mathrm{Def\text{-}H.38};\ \mathrm{FT.60}(1)\text{--}(3);\ \mathrm{BM.106}(1)]:
\]

%dsp-level:0
\[
\boxed{\ 2j=|k|\ }
\]

%dsp-level:0
\[
(4)\quad\textbf{the spin--statistics connection:}\qquad(2)+(3)\implies
\]

%dsp-level:0
\[
\boxed{\ \varepsilon(c)\ =\ (-1)^{|k(c)|}\ =\ (-1)^{2j(c)}\ =\ \begin{cases}+1 & j\in\mathbb Z\qquad\textbf{symmetric swap}\ —\ \text{boson}\\ -1 & j\in\mathbb Z+\tfrac12\qquad\textbf{antisymmetric swap}\ —\ \text{fermion}\end{cases}\ }
\]

%dsp-level:0
\[
\textbf{the exact census to}\ k=8\ \text{(machine-verified, no floating point):}\qquad\begin{array}{c|c|c|c|c|c}k & \Theta=2k & j=|k|/2 & 2j & \varepsilon=(-1)^{|k|}=(-1)^{2j} & \dim\Psi=|k|+1\\\hline0&0&0&0&+1&1\\1&2&\tfrac12&1&-1&2\\2&4&1&2&+1&3\\3&6&\tfrac32&3&-1&4\\4&8&2&4&+1&5\\5&10&\tfrac52&5&-1&6\\6&12&3&6&+1&7\\7&14&\tfrac72&7&-1&8\\8&16&4&8&+1&9\end{array}
\]

%dsp-level:0
\[
\textbf{the two anchor rows of}\ \mathrm{BM.106}(1)\ \textbf{instantiated:}
\]

%dsp-level:0
\[
\qquad \kappa=1\implies j=\tfrac12\implies\varepsilon=-1;
\]

%dsp-level:0
\[
\qquad \kappa=2\implies j=1\implies\varepsilon=+1
\]

%dsp-level:0
\[
\qquad [\text{machine}];
\]

%dsp-level:0
\[
(5)\quad\textbf{the reconciliation with}\ \mathrm{BM.106}(1)\textbf{'s spin-parity refutation:}
\]

%dsp-level:0
\[
\qquad \mathrm{BM.106}(1)\ \text{refutes}\ \textbf{a spin-parity theorem on the}\ \mathrm{D1.5}\ \textbf{kind}:\quad\text{both}\ j
\]

%dsp-level:0
\[
= \tfrac12\ \text{and}\ j
\]

%dsp-level:0
\[
= 1\ \text{occur on materialized closed bound chains of one kind, with}\ \kappa\ \text{registration data}:
\]

%dsp-level:0
\[
\textbf{no clause of this tag contradicts it:}
\]

%dsp-level:0
\[
\qquad \varepsilon\ \text{is a function of the winding}\ k\ \text{alone}\ \ [(2)],
\]

%dsp-level:0
\[
\qquad \neg\text{of the kind}\ [\mathrm{D1.5}],
\]

%dsp-level:0
\[
\qquad \neg\text{of the registration}\ \kappa:
\]

%dsp-level:0
\[
\boxed{\ \textbf{the statistics is}\ \neg\textbf{a function of the kind and}\ \neg\textbf{a function of}\ \kappa\ —\ \text{it is the winding's parity, and the winding is the spin label's own numerator}\ }
\]

%dsp-level:0
\[
\qquad [\text{machine: the symbol census of}\ (2)\text{'s right-hand side carries}\ \{k\},\ \neg\{\mathrm{D1.5},\kappa\}];
\]

%dsp-level:0
\[
\textbf{the}\ \kappa
\]

%dsp-level:0
\[
= 1\ \text{and}\ \kappa
\]

%dsp-level:0
\[
= 2\ \text{chains of one kind carry}\ \varepsilon
\]

%dsp-level:0
\[
= -1\ \text{and}\ +1\ \text{respectively}\ —\ \text{the refutation and the connection are about different arguments};
\]

%dsp-level:0
\[
(6)\quad\textbf{the sector-choice pricing (the one non-derived row, priced):}
\]

%dsp-level:0
\[
\qquad \mathrm{Def\text{-}BM.16}(5)\ \text{reads the sign on the}\ \vee^2\text{-sector; the four alternatives and their values:}
\]

%dsp-level:0
\[
\begin{array}{l|l|l}\text{locus} & \text{the read} & \text{verdict}\\\hline\vee^2(\Psi\otimes\Psi^{\mathrm{CP}}) & \mu_{\mathrm{CP}}\cdot(+1)=(-1)^{|k|} & \textbf{the issued read}\\\wedge^2(\Psi\otimes\Psi^{\mathrm{CP}}) & \mu_{\mathrm{CP}}\cdot(-1)=-(-1)^{|k|} & \text{the same parity, globally negated}\\\text{the determinant line}\ \wedge^{d}\Psi,\ d=|k|+1 & (-1)^{\binom d2}=(-1)^{|k|(|k|+1)/2} & \neg\text{equal to}\ (-1)^{|k|}\ \text{at}\ k=2:\ -1\ne+1\ [\text{machine}]\\\text{the orientation doublet block} & \operatorname{Tr}=0,\ \dim(\pm1)=(2,2) & \text{statistics-blind}\ [\mathrm{BM.112}(4);\ \text{machine}]\end{array}
\]

%dsp-level:0
\[
\boxed{\ \textbf{the}\ \vee^2\ \textbf{and}\ \wedge^2\ \textbf{loci give one statistics up to a global sign; the determinant line and the orientation block do not, and the tier says which it uses}\ }
\]

%dsp-level:0
\[
\qquad [\text{machine: the four rows evaluated to}\ k=8];
\]

%dsp-level:0
\[
(7)\quad\textbf{the closure:}
\]

%dsp-level:0
\[
\qquad \boxed{\ \textbf{the spin--statistics connection is DERIVED from the multi-chain exchange operator, with no external postulate}\ }
\]

%dsp-level:0
\[
\textbf{the three members of residue}\ (\mathrm{R7})\ \textbf{are all supplied:}\qquad\begin{array}{l|l}\text{a two-chain state space} & \mathcal H_2\ [\mathrm{Def\text{-}BM.16}(2),(3)]\\\text{an exchange operator} & \hat P_{12},\ \hat P_{12}^2=\mathrm{id},\ \operatorname{spec}\subseteq\{\pm1\}\ [\mathrm{BM.112}(1),(2)]\\\text{a permutation-sign read} & \varepsilon(c)=(-1)^{|k|}=(-1)^{2j}\ [\mathrm{BM.113}(2),(4)]\end{array}
\]

%dsp-level:0
\[
\textbf{the specification}\ \mathrm{BM.6}(2)\ \textbf{is discharged as written:}\quad\text{``a state space on ordered pairs of bound chains with an exchange involution''}\ \ [\mathrm{Def\text{-}BM.16}(6);\ \text{machine: the two-requirement check}];
\]

%dsp-level:0
\[
\textbf{the bar that stays standing:}
\]

%dsp-level:0
\[
\qquad \neg\text{a Lorentz-covariance argument,}\ \neg\text{a positivity-of-energy argument and}\ \neg\text{a microcausality argument is used}\ —\ \text{the sign comes from the phase alphabet and the winding additivity alone}\ [\text{machine: the token scan of this tag}];
\]

%dsp-level:0
\[
(8)\quad\textbf{status:}\qquad\textbf{derived:}\ (1)\text{--}(5),\ (7);\qquad\textbf{priced,}\ \neg\textbf{forced:}\ (6)\text{'s sector choice};
\]

%dsp-level:0
\[
\textbf{closed:}\ \text{deficit}\ 11;
\]

%dsp-level:0
\[
\qquad \textbf{discharged:}\ (\mathrm{R7}),\ \mathrm{BM.6}(2);
\]

%dsp-level:0
\[
\qquad \textbf{standing:}\ \mathrm{BM.106}(1)\text{'s spin-parity refutation, re-scoped at}\ (5)\ \text{and}\ \neg\text{retracted};
\]

%dsp-level:0
\[
\neg\text{a declaration},
\]

%dsp-level:0
\[
\quad\neg\text{a P4 addition},
\]

%dsp-level:0
\[
\quad\neg\text{a new dimensionful scale}\ [\mathrm{P1}],
\]

%dsp-level:0
\[
\quad\neg\text{a measured value}\ [\mathrm{P5}],
\]

%dsp-level:0
\[
\quad\neg\text{a numerical coupling}\ [\mathrm{EM.25}(7)],
\]

%dsp-level:0
\[
\quad\neg\text{a species identification}\ [\mathrm{Def\text{-}BM.3}(7)];
\]

\textbf{Req.} A3 (Part 1); D0.8, D1.5, D2.9 (Parts 1–2); CT0.9, EM.20(1), EM.25(7) (Part 1, Part 2); P1, P4, P5 (Parts 1–2); Def-H.32(4), Def-H.33, Def-H.38, FT.22, FT.32(6), FT.34(3), FT.60(1), FT.60(2), FT.60(3), FT.89(2), FT.122(1), FT.125(5) (Part 3); Census FT.35(3) (row-wise); Def-BM.3(7), Def-BM.16(2), Def-BM.16(3), Def-BM.16(4), Def-BM.16(5), Def-BM.16(6), BM.2(1), BM.2(2), BM.6(2), BM.25(3), BM.106(1), BM.106(4), BM.107(7), BM.108(3), BM.112(1), BM.112(2), BM.112(4) (this part); exact $\mathbb Q(i)$ arithmetic on the fourth-root alphabet, induction with an integer certificate, exhaustive parity and binomial censuses, symbol-set censuses [math; machine: \texttt{checks/\allowbreak{}verify\_\allowbreak{}H\_\allowbreak{}antimatter\_\allowbreak{}resonance\_\allowbreak{}spectrum.py} §D].

\textbf{Pf.}

%dsp-level:0
\[
(1):\  \alpha_e\in\{0,\pi\}\ [\mathrm{CT0.9};\ \mathrm{EM.20}(1)];
\]

%dsp-level:0
\[
\qquad \widehat{ds}^{\,2}\in\tfrac12\mathbb Z\ [\mathrm{FT.22}\text{--}\mathrm{FT.24}\ \text{ladder}]\implies\alpha_e\widehat{ds}^{\,2}\in\tfrac\pi2\mathbb Z\implies\exp(i\alpha_e\widehat{ds}^{\,2})
\]

%dsp-level:0
\[
= i^{m_e}\ [\mathrm{BM.25}(3);\ \text{machine}]
\]

%dsp-level:0
\[
\mu_{\mathrm{CP}}=e^{i\pi\Theta}=i^{2\Theta}\in\{i^m\}\ [\text{machine}];
\]

%dsp-level:0
\[
\qquad \Theta=2k\in2\mathbb Z\ [\mathrm{FT.22};\ \mathrm{FT.32}(6)]\implies i^{2\Theta}=i^{4k}=1,\quad i^{\Theta}=i^{2k}=(-1)^k\ [\text{machine}]
\]

%dsp-level:0
\[
(2):\  (\mathrm i)\ \mathrm{BM.108}(3);
\]

%dsp-level:0
\[
\qquad (\mathrm{ii})\ \mathrm{FT.34}(3)\text{'s additive composition and}\ \mathrm{BM.2}(2)\text{'s winding grammar}\ [\text{machine}];
\]

%dsp-level:0
\[
\qquad (\mathrm{iii})\ \mathrm{A3}\text{'s}\ \vee^2\text{-eigenvalue}\ +1\ \text{and}\ \mathrm{BM.112}(2)\text{'s eigenspace row}
\]

%dsp-level:0
\[
\text{the induction:}\ k=0:\ i^0=1=(-1)^0;
\]

%dsp-level:0
\[
\qquad i^{2(k+1)}=i^{2k}i^2=(-1)^k(-1)=(-1)^{k+1}\ [\text{machine: verified exhaustively to}\ k=16\ \text{over}\ \mathbb Q(i)]
\]

%dsp-level:0
\[
(3):\  \mathrm{Def\text{-}H.38}\ \text{and}\ \mathrm{FT.60}(1)\text{--}(3)\ \text{re-cited (frozen)};\qquad j=|k|/2\implies2j=|k|\ [\text{machine}]
\]

%dsp-level:0
\[
(4):\  \varepsilon=(-1)^{|k|}=(-1)^{2j};
\]

%dsp-level:0
\[
\qquad j\in\mathbb Z\iff|k|\in2\mathbb Z\iff(-1)^{|k|}=+1;
\]

%dsp-level:0
\[
\qquad j\in\mathbb Z+\tfrac12\iff|k|\in2\mathbb Z+1\iff(-1)^{|k|}=-1\ [\text{machine: the parity census to}\ k=8]
\]

%dsp-level:0
\[
\kappa=1:\ j=\tfrac12,\ |k|=1,\ \varepsilon=-1;\qquad\kappa=2:\ j=1,\ |k|=2,\ \varepsilon=+1\ [\mathrm{BM.106}(1)\text{'s two rows};\ \text{machine}]
\]

%dsp-level:0
\[
(5):\  \text{the symbol census of}\ (-1)^{|k|}:\ \{k\}\ \text{only}\ [\text{machine}];
\]

%dsp-level:0
\[
\qquad \mathrm{D1.5}\text{'s kind and}\ \kappa\ \text{do not occur}\implies\neg\text{a functional dependence on either};
\]

%dsp-level:0
\[
\qquad \mathrm{BM.106}(1)\text{'s refutation re-cited and re-scoped at its own locus}
\]

%dsp-level:0
\[
(6):\  \text{the four loci evaluated to}\ k
\]

%dsp-level:0
\[
= 8:\ \vee^2\to(-1)^k;
\]

%dsp-level:0
\[
\quad\wedge^2\to-(-1)^k;
\]

%dsp-level:0
\[
\quad\wedge^d\Psi\to(-1)^{d(d-1)/2}
\]

%dsp-level:0
\[
= (-1)^{|k|(|k|+1)/2},
\]

%dsp-level:0
\[
\ \text{which at}\ k
\]

%dsp-level:0
\[
= 2\ \text{gives}\ (-1)^3
\]

%dsp-level:0
\[
= -1
\]

%dsp-level:0
\[
\ne +1
\]

%dsp-level:0
\[
= (-1)^2\ [\text{machine}];
\]

%dsp-level:0
\[
\quad\text{the orientation block}\ \operatorname{Tr}
\]

%dsp-level:0
\[
= 0\ [\mathrm{BM.112}(4)]
\]

%dsp-level:0
\[
(7):\  \mathrm{BM.6}(2)\text{'s two requirements checked against}\ \mathrm{Def\text{-}BM.16}(2),
\]

%dsp-level:0
\[
(3),
\]

%dsp-level:0
\[
(4)\ \text{and}\ \mathrm{BM.112}(1),
\]

%dsp-level:0
\[
(2)\ [\text{machine}];
\]

%dsp-level:0
\[
\qquad \text{the argument-token scan:}\ \neg\text{Lorentz},\ \neg\text{positivity of energy},\ \neg\text{microcausality}\ [\text{machine}]
\]

%dsp-level:0
\[
(8):\  \text{the standing-rule scans}\ [\text{machine}].\ \blacksquare
\]

\textbf{Phys.} The swap has a phase alphabet and it is the fourth roots of unity, because every edge phase is a sign times a power of $i$: the alignment takes two values and the line element takes half-integers. The CP transport multiplies by $e^{i\pi\Theta}$, and on the materialized locus the winding is an even integer, so that phase is exactly one while its half-turn $i^\Theta$ is exactly $(-1)^k$. That is the whole mechanism. The exchange sign is the winding's parity; the spin label is half the winding; so the exchange sign is $(-1)^{2j}$, plus one for integer spin and minus one for half-integer spin, and the connection is a theorem with a two-line induction certificate behind it. It does not contradict the earlier refutation, and the tier is explicit about why: what was refuted was a spin-parity theorem on the bound-or-transit kind, and both parities do occur on one kind at the floor, because the statistics depends on the winding and not on the kind, nor on the registration datum. The one thing that is chosen rather than derived is which sector of the two-chain space the sign is read on, and the tier prices the choice: the symmetric and antisymmetric sectors give the same statistics up to a global sign, while the determinant line and the orientation block give different answers or none, so the read is stated with its locus attached. What is not used is any of the three classical routes — no Lorentz covariance, no positivity of energy, no microcausality — and deficit eleven closes on the phase alphabet and the additivity of the winding, with the residue the previous tier named now discharged at two chains and the three-chain space left open where it was.

\end{theorem}

\begin{definition}[{\texorpdfstring{Def-BM.17 — The local resonance and its two threshold functions}{Def-BM.17 - The local resonance and its two threshold functions}}]\label{tag:Def-BM-17}
\textbf{Def.} A resonance $:=$ an eigenvalue of a multi-chain ensemble together with its two derived local functions. A construction on $\mathrm{Def\text{-}BM.14}$'s six-state assembly, $\mathrm{BM.13}$'s width ladder, $\mathrm{BM.34}$'s ADM split, $\mathrm{BM.44}$'s per-tick weight and $\mathrm{BM.88}$'s coupling profile; no new axiom, no new datum, no new parameter, no cosmological coordinate.

%dsp-level:0
\[
(1)\quad\textbf{the resonance tuple:}\qquad\boxed{\ \mathcal R\ :=\ \bigl(n,\ \sigma,\ Q,\ \mathbf s,\ k\bigr)\ }
\]

%dsp-level:0
\[
n\in\{1,2,3\}\ \ \text{the chain count}\ [\mathrm{CT3.17a}(1):\ \text{a sealed support partitions into}\ 2\text{'s and}\ 3\text{'s}];
\]

%dsp-level:0
\[
\qquad \sigma\in\{\delta,\varepsilon\}\ \ \text{the junction grammar}\ [\mathrm{D3.9};\ \mathrm{Def\text{-}BM.1}];
\]

%dsp-level:0
\[
Q\in\{-n,\dots,+n\}\ \ \text{the charge class}\ [\mathrm{Def\text{-}BM.8}(1);\ \mathrm{BM.103}(5)];
\]

%dsp-level:0
\[
\qquad \mathbf s=(s_1,\dots,s_n)\in\mathcal S^n\ \ \text{the constituent states}\ [\mathrm{Def\text{-}BM.14}(2):\ |\mathcal S|=6];
\]

%dsp-level:0
\[
k\in\mathbb Z_{\ge0}\ \ \text{the winding}\ [\mathrm{FT.22};\ \mathrm{Def\text{-}H.38}];
\]

%dsp-level:0
\[
\qquad \textbf{constitutive clause:}\ \ \text{the tuple is derived structural data}\ [\mathrm{Def\text{-}BM.2}\ \text{pattern}];
\]

%dsp-level:0
\[
\qquad \textbf{no entry is a species name}\ [\mathrm{Def\text{-}BM.3}(7)];
\]

%dsp-level:0
\[
(2)\quad\textbf{the internal count and the informational mass:}
\]

%dsp-level:0
\[
\qquad E_{\mathrm{int}}(\mathcal R)=\sum_{a=1}^{n}E_{\mathrm{int}}(s_a),
\]

%dsp-level:0
\[
\qquad E_{\mathrm{int}}(s)\in\{2,\ 6,\ 12,\ 24,\ 60\}\ \ \text{on the computed cells}
\]

%dsp-level:0
\[
\qquad [\mathrm{D0.8};\ \mathrm{BM.13}(2);\ \mathrm{FT.28};\ \text{Census FT.35}(2),\ \text{FT.42}(5)\ \text{row-wise}]:
\]

%dsp-level:0
\[
m_{\mathrm{hid}}(\mathcal R)=\frac{E_{\mathrm{int}}(\mathcal R)\ln2}{\mathfrak h}
\]

%dsp-level:0
\[
\qquad [\mathrm{FT.44}(4)];
\]

%dsp-level:0
\[
\qquad m^2_{\mathrm{pole}}(\mathcal R)=\hat\mu'^2\,\lambda^{(n)}(\mathbf s),
\]

%dsp-level:0
\[
\qquad \lambda^{(n)}(\mathbf s):=\sum_{a=1}^{n}D_S(s_a)
\]

%dsp-level:0
\[
\qquad [\mathrm{Def\text{-}BM.14}(3);\ \mathrm{BM.111}(2)];
\]

%dsp-level:0
\[
(3)\quad\textbf{the materialization threshold:}
\]

%dsp-level:0
\[
\qquad \textbf{constitutive clause:}\ \ \text{a resonance materializes when the locally available rest energy reaches its own scale},
\]

%dsp-level:0
\[
E_{\mathrm{rest}}=N\,m_{\mathrm{hid}}(\mathcal R)=\frac{N\,E_{\mathrm{int}}(\mathcal R)\ln2}{\mathfrak h}
\]

%dsp-level:0
\[
\qquad [\mathrm{BM.34}(4),(7):\ E=N\sqrt{m^2+p^2}-N^ip_i\ \text{at}\ p=0,\ N^i=0;\ \mathrm{BM.88}(2)(\mathrm{iii}):\ \text{the lapse}]:
\]

%dsp-level:0
\[
\boxed{\ \textbf{the materialization condition:}\qquad N\,E\ \ge\ \mathcal T_0(\mathcal R)\ :=\ \frac{E_{\mathrm{int}}(\mathcal R)\,\mathfrak h}{\ln2}\ }
\]

%dsp-level:0
\[
\qquad [E\ \text{the local energy scale};\ \text{machine: one Laurent monomial}];
\]

%dsp-level:0
\[
\mathcal T(\mathcal R;\,N)\ :=\ \frac{\mathcal T_0(\mathcal R)}{N}\ =\ \frac{E_{\mathrm{int}}(\mathcal R)\,\mathfrak h}{N\,\ln2}
\]

%dsp-level:0
\[
\qquad [\text{the threshold as a function of the lapse alone};\ \text{machine}];
\]

%dsp-level:0
\[
\textbf{the curvature read (exact, via the lock):}
\]

%dsp-level:0
\[
\qquad G=2\pi E_{\mathrm{int}}^2\ln2\implies E_{\mathrm{int}}=\sqrt{\frac{G}{2\pi\ln2}}\implies\boxed{\ \mathcal T(\mathcal R;\,G,N)\ =\ \frac{\mathfrak h}{N\,\ln2}\sqrt{\frac{G}{2\pi\ln2}}\ }
\]

%dsp-level:0
\[
\qquad [\mathrm{BM.16}(1);\ \mathrm{BM.88}(3);\ \text{machine}];
\]

%dsp-level:0
\[
\textbf{the horizon read (comparison scale, quarantined):}
\]

%dsp-level:0
\[
\qquad r_s=\frac{\mathfrak h E_{\mathrm{int}}}2=2G_Nm_{\mathrm{hid}}\implies\boxed{\ \mathcal T(\mathcal R;\,r_s,N)\ =\ \frac{2\,r_s}{N}\ }
\]

%dsp-level:0
\[
\qquad [\mathrm{BM.15}(4);\ \mathrm{CT2.9};\ \mathrm{P3}:\ G_N\ \text{is a named combination for external comparison};\ \neg\text{a horizon identification}\ [\mathrm{BM.15}\ \text{comparison discipline}]];
\]

%dsp-level:0
\[
(4)\quad\textbf{the decay width:}
\]

%dsp-level:0
\[
\qquad \textbf{constitutive clause:}\ \ \text{the width is a weight per tick,}\ [\Gamma]
\]

%dsp-level:0
\[
= \mathfrak q^0,
\]

%dsp-level:0
\[
\ \neg\text{a rate per external second}\ [\mathrm{BM.44}(4),(5);\ \mathrm{D1.8}]:
\]

%dsp-level:0
\[
\boxed{\ \Gamma(\mathcal R;\,E,N,G,\varepsilon_{\mathrm{EM}})\ :=\ \Gamma_{\mathrm{tr}}(\varepsilon_q)\cdot f\!\Bigl(\frac{r}{\Delta x_{\mathrm B}(E)}\Bigr)\cdot\frac{\varepsilon_{\mathrm{eff}}}{\varepsilon_{\mathrm{EM}}}\cdot\mathcal N(\mathcal R)\cdot\chi_{\{\,NE\ \ge\ \mathcal T(\mathcal R;N)\,\}}\ }
\]

%dsp-level:0
\[
\Gamma_{\mathrm{tr}}=\frac{x_+x_-}{1-x_+x_-}\in\Bigl[0,\tfrac13\Bigr]
\]

%dsp-level:0
\[
\qquad [\mathrm{BM.44}(2);\ \text{frozen}];
\]

%dsp-level:0
\[
\qquad f=1-\frac1k,\quad k=\frac{r}{\Delta x_{\mathrm B}(E)}\ge1
\]

%dsp-level:0
\[
\qquad [\mathrm{BM.88}(2)(\mathrm{ii})];
\]

%dsp-level:0
\[
\varepsilon_{\mathrm{eff}}=Nf\,\mathsf g(G)\,\varepsilon_{\mathrm{EM}}\implies\frac{\varepsilon_{\mathrm{eff}}}{\varepsilon_{\mathrm{EM}}}=N\,f\,\mathsf g(G)
\]

%dsp-level:0
\[
\qquad [\mathrm{BM.88}(4);\ \mathrm{BM.101}(2)];
\]

%dsp-level:0
\[
\qquad \mathsf g\in\{1,\ 1+G,\ \tfrac1{1+G}\}\ \ \text{priced,}\ \neg\text{adopted}
\]

%dsp-level:0
\[
\qquad [\mathrm{BM.88}(4);\ \mathrm{BM.104}(5)];
\]

%dsp-level:0
\[
\mathcal N(\mathcal R)\ :=\ \#\bigl\{\mathcal R'\ \text{in the same}\ (n,\sigma,Q)\ \text{class with}\ \lambda^{(n)}(\mathcal R')<\lambda^{(n)}(\mathcal R)\bigr\}
\]

%dsp-level:0
\[
\qquad [\text{the lighter-resonance channel count};\ \mathrm{BM.41}(3)\ \text{pattern}];
\]

%dsp-level:0
\[
\chi_{\{\cdot\}}\in\{0,1\}
\]

%dsp-level:0
\[
\qquad [\text{the exact indicator}];
\]

%dsp-level:0
\[
\qquad \textbf{non-derived clause:}\ \ \text{the product form}\ —\ \text{three derived factors and one count, multiplied};
\]

%dsp-level:0
\[
(5)\quad\textbf{the derived rows of the width (recorded at the definition site):}
\]

%dsp-level:0
\[
\qquad \Gamma_{\mathrm{tr}}(0)=0\ \ \text{exactly}\ [\mathrm{BM.23}(8);\ \mathrm{BM.44}(2)]\implies\boxed{\ \Gamma(\mathcal R;\,\varepsilon_q{=}0)=0\ \ \forall\mathcal R,\ \forall(E,N,G,\varepsilon_{\mathrm{EM}})\ }
\]

%dsp-level:0
\[
f(1)=0\ [\mathrm{BM.88}(5)]\implies\boxed{\ \Gamma(\mathcal R;\,r=\Delta x_{\mathrm B})=0\ \ \text{at the member's own envelope}\ }
\]

%dsp-level:0
\[
\mathcal N(\mathcal R)
\]

%dsp-level:0
\[
= 0\ \ \text{for the class's lightest resonance}\implies\boxed{\ \textbf{the ground state of every}\ (n,\sigma,Q)\ \textbf{class has width exactly zero}\ }
\]

%dsp-level:0
\[
\Gamma\le\tfrac13\cdot1\cdot N\,\mathsf g\cdot\mathcal N
\]

%dsp-level:0
\[
\qquad [\text{machine:}\ f<1,\ \varepsilon_{\mathrm{eff}}/\varepsilon_{\mathrm{EM}}=Nf\mathsf g<N\mathsf g];
\]

%dsp-level:0
\[
\qquad \tau_{\mathrm{life}}:=\frac1\Gamma\ge3\ \ \text{ticks at}\ \mathcal N=1,\ f\to1,\ N\mathsf g\to1
\]

%dsp-level:0
\[
\qquad [\mathrm{BM.44}(3);\ \text{machine}];
\]

%dsp-level:0
\[
(6)\quad\textbf{the standing orders, as scope clauses:}
\]

%dsp-level:0
\[
\qquad (\mathrm i)\ \ \neg\textbf{a cosmological coordinate exists:}\ \text{the ledger knows the local tick}\ \tau\ [\mathrm{A6};\ \mathrm{D1.8}],
\]

%dsp-level:0
\[
\ \text{the local curvature}\ G\ [\mathrm{BM.88}(3)],
\]

%dsp-level:0
\[
\ \text{the local lapse}\ N\ [\mathrm{BM.88}(2)(\mathrm{iii})]\ \text{and the local energy scale}\ E\ [\mathrm{BM.88}(2)(\mathrm i)];
\]

%dsp-level:0
\[
\ \neg\text{an age and}\ \neg\text{a global time};
\]

%dsp-level:0
\[
(\mathrm{ii})\quad\neg\textbf{a binary stable/unstable category:}
\]

%dsp-level:0
\[
\qquad \textbf{every eigenvalue of every multi-chain ensemble is a resonance};
\]

%dsp-level:0
\[
\qquad \text{``stable''}\ :=\ \Gamma\ \text{evaluates to}\ 0\ \text{at the ambient calibration point}\ [(5)];
\]

%dsp-level:0
\[
\qquad \text{``unstable'' or ``heavy''}\ :=\ \text{the local parameters cross}\ \mathcal T\ \text{or open}\ \mathcal N\ [(3),(4)];
\]

%dsp-level:0
\[
(\mathrm{iii})\quad\neg\textbf{a permanently forbidden state:}
\]

%dsp-level:0
\[
\qquad \mathcal T(\mathcal R;N)\in(0,\infty)\ \ \forall N\in(0,\infty)\implies\textbf{the threshold locus is never empty}\ [\mathrm{BM.116}(4)];
\]

%dsp-level:0
\[
(7)\quad\textbf{status and scope:}
\]

%dsp-level:0
\[
\qquad \textbf{derived:}\ (2),\ (3)\text{'s three forms},\ (5)\text{'s five rows};
\]

%dsp-level:0
\[
\qquad \textbf{constitutive:}\ (1)\text{'s tuple},\ (3)\text{'s and}\ (4)\text{'s definitional forms};
\]

%dsp-level:0
\[
\qquad \textbf{non-derived:}\ (4)\text{'s product form};
\]

%dsp-level:0
\[
\textbf{the scopes inherited at every use:}\ \ \mathrm{CT2.6}\text{'s continuum-limit hypotheses for the chart-side reads}\ [\mathrm{BM.34}\ \mathrm{Rmk};\ \mathrm{Def\text{-}BM.4}(5)\ \text{envelope discipline}];
\]

%dsp-level:0
\[
\qquad \mathrm{BM.8}(2)\text{'s stationarity hypothesis for the chart evolution}\ [\mathrm{BM.44}(3)];
\]

%dsp-level:0
\[
\neg\text{a P4 addition},
\]

%dsp-level:0
\[
\quad\neg\text{a new dimensionful scale}\ [\mathrm{P1}],
\]

%dsp-level:0
\[
\quad\neg\text{a measured value}\ [\mathrm{P5}],
\]

%dsp-level:0
\[
\quad\neg\text{a numerical coupling}\ [\mathrm{EM.25}(7)],
\]

%dsp-level:0
\[
\quad\neg\text{a lifetime in external seconds}\ [\mathrm{BM.44}(5)];
\]

%dsp-level:0
\[
\textbf{beyond this tier:}\ \ \mathsf g(G)\ [\mathrm{BM.104}(5)],
\]

%dsp-level:0
\[
\quad e_0\ [\mathrm{EM.25}(7)],
\]

%dsp-level:0
\[
\quad N\ \text{as solution data}\ [\mathrm{CT2.7};\ \mathrm{BM.34}(2)],\quad\rho_{\mathrm B}\ [\mathrm{A\text{-}BM.9}],
\]

%dsp-level:0
\[
\quad (\mathrm{R3})\ \text{the flavour composition},
\]

%dsp-level:0
\[
\quad (\mathrm{R8})\ \text{the}\ 2\text{-chain sector masses}.
\]

\textbf{Rmk (why "resonance" and not "particle").} Clause (6)(ii) is the tier's standing order written as a scope clause, and it is a statement about the mathematics rather than a naming preference: the ensemble's eigenvalues are indexed by the tuple of (1), each carries a threshold function and a width function, and both are exact symbolic functions of the four local parameters. Nothing in the construction distinguishes an eigenvalue whose width happens to vanish at one calibration point from one whose width happens not to; the difference is a value of a function, not a category. The word \textit{particle} would import a binary the ledger does not carry, and the word \textit{resonance} imports nothing: it names an eigenvalue with two functions attached.

\textbf{Rmk (the horizon row's quarantine).} Clause (3)'s third form is written with $r_s$, which BM.15(4) builds from $G_N$ — a named combination of the single scale for external comparison, not an input [P3]. The row is a change of variable inside the corpus's own monomial and it is quarantined exactly as BM.15's comparison discipline requires: no collapse reading, no horizon identification, and no claim that any resonance sits at any measured radius.

\textbf{Req.} A6, D0.8, D1.8, D2.9 (Parts 1–2); CT2.6, CT2.7, CT2.9, EM.25(7) (Part 2); P1, P3, P4, P5 (Parts 1–2); Def-H.20, Def-H.38, FT.22, FT.28, FT.44(4) (Part 3, Census); Census FT.35(2), FT.42(5) (row-wise); Def-BM.1, Def-BM.2, Def-BM.3(7), Def-BM.4(1), Def-BM.4(5), Def-BM.8(1), Def-BM.14(2), Def-BM.14(3), BM.8(2), BM.13(2), BM.15(4), BM.16(1), BM.23(8), BM.34(2), BM.34(4), BM.34(7), BM.41(3), BM.44(2), BM.44(3), BM.44(4), BM.44(5), BM.88(2), BM.88(3), BM.88(4), BM.88(5), BM.101(2), BM.103(5), BM.104(5), BM.111(2), A-BM.9 (this part); exact Laurent-monomial arithmetic with $\mathfrak q$-exponent censuses, exact rational function rows [math; machine: \texttt{checks/\allowbreak{}verify\_\allowbreak{}H\_\allowbreak{}antimatter\_\allowbreak{}resonance\_\allowbreak{}spectrum.py} §E].

\textbf{Pf.}

%dsp-level:0
\[
\text{Definition block: the constitutive clauses are}\ (1)\text{'s tuple and the definitional forms of}\ (3),
\]

%dsp-level:0
\[
(4);
\]

%dsp-level:0
\[
\ \text{the non-derived clause is}\ (4)\text{'s product form};
\]

%dsp-level:0
\[
\ (2),
\]

%dsp-level:0
\[
(3)\text{'s three exact forms and}\ (5)\text{'s five rows are derived and machine-checked};
\]

%dsp-level:0
\[
\ (6),
\]

%dsp-level:0
\[
(7)\ \text{are scope records}.\ \blacksquare
\]

\textbf{Phys.} A resonance is a five-entry tuple — how many chains, which seal, which charge class, which six states, which winding — and it carries two functions of the local parameters. The first says when it appears: the rest energy the metric's own ADM split gives is the lapse times the mass, so the resonance materializes when the local scale times the lapse reaches its edge count times the single scale over the log of two, and the same condition read through the calibration lock is a square root of the local curvature, or read through the comparison radius it is twice that radius over the lapse. The second says how fast it goes: the frozen per-tick transition weight, times the envelope fraction that vanishes at the member's own width, times the effective coupling over the bare one, times the count of strictly lighter resonances in its class, times the exact indicator that it is materialized at all. Both functions vanish for whole regions rather than for special objects — at the sharp crystallization threshold nothing decays, at the envelope nothing decays, and the lightest member of every class has nothing lighter to go to, so its width is zero by count. That is the entire content of "stable": a function evaluating to zero where it is being evaluated. Two scope clauses carry the tier's standing orders: no cosmological coordinate exists, because the ledger owns a tick, a curvature, a lapse and a scale and nothing global, and no binary category exists, because every eigenvalue is a resonance and the adjectives are values. The third says nothing is forbidden, since every threshold is a finite positive number at every positive lapse, so every locus is non-empty.

\end{definition}

\begin{theorem}[{\texorpdfstring{BM.114 — The resonance spectrum: the two characteristic polynomials, the orbit law, and the one exact collision}{BM.114 - The resonance spectrum: the two characteristic polynomials, the orbit law, and the one exact collision}}]\label{tag:BM-114}
\textbf{Thm.}

%dsp-level:0
\[
(1)\quad\textbf{the transport-register characteristic polynomial at every chain count:}
\]

%dsp-level:0
\[
\qquad W_n(Q)=\bigl[z^Q\bigr]\bigl(2+z+z^{-1}\bigr)^{n},
\]

%dsp-level:0
\[
\qquad 2+z+z^{-1}=z^{-1}(z+1)^2
\]

%dsp-level:0
\[
\qquad [\mathrm{Def\text{-}BM.13}(3);\ \mathrm{BM.103}(3)]:
\]

%dsp-level:0
\[
\boxed{\ \bigl(2+z+z^{-1}\bigr)^{n}\ =\ \frac{(z+1)^{2n}}{z^{n}}\ \ \text{identically in}\ \mathbb Z[z,z^{-1}]\implies W_n(Q)=\binom{2n}{Q+n}\ \ \text{exactly, on all}\ 2n+1\ \text{classes}\ }
\]

%dsp-level:0
\[
[\text{machine: the Laurent identity and the exhaustive}\ 4^n\ \text{end-sign census at}\ n=1,2,3];
\]

%dsp-level:0
\[
\qquad \sum_{Q=-n}^{n}W_n(Q)=2^{2n}=4^{n}
\]

%dsp-level:0
\[
\qquad [\text{machine}];
\]

%dsp-level:0
\[
\textbf{the single root:}
\]

%dsp-level:0
\[
\qquad z
\]

%dsp-level:0
\[
= -1\ \ \text{of multiplicity}\ 2n\ \ \text{(the}\ 2n\ \text{directed edges), and the}\ 2n+1\ \text{class weights are its binomial coefficients}
\]

%dsp-level:0
\[
\qquad [\mathrm{BM.103}(3)\ \text{transported to every}\ n;\ \text{machine}];
\]

%dsp-level:0
\[
\textbf{the three formed levels:}\qquad\begin{array}{c|c|c|c}n & \text{seal} & W_n(Q) & \text{sum}\\\hline1 & \neg\text{sealed (a lone attachment is never sealed)} & (1,\ 2,\ 1) & 4=4^1\\2 & \delta\ \text{(meson)} & (1,\ 4,\ \mathbf6,\ 4,\ 1) & 16=4^2\\3 & \varepsilon\ \text{(baryon)} & (1,\ 6,\ 15,\ \mathbf{20},\ 15,\ 6,\ 1) & 64=4^3\end{array}\qquad[\text{machine: exact binomials}];
\]

%dsp-level:0
\[
\boxed{\ \textbf{the}\ n{=}1\ \textbf{row is formed but}\ \neg\textbf{sealed:}\quad g(1)=0\implies\mathcal C=0\implies\text{the colour block vanishes}\ }
\]

%dsp-level:0
\[
\qquad [\mathrm{CT3.17a}(1),(2);\ \mathrm{FT.19};\ \text{machine}];
\]

%dsp-level:0
\[
(2)\quad\textbf{the}\ \varepsilon\textbf{-seal's antisymmetry and the}\ \delta\textbf{-seal's symmetry (the only two antisymmetries the corpus owns):}
\]

%dsp-level:0
\[
\qquad B=\varepsilon_{ijk}q^iq^jq^k\ \ \text{is antisymmetric in its three colour slots};
\]

%dsp-level:0
\[
\qquad M=\delta^i{}_jc^j\bar c_i\ \ \text{is symmetric in its two}
\]

%dsp-level:0
\[
\qquad [\mathrm{D3.9};\ \mathrm{CT3.15};\ \mathrm{CT3.17};\ \mathrm{BM.106}(5)]:
\]

%dsp-level:0
\[
\textbf{colour is}\ \neg\textbf{observable by construction}
\]

%dsp-level:0
\[
\qquad [\mathrm{CT3.22}]\implies\boxed{\ \textbf{the colour slot symmetry is}\ \neg\textbf{the statistics}\ —\ \text{the statistics is}\ \varepsilon(c)\ \text{on the winding carrier}\ [\mathrm{BM.113}(2)]\ }
\]

%dsp-level:0
\[
\textbf{the transport assembly is sign-free at every}\ n:
\]

%dsp-level:0
\[
\qquad W_n(Q)=\binom{2n}{Q+n}>0\ \ \text{on all}\ 2n+1\ \text{classes}
\]

%dsp-level:0
\[
\qquad [\mathrm{BM.103}(3);\ \text{machine}];
\]

%dsp-level:0
\[
(3)\quad\textbf{the defect-register eigenvalues (the resonance ladder):}
\]

%dsp-level:0
\[
\qquad \lambda^{(n)}(\mathbf s)=\sum_{a=1}^{n}D_S(s_a),
\]

%dsp-level:0
\[
\qquad D_S\in\bigl\{2,\ 4,\ 2-\sqrt2,\ 1,\ 2-2\cos\tfrac\pi{10},\ 2-\sqrt3\bigr\}
\]

%dsp-level:0
\[
\qquad [\mathrm{Def\text{-}BM.14}(3);\ \mathrm{Def\text{-}BM.17}(2)]:
\]

%dsp-level:0
\[
m^2_{\mathrm{pole}}(\mathcal R)=\hat\mu'^2\,\lambda^{(n)}(\mathbf s)\qquad[\mathrm{Def\text{-}H.37}(2);\ \mathrm{FT.54}(4);\ \text{machine}];
\]

%dsp-level:0
\[
\textbf{the exact eigenvalue field:}
\]

%dsp-level:0
\[
\qquad \mathfrak u:=2\cos\tfrac\pi{10},
\]

%dsp-level:0
\[
\qquad \mathfrak u^4-5\mathfrak u^2+5=0,
\]

%dsp-level:0
\[
\qquad \sqrt5=2\mathfrak u^2-5\in\mathbb Q(\mathfrak u),
\]

%dsp-level:0
\[
\qquad 2-2\cos\tfrac\pi{10}=2-\mathfrak u:
\]

%dsp-level:0
\[
\boxed{\ \mathbb M\ :=\ \mathbb Q\bigl(\sqrt2,\ \sqrt3,\ \mathfrak u\bigr),\qquad[\mathbb M:\mathbb Q]=2\cdot2\cdot4=16,\qquad\textbf{basis}\ \bigl\{\sqrt2^{\,i}\sqrt3^{\,j}\,\mathfrak u^{\,\ell}\bigr\}_{i,j\in\{0,1\},\ \ell\in\{0,1,2,3\}}\ }
\]

%dsp-level:0
\[
\qquad [\text{machine: the three defining relations verified exactly, and}\ (\sqrt2)^2{=}2,(\sqrt3)^2{=}3,(2\mathfrak u^2{-}5)^2{=}5];
\]

%dsp-level:0
\[
\textbf{every resonance eigenvalue lies in}\ \mathbb M,
\]

%dsp-level:0
\[
\ \text{and equality of two eigenvalues is decidable by basis-coordinate equality}
\]

%dsp-level:0
\[
\qquad [\text{machine:}\ \mathbb M\text{'s basis is linearly independent over}\ \mathbb Q];
\]

%dsp-level:0
\[
(4)\quad\textbf{the orbit law (the multiplicity is the stabilizer index):}
\]

%dsp-level:0
\[
\qquad S_n\ \text{acts on}\ \mathcal S^n\ \text{by permuting the constituent slots;}\ \ \lambda^{(n)}\ \textbf{is}\ S_n\textbf{-invariant}\ [\text{machine: a sum}]:
\]

%dsp-level:0
\[
\#\bigl(\text{orbits}\bigr)=\binom{6+n-1}{n},
\]

%dsp-level:0
\[
\qquad |\text{orbit of}\ \mathbf s|=\frac{n!}{|\operatorname{Stab}(\mathbf s)|},
\]

%dsp-level:0
\[
\qquad \sum_{\text{orbits}}|\text{orbit}|=6^{n}
\]

%dsp-level:0
\[
\qquad [\text{machine: exact at}\ n=1,2,3];
\]

%dsp-level:0
\[
\textbf{the three levels:}\qquad\begin{array}{c|c|c|c|c}n & \#\text{orbits} & \text{orbit sizes} & \sum & \text{multiplicity census}\\\hline1 & \binom61=6 & 1 & 6 & \{1^6\}\\2 & \binom72=21 & 1\ (\text{6 diagonal}),\ 2\ (\text{15 off-diagonal}) & 36 & \{1^6,\ 2^{15}\}\\3 & \binom83=56 & 1\ (\text{6}),\ 3\ (\text{30}),\ 6\ (\text{20}) & 216 & \text{see}\ (5)\end{array}
\]

%dsp-level:0
\[
[\text{machine: exhaustive over the}\ 6^n\ \text{ordered tuples}];\qquad\boxed{\ 15=\binom62=\dim\wedge^2\mathbb C^6,\qquad20=\binom63\ }
\]

%dsp-level:0
\[
\textbf{the count is a multiset count,}\ \neg\textbf{a representation dimension:}
\]

%dsp-level:0
\[
\qquad \neg\exists\ \text{a flavour}\ \mathsf{SU}(6)\ \text{in the inventory}\ [\mathrm{BM.103}(7);\ \text{the flavour symmetry is}\ C_3\ [\mathrm{FT.111}(2);\ \mathrm{FT.142}(5)]];
\]

%dsp-level:0
\[
(5)\quad\textbf{the three-chain ladder in full:}
\]

%dsp-level:0
\[
\qquad 56\ \text{orbits}\ \longrightarrow\ \boxed{\ 55\ \textbf{distinct eigenvalues}\ }
\]

%dsp-level:0
\[
\qquad [\text{machine: exact equality in}\ \mathbb M]:
\]

%dsp-level:0
\[
\textbf{exactly one collision, and it is the powers-of-two triple:}
\]

%dsp-level:0
\[
\qquad \lambda^{(3)}(s_1,s_1,s_1)=2+2+2=6=4+1+1=\lambda^{(3)}(s_2,s_4,s_4),
\]

%dsp-level:0
\[
\qquad \text{multiplicity}\ 1+3=4
\]

%dsp-level:0
\[
[\text{machine: the exact rational identity};\ \mathrm{BM.82}(4)\text{'s collision census re-cited at the constituent level: three of the six defect values are the powers}\ 1,2,4];
\]

%dsp-level:0
\[
\textbf{the multiplicity census of the 55:}
\]

%dsp-level:0
\[
\qquad \boxed{\ \{1^{5},\ 3^{29},\ 4^{1},\ 6^{20}\}\ }
\]

%dsp-level:0
\[
\qquad \text{sum}\ =\ 5+87+4+120=216=6^3
\]

%dsp-level:0
\[
\qquad [\text{machine: exact}];
\]

%dsp-level:0
\[
\textbf{the multiplicity is always a divisor of}\ n!=6:
\]

%dsp-level:0
\[
\qquad \{1,3,4,6\}\subset\{1,2,3,6\}\cup\{4\},
\]

%dsp-level:0
\[
\qquad \text{and the}\ 4\ \text{is the collision's}\ 1+3\ \ [\text{machine}];
\]

%dsp-level:0
\[
(6)\quad\textbf{the minimal-polynomial degree census (the subfield ladder of}\ \mathbb M\text{):}\qquad\begin{array}{c|c|c}n & \text{degree census over the distinct eigenvalues} & \max\\\hline1 & \{1^{3},\ 2^{2},\ 4^{1}\} & 4\\2 & \{1^{6},\ 2^{8},\ 4^{5},\ 8^{2}\} & 8\\3 & \{1^{9},\ 2^{20},\ 4^{15},\ 8^{10},\ 16^{1}\} & 16=[\mathbb M:\mathbb Q]\end{array}
\]

%dsp-level:0
\[
[\text{machine: each minimal polynomial computed exactly by linear dependence in the}\ 16\text{-dimensional}\ \mathbb Q\text{-space}\ \mathbb M,\ \text{then verified by exact evaluation at its own root}];
\]

%dsp-level:0
\[
\textbf{the pinned rows:}
\]

%dsp-level:0
\[
\qquad D_S(s_3)=2-\sqrt2:\ x^2-4x+2;
\]

%dsp-level:0
\[
\qquad D_S(s_5)=2-\mathfrak u:\ x^4-8x^3+19x^2-12x+1;
\]

%dsp-level:0
\[
\qquad D_S(s_6)=2-\sqrt3:\ x^2-4x+1
\]

%dsp-level:0
\[
\qquad [\mathrm{Def\text{-}BM.14}(3);\ \text{machine re-verified}];
\]

%dsp-level:0
\[
\textbf{the unique degree-16 resonance eigenvalue:}
\]

%dsp-level:0
\[
\qquad \lambda^{(3)}(s_3,s_5,s_6)=(2-\sqrt2)+(2-\mathfrak u)+(2-\sqrt3)=6-\sqrt2-\mathfrak u-\sqrt3:
\]

%dsp-level:0
\[
\qquad \textbf{it generates}\ \mathbb M\ \text{over}\ \mathbb Q
\]

%dsp-level:0
\[
[\text{machine: its minimal polynomial has degree exactly}\ 16;\ \text{the leading rows are}\ x^{16}-96x^{15}+4260x^{14}-115920x^{13}+\cdots+249713695921];
\]

%dsp-level:0
\[
\textbf{the minimal eigenvalue at each level is a quartic:}
\]

%dsp-level:0
\[
\qquad n\,\bigl(2-\mathfrak u\bigr):
\]

%dsp-level:0
\[
\qquad n=1:\ x^4-8x^3+19x^2-12x+1;
\]

%dsp-level:0
\[
\qquad n=2:\ x^4-16x^3+76x^2-96x+16;
\]

%dsp-level:0
\[
\qquad n=3:\ x^4-24x^3+171x^2-324x+81
\]

%dsp-level:0
\[
[\text{machine: the three quartics computed exactly and each verified at its own root}];
\]

%dsp-level:0
\[
(7)\quad\textbf{the ladder is a resonance ladder, not a stability ladder:}
\]

%dsp-level:0
\[
\qquad \neg\exists\ \text{a derived eigenvalue stability condition}\ [\mathrm{BM.103}(6):\ \text{the three grounds re-cited}\ —\ (\mathrm{M1})/(\mathrm{M2})\ \text{undeclared},\ \mathrm{FT.66}(4)\text{'s index-zero refutation, and the count}\to\text{strength bar}]:
\]

%dsp-level:0
\[
\boxed{\ \textbf{what orders the ladder is the eigenvalue's own value, and what decides its appearance is}\ \mathcal T(\mathcal R;N),\ \neg\textbf{a stability class}\ }
\]

%dsp-level:0
\[
\qquad [\mathrm{Def\text{-}BM.17}(3),(6)(\mathrm{ii});\ \mathrm{BM.115};\ \mathrm{BM.116}];
\]

%dsp-level:0
\[
(8)\quad\textbf{status:}\qquad\textbf{derived:}\ (1)\text{--}(7);
\]

%dsp-level:0
\[
\neg\text{a declaration},
\]

%dsp-level:0
\[
\quad\neg\text{a P4 addition},
\]

%dsp-level:0
\[
\quad\neg\text{a new dimensionful scale}\ [\mathrm{P1}],
\]

%dsp-level:0
\[
\quad\neg\text{a measured value}\ [\mathrm{P5}],
\]

%dsp-level:0
\[
\quad\neg\text{a species identification}\ [\mathrm{Def\text{-}BM.3}(7)],
\]

%dsp-level:0
\[
\quad\neg\text{a currency conversion}\ [\mathrm{BM.83}(3)];
\]

\textbf{Req.} D2.9, D3.9 (Parts 1–2); CT3.15, CT3.17, CT3.17a(1), CT3.17a(2), CT3.22 (Part 2); P1, P4, P5 (Parts 1–2); Def-H.37(2), FT.19, FT.54(4), FT.111(2), FT.142(5) (Part 3); Def-BM.3(7), Def-BM.13(3), Def-BM.14(3), Def-BM.17(2), BM.66(6), BM.82(4), BM.83(3), BM.103(3), BM.103(6), BM.103(7), BM.106(5), BM.111(2), BM.113(2) (this part); Laurent-polynomial identities over $\mathbb Z[z,z^{-1}]$, exhaustive binomial and multiset-orbit enumeration, exact linear algebra over $\mathbb Q$ in a $16$-dimensional algebra, exact real ordering by a rigorous dyadic interval embedding [math; machine: \texttt{checks/\allowbreak{}verify\_\allowbreak{}H\_\allowbreak{}antimatter\_\allowbreak{}resonance\_\allowbreak{}spectrum.py} §F].

\textbf{Pf.}

%dsp-level:0
\[
(1):\  2+z+z^{-1}=z^{-1}(z^2+2z+1)=z^{-1}(z+1)^2\ [\text{machine: one Laurent identity}];
\]

%dsp-level:0
\[
\qquad \bigl(z^{-1}(z+1)^2\bigr)^n=z^{-n}(z+1)^{2n}\ [\text{machine}]
\]

%dsp-level:0
\[
[z^Q]z^{-n}(z+1)^{2n}
\]

%dsp-level:0
\[
= [z^{Q+n}](z+1)^{2n}
\]

%dsp-level:0
\[
= \binom{2n}{Q+n}\ [\text{machine: against the exhaustive}\ 4^n\ \text{census at}\ n=1,2,3];
\]

%dsp-level:0
\[
\qquad \sum_Q\binom{2n}{Q+n}=2^{2n}\ [\text{machine}]
\]

%dsp-level:0
\[
g(1)=0\ [\mathrm{CT3.17a}(2)];\qquad|\mathrm{Colour}(\Sigma)|=1\implies\mathcal C=0\implies\text{the block vanishes}\ [\mathrm{CT3.17a}(1);\ \mathrm{FT.19}]
\]

%dsp-level:0
\[
(2):\  \mathrm{D3.9},\ \mathrm{CT3.15},\ \mathrm{CT3.17}\ \text{re-cited};
\]

%dsp-level:0
\[
\qquad \mathrm{CT3.22}\text{'s non-observability re-cited};
\]

%dsp-level:0
\[
\qquad W_n(Q)=\binom{2n}{Q+n}>0\ [\text{machine: positivity of binomials}]
\]

%dsp-level:0
\[
(3):\  \mathfrak u
\]

%dsp-level:0
\[
= 2\cos\tfrac\pi{10}:\quad\cos^2\tfrac\pi{10}
\]

%dsp-level:0
\[
= \tfrac{5+\sqrt5}8\implies\mathfrak u^2
\]

%dsp-level:0
\[
= \tfrac{5+\sqrt5}2\implies\mathfrak u^4
\]

%dsp-level:0
\[
= \tfrac{15+5\sqrt5}2\implies\mathfrak u^4-5\mathfrak u^2+5
\]

%dsp-level:0
\[
= \tfrac{15+5\sqrt5-25-5\sqrt5}2+5
\]

%dsp-level:0
\[
= 0\ [\text{machine: exact in}\ \mathbb Q(\sqrt5)]
\]

%dsp-level:0
\[
\sqrt5
\]

%dsp-level:0
\[
= 2\mathfrak u^2-5\ [\text{machine:}\ (2\mathfrak u^2-5)^2=4\cdot\tfrac{5+\sqrt5}2\cdot2-20\mathfrak u^2+25\ \text{reduced by}\ \mathfrak u^4=5\mathfrak u^2-5\ \text{to}\ 5];
\]

%dsp-level:0
\[
\qquad [\mathbb Q(\mathfrak u):\mathbb Q]=4\ [\mathfrak u^4-5\mathfrak u^2+5\ \text{irreducible:}\ \neg\text{a rational root and}\ \mathfrak u^2=\tfrac{5+\sqrt5}2\notin\mathbb Q]
\]

%dsp-level:0
\[
[\mathbb M:\mathbb Q]
\]

%dsp-level:0
\[
= [\mathbb M:\mathbb Q(\mathfrak u)]\cdot[\mathbb Q(\mathfrak u):\mathbb Q]
\]

%dsp-level:0
\[
= 4\cdot4
\]

%dsp-level:0
\[
= 16\ [\sqrt2,\sqrt3\notin\mathbb Q(\mathfrak u)\ \text{by the degree census of}\ (6);\ \text{machine: the}\ 16\ \text{basis elements are}\ \mathbb Q\text{-linearly independent}]
\]

%dsp-level:0
\[
(4):\  \lambda^{(n)}(\pi\cdot\mathbf s)=\sum_aD_S(s_{\pi(a)})=\sum_aD_S(s_a)=\lambda^{(n)}(\mathbf s)\ [\text{machine}];
\]

%dsp-level:0
\[
\qquad \#\text{orbits}=\binom{6+n-1}n\ [\text{multiset count};\ \text{machine: exhaustive at}\ n=1,2,3]
\]

%dsp-level:0
\[
|\text{orbit}|=n!/|\operatorname{Stab}|\ \text{with}\ |\operatorname{Stab}|=\prod_{\text{distinct values}}(\text{multiplicity})!\ [\text{machine}];
\]

%dsp-level:0
\[
\qquad \text{the size sums}\ 6,36,216\ [\text{machine}]
\]

%dsp-level:0
\[
\binom62=15=\dim\wedge^2\mathbb C^6=\tfrac{6\cdot5}2;\qquad\binom63=20\ [\text{machine}]
\]

%dsp-level:0
\[
(5):\  2+2+2=6\ \text{and}\ 4+1+1=6\ \text{in}\ \mathbb Q\subset\mathbb M\ [\text{machine: basis-coordinate equality}];
\]

%dsp-level:0
\[
\qquad \text{the exhaustive census finds no second collision among the}\ 56\ [\text{machine}]
\]

%dsp-level:0
\[
\text{the multiplicity census}\ \{1^5,3^{29},4^1,6^{20}\}:\quad5+3\cdot29+4+6\cdot20=5+87+4+120=216\ [\text{machine}]
\]

%dsp-level:0
\[
(6):\  \text{for each eigenvalue}\ x,
\]

%dsp-level:0
\[
\ \text{the smallest}\ d\ \text{with}\ \{1,x,\dots,x^d\}\ \mathbb Q\text{-dependent in the}\ 16\text{-dimensional}\ \mathbb M\ \text{gives the degree, and the dependency gives the monic polynomial}\ [\text{machine: exact Fraction elimination}];
\]

%dsp-level:0
\[
\ \text{each is verified by exact evaluation at its own root}\ [\text{machine}]
\]

%dsp-level:0
\[
x^2-4x+2\ \text{at}\ 2-\sqrt2:\ (2-\sqrt2)^2-4(2-\sqrt2)+2=(6-4\sqrt2)-(8-4\sqrt2)+2=0\ [\text{machine}];
\]

%dsp-level:0
\[
\qquad x^2-4x+1\ \text{at}\ 2-\sqrt3:\ (7-4\sqrt3)-(8-4\sqrt3)+1=0\ [\text{machine}]
\]

%dsp-level:0
\[
(2-\mathfrak u):\ (2-x)^4-5(2-x)^2+5=x^4-8x^3+19x^2-12x+1\ [\text{machine: the exact expansion}];
\]

%dsp-level:0
\[
\qquad \text{the three minimal quartics and the degree-16 polynomial computed and root-verified}\ [\text{machine}]
\]

%dsp-level:0
\[
(7):\  \mathrm{BM.103}(6)\text{'s three grounds re-cited};\qquad(8):\  \text{the standing-rule scans}\ [\text{machine}].\ \blacksquare
\]

\textbf{Phys.} Two characteristic polynomials come out, one per currency, and neither is fitted. In the transport register the generating function is the same Laurent monomial at every chain count — one plus z plus one over z, raised to n, which is one plus z to the two n over z to the n — so the class weight is a binomial coefficient, the sum is four to the n, and the whole spectrum is one root at minus one with multiplicity twice the chain count. The first level is formed but not sealed, since a lone attachment is never sealed, which is why the ladder starts at two. In the defect register the eigenvalues are sums of the six branch values, and they all live in one exact field of degree sixteen over the rationals built from root two, root three and twice the cosine of eighteen degrees, whose own relation makes root five a polynomial in it — so equality of two eigenvalues is a decidable integer question rather than a numerical one. The orbit law is clean: the symmetric group permutes the constituent slots, the eigenvalue is invariant, the orbit count is a multiset binomial, and the multiplicity of each eigenvalue is the orbit size, a divisor of n factorial. Fifteen and twenty appear, and they are exactly the wedge-square dimension and the three-subset count of the six states — but the tier says plainly that these are counts and not representation dimensions, because the flavour symmetry the corpus derives is a cyclic three and not a special unitary six. The three-chain ladder has fifty-six orbits and fifty-five eigenvalues, and the single collision is the one the constituent census had already found: two plus two plus two equals four plus one plus one, because three of the six defect values are the powers one, two and four. The minimal-polynomial degrees climb one, two, four, eight, sixteen with the number of distinct irrational ingredients, and exactly one three-chain resonance — the sum of the root-two, cosine and root-three branches — generates the whole field.

\end{theorem}

\begin{theorem}[{\texorpdfstring{BM.115 — The materialization threshold: the exact function of $E$, $N$ and $G$, and the never-empty locus}{BM.115 - The materialization threshold: the exact function of , and , and the never-empty locus}}]\label{tag:BM-115}
\textbf{Thm.}

%dsp-level:0
\[
(1)\quad\textbf{the threshold function:}
\]

%dsp-level:0
\[
\qquad \mathcal T_0(\mathcal R)=\frac{E_{\mathrm{int}}(\mathcal R)\,\mathfrak h}{\ln2},
\]

%dsp-level:0
\[
\qquad \mathcal T(\mathcal R;N)=\frac{\mathcal T_0(\mathcal R)}N,
\]

%dsp-level:0
\[
\qquad \textbf{materialization}\iff NE\ge\mathcal T_0(\mathcal R)
\]

%dsp-level:0
\[
\qquad [\mathrm{Def\text{-}BM.17}(3);\ \text{machine}]:
\]

%dsp-level:0
\[
[\mathcal T_0]=[\mathfrak h]=\mathfrak q^{-1/2}\cdot\mathfrak q^{+1/2}\cdot\mathfrak q^{0}=\ \ \text{an energy:}\quad[E]=\mathfrak q^{+1/2},
\]

%dsp-level:0
\[
\qquad [\mathcal T]=\mathfrak q^{+1/2}
\]

%dsp-level:0
\[
\qquad [\mathrm{CT2.1};\ \mathrm{BM.102}(2);\ \text{machine: the exponent arithmetic}];
\]

%dsp-level:0
\[
\mathcal T\ \textbf{is strictly decreasing in the lapse:}
\]

%dsp-level:0
\[
\qquad N_1<N_2\implies\mathcal T(\mathcal R;N_1)>\mathcal T(\mathcal R;N_2)>0
\]

%dsp-level:0
\[
\qquad [\text{machine: exact rationals}];
\]

%dsp-level:0
\[
\boxed{\ \textbf{a deeper lapse lowers every threshold at once}\ —\ \text{the whole ladder moves as one function of}\ N\ }
\]

%dsp-level:0
\[
(2)\quad\textbf{the two exact threshold ladders:}
\]

%dsp-level:0
\[
\qquad \textbf{(i) the computed-cell ladder}\quad E_{\mathrm{int}}(s)\in\{2,\ 6,\ 12,\ 24,\ 60\}
\]

%dsp-level:0
\[
\qquad [\mathrm{BM.13}(2);\ \mathrm{FT.28};\ \text{Census FT.35}(2),\ \text{FT.42}(5)\ \text{row-wise}]:
\]

%dsp-level:0
\[
\begin{array}{c|c|l|c}n & \#\{E_{\mathrm{int}}\} & \text{the exact ladder} & \text{span ratio}\\\hline1 & 5 & 2,\ 6,\ 12,\ 24,\ 60 & 30\\2 & 15 & 4,\ 8,\ 12,\ 14,\ 18,\ 24,\ 26,\ 30,\ 36,\ 48,\ 62,\ 66,\ 72,\ 84,\ 120 & 30\\3 & 33 & 6,\ 10,\ 14,\ 16,\ 18,\ 20,\ 24,\ 26,\ 28,\ 30,\ 32,\ 36,\ 38,\ 42,\ 48,\ 50,\ 54,\ 60,\ 64,\ 68,\ 72,\ 74,\ 78,\ 84,\ 86,\ 90,\ 96,\ 108,\ 122,\ 126,\ 132,\ 144,\ 180 & 30\end{array}
\]

%dsp-level:0
\[
[\text{machine: the exact sum sets with repetitions}];
\]

%dsp-level:0
\[
\qquad \boxed{\ \min E_{\mathrm{int}}^{(n)}=2n\ \ \text{(the}\ n\ \text{digon floors)},\qquad\max E_{\mathrm{int}}^{(n)}=60n,\qquad\frac{\max}{\min}=30\ \ \textbf{at every}\ n\ }
\]

%dsp-level:0
\[
\textbf{(ii) the six-state ladder}\quad E_{\mathrm{int}}(s)=2E(s)\in\{12,\ 24,\ 60\}\ \ \text{on the five Platonic solids}
\]

%dsp-level:0
\[
\qquad [\mathrm{Def\text{-}BM.14}(1);\ \mathrm{BM.60}(1):\ 12:24:60=1:2:5]:
\]

%dsp-level:0
\[
\begin{array}{c|c|l|c}n & \#\{E_{\mathrm{int}}\} & \text{the exact ladder} & \text{span ratio}\\\hline1 & 3 & 12,\ 24,\ 60 & 5\\2 & 6 & 24,\ 36,\ 48,\ 72,\ 84,\ 120 & 5\\3 & 10 & 36,\ 48,\ 60,\ 72,\ 84,\ 96,\ 108,\ 132,\ 144,\ 180 & 5\end{array}
\]

%dsp-level:0
\[
[\text{machine: the exact sum sets}];
\]

%dsp-level:0
\[
\qquad \boxed{\ \min=12n,\qquad\max=60n,\qquad\frac{\max}{\min}=5=\frac{60}{12}\ \ \textbf{at every}\ n\ }
\]

%dsp-level:0
\[
\qquad [\text{the floor ladder's own ratio}\ 1:2:5\ \text{re-cited}];
\]

%dsp-level:0
\[
\textbf{the two ladders are nested, and the nesting is exact:}
\]

%dsp-level:0
\[
\qquad \{12,24,60\}\subset\{2,6,12,24,60\}\implies\text{every six-state threshold is a computed-cell threshold};
\]

%dsp-level:0
\[
\textbf{the two extra cells}\ 2\ \text{and}\ 6\ \text{are the digon floor}\ [\mathrm{FT.28}]\ \text{and the baryon junction}\ [\text{Census FT.35}(2)]:
\]

%dsp-level:0
\[
\qquad \text{they lower the minimum from}\ 12n\ \text{to}\ 2n\ \text{and raise the span from}\ 5\ \text{to}\ 30\ [\text{machine}];
\]

%dsp-level:0
\[
\textbf{the meson and baryon rows on the computed-cell ladder:}
\]

%dsp-level:0
\[
\qquad m_{\mathrm{hid}}^{(2)}\in\Bigl[\frac{4\ln2}{\mathfrak h},\ \frac{120\ln2}{\mathfrak h}\Bigr],
\]

%dsp-level:0
\[
\qquad m_{\mathrm{hid}}^{(3)}\in\Bigl[\frac{6\ln2}{\mathfrak h},\ \frac{180\ln2}{\mathfrak h}\Bigr]
\]

%dsp-level:0
\[
\qquad [\mathrm{BM.111}(5);\ \text{machine}];
\]

%dsp-level:0
\[
\textbf{and on the six-state ladder:}
\]

%dsp-level:0
\[
\qquad m_{\mathrm{hid}}^{(2)}\in\Bigl[\frac{24\ln2}{\mathfrak h},\ \frac{120\ln2}{\mathfrak h}\Bigr],
\]

%dsp-level:0
\[
\qquad m_{\mathrm{hid}}^{(3)}\in\Bigl[\frac{36\ln2}{\mathfrak h},\ \frac{180\ln2}{\mathfrak h}\Bigr]
\]

%dsp-level:0
\[
\qquad [\text{machine}];
\]

%dsp-level:0
\[
(3)\quad\textbf{the curvature read:}
\]

%dsp-level:0
\[
\qquad G=2\pi E_{\mathrm{int}}^2\ln2\in\mathfrak q^0,
\]

%dsp-level:0
\[
\qquad G\ \ \mathfrak h\text{-free},\ \Lambda\text{-free},\ \mathrm{P4}\text{-free}
\]

%dsp-level:0
\[
\qquad [\mathrm{BM.16}(1);\ \mathrm{BM.88}(3);\ \text{machine: the three freedom censuses re-cited}]:
\]

%dsp-level:0
\[
\boxed{\ E_{\mathrm{int}}=\sqrt{\frac{G}{2\pi\ln2}}\implies\mathcal T(\mathcal R;\,G,N)=\frac{\mathfrak h}{N\ln2}\sqrt{\frac{G}{2\pi\ln2}}\ }
\]

%dsp-level:0
\[
\qquad [\text{machine: the exact substitution}];
\]

%dsp-level:0
\[
G^{(n)}\in\bigl[2\pi\ln2\cdot4n^2,\ \ 2\pi\ln2\cdot3600n^2\bigr],
\]

%dsp-level:0
\[
\qquad \frac{G_{\max}}{G_{\min}}=900=30^2\ \ \textbf{at every}\ n
\]

%dsp-level:0
\[
\qquad [\text{machine: exact integers}];
\]

%dsp-level:0
\[
\textbf{the curvature is a derived function of the count, so the threshold is a derived function of the local curvature}\ —\ \neg\text{an independent datum}\ [\mathrm{BM.88}(3);\ \text{machine}];
\]

%dsp-level:0
\[
(4)\quad\textbf{the horizon read (comparison scale, quarantined):}
\]

%dsp-level:0
\[
\qquad r_s=2G_Nm_{\mathrm{hid}}=\tfrac12\mathfrak h E_{\mathrm{int}},
\]

%dsp-level:0
\[
\qquad \Delta x_{\mathrm B}\cdot r_s=2G_N=\frac{\mathfrak h^2}{2\ln2}
\]

%dsp-level:0
\[
\qquad [\mathrm{BM.15}(4);\ \mathrm{CT2.9};\ \mathrm{P3}]:
\]

%dsp-level:0
\[
\boxed{\ \mathcal T(\mathcal R;\,r_s,N)=\frac{2\,r_s}{N}\ }
\]

%dsp-level:0
\[
\qquad [\text{machine:}\ \tfrac{\mathfrak h}{N\ln2}\cdot\tfrac{2r_s}{\mathfrak h}=\tfrac{2r_s}N];
\]

%dsp-level:0
\[
\qquad r_s^{(n)}\in[n\mathfrak h,\ 30n\mathfrak h],
\]

%dsp-level:0
\[
\qquad \frac{r_{s,\max}}{r_{s,\min}}=30
\]

%dsp-level:0
\[
\qquad [\text{machine}];
\]

%dsp-level:0
\[
\textbf{the comparison discipline:}\quad\neg\text{a collapse reading},
\]

%dsp-level:0
\[
\quad\neg\text{a horizon identification},
\]

%dsp-level:0
\[
\quad\neg\text{a measured radius}\ [\mathrm{BM.15}(4)\ \text{comparison discipline};\ \mathrm{P5};\ \mathrm{Def\text{-}BM.3}(6)];
\]

%dsp-level:0
\[
\textbf{what the row is:}\quad\text{a change of variable inside the corpus's own monomial, which exhibits the threshold as}\ \textbf{twice the comparison radius over the lapse};
\]

%dsp-level:0
\[
(5)\quad\textbf{the never-empty locus (no forbidden state):}
\]

%dsp-level:0
\[
\qquad \forall\mathcal R,\ \forall N\in(0,\infty):
\]

%dsp-level:0
\[
\qquad \mathcal T(\mathcal R;N)\in(0,\infty)
\]

%dsp-level:0
\[
\qquad [E_{\mathrm{int}}\in\mathbb Z_{\ge2},\ \mathfrak h>0,\ \ln2>0\ [\mathrm{BM.13}(3):\ \tfrac12<\ln2<1]]:
\]

%dsp-level:0
\[
\{E>0:\ NE\ge\mathcal T_0(\mathcal R)\}=\Bigl[\frac{\mathcal T_0(\mathcal R)}N,\ \infty\Bigr)\ne\varnothing
\]

%dsp-level:0
\[
\qquad [\text{machine: the half-line is nonempty at every admissible}\ N];
\]

%dsp-level:0
\[
\boxed{\ \neg\exists\ \textbf{a permanently forbidden resonance: every eigenvalue's threshold locus is a nonempty half-line in the local energy scale}\ }
\]

%dsp-level:0
\[
\textbf{the two named triples are neither formed nor excluded:}
\]

%dsp-level:0
\[
\qquad \neg\exists\ \text{a triple-selection rule in the corpus}\ [\mathrm{BM.103}(8):\ (\mathrm{R3})]:
\]

%dsp-level:0
\[
\textbf{but both are admissible threshold-wise:}\quad\text{any}\ \mathbf s\in\mathcal S^3\ \text{has}\ E_{\mathrm{int}}(\mathbf s)
\]

%dsp-level:0
\[
\in [6,180]\subset\mathbb Z_{>0}\ \text{and hence a nonempty locus}\ [(5)\text{'s first row};\ \text{machine: all}\ 216\ \text{ordered triples}];
\]

%dsp-level:0
\[
\boxed{\ \textbf{there are no forbidden compositions; there are compositions whose thresholds require a higher local scale}\ }
\]

%dsp-level:0
\[
(6)\quad\textbf{the accelerator and horizon regimes are evaluations, not mechanisms:}
\]

%dsp-level:0
\[
\qquad \textbf{a high local scale}\ E\gg E_1:\quad NE\ge\mathcal T_0(\mathcal R)\ \text{holds for every}\ \mathcal R\ \text{with}\ E_{\mathrm{int}}(\mathcal R)\le NE\ln2/\mathfrak h
\]

%dsp-level:0
\[
\qquad [\text{machine}]:
\]

%dsp-level:0
\[
\textbf{a small lapse}\ N\ll1\ \text{(the deep-gravity read):}\quad\mathcal T(\mathcal R;N)
\]

%dsp-level:0
\[
= \mathcal T_0/N\to\infty\implies\textbf{thresholds rise and fewer resonances materialize}
\]

%dsp-level:0
\[
\qquad [\text{machine: monotonicity from}\ (1)];
\]

%dsp-level:0
\[
\textbf{a large comparison radius}\ r_s:\quad\mathcal T=2r_s/N\ \text{grows linearly}\qquad[(4);\ \text{machine}];
\]

%dsp-level:0
\[
\boxed{\ \textbf{both regimes are the same evaluation of the same function at different local parameters}\ —\ \neg\text{two mechanisms,}\ \neg\text{a new object}\ }
\]

%dsp-level:0
\[
\neg\text{a cosmological coordinate is used:}\ \text{the rows are stated in}\ (E,\ N,\ G,\ r_s)\ \text{and}\ \neg\text{in any global time}\ [\mathrm{Def\text{-}BM.17}(6)(\mathrm i);\ \mathrm{A6};\ \mathrm{D1.8}];
\]

%dsp-level:0
\[
(7)\quad\textbf{the crossing strength re-cited and}\ \neg\textbf{re-evaluated:}
\]

%dsp-level:0
\[
\qquad \varepsilon_\ast(w)=\frac{\ln(4/3)}{3w\ln2}=\frac{2-\log_23}{3w},
\]

%dsp-level:0
\[
\qquad \tfrac18<\varepsilon_\ast<\tfrac5{12}
\]

%dsp-level:0
\[
\qquad [\mathrm{BM.101}(3);\ \text{machine: the enclosure re-certified on}\ 2^{13}>3^8,\ 3^{12}>2^{19}]:
\]

%dsp-level:0
\[
\textbf{its role here:}\quad\text{the crossing locus}\ \mathcal L_\ast\ \text{is where the mass splitting changes sign, and it is a level set of the same profile}\ \varepsilon_{\mathrm{eff}}\ \text{that the width function carries}\ [\mathrm{BM.101}(4);\ \mathrm{Def\text{-}BM.17}(4)]:
\]

%dsp-level:0
\[
\neg\text{a numerical coupling enters},
\]

%dsp-level:0
\[
\qquad \neg\text{a measured value is compared}
\]

%dsp-level:0
\[
\qquad [\mathrm{P5};\ \mathrm{EM.25}(7);\ \text{machine: the token scan}];
\]

%dsp-level:0
\[
(8)\quad\textbf{status:}\qquad\textbf{derived:}\ (1)\text{--}(6);\qquad\textbf{re-cited,}\ \neg\textbf{re-evaluated:}\ (7);
\]

%dsp-level:0
\[
\neg\text{a declaration},
\]

%dsp-level:0
\[
\quad\neg\text{a P4 addition},
\]

%dsp-level:0
\[
\quad\neg\text{a new dimensionful scale}\ [\mathrm{P1}],
\]

%dsp-level:0
\[
\quad\neg\text{a measured value}\ [\mathrm{P5}],
\]

%dsp-level:0
\[
\quad\neg\text{a species identification}\ [\mathrm{Def\text{-}BM.3}(7)];
\]

\textbf{Req.} D0.8, D1.8, D2.9 (Parts 1–2); A6, CT2.1, CT2.9, EM.25(7) (Part 1, Part 2); P1, P3, P4, P5 (Parts 1–2); FT.28, FT.44(4) (Part 3, Census); Census FT.35(2), FT.42(5) (row-wise); Def-BM.3(6), Def-BM.3(7), Def-BM.4(1), Def-BM.17(3), Def-BM.17(4), BM.13(2), BM.13(3), BM.15(4), BM.16(1), BM.88(3), BM.101(3), BM.101(4), BM.102(2), BM.103(8), BM.111(5) (this part); exact integer sum-set enumeration, exact rational monotonicity, Laurent-monomial exponent censuses, integer power certificates [math; machine: \texttt{checks/\allowbreak{}verify\_\allowbreak{}H\_\allowbreak{}antimatter\_\allowbreak{}resonance\_\allowbreak{}spectrum.py} §E].

\textbf{Pf.}

%dsp-level:0
\[
(1):\  \mathcal T_0=E_{\mathrm{int}}\mathfrak h/\ln2\ [\mathrm{Def\text{-}BM.17}(3);\ \text{machine}];
\]

%dsp-level:0
\[
\qquad [\mathcal T_0]=[\mathfrak h]=\mathfrak q^{-1/2}\ \text{and}\ [E_{\mathrm{int}}]=[\ln2]=\mathfrak q^0\implies[\mathcal T_0]=\mathfrak q^{-1/2}\cdot\mathfrak q^{0}
\]

%dsp-level:0
\[
\text{with}\ E_{\mathrm{int}}=\mathfrak h^{-1}E\ \text{dimensionless:}
\]

%dsp-level:0
\[
\quad [\mathcal T_0]=[\mathfrak h]\cdot[E]/[\mathfrak h]=\mathfrak q^{+1/2}\ [\mathrm{CT2.1};\ \mathrm{BM.102}(2);\ \text{machine: the exponent arithmetic}];
\]

%dsp-level:0
\[
\qquad N_1<N_2\implies\mathcal T_0/N_1>\mathcal T_0/N_2\ [\text{machine: exact rationals}]
\]

%dsp-level:0
\[
(2):\  \text{the three sum sets computed exhaustively from}\ \{2,6,12,24,60\}\ \text{with repetitions,}\ n
\]

%dsp-level:0
\[
= 1,
\]

%dsp-level:0
\[
2,\qquad 3:
\]

%dsp-level:0
\[
\quad 5,\ 15,\ 33\ \text{distinct values};
\]

%dsp-level:0
\[
\qquad \min=2n,\ \max=60n,\ 60n/2n=30\ [\text{machine: exact integers}]
\]

%dsp-level:0
\[
(3):\  G=2\pi E_{\mathrm{int}}^2\ln2\implies E_{\mathrm{int}}=\sqrt{G/(2\pi\ln2)}\ [\text{machine: the positive branch at}\ E_{\mathrm{int}}>0];
\]

%dsp-level:0
\[
\qquad \text{substitution into}\ (1)\ [\text{machine}]
\]

%dsp-level:0
\[
G_{\min}=2\pi\ln2\cdot(2n)^2=4n^2\cdot2\pi\ln2,
\]

%dsp-level:0
\[
\quad G_{\max}=2\pi\ln2\cdot(60n)^2=3600n^2\cdot2\pi\ln2,
\]

%dsp-level:0
\[
\quad 3600/4=900=30^2\ [\text{machine}]
\]

%dsp-level:0
\[
(4):\  r_s
\]

%dsp-level:0
\[
= \tfrac12\mathfrak h E_{\mathrm{int}}
\]

%dsp-level:0
\[
\implies E_{\mathrm{int}}
\]

%dsp-level:0
\[
= 2r_s/\mathfrak h\implies\mathcal T_0
\]

%dsp-level:0
\[
= \tfrac{\mathfrak h}{\ln2}\cdot\tfrac{2r_s}{\mathfrak h}
\]

%dsp-level:0
\[
= \tfrac{2r_s}{\ln2}\cdot\ln2\ [\text{machine}];
\]

%dsp-level:0
\[
\qquad \mathcal T=\mathcal T_0/N\ \text{with}\ \mathcal T_0=2r_s\ \text{in energy units}\ [\text{machine}]
\]

%dsp-level:0
\[
r_s\in[\tfrac12\mathfrak h\cdot2n,\ \tfrac12\mathfrak h\cdot60n]=[n\mathfrak h,\ 30n\mathfrak h]\ [\text{machine}]
\]

%dsp-level:0
\[
(5):\  E_{\mathrm{int}}\ge2\ [\mathrm{FT.28}]\implies\mathcal T_0\ge2\mathfrak h/\ln2>0\ [\tfrac12<\ln2<1\ [\mathrm{BM.13}(3)];\ \text{machine}];
\]

%dsp-level:0
\[
\qquad \text{the half-line}\ [\mathcal T_0/N,\infty)\ \text{is nonempty}\ \forall N>0\ [\text{machine}]
\]

%dsp-level:0
\[
\text{all}\ 216\ \text{ordered triples}: E_{\mathrm{int}}\in[6,180]\subset\mathbb Z_{>0}\ [\text{machine: exhaustive}];
\]

%dsp-level:0
\[
\qquad \neg\text{a selection rule exists}\ [\mathrm{BM.103}(8)]
\]

%dsp-level:0
\[
(6):\  NE\ge\mathcal T_0\iff E_{\mathrm{int}}\le NE\ln2/\mathfrak h\ [\text{machine}];
\]

%dsp-level:0
\[
\qquad N\downarrow\implies\mathcal T\uparrow\ [(1)];
\]

%dsp-level:0
\[
\qquad (7):\  \mathrm{BM.101}(3),
\]

%dsp-level:0
\[
(4)\ \text{re-cited, the enclosure re-certified on the two integer powers}\ [\text{machine}];
\]

%dsp-level:0
\[
\qquad (8):\  \text{the standing-rule scans}\ [\text{machine}].\ \blacksquare
\]

\textbf{Phys.} The threshold is one Laurent monomial and it has three exact readings. Read in the local scale it is the resonance's edge count times the single free scale over the log of two, divided by the lapse, so a deeper lapse lowers every threshold in the ladder at once and the whole spectrum moves as one function. Read in the local curvature it is a square root, because the calibration lock makes the curvature a derived quadratic function of the same count — and the lock is free of the scale, free of the cosmological term and free of every declared parameter, so the curvature reading adds no datum. Read in the comparison radius it is twice the radius over the lapse, and that row is quarantined exactly as the corpus quarantines every $G_N$-built quantity: no collapse reading, no horizon identification, no measured radius. The ladders themselves are exact integer sets — five values for one chain, fifteen for a meson, thirty-three for a baryon — running from twice the chain count to sixty times it, so the span ratio is exactly thirty at every level and the curvature span is exactly nine hundred. Then the row that answers the question about forbidden states: every threshold is a finite positive number at every positive lapse, so every locus is a nonempty half-line, and all two hundred sixteen ordered constituent triples sit inside it. Nothing is excluded; what differs is where in the local scale the half-line starts. Both regimes the question names — a collider and a deep-gravity region — are the same evaluation of the same function at different local parameters, one raising the scale and one lowering the lapse, and neither needs a mechanism, a cosmological clock, or a new object.

\end{theorem}

\begin{theorem}[{\texorpdfstring{BM.116 — The decay width: the exact function of the local parameters, its three exact zeros, and the lifetime floor}{BM.116 - The decay width: the exact function of the local parameters, its three exact zeros, and the lifetime floor}}]\label{tag:BM-116}
\textbf{Thm.}

%dsp-level:0
\[
(1)\quad\textbf{the width function:}
\]

%dsp-level:0
\[
\qquad \boxed{\ \Gamma(\mathcal R;\,E,N,G,\varepsilon_{\mathrm{EM}})\ =\ \underbrace{\frac{x_+x_-}{1-x_+x_-}}_{\Gamma_{\mathrm{tr}}(\varepsilon_q)}\ \cdot\ \underbrace{\Bigl(1-\frac{\Delta x_{\mathrm B}(E)}{r}\Bigr)}_{f}\ \cdot\ \underbrace{N\,\mathsf g(G)}_{\varepsilon_{\mathrm{eff}}/\varepsilon_{\mathrm{EM}}\,\cdot\,f^{-1}\cdot f}\ \cdot\ \underbrace{\mathcal N(\mathcal R)}_{\text{the channel count}}\ \cdot\ \underbrace{\chi_{\{NE\ge\mathcal T_0(\mathcal R)\}}}_{\text{materialization}}\ }
\]

%dsp-level:0
\[
\textbf{assembled from four frozen rows and one count:}
\]

%dsp-level:0
\[
\qquad \Gamma_{\mathrm{tr}}=\tfrac{x_+x_-}{1-x_+x_-}\in[0,\tfrac13]\ [\mathrm{BM.44}(2)];
\]

%dsp-level:0
\[
\qquad f=1-\tfrac1k,\ k=\tfrac{r}{\Delta x_{\mathrm B}(E)}\ge1\ [\mathrm{BM.88}(2)(\mathrm{ii})];
\]

%dsp-level:0
\[
\frac{\varepsilon_{\mathrm{eff}}}{\varepsilon_{\mathrm{EM}}}=Nf\,\mathsf g(G)\ [\mathrm{BM.88}(4);\ \mathrm{BM.101}(2)];
\]

%dsp-level:0
\[
\qquad \mathcal N(\mathcal R)=\#\{\mathcal R'\ \text{in the same}\ (n,\sigma,Q)\ \text{class}:\ \lambda^{(n)}(\mathcal R')<\lambda^{(n)}(\mathcal R)\}\ [\mathrm{Def\text{-}BM.17}(4)];
\]

%dsp-level:0
\[
[\Gamma]=\mathfrak q^0,\qquad[\tau_{\mathrm{life}}]=\mathfrak q^0\ \ \text{ticks}\qquad[\mathrm{BM.44}(4);\ \mathrm{D1.8};\ \mathrm{BM.102}(2);\ \text{machine: the exponent census}];
\]

%dsp-level:0
\[
(2)\quad\textbf{the flicker factor and its exact range:}
\]

%dsp-level:0
\[
\qquad x_+x_-=4\varepsilon_q(1-4\varepsilon_q)=\frac14-16\Bigl(\varepsilon_q-\frac18\Bigr)^{\!2}\le\frac14,
\]

%dsp-level:0
\[
\qquad \text{equality}\iff\varepsilon_q=\frac18
\]

%dsp-level:0
\[
\qquad [\mathrm{BM.19}(2);\ \mathrm{BM.44}(1);\ \text{machine: the SOS completion and the}\ 256\text{-point grid}]:
\]

%dsp-level:0
\[
\boxed{\ \Gamma_{\mathrm{tr}}(0)=0,\qquad\Gamma_{\mathrm{tr}}\bigl(\tfrac18\bigr)=\frac{1/4}{3/4}=\frac13,\qquad\Gamma_{\mathrm{tr}}\in\Bigl[0,\tfrac13\Bigr]\ \ \text{on}\ \ \varepsilon_q\in\Bigl[0,\tfrac18\Bigr],\qquad\Gamma_{\mathrm{tr}}\ \text{strictly increasing there}\ }
\]

%dsp-level:0
\[
[\mathrm{BM.44}(2);\ \text{machine:}\ \Gamma_{\mathrm{tr}}(\varepsilon_2)-\Gamma_{\mathrm{tr}}(\varepsilon_1)>0\ \text{via}\ 4(\varepsilon_2-\varepsilon_1)(1-4(\varepsilon_1+\varepsilon_2))>0];
\]

%dsp-level:0
\[
(3)\quad\textbf{the three exact zeros (the derived ``stable'' loci):}
\]

%dsp-level:0
\[
\qquad \boxed{\ (\mathrm Z1)\quad\varepsilon_q=0\implies\Gamma\equiv0\ \ \text{for every resonance at every local parameter}\ }
\]

%dsp-level:0
\[
\qquad [\mathrm{BM.23}(8);\ \mathrm{BM.19}(3);\ \text{machine}]:
\]

%dsp-level:0
\[
\textbf{the mechanism:}
\]

%dsp-level:0
\[
\quad x_-
\]

%dsp-level:0
\[
= 4\varepsilon_q
\]

%dsp-level:0
\[
= 0\implies\neg\exists\ \text{a dissolution move}\implies\text{the ensemble collapses to its crystallized member}\implies\neg\text{a transition weight beyond the identity};
\]

%dsp-level:0
\[
(\mathrm Z2)\quad r=\Delta x_{\mathrm B}(E)\implies k=1\implies f=0\implies\Gamma=0
\]

%dsp-level:0
\[
\qquad [\mathrm{BM.88}(5);\ \text{machine}]:
\]

%dsp-level:0
\[
\qquad \textbf{at the member's own Born envelope the width vanishes};
\]

%dsp-level:0
\[
(\mathrm Z3)\quad\mathcal N(\mathcal R)=0\implies\Gamma=0
\]

%dsp-level:0
\[
\qquad [\mathrm{Def\text{-}BM.17}(5);\ \text{machine}]:
\]

%dsp-level:0
\[
\qquad \boxed{\ \textbf{the lightest resonance of every}\ (n,\sigma,Q)\ \textbf{class has width exactly zero, at every local parameter}\ }
\]

%dsp-level:0
\[
\textbf{the count of zero-width resonances at the three formed levels:}\qquad\begin{array}{c|c|c}\text{level} & \#\{\mathcal N=0\} & \text{locus}\\\hline n=2,\ \sigma=\delta & 3 & \text{one per charge class}\ Q\in\{0,\pm1\}\ [\mathrm{BM.110}(3)]\\n=3,\ \sigma=\varepsilon & 7 & \text{one per charge class}\ Q\in\{-3,\dots,+3\}\ [\mathrm{BM.114}(1)]\end{array}\qquad[\text{machine: exact}];
\]

%dsp-level:0
\[
(4)\quad\textbf{the fourth zero is the threshold's, and it is not a width:}
\]

%dsp-level:0
\[
\qquad NE<\mathcal T_0(\mathcal R)\implies\chi=0\implies\Gamma=0
\]

%dsp-level:0
\[
\qquad [\text{machine}]:
\]

%dsp-level:0
\[
\boxed{\ \textbf{an unmaterialized resonance has no width because it has no ensemble}\ —\ \text{the zero is the indicator's,}\ \neg\text{a decay}\ }
\]

%dsp-level:0
\[
\qquad [\mathrm{Def\text{-}BM.17}(3);\ \mathrm{BM.94}(5):\ \text{the materialization restriction}];
\]

%dsp-level:0
\[
\textbf{the distinction the tier keeps:}\quad\Gamma=0\ \text{from}\ (\mathrm Z1),(Z2),(Z3)\ \textbf{is a stable resonance};
\]

%dsp-level:0
\[
\qquad \chi=0\ \textbf{is an absent one};
\]

%dsp-level:0
\[
\qquad \neg\text{a binary category covers both}\ [\mathrm{Def\text{-}BM.17}(6)(\mathrm{ii})];
\]

%dsp-level:0
\[
(5)\quad\textbf{the bound and the lifetime floor:}\qquad f<1\ \forall k<\infty\ [\mathrm{BM.88}(5);\ \text{machine}]\implies
\]

%dsp-level:0
\[
\boxed{\ 0\ \le\ \Gamma(\mathcal R)\ \le\ \frac13\cdot N\,\mathsf g(G)\cdot\mathcal N(\mathcal R)\ }\qquad[\text{machine: the three factor bounds composed}];
\]

%dsp-level:0
\[
\mathcal N=1,\ f\to1,\ N\mathsf g\to1:
\]

%dsp-level:0
\[
\qquad \Gamma\to\Gamma_{\mathrm{tr}}\le\tfrac13\implies\boxed{\ \tau_{\mathrm{life}}:=\frac1\Gamma\ \ge\ 3\ \ \textbf{ticks}\ }
\]

%dsp-level:0
\[
\qquad [\mathrm{BM.44}(3);\ \text{machine: exact reciprocal}];
\]

%dsp-level:0
\[
\varepsilon_q=0:\qquad\Gamma=0\implies\tau_{\mathrm{life}}\ \textbf{has no finite value}\qquad[\mathrm{BM.44}(3);\ \text{machine}];
\]

%dsp-level:0
\[
\neg\text{a lifetime in external seconds is derived and}\ \neg\text{a lifetime is compared}
\]

%dsp-level:0
\[
\qquad [\mathrm{BM.44}(5);\ \mathrm{P5};\ \text{machine: zero external unit tokens}];
\]

%dsp-level:0
\[
(6)\quad\textbf{the energy dependence (the width grows with the local scale):}
\]

%dsp-level:0
\[
\qquad \text{at fixed}\ r,
\]

%dsp-level:0
\[
\qquad k(E)=\frac{r\,E_{\mathrm{int}}(E)\ln2}{\mathfrak h},
\]

%dsp-level:0
\[
\qquad f=1-\frac{\mathfrak h}{r\,E_{\mathrm{int}}(E)\ln2}
\]

%dsp-level:0
\[
\qquad [\mathrm{BM.88}(6);\ \text{machine}]:
\]

%dsp-level:0
\[
\boxed{\ E\uparrow\implies k\uparrow\implies f\uparrow\implies\Gamma\uparrow\ }
\]

%dsp-level:0
\[
\qquad [\mathrm{BM.88}(6)\text{'s three monotonicity rows re-cited};\ \text{machine: positive coefficients}];
\]

%dsp-level:0
\[
\textbf{the calibration-scale reading:}
\]

%dsp-level:0
\[
\qquad \text{``low energy''}\ :=\ E\ \text{near}\ E_1=\mathfrak h/(r\ln2),\ \text{where}\ f(1)=0;
\]

%dsp-level:0
\[
\qquad \text{``high energy''}\ :=\ E\gg E_1,\ \text{where}\ f\to1
\]

%dsp-level:0
\[
\qquad [\mathrm{BM.101}(5);\ \mathrm{BM.88}(2)(\mathrm i);\ \text{machine}]:
\]

%dsp-level:0
\[
\boxed{\ \textbf{at the low calibration point every width vanishes through}\ f\ \textbf{and through}\ \chi\ \textbf{for the unmaterialized ones; at the high one the widths open in the order of}\ \mathcal N\ \textbf{and}\ \mathcal T_0\ }
\]

%dsp-level:0
\[
\textbf{and both are evaluations of one function,}\ \neg\textbf{two regimes of two theories}
\]

%dsp-level:0
\[
\qquad [\mathrm{Def\text{-}BM.17}(6)(\mathrm{ii});\ \text{machine}];
\]

%dsp-level:0
\[
(7)\quad\textbf{the two-sided bar and the width:}
\]

%dsp-level:0
\[
\qquad N\mathsf g\,\varepsilon_{\mathrm{EM}}\le\tfrac18\implies\Delta m>0;
\]

%dsp-level:0
\[
\qquad \ge\tfrac5{12}\implies\Delta m<0
\]

%dsp-level:0
\[
\qquad [\mathrm{BM.99}(3);\ \mathrm{BM.101}(3):\ \text{the same enclosure}]:
\]

%dsp-level:0
\[
\textbf{the consequence for the width:}\quad\frac{\varepsilon_{\mathrm{eff}}}{\varepsilon_{\mathrm{EM}}}
\]

%dsp-level:0
\[
= Nf\,
\]

%dsp-level:0
\[
\mathsf g
\]

%dsp-level:0
\[
\le N\mathsf g\implies\boxed{\ \text{the bar's band}\ \bigl[\tfrac18,\tfrac5{12}\bigr]\ \text{bounds the width's coupling factor exactly}\ }
\]

%dsp-level:0
\[
\textbf{and the sign of the mass splitting flips inside it}\ —\ \text{one function, two readings}\ [\mathrm{BM.101}(1),(2);\ \text{machine}];
\]

%dsp-level:0
\[
(8)\quad\textbf{status:}\qquad\textbf{derived:}\ (1)\text{--}(7);
\]

%dsp-level:0
\[
\textbf{non-derived and named:}\ \ (1)\text{'s product form}\ [\mathrm{Def\text{-}BM.17}(4)\ \text{non-derived clause}];
\]

%dsp-level:0
\[
\qquad \textbf{priced,}\ \neg\textbf{adopted:}\ \ \mathsf g\in\{1,1+G,\tfrac1{1+G}\}\ [\mathrm{BM.88}(4);\ \mathrm{BM.104}(5)];
\]

%dsp-level:0
\[
\neg\text{a declaration},
\]

%dsp-level:0
\[
\quad\neg\text{a P4 addition},
\]

%dsp-level:0
\[
\quad\neg\text{a new dimensionful scale}\ [\mathrm{P1}],
\]

%dsp-level:0
\[
\quad\neg\text{a measured value}\ [\mathrm{P5}],
\]

%dsp-level:0
\[
\quad\neg\text{a numerical coupling}\ [\mathrm{EM.25}(7)],
\]

%dsp-level:0
\[
\quad\neg\text{a species identification}\ [\mathrm{Def\text{-}BM.3}(7)];
\]

\textbf{Req.} D0.8, D1.8, D2.9 (Parts 1–2); CT0.8, EM.25(7) (Part 1, Part 2); P1, P4, P5 (Parts 1–2); Def-H.20, FT.44(4) (Part 3); Def-BM.3(7), Def-BM.4(1), Def-BM.17(3), Def-BM.17(4), Def-BM.17(5), Def-BM.17(6), BM.19(2), BM.19(3), BM.23(8), BM.44(1), BM.44(2), BM.44(3), BM.44(4), BM.44(5), BM.88(2), BM.88(4), BM.88(5), BM.88(6), BM.94(5), BM.99(3), BM.101(1), BM.101(2), BM.101(3), BM.101(5), BM.102(2), BM.104(5), BM.110(3), BM.114(1), BM.115(1) (this part); exact rational resolvent identities, the SOS parabola completion, exhaustive rational grid censuses, exact monotonicity by polynomial sign, Laurent-monomial exponent censuses [math; machine: \texttt{checks/\allowbreak{}verify\_\allowbreak{}H\_\allowbreak{}antimatter\_\allowbreak{}resonance\_\allowbreak{}spectrum.py} §E].

\textbf{Pf.}

%dsp-level:0
\[
(1):\  \text{the four frozen rows substituted into}\ \mathrm{Def\text{-}BM.17}(4)\text{'s product}\ [\text{machine: the monomial assembly}];
\]

%dsp-level:0
\[
\qquad \frac{\varepsilon_{\mathrm{eff}}}{\varepsilon_{\mathrm{EM}}}=\frac{Nf\,\mathsf g\,\varepsilon_{\mathrm{EM}}}{\varepsilon_{\mathrm{EM}}}=Nf\,\mathsf g\ [\text{machine}];
\]

%dsp-level:0
\[
\qquad [\Gamma]=\mathfrak q^0\ [\text{every factor dimensionless};\ \text{machine}]
\]

%dsp-level:0
\[
(2):\  \mathrm{BM.44}(1),
\]

%dsp-level:0
\[
(2)\text{'s rows re-cited}\ [\text{machine:}\ 4\varepsilon_q(1-4\varepsilon_q)=\tfrac14-16(\varepsilon_q-\tfrac18)^2\ \text{by expansion};\ \text{the}\ 256\text{-point grid census with equality only at}\ k=32]
\]

%dsp-level:0
\[
\Gamma_{\mathrm{tr}}(\tfrac18)=\tfrac{1/4}{1-1/4}=\tfrac{1/4}{3/4}=\tfrac13\ [\text{machine: exact rationals}];
\]

%dsp-level:0
\[
\qquad \text{monotonicity via}\ 4(\varepsilon_2-\varepsilon_1)(1-4(\varepsilon_1+\varepsilon_2))>0\ \text{on}\ \varepsilon_1<\varepsilon_2\le\tfrac18\ [\text{machine: the polynomial identity and a rational witness grid}]
\]

%dsp-level:0
\[
(3):\  (\mathrm Z1)\ x_-=4\cdot0=0\implies x_+x_-=0\implies\Gamma_{\mathrm{tr}}=0/(1-0)=0\ [\text{machine}];
\]

%dsp-level:0
\[
\qquad (\mathrm Z2)\ k=1\implies f=1-1=0\ [\mathrm{BM.88}(5);\ \text{machine}];
\]

%dsp-level:0
\[
\qquad (\mathrm Z3)\ \mathcal N=0\implies\text{the product is}\ 0\ [\text{machine}]
\]

%dsp-level:0
\[
\#\{\mathcal N=0\}:\ \text{at}\ n
\]

%dsp-level:0
\[
= 2\ \text{the three classes}\ Q\in\{0,\pm1\}\ \text{each have a lightest member}\ (Q=\pm2\ \text{has one member and is its own lightest, counted once each});
\]

%dsp-level:0
\[
\ \text{the exact count over the}\ \lambda^{(2)}\ \text{ladder of}\ \mathrm{BM.114}(3)\ \text{is}\ 3\ \text{on the sealed classes}\ [\text{machine}]
\]

%dsp-level:0
\[
\text{at}\ n=3\ \text{the seven classes}\ Q=-3,\dots,+3\ \text{each contribute one}\ [\text{machine}]
\]

%dsp-level:0
\[
(4):\  \chi_{\{NE<\mathcal T_0\}}=0\implies\Gamma=0\ [\text{machine}];
\]

%dsp-level:0
\[
\qquad \mathrm{BM.94}(5)\text{'s materialization restriction re-cited}:\ E\ge2\ \text{supports no closed chain below it}
\]

%dsp-level:0
\[
(5):\  f=1-\tfrac1k<1\iff\tfrac1k>0\iff k<\infty\ [\mathrm{BM.88}(5);\ \text{machine}];
\]

%dsp-level:0
\[
\qquad \Gamma_{\mathrm{tr}}\le\tfrac13\ [(2)];
\]

%dsp-level:0
\[
\qquad \text{the product bound}\ [\text{machine}]
\]

%dsp-level:0
\[
\mathcal N=1,\ f\to1,\ N\mathsf g\to1:\ \Gamma\to\tfrac13\implies\tau_{\mathrm{life}}\to3\ [\mathrm{BM.44}(3);\ \text{machine: the exact reciprocal}];
\]

%dsp-level:0
\[
\qquad \Gamma=0\implies1/\Gamma\ \text{has no finite value}\ [\text{machine}]
\]

%dsp-level:0
\[
(6):\  \mathrm{BM.88}(6)\text{'s three monotonicity rows re-cited};
\]

%dsp-level:0
\[
\qquad f\ \text{is a factor of}\ \Gamma\ \text{with all other factors}\
\]

%dsp-level:0
\[
\ge 0\implies\Gamma\ \text{increases with}\ f\ [\text{machine: positive-coefficient composition}]
\]

%dsp-level:0
\[
(7):\  \mathrm{BM.99}(3),
\]

%dsp-level:0
\[
\ \mathrm{BM.101}(3)\ \text{re-cited, the enclosure re-certified on}\ 2^{13}>3^8\ \text{and}\ 3^{12}>2^{19}\ [\text{machine}];
\]

%dsp-level:0
\[
\qquad f\le1\implies Nf\,\mathsf g\le N\mathsf g\ [\text{machine}];
\]

%dsp-level:0
\[
\qquad (8):\  \text{the standing-rule scans}\ [\text{machine}].\ \blacksquare
\]

\textbf{Phys.} The width is a product of four things the corpus had already frozen and one thing it can count. The flicker factor is the transition content of the Janus resolvent at unit fugacity, a rational function of the single crystallization threshold that climbs from zero to exactly a third at the threshold's midpoint and is strictly increasing all the way. The envelope factor is one minus one over k, which vanishes at the member's own Born width. The coupling factor is the effective coupling over the bare one, the lapse times the envelope fraction times the priced curvature composition. The count is how many strictly lighter resonances sit in the same chain-count, seal and charge class. And the indicator says whether the resonance is materialized at all. Three of those give exact zeros with distinct mechanisms: no flicker means no dissolution move and hence no transition weight; sitting at your own envelope means the enclosed flux fraction is zero; and being the lightest in your class means there is nothing to go to, so the count is zero and the width is zero by arithmetic. The fourth zero is different in kind and the tier keeps it apart — an unmaterialized resonance has no width because it has no ensemble, which is absence rather than stability. The bound follows by composing the factor bounds, and the lifetime floor is three ticks at the extreme, exactly as the vertex tier derived, with no second anywhere in it. Energy enters monotonically: raising the local scale raises the envelope ratio, raises the enclosed fraction, and raises every width, so "low energy" and "high energy" are two points on one function where the widths are respectively all zero and open in the order of their thresholds and their counts. The two-sided bar on the mass splitting is the same product seen from the other side, and its band is exactly the bound on the width's coupling factor.

\end{theorem}

\begin{theorem}[{\texorpdfstring{BM.117 — The groupings, the ceilings, and the octet adjudication}{BM.117 - The groupings, the ceilings, and the octet adjudication}}]\label{tag:BM-117}
\textbf{Thm.}

%dsp-level:0
\[
(1)\quad\textbf{the groupings that are derived:}\qquad\begin{array}{l|l|l}\text{grouping} & \text{the exact value} & \text{locus}\\\hline\text{the charge classes at}\ n & 2n+1\ \text{classes with the binomial weights}\ \binom{2n}{Q+n} & [\mathrm{BM.114}(1)]\\\text{the sealing grammar} & g(N)=\sum_{2a+3b=N}\frac{N!}{a!\,b!\,2^a6^b}=(0,1,1,3,10,25,105)\ \text{for}\ N=1..7 & [\mathrm{CT3.17a}(2);\ \mathrm{Def\text{-}BM.9}]\\\text{the multiset orbits at}\ n & \binom{6+n-1}n=(6,\ 21,\ 56) & [\mathrm{BM.114}(4)]\\\text{the all-distinct orbits} & \binom6n=(6,\ 15,\ 20)\ \ \text{with multiplicity}\ n! & [\mathrm{BM.114}(4);\ \text{machine}]\\\text{the generation floors} & 3,\ \text{ordered exactly by}\ 2E=12<24<60 & [\mathrm{Def\text{-}BM.14}(1),(6);\ \mathrm{BM.60}(1)]\\\text{the branches} & 2\ \text{per floor, ordered exactly by}\ D_S(\mathcal F_g,-)<D_S(\mathcal F_g,+) & [\mathrm{Def\text{-}BM.14}(4)]\\\text{the flavour symmetry} & C_3,\ \text{its orbits the three channel lines} & [\mathrm{FT.111}(2);\ \mathrm{FT.158}(1);\ \mathrm{FT.142}(5)]\\\text{the defect field} & \mathbb M,\ [\mathbb M:\mathbb Q]=16,\ \text{with the subfield ladder}\ 1,2,4,8,16 & [\mathrm{BM.114}(3),(6)]\end{array}
\]

%dsp-level:0
\[
(2)\quad\textbf{the grouping the spectrum actually carries:}
\]

%dsp-level:0
\[
\qquad \textbf{the three-level grouping of the}\ n\textbf{-chain ladder is}\ (\text{floor},\ \text{branch},\ \text{charge class})\ \text{with the orbit multiplicities of}\ \mathrm{BM.114}(4):
\]

%dsp-level:0
\[
\boxed{\ \textbf{at}\ n=3:\quad56\ \text{orbits}\ =\ 6\ \text{of size}\ 1\ +\ 30\ \text{of size}\ 3\ +\ 20\ \text{of size}\ 6,\qquad55\ \text{eigenvalues after the one collision}\ }
\]

%dsp-level:0
\[
\textbf{the}\ 20\ \textbf{all-distinct orbits are the}\ \varepsilon\textbf{-seal's fully antisymmetric slot pattern realized on distinct states:}
\]

%dsp-level:0
\[
\qquad \varepsilon_{ijk}\ \text{seals three fundamentals antisymmetrically}\ [\mathrm{D3.9};\ \mathrm{CT3.15}],
\]

%dsp-level:0
\[
\ \text{and}\ S_3\ \text{acts freely exactly on the all-distinct tuples}\ [\text{machine}]:
\]

%dsp-level:0
\[
\boxed{\ 20=\binom63\ \ \textbf{is derived as an orbit count and is}\ \neg\textbf{an}\ \mathsf{SU}(3)\ \textbf{decuplet dimension}\ }
\]

%dsp-level:0
\[
\qquad [\neg\exists\ \text{a flavour}\ \mathsf{SU}(3)\ \text{in the inventory}\ [\mathrm{BM.103}(7)];\ \text{machine: the symbol census}];
\]

%dsp-level:0
\[
(3)\quad\textbf{the octet adjudication (re-certified, and what replaces it):}
\]

%dsp-level:0
\[
\qquad \textbf{the derived}\ 3\ \textbf{is the colour seal's fundamental}\ [\mathrm{CT3.15};\ \mathrm{CT3.17};\ \mathrm{D3.9}];
\]

%dsp-level:0
\[
\qquad \textbf{the derived}\ 8\ \textbf{is a mask-capacity ratio, barred as a coupling}\ [\mathrm{BM.36}(2),(3);\ \mathrm{BM.61}(8)(\mathrm{O1})];
\]

%dsp-level:0
\[
\textbf{the derived}\ 10\ \textbf{is the grammar count}\ g(5)=10\ [\mathrm{CT3.17a}(2);\ \mathrm{Def\text{-}BM.9};\ \mathrm{BM.31}];
\]

%dsp-level:0
\[
\qquad \textbf{the derived}\ 15\ \textbf{and}\ 20\ \textbf{are the multiset orbit counts of}\ (1);
\]

%dsp-level:0
\[
\boxed{\ \textbf{the baryon octet and the meson octet are}\ \neg\textbf{derived, and the groupings that are derived are the binomial charge classes, the multiset orbits, the}\ C_3\ \textbf{triplet and the grammar counts}\ }
\]

%dsp-level:0
\[
[\mathrm{BM.103}(7)\ \text{re-cited and extended with the resonance census};\ \text{machine:}\ 8\notin\{2n+1,\ \binom{6+n-1}n,\ \binom6n,\ g(N)\}\ \text{at}\ n\le3,\ N\le7];
\]

%dsp-level:0
\[
(4)\quad\textbf{the ceilings:}
\]

%dsp-level:0
\[
\qquad \textbf{the per-constituent pole ceiling}\quad D_S(s)\le4\implies m^2_{\mathrm{pole}}(s)\le4\hat\mu'^2
\]

%dsp-level:0
\[
\qquad [\mathrm{FT.76}(3);\ \mathrm{Def\text{-}BM.14}(3):\ \max D_S=4\ \text{at}\ s_2;\ \text{machine}]:
\]

%dsp-level:0
\[
\textbf{the}\ n\textbf{-chain resonance ceiling is its}\ n\textbf{-fold sum:}
\]

%dsp-level:0
\[
\qquad \boxed{\ \lambda^{(n)}_{\max}=4n,\qquad m^2_{\mathrm{pole}}(\mathcal R)\le4n\,\hat\mu'^2,\qquad\text{attained exactly by}\ \mathbf s=(s_2,\dots,s_2)\ }
\]

%dsp-level:0
\[
[\text{machine:}\ 4n\ \text{is in the ladder at every}\ n=1,2,3,\ \text{with multiplicity}\ 1];
\]

%dsp-level:0
\[
\textbf{the locus discipline:}
\]

%dsp-level:0
\[
\qquad \mathrm{FT.76}(3)\text{'s ceiling is a property of a single chain's pole-side read and}\ \mathrm{FT.79}(2)\text{'s of a single chain's colour-side read}\ [\mathrm{BM.103}(6);\ \mathrm{FT.84}(2):\ \text{the two-ceiling synthesis}]:
\]

%dsp-level:0
\[
\boxed{\ \textbf{an}\ n\textbf{-chain composite eigenvalue is a sum of}\ n\ \textbf{single-chain reads and is}\ \neg\textbf{a single-chain read}\ —\ \neg\text{a ceiling is violated and}\ \neg\text{a ceiling is extended to the composite as such}\ }
\]

%dsp-level:0
\[
\textbf{the informational ceiling:}
\]

%dsp-level:0
\[
\qquad E_{\mathrm{int}}^{(n)}\in[12n,\ 60n]\ \text{on the six-state ladder}\implies m_{\mathrm{hid}}^{(n)}\in\Bigl[\frac{12n\ln2}{\mathfrak h},\ \frac{60n\ln2}{\mathfrak h}\Bigr]
\]

%dsp-level:0
\[
\qquad [\mathrm{BM.115}(2)(\mathrm{ii});\ \text{machine}];
\]

%dsp-level:0
\[
\textbf{the two currencies stay incommensurate:}
\]

%dsp-level:0
\[
\qquad \neg\exists\ \text{a derived ratio between}\ \lambda^{(n)}\ \text{and}\ E_{\mathrm{int}}^{(n)}\ —\ \text{the single residue is}\ \rho_{\mathrm B}
\]

%dsp-level:0
\[
\qquad [\mathrm{BM.82}(1),(5),(7);\ \mathrm{BM.69}(4);\ \mathrm{A\text{-}BM.9}];
\]

%dsp-level:0
\[
(5)\quad\textbf{the deficit-12 row is untouched:}
\]

%dsp-level:0
\[
\qquad \text{the two mass ceilings stand as properties of the reads they were proved for}\ [\mathrm{FT.76}(3);\ \mathrm{FT.79}(2);\ \mathrm{FT.97}(5):\ \text{the three-ceiling synthesis}]:
\]

%dsp-level:0
\[
\textbf{the resonance ladder adds no fourth ceiling and lifts none:}
\]

%dsp-level:0
\[
\qquad \lambda^{(n)}\ \text{is a sum of ceiling-bounded terms and its own bound}\ 4n\ \text{is the sum's}\ [(4)];
\]

%dsp-level:0
\[
\neg\text{a measured mass, hierarchy or range is compared}
\]

%dsp-level:0
\[
\qquad [\mathrm{P5};\ \mathrm{FT.84}(5\mathrm b);\ \text{machine: zero external tokens}];
\]

%dsp-level:0
\[
(6)\quad\textbf{the no-forbidden-state row, at the grouping level:}
\]

%dsp-level:0
\[
\qquad \textbf{every one of the}\ 6^n\ \textbf{ordered constituent tuples is a formed resonance with a nonempty threshold locus:}
\]

%dsp-level:0
\[
\#\{\text{ordered tuples at}\ n\}=6^n=(6,\ 36,\ 216),
\]

%dsp-level:0
\[
\qquad \#\{\text{distinct eigenvalues}\}=(6,\ 21,\ 55),
\]

%dsp-level:0
\[
\qquad \#\{\text{tuples with an empty locus}\}=0
\]

%dsp-level:0
\[
\qquad [\mathrm{BM.115}(5);\ \text{machine: exhaustive}]:
\]

%dsp-level:0
\[
\boxed{\ \textbf{there is no excluded composition at any level}\ —\ \text{the two named heaviest-floor triples are unread,}\ \neg\text{unreachable}\ }
\]

%dsp-level:0
\[
\qquad [\mathrm{BM.103}(8)];
\]

%dsp-level:0
\[
(7)\quad\textbf{status:}\qquad\textbf{derived:}\ (1),\ (2),\ (4),\ (6);\qquad\textbf{re-certified and extended:}\ (3);\qquad\textbf{untouched:}\ (5);
\]

%dsp-level:0
\[
\neg\text{a declaration},
\]

%dsp-level:0
\[
\quad\neg\text{a P4 addition},
\]

%dsp-level:0
\[
\quad\neg\text{a new dimensionful scale}\ [\mathrm{P1}],
\]

%dsp-level:0
\[
\quad\neg\text{a measured value}\ [\mathrm{P5}],
\]

%dsp-level:0
\[
\quad\neg\text{a currency conversion}\ [\mathrm{BM.83}(3)],
\]

%dsp-level:0
\[
\quad\neg\text{a species identification}\ [\mathrm{Def\text{-}BM.3}(7)];
\]

%dsp-level:0
\[
\alpha_B:\quad\neg\text{added},\quad\neg\text{derived},\quad\text{barred on}\ \mathrm{BM.61}(8)\text{'s four grounds}\ [\text{machine: the token scan of this tier}];
\]

\textbf{Req.} D2.9, D3.9 (Parts 1–2); CT3.15, CT3.17, CT3.17a(2) (Part 2); P1, P4, P5 (Parts 1–2); FT.76(3), FT.79(2), FT.84(2), FT.84(5), FT.97(5), FT.111(2), FT.142(5), FT.158(1) (Part 3); Def-BM.3(7), Def-BM.9, Def-BM.14(1), Def-BM.14(3), Def-BM.14(4), Def-BM.14(6), BM.31, BM.36(2), BM.36(3), BM.60(1), BM.61(8), BM.69(4), BM.82(1), BM.82(5), BM.82(7), BM.83(3), BM.103(7), BM.103(8), BM.114(1), BM.114(3), BM.114(4), BM.114(6), BM.115(2), BM.115(5), A-BM.9 (this part); exact integer orbit and grammar enumeration, exhaustive tuple censuses, exact radical comparison in $\mathbb M$ [math; machine: \texttt{checks/\allowbreak{}verify\_\allowbreak{}H\_\allowbreak{}antimatter\_\allowbreak{}resonance\_\allowbreak{}spectrum.py} §F].

\textbf{Pf.}

%dsp-level:0
\[
(1):\  \text{the eight rows re-cited from their own tags and their exact values re-computed}\ [\text{machine}];
\]

%dsp-level:0
\[
\qquad g(N)\ \text{for}\ N=1..7\ \text{by the labelled set-partition formula}\ [\mathrm{CT3.17a}(2);\ \text{machine}]
\]

%dsp-level:0
\[
(2):\  6+30+20=56,
\]

%dsp-level:0
\[
\qquad 6\cdot1+30\cdot3+20\cdot6=6+90+120=216=6^3\ [\text{machine}];
\]

%dsp-level:0
\[
\qquad S_3\ \text{acts freely on the all-distinct tuples}: |\operatorname{Stab}|
\]

%dsp-level:0
\[
= 1
\]

%dsp-level:0
\[
\implies |\text{orbit}|
\]

%dsp-level:0
\[
= 3!
\]

%dsp-level:0
\[
= 6\ \text{and}\ \#\text{orbits}
\]

%dsp-level:0
\[
= \tfrac{6\cdot5\cdot4}{6}
\]

%dsp-level:0
\[
= 20\ [\text{machine}]
\]

%dsp-level:0
\[
\varepsilon_{ijk}\ \text{antisymmetric in three slots}\ [\mathrm{D3.9};\ \mathrm{CT3.15}];
\]

%dsp-level:0
\[
\qquad \neg\exists\ \text{a flavour}\ \mathsf{SU}(3)\ [\mathrm{BM.103}(7);\ \text{machine: the symbol census}]\implies20\ \text{is an orbit count}
\]

%dsp-level:0
\[
(3):\  \mathrm{BM.103}(7)\text{'s four rows re-cited};
\]

%dsp-level:0
\[
\qquad 8\notin\{2n+1\}_{n\le3}\cup\{\binom{6+n-1}n\}_{n\le3}\cup\{\binom6n\}_{n\le3}\cup\{g(N)\}_{N\le7}\ [\text{machine: exhaustive membership test over}\ \{3,5,7\}\cup\{6,21,56\}\cup\{6,15,20\}\cup\{0,1,1,3,10,25,105\}]
\]

%dsp-level:0
\[
(4):\  \max_sD_S(s)=4\ \text{at}\ s_2\ [\mathrm{Def\text{-}BM.14}(3);\ \text{machine: exact comparison in}\ \mathbb M];
\]

%dsp-level:0
\[
\qquad \lambda^{(n)}_{\max}=n\cdot4=4n\ \text{and it is attained by}\ (s_2,\dots,s_2)\ \text{with multiplicity}\ 1\ [\mathrm{BM.114}(3);\ \text{machine}]
\]

%dsp-level:0
\[
\mathrm{FT.76}(3),\ \mathrm{FT.79}(2),\ \mathrm{FT.84}(2)\ \text{re-cited at their own loci};
\]

%dsp-level:0
\[
\qquad E_{\mathrm{int}}^{(n)}
\]

%dsp-level:0
\[
\in [12n,60n]\ \text{on the six-state ladder and}\ [2n,60n]\ \text{on the computed-cell ladder}\ [\mathrm{BM.115}(2)(\mathrm i),(\mathrm{ii})];
\]

%dsp-level:0
\[
\qquad \mathrm{BM.82}(1),(5),(7)\ \text{and}\ \mathrm{BM.69}(4)\text{'s four grounds re-cited}
\]

%dsp-level:0
\[
(5):\  \mathrm{FT.97}(5)\text{'s three-ceiling synthesis re-cited};
\]

%dsp-level:0
\[
\qquad \lambda^{(n)}\le4n\ \text{is the sum's own bound}\ [(4);\ \text{machine}];
\]

%dsp-level:0
\[
\qquad \text{the external-token scan}\ [\text{machine}]
\]

%dsp-level:0
\[
(6):\  6^n=(6,36,216);
\]

%dsp-level:0
\[
\qquad \#\text{distinct}=(6,21,55)\ [\mathrm{BM.114}(3),(5);\ \text{machine}];
\]

%dsp-level:0
\[
\qquad \#\{\text{empty loci}\}=0\ [\mathrm{BM.115}(5);\ \text{machine: all}\ 216\ \text{triples carry}\ E_{\mathrm{int}}\in[6,180]\subset\mathbb Z_{>0}]
\]

%dsp-level:0
\[
(7):\  \text{the standing-rule scans and the barred-object token scan}\ [\text{machine}].\ \blacksquare
\]

\textbf{Phys.} Asked to group the spectrum, the corpus gives what it has and no more. Eight groupings are derived and each has an exact value: the binomial charge classes, the sealing grammar's counts, the multiset orbits, the all-distinct orbits, the three ordered generation floors, the two ordered branches, the cyclic three of the flavour channel, and the subfield ladder of the eigenvalue field. The one the resonance spectrum actually carries is the three-level grouping by floor, branch and charge class with the orbit multiplicities, and at three chains that is six orbits of size one, thirty of size three and twenty of size six — fifty-six orbits giving fifty-five eigenvalues after the single powers-of-two collision. Twenty appears, and it is exactly the count of ways to pick three distinct states from six, which is also where the epsilon seal's free symmetric-group action lives; the tier states plainly that this is an orbit count and not a decuplet dimension, because the flavour symmetry the corpus derives is a cyclic three and there is no special unitary three anywhere in the inventory. The octet adjudication is re-certified and extended: the derived three is colour's, the derived eight is a barred capacity ratio, the derived ten is a grammar count, and eight does not occur among any of the four grouping families at any formed level. The ceilings behave exactly as they should — the single-chain pole ceiling is four, the composite ceiling is four times the chain count and is attained by the heaviest branch at every level, and a composite eigenvalue is a sum of single-chain reads rather than a single-chain read, so no ceiling is violated and none is extended. And the count that answers the forbidden-state question a second time is the simplest one: all two hundred sixteen ordered constituent triples are formed resonances with nonempty threshold loci, so the two heaviest compositions anyone might name are unread rather than unreachable.

\end{theorem}

\medskip\noindent\textit{\texorpdfstring{BM.118 — The status of the Antimatter and Resonance Tier: adjudication, scope, and the open items}{BM.118 - The status of the Antimatter and Resonance Tier: adjudication, scope, and the open items}.}\ Structural-constraints record: see Appendix B, Section B.1.31 of the Addendum.\par\medskip

\phantomsection
\section[Tier BM.22 — The Annihilation Channel, Concrete Spin, and…]{\texorpdfstring{Tier BM.22 — The Annihilation Channel, Concrete Spin, and Fermion Identification: the Janus Tunnel, the Zig-Zag World Line, and the Spin-Dependent Transport Defect}{Tier BM.22 - The Annihilation Channel, Concrete Spin, and Fermion Identification: the Janus Tunnel, the Zig-Zag World Line, and the Spin-Dependent Transport Defect}}

\textbf{Locus.} This tier takes the residue $(\mathrm{R9})$ that $\mathrm{BM.118}(7)$ named — \textit{the annihilation channel: an inter-sector transition operator, $\neg$in the inventory} — and constructs the channel, then follows the two chains that fall out of it: the concrete spin and the composite statistics. The crossing is not a velocity and the tier says so first: $c=1$ is a convention, the dispersion gives a strict subluminal deficit $1-v=m^2/(E(E+p))>0$ at every finite momentum, and the boundary the corpus actually derives between the materialized and the fluctuation content is the \textbf{Janus real/virtual register}, whose $\mathcal V$ sector is the maximally raw-entropic state $\rho_\circ=\tfrac14I$ with no fermion terminus at $|E_{\mathcal R}|=0$ ($\mathrm{BM.119}$). On that register the crossing operator is derived, not posited: the two moves $\hat J_\pm$ of $\mathrm{D1.2}$, the master generator $Q=x_+(\hat J_+-P_{\mathcal V})+x_-(\hat J_--P_{\mathcal R})$, its one-dimensional kernel giving the stationary measure, its spectrum $\{0,-(x_++x_-)\}$, and detailed balance as a theorem ($\mathrm{BM.120}$). The zig-zag is then the resolvent's own evenness — $(\hat U^{2m})_{\varnothing\varnothing}=(x_+x_-)^m$ with every odd term zero — so each term is a create-propagate-annihilate round trip whose legs all run forward in the ledger order, and the "backward leg" of a diagram is here the CP-conjugate \textit{configuration} on a forward leg, which is why no inverse tick is needed ($\mathrm{BM.121}$). The pair channel is derived with its charge balance, its exact threshold, its weight and its three CP legs — the nilpotent weak step for the charge inversion, the CP transport for the parity inversion, the local scale for the energy — while the two-photon supply is barred on four named grounds, so the channel closes and the amplitude does not ($\mathrm{BM.122}$). Annihilation is the reverse move, its content returning to the $\mathcal V$ sector at the entropy maximum rather than into radiation legs, and $(\mathrm{R9})$ closes at the channel level with the cross-section residue narrowing to the momentum-space resolvent and the gauge coupling of the charge read ($\mathrm{BM.123}$). On the spin side the tier separates two angles that a single-symbol reading would fuse: the winding $\Theta$ and the frame-holonomy angle $\Theta_R$ are provably not proportional, on a dimension mismatch, so the spin-dependent transport defect is a \textbf{multiplet ladder} $\Delta(m)=4\sin^2(m\Theta_R/4)$ over $m\in\{k,k-2,\dots,-k\}$ and not a number $\Delta(j)$ ($\mathrm{BM.124}$, $\mathrm{BM.125}$). A spin-dependent crossing weight is refuted on the four standing grounds, and what survives is a \textit{domain} correlation rather than a prohibition ($\mathrm{BM.126}$). The winding parity is then a ledger read, separated from the $\mathrm{D1.5}$ kind read that $\mathrm{BM.106}(1)$ already showed cannot carry a spin-parity theorem ($\mathrm{BM.127}$). Finally the three-chain exchange structure is built — $\mathcal H_3$, the $S_3$ action with its exact character formula $d^{\,c(\sigma)}$, and the vanishing of $\wedge^3(\mathbb C^2)$ against the one-dimensionality of $\wedge^3(\mathbb C^3)$, which is why the baryon's antisymmetry is colour's and the statistics is the winding's — giving the multiplicative exchange law $\varepsilon=\prod_a\varepsilon(c_a)$, the meson's unconditional boson sign, and the baryon's sign as an exact function of its three constituent windings ($\mathrm{BM.128}$, $\mathrm{BM.129}$). $\mathrm{BM.130}$ adjudicates.

\textbf{Standing rows.} $\alpha_B$ stays barred on $\mathrm{BM.61}(8)$'s four grounds; §Cal-BM stays empty except for $\mathrm{A\text{-}BM.9}$'s logged slot with its value $\neg$issued. No right-ray matter line, no seesaw-type mechanism, no Yukawa-type object and no external mixing matrix is suggested, added or required $[\mathrm{EW.15};\ \mathrm{FT.145}(3);\ \mathrm{BM.96}(4),(6)\ \text{re-cited},\ \neg\text{extended}]$. No clause of $\mathrm{BM.1}$–$\mathrm{BM.118}$ is retracted, replaced or renumbered: the rows this tier touches are re-scoped at their own loci — $\mathrm{BM.118}(7)$'s residue $(\mathrm{R9})$, $\mathrm{BM.112}(7)$'s and $\mathrm{Def\text{-}BM.16}(7)$'s residue $(\mathrm{R7})$ at $n\ge3$, $\mathrm{BM.109}(2)$'s pair-weight naming, $\mathrm{BM.114}(2)$'s colour-antisymmetry row, $\mathrm{BM.27}(7)$, $\mathrm{BM.29}(5)$ and $\mathrm{Def\text{-}BM.15}(5)$'s no-annihilation rows, $\mathrm{BM.6}(6)$'s resolvent specification and $\mathrm{BM.6}(2)$'s discharged specification — and each rescope is recorded in $\mathrm{BM.130}(5)$, in $\mathrm{BM.77}(7)$'s pattern. No measured coupling, mass, width, lifetime, cross-section or ratio enters any derivation $[\mathrm{P5}]$. No cosmological coordinate is used: a particle accelerator and a deep-gravity region are both local regions with high $(E,N,G)$, and every threshold and width is a function of those four local parameters alone $[\mathrm{Def\text{-}BM.17}(6)(\mathrm i);\ \mathrm{A6};\ \mathrm{D1.8}]$. "Low energy" and "high energy" are calibration points on the derived functions and $\neg$regimes of two theories $[\mathrm{BM.116}(6);\ \mathrm{Def\text{-}BM.17}(6)(\mathrm{ii})]$. No black-hole physics is derived: the comparison radius $r_s$ enters only through $\mathrm{BM.15}(4)$'s quarantined change of variable, with $\neg$a collapse reading and $\neg$a horizon identification $[\mathrm{P3};\ \mathrm{BM.115}(4)]$. Species names are confined to the identification layer $[\mathrm{Def\text{-}BM.3}(7)]$. Topology labels are internal bookkeeping $[\mathrm{Def\text{-}H.32}(4)]$. Amplitude $\ne$ mass $[\mathrm{FT.79}(5)]$. Rung $\ne$ scale $[\mathrm{D0.8}]$. No currency is converted $[\mathrm{BM.83}(3);\ \mathrm{BM.102}(6)]$. No theorem of this tier imports a HistLog disposition as premise: Q9, Q11, items 7, 13, 17, 20 and deficits $1'$, $3$, $4$, $6$, $11$, $12$ are cited as pointers. One dimension only: units are powers of $\mathfrak q$ $[\mathrm{D2.9};\ \mathrm{CT2.1}]$.

\textbf{Local notation, fixed for the whole tier.} $\mathcal J:=\operatorname{span}_{\mathbb C}\{|\mathcal V\rangle,\ |\mathcal R\rangle\}\cong\mathbb C^2$ the per-pair Janus register in the (virtual, real) order of $\mathrm{FT.108}(4)$; $P_{\mathcal V},P_{\mathcal R}$ its two orthogonal projectors; $\hat J_+:=|\mathcal R\rangle\langle\mathcal V|$ the crystallization move and $\hat J_-:=|\mathcal V\rangle\langle\mathcal R|$ the dissolution move $[\mathrm{D1.2}]$; $Q$ the master generator over $\mathbb Z[x_+,x_-]$ $[\mathrm{FT.108}(4)]$; $\varepsilon_q\in[0,\tfrac14)$, $x_+=1-4\varepsilon_q$, $x_-=4\varepsilon_q$, $p=x_+/(x_++x_-)=1-4\varepsilon_q$ $[\mathrm{A4};\ \mathrm{D0.4};\ \mathrm{CT0.8};\ \mathrm{BM.19}(2)]$; $w$ the formal fugacity and $\hat U$ the single-channel step $[\mathrm{FT.108}(5);\ \mathrm{Def\text{-}H.32}(3)]$; $\Theta(c)=\sum_{e\in E(c)}\widehat{ds}^{\,2}(e)\in2\mathbb Z$ the winding, $k=\Theta/2\in\mathbb Z_{\ge0}$, $j=|k|/2$, $\Psi(c)=\mathrm{Sym}^{|k|}(\mathcal H_{\mathsf v})$ with $\mathcal H_{\mathsf v}\cong\mathbb C^2$ $[\mathrm{FT.22};\ \mathrm{Def\text{-}H.38};\ \mathrm{EW.32}]$; $\Theta_R$ the frame-holonomy angle with $[\Theta_R]=\mathfrak q^0$ $[\mathrm{Def\text{-}H.21}(5),(6);\ \mathrm{FT.55}(1)]$; $\phi:=\Theta_R/2$; $T_m$ the Chebyshev polynomial of the first kind; $\Delta(m):=4\sin^2(m\Theta_R/4)=2-2\cos(m\phi)$ the transport defect at multiplet weight $m$ $[\mathrm{FT.60}(5);\ \mathrm{FT.54}(4);\ \mathrm{EW.33}]$; $\mathcal H_\omega(c)$, $\hat\omega$, $\hat\chi$, $P_{L/R}$, $|L\rangle$, $|R\rangle$ as in $\mathrm{Def\text{-}H.33}$; $\widehat V=E_+\otimes P_L$ the nilpotent step $[\mathrm{BM.76}(1)]$; $\mu_{\mathrm{CP}}=e^{i\pi\Theta}$ and $\varepsilon(c)=(-1)^{|k(c)|}=(-1)^{2j(c)}$ $[\mathrm{BM.108}(3);\ \mathrm{BM.113}(2),(4)]$; $\mathcal H_2$, $\hat P_{12}$ as in $\mathrm{Def\text{-}BM.16}$; $d:=\Theta(P)\ge0$ a connector winding $[\mathrm{BM.2}(2);\ \mathrm{FT.34}(3)]$; $S_\sigma$ the sign of a permutation $\sigma$; $c(\sigma)$ its cycle count; $\mathcal T_0$, $\mathcal T$, $\Gamma$, $\mathcal N$, $\chi_{\{\cdot\}}$ as in $\mathrm{Def\text{-}BM.17}$; $[\cdot]$ the $\mathfrak q$-exponent map $[\mathrm{CT2.1};\ \mathrm{D2.9}]$.

\begin{definition}[{\texorpdfstring{Def-BM.18 — The Janus register and its crossing operators}{Def-BM.18 - The Janus register and its crossing operators}}]\label{tag:Def-BM-18}
\textbf{Def.} A construction on $\mathrm{D1.1}$–$\mathrm{D1.2}$'s real/virtual distinction, $\mathrm{FT.108}(4)$'s master generator and $\mathrm{Def\text{-}BM.16}$'s two-chain space; no new axiom, no new datum, no new parameter, no velocity operator, no position operator, no external scattering object.

%dsp-level:0
\[
(1)\quad\textbf{the register:}
\]

%dsp-level:0
\[
\qquad \boxed{\ \mathcal J\ :=\ \operatorname{span}_{\mathbb C}\bigl\{|\mathcal V\rangle,\ |\mathcal R\rangle\bigr\}\ \cong\ \mathbb C^2\ }
\]

%dsp-level:0
\[
\qquad \text{in the}\ (\text{virtual},\ \text{real})\ \text{order of}\ \mathrm{FT.108}(4);
\]

%dsp-level:0
\[
e\in\mathcal R\iff\bigl|\Pi_{\mathcal V}\rho\bigr|\le\varepsilon_q\ \vee\ \bigl|\Pi_{\mathcal V}\rho\bigr|\ge1-\varepsilon_q;\qquad\neg\text{registration}\implies e\in\mathcal V\qquad[\mathrm{A4};\ \mathrm{D1.1};\ \mathrm{CT0.1}]:
\]

%dsp-level:0
\[
\textbf{constitutive clause:}\ \ \mathcal J\ \text{is the two-state register of one pair across the}\ \mathrm{A4}\ \text{criterion};
\]

%dsp-level:0
\[
\qquad \textbf{non-derived clause:}\ \ \text{the basis order}\ (\mathcal V,\mathcal R);
\]

%dsp-level:0
\[
(2)\quad\textbf{the two crossing operators:}
\]

%dsp-level:0
\[
\qquad \boxed{\ \hat J_+\ :=\ |\mathcal R\rangle\langle\mathcal V|\ \ (\text{crystallization},\ J^+:\ n_{\mathcal R}:0\to1),\qquad\hat J_-\ :=\ |\mathcal V\rangle\langle\mathcal R|\ \ (\text{dissolution},\ J^-:\ n_{\mathcal R}:1\to0)\ }
\]

%dsp-level:0
\[
[\mathrm{D1.2}:\ \text{a seal is the transition of one pair across the}\ \mathrm{A4}\ \text{criterion, in either branch}];
\]

%dsp-level:0
\[
\qquad P_{\mathcal V}:=|\mathcal V\rangle\langle\mathcal V|,\quad P_{\mathcal R}:=|\mathcal R\rangle\langle\mathcal R|,\quad P_{\mathcal V}+P_{\mathcal R}=\mathrm{id}
\]

%dsp-level:0
\[
\qquad [\text{machine}];
\]

%dsp-level:0
\[
(3)\quad\textbf{the master generator:}\qquad\boxed{\ Q\ :=\ \begin{pmatrix}-x_+&x_+\\ x_-&-x_-\end{pmatrix}\ \ \text{over}\ \ \mathbb Z[x_+,x_-]\ }\qquad[\mathrm{FT.108}(4):\ \text{the per-pair Janus dynamics}]:
\]

%dsp-level:0
\[
\textbf{the decomposition (derived, recorded here):}
\]

%dsp-level:0
\[
\qquad Q\ =\ x_+\bigl(\hat J_+-P_{\mathcal V}\bigr)\ +\ x_-\bigl(\hat J_--P_{\mathcal R}\bigr)
\]

%dsp-level:0
\[
\qquad [\text{machine: entrywise at six rational witnesses}];
\]

%dsp-level:0
\[
\textbf{constitutive clause:}\ \ Q\ \text{is}\ \mathrm{FT.108}(4)\text{'s own generator,}\ \neg\text{a new dynamics};
\]

%dsp-level:0
\[
\qquad \textbf{the row convention:}\ \ Q\ \text{acts on row vectors, so}\ Q_{ij}\ \text{is the rate from}\ i\ \text{to}\ j;
\]

%dsp-level:0
\[
(4)\quad\textbf{the crossing space of a chain:}
\]

%dsp-level:0
\[
\qquad \boxed{\ \mathcal T_{\mathcal J}(c)\ :=\ \mathcal J\ \otimes\ \mathcal H_\omega(c)\ \otimes\ \Psi(c),\qquad\dim\mathcal T_{\mathcal J}(c)\ =\ 4\,\bigl(|k(c)|+1\bigr)\ }
\]

%dsp-level:0
\[
[\mathrm{Def\text{-}H.33};\ \mathrm{Def\text{-}H.38};\ \mathrm{Def\text{-}BM.16}(1);\ \text{machine: exact integers}];
\]

%dsp-level:0
\[
\qquad \textbf{the pair crossing space:}\quad\mathcal T_{\mathcal J}^{\,2}(c):=\mathcal T_{\mathcal J}(c)\otimes\mathcal T_{\mathcal J}(c^{\mathrm{CP}}),
\]

%dsp-level:0
\[
\qquad \dim=16\,\bigl(|k|+1\bigr)^2;
\]

%dsp-level:0
\[
(5)\quad\textbf{the channel clauses:}
\]

%dsp-level:0
\[
\qquad \textbf{a pair creation}\ :
\]

%dsp-level:0
\[
= \ \hat J_+\ \text{on two registers, with the two members identified as a CP-pair}\ [\mathrm{Def\text{-}BM.15}(3)];
\]

%dsp-level:0
\[
\qquad \textbf{weight}\ x_+^2\ \text{for the two one-way moves and}\ x_+x_-\ \textbf{for one round trip};
\]

%dsp-level:0
\[
\textbf{an annihilation}\ :=\ \hat J_-\ \text{on the same two registers};\qquad\textbf{weight}\ x_-^2\ \text{one-way,}\ x_+x_-\ \text{per round trip};
\]

%dsp-level:0
\[
\textbf{a zig-zag}\ :
\]

%dsp-level:0
\[
= \ \text{a finite alternating sequence of}\ \hat J_+\ \text{and}\ \hat J_-\ \text{moves, every step recorded in the ledger order}\ [\mathrm{A6};\ \mathrm{D1.8}];
\]

%dsp-level:0
\[
\textbf{constitutive clause:}\ \ \text{the three words name move sequences on}\ \mathcal J,
\]

%dsp-level:0
\[
\ \neg\text{world lines in a spacetime and}\ \neg\text{asymptotic scattering states};
\]

%dsp-level:0
\[
(6)\quad\textbf{the bars (standing):}
\]

%dsp-level:0
\[
\qquad \neg\exists\ \text{a velocity operator, a position operator or a momentum-space state on}\ \mathcal J\ [\mathrm{BM.10}(1):\ \neg[\hat x,\hat p]=i\ \text{in the finite register};\ \mathrm{BM.6}(6):\ \text{the momentum-space}\ \Sigma\text{-resolvent is an open interface}];
\]

%dsp-level:0
\[
\neg\exists\ \text{a cross-section, an S-matrix element or a Lorentz-invariant phase space}\ [\mathrm{BM.6}(6);\ \mathrm{BM.118}(7)(\mathrm{R9})];
\]

%dsp-level:0
\[
\neg\exists\ \text{a direct}\ \mathrm{EM}^+\to\mathrm{EM}^-\ \text{operator: the sectors are amplitude-equal configuration orbits and the channel that connects them is mediated by}\ \mathcal J\ \text{and}\ \mathrm{CP}\ [\mathrm{BM.25}(4);\ \mathrm{BM.27}(6),(7);\ \mathrm{BM.29}(5)];
\]

%dsp-level:0
\[
\textbf{the kind bar:}
\]

%dsp-level:0
\[
\qquad \mathcal J\ \text{is}\ \neg\mathcal H_\omega\ \text{and}\ \neg\text{one of the four single-object}\ \mathbb Z_2\ \text{gradings}\ [\mathrm{BM.106}(4);\ \mathrm{BM.112}(5);\ \mathrm{FT.125}(5);\ \mathrm{FT.89}(2)]:
\]

%dsp-level:0
\[
\boxed{\ \textbf{the real/virtual distinction is the Janus register; the orientation distinction is}\ \mathcal H_\omega\textbf{; identifying them is a barred kind confusion}\ }
\]

%dsp-level:0
\[
(7)\quad\textbf{status and scope:}
\]

%dsp-level:0
\[
\qquad \textbf{derived:}\ (3)\text{'s decomposition},\ (4)\text{'s dimensions};
\]

%dsp-level:0
\[
\qquad \textbf{constitutive:}\ (1)\text{'s register},\ (2)\text{'s two moves},\ (5)\text{'s three names};
\]

%dsp-level:0
\[
\textbf{the scope inherited at every use:}\ \ \mathrm{Def\text{-}H.32}\text{'s}\ (\mathrm{S1})\text{--}(\mathrm{S4});
\]

%dsp-level:0
\[
\qquad \mathrm{Def\text{-}H.38}\text{'s nonnegative cone and its}\ k\ge3\ \text{aspect bar};
\]

%dsp-level:0
\[
\qquad \mathrm{BM.8}(2)\text{'s stationarity hypothesis for any chart evolution};
\]

%dsp-level:0
\[
\neg\text{a P4 addition},
\]

%dsp-level:0
\[
\quad\neg\text{a new dimensionful scale}\ [\mathrm{P1}],
\]

%dsp-level:0
\[
\quad\neg\text{a measured value}\ [\mathrm{P5}],
\]

%dsp-level:0
\[
\quad\neg\text{a species identification}\ [\mathrm{Def\text{-}BM.3}(7)],
\]

%dsp-level:0
\[
\quad\neg\text{a rate datum}\ [\mathrm{Cal\text{-}F}\ \text{carries no rate};\ \mathrm{BM.44}(5)];
\]

%dsp-level:0
\[
\textbf{beyond this tier:}\ \ \text{the momentum-space}\ \Sigma\text{-resolvent}\ [\mathrm{BM.6}(6)],
\]

%dsp-level:0
\[
\qquad \text{the gauge coupling of the charge read}\ [\mathrm{Def\text{-}BM.8}(8)].
\]

\textbf{Rmk (what the register is and is not).} $\mathcal J$ is the two-state register of \textit{one pair} across the crystallization criterion — the object D1.2 named and FT.108 derived a stationary measure for. It is not a spacetime, not a world line, and not a Hilbert space of asymptotic states. Clause (6)'s kind bar is the load-bearing one: the corpus has an orientation doublet $\mathcal H_\omega$ per chain, and it would be a category error to read the real/virtual distinction as that doublet's two rays, exactly as it would be to read a single-chain grade as a permutation sign. The zig-zag of clause (5) is a word in two moves, recorded forward, and its weight is a monomial in $x_+$ and $x_-$.

\textbf{Req.} A4, A6, D0.8, D1.1, D1.2, D1.8 (Part 1); CT0.1, CT0.8, CT2.1, D2.9 (Part 1, Part 2); P1, P3, P4, P5 (Parts 1–2); Def-H.21(5), Def-H.21(6), Def-H.32, Def-H.33, Def-H.38, FT.22, FT.89(2), FT.108(4), FT.108(5), FT.125(5) (Part 3); Def-BM.3(7), Def-BM.15(3), Def-BM.16(1), BM.2(2), BM.6(6), BM.8(2), BM.10(1), BM.25(4), BM.27(6), BM.27(7), BM.29(5), BM.44(5), BM.106(4), BM.112(5), BM.118(7) (this part); exact $2\times2$ matrix algebra over $\mathbb Z[x_+,x_-]$ and over $\mathbb Q(i)$, exact integer dimension arithmetic [math; machine: \texttt{checks/\allowbreak{}verify\_\allowbreak{}H\_\allowbreak{}janus\_\allowbreak{}tunnel\_\allowbreak{}concrete\_\allowbreak{}spin.py} §A].

\textbf{Pf.}

%dsp-level:0
\[
\text{Definition block: the constitutive clauses are}\ (1)\text{'s register},
\]

%dsp-level:0
\[
\ (2)\text{'s two moves and}\ (5)\text{'s three names};
\]

%dsp-level:0
\[
\ \text{the non-derived clauses are the basis order of}\ (1)\ \text{and the row convention of}\ (3);
\]

%dsp-level:0
\[
\ (3)\text{'s decomposition and}\ (4)\text{'s dimensions are derived and machine-checked};
\]

%dsp-level:0
\[
\ (6),
\]

%dsp-level:0
\[
(7)\ \text{are bars and scope records}.\ \blacksquare
\]

\textbf{Phys.} One pair, crossing a criterion, has two states and two moves, and the corpus already derived both: D1.2 named the seal as a transition across the A4 criterion in either branch, and FT.108 wrote down the master generator whose kernel is one-dimensional and gives the stationary measure. This definition makes the two moves into operators on a two-dimensional register, decomposes the generator as a sum of jump-minus-loss terms with the two derived weights, and tensors the register with the chain's own orientation doublet and winding carrier to give a crossing space of exact finite dimension. Three words are then defined and no more: a pair creation is the crystallization move on two registers whose members are identified as a CP-pair, an annihilation is the dissolution move on the same two, and a zig-zag is a finite alternating word in the two moves, every step recorded in ledger order. What the definition refuses is the vocabulary that would overreach — no velocity operator, no position operator, no momentum-space state, no cross-section, no S-matrix, and above all no identification of the real/virtual register with the orientation doublet, which is a different two-dimensional object doing different work.

\end{definition}

\medskip\noindent\textit{\texorpdfstring{BM.119 — The c-threshold adjudication: no velocity boundary, the exact subluminal deficit, and the register boundary that is derived}{BM.119 - The c-threshold adjudication: no velocity boundary, the exact subluminal deficit, and the register boundary that is derived}.}\ Structural-constraints record: see Appendix B, Section B.1.32 of the Addendum.\par\medskip

\begin{theorem}[{\texorpdfstring{BM.120 — The Janus tunnel: the crossing operators, the exact spectrum, and detailed balance}{BM.120 - The Janus tunnel: the crossing operators, the exact spectrum, and detailed balance}}]\label{tag:BM-120}
\textbf{Thm.}

%dsp-level:0
\[
(1)\quad\textbf{the move algebra (exact, over}\ \mathbb Q\text{):}
\]

%dsp-level:0
\[
\qquad \hat J_+^2=0,
\]

%dsp-level:0
\[
\qquad \hat J_-^2=0,
\]

%dsp-level:0
\[
\qquad \hat J_+\hat J_-=P_{\mathcal V},
\]

%dsp-level:0
\[
\qquad \hat J_-\hat J_+=P_{\mathcal R},
\]

%dsp-level:0
\[
\qquad \{\hat J_+,\hat J_-\}=\mathrm{id}
\]

%dsp-level:0
\[
[\text{machine: all five identities entrywise on the}\ 2\times2\ \text{register}];
\]

%dsp-level:0
\[
\qquad \boxed{\ \textbf{each move is nilpotent and the two anticommute to the identity}\ }
\]

%dsp-level:0
\[
\textbf{the reading:}\quad\hat J_+\hat J_-=P_{\mathcal V}\ \text{is the round trip}\ \mathcal V\to\mathcal R\to\mathcal V;
\]

%dsp-level:0
\[
\qquad \hat J_-\hat J_+=P_{\mathcal R}\ \text{is the round trip}\ \mathcal R\to\mathcal V\to\mathcal R;
\]

%dsp-level:0
\[
\neg\exists\ \text{a doubled move in the same direction}\ —\ \textbf{the}\ \mathrm{A4}\ \textbf{criterion is binary, and the nilpotency is its operator form}\ [\mathrm{D1.2};\ \mathrm{BM.23}(2)\text{'s overflow pattern}];
\]

%dsp-level:0
\[
(2)\quad\textbf{the generator decomposes as jump minus loss:}
\]

%dsp-level:0
\[
\qquad \boxed{\ Q\ =\ x_+\bigl(\hat J_+-P_{\mathcal V}\bigr)\ +\ x_-\bigl(\hat J_--P_{\mathcal R}\bigr)\ }
\]

%dsp-level:0
\[
\qquad [\mathrm{FT.108}(4);\ \text{machine: entrywise at six rational witnesses}\ (x_+,x_-)\in\{(\tfrac12,\tfrac12),(\tfrac34,\tfrac14),(1,0),(0,1),(\tfrac23,\tfrac13),(\tfrac78,\tfrac18)\}];
\]

%dsp-level:0
\[
\textbf{the row-sum law:}
\]

%dsp-level:0
\[
\qquad Q\ \text{acts on row vectors and every row sums to zero:}\quad(-x_+)+x_+=0,
\]

%dsp-level:0
\[
\qquad x_-+(-x_-)=0
\]

%dsp-level:0
\[
\qquad [\text{machine}]:
\]

%dsp-level:0
\[
\therefore\ Q\ \textbf{is a master generator and}\ \neg\textbf{a unitary}\ —\ \text{the crossing is a stochastic ledger move,}\ \neg\text{a rotation of a state};
\]

%dsp-level:0
\[
(3)\quad\textbf{the kernel, the stationary measure and detailed balance:}
\]

%dsp-level:0
\[
\qquad (x_-,\ x_+)\,Q\ =\ \bigl(x_-(-x_+)+x_+x_-,\ \ x_-x_++x_+(-x_-)\bigr)\ =\ (0,\ 0)
\]

%dsp-level:0
\[
\qquad [\text{machine: exact in}\ \mathbb Z[x_+,x_-]];
\]

%dsp-level:0
\[
\det Q\ =\ (-x_+)(-x_-)-x_+x_-\ =\ 0,\qquad Q\ne0\implies\operatorname{rk}Q=1,\qquad\dim\ker Q=1\qquad[\mathrm{FT.108}(4);\ \text{machine}];
\]

%dsp-level:0
\[
\boxed{\ p\ =\ \frac{x_+}{x_++x_-}\ =\ 1-4\varepsilon_q\ \ \textbf{identically}\ }
\]

%dsp-level:0
\[
\qquad [\mathrm{BM.19}(2);\ \text{machine: exact at seven rational thresholds}];
\]

%dsp-level:0
\[
\textbf{detailed balance is derived, not assumed:}
\]

%dsp-level:0
\[
\qquad p_{\mathcal V}\,x_+\ =\ p_{\mathcal R}\,x_-\iff\frac{x_-}{x_++x_-}\,x_+=\frac{x_+}{x_++x_-}\,x_-\iff x_-x_+=x_+x_-
\]

%dsp-level:0
\[
\qquad [\mathrm{FT.108}(4);\ \text{machine}];
\]

%dsp-level:0
\[
(4)\quad\textbf{the exact spectrum:}
\]

%dsp-level:0
\[
\qquad \operatorname{tr}Q=-(x_++x_-),
\]

%dsp-level:0
\[
\qquad \det Q=0\implies\boxed{\ \operatorname{spec}Q=\bigl\{0,\ -(x_++x_-)\bigr\}\ }
\]

%dsp-level:0
\[
\qquad [\text{machine: the characteristic polynomial}\ \lambda^2+(x_++x_-)\lambda];
\]

%dsp-level:0
\[
\textbf{at the specialization}\ x_++x_-=1:\qquad\operatorname{spec}Q=\{0,\ -1\},\qquad\operatorname{tr}Q=-1,\qquad\operatorname{rk}Q=1\qquad[\text{machine}];
\]

%dsp-level:0
\[
\boxed{\ \textbf{the zero eigenvalue is the stationary mode and}\ -1\ \textbf{is the relaxation mode at unit total weight}\ }
\]

%dsp-level:0
\[
\qquad [\text{machine: the two eigenvectors}\ (x_-,x_+)\ \text{and}\ (1,-1)];
\]

%dsp-level:0
\[
(5)\quad\textbf{the tunnel operator and its weights:}
\]

%dsp-level:0
\[
\qquad \textbf{one-way creation}\ \hat J_+:\ \text{weight}\ x_+=1-4\varepsilon_q\in(0,1];
\]

%dsp-level:0
\[
\qquad \textbf{one-way annihilation}\ \hat J_-:\ \text{weight}\ x_-=4\varepsilon_q\in[0,1);
\]

%dsp-level:0
\[
\textbf{the round trip:}
\]

%dsp-level:0
\[
\qquad \boxed{\ \mathcal W_{\text{round}}\ =\ x_+x_-\ =\ 4\varepsilon_q(1-4\varepsilon_q)\ =\ \frac14-16\Bigl(\varepsilon_q-\frac18\Bigr)^{\!2}\ \le\ \frac14\ }
\]

%dsp-level:0
\[
\qquad [\mathrm{BM.109}(2);\ \mathrm{BM.44}(1);\ \text{machine: the SOS completion and the}\ 256\text{-point grid}];
\]

%dsp-level:0
\[
\textbf{the detailed-balance point is the maximum:}
\]

%dsp-level:0
\[
\qquad x_+=x_-\iff\varepsilon_q=\frac18,
\]

%dsp-level:0
\[
\qquad \text{and}
\]

%dsp-level:0
\[
\qquad x_+x_-\le\Bigl(\frac{x_++x_-}2\Bigr)^{\!2}=\frac14\ \ \text{with equality}\iff x_+=x_-
\]

%dsp-level:0
\[
[\text{machine: the AM--GM row on exact rationals, equality locus verified}];
\]

%dsp-level:0
\[
\qquad \boxed{\ \textbf{the round-trip weight is maximal exactly where the two moves balance}\ }
\]

%dsp-level:0
\[
(6)\quad\textbf{the type correction (the crossing acts on}\ \mathcal J\textbf{, not on}\ \mathcal H_\omega\textbf{):}
\]

%dsp-level:0
\[
\qquad \mathcal H_\omega(c)=\operatorname{span}\{|c\rangle,|\omega(c)\rangle\}\ \text{is one chain's orientation doublet}\ [\mathrm{Def\text{-}H.33}];
\]

%dsp-level:0
\[
\qquad \mathcal J=\operatorname{span}\{|\mathcal V\rangle,|\mathcal R\rangle\}\ \text{is one pair's registration register}\ [\mathrm{Def\text{-}BM.18}(1)]:
\]

%dsp-level:0
\[
\textbf{the two are different objects with different involutions:}
\]

%dsp-level:0
\[
\qquad \hat\omega^2=\mathrm{id}\ \text{with}\ \operatorname{spec}\{\pm1\}\ \text{and}\ \operatorname{tr}\hat\omega=0\ [\mathrm{FT.122}(1)];
\]

%dsp-level:0
\[
\qquad \hat J_\pm^2=0\ \text{with}\ \operatorname{tr}\hat J_\pm=0\ \text{and}\ \neg\text{diagonalizable}\ [(1)]:
\]

%dsp-level:0
\[
\boxed{\ \neg\exists\ \text{an isomorphism}\ \mathcal J\cong\mathcal H_\omega\ \textbf{that respects the two operator algebras}\ —\ \hat\omega\ \text{is an involution and}\ \hat J_\pm\ \text{are nilpotent}\ }
\]

%dsp-level:0
\[
\qquad [\text{machine:}\ \hat\omega^2=\mathrm{id}\ne0\ \text{against}\ \hat J_\pm^2=0];
\]

%dsp-level:0
\[
\therefore\ \textbf{the crossing operator's type is}\quad\hat J_\pm\otimes\mathrm{id}\otimes\mathrm{id}\ \ \text{on}\ \ \mathcal T_{\mathcal J}(c)=\mathcal J\otimes\mathcal H_\omega(c)\otimes\Psi(c)
\]

%dsp-level:0
\[
\qquad [\mathrm{Def\text{-}BM.18}(4);\ \text{machine}];
\]

%dsp-level:0
\[
(7)\quad\textbf{the forward-only row:}
\]

%dsp-level:0
\[
\qquad \textbf{the move system is bidirectional as a set of moves and has a stationary measure}\ [(3)];
\]

%dsp-level:0
\[
\qquad \textbf{the tick is not bidirectional:}\quad\neg\exists\,\tau^{-1}\ [\mathrm{A6}];
\]

%dsp-level:0
\[
\bigl|\mathfrak L(\tau_{n+1})\bigr|=\bigl|\mathfrak L(\tau_n)\bigr|+n_{\mathrm{real\_out}}(\tau_n)\ge\bigl|\mathfrak L(\tau_n)\bigr|\ [\mathrm{A6};\ \mathrm{A7}];\qquad V'=\bigl(1-\tfrac\gamma4\bigr)^2V\le V\ [\mathrm{L2}];
\]

%dsp-level:0
\[
\boxed{\ \textbf{the ensemble sums both moves; the ledger records only growth}\ —\ \text{the formal mirror exists and the dynamical reversal does not, on the Janus register exactly as on the ledger}\ }
\]

%dsp-level:0
\[
\qquad [\mathrm{BM.25}(2);\ \mathrm{Def\text{-}BM.7}(2);\ \mathrm{FT.144}(4)];
\]

%dsp-level:0
\[
(8)\quad\textbf{status:}\qquad\textbf{derived:}\ (1)\text{--}(7);
\]

%dsp-level:0
\[
\neg\text{a declaration},
\]

%dsp-level:0
\[
\quad\neg\text{a P4 addition},
\]

%dsp-level:0
\[
\quad\neg\text{a new dimensionful scale}\ [\mathrm{P1}],
\]

%dsp-level:0
\[
\quad\neg\text{a measured value}\ [\mathrm{P5}],
\]

%dsp-level:0
\[
\quad\neg\text{a numerical coupling}\ [\mathrm{EM.25}(7)],
\]

%dsp-level:0
\[
\quad\neg\text{a rate datum in external units}\ [\mathrm{BM.44}(5)],
\]

%dsp-level:0
\[
\quad\neg\text{a species identification}\ [\mathrm{Def\text{-}BM.3}(7)];
\]

\textbf{Req.} A4, A6, A7, D0.4, D1.2, D1.8 (Part 1); CT0.8, EM.25(7) (Part 1, Part 2); P1, P3, P4, P5 (Parts 1–2); L2 (Part 1); Def-H.32, Def-H.33, FT.108(4), FT.108(5), FT.122(1), FT.144(4) (Part 3); Def-BM.3(7), Def-BM.7(2), Def-BM.18(1), Def-BM.18(4), BM.19(2), BM.23(2), BM.25(2), BM.44(1), BM.44(5), BM.109(2), BM.113(2) (this part); exact $2\times2$ matrix algebra over $\mathbb Z[x_+,x_-]$ and $\mathbb Q$, the SOS parabola completion, AM–GM on exact rationals with its equality locus, exhaustive rational grid censuses [math; machine: \texttt{checks/\allowbreak{}verify\_\allowbreak{}H\_\allowbreak{}janus\_\allowbreak{}tunnel\_\allowbreak{}concrete\_\allowbreak{}spin.py} §A].

\textbf{Pf.}

%dsp-level:0
\[
(1):\ \hat J_+=\left(\begin{smallmatrix}0&1\\0&0\end{smallmatrix}\right),\ \hat J_-=\left(\begin{smallmatrix}0&0\\1&0\end{smallmatrix}\right)\ \text{in the}\ (\mathcal V,\mathcal R)\ \text{row convention}\ [\mathrm{Def\text{-}BM.18}(2)]:
\]

%dsp-level:0
\[
\hat J_+^2=\left(\begin{smallmatrix}0&1\\0&0\end{smallmatrix}\right)^{\!2}=0,\qquad\hat J_-^2=\left(\begin{smallmatrix}0&0\\1&0\end{smallmatrix}\right)^{\!2}=0;\qquad\hat J_+\hat J_-=\left(\begin{smallmatrix}1&0\\0&0\end{smallmatrix}\right)=P_{\mathcal V},\qquad\hat J_-\hat J_+=\left(\begin{smallmatrix}0&0\\0&1\end{smallmatrix}\right)=P_{\mathcal R};
\]

%dsp-level:0
\[
\hat J_+\hat J_-+\hat J_-\hat J_+=P_{\mathcal V}+P_{\mathcal R}=\mathrm{id}\ [\text{machine: all five identities entrywise over}\ \mathbb Q]
\]

%dsp-level:0
\[
(2):\ x_+(\hat J_+-P_{\mathcal V})+x_-(\hat J_--P_{\mathcal R})\ \text{expanded entrywise gives}\ \left(\begin{smallmatrix}-x_+&x_+\\x_-&-x_-\end{smallmatrix}\right)=Q\ [\text{machine: six rational witnesses}];\qquad\text{the two row sums}\ [\text{machine}]
\]

%dsp-level:0
\[
(3):\ (x_-,x_+)Q=(0,0)\ [\text{machine: exact in}\ \mathbb Z[x_+,x_-]];
\]

%dsp-level:0
\[
\qquad \det Q=x_+x_--x_+x_-=0;
\]

%dsp-level:0
\[
\qquad Q\ne0\ \text{at}\ x_+x_->0\implies\operatorname{rk}=1\ [\text{machine}]
\]

%dsp-level:0
\[
p=\tfrac{x_+}{x_++x_-}=\tfrac{1-4\varepsilon_q}{1-4\varepsilon_q+4\varepsilon_q}=1-4\varepsilon_q\ [\text{machine: exact at seven thresholds}];
\]

%dsp-level:0
\[
\qquad \text{detailed balance}: p_{\mathcal V}x_+=\tfrac{x_-x_+}{x_++x_-}=p_{\mathcal R}x_-\ [\text{machine}]
\]

%dsp-level:0
\[
(4):\ \lambda^2-(\operatorname{tr}Q)\lambda+\det Q=\lambda^2+(x_++x_-)\lambda=\lambda(\lambda+x_++x_-)\ [\text{machine}];
\]

%dsp-level:0
\[
\qquad \text{at}\ x_++x_-=1:\ \{0,-1\}\ [\text{machine}]
\]

%dsp-level:0
\[
Q(x_-,x_+)
\]

%dsp-level:0
\[
= (0,0)\ \text{and}\ Q(1,-1)
\]

%dsp-level:0
\[
= (-(x_++x_-),\ (x_++x_-))
\]

%dsp-level:0
\[
= -(x_++x_-)(1,-1)\ [\text{machine: the two eigenvectors}]
\]

%dsp-level:0
\[
(5):\ x_+x_-=4\varepsilon_q(1-4\varepsilon_q)=\tfrac14-16(\varepsilon_q-\tfrac18)^2\ [\text{machine: expansion}];
\]

%dsp-level:0
\[
\qquad \text{the}\ 256\text{-point grid:}\ \le\tfrac14,\ =\tfrac14\ \text{only at}\ \varepsilon_q=\tfrac18\ [\text{machine}]
\]

%dsp-level:0
\[
\text{AM--GM}:\ x_+x_-\le\bigl(\tfrac{x_++x_-}2\bigr)^2
\]

%dsp-level:0
\[
\iff 0
\]

%dsp-level:0
\[
\le (x_+-x_-)^2\ [\text{machine: the square is nonnegative on exact rationals}],
\]

%dsp-level:0
\[
\ \text{equality}
\]

%dsp-level:0
\[
\iff x_+
\]

%dsp-level:0
\[
= x_-
\]

%dsp-level:0
\[
\iff 1-4\varepsilon_q
\]

%dsp-level:0
\[
= 4\varepsilon_q\iff\varepsilon_q
\]

%dsp-level:0
\[
= \tfrac18\ [\text{machine}]
\]

%dsp-level:0
\[
(6):\ \hat\omega=\left(\begin{smallmatrix}0&1\\1&0\end{smallmatrix}\right),\ \hat\omega^2=\mathrm{id},\ \operatorname{tr}\hat\omega=0,\ \operatorname{spec}=\{\pm1\}\ [\mathrm{Def\text{-}H.33};\ \mathrm{FT.122}(1);\ \text{machine}];\qquad\hat J_\pm^2=0\ne\mathrm{id}\ [(1);\ \text{machine}]\implies\neg\text{an algebra isomorphism}
\]

%dsp-level:0
\[
(7):\ \mathrm{A6},\ \mathrm{A7},\ \mathrm{L2},\ \mathrm{BM.25}(2),\ \mathrm{FT.144}(4)\ \text{rows re-cited};
\]

%dsp-level:0
\[
\qquad (8):\ \text{the standing-rule scans}\ [\text{machine}].\ \blacksquare
\]

\textbf{Phys.} The crossing has an algebra and it is small enough to write down completely. Each move squares to zero, because the crystallization criterion is binary and a pair cannot crystallize twice without dissolving in between; the two moves anticommute to the identity, and their two products are exactly the two round trips. The generator is the sum of a jump term and a loss term for each move, its rows sum to zero, and it is therefore a master generator rather than a unitary — the crossing is a stochastic ledger move, not a rotation of a state. Its kernel is the line spanned by the two weights, which is where the stationary measure comes from, and detailed balance is an identity rather than an assumption: the flux out of the virtual state equals the flux out of the real one at that measure, because both are the same product. The spectrum is zero and minus the total weight, so at unit total weight there is one stationary mode and one relaxation mode with rate exactly one. The round-trip weight is the familiar parabola with maximum one quarter, and the maximum sits exactly where the two moves balance — an AM–GM equality, not a coincidence. Two clarifications keep the construction honest. The register is not the orientation doublet: one carries an involution, the other carries nilpotents, and no isomorphism respects both algebras, so the crossing operator acts on the register tensored with the chain's own spaces. And the bidirectionality is in the move set, not in the tick: the ensemble sums both moves and has a stationary measure, while the ledger only grows and the reversed history remains inadmissible at its first event.

\end{theorem}

\begin{theorem}[{\texorpdfstring{BM.121 — The zig-zag world line: the even resolvent, the round-trip weights, and no backward leg}{BM.121 - The zig-zag world line: the even resolvent, the round-trip weights, and no backward leg}}]\label{tag:BM-121}
\textbf{Thm.}

%dsp-level:0
\[
(1)\quad\textbf{the vacuum-to-vacuum matrix elements are even:}
\]

%dsp-level:0
\[
\qquad \bigl(\hat U^{2m}\bigr)_{\varnothing\varnothing}=(x_+x_-)^{m},
\]

%dsp-level:0
\[
\qquad \bigl(\hat U^{2m+1}\bigr)_{\varnothing\varnothing}=0
\]

%dsp-level:0
\[
\qquad [\mathrm{FT.108}(5)\ \mathrm{Pf};\ \text{machine: coefficient-wise in}\ \mathbb Q[w,x_+,x_-]/(w^{12})]:
\]

%dsp-level:0
\[
\textbf{the mechanism:}\quad\text{a return to}\ \varnothing\ \text{requires an equal number of}\ \hat J_+\ \text{and}\ \hat J_-\ \text{moves, and}\ \hat J_\pm^2
\]

%dsp-level:0
\[
= 0\ \text{forces strict alternation}\ [\mathrm{BM.120}(1);\ \text{machine}];
\]

%dsp-level:0
\[
(2)\quad\textbf{the resolvent is the zig-zag sum:}
\]

%dsp-level:0
\[
\qquad \boxed{\ \sum_{n\ge0}\langle\varnothing|\hat U^{n}|\varnothing\rangle\,w^{n}\ =\ \sum_{m\ge0}\bigl(w^2x_+x_-\bigr)^{m}\ =\ \frac{1}{1-w^2x_+x_-}\ \ \text{in}\ \mathbb Q[w,x_+,x_-]/(w^{12})\ }
\]

%dsp-level:0
\[
[\mathrm{FT.108}(5);\ \mathrm{BM.23}(7);\ \mathrm{BM.44}(1);\ \text{machine:}\ (1-w^2x_+x_-)\cdot\sum\equiv1\ \text{coefficient-wise at five exact values of}\ x_+x_-\in\{0,\tfrac3{16},\tfrac14,\tfrac12,1\}];
\]

%dsp-level:0
\[
\textbf{every odd coefficient is exactly zero:}
\]

%dsp-level:0
\[
\qquad [w^{2m+1}]\,\frac1{1-w^2x_+x_-}=0\ \ \forall m
\]

%dsp-level:0
\[
\qquad [\text{machine: the coefficient census to}\ w^{11}];
\]

%dsp-level:0
\[
(3)\quad\textbf{the zig-zag is a sum over round trips, and every leg is forward:}
\]

%dsp-level:0
\[
\qquad \textbf{the}\ m\textbf{-th term}\ (w^{2}x_+x_-)^{m}\ \textbf{is}\ m\ \textbf{create--propagate--annihilate round trips}:
\]

%dsp-level:0
\[
\varnothing\ \xrightarrow{\ \hat J_+\ }\ \{\mathcal R\}\ \xrightarrow{\ \hat J_-\ }\ \varnothing\ \xrightarrow{\ \hat J_+\ }\ \{\mathcal R\}\ \xrightarrow{\ \hat J_-\ }\ \varnothing\ \cdots
\]

%dsp-level:0
\[
\qquad \text{with each arrow a step of the ledger order}\ \tau_n\to\tau_{n+1}
\]

%dsp-level:0
\[
\qquad [\mathrm{A6};\ \mathrm{D1.8}]:
\]

%dsp-level:0
\[
\boxed{\ \textbf{zig-zag}\ :=\ \text{the}\ m\textbf{-round-trip word in}\ \{\hat J_+,\hat J_-\},\ \ \text{weight}\ (x_+x_-)^{m},\ \ \text{with}\ \neg\textbf{a backward leg anywhere in it}\ }
\]

%dsp-level:0
\[
(4)\quad\textbf{the exact weight ladder:}\qquad\mathcal W_m=(x_+x_-)^{m}=\bigl[4\varepsilon_q(1-4\varepsilon_q)\bigr]^{m}\qquad[\text{machine}]:
\]

%dsp-level:0
\[
\begin{array}{c|c|c|c}\varepsilon_q & x_+x_- & \mathcal W_1 & \mathcal W_3\\\hline0 & 0 & 0 & 0\\\tfrac1{16} & \tfrac3{16} & \tfrac3{16} & \tfrac{27}{4096}\\\tfrac18 & \tfrac14 & \tfrac14 & \tfrac1{64}\\\tfrac3{16} & \tfrac3{16} & \tfrac3{16} & \tfrac{27}{4096}\\\tfrac14 & 0 & 0 & 0\end{array}\qquad[\text{machine: exact rationals}];
\]

%dsp-level:0
\[
\boxed{\ \mathcal W_m\ \textbf{is maximal at}\ \varepsilon_q=\tfrac18\ \textbf{for every}\ m,\ \textbf{and}\ \mathcal W_m\le4^{-m}\ \textbf{with equality exactly there}\ }
\]

%dsp-level:0
\[
\qquad [\mathrm{BM.120}(5);\ \text{machine: the monotone power of the parabola}];
\]

%dsp-level:0
\[
\textbf{the two endpoints annihilate the whole ladder:}
\]

%dsp-level:0
\[
\qquad \varepsilon_q=0\implies\mathcal W_m=0\ \forall m\ge1;
\]

%dsp-level:0
\[
\qquad \varepsilon_q\to\tfrac14^-\implies\mathcal W_m\to0\ \forall m\ge1
\]

%dsp-level:0
\[
\qquad [\mathrm{BM.109}(2);\ \text{machine}];
\]

%dsp-level:0
\[
(5)\quad\textbf{the pole and its enclosure:}
\]

%dsp-level:0
\[
\qquad w^2x_+x_-=1\implies|w|=\frac1{\sqrt{x_+x_-}}\ge2,
\]

%dsp-level:0
\[
\qquad \text{equality}\iff\varepsilon_q=\frac18
\]

%dsp-level:0
\[
\qquad [x_+x_-\le\tfrac14;\ \mathrm{BM.23}(7);\ \mathrm{BM.44}(1);\ \text{machine}]:
\]

%dsp-level:0
\[
\boxed{\ \textbf{the zig-zag series converges for}\ |w|<2\ \textbf{at every admissible threshold, and its pole never enters the unit disc}\ }
\]

%dsp-level:0
\[
(6)\quad\textbf{no backward leg, on three independent grounds:}\qquad(\mathrm i)\ \ \neg\exists\,\tau^{-1}\qquad[\mathrm{A6}];
\]

%dsp-level:0
\[
(\mathrm{ii})\quad\text{the pure order reversal of every completed nonempty admissible history leads with}\ J^-\ \text{and is inadmissible at its first event}
\]

%dsp-level:0
\[
\qquad [\mathrm{FT.144}(4)(\mathrm{ii});\ \text{machine: exhaustive}];
\]

%dsp-level:0
\[
(\mathrm{iii})\quad\Theta_{\mathfrak L}\ \textbf{fixes every completed alternating word pointwise:}
\]

%dsp-level:0
\[
\qquad \Theta_{\mathfrak L}\bigl((J^+J^-)^{m}\bigr)=(J^+J^-)^{m}\ \ \forall m
\]

%dsp-level:0
\[
\qquad [\mathrm{BM.25}(2);\ \text{machine: all}\ m\le6]:
\]

%dsp-level:0
\[
\boxed{\ \textbf{the reversed reading of a zig-zag is the same word}\ —\ \text{so the diagram's backward leg needs no inverse tick, because the corpus's own mirror already fixes it}\ }
\]

%dsp-level:0
\[
(7)\quad\textbf{the CP reading of the second leg:}
\]

%dsp-level:0
\[
\qquad \text{the member created by}\ \hat J_+\ \text{and the member annihilated by}\ \hat J_-\ \text{in one round trip are a CP-pair:}\quad\{c,\ c^{\mathrm{CP}}\},
\]

%dsp-level:0
\[
\qquad Q=q+(-q)=0
\]

%dsp-level:0
\[
\qquad [\mathrm{Def\text{-}BM.15}(3);\ \mathrm{BM.109}(1);\ \text{machine}]:
\]

%dsp-level:0
\[
\textbf{the second leg carries the conjugate configuration forward,}\ \neg\textbf{the original configuration backward:}
\]

%dsp-level:0
\[
\qquad |w(c^{\mathrm{CP}})|^2=|w(c)|^2,
\]

%dsp-level:0
\[
\qquad w(c^{\mathrm{CP}})=\mu_{\mathrm{CP}}\,w(c),
\]

%dsp-level:0
\[
\qquad \mu_{\mathrm{CP}}=e^{i\pi\Theta}=1\ \ \text{on}\ \Theta\in2\mathbb Z
\]

%dsp-level:0
\[
[\mathrm{BM.108}(2),(3);\ \mathrm{EM.24}(5);\ \text{machine}];
\]

%dsp-level:0
\[
\qquad \boxed{\ \textbf{the zig-zag is}\ \text{chain}\to\text{CP-conjugate chain, both forward}\ }
\]

%dsp-level:0
\[
\qquad [\mathrm{BM.129}(1)\text{'s boson row is the same statement read on the winding}];
\]

%dsp-level:0
\[
(8)\quad\textbf{the ledger-count row:}\qquad\bigl|\mathfrak L(\tau_{n+1})\bigr|-\bigl|\mathfrak L(\tau_n)\bigr|=n_{\mathrm{real\_out}}(\tau_n)\ge0\qquad[\mathrm{A6}]:
\]

%dsp-level:0
\[
\text{a round trip returns}\ n_{\mathcal R}\ \text{to its starting value while the record grows by the events it logged}
\]

%dsp-level:0
\[
\qquad [\text{machine:}\ \hat J_+\hat J_-=P_{\mathcal V}\ \text{restores the register,}\ \mathrm{A6}\ \text{keeps the record}]:
\]

%dsp-level:0
\[
\boxed{\ \textbf{the zig-zag is closed on the register and non-decreasing on the record}\ —\ \text{which is exactly what a forward-only ledger can support}\ }
\]

%dsp-level:0
\[
(9)\quad\textbf{status:}\qquad\textbf{derived:}\ (1)\text{--}(8);
\]

%dsp-level:0
\[
\textbf{re-scoped at its own locus,}\ \neg\textbf{retracted:}\quad\mathrm{BM.109}(2)\text{'s}\ \mathcal W_{\text{pair}}
\]

%dsp-level:0
\[
= x_+x_-\ \textbf{is the round-trip weight}\ —\ \text{the one-way weights are}\ x_+\ \text{and}\ x_-\ \text{separately}\ [(4);\ \mathrm{BM.120}(5)];
\]

%dsp-level:0
\[
\neg\text{a declaration},
\]

%dsp-level:0
\[
\quad\neg\text{a P4 addition},
\]

%dsp-level:0
\[
\quad\neg\text{a new dimensionful scale}\ [\mathrm{P1}],
\]

%dsp-level:0
\[
\quad\neg\text{a measured value}\ [\mathrm{P5}],
\]

%dsp-level:0
\[
\quad\neg\text{a species identification}\ [\mathrm{Def\text{-}BM.3}(7)],
\]

%dsp-level:0
\[
\quad\neg\text{a cosmological coordinate}\ [\mathrm{Def\text{-}BM.17}(6)(\mathrm i)];
\]

%dsp-level:0
\[
\neg\text{a world line in a spacetime is constructed: the zig-zag is a word in two moves on a finite register}\ [\mathrm{Def\text{-}BM.18}(5),(6)];
\]

\textbf{Req.} A6, A7, D1.2, D1.8 (Part 1); P1, P4, P5 (Parts 1–2); Def-H.32, FT.108(5), FT.144(4) (Part 3); Def-BM.3(7), Def-BM.15(3), Def-BM.17(6), Def-BM.18(5), Def-BM.18(6), BM.23(7), BM.25(2), BM.44(1), BM.108(2), BM.108(3), BM.109(1), BM.109(2), BM.120(1), BM.120(5), BM.129(1), EM.24(5) (this part); exact truncated power-series arithmetic in $\mathbb Q[w,x_+,x_-]/(w^{12})$, exact rational weight ladders, exhaustive ledger-record enumeration [math; machine: \texttt{checks/\allowbreak{}verify\_\allowbreak{}H\_\allowbreak{}janus\_\allowbreak{}tunnel\_\allowbreak{}concrete\_\allowbreak{}spin.py} §B].

\textbf{Pf.}

%dsp-level:0
\[
(1):\ \mathrm{FT.108}(5)\text{'s Pf row re-cited};
\]

%dsp-level:0
\[
\qquad \text{strict alternation}: \hat J_+^2
\]

%dsp-level:0
\[
= \hat J_-^2
\]

%dsp-level:0
\[
= 0\ [\mathrm{BM.120}(1);\ \text{machine}]\implies\text{a word returning to}\ \varnothing\ \text{has equal counts of the two moves}
\]

%dsp-level:0
\[
(2):\ (1-w^2x_+x_-)\sum_{m=0}^{5}(w^2x_+x_-)^m
\]

%dsp-level:0
\[
= 1-(w^2x_+x_-)^6
\]

%dsp-level:0
\[
\equiv 1\ \bmod\ w^{12}\ [\text{machine: coefficient-wise at five exact values of}\ x_+x_-];
\]

%dsp-level:0
\[
\qquad [w^{2m+1}]=0\ \forall m\ [\text{machine: the census to}\ w^{11}]
\]

%dsp-level:0
\[
(3):\ \text{each factor}\ w^2x_+x_-\ \text{is one}\ \hat J_+\hat J_-\ \text{pair}\ [(1);\ \mathrm{Def\text{-}BM.18}(5)];
\]

%dsp-level:0
\[
\qquad \text{each move is a ledger step}\ [\mathrm{D1.2};\ \mathrm{A6}]
\]

%dsp-level:0
\[
(4):\ \mathcal W_m=(x_+x_-)^m\ [\text{machine: exact rationals at the five thresholds}];
\]

%dsp-level:0
\[
\qquad x_+x_-\le\tfrac14\ [\mathrm{BM.120}(5)]\implies(x_+x_-)^m\le4^{-m}\ \text{with equality}\iff x_+=x_-\ [\text{machine: monotone power}]
\]

%dsp-level:0
\[
\varepsilon_q=0\implies x_+x_-=0\implies0^m=0\ \forall m\ge1;\qquad\varepsilon_q\to\tfrac14^-\implies x_+\to0\implies x_+x_-\to0\ [\text{machine}]
\]

%dsp-level:0
\[
(5):\ |w|^2x_+x_-=1\implies|w|^2=1/(x_+x_-)\ge4\implies|w|\ge2\ [\text{machine:}\ x_+x_-\le\tfrac14];
\]

%dsp-level:0
\[
\qquad \text{equality at}\ \varepsilon_q=\tfrac18\ [(4)]
\]

%dsp-level:0
\[
(6):\ (\mathrm i)\ \mathrm{A6}\ \text{verbatim};
\]

%dsp-level:0
\[
\qquad (\mathrm{ii})\ \mathrm{FT.144}(4)(\mathrm{ii})\ \text{re-cited}\ [\text{machine: exhaustive}];
\]

%dsp-level:0
\[
\qquad (\mathrm{iii})\ \Theta_{\mathfrak L}((J^+J^-)^m)=(J^-J^+)^m\ \text{reversed}\ =(J^+J^-)^m\ [\mathrm{BM.25}(2);\ \text{machine: all}\ m\le6]
\]

%dsp-level:0
\[
(7):\ \mathrm{Def\text{-}BM.15}(3),\ \mathrm{BM.109}(1)\ \text{re-cited};\qquad\mu_{\mathrm{CP}}=e^{i\pi\Theta}=1\ \text{at}\ \Theta\in2\mathbb Z\ [\mathrm{BM.108}(3);\ \text{machine}]
\]

%dsp-level:0
\[
(8):\ \mathrm{A6}\text{'s increment row re-cited};
\]

%dsp-level:0
\[
\qquad \hat J_+\hat J_-=P_{\mathcal V}\ \text{restores the register}\ [\mathrm{BM.120}(1);\ \text{machine}];
\]

%dsp-level:0
\[
\qquad (9):\ \text{the standing-rule scans}\ [\text{machine}].\ \blacksquare
\]

\textbf{Phys.} The zig-zag is not drawn, it is summed, and the sum turns out to be even. A return to the vacuum needs as many crystallizations as dissolutions, and since neither move can be applied twice in a row the words strictly alternate — so every odd matrix element vanishes identically and the vacuum-to-vacuum resolvent is a geometric series in the square of the fugacity with ratio the round-trip weight. That is the whole content of the zig-zag: term by term it is a create-propagate-annihilate excursion, weighted by the parabola to the power of the number of excursions, maximal at one eighth of the threshold and exactly zero at both endpoints, so the entire ladder is annihilated by a non-flickering ledger. The pole sits at fugacity two or beyond and never enters the unit disc. Three independent rows say why no leg runs backward: the tick has no inverse, the pure reversal of a completed history is inadmissible at its first event, and — the sharpest of the three — the corpus's own record involution fixes every completed alternating word pointwise, so the reversed reading of a zig-zag is literally the same word. What looks like a backward-propagating member in a diagram is here the CP-conjugate configuration carried forward, with equal modulus and a phase that is exactly one on the materialized locus. And the ledger keeps its own accounting: a round trip restores the register while the record grows by the events it logged, which is precisely the one thing a forward-only ledger can support.

\end{theorem}

\begin{theorem}[{\texorpdfstring{BM.122 — The pair channel: charge balance, the exact threshold, the CP structure, and the two-photon bar}{BM.122 - The pair channel: charge balance, the exact threshold, the CP structure, and the two-photon bar}}]\label{tag:BM-122}
\textbf{Thm.}

%dsp-level:0
\[
(1)\quad\textbf{the channel is charge-balanced exactly:}
\]

%dsp-level:0
\[
\qquad Q_{\text{in}}=0\ \ (\text{the}\ \mathcal V\ \text{sector carries no registered charge}\ [\mathrm{D1.1};\ \mathrm{FT.144}(4)]),
\]

%dsp-level:0
\[
\qquad Q_{\text{out}}=q+(-q)=0
\]

%dsp-level:0
\[
\qquad [\mathrm{Def\text{-}BM.15}(3);\ \mathrm{BM.109}(1);\ \text{machine: exhaustive over}\ q\in\{-1,0,+1\}]:
\]

%dsp-level:0
\[
\boxed{\ Q_{\mathrm{EM}}=\sum_{c\in U}q_{\mathrm{EM}}(c)=0\ \ \text{survives the channel, in every valid universe}\ }
\]

%dsp-level:0
\[
\qquad [\mathrm{EM.11};\ \mathrm{D2.11};\ \mathrm{BM.23}(4)];
\]

%dsp-level:0
\[
\textbf{the compensator is unique:}\qquad\#\{q\in\{-1,0,+1\}:\ q=-q_{\text{in}}\}=1\qquad[\mathrm{BM.23}(4);\ \mathrm{BM.108}(5);\ \text{machine}];
\]

%dsp-level:0
\[
(2)\quad\textbf{the exact threshold:}
\]

%dsp-level:0
\[
\qquad \mathcal T_0(\text{pair})=\frac{E^{\text{pair}}_{\mathrm{int}}\,\mathfrak h}{\ln2}=\frac{2E_{\mathrm{int}}(c)\,\mathfrak h}{\ln2}
\]

%dsp-level:0
\[
\qquad [\mathrm{Def\text{-}BM.17}(3);\ \mathrm{BM.109}(1)]:
\]

%dsp-level:0
\[
\textbf{at the minimal materialized floor}\ E_{\mathrm{int}}(c)=2\ [\mathrm{FT.28};\ \mathrm{BM.13}(3)]:
\]

%dsp-level:0
\[
\qquad \boxed{\ \mathcal T_0(\text{pair})\ =\ \frac{4\,\mathfrak h}{\ln2},\qquad\textbf{materialization}\iff NE\ \ge\ \frac{4\mathfrak h}{\ln2}\ }
\]

%dsp-level:0
\[
[\text{machine: one Laurent monomial};\ \mathrm{BM.115}(2)(\mathrm i)\text{'s}\ n{=}2\ \textbf{minimum, re-derived}];
\]

%dsp-level:0
\[
\qquad \mathcal T(\text{pair};N)=\frac{4\mathfrak h}{N\ln2},
\]

%dsp-level:0
\[
\qquad \mathcal T(\text{pair};r_s,N)=\frac{2r_s}N\ \ \text{at}\ r_s=2\mathfrak h
\]

%dsp-level:0
\[
\qquad [\mathrm{BM.115}(1),(4);\ \text{machine}];
\]

%dsp-level:0
\[
\textbf{the general ladder:}
\]

%dsp-level:0
\[
\qquad \mathcal T_0(\text{pair})\in\Bigl\{\frac{4\mathfrak h}{\ln2},\ \dots,\ \frac{120\mathfrak h}{\ln2}\Bigr\}\ \ \text{on the computed cells, span ratio exactly}\ 30
\]

%dsp-level:0
\[
\qquad [\mathrm{BM.111}(5);\ \mathrm{BM.115}(2)(\mathrm i);\ \text{machine}];
\]

%dsp-level:0
\[
(3)\quad\textbf{the channel weight:}
\]

%dsp-level:0
\[
\qquad \boxed{\ \mathcal W_{\text{channel}}\ =\ x_+^{\,2}\cdot\chi_{\{NE\ge\mathcal T_0\}}\ =\ (1-4\varepsilon_q)^2\cdot\chi_{\{NE\ge\mathcal T_0\}}\ }
\]

%dsp-level:0
\[
\qquad [\text{two one-way crystallization moves};\ \mathrm{Def\text{-}BM.18}(5);\ \text{machine}]:
\]

%dsp-level:0
\[
\textbf{the round-trip completion:}
\]

%dsp-level:0
\[
\qquad \mathcal W_{\text{round}}=x_+x_-=4\varepsilon_q(1-4\varepsilon_q)\le\tfrac14
\]

%dsp-level:0
\[
\qquad [\mathrm{BM.109}(2);\ \mathrm{BM.120}(5);\ \text{machine}];
\]

%dsp-level:0
\[
\boxed{\ \mathcal W_{\text{channel}}>0\iff\varepsilon_q<\tfrac14\ \wedge\ NE\ge\mathcal T_0;\qquad\mathcal W_{\text{channel}}=1\iff\varepsilon_q=0\ \wedge\ NE\ge\mathcal T_0\ }
\]

%dsp-level:0
\[
\textbf{the two limits are opposite and both exact:}\quad\text{at}\ \varepsilon_q
\]

%dsp-level:0
\[
= 0\ \textbf{creation is certain and annihilation impossible}\ (x_+=1,\ x_-=0);
\]

%dsp-level:0
\[
\quad\text{at}\ \varepsilon_q\to\tfrac14^-\ \textbf{creation vanishes}\ (x_+\to0)
\]

%dsp-level:0
\[
\qquad [\text{machine}];
\]

%dsp-level:0
\[
(4)\quad\textbf{the CP structure of the channel — the three derived legs:}\qquad\begin{array}{l|l|l}\text{the inversion} & \text{the derived operator} & \text{the exact rows}\\\hline\text{charge} & \widehat V=E_+\otimes P_L\ \ [\mathrm{BM.76}(1)] & \widehat V^2=0,\ \operatorname{rk}\widehat V=1,\ \dim\ker\widehat V=3,\ [W_3\otimes I_\omega,\widehat V]=\widehat V\ \ [\mathrm{BM.76}(2)\text{--}(4);\ \text{machine over}\ \mathbb Q(i)]\\\text{parity/orientation} & \mu_{\mathrm{CP}}=e^{i\pi\Theta}\ \text{and}\ \hat\omega\ \ [\mathrm{BM.108}(3);\ \mathrm{FT.122}(4)] & \hat\omega\hat\chi\hat\omega^{-1}=-\hat\chi,\ \hat\omega P_L\hat\omega^{-1}=P_R,\ \mu_{\mathrm{CP}}=1\ \text{on}\ \Theta\in2\mathbb Z\ \ [\text{machine}]\\\text{energy} & \text{the local scale against}\ \mathcal T_0\ \ [\mathrm{Def\text{-}BM.17}(3)] & NE\ge\mathcal T_0,\ \ \mathcal T_0=E^{\text{pair}}_{\mathrm{int}}\mathfrak h/\ln2\ \ [\mathrm{BM.115}(1);\ \text{machine}]\end{array}
\]

%dsp-level:0
\[
\textbf{the image vertex:}
\]

%dsp-level:0
\[
\qquad (I\otimes\hat\omega)\,\widehat V\,(I\otimes\hat\omega)^{-1}=E_+\otimes P_R
\]

%dsp-level:0
\[
\qquad [\mathrm{BM.76}(5);\ \text{machine}]:
\]

%dsp-level:0
\[
\qquad \boxed{\ \textbf{the step fires from the left ray and its image from the right ray}\ —\ \text{the two members of the pair occupy opposite rays}\ }
\]

%dsp-level:0
\[
[\mathrm{BM.73}(3);\ \mathrm{BM.74}:\ \text{the two emitted members of a vertex occupy opposite rays};\ \mathrm{A\text{-}BM.7}(b);\ \text{machine}];
\]

%dsp-level:0
\[
(5)\quad\textbf{the nilpotency is the channel's selection rule:}
\]

%dsp-level:0
\[
\qquad \widehat V^2=0\implies\boxed{\ \neg\exists\ \text{an iterated charged step inside one doublet}\ }
\]

%dsp-level:0
\[
\qquad [\mathrm{BM.76}(2);\ \mathrm{BM.23}(3)\text{'s ternary overflow}]:
\]

%dsp-level:0
\[
\textbf{the consequence for the channel:}\quad\text{one charge inversion per member, exactly}\ —\ \text{the pair is}\ (q,\ -q)\ \text{and}\ \neg(q,\ q)\ \text{or}\ (q,\ q')\ \text{with}\ |q-q'|\ge2
\]

%dsp-level:0
\[
\qquad [\text{machine: the}\ 3^2=9\ \text{charge pairs against the balanced subset}];
\]

%dsp-level:0
\[
\#\{(q_1,q_2)\in\{-1,0,+1\}^2:\ q_1+q_2=0\}=3,\qquad\#\{(q_1,q_2):\ q_1+q_2\ne0\}=6\qquad[\text{machine}];
\]

%dsp-level:0
\[
(6)\quad\textbf{the two-photon bar (four independent grounds):}\qquad\textbf{the supply leg of}\ \gamma+\gamma\to\text{pair}\ \textbf{is}\ \neg\textbf{constructed:}
\]

%dsp-level:0
\[
(\mathrm O1)\quad\textbf{no gauged chiral coupling:}\ \text{the gauge blocks act on the}\ 16\text{-dimensional pair-channel register and on slot algebras;}\ \neg\text{a block couples to a chiral matter multiplet}
\]

%dsp-level:0
\[
\qquad [\mathrm{FT.80}(7)(\mathrm a);\ \mathrm{CT0.12};\ \mathrm{CT3.20}];
\]

%dsp-level:0
\[
(\mathrm O2)\quad\textbf{the transit alternative is not an inventory candidate:}\ \text{the gauge coupling of the charge read is deficit}\ 1'\ \text{content}
\]

%dsp-level:0
\[
\qquad [\mathrm{BM.41}(3);\ \mathrm{Def\text{-}BM.8}(8);\ \mathrm{FT.125}(4)];
\]

%dsp-level:0
\[
(\mathrm O3)\quad\textbf{no momentum-space resolvent:}\ \text{plane-wave boundary-channel states and the}\ \Sigma\text{-resolvent are an open interface, blocked pending their own declaration}
\]

%dsp-level:0
\[
\qquad [\mathrm{BM.6}(6)];
\]

%dsp-level:0
\[
(\mathrm O4)\quad\textbf{the transit object carries no winding carrier:}\ \text{the}\ \mathrm{D3.7}\ \text{trivial chain is excluded from}\ \mathrm{Def\text{-}H.38}\text{'s scope, so}\ \neg\text{a}\ j
\]

%dsp-level:0
\[
= |k|/2\ \text{read and}\ \neg\text{an}\ \varepsilon\ \text{read attach to it; its polarization count is}\ N_{\mathrm{pol}}
\]

%dsp-level:0
\[
= 2\ \text{and its helicity is a transit-side read with no identification}
\]

%dsp-level:0
\[
\qquad [\mathrm{D3.7};\ \mathrm{EM.8};\ \mathrm{Def\text{-}H.38}\ \text{Scope}];
\]

%dsp-level:0
\[
\boxed{\ \neg\exists\ \text{a two-transit-quantum supply vertex in the inventory}\implies\textbf{the pair channel's amplitude is}\ \neg\textbf{formed, while the channel itself is}\ }
\]

%dsp-level:0
\[
(7)\quad\textbf{what is formed and what is not:}\qquad\begin{array}{l|l}\textbf{formed} & \neg\textbf{formed}\\\hline\text{the charge balance}\ (1) & \text{a cross-section}\ [\mathrm{BM.6}(6);\ \mathrm{Def\text{-}BM.18}(6)]\\\text{the exact threshold}\ \mathcal T_0\ (2) & \text{an S-matrix element}\ [\mathrm{BM.6}(6)]\\\text{the channel weight}\ x_+^2\chi\ (3) & \text{a Lorentz-invariant phase space}\ [\mathrm{Def\text{-}BM.18}(6)]\\\text{the three CP legs}\ (4) & \text{the two-quanta supply vertex}\ (6)\\\text{the selection rule}\ \widehat V^2=0\ (5) & \text{a measured rate or branching fraction}\ [\mathrm{P5};\ \mathrm{BM.44}(5)]\end{array}
\]

%dsp-level:0
\[
(8)\quad\textbf{status:}\qquad\textbf{derived:}\ (1)\text{--}(5);\qquad\textbf{refuted as formed, with four grounds:}\ (6);
\]

%dsp-level:0
\[
\neg\text{a declaration},
\]

%dsp-level:0
\[
\quad\neg\text{a P4 addition},
\]

%dsp-level:0
\[
\quad\neg\text{a new dimensionful scale}\ [\mathrm{P1}],
\]

%dsp-level:0
\[
\quad\neg\text{a measured value}\ [\mathrm{P5}],
\]

%dsp-level:0
\[
\quad\neg\text{a numerical coupling}\ [\mathrm{EM.25}(7)],
\]

%dsp-level:0
\[
\quad\neg\text{a species identification}\ [\mathrm{Def\text{-}BM.3}(7)];
\]

%dsp-level:0
\[
\neg\text{a right-ray matter line is added: the image vertex of}\ (4)\ \text{is the}\ \hat\omega\text{-image of the same single-arrow locus, in the representation}\ \{0,-1\}\ [\mathrm{BM.96}(4),(6);\ \mathrm{BM.73}(2)(b);\ \mathrm{A\text{-}BM.7}];
\]

\textbf{Req.} A4, A6, D0.8, D1.1, D1.2, D1.8, D2.11, D3.7 (Parts 1–2); CT0.12, CT3.20, EM.8, EM.11, EM.24(5), EM.25(7) (Part 1, Part 2); P1, P4, P5 (Parts 1–2); Def-H.4, Def-H.38, FT.28, FT.44(4), FT.80(7), FT.122(4), FT.125(4), FT.144(4) (Part 3); Def-BM.3(7), Def-BM.8(8), Def-BM.15(3), Def-BM.17(3), Def-BM.18(5), Def-BM.18(6), BM.6(6), BM.13(3), BM.23(3), BM.23(4), BM.41(3), BM.44(5), BM.73(2), BM.73(3), BM.74, BM.76(1), BM.76(2), BM.76(3), BM.76(4), BM.76(5), BM.108(3), BM.108(5), BM.109(1), BM.109(2), BM.111(5), BM.115(1), BM.115(2), BM.115(4), BM.120(5), A-BM.7 (this part); exact $\mathbb Q(i)$ matrix algebra, exhaustive ternary charge-pair enumeration, Laurent-monomial assembly, exact rational weight ladders [math; machine: \texttt{checks/\allowbreak{}verify\_\allowbreak{}H\_\allowbreak{}janus\_\allowbreak{}tunnel\_\allowbreak{}concrete\_\allowbreak{}spin.py} §B].

\textbf{Pf.}

%dsp-level:0
\[
(1):\ q+(-q)=0\ [\text{machine: three cases}];\qquad\mathrm{EM.11},\ \mathrm{D2.11}\ \text{re-cited};\qquad\#\{q:q=-q_{\text{in}}\}=1\ [\text{machine}]
\]

%dsp-level:0
\[
(2):\ E^{\text{pair}}_{\mathrm{int}}=2E_{\mathrm{int}}(c)\ [\mathrm{BM.109}(1)];
\]

%dsp-level:0
\[
\qquad \mathcal T_0=E^{\text{pair}}_{\mathrm{int}}\mathfrak h/\ln2\ [\mathrm{Def\text{-}BM.17}(3)];
\]

%dsp-level:0
\[
\qquad E_{\mathrm{int}}=2\implies\mathcal T_0=4\mathfrak h/\ln2\ [\text{machine}]
\]

%dsp-level:0
\[
\mathcal T=2r_s/N\ \text{at}\ r_s=\tfrac12\mathfrak h E^{\text{pair}}_{\mathrm{int}}=2\mathfrak h\ [\mathrm{BM.115}(4);\ \text{machine}];
\]

%dsp-level:0
\[
\qquad \text{the ladder and its span ratio}\ 30\ [\mathrm{BM.115}(2)(\mathrm i);\ \text{machine}]
\]

%dsp-level:0
\[
(3):\ \text{two one-way crystallization moves, each of weight}\ x_+\ [\mathrm{Def\text{-}BM.18}(5);\ \mathrm{BM.120}(5)];
\]

%dsp-level:0
\[
\qquad x_+^2=(1-4\varepsilon_q)^2\ [\text{machine}]
\]

%dsp-level:0
\[
\varepsilon_q=0\implies x_+=1,\ x_-=0;
\]

%dsp-level:0
\[
\qquad \varepsilon_q\to\tfrac14^-\implies x_+\to0\ [\text{machine: exact at the two endpoints of the admissible range}\ [\mathrm{CT0.8}]]
\]

%dsp-level:0
\[
(4):\ \widehat V=E_+\otimes P_L;
\]

%dsp-level:0
\[
\quad E_+^2=0\implies\widehat V^2=0;\quad\operatorname{rk}\widehat V=\operatorname{rk}E_+\cdot\operatorname{rk}P_L=1;
\]

%dsp-level:0
\[
\quad [W_3\otimes I_\omega,\widehat V]=\widehat V\ \text{on}\ [\sigma_3,E_+]=2E_+,\ W_3=\tfrac12\sigma_3\ [\mathrm{BM.76}(1)\text{--}(4);\ \text{machine over}\ \mathbb Q(i)]
\]

%dsp-level:0
\[
(I\otimes\hat\omega)\widehat V(I\otimes\hat\omega)^{-1}=E_+\otimes\hat\omega P_L\hat\omega^{-1}=E_+\otimes P_R\ [\mathrm{FT.122}(4);\ \mathrm{BM.76}(5);\ \text{machine}]
\]

%dsp-level:0
\[
(5):\ \widehat V^2=0\ [\mathrm{BM.76}(2)];
\]

%dsp-level:0
\[
\qquad \text{the balanced subset of the}\ 9\ \text{charge pairs is}\ \{(-1,+1),(0,0),(+1,-1)\},
\]

%dsp-level:0
\[
\ \#\
\]

%dsp-level:0
\[
= 3\ [\text{machine: exhaustive}]
\]

%dsp-level:0
\[
(6):\ \mathrm{FT.80}(7)(\mathrm a),
\]

%dsp-level:0
\[
\ \mathrm{BM.41}(3),
\]

%dsp-level:0
\[
\ \mathrm{Def\text{-}BM.8}(8),
\]

%dsp-level:0
\[
\ \mathrm{BM.6}(6),
\]

%dsp-level:0
\[
\ \mathrm{D3.7},
\]

%dsp-level:0
\[
\ \mathrm{EM.8},
\]

%dsp-level:0
\[
\ \mathrm{Def\text{-}H.38}\ \text{Scope rows re-cited}\ [\text{machine: the four-ground census}];
\]

%dsp-level:0
\[
\qquad N_{\mathrm{pol}}=(4-1)-1=2\ [\mathrm{EM.8};\ \text{machine}]
\]

%dsp-level:0
\[
(7),(8):\ \text{the two-column assembly and the standing-rule scans}\ [\text{machine}].\ \blacksquare
\]

\textbf{Phys.} The channel that opens is exactly as constrained as the ledger can make it. Charge balances because the pair is a chain and its conjugate, the compensating value is unique in the ternary alphabet, and the nilpotent weak step allows precisely one charge inversion per member — so of the nine ordered charge pairs only the three balanced ones are formed. The threshold is one Laurent monomial: twice the constituent edge count times the single scale over the log of two, which at the digon floor is four units, the very minimum of the two-chain ladder the resonance tier already computed. The weight is the square of the crystallization weight times the indicator that the local scale has reached that threshold, and its two limits are exact opposites — at the sharp threshold creation is certain and annihilation impossible, while near full dissolution creation itself vanishes. The CP structure decomposes into three derived legs: a nilpotent step operator that fires from the occupied ray alone and whose orientation image fires from the other, the CP transport phase that is exactly one on the materialized locus, and the local scale against the threshold. What does not exist is the supply vertex. Four independent grounds bar it: nothing couples a gauge block to a chiral matter multiplet, the transit alternative is deficit-one-prime content, there is no momentum-space resolvent, and the transit object is excluded from the winding carrier's scope so neither a spin label nor a statistics sign attaches to it. So the channel is derived and the amplitude is not, and the tier records both halves rather than promoting one into the other.

\end{theorem}

\begin{theorem}[{\texorpdfstring{BM.123 — The annihilation channel: the reverse move, the entropy destination, and the closure of the cross-section residue}{BM.123 - The annihilation channel: the reverse move, the entropy destination, and the closure of the cross-section residue}}]\label{tag:BM-123}
\textbf{Thm.}

%dsp-level:0
\[
(1)\quad\textbf{the reverse move is the same algebra:}
\]

%dsp-level:0
\[
\qquad \text{annihilation}\ :=\ \hat J_-\ \text{on each of the pair's two registers},
\]

%dsp-level:0
\[
\qquad \text{weight}\ x_-^{\,2}=(4\varepsilon_q)^2
\]

%dsp-level:0
\[
\qquad [\mathrm{Def\text{-}BM.18}(5);\ \mathrm{BM.120}(5);\ \text{machine}]:
\]

%dsp-level:0
\[
\boxed{\ \mathcal W_{\text{ann}}\ =\ x_-^{\,2}\cdot\chi_{\{\text{the pair is materialized}\}}\ =\ 16\varepsilon_q^2\cdot\chi_{\{NE\ge\mathcal T_0\}}\ }
\]

%dsp-level:0
\[
\qquad [\text{machine: one monomial in}\ \varepsilon_q];
\]

%dsp-level:0
\[
\textbf{the charge balance:}
\]

%dsp-level:0
\[
\qquad Q_{\text{in}}=q+(-q)=0,
\]

%dsp-level:0
\[
\qquad Q_{\text{out}}=0\ \ (\text{the}\ \mathcal V\ \text{sector is unregistered})
\]

%dsp-level:0
\[
\qquad [\mathrm{D1.1};\ \mathrm{EM.11};\ \text{machine}]:
\]

%dsp-level:0
\[
\qquad \boxed{\ \text{the balance holds in both directions}\ }
\]

%dsp-level:0
\[
(2)\quad\textbf{the two directions and their exact ratio:}
\]

%dsp-level:0
\[
\qquad \frac{\mathcal W_{\text{ann}}}{\mathcal W_{\text{channel}}}\ =\ \frac{x_-^2}{x_+^2}\ =\ \Bigl(\frac{4\varepsilon_q}{1-4\varepsilon_q}\Bigr)^{\!2}
\]

%dsp-level:0
\[
\qquad [\text{machine: exact rational function}]:
\]

%dsp-level:0
\[
\boxed{\ \frac{\mathcal W_{\text{ann}}}{\mathcal W_{\text{channel}}}=1\iff\varepsilon_q=\frac18,\qquad<1\iff\varepsilon_q<\frac18,\qquad>1\iff\varepsilon_q>\frac18\ }
\]

%dsp-level:0
\[
[\text{machine: the sign census on a}\ 256\text{-point rational grid}];
\]

%dsp-level:0
\[
\qquad \textbf{the ratio is the detailed-balance indicator:}\ \text{equality holds exactly at}\ \mathrm{BM.120}(5)\text{'s maximum of the round-trip weight};
\]

%dsp-level:0
\[
\textbf{the two endpoint rows:}
\]

%dsp-level:0
\[
\qquad \varepsilon_q=0:\ \ \frac{\mathcal W_{\text{ann}}}{\mathcal W_{\text{channel}}}=0\ \ (\text{nothing annihilates});
\]

%dsp-level:0
\[
\qquad \varepsilon_q\to\tfrac14^-:\ \ \mathcal W_{\text{channel}}\to0\ \ \text{and the ratio diverges}
\]

%dsp-level:0
\[
\qquad [\text{machine}];
\]

%dsp-level:0
\[
(3)\quad\textbf{the round trip closes the two directions:}
\]

%dsp-level:0
\[
\qquad \mathcal W_{\text{channel}}\cdot\mathcal W_{\text{ann}}\ \text{is}\ \neg\text{the round-trip weight};
\]

%dsp-level:0
\[
\qquad \textbf{the round trip is one creation and one annihilation:}
\]

%dsp-level:0
\[
\boxed{\ \mathcal W_{\text{round}}\ =\ x_+x_-\ =\ 4\varepsilon_q(1-4\varepsilon_q)\ \le\ \frac14\ }\qquad[\mathrm{BM.109}(2);\ \mathrm{BM.120}(5);\ \mathrm{BM.121}(4);\ \text{machine}];
\]

%dsp-level:0
\[
\textbf{the resolvent that sums them:}
\]

%dsp-level:0
\[
\qquad \frac1{1-w^2x_+x_-},
\]

%dsp-level:0
\[
\qquad \text{pole at}\ |w|\ge2
\]

%dsp-level:0
\[
\qquad [\mathrm{BM.121}(2),(5);\ \mathrm{FT.108}(5);\ \text{machine}];
\]

%dsp-level:0
\[
(4)\quad\textbf{where the content goes (the entropy destination):}
\]

%dsp-level:0
\[
\qquad \text{a dissolved pair returns to the}\ \mathcal V\ \text{sector:}\quad n_{\mathcal R}=0\implies\rho_{pq}=\rho_\circ=\tfrac14I,
\]

%dsp-level:0
\[
\qquad S_{\mathrm{raw}}(\rho_\circ)=\ln4=\max,
\]

%dsp-level:0
\[
\qquad \arg\max=\bigl\{\tfrac14I\bigr\}
\]

%dsp-level:0
\[
[\mathrm{FT.144}(4);\ \mathrm{CT0.1};\ \mathrm{CT0.2};\ \mathrm{A5};\ \text{machine}];\qquad E_{\mathrm{int}}(R)=|R|,\qquad S_{\mathfrak L}=n_{\mathcal R}\ln2\qquad[\mathrm{Def\text{-}BM.5}(4);\ \mathrm{CT1.29}]:
\]

%dsp-level:0
\[
\boxed{\ \textbf{annihilation returns the pair's registered content to the maximally raw-entropic virtual sector}\ —\ \neg\text{a radiation leg carries it, because}\ \neg\text{a supply vertex exists}\ [\mathrm{BM.122}(6)]\ }
\]

%dsp-level:0
\[
\textbf{the ledger row:}\qquad|E_{\mathcal R}|=0\implies\neg\exists\ \text{a fermion terminus}\qquad[\mathrm{D1.4};\ \mathrm{Def\text{-}H.4};\ \mathrm{FT.144}(4);\ \text{machine}];
\]

%dsp-level:0
\[
(5)\quad\textbf{the inter-sector row, re-scoped and not retracted:}
\]

%dsp-level:0
\[
\qquad \mathrm{BM.27}(7),
\]

%dsp-level:0
\[
\ \mathrm{BM.29}(5)\ \text{and}\ \mathrm{Def\text{-}BM.15}(5)\ \text{record}\ \neg\text{an annihilation channel,}\ \neg\text{a pair-production channel and}\ \neg\text{a cross-section between}\ \mathrm{EM}^+\ \text{and}\ \mathrm{EM}^-:
\]

%dsp-level:0
\[
\textbf{the frozen content stands:}
\]

%dsp-level:0
\[
\qquad \neg\exists\ \textbf{a direct}\ \mathrm{EM}^+\to\mathrm{EM}^-\ \textbf{operator};
\]

%dsp-level:0
\[
\qquad \text{the two sectors remain amplitude-equal configuration orbits of one involution}\ [\mathrm{BM.25}(4);\ \mathrm{BM.27}(6);\ |w(c^{\mathrm{CP}})|^2=|w(c)|^2];
\]

%dsp-level:0
\[
\textbf{what is constructed:}
\]

%dsp-level:0
\[
\qquad \boxed{\ \text{a}\ \textbf{mediated}\ \text{channel}\ —\ \mathcal J\text{'s two moves composed with the CP identification of}\ \mathrm{Def\text{-}BM.15}(2)\ —\ \text{connects the two sectors through the}\ \mathcal V\ \textbf{sector, not through a direct operator}\ }
\]

%dsp-level:0
\[
\textbf{the composition, exactly:}
\]

%dsp-level:0
\[
\qquad \varnothing\ \xrightarrow{\ \hat J_+\otimes\hat J_+\ }\ \{c,\ c^{\mathrm{CP}}\}\ \xrightarrow{\ \hat J_-\otimes\hat J_-\ }\ \varnothing,
\]

%dsp-level:0
\[
\qquad \text{weight}\ x_+^2x_-^2\ \ \text{for the full out-and-back on two registers}
\]

%dsp-level:0
\[
[\text{machine: the monomial}];\qquad\neg\text{a cross-section is attached to it}\ [\mathrm{Def\text{-}BM.18}(6);\ \mathrm{BM.6}(6)];
\]

%dsp-level:0
\[
(6)\quad\textbf{the closure of the cross-section residue:}
\]

%dsp-level:0
\[
\qquad \mathrm{BM.118}(7)\ \text{named}\ (\mathrm{R9})\ \text{as}\ \textit{the annihilation channel: an inter-sector transition operator between}\ \mathrm{EM}^+\ \text{and}\ \mathrm{EM}^-,
\]

%dsp-level:0
\[
\ \neg\text{in the inventory}:
\]

%dsp-level:0
\[
\boxed{\ (\mathrm{R9})\ \ \textbf{is CLOSED at the channel level:}\quad\text{a mediated annihilation channel is constructed, with an exact weight}\ x_-^2\chi,\ \text{an exact threshold}\ \mathcal T_0,\ \text{an exact charge balance and an exact entropy destination}\ }
\]

%dsp-level:0
\[
\textbf{the residue narrows and does not vanish:}\qquad\begin{array}{l|l}(\mathrm{R9}')\ \ \text{the annihilation cross-section} & \text{requires the momentum-space}\ \Sigma\text{-resolvent}\ [\mathrm{BM.6}(6)]\ \textbf{and}\ \text{the gauge coupling of the charge read}\ [\mathrm{Def\text{-}BM.8}(8);\ \mathrm{BM.122}(6)]\\(\mathrm{R9}'')\ \ \text{the radiation legs} & \neg\text{a transit-quantum supply vertex exists}\ [\mathrm{BM.122}(6)(\mathrm O1)\text{--}(\mathrm O4)]\end{array}
\]

%dsp-level:0
\[
(7)\quad\textbf{the arrow rows stand:}
\]

%dsp-level:0
\[
\qquad \text{the channel is forward-only:}\ \neg\exists\,\tau^{-1}\ [\mathrm{A6}];
\]

%dsp-level:0
\[
\qquad \neg\exists\ \text{an inverse decay or an inverse creation in the ledger}\ [\mathrm{BM.23}(9);\ \mathrm{BM.44}(6);\ \mathrm{FT.144}(2)]:
\]

%dsp-level:0
\[
\textbf{the reversibility is in the move set and the stationary measure,}\ \neg\textbf{in the tick}
\]

%dsp-level:0
\[
\qquad [\mathrm{BM.120}(3),(7);\ \text{machine}];
\]

%dsp-level:0
\[
\Gamma(\text{pair})=\Gamma_{\mathrm{tr}}(\varepsilon_q)\cdot f\cdot\frac{\varepsilon_{\mathrm{eff}}}{\varepsilon_{\mathrm{EM}}}\cdot\mathcal N\cdot\chi_{\{NE\ge\mathcal T_0\}}
\]

%dsp-level:0
\[
\qquad [\mathrm{Def\text{-}BM.17}(4);\ \mathrm{BM.116}(1):\ \text{the width function applies to the pair resonance unchanged}];
\]

%dsp-level:0
\[
(8)\quad\textbf{status:}
\]

%dsp-level:0
\[
\qquad \textbf{derived:}\ (1)\text{--}(5),\ (7);
\]

%dsp-level:0
\[
\qquad \textbf{closed at the channel level:}\ (6);
\]

%dsp-level:0
\[
\qquad \textbf{named and}\ \neg\textbf{filled:}\ (\mathrm{R9}'),\ (\mathrm{R9}'');
\]

%dsp-level:0
\[
\textbf{re-scoped at their own loci,}\ \neg\textbf{retracted:}\quad\mathrm{BM.27}(7),
\]

%dsp-level:0
\[
\ \mathrm{BM.29}(5),
\]

%dsp-level:0
\[
\ \mathrm{Def\text{-}BM.15}(5),
\]

%dsp-level:0
\[
\ \mathrm{BM.108}(7),
\]

%dsp-level:0
\[
\ \mathrm{BM.109}(5)\ \text{(each frozen on its own locus:}\ \neg\text{a}\ \textbf{direct}\ \text{inter-sector operator exists)};
\]

%dsp-level:0
\[
\neg\text{a declaration},
\]

%dsp-level:0
\[
\quad\neg\text{a P4 addition},
\]

%dsp-level:0
\[
\quad\neg\text{a new dimensionful scale}\ [\mathrm{P1}],
\]

%dsp-level:0
\[
\quad\neg\text{a measured value}\ [\mathrm{P5}],
\]

%dsp-level:0
\[
\quad\neg\text{a numerical coupling}\ [\mathrm{EM.25}(7)],
\]

%dsp-level:0
\[
\quad\neg\text{a lifetime or cross-section in external units}\ [\mathrm{BM.44}(5)],
\]

%dsp-level:0
\[
\quad\neg\text{a species identification}\ [\mathrm{Def\text{-}BM.3}(7)];
\]

\textbf{Req.} A5, A6, D0.8, D1.1, D1.4, D1.8, D2.11 (Part 1); CT0.1, CT0.2, CT0.8, CT1.29, EM.11, EM.25(7) (Part 1, Part 2); P1, P4, P5 (Parts 1–2); Def-H.4, Def-H.32, FT.108(5), FT.144(2), FT.144(4) (Part 3); Def-BM.3(7), Def-BM.5(4), Def-BM.8(8), Def-BM.15(2), Def-BM.15(5), Def-BM.17(3), Def-BM.17(4), Def-BM.18(5), Def-BM.18(6), BM.6(6), BM.23(9), BM.25(4), BM.27(6), BM.27(7), BM.29(5), BM.44(5), BM.44(6), BM.108(7), BM.109(2), BM.109(5), BM.116(1), BM.118(7), BM.120(3), BM.120(5), BM.120(7), BM.121(2), BM.121(4), BM.121(5), BM.122(6) (this part); exact rational function rows, exhaustive rational grid sign censuses, Laurent-monomial assembly [math; machine: \texttt{checks/\allowbreak{}verify\_\allowbreak{}H\_\allowbreak{}janus\_\allowbreak{}tunnel\_\allowbreak{}concrete\_\allowbreak{}spin.py} §B].

\textbf{Pf.}

%dsp-level:0
\[
(1):\ \hat J_-\ \text{on two registers with weight}\ x_-\ \text{each}\ [\mathrm{Def\text{-}BM.18}(5);\ \mathrm{BM.120}(5)];
\]

%dsp-level:0
\[
\qquad x_-^2=(4\varepsilon_q)^2=16\varepsilon_q^2\ [\text{machine}]
\]

%dsp-level:0
\[
Q_{\text{in}}=q+(-q)=0\ [\mathrm{BM.109}(1);\ \text{machine}];\qquad Q_{\text{out}}=0\ [\mathrm{D1.1}:\ \text{virtual pairs are unregistered}]
\]

%dsp-level:0
\[
(2):\ \tfrac{x_-^2}{x_+^2}=\bigl(\tfrac{4\varepsilon_q}{1-4\varepsilon_q}\bigr)^2;
\]

%dsp-level:0
\[
\qquad =1\iff4\varepsilon_q=1-4\varepsilon_q\iff\varepsilon_q=\tfrac18\ [\text{machine: the cleared linear equation}]
\]

%dsp-level:0
\[
\text{the}\ 256\text{-point grid census of the sign}\ [\text{machine}];
\]

%dsp-level:0
\[
\qquad \varepsilon_q=0\implies x_-=0\implies\text{the ratio is}\ 0;
\]

%dsp-level:0
\[
\qquad \varepsilon_q\to\tfrac14^-\implies x_+\to0\implies\text{the ratio diverges}\ [\text{machine}]
\]

%dsp-level:0
\[
(3):\ x_+x_-=4\varepsilon_q(1-4\varepsilon_q)\le\tfrac14\ [\mathrm{BM.120}(5);\ \text{machine}];\qquad x_+^2x_-^2=(x_+x_-)^2\le\tfrac1{16}\ [\text{machine}]
\]

%dsp-level:0
\[
(4):\ \mathrm{FT.144}(4)\text{'s}\ \mathcal V\text{-sector rows re-cited (frozen)};
\]

%dsp-level:0
\[
\qquad S_{\mathfrak L}
\]

%dsp-level:0
\[
= n_{\mathcal R}\ln2\ [\mathrm{CT1.29}]\ \text{and}\ n_{\mathcal R}\to0\ \text{under}\ \hat J_-\ [\mathrm{D1.2}]\implies\text{the registered content leaves while the raw entropy reaches its maximum}\ [\mathrm{CT0.2};\ \mathrm{A5}]
\]

%dsp-level:0
\[
(5):\ \mathrm{BM.25}(4),\ \mathrm{BM.27}(6)\ \text{re-cited};
\]

%dsp-level:0
\[
\qquad |w(c^{\mathrm{CP}})|^2=|w(c)|^2\ [\mathrm{EM.24}(5);\ \mathrm{BM.108}(2)];
\]

%dsp-level:0
\[
\qquad \text{the mediated composition}\ \hat J_+\otimes\hat J_+\ \text{then}\ \hat J_-\otimes\hat J_-\ \text{with weight}\ x_+^2x_-^2\ [\text{machine: the monomial}]
\]

%dsp-level:0
\[
(6):\ \mathrm{BM.118}(7)\text{'s}\ (\mathrm{R9})\ \text{row re-cited};
\]

%dsp-level:0
\[
\qquad \text{the two narrowed residues cite their own tags}\ [\mathrm{BM.6}(6);\ \mathrm{BM.122}(6)]
\]

%dsp-level:0
\[
(7):\ \mathrm{A6},\ \mathrm{BM.23}(9),\ \mathrm{BM.44}(6),\ \mathrm{FT.144}(2)\ \text{re-cited};
\]

%dsp-level:0
\[
\qquad (8):\ \text{the rescope census and the standing-rule scans}\ [\text{machine}].\ \blacksquare
\]

\textbf{Phys.} Annihilation is the same algebra read the other way, and it costs the square of the dissolution weight. The two directions have an exact ratio — four epsilon over one minus four epsilon, squared — which is one precisely at the threshold value one eighth, below one before it and above one after, so the ratio is nothing other than the detailed-balance indicator and it peaks where the round-trip weight does. At the sharp threshold the ratio is zero: nothing annihilates, which is the same statement as the no-flicker theorem seen from the other side. The round trip that joins the two directions carries the familiar parabola, and the resolvent that sums all round trips has its pole at fugacity two or beyond. What happens to the content is exact and slightly austere: a dissolved pair returns to the virtual sector at one quarter of the identity, the unique maximizer of the raw entropy, and the ledger's own entropy read goes to zero with the registered count — so nothing is carried off by radiation, because no supply vertex exists to carry it. That absence is what keeps the frozen rows honest: there is still no direct operator between the two charge sectors and they remain amplitude-equal orbits of one involution, and what the tier adds is a mediated channel through the virtual sector, composed of the two Janus moves and the CP identification. Residue (R9) closes at exactly that level — a channel with a weight, a threshold, a balance and a destination — and narrows into two named successors, the cross-section which needs a momentum-space resolvent and a gauge coupling of the charge read, and the radiation legs which need a transit supply vertex. The arrow is untouched throughout: the reversibility lives in the move set and its stationary measure, never in the tick.

\end{theorem}

\begin{definition}[{\texorpdfstring{Def-BM.19 — The spin-dependent transport defect (the multiplet-ladder read)}{Def-BM.19 - The spin-dependent transport defect (the multiplet-ladder read)}}]\label{tag:Def-BM-19}
\textbf{Def.} A read specification on $\mathrm{Def\text{-}H.37}(2)$'s holonomy defect and $\mathrm{FT.60}(5)$'s multiplet ladder; no new axiom, no new datum, no new parameter, and no identification of the two angles.

%dsp-level:0
\[
(1)\quad\textbf{the two angles, kept apart:}\qquad\begin{array}{l|l|l}\text{angle} & \text{definition} & \text{dimension}\\\hline\Theta(c)=\sum_{e\in E(c)}\widehat{ds}^{\,2}(e) & \text{the winding of a materialized closed chain}\ [\mathrm{FT.22};\ \mathrm{EM.24}(5)] & [\Theta]=\mathfrak q^{+1}\ [\mathrm{FT.55}(1)\ \mathrm{Pf}:\ \Theta\mathfrak h^2=\oint ds^2]\\\Theta_R(c) & \text{the frame-holonomy angle of}\ \mathrm{Def\text{-}H.21}\text{'s spin connection}\ [\mathrm{Def\text{-}H.21}(5),(6)] & [\Theta_R]=\mathfrak q^0\ [\mathrm{Def\text{-}H.21}(6):\ K\,dA\ \text{classical}]\end{array}
\]

%dsp-level:0
\[
k=\frac\Theta2\in\mathbb Z_{\ge0},
\]

%dsp-level:0
\[
\qquad j=\frac{|k|}2,
\]

%dsp-level:0
\[
\qquad [k]=\mathfrak q^{+1},
\]

%dsp-level:0
\[
\qquad [j]=\mathfrak q^{+1}
\]

%dsp-level:0
\[
\qquad [\mathrm{FT.22};\ \mathrm{Def\text{-}H.38};\ \mathrm{FT.55}(1)\ \mathrm{Pf};\ \text{machine: the exponent arithmetic}];
\]

%dsp-level:0
\[
(2)\quad\textbf{the no-go (re-cited, and constitutive for this definition):}
\]

%dsp-level:0
\[
\qquad \boxed{\ \text{the identification}\ \Theta_R(c)=2\pi k(c)\ \ \textbf{is ruled out}\ }
\]

%dsp-level:0
\[
\qquad [\mathrm{FT.55}(1):\ \text{the calibration no-go}]:
\]

%dsp-level:0
\[
\textbf{the ground is dimensional and exact:}
\]

%dsp-level:0
\[
\qquad [\Theta_R]=\mathfrak q^0\ \ \text{against}\ \ [2\pi k]=\mathfrak q^{+1}
\]

%dsp-level:0
\[
\qquad [\text{machine: the two exponent rows disagree}];
\]

%dsp-level:0
\[
\textbf{the corpus's own scope row:}
\]

%dsp-level:0
\[
\qquad \text{under}\ \mathrm{Def\text{-}H.37}(2)\ \text{the defect ladder reads in}\ \Theta_R\ \text{at fixed weight}\ m\ \text{with}\ \neg\text{a}\ k\text{--}\Theta_R\ \text{proportionality anywhere}
\]

%dsp-level:0
\[
\qquad [\mathrm{FT.60}\ \mathrm{Rmk}];
\]

%dsp-level:0
\[
(3)\quad\textbf{the multiplet ladder:}\qquad m\in\{k,\ k-2,\ \dots,\ -k\},\qquad\#m=k+1=2j+1\qquad[\mathrm{FT.60}(5);\ \mathrm{Def\text{-}H.38}]:
\]

%dsp-level:0
\[
\boxed{\ \Delta(m)\ :=\ \bigl\|(\mathcal U^{(k)}-1)\psi_m\bigr\|^2\ =\ \bigl|e^{im\phi}-1\bigr|^2\ =\ 2-2\cos(m\phi)\ =\ 4\sin^2\frac{m\phi}2\ }
\]

%dsp-level:0
\[
\qquad \phi:=\frac{\Theta_R}2
\]

%dsp-level:0
\[
\qquad [\mathrm{FT.60}(5);\ \mathrm{FT.54}(4);\ \mathrm{EW.33}\ \text{pattern};\ \mathrm{Def\text{-}H.21}(6)];
\]

%dsp-level:0
\[
\textbf{constitutive clause:}\ \ \text{the spin-dependence enters through the multiplet weight}\ m,
\]

%dsp-level:0
\[
\ \neg\text{through the winding}\ \Theta\ \text{and}\ \neg\text{through a single value per}\ j;
\]

%dsp-level:0
\[
(4)\quad\textbf{the exact algebraic form:}
\]

%dsp-level:0
\[
\qquad \cos(m\phi)=T_m(\cos\phi),
\]

%dsp-level:0
\[
\qquad T_{m+1}=2x\,T_m-T_{m-1},
\]

%dsp-level:0
\[
\qquad T_0=1,\ T_1=x
\]

%dsp-level:0
\[
\qquad [\text{math: Chebyshev}]:
\]

%dsp-level:0
\[
\boxed{\ \Delta(m)\ =\ 2-2\,T_{|m|}\bigl(\cos\tfrac{\Theta_R}2\bigr)\ }
\]

%dsp-level:0
\[
\textbf{the state-level specialization:}
\]

%dsp-level:0
\[
\qquad \Theta_R=\delta_v\ \ \text{at the vertex-figure locus}\implies\Delta=D_S\in\bigl\{2,\ 4,\ 2-\sqrt2,\ 1,\ 2-2\cos\tfrac\pi{10},\ 2-\sqrt3\bigr\}
\]

%dsp-level:0
\[
\qquad [\mathrm{Def\text{-}BM.14}(3);\ \mathrm{BM.82}(1);\ \mathrm{BM.114}(3)];
\]

%dsp-level:0
\[
(5)\quad\textbf{the standing zero:}
\]

%dsp-level:0
\[
\qquad k\ \text{even}\implies0\in\{k,k-2,\dots,-k\}\implies\boxed{\ \Delta(0)=2-2\cos0=0\ \ \text{exactly, at every value of}\ \Theta_R\ }
\]

%dsp-level:0
\[
[\mathrm{FT.60}(5):\ \text{every even}\ k\ \text{carries a defect-massless weight};\ \text{machine}];
\]

%dsp-level:0
\[
\qquad k\ \text{odd}\implies0\notin\{k,k-2,\dots,-k\}\implies\neg\text{a standing zero}\ [\text{machine}];
\]

%dsp-level:0
\[
(6)\quad\textbf{the mass ladder and its ceiling:}
\]

%dsp-level:0
\[
\qquad m^2_{\mathrm{pole}}(m)=\hat\mu'^2\,\Delta(m),
\]

%dsp-level:0
\[
\qquad 0\le\Delta(m)\le4\implies\boxed{\ 0\le m^2_{\mathrm{pole}}(m)\le4\hat\mu'^2\ }
\]

%dsp-level:0
\[
\qquad [\mathrm{Def\text{-}H.37}(2);\ \mathrm{FT.54}(1);\ \mathrm{FT.76}(3);\ \text{machine}];
\]

%dsp-level:0
\[
\textbf{the two}\ \mathrm{Def\text{-}H.37}\ \textbf{clauses and their matching:}
\]

%dsp-level:0
\[
\qquad (2)\ \text{implies}\ (1)\ \text{at leading order under}
\]

%dsp-level:0
\[
\qquad \boxed{\ \hat\mu'^2\,\frac{\Theta_R^2}4\ =\ \hat\mu^2\,\ell_c^2\,G(c)\ }
\]

%dsp-level:0
\[
\qquad [\mathrm{Def\text{-}H.37}:\ \text{recorded, not imposed}];
\]

%dsp-level:0
\[
(7)\quad\textbf{the residues inherited, not filled:}
\]

%dsp-level:0
\[
\qquad (\mathrm{R2})\ \ \Theta_R
\]

%dsp-level:0
\[
= \delta_v\ \text{is read at the vertex-figure locus, whose ledger realization}\ —\ \text{a clean-enclosure loop through conical vertices}\ —\ \text{is}\ \neg\text{constructed}
\]

%dsp-level:0
\[
\qquad [\mathrm{BM.67}(9),(10);\ \mathrm{BM.82}(6)(\mathrm{R2});\ \text{Census FT.63}(4)\ \mathrm{Rmk}];
\]

%dsp-level:0
\[
(\mathrm{R1})\ \ \rho_{\mathrm B}
\]

%dsp-level:0
\[
= \hat\mu'^2\mathfrak h^2/(\ln2)^2\ \text{is the single cross-currency residue, its slot declared and its value}\ \neg\text{issued}
\]

%dsp-level:0
\[
\qquad [\mathrm{A\text{-}BM.9};\ \mathrm{BM.82}(1),(6)];
\]

%dsp-level:0
\[
(8)\quad\textbf{status and scope:}
\]

%dsp-level:0
\[
\qquad \textbf{derived:}\ (1)\text{'s dimensions},\ (4)\text{'s recursion and specialization},\ (5),\ (6)\text{'s ceiling};
\]

%dsp-level:0
\[
\qquad \textbf{constitutive:}\ (3)\text{'s locus of the spin-dependence};
\]

%dsp-level:0
\[
\qquad \textbf{re-cited:}\ (2);
\]

%dsp-level:0
\[
\textbf{the scope inherited at every use:}\ \ \mathrm{Def\text{-}H.38}\text{'s nonnegative cone and its}\ k
\]

%dsp-level:0
\[
\ge 3\ \text{aspect bar}\ [\neg\text{an aspect claim is made here}];
\]

%dsp-level:0
\[
\qquad \mathrm{Def\text{-}H.37}\text{'s two P4-class residues}\ \hat\mu^2,\ \hat\mu'^2\ [\mathrm{Cal\text{-}F}];
\]

%dsp-level:0
\[
\textbf{the bar:}
\]

%dsp-level:0
\[
\qquad \neg\text{a product of the two}\ \mathrm{Def\text{-}H.37}\ \text{clauses' factors is a derived mass}\ —\ \text{the formed rows are the two clauses and their matching identity}\ (6);
\]

%dsp-level:0
\[
\qquad \neg\text{a}\ k\text{--}\Theta_R\ \text{proportionality is asserted anywhere};
\]

%dsp-level:0
\[
\neg\text{a P4 addition},
\]

%dsp-level:0
\[
\quad\neg\text{a new dimensionful scale}\ [\mathrm{P1}],
\]

%dsp-level:0
\[
\quad\neg\text{a measured value}\ [\mathrm{P5}],
\]

%dsp-level:0
\[
\quad\neg\text{a species identification}\ [\mathrm{Def\text{-}BM.3}(7)],
\]

%dsp-level:0
\[
\quad\neg\text{a currency conversion}\ [\mathrm{BM.83}(3)].
\]

\textbf{Rmk (why the defect is a ladder and not a number).} A single-valued $\Delta(j)$ would require the frame-holonomy angle to be a function of the winding, and that identification is ruled out on a dimension count: the winding carries the free scale with weight minus two while the holonomy angle is dimensionless. What the corpus derives instead is a ladder — the defect read evaluated at each weight of the spin multiplet — so the spin fixes how many values there are ($2j+1$ of them) and which one is exactly zero when $k$ is even, while $\Theta_R$ remains a separate datum whose ledger realization is residue (R2). Reading the ladder as a number is the same category error as reading a count as a strength, and the definition bars it explicitly.

\textbf{Req.} D2.9 (Parts 1–2); CT2.1 (Part 2); P1, P4, P5 (Parts 1–2); Def-H.21(5), Def-H.21(6), Def-H.37, Def-H.38, FT.22, FT.54(1), FT.54(4), FT.55(1), FT.60(5), FT.76(3), EM.24(5) (Part 3); EW.33 (Part 2); Census FT.63(4) (row-wise); Def-BM.3(7), Def-BM.14(3), BM.67(9), BM.67(10), BM.82(1), BM.82(6), BM.83(3), BM.114(3), A-BM.9 (this part); the Chebyshev recursion over $\mathbb Z[x]$, exact radical identities, $\mathfrak q$-exponent arithmetic [math; machine: \texttt{checks/\allowbreak{}verify\_\allowbreak{}H\_\allowbreak{}janus\_\allowbreak{}tunnel\_\allowbreak{}concrete\_\allowbreak{}spin.py} §C].

\textbf{Pf.}

%dsp-level:0
\[
\text{Definition block: the constitutive clause is}\ (3)\text{'s locus of the spin-dependence};
\]

%dsp-level:0
\[
\ \text{the non-derived clauses are}\ \neg\text{present}\ —\ (1)\text{'s dimensions},
\]

%dsp-level:0
\[
\ (2)\text{'s no-go},
\]

%dsp-level:0
\[
\ (4)\text{'s recursion},
\]

%dsp-level:0
\[
\ (5)\text{'s zero and}\ (6)\text{'s ceiling are all derived or re-cited};
\]

%dsp-level:0
\[
\ (7),
\]

%dsp-level:0
\[
(8)\ \text{are residue and scope records}.\ \blacksquare
\]

\textbf{Phys.} Two angles appear in this corpus and they are not the same object: the winding, which counts the line element around a closed chain and carries the free scale, and the frame-holonomy angle, which is dimensionless. Their identification is ruled out by that dimension mismatch alone, so a transport defect that depended on the winding through the holonomy formula cannot be formed. What is formed is better: a ladder. The spin fixes the multiplet, the multiplet has two-j-plus-one weights spaced by two, and the defect read at each weight is two minus two cosines of that weight times half the holonomy angle — equivalently four sine-squared of a quarter, and equivalently an integer polynomial in the cosine through the Chebyshev recursion. Two rows fall out exactly. Whenever the winding is even the ladder contains the weight zero, and the defect at weight zero is identically zero at every curvature flux, which is a standing massless state in the middle of the ladder rather than a tuning. And the whole ladder is bounded by four times the defect scale, the ceiling the pole-side read already carried. The state-level defect values are the ladder read at the vertex-figure angle, which is where the six exact algebraic numbers come from — and that locus, together with the single cross-currency ratio, stays a named residue rather than a filled one.

\end{definition}

\begin{theorem}[{\texorpdfstring{BM.124 — The spin–winding relation and the multiplet defect ladder}{BM.124 - The spin-winding relation and the multiplet defect ladder}}]\label{tag:BM-124}
\textbf{Thm.}

%dsp-level:0
\[
(1)\quad\textbf{the exact spin–winding relation:}
\]

%dsp-level:0
\[
\qquad k=\frac\Theta2\in\mathbb Z_{\ge0}\ \ \text{on shell},
\]

%dsp-level:0
\[
\qquad j=\frac{|k|}2,
\]

%dsp-level:0
\[
\qquad \Psi(c)=\mathrm{Sym}^{|k|}(\mathcal H_{\mathsf v}),
\]

%dsp-level:0
\[
\qquad \dim\Psi=|k|+1
\]

%dsp-level:0
\[
\qquad [\mathrm{FT.22};\ \mathrm{FT.25}(3);\ \mathrm{Def\text{-}H.38};\ \mathrm{EW.32}]:
\]

%dsp-level:0
\[
\boxed{\ \Theta\ =\ 2k\ =\ 4j\qquad\text{on the nonnegative cone},\qquad\Theta\in2\mathbb Z\ \ \text{for every materialized closed chain}\ }
\]

%dsp-level:0
\[
[\text{machine: the exact relation at every}\ k\in\{0,\dots,8\}];\qquad\textbf{the inverse read:}\qquad\boxed{\ j\ =\ \frac{|\Theta|}4\ }
\]

%dsp-level:0
\[
(2)\quad\textbf{the multiplet:}
\]

%dsp-level:0
\[
\qquad m\in\{k,\ k-2,\ \dots,\ -k\},
\]

%dsp-level:0
\[
\qquad \boxed{\ \#m\ =\ k+1\ =\ 2j+1\ }
\]

%dsp-level:0
\[
\qquad [\mathrm{FT.60}(5);\ \mathrm{Def\text{-}H.38};\ \text{machine: exact at every}\ k\le8];
\]

%dsp-level:0
\[
\textbf{the Casimir:}
\]

%dsp-level:0
\[
\qquad \frac k2\Bigl(\frac k2+1\Bigr)=j(j+1)
\]

%dsp-level:0
\[
\qquad [\mathrm{FT.60}(3);\ \text{machine}];
\]

%dsp-level:0
\[
\qquad \textbf{the stretched total:}
\]

%dsp-level:0
\[
\qquad m_{\max}=k=2j
\]

%dsp-level:0
\[
\qquad [\mathrm{FT.60}(6):\ \text{maximal-weight CG on the cone}];
\]

%dsp-level:0
\[
(3)\quad\textbf{the spin-dependent transport defect is the ladder, and the ladder is exact:}
\]

%dsp-level:0
\[
\qquad \boxed{\ \Delta(m)\ =\ 4\sin^2\frac{m\Theta_R}4\ =\ 2-2\cos\frac{m\Theta_R}2\ =\ 2-2\,T_{|m|}\Bigl(\cos\frac{\Theta_R}2\Bigr)\ }
\]

%dsp-level:0
\[
\qquad [\mathrm{Def\text{-}BM.19}(3),(4);\ \mathrm{FT.60}(5);\ \mathrm{FT.54}(4)]:
\]

%dsp-level:0
\[
\textbf{the Chebyshev recursion verified exactly:}
\]

%dsp-level:0
\[
\qquad T_0=1,\quad T_1=x,\quad T_{m+1}=2x\,T_m-T_{m-1}
\]

%dsp-level:0
\[
\qquad [\text{machine: the integer-polynomial recursion to}\ m=8];
\]

%dsp-level:0
\[
\textbf{the exact witness:}
\]

%dsp-level:0
\[
\qquad T_m\bigl(\tfrac12\bigr)=\cos\frac{m\pi}3\ \ \text{on the six-cycle}\ \bigl(1,\ \tfrac12,\ -\tfrac12,\ -1,\ -\tfrac12,\ \tfrac12\bigr)
\]

%dsp-level:0
\[
\qquad [\text{machine: all}\ m\le8];
\]

%dsp-level:0
\[
(4)\quad\textbf{the refutation of a single-valued}\ \Delta(j):\qquad\#\{m\}=2j+1\ge2\ \ \text{for every}\ j\ge\tfrac12\qquad[(2);\ \text{machine}]:
\]

%dsp-level:0
\[
\boxed{\ \textbf{the spin determines the ladder's length,}\ \neg\textbf{a single defect value}\ —\ \Delta(j)\ \textbf{as one number is}\ \neg\textbf{formed}\ }
\]

%dsp-level:0
\[
\textbf{the two anchors:}\qquad\begin{array}{c|c|c|c}k & j & m\ \text{values} & \Delta(m)\\\hline1 & \tfrac12 & \{+1,\ -1\} & 2-2\cos\tfrac{\Theta_R}2\ \ \text{(both weights agree, by evenness)}\\2 & 1 & \{+2,\ 0,\ -2\} & 2-2\cos\Theta_R,\quad\mathbf0,\quad2-2\cos\Theta_R\end{array}\qquad[\mathrm{FT.60}(5);\ \mathrm{BM.106}(1)\text{'s two floor rows};\ \text{machine}];
\]

%dsp-level:0
\[
\boxed{\ \textbf{at}\ j=\tfrac12\ \textbf{the ladder is state-independent and}\ \neg\textbf{split; at}\ j=1\ \textbf{it splits into three with an exact zero in the middle}\ }
\]

%dsp-level:0
\[
\qquad [\mathrm{FT.60}(5):\ k{=}1\ \text{reproduces}\ \mathrm{FT.54}(4)\ \text{and its state-independence is special to the seed}];
\]

%dsp-level:0
\[
(5)\quad\textbf{the two angles are}\ \neg\textbf{proportional, on a dimension count:}
\]

%dsp-level:0
\[
\qquad [k]=\mathfrak q^{+1},
\]

%dsp-level:0
\[
\qquad [\Theta_R]=\mathfrak q^0,
\]

%dsp-level:0
\[
\qquad [2\pi]=\mathfrak q^0
\]

%dsp-level:0
\[
\qquad [\mathrm{FT.55}(1)\ \mathrm{Pf};\ \mathrm{Def\text{-}H.21}(6);\ \text{machine: the exponent rows}]:
\]

%dsp-level:0
\[
\boxed{\ \Theta_R=2\pi k\ \ \textbf{would require}\ \ \mathfrak q^0=\mathfrak q^{+1}:\ \ \textbf{dimensionally impossible, with no measurement used}\ }
\]

%dsp-level:0
\[
\qquad [\mathrm{FT.55}(1):\ \text{the calibration no-go re-cited}];
\]

%dsp-level:0
\[
\textbf{the consequence:}
\]

%dsp-level:0
\[
\qquad \Delta\ \text{is}\ \neg\text{a function of}\ \Theta,\ \text{hence}\ \neg\text{a function of}\ j\ \text{alone};
\]

%dsp-level:0
\[
\qquad \Delta\ \text{is a function of the pair}\ (m,\ \Theta_R)\ \text{with}\ m\ \text{spin-determined and}\ \Theta_R\ \text{independent}\ [\mathrm{Def\text{-}BM.19}(3),(6)];
\]

%dsp-level:0
\[
\textbf{the corpus's own scope row:}
\]

%dsp-level:0
\[
\qquad \neg\text{a}\ k\text{--}\Theta_R\ \text{proportionality anywhere}\ [\mathrm{FT.60}\ \mathrm{Rmk};\ \mathrm{FT.55}(4):\ \text{the winding enters as a copy count,}\ \neg\text{as an angle}];
\]

%dsp-level:0
\[
(6)\quad\textbf{the standing zero:}\qquad k\ \text{even}\iff0\in\{k,k-2,\dots,-k\}\qquad[\text{machine: exhaustive to}\ k=8]:
\]

%dsp-level:0
\[
\boxed{\ \Delta(0)=2-2T_0(\cos\tfrac{\Theta_R}2)=2-2=0\ \ \textbf{exactly, at every}\ \Theta_R,\ \textbf{for every even}\ k\ }\qquad[\mathrm{FT.60}(5);\ \text{machine}];
\]

%dsp-level:0
\[
k\ \text{odd}\implies m\in\{\pm1,\pm3,\dots,\pm k\}\implies\neg\exists\,m=0\implies\neg\text{a standing zero}\qquad[\text{machine}];
\]

%dsp-level:0
\[
\textbf{the parity row that follows:}
\]

%dsp-level:0
\[
\qquad \boxed{\ \textbf{the defect ladder carries an exact massless weight precisely when the exchange sign is}\ +1\ }
\]

%dsp-level:0
\[
\qquad [k\ \text{even}\iff\varepsilon=(-1)^k=+1\ [\mathrm{BM.113}(4)];\ \text{machine: the two parities coincide}];
\]

%dsp-level:0
\[
(7)\quad\textbf{the state-level specialization:}
\]

%dsp-level:0
\[
\qquad \Theta_R=\delta_v=2\pi-d\,\frac{(L-2)\pi}L,
\]

%dsp-level:0
\[
\qquad D_S=4\sin^2\frac{\delta_v}4
\]

%dsp-level:0
\[
\qquad [\mathrm{Def\text{-}BM.14}(3);\ \text{Census FT.63}(4),(6);\ \mathrm{BM.82}(1)]:
\]

%dsp-level:0
\[
\textbf{the six values are the ladder read at six vertex-figure angles,}\ \neg\textbf{six spin values:}
\]

%dsp-level:0
\[
\qquad D_S\in\bigl\{2,\ 4,\ 2-\sqrt2,\ 1,\ 2-2\cos\tfrac\pi{10},\ 2-\sqrt3\bigr\},
\]

%dsp-level:0
\[
\qquad \text{four fields, six minimal polynomials}
\]

%dsp-level:0
\[
[\mathrm{Def\text{-}BM.14}(3);\ \mathrm{BM.114}(3);\ \text{machine re-verified}];
\]

%dsp-level:0
\[
\qquad \textbf{the index is}\ (\mathcal F_g,\beta),\ \neg\ j
\]

%dsp-level:0
\[
\qquad [\mathrm{BM.90}(5):\ \text{the branch-blindness theorem};\ \text{machine: the symbol census carries}\ \neg j];
\]

%dsp-level:0
\[
(8)\quad\textbf{status:}
\]

%dsp-level:0
\[
\qquad \textbf{derived:}\ (1),\ (2),\ (3),\ (6),\ (7);
\]

%dsp-level:0
\[
\qquad \textbf{refuted as formed:}\ (4)\text{'s single-valued}\ \Delta(j),\ (5)\text{'s}\ \Theta_R=2\pi k;
\]

%dsp-level:0
\[
\neg\text{a declaration},
\]

%dsp-level:0
\[
\quad\neg\text{a P4 addition},
\]

%dsp-level:0
\[
\quad\neg\text{a new dimensionful scale}\ [\mathrm{P1}],
\]

%dsp-level:0
\[
\quad\neg\text{a measured value}\ [\mathrm{P5}],
\]

%dsp-level:0
\[
\quad\neg\text{a numerical coupling}\ [\mathrm{EM.25}(7)],
\]

%dsp-level:0
\[
\quad\neg\text{a species identification}\ [\mathrm{Def\text{-}BM.3}(7)];
\]

\textbf{Req.} D2.9 (Parts 1–2); CT2.1, EW.32 (Part 2); P1, P4, P5 (Parts 1–2); Def-H.21(6), Def-H.37(2), Def-H.38, FT.22, FT.25(3), FT.54(4), FT.55(1), FT.55(4), FT.60(3), FT.60(5), FT.60(6) (Part 3); Census FT.63(4), FT.63(6) (row-wise); Def-BM.3(7), Def-BM.14(3), Def-BM.19(3), Def-BM.19(4), Def-BM.19(6), BM.82(1), BM.90(5), BM.106(1), BM.113(4), BM.114(3) (this part); the Chebyshev recursion over $\mathbb Z[x]$ with an exact rational witness, $\mathfrak q$-exponent arithmetic, exhaustive parity censuses [math; machine: \texttt{checks/\allowbreak{}verify\_\allowbreak{}H\_\allowbreak{}janus\_\allowbreak{}tunnel\_\allowbreak{}concrete\_\allowbreak{}spin.py} §C].

\textbf{Pf.}

%dsp-level:0
\[
(1):\ k=\Theta/2\ [\mathrm{FT.22};\ \mathrm{FT.25}(3)];
\]

%dsp-level:0
\[
\qquad j=|k|/2\ [\mathrm{Def\text{-}H.38}]\implies\Theta=2k=2\cdot2j=4j\ [\text{machine: exact at every}\ k\le8];
\]

%dsp-level:0
\[
\qquad \Theta\in2\mathbb Z\ [\mathrm{FT.22}]
\]

%dsp-level:0
\[
(2):\ \#\{k,k-2,\dots,-k\}=k+1\ [\text{machine: the arithmetic-progression count}];
\]

%dsp-level:0
\[
\qquad k+1=2\cdot\tfrac k2+1=2j+1\ [\text{machine}];
\]

%dsp-level:0
\[
\qquad \tfrac k2(\tfrac k2+1)=j(j+1)\ \text{at}\ j=k/2\ [\text{machine}]
\]

%dsp-level:0
\[
(3):\ \|(U^{(k)}-1)\psi_m\|^2
\]

%dsp-level:0
\[
= |e^{im\phi}-1|^2
\]

%dsp-level:0
\[
= (e^{im\phi}-1)(e^{-im\phi}-1)
\]

%dsp-level:0
\[
= 2-e^{im\phi}-e^{-im\phi}
\]

%dsp-level:0
\[
= 2-2\cos(m\phi)\ [\mathrm{FT.60}(5);\ \text{machine over}\ \mathbb Q(i)]
\]

%dsp-level:0
\[
2-2\cos(m\phi)=4\sin^2\tfrac{m\phi}2\ \text{at}\ \phi=\Theta_R/2\ [\text{machine: the half-angle identity}];
\]

%dsp-level:0
\[
\qquad \cos(m\phi)=T_m(\cos\phi)\ \text{with}\ T_{m+1}=2xT_m-T_{m-1}\ [\text{machine: the recursion to}\ m=8]
\]

%dsp-level:0
\[
T_m(\tfrac12)=\cos\tfrac{m\pi}3:
\]

%dsp-level:0
\[
\quad m
\]

%dsp-level:0
\[
= 0..8\mapsto\bigl(1,\tfrac12,-\tfrac12,-1,-\tfrac12,\tfrac12,1,\tfrac12,-\tfrac12\bigr)\ [\text{machine: exact rationals on the six-cycle}]
\]

%dsp-level:0
\[
(4):\ \#\{m\}=k+1\ge2\ \text{for}\ k\ge1\ [\text{machine}];
\]

%dsp-level:0
\[
\qquad k=1:\ m\in\{+1,-1\},\ \Delta(\pm1)=2-2\cos(\pm\tfrac{\Theta_R}2)=2-2\cos\tfrac{\Theta_R}2\ \text{by evenness}\ [\text{machine}]
\]

%dsp-level:0
\[
k=2:\ m\in\{+2,0,-2\},\ \Delta(\pm2)=2-2\cos\Theta_R,\ \Delta(0)=0\ [\text{machine}]
\]

%dsp-level:0
\[
(5):\ \Theta\mathfrak h^2=\oint ds^2\implies[\Theta]=[\mathfrak h^{-2}]=\mathfrak q^{+1}\implies[k]=\mathfrak q^{+1}\ [\mathrm{FT.55}(1)\ \mathrm{Pf};\ \text{machine}];
\]

%dsp-level:0
\[
\qquad [\Theta_R]=\mathfrak h^0=\mathfrak q^0\ [\mathrm{Def\text{-}H.21}(6)];
\]

%dsp-level:0
\[
\qquad \mathfrak q^0\ne\mathfrak q^{+1}\ [\text{machine}]
\]

%dsp-level:0
\[
(6):\ k\ \text{even}\iff\exists r:\ k-2r=0\iff r=k/2\in\mathbb Z\ [\text{machine: exhaustive to}\ k=8];
\]

%dsp-level:0
\[
\qquad \Delta(0)=2-2T_0=2-2\cdot1=0\ [\text{machine}]
\]

%dsp-level:0
\[
\varepsilon=(-1)^k\ [\mathrm{BM.113}(4)];
\]

%dsp-level:0
\[
\qquad k\ \text{even}\iff\varepsilon=+1\iff\Delta(0)=0\ \text{is in the ladder}\ [\text{machine: the two parity censuses coincide}]
\]

%dsp-level:0
\[
(7):\ \mathrm{Def\text{-}BM.14}(3)\text{'s table re-cited};
\]

%dsp-level:0
\[
\qquad \text{the symbol census of}\ D_S\ \text{carries}\ \{\mathcal F_g,\beta,\delta_v,L,d\}\ \text{and}\ \neg\{j,k\}\ [\mathrm{BM.90}(5);\ \text{machine}];
\]

%dsp-level:0
\[
\qquad (8):\ \text{the standing-rule scans}\ [\text{machine}].\ \blacksquare
\]

\textbf{Phys.} The spin and the winding are the same integer read twice, and the relation is exact rather than fitted: the winding is twice the copy number and the spin is half of it, so the winding is four times the spin on the cone, and it is always an even integer for a materialized closed chain. What the spin then controls is not one mass but a ladder — two-j-plus-one weights spaced by two — and the defect at each weight is two minus two cosines of the weight times half the holonomy angle, which is also four sine-squared of a quarter and also an integer polynomial in the cosine by the Chebyshev recursion. At the two anchor windings the difference is visible: at spin one half the ladder has two weights that agree by evenness and nothing splits, while at spin one it splits into three with an exact zero in the middle. That zero is not a tuning but a parity theorem — the ladder contains the weight zero precisely when the winding is even, which is precisely when the exchange sign is plus one, so the defect ladder carries a massless weight exactly on the bosonic side. What the tier refuses is the fusion of the two angles: the winding carries the free scale and the holonomy angle does not, so identifying them would equate two different dimensions, and that is the whole content of the calibration no-go. The six exact defect numbers are therefore the ladder read at six vertex-figure angles and indexed by floor and branch, not by spin, which is why a spin-dependent retention probability was refuted earlier and stays refuted.

\end{theorem}

\begin{theorem}[{\texorpdfstring{BM.125 — The spin-dependent mass: the exact ladder, the two clauses and their matching, and the ceiling}{BM.125 - The spin-dependent mass: the exact ladder, the two clauses and their matching, and the ceiling}}]\label{tag:BM-125}
\textbf{Thm.}

%dsp-level:0
\[
(1)\quad\textbf{the mass ladder:}
\]

%dsp-level:0
\[
\qquad \boxed{\ m^2_{\mathrm{pole}}(m)\ =\ \hat\mu'^2\,\Delta(m)\ =\ 4\hat\mu'^2\sin^2\frac{m\Theta_R}4,\qquad m\in\{k,k-2,\dots,-k\},\qquad k=2j\ }
\]

%dsp-level:0
\[
\qquad [\mathrm{Def\text{-}H.37}(2);\ \mathrm{Def\text{-}BM.19}(3),(6);\ \mathrm{FT.54}(4)]:
\]

%dsp-level:0
\[
\textbf{the ladder's length is the spin's:}
\]

%dsp-level:0
\[
\qquad \#\{m^2_{\mathrm{pole}}\ \text{values}\}\ \le\ 2j+1
\]

%dsp-level:0
\[
\qquad [\text{equality when the}\ \Delta(m)\ \text{are distinct};\ \text{machine}];
\]

%dsp-level:0
\[
(2)\quad\textbf{the exact range and the ceiling:}\qquad0\le4\sin^2\frac{m\Theta_R}4\le4\qquad[\text{machine: the sine-square bound}]:
\]

%dsp-level:0
\[
\boxed{\ 0\ \le\ m^2_{\mathrm{pole}}(m)\ \le\ 4\hat\mu'^2\ }\qquad[\mathrm{FT.76}(3):\ \text{the pole-side ceiling re-cited};\ \text{machine}];
\]

%dsp-level:0
\[
\textbf{the ceiling is attained:}
\]

%dsp-level:0
\[
\qquad \Delta=4\iff\sin^2\frac{m\Theta_R}4=1\iff m\Theta_R\equiv2\pi\ (\mathrm{mod}\ 4\pi)
\]

%dsp-level:0
\[
\qquad [\text{machine: the exact locus}];
\]

%dsp-level:0
\[
\textbf{the ladder's own ceiling is the composite's:}
\]

%dsp-level:0
\[
\qquad \lambda^{(n)}_{\max}=4n\ \ \text{is the sum of}\ n\ \text{attained ceilings}
\]

%dsp-level:0
\[
\qquad [\mathrm{BM.117}(4);\ \text{machine}];
\]

%dsp-level:0
\[
(3)\quad\textbf{the two}\ \mathrm{Def\text{-}H.37}\ \textbf{clauses and their matching identity:}
\]

%dsp-level:0
\[
\qquad (1)\ \ m^2_{\mathrm{pole}}=\hat\mu^2\ell_c^2G(c),
\]

%dsp-level:0
\[
\qquad G(c)=\sum_{A,B}c_Ac_B\operatorname{Re}\langle\mathcal G_A\mathsf v,\mathcal G_B\mathsf v\rangle
\]

%dsp-level:0
\[
\qquad [\mathrm{EW.33}\ \text{Gram data}];
\]

%dsp-level:0
\[
(2)\ \ m^2_{\mathrm{pole}}=\hat\mu'^2\cdot4\sin^2\frac{\Theta_R}4,
\]

%dsp-level:0
\[
\qquad [\hat\mu'^2]=\mathfrak q,
\]

%dsp-level:0
\[
\qquad [\hat\mu^2]=\mathfrak q^2
\]

%dsp-level:0
\[
\qquad [\mathrm{Def\text{-}H.37}(1),(2);\ \mathrm{EW.33}\ \text{packaging}];
\]

%dsp-level:0
\[
\boxed{\ (2)\ \text{implies}\ (1)\ \text{at leading order under the matching}\ \ \hat\mu'^2\,\frac{\Theta_R^2}4\ =\ \hat\mu^2\,\ell_c^2\,G(c)\ }
\]

%dsp-level:0
\[
\qquad [\mathrm{Def\text{-}H.37}:\ \text{recorded, not imposed};\ \mathrm{FT.54}(2);\ \mathrm{FT.33}];
\]

%dsp-level:0
\[
\textbf{the dimension check on the matching:}
\]

%dsp-level:0
\[
\qquad [\hat\mu'^2\Theta_R^2]=\mathfrak q\cdot\mathfrak q^0=\mathfrak q,
\]

%dsp-level:0
\[
\qquad [\hat\mu^2\ell_c^2G(c)]=\mathfrak q^2\cdot[\ell_c^2G(c)]=\mathfrak q\implies[\ell_c^2G(c)]=\mathfrak q^{-1}
\]

%dsp-level:0
\[
\qquad [\text{machine: the exponent arithmetic}];
\]

%dsp-level:0
\[
(4)\quad\textbf{the refutation of the directive's product form:}
\]

%dsp-level:0
\[
\qquad m(j)=\hat\mu'\,\ell_c\sqrt{G(c)\,\Delta(j)}\ \ \text{would require:}
\]

%dsp-level:0
\[
(\mathrm i)\quad\text{a single}\ \Delta(j)\ —\ \text{refuted: the defect is a ladder of}\ 2j+1\ \text{values}\ [\mathrm{BM.124}(4)];
\]

%dsp-level:0
\[
(\mathrm{ii})\quad\text{a product of clause}\ (1)\text{'s}\ \ell_c^2G(c)\ \text{with clause}\ (2)\text{'s}\ \Delta\ —\ \text{refuted: the two clauses are two reads of one mass, related by}\ (3)\text{'s matching,}\ \neg\text{two multiplicative factors};
\]

%dsp-level:0
\[
(\mathrm{iii})\quad\text{a dimension match}\ —\ \text{refuted:}\ [\hat\mu'\ell_c\sqrt{G\Delta}]
\]

%dsp-level:0
\[
= \mathfrak q^{1/2}\cdot[\ell_c]\cdot\mathfrak q^{-1/2}\ \text{against}\ [m]
\]

%dsp-level:0
\[
= \mathfrak q^{1/2},
\]

%dsp-level:0
\[
\ \text{so}\ [\ell_c]\ \text{would have to be}\ \mathfrak q^0\ \text{while}\ (3)\text{'s matching gives}\ [\ell_c^2G]\qquad = \mathfrak q^{-1}
\]

%dsp-level:0
\[
\qquad [\text{machine: the two exponent rows disagree}];
\]

%dsp-level:0
\[
\boxed{\ \textbf{the formed rows are the two clauses separately and their matching identity}\ (3)\ \textbf{— the product form is}\ \neg\textbf{formed on three independent grounds}\ }
\]

%dsp-level:0
\[
(5)\quad\textbf{the standing zero is a massless weight, not a massless species:}
\]

%dsp-level:0
\[
\qquad k\ \text{even}\implies m=0\in\text{the ladder}\implies m^2_{\mathrm{pole}}(0)=0\ \ \text{exactly at every}\ \Theta_R
\]

%dsp-level:0
\[
\qquad [\mathrm{BM.124}(6);\ \mathrm{FT.60}(5);\ \text{machine}]:
\]

%dsp-level:0
\[
\textbf{the scope bars:}
\]

%dsp-level:0
\[
\qquad \neg\text{a species identification attaches to the zero weight}\ [\mathrm{Def\text{-}BM.3}(7);\ \mathrm{P5}];
\]

%dsp-level:0
\[
\qquad \neg\text{a massless}\ \nu\ \text{line is formed}\ —\ \text{the rung-one}\ \nu\ \text{row has}\ m_\nu^{\mathrm{rest}}>0\ \text{strictly on two SOS sandwiches}\ [\mathrm{BM.96}(3),(4)\ \text{re-cited},\ \neg\text{extended}];
\]

%dsp-level:0
\[
\neg\text{a right-ray matter line,}\ \neg\text{a seesaw-type mechanism and}\ \neg\text{a Yukawa-type object is added}\ [\mathrm{BM.96}(4),(6);\ \mathrm{EW.15};\ \mathrm{FT.145}(3)];
\]

%dsp-level:0
\[
(6)\quad\textbf{the residue rows:}
\]

%dsp-level:0
\[
\qquad (\mathrm{R2})\ \ \Theta_R
\]

%dsp-level:0
\[
= \delta_v\ \text{at the vertex-figure locus, whose ledger realization is}\ \neg\text{constructed}\ [\mathrm{BM.67}(9),(10);\ \mathrm{BM.82}(6)(\mathrm{R2})]\ \text{— inherited at every use of}\ (1);
\]

%dsp-level:0
\[
(\mathrm{R1})\ \ \rho_{\mathrm B}=\hat\mu'^2\mathfrak h^2/(\ln2)^2\ \text{is the single cross-currency residue:}
\]

%dsp-level:0
\[
\qquad \frac{m^2_{\mathrm{pole}}}{m^2_{\mathrm{hid}}}=\rho_{\mathrm B}\mathcal R_s,
\]

%dsp-level:0
\[
\qquad \mathcal R_s=\frac{D_S}{4E^2}\in\overline{\mathbb Q}
\]

%dsp-level:0
\[
\qquad [\mathrm{BM.82}(1),(2);\ \mathrm{A\text{-}BM.9}:\ \text{the slot declared, the value}\ \neg\text{issued}];
\]

%dsp-level:0
\[
\neg\text{a currency is converted:}\ [\Delta]
\]

%dsp-level:0
\[
= \mathfrak q^0\ \text{and}\ [\hat\mu'^2]
\]

%dsp-level:0
\[
= \mathfrak q\ \text{against}\ [m^2_{\mathrm{hid}}]
\]

%dsp-level:0
\[
= \mathfrak q^{+1}\ \text{and}\ [\ln(4/3)]\qquad = \mathfrak q^0
\]

%dsp-level:0
\[
\qquad [\mathrm{BM.102}(2);\ \mathrm{BM.83}(3);\ \text{machine: the four exponent rows}];
\]

%dsp-level:0
\[
(7)\quad\textbf{status:}\qquad\textbf{derived:}\ (1),\ (2),\ (3),\ (5)\text{'s zero row};\qquad\textbf{refuted as formed, on three grounds:}\ (4);
\]

%dsp-level:0
\[
\textbf{inherited,}\ \neg\textbf{filled:}\ (6)\text{'s}\ (\mathrm{R1}),\ (\mathrm{R2});
\]

%dsp-level:0
\[
\neg\text{a declaration},
\]

%dsp-level:0
\[
\quad\neg\text{a P4 addition},
\]

%dsp-level:0
\[
\quad\neg\text{a new dimensionful scale}\ [\mathrm{P1}],
\]

%dsp-level:0
\[
\quad\neg\text{a measured value}\ [\mathrm{P5}],
\]

%dsp-level:0
\[
\quad\neg\text{a numerical coupling}\ [\mathrm{EM.25}(7)],
\]

%dsp-level:0
\[
\quad\neg\text{a species identification}\ [\mathrm{Def\text{-}BM.3}(7)];
\]

\textbf{Req.} D2.9 (Parts 1–2); CT2.1, EW.15, EW.33 (Part 2); P1, P4, P5 (Parts 1–2); Def-H.37(1), Def-H.37(2), FT.33, FT.54(2), FT.54(4), FT.60(5), FT.76(3), FT.145(3) (Part 3); Def-BM.3(7), Def-BM.14(3), Def-BM.19(3), Def-BM.19(6), BM.67(9), BM.67(10), BM.82(1), BM.82(2), BM.82(6), BM.83(3), BM.96(3), BM.96(4), BM.96(6), BM.102(2), BM.117(4), BM.124(4), BM.124(6), A-BM.9 (this part); exact radical and trigonometric identities, $\mathfrak q$-exponent arithmetic, the sine-square bound on exact rationals [math; machine: \texttt{checks/\allowbreak{}verify\_\allowbreak{}H\_\allowbreak{}janus\_\allowbreak{}tunnel\_\allowbreak{}concrete\_\allowbreak{}spin.py} §C].

\textbf{Pf.}

%dsp-level:0
\[
(1):\ \mathrm{Def\text{-}H.37}(2)\text{'s}\ m^2_{\mathrm{pole}}
\]

%dsp-level:0
\[
= \hat\mu'^2\cdot4\sin^2(\Theta_R/4)\ \text{read at the multiplet weight}\ m\Theta_R
\]

%dsp-level:0
\[
\mapsto m\Theta_R\ [\mathrm{FT.60}(5);\ \mathrm{FT.54}(4)\text{'s exact SU(2) defect}];
\]

%dsp-level:0
\[
\qquad \#\{m\}=k+1=2j+1\ [\mathrm{BM.124}(2)]
\]

%dsp-level:0
\[
(2):\ 0\le\sin^2\le1\implies0\le4\sin^2\le4\ [\text{machine}];
\]

%dsp-level:0
\[
\qquad \mathrm{FT.76}(3)\text{'s ceiling re-cited};
\]

%dsp-level:0
\[
\qquad \sin^2\tfrac{m\Theta_R}4=1\iff\tfrac{m\Theta_R}4\equiv\tfrac\pi2\ (\mathrm{mod}\ \pi)\iff m\Theta_R\equiv2\pi\ (\mathrm{mod}\ 4\pi)\ [\text{machine}]
\]

%dsp-level:0
\[
\lambda^{(n)}=\sum_aD_S(s_a)\le\sum_a4=4n\ [\mathrm{BM.114}(3);\ \mathrm{BM.117}(4);\ \text{machine}]
\]

%dsp-level:0
\[
(3):\ \mathrm{Def\text{-}H.37}(1),(2)\ \text{re-cited};
\]

%dsp-level:0
\[
\qquad 4\sin^2\tfrac{\Theta_R}4
\]

%dsp-level:0
\[
= \Theta_R^2+O(\Theta_R^4)\implies\hat\mu'^2\Theta_R^2/4
\]

%dsp-level:0
\[
= \hat\mu^2\ell_c^2G(c)\ \text{at leading order}\ [\mathrm{FT.54}(2);\ \mathrm{FT.33}]
\]

%dsp-level:0
\[
[\hat\mu'^2\Theta_R^2]=\mathfrak q\cdot\mathfrak q^0=\mathfrak q;
\]

%dsp-level:0
\[
\qquad [\hat\mu^2]=\mathfrak q^2\implies[\ell_c^2G]=\mathfrak q^{-1}\ [\text{machine: the exponent arithmetic}]
\]

%dsp-level:0
\[
(4):\ (\mathrm i)\ \mathrm{BM.124}(4);
\]

%dsp-level:0
\[
\qquad (\mathrm{ii})\ (3)\text{'s matching is an equality of two reads of one mass, so their product has}\ [m^4]\ \text{not}\ [m^2]\ [\text{machine: the exponent row}];
\]

%dsp-level:0
\[
\qquad (\mathrm{iii})\ \text{the two}\ [\ell_c]\ \text{rows disagree}\ [\text{machine}]
\]

%dsp-level:0
\[
(5):\ \Delta(0)=0\implies m^2_{\mathrm{pole}}(0)=\hat\mu'^2\cdot0=0\ [\text{machine}];
\]

%dsp-level:0
\[
\qquad \mathrm{BM.96}(3),(4)\text{'s strict-positivity rows re-cited at their own locus}\ [\neg\text{extended}]
\]

%dsp-level:0
\[
(6):\ \mathrm{BM.82}(1),(2)\ \text{re-cited};
\]

%dsp-level:0
\[
\qquad [\Delta]=\mathfrak q^0,\ [\hat\mu'^2]=\mathfrak q,\ [m^2_{\mathrm{hid}}]=\mathfrak q^{+1}\ [\mathrm{BM.102}(2);\ \text{machine}];
\]

%dsp-level:0
\[
\qquad (7):\ \text{the standing-rule scans}\ [\text{machine}].\ \blacksquare
\]

\textbf{Phys.} The mass ladder is the defect ladder times one scale, so it inherits everything the ladder has: two-j-plus-one values, an exact zero at the middle weight whenever the winding is even, and a hard ceiling at four times the defect scale that the pole-side read already carried. The composite ceiling is just the sum of attained constituent ceilings, which is why the resonance tier's four-n bound is not a new ceiling but the same one counted. The two clauses of the defect-transfer interface are then kept apart rather than multiplied: one reads the mass through the Gram data of the weak transport with a length factor, the other through the holonomy angle with a defect scale, and they meet in a single matching identity whose dimension check forces the length-times-Gram factor to carry weight minus one. The product form the question proposes fails on three separate counts — there is no single defect value to put in it, multiplying two reads of one mass gives the wrong power, and the dimensions do not close — so the tier supplies the matching identity instead. The massless middle weight is a property of a ladder and not a species claim: the rung-one neutral line keeps its strictly positive rest mass on two sum-of-squares sandwiches, and no right-ray line, no seesaw-type mechanism and no Yukawa-type object is added anywhere. Two residues are inherited and left unfilled, the vertex-figure locus and the single cross-currency ratio, and no currency is converted between them.

\end{theorem}

\medskip\noindent\textit{\texorpdfstring{BM.126 — The spin-dependence of the Janus crossing: refuted on the four standing grounds, with the derived domain row}{BM.126 - The spin-dependence of the Janus crossing: refuted on the four standing grounds, with the derived domain row}.}\ Structural-constraints record: see Appendix B, Section B.2.7 of the Addendum.\par\medskip

\begin{theorem}[{\texorpdfstring{BM.127 — The winding-parity identification: two reads separated}{BM.127 - The winding-parity identification: two reads separated}}]\label{tag:BM-127}
\textbf{Thm.}

%dsp-level:0
\[
(1)\quad\textbf{the statistics read:}\qquad\boxed{\ \varepsilon(c)\ =\ (-1)^{|k(c)|}\ =\ (-1)^{2j(c)}\ =\ \begin{cases}+1 & j\in\mathbb Z\\ -1 & j\in\mathbb Z+\tfrac12\end{cases}\ }\qquad[\mathrm{BM.113}(2),(4);\ \text{machine: the census to}\ k=8];
\]

%dsp-level:0
\[
(2)\quad\textbf{the winding is a ledger datum:}
\]

%dsp-level:0
\[
\qquad \Theta(c)=\sum_{e\in E(c)}\widehat{ds}^{\,2}(e),
\]

%dsp-level:0
\[
\qquad \widehat{ds}^{\,2}(e)=\kappa_e\ \ \text{is registration data determined at crystallization}
\]

%dsp-level:0
\[
\qquad [\mathrm{FT.22};\ \mathrm{FT.24};\ \text{Census FT.28}(2);\ \text{A §0.2a}]:
\]

%dsp-level:0
\[
\Theta_{\mathrm{int}}=2\kappa\in2\mathbb Z,\qquad k=\frac\Theta2=\kappa\in\mathbb Z\setminus\{0\}\qquad[\mathrm{BM.106}(1);\ \text{machine}];
\]

%dsp-level:0
\[
\boxed{\ \textbf{the parity of}\ k\ \textbf{is a read of registration data}\ —\ \neg\text{an external postulate,}\ \neg\text{a declared assignment and}\ \neg\text{a measurement}\ }
\]

%dsp-level:0
\[
\qquad [\mathrm{P5};\ \text{machine: the symbol census of}\ (-1)^{|k|}\ \text{carries}\ \{k\}\ \text{alone}];
\]

%dsp-level:0
\[
(3)\quad\textbf{the kind read is a different object:}
\]

%dsp-level:0
\[
\qquad \text{bound}\ :=\ \text{closed self-contained chain, content fermionic};
\]

%dsp-level:0
\[
\qquad \text{transit}\ :=\ \text{propagating open chain, content bosonic}
\]

%dsp-level:0
\[
\qquad [\mathrm{D1.5}:\ \textbf{a definitional label}]:
\]

%dsp-level:0
\[
\textbf{the dichotomy is the closed/open read}\qquad[\mathrm{FT.80}(5);\ \mathrm{D1.4}];
\]

%dsp-level:0
\[
(4)\quad\textbf{the two reads are}\ \neg\textbf{equated, and the floor census is the proof:}
\]

%dsp-level:0
\[
\qquad \kappa=1\implies k=1\implies j=\tfrac12\implies\varepsilon=-1;
\]

%dsp-level:0
\[
\qquad \kappa=2\implies k=2\implies j=1\implies\varepsilon=+1
\]

%dsp-level:0
\[
\qquad [\mathrm{BM.106}(1);\ \text{machine}]:
\]

%dsp-level:0
\[
\boxed{\ \textbf{both exchange signs occur on materialized closed bound chains of one}\ \mathrm{D1.5}\ \textbf{kind}\ —\ \text{so the kind label is}\ \neg\textbf{a statistics and the statistics is}\ \neg\textbf{a kind}\ }
\]

%dsp-level:0
\[
\textbf{the exact separation:}\qquad\begin{array}{l|l|l}\text{read} & \text{its argument} & \text{its values}\\\hline\text{the statistics}\ \varepsilon & k\ \text{(an integer registration datum)} & \{+1,\ -1\},\ \text{sharp per chain}\\\text{the}\ \mathrm{D1.5}\ \text{kind} & \text{closed against open}\ [\mathrm{D1.4}] & \{\text{bound},\ \text{transit}\},\ \text{both signs on one kind}\\\text{the registration}\ \kappa & \widehat{ds}^{\,2}(e)\ \text{per edge} & \mathbb Z\setminus\{0\},\ \neg\text{a statistics}\end{array}
\]

%dsp-level:0
\[
(5)\quad\textbf{the naming discipline:}
\]

%dsp-level:0
\[
\qquad \textbf{``fermion''/``boson'' occur in this corpus in two distinct uses, and the tier keeps them apart:}
\]

%dsp-level:0
\[
(\mathrm i)\quad\textbf{as statistics labels:}\ \ \text{reads of}\ \varepsilon\ [(1)];
\]

%dsp-level:0
\[
\qquad (\mathrm{ii})\quad\textbf{as content labels:}\ \ \mathrm{D1.5}\text{'s bound/transit vocabulary}\ [(3)];
\]

%dsp-level:0
\[
\boxed{\ \neg\text{a species name is attached to either use}\ —\ \text{the identification layer is}\ \mathrm{Def\text{-}BM.3},\ \mathrm{Def\text{-}BM.6}\ \text{and}\ \mathrm{Def\text{-}BM.14},\ \text{and}\ \neg\text{a clause of this tier names a particle}\ }
\]

%dsp-level:0
\[
\qquad [\mathrm{Def\text{-}BM.3}(7);\ \text{machine: the species-token scan}];
\]

%dsp-level:0
\[
(6)\quad\textbf{what is complete and what is open:}
\]

%dsp-level:0
\[
\qquad \textbf{complete:}\ \varepsilon\ \text{is an exact function of a ledger datum, with no external input}\ [(1),(2)];
\]

%dsp-level:0
\[
\textbf{beyond this tier:}\qquad\begin{array}{l|l}\text{the winding of a named sealed class} & (\mathrm{R3})\ /\ \text{deficit}\ 1':\ \text{which three of the six states a class occupies, and which winding each constituent carries}\ [\mathrm{BM.28}(7);\ \mathrm{BM.103}(8);\ \mathrm{BM.117}(6)]\\\text{the connector winding}\ d=\Theta(P) & \mathrm{BM.2}(2):\ d\ \text{is registration data; an open connector carries}\ S=0\ \text{and hence}\ d=0\ [\mathrm{Def\text{-}H.32}(4);\ \mathrm{BM.128}(7)]\\\text{the}\ k\ge3\ \text{aspect claim} & \mathrm{Def\text{-}H.38}\ \text{Scope: carriers for}\ k\ge3\ \text{exceed the registered pair-channel inventory}\ [\mathrm{CT0.12}]\end{array}
\]

%dsp-level:0
\[
\boxed{\ \textbf{the identification is complete as a function and open as a value per named class}\ }
\]

%dsp-level:0
\[
(7)\quad\textbf{status:}\qquad\textbf{derived:}\ (1)\text{--}(5),\ (6)\text{'s completeness row};
\]

%dsp-level:0
\[
\textbf{re-scoped at its own locus,}\ \neg\textbf{retracted:}\quad\mathrm{BM.106}(1)\text{'s spin-parity refutation stands and is what makes}\ (4)\ \text{a separation rather than a contradiction};
\]

%dsp-level:0
\[
\neg\text{a declaration},
\]

%dsp-level:0
\[
\quad\neg\text{a P4 addition},
\]

%dsp-level:0
\[
\quad\neg\text{a new dimensionful scale}\ [\mathrm{P1}],
\]

%dsp-level:0
\[
\quad\neg\text{a measured value}\ [\mathrm{P5}],
\]

%dsp-level:0
\[
\quad\neg\text{a species identification}\ [\mathrm{Def\text{-}BM.3}(7)];
\]

\textbf{Req.} D0.8, D1.4, D1.5 (Parts 1–2); CT0.12 (Part 1); P1, P4, P5 (Parts 1–2); Def-H.32(4), Def-H.38, FT.22, FT.24, FT.80(5) (Part 3); Census FT.28(2) (row-wise); Def-BM.3(7), Def-BM.6, Def-BM.14, BM.2(2), BM.28(7), BM.103(8), BM.106(1), BM.113(2), BM.113(4), BM.117(6), BM.128(7) (this part); exhaustive parity and registration censuses, symbol-set censuses [math; machine: \texttt{checks/\allowbreak{}verify\_\allowbreak{}H\_\allowbreak{}janus\_\allowbreak{}tunnel\_\allowbreak{}concrete\_\allowbreak{}spin.py} §D].

\textbf{Pf.}

%dsp-level:0
\[
(1):\ \mathrm{BM.113}(2),(4)\ \text{re-cited (frozen)};\qquad\text{the parity census to}\ k=8\ [\text{machine}]
\]

%dsp-level:0
\[
(2):\ \Theta=\sum_e\widehat{ds}^{\,2}(e)\ [\mathrm{FT.22}];
\]

%dsp-level:0
\[
\qquad \widehat{ds}^{\,2}(e)=\kappa_e\ \text{registration data}\ [\mathrm{FT.24};\ \text{Census FT.28}(2);\ \text{A §0.2a}];
\]

%dsp-level:0
\[
\qquad \Theta_{\mathrm{int}}=2\kappa\implies k=\kappa\ [\mathrm{BM.106}(1);\ \text{machine}]
\]

%dsp-level:0
\[
\text{the symbol census of}\ (-1)^{|k|}:\ \{k\}\ [\text{machine}]
\]

%dsp-level:0
\[
(3):\ \mathrm{D1.5},\ \mathrm{D1.4},\ \mathrm{FT.80}(5)\ \text{re-cited}\ [\text{a definitional label and the closed/open read}]
\]

%dsp-level:0
\[
(4):\ \kappa=1\implies k=1,\ j=\tfrac12,\ \varepsilon=(-1)^1=-1;
\]

%dsp-level:0
\[
\qquad \kappa=2\implies k=2,\ j=1,\ \varepsilon=(-1)^2=+1\ [\mathrm{BM.106}(1);\ \text{machine: the two exact rows}]
\]

%dsp-level:0
\[
\text{both chains are materialized closed bound chains of one}\ \mathrm{D1.5}\ \text{kind}\ [\mathrm{Def\text{-}H.4};\ \mathrm{BM.106}(1)]\implies\text{the kind does not determine the sign}
\]

%dsp-level:0
\[
(5):\ \text{the two uses enumerated};
\]

%dsp-level:0
\[
\qquad \mathrm{Def\text{-}BM.3}(7)\text{'s scope row re-cited};
\]

%dsp-level:0
\[
\qquad \text{the species-token scan of this tier}\ [\text{machine}]
\]

%dsp-level:0
\[
(6):\ \text{the three open rows cite their own tags}\ [\text{machine}];
\]

%dsp-level:0
\[
\qquad (7):\ \text{the standing-rule scans}\ [\text{machine}].\ \blacksquare
\]

\textbf{Phys.} The statistics sign is a parity of a registration integer, and that is the whole content of the identification: the winding is a sum of per-edge line elements, each of which is registration data fixed at crystallization, so its parity is a ledger read with no postulate, no declared assignment and no measurement behind it. What the tier is careful about is that this read is not the same object as the bound-or-transit label the corpus has carried since its definition layer, and the proof that they differ is the floor census: at registration one a chain has spin one half and sign minus one, at registration two it has spin one and sign plus one, and both are materialized closed bound chains of the same kind. So the kind label cannot be a statistics and the statistics cannot be a kind, and the earlier refutation of a spin-parity theorem on the kind is exactly what makes this a clean separation instead of a contradiction. The two words fermion and boson are then used in two disciplined ways — as reads of the exchange sign and as the content vocabulary — with no species name attached to either. What is complete is the function; what stays open is its value on any named sealed class, because which windings a class occupies is the same missing datum as the flavour composition, together with the connector winding and the aspect bar at high winding.

\end{theorem}

\begin{theorem}[{\texorpdfstring{BM.128 — The three-chain exchange structure: the space, the $S_3$ action, the character formula, and the vanishing wedge}{BM.128 - The three-chain exchange structure: the space, the action, the character formula, and the vanishing wedge}}]\label{tag:BM-128}
\textbf{Thm.}

%dsp-level:0
\[
(1)\quad\textbf{the three-chain state space:}
\]

%dsp-level:0
\[
\qquad \boxed{\ \mathcal H_3(c_1,c_2,c_3)\ :=\ \bigotimes_{a=1}^{3}\Bigl(\mathcal H_\omega(c_a)\otimes\Psi(c_a)\Bigr),\qquad\dim\mathcal H_3\ =\ 8\prod_{a=1}^{3}\bigl(|k_a|+1\bigr)\ }
\]

%dsp-level:0
\[
[\mathrm{Def\text{-}BM.16}(2)\ \text{extended to}\ n=3;\ \text{machine: exact integers}];\qquad\textbf{the anchors:}\qquad\begin{array}{c|c|c}k=(k_1,k_2,k_3) & \prod(|k_a|+1) & \dim\mathcal H_3\\\hline(1,1,1) & 8 & 64\\(2,2,2) & 27 & 216\\(1,1,2) & 12 & 96\\(1,2,3) & 24 & 192\end{array}
\]

%dsp-level:0
\[
(2)\quad\textbf{the}\ S_3\ \textbf{action:}
\]

%dsp-level:0
\[
\qquad \hat P_\sigma\ :=\ \text{the factor permutation by}\ \sigma\in S_3,
\]

%dsp-level:0
\[
\qquad |S_3|=6,
\]

%dsp-level:0
\[
\qquad \boxed{\ \hat P_\sigma\hat P_\tau=\hat P_{\sigma\tau}\ \ \forall\sigma,\tau\in S_3\ }
\]

%dsp-level:0
\[
\qquad [\text{machine: all}\ 36\ \text{ordered pairs, entrywise over}\ \mathbb Q(i)]:
\]

%dsp-level:0
\[
\textbf{the orders:}\qquad\operatorname{ord}(\hat P_\sigma)=\operatorname{ord}(\sigma)=\begin{cases}1 & \sigma=\mathrm{id}\\ 2 & \sigma\ \text{a transposition}\ (3\ \text{of them})\\ 3 & \sigma\ \text{a}\ 3\text{-cycle}\ (2\ \text{of them})\end{cases}\qquad[\text{machine: exact matrix powers}];
\]

%dsp-level:0
\[
\textbf{the sign representation:}
\]

%dsp-level:0
\[
\qquad \operatorname{sgn}:S_3\to\{\pm1\},
\]

%dsp-level:0
\[
\qquad \ker(\operatorname{sgn})=A_3,\ |A_3|=3,
\]

%dsp-level:0
\[
\qquad S_3/A_3\cong\mathbb Z_2
\]

%dsp-level:0
\[
\qquad [\text{machine: the inversion census}\ (0,1,1,2,2,3)];
\]

%dsp-level:0
\[
(3)\quad\textbf{the character formula (exact):}
\]

%dsp-level:0
\[
\qquad \boxed{\ \chi(\sigma)\ :=\ \operatorname{Tr}\hat P_\sigma\big|_{(\mathbb C^d)^{\otimes n}}\ =\ d^{\,c(\sigma)}\ }
\]

%dsp-level:0
\[
\qquad c(\sigma)=\text{the cycle count of}\ \sigma
\]

%dsp-level:0
\[
[\text{machine: verified on all}\ 2!\ \text{permutations at}\ d\in\{2,3\}\ \text{and all}\ 3!\ \text{at}\ d\in\{2,3\}];
\]

%dsp-level:0
\[
\qquad \textbf{the three rows at}\ n=3,\ d=2:
\]

%dsp-level:0
\[
\qquad \chi(\mathrm{id})=2^3=8,
\]

%dsp-level:0
\[
\qquad \chi(\text{transposition})=2^2=4,
\]

%dsp-level:0
\[
\qquad \chi(3\text{-cycle})=2^1=2;
\]

%dsp-level:0
\[
\textbf{the consequence for the exchange sign:}
\]

%dsp-level:0
\[
\qquad \text{a transposition's trace is}\ d^{\,n-1},\ \neg d^{\,n}:
\]

%dsp-level:0
\[
\qquad \boxed{\ \textbf{the flip is not the identity on the}\ n\textbf{-fold space for}\ n\ge2\ }
\]

%dsp-level:0
\[
\qquad [\text{machine}];
\]

%dsp-level:0
\[
(4)\quad\textbf{the CP-twisted operator:}
\]

%dsp-level:0
\[
\qquad \hat P^{\mathrm{CP}}_\sigma\ :=\ \hat\mu_{\mathrm{CP}}\circ\hat P_\sigma\ \ \text{on}\ \mathcal H_3,
\]

%dsp-level:0
\[
\qquad \hat\mu_{\mathrm{CP}}=\mu_{\mathrm{CP}}\,\mathrm{id},
\]

%dsp-level:0
\[
\qquad \mu_{\mathrm{CP}}=e^{i\pi\Theta}=1\ \ \text{on}\ \Theta\in2\mathbb Z
\]

%dsp-level:0
\[
\qquad [\mathrm{BM.112}(1);\ \mathrm{BM.108}(3)]:
\]

%dsp-level:0
\[
\boxed{\ \bigl(\hat P^{\mathrm{CP}}_{(12)}\bigr)^2=\mathrm{id}\ \ \text{unconditionally},\qquad\operatorname{spec}\hat P^{\mathrm{CP}}_{(12)}\subseteq\{+1,-1\},\qquad\operatorname{Tr}\hat P^{\mathrm{CP}}_{(12)}=\mu_{\mathrm{CP}}\,d^{\,c((12))}\prod_{a}(|k_a|+1)\ }
\]

%dsp-level:0
\[
[\mathrm{BM.112}(1),(2)\ \text{transported to}\ n=3;\ \text{machine}];
\]

%dsp-level:0
\[
(5)\quad\textbf{the vanishing wedge (the structural row of this tag):}
\]

%dsp-level:0
\[
\qquad \dim\wedge^{3}(\mathbb C^d)=\binom d3:
\]

%dsp-level:0
\[
\qquad \boxed{\ \dim\wedge^{3}(\mathbb C^2)=\binom23=0,\qquad\dim\wedge^{3}(\mathbb C^3)=\binom33=1\ }
\]

%dsp-level:0
\[
[\text{machine: the explicit antisymmetrizer}\ \tfrac1{3!}\sum_{\sigma}\operatorname{sgn}(\sigma)\hat P_\sigma\ \text{has rank}\ 0\ \text{on}\ (\mathbb C^2)^{\otimes3}\ \text{and rank}\ 1\ \text{on}\ (\mathbb C^3)^{\otimes3},\ \text{computed exactly over}\ \mathbb Q(i)]:
\]

%dsp-level:0
\[
\textbf{the reading:}
\]

%dsp-level:0
\[
\qquad \boxed{\ \textbf{three copies of the}\ 2\textbf{-dimensional seed carrier admit no totally antisymmetric carrier state, while three colour slots admit exactly one}\ —\ \text{the}\ \varepsilon\text{-seal}\ }
\]

%dsp-level:0
\[
(6)\quad\textbf{the sharpening of}\ \mathrm{BM.114}(2):
\]

%dsp-level:0
\[
\qquad B=\varepsilon_{ijk}q^iq^jq^k\ \ \text{is antisymmetric in the colour slots}\ [\mathrm{D3.9};\ \mathrm{CT3.15};\ \mathrm{CT3.17}]:
\]

%dsp-level:0
\[
\textbf{the two antisymmetries are different objects with different dimensions:}\qquad\begin{array}{l|l|l}\text{antisymmetry} & \text{its space} & \text{its dimension}\\\hline\text{the colour seal} & \wedge^3(\mathbb C^3_{\text{colour}}) & 1\ \ [\text{machine}]\\\text{the carrier wedge} & \wedge^3(\mathbb C^2_{\mathsf v}) & 0\ \ [\text{machine}]\\\text{the statistics sign} & \text{the winding parity on}\ \mathcal H_3 & \{+1,-1\}\ \ [\mathrm{BM.113};\ (7)]\end{array}
\]

%dsp-level:0
\[
\boxed{\ \textbf{the baryon's total antisymmetry is carried by colour, which is}\ \neg\textbf{observable}\ [\mathrm{CT3.22}],\ \textbf{and the statistics is carried by the winding}\ —\ \neg\textbf{one is the other}\ }
\]

%dsp-level:0
\[
(7)\quad\textbf{the multiplicative exchange law:}
\]

%dsp-level:0
\[
\qquad k(\Gamma)=k(c_1)+k(c_2)+d_{12}
\]

%dsp-level:0
\[
\qquad [\mathrm{FT.34}(3);\ \mathrm{BM.2}(2):\ \text{the additive winding law}]:
\]

%dsp-level:0
\[
\textbf{the connector row:}
\]

%dsp-level:0
\[
\qquad d_{12}=\Theta(P)\ge0\ \ \text{for a closed connector};
\]

%dsp-level:0
\[
\qquad \textbf{an open chain carries}\ S=0\implies d=0
\]

%dsp-level:0
\[
\qquad [\mathrm{Def\text{-}H.32}(4);\ \text{machine}];
\]

%dsp-level:0
\[
\textbf{at a sealed junction with an open connector:}
\]

%dsp-level:0
\[
\qquad \boxed{\ k_\Gamma=\sum_{a=1}^{n}k_a\implies\varepsilon(\Gamma)=(-1)^{\left|\sum_ak_a\right|}=\prod_{a=1}^{n}(-1)^{|k_a|}=\prod_{a=1}^{n}\varepsilon(c_a)\ }
\]

%dsp-level:0
\[
[\text{machine: the multiplicativity verified on all}\ 5^3=125\ \text{triples at}\ n=3\ \text{and all}\ 5^2=25\ \text{pairs at}\ n=2];
\]

%dsp-level:0
\[
(8)\quad\textbf{the closure of}\ (\mathrm{R7})\ \textbf{at}\ n=3:\qquad\mathrm{BM.112}(7)\ \text{and}\ \mathrm{Def\text{-}BM.16}(7)\ \text{left}\ \mathcal H_3\ \text{and its}\ S_3\ \text{action}\ \neg\text{constructed}:
\]

%dsp-level:0
\[
\boxed{\ (\mathrm{R7})\ \ \textbf{is discharged at}\ n=3:\quad\text{the space}\ (1),\ \text{the group action with its character}\ (2),(3),\ \text{the involution}\ (4)\ \text{and the permutation-sign read}\ (7)\ \textbf{are all constructed}\ }
\]

%dsp-level:0
\[
\textbf{the residue narrows:}
\]

%dsp-level:0
\[
\qquad (\mathrm{R7}')\ \ \mathcal H_n\ \text{and its}\ S_n\ \text{action for}\ n\ge4\ \ \text{are}\ \neg\textbf{constructed}
\]

%dsp-level:0
\[
\qquad [\text{the sealing grammar admits}\ n\ge4\ \text{supports where}\ \delta\text{- and}\ \varepsilon\text{-groupings mix}\ [\mathrm{BM.31}(2);\ \mathrm{BM.110}(5)]];
\]

%dsp-level:0
\[
(9)\quad\textbf{status:}\qquad\textbf{derived:}\ (1)\text{--}(7);\qquad\textbf{closed:}\ (\mathrm{R7})\ \text{at}\ n=3;\qquad\textbf{named:}\ (\mathrm{R7}');
\]

%dsp-level:0
\[
\textbf{re-scoped at its own locus,}\ \neg\textbf{retracted:}\quad\mathrm{BM.114}(2)\text{'s colour-antisymmetry row, sharpened by}\ (6);
\]

%dsp-level:0
\[
\neg\text{a declaration},
\]

%dsp-level:0
\[
\quad\neg\text{a P4 addition},
\]

%dsp-level:0
\[
\quad\neg\text{a new dimensionful scale}\ [\mathrm{P1}],
\]

%dsp-level:0
\[
\quad\neg\text{a measured value}\ [\mathrm{P5}],
\]

%dsp-level:0
\[
\quad\neg\text{a species identification}\ [\mathrm{Def\text{-}BM.3}(7)];
\]

%dsp-level:0
\[
\neg\text{a flavour}\ \mathsf{SU}(3)\ \text{or}\ \mathsf{SU}(6)\ \text{is constructed:}\ S_3\ \text{permutes}\ \textbf{constituent slots},
\]

%dsp-level:0
\[
\ \neg\text{flavour states}\ [\mathrm{BM.117}(2),(3)];
\]

\textbf{Req.} A3, D1.4 (Part 1); CT3.15, CT3.17, CT3.22, D3.9 (Part 2); P1, P4, P5 (Parts 1–2); Def-H.32(4), Def-H.33, Def-H.38, FT.34(3), FT.108(4) (Part 3); Def-BM.3(7), Def-BM.16(2), Def-BM.16(7), BM.2(2), BM.31(2), BM.108(3), BM.110(5), BM.112(1), BM.112(2), BM.112(7), BM.113, BM.114(2), BM.117(2), BM.117(3) (this part); exact matrix algebra over $\mathbb Q(i)$, exact rank by Fraction elimination, permutation-group enumeration, exhaustive multiplicativity censuses [math; machine: \texttt{checks/\allowbreak{}verify\_\allowbreak{}H\_\allowbreak{}janus\_\allowbreak{}tunnel\_\allowbreak{}concrete\_\allowbreak{}spin.py} §E].

\textbf{Pf.}

%dsp-level:0
\[
(1):\ \dim\mathcal H_\omega
\]

%dsp-level:0
\[
= 2\ \text{and}\ \dim\Psi
\]

%dsp-level:0
\[
= |k|+1\ \text{per constituent}\ [\mathrm{Def\text{-}BM.16}(2);\ \mathrm{Def\text{-}H.38}]\implies\dim\mathcal H_3
\]

%dsp-level:0
\[
= 2^3\prod_a(|k_a|+1)\ [\text{machine: the four anchor rows}]
\]

%dsp-level:0
\[
(2):\ \hat P_\sigma\ \text{built on the}\ d^3\ \text{basis}\ \{e_{i_1}\otimes e_{i_2}\otimes e_{i_3}\}\ \text{by}\ e_{i_1}\otimes e_{i_2}\otimes e_{i_3}
\]

%dsp-level:0
\[
\mapsto e_{i_{\sigma^{-1}(1)}}\otimes\cdots\ [\text{machine}];
\]

%dsp-level:0
\[
\qquad \hat P_\sigma\hat P_\tau=\hat P_{\sigma\tau}\ \text{on all}\ 36\ \text{pairs}\ [\text{machine: entrywise over}\ \mathbb Q(i)]
\]

%dsp-level:0
\[
\text{the order census}\ (1,2,2,3,3,2)\ \text{and the inversion census}\ (0,1,1,2,2,3)\ [\text{machine}];
\]

%dsp-level:0
\[
\qquad |A_3|=3,\ 6/3=2\ [\text{machine}]
\]

%dsp-level:0
\[
(3):\ \text{the basis element}\ e_{i_1}\otimes e_{i_2}\otimes e_{i_3}\ \text{is fixed by}\ \hat P_\sigma
\]

%dsp-level:0
\[
\iff i\ \text{is constant on each cycle of}\ \sigma\implies\operatorname{Tr}
\]

%dsp-level:0
\[
= d^{c(\sigma)}\ [\text{machine: verified on all permutations at}\ (n,d)\in\{2,3\}\times\{2,3\}]
\]

%dsp-level:0
\[
n=3,\ d=2:\ c(\mathrm{id})=3\implies8;\quad c(\text{transposition})=2\implies4;\quad c(3\text{-cycle})=1\implies2\ [\text{machine}]
\]

%dsp-level:0
\[
(4):\ \hat\mu_{\mathrm{CP}}\ \text{is central}\implies(\hat\mu_{\mathrm{CP}}\hat P_{(12)})^2=\mu_{\mathrm{CP}}\mu_{\mathrm{CP}}^\ast\hat P_{(12)}^2=\mathrm{id}\ [\mathrm{BM.112}(1);\ \text{machine}];
\]

%dsp-level:0
\[
\qquad \operatorname{Tr}=\mu_{\mathrm{CP}}\chi((12))\prod_a(|k_a|+1)\ [(3);\ \text{machine}]
\]

%dsp-level:0
\[
(5):\ \dim\wedge^3(\mathbb C^d)=\binom d3\ [\text{math}];
\]

%dsp-level:0
\[
\qquad \binom23=0,\ \binom33=1\ [\text{machine}];
\]

%dsp-level:0
\[
\qquad \text{the explicit antisymmetrizer}\ \tfrac1{6}\sum_\sigma\operatorname{sgn}(\sigma)\hat P_\sigma\ \text{has rank}\ 0\ \text{at}\ d
\]

%dsp-level:0
\[
= 2\ \text{and rank}\ 1\ \text{at}\ d
\]

%dsp-level:0
\[
= 3\ [\text{machine: exact Fraction elimination on the}\ 8\times8\ \text{and}\ 27\times27\ \text{matrices}]
\]

%dsp-level:0
\[
(6):\ \mathrm{D3.9},\ \mathrm{CT3.15},\ \mathrm{CT3.17},\ \mathrm{CT3.22}\ \text{re-cited};\qquad\text{the three-row census}\ [(5);\ \mathrm{BM.113}(4);\ \text{machine}]
\]

%dsp-level:0
\[
(7):\ k(\Gamma)=\sum_ak_a+d\ [\mathrm{FT.34}(3);\ \mathrm{BM.2}(2)];\qquad\text{an open chain carries}\ S=0\ [\mathrm{Def\text{-}H.32}(4)]\implies d=0;
\]

%dsp-level:0
\[
(-1)^{|k_1+k_2+k_3|}
\]

%dsp-level:0
\[
= (-1)^{k_1+k_2+k_3}\ \text{on}\ k_a
\]

%dsp-level:0
\[
\ge 0\ [\mathrm{Def\text{-}H.38}\ \text{Scope}]
\]

%dsp-level:0
\[
= (-1)^{k_1}(-1)^{k_2}(-1)^{k_3}\ [\text{machine: all}\ 125\ \text{triples in}\ \{0,\dots,4\}^3\ \text{and all}\ 25\ \text{pairs}]
\]

%dsp-level:0
\[
(8):\ \mathrm{BM.112}(7),\ \mathrm{Def\text{-}BM.16}(7)\ \text{re-cited};
\]

%dsp-level:0
\[
\qquad \text{the four members discharged against}\ \mathrm{BM.106}(7)\text{'s}\ (\mathrm{R7})\ [\text{machine: the row-by-row check}];
\]

%dsp-level:0
\[
\qquad (9):\ \text{the standing-rule scans}\ [\text{machine}].\ \blacksquare
\]

\textbf{Phys.} Three chains now have a space, a group and a character, and the three fit together exactly. The space is the tensor product of the three orientation doublets and the three winding carriers, eight times the product of the carrier dimensions, which at the two anchors is sixty-four and two hundred sixteen. The symmetric group on three letters acts by permuting factors, the operators multiply as the permutations do on all thirty-six ordered pairs, and the traces obey one clean formula: the dimension to the number of cycles, giving eight, four and two at three chains over a two-dimensional carrier. Twisting by the CP phase changes nothing on the materialized locus, where the phase is one, and the transposition remains an unconditional involution. The structural surprise is the wedge: three copies of the two-dimensional seed carrier have no totally antisymmetric state at all, because the binomial coefficient of three taken from two is zero, while three colour slots have exactly one, and an explicit antisymmetrizer confirms both ranks over the exact complex field. That is why the baryon's antisymmetry is colour's and not the carrier's — and colour is unobservable by construction, so the statistics has to live somewhere else, which is the winding. The winding then supplies the law the tier was after: at a sealed junction with an open connector the total winding is the sum of the constituent windings, so the exchange sign of a composite is the product of its constituents' signs, verified on every triple in a five-by-five-by-five census. That closes the exchange-structure residue at three chains, and narrows it to four and beyond, where the sealing grammar starts mixing delta and epsilon groupings.

\end{theorem}

\begin{theorem}[{\texorpdfstring{BM.129 — The meson is a boson and the baryon's sign is the product: the composite identification}{BM.129 - The meson is a boson and the baryon's sign is the product: the composite identification}}]\label{tag:BM-129}
\textbf{Thm.}

%dsp-level:0
\[
(1)\quad\textbf{the meson:}
\]

%dsp-level:0
\[
\qquad \text{a}\ \delta\text{-sealed CP-pair with an open connector:}\quad k_M=k(c)+k(c^{\mathrm{CP}})+d=k+k+0=2k
\]

%dsp-level:0
\[
\qquad [\mathrm{BM.128}(7);\ \mathrm{BM.113}(2)(\mathrm{ii}):\ \Theta(c^{\mathrm{CP}})=\Theta(c)\ \text{on the materialized locus}]:
\]

%dsp-level:0
\[
\boxed{\ \varepsilon_M\ =\ (-1)^{|2k|}\ =\ \bigl((-1)^2\bigr)^{k}\ =\ +1\qquad\textbf{for every constituent winding}\ k\ge0\ }
\]

%dsp-level:0
\[
\qquad [\text{machine: verified at every}\ k\in\{0,\dots,8\}];
\]

%dsp-level:0
\[
j_M=\frac{|k_M|}2=\frac{2k}2=k\in\mathbb Z
\]

%dsp-level:0
\[
\qquad [\text{machine}]:
\]

%dsp-level:0
\[
\qquad \boxed{\ \textbf{the meson's total spin is an integer and its exchange sign is}\ +1\ —\ \textbf{it is a boson, unconditionally}\ }
\]

%dsp-level:0
\[
(2)\quad\textbf{the robustness row (the orientation-sign reading):}
\]

%dsp-level:0
\[
\qquad \text{the alternative reading}\ k(c^{\mathrm{CP}})=-k(c)\ \ [\hat\omega:\ k\mapsto-k\ \text{on the carrier};\ \mathrm{BM.73}(4);\ \mathrm{A\text{-}BM.7}(b)]:
\]

%dsp-level:0
\[
k_M=k+(-k)+0=0\implies j_M=0,\qquad\varepsilon_M=(-1)^0=+1\qquad[\text{machine}];
\]

%dsp-level:0
\[
\boxed{\ \textbf{both readings give}\ \varepsilon_M=+1\ —\ \textbf{the meson's statistics is robust under the unresolved orientation-sign row, while its total spin is not}\ }
\]

%dsp-level:0
\[
\qquad [j_M=k\ \text{against}\ j_M=0;\ \text{the spin value is}\ \neg\text{fixed and the statistics is}];
\]

%dsp-level:0
\[
(3)\quad\textbf{the closed-connector variant:}
\]

%dsp-level:0
\[
\qquad k_M=2k+d\implies\boxed{\ \varepsilon_M=(-1)^{d}\ }
\]

%dsp-level:0
\[
\qquad [\text{machine: the two constituent windings cancel modulo two}]:
\]

%dsp-level:0
\[
\textbf{the meson's statistics is then the connector's parity alone:}
\]

%dsp-level:0
\[
\qquad d\ \text{even}\implies+1,
\]

%dsp-level:0
\[
\qquad d\ \text{odd}\implies-1
\]

%dsp-level:0
\[
\qquad [\text{machine:}\ d\in\{0,1,2,3,4\}];
\]

%dsp-level:0
\[
d=\Theta(P)\ge0\ \ \text{is registration data}
\]

%dsp-level:0
\[
\qquad [\mathrm{BM.2}(2);\ \mathrm{FT.34}(3)];
\]

%dsp-level:0
\[
\qquad \textbf{the computed cells:}\ \ \text{an open connector carries}\ S=0\implies d=0\implies\varepsilon_M=+1
\]

%dsp-level:0
\[
\qquad [\mathrm{Def\text{-}H.32}(4);\ \text{machine}];
\]

%dsp-level:0
\[
(4)\quad\textbf{the baryon:}\qquad\text{an}\ \varepsilon\text{-sealed 3-juncture with an open connector:}\quad k_B=k_1+k_2+k_3\qquad[\mathrm{BM.128}(7)]:
\]

%dsp-level:0
\[
\boxed{\ \varepsilon_B\ =\ (-1)^{k_1+k_2+k_3}\ =\ \varepsilon(c_1)\,\varepsilon(c_2)\,\varepsilon(c_3),\qquad j_B=\frac{|k_1+k_2+k_3|}2\ }
\]

%dsp-level:0
\[
\textbf{the exact criterion:}
\]

%dsp-level:0
\[
\qquad \boxed{\ \varepsilon_B=-1\iff\#\{a:\ k_a\ \text{odd}\}\ \text{is odd}\ }
\]

%dsp-level:0
\[
\qquad [\text{machine: all}\ 125\ \text{triples in}\ \{0,\dots,4\}^3;\ 62\ \text{fermionic and}\ 63\ \text{bosonic}];
\]

%dsp-level:0
\[
(5)\quad\textbf{the anchor and the computed floors:}
\]

%dsp-level:0
\[
\qquad \textbf{all three constituents at the minimal odd winding}\ k_a=1:
\]

%dsp-level:0
\[
\qquad k_B=3,
\]

%dsp-level:0
\[
\qquad j_B=\tfrac32,
\]

%dsp-level:0
\[
\qquad \varepsilon_B=-1
\]

%dsp-level:0
\[
\qquad [\text{machine}]:
\]

%dsp-level:0
\[
\textbf{the digon floor:}\quad L=2\implies\kappa\equiv1,\ k=L/2=1\ \ (\text{odd})
\]

%dsp-level:0
\[
\qquad [\mathrm{FT.32}(6);\ \mathrm{FT.28};\ \text{machine}]:
\]

%dsp-level:0
\[
\qquad \boxed{\ \textbf{the minimal materialized constituent carries}\ k=1,\ \text{so the three-floor baryon is fermionic}\ }
\]

%dsp-level:0
\[
\textbf{the cycle ladder's parity:}\qquad L\ \text{even}\implies k=L/2;\qquad L\ \text{odd}\implies k=L\qquad[\mathrm{FT.32}(6);\ \text{machine}]:
\]

%dsp-level:0
\[
\begin{array}{c|c|c|c}L & k & k\ \text{parity} & \varepsilon\\\hline2 & 1 & \text{odd} & -1\\3 & 3 & \text{odd} & -1\\4 & 2 & \text{even} & +1\\5 & 5 & \text{odd} & -1\\6 & 3 & \text{odd} & -1\\8 & 4 & \text{even} & +1\end{array}\qquad[\text{machine: exact}];
\]

%dsp-level:0
\[
(6)\quad\textbf{what is derived and what awaits a datum:}
\]

%dsp-level:0
\[
\qquad \textbf{derived:}\ \varepsilon_B\ \text{as an exact function of the three constituent windings}\ [(4)];
\]

%dsp-level:0
\[
\ \text{the multiplicative law}\ [\mathrm{BM.128}(7)];
\]

%dsp-level:0
\[
\ \text{the criterion}\ [(4)];
\]

%dsp-level:0
\[
\textbf{awaits}\ (\mathrm{R3})/\text{deficit}\ 1':
\]

%dsp-level:0
\[
\qquad \boxed{\ \textbf{which windings a named sealed class occupies is}\ \neg\textbf{derived}\ —\ \text{so the baryon's statistics is complete as a function and open as a value per class}\ }
\]

%dsp-level:0
\[
[\mathrm{BM.28}(7);\ \mathrm{BM.103}(8);\ \mathrm{BM.117}(6);\ \mathrm{BM.127}(6);\ \text{machine: the symbol census of the criterion carries}\ \{k_1,k_2,k_3\}\ \text{and}\ \neg\text{a class name}];
\]

%dsp-level:0
\[
(7)\quad\textbf{the colour row and the octet row:}
\]

%dsp-level:0
\[
\qquad \textbf{the}\ \varepsilon\textbf{-seal's antisymmetry is colour's, one-dimensional, and}\ \neg\textbf{observable:}
\]

%dsp-level:0
\[
\qquad \dim\wedge^3(\mathbb C^3_{\text{colour}})=1,
\]

%dsp-level:0
\[
\qquad \dim\wedge^3(\mathbb C^2_{\mathsf v})=0
\]

%dsp-level:0
\[
\qquad [\mathrm{BM.128}(5),(6);\ \mathrm{CT3.22}];
\]

%dsp-level:0
\[
\textbf{no flavour multiplet is formed:}
\]

%dsp-level:0
\[
\qquad S_3\ \text{permutes constituent slots},\ \neg\text{flavour states};
\]

%dsp-level:0
\[
\qquad \neg\exists\ \text{a flavour}\ \mathsf{SU}(3);
\]

%dsp-level:0
\[
\qquad \text{the octet and decuplet groupings are}\ \neg\text{derived}
\]

%dsp-level:0
\[
\qquad [\mathrm{BM.117}(2),(3);\ \mathrm{BM.114}(4)];
\]

%dsp-level:0
\[
(8)\quad\textbf{status:}\qquad\textbf{derived:}\ (1)\text{--}(5),\ (7);\qquad\textbf{open, awaiting}\ (\mathrm{R3}):\ (6)\text{'s per-class value};
\]

%dsp-level:0
\[
\textbf{re-scoped at its own locus,}\ \neg\textbf{retracted:}\quad\mathrm{BM.117}(2),
\]

%dsp-level:0
\[
(3)\text{'s orbit-count rows stand and are}\ \neg\text{promoted to multiplets};
\]

%dsp-level:0
\[
\neg\text{a declaration},
\]

%dsp-level:0
\[
\quad\neg\text{a P4 addition},
\]

%dsp-level:0
\[
\quad\neg\text{a new dimensionful scale}\ [\mathrm{P1}],
\]

%dsp-level:0
\[
\quad\neg\text{a measured value}\ [\mathrm{P5}],
\]

%dsp-level:0
\[
\quad\neg\text{a numerical coupling}\ [\mathrm{EM.25}(7)],
\]

%dsp-level:0
\[
\quad\neg\text{a species identification}\ [\mathrm{Def\text{-}BM.3}(7)];
\]

\textbf{Req.} A3, D1.4 (Part 1); CT3.22 (Part 2); P1, P4, P5 (Parts 1–2); Def-H.32(4), Def-H.38, FT.28, FT.32(6), FT.34(3) (Part 3); Def-BM.3(7), Def-BM.16(2), BM.2(2), BM.28(7), BM.73(4), BM.103(8), BM.113(2), BM.114(4), BM.117(2), BM.117(3), BM.117(6), BM.127(6), BM.128(5), BM.128(6), BM.128(7), A-BM.7 (this part); exhaustive parity and triple censuses, exact integer winding arithmetic, symbol-set censuses [math; machine: \texttt{checks/\allowbreak{}verify\_\allowbreak{}H\_\allowbreak{}janus\_\allowbreak{}tunnel\_\allowbreak{}concrete\_\allowbreak{}spin.py} §E].

\textbf{Pf.}

%dsp-level:0
\[
(1):\ k_M=k(c)+k(c^{\mathrm{CP}})+d\ [\mathrm{BM.128}(7)];
\]

%dsp-level:0
\[
\qquad \Theta(c^{\mathrm{CP}})=\Theta(c)\ \text{on the materialized locus}\ [\mathrm{BM.113}(2)(\mathrm{ii})]\implies k(c^{\mathrm{CP}})=k(c)=k;
\]

%dsp-level:0
\[
\qquad d=0\ \text{for an open connector}\ [\mathrm{Def\text{-}H.32}(4)]
\]

%dsp-level:0
\[
k_M=2k\implies\varepsilon_M=(-1)^{|2k|}=((-1)^2)^k=1\ [\text{machine: every}\ k\in\{0,\dots,8\}];
\]

%dsp-level:0
\[
\qquad j_M=|2k|/2=k\in\mathbb Z\ [\text{machine}]
\]

%dsp-level:0
\[
(2):\ k(c^{\mathrm{CP}})=-k(c)\ [\hat\omega:\ k\mapsto-k\ [\mathrm{BM.73}(4)]]\implies k_M=0\implies\varepsilon_M=(-1)^0=+1,\ j_M=0\ [\text{machine}];
\]

%dsp-level:0
\[
\qquad \text{the two readings agree on}\ \varepsilon_M\ \text{and differ on}\ j_M\ [\text{machine}]
\]

%dsp-level:0
\[
(3):\ k_M=2k+d\implies(-1)^{|2k+d|}=(-1)^{2k+d}=(-1)^{2k}(-1)^d=(-1)^d\ [\text{machine:}\ d\in\{0,\dots,4\}]
\]

%dsp-level:0
\[
(4):\ k_B=k_1+k_2+k_3\ [\mathrm{BM.128}(7)];
\]

%dsp-level:0
\[
\qquad (-1)^{k_1+k_2+k_3}=(-1)^{k_1}(-1)^{k_2}(-1)^{k_3}\ [\text{machine: all}\ 125\ \text{triples}];
\]

%dsp-level:0
\[
\qquad \#\{a:k_a\ \text{odd}\}\ \text{odd}\iff\varepsilon_B=-1\ [\text{machine:}\ 62\ \text{of}\ 125]
\]

%dsp-level:0
\[
(5):\ k_a=1\ \forall a\implies k_B=3,\ j_B=3/2,\ \varepsilon_B=(-1)^3=-1\ [\text{machine}];
\]

%dsp-level:0
\[
\qquad L=2\implies\kappa\equiv1,\ k=L/2=1\ [\mathrm{FT.32}(6);\ \mathrm{FT.28}];
\]

%dsp-level:0
\[
\qquad \text{the six-row ladder}\ [\text{machine}]
\]

%dsp-level:0
\[
(6):\ \mathrm{BM.28}(7),\ \mathrm{BM.103}(8),\ \mathrm{BM.117}(6),\ \mathrm{BM.127}(6)\ \text{re-cited};
\]

%dsp-level:0
\[
\qquad \text{the symbol census of}\ (-1)^{k_1+k_2+k_3}\ \text{carries}\ \{k_1,k_2,k_3\}\ \text{alone}\ [\text{machine}]
\]

%dsp-level:0
\[
(7):\ \mathrm{BM.128}(5),(6)\ \text{and}\ \mathrm{CT3.22}\ \text{re-cited};
\]

%dsp-level:0
\[
\qquad \mathrm{BM.117}(2),(3)\text{'s octet rows re-cited}\ [\neg\text{promoted}];
\]

%dsp-level:0
\[
\qquad (8):\ \text{the standing-rule scans}\ [\text{machine}].\ \blacksquare
\]

\textbf{Phys.} Two composites, two clean answers, and one honest gap. The meson is a boson and it is one unconditionally: its constituents are a chain and its conjugate, so their windings are equal and their sum is even, which makes the total spin an integer and the exchange sign plus one — and the conclusion survives the corpus's one unresolved orientation-sign convention, since the alternative reading puts the total winding at zero and the sign at plus one as well, differing only in the spin value. With a closed connector the answer changes to the connector's own parity, because the two constituent windings cancel modulo two, and on the computed cells the connector is open so the boson sign is what holds. The baryon is not unconditional and the tier does not pretend otherwise: its sign is the product of its three constituents' signs, so it is fermionic exactly when an odd number of them carry odd winding, and at the minimal materialized floor the winding is one — odd — which makes the three-floor baryon fermionic with total spin three halves. Sixty-two of the one hundred twenty-five winding triples in a small census are fermionic and sixty-three bosonic, so the criterion genuinely discriminates. What is missing is not the rule but its input: which windings a named sealed class occupies is the same unread datum as the flavour composition, so the baryon's statistics is complete as a function and open as a value per class. And the antisymmetry that makes a baryon a baryon is colour's — one-dimensional over three colour slots against a carrier wedge that vanishes at dimension two — unobservable by construction and not the statistics, with no flavour multiplet formed anywhere in the argument.

\end{theorem}

\medskip\noindent\textit{\texorpdfstring{BM.130 — The status of the Annihilation, Concrete Spin and Fermion Identification Tier: adjudication, scope, and the open items}{BM.130 - The status of the Annihilation, Concrete Spin and Fermion Identification Tier: adjudication, scope, and the open items}.}\ Structural-constraints record: see Appendix B, Section B.1.33 of the Addendum.\par\medskip

\phantomsection
\section[Tier BM.23 — The Closing Tier: the Haar–Poisson Duality and…]{\texorpdfstring{Tier BM.23 — The Closing Tier: the Haar–Poisson Duality and the Lapse, the Selected Curvature Composition, the Momentum-Space Resolvent, the Charge and Statistics Censuses, the Two-Chain Ensemble, the Seed-Tensoring Mechanism, the Convention Declarations, the Ceiling Study, the Jensen Residue, the Pair Amplitude, and the Era Closure}{Tier BM.23 - The Closing Tier: the Haar-Poisson Duality and the Lapse, the Selected Curvature Composition, the Momentum-Space Resolvent, the Charge and Statistics Censuses, the Two-Chain Ensemble, the Seed-Tensoring Mechanism, the Convention Declarations, the Ceiling Study, the Jensen Residue, the Pair Amplitude, and the Era Closure}}

\begin{theorem}[{\texorpdfstring{BM.131 — The Haar–Poisson duality: the twirl fixed point, the complementary factorization, and the proper-time contraction}{BM.131 - The Haar-Poisson duality: the twirl fixed point, the complementary factorization, and the proper-time contraction}}]\label{tag:BM-131}
\textbf{Thm.}

%dsp-level:0
\[
(1)\quad\textbf{the two evolutions, from the axioms:}
\]

%dsp-level:0
\[
\qquad \textbf{the Haar side:}\quad\mu_\star=\mu_{\mathrm{Haar}}\bigl(\mathrm{SU}(2)\bigr),
\]

%dsp-level:0
\[
\qquad \mathcal T(\rho)=U\rho U^\dagger,
\]

%dsp-level:0
\[
\qquad U\sim\mu_\star,
\]

%dsp-level:0
\[
\qquad U\notin\mathfrak L
\]

%dsp-level:0
\[
\qquad [\mathrm{A6}];
\]

%dsp-level:0
\[
\textbf{the Poisson side:}
\]

%dsp-level:0
\[
\qquad K:=-\log\rho
\]

%dsp-level:0
\[
\qquad [\mathrm{L1}:\ \text{the modular Hamiltonian}],
\]

%dsp-level:0
\[
\qquad \tilde{\mathcal T}_\gamma(\rho)=\frac{\rho\,e^{-\gamma\rho}}{\mathfrak T\bigl(\rho e^{-\gamma\rho}\bigr)},
\]

%dsp-level:0
\[
\qquad \gamma=1-4\varepsilon_q\in(0,1]
\]

%dsp-level:0
\[
\qquad [\mathrm{D0.3};\ \mathrm{L0}:\ \text{one tilt step per A4 event}];
\]

%dsp-level:0
\[
\textbf{the ledger increment factorizes:}
\]

%dsp-level:0
\[
\qquad \rho\ \longmapsto\ U\,\tilde{\mathcal T}_\gamma(\rho)\,U^\dagger
\]

%dsp-level:0
\[
\qquad [\mathrm{A6}:\ \tau_n\to\tau_{n+1};\ \mathrm{L0}:\ \text{the crystallization--tilt coupling, declared at the point of use}];
\]

%dsp-level:0
\[
(2)\quad\textbf{the exact K-update (the Poisson evolution):}
\]

%dsp-level:0
\[
\qquad \lambda_i'=\frac{\lambda_ie^{-\gamma\lambda_i}}{Z_\gamma},
\]

%dsp-level:0
\[
\qquad Z_\gamma=\textstyle\sum_j\lambda_je^{-\gamma\lambda_j}
\]

%dsp-level:0
\[
\qquad [\mathrm{L1}(3)\ \mathrm{Pf}:\ \text{the eigenvalue flow}]:
\]

%dsp-level:0
\[
\boxed{\ \delta K\ :=\ K'-K\ =\ \gamma\,\rho\ +\ \log Z_\gamma\cdot I\ }
\]

%dsp-level:0
\[
\qquad [\log\lambda_i'=\log\lambda_i-\gamma\lambda_i-\log Z_\gamma\ \text{exactly};\ \text{machine: the formal-series identity to order}\ 12];
\]

%dsp-level:0
\[
\gamma=1\ (\text{the sharp point}\ \varepsilon_q=0\ [\mathrm{D0.4}]):\qquad\delta K=\rho+\log Z\cdot I\qquad[\text{machine}];
\]

%dsp-level:0
\[
\textbf{the vacuum is the Poisson fixed point:}
\]

%dsp-level:0
\[
\qquad \rho=\rho_\circ=\tfrac14I:\quad Z_\gamma=e^{-\gamma/4},
\]

%dsp-level:0
\[
\qquad \lambda_i'=\tfrac14,
\]

%dsp-level:0
\[
\qquad \boxed{\ \delta K=0\ \ \text{exactly}\ }
\]

%dsp-level:0
\[
\qquad [\mathrm{L1}(1);\ \text{machine:}\ e^{-\gamma/4}\cdot e^{+\gamma/4}=1\ \text{coefficient-wise in}\ \mathbb Q[[\gamma]]\ \text{to order}\ 12];
\]

%dsp-level:0
\[
(3)\quad\textbf{the Haar invariances (the measure side):}
\]

%dsp-level:0
\[
\qquad V\bigl(\tilde K(\mathcal T(\rho))\bigr)=V\bigl(\tilde K(\rho)\bigr)\ \ \text{for every draw}
\]

%dsp-level:0
\[
\qquad [\mathrm{L1}(2);\ \text{machine: the spectral invariance}\ \operatorname{Tr}\bigl((U\rho U^\dagger)^n\bigr)=\operatorname{Tr}(\rho^n),\ n=1,2,3,\ \text{on}\ 40\ \text{exact rational}\ \mathrm{SU}(2)\ \text{points and generic rational states}];
\]

%dsp-level:0
\[
K\mapsto UKU^\dagger\implies\operatorname{spec}K\ \text{invariant};\qquad\text{the draws are unrecorded}\ [\mathrm{A6}:\ U\notin\mathfrak L];\qquad\neg\exists\,\tau^{-1}\qquad[\mathrm{A6}];
\]

%dsp-level:0
\[
\textbf{the stationarity of the measure:}
\]

%dsp-level:0
\[
\qquad \text{conjugation by a fixed}\ U_0\ \text{pushes}\ \mu_\star\ \text{to}\ \mu_\star
\]

%dsp-level:0
\[
\qquad [\text{math: the two-sided invariance and uniqueness of the Haar measure on a compact group}];
\]

%dsp-level:0
\[
(4)\quad\textbf{the twirl theorem (the duality's fixed locus):}
\]

%dsp-level:0
\[
\qquad \mathcal H_p\otimes\mathcal H_q=\vee^2\oplus\wedge^2
\]

%dsp-level:0
\[
\qquad [\mathrm{A3}],
\]

%dsp-level:0
\[
\qquad \Pi_{\mathcal R}=\tfrac12(I+F),
\]

%dsp-level:0
\[
\qquad \Pi_{\mathcal V}=\tfrac12(I-F),
\]

%dsp-level:0
\[
\qquad \operatorname{rk}\Pi_{\mathcal R}=3,
\]

%dsp-level:0
\[
\qquad \operatorname{rk}\Pi_{\mathcal V}=1
\]

%dsp-level:0
\[
\qquad [\mathrm{FT.107}(2);\ \text{machine: the flip}\ F^2=I]:
\]

%dsp-level:0
\[
\boxed{\ \int_{\mathrm{SU}(2)}\!\bigl(U\otimes U\bigr)\rho\,\bigl(U\otimes U\bigr)^\dagger d\mu_\star(U)\ =\ \frac{\mathfrak T\bigl(\Pi_{\mathcal R}\rho\bigr)}3\,\Pi_{\mathcal R}\ +\ \mathfrak T\bigl(\Pi_{\mathcal V}\rho\bigr)\,\Pi_{\mathcal V}\ }
\]

%dsp-level:0
\[
\qquad [\text{math: Schur's lemma on the}\ 3\oplus1\ \text{decomposition};\ \text{machine: idempotence, the block-trace preservation, and the equivariance}\ \mathrm{tw}\bigl((U\otimes U)\rho(U\otimes U)^\dagger\bigr)=(U\otimes U)\,\mathrm{tw}(\rho)\,(U\otimes U)^\dagger\ \text{on all}\ 40\ \text{exact rational points}];
\]

%dsp-level:0
\[
\textbf{the twirl fixed locus is the block-maximal-mixed family}\ \bigl\{\tfrac x3\Pi_{\mathcal R}+(1-x)\Pi_{\mathcal V}:\ x\in[0,1]\bigr\}
\]

%dsp-level:0
\[
\qquad [\text{machine}];
\]

%dsp-level:0
\[
\textbf{the unique common fixed point of both sides:}
\]

%dsp-level:0
\[
\qquad \lambda_i'=\lambda_i\ \forall i\iff e^{-\gamma\lambda_i}=Z_\gamma\ \forall i\iff\lambda_i\ \text{all equal}\ (\gamma>0)\iff\rho=\tfrac14I;
\]

%dsp-level:0
\[
\boxed{\ \Bigl\{\text{Poisson fixed points}\Bigr\}\cap\Bigl\{\text{Haar-twirl fixed points}\Bigr\}\ =\ \Bigl\{\tfrac14I\Bigr\}\ }
\]

%dsp-level:0
\[
\qquad [\text{machine:}\ x/3=1-x\iff x=\tfrac34\ \text{gives}\ \rho=\tfrac14I\ \text{on the block family};\ \text{the tilt fixed point is}\ \tfrac14I\ [\mathrm{L1}(1),(5)]];
\]

%dsp-level:0
\[
(5)\quad\textbf{the complementary factorization (what the record sees):}
\]

%dsp-level:0
\[
\qquad V\ \text{is invariant under the Haar factor}\ [(3)]\ \text{and strictly contracting under the Poisson factor:}
\]

%dsp-level:0
\[
\qquad V'=\Bigl(1-\frac\gamma4\Bigr)^{\!2}V+O(\delta\rho^3)
\]

%dsp-level:0
\[
\qquad [\mathrm{L2}];
\]

%dsp-level:0
\[
\Bigl(1-\frac\gamma4\Bigr)^{\!2}+\frac{\gamma(8-\gamma)}{16}=1
\]

%dsp-level:0
\[
\qquad [\mathrm{BM.51}(2);\ \text{machine: exact in}\ \mathbb Q[\gamma]]\implies V_n=\Bigl(1-\frac\gamma4\Bigr)^{2n}V_0+O
\]

%dsp-level:0
\[
\qquad [\mathrm{BM.51}(3)];
\]

%dsp-level:0
\[
\boxed{\ \textbf{the unrecorded Haar draws drop out of the record exactly:}\ \ V_n\ \text{depends on the tilt count alone}\ }
\]

%dsp-level:0
\[
\qquad [\mathrm{L1}(2)+\mathrm{L2};\ \text{machine}];
\]

%dsp-level:0
\[
\textbf{the reading:}
\]

%dsp-level:0
\[
\qquad \text{the measure side is the exact gauge freedom of the record};
\]

%dsp-level:0
\[
\qquad \text{the Poisson side is the recorded dynamics};
\]

%dsp-level:0
\[
\qquad \text{the duality fixed point is the A5 vacuum}\ \rho_\circ\ [\mathrm{A5}];
\]

%dsp-level:0
\[
(6)\quad\textbf{the finite-group companion (the same twirl at the carrier level):}
\]

%dsp-level:0
\[
\qquad \Pi_k:=\frac1{k!}\sum_{\sigma\in S_k}\hat P_\sigma
\]

%dsp-level:0
\[
\qquad [\text{the averaging over the finite group}\ S_k\ \text{is its Haar measure}]:
\]

%dsp-level:0
\[
\Pi_k^2=\Pi_k,
\]

%dsp-level:0
\[
\qquad \operatorname{rk}\Pi_k=k+1,
\]

%dsp-level:0
\[
\qquad \operatorname{im}\Pi_k=\mathrm{Sym}^k\bigl(\mathbb C^2\bigr)
\]

%dsp-level:0
\[
\qquad [\text{machine: exact rank over}\ \mathbb Q\ \text{for}\ k=1..6;\ \text{forward pointer: the mechanism is BM.138}];
\]

%dsp-level:0
\[
(7)\quad\textbf{the proper-time contraction (the metric read):}
\]

%dsp-level:0
\[
\qquad ds^2=-N^2d\tau^2+g_{ij}\bigl(dx^i+N^id\tau\bigr)\bigl(dx^j+N^jd\tau\bigr)
\]

%dsp-level:0
\[
\qquad [\mathrm{BM.34}(2);\ \text{machine: the line-element identity}];
\]

%dsp-level:0
\[
\text{at rest}\ \bigl(dx^i=0\bigr):\quad ds=N\,d\tau;
\]

%dsp-level:0
\[
\qquad \text{the tick is the chart coordinate}\ x^0=\tau_n
\]

%dsp-level:0
\[
\qquad [\mathrm{D1.8}];
\]

%dsp-level:0
\[
\qquad N=N(k,E)\ \text{is derived in}\ \mathrm{BM.132};
\]

%dsp-level:0
\[
\boxed{\ V(s)\ =\ \Bigl(1-\frac\gamma4\Bigr)^{2s/N}V_0\ }
\]

%dsp-level:0
\[
\qquad [\text{the reparametrization of}\ (5)\text{'s law by}\ s=Nn;\ \text{machine: at}\ \gamma=1,\ N=1/m,\ m\in\{1,2,3,4\}:\ \bigl(\tfrac9{16}\bigr)^{m}\ \text{exact integer powers, strictly decreasing in}\ m];
\]

%dsp-level:0
\[
\textbf{the two deficit species stay separated:}
\]

%dsp-level:0
\[
\qquad \frac{\gamma(8-\gamma)}{16}\ \ \text{per tick (informational)}
\]

%dsp-level:0
\[
\qquad \text{against}
\]

%dsp-level:0
\[
\qquad 1-N\ \ \text{(metric)};
\]

%dsp-level:0
\[
\neg\text{a numerical identification of the two};
\]

%dsp-level:0
\[
\qquad \text{they join only through}\ (7)\text{'s exponent}
\]

%dsp-level:0
\[
\qquad [\mathrm{BM.37}(2):\ \text{disjoint alphabets}];
\]

%dsp-level:0
\[
\textbf{the scope row:}
\]

%dsp-level:0
\[
\qquad (7)\ \text{is a chart-side reparametrization of a fixed law,}\ \neg\text{a dynamical coupling:}
\]

%dsp-level:0
\[
\qquad \neg\text{a feedback of the contraction into the metric is formed}\ [\mathrm{BM.104}(4)\text{'s freedom-census pattern}];
\]

%dsp-level:0
\[
\text{the tick-rate/field coupling is re-scoped by it: the redshifted contraction read is formed at the point of use, inheriting}\ \mathrm{CT2.6}\text{'s hypotheses}\ [\text{A §0.2a}];
\]

%dsp-level:0
\[
\qquad \text{its absolute-frequency datum remains calibration-species}\ [\mathrm{CT0.3}:\ \tau_c=1\ \text{tick}];
\]

%dsp-level:0
\[
(8)\quad\textbf{status:}
\]

%dsp-level:0
\[
\qquad \textbf{derived:}\ (1)\text{--}(7);
\]

%dsp-level:0
\[
\qquad \text{the duality is the pair}\ (\mu_\star,\ K)\ \text{with the three exact rows: the Haar invariance of}\ V,
\]

%dsp-level:0
\[
\ \text{the Poisson contraction of}\ V,
\]

%dsp-level:0
\[
\ \text{the unique common fixed point}\ \tfrac14I;
\]

%dsp-level:0
\[
\neg\text{a declaration},\qquad\neg\text{a P4 addition},\qquad\neg\text{a new dimensionful scale}\ [\mathrm{P1}],\qquad\neg\text{a measured value}\ [\mathrm{P5}],\qquad
\]

%dsp-level:0
\[
\textbf{the lapse gap is closed by}\ \mathrm{BM.132};
\]

%dsp-level:0
\[
\qquad \text{the row}\ \mathsf g,
\]

%dsp-level:0
\[
\ e_0,
\]

%dsp-level:0
\[
\ N\ \text{of}\ \mathrm{BM.130}(7)\ \text{is re-scoped:}\ N\ \text{derived}\ [\mathrm{BM.132}],
\]

%dsp-level:0
\[
\ \ \mathsf g\ \text{selected}\ [\mathrm{BM.133}],
\]

%dsp-level:0
\[
\ \ e_0\ \text{consolidated as calibration}\ [\mathrm{BM.143}(4)].
\]

\textbf{Req.} A3, A5, A6 (Part 1); D0.3, D0.4, D1.8 (Part 1); CT0.3, CT0.8 (Part 1); L0, L1, L2 (Part 1); CT2.6 (Part 2); P1, P5 (Parts 1–2); FT.107(2) (Part 3); Def-BM.10, BM.34(2), BM.37(2), BM.51(2), BM.51(3), BM.104(4) (this part); Schur's lemma on the $3\oplus1$ decomposition, the Haar-measure invariance and uniqueness on a compact group, formal power-series arithmetic over $\mathbb Q$ with the exact $\exp/\log$ pair, exact rational $\mathrm{SU}(2)$ points from Pythagorean quadruples, exact rank by Fraction elimination, Zariski-density substitution [math; machine: \texttt{checks/\allowbreak{}verify\_\allowbreak{}H\_\allowbreak{}bm\_\allowbreak{}closing\_\allowbreak{}tier.py} §A].

\textbf{Pf.}

%dsp-level:0
\[
(1):\ \mathrm{A6}\ \text{supplies}\ \mu_\star,\ \mathcal T,\ U\notin\mathfrak L;
\]

%dsp-level:0
\[
\qquad \mathrm{D0.3}\ \text{supplies}\ \tilde{\mathcal T}_\gamma;
\]

%dsp-level:0
\[
\qquad \mathrm{L0}\ \text{supplies the one-tilt-per-event coupling (declared scope, inherited at the point of use)};
\]

%dsp-level:0
\[
\qquad \text{the factorization is the composition of the two axiom-level moves}
\]

%dsp-level:0
\[
(2):\ \rho\ \text{faithful}\implies\rho,\ \tilde{\mathcal T}_\gamma(\rho)\ \text{simultaneously diagonalizable}\ [\mathrm{L1}(3)\ \mathrm{Pf}]:
\]

%dsp-level:0
\[
\quad K'
\]

%dsp-level:0
\[
= -\log\rho'
\]

%dsp-level:0
\[
= -\log\rho+\gamma\rho+\log Z_\gamma\cdot I\ \ \text{as an operator identity on the eigenbasis}\ [\text{machine: the}\ \mathbb Q[[t,\gamma]]\ \text{series identity}\ \log\bigl(\lambda e^{-\gamma\lambda}/Z_\gamma\bigr)+\gamma\lambda+\log Z_\gamma-\log\lambda\equiv0\ \text{to order}\ 12];
\]

%dsp-level:0
\[
\rho=\tfrac14I:\quad Z_\gamma=4\cdot\tfrac14e^{-\gamma/4}=e^{-\gamma/4};
\]

%dsp-level:0
\[
\qquad \lambda'=\tfrac14e^{-\gamma/4}/e^{-\gamma/4}=\tfrac14;
\]

%dsp-level:0
\[
\qquad \delta K=\gamma\cdot\tfrac14I+\log\bigl(e^{-\gamma/4}\bigr)I=0\ \ [\text{machine:}\ \exp(-\gamma/4)\exp(\gamma/4)=1\ \text{coefficient-wise}]
\]

%dsp-level:0
\[
(3):\ \mathrm{L1}(2);\qquad K(\mathcal T\rho)=UKU^\dagger\ [\mathrm{L1}(2)\ \mathrm{Pf}:\ \text{functional-calculus covariance}]\implies\operatorname{spec}K\ \text{invariant};\qquad\text{the stationarity of}\ \mu_\star\ \text{under conjugation is the two-sided invariance}\ [\text{math}];\qquad\text{the machine rows:}\ \operatorname{Tr}\bigl((U\rho U^\dagger)^n\bigr)=\operatorname{Tr}(\rho^n),\ n=1,2,3,\ \text{on}\ 40\ \text{exact rational points}\ U=\bigl(\begin{smallmatrix}a+bi&c+di\\-c+di&a-bi\end{smallmatrix}\bigr)/n,\ a^2{+}b^2{+}c^2{+}d^2{=}n^2,\ \text{and generic rational Hermitian}\ \rho
\]

%dsp-level:0
\[
(4):\ \mathbb C^2\otimes\mathbb C^2=\vee^2\oplus\wedge^2\ [\mathrm{A3}]:\quad U\otimes U\ \text{acts irreducibly on each block};
\]

%dsp-level:0
\[
\qquad \text{the twirl of any intertwiner-average projects onto the commutant}\
\]

%dsp-level:0
\[
= \ \operatorname{span}\{\Pi_{\mathcal R},\Pi_{\mathcal V}\}\ [\text{math: Schur}]:\quad\text{the block weights}\ \mathfrak T(\Pi_{\mathcal R}\rho),
\]

%dsp-level:0
\[
\ \mathfrak T(\Pi_{\mathcal V}\rho)\ \text{are the only invariants};
\]

%dsp-level:0
\[
\text{machine:}\ \Pi_{\mathcal R}+\Pi_{\mathcal V}=I,\ \Pi_{\mathcal R}\Pi_{\mathcal V}=0,\ \operatorname{rk}=(3,1);
\]

%dsp-level:0
\[
\qquad \mathrm{tw}\circ\mathrm{tw}=\mathrm{tw};
\]

%dsp-level:0
\[
\qquad \text{the block traces of}\ \mathrm{tw}(\rho)\ \text{equal those of}\ \rho;
\]

%dsp-level:0
\[
\qquad \text{the equivariance on all}\ 40\ \text{rational points};
\]

%dsp-level:0
\[
\qquad x/3=1-x\iff x=\tfrac34\implies\tfrac{x}3\Pi_{\mathcal R}+(1-x)\Pi_{\mathcal V}=\tfrac14I
\]

%dsp-level:0
\[
(5):\ V\bigl(\tilde K(U\tilde{\mathcal T}_\gamma(\rho)U^\dagger)\bigr)=V\bigl(\tilde K(\tilde{\mathcal T}_\gamma(\rho))\bigr)\ [(3)]=\bigl(1-\tfrac\gamma4\bigr)^2V(\tilde K(\rho))+O\ [\mathrm{L2}];
\]

%dsp-level:0
\[
\qquad \text{iteration:}\ \mathrm{BM.51}(3);
\]

%dsp-level:0
\[
\qquad \text{the record}\ V_n\ \text{carries}\ \neg\text{a draw label}\ [\mathrm{A6}:\ U\notin\mathfrak L]
\]

%dsp-level:0
\[
(6):\ \Pi_k^2
\]

%dsp-level:0
\[
= \frac1{(k!)^2}\sum_{\sigma,\tau}\hat P_{\sigma\tau}
\]

%dsp-level:0
\[
= \frac1{k!}\sum_{\tau'}\hat P_{\tau'}
\]

%dsp-level:0
\[
= \Pi_k\ [\mathrm{BM.128}(2)\text{'s product law transported};\ \text{machine: entrywise over}\ \mathbb Q];
\]

%dsp-level:0
\[
\qquad \operatorname{rk}\Pi_k=k+1\ [\text{machine: exact Fraction elimination for}\ k=1..6:\ \text{ranks}\ (2,3,4,5,6,7)];
\]

%dsp-level:0
\[
\qquad \operatorname{im}\Pi_k=\mathrm{Sym}^k\ [\mathrm{Def\text{-}H.38}\ \text{carrier}]
\]

%dsp-level:0
\[
(7):\ ds^2\big|_{dx^i=0}=-N^2d\tau^2\implies ds=N\,d\tau\ [\mathrm{BM.34}(2);\ \text{machine}];
\]

%dsp-level:0
\[
\qquad s=Nn\implies V_n=\bigl(1-\tfrac\gamma4\bigr)^{2n}V_0=\bigl(1-\tfrac\gamma4\bigr)^{2s/N}V_0\ [\text{substitution}];
\]

%dsp-level:0
\[
\qquad \text{the exact powers at}\ \gamma
\]

%dsp-level:0
\[
= 1:\ \bigl(\tfrac9{16}\bigr)^m,
\]

%dsp-level:0
\[
\ m
\]

%dsp-level:0
\[
= 1..4\ [\text{machine: strictly decreasing, exact integers}\ 9^m/16^m];
\]

%dsp-level:0
\[
\text{the species row:}\ \gamma(8-\gamma)/16\in\mathbb Q\ \text{against}\ 1-N\in\mathfrak q^0\ \text{chart data}:\ \neg\text{a derived equation relates them}\ [\mathrm{BM.37}(2)\ \text{pattern}];
\]

%dsp-level:0
\[
\qquad (8):\ \text{assembly}.\ \blacksquare
\]

\textbf{Phys.} The corpus owns exactly two state evolutions, and they are dual in the strict sense: an unrecorded random conjugation drawn from the invariant Haar measure, and a recorded modular update generated by $K=-\log\rho$ whose exact step is $\delta K=\gamma\rho+\log Z_\gamma\cdot I$. The measure side moves states inside the level sets of the Lyapunov function and preserves the modular spectrum; the dynamics side cuts strictly across the levels toward the vacuum. The Haar average of any pair state is its block-maximal-mixed projection, and the tilt's unique fixed point is the maximally mixed state itself, so the two sides share exactly one fixed point: the A5 vacuum. Because the draws are unrecorded, they drop out of every ledger read, and the contraction law the record obeys is purely Poisson-side — the measure is the record's exact gauge freedom. At the finite-group level the same twirl builds the spin carriers: averaging over $S_k$ is the symmetrizer, whose image is $\mathrm{Sym}^k$ of dimension $k+1$. Finally the metric joins: the lapse converts chart ticks to proper time along a resting worldline, so the same exact contraction read per proper time carries the lapse in the exponent, with the informational and metric deficit species kept apart and joined only by that exponent. No feedback is formed, no scale added, nothing measured.

\end{theorem}

\begin{theorem}[{\texorpdfstring{BM.132 — The lapse: the derived function of the envelope ratio and the curvature, its formed domain, and the envelope row}{BM.132 - The lapse: the derived function of the envelope ratio and the curvature, its formed domain, and the envelope row}}]\label{tag:BM-132}
\textbf{Thm.}

%dsp-level:0
\[
(1)\quad\textbf{the source and the potential (re-derived at the point of use):}
\]

%dsp-level:0
\[
\qquad T_{\mu\nu}=\mathrm{diag}(\rho,0,0,0)\ \ \text{at rest}
\]

%dsp-level:0
\[
\qquad [\mathrm{BM.9}(4);\ \mathrm{CT1.22}];
\]

%dsp-level:0
\[
G_{\mu\nu}+\Lambda g_{\mu\nu}=\kappa^2T_{\mu\nu}
\]

%dsp-level:0
\[
\qquad [\mathrm{CT1.23}];
\]

%dsp-level:0
\[
\qquad g_{00}=-(1+2\Phi),
\]

%dsp-level:0
\[
\qquad G^{(1)}_{00}=2\nabla^2\Phi,
\]

%dsp-level:0
\[
\qquad \nabla^2\Phi=\frac{\kappa^2}2\rho=4\pi G_N\rho
\]

%dsp-level:0
\[
\qquad [\mathrm{BM.43}(4);\ \text{machine re-run}];
\]

%dsp-level:0
\[
\Phi(r)=-\frac{G_Nm}{r}\ \ \text{outside the source support},
\]

%dsp-level:0
\[
\qquad G_Nm_{\mathrm{hid}}=\frac{\mathfrak hE_{\mathrm{int}}}4,
\]

%dsp-level:0
\[
\qquad r_s=2G_Nm_{\mathrm{hid}}=\frac{\mathfrak hE_{\mathrm{int}}}2
\]

%dsp-level:0
\[
\qquad [\mathrm{BM.43}(4);\ \mathrm{P3}:\ G_N\ \text{a named combination, quarantined}];
\]

%dsp-level:0
\[
(2)\quad\textbf{the ADM lapse at zero shift:}\qquad N^2=\Omega^2+\gamma^{ij}N_iN_j,\qquad w=0\implies N^i=0,\quad N=\Omega\qquad[\mathrm{BM.34}(1);\ \mathrm{CT1.12}];
\]

%dsp-level:0
\[
\text{the linearized dictionary}\ \Omega^2=1+2\Phi
\]

%dsp-level:0
\[
\qquad [\mathrm{CT1.12}:\ \Omega=e^{\alpha_\Omega a_{00}};\ \text{the conformal mode reads the potential};\ a_{00}=\frac1{2\alpha_\Omega}\ln(1+2\Phi)\ \text{on the formed locus}];
\]

%dsp-level:0
\[
\boxed{\ N(r,E)^2\ =\ 1+2\Phi\ =\ 1-\frac{2G_Nm}{r}\ =\ 1-\frac{r_s}{r}\ =\ 1-\frac{G(E)}{4\pi\,k}\ },
\]

%dsp-level:0
\[
\qquad k:=\frac{r}{\Delta x_{\mathrm B}(E)},
\]

%dsp-level:0
\[
\qquad G(E)=2\pi E_{\mathrm{int}}(E)^2\ln2
\]

%dsp-level:0
\[
\qquad [\mathrm{BM.88}(2)(\mathrm i),(\mathrm{iii});\ \mathrm{BM.43}(4);\ \text{machine: one identity in}\ \mathbb Q[E_{\mathrm{int}},\pi,\ln2,k]];
\]

%dsp-level:0
\[
\textbf{the lapse is a derived function of the envelope ratio and the derived curvature argument:}
\]

%dsp-level:0
\[
\qquad N=N(k,E)\in\mathfrak q^0
\]

%dsp-level:0
\[
\qquad [\mathrm{CT2.1};\ \text{machine: the exponent census}];
\]

%dsp-level:0
\[
(3)\quad\textbf{the formed domain:}
\]

%dsp-level:0
\[
\qquad N^2\ge0\iff k\ \ge\ \frac{G(E)}{4\pi}\ =\ \frac{E_{\mathrm{int}}^2\ln2}2\ =\ \frac{r_s}{\Delta x_{\mathrm B}}
\]

%dsp-level:0
\[
\qquad [\text{machine: the monomial}];
\]

%dsp-level:0
\[
\textbf{the envelope row:}
\]

%dsp-level:0
\[
\qquad \text{on every computed cell}\ E_{\mathrm{int}}\in\{2,6,12,24,60\}\ [\mathrm{BM.13}(2);\ \mathrm{Census\ FT.35}(2),\ \mathrm{FT.42}(5)\ \text{row-wise}]:
\]

%dsp-level:0
\[
\frac{G}{4\pi}=\Bigl(2,\ 18,\ 72,\ 288,\ 1800\Bigr)\ln2\ >\ 1
\]

%dsp-level:0
\[
\qquad \Bigl[\ln2>\frac12\iff4>e:\quad e<\frac{876809}{322560}<4;\ \text{machine: the exact rational enclosure of}\ e\Bigr];
\]

%dsp-level:0
\[
\boxed{\ 1+2\Phi\bigl(\Delta x_{\mathrm B}\bigr)=1-\frac{E_{\mathrm{int}}^2\ln2}2\ <\ 0\ \ \text{on every computed cell}\ }
\]

%dsp-level:0
\[
\qquad [\mathrm{BM.43}(4):\ \Phi(\Delta x_{\mathrm B})=-E_{\mathrm{int}}^2\ln2/4;\ \text{machine}];
\]

%dsp-level:0
\[
\textbf{the interior row (the uniform-ball read):}
\]

%dsp-level:0
\[
\qquad \Phi_{\mathrm{int}}(r)=-\frac{G_Nm}{2R}\Bigl(3-\frac{r^2}{R^2}\Bigr),
\]

%dsp-level:0
\[
\qquad \nabla^2\Phi_{\mathrm{int}}=4\pi G_N\rho,
\]

%dsp-level:0
\[
\qquad \Phi_{\mathrm{int}}(R)=\Phi(R)
\]

%dsp-level:0
\[
\qquad [\text{machine: the harmonicity and matching identities in}\ \mathbb Q[x,y,z,r,R]];
\]

%dsp-level:0
\[
1+2\Phi_{\mathrm{int}}(0)=1-\frac{3E_{\mathrm{int}}^2\ln2}4<0\ \ \text{on every computed cell}
\]

%dsp-level:0
\[
\qquad \Bigl[E_{\mathrm{int}}=2:\ 3\ln2>1\ \text{on}\ \ln2>\frac13:\ e^3<8;\ \text{machine}\Bigr];
\]

%dsp-level:0
\[
\boxed{\ \textbf{the self-lapse of a sealed composite is}\ \neg\textbf{formed at its own Born width: the envelope lies strictly inside the domain boundary on every computed cell}\ }
\]

%dsp-level:0
\[
(4)\quad\textbf{the monotonicity, the limits, and the gauge agreements:}
\]

%dsp-level:0
\[
\qquad \frac{\partial N^2}{\partial k}=\frac{G}{4\pi k^2}>0:\quad N\ \text{strictly increasing in}\ k\ \text{on the formed domain};
\]

%dsp-level:0
\[
\qquad N\ \text{strictly decreasing in}\ E_{\mathrm{int}}
\]

%dsp-level:0
\[
\qquad [\text{machine: sign census on exact rationals}];
\]

%dsp-level:0
\[
N=0\ \ \text{at the boundary}\ \ k=\frac{G}{4\pi};\qquad N\to1^-\ \ \text{as}\ \ k\to\infty;\qquad N\in(0,1]\ \ \text{on the formed locus}\qquad[\text{machine}];
\]

%dsp-level:0
\[
\textbf{the vacuum gauge agreement:}
\]

%dsp-level:0
\[
\qquad (a_{00},q_{ij},w_i)=0\implies N=1
\]

%dsp-level:0
\[
\qquad [\mathrm{BM.34}(2);\ \mathrm{BM.88}(2)(\mathrm{iii});\ \text{machine}]
\]

%dsp-level:0
\[
\qquad \text{and}
\]

%dsp-level:0
\[
\qquad \lim_{k\to\infty}N(k,E)=1;
\]

%dsp-level:0
\[
\textbf{T-evenness:}
\]

%dsp-level:0
\[
\qquad N\ \text{carries}\ w\ \text{only through}\ \gamma^{ij}N_iN_j=N^2-\Omega^2;\quad\text{at zero shift}\ N\ \text{is}\ w\text{-free}
\]

%dsp-level:0
\[
\qquad [\mathrm{BM.34}(3);\ \text{machine}];
\]

%dsp-level:0
\[
(5)\quad\textbf{the three substitutions into derived objects:}
\]

%dsp-level:0
\[
\qquad \textbf{(i) the dispersion:}\quad E=N(k,E)\sqrt{m^2+\gamma^{ij}p_ip_j}-N^ip_i
\]

%dsp-level:0
\[
\qquad [\mathrm{BM.34}(4);\ \text{the lapse entry now carries}\ (2)\text{'s function}];
\]

%dsp-level:0
\[
\textbf{(ii) the threshold:}
\]

%dsp-level:0
\[
\qquad \mathcal T(k,E)=\frac{\mathcal T_0}{N(k,E)}=\frac{E_{\mathrm{int}}\,\mathfrak h}{\ln2}\Bigl(1-\frac{E_{\mathrm{int}}^2\ln2}{2k}\Bigr)^{-1/2}
\]

%dsp-level:0
\[
\qquad [\mathrm{BM.115}(1);\ \mathrm{Def\text{-}BM.17}(3);\ \text{machine: the substitution}]:
\]

%dsp-level:0
\[
\mathcal T\to\infty\ \ \text{at the domain boundary};
\]

%dsp-level:0
\[
\qquad \mathcal T\to\mathcal T_0\ \ \text{as}\ \ k\to\infty;
\]

%dsp-level:0
\[
\qquad \textbf{a deeper lapse lowers every threshold at once, and the boundary locus raises every threshold without bound}\ [\mathrm{BM.115}(1)\ \text{re-cited}];
\]

%dsp-level:0
\[
\textbf{(iii) the coupling profile:}
\]

%dsp-level:0
\[
\qquad e(r,E)=e_0\sqrt{N(k,E)\cdot f(k)\cdot\mathsf g\bigl(G(E)\bigr)}
\]

%dsp-level:0
\[
\qquad [\mathrm{BM.88}(4);\ \mathsf g\ \text{carried symbolically here, selected in}\ \mathrm{BM.133}];
\]

%dsp-level:0
\[
(6)\quad\textbf{the bars:}
\]

%dsp-level:0
\[
\qquad r_s\ \text{enters as}\ \mathrm{P3}\text{'s named combination, quarantined:}\ \neg\text{a collapse reading,}\ \neg\text{a horizon identification,}\ \neg\text{black-hole physics}\ [\mathrm{P3};\ \mathrm{BM.15}(4);\ \mathrm{BM.115}(4);\ \mathrm{BM.130}(6)\ \text{discipline inherited}];
\]

%dsp-level:0
\[
\neg\text{a cosmological coordinate}\ [\mathrm{Def\text{-}BM.17}(6)(\mathrm i)];
\]

%dsp-level:0
\[
\qquad \neg\text{a measured curvature, radius or redshift}\ [\mathrm{P5}];
\]

%dsp-level:0
\[
\qquad \text{the weak-field rows inherit}\ \mathrm{CT2.6}\text{'s hypotheses at the point of use}\ [\text{A §0.2a};\ \mathrm{BM.43}(4)\ \text{scope}];
\]

%dsp-level:0
\[
\text{the full non-linear joint solution remains}\ \neg\text{constructed}\ [\mathrm{BM.43}(5)\ \text{re-cited,}\ \neg\text{extended}];
\]

%dsp-level:0
\[
(7)\quad\textbf{status:}
\]

%dsp-level:0
\[
\qquad \textbf{derived:}\ (1)\text{--}(5);
\]

%dsp-level:0
\[
\qquad \boxed{\ \textbf{the lapse gap is closed:}\ \ N(k,E)=\sqrt{1-\frac{G(E)}{4\pi k}}\ \ \text{on the formed domain}\ \ k\ge\frac{G(E)}{4\pi}\ };
\]

%dsp-level:0
\[
\qquad \text{the row}\ N\ \text{as solution data}\ [\mathrm{BM.88}(8);\ \mathrm{BM.93}(6);\ \mathrm{BM.100}(6);\ \mathrm{BM.107}(7);\ \mathrm{BM.118}(7);\ \mathrm{BM.130}(7)]\ \textbf{is discharged}.
\]

\textbf{Req.} A8 (Part 1); D2.1 (Part 1); CT1.12, CT1.22, CT1.23, CT2.1, CT2.6, CT2.9 (Part 2); P1, P3, P5 (Parts 1–2); Def-BM.4, Def-BM.10, Def-BM.17(3), Def-BM.17(6), BM.9(4), BM.13(2), BM.15(4), BM.34(1), BM.34(2), BM.34(3), BM.34(4), BM.43(4), BM.43(5), BM.88(2), BM.88(3), BM.88(4), BM.88(8), BM.115(1), BM.115(4), BM.130(6), BM.130(7) (this part); Census FT.35(2), FT.42(5) (row-wise); the exact rational enclosure of $e$ by the factorial series, the harmonicity of $1/r$ and the interior-ball potential in quotient rings, Laurent and monomial identities, sign censuses on exact rationals [math; machine: \texttt{checks/\allowbreak{}verify\_\allowbreak{}H\_\allowbreak{}bm\_\allowbreak{}closing\_\allowbreak{}tier.py} §A–§B].

\textbf{Pf.}

%dsp-level:0
\[
(1):\ \mathrm{BM.43}(4)\text{'s rows re-derived:}\ G^{(1)}_{00}
\]

%dsp-level:0
\[
= 2\nabla^2\Phi\ [\text{machine: the jet-ring Christoffel}\to\text{Ricci}\to\text{Einstein chain}];
\]

%dsp-level:0
\[
\qquad \nabla^2(1/r)=0\ [\text{machine: the quotient-ring identity}];
\]

%dsp-level:0
\[
\qquad G_Nm_{\mathrm{hid}}=\frac{\mathfrak h^2}{4\ln2}\cdot\frac{E_{\mathrm{int}}\ln2}{\mathfrak h}=\frac{\mathfrak hE_{\mathrm{int}}}4\ [\mathrm{P3};\ \mathrm{FT.44}(4);\ \text{machine}]
\]

%dsp-level:0
\[
(2):\ w=0\implies N_i=\Omega^2\varkappa\beta w_i=0\implies N^2=\Omega^2\ [\mathrm{BM.34}(1)];
\]

%dsp-level:0
\[
\qquad \text{the linearized dictionary:}\ g_{00}
\]

%dsp-level:0
\[
= -\Omega^2
\]

%dsp-level:0
\[
= -(1+2\Phi)\ [\mathrm{CT1.12};\ \mathrm{BM.34}(2)\text{'s line element at zero shift}];
\]

%dsp-level:0
\[
\qquad \text{the substitution:}\ 1-\frac{2G_Nm}r=1-\frac{r_s}r;
\]

%dsp-level:0
\[
\qquad \frac{r_s}{\Delta x_{\mathrm B}}
\]

%dsp-level:0
\[
= \frac{\mathfrak hE_{\mathrm{int}}/2}{\mathfrak h/(E_{\mathrm{int}}\ln2)}
\]

%dsp-level:0
\[
= \frac{E_{\mathrm{int}}^2\ln2}2
\]

%dsp-level:0
\[
= \frac{G}{4\pi}\ [\mathrm{BM.88}(3);\ \text{machine: one monomial}]
\]

%dsp-level:0
\[
(3):\ N^2\ge0\iff\frac{G}{4\pi k}\le1;
\]

%dsp-level:0
\[
\qquad E_{\mathrm{int}}\ge2\ [\mathrm{BM.13}(3)]\implies\frac{G}{4\pi}\ge2\ln2>1\ [\text{machine:}\ e<\tfrac{876809}{322560}<4\implies\ln4>1];
\]

%dsp-level:0
\[
\qquad \text{the interior row: the matching identities and}\ 3E_{\mathrm{int}}^2\ln2>4\ \text{at}\ E_{\mathrm{int}}=2\ [\text{machine:}\ e^3<8]
\]

%dsp-level:0
\[
(4):\ \text{the derivative sign and the endpoint values on exact rationals}\ [\text{machine}];
\]

%dsp-level:0
\[
\qquad (5):\ \text{substitution into}\ \mathrm{BM.34}(4),\ \mathrm{BM.115}(1),\ \mathrm{BM.88}(4)\ [\text{machine: the monomial rows}];
\]

%dsp-level:0
\[
\qquad (6),(7):\ \text{assembly and the rescope census}.\ \blacksquare
\]

\textbf{Phys.} The lapse was solution data; it is now a function. The composite's own dust source runs through the derived field equations to the weak-field potential, the ADM split reads the potential as the zero-shift lapse, and the curvature lock turns the answer into one dimensionless monomial: $N^2=1-G/(4\pi k)$ on the envelope ratio. The function comes with its own formed domain, and the domain has a sharp edge: on every computed cell the boundary $k=G/(4\pi)$ lies strictly beyond the Born width, so a sealed composite's self-lapse is not formed at its own envelope — the interior read only deepens the failure. The vacuum gauge is recovered exactly at infinity, the threshold ladder inherits the lapse and diverges at the boundary, and the running coupling's energy factor is now explicit. No black-hole reading is attached to the boundary, the Schwarzschild combination stays quarantined, and no measured number enters.

\end{theorem}

\begin{theorem}[{\texorpdfstring{BM.133 — The selected curvature composition: two derived criteria, the unique survivor, and the completed running coupling}{BM.133 - The selected curvature composition: two derived criteria, the unique survivor, and the completed running coupling}}]\label{tag:BM-133}
\textbf{Thm.}

%dsp-level:0
\[
(1)\quad\textbf{the three priced rows (re-cited, frozen):}\qquad\mathsf g\in\Bigl\{1,\quad1+G,\quad\frac1{1+G}\Bigr\}\qquad[\mathrm{BM.88}(4);\ \mathrm{BM.104}(2)]:
\]

%dsp-level:0
\[
\text{the back-reaction fixed-point rows:}
\]

%dsp-level:0
\[
\qquad (\mathrm i)\ \ \mathsf g\equiv1;
\]

%dsp-level:0
\[
\qquad (\mathrm{ii})\ \ \mathsf g_{\mathrm{fp}}=\frac{1+G}{1-aG},\ \ \text{admissible}\iff aG<1;
\]

%dsp-level:0
\[
\qquad (\mathrm{iii})\ \ aG\,\mathsf g^2+(1+G)\,\mathsf g-1=0,\ \ \mathsf g_{\mathrm{fp}}=\frac{\sqrt{(1+G)^2+4aG}-(1+G)}{2aG}
\]

%dsp-level:0
\[
\qquad [\mathrm{BM.104}(2);\ \text{machine re-run}];
\]

%dsp-level:0
\[
a=\tfrac12Nw\,\varepsilon_{\mathrm{EM}}Q^2\ge0,\qquad G>0\qquad[\mathrm{BM.104}(1),(3);\ \mathrm{BM.88}(3)];
\]

%dsp-level:0
\[
(2)\quad\textbf{criterion}\ (\mathrm S1)\ \textbf{— the Lyapunov contractivity of the back-reaction map:}
\]

%dsp-level:0
\[
\qquad F_{(\mathrm i)}(\mathsf g)=1;
\]

%dsp-level:0
\[
\qquad F_{(\mathrm{ii})}(\mathsf g)=1+G+aG\,\mathsf g;
\]

%dsp-level:0
\[
\qquad F_{(\mathrm{iii})}(\mathsf g)=\frac1{1+G+aG\,\mathsf g}
\]

%dsp-level:0
\[
\qquad [\mathrm{BM.104}(1)\text{'s loop read as a map}]:
\]

%dsp-level:0
\[
F'_{(\mathrm i)}=0;
\]

%dsp-level:0
\[
\qquad F'_{(\mathrm{ii})}=aG:\quad\text{contractive}\iff aG<1\ —\ \text{conditional, with the pole at}\ aG=1\ \text{where}\ \mathsf g_{\mathrm{fp}}\ \text{diverges}
\]

%dsp-level:0
\[
\qquad [\mathrm{BM.104}(2)(\mathrm{ii}),\ (3);\ \text{machine}];
\]

%dsp-level:0
\[
F'_{(\mathrm{iii})}(\mathsf g)=-\frac{aG}{(1+G+aG\,\mathsf g)^2}:
\]

%dsp-level:0
\[
\qquad \text{at the fixed point}\ 1+G+aG\,\mathsf g_{\mathrm{fp}}=\frac1{\mathsf g_{\mathrm{fp}}}\implies\boxed{\ \bigl|F'_{(\mathrm{iii})}\bigr|=aG\,\mathsf g_{\mathrm{fp}}^2=1-(1+G)\,\mathsf g_{\mathrm{fp}}\ }
\]

%dsp-level:0
\[
\textbf{the unconditional contraction:}
\]

%dsp-level:0
\[
\qquad 0<\bigl|F'_{(\mathrm{iii})}\bigr|<1\iff0<\mathsf g_{\mathrm{fp}}<\frac1{1+G}
\]

%dsp-level:0
\[
\qquad \text{and}
\]

%dsp-level:0
\[
\qquad (1+G)^2\bigl((1+G)^2+4aG\bigr)<\bigl((1+G)^2+2aG\bigr)^2\iff0<4a^2G^2
\]

%dsp-level:0
\[
\qquad [\text{machine: the squaring argument, exact in}\ \mathbb Q[G,a];\ \text{both sides positive on the formed locus}];
\]

%dsp-level:0
\[
\boxed{\ \textbf{row}\ (\mathrm{iii})\ \textbf{is unconditionally Lyapunov-contractive on the whole formed locus; row}\ (\mathrm{ii})\ \textbf{is not; row}\ (\mathrm i)\ \textbf{is contractive but}\ G\textbf{-inert}\ }
\]

%dsp-level:0
\[
\qquad [\mathrm{L1}\text{--}\mathrm{L2}\ \text{contraction discipline};\ \mathrm{BM.51}(2)\ \text{pattern}];
\]

%dsp-level:0
\[
(3)\quad\textbf{criterion}\ (\mathrm S2)\ \textbf{— the minimal lapse-domain loss:}
\]

%dsp-level:0
\[
\qquad \text{the corrected curvature}\ G_{\mathrm{corr}}=G(1+a\,\mathsf g)\ [\mathrm{BM.104}(1)]\ \text{moves the lapse boundary:}
\]

%dsp-level:0
\[
\qquad N_{\mathrm{corr}}(k)^2=1-\frac{G(1+a\,\mathsf g)}{4\pi k}\ \ \text{formed}\iff k\ge\frac{G(1+a\,\mathsf g)}{4\pi}
\]

%dsp-level:0
\[
\qquad [\mathrm{BM.132}(2),(3)];
\]

%dsp-level:0
\[
\text{every positive composition shrinks the formed domain}\ (a\,\mathsf g>0);
\]

%dsp-level:0
\[
\qquad \textbf{the loss is exactly ordered:}
\]

%dsp-level:0
\[
\qquad \boxed{\ a\,\mathsf g_{\mathrm{fp}}^{(\mathrm{iii})}\ <\ \frac{a}{1+G}\ <\ a\ =\ a\,\mathsf g^{(\mathrm i)}\ <\ a\,(1+G)\ =\ a\,\mathsf g^{(\mathrm{ii})}\ }
\]

%dsp-level:0
\[
\qquad [\text{machine: the strict chain on}\ a,G>0];
\]

%dsp-level:0
\[
\textbf{row}\ (\mathrm{ii})\ \textbf{at the pole loses the whole domain:}
\]

%dsp-level:0
\[
\qquad aG\to1:\quad\mathsf g_{\mathrm{fp}}\to\infty,\quad G_{\mathrm{corr}}\to\infty,\quad\{k:\ N_{\mathrm{corr}}\ \text{formed}\}\to\varnothing
\]

%dsp-level:0
\[
\qquad [\mathrm{BM.104}(2)(\mathrm{ii}),\ (3);\ \text{machine}];
\]

%dsp-level:0
\[
\boxed{\ \textbf{row}\ (\mathrm{iii})\ \textbf{loses the least formed lapse domain of the three, uniformly in}\ (a,G)\ }
\]

%dsp-level:0
\[
(4)\quad\textbf{the unique survivor:}
\]

%dsp-level:0
\[
\qquad (\mathrm S1)\wedge(\mathrm S2):\quad\text{row}\ (\mathrm{ii})\ \text{fails both (conditionality, maximal loss with a total collapse at the pole)};
\]

%dsp-level:0
\[
\qquad \text{row}\ (\mathrm i)\ \text{passes}\ (\mathrm S1)\ \text{and fails}\ (\mathrm S2)\ \text{and is}\ G\text{-inert}\ (2);
\]

%dsp-level:0
\[
\boxed{\ \mathsf g(G)\ =\ \frac1{1+G}\ }
\]

%dsp-level:0
\[
\qquad \textbf{selected by two derived criteria;}
\]

%dsp-level:0
\[
\qquad \mathsf g\in\mathfrak q^0,
\]

%dsp-level:0
\[
\qquad 0<\mathsf g\le1,
\]

%dsp-level:0
\[
\qquad \mathsf g\ \text{strictly decreasing in}\ G,
\]

%dsp-level:0
\[
\qquad \mathsf g(0)=1\ \ \text{the common vacuum-normalization row of all three}\ [\mathrm{BM.104}(2);\ \text{machine}];
\]

%dsp-level:0
\[
\textbf{the exact values at the computed cells:}
\]

%dsp-level:0
\[
\qquad \mathsf g=\frac1{1+2\pi E_{\mathrm{int}}^2\ln2},
\]

%dsp-level:0
\[
\qquad E_{\mathrm{int}}\in\{2,6,12,24,60\}:\quad\mathsf g\in\Bigl\{\frac1{1+8\pi\ln2},\ \frac1{1+72\pi\ln2},\ \frac1{1+288\pi\ln2},\ \frac1{1+1152\pi\ln2},\ \frac1{1+7200\pi\ln2}\Bigr\}
\]

%dsp-level:0
\[
\qquad [\text{machine: the monomial substitutions, carried exactly}];
\]

%dsp-level:0
\[
\textbf{the observable consequence (re-cited):}
\]

%dsp-level:0
\[
\qquad \mathrm{BM.101}(5)\text{'s crossing count on the selected row is two, one or none}\ —\ \text{the composition is observable in the energy-scale graph's topology, and the selection is now made}\ [\mathrm{BM.104}(6)\ \text{re-scoped}];
\]

%dsp-level:0
\[
(5)\quad\textbf{the}\ \mathrm{BM.104}\ \textbf{bars are respected:}
\]

%dsp-level:0
\[
\qquad \textbf{the loop stays excluded:}\ G_{\mathrm{corr}}\ \neg\text{re-enters}\ G;
\]

%dsp-level:0
\[
\ \text{the lock's argument remains the informational mass}\ [\mathrm{BM.104}(4)\ \text{frozen};\ \mathrm{BM.88}(3);\ \text{machine: the alphabet census}\ \{E_{\mathrm{int}},\ln2,\pi\}\cap\{N,w,\varepsilon_{\mathrm{EM}},Q\}=\varnothing\ \text{re-run}];
\]

%dsp-level:0
\[
\textbf{the four grounds stand and none is used:}
\]

%dsp-level:0
\[
\qquad (\mathrm i)\ \neg\text{a count}\to\text{strength map enters}\ [\mathrm{BM.36}(3);\ \mathrm{BM.61}(8)(\mathrm{O1})];
\]

%dsp-level:0
\[
\qquad (\mathrm{ii})\ \neg\text{a defect-conditioned weighting}\ [\mathrm{BM.90}(3)];
\]

%dsp-level:0
\[
\qquad (\mathrm{iii})\ \text{consistency is}\ \neg\text{used as sufficient}\ —\ \text{the selection runs on}\ (\mathrm S1),
\]

%dsp-level:0
\[
(\mathrm S2),
\]

%dsp-level:0
\[
\ \text{which are stability and metric-formedness criteria, not the fixed-point equation};
\]

%dsp-level:0
\[
\qquad (\mathrm{iv})\ \neg\text{a P4 datum: the survivor is unique, so no choice remains}\ [\mathrm{BM.104}(5)(\mathrm{iv})\ \text{discharged by uniqueness}];
\]

%dsp-level:0
\[
\mathrm{BM.104}(5)\text{'s}\ \neg\exists\text{-row is re-scoped to the four grounds it cites,}\ \neg\text{retracted};
\]

%dsp-level:0
\[
\qquad \mathrm{BM.104}(6)\text{'s discrimination row is completed by the selection};
\]

%dsp-level:0
\[
(6)\quad\textbf{the completed running coupling:}
\]

%dsp-level:0
\[
\qquad \boxed{\ e(r,E)\ =\ e_0\,\sqrt{\,N(k,E)\cdot f(k)\cdot\mathsf g\bigl(G(E)\bigr)\,}\ =\ e_0\left[\frac{\bigl(1-\frac{G}{4\pi k}\bigr)^{1/2}\bigl(1-\frac1k\bigr)}{1+G}\right]^{1/2}\ }
\]

%dsp-level:0
\[
\qquad [\mathrm{BM.88}(4);\ \mathrm{BM.132}(2);\ \text{machine: the assembled monomial}];
\]

%dsp-level:0
\[
\textbf{all three factors are derived:}
\]

%dsp-level:0
\[
\qquad N(k,E)\ [\mathrm{BM.132}];
\]

%dsp-level:0
\[
\qquad f(k)=1-\frac1k\ [\mathrm{BM.88}(2)(\mathrm{ii});\ \mathrm{BM.47}(6)];
\]

%dsp-level:0
\[
\qquad \mathsf g(G)=\frac1{1+G}\ [(4)];
\]

%dsp-level:0
\[
\qquad \textbf{only}\ e_0\ \textbf{remains calibration}\ [\mathrm{EM.25}(7);\ \mathrm{BM.143}(4)];
\]

%dsp-level:0
\[
\textbf{the formed locus:}
\]

%dsp-level:0
\[
\qquad k\ \ge\ \frac{G}{4\pi}\ \ \wedge\ \ k>1:
\]

%dsp-level:0
\[
\qquad \text{both factors strictly positive on its interior};
\]

%dsp-level:0
\[
\qquad e\to0\ \ \text{at the lapse boundary}\ (N\to0);
\]

%dsp-level:0
\[
\qquad e\to\frac{e_0}{\sqrt{1+G}}\ \ \text{as}\ \ k\to\infty
\]

%dsp-level:0
\[
\qquad [\text{machine: the exact limit rows}];
\]

%dsp-level:0
\[
\textbf{the monotonicity:}
\]

%dsp-level:0
\[
\qquad k\mapsto N(k)f(k)\ \text{strictly increasing on the formed locus}
\]

%dsp-level:0
\[
\qquad [\text{both factors positive and strictly increasing}\ [\mathrm{BM.88}(2)(\mathrm{ii});\ (4)];\ \text{machine: the derivative sign on exact rationals}];
\]

%dsp-level:0
\[
\mathrm{BM.88}(5)\text{'s envelope zero}\ f(1)
\]

%dsp-level:0
\[
= 0\ \textbf{is re-scoped:}\ \text{it is a vacuum-gauge read}\ (N{=}1);
\]

%dsp-level:0
\[
\ \text{on the derived self-lapse the envelope locus itself is}\ \neg\text{formed}\ [\mathrm{BM.132}(3)];
\]

%dsp-level:0
\[
\qquad \text{the formed coupling is strictly positive on its whole formed locus and vanishes only at the boundary};
\]

%dsp-level:0
\[
(7)\quad\textbf{status:}
\]

%dsp-level:0
\[
\qquad \textbf{derived:}\ (1)\text{--}(6);
\]

%dsp-level:0
\[
\qquad \textbf{the curvature composition is closed:}\ \mathsf g(G)=\frac1{1+G};
\]

%dsp-level:0
\[
\qquad \text{the row}\ \mathsf g\ \text{of}\ \mathrm{BM.130}(7)\ \textbf{is discharged};
\]

%dsp-level:0
\[
\neg\text{a declaration}\ (\text{the survivor is unique}),
\]

%dsp-level:0
\[
\qquad \neg\text{a P4 addition},
\]

%dsp-level:0
\[
\qquad \neg\text{a new dimensionful scale}\ [\mathrm{P1}],
\]

%dsp-level:0
\[
\qquad \neg\text{a measured coupling}\ [\mathrm{P5}].
\]

\textbf{Req.} A4 (Part 1); D0.3, D0.4 (Part 1); CT2.1 (Part 2); P1, P4, P5 (Parts 1–2); L1, L2 (Part 1); EM.25(7) (Part 2); Def-BM.12, BM.36(3), BM.47(6), BM.51(2), BM.61(8), BM.88(2), BM.88(3), BM.88(4), BM.88(5), BM.90(3), BM.101(5), BM.104, BM.130(7), BM.132(2), BM.132(3) (this part); the squaring argument in $\mathbb Q[G,a]$, the strict-order chain on positive monomials, exact rational limit and derivative sign rows, the radical identity $(1+G)\mathsf g_{\mathrm{fp}}<1$ in $\mathbb Q[G,a,s]/(s^2-(1+G)^2-4aG)$ [math; machine: \texttt{checks/\allowbreak{}verify\_\allowbreak{}H\_\allowbreak{}bm\_\allowbreak{}closing\_\allowbreak{}tier.py} §B].

\textbf{Pf.}

%dsp-level:0
\[
(1):\ \mathrm{BM.104}(2)\text{'s rows re-verified by substitution}\ [\text{machine}];
\]

%dsp-level:0
\[
\qquad (2):\ F'_{(\mathrm{iii})}\ \text{at the fixed point:}\ 1+G+aG\mathsf g_{\mathrm{fp}}
\]

%dsp-level:0
\[
= 1/\mathsf g_{\mathrm{fp}}\ [\text{the fixed-point equation rearranged}]
\]

%dsp-level:0
\[
\implies |F'|
\]

%dsp-level:0
\[
= aG\mathsf g_{\mathrm{fp}}^2;
\]

%dsp-level:0
\[
\qquad aG\mathsf g_{\mathrm{fp}}^2=1-(1+G)\mathsf g_{\mathrm{fp}}\ [\text{machine: the equation}\ aG\mathsf g^2+(1+G)\mathsf g-1=0\ \text{at}\ \mathsf g=\mathsf g_{\mathrm{fp}}];
\]

%dsp-level:0
\[
\mathsf g_{\mathrm{fp}}<\tfrac1{1+G}:
\]

%dsp-level:0
\[
\quad (1+G)\sqrt{(1+G)^2+4aG}<(1+G)^2+2aG
\]

%dsp-level:0
\[
\iff (1+G)^2\bigl((1+G)^2+4aG\bigr)<\bigl((1+G)^2+2aG\bigr)^2
\]

%dsp-level:0
\[
\iff 0<4a^2G^2\ [\text{machine: expansion in}\ \mathbb Q[G,a]]
\]

%dsp-level:0
\[
(3):\ \text{the domain-boundary monomial}\ \frac{G(1+a\mathsf g)}{4\pi}\ \text{is strictly increasing in}\ \mathsf g>0;
\]

%dsp-level:0
\[
\qquad \text{the chain}\ a\mathsf g_{\mathrm{fp}}<\frac a{1+G}<a<a(1+G)\ \text{on}\ (2)\text{'s inequality and}\ 1+G>1\ [\text{machine}];
\]

%dsp-level:0
\[
\qquad \text{the pole row:}\ (1-aG)\mathsf g_{\mathrm{fp}}=1+G\ \text{and}\ aG=1\implies0=1+G\ne0\ [\text{machine}]
\]

%dsp-level:0
\[
(4):\ \text{the census of}\ (2),(3)\ \text{over the three rows};
\]

%dsp-level:0
\[
\qquad \text{the cell values:}\ G=2\pi E_{\mathrm{int}}^2\ln2\ \text{substituted}\ [\text{machine: monomials}];
\]

%dsp-level:0
\[
\qquad (5):\ \text{the alphabet census re-run and the four grounds re-cited at their loci}\ [\text{machine}];
\]

%dsp-level:0
\[
(6):\ e^2=e_0^2\,N\,f\,\mathsf g\ \text{with}\ (4)\text{'s}\ \mathsf g\ \text{and}\ \mathrm{BM.132}(2)\text{'s}\ N\ [\text{machine: the assembled product}];
\]

%dsp-level:0
\[
\qquad \frac{d}{dk}\bigl[Nf\bigr]>0\ \text{on}\ k>\max\bigl(1,\frac G{4\pi}\bigr):\ \text{both factors positive increasing}\ [\text{machine: rational grid and product rule sign census}];
\]

%dsp-level:0
\[
\qquad (7):\ \text{assembly}.\ \blacksquare
\]

\textbf{Phys.} Three curvature compositions were priced and none was selected, because consistency alone is satisfied three times over. Two derived criteria break the tie. The first is the corpus's own contraction discipline: the back-reaction map of row (iii) is a strict contraction at its fixed point everywhere on the formed locus, row (ii) is contractive only below its pole and collapses there, row (i) is inert. The second is the metric's: every composition inflates the corrected horizon locus and shrinks the formed domain of the lapse, and row (iii) shrinks it least, uniformly. The unique survivor is $\mathsf g=1/(1+G)$ — curvature-diluted, vacuum-normalized, strictly decreasing. The back-reaction loop stays excluded: the lock's argument remains the informational mass, and the selection consumes stability and formedness, never a measured value or a free parameter. With the lapse derived and the composition selected, the running coupling is complete as a function of $(k,E)$ — every factor derived, every zero located, only $e_0$ left to calibration.

\end{theorem}

\begin{definition}[{\texorpdfstring{Def-BM.20 — The momentum-space $\Sigma$-resolvent (the construction BM.6(6) specifies)}{Def-BM.20 - The momentum-space -resolvent (the construction BM.6(6) specifies)}}]\label{tag:Def-BM-20}
\textbf{Def.} A chart-side interface definition on $\mathrm{BM.34}$'s ADM split and $\mathrm{CT2.6}$'s charts; the non-derived clauses are stated exactly and priced inside Part 4; no new axiom, no new datum, no new dimensionful scale. Logged as §II.4 addendum A-BM.12 before use.

%dsp-level:0
\[
(1)\quad\textbf{the boundary-channel states:}
\]

%dsp-level:0
\[
\qquad |p\rangle\ :=\ \text{the chart plane-wave state of covariant momentum label}\ p_\mu\ \text{on}\ \mathrm{CT2.6}\text{'s arena}
\]

%dsp-level:0
\[
\qquad [\text{math: plane waves on the chart}];
\]

%dsp-level:0
\[
\textbf{non-derived clause (the embedding):}
\]

%dsp-level:0
\[
\qquad \text{the identification of the boundary-channel register's momentum label with the chart's}\ p_\mu\ \text{— an}\ \iota_\Sigma\text{-type embedding at the momentum level}\ [\mathrm{BM.6}(6)];
\]

%dsp-level:0
\[
\textbf{scope:}
\]

%dsp-level:0
\[
\qquad \mathrm{CT2.6}\text{'s continuum-limit hypotheses are inherited at every point of use}\ [\text{A §0.2a}];
\]

%dsp-level:0
\[
\qquad \text{the states are chart objects,}\ \neg\text{ledger objects};
\]

%dsp-level:0
\[
\qquad \neg\text{a canonical pair is imported into the finite register}\ [\mathrm{BM.10}(1)\ \text{stands}];
\]

%dsp-level:0
\[
(2)\quad\textbf{the resolvent:}\qquad\boxed{\ R(z;\mathbf p)\ :=\ \frac1{\bigl(p_0-N^ip_i\bigr)^2-N^2\bigl(m^2+\gamma^{ij}p_ip_j\bigr)}\bigg|_{p_0=-z}\ }
\]

%dsp-level:0
\[
\textbf{constitutive clause:}
\]

%dsp-level:0
\[
\qquad \text{the denominator is the derived mass-shell operator}\ —\ \text{its zero locus is the dispersion relation}\ [\mathrm{BM.34}(4)];
\]

%dsp-level:0
\[
\qquad m=m_{\mathrm{hid}}(\mathcal R)\ \text{the composite's dispersion mass}\ [\mathrm{Def\text{-}BM.10};\ \mathrm{BM.34}(7)];
\]

%dsp-level:0
\[
(3)\quad\textbf{the Fourier convention (declared):}
\]

%dsp-level:0
\[
\qquad \text{the chart measure carries}\ \frac{d^3p}{(2\pi)^3}\ \text{per spatial dimension};
\]

%dsp-level:0
\[
\qquad \textbf{cost:}\ \neg\text{a scale,}\ \neg\text{a P4 member}\ —\ \text{any other normalization rescales every amplitude read off it by a declared constant, so the convention is a naming, not a datum}\ [\mathrm{D2.9}:\ \hbar=1];
\]

%dsp-level:0
\[
(4)\quad\textbf{the vertex normalization (constitutive minimality):}
\]

%dsp-level:0
\[
\qquad \mathcal V=1:
\]

%dsp-level:0
\[
\qquad \textbf{the reduction clause:}\ \ \text{at}\ e
\]

%dsp-level:0
\[
= 1,
\]

%dsp-level:0
\[
\ q
\]

%dsp-level:0
\[
= 1\ \text{every amplitude of}\ \mathrm{BM.142}\ \text{reduces to the register's channel weight}\ [\mathrm{BM.122}(3);\ \mathrm{BM.123}(1)]:
\]

%dsp-level:0
\[
\textbf{price:}
\]

%dsp-level:0
\[
\qquad \text{the unique normalization with this reduction}\ —\ \text{any other choice is an}\ e_0\text{-species rescaling}\ [\mathrm{EM.25}(7);\ \text{calibration,}\ \neg\text{structure}];
\]

%dsp-level:0
\[
(5)\quad\textbf{the bars:}
\]

%dsp-level:0
\[
\qquad \neg\text{an S-matrix unitarity claim}\ (\text{the formed reads are tree-level symbolic});
\]

%dsp-level:0
\[
\qquad \text{amplitude}\ne\text{mass}\ [\mathrm{FT.79}(5)];
\]

%dsp-level:0
\[
\qquad \neg\text{a numerical cross-section,}\ \neg\text{a measured comparison}\ [\mathrm{P5}];
\]

%dsp-level:0
\[
\neg\text{an}\ i\varepsilon\text{-prescription import:}\ R\ \text{is read meromorphically, sign-free off the poles}\ [\mathrm{FT.26}(1)\ \text{pattern}];
\]

%dsp-level:0
\[
\qquad \text{the}\ \aleph\text{-domain at the junction couplings is}\ \mathrm{E2}\ [\text{A-BM.10},\ \text{issued in}\ \mathrm{BM.139}];
\]

%dsp-level:0
\[
(6)\quad\textbf{status:}
\]

%dsp-level:0
\[
\qquad \textbf{constructed}\ —\ \mathrm{BM.6}(6)\text{'s specification is discharged;}
\]

%dsp-level:0
\[
\qquad \text{the resolvent specification is discharged by}\ \mathrm{BM.134};
\]

%dsp-level:0
\[
\qquad \text{the two derived clauses}\ (2),
\]

%dsp-level:0
\[
(4)\ \text{carry no new freedom beyond}\ (1),
\]

%dsp-level:0
\[
(3)\text{'s declared identifications}.
\]

\textbf{Req.} D2.5, D2.9 (Parts 1–2); CT2.6 (Part 2); P4, P5 (Parts 1–2); FT.26(1), FT.79(5) (Part 3); Def-BM.10, Def-BM.17, BM.6(6), BM.10(1), BM.34(4), BM.34(7), BM.122(3), BM.123(1), EM.25(7) (this part or Part 2); A-BM.12 (§II.4, logged before use); the plane-wave basis on a chart [math].

\textbf{Pf.} Specification content: each clause states its requirement and cites its anchors; the non-derived clauses (1) and (3) are declared at the point of use and logged in §II.4; the constitutive clauses (2) and (4) are constructions on derived objects, each priced in its own clause; anchor existence machine-checked [\texttt{checks/\allowbreak{}verify\_\allowbreak{}H\_\allowbreak{}bm\_\allowbreak{}closing\_\allowbreak{}tier.py} §C: the document pins]. $\blacksquare$

\textbf{Phys.} The momentum-space resolvent the corpus named as an open interface is now a definition with exactly two declared ingredients — the boundary-channel embedding and the Fourier normalization — and everything else forced: the denominator is the derived mass-shell operator, so the pole locus is the dispersion relation, and the vertex normalization is pinned by its reduction to the register weight at unit coupling. The bars travel with it: no unitarity claim, no numerical cross-section, no measured comparison, and no Feynman prescription imported from outside.

\end{definition}

\begin{theorem}[{\texorpdfstring{BM.134 — The resolvent's derived content: the form uniqueness, the pole reads, the on-shell measure, and the two-body phase space}{BM.134 - The resolvent's derived content: the form uniqueness, the pole reads, the on-shell measure, and the two-body phase space}}]\label{tag:BM-134}
\textbf{Thm.}

%dsp-level:0
\[
(1)\quad\textbf{the pole locus is the derived dispersion:}
\]

%dsp-level:0
\[
\qquad R(z;\mathbf p)^{-1}=0\iff\bigl(p_0-N^ip_i\bigr)^2=N^2\bigl(m^2+\gamma^{ij}p_ip_j\bigr)\iff g^{\mu\nu}p_\mu p_\nu=-m^2
\]

%dsp-level:0
\[
\qquad [\mathrm{BM.34}(4);\ \text{machine: the master polynomial identity re-derived, cleared by}\ (\det\gamma)^2,\ \text{with three exact rational ADM witnesses}];
\]

%dsp-level:0
\[
\textbf{the branches:}
\]

%dsp-level:0
\[
\qquad p_0=N^ip_i\pm N\sqrt{m^2+p^2};
\]

%dsp-level:0
\[
\qquad E=-p_0=N\sqrt{m^2+p^2}-N^ip_i\ \ \text{on the positive branch}
\]

%dsp-level:0
\[
\qquad [\mathrm{BM.34}(4);\ \mathrm{BM.119}(1)];
\]

%dsp-level:0
\[
\textbf{the rest-pole read:}\qquad\mathbf p=0,\ w=0:\quad p_0=\pm Nm;\qquad E=Nm\qquad[\mathrm{BM.34}(7);\ \text{machine}];
\]

%dsp-level:0
\[
\textbf{the pole displacement:}
\]

%dsp-level:0
\[
\qquad \text{in the vacuum gauge}\ (N{=}1,\ w{=}0):\quad R=\frac1{p_0^2-m^2-\mathbf p^2}:
\]

%dsp-level:0
\[
\qquad \boxed{\ m^2\ \ \text{is exactly the pole displacement in}\ p_0^2\ \text{against the massless pole}\ }
\]

%dsp-level:0
\[
\qquad [\mathrm{D2.5}\ \text{vocabulary};\ \text{machine: the cleared identity}];
\]

%dsp-level:0
\[
(2)\quad\textbf{the form uniqueness:}
\]

%dsp-level:0
\[
\qquad \text{among rational forms}\ R'=\bigl(\alpha\,(p_0-N^ip_i)^2+\beta\,N^2(m^2+p^2)+\delta\bigr)^{-1}\ \text{over}\ \mathbb Q(\alpha,\beta,\delta):
\]

%dsp-level:0
\[
\text{the five conditions}
\]

%dsp-level:0
\[
\qquad (\mathrm i)\ \ \text{the pole locus is}\ \mathrm{BM.34}(4)\text{'s mass shell};
\]

%dsp-level:0
\[
\qquad (\mathrm{ii})\ \ \text{the vacuum-gauge reduction}\ R'\to\bigl(p_0^2-m^2-p^2\bigr)^{-1};
\]

%dsp-level:0
\[
(\mathrm{iii})\ \ \text{T-evenness under}\ (N^i,p_i)\mapsto(-N^i,-p_i)\ [\mathrm{BM.34}(4)];
\]

%dsp-level:0
\[
\qquad (\mathrm{iv})\ \ \text{hermitian sign-freeness off the poles}\ [\mathrm{FT.26}(1)\ \text{pattern}];
\]

%dsp-level:0
\[
\qquad (\mathrm v)\ \ \text{simple poles (minimality, no squared denominators)};
\]

%dsp-level:0
\[
\boxed{\ \text{force}\ \ (\alpha,\beta,\delta)=(1,-1,0)\ \ \text{exactly}\ }
\]

%dsp-level:0
\[
\qquad [\text{machine: the exact linear solve over}\ \mathbb Q\ \text{on rational ADM witnesses;}\ (\mathrm i)\implies\alpha=-\beta,\ \delta=0;\ (\mathrm{ii})\implies\alpha=1];
\]

%dsp-level:0
\[
\textbf{the declared content is exactly}\ \mathrm{Def\text{-}BM.20}(1),(3)\text{'s two identifications:}
\]

%dsp-level:0
\[
\qquad \text{every other row is forced};
\]

%dsp-level:0
\[
(3)\quad\textbf{the on-shell measure (the phase space from the resolvent and the chart measure):}
\]

%dsp-level:0
\[
\qquad \delta\bigl((p_0-N^ip_i)^2-N^2E_{\mathbf p}^2\bigr)=\frac{\delta\bigl(p_0-N^ip_i+NE_{\mathbf p}\bigr)+\delta\bigl(p_0-N^ip_i-NE_{\mathbf p}\bigr)}{2N\,E_{\mathbf p}},
\]

%dsp-level:0
\[
\qquad E_{\mathbf p}:=\sqrt{m^2+\gamma^{ij}p_ip_j}
\]

%dsp-level:0
\[
\qquad [\text{math: the delta composition rule};\ \text{machine: the Jacobian coefficient at five rational witnesses}];
\]

%dsp-level:0
\[
\boxed{\ d\mu_{\mathrm{shell}}\ =\ \frac{d^3p}{2\,N\sqrt{m^2+\gamma^{ij}p_ip_j}}\ =\ \frac{d^3p}{2E}\ }
\]

%dsp-level:0
\[
\qquad [\text{the positive branch};\ \mathrm{Def\text{-}BM.20}(3)\text{'s}\ (2\pi)^3\ \text{carried with it}];
\]

%dsp-level:0
\[
\textbf{the vacuum-gauge reduction:}
\]

%dsp-level:0
\[
\qquad N=1,\ \gamma=\delta:\quad d\mu_{\mathrm{shell}}=\frac{d^3p}{2\sqrt{m^2+p^2}}
\]

%dsp-level:0
\[
\qquad [\text{machine}];
\]

%dsp-level:0
\[
\qquad [d\mu_{\mathrm{shell}}]=\mathfrak q^{+1}
\]

%dsp-level:0
\[
\qquad [\mathrm{CT2.1};\ \text{machine: the exponent census}];
\]

%dsp-level:0
\[
(4)\quad\textbf{the two-body phase space and the fluxes (CM frame, zero shift):}
\]

%dsp-level:0
\[
\qquad z=E_1+E_2;
\]

%dsp-level:0
\[
\qquad E_i=N\sqrt{m_i^2+p^2};
\]

%dsp-level:0
\[
\qquad p_\ast=\frac1{2N}\sqrt{z^2-4N^2m^2}\ \ \text{at}\ m_1=m_2=m
\]

%dsp-level:0
\[
\qquad [\text{machine: the exact solve}];
\]

%dsp-level:0
\[
\boxed{\ \Phi_2(z;m,N)\ =\ \frac{p_\ast}{4\pi N^2\,z}\ =\ \frac{\sqrt{z^2-4N^2m^2}}{8\pi N^3\,z}\ }
\]

%dsp-level:0
\[
\qquad [\text{math: the delta reduction of}\ \int d\mu_1d\mu_2\ \text{with}\ (2\pi)^4\delta^4\ \text{and}\ \mathrm{Def\text{-}BM.20}(3)\text{'s convention};\ \text{machine: the coefficient algebra at five rational witnesses}];
\]

%dsp-level:0
\[
\textbf{the massless final states:}
\]

%dsp-level:0
\[
\qquad \Phi_2^{(\gamma\gamma)}=\frac1{8\pi N^3}
\]

%dsp-level:0
\[
\qquad [\text{machine}];
\]

%dsp-level:0
\[
\qquad \textbf{the fluxes:}
\]

%dsp-level:0
\[
\qquad F_{\gamma\gamma}=\frac{2z^2}{N^2}\ \ (\text{two massless beams}),
\]

%dsp-level:0
\[
\qquad F_{\mathrm{pair}}=\frac{2z\sqrt{z^2-4N^2m^2}}{N^2}\ \ (\text{two massive beams})
\]

%dsp-level:0
\[
[\text{machine:}\ p_1\!\cdot\!p_2=-\frac{z^2}{2N^2}\ \text{massless and}\ m^2-\frac{z^2}{2N^2}\ \text{massive at zero shift};\ F=4\sqrt{(p_1\!\cdot\!p_2)^2-m_1^2m_2^2}];
\]

%dsp-level:0
\[
\textbf{the formed domain:}
\]

%dsp-level:0
\[
\qquad \Phi_2\ \text{formed}\iff z\ge2Nm;
\]

%dsp-level:0
\[
\qquad \Phi_2=0\ \ \text{exactly at threshold};
\]

%dsp-level:0
\[
\qquad [\Phi_2]=\mathfrak q^0,
\]

%dsp-level:0
\[
\qquad [F]=\mathfrak q^{+1}
\]

%dsp-level:0
\[
\qquad [\text{machine: the exponent census}];
\]

%dsp-level:0
\[
(5)\quad\textbf{the two resolvents, two species:}
\]

%dsp-level:0
\[
\qquad \text{the chart resolvent}\ R(z;\mathbf p)\ [(1)]\ \text{and the register's zig-zag resolvent}\ \frac1{1-w^2x_+x_-}\ [\mathrm{BM.121}(2);\ \mathrm{BM.23}(7)]\ \textbf{are}\ \neg\text{identified}:
\]

%dsp-level:0
\[
\text{the first is meromorphic in the chart energy with poles at the dispersion branches};
\]

%dsp-level:0
\[
\qquad \text{the second is a formal even series in the ledger weight with its pole at}\ |w|\ge2\ [\mathrm{BM.121}(5)];
\]

%dsp-level:0
\[
\boxed{\ \text{they meet only through the threshold indicator}\ \chi_{\{NE\ge\mathcal T_0\}}\ }
\]

%dsp-level:0
\[
\qquad [\mathrm{Def\text{-}BM.17}(3);\ \mathrm{BM.122}(3);\ \text{the currency separation stands}];
\]

%dsp-level:0
\[
(6)\quad\textbf{the junction domain:}
\]

%dsp-level:0
\[
\qquad \text{the}\ \aleph\text{-reads at the sealed junctions where the boundary-channel states couple run in the}\ \mathrm{E2}\ \text{domain}\ [\text{A-BM.10},\ \text{issued in}\ \mathrm{BM.139}(2);\ \text{forward pointer}];
\]

%dsp-level:0
\[
(7)\quad\textbf{the closure row:}
\]

%dsp-level:0
\[
\qquad \boxed{\ \textbf{the momentum-space}\ \Sigma\textbf{-resolvent is constructed}\ }
\]

%dsp-level:0
\[
\qquad [\mathrm{Def\text{-}BM.20};\ \text{machine: the discharge census against}\ \mathrm{BM.6}(6)\text{'s three requirements — plane-wave boundary-channel states}\ (1),\ Z(\mathcal R,\Sigma;k^2)\text{-type two-point read}\ (2),\ m^2\ \text{as the pole displacement}\ (1)];
\]

%dsp-level:0
\[
\textbf{what it unblocks:}
\]

%dsp-level:0
\[
\qquad \text{the pair amplitude and the annihilation cross-section}\ [\mathrm{BM.142}];
\]

%dsp-level:0
\[
\qquad \textbf{what stays barred:}
\]

%dsp-level:0
\[
\qquad \neg\text{an S-matrix element beyond the formed amplitudes},\quad\neg\text{a measured cross-section}\ [\mathrm{P5}].
\]

\textbf{Req.} D2.5, D2.9 (Parts 1–2); CT2.1, CT2.6 (Part 2); P5 (Parts 1–2); FT.26(1), FT.79(5) (Part 3); Def-BM.10, Def-BM.17(3), Def-BM.20, BM.6(6), BM.10(1), BM.23(7), BM.34(1), BM.34(4), BM.34(7), BM.119(1), BM.121(2), BM.121(5), BM.122(3), BM.132(2) (this part); A-BM.10, A-BM.12 (§II.4, logged before use; forward pointer for A-BM.10's issuance row); the delta-composition rule, the exact linear solve on rational witnesses, the two-body phase-space integral in the declared convention, Laurent exponent censuses [math; machine: \texttt{checks/\allowbreak{}verify\_\allowbreak{}H\_\allowbreak{}bm\_\allowbreak{}closing\_\allowbreak{}tier.py} §C].

\textbf{Pf.}

%dsp-level:0
\[
(1):\ \mathrm{BM.34}(4)\text{'s master identity re-derived entrywise on three rational ADM witnesses}\ (\Omega^2,\varkappa,\beta,w_i,q_{ij}\ \text{rational,}\ q\ \text{off-diagonal});
\]

%dsp-level:0
\[
\qquad \text{the branch solve and the rest substitution}\ [\text{machine}];
\]

%dsp-level:0
\[
\qquad \text{the vacuum gauge:}\ N=1,\ w=0,\ \gamma=\delta\implies R=(p_0^2-m^2-p^2)^{-1}\ [\text{machine}]
\]

%dsp-level:0
\[
(2):\ \text{the conditions as linear constraints on}\ (\alpha,\beta,\delta):
\]

%dsp-level:0
\[
\quad (\mathrm i)\ \text{the zero set coincides with the mass shell}\implies\alpha(p_0-N^ip_i)^2+\beta N^2(m^2+p^2)+\delta\ \text{is proportional to}\ (p_0-N^ip_i)^2-N^2(m^2+p^2)\implies\delta
\]

%dsp-level:0
\[
= 0,
\]

%dsp-level:0
\[
\ \beta\qquad = -\alpha;
\]

%dsp-level:0
\[
\qquad (\mathrm{ii})\ \text{fixes}\ \alpha=1\ [\text{machine: the exact solve over}\ \mathbb Q];
\]

%dsp-level:0
\[
\qquad (\mathrm{iii}),(iv),(\mathrm v)\ \text{are then automatic}\ [\text{machine: the substitution rows}]
\]

%dsp-level:0
\[
(3):\ \delta(g(p_0))=\sum_{\text{roots}}\delta(p_0-r_i)/|g'(r_i)|,
\]

%dsp-level:0
\[
\qquad g'(p_0)=2(p_0-N^ip_i)=\pm2NE_{\mathbf p}\ \ \text{at the roots}\ [\text{math}];
\]

%dsp-level:0
\[
\qquad \text{machine: the coefficient at}\ (N,\mathbf p,m)\in\{(1,\cdot,1),(2,\cdot,1),(1,\cdot,2),(3,\cdot,2),(2,\cdot,3)\}\ \text{rational grids};
\]

%dsp-level:0
\[
\qquad \text{the exponent census}\ [\mathrm{CT2.1}:\ [d^3p]=\mathfrak q^{+3/2},\ [E]=\mathfrak q^{+1/2}]
\]

%dsp-level:0
\[
(4):\ \int\frac{d^3p}{(2\pi)^3}\frac{(2\pi)^4\delta(z-E_1-E_2)}{4E_1E_2}\bigg|_{\mathbf p_2=-\mathbf p_1}
\]

%dsp-level:0
\[
= \frac1{4(2\pi)^2}\int d\Omega\,
\]

%dsp-level:0
\[
\frac{p^2\,dp}{E_1E_2}\delta(z-E_1-E_2)
\]

%dsp-level:0
\[
= \frac{4\pi}{16\pi^2}\cdot\frac{p_\ast}{N^2z}\ [\text{machine: the Jacobian}\ \tfrac{d(E_1+E_2)}{dp}=\tfrac{2N^2p_\ast}z\ \text{at five rational witnesses}];
\]

%dsp-level:0
\[
\qquad \text{the fluxes from}\ p_1\!\cdot\!p_2\ \text{at zero shift}\ [\text{machine}]
\]

%dsp-level:0
\[
(5),
\]

%dsp-level:0
\[
(6),
\]

%dsp-level:0
\[
(7):\ \text{the species census, the domain row, and the discharge census against}\ \mathrm{BM.6}(6)\text{'s specification text}\ [\text{machine: the document pins}].\ \blacksquare
\]

\textbf{Phys.} The resolvent's declared surface is two naming choices; everything else is forced. Its poles are exactly the derived dispersion's two branches, its rest pole reads the composite's mass as the displacement against the massless pole, and among all rational forms of its degree the five derived conditions pin it uniquely. The phase space falls out of the same mass shell: the on-shell measure is the delta reduction of the chart measure, the two-body phase space and both fluxes are exact functions of $(z,m,N)$, formed precisely above threshold. The chart resolvent and the register's zig-zag resolvent stay separate species that meet only at the threshold indicator, and the item the corpus carried since the fermion era — the momentum-space $\Sigma$-resolvent — closes as constructed.

\end{theorem}

\begin{theorem}[{\texorpdfstring{BM.135 — The multi-chain additive charge census: the baryon and meson rows, the stretched states, and the gauged read}{BM.135 - The multi-chain additive charge census: the baryon and meson rows, the stretched states, and the gauged read}}]\label{tag:BM-135}
\textbf{Thm.}

%dsp-level:0
\[
(1)\quad\textbf{the read, re-certified:}
\]

%dsp-level:0
\[
\qquad Q(c_1,\dots,c_k)=\sum_{i=1}^{k}q_{\mathrm{core}}(c_i),
\]

%dsp-level:0
\[
\qquad q_{\mathrm{core}}(c)=\tfrac12\bigl(\sigma_{s(c)}-\sigma_{t(c)}\bigr)\in\{-1,0,+1\}
\]

%dsp-level:0
\[
\qquad [\mathrm{Def\text{-}BM.8}(1);\ \mathrm{FT.40};\ \mathrm{EM.9};\ \mathrm{D3.4}];
\]

%dsp-level:0
\[
\textbf{the three specification clauses stand discharged:}
\]

%dsp-level:0
\[
\qquad Q|_{\text{single chain}}=q_{\mathrm{core}};
\]

%dsp-level:0
\[
\qquad Q(\mathcal X\sqcup\mathcal Y)=Q(\mathcal X)+Q(\mathcal Y)\ \ \text{on disjoint sealings};
\]

%dsp-level:0
\[
\qquad Q\ \text{defined only on sealed supports}\ \bigl(\mathcal C=1\bigr)
\]

%dsp-level:0
\[
\qquad [\mathrm{Def\text{-}BM.8}(2)\text{--}(4);\ \mathrm{BM.6}(1);\ \mathrm{FT.19}];
\]

%dsp-level:0
\[
(2)\quad\textbf{the baryon charge (the}\ \varepsilon\textbf{-sealed 3-juncture):}
\]

%dsp-level:0
\[
\qquad Q_B=q_1+q_2+q_3\in\{-3,\dots,+3\},
\]

%dsp-level:0
\[
\qquad W_3(Q)=\binom6{Q+3}=(1,\ 6,\ 15,\ \mathbf{20},\ 15,\ 6,\ 1)
\]

%dsp-level:0
\[
\qquad [\mathrm{Def\text{-}BM.8}(5),(6);\ \mathrm{BM.28}(2);\ \mathrm{BM.114}(1);\ \text{machine: the exhaustive}\ 4^3{=}64\ \text{end-sign census re-run}];
\]

%dsp-level:0
\[
\textbf{the named-class rows:}
\]

%dsp-level:0
\[
\qquad \text{the neutral class:}\ Q_B=0,
\]

%dsp-level:0
\[
\qquad E_{\mathrm{int}}^{(n)}=6
\]

%dsp-level:0
\[
\qquad [\mathrm{BM.17};\ \mathrm{Def\text{-}BM.3}(2);\ \mathrm{Census\ FT.35}(2)];
\]

%dsp-level:0
\[
\qquad \text{the unit class:}\ Q_B=+1
\]

%dsp-level:0
\[
\qquad [\mathrm{BM.26};\ \mathrm{Def\text{-}BM.3}(3)];
\]

%dsp-level:0
\[
\qquad \text{the CPT image:}\ Q_B=-1
\]

%dsp-level:0
\[
\qquad [\mathrm{BM.27}];
\]

%dsp-level:0
\[
\textbf{the transport profile of the neutral class:}
\]

%dsp-level:0
\[
\qquad m_q=(2,1,1)
\]

%dsp-level:0
\[
\qquad [\mathrm{Def\text{-}BM.13}(1);\ \mathrm{FT.106}(3\mathrm a)];
\]

%dsp-level:0
\[
\qquad \frac{W(0)}{W(\pm1)}=\frac{20}{15}=\frac43
\]

%dsp-level:0
\[
\qquad [\mathrm{BM.28}(2);\ \text{machine}];
\]

%dsp-level:0
\[
(3)\quad\textbf{the meson charge (the}\ \delta\textbf{-sealed 2-juncture):}
\]

%dsp-level:0
\[
\qquad Q_M=q_1+q_2\in\{-2,\dots,+2\},
\]

%dsp-level:0
\[
\qquad W_2(Q)=\binom4{Q+2}=(1,\ 4,\ \mathbf6,\ 4,\ 1)
\]

%dsp-level:0
\[
\qquad [\mathrm{BM.110}(3);\ \mathrm{BM.114}(1);\ \text{machine: the}\ 4^2\ \text{census re-run}];
\]

%dsp-level:0
\[
\textbf{the CP-pair rows:}\qquad Q_M\bigl(c,c^{\mathrm{CP}}\bigr)=q+(-q)=0\ \ \text{for all three}\ q\qquad[\mathrm{Def\text{-}BM.15}(3);\ \text{machine: exhaustive}];
\]

%dsp-level:0
\[
\textbf{the CP-pairs saturate the neutral class:}
\]

%dsp-level:0
\[
\qquad 1+4+1=6=W_2(0),
\]

%dsp-level:0
\[
\qquad \#\{Q{=}0\ \text{tuples outside the CP-pair forms}\}=0
\]

%dsp-level:0
\[
\qquad [\mathrm{BM.110}(3);\ \text{machine re-run}];
\]

%dsp-level:0
\[
(4)\quad\textbf{the stretched and unit rows (both charge-read questions adjudicated):}
\]

%dsp-level:0
\[
\qquad \boxed{\ \#|Q|=2\ \ \text{states:}\ \ 12\ \ \text{at}\ k=3,\quad2\ \ \text{at}\ k=2;\qquad\#|Q|=3\ \ \text{states:}\ \ 2\ \ \text{at}\ k=3\ }
\]

%dsp-level:0
\[
\qquad [W_3(\pm2)=6,\ W_3(\pm3)=1,\ W_2(\pm2)=1;\ \text{machine: exhaustive}];
\]

%dsp-level:0
\[
\textbf{the multi-chain charge read is derived:}
\]

%dsp-level:0
\[
\qquad |Q|=2\ \ \text{and}\ \ |Q|=3\ \ \text{observable integer states exist on sealed supports of the additive read}
\]

%dsp-level:0
\[
\qquad [\text{the requirement is discharged};\ \mathrm{FT.41}(\mathrm a):\ N_{\mathrm{charge}}=2\ \text{end registers per chain re-cited}];
\]

%dsp-level:0
\[
\textbf{the ternary-read row stands re-certified:}
\]

%dsp-level:0
\[
\qquad \text{the read is integer-valued at every}\ k:\quad\tfrac13,\ \tfrac23\ \ \text{are structurally unreachable}
\]

%dsp-level:0
\[
\qquad [\mathrm{Def\text{-}BM.8}(5);\ \mathrm{BM.28}(7);\ \text{machine: the sum-closure scan over}\ \{-1,0,1\}^k,\ k=1..4];
\]

%dsp-level:0
\[
(5)\quad\textbf{the gauged read:}
\]

%dsp-level:0
\[
\qquad \partial_\mu F^{\mu\nu}=e\,J_Q^\nu,
\]

%dsp-level:0
\[
\qquad \text{the}\ \mathsf A\text{-coefficient}\ e\,Q,
\]

%dsp-level:0
\[
\qquad e=\frac{gg'}{\sqrt{g^2+g'^2}}
\]

%dsp-level:0
\[
\qquad [\mathrm{Def\text{-}BM.12}(1)\text{--}(3);\ \mathrm{EW.35};\ \mathrm{EM.15}];
\]

%dsp-level:0
\[
\textbf{the running entry:}
\]

%dsp-level:0
\[
\qquad e=e(r,E)\ \ \text{with all three factors derived}
\]

%dsp-level:0
\[
\qquad [\mathrm{BM.88}(4);\ \mathrm{BM.132}(2);\ \mathrm{BM.133}(6)];
\]

%dsp-level:0
\[
\mathrm{Def\text{-}BM.8}(8)\text{'s}\ \neg\text{constructed row (the coupling of}\ Q\ \text{to a gauge field)}\ \textbf{is discharged by}\ \mathrm{Def\text{-}BM.12}(2),
\]

%dsp-level:0
\[
(3)\ \text{and re-scoped,}\ \neg\text{retracted};
\]

%dsp-level:0
\[
(6)\quad\textbf{additivity at general}\ n\ \textbf{and the grammar row:}
\]

%dsp-level:0
\[
\qquad \text{a sealed support partitions into blocks of size}\ 2\ \text{and}\ 3
\]

%dsp-level:0
\[
\qquad [\mathrm{CT3.17a}(1)];
\]

%dsp-level:0
\[
\qquad Q(\Gamma)=\sum_{\text{blocks}}Q(\text{block})
\]

%dsp-level:0
\[
\qquad [\text{machine: exhaustive at}\ n\le4\ \text{on ternary alphabets}];
\]

%dsp-level:0
\[
\text{the}\ \delta/\varepsilon\text{-mixing from}\ N
\]

%dsp-level:0
\[
= 4\ [\mathrm{BM.31}(2);\ \mathrm{BM.33}]\ \text{governs which supports are sealed,}\ \neg\text{the additivity of the read on sealed blocks};
\]

%dsp-level:0
\[
(7)\quad\textbf{the closure row:}
\]

%dsp-level:0
\[
\qquad \boxed{\ \textbf{the additive multi-chain charge read exists, is constructed, and is gauged on the chart side}\ }
\]

%dsp-level:0
\[
\qquad [\mathrm{Def\text{-}BM.8};\ (1)\text{--}(6);\ \text{machine: the three-clause discharge census against}\ \mathrm{BM.6}(1)\text{'s specification}];
\]

%dsp-level:0
\[
\text{the ternary-read row is complete: no residue at any chain count};
\]

%dsp-level:0
\[
\qquad \neg\text{a species name attaches to any charge class beyond the identification layer}\ [\mathrm{Def\text{-}BM.3}(7)];
\]

%dsp-level:0
\[
\qquad \neg\text{a measured charge enters}\ [\mathrm{P5}].
\]

\textbf{Req.} D2.14, D3.4, D3.9 (Parts 1–2); CT3.17a, CT3.22, EM.9, EM.11, EM.15, EW.35 (Part 2); P2, P5 (Parts 1–2); FT.19, FT.40, FT.41, FT.106(3) (Part 3); Def-BM.3, Def-BM.8, Def-BM.12, Def-BM.13, Def-BM.15, BM.6(1), BM.17, BM.26, BM.27, BM.28(2), BM.28(7), BM.31(2), BM.33, BM.88(4), BM.110(3), BM.114(1), BM.132(2), BM.133(6) (this part); Census FT.35(2) (row-wise); exhaustive ternary enumeration, generating-function convolution, sum-closure scans [math; machine: \texttt{checks/\allowbreak{}verify\_\allowbreak{}H\_\allowbreak{}bm\_\allowbreak{}closing\_\allowbreak{}tier.py} §D].

\textbf{Pf.}

%dsp-level:0
\[
(1)\text{--}(3):\ \mathrm{Def\text{-}BM.8}\text{'s construction re-certified clause by clause and the two censuses re-run exhaustively}\ [\text{machine}];
\]

%dsp-level:0
\[
\qquad \text{the named-class rows consumed from their own tags row-wise}\ [\mathrm{Def\text{-}BM.3}(7)\text{'s discipline}]
\]

%dsp-level:0
\[
(4):\ \text{the counts from}\ W_3,W_2\ \text{directly}\ [\text{machine}];
\]

%dsp-level:0
\[
\qquad \text{the unreachability scan:}\ \textstyle\sum_iq_i\in\mathbb Z\ \forall q_i\in\{-1,0,1\}\ \text{at}\ k\le4\ [\text{machine}]
\]

%dsp-level:0
\[
(5):\ \mathrm{Def\text{-}BM.12}(1)\text{--}(3)\text{'s rows re-cited with the running coupling substituted}\ [\mathrm{BM.133}(6)];
\]

%dsp-level:0
\[
\qquad (6):\ \text{the block sum on ternary alphabets, exhaustive at}\ n\le4\ [\text{machine}];
\]

%dsp-level:0
\[
\qquad (7):\ \text{the discharge census}.\ \blacksquare
\]

\textbf{Phys.} The additive charge read the era opened as its first construction spec is re-certified, censused and gauged. Baryons carry seven integer charge classes with the frozen weight row and the neutral-to-unit ratio four thirds; mesons carry five, and their CP pairs saturate the neutral class exactly. The stretched states settle two ledger rows at once: observable charge two exists — twelve states at three chains, two at two — so the integer charge read is derived, while the fractional alphabet stays structurally unreachable, so the ternary read's consistent no-go stands. The chart-side gauge coupling turns the read into a current with a fully derived running coefficient, and the question the corpus carried since the fermion era — whether an additive multi-chain charge read exists — closes as constructed.

\end{theorem}

\begin{theorem}[{\texorpdfstring{BM.136 — The $S_3$ action is the colour 3-vertex: the equivariance theorem, the exchange structure at every $n$, and the statistics per class}{BM.136 - The action is the colour 3-vertex: the equivariance theorem, the exchange structure at every , and the statistics per class}}]\label{tag:BM-136}
\textbf{Thm.}

%dsp-level:0
\[
(1)\quad\textbf{the colour 3-vertex and its equivariance:}\qquad B=\varepsilon_{ijk}\,c_1^ic_2^jc_3^k\qquad[\mathrm{D3.9};\ \mathrm{Def\text{-}BM.1};\ \mathrm{CT3.15}];
\]

%dsp-level:0
\[
\boxed{\ \varepsilon_{\sigma(i)\sigma(j)\sigma(k)}=\operatorname{sgn}(\sigma)\,\varepsilon_{ijk}\ \ \forall\sigma\in S_3\ }
\]

%dsp-level:0
\[
\qquad [\text{machine: exhaustive over the}\ 6\times27\ \text{index pairs}];
\]

%dsp-level:0
\[
\dim\wedge^3(\mathbb C^3)=\binom33=1:
\]

%dsp-level:0
\[
\qquad \textbf{the}\ \varepsilon\textbf{-seal is the unique invariant line, and}\ S_3\ \textbf{acts on it by the sign representation}
\]

%dsp-level:0
\[
\qquad [\mathrm{BM.128}(5);\ \text{machine: the antisymmetrizer rank re-run}];
\]

%dsp-level:0
\[
\text{colour is}\ \neg\text{observable by construction}
\]

%dsp-level:0
\[
\qquad [\mathrm{CT3.22}];
\]

%dsp-level:0
\[
\qquad \text{the colour generators act as}\ \mathfrak{su}(3)\ \text{on the slots}
\]

%dsp-level:0
\[
\qquad [\mathrm{CT3.8};\ \mathrm{FT.82}(3)];
\]

%dsp-level:0
\[
(2)\quad\textbf{the intertwiner theorem (the}\ S_3\ \textbf{action is literally the colour 3-vertex):}
\]

%dsp-level:0
\[
\qquad \text{the sealed baryon is the contraction}\ \langle\varepsilon|\,
\]

%dsp-level:0
\[
:\ (\mathbb C^3)^{\otimes3}\otimes\mathcal H_3\to\mathbb C\ \ \text{with}\ \hat P_\sigma\ \text{the slot permutation}\ [\mathrm{BM.128}(2)]:
\]

%dsp-level:0
\[
\boxed{\ \bigl(\hat P_\sigma\otimes D_{\mathrm{col}}(\sigma)\bigr)\,|\varepsilon\rangle\otimes|\psi\rangle\ \ \text{is}\ \ S_3\textbf{-invariant exactly when the slot read compensates the colour sign:}\ \ \hat P_\sigma\big|_{\text{sealed}}=\operatorname{sgn}(\sigma)\ }
\]

%dsp-level:0
\[
[\text{machine: the explicit}\ 27\text{-dimensional colour tensor}\ |\varepsilon\rangle=\sum_{abc}\varepsilon_{abc}|abc\rangle\ \text{and all six}\ \sigma\in S_3;\ D_{\mathrm{col}}(\sigma)|\varepsilon\rangle=\operatorname{sgn}(\sigma)|\varepsilon\rangle\ \text{entrywise}];
\]

%dsp-level:0
\[
\textbf{the reading:}
\]

%dsp-level:0
\[
\qquad \text{the sealing grammar's antisymmetry and the slot space's permutation action are the two halves of one equivariant map}\ —\ \text{the}\ S_3\ \text{action on the 3-chain state space is the colour 3-vertex read on slots};
\]

%dsp-level:0
\[
(3)\quad\textbf{the exchange structure at every}\ n\ \textbf{(residue}\ (\mathrm{R7}')\ \textbf{closed):}
\]

%dsp-level:0
\[
\qquad \boxed{\ \mathcal H_n\ :=\ \bigotimes_{a=1}^{n}\Bigl(\mathcal H_\omega(c_a)\otimes\Psi(c_a)\Bigr),\qquad\dim\mathcal H_n=2^n\prod_{a=1}^{n}\bigl(|k_a|+1)\ }\qquad[\mathrm{BM.128}(1)\ \text{extended uniformly}];
\]

%dsp-level:0
\[
\hat P_\sigma\hat P_\tau=\hat P_{\sigma\tau}
\]

%dsp-level:0
\[
\qquad \forall\sigma,\tau\in S_n
\]

%dsp-level:0
\[
\qquad [\mathrm{BM.128}(2)\text{'s convention, uniform in}\ n;\ \text{machine: all}\ 576\ \text{ordered pairs at}\ n=4,\ d=2,\ \text{entrywise over}\ \mathbb Q];
\]

%dsp-level:0
\[
\boxed{\ \chi(\sigma)=\operatorname{Tr}\hat P_\sigma\big|_{(\mathbb C^d)^{\otimes n}}=d^{\,c(\sigma)}\ \ \text{for every}\ n\ }
\]

%dsp-level:0
\[
\qquad [\mathrm{BM.128}(3)\ \text{extended};\ \text{machine: all}\ 24\ \text{permutations at}\ (n,d)=(4,2)\ \text{and}\ (4,3):\ \text{the five cycle types give}\ (16,8,4,4,2)\ \text{at}\ d=2\ \text{and}\ (81,27,9,9,3)\ \text{at}\ d=3];
\]

%dsp-level:0
\[
\varepsilon(\Gamma)=\prod_{a=1}^{n}\varepsilon(c_a)
\]

%dsp-level:0
\[
\qquad \text{at an open connector}
\]

%dsp-level:0
\[
\qquad [\mathrm{BM.128}(7)\ \text{extended};\ \text{machine: all}\ 625\ \text{quadruples in}\ \{0,\dots,4\}^4:\ \ 312\ \text{fermionic},\ 313\ \text{bosonic}];
\]

%dsp-level:0
\[
\textbf{the wedge rows:}
\]

%dsp-level:0
\[
\qquad \dim\wedge^n(\mathbb C^2)=\binom2n=0\ \ \forall n\ge3;
\]

%dsp-level:0
\[
\qquad \dim\wedge^3(\mathbb C^3)=1;
\]

%dsp-level:0
\[
\qquad \dim\operatorname{Inv}\bigl((\mathbb C^3)^{\otimes4}\bigr)=0
\]

%dsp-level:0
\[
\qquad [\text{machine: the antisymmetrizer ranks}\ (2,4),(3,4),(2,3),(3,3)\ \text{exact};\ \text{the weight census: no zero-weight vector in}\ (\mathbb C^3)^{\otimes4}];
\]

%dsp-level:0
\[
\boxed{\ \textbf{four colour slots admit no singlet: the}\ N\ge4\ \textbf{sealings must mix}\ \delta\ \textbf{and}\ \varepsilon\ \textbf{blocks}\ }
\]

%dsp-level:0
\[
\qquad [\mathrm{BM.31}(2);\ \mathrm{BM.33};\ \mathrm{CT3.17a}(2):\ g(4)=10\ \text{groupings};\ \text{machine: the census of}\ \{2,3\}\text{-partitions at}\ N=4];
\]

%dsp-level:0
\[
\textbf{the grammar rescope row:}
\]

%dsp-level:0
\[
\qquad \text{the}\ \delta/\varepsilon\text{-mixing governs which}\ n\text{-chain configurations are sealed (observability}\ [\mathrm{BM.1}]),
\]

%dsp-level:0
\[
\ \neg\text{the exchange structure, which is uniform in}\ n;
\]

%dsp-level:0
\[
(4)\quad\textbf{the winding assignment per class:}
\]

%dsp-level:0
\[
\qquad \textbf{(a) the computed floor:}\ \ \text{the minimal materialized constituent is the digon with}\ k=1
\]

%dsp-level:0
\[
\qquad [\mathrm{FT.28};\ \mathrm{FT.32}(6);\ \mathrm{FT.107}(1);\ \mathrm{BM.129}(5)]:
\]

%dsp-level:0
\[
\boxed{\ (k_1,k_2,k_3)=(1,1,1):\qquad\varepsilon_B=(-1)^3=-1,\qquad j_B=\frac32\ }\qquad[\mathrm{BM.129}(4),(5);\ \text{machine}];
\]

%dsp-level:0
\[
\textbf{the neutral and unit classes share the floor:}
\]

%dsp-level:0
\[
\qquad E_{\mathrm{int}}
\]

%dsp-level:0
\[
= 6\ \text{on both}\ [\mathrm{BM.17};\ \mathrm{BM.26};\ \mathrm{BM.28}(1)]\implies\textbf{both named baryon classes are fermionic at the computed floor, derived}\ ;
\]

%dsp-level:0
\[
\mathrm{BM.129}(6),
\]

%dsp-level:0
\[
\ \mathrm{BM.130}(1)\text{'s ``value per named class awaits}\ (\mathrm{R3})\text{'' row is discharged at the floor locus:}\ \varepsilon_B\ \text{depends on the winding alone, and the winding at the floor is derived};
\]

%dsp-level:0
\[
\textbf{(b) the six-state classes under the declared corner convention:}
\]

%dsp-level:0
\[
\qquad k_{vf}=(3,3,2,3,5)\ \ \text{on}\ (\text{tetra},\ \text{cube},\ \text{octa},\ \text{dodeca},\ \text{icosa})
\]

%dsp-level:0
\[
\qquad [\text{A-BM.11 (variant B), issued in}\ \mathrm{BM.139}(4);\ \mathrm{FT.76}(1);\ \mathrm{Def\text{-}BM.14}(1),(2):\ \text{the floor map}\ 2E\in\{12,24,60\}\ \text{gives}\ k_{vf}=(3,3\,|\,2,2\,|\,3,5)\ \text{on the six states}]:
\]

%dsp-level:0
\[
\varepsilon_B(\mathbf s)=(-1)^{k_{vf}(s_1)+k_{vf}(s_2)+k_{vf}(s_3)}:
\]

%dsp-level:0
\[
\qquad \boxed{\ \#\{\text{fermionic}\}=112,\qquad\#\{\text{bosonic}\}=104\ }
\]

%dsp-level:0
\[
\qquad [\text{machine: the exhaustive}\ 6^3{=}216\ \text{census}];
\]

%dsp-level:0
\[
\textbf{the variant dependence is exact:}
\]

%dsp-level:0
\[
\qquad \text{variant A}\ \bigl(k_f=(3,2\,|\,3,3\,|\,5,3)\bigr):\quad140\ \ \text{fermionic},\ 76\ \ \text{bosonic}
\]

%dsp-level:0
\[
\qquad [\text{machine}]\implies\textbf{the census is convention-dependent and the declaration}\ \text{A-BM.11}\ \textbf{pins it};
\]

%dsp-level:0
\[
(5)\quad\textbf{the bars:}
\]

%dsp-level:0
\[
\qquad \neg\text{a flavour}\ \mathsf{SU}(3),\ \ \neg\text{an}\ \mathsf{SU}(6):
\]

%dsp-level:0
\[
\quad S_n\ \text{permutes constituent slots}\ [\mathrm{BM.117}(3);\ \mathrm{BM.128}(9);\ \mathrm{BM.129}(7)\ \text{re-cited}];
\]

%dsp-level:0
\[
\textbf{the flavour occupancy remains}\ (\mathrm{R3})\textbf{-species:}
\]

%dsp-level:0
\[
\qquad \text{which three of the six states a named class occupies is the generation-selection question}\ [\mathrm{BM.60}(1):\ \text{three lines derived, three selected}\ \neg\text{derived}]\ —\ \textbf{named for the Part-5 ledger,}\ \neg\text{a baryomesonic residue};
\]

%dsp-level:0
\[
\neg\text{a species name on any row}\ [\mathrm{Def\text{-}BM.3}(7)];\qquad\neg\text{a measured value}\ [\mathrm{P5}];
\]

%dsp-level:0
\[
(6)\quad\textbf{status:}
\]

%dsp-level:0
\[
\qquad \textbf{derived:}\ (1)\text{--}(4);
\]

%dsp-level:0
\[
\qquad \boxed{\ (\mathrm{R7}')\ \textbf{is closed:}\ \mathcal H_n,\ \hat P_\sigma,\ \chi(\sigma)=d^{\,c(\sigma)},\ \varepsilon(\Gamma)=\prod_a\varepsilon(c_a)\ \ \text{at every}\ n};
\]

%dsp-level:0
\[
\qquad \textbf{the}\ S_3\textbf{-action gap is closed: the action is the colour vertex's slot read};
\]

%dsp-level:0
\[
\textbf{the winding assignment per class is derived at both loci (floor and six-state);}
\]

%dsp-level:0
\[
\qquad \text{the statistics function is complete and its values are derived per class}\ [\mathrm{BM.126};\ \mathrm{BM.129}(4)].
\]

\textbf{Req.} A3 (Part 1); D2.14, D3.9 (Parts 1–2); CT3.8, CT3.15, CT3.17a, CT3.22 (Part 2); P5 (Parts 1–2); FT.28, FT.32(6), FT.76(1), FT.82(3), FT.107(1) (Part 3, Census); Def-BM.1, Def-BM.3(7), Def-BM.14, BM.1, BM.17, BM.26, BM.28(1), BM.31(2), BM.33, BM.60(1), BM.117(3), BM.126, BM.128, BM.129 (this part); A-BM.11 (§II.4, logged before use; forward pointer for its issuance row); permutation-module enumeration, antisymmetrizer ranks by exact Fraction elimination, the weight-zero census of $(\mathbb C^3)^{\otimes4}$, exhaustive parity censuses [math; machine: \texttt{checks/\allowbreak{}verify\_\allowbreak{}H\_\allowbreak{}bm\_\allowbreak{}closing\_\allowbreak{}tier.py} §E].

\textbf{Pf.}

%dsp-level:0
\[
(1):\ \text{the}\ 6\times27\ \text{equivariance census and the rank-1 antisymmetrizer}\ [\text{machine}];
\]

%dsp-level:0
\[
\qquad (2):\ D_{\mathrm{col}}(\sigma)|\varepsilon\rangle=\operatorname{sgn}(\sigma)|\varepsilon\rangle\ \text{entrywise on the explicit tensor};
\]

%dsp-level:0
\[
\qquad \text{the invariance of}\ |\varepsilon\rangle\otimes|\psi\rangle\ \text{under the diagonal action with}\ \hat P_\sigma\ \text{on slots}\ [\text{machine}]
\]

%dsp-level:0
\[
(3):\ \hat P_\sigma\ \text{built as the factor permutation on}\ (\mathbb C^d)^{\otimes n}:\ \ \text{the product law on all}\ 576\ \text{pairs at}\ n
\]

%dsp-level:0
\[
= 4\ \text{in}\ \mathrm{BM.128}(2)\text{'s composition convention}\ [\text{machine: the pushforward factor action}];
\]

%dsp-level:0
\[
\qquad \chi(\sigma)=d^{c(\sigma)}\ \text{on all}\ 24\ \text{permutations at}\ d\in\{2,3\}\ [\text{machine}];
\]

%dsp-level:0
\[
\qquad \text{the antisymmetrizer ranks}\ (0,0,0,1)\ \text{at}\ ((2,4),(3,4),(2,3),(3,3))\ [\text{machine}];
\]

%dsp-level:0
\[
\qquad \dim\operatorname{Inv}((\mathbb C^3)^{\otimes4})
\]

%dsp-level:0
\[
= 0:\ \text{the weights of}\ \mathbb C^3\ \text{are}\ \{(1,0),(0,1),(-1,-1)\}:\ \ a(1,0)+b(0,1)+c(-1,-1)
\]

%dsp-level:0
\[
= (0,0)\ \wedge\ a+b+c
\]

%dsp-level:0
\[
= 4
\]

%dsp-level:0
\[
\implies a
\]

%dsp-level:0
\[
= b
\]

%dsp-level:0
\[
= c
\]

%dsp-level:0
\[
= 4/3\notin\mathbb Z\ [\text{machine: the exhaustive composition census}];
\]

%dsp-level:0
\[
\qquad \text{the parity census on}\ 625\ \text{quadruples}\ [\text{machine}]
\]

%dsp-level:0
\[
(4):\ \text{the floor row from}\ \mathrm{BM.129}(5)\text{'s digon line and}\ \varepsilon_B=(-1)^{k_1+k_2+k_3}\ [\mathrm{BM.129}(4)];
\]

%dsp-level:0
\[
\qquad \text{the six-state censuses:}\ 112/104\ \text{and}\ 140/76\ \text{by exhaustive enumeration of the parity patterns}\ (4\ \text{odd},\ 2\ \text{even})\ \text{and}\ (5\ \text{odd},\ 1\ \text{even})\ [\text{machine}];
\]

%dsp-level:0
\[
\qquad (5),(6):\ \text{assembly}.\ \blacksquare
\]

\textbf{Phys.} The permutation action on the three-chain state space is not a bookkeeping device: it is the colour 3-vertex read on slots. The epsilon tensor is the unique invariant line of three colour slots, it transforms by the sign of the permutation, and the sealed baryon is exactly the intertwiner that pairs the two actions — so the antisymmetry lives in colour, unobservable by construction, while the observable statistics lives in the winding. The exchange structure extends uniformly to every chain count with the character formula intact, and the four-slot colour census explains the grammar: no four-fold singlet exists, which is why seals mix delta and epsilon blocks from four attachments on. The winding assignment per class is derived at both loci — the digon floor gives both named baryon classes the fermion sign and spin three halves, and the six-state census under the declared vertex-figure convention gives 112 fermionic against 104 bosonic classes, a census that provably depends on the corner convention and is pinned by its declaration. What remains open is not baryomesonic: which states a named class occupies is the generation-selection question, and it leaves for the next era.

\end{theorem}

\begin{theorem}[{\texorpdfstring{BM.137 — The two-chain ensemble measure, the meson's formed binding, the sector-mass pincer, and the two-chain spectrum}{BM.137 - The two-chain ensemble measure, the meson's formed binding, the sector-mass pincer, and the two-chain spectrum}}]\label{tag:BM-137}
\textbf{Thm.}

%dsp-level:0
\[
(1)\quad\textbf{the general pair-specialization law:}
\]

%dsp-level:0
\[
\qquad \text{an}\ n\text{-pair floor crystallizes each pair independently with the stationary weight}\ p=1-4\varepsilon_q
\]

%dsp-level:0
\[
\qquad [\mathrm{BM.19}(2);\ \mathrm{BM.94}(4)\text{'s marked specialization, extended uniformly}]:
\]

%dsp-level:0
\[
P(E{=}2j)=\binom njp^ju^{n-j},
\]

%dsp-level:0
\[
\qquad u:=1-p=4\varepsilon_q;
\]

%dsp-level:0
\[
\qquad \text{the materialization restriction}\ E\ge2\ \ \text{is derived}
\]

%dsp-level:0
\[
\qquad [\mathrm{BM.94}(5);\ \mathrm{FT.144}(4);\ \mathrm{Def\text{-}H.4}];
\]

%dsp-level:0
\[
\boxed{\ D_n\ :=\ \frac{1-u^n}{1-u}\ =\ 1+u+\cdots+u^{n-1},\qquad w_{\mathrm{mat}}^{(n)}=\frac1{D_n},\qquad\langle E\rangle^{(n)}_{\mathrm{mat}}=\frac{2n}{D_n}\in(2,\ 2n]\ }
\]

%dsp-level:0
\[
\qquad [\text{machine: the polynomial identity}\ (1-u)D_n=1-u^n\ \text{for}\ n=1..8;\ \text{the binomial moment algebra exact in}\ \mathbb Q[p]];
\]

%dsp-level:0
\[
\operatorname{Var}^{(n)}_{\mathrm{mat}}(E)=\frac{4n\bigl[(1+(n-1)p)D_n-n\bigr]}{D_n^2};
\]

%dsp-level:0
\[
\qquad n=3:\quad\frac{12\,p(1-p)(3-2p)}{D^2}
\]

%dsp-level:0
\[
\qquad [\mathrm{BM.94}(6)\ \textbf{reproduced exactly};\ \text{machine}];
\]

%dsp-level:0
\[
n=1:\quad w_{\mathrm{mat}}^{(1)}=1:
\]

%dsp-level:0
\[
\qquad \textbf{a single materialized pair has no flicker: the restriction}\ E
\]

%dsp-level:0
\[
\ge 2\ \textbf{forces its crystallized member}
\]

%dsp-level:0
\[
\qquad [\text{machine}];
\]

%dsp-level:0
\[
(2)\quad\textbf{the two-chain ensemble measure (constructed):}
\]

%dsp-level:0
\[
\qquad \text{the meson floor}\ E^{cryst}=2E_{\mathrm{int}}(c)=4\ \ \text{is two pairs}
\]

%dsp-level:0
\[
\qquad [\mathrm{BM.111}(1);\ \text{the digon floor}\ E_{\mathrm{int}}(c)=2\ [\mathrm{FT.28};\ \mathrm{BM.13}(3)]]:
\]

%dsp-level:0
\[
\boxed{\ P_{\mathrm{mat}}(E{=}2)=\frac{2u}{1+u},\qquad P_{\mathrm{mat}}(E{=}4)=\frac{p}{1+u},\qquad\langle E\rangle^{(2)}_{\mathrm{mat}}=\frac4{1+u}\in(2,4],\qquad w_{\mathrm{mat}}^{(2)}=\frac1{1+u}\in\Bigl(\tfrac12,\ 1\Bigr]\ }
\]

%dsp-level:0
\[
\operatorname{Var}^{(2)}_{\mathrm{mat}}(E)=\frac{8p(1-p)}{(2-p)^2};
\]

%dsp-level:0
\[
\qquad \frac{\operatorname{Var}^{(2)}}{\langle E\rangle^{(2)2}}=\frac{p(1-p)}2\le\frac18,
\]

%dsp-level:0
\[
\qquad \frac{\sigma^{(2)}}{\langle E\rangle^{(2)}}\le\frac1{2\sqrt2}<\frac12
\]

%dsp-level:0
\[
\qquad [\text{machine: exact rational functions};\ \tfrac18-\tfrac{p(1-p)}2=\tfrac{(2p-1)^2}{8}\ge0];
\]

%dsp-level:0
\[
\textbf{the two-bit floor survives the averaging:}
\]

%dsp-level:0
\[
\qquad E\ge2\ \ \text{on the whole materialized support}\implies\langle m_{\mathrm{hid}}\rangle^{(2)}_{\mathrm{mat}}>\frac{2\ln2}{\mathfrak h}
\]

%dsp-level:0
\[
\qquad [\mathrm{BM.70}(2);\ \mathrm{BM.94}(8)\ \text{pattern};\ \text{machine}];
\]

%dsp-level:0
\[
\textbf{the no-flicker collapse:}
\]

%dsp-level:0
\[
\qquad \varepsilon_q=0\implies u=0\implies P_{\mathrm{mat}}(E{=}4)=1,
\]

%dsp-level:0
\[
\qquad w_{\mathrm{mat}}^{(2)}=1
\]

%dsp-level:0
\[
\qquad [\mathrm{BM.19}(3);\ \mathrm{Def\text{-}BM.5}(5);\ \text{machine}];
\]

%dsp-level:0
\[
(3)\quad\textbf{the meson's ensemble binding (formed, nonzero — the}\ (\mathrm{R8})\ \textbf{row):}
\]

%dsp-level:0
\[
\qquad \mathcal B^{(n)}_{\mathrm{ens}}\ :=\ \bigl\langle m_{\mathrm{hid}}\bigr\rangle^{(n)}_{\mathrm{mat}}-\sum_{a=1}^n\bigl\langle m_{\mathrm{hid}}\bigr\rangle^{(1)}_{\mathrm{mat}}\ =\ \Bigl(\frac{2n}{D_n}-2n\Bigr)\frac{\ln2}{\mathfrak h}
\]

%dsp-level:0
\[
\qquad [\mathrm{FT.44}(4);\ (1)];
\]

%dsp-level:0
\[
\boxed{\ \mathcal B^{(n)}_{\mathrm{ens}}\ =\ -\,\frac{2n\,u\,D_{n-1}}{D_n}\cdot\frac{\ln2}{\mathfrak h}\ }
\]

%dsp-level:0
\[
\qquad [D_n-1=uD_{n-1};\ \text{machine: the polynomial identity for}\ n=2,3,4];
\]

%dsp-level:0
\[
n=2:
\]

%dsp-level:0
\[
\qquad \boxed{\ \mathcal B^{(2)}_{\mathrm{ens}}=-\frac{4u}{1+u}\cdot\frac{\ln2}{\mathfrak h}=-\frac{16\,\varepsilon_q\ln2}{(1+4\varepsilon_q)\,\mathfrak h}\ },
\]

%dsp-level:0
\[
\qquad \mathcal B^{(2)}_{\mathrm{ens}}<0\ \ \forall\varepsilon_q>0,
\]

%dsp-level:0
\[
\qquad \mathcal B^{(2)}_{\mathrm{ens}}=0\iff\varepsilon_q=0,
\]

%dsp-level:0
\[
\qquad -\frac{2\ln2}{\mathfrak h}<\mathcal B^{(2)}_{\mathrm{ens}}\le0
\]

%dsp-level:0
\[
\qquad [\text{machine: the sign and range census on}\ u\in[0,1)];
\]

%dsp-level:0
\[
\textbf{the binding is the flicker's, not the count's, at}\ n=2\ \textbf{as at}\ n=3
\]

%dsp-level:0
\[
\qquad [\mathrm{BM.21}(3)\ \text{pattern}];
\]

%dsp-level:0
\[
\qquad \text{the two vanishing routes stand untouched:}\ \mathcal B^{\text{meson}}_{\mathrm{info}}=\mathcal B^{\text{meson}}_{\mathrm{pole}}=0
\]

%dsp-level:0
\[
\qquad [\mathrm{BM.111}(1),(2)\ \text{re-cited,}\ \neg\text{retracted}];
\]

%dsp-level:0
\[
(4)\quad\textbf{the sector-mass pincer at budget}\ B=2\ \textbf{(the}\ 2\textbf{-chain sector masses, computed):}
\]

%dsp-level:0
\[
\qquad \mathcal G_2\ \text{carries one chain-bearing skeleton, the digon}\ K_2
\]

%dsp-level:0
\[
\qquad [\mathrm{Def\text{-}H.32}(1);\ \mathrm{FT.105}:\ \text{the connected simple graphs on}\ \le2\ \text{nodes}]:
\]

%dsp-level:0
\[
\textbf{the pincer:}
\]

%dsp-level:0
\[
\qquad \text{a closed chain requires}\ |R|\ge2
\]

%dsp-level:0
\[
\qquad [\mathrm{Def\text{-}H.4}:\ \delta h(c)\ge2]
\]

%dsp-level:0
\[
\qquad \text{and the budget gives}\ |R|\le2E(K_2)=2
\]

%dsp-level:0
\[
\qquad [\mathrm{Def\text{-}H.32}(1)]:
\]

%dsp-level:0
\[
\implies|R|=2E\implies n_{\mathcal V}=2E-|R|=0\implies\Delta S=\tfrac32n_{\mathcal V}=0\qquad[\mathrm{Def\text{-}H.32}(6);\ \mathrm{FT.107}(2)];
\]

%dsp-level:0
\[
\boxed{\ m_2^2(q,S)=0\ \ \text{on every formed sector, and}\ Z_2^{\mathrm{norm}}(q,S)\ne0\ \ \text{only at}\ S=1\ \ \text{with}\ Z_2^{\mathrm{norm}}=(4,\ 2,\ 2)\ \ \text{at}\ q=(0,+1,-1)\ }
\]

%dsp-level:0
\[
\qquad [\text{machine: the exhaustive enumeration of}\ \mathcal G_2\ \text{—}\ 4\ \text{patterns on}\ K_2,\ \text{the}\ \sigma\text{-resolved instance census, the closed form}\ m_q2^{|R|-2}\ \text{of}\ \mathrm{FT.106}(3)\ \text{cross-checked}];
\]

%dsp-level:0
\[
\textbf{the structural reason:}
\]

%dsp-level:0
\[
\qquad \boxed{\ \textbf{at budget two, chain-bearing content is fully crystallized: the sector ensemble has no flicker room}\ };
\]

%dsp-level:0
\[
\qquad \text{the sector-species binding is}\ 0-0-0=0\ \ \text{exactly};
\]

%dsp-level:0
\[
\textbf{the two ensemble reads are separated species:}
\]

%dsp-level:0
\[
\qquad \text{the sector/defect read gives}\ 0;\ \text{the informational read gives}\ (3)\text{'s strict negative};
\]

%dsp-level:0
\[
\qquad \neg\text{a derived equation relates them}\ —\ \text{that is}\ (\mathrm{R1})\text{'s bridge}
\]

%dsp-level:0
\[
\qquad [\mathrm{BM.37}(2);\ \mathrm{BM.98};\ \mathrm{A\text{-}BM.9}];
\]

%dsp-level:0
\[
(5)\quad\textbf{the two-chain spectrum (the resonance rows at}\ n=2\textbf{):}
\]

%dsp-level:0
\[
\qquad W_2(Q)=(1,4,6,4,1),
\]

%dsp-level:0
\[
\qquad \sum_QW_2=16=4^2
\]

%dsp-level:0
\[
\qquad [\mathrm{BM.114}(1);\ \text{machine re-run}];
\]

%dsp-level:0
\[
\lambda^{(2)}(s_1,s_2)=D_S(s_1)+D_S(s_2),
\]

%dsp-level:0
\[
\qquad D_S\in\bigl\{2,\ 4,\ 2-\sqrt2,\ 1,\ 2-\mathfrak u,\ 2-\sqrt3\bigr\},
\]

%dsp-level:0
\[
\qquad \mathfrak u=2\cos\tfrac\pi{10}
\]

%dsp-level:0
\[
\qquad [\mathrm{Def\text{-}BM.14}(3);\ \mathrm{BM.111}(2);\ \mathrm{BM.114}(3)]:
\]

%dsp-level:0
\[
\boxed{\ \text{the}\ 21\ \text{multiset orbits give}\ 21\ \text{distinct eigenvalues with}\ \neg\text{a collision}\ }
\]

%dsp-level:0
\[
\qquad [\text{machine: the exact equality census in}\ \mathbb M;\ \text{contrast: one collision at}\ n=3\ [\mathrm{BM.114}(5)]];
\]

%dsp-level:0
\[
\lambda^{(2)}_{\min}=4-2\mathfrak u=2\bigl(2-2\cos\tfrac\pi{10}\bigr),
\]

%dsp-level:0
\[
\qquad \lambda^{(2)}_{\max}=8=4n
\]

%dsp-level:0
\[
\qquad [\mathrm{BM.117}(4)\ \text{re-certified at}\ n=2;\ \text{machine: the exact extrema in}\ \mathbb M];
\]

%dsp-level:0
\[
m^2_{\mathrm{pole}}(\mathcal R)=\hat\mu'^2\lambda^{(2)}\in\bigl[(4-2\mathfrak u)\,\hat\mu'^2,\ 8\hat\mu'^2\bigr];
\]

%dsp-level:0
\[
\qquad m^{meson}_{\mathrm{hid}}=\frac{2E_{\mathrm{int}}(c)\ln2}{\mathfrak h};
\]

%dsp-level:0
\[
\qquad \bigl\langle m_{\mathrm{hid}}\bigr\rangle^{(2)}_{\mathrm{mat}}=\frac{4\ln2}{(1+u)\,\mathfrak h}\ \ \text{at the digon floor}
\]

%dsp-level:0
\[
\qquad [\mathrm{Def\text{-}BM.17}(2);\ (2);\ \text{machine}];
\]

%dsp-level:0
\[
\varepsilon_M=+1\ \ \text{unconditionally},\qquad j_M=k\in\mathbb Z\qquad[\mathrm{BM.129}(1),(2)\ \text{re-cited}];
\]

%dsp-level:0
\[
(6)\quad\textbf{the octet adjudication stands, and the spectrum's shape is the derived one:}
\]

%dsp-level:0
\[
\qquad 8\ \ \text{occurs in no grouping family at any formed level}
\]

%dsp-level:0
\[
\qquad [\mathrm{BM.117}(3);\ \text{machine: the scan re-run over}\ W_n,\ g(N),\ \binom{6+n-1}n,\ \binom6n,\ n\le3,\ N\le7];
\]

%dsp-level:0
\[
\textbf{the meson spectrum is the eigenvalue ladder}\ (5)\ \textbf{with its charge census, its statistics and its formed binding}\ —\ \neg\text{an octet,}\ \neg\text{a nonet},
\]

%dsp-level:0
\[
\ \neg\text{a flavour multiplet};
\]

%dsp-level:0
\[
\textbf{the naming row:}
\]

%dsp-level:0
\[
\qquad \text{species names for meson classes remain confined to the identification layer}\ [\mathrm{Def\text{-}BM.3}(7);\ \mathrm{Def\text{-}BM.6}(7)];
\]

%dsp-level:0
\[
\qquad \neg\text{a measured mass, width or mixing enters}\ [\mathrm{P5}];
\]

%dsp-level:0
\[
(7)\quad\textbf{status:}
\]

%dsp-level:0
\[
\qquad \textbf{derived:}\ (1)\text{--}(6);
\]

%dsp-level:0
\[
\qquad \boxed{\ (\mathrm{R8})\ \textbf{is closed: the}\ 2\textbf{-chain sector masses are computed (identically zero, with the pincer), the two-chain ensemble measure is constructed, and the meson's ensemble binding is formed and strictly negative at every nonzero threshold}\ }
\]

%dsp-level:0
\[
[\mathrm{BM.111}(3)\text{'s}\ \neg\text{formed row discharged}];
\]

%dsp-level:0
\[
\qquad \neg\text{a declaration},
\]

%dsp-level:0
\[
\qquad \neg\text{a P4 addition},
\]

%dsp-level:0
\[
\qquad \neg\text{a new dimensionful scale}\ [\mathrm{P1}],
\]

%dsp-level:0
\[
\qquad \neg\text{a measured value}\ [\mathrm{P5}].
\]

\textbf{Req.} A4 (Part 1); D0.4, D0.8 (Part 1); CT0.8, CT1.29 (Part 1); CT2.8 (Part 2); P1, P5 (Parts 1–2); Def-H.4, Def-H.32, FT.28, FT.32(6), FT.44(4), FT.105, FT.106(3), FT.107(1), FT.107(2), FT.144(4) (Part 3); Def-BM.5, Def-BM.14(3), Def-BM.17(2), BM.13(3), BM.19(2), BM.19(3), BM.21(3), BM.37(2), BM.70(2), BM.94(4), BM.94(5), BM.94(6), BM.94(8), BM.98, BM.111, BM.114(1), BM.114(3), BM.114(5), BM.117(3), BM.117(4), BM.129(1), BM.129(2) (this part); Census FT.35(2) (row-wise); the A-BM.9 declaration; binomial moment algebra over $\mathbb Q[p]$, the degree-16 field $\mathbb M$ with its exact equality census, exhaustive configuration enumeration at budget two, polynomial identities [math; machine: \texttt{checks/\allowbreak{}verify\_\allowbreak{}H\_\allowbreak{}bm\_\allowbreak{}closing\_\allowbreak{}tier.py} §F].

\textbf{Pf.}

%dsp-level:0
\[
(1):\ P(E{=}2j)\ \text{is the}\ n\text{-pair binomial}\ [\mathrm{BM.94}(4)\text{'s specialization at}\ n];
\]

%dsp-level:0
\[
\qquad \text{conditioning on}\ j\ge1:\ \ 1-u^n=pD_n;
\]

%dsp-level:0
\[
\qquad \langle E\rangle^{(n)}_{\mathrm{mat}}=2\cdot\frac{np}{1-u^n}=\frac{2n}{D_n}\ [\text{machine: exact in}\ \mathbb Q[p]];
\]

%dsp-level:0
\[
\langle E^2\rangle^{(n)}_{\mathrm{mat}}=4\cdot\frac{np(1+(n-1)p)}{1-u^n}=\frac{4n(1+(n-1)p)}{pD_n}\ [\text{machine}];
\]

%dsp-level:0
\[
\qquad \operatorname{Var}=\langle E^2\rangle-\langle E\rangle^2\ [\text{machine}];
\]

%dsp-level:0
\[
\qquad n=3\ \text{row against}\ \mathrm{BM.94}(6)\text{'s frozen expression}\ [\text{machine: identical rational functions}];
\]

%dsp-level:0
\[
\qquad n=1:\ D_1=1,\ P(E{=}2)=1\ [\text{machine}]
\]

%dsp-level:0
\[
(2):\ n=2\ \text{substitution:}\ D_2=1+u=2-p;
\]

%dsp-level:0
\[
\qquad P_{\mathrm{mat}}(2)=\frac{2pu}{pD_2}=\frac{2u}{1+u},\quad P_{\mathrm{mat}}(4)=\frac{p^2}{pD_2}=\frac p{1+u}\ [\text{machine}];
\]

%dsp-level:0
\[
\qquad \text{the variance and the relative-width SOS certificate}\ [\text{machine}]
\]

%dsp-level:0
\[
(3):\ \mathcal B^{(n)}
\]

%dsp-level:0
\[
= \bigl(\frac{2n}{D_n}-2n\bigr)\frac{\ln2}{\mathfrak h}
\]

%dsp-level:0
\[
= -2n\frac{D_n-1}{D_n}\frac{\ln2}{\mathfrak h}
\]

%dsp-level:0
\[
= -\frac{2nuD_{n-1}}{D_n}\frac{\ln2}{\mathfrak h}\ [D_n-1=uD_{n-1};\ \text{machine}];
\]

%dsp-level:0
\[
\qquad \text{the sign, zero and range rows on}\ u\in[0,1)\ [\text{machine: rational grid and the factor census}]
\]

%dsp-level:0
\[
(4):\ \mathcal G_2=\{\bullet,\ K_2\}\ [\mathrm{FT.105}\text{'s census at}\ V\le2];
\]

%dsp-level:0
\[
\qquad \text{the}\ 4\ \text{patterns of}\ K_2\ \text{enumerated:}\ |R|\in\{0,1,1,2\};
\]

%dsp-level:0
\[
\qquad \text{closed chains exist only at}\ |R|=2\ (\text{the digon},\ k=1,\ S=1,\ 2\ \text{bases})\ [\text{machine}];
\]

%dsp-level:0
\[
\qquad \Delta S=\tfrac32(2-2)=0;
\]

%dsp-level:0
\[
\qquad Z_2^{\mathrm{norm}}(q,1)
\]

%dsp-level:0
\[
= 2\ \text{instances}\times m_q2^{0}
\]

%dsp-level:0
\[
= (4,2,2)\ [\mathrm{FT.106}(3)\text{'s closed form};\ \text{machine: the}\ \sigma\text{-resolved census agrees}];
\]

%dsp-level:0
\[
\qquad S\ne1:\ \neg\text{instances}\implies Z=0
\]

%dsp-level:0
\[
(5):\ \text{the}\ 21\ \text{sums computed in}\ \mathbb M\ \text{with the basis-equality census}\ [\text{machine}];
\]

%dsp-level:0
\[
\qquad \text{the extrema by the exact order of}\ \mathbb M\text{'s positive cone on the six}\ D_S\ [\text{machine: dyadic interval embedding}];
\]

%dsp-level:0
\[
\qquad (6):\ \text{the grouping scan re-run}\ [\text{machine}];
\]

%dsp-level:0
\[
\qquad (7):\ \text{assembly}.\ \blacksquare
\]

\textbf{Phys.} The meson gets its ensemble. The flicker measure generalizes from the marked three-pair specialization to a law at every $n$: the materialized weight is one over the geometric sum $1+u+\cdots+u^{n-1}$, the mean count is $2n$ over that sum, and at two chains the measure is exact — a mean between two and four, a relative width bounded by one over two root two, and the two-bit floor intact. The binding that residue (R8) named is formed and nonzero: strictly negative at every nonzero threshold, zero exactly at the sharp point, the same flicker-not-count structure the baryon showed. The sector route tells the complementary story: at budget two a chain-bearing configuration has no virtual room at all, so the two-chain sector masses vanish identically — a pincer, not a coincidence — and the two ensemble reads stay separated currencies pending the bridge slot. The spectrum completes the row: five charge classes with weights one-four-six-four-one, twenty-one pole eigenvalues with no collision at all, bounded by four-$n$, bosonic unconditionally. The octet remains unformed and the names remain where the identification layer keeps them.

\end{theorem}

\begin{theorem}[{\texorpdfstring{BM.138 — The seed-tensoring mechanism: the symmetrizer as a finite-group twirl, the transport and composition closure, and the two residues adjudicated}{BM.138 - The seed-tensoring mechanism: the symmetrizer as a finite-group twirl, the transport and composition closure, and the two residues adjudicated}}]\label{tag:BM-138}
\textbf{Thm.}

%dsp-level:0
\[
(1)\quad\textbf{the mechanism, derived:}
\]

%dsp-level:0
\[
\qquad \mathcal H_{\mathsf v}\cong\mathbb C^2\ \ \text{the transported seed}
\]

%dsp-level:0
\[
\qquad [\mathrm{EW.32};\ \mathrm{A2}:\ \text{the channel capacity is one qubit}];
\]

%dsp-level:0
\[
\boxed{\ \Psi(c)\ =\ \mathrm{Sym}^{|k(c)|}\bigl(\mathcal H_{\mathsf v}\bigr)\ =\ \Pi_{|k|}\Bigl(\mathcal H_{\mathsf v}^{\otimes|k|}\Bigr),\qquad\Pi_k:=\frac1{k!}\sum_{\sigma\in S_k}\hat P_\sigma\ }
\]

%dsp-level:0
\[
\qquad [\mathrm{Def\text{-}H.38};\ \mathrm{BM.128}(2)\text{'s permutation operators on identical factors}];
\]

%dsp-level:0
\[
\Pi_k^2=\Pi_k,
\]

%dsp-level:0
\[
\qquad \operatorname{rk}\Pi_k=k+1=\dim\mathrm{Sym}^k(\mathbb C^2),
\]

%dsp-level:0
\[
\qquad \Pi_k\hat P_\tau=\hat P_\tau\Pi_k=\Pi_k\ \ \forall\tau\in S_k
\]

%dsp-level:0
\[
\qquad [\text{machine: idempotence, exact rank over}\ \mathbb Q,\ \text{and the commutation for}\ k=1..6];
\]

%dsp-level:0
\[
\boxed{\ \textbf{the symmetrizer is the Haar twirl of the finite group}\ S_k\ }
\]

%dsp-level:0
\[
\qquad [\mathrm{BM.131}(6);\ \text{math: the averaging over a finite group is its Haar measure}];
\]

%dsp-level:0
\[
\textbf{the seed-tensoring read:}
\]

%dsp-level:0
\[
\qquad \text{the carrier is the}\ k\text{-fold tensor power of the seed, projected onto the invariants of the group average}\ —\ \textbf{tensoring with itself, then twirling};
\]

%dsp-level:0
\[
(2)\quad\textbf{the diagonal transport:}\qquad\mathcal U_c^{\otimes k}\big|_{\mathrm{Sym}^k}\qquad[\mathrm{Def\text{-}H.38};\ \mathrm{EW.33}:\ \text{the wrapping transport}];
\]

%dsp-level:0
\[
\bigl[\mathcal U^{\otimes k},\ \Pi_k\bigr]=0
\]

%dsp-level:0
\[
\qquad [\text{machine: the infinitesimal rows}\ \bigl[\textstyle\sum_a W_a^{(i)},\Pi_k\bigr]=0\ \text{on the}\ \mathfrak{su}(2)\ \text{generator triple, exact over}\ \mathbb Q(i),\ k\le6];
\]

%dsp-level:0
\[
\textbf{the weight ladder from the twirl:}
\]

%dsp-level:0
\[
\qquad \operatorname{Tr}\Bigl(\Pi_k\,D(z)^{\otimes k}\Bigr)=\sum_{m\in\{k,k-2,\dots,-k\}}z^{m},
\]

%dsp-level:0
\[
\qquad D(z)=\mathrm{diag}(z,z^{-1})
\]

%dsp-level:0
\[
\qquad [\mathrm{FT.60}(4)\ \text{re-derived};\ \text{machine: the Laurent identity at five rational}\ z\ \text{for}\ k\le8];
\]

%dsp-level:0
\[
\text{the Casimir row}\ \tfrac k2\bigl(\tfrac k2+1\bigr)\qquad[\mathrm{FT.60}(3)\ \text{re-cited}];\qquad j=\frac{|k|}2\qquad[\mathrm{Def\text{-}H.38};\ \mathrm{FT.55}(2)];
\]

%dsp-level:0
\[
(3)\quad\textbf{the composition closure (the junctions tensor by the same mechanism):}
\]

%dsp-level:0
\[
\qquad \mathrm{Sym}^{k_1}\otimes\mathrm{Sym}^{k_2}\supseteq\mathrm{Sym}^{k_1+k_2}\ \ \text{as the maximal-weight line on the cone}
\]

%dsp-level:0
\[
\qquad [\mathrm{FT.60}(6);\ \text{machine: the top weight}\ z^{k_1+k_2}\ \text{occurs with multiplicity}\ 1\ \text{in the character product, Laurent-exact}];
\]

%dsp-level:0
\[
\Psi(\Gamma)=\mathrm{Sym}^{k_1+k_2+d_{12}}\subseteq\mathrm{Sym}^{k_1}\otimes\mathrm{Sym}^{k_2}\otimes\mathrm{Sym}^{d_{12}}\qquad[\mathrm{BM.2}(2);\ \mathrm{FT.34}(3):\ k(\Gamma)=k_1+k_2+d_{12}];
\]

%dsp-level:0
\[
\boxed{\ \textbf{the connector is a carrier too:}\ \ \mathrm{Sym}^{d_{12}}(\mathcal H_{\mathsf v})\ \ \text{is built by the same seed-tensoring, so every junction composition closes inside the mechanism}\ }
\]

%dsp-level:0
\[
(4)\quad\textbf{the defect-ladder coherence:}
\]

%dsp-level:0
\[
\qquad \text{the weights}\ m\in\{k,k{-}2,\dots,-k\}\ \text{of}\ (2)\ \text{are the ladder}\ \Delta(m)
\]

%dsp-level:0
\[
= 4\sin^2\frac{m\Theta_R}4\ \text{reads on the tensorized carrier's own weight basis}
\]

%dsp-level:0
\[
\qquad [\mathrm{Def\text{-}BM.19}(3);\ \mathrm{BM.124}(3);\ \mathrm{Def\text{-}H.37}(2)]:
\]

%dsp-level:0
\[
\text{the mechanism is what makes the carrier and the defect read the same multiplet}\ —\ \neg\text{a new coupling,}\ \neg\text{a proportionality}\ \Theta_R
\]

%dsp-level:0
\[
= 2\pi k\ [\mathrm{BM.124}(5)\text{'s dimensional bar re-cited}];
\]

%dsp-level:0
\[
(5)\quad\textbf{residue (i) adjudicated — the coupling datum:}
\]

%dsp-level:0
\[
\qquad \operatorname{copies}(c)=k(c)\ \ \text{is}\ \mathrm{Def\text{-}H.38}\text{'s constitutive clause};
\]

%dsp-level:0
\[
\textbf{the derivation route is blocked by a frozen minimality theorem:}
\]

%dsp-level:0
\[
\qquad \text{the uniform-minimal triangle's winding}\ 3\ \text{has no three materialized winding-}1\ \text{subloops, and admissible decompositions are basis-relative}
\]

%dsp-level:0
\[
\qquad [\mathrm{FT.60}(7\mathrm a)\text{--}(7\mathrm b)];
\]

%dsp-level:0
\[
\boxed{\ \textbf{the mechanism is derived and the coupling integer stays a priced interface:}\ \ \text{what remained was the construction half, and it is complete}\ };
\]

%dsp-level:0
\[
(6)\quad\textbf{residue (ii) adjudicated — the}\ k\ge3\ \textbf{carrier inventory:}
\]

%dsp-level:0
\[
\qquad \mathrm{Sym}^{k\ge3}\ (\dim\ge4)\ \text{has no registered aspect block}\ [\mathrm{CT0.12}:\ \text{blocks}\ (3,3,5,3,1,1)];
\]

%dsp-level:0
\[
\textbf{the adjudication:}
\]

%dsp-level:0
\[
\qquad \text{the carrier is a chain-layer object,}\ \neg\text{a force aspect:}\ \ \text{every clause of this part consumes}\ \Psi\ \text{through}\ \bigl(\dim,\ \text{weights},\ \varepsilon,\ j,\ \text{Casimir},\ \chi\bigr)\ \text{alone}
\]

%dsp-level:0
\[
\qquad [\text{machine: the inventory scan of}\ \mathrm{Def\text{-}H.38}\text{'s consumption sites in Part 4}];
\]

%dsp-level:0
\[
\neg\text{an aspect claim at}\ k\ge3\ \ \text{is made or needed}\ [\mathrm{Def\text{-}H.38}\ \text{Scope re-cited}];
\]

%dsp-level:0
\[
\qquad \textbf{the inventory-extension question is void for this era};
\]

%dsp-level:0
\[
(7)\quad\textbf{status:}
\]

%dsp-level:0
\[
\qquad \textbf{derived:}\ (1)\text{--}(4);
\]

%dsp-level:0
\[
\qquad \textbf{adjudicated:}\ (5),(6);
\]

%dsp-level:0
\[
\qquad \boxed{\ \textbf{the seed-tensoring mechanism is derived}\ };
\]

%dsp-level:0
\[
\qquad \neg\text{a declaration},
\]

%dsp-level:0
\[
\qquad \neg\text{a P4 addition},
\]

%dsp-level:0
\[
\qquad \neg\text{a measured value}\ [\mathrm{P5}].
\]

\textbf{Req.} A2 (Part 1); D1.4 (Part 1); CT0.12 (Part 1); P5 (Parts 1–2); EW.32, EW.33 (Part 2); Def-H.37(2), Def-H.38, FT.34(3), FT.55(2), FT.60(3), FT.60(4), FT.60(6), FT.60(7) (Part 3); Def-BM.19(3), BM.2(2), BM.124(3), BM.124(5), BM.128(2), BM.131(6) (this part); exact rank by Fraction elimination, Laurent-polynomial character identities, the $\mathfrak{su}(2)$ infinitesimal commutation over $\mathbb Q(i)$, finite-group averaging [math; machine: \texttt{checks/\allowbreak{}verify\_\allowbreak{}H\_\allowbreak{}bm\_\allowbreak{}closing\_\allowbreak{}tier.py} §G].

\textbf{Pf.}

%dsp-level:0
\[
(1):\ \Pi_k^2
\]

%dsp-level:0
\[
= \Pi_k\ \text{from}\ \hat P_\sigma\hat P_\tau
\]

%dsp-level:0
\[
= \hat P_{\sigma\tau}\ \text{and the group-average rearrangement}\ [\mathrm{BM.128}(2)\ \text{transported to identical factors};\ \text{machine}];
\]

%dsp-level:0
\[
\qquad \operatorname{rk}\Pi_k:\ \text{exact elimination on the}\ 2^k\text{-dimensional symmetrizer for}\ k
\]

%dsp-level:0
\[
= 1..6:\ \text{ranks}\ (2,3,4,5,6,7)\ [\text{machine}]
\]

%dsp-level:0
\[
(2):\ \bigl[\hat P_\sigma,\mathcal U^{\otimes k}\bigr]
\]

%dsp-level:0
\[
= 0\ \text{for the diagonal action}\ [\text{math: the permutation and the diagonal tensor power commute}];
\]

%dsp-level:0
\[
\qquad \text{the infinitesimal rows on the generator triple}\ W_a
\]

%dsp-level:0
\[
= \tfrac12\sigma_a:\ \ \bigl[\textstyle\sum_i I^{\otimes(i-1)}\otimes W_a\otimes I^{\otimes(k-i)},\ \Pi_k\bigr]
\]

%dsp-level:0
\[
= 0\ [\text{machine over}\ \mathbb Q(i)];
\]

%dsp-level:0
\[
\qquad \text{the character:}\ \operatorname{Tr}(\Pi_kD^{\otimes k})
\]

%dsp-level:0
\[
= \frac1{k!}\sum_\sigma\prod_{\text{cycles}}(z^{L}+z^{-L})
\]

%dsp-level:0
\[
= \sum_mz^m\ [\text{machine: five rational}\ z\ \text{values,}\ k\le8]
\]

%dsp-level:0
\[
(3):\ \text{the character product}\ \bigl(\sum_iz^{k_1-2i}\bigr)\bigl(\sum_jz^{k_2-2j}\bigr)\ \text{carries}\ z^{k_1+k_2}\ \text{with coefficient}\ 1\ [\text{machine: Laurent multiplication}];
\]

%dsp-level:0
\[
\qquad \text{the connector row:}\ \mathrm{Sym}^{d_{12}}\ \text{is}\ (1)\text{'s construction at}\ k=d_{12}\ [\mathrm{BM.2}(2)]
\]

%dsp-level:0
\[
(4):\ \text{the ladder reads on the weight basis}\ \{\psi_m\}\ \text{of}\ (2)\ [\mathrm{Def\text{-}BM.19}(3);\ \mathrm{BM.124}(3)\ \text{re-cited}];
\]

%dsp-level:0
\[
\qquad (5):\ \mathrm{FT.60}(7)\text{'s two obstruction rows re-cited}\ —\ \text{the irreducibility is frozen};
\]

%dsp-level:0
\[
\qquad (6):\ \text{the consumption-site scan}\ [\text{machine: every}\ \Psi\text{-use in Part 4 is through the six listed reads}];
\]

%dsp-level:0
\[
\qquad (7):\ \text{assembly}.\ \blacksquare
\]

\textbf{Phys.} The carrier's construction is no longer a black box with a declared integer on top: the mechanism underneath the integer is derived. The symmetric power is the image of the group average over the $k$-fold tensor power of the seed — the same twirl structure the Haar–Poisson duality runs on $\mathrm{SU}(2)$, here on the finite permutation group — with rank $k+1$ exact, diagonal transport commuting with the average, and the weight ladder falling out of a Laurent identity. Junctions compose by the same machine: the connector itself carries a symmetric-power carrier, so every glued chain stays inside the mechanism. What remains declared is exactly one integer per chain, and a frozen minimality theorem says no construction below it exists; what remains questioned — an aspect block for carriers of dimension four and up — is adjudicated void for this era, because no clause consumes a carrier except through its dimension, weights, sign, spin, Casimir and character.

\end{theorem}

\begin{theorem}[{\texorpdfstring{BM.139 — The two convention declarations and the sector locus}{BM.139 - The two convention declarations and the sector locus}}]\label{tag:BM-139}
\textbf{Thm.}

%dsp-level:0
\[
(1)\quad\textbf{the}\ \aleph\textbf{-domain conflict (exactly as computed):}
\]

%dsp-level:0
\[
\qquad \text{three domains for the vertex-amplitude exponent}\ \aleph\ \text{entering}\ w=\mathcal J\prod_{e\in R}G(e)(-1)^{\aleph}
\]

%dsp-level:0
\[
\qquad [\mathrm{Def\text{-}H.27};\ \mathrm{A8};\ \mathrm{D3.3}]:
\]

%dsp-level:0
\[
\text{literal}\ (\mathrm{S2},\ \text{frozen});\qquad\mathrm{E1}\ (\text{chain-dominant});\qquad\mathrm{E2}\ (\text{literal-forced})\qquad[\mathrm{Def\text{-}H.27}(\mathrm a)\text{--}(\mathrm c)];
\]

%dsp-level:0
\[
\textbf{pairwise inequivalent, with the exact witnesses:}
\]

%dsp-level:0
\[
\qquad \text{the locked single-vertex star}\ (n_{\mathrm{in}}{+}n_{\mathrm{out}}=n\in\{4,5,6\},\ \text{balanced}):
\]

%dsp-level:0
\[
\qquad \text{literal}\ 0;
\]

%dsp-level:0
\[
\qquad \mathrm{E2}=\mathrm{E1}=\sum_{a\ge2}\binom{n_{\mathrm{in}}}a\binom{n_{\mathrm{out}}}aa!\,4^a=(32,\ 96,\ 672)
\]

%dsp-level:0
\[
\qquad [\mathrm{FT.75}(1);\ \text{machine: the binomial sums re-derived exactly}];
\]

%dsp-level:0
\[
\text{the locked digon hub with}\ a\ \text{digons}:
\]

%dsp-level:0
\[
\qquad \text{literal}\ 4^a;
\]

%dsp-level:0
\[
\qquad \mathrm{E2}=0;
\]

%dsp-level:0
\[
\qquad \mathrm{E1}=(!a)\,4^a\ \ \text{with}\ (!a)=(1,2,9,44)\ \ \text{at}\ a=2..5
\]

%dsp-level:0
\[
\qquad [\mathrm{FT.75}(1);\ \text{machine: the derangement census}];
\]

%dsp-level:0
\[
\boxed{\ \text{no further computation settles the choice from the frozen corpus: it is a declaration species}\ }\qquad;
\]

%dsp-level:0
\[
(2)\quad\textbf{A-BM.10 issued (the}\ \aleph\textbf{-domain):}
\]

%dsp-level:0
\[
\qquad \textbf{the baryomesonic era runs:}
\]

%dsp-level:0
\[
\qquad (\mathrm a)\ \ \text{the topology-blind ensembles on}\ \mathrm{Def\text{-}H.32}\text{'s inherited literal domain}\ (\mathrm{S2})\ —\ \text{unchanged, with every frozen ensemble row standing in its own domain}\ [\mathrm{Def\text{-}H.32}\ \text{Scopes};\ \mathrm{FT.74}\ \text{status row}];
\]

%dsp-level:0
\[
(\mathrm b)\ \ \text{the sealed-junction amplitude reads}\ —\ \text{where the chain-local pass-through usage of the telescoping and EM rows is the forced datum}\ [\mathrm{FT.21}(\mathrm i);\ \mathrm{FT.62}(5);\ \mathrm{BM.7}\ \text{species}]\ —\ \text{on the}\ \mathrm{E2}\ \text{domain}\ [\mathrm{Def\text{-}H.27}(\mathrm c)];
\]

%dsp-level:0
\[
(\mathrm c)\ \ \ \mathrm{E1}\ \textbf{recorded as a computed variant,}\ \neg\textbf{adopted};
\]

%dsp-level:0
\[
\textbf{the grounds:}
\]

%dsp-level:0
\[
\qquad \mathrm{E2}\ \text{retains the literal}\ \mathrm{D3.3}\ \text{read at every type-}(1,1)\ \text{vertex, so the frozen literal-domain records carry over unmodified there};
\]

%dsp-level:0
\[
\qquad \text{the chain choice at type-}(1,1)\ \text{is}\ \aleph\text{-gauge}\ [\mathrm{Def\text{-}H.27}(\mathrm c):\ \text{the datum is type-determined}]\ \text{and}\ \mathrm{E2}\ \text{fixes that gauge by vertex type, while}\ \mathrm{E1}\ \text{leaves it to the decomposition};
\]

%dsp-level:0
\[
\textbf{the derived consequences:}
\]

%dsp-level:0
\[
\qquad \text{the locked star is revived with the closed forms}\ (32,96,672);
\]

%dsp-level:0
\[
\qquad \text{the locked digon hub cancels in}\ \mathrm{E2}\ —\ \text{the bound-content cancellation at hubs is preserved, sharpening}\ \mathrm{FT.67}(3\mathrm a);
\]

%dsp-level:0
\[
\qquad \kappa\text{-invariance and the}\ (-1)^{\sum\kappa}\ \text{scaling carry over verbatim}\ [\mathrm{Def\text{-}H.27}\ \text{Rmk}];
\]

%dsp-level:0
\[
(3)\quad\textbf{the corner-convention conflict:}
\]

%dsp-level:0
\[
\qquad \text{the four}\ \Theta_R\text{-readings are computed:}
\]

%dsp-level:0
\[
\qquad \text{the vertex figure}\ \delta_v\ (\text{the only non-trivial spectrum});
\]

%dsp-level:0
\[
\qquad \text{the flat face}\ 0\ \text{and the global sweep}\ V\delta_v=4\pi\ (\text{both identically massless});
\]

%dsp-level:0
\[
\qquad \text{the winding pairing inside the first:}\ \text{variant A}\ (\delta_v,k_f),\ \text{variant B}\ (\delta_v,k_{vf})
\]

%dsp-level:0
\[
\qquad [\mathrm{FT.76}(1);\ \mathrm{Census\ FT.76}(1)\text{--}(2);\ \text{machine:}\ 2-2T_m(1)=0\ \forall m\ \text{in}\ \mathbb Z[x]\ \text{for the two massless readings}];
\]

%dsp-level:0
\[
k_f=(3,2,3,5,3),
\]

%dsp-level:0
\[
\qquad k_{vf}=(3,3,2,3,5)
\]

%dsp-level:0
\[
\qquad [\mathrm{FT.76}(1);\ \text{machine: both re-derived from the minimality rule}\ k=L\kappa_{\min}/2\ \text{on the face and vertex-figure loop lengths}\ L_f=(3,4,3,5,3),\ L_{vf}=(3,3,4,3,5)];
\]

%dsp-level:0
\[
\textbf{the variants are materially different:}
\]

%dsp-level:0
\[
\qquad \text{the multiplet tables differ}\ [\mathrm{Census\ FT.76}(2)];
\]

%dsp-level:0
\[
\qquad \textbf{the per-class statistics censuses differ:}\ \ 112/104\ \ \text{in B against}\ 140/76\ \ \text{in A}
\]

%dsp-level:0
\[
\qquad [\mathrm{BM.136}(4);\ \text{machine}];
\]

%dsp-level:0
\[
(4)\quad\textbf{A-BM.11 issued (the defect corner convention):}
\]

%dsp-level:0
\[
\qquad \boxed{\ \textbf{variant B — the vertex-figure winding}\ k_{vf}\ \text{with}\ \Theta_R=\delta_v\ }
\]

%dsp-level:0
\[
\qquad [\mathrm{BM.6}(4)\text{'s specification discharged}];
\]

%dsp-level:0
\[
\textbf{the grounds:}
\]

%dsp-level:0
\[
\qquad (\mathrm i)\ \ \text{loop consistency: the}\ \Theta_R
\]

%dsp-level:0
\[
= \delta_v\ \text{read is the vertex-figure loop's own defect, and pairing it with the face winding would mix two loci}\ [\mathrm{FT.76}(1):\ \text{variant B is the loop-consistent reading}];
\]

%dsp-level:0
\[
(\mathrm{ii})\ \ \text{the whole defect ladder of this part runs at the vertex-figure locus}\ [\mathrm{Def\text{-}BM.14}(3);\ \mathrm{Def\text{-}BM.19}(3);\ \mathrm{BM.79}(1);\ \mathrm{BM.124}(3);\ \mathrm{BM.125}];
\]

%dsp-level:0
\[
\qquad (\mathrm{iii})\ \ \text{the ceiling}\ \sqrt{4/D_S}\ \text{is winding-independent and untouched}\ [\mathrm{FT.76}(3)];
\]

%dsp-level:0
\[
\textbf{the effects:}
\]

%dsp-level:0
\[
\qquad \mathrm{FT.76},
\]

%dsp-level:0
\[
\ \mathrm{FT.77}\text{'s conditional clauses and this part's ladders convert from convention-tagged to declared data at their points of use},
\]

%dsp-level:0
\[
\ \neg\text{retracted};
\]

%dsp-level:0
\[
\qquad \text{variant A's tables remain computed data, recorded and}\ \neg\text{adopted};
\]

%dsp-level:0
\[
\qquad \text{the winding assignment per class is pinned}\ [\mathrm{BM.136}(4)];
\]

%dsp-level:0
\[
(5)\quad\textbf{A-BM.13 issued (the statistics locus):}
\]

%dsp-level:0
\[
\qquad \boxed{\ \textbf{the}\ \vee^2\textbf{-sector is the standing locus of the sign read}\ }
\]

%dsp-level:0
\[
\qquad [\mathrm{Def\text{-}BM.16}(5)\text{'s non-derived clause converted to declared data}];
\]

%dsp-level:0
\[
\textbf{the grounds (the four priced rows):}
\]

%dsp-level:0
\[
\qquad \vee^2\ \text{and}\ \wedge^2\ \text{give one statistics up to a global sign};
\]

%dsp-level:0
\[
\ \text{the determinant line and the orientation block do not}\ [\mathrm{BM.113}(6);\ \text{machine re-run to}\ k=8];
\]

%dsp-level:0
\[
\vee^2\ \text{is the flip's}\ +1\text{-eigensector on the CP pair}\ [\mathrm{A3}]\ \text{and matches the carrier's own symmetry}\ —\ \wedge^3(\mathbb C^2)
\]

%dsp-level:0
\[
= 0\ \text{leaves the carrier no antisymmetric locus at}\ n\qquad = 3
\]

%dsp-level:0
\[
\qquad [\mathrm{BM.128}(5)];
\]

%dsp-level:0
\[
\textbf{A-BM.12 (the resolvent interface) is issued at}\ \mathrm{Def\text{-}BM.20}\ \text{and recorded here in the ledger census};
\]

%dsp-level:0
\[
(6)\quad\textbf{the convention question closed — the census:}
\]

%dsp-level:0
\[
\qquad (\mathrm d)\ \ \text{the}\ \aleph\text{-domain: declared}\ [(2)];
\]

%dsp-level:0
\[
\qquad (\mathrm e)\ \ \text{the corner convention: declared}\ [(4)];
\]

%dsp-level:0
\[
(\mathrm b)\ \ \text{the exchange block}\ \mathcal K^{(\mathrm{exch})}:\ \textbf{adjudicated void for this era}\ —\ \text{the chirality read has its locus in the carrier rays}\ [\mathrm{A\text{-}BM.7};\ \mathrm{BM.73}],
\]

%dsp-level:0
\[
\ \text{and no}\ k
\]

%dsp-level:0
\[
\ge 3\ \text{aspect block is consumed}\ [\mathrm{BM.138}(6)];
\]

%dsp-level:0
\[
\qquad \neg\text{a declared exchange block is introduced};
\]

%dsp-level:0
\[
(\mathrm a),
\]

%dsp-level:0
\[
(\mathrm c)\ \ \text{the classification program's remaining cells}\ (\text{multi-solid chains, the single-vertex ladder, separated-contact hosting maps})\ \textbf{are named for the Part-5 cell program}\ —\ \neg\text{era objectives};
\]

%dsp-level:0
\[
\boxed{\ \textbf{both conventions are declared, the void residue is adjudicated, and the cell remainder is transferred}\ }
\]

%dsp-level:0
\[
(7)\quad\textbf{the declaration ledger:}
\]

%dsp-level:0
\[
\qquad \textbf{four addenda issued by this tier:}\ \ \text{A-BM.10}\ (\aleph\text{-domain}),
\]

%dsp-level:0
\[
\ \ \text{A-BM.11}\ (\text{corner convention}),
\]

%dsp-level:0
\[
\ \ \text{A-BM.12}\ (\text{resolvent interface}),
\]

%dsp-level:0
\[
\ \ \text{A-BM.13}\ (\text{statistics locus});
\]

%dsp-level:0
\[
\textbf{all four logged in §II.4 before use};
\]

%dsp-level:0
\[
\qquad \neg\text{a P4 member added by any of them};
\]

%dsp-level:0
\[
\qquad \neg\text{a dimensionful scale}\ [\mathrm{P1}];
\]

%dsp-level:0
\[
\qquad \neg\text{a measured value}\ [\mathrm{P5}];
\]

%dsp-level:0
\[
\qquad \text{§Cal-BM gains no parameter}\ [\mathrm{BM.143}(4)\text{'s consolidation is a register,}\ \neg\text{an addition}].
\]

\textbf{Req.} A8 (Part 1); D3.3 (Part 1); CT0.9 (Part 1); P1, P5 (Parts 1–2); Def-H.27, Def-H.32, FT.21, FT.62(5), FT.67(1), FT.67(3), FT.74, FT.75, FT.76(1), FT.76(3), FT.77 (Part 3); A3 (Part 1); Def-BM.14(3), Def-BM.16(5), Def-BM.19(3), Def-BM.20, BM.6(4), BM.7, BM.73, BM.79(1), BM.113(6), BM.124(3), BM.125, BM.128(5), BM.136(4), BM.138(6), A-BM.7 (this part); Census FT.76(1), FT.76(2) (row-wise); A-BM.10, A-BM.11, A-BM.13 (§II.4, logged before use); exact binomial and derangement enumeration, the Chebyshev identity $T_m(1)=1$ over $\mathbb Z[x]$, exhaustive parity censuses [math; machine: \texttt{checks/\allowbreak{}verify\_\allowbreak{}H\_\allowbreak{}bm\_\allowbreak{}closing\_\allowbreak{}tier.py} §H].

\textbf{Pf.}

%dsp-level:0
\[
(1):\ \text{the star sums}\ \sum_{a\ge2}\binom{n_{\mathrm{in}}}a\binom{n_{\mathrm{out}}}aa!4^a\ \text{evaluated exactly at}\ (2,2),
\]

%dsp-level:0
\[
(2,3),
\]

%dsp-level:0
\[
(3,3):\ (32,96,672)\ [\text{machine}];
\]

%dsp-level:0
\[
\qquad \text{the derangements}\ (!a)
\]

%dsp-level:0
\[
= \sum_{j\le a}(-1)^ja!/j!
\]

%dsp-level:0
\[
= (1,2,9,44)\ \text{and the hub values}\ (16,128,2304,45056)\ [\text{machine}]
\]

%dsp-level:0
\[
(3):\ k_f,
\]

%dsp-level:0
\[
k_{vf}\ \text{from}\ k
\]

%dsp-level:0
\[
= L\kappa_{\min}/2\ \text{with}\ \kappa_{\min}
\]

%dsp-level:0
\[
= 2\ (L\ \text{odd}),
\]

%dsp-level:0
\[
\ 1\ (L\ \text{even})\ [\mathrm{FT.32}(6);\ \mathrm{FT.45}]:\ \ L_f
\]

%dsp-level:0
\[
= (3,4,3,5,3)
\]

%dsp-level:0
\[
\mapsto (3,2,3,5,3);
\]

%dsp-level:0
\[
\quad L_{vf}=(3,3,4,3,5)\mapsto(3,3,2,3,5)\ [\text{machine}];
\]

%dsp-level:0
\[
\qquad \text{the massless readings:}\ 2-2T_m(1)=0\ \text{and}\ 2-2T_m(\cos2\pi)=0\ \forall m\ [\text{machine over}\ \mathbb Z[x]]
\]

%dsp-level:0
\[
(2),
\]

%dsp-level:0
\[
(4)\text{--}(7):\ \text{the declaration rows cite their grounds and their loci; the §II.4 entries carry the registry grammar}\ [\text{machine: the logged-before-use document pins}].\ \blacksquare
\]

\textbf{Phys.} Two conflicts the corpus could compute but not settle are settled by declaration, each with its grounds on the record. The amplitude-domain question — literal against chain-local — resolves pointwise: the ensembles keep the frozen literal domain, and the sealed junctions, where the chain-local pass-through reads were always the forced datum, run the literal-forced domain E2, which is canonical because it fixes the gauge by vertex type rather than by a choice of decomposition; E1 stays computed and unadopted. The corner convention resolves to the vertex-figure pairing, the loop-consistent reading, which is where this part's whole defect ladder already lives — and the declaration is load-bearing, because the per-class statistics census provably differs between the variants.

\end{theorem}

\medskip\noindent\textit{\texorpdfstring{BM.140 — The ceiling study: the Platonic defect census, the dodecahedron selection theorem, the five-currency range table, and the adjudication of deficit 12}{BM.140 - The ceiling study: the Platonic defect census, the dodecahedron selection theorem, the five-currency range table, and the adjudication of deficit 12}.}\ Structural-constraints record: see Appendix B, Section B.2.8 of the Addendum.\par\medskip

\begin{theorem}[{\texorpdfstring{BM.141 — The Jensen residue: the exact closed form, its range, its zero locus, the two-chain companion, and the width adjudication}{BM.141 - The Jensen residue: the exact closed form, its range, its zero locus, the two-chain companion, and the width adjudication}}]\label{tag:BM-141}
\textbf{Thm.}

%dsp-level:0
\[
(1)\quad\textbf{the residue computed:}
\]

%dsp-level:0
\[
\qquad J_3(p)\ :=\ \langle E\rangle_{\mathrm{mat}}\bigl\langle E^{-1}\bigr\rangle_{\mathrm{mat}}-1\ =\ \frac{18-27p+11p^2}{2D^2}-1,
\]

%dsp-level:0
\[
\qquad D=3-3p+p^2,
\]

%dsp-level:0
\[
\qquad p=1-4\varepsilon_q
\]

%dsp-level:0
\[
\qquad [\mathrm{BM.94}(6),(8)\ \text{re-cited}];
\]

%dsp-level:0
\[
\boxed{\ J_3(p)\ =\ \frac{p\,(1-p)\,\bigl(9-10p+2p^2\bigr)}{2D^2}\ }
\]

%dsp-level:0
\[
\qquad [\text{machine: the polynomial identity}\ 18-27p+11p^2-2D^2=p(1-p)(9-10p+2p^2)\ \text{exact in}\ \mathbb Q[p]];
\]

%dsp-level:0
\[
(2)\quad\textbf{the sign, the zeros, and the range:}
\]

%dsp-level:0
\[
\qquad J_3\ge0\ \ \text{on}\ p\in(0,1]
\]

%dsp-level:0
\[
\qquad [\mathrm{BM.94}(8)\text{'s Cauchy--Schwarz row re-derived};\ \text{machine:}\ 9-10p+2p^2>0\ \text{on}\ (0,1]\ \text{— its roots}\ \frac{5\pm\sqrt7}2\ \text{both exceed}\ 1\ \text{in}\ \mathbb Q(\sqrt7):\ \frac{5-\sqrt7}2>1\iff3>\sqrt7\iff9>7];
\]

%dsp-level:0
\[
\boxed{\ J_3=0\iff p=1\iff\varepsilon_q=0\ }
\]

%dsp-level:0
\[
\qquad [\text{machine: the factor census}];
\]

%dsp-level:0
\[
\qquad J_3\bigl(\tfrac12\bigr)=\frac9{49}
\]

%dsp-level:0
\[
\qquad [\mathrm{BM.94}(8):\ \tfrac{58}{49}-1;\ \text{machine}];
\]

%dsp-level:0
\[
J_3<\frac15\ \ \text{on}\ (0,1)
\]

%dsp-level:0
\[
\qquad [12p^4-72p^3+125p^2-81p+18>0;\ \mathrm{BM.94}(8)\ \text{re-derived};\ \text{machine: exact rational bisection, no roots on}\ (0,1)];
\]

%dsp-level:0
\[
\boxed{\ \sup_{(0,1)}J_3=J_3(p_\ast)\in\Bigl(\frac{93}{500},\ \frac3{16}\Bigr),\qquad p_\ast\in\Bigl(\frac{288131}{524288},\ \frac{576263}{1048576}\Bigr)\ }
\]

%dsp-level:0
\[
\qquad [\text{machine: the unique critical point by Sturm isolation of}\ -44p^5+294p^4-792p^3+1062p^2-684p+162;\qquad J_3<\tfrac3{16}\ \text{strict:}\ 38p^4-228p^3+394p^2-252p+54>0\ \text{on}\ (0,1),\ \text{no roots}];
\]

%dsp-level:0
\[
(3)\quad\textbf{the two-chain companion:}
\]

%dsp-level:0
\[
\qquad \langle E\rangle^{(2)}_{\mathrm{mat}}=\frac4{1+u},
\]

%dsp-level:0
\[
\qquad \bigl\langle E^{-1}\bigr\rangle^{(2)}_{\mathrm{mat}}=\frac{3u+1}{4(1+u)}
\]

%dsp-level:0
\[
\qquad [\mathrm{BM.137}(2);\ \text{machine}];
\]

%dsp-level:0
\[
\boxed{\ J_2(p)\ =\ \frac{4-3p}{(2-p)^2}-1\ =\ \frac{p(1-p)}{(2-p)^2}\ =\ \frac{u(1-u)}{(1+u)^2}\ }
\]

%dsp-level:0
\[
\qquad [\text{machine: the identity exact in}\ \mathbb Q[p]];
\]

%dsp-level:0
\[
\boxed{\ J_2\le\frac18,\qquad\text{equality}\iff p=\frac23\iff\varepsilon_q=\frac1{12}\ }
\]

%dsp-level:0
\[
\qquad [\text{machine: the SOS certificate}\ \frac{(1+u)^2}8-u(1-u)=\frac{(3u-1)^2}8\ge0;\ \text{the stationary point}\ 3p^2-8p+4=(3p-2)(p-2)];
\]

%dsp-level:0
\[
J_2\bigl(\tfrac12\bigr)=\frac19<\frac9{49}=J_3\bigl(\tfrac12\bigr);
\]

%dsp-level:0
\[
\qquad \frac18<\frac15;
\]

%dsp-level:0
\[
\qquad \textbf{the two-chain residue is strictly tighter than the three-chain bar}
\]

%dsp-level:0
\[
\qquad [\text{machine}];
\]

%dsp-level:0
\[
(4)\quad\textbf{the exact crossing of the two residues:}
\]

%dsp-level:0
\[
\qquad J_3-J_2=\frac{p\,(p-1)\,\bigl(6p^3-27p^2+40p-18\bigr)}{2D^2(2-p)^2}
\]

%dsp-level:0
\[
\qquad [\text{machine: the factorization exact in}\ \mathbb Q[p]];
\]

%dsp-level:0
\[
6p^3-27p^2+40p-18\ \ \text{is strictly increasing on}\ (0,1)\ \bigl[\text{its derivative}\ 18p^2-54p+40\ \text{has roots}\ \tfrac43,\tfrac53,\ \text{both outside}\bigr]\ \text{with the single root}\ p_c\in\Bigl(\frac{108267}{131072},\ \frac{866137}{1048576}\Bigr)
\]

%dsp-level:0
\[
\qquad [\text{machine: Sturm isolation}];
\]

%dsp-level:0
\[
\boxed{\ J_3>J_2\iff p<p_c;\qquad J_3<J_2\iff p>p_c\ }\qquad[\text{machine: the sign census on exact rationals}];
\]

%dsp-level:0
\[
(5)\quad\textbf{the width adjudication (the two formed reads of}\ \Delta x_{\mathrm B}\textbf{):}
\]

%dsp-level:0
\[
\qquad \Delta x_{\mathrm B}^{\mathrm{eff}}=\frac{\mathfrak h}{\langle E\rangle_{\mathrm{mat}}\ln2}=D\,\Delta x_{\mathrm B}^{cryst};
\]

%dsp-level:0
\[
\qquad \frac{\mathfrak h}{\ln2}\bigl\langle E^{-1}\bigr\rangle_{\mathrm{mat}}\ \ge\ \Delta x_{\mathrm B}^{\mathrm{eff}}
\]

%dsp-level:0
\[
\qquad [\mathrm{BM.94}(8)\ \text{re-cited}];
\]

%dsp-level:0
\[
\boxed{\ \textbf{the relative gap of the two reads is exactly the Jensen residue:}\qquad\Bigl(\bigl\langle E^{-1}\bigr\rangle_{\mathrm{mat}}-\frac1{\langle E\rangle_{\mathrm{mat}}}\Bigr)\langle E\rangle_{\mathrm{mat}}=J_3\ }
\]

%dsp-level:0
\[
\qquad [\text{machine: the identity}];
\]

%dsp-level:0
\[
\textbf{the adjudication:}
\]

%dsp-level:0
\[
\qquad \text{the tier's evaluations at the width of the mean}\ [\mathrm{BM.94}(8);\ \mathrm{BM.98}(5)]\ \textbf{stand, with the difference now an exact computed function instead of a residue};
\]

%dsp-level:0
\[
(6)\quad\textbf{status:}
\]

%dsp-level:0
\[
\qquad \textbf{derived:}\ (1)\text{--}(5);
\]

%dsp-level:0
\[
\qquad \boxed{\ (\mathrm{R6})\ \textbf{is closed: the Jensen residue is computed as an exact rational function with its zero locus, its exact supremum enclosure and its two-chain companion}\ }
\]

%dsp-level:0
\[
[\mathrm{BM.94}(8),\ \mathrm{BM.98}(5)\text{'s residue rows discharged}];
\]

%dsp-level:0
\[
\qquad \neg\text{a declaration},
\]

%dsp-level:0
\[
\qquad \neg\text{a P4 addition},
\]

%dsp-level:0
\[
\qquad \neg\text{a measured value}\ [\mathrm{P5}].
\]

\textbf{Req.} A4 (Part 1); D0.4 (Part 1); CT0.8 (Part 1); P5 (Parts 1–2); BM.94(6), BM.94(8), BM.98(5), BM.137(2) (this part); polynomial factorization over $\mathbb Q$, Sturm isolation and exact rational bisection, SOS certificates, the Cauchy–Schwarz equality locus [math; machine: \texttt{checks/\allowbreak{}verify\_\allowbreak{}H\_\allowbreak{}bm\_\allowbreak{}closing\_\allowbreak{}tier.py} §J].

\textbf{Pf.}

%dsp-level:0
\[
(1):\ 18-27p+11p^2-2(3-3p+p^2)^2=p(1-p)(9-10p+2p^2)\ [\text{machine: expansion in}\ \mathbb Q[p]];
\]

%dsp-level:0
\[
\qquad (2):\ \text{the root census of}\ 9-10p+2p^2\ \text{in}\ \mathbb Q(\sqrt7);
\]

%dsp-level:0
\[
\qquad \text{the bisection of}\ 12p^4-72p^3+125p^2-81p+18\ \text{and}\ 38p^4-228p^3+394p^2-252p+54\ \text{on}\ (0,1)\ \text{at exact rational endpoints}\ [\text{machine}]
\]

%dsp-level:0
\[
(3):\ \bigl\langle E^{-1}\bigr\rangle^{(2)}_{\mathrm{mat}}=\frac{\frac12\cdot2u+\frac14p}{1+u}\ \text{normalized}\ [\mathrm{BM.137}(2)\text{'s measure};\ \text{machine}];
\]

%dsp-level:0
\[
\qquad J_2=\frac4{1+u}\cdot\frac{3u+1}{4(1+u)}-1=\frac{3u+1-(1+u)^2}{(1+u)^2}=\frac{u(1-u)}{(1+u)^2}\ [\text{machine}];
\]

%dsp-level:0
\[
\qquad \text{the SOS certificate and the stationary factorization}\ [\text{machine}]
\]

%dsp-level:0
\[
(4):\ J_3-J_2\ \text{brought to the common denominator}\ 2D^2(2-p)^2\ \text{and factored}\ [\text{machine}];
\]

%dsp-level:0
\[
\qquad (5):\ \bigl\langle E^{-1}\bigr\rangle\langle E\rangle-\langle E\rangle/\langle E\rangle=J_3\ [\text{machine}];
\]

%dsp-level:0
\[
\qquad (6):\ \text{assembly}.\ \blacksquare
\]

\textbf{Phys.} The last pure-mathematics residue of the era is discharged: the Jensen residue of the materialized flicker distribution is one exact rational function, $p(1-p)(9-10p+2p^2)$ over twice the squared geometric sum, vanishing precisely at the sharp crystallization point, capped strictly below one fifth with its supremum enclosed between two rationals and attained at a uniquely isolated critical point. The two-chain companion is tighter — bounded by one eighth with a sum-of-squares certificate and equality at threshold over twelve — and the two residues cross exactly once, at an isolable rational-bracketed point. The gap between the two formed width reads of the Born envelope is precisely this residue, so the era's evaluations at the width of the mean stand with their correction now computed rather than named.

\end{theorem}

\begin{theorem}[{\texorpdfstring{BM.142 — The supply vertex, the pair amplitude, and the annihilation cross-section: the four grounds resolved and residues $(\mathrm{R9}')$, $(\mathrm{R9}'')$ closed}{BM.142 - The supply vertex, the pair amplitude, and the annihilation cross-section: the four grounds resolved and residues , closed}}]\label{tag:BM-142}
\textbf{Thm.}

%dsp-level:0
\[
(1)\quad\textbf{the four grounds of}\ \mathrm{BM.122}(6)\ \textbf{resolved, one by one:}\qquad\begin{array}{l|l|l}\text{ground} & \text{the resolution} & \text{the row}\\\hline(\mathrm O1)\ \neg\text{a gauged chiral coupling} & \textbf{dissolved for the EM legs:}\ \text{the gauged coupling of the charge read is constructed chart-side} & [\mathrm{Def\text{-}BM.12}(1)\text{--}(3):\ \partial_\mu F^{\mu\nu}=eJ_Q^\nu;\ \mathrm{EM.15}]\\\text{} & \text{the W-species chiral gauging remains}\ \neg\text{constructed}\ —\ \text{barred for weak legs, recorded} & [\mathrm{BM.76};\ \mathrm{A\text{-}BM.7}\ \text{scope}]\\\hline(\mathrm O2)\ \text{deficit}\ 1' & \textbf{discharged:}\ \text{the assignment rule is issued and the attachment derived} & [\mathrm{A\text{-}BM.7};\ \mathrm{BM.73}(1),(2);\ \mathrm{BM.75}(3)]\\\hline(\mathrm O3)\ \neg\text{a momentum resolvent} & \textbf{discharged:}\ \text{the resolvent and the phase space are constructed} & [\mathrm{Def\text{-}BM.20};\ \mathrm{BM.134}]\\\hline(\mathrm O4)\ \text{the transit leg carries no winding carrier} & \textbf{re-scoped:}\ \text{the exclusion stands, but the vertex needs}\ \neg\text{an}\ \varepsilon/j\ \text{read of the legs} & [\mathrm{Def\text{-}H.38}\ \text{Scope frozen}];\ \text{the legs couple by}\ (q,\ \text{helicity}):\ N_{\mathrm{pol}}=2\ [\mathrm{P2};\ \mathrm{EM.8}]\\\text{} & \text{and the register's moves apply to transit legs} & [\mathrm{BM.126}(5)\ \text{domain row}]\end{array}
\]

%dsp-level:0
\[
(2)\quad\textbf{the supply vertex (constructed):}
\]

%dsp-level:0
\[
\qquad \boxed{\ \widehat{\mathcal V}_{\mathrm{sup}}\ :=\ e(k,E)^2\cdot\bigl(\hat J_+\otimes\hat J_+\bigr)\ }
\]

%dsp-level:0
\[
\qquad \text{on}\ \mathcal J\otimes\mathcal J\otimes\bigl(\text{the two transit helicity legs}\bigr)
\]

%dsp-level:0
\[
\qquad [\mathrm{Def\text{-}BM.18}(2);\ \mathrm{Def\text{-}BM.12}(2);\ \mathrm{Def\text{-}BM.20}(4)];
\]

%dsp-level:0
\[
\bigl(\hat J_+\otimes\hat J_+\bigr)^2=0
\]

%dsp-level:0
\[
\qquad [\hat J_+^2=0;\ \mathrm{BM.120}(1);\ \text{machine: the}\ 4\times4\ \text{identity over}\ \mathbb Z]:
\]

%dsp-level:0
\[
\qquad \boxed{\ \textbf{the vertex is nilpotent: the register cannot supply a pair twice}\ };
\]

%dsp-level:0
\[
\textbf{the vertex rows:}
\]

%dsp-level:0
\[
\qquad \text{charge balance}\ Q_{\mathrm{in}}=Q_{\mathrm{out}}=0\ [\mathrm{BM.122}(1)];
\]

%dsp-level:0
\[
\qquad \text{the threshold gate}\ \chi_{\{NE\ge\mathcal T_0\}}\ [\mathrm{BM.122}(2);\ \mathrm{Def\text{-}BM.17}(3)];
\]

%dsp-level:0
\[
\qquad \text{the helicity census}\ 2\times2
\]

%dsp-level:0
\[
= 4\ \text{combinations, equally weighted by the metric's charge-blindness}\ [\mathrm{P2}:\ N_{\gamma\text{-pol}}=2;\ \mathrm{EM.8};\ \mathrm{BM.58}];
\]

%dsp-level:0
\[
(3)\quad\textbf{the pair amplitude (formed, symbolic):}
\]

%dsp-level:0
\[
\qquad \boxed{\ \mathcal M_{\mathrm{pair}}(z;k,E)\ =\ e(k,E)^2\,q^2\,x_+^2\,\chi\ =\ e_0^2\,q^2x_+^2\,\chi\cdot\frac{N(k,E)\,f(k)}{1+G(E)}\ }
\]

%dsp-level:0
\[
\qquad [\mathrm{Def\text{-}BM.20}(4)\text{'s}\ \mathcal V=1;\ \mathrm{BM.133}(6);\ \mathrm{BM.132}(2);\ \text{machine: the monomial census}];
\]

%dsp-level:0
\[
[\mathcal M]=\mathfrak q^0,\qquad[e]=\mathfrak q^0,\qquad[x_\pm]=\mathfrak q^0\qquad[\mathrm{CT2.1};\ \mathrm{BM.102}(2);\ \text{machine: the exponent census}];
\]

%dsp-level:0
\[
\textbf{the reduction clause:}
\]

%dsp-level:0
\[
\qquad e=1,\ q=1:\quad\mathcal M_{\mathrm{pair}}=\mathcal W_{\text{channel}}=x_+^2\chi
\]

%dsp-level:0
\[
\qquad [\mathrm{BM.122}(3);\ \mathrm{Def\text{-}BM.20}(4);\ \text{machine}];
\]

%dsp-level:0
\[
\textbf{the exact zeros:}
\]

%dsp-level:0
\[
\qquad \mathcal M=0\iff e_0=0\ \vee\ q=0\ \vee\ \varepsilon_q=\tfrac14\ \vee\ \chi=0\ \vee\ N=0\ \vee\ f=0
\]

%dsp-level:0
\[
\qquad [\text{machine: the factor census}];
\]

%dsp-level:0
\[
(4)\quad\textbf{the pair cross-section (the formed chart read):}
\]

%dsp-level:0
\[
\qquad \sigma_{\mathrm{pair}}:=\frac{|\mathcal M_{\mathrm{pair}}|^2\,\Phi_2\bigl(z;\,2m_{\mathrm{hid}},N\bigr)}{F_{\gamma\gamma}}
\]

%dsp-level:0
\[
\qquad [\mathrm{BM.134}(4);\ \text{the final state is the massive pair}]:
\]

%dsp-level:0
\[
\boxed{\ \sigma_{\mathrm{pair}}(z;k,E)\ =\ \frac{e_0^4\,q^4x_+^4\,\chi\ N\,f^2}{16\pi\,z^3\,(1+G)^2}\ \sqrt{z^2-4N^2\bigl(2m_{\mathrm{hid}}\bigr)^2}\ }
\]

%dsp-level:0
\[
\qquad [\text{machine: the coefficient algebra exact};\ m_{\mathrm{pair}}=2m_{\mathrm{hid}}\ [\mathrm{BM.111}(1)]];
\]

%dsp-level:0
\[
[\sigma_{\mathrm{pair}}]=\mathfrak q^{-1}\ =\ \text{the area species of}\ (\Delta x_{\mathrm B})^2\qquad[\text{machine: the exponent census}];
\]

%dsp-level:0
\[
\textbf{the threshold rows:}
\]

%dsp-level:0
\[
\qquad \sigma_{\mathrm{pair}}=0\ \ \text{exactly at}\ \ z=2N\cdot2m_{\mathrm{hid}}\ \ (\text{the square-root onset});
\]

%dsp-level:0
\[
\qquad \sigma_{\mathrm{pair}}>0\ \ \text{strictly above it};
\]

%dsp-level:0
\[
\qquad \text{formed on}\ \ z\ge4Nm_{\mathrm{hid}}\ \wedge\ k\ge\frac G{4\pi}\ \wedge\ NE\ge\mathcal T_0
\]

%dsp-level:0
\[
\qquad [\text{machine}];
\]

%dsp-level:0
\[
(5)\quad\textbf{the annihilation cross-section:}
\]

%dsp-level:0
\[
\qquad \mathcal M_{\mathrm{ann}}=e(k,E)^2q^2x_-^2\chi
\]

%dsp-level:0
\[
\qquad [\mathrm{BM.123}(1)\text{'s reverse move}\ \hat J_-\otimes\hat J_-;\ \text{the same vertex normalization};\ m\bigl(c^{\mathrm{CP}}\bigr)=m(c)\ [\mathrm{BM.27};\ \mathrm{BM.108}(3)]];
\]

%dsp-level:0
\[
\sigma_{\mathrm{ann}}:=\frac{|\mathcal M_{\mathrm{ann}}|^2\,\Phi_2^{(\gamma\gamma)}}{F_{\mathrm{pair}}}\qquad[\text{the final state massless},\ \text{the initial state massive}]:
\]

%dsp-level:0
\[
\boxed{\ \sigma_{\mathrm{ann}}(z;k,E)\ =\ \frac{e_0^4\,q^4x_-^4\,\chi\ N\,f^2}{16\pi\,z\,(1+G)^2\,\sqrt{z^2-4N^2\bigl(2m_{\mathrm{hid}}\bigr)^2}}\ }
\]

%dsp-level:0
\[
\qquad [\text{machine: the coefficient algebra exact};\ [\sigma_{\mathrm{ann}}]=\mathfrak q^{-1}];
\]

%dsp-level:0
\[
\textbf{the threshold behaviour:}
\]

%dsp-level:0
\[
\qquad \sigma_{\mathrm{ann}}\ \ \text{has the exact square-root pole at}\ \ z=4Nm_{\mathrm{hid}};
\]

%dsp-level:0
\[
\qquad \sigma_{\mathrm{ann}}\cdot F_{\mathrm{pair}}\ \ \text{is regular there}
\]

%dsp-level:0
\[
\qquad [\text{machine: the product}\ =\ \frac{e_0^4q^4x_-^4\chi\,f^2}{8\pi\,N\,(1+G)^2}\ \text{regular in the formed factors}];
\]

%dsp-level:0
\[
(6)\quad\textbf{the ratio and the detailed-balance elevation:}
\]

%dsp-level:0
\[
\qquad \boxed{\ \frac{\sigma_{\mathrm{ann}}}{\sigma_{\mathrm{pair}}}\ =\ \Bigl(\frac{x_-}{x_+}\Bigr)^{4}\cdot\frac{z^2}{z^2-4N^2(2m_{\mathrm{hid}})^2}\ }
\]

%dsp-level:0
\[
\qquad [\text{machine: the symbolic division exact}];
\]

%dsp-level:0
\[
\textbf{the amplitude-level detailed balance:}
\]

%dsp-level:0
\[
\qquad |\mathcal M_{\mathrm{ann}}|=|\mathcal M_{\mathrm{pair}}|\iff\varepsilon_q=\frac18
\]

%dsp-level:0
\[
\qquad [\mathrm{BM.123}(2)\text{'s indicator elevated to the amplitude};\ \text{machine}];
\]

%dsp-level:0
\[
\text{at}\ \varepsilon_q
\]

%dsp-level:0
\[
= \frac18\ \text{the cross-section ratio is the pure kinematic factor}\ \frac{z^2}{z^2-16N^2m_{\mathrm{hid}}^2}>1\ \ \text{strictly above threshold}
\]

%dsp-level:0
\[
\qquad [\text{machine}];
\]

%dsp-level:0
\[
(7)\quad\textbf{the width row and its species bar:}
\]

%dsp-level:0
\[
\qquad \text{the chart-frequency width}\ \Gamma^{\mathrm{chart}}
\]

%dsp-level:0
\[
= \frac{|\mathcal M|^2\Phi_2}{2z}\ \ \text{carries}\ [\Gamma^{\mathrm{chart}}]
\]

%dsp-level:0
\[
= \mathfrak q^{-1/2}\ \ \text{at}\ \mathcal V
\]

%dsp-level:0
\[
= 1\ —\ \text{a time,}\ \neg\text{a rate}:
\]

%dsp-level:0
\[
\text{the mass-dressed normalization that would convert it is}\ \neg\text{declared}\ [\mathrm{Def\text{-}BM.20}(4)\text{'s minimality}]\implies\boxed{\ \textbf{the formed width read stays the ledger's per-tick species}\ }\ [\mathrm{BM.116}(1);\ \mathrm{BM.44}(4),(5):\ [\Gamma]=\mathfrak q^0\ \text{ticks}^{-1};\ \tau_{\mathrm{life}}\ge3\ \text{ticks re-cited}];
\]

%dsp-level:0
\[
(8)\quad\textbf{the closure rows:}
\]

%dsp-level:0
\[
\qquad \boxed{\ (\mathrm{R9}')\ \textbf{is closed: the annihilation cross-section is constructed as an exact symbolic chart-side read on the formed domain}\ }
\]

%dsp-level:0
\[
\boxed{\ (\mathrm{R9}'')\ \textbf{is closed at the vertex level: the two-transit-quantum supply vertex is constructed, nilpotent, charge-balanced and threshold-gated}\ }
\]

%dsp-level:0
\[
\qquad [\mathrm{BM.118}(7),\ \mathrm{BM.130}(7)\text{'s rows discharged}];
\]

%dsp-level:0
\[
\mathrm{BM.130}(3)\text{'s}\ \neg\text{cross-section row is discharged by construction and re-scoped,}\ \neg\text{retracted};
\]

%dsp-level:0
\[
\qquad \mathrm{BM.122}(6)\text{'s four grounds carry}\ (1)\text{'s per-row dispositions};
\]

%dsp-level:0
\[
(9)\quad\textbf{the bars:}
\]

%dsp-level:0
\[
\qquad \neg\text{a numerical cross-section, rate or branching fraction is evaluated}\ [\mathrm{P5};\ \mathrm{BM.44}(5)];
\]

%dsp-level:0
\[
\qquad \neg\text{an S-matrix unitarity claim}\ [\mathrm{Def\text{-}BM.20}(5)];
\]

%dsp-level:0
\[
\qquad \text{amplitude}\ne\text{mass}\ [\mathrm{FT.79}(5)];
\]

%dsp-level:0
\[
e_0\ \ \text{remains calibration}\ [\mathrm{EM.25}(7);\ \mathrm{BM.143}(4)];
\]

%dsp-level:0
\[
\qquad \neg\text{a measured external constant is identified with}\ \sigma_{\mathrm{pair}},\ \sigma_{\mathrm{ann}}\ \text{or any factor}\ [\mathrm{P5}];
\]

%dsp-level:0
\[
\qquad \neg\text{a species name enters any clause}\ [\mathrm{Def\text{-}BM.3}(7)].
\]

\textbf{Req.} A4, A8 (Part 1); D0.4, D1.2 (Part 1); CT1.22, CT2.1 (Part 2); EM.8, EM.15, EM.25(7) (Part 2); P2, P5 (Parts 1–2); Def-H.38 (Part 3); FT.79(5) (Part 3); Def-BM.12, Def-BM.17(3), Def-BM.18(2), Def-BM.20, BM.27, BM.44(4), BM.44(5), BM.58, BM.73(1), BM.73(2), BM.75(3), BM.76, BM.102(2), BM.108(3), BM.111(1), BM.116(1), BM.118(7), BM.120(1), BM.122, BM.123, BM.126(5), BM.130(3), BM.130(7), BM.132(2), BM.133(6), BM.134(4) (this part); the declaration addendum A-BM.7 (pointers); exact matrix identities over $\mathbb Z$, monomial censuses in symbolic rings, the coefficient algebra of the two-body reads [math; machine: \texttt{checks/\allowbreak{}verify\_\allowbreak{}H\_\allowbreak{}bm\_\allowbreak{}closing\_\allowbreak{}tier.py} §K].

\textbf{Pf.}

%dsp-level:0
\[
(1):\ \text{the four rows cite their discharging tags; the W-species bar and the}\ \mathrm{E2}\text{-domain inheritance recorded}\ [\mathrm{Def\text{-}BM.20}(5);\ \text{machine: the document pins}]
\]

%dsp-level:0
\[
(2):\ \bigl(\hat J_+\otimes\hat J_+\bigr)^2=\hat J_+^2\otimes\hat J_+^2=0\otimes0=0\ [\mathrm{BM.120}(1);\ \text{machine: the}\ 4\times4\ \text{product over}\ \mathbb Z];
\]

%dsp-level:0
\[
\qquad \text{the helicity census}\ 2^2=4\ [\mathrm{P2};\ \text{machine}]
\]

%dsp-level:0
\[
(3),(4):\ \mathcal M\ \text{assembled from}\ e^2=e_0^2Nf\mathsf g=e_0^2Nf/(1+G)\ [\mathrm{BM.133}(6);\ \text{machine}];
\]

%dsp-level:0
\[
\qquad \sigma_{\mathrm{pair}}
\]

%dsp-level:0
\[
= \frac{|\mathcal M|^2}{F_{\gamma\gamma}}\Phi_2
\]

%dsp-level:0
\[
= \frac{e_0^4q^4x_+^4\chi N^2f^2}{(1+G)^2}\cdot\frac{\sqrt{z^2-16N^2m^2}}{8\pi N^3z}\cdot\frac{N^2}{2z^2}\ [\mathrm{BM.134}(4);\ \text{machine: the coefficient algebra}];
\]

%dsp-level:0
\[
\text{the exponent census}\ [\sigma]=\mathfrak q^{-1}\ [\mathrm{CT2.1}]
\]

%dsp-level:0
\[
(5),(6):\ \sigma_{\mathrm{ann}}=\frac{e_0^4q^4x_-^4\chi N^2f^2}{(1+G)^2}\cdot\frac1{8\pi N^3}\cdot\frac{N^2}{2z\sqrt{z^2-16N^2m^2}}\ [\text{machine}];
\]

%dsp-level:0
\[
\qquad \text{the ratio by symbolic division}\ [\text{machine}];
\]

%dsp-level:0
\[
\qquad \text{the detailed-balance locus}\ x_-=x_+\iff\varepsilon_q=\tfrac18\ [\mathrm{BM.120}(5);\ \mathrm{BM.123}(2)]
\]

%dsp-level:0
\[
(7):\ \bigl[\tfrac{|\mathcal M|^2\Phi_2}{2z}\bigr]=\mathfrak q^0\cdot\mathfrak q^0\cdot\mathfrak q^{-1/2}=\mathfrak q^{-1/2}\ [\text{machine: the exponent census}];
\]

%dsp-level:0
\[
\qquad (8),(9):\ \text{assembly and the token scan}.\ \blacksquare
\]

\textbf{Phys.} The two-photon bar falls ground by ground: the electromagnetic gauging of the charge read existed chart-side, the assignment rule was issued, the resolvent is now constructed, and the transit legs couple by charge and helicity — the one exclusion that stands, the winding carrier of a trivial chain, simply never entered the vertex. What is built is minimal and honest: a nilpotent supply vertex that cannot fire twice, an amplitude that reduces to the register's own channel weight at unit coupling, and two cross-sections that are exact symbolic functions of the envelope ratio, the local energy, the flicker threshold and the calibration coupling — with an exact square-root onset for production, an exact square-root pole for annihilation whose flux product stays regular, and a ratio that carries the detailed-balance point of the register elevated to the amplitude level. The chart-frequency width is refused on dimensional grounds rather than fudged: the per-tick ledger species remains the corpus's width read. Both residues the annihilation tier named are closed, and nothing measured entered any row.

\end{theorem}

\medskip\noindent\textit{\texorpdfstring{BM.143 — The house cleaning: the void residues closed, the phantom gaps reclassified, the calibration register consolidated, and the Part-5 requirements named}{BM.143 - The house cleaning: the void residues closed, the phantom gaps reclassified, the calibration register consolidated, and the Part-5 requirements named}.}\ Structural-constraints record: see Appendix B, Section B.4.2 of the Addendum.\par\medskip

\medskip\noindent\textit{\texorpdfstring{BM.144 — The status of the Closing Tier and the closure of the Baryomesonic era}{BM.144 - The status of the Closing Tier and the closure of the Baryomesonic era}.}\ Structural-constraints record: see Appendix B, Section B.1.34 of the Addendum.\par\medskip

\phantomsection
\section[Tier BM.24 — The Continuum-Limit Tier: the Force Operator…]{\texorpdfstring{Tier BM.24 — The Continuum-Limit Tier: the Force Operator over the Graph, the Continuum-Limit Function, the Metric–Energy–Continuum Map, the Breaking Threshold, the Locked-Sector Complement, the Universality, the Reclassification of the Continuum Conditions, and the Closure of the Foundational Era}{Tier BM.24 - The Continuum-Limit Tier: the Force Operator over the Graph, the Continuum-Limit Function, the Metric-Energy-Continuum Map, the Breaking Threshold, the Locked-Sector Complement, the Universality, the Reclassification of the Continuum Conditions, and the Closure of the Foundational Era}}

\begin{theorem}[{\texorpdfstring{BM.145 — The force operator acting over the graph: the generator census, the mask ladder, and the derived action}{BM.145 - The force operator acting over the graph: the generator census, the mask ladder, and the derived action}}]\label{tag:BM-145}
\textbf{Thm.}

%dsp-level:0
\[
(1)\quad\textbf{the complete operator (re-cited, exact):}
\]

%dsp-level:0
\[
\qquad \mathbf F=F_{(+,\mathcal R)}\oplus F_{(-,\mathcal R)}\oplus F_{(+,\mathcal V)}\oplus F_{(-,\mathcal V)}\oplus F_{\mathrm{st}}\ \ \text{always applies in full}
\]

%dsp-level:0
\[
\qquad [\mathrm{A8};\ \mathrm{BM.52}(1)];
\]

%dsp-level:0
\[
F_{(+,\mathcal R)}=F_G\oplus F_H\oplus F_{D^+};
\]

%dsp-level:0
\[
\qquad F_{(-,\mathcal R)}=F_{W^+}\oplus F_{W^-}\oplus F_Z;
\]

%dsp-level:0
\[
\qquad F_{(+,\mathcal V)}=F_{J^+};
\]

%dsp-level:0
\[
\qquad F_{(-,\mathcal V)}=F_{J^-};
\]

%dsp-level:0
\[
\qquad F_{\mathrm{st}}=F_{\mathrm{EM}}\oplus F_{\mathrm{strong}}\oplus F_{D^-}
\]

%dsp-level:0
\[
\qquad [\mathrm{A8}];
\]

%dsp-level:0
\[
\textbf{the pair-arena block census:}\qquad\begin{array}{c|c|c}\text{aspect} & \text{the block} & \dim_{\mathbb R}\\\hline F_{J^+} & \mathcal E_{-+}=\operatorname{span}\{J_i\} & 3\\ F_{J^-} & \mathcal E_{--}=\operatorname{span}\{K_i\} & 3\\ F_G & \operatorname{span}\{\hat E^{ij}\}\ \text{traceless} & 5\\ F_W & \mathcal E_{+-}=\operatorname{span}\{W_i\} & 3\\ F_H & \operatorname{span}\{\Sigma\} & 1\\ F_{D^+} & \operatorname{span}\{E_0\} & 1\\\hline\text{total} & \mathfrak H(\mathbb C^4) & 16=C_{\mathrm{cap}}\end{array}\qquad[\mathrm{CT0.12};\ \mathrm{CT1.1};\ \mathrm{P2}:\ \text{machine:}\ 3+3+5+3+1+1=16];
\]

%dsp-level:0
\[
\textbf{the stack aspects act on the vertex/slot registers,}\ \neg\textbf{on the pair channel:}
\]

%dsp-level:0
\[
\qquad F_{\mathrm{EM}},\ F_{\mathrm{strong}},\ F_{D^-}
\]

%dsp-level:0
\[
\qquad [\mathrm{CT0.12}\ \mathrm{Rmk};\ \mathrm{D3.1};\ \mathrm{D4.1}];
\]

%dsp-level:0
\[
(2)\quad\textbf{the activation and the mask ladder:}
\]

%dsp-level:0
\[
\qquad \mathcal A(0)
\]

%dsp-level:0
\[
= \{J\};
\]

%dsp-level:0
\[
\quad\mathcal A(1)
\]

%dsp-level:0
\[
= \{J,G,W,H,D^+\};
\]

%dsp-level:0
\[
\quad\mathcal A(2)
\]

%dsp-level:0
\[
= \mathcal A(1)\cup\{\mathrm{EM}\};
\]

%dsp-level:0
\[
\quad\mathcal A(3)
\]

%dsp-level:0
\[
= \mathcal A(2)\cup\{\mathrm{strong}\};
\]

%dsp-level:0
\[
\quad\mathcal A(n\ge4)
\]

%dsp-level:0
\[
= \mathcal A(3)\cup\{D^-\}
\]

%dsp-level:0
\[
\qquad [\mathrm{CT0.13}];
\]

%dsp-level:0
\[
\mathbf F_M^{(n)}=\bigoplus_{A\in\mathcal A(n)}F_A\big|_{\mathcal K^{(A)}}
\]

%dsp-level:0
\[
\qquad [\mathrm{CT2.4}:\ \text{masking never mixes forces}];
\]

%dsp-level:0
\[
\qquad \dim\mathcal K_n=(6,\ 16,\ 17,\ 25)\ \ \text{strict}
\]

%dsp-level:0
\[
\qquad [\mathrm{CT0.14};\ \mathrm{BM.36}(1)];
\]

%dsp-level:0
\[
\textbf{the increments are derived:}
\]

%dsp-level:0
\[
\qquad 1\to2:\ +1=\dim Z\bigl(\iota_2(\mathfrak u(2))\bigr)=\dim\mathcal K^{(\mathrm{EM})}
\]

%dsp-level:0
\[
\qquad [\mathrm{D2.12}];
\]

%dsp-level:0
\[
\qquad 2\to3:\ +8=\dim\mathfrak{su}(3)=\dim\mathcal K^{(\mathrm{strong})}
\]

%dsp-level:0
\[
\qquad [\mathrm{CT3.8}];
\]

%dsp-level:0
\[
\text{the observable bound}\ \dim\mathcal K_\Sigma\le1+5+1+3+8=18\qquad[\mathrm{FT.26}(3);\ \text{machine: the integer census}];
\]

%dsp-level:0
\[
(3)\quad\textbf{the action over the graph (the composition, every step derived):}
\]

%dsp-level:0
\[
\qquad \boxed{\ \hat{\mathcal F}:\ \bigl(\Gamma_{\mathcal R},\ \text{registrations}\bigr)\ \longmapsto\ \bigl(M,\ g_{\mu\nu},\ T_{\mu\nu},\ \operatorname{spec}\bigr)\ }
\]

%dsp-level:0
\[
\begin{array}{l|l|l}\text{step} & \text{the map} & \text{the locus}\\\hline\text{the weights} & (\Gamma,\sigma_e,\kappa_e,\alpha_e)\mapsto w(C)=\mathcal J(C)\prod_{e\in R}G(e)\,(-1)^{\aleph},\quad G(e)=\sigma_e e^{i\pi\kappa_e} & [\mathrm{A8};\ \mathrm{D3.5};\ \mathrm{D3.6};\ \mathrm{FT.9}]\\\hline\text{the domain} & \aleph\ \text{literal on the ensembles,}\ \mathrm{E2}\ \text{at the sealed junctions} & [\mathrm{Def\text{-}H.32}\ \text{Scopes};\ \text{A-BM.10};\ \mathrm{BM.139}(2)]\\\hline\text{the dictionary} & C\mapsto(a_{00},q_{ij},w_i)\mapsto g_{\mu\nu}=\Omega^2\bar g_{\mu\nu},\quad\Omega=e^{\alpha_\Omega a_{00}} & [\mathrm{D2.1};\ \mathrm{CT1.12}]\\\hline\text{the grouping} & Z[g]=\sum_{C:\,g(C)=g}|w(C)|\,e^{i\Phi(C)} & [\mathrm{CT1.20}]\\\hline\text{the action} & S_{\mathrm{eff}}=-i\log Z;\quad S_{\mathrm{geom}}=\frac1{2\kappa^2}\int\sqrt{-g}\,(R-2\Lambda) & [\mathrm{CT1.20};\ \mathrm{CT1.21};\ \mathrm{CT2.7}]\\\hline\text{the spectrum} & Z(\mathcal R,n)=\langle\mathrm{out}|(1-\mathbf F_M^{(n)})^{-1}|\mathrm{in}\rangle=\sum_a\frac{\langle\mathrm{out}|P_a|\mathrm{in}\rangle}{1-\lambda_a};\quad\text{poles}\iff1\in\operatorname{spec}\mathbf F_M^{(n)} & [\mathrm{CT2.5}]\\\hline\text{the source} & T_{\mu\nu}=-\frac2{\sqrt{-g}}\frac{\delta S_{\mathrm{force}}}{\delta g^{\mu\nu}}=T^{\mathrm{bd}}_{\mu\nu}+T^{\mathrm{tr}}_{\mu\nu} & [\mathrm{CT1.22}]\\\hline\text{the field equations} & G_{\mu\nu}+\Lambda g_{\mu\nu}=\kappa^2T_{\mu\nu} & [\mathrm{CT1.23}]\\\hline\text{the target} & M=\Gamma_{\mathcal R}/\!\!\sim\ \ \text{a smooth 4-manifold with the charts}\ (\tau_n,h_1,h_2,h_3) & [\mathrm{CT2.6};\ \mathrm{D1.8}]\end{array}
\]

%dsp-level:0
\[
\textbf{the composition is the action:}
\]

%dsp-level:0
\[
\qquad \text{no step adds a datum}\ —\ \text{each row is a derived map of the cited tag}
\]

%dsp-level:0
\[
\qquad [\text{machine: the locus census, row by row}];
\]

%dsp-level:0
\[
(4)\quad\textbf{the state-side action (the second face, separated by species):}
\]

%dsp-level:0
\[
\qquad \text{the Haar draws and the Gibbs tilt act on the pair states,}\ \neg\text{on the arena blocks:}
\]

%dsp-level:0
\[
\qquad \mathcal T(\rho)=U\rho U^\dagger,\quad\tilde{\mathcal T}_\gamma(\rho)=\rho e^{-\gamma\rho}/Z_\gamma
\]

%dsp-level:0
\[
\qquad [\mathrm{A6};\ \mathrm{D0.3};\ \mathrm{L0}];
\]

%dsp-level:0
\[
V'=\Bigl(1-\frac\gamma4\Bigr)^2V+O(\delta\rho^3),\qquad\mu'=\Bigl(1-\frac\gamma4\Bigr)\mu+O(\mu^2)\qquad[\mathrm{L1};\ \mathrm{L2}];\qquad\text{the duality rows}\ [\mathrm{BM.131}];
\]

%dsp-level:0
\[
\textbf{the two actions meet only through the dictionary:}
\]

%dsp-level:0
\[
\qquad \text{the sector coefficients of}\ \mathrm{D2.1}\ \text{are the pair-state components}\ [\mathrm{CT1.1};\ \mathrm{L3}(1)];
\]

%dsp-level:0
\[
\qquad \text{the Haar face preserves their spectrum and the tilt face contracts it}\ [\mathrm{BM.131}(3);\ \mathrm{L2}]:
\]

%dsp-level:0
\[
\boxed{\ \text{both actions preserve the derived coefficient bounds of}\ \mathrm{BM.151}(2)\ (\text{forward pointer})\ —\ \text{the state-side dynamics never leaves the formed locus}\ }
\]

%dsp-level:0
\[
(5)\quad\textbf{the locked face of the operator:}
\]

%dsp-level:0
\[
\qquad F_{D^-}=\bigoplus_{n\ge4}F_n
\]

%dsp-level:0
\[
\qquad [\mathrm{D4.1};\ \mathrm{A8}];
\]

%dsp-level:0
\[
\qquad D_{\mathrm{tick}}(n)=2^{n+1},
\]

%dsp-level:0
\[
\qquad O(n)=\max\bigl(0,\ 2^{n+1}-16\bigr),
\]

%dsp-level:0
\[
\qquad \text{the last admissible rung is}\ 3
\]

%dsp-level:0
\[
\qquad [\mathrm{CT0.5};\ \mathrm{CT0.6};\ \mathrm{CT4.1}];
\]

%dsp-level:0
\[
\text{machine:}\ O=(0,0,0,0,16,48,112,240,496)\ \ \text{at}\ n=(0,1,2,3,4,5,6,7,8);
\]

%dsp-level:0
\[
\qquad O(4)=16=C_{\mathrm{cap}}\ \ \text{exactly one further channel quantum per tick}
\]

%dsp-level:0
\[
\qquad [\mathrm{CT4.1}];
\]

%dsp-level:0
\[
(6)\quad\textbf{status:}
\]

%dsp-level:0
\[
\qquad \textbf{derived:}\ (1)\text{--}(5);
\]

%dsp-level:0
\[
\qquad \text{the census is complete:}\ \neg\text{an aspect outside}\ \mathrm{A8}\text{'s list,}\ \neg\text{a block outside}\ \mathrm{CT0.12}\text{'s tiling,}\ \neg\text{a rung beyond}\ 3\ \text{in the accessible mask}\ [\mathrm{FT.83}(2)\text{--}(3):\ \text{the}\ n\ge4\ \text{no-go}];
\]

%dsp-level:0
\[
\neg\text{a declaration},\qquad\neg\text{a P4 addition},\qquad\neg\text{a new dimensionful scale}\ [\mathrm{P1}],\qquad\neg\text{a measured value}\ [\mathrm{P5}].
\]

\textbf{Req.} A2, A7, A8 (Part 1); C8 (Part 1); D0.3, D0.8, D0.9, D1.2, D1.8, D2.1, D2.3, D2.12, D3.1, D3.5, D3.6, D4.1 (Parts 1–2); CT0.5, CT0.6, CT0.11, CT0.12, CT0.13, CT0.14, CT1.1, CT1.2, L0, L1, L2, L3(1) (Part 1); CT1.12, CT1.20, CT1.21, CT2.4, CT2.5, CT2.6, CT2.7, CT3.8, CT4.1 (Part 2); P1, P2, P5 (Parts 1–2); Def-H.32, FT.9, FT.26(3), FT.83(2), FT.83(3) (Part 3); BM.36(1), BM.52(1), BM.131(3), BM.139(2) (this part); A-BM.10 (§II.4, logged before use); exact integer censuses [math; machine: \texttt{checks/\allowbreak{}verify\_\allowbreak{}H\_\allowbreak{}continuum\_\allowbreak{}limit\_\allowbreak{}era1\_\allowbreak{}closure.py} §A]; cf. BM.151 (forward pointer: the coefficient bounds).

\textbf{Pf.}

%dsp-level:0
\[
(1),
\]

%dsp-level:0
\[
(2):\ \mathrm{A8},
\]

%dsp-level:0
\[
\ \mathrm{CT0.12},
\]

%dsp-level:0
\[
\ \mathrm{CT0.13},
\]

%dsp-level:0
\[
\ \mathrm{CT2.4},
\]

%dsp-level:0
\[
\ \mathrm{CT0.14},
\]

%dsp-level:0
\[
\ \mathrm{BM.36}(1)\text{'s rows re-assembled}\ [\text{machine: the dimension sums}\ 3+3+5+3+1+1=16,\ 16+1=17,\ 17+8=25,\ 1+5+1+3+8=18]
\]

%dsp-level:0
\[
(3):\ \text{each row is the cited tag's own derived map; the composition is read off the shared domains and codomains}\ [\text{machine: the locus census — every row's tag exists and carries the cited clause}];
\]

%dsp-level:0
\[
\qquad \text{the}\ \aleph\text{-domain row is A-BM.10's pointwise declaration}\ [\mathrm{BM.139}(2)]
\]

%dsp-level:0
\[
(4):\ \mathrm{A6},\ \mathrm{D0.3},\ \mathrm{L0}\text{'s rows re-cited};
\]

%dsp-level:0
\[
\qquad \text{the coefficient components:}\ \delta\rho=\tfrac14\sum_ar_aG_a,\ r_a=\mathfrak T(\delta\rho\,G_a)\ [\mathrm{L3}(1);\ \mathrm{CT1.1}];
\]

%dsp-level:0
\[
\qquad \text{the preservation rows:}\ [\mathrm{BM.131}(3);\ \mathrm{L2}]\ [\text{machine: the contraction factor}\ (1-\tfrac\gamma4)\le1\ \text{on}\ \gamma\in(0,1]]
\]

%dsp-level:0
\[
(5):\ O(n)=\max(0,2^{n+1}-16)\ \text{evaluated at}\ n=0..8\ [\text{machine: exact integers}];\qquad(6):\ \text{assembly}.\ \blacksquare
\]

\textbf{Phys.} The force operator is not a list: it is one object with three exact faces. Its algebra face tiles the sixteen dimensions of the pair arena with the six pair aspects and stacks the remaining three on the vertex registers; its mask face climbs the rung ladder six, sixteen, seventeen, twenty-five with the two derived increments; its dynamical face is the composed action that carries a registered graph to weights, to the grouped sum, to the effective action, to the masked resolvent's spectrum, to the source and the field equations, to the smooth target manifold. Beside it runs the state-side action — the Haar draws and the tilt — which touches the same data only through the dictionary, preserves the coefficient bounds, and never leaves the formed locus. The locked face closes the census: beyond rung three the tick demand outruns the channel capacity by an exact power of two, and that excess is the operator's own dark-saturation summand. Nothing is added, nothing declared, nothing measured.

\end{theorem}

\begin{theorem}[{\texorpdfstring{BM.146 — The continuum-limit function: the two variables, the five contributors, and the exact functional form}{BM.146 - The continuum-limit function: the two variables, the five contributors, and the exact functional form}}]\label{tag:BM-146}
\textbf{Thm.}

%dsp-level:0
\[
(1)\quad\textbf{the variables and the notation row:}
\]

%dsp-level:0
\[
\qquad \text{a cell's content is}\ E_{\mathrm{int}}=|R|=n_{\mathcal R}
\]

%dsp-level:0
\[
\qquad [\mathrm{D0.8};\ \mathrm{Def\text{-}BM.5}(4)];
\]

%dsp-level:0
\[
\qquad \text{its ledger flux is}\ S_{\mathfrak L}=E_{\mathrm{int}}\ln2
\]

%dsp-level:0
\[
\qquad [\mathrm{CT1.29};\ \mathrm{C9}];
\]

%dsp-level:0
\[
\text{its chart span enters through the envelope ratio}
\]

%dsp-level:0
\[
\qquad k=\frac{r}{\Delta x_{\mathrm B}(E)}=\frac{r\,E_{\mathrm{int}}\ln2}{\mathfrak h}
\]

%dsp-level:0
\[
\qquad [\mathrm{BM.88}(2)(\mathrm i);\ \mathrm{Def\text{-}BM.4}(1)];
\]

%dsp-level:0
\[
\boxed{\ \textbf{the notation row:}\ \ \mathfrak q\ \ \textbf{remains the dimension bookkeeper}\ [\mathrm{CT2.1}];\ \ \text{the flux variable is the ledger's own}\ S_{\mathfrak L}\ —\ \text{count and flux are one datum up to}\ \ln2\ }
\]

%dsp-level:0
\[
(2)\quad\textbf{the five major contributors (each derived at its own locus):}\qquad\begin{array}{l|l|l}\text{contributor} & \text{the exact row} & \text{the locus}\\\hline\text{the Born width} & \Delta x_{\mathrm B}=\mathfrak h/(E_{\mathrm{int}}\ln2),\quad\Delta x_{\mathrm B}\cdot m_{\mathrm{hid}}=1 & [\mathrm{Def\text{-}BM.4}(1);\ \mathrm{BM.13}(1)]\\\hline\text{the flicker density} & w^{(n)}_{\mathrm{mat}}=\frac1{D_n},\quad\langle E\rangle^{(n)}_{\mathrm{mat}}=\frac{2n}{D_n},\quad D_n=\frac{1-u^n}{1-u} & [\mathrm{BM.137}(1);\ \mathrm{BM.94};\ \mathrm{BM.105}]\\\hline\text{the Janus boundary} & Q=x_+(\hat J_+-P_{\mathcal V})+x_-(\hat J_--P_{\mathcal R}),\quad\operatorname{spec}Q=\{0,-1\},\quad p=1-4\varepsilon_q & [\mathrm{Def\text{-}BM.18};\ \mathrm{BM.120};\ \mathrm{BM.19}(2)]\\\hline\text{the Haar--Poisson lapse} & N(k,E)=\sqrt{1-\frac{G(E)}{4\pi k}}=\sqrt{1-\frac{\mathfrak hE_{\mathrm{int}}}{2r}} & [\mathrm{BM.131};\ \mathrm{BM.132}(2)]\\\hline\text{the locked sector} & O(n)=\max(0,2^{n+1}-16),\quad\langle\rho_{D^-}\rangle_\Gamma(p)=\sum_vO(n_v)\bigl(1-f_{\mathrm{adm}}(n_v,p)\bigr) & [\mathrm{CT4.1};\ \mathrm{CT4.5a}(1)]\end{array}
\]

%dsp-level:0
\[
(3)\quad\textbf{the continuum-limit function (the exact functional form):}
\]

%dsp-level:0
\[
\qquad \text{the chart read of a cell is the tuple}
\]

%dsp-level:0
\[
\boxed{\ \mathcal C(E_{\mathrm{int}},r)\ =\ \Bigl(g_{\mu\nu}(E_{\mathrm{int}},r),\ \ N(k,E),\ \ \mathcal T(k,E),\ \ w^{(n)}_{\mathrm{mat}},\ \ \Gamma(E,r)\Bigr)\ }
\]

%dsp-level:0
\[
g_{\mu\nu}=\eta_{\mu\nu}+\frac{\mathfrak hE_{\mathrm{int}}}{2r}\,\mathbb 1_4+O\Bigl(\frac{\mathfrak h^2E_{\mathrm{int}}^2}{r^2}\Bigr)
\]

%dsp-level:0
\[
\qquad [\mathrm{BM.43}(4);\ \mathrm{BM.132}(2)];
\]

%dsp-level:0
\[
\qquad N=\sqrt{1-\frac{\mathfrak hE_{\mathrm{int}}}{2r}}
\]

%dsp-level:0
\[
\qquad [\mathrm{BM.132}(2)];
\]

%dsp-level:0
\[
\mathcal T=\frac{\mathcal T_0}N=\frac{E_{\mathrm{int}}\,\mathfrak h}{\ln2}\Bigl(1-\frac{\mathfrak hE_{\mathrm{int}}}{2r}\Bigr)^{-1/2}
\]

%dsp-level:0
\[
\qquad [\mathrm{BM.115}(1);\ \mathrm{BM.132}(5)];
\]

%dsp-level:0
\[
\qquad \Gamma=\Gamma_{\mathrm{tr}}\,f\,N\,\mathsf g\,\mathcal N\,\chi
\]

%dsp-level:0
\[
\qquad [\mathrm{BM.116}(1);\ \mathrm{BM.133}(4)];
\]

%dsp-level:0
\[
\textbf{every component is an exact closed form in}\ \bigl(E_{\mathrm{int}},\ r,\ \mathfrak h,\ \ln2,\ \varepsilon_q\bigr)
\]

%dsp-level:0
\[
\qquad [\text{machine: the monomial census}];
\]

%dsp-level:0
\[
(4)\quad\textbf{the formed domain (the exact conditions):}
\]

%dsp-level:0
\[
\qquad \boxed{\ \mathcal C\ \ \text{is formed}\iff E_{\mathrm{int}}\ge2\ \ \wedge\ \ \varkappa\|q\|_{\mathrm{op}}<1\ \ \wedge\ \ E_{\mathrm{int}}\le\frac{2r}{\mathfrak h}\ \ \wedge\ \ n(v)\le3\ \ \text{on the accessible vertices}\ }
\]

%dsp-level:0
\[
(\mathrm C1)\ E_{\mathrm{int}}\ge2:\ \ \text{the sub-hop/materialization floor}\qquad[\mathrm{Census\ FT.28}(1);\ \mathrm{Def\text{-}H.4};\ \mathrm{BM.13}(3)];
\]

%dsp-level:0
\[
(\mathrm C2)\ \varkappa\|q\|_{\mathrm{op}}<1:\ \ \text{the metric's nondegeneracy}\qquad[\mathrm{CT1.12};\ \mathrm{CT2.6}];
\]

%dsp-level:0
\[
(\mathrm C3)\ E_{\mathrm{int}}\le2r/\mathfrak h:\ \ \text{the lapse-formedness condition, in the capacity form}
\]

%dsp-level:0
\[
\qquad [\mathrm{BM.132}(3);\ \mathrm{BM.148}(1);\ \text{machine: the cross-multiplied identity}\ r\ge\mathfrak hE/2];
\]

%dsp-level:0
\[
(\mathrm C4)\ n(v)\le3:\ \ \text{the accessibility bound}\ —\ \text{oversubscribed vertices pass their excess to}\ D^-
\]

%dsp-level:0
\[
\qquad [\mathrm{CT0.11};\ \mathrm{CT4.1};\ \mathrm{D1.11}];
\]

%dsp-level:0
\[
(5)\quad\textbf{the limit}\ \mathfrak h\to0\ \textbf{(the derived sense in which the limit exists):}
\]

%dsp-level:0
\[
\qquad \text{at fixed}\ (E_{\mathrm{int}},r):
\]

%dsp-level:0
\[
\qquad E_{\mathrm{crit}}(r)=\frac{2r}{\mathfrak h}\to\infty
\]

%dsp-level:0
\[
\qquad [\text{machine: the Archimedean row on exact rationals}];
\]

%dsp-level:0
\[
\implies\text{every finite-content cell becomes chart-formed:}
\]

%dsp-level:0
\[
\qquad \forall E_{\mathrm{int}}\ \exists\,\mathfrak h_0:\ \forall\mathfrak h<\mathfrak h_0:\ (\mathrm C3)\ \text{holds}
\]

%dsp-level:0
\[
\qquad [\mathfrak h_0=2r/E_{\mathrm{int}}\ \text{explicit}];
\]

%dsp-level:0
\[
\textbf{the refinement row:}
\]

%dsp-level:0
\[
\qquad \text{a cell of span}\ s\ \text{carries}\ k=\frac{sE_{\mathrm{int}}\ln2}{\mathfrak h}\ \ \text{envelopes, each with}\ \ge2\ \ \text{directed real edges}
\]

%dsp-level:0
\[
\qquad [\mathrm{Def\text{-}H.4};\ \mathrm{Census\ FT.28}(1);\ \text{machine: the monomial}]:
\]

%dsp-level:0
\[
\boxed{\ \#\{\text{real edges per cell}\}\ \ge\ \frac{2sE_{\mathrm{int}}\ln2}{\mathfrak h}\ \longrightarrow\ \infty\ \ \text{as}\ \ \mathfrak h\to0\ }\qquad[\text{the witness of}\ \mathrm{BM.151}(1)];
\]

%dsp-level:0
\[
(6)\quad\textbf{the universality preview:}
\]

%dsp-level:0
\[
\qquad \mathcal C\ \text{carries}\ \neg\text{an aspect label:}\ \ \text{the force content enters through}\ \bigl(T_{\mu\nu}\ \text{channel},\ \text{mask slot},\ \text{coupling profile}\bigr)\ \ \text{alone}
\]

%dsp-level:0
\[
\qquad [\mathrm{CT1.22};\ \mathrm{CT2.4};\ \mathrm{BM.150}(3)\ \text{forward pointer}];
\]

%dsp-level:0
\[
(7)\quad\textbf{status:}
\]

%dsp-level:0
\[
\qquad \textbf{derived:}\ (1)\text{--}(6);
\]

%dsp-level:0
\[
\qquad \neg\text{a declaration},
\]

%dsp-level:0
\[
\qquad \neg\text{a P4 addition},
\]

%dsp-level:0
\[
\qquad \neg\text{a cosmological coordinate}\ [\mathrm{Def\text{-}BM.17}(6)(\mathrm i)],
\]

%dsp-level:0
\[
\qquad \neg\text{a measured value}\ [\mathrm{P5}].
\]

\textbf{Req.} A8 (Part 1); C9, D0.4, D0.8 (Part 1); CT0.11, CT1.29, CT2.1 (Parts 1–2); CT1.12, CT2.6, CT4.1, CT4.5a(1) (Part 2); P1, P5 (Parts 1–2); Def-H.4, Def-H.32 (Part 3); Def-BM.4(1), Def-BM.17(6), Def-BM.18, BM.13(1), BM.13(3), BM.19(2), BM.88(2), BM.94, BM.105, BM.115(1), BM.116(1), BM.120, BM.131, BM.132(2), BM.132(3), BM.132(5), BM.133(4), BM.137(1) (this part); Census FT.28(1) (row-wise); monomial identities, exact rational limit rows [math; machine: \texttt{checks/\allowbreak{}verify\_\allowbreak{}H\_\allowbreak{}continuum\_\allowbreak{}limit\_\allowbreak{}era1\_\allowbreak{}closure.py} §B].

\textbf{Pf.}

%dsp-level:0
\[
(1):\ \text{the count/flux identity}\ S_{\mathfrak L}=n_{\mathcal R}\ln2=E_{\mathrm{int}}\ln2\ [\mathrm{CT1.29};\ \mathrm{D0.8};\ \text{machine: one monomial}];
\]

%dsp-level:0
\[
\qquad \text{the envelope-ratio substitution}\ k=rE\ln2/\mathfrak h\ [\mathrm{Def\text{-}BM.4}(1);\ \text{machine}]
\]

%dsp-level:0
\[
(2):\ \text{the five rows re-cited at their loci}\ [\text{machine: the citation census}];
\]

%dsp-level:0
\[
\qquad (3):\ \text{the substitutions}\ G
\]

%dsp-level:0
\[
= 2\pi E^2\ln2,
\]

%dsp-level:0
\[
\ \Delta x_{\mathrm B}
\]

%dsp-level:0
\[
= \mathfrak h/(E\ln2):\quad\frac{G}{4\pi k}
\]

%dsp-level:0
\[
= \frac{2\pi E^2\ln2}{4\pi}\cdot\frac{\mathfrak h}{rE\ln2}
\]

%dsp-level:0
\[
= \frac{\mathfrak hE}{2r}\ [\text{machine: the rational-function identity}]
\]

%dsp-level:0
\[
(4):\ (\mathrm C1)\text{--}(\mathrm C4)\ \text{re-cited from their own tags};
\]

%dsp-level:0
\[
\qquad (\mathrm C3)\ \text{equivalence:}\ 1-\frac{\mathfrak hE}{2r}
\]

%dsp-level:0
\[
\ge 0
\]

%dsp-level:0
\[
\iff r\ge\frac{\mathfrak hE}2
\]

%dsp-level:0
\[
\iff E\le\frac{2r}{\mathfrak h}\ [\text{machine: cross-multiplication on positive denominators}]
\]

%dsp-level:0
\[
(5):\ \mathfrak h_0=2r/E:\ \ \mathfrak h<\mathfrak h_0\implies E<2r/\mathfrak h\ [\text{machine: rational witnesses}];
\]

%dsp-level:0
\[
\qquad \text{the edge-count monomial}\ 2k=2sE\ln2/\mathfrak h\ [\text{machine}];
\]

%dsp-level:0
\[
\qquad (6),(7):\ \text{assembly}.\ \blacksquare
\]

\textbf{Phys.} The continuum limit stops being a slogan and becomes a function. It takes two ledger variables — the cell's content and its chart span — and returns five exact closed forms: the linearized metric, the lapse, the materialization threshold, the flicker density, and the width. Its domain is four derived conditions, and the third of them is a capacity: a cell of content $E$ is chart-readable out to radius $r$ exactly when $E\le 2r/\mathfrak h$. The limit $\mathfrak h\to0$ is the statement that this capacity outruns every finite content, while the cell's real-edge count grows without bound — the two halves of the continuum condition, both witnessed inside the ledger. The function carries no aspect label: force identity lives in the source, the mask and the coupling profile, never in the limit itself.

\end{theorem}

\begin{theorem}[{\texorpdfstring{BM.147 — The metric–energy–continuum map: the dictionary bijection, the closed form of the source read, and the target}{BM.147 - The metric-energy-continuum map: the dictionary bijection, the closed form of the source read, and the target}}]\label{tag:BM-147}
\textbf{Thm.}

%dsp-level:0
\[
(1)\quad\textbf{the metric is a map (the dictionary):}\qquad\bigl(a_{00},\ q_{ij},\ w_i\bigr)\in\mathcal E_{++}\oplus\mathcal E_{+-}\ \longmapsto\ g_{\mu\nu}=\Omega^2\bar g_{\mu\nu},\qquad\Omega=e^{\alpha_\Omega a_{00}},\qquad\bar g=\begin{pmatrix}-1&\varkappa N_j\\ \varkappa N_i&\delta_{ij}+\varkappa q_{ij}\end{pmatrix},\qquad N_i=\beta w_i\cdot\Omega^2\ [\mathrm{CT1.12};\ \mathrm{BM.34}(1)];
\]

%dsp-level:0
\[
7+3=10=\tfrac{4\cdot5}2:
\]

%dsp-level:0
\[
\qquad \boxed{\ \text{the dictionary is a bijection between the ten sector coefficients and the ten metric components}\ }
\]

%dsp-level:0
\[
\qquad [\mathrm{CT1.12};\ \text{machine: the dimension count}];
\]

%dsp-level:0
\[
\operatorname{sig}(g)=(-,+,+,+)\iff\varkappa\|q\|_{\mathrm{op}}<1\ \ \text{on the perturbative regime}
\]

%dsp-level:0
\[
\qquad [\mathrm{CT1.12}];
\]

%dsp-level:0
\[
\qquad (a_{00},q,w)=0\longmapsto\text{the flat metric}
\]

%dsp-level:0
\[
\qquad [\mathrm{BM.34}(2)];
\]

%dsp-level:0
\[
(2)\quad\textbf{the energy determines the map (the closed form on the dust source):}
\]

%dsp-level:0
\[
\qquad T_{\mu\nu}=\mathrm{diag}(\rho,0,0,0),
\]

%dsp-level:0
\[
\qquad \rho=\frac{m}{V_{\mathrm{eff}}},
\]

%dsp-level:0
\[
\qquad m=\frac{E_{\mathrm{int}}\ln2}{\mathfrak h}
\]

%dsp-level:0
\[
\qquad [\mathrm{BM.9}(4);\ \mathrm{BM.14};\ \mathrm{FT.44}(4)];
\]

%dsp-level:0
\[
G_{\mu\nu}+\Lambda g_{\mu\nu}=\kappa^2T_{\mu\nu}
\]

%dsp-level:0
\[
\qquad [\mathrm{CT1.23}];
\]

%dsp-level:0
\[
\qquad \nabla^2\Phi=4\pi G_N\rho,
\]

%dsp-level:0
\[
\qquad \Phi=-\frac{G_Nm}r,
\]

%dsp-level:0
\[
\qquad G_N=\frac{\mathfrak h^2}{4\ln2},
\]

%dsp-level:0
\[
\qquad G_Nm=\frac{\mathfrak hE_{\mathrm{int}}}4
\]

%dsp-level:0
\[
\qquad [\mathrm{BM.43}(4);\ \mathrm{P3};\ \mathrm{CT2.9}];
\]

%dsp-level:0
\[
\boxed{\ g_{00}=-\Bigl(1-\frac{\mathfrak hE_{\mathrm{int}}}{2r}\Bigr),\qquad g_{ij}=\Bigl(1+\frac{\mathfrak hE_{\mathrm{int}}}{2r}\Bigr)\delta_{ij},\qquad g_{0i}=0\ }
\]

%dsp-level:0
\[
\qquad [\text{machine: the substitution}\ r_s=\mathfrak hE_{\mathrm{int}}/2\ \text{in}\ \mathrm{BM.43}(4)\text{'s weak-field rows}];
\]

%dsp-level:0
\[
\textbf{the pure-trace row:}
\]

%dsp-level:0
\[
\qquad g_{\mu\nu}=\eta_{\mu\nu}+\frac{\mathfrak hE_{\mathrm{int}}}{2r}\,\mathbb 1_4\ +\ O\Bigl(\frac{\mathfrak h^2E^2}{r^2}\Bigr):
\]

%dsp-level:0
\[
\qquad \boxed{\ \textbf{the linearized self-field of the dust cell is a pure-trace (conformal) perturbation}\ }
\]

%dsp-level:0
\[
\text{its whole content sits in the conformal mode:}
\]

%dsp-level:0
\[
\qquad \Omega^2=1-\frac{\mathfrak hE_{\mathrm{int}}}{2r},
\]

%dsp-level:0
\[
\qquad a_{00}=\frac1{2\alpha_\Omega}\ln\Bigl(1-\frac{\mathfrak hE_{\mathrm{int}}}{2r}\Bigr),
\]

%dsp-level:0
\[
\qquad \hat q=0,
\]

%dsp-level:0
\[
\qquad w=0
\]

%dsp-level:0
\[
\qquad [\text{machine: the exact solve of}\ \gamma_{ij}=\Omega^2(\delta_{ij}+\varkappa q_{ij})\ \text{against}\ g_{ij}];
\]

%dsp-level:0
\[
\textbf{the energy–lapse identity:}
\]

%dsp-level:0
\[
\qquad N^2=\Omega^2=1-\frac{\mathfrak hE_{\mathrm{int}}}{2r}=1-\frac{G(E)}{4\pi k}
\]

%dsp-level:0
\[
\qquad [\mathrm{BM.132}(2);\ \text{machine: the rational-function identity of}\ (3)\text{'s rows}];
\]

%dsp-level:0
\[
(3)\quad\textbf{the inverse read (the energy from the metric):}
\]

%dsp-level:0
\[
\qquad R=\kappa^2(\rho-3p)+4\Lambda,
\]

%dsp-level:0
\[
\qquad p=0,\ \Lambda=0:
\]

%dsp-level:0
\[
\qquad RV=\kappa^2m
\]

%dsp-level:0
\[
\qquad [\mathrm{BM.15}(2);\ \mathrm{BM.20}(1)]:
\]

%dsp-level:0
\[
\boxed{\ m=\frac1{\kappa^2}\int_\Sigma R\,dV\ }
\]

%dsp-level:0
\[
\qquad [\mathrm{BM.20}(1)];
\]

%dsp-level:0
\[
\qquad \kappa^2\rho=\frac{2\pi\,\mathfrak h\,E_{\mathrm{int}}}{V},
\]

%dsp-level:0
\[
\qquad (\Delta x_{\mathrm B})^2\kappa^2\rho=2\pi E_{\mathrm{int}}^2\ln2=G
\]

%dsp-level:0
\[
\qquad [\mathrm{BM.9}(5);\ \mathrm{BM.16}(1);\ \text{machine: monomials}];
\]

%dsp-level:0
\[
\textbf{the map is invertible on-shell at the leading reads:}
\]

%dsp-level:0
\[
\qquad E_{\mathrm{int}}=\frac{\kappa^2\rho V}{2\pi\mathfrak h}=\sqrt{\frac{G}{2\pi\ln2}}
\]

%dsp-level:0
\[
\qquad [\text{machine: the two exact inverses}];
\]

%dsp-level:0
\[
(4)\quad\textbf{the continuum is the target:}
\]

%dsp-level:0
\[
\qquad M=\Gamma_{\mathcal R}/\!\!\sim,
\]

%dsp-level:0
\[
\qquad \chi([v])=(\tau_n,h_1,h_2,h_3),
\]

%dsp-level:0
\[
\qquad \text{transitions smooth with}\ |\det J|=1
\]

%dsp-level:0
\[
\qquad [\mathrm{CT2.6};\ \mathrm{D1.8}];
\]

%dsp-level:0
\[
\textbf{the target-space rows:}
\]

%dsp-level:0
\[
\qquad \text{the image of the dictionary is the Lorentzian metrics of signature}\ (-,+,+,+)\ \text{on the perturbative regime}\ [\mathrm{CT1.12}];
\]

%dsp-level:0
\[
\qquad \text{the dispersion target}\ E=N\sqrt{m^2+p^2}-N^ip_i\ [\mathrm{BM.34}(4)];
\]

%dsp-level:0
\[
\textbf{the map's full statement:}
\]

%dsp-level:0
\[
\qquad \boxed{\ \text{metric}\ =\ \text{dictionary}\circ\ \text{field-equation solve}\circ\ \text{source read}\ (E_{\mathrm{int}},r)\ \longmapsto\ \text{continuum}\ }
\]

%dsp-level:0
\[
\qquad [(1)+(2)+(4);\ \text{the nonlinear joint solve remains}\ \neg\text{constructed}\ [\mathrm{BM.43}(5)\ \text{re-cited}]];
\]

%dsp-level:0
\[
(5)\quad\textbf{status:}
\]

%dsp-level:0
\[
\qquad \textbf{derived:}\ (1)\text{--}(4);
\]

%dsp-level:0
\[
\qquad \neg\text{a declaration},
\]

%dsp-level:0
\[
\qquad \neg\text{a P4 addition},
\]

%dsp-level:0
\[
\qquad r_s\ \text{quarantined as}\ \mathrm{P3}\text{'s named combination},
\]

%dsp-level:0
\[
\qquad \neg\text{a measured value}\ [\mathrm{P5}].
\]

\textbf{Req.} A8 (Part 1); D1.8, D2.1 (Part 1); CT1.12, CT1.23, CT2.6, CT2.9 (Part 2); P3, P5 (Parts 1–2); FT.44(4) (Census); Def-BM.4, BM.9(4), BM.9(5), BM.14, BM.15(2), BM.16(1), BM.20(1), BM.34(1), BM.34(2), BM.34(4), BM.43(4), BM.43(5), BM.115(1), BM.132(2) (this part); the exact weak-field solve, monomial and rational-function identities [math; machine: \texttt{checks/\allowbreak{}verify\_\allowbreak{}H\_\allowbreak{}continuum\_\allowbreak{}limit\_\allowbreak{}era1\_\allowbreak{}closure.py} §C].

\textbf{Pf.}

%dsp-level:0
\[
(1):\ \mathrm{CT1.12}\text{'s bijection row and signature criterion re-cited}\ [\text{machine: the dimension count}];
\]

%dsp-level:0
\[
\qquad (2):\ \Phi=-G_Nm/r\ \text{with}\ G_Nm=\frac{\mathfrak h^2}{4\ln2}\cdot\frac{E\ln2}{\mathfrak h}=\frac{\mathfrak hE}4\ [\mathrm{P3};\ \mathrm{FT.44}(4);\ \text{machine}];
\]

%dsp-level:0
\[
\qquad g_{00}=-(1+2\Phi)=-\bigl(1-\frac{\mathfrak hE}{2r}\bigr),\quad g_{ij}=(1-2\Phi)\delta_{ij}=\bigl(1+\frac{\mathfrak hE}{2r}\bigr)\delta_{ij}\ [\mathrm{BM.43}(4);\ \text{machine}]
\]

%dsp-level:0
\[
\text{the pure-trace row:}\ h_{00}=+\frac{\mathfrak hE}{2r}=h_{11}=h_{22}=h_{33}\ [\text{machine: the component census}];
\]

%dsp-level:0
\[
\qquad \text{the conformal solve:}\ \Omega^2\bigl(\delta_{ij}+\varkappa q_{ij}\bigr)
\]

%dsp-level:0
\[
= \bigl(1+\frac{\mathfrak hE}{2r}\bigr)\delta_{ij}\ \wedge\ \Omega^2
\]

%dsp-level:0
\[
= 1-\frac{\mathfrak hE}{2r}:\quad\varkappa q_{ij}
\]

%dsp-level:0
\[
= \Bigl(\frac{1+\mathfrak hE/(2r)}{1-\mathfrak hE/(2r)}-1\Bigr)\delta_{ij}
\]

%dsp-level:0
\[
= \frac{\mathfrak hE/r}{1-\mathfrak hE/(2r)}\,
\]

%dsp-level:0
\[
\delta_{ij}\ [\text{machine: the rational-function identity}]
\]

%dsp-level:0
\[
(3):\ \mathrm{BM.15}(2),\ \mathrm{BM.20}(1),\ \mathrm{BM.9}(5),\ \mathrm{BM.16}(1)\text{'s monomials re-verified}\ [\text{machine}];
\]

%dsp-level:0
\[
\qquad \text{the two inverses:}\ \frac{\kappa^2}{2\pi\mathfrak h}\cdot\frac{2\pi\mathfrak hE}{V}\cdot V=E;
\]

%dsp-level:0
\[
\qquad \Bigl(\sqrt{\frac{G}{2\pi\ln2}}\Bigr)^2\cdot2\pi\ln2=G\ [\text{machine}]
\]

%dsp-level:0
\[
(4),(5):\ \text{assembly and the scope re-citations}.\ \blacksquare
\]

\textbf{Phys.} The three names in the directive are one chain. The metric is the dictionary: ten sector coefficients of the ledger's pair states, bijectively repackaged as ten metric components. The energy determines the map: run the composite's own dust source through the field equations and the linearized solve closes in one monomial — the perturbation is pure trace, entirely conformal, the lapse and the conformal mode carrying the same exact function of content over radius. The map inverts on-shell: the curvature flux reads the content back. The continuum is the target: the quotient manifold with its tick-and-hop charts, whose metric image is exactly the Lorentzian signatures on the perturbative regime. The full non-linear solve stays unconstructed, and nothing measured enters anywhere.

\end{theorem}

\begin{theorem}[{\texorpdfstring{BM.148 — The breaking threshold: the exact capacity of a chart cell}{BM.148 - The breaking threshold: the exact capacity of a chart cell}}]\label{tag:BM-148}
\textbf{Thm.}

%dsp-level:0
\[
(1)\quad\textbf{the threshold (three exact forms of one condition):}
\]

%dsp-level:0
\[
\qquad \text{the chart read of a cell is formed}\iff r\ \ge\ r_s(E)=\frac{\mathfrak hE_{\mathrm{int}}}2
\]

%dsp-level:0
\[
\qquad [\mathrm{BM.132}(3);\ \mathrm{BM.147}(2)]:
\]

%dsp-level:0
\[
\boxed{\ E_{\mathrm{crit}}(r)=\frac{2r}{\mathfrak h}\ }\qquad\boxed{\ S_{\mathrm{crit}}(r)=\frac{2r\ln2}{\mathfrak h}\ }\qquad\boxed{\ k_{\mathrm{crit}}(E)=\frac{E_{\mathrm{int}}^2\ln2}2=\frac{G(E)}{4\pi}\ }
\]

%dsp-level:0
\[
\textbf{the flux form:}
\]

%dsp-level:0
\[
\qquad \text{the map breaks}\iff S_{\mathfrak L}(\text{cell})>S_{\mathrm{crit}}(r)
\]

%dsp-level:0
\[
\qquad [S_{\mathfrak L}=E_{\mathrm{int}}\ln2;\ \mathrm{CT1.29};\ \text{machine: the cross-multiplied identity}];
\]

%dsp-level:0
\[
(2)\quad\textbf{the dimension census:}
\]

%dsp-level:0
\[
\qquad [E_{\mathrm{crit}}]=\frac{\mathfrak q^{-1/2}}{\mathfrak q^{-1/2}}=\mathfrak q^0\ (\text{a count});
\]

%dsp-level:0
\[
\qquad [S_{\mathrm{crit}}]=\mathfrak q^0\ (\text{nats});
\]

%dsp-level:0
\[
\qquad [r_{\mathrm{crit}}]=\mathfrak q^{-1/2}\ (\text{a length})
\]

%dsp-level:0
\[
\qquad [\mathrm{CT2.1};\ \text{machine: the exponent arithmetic}];
\]

%dsp-level:0
\[
(3)\quad\textbf{the integer row:}
\]

%dsp-level:0
\[
\qquad E_{\mathrm{int}}\in\mathbb Z\implies\text{the formed contents at radius}\ r\ \text{are}\ \{2,\dots,\lfloor2r/\mathfrak h\rfloor\}
\]

%dsp-level:0
\[
\qquad [\mathrm{BM.13}(3):\ E\ge2;\ \text{machine}];
\]

%dsp-level:0
\[
\textbf{the one-hop capacity:}
\]

%dsp-level:0
\[
\qquad r=\mathfrak h:\quad E_{\mathrm{crit}}(\mathfrak h)=2:
\]

%dsp-level:0
\[
\qquad \boxed{\ \text{a chart cell of one hop carries exactly the digon floor and nothing heavier}\ }
\]

%dsp-level:0
\[
\qquad [\mathrm{Census\ FT.28};\ \text{machine}];
\]

%dsp-level:0
\[
\textbf{the computed cells:}
\]

%dsp-level:0
\[
\qquad r_s=\Bigl(1,\ 3,\ 6,\ 12,\ 30\Bigr)\,\mathfrak h
\]

%dsp-level:0
\[
\qquad \text{at}\ E_{\mathrm{int}}\in\{2,6,12,24,60\}
\]

%dsp-level:0
\[
\qquad [\mathrm{BM.43}(1);\ \text{machine: the integer census}];
\]

%dsp-level:0
\[
(4)\quad\textbf{the envelope corollary (the second route to}\ \mathrm{BM.132}(3)\text{'s row):}
\]

%dsp-level:0
\[
\qquad r=\Delta x_{\mathrm B}(E)=\frac{\mathfrak h}{E\ln2}:
\]

%dsp-level:0
\[
\qquad E\le E_{\mathrm{crit}}\bigl(\Delta x_{\mathrm B}(E)\bigr)=\frac2{E\ln2}\iff E^2\ln2\le2:
\]

%dsp-level:0
\[
E_{\mathrm{int}}\ge2\implies E^2\ln2\ge4\ln2=\ln16>2
\]

%dsp-level:0
\[
\qquad \Bigl[\ln16>2\iff16>e^2:\quad e^2<\Bigl(\frac{876809}{322560}\Bigr)^2<8<16;\ \text{machine: the exact enclosure}\Bigr];
\]

%dsp-level:0
\[
\boxed{\ \textbf{no materialized cell is chart-formed at its own envelope:}\ \ \text{the capacity}\ \frac2{E\ln2}\ \text{is below every admissible}\ E\ }
\]

%dsp-level:0
\[
(5)\quad\textbf{the monotonicity and the limit:}
\]

%dsp-level:0
\[
\qquad E_{\mathrm{crit}}(r)\ \text{strictly increasing in}\ r;
\]

%dsp-level:0
\[
\qquad S_{\mathrm{crit}}(r)\ \text{strictly increasing};
\]

%dsp-level:0
\[
\qquad r_{\mathrm{crit}}(E)\ \text{strictly increasing in}\ E;
\]

%dsp-level:0
\[
\qquad k_{\mathrm{crit}}(E)\ \text{strictly increasing and quadratic}
\]

%dsp-level:0
\[
\qquad [\text{machine: the sign census on exact rationals}];
\]

%dsp-level:0
\[
\lim_{\mathfrak h\to0}E_{\mathrm{crit}}(r)=\infty\ \ \text{at fixed}\ r
\]

%dsp-level:0
\[
\qquad [\mathrm{BM.146}(5);\ \text{machine}];
\]

%dsp-level:0
\[
\qquad \boxed{\ \textbf{the breaking locus recedes without bound in the continuum limit}\ }
\]

%dsp-level:0
\[
(6)\quad\textbf{the quarantine rows:}
\]

%dsp-level:0
\[
\qquad r_s\ \ \text{enters only as}\ \mathrm{P3}\text{'s named combination:}\ \neg\text{a collapse reading,}\ \neg\text{a horizon identification,}\ \neg\text{black-hole physics}\ [\mathrm{P3};\ \mathrm{BM.132}(6);\ \mathrm{BM.115}(4)\ \text{discipline}];
\]

%dsp-level:0
\[
\text{the breaking locus is a domain boundary of the chart reads}\ —\ \text{the ledger side of the same boundary is}\ \mathrm{BM.149};
\]

%dsp-level:0
\[
(7)\quad\textbf{status:}\qquad\textbf{derived:}\ (1)\text{--}(6);\qquad\neg\text{a declaration},\qquad\neg\text{a P4 addition},\qquad\neg\text{a numerical}\ \mathfrak h\ [\mathrm{P1};\ \mathrm{P5}].
\]

\textbf{Req.} D0.8 (Part 1); CT1.29, CT2.1 (Parts 1–2); P1, P3, P5 (Parts 1–2); Def-BM.4, BM.13(3), BM.43(1), BM.88(3), BM.115(4), BM.132(3), BM.132(6), BM.146(5), BM.147(2) (this part); Census FT.28 (row-wise); the exact rational enclosure of $e^2$, cross-multiplied identities, integer censuses [math; machine: \texttt{checks/\allowbreak{}verify\_\allowbreak{}H\_\allowbreak{}continuum\_\allowbreak{}limit\_\allowbreak{}era1\_\allowbreak{}closure.py} §D].

\textbf{Pf.}

%dsp-level:0
\[
(1):\ r\ge\frac{\mathfrak hE}2
\]

%dsp-level:0
\[
\iff E\le\frac{2r}{\mathfrak h}
\]

%dsp-level:0
\[
\iff E\ln2\le\frac{2r\ln2}{\mathfrak h}\ [\text{machine: cross-multiplication on positive denominators}];
\]

%dsp-level:0
\[
\qquad k_{\mathrm{crit}}:\ \frac{G}{4\pi}=\frac{2\pi E^2\ln2}{4\pi}=\frac{E^2\ln2}2\ [\mathrm{BM.88}(3);\ \text{machine}]
\]

%dsp-level:0
\[
(2):\ [\mathfrak h]=\mathfrak q^{-1/2}\ [\mathrm{P1};\ \mathrm{CT2.1}]:\ \text{the three exponent rows}\ [\text{machine}];
\]

%dsp-level:0
\[
\qquad (3):\ r_s=\mathfrak hE/2\ \text{at}\ E\in\{2,6,12,24,60\}:\ (1,3,6,12,30)\,\mathfrak h\ [\text{machine: integers}]
\]

%dsp-level:0
\[
(4):\ E^2\ln2\le2\ \text{has no integer solution}\ \ge2:
\]

%dsp-level:0
\[
\quad 4\ln2=\ln16>2\iff16>e^2\ [\mathrm{BM.132}(3)\text{'s enclosure re-run};\ \text{machine}];
\]

%dsp-level:0
\[
\qquad (5):\ \text{the derivative signs on exact rationals};
\]

%dsp-level:0
\[
\qquad (6),(7):\ \text{assembly}.\ \blacksquare
\]

\textbf{Phys.} The continuum has a capacity, and the capacity is exact: a chart cell of radius $r$ carries at most $2r/\mathfrak h$ of internal count, equivalently $2r\ln2/\mathfrak h$ nats of ledger flux, before its own metric read leaves the formed domain. At one hop of radius the capacity is precisely the digon floor — the lightest materialized cell there is also the heaviest admissible one. No materialized cell fits inside its own envelope: the capacity at the Born width is below every admissible content, a second, purely counting route to the envelope row. And in the continuum limit the breaking locus recedes without bound, which is the exact sense in which the limit exists. The Schwarzschild combination stays quarantined; the boundary is a domain edge of the chart reads, and its ledger face is the next theorem.

\end{theorem}

\begin{theorem}[{\texorpdfstring{BM.149 — The locked sector at the breaking locus: the two faces of one capacity, the absorption, the membrane, and the complement theorem}{BM.149 - The locked sector at the breaking locus: the two faces of one capacity, the absorption, the membrane, and the complement theorem}}]\label{tag:BM-149}
\textbf{Thm.}

%dsp-level:0
\[
(1)\quad\textbf{the two faces of one capacity:}\qquad\begin{array}{l|l|l}\text{the face} & \text{the capacity} & \text{the locus}\\\hline\text{chart-side} & E_{\mathrm{int}}\le2r/\mathfrak h:\ \ \text{the metric resolution of a cell} & [\mathrm{BM.148}(1)]\\\hline\text{ledger-side} & k_v\le3\ \ \text{per accessible vertex};\quad D_{\mathrm{tick}}(n)\le C_{\mathrm{cap}}=16\ \ \text{per tick} & [\mathrm{CT0.11};\ \mathrm{CT0.5};\ \mathrm{CT0.6};\ \mathrm{D1.11}]\end{array}
\]

%dsp-level:0
\[
\textbf{the excess rows:}
\]

%dsp-level:0
\[
\qquad \text{chart-side excess: the read is}\ \neg\text{formed}\ [\mathrm{BM.132}(3)];
\]

%dsp-level:0
\[
\qquad \text{ledger-side excess:}\ O(n)=\max(0,2^{n+1}-16)>0\ \ \text{at}\ n\ge4\ \ \text{falls to}\ D^-\ [\mathrm{CT4.1};\ \mathrm{D4.1}];
\]

%dsp-level:0
\[
(2)\quad\textbf{the absorption (exact rows):}
\]

%dsp-level:0
\[
\qquad \text{rate}\ \propto O(n)
\]

%dsp-level:0
\[
\qquad [\mathrm{D4.2}];
\]

%dsp-level:0
\[
\qquad O(4)=16=C_{\mathrm{cap}}:\ \ \text{exactly one further channel quantum per tick}
\]

%dsp-level:0
\[
\qquad [\mathrm{CT4.1}];
\]

%dsp-level:0
\[
\langle\rho_{D^-}\rangle_\Gamma(p)=\sum_vO(n_v)\,\bigl(1-f_{\mathrm{adm}}(n_v,p)\bigr),
\]

%dsp-level:0
\[
\qquad f_{\mathrm{adm}}(n,p)=\sum_{k\le3}\binom nkp^k(1-p)^{n-k}
\]

%dsp-level:0
\[
\qquad [\mathrm{CT4.5a}(1):\ \text{an exact polynomial in}\ p\ \text{with integer coefficients}];
\]

%dsp-level:0
\[
\textbf{the anchors re-derived:}
\]

%dsp-level:0
\[
\qquad \langle\rho_{D^-}\rangle_\Gamma(1)=\bigl(448,\ 896,\ 2976,\ 2240,\ 24384\bigr)
\]

%dsp-level:0
\[
\qquad \text{on the five Platonic skeletons}
\]

%dsp-level:0
\[
\qquad [\mathrm{CT4.5a}(2);\ \text{machine:}\ 4\,O(6),\ 8\,O(6),\ 6\,O(8),\ 20\,O(6),\ 12\,O(10)\ \text{with}\ O(m)=2^{m+1}-16];
\]

%dsp-level:0
\[
(3)\quad\textbf{the one-way membrane (re-derived):}
\]

%dsp-level:0
\[
\qquad \forall n\le3:\quad\mathcal P_M^{(n)}F_{D^-}\mathcal P_M^{(n)}=0;
\]

%dsp-level:0
\[
\qquad \neg\exists\ \text{a pole of}\ Z(\mathcal R,n)\ \ \text{carrying locked content}
\]

%dsp-level:0
\[
\qquad [\mathrm{CT4.2}];
\]

%dsp-level:0
\[
\text{locked content: physical norm}\ 0\ \ \text{at accessible rungs},
\]

%dsp-level:0
\[
\qquad w(C)\ne0\ \ \text{in the unmasked A8 sum},
\]

%dsp-level:0
\[
\qquad \neg\exists\ \text{a deleted entry}
\]

%dsp-level:0
\[
\qquad [\mathrm{CT4.3};\ \mathrm{CT2.8}];
\]

%dsp-level:0
\[
(4)\quad\textbf{the maximal locked entropy:}
\]

%dsp-level:0
\[
\qquad \text{the entropy triangle's third vertex:}
\]

%dsp-level:0
\[
\qquad D^-:\ \bigl(S_{\mathrm{raw}},S_{\mathfrak L}\bigr)=\bigl(\min_{\mathrm{acc}},\ \max\wedge\neg\text{accessible}\bigr)
\]

%dsp-level:0
\[
\qquad [\mathrm{CT4.7}];
\]

%dsp-level:0
\[
S_{\mathfrak L}=n_{\mathcal R}\ln2\ \ \text{is retained in full on absorption}
\]

%dsp-level:0
\[
\qquad [\mathrm{CT4.3}:\ \text{absorbed}\subseteq\mathfrak L;\ \mathrm{CT1.29}];
\]

%dsp-level:0
\[
\qquad \text{the accessible per-vertex content is bounded by}\ n(v)\ [\mathrm{D1.11}];
\]

%dsp-level:0
\[
(5)\quad\textbf{the dynamic emission:}
\]

%dsp-level:0
\[
\qquad \text{at the dynamic level the move system carries the emission and dissolution moves}\ N^\pm(v,S)\ \ \text{with formal weights}\ y_d,\ z_dx_-^{r(v)}
\]

%dsp-level:0
\[
\qquad [\mathrm{Def\text{-}H.32}(7);\ \mathrm{FT.108}];
\]

%dsp-level:0
\[
\textbf{the scope row:}
\]

%dsp-level:0
\[
\qquad \text{emission is a ledger move,}\ \neg\text{an accessible-mask return:}\ \ \text{re-entry to the accessible sector runs only through the A4 crossing and the register's moves}\ \hat J_\pm
\]

%dsp-level:0
\[
\qquad [\mathrm{CT4.2};\ \mathrm{Def\text{-}BM.18};\ \mathrm{BM.120};\ \text{machine: the citation census}];
\]

%dsp-level:0
\[
(6)\quad\textbf{the two scales:}
\]

%dsp-level:0
\[
\qquad \text{the locked sector gravitates as pressureless dust:}\ \ T^{(D^-)}_{\mu\nu}=\rho_{D^-}u_\mu u_\nu,
\]

%dsp-level:0
\[
\quad w_{D^-}=0,\quad\rho_{D^-}\ge0\ \bigl(=0\iff\neg\exists v:\ n(v)\ge4\bigr)
\]

%dsp-level:0
\[
\qquad [\mathrm{CT4.5}];
\]

%dsp-level:0
\[
\text{at the local scale the same sector is the membrane of}\ (3)
\]

%dsp-level:0
\[
\qquad [\mathrm{CT4.2}];
\]

%dsp-level:0
\[
\qquad \text{the junctions at}\ n\ge4\ \ \text{carry its lock content: retained,}\ \neg\text{accessible}
\]

%dsp-level:0
\[
\qquad [\mathrm{BM.2}(4);\ \mathrm{D1.11};\ \mathrm{Census\ FT.35}:\ \text{the first}\ D^-\text{-carrying floor}];
\]

%dsp-level:0
\[
(7)\quad\textbf{the complement theorem (the derived row of this tag):}
\]

%dsp-level:0
\[
\qquad \text{for a cell whose content exceeds the chart capacity}\ \bigl(E_{\mathrm{int}}>2r/\mathfrak h\bigr):
\]

%dsp-level:0
\[
(\mathrm i)\quad\text{the chart read is}\ \neg\text{formed}\ [\mathrm{BM.148}(1);\ \mathrm{BM.132}(3)];
\]

%dsp-level:0
\[
\qquad (\mathrm{ii})\quad\text{the ledger read stays exact and non-decreasing:}\ S_{\mathfrak L}=E_{\mathrm{int}}\ln2,
\]

%dsp-level:0
\[
\qquad |\mathfrak L(\tau_{n+1})|\ge|\mathfrak L(\tau_n)|
\]

%dsp-level:0
\[
\qquad [\mathrm{CT1.29};\ \mathrm{A6}];
\]

%dsp-level:0
\[
(\mathrm{iii})\quad\text{the accessible content is membrane-governed:}\ \neg\exists\ \text{an accessible mask returning locked content}\ [\mathrm{CT4.2};\ \mathrm{CT4.3}];
\]

%dsp-level:0
\[
\boxed{\ \textbf{what the metric cannot read, the ledger still retains:}\ \ \text{the breaking of the continuum map is}\ \neg\text{a loss of content}\ —\ \text{the content passes to the locked species, chart-side unreadable and ledger-side exact}\ }
\]

%dsp-level:0
\[
(8)\quad\textbf{status:}
\]

%dsp-level:0
\[
\qquad \textbf{derived:}\ (1)\text{--}(7);
\]

%dsp-level:0
\[
\qquad \neg\text{a declaration},
\]

%dsp-level:0
\[
\qquad \neg\text{a P4 addition},
\]

%dsp-level:0
\[
\qquad \neg\text{black-hole physics}\ [\mathrm{P3}\ \text{discipline};\ \mathrm{BM.148}(6)],
\]

%dsp-level:0
\[
\qquad \neg\text{a measured value}\ [\mathrm{P5}].
\]

\textbf{Req.} A6, A8 (Part 1); D0.9, D1.11, D4.1, D4.2 (Parts 1–2); CT0.5, CT0.6, CT0.11, CT1.29 (Part 1); CT2.8, CT4.1, CT4.2, CT4.3, CT4.5, CT4.5a, CT4.7 (Part 2); P3, P5 (Parts 1–2); Def-H.32(7), FT.108 (Part 3); Def-BM.18, BM.2(4), BM.120, BM.132(3), BM.148(1), BM.148(6) (this part); Census FT.35 (row-wise); exact integer and polynomial censuses [math; machine: \texttt{checks/\allowbreak{}verify\_\allowbreak{}H\_\allowbreak{}continuum\_\allowbreak{}limit\_\allowbreak{}era1\_\allowbreak{}closure.py} §E].

\textbf{Pf.}

%dsp-level:0
\[
(1):\ \text{the two capacity rows re-cited from their own tags};
\]

%dsp-level:0
\[
\qquad (2):\ O(n)=\max(0,2^{n+1}-16)\ \text{at}\ n=4..10\ [\text{machine: exact integers}];
\]

%dsp-level:0
\[
\qquad \text{the anchors:}\ O(6)=112,\ O(8)=496,\ O(10)=2032:
\]

%dsp-level:0
\[
\quad 4\cdot112=448,\ 8\cdot112=896,\ 6\cdot496=2976,\ 20\cdot112=2240,\ 12\cdot2032=24384\ [\text{machine}]
\]

%dsp-level:0
\[
(3),(4):\ \mathrm{CT4.2},\ \mathrm{CT4.3},\ \mathrm{CT4.7}\text{'s rows re-derived at their loci}\ [\text{machine: the citation pins}];
\]

%dsp-level:0
\[
\qquad (5):\ \mathrm{Def\text{-}H.32}(7)\text{'s moves with}\ \mathrm{CT4.2}\text{'s mask row}
\]

%dsp-level:0
\[
(6):\ \mathrm{CT4.5}\text{'s dust rows re-cited};
\]

%dsp-level:0
\[
\qquad (7):\ (\mathrm i)\text{--}(\mathrm{iii})\ \text{conjoined on the common domain}\ E>2r/\mathfrak h\ [\text{machine: the domain identity}];
\]

%dsp-level:0
\[
\qquad (8):\ \text{assembly}.\ \blacksquare
\]

\textbf{Phys.} The breaking of the continuum map is not a hole in the ledger — it is the locked sector's entrance. The same capacity theorem has two faces: the chart can resolve a cell only out to its content in hop units, and a vertex can expose only three incident quanta to an accessible mask; the excess on the first face is unreadable geometry, on the second it is absorption into the dark-saturation summand at an exact power-of-two rate, with the Platonic anchor values recomputed to the integer. The membrane is one-way: no accessible mask returns the content, no accessible pole carries it, yet nothing is deleted — the record survives with its full ledger entropy, at the maximal-locked vertex of the entropy triangle. Emission exists as a ledger move, but re-entry runs only through the A4 crossing. And the complement theorem states the whole design in one row: what the metric cannot read, the ledger still retains.

\end{theorem}

\begin{theorem}[{\texorpdfstring{BM.150 — The universality of the continuum limit: one operator, one arena, one action, three fingerprints}{BM.150 - The universality of the continuum limit: one operator, one arena, one action, three fingerprints}}]\label{tag:BM-150}
\textbf{Thm.}

%dsp-level:0
\[
(1)\quad\textbf{one operator:}
\]

%dsp-level:0
\[
\qquad \mathbf F\ \ \text{always applies in full, and its blocks are exhausted:}
\]

%dsp-level:0
\[
\qquad 3+3+5+3+1+1=16=C_{\mathrm{cap}}\ \ \text{on the pair arena};
\]

%dsp-level:0
\[
\qquad \mathrm{EM},\ \mathrm{strong},\ D^-\ \ \text{on the stack registers}
\]

%dsp-level:0
\[
\qquad [\mathrm{A8};\ \mathrm{BM.52}(1);\ \mathrm{CT0.12}];
\]

%dsp-level:0
\[
\neg\exists\ \text{an aspect outside the census}\ [\mathrm{CT0.13};\ \mathrm{FT.83}(2)\text{--}(3)];
\]

%dsp-level:0
\[
\qquad \neg\exists\ \text{a further block inside}\ \mathfrak H(\mathbb C^4)\ [\mathrm{CT0.12}:\ \text{the tiling is exact}];
\]

%dsp-level:0
\[
(2)\quad\textbf{one arena and one action:}
\]

%dsp-level:0
\[
\qquad M,\ \ \text{the charts}\ (\tau_n,h_1,h_2,h_3),\ \ \text{and the dictionary are aspect-blind:}
\]

%dsp-level:0
\[
\qquad \text{the weights depend on}\ (\Gamma,\ \text{registrations})\ \ \text{only}\implies\text{relabeling-invariant}\implies\text{diffeomorphism-invariant}
\]

%dsp-level:0
\[
\qquad [\mathrm{CT2.6};\ \mathrm{CT2.7}\ \mathrm{Pf}];
\]

%dsp-level:0
\[
\textbf{the Lovelock closure:}
\]

%dsp-level:0
\[
\qquad S_{\mathrm{geom}}=\frac1{2\kappa^2}\int\sqrt{-g}\,(R-2\Lambda)\ \ \text{with}\ \neg\exists\ \text{additional local higher-curvature terms in}\ D=4
\]

%dsp-level:0
\[
\qquad [\mathrm{CT2.7}:\ \text{the target dynamics is unique and force-independent}];
\]

%dsp-level:0
\[
(3)\quad\textbf{three fingerprints carry the force identity:}
\]

%dsp-level:0
\[
\qquad \boxed{\ \text{aspect content enters the continuum exactly through}\ \ \bigl(T_{\mu\nu}\text{'s channel},\ \ \text{the mask slot},\ \ \text{the coupling profile}\bigr)\ }
\]

%dsp-level:0
\[
(\mathrm i)\quad T_{\mu\nu}=T^{\mathrm{bd}}_{\mu\nu}+T^{\mathrm{tr}}_{\mu\nu}:
\]

%dsp-level:0
\[
\qquad \text{the bound channel}\ (\rho+p)u_\mu u_\nu+pg_{\mu\nu}\ \ \text{against the transit channel}\ \ \sum_a\bigl(F^{(a)}_{\mu\alpha}F^{(a)\ \alpha}_{\ \ \nu}-\tfrac14g_{\mu\nu}F^2\bigr)
\]

%dsp-level:0
\[
\qquad [\mathrm{CT1.22};\ \mathrm{D1.5}];
\]

%dsp-level:0
\[
(\mathrm{ii})\quad\text{the mask slot:}\ \ \mathbf F_M^{(n)}=\bigoplus_{A\in\mathcal A(n)}F_A|_{\mathcal K^{(A)}},
\]

%dsp-level:0
\[
\qquad \dim\mathcal K_n=(6,16,17,25)
\]

%dsp-level:0
\[
\qquad [\mathrm{CT2.4};\ \mathrm{CT0.14};\ \mathrm{BM.36}(1)];
\]

%dsp-level:0
\[
(\mathrm{iii})\quad\text{the coupling profile:}\ \ e(r,E)=e_0\bigl[Nf/(1+G)\bigr]^{1/2},
\]

%dsp-level:0
\[
\qquad (g,g')\ \ \text{family forms},
\]

%dsp-level:0
\[
\qquad \text{the channel constants}
\]

%dsp-level:0
\[
\qquad [\mathrm{BM.133}(6);\ \mathrm{BM.39}(1);\ \mathrm{BM.80}(1)];
\]

%dsp-level:0
\[
(4)\quad\textbf{the per-force rows (each a read of the same map):}\qquad\begin{array}{l|l|l}\text{aspect} & \text{its continuum read} & \text{the locus}\\\hline\text{gravity} & \text{the field equations themselves:}\ G_{\mu\nu}+\Lambda g_{\mu\nu}=\kappa^2T_{\mu\nu} & [\mathrm{CT1.23}]\\\hline\text{the uniform dark aspect} & F_{D^+}=\Lambda\mathbb I:\ \ \text{excluded from}\ Z_{\mathrm{force}},\ \ \text{carried by}\ \Lambda\ \ \text{in}\ S_{\mathrm{geom}} & [\mathrm{A8};\ \mathrm{D4.3};\ \mathrm{CT1.22}]\\\hline\text{EM} & \partial_\mu F^{\mu\nu}=eJ_Q^\nu;\ \ \text{the Coulomb read outside the support} & [\mathrm{Def\text{-}BM.12}(1)\text{--}(3);\ \mathrm{EM.15}]\\\hline\text{weak} & \text{the shift block:}\ \ N_i=\Omega^2\varkappa\beta w_i\ —\ \ \text{the weak aspect is the metric's own}\ g_{0i} & [\mathrm{FT.157}(1);\ \mathrm{CT1.12};\ \mathrm{BM.34}(3)]\\\hline\text{colour} & \text{sealed content only: the trace law}\ \mathfrak T(H)=3a;\ \ \neg\text{a triplet pole, physical norm}\ 0 & [\mathrm{CT3.22};\ \mathrm{BM.1}(2);\ \mathrm{FT.117}]\\\hline\text{Janus} & \partial_\mu J^\mu=0\ \ \text{with}\ \ J^0=\text{the seal density},\ J^i=\text{the hop current on the}\ w_i\ \text{carrier} & [\mathrm{BM.7};\ \mathrm{BM.11}(1)\text{--}(3)]\\\hline\text{the locked sector} & T^{(D^-)}_{\mu\nu}=\rho_{D^-}u_\mu u_\nu,\quad w_{D^-}=0 & [\mathrm{CT4.5}]\end{array}
\]

%dsp-level:0
\[
(5)\quad\textbf{the pp-wave witness (the consistency check that names no force):}
\]

%dsp-level:0
\[
\qquad \text{the nonlinear contracted Bianchi identity}\ \nabla_\mu G^{\mu\nu}
\]

%dsp-level:0
\[
= 0\ \ \text{holds on the exact rational pp-wave:}\ \ \det g
\]

%dsp-level:0
\[
= -1\ \ \text{identically, metric compatibility on all}\ 64\ \ \text{components, the}\ \Lambda\ \ \text{row}\ \nabla_\mu(\Lambda g^{\mu\nu})\qquad = 0
\]

%dsp-level:0
\[
\qquad [\mathrm{BM.11}(4);\ \text{machine re-cited}];
\]

%dsp-level:0
\[
\textbf{the witness is source-free:}
\]

%dsp-level:0
\[
\qquad \text{the identity is a property of the dictionary image,}\ \neg\text{of any aspect's content}\ —\ \text{it holds whatever filled}\ T_{\mu\nu}
\]

%dsp-level:0
\[
\qquad [\mathrm{CT1.23}\ \mathrm{Pf}\text{'s}\ \delta\text{-identity structure};\ \text{machine: the citation census}];
\]

%dsp-level:0
\[
(6)\quad\textbf{the state-side row:}
\]

%dsp-level:0
\[
\qquad \text{the Lyapunov contraction splits no class:}\ \ \Delta m_{\mathrm{Lyapunov}}=0\ \ \text{on every read species at once}
\]

%dsp-level:0
\[
\qquad [\mathrm{BM.51}(6):\ \text{every input is class-shared}];
\]

%dsp-level:0
\[
\qquad \text{the duality faces are aspect-blind}\ [\mathrm{BM.131}(3),(5)];
\]

%dsp-level:0
\[
(7)\quad\textbf{the universality theorem:}
\]

%dsp-level:0
\[
\qquad \boxed{\ \textbf{the continuum limit is one map with one domain law:}\ \ \text{force identity enters exactly through}\ (3)\text{'s three fingerprints and never through the arena, the action, the domain law, or the limit}\ }
\]

%dsp-level:0
\[
\textbf{the per-force consequence:}
\]

%dsp-level:0
\[
\qquad \text{the colour, weak, EM and gravitational reads share}\ \mathcal C\text{'s formed domain}\ [\mathrm{BM.146}(4)]\ \text{and}\ \mathrm{CT2.7}\text{'s unique action}\ —\ \text{what differs is the source channel, the mask slot, and the coupling profile, all three derived}
\]

%dsp-level:0
\[
\qquad [(4)\text{'s table}];
\]

%dsp-level:0
\[
(8)\quad\textbf{status:}
\]

%dsp-level:0
\[
\qquad \textbf{derived:}\ (1)\text{--}(7);
\]

%dsp-level:0
\[
\qquad \neg\text{a declaration},
\]

%dsp-level:0
\[
\qquad \neg\text{a P4 addition},
\]

%dsp-level:0
\[
\qquad \neg\text{a measured value}\ [\mathrm{P5}],
\]

%dsp-level:0
\[
\qquad \neg\text{an identification of any aspect with an external force law}\ [\mathrm{P5};\ \text{the names are the corpus's own}].
\]

\textbf{Req.} A8 (Part 1); D1.5, D4.3 (Parts 1–2); CT0.13, CT0.14, CT1.22, CT1.23, CT2.4, CT2.6, CT2.7, CT3.22, CT4.5 (Part 2); P5 (Parts 1–2); EM.15 (Part 2); FT.117, FT.157(1), FT.83(2), FT.83(3) (Part 3); Def-BM.12, BM.1(2), BM.7, BM.11(1), BM.11(3), BM.11(4), BM.34(3), BM.36(1), BM.39(1), BM.51(6), BM.52(1), BM.80(1), BM.131(3), BM.131(5), BM.133(6), BM.146(4) (this part); the citation censuses [machine: \texttt{checks/\allowbreak{}verify\_\allowbreak{}H\_\allowbreak{}continuum\_\allowbreak{}limit\_\allowbreak{}era1\_\allowbreak{}closure.py} §F].

\textbf{Pf.}

%dsp-level:0
\[
(1),(2):\ \text{the censuses re-run and}\ \mathrm{CT2.7}\text{'s uniqueness re-cited}\ [\text{machine: the integer sums}];
\]

%dsp-level:0
\[
\qquad (3):\ \text{the three fingerprints assembled from}\ \mathrm{CT1.22},
\]

%dsp-level:0
\[
\ \mathrm{CT2.4},
\]

%dsp-level:0
\[
\ \mathrm{BM.133}(6)\ \text{and the family forms}\ [\text{machine: the locus census}]
\]

%dsp-level:0
\[
(4):\ \text{each row re-cited from its own derived tag}\ [\text{machine: the table's citation pins}];
\]

%dsp-level:0
\[
\qquad (5):\ \mathrm{BM.11}(4)\text{'s witness rows re-cited; the source-freeness is the}\ \delta\text{-identity structure of}\ \mathrm{CT1.23}\text{'s Pf}\ [\text{machine: the citation census}]
\]

%dsp-level:0
\[
(6):\ \mathrm{BM.51}(6)\text{'s class-degenerate zero re-cited};\qquad(7),(8):\ \text{assembly}.\ \blacksquare
\]

\textbf{Phys.} Universality is not an assumption about forces agreeing — it is a structural census. There is one operator, its blocks tile the arena exactly, and the stack aspects sit on their own registers. There is one target manifold, one dictionary, and one geometric action, unique by the Lovelock count, with no term that could carry a force label. Force identity has exactly three homes: which channel of the stress tensor the aspect fills, which slot of the mask ladder it occupies, and which coupling profile it runs. The per-force table reads all seven aspects through the same map, and the pp-wave witness certifies the identity that no source ever signs: the contracted Bianchi row holds because of what the dictionary image is, not because of what fed it. The state-side dynamics agrees: the contraction splits no class. The continuum limit is the same limit for every force, and the differences are derived fingerprints, not different limits.

\end{theorem}

\begin{theorem}[{\texorpdfstring{BM.151 — The derivation of CT2.6(ii)–(iii): the unbounded ladder, the spectral bound, and the empty extra-axiomatic list}{BM.151 - The derivation of CT2.6(ii)-(iii): the unbounded ladder, the spectral bound, and the empty extra-axiomatic list}}]\label{tag:BM-151}
\textbf{Thm.}

%dsp-level:0
\[
(1)\quad\textbf{condition (ii) is a theorem of the ledger's own configuration space:}
\]

%dsp-level:0
\[
\qquad \text{(ii):}\ \ \text{coarse-graining cells contain unboundedly many real edges as}\ \mathfrak h\to0
\]

%dsp-level:0
\[
\qquad [\mathrm{CT2.6}\ \text{statement}];
\]

%dsp-level:0
\[
(\mathrm a)\quad\textbf{the configuration space:}
\]

%dsp-level:0
\[
\qquad \mathrm{A7}\ \text{admits}\ \Gamma=(V,E)\ \ \text{at every scale};
\]

%dsp-level:0
\[
\qquad R\subseteq E_{\mathrm{dir}}\ \ \text{is unconstrained in size};
\]

%dsp-level:0
\[
\qquad E_{\mathrm{int}}=|R|
\]

%dsp-level:0
\[
\qquad [\mathrm{A7};\ \mathrm{D0.10};\ \mathrm{D0.8};\ \mathrm{Def\text{-}H.32}(1)\text{--}(2)];
\]

%dsp-level:0
\[
(\mathrm b)\quad\textbf{the witness ladder:}
\]

%dsp-level:0
\[
\qquad \text{the cycle floors}\ \Sigma_L\ \ \text{exist at every}\ L\ \ \text{with}\ E_{\mathrm{int}}=L
\]

%dsp-level:0
\[
\qquad [\mathrm{FT.32}];
\]

%dsp-level:0
\[
\qquad \text{the polygon family is unbounded with injective mass/rank/null reads and}\ \neg\text{a combinatorial or geometric ceiling}
\]

%dsp-level:0
\[
\qquad [\mathrm{FT.45};\ \mathrm{FT.46}];
\]

%dsp-level:0
\[
\boxed{\ \forall n\ \ \exists\ \ \text{an admissible cell configuration with}\ E_{\mathrm{int}}\ge n\ }
\]

%dsp-level:0
\[
\qquad [\text{machine: the constructive witness}\ R=E_{\mathrm{dir}}(\Sigma_L)\ \text{at}\ L=2..8];
\]

%dsp-level:0
\[
(\mathrm c)\quad\textbf{the refinement row:}
\]

%dsp-level:0
\[
\qquad \text{at fixed chart span}\ s:
\]

%dsp-level:0
\[
\qquad k=\frac{s}{\Delta x_{\mathrm B}}=\frac{sE_{\mathrm{int}}\ln2}{\mathfrak h}\ \ \text{envelopes per cell}
\]

%dsp-level:0
\[
\qquad [\mathrm{Def\text{-}BM.4}(1);\ \mathrm{FT.44}(4)];
\]

%dsp-level:0
\[
\qquad \text{every materialized chain carries}\ \ge2\ \ \text{directed real edges}
\]

%dsp-level:0
\[
\qquad [\mathrm{Def\text{-}H.4};\ \mathrm{Census\ FT.28}(1)]:
\]

%dsp-level:0
\[
\#\{E_{\mathcal R}\ \text{per cell}\}\ \ge\ \frac{2sE_{\mathrm{int}}\ln2}{\mathfrak h}\ \xrightarrow[\ \mathfrak h\to0\ ]{}\ \infty
\]

%dsp-level:0
\[
\qquad [\text{machine: the Archimedean row on exact rationals}];
\]

%dsp-level:0
\[
(\mathrm d)\quad\textbf{the scope audit:}
\]

%dsp-level:0
\[
\qquad \text{the witness rows}\ (\mathrm a)\text{--}(\mathrm c)\ \text{are combinatorial and card content:}\ \neg\text{one of them uses}\ \mathrm{CT2.6}\text{'s conclusion}\ (M\ \text{a manifold})\ \text{as premise};
\]

%dsp-level:0
\[
\text{the}\ \mathrm{CT2.6}\text{-citations of}\ \mathrm{Def\text{-}BM.4},
\]

%dsp-level:0
\[
\ \mathrm{BM.13},
\]

%dsp-level:0
\[
\ \mathrm{FT.47}\ \text{are chart-arena inheritances}\ (\text{scope rows}),
\]

%dsp-level:0
\[
\ \neg\text{uses of}\ (\mathrm{ii})/(\mathrm{iii})\ \text{as premises}
\]

%dsp-level:0
\[
\qquad [\text{machine: the citation-pattern scan}];
\]

%dsp-level:0
\[
\text{the materialization hypotheses}\ (\mathrm{M1}/\mathrm{M2}\ \text{species})\ \text{remain declared spectral scopes at their own loci, untouched and distinct}
\]

%dsp-level:0
\[
\qquad [\text{A §0.2a integration note};\ \mathrm{Def\text{-}H.20}];
\]

%dsp-level:0
\[
(2)\quad\textbf{condition (iii) is a theorem of the pair-state spectral bound:}
\]

%dsp-level:0
\[
\qquad \text{(iii):}\ \ \text{sector coefficient fields}\ [\mathrm{D2.1}]\ \ \text{of bounded variation across neighbouring cells}
\]

%dsp-level:0
\[
\qquad [\mathrm{CT2.6}\ \text{statement}];
\]

%dsp-level:0
\[
(\mathrm a)\quad\textbf{the coefficients are state components:}
\]

%dsp-level:0
\[
\qquad \rho=\tfrac14I+\tfrac14\sum_ar_aG_a,
\]

%dsp-level:0
\[
\qquad r_a=\mathfrak T(\delta\rho\,G_a)\in\mathbb R,
\]

%dsp-level:0
\[
\qquad |r|^2=4\,\mathfrak T(\delta\rho^2)
\]

%dsp-level:0
\[
\qquad [\mathrm{L3}(1);\ \mathrm{CT1.1}:\ \mathfrak T(G_aG_b)=4\delta_{ab}];
\]

%dsp-level:0
\[
(\mathrm b)\quad\textbf{the spectral bound:}
\]

%dsp-level:0
\[
\qquad \operatorname{spec}(\delta\rho)\subset\bigl[-\tfrac14,\tfrac34\bigr],
\]

%dsp-level:0
\[
\qquad \textstyle\sum_i\mu_i=0
\]

%dsp-level:0
\[
\qquad [\mathrm{A1}:\ \rho\succeq0,\ \mathfrak T\rho=1;\ \mathrm{A5};\ \mathrm{L1}]:
\]

%dsp-level:0
\[
\max\ \textstyle\sum_i\mu_i^2=\tfrac34\ \ \text{attained exactly at the vertex orbit}\ \bigl(\tfrac34,-\tfrac14,-\tfrac14,-\tfrac14\bigr)
\]

%dsp-level:0
\[
\qquad [\text{machine: the vertex enumeration of the polytope}\ \{\sum\mu=0\}\cap[-\tfrac14,\tfrac34]^4\ \text{and convexity of}\ \sum\mu^2];
\]

%dsp-level:0
\[
\boxed{\ |r_a|\ \le\ |r|\ =\ 2\sqrt{\textstyle\sum_i\mu_i^2}\ \le\ 2\sqrt{\tfrac34}\ =\ \sqrt3\ }\qquad[\text{machine:}\ |r|^2\le3\ \text{exact}];
\]

%dsp-level:0
\[
\textbf{the component chart:}
\]

%dsp-level:0
\[
\qquad \rho=\tfrac14\sum_{\mu\nu}c_{\mu\nu}\,\sigma_\mu\otimes\sigma_\nu,
\]

%dsp-level:0
\[
\qquad c_{\mu\nu}=\mathfrak T\bigl(\rho\,\sigma_\mu\otimes\sigma_\nu\bigr),
\]

%dsp-level:0
\[
\qquad |c_{\mu\nu}|\le\bigl\|\sigma_\mu\otimes\sigma_\nu\bigr\|_{\mathrm{op}}=1
\]

%dsp-level:0
\[
\qquad [\text{math: the state-functional bound};\ \text{machine: exact rational density-matrix samples}];
\]

%dsp-level:0
\[
\textbf{the}\ \mathrm{D2.1}\ \textbf{row:}
\]

%dsp-level:0
\[
\qquad \text{the sector coefficients}\ (a_{00},q_{ij},w_i)\ \text{are fixed rational projections of either chart}\ [\mathrm{D2.1}:\ \text{the basis row};\ \mathrm{CT1.2}]:
\]

%dsp-level:0
\[
\qquad \text{uniformly bounded by explicit rational constants,}\ \neg\text{free data};
\]

%dsp-level:0
\[
(\mathrm c)\quad\textbf{the fields are averages:}
\]

%dsp-level:0
\[
\qquad \text{coefficients}=\text{local ledger averages}
\]

%dsp-level:0
\[
\qquad [\mathrm{CT2.6}\ \mathrm{Pf}\ \text{row}];
\]

%dsp-level:0
\[
\qquad \text{convex combinations preserve the ball:}\ \|\text{field}\|\le\sqrt3\ \ \text{per coefficient};
\]

%dsp-level:0
\[
(\mathrm d)\quad\textbf{the variation bound:}
\]

%dsp-level:0
\[
\qquad \boxed{\ |\Delta\,\text{coefficient}|\ \le\ 2\sqrt3\ \ \text{across any two cells},\qquad\mathrm{TV}(m\ \ \text{neighbouring cells})\ \le\ 2\sqrt3\,m\ }
\]

%dsp-level:0
\[
\qquad [\text{machine:}\ (2\sqrt3)^2=12\ \text{exact};\ \text{the triangle inequality}];
\]

%dsp-level:0
\[
(\mathrm e)\quad\textbf{the tighter formed-locus census:}
\]

%dsp-level:0
\[
\qquad \varkappa\|q\|_{\mathrm{op}}<1\ [\mathrm{CT1.12}];
\]

%dsp-level:0
\[
\qquad N\in(0,1]\ [\mathrm{BM.132}(4)];
\]

%dsp-level:0
\[
\qquad w^{(n)}_{\mathrm{mat}}\in\bigl(\tfrac1n,1\bigr]\ [\mathrm{BM.137}(1)];
\]

%dsp-level:0
\[
\qquad p\in(0,1]\ [\mathrm{D0.4};\ \mathrm{BM.19}(2)];
\]

%dsp-level:0
\[
f\in[0,1)\ [\mathrm{BM.88}(2)(\mathrm{ii})];
\]

%dsp-level:0
\[
\qquad \mathsf g\in(0,1]\ [\mathrm{BM.133}(4)];
\]

%dsp-level:0
\[
\qquad \Gamma_{\mathrm{tr}}\in\bigl[0,\tfrac13\bigr]\ [\mathrm{BM.44}(2);\ \mathrm{BM.116}(2)];
\]

%dsp-level:0
\[
\qquad w_{\mathrm{mat}}\in\bigl[\tfrac13,1\bigr]\ \ \text{measure-free}\ [\mathrm{BM.105}(2)\text{--}(4)];
\]

%dsp-level:0
\[
\textbf{every monotone row's total variation equals its range:}
\]

%dsp-level:0
\[
\qquad \text{the exact variation widths are}\ \bigl(1,\ 1-\tfrac1n,\ 1,\ 1,\ 1,\ \tfrac13,\ \tfrac23\bigr)\ \ \text{on}\ \bigl(N,\ w^{(n)}_{\mathrm{mat}},\ p,\ f,\ \mathsf g,\ \Gamma_{\mathrm{tr}},\ w_{\mathrm{mat}}\bigr)
\]

%dsp-level:0
\[
\qquad [\text{machine: the monotone-BV identity on exact rationals}];
\]

%dsp-level:0
\[
(\mathrm f)\quad\textbf{the dynamics row:}
\]

%dsp-level:0
\[
\qquad \text{the Haar face preserves the coefficient spectrum and the tilt face contracts it:}
\]

%dsp-level:0
\[
\qquad \mu'=\bigl(1-\tfrac\gamma4\bigr)\mu+O(\mu^2),
\]

%dsp-level:0
\[
\qquad \bigl(1-\tfrac\gamma4\bigr)\le1
\]

%dsp-level:0
\[
\qquad [\mathrm{BM.131}(3);\ \mathrm{L2};\ \text{machine}]:
\]

%dsp-level:0
\[
\text{the bound ball}\ \{|\mu|\le\tfrac34\}\ \ \text{is invariant under both faces}\ —\ \text{the recorded dynamics never violates}\ (\mathrm b);
\]

%dsp-level:0
\[
(3)\quad\textbf{the derivation:}
\]

%dsp-level:0
\[
\qquad \boxed{\ \mathrm{CT2.6}(\mathrm{ii}),(\mathrm{iii})\ \ \textbf{are theorems at the stated loci:}\ \ (\mathrm{ii})\ \text{by}\ (1),\qquad(\mathrm{iii})\ \text{by}\ (2)\ }
\]

%dsp-level:0
\[
\textbf{the extra-axiomatic list of}\ \text{A §0.2a}\ \textbf{is empty};
\]

%dsp-level:0
\[
\qquad \mathrm{CT2.6}\text{'s statement and Pf stand unmodified in their conclusions}\ —\ \text{the conditions are now supplied by derivation rather than by declaration};
\]

%dsp-level:0
\[
\textbf{the in-place amendments:}
\]

%dsp-level:0
\[
\qquad \text{A §0.2a's list sentence and}\ \mathrm{CT2.6}\text{'s two condition labels carry the derivation markers}\ [\text{the FT.62/FT.66 patch pattern}];
\]

%dsp-level:0
\[
\qquad \text{the point-of-use declarations across the corpus stand as scope-inheritance markers over the same loci};
\]

%dsp-level:0
\[
(4)\quad\textbf{the packaging:}
\]

%dsp-level:0
\[
\qquad \neg\text{a declaration is issued by this tag};
\]

%dsp-level:0
\[
\qquad \neg\text{a P4 member};
\]

%dsp-level:0
\[
\qquad \mathfrak h\ \ \text{remains the single dimensionful scale}\ [\mathrm{P1}];
\]

%dsp-level:0
\[
\qquad \neg\text{a measured value}\ [\mathrm{P5}];
\]

%dsp-level:0
\[
\text{no tag loses content:}
\]

%dsp-level:0
\[
\qquad \text{every tag inheriting}\ \mathrm{CT2.6}\text{'s hypotheses inherits the same conditions, now derived at the same scope}\ [\text{machine: the inheritance census}];
\]

%dsp-level:0
\[
(5)\quad\textbf{the gap row:}
\]

%dsp-level:0
\[
\qquad \textbf{the continuum-limit gap}\ —\ \text{the extra-axiomatic entry of}\ \text{A §0.2a},
\]

%dsp-level:0
\[
\ \text{the corpus's last non-axiomatic import}\ —\ \textbf{is closed};
\]

%dsp-level:0
\[
\text{the finite-}\mathfrak h\ \ \text{residue is distinct and stands}\ [\text{the clean-bifurcation loci};\ \mathrm{FT.66}(3);\ \mathrm{BM.143}(7)];
\]

%dsp-level:0
\[
(6)\quad\textbf{status:}
\]

%dsp-level:0
\[
\qquad \textbf{derived:}\ (1),(2);
\]

%dsp-level:0
\[
\qquad \textbf{adjudicated:}\ (3)\text{--}(5);
\]

%dsp-level:0
\[
\qquad \boxed{\ \textbf{the axiom surface is now the entire non-derived surface of the corpus}\ }.
\]

\textbf{Req.} A1, A5, A7 (Part 1); D0.4, D0.8, D0.10 (Part 1); CT1.1, CT1.2, CT1.10, CT1.29, L1, L2 (Part 1); CT1.12, CT2.6, D2.1 (Part 2); P1, P5 (Parts 1–2); Def-H.4, Def-H.20, Def-H.32(1), Def-H.32(2), FT.32, FT.44(4), FT.45, FT.46, FT.66(3), L3(1) (Part 3); Def-BM.4(1), BM.13, BM.19(2), BM.44(2), BM.47(6), BM.88(2), BM.105, BM.116(2), BM.131(3), BM.132(4), BM.133(4), BM.137(1), BM.143(7) (this part); Census FT.28(1) (row-wise); vertex enumeration of a rational polytope, convexity of the quadratic form, the monotone-BV identity, Archimedean rows on exact rationals [math; machine: \texttt{checks/\allowbreak{}verify\_\allowbreak{}H\_\allowbreak{}continuum\_\allowbreak{}limit\_\allowbreak{}era1\_\allowbreak{}closure.py} §G].

\textbf{Pf.}

%dsp-level:0
\[
(1)(\mathrm b):\ \Sigma_L\ \text{at}\ L=2..8:\ \ \text{the unidirectional}\ L\text{-cycle with}\ R=E_{\mathrm{dir}}(\Sigma_L):
\]

%dsp-level:0
\[
\quad |R|=L\ [\mathrm{FT.32};\ \text{machine: the constructive witness on the cycle skeletons}];
\]

%dsp-level:0
\[
\qquad (\mathrm c):\ \frac{2sE\ln2}{\mathfrak h}
\]

%dsp-level:0
\[
\ge T\iff\mathfrak h\le\frac{2sE\ln2}T:\ \ \text{witness}\ \mathfrak h
\]

%dsp-level:0
\[
= \frac{2sE\ln2}{T+1}\ [\text{machine: exact rationals at}\ T=10^6]
\]

%dsp-level:0
\[
(2)(\mathrm b):\ \text{the polytope}\ \{\mu\in[-\tfrac14,\tfrac34]^4:\ \sum\mu_i=0\}\ \text{has the vertex orbit}\ \binom41\bigl(\tfrac34,-\tfrac14,-\tfrac14,-\tfrac14\bigr):\ \ \text{any vertex fixes three coordinates at bounds; two}\ +\tfrac34\text{'s force the fourth to}\ -\tfrac32<-\tfrac14\ [\text{machine: the exhaustive bound-pattern enumeration}];
\]

%dsp-level:0
\[
\textstyle\sum\mu^2\ \text{convex}\implies\max\ \text{at a vertex}:\ \ \tfrac9{16}+\tfrac3{16}=\tfrac34\ [\text{machine}];
\]

%dsp-level:0
\[
\qquad |r|^2=4\sum\mu_i^2\le3\ [\mathrm{L3}(1)];
\]

%dsp-level:0
\[
\qquad |r_a|\le|r|\le\sqrt3\ [\text{Cauchy--Schwarz on components};\ \text{machine:}\ r_a^2\le|r|^2]
\]

%dsp-level:0
\[
(2)(\mathrm d):\ |\Delta|\le|r_a|+|r_a'|\le2\sqrt3;
\]

%dsp-level:0
\[
\qquad (2\sqrt3)^2=12\ [\text{machine}];
\]

%dsp-level:0
\[
\qquad (2)(\mathrm e):\ \text{the range rows re-cited; the monotone-BV identity:}\ \mathrm{TV}
\]

%dsp-level:0
\[
= \sup-\inf\ \text{on monotone functions}\ [\text{machine: rational grids}]
\]

%dsp-level:0
\[
(1)(\mathrm d),
\]

%dsp-level:0
\[
(3)\text{--}(6):\ \text{the citation-pattern scan, the amendment pins, and the inheritance census}\ [\text{machine}].\ \blacksquare
\]

\textbf{Phys.} The corpus's last non-axiomatic import is dissolved, and both halves dissolve into structure the ledger already owned. The richness condition was never an outside gift: the ledger's own configuration space carries cells of every finite content — the cycle ladder witnesses each level constructively — and refining the hop length at fixed span drives every cell's real-edge count past any bound. The regularity condition was never an outside gift either: the sector coefficients are components of a positive, unit-trace pair state, whose spectrum lives in an exact interval; convexity pins the quadratic extremum at a single vertex orbit, the bound is the square root of three, fields are averages of bounded data, and every weight family of the corpus carries its own exact range, so total variation equals range width everywhere. Neither derivation touches the manifold conclusion it licenses, and the scope citations that once looked like dependencies are audited as arena inheritances. The extra-axiomatic list is empty; the axiom surface is the entire non-derived surface; the conditions remain in force at every point of use, now as theorems.

\end{theorem}

\medskip\noindent\textit{\texorpdfstring{BM.152 — The status of the Continuum-Limit Tier and the closure of the foundational era}{BM.152 - The status of the Continuum-Limit Tier and the closure of the foundational era}.}\ Structural-constraints record: see Appendix B, Section B.1.35 of the Addendum.\par\medskip

\phantomsection
\section[Closure exports of Part 4 (corpus terminus)]{\texorpdfstring{Closure exports of Part 4 (corpus terminus)}{Closure exports of Part 4 (corpus terminus)}}

%dsp-level:0
\[
\begin{array}{lll}
\mathrm{Def\text{-}BM.1}\ \text{sealed composites};\ \text{validity/read clauses} & \mathrm{Def\text{-}BM.2}\ \text{composite card grammar} & \mathrm{Def\text{-}BM.3}\ \text{the identification layer (declared)}\\
\mathrm{BM.1}\ \text{observability switch};\ \text{trace-only read} & \mathrm{BM.2}\ 4\mathbb Z\ \text{junction phase};\ \text{additive winding} & \mathrm{BM.3}\ C_3\ \text{forced at the seal};\ \delta_q,\ 3\delta_q=\delta_\ell\\
\mathrm{BM.4}\ \text{colour/weak invariance};\ \text{parity passage} & \mathrm{BM.5}\ 25\ \text{cells},\ 6\ \text{families};\ \text{conditioning law} & \mathrm{BM.6}\ 8\ \text{open construction specs}\\
\mathrm{BM.7}\ \#J^+{=}\#J^-;\ \text{telescoping};\ \partial^2{=}0;\ \neg U(1)\ \text{on}\ q_{\mathrm{EM}} & \mathrm{BM.8}\ \text{hop}\to\text{translation};\ C_3\to SO(3);\ \text{tick obstruction} & \mathrm{BM.9}\ T_{\mu\nu};\ \text{sign reconciliation};\ \text{dust source}\\
\mathrm{BM.10}\ \neg[\hat x,\hat p]{=}i\ \text{in}\ 16;\ \Delta x\Delta p\ge\tfrac12;\ \text{MT},\ \text{CN} & \mathrm{BM.11}\ \partial_\mu J^\mu{=}0;\ \nabla_\mu T^{\mu\nu}{=}0\ \text{two routes} & \mathrm{BM.12}\ \text{the Noether adjudication};\ \neg\text{Cal-BM parameter}\\
\mathrm{Def\text{-}BM.4}\ \Delta x_{\mathrm B}{=}1/m_{\mathrm{hid}};\ V_{\mathrm{eff}}{=}(\Delta x_{\mathrm B})^3 & \mathrm{BM.13}\ \Delta x_{\mathrm B}{=}\mathfrak h/(E\ln2);\ \text{sub-hop}\ E\ge2 & \mathrm{BM.14}\ \rho{=}(E\ln2)^4/\mathfrak h^4;\ [\rho]{=}\mathfrak q^{+2}\\
\mathrm{BM.15}\ R{=}\kappa^2(\rho{-}3p){+}4\Lambda;\ \kappa^2\rho{=}2\pi E^4(\ln2)^3/\mathfrak h^2 & \mathrm{BM.16}\ (\Delta x_{\mathrm B})^2\kappa^2\rho{=}2\pi E^2\ln2;\ \Delta x_{\mathrm B}r_s{=}2G_N & \mathrm{BM.17}\ E_{\mathrm{int}}^{(n)}{=}6:\ \text{the neutron row}\\
\mathrm{BM.18}\ \text{the calibration adjudication} & \neg\text{a new dimensionful scale}\ [\mathrm{P1}] & \neg\text{a measured value}\ [\mathrm{P5}]\\
\mathrm{Def\text{-}BM.5}\ \mathcal E(\mathcal X);\ p{=}x_+/(x_+{+}x_-);\ E_{\mathrm{int}}(R){=}|R| & \mathrm{Def\text{-}BM.6}\ \text{the family naming rows};\ Q{=}\tfrac23 & \mathrm{BM.19}\ p{=}1{-}4\varepsilon_q;\ \Gamma{=}0\ \text{at}\ \varepsilon_q{=}0\\
\mathrm{BM.20}\ RV{=}\kappa^2m;\ \langle m\rangle{=}\text{the averaged flux};\ m^2_B(q,S) & \mathrm{BM.21}\ \mathcal B{\ne}0\ \text{in the ensemble},\ 0\ \text{in the count} & \mathrm{BM.22}\ F_{LR}{=}qI_2;\ 2|q|\in\{0,2\}\\
\mathrm{BM.23}\ \Delta q{=}\pm1;\ \text{the forced carrier};\ \text{the resolvent rate} & \mathrm{BM.24}\ \text{the flicker adjudication} & \neg\text{Cal-BM parameter};\ \neg\text{a P4 addition}\\
\mathrm{Def\text{-}BM.7}\ \mathrm{EM}^\pm;\ \text{the EM scar};\ \text{the formal-mirror scope} & \mathrm{BM.25}\ \mathrm{CPT}^2{=}\mathrm{id};\ w{\mapsto}w^\ast;\ \lambda\in\{\pm1,\pm i\} & \mathrm{BM.26}\ \text{the proton row};\ 2|q|{=}2\\
\mathrm{BM.27}\ \text{the antiproton row};\ \text{one orbit of size two} & \mathrm{BM.28}\ W(0){=}20,\ W({\pm}1){=}15;\ \text{three mass routes} & \mathrm{BM.29}\ \text{the sector adjudication};\ \neg\text{inter-sector dynamics}\\
\mathrm{Def\text{-}BM.8}\ Q{=}\textstyle\sum_iq_{\mathrm{core}}(c_i);\ \text{BM.6}(1)\ \text{discharged} & \text{(blank)} & \text{(blank)}\\
\mathrm{BM.30}\ \text{ratios}\ 1,\ \tfrac43,\ \tfrac{58}{13};\ \neg\text{a route adopted} & \neg\text{a fractional charge}\ [\mathrm{BM.28}(7)] & \neg\text{Cal-BM parameter};\ \neg\text{a measured value}\ [\mathrm{P5}]\\
\mathrm{Def\text{-}BM.9}\ \mathcal E_{\mathrm{ext}};\ \Pi_{\mathcal C};\ g(N){=}(1,1,3,10,25,105) & \mathrm{BM.31}\ g(3){=}1,\ g_\delta(3){=}0;\ N{=}3\ \text{protected} & \mathrm{BM.32}\ Z_\delta{=}0\ \text{by filtration};\ \langle m\rangle_{\mathrm{ext}}{=}\langle m\rangle_\varepsilon\\
\mathrm{BM.33}\ \text{the test's adjudication} & \text{the protection fails at}\ N\ge4 & \neg\text{a new declaration}\\
\mathrm{Def\text{-}BM.10}\ E{=}N\sqrt{m^2{+}\gamma^{ij}p_ip_j}{-}N^ip_i;\ \text{chart-side} & \mathrm{BM.34}\ (p_0{-}N^ip_i)^2{=}N^2(m^2{+}p^2);\ \mathrm{KE}{=}\tfrac{Np^2}{2m}{+}O(p^4);\ E^2\ge\tfrac74m_{\mathrm{hid}}^2 & \mathrm{BM.35}\ m_{\mathrm{rest}}{=}\Delta S_{\mathfrak L}/\mathfrak h{=}m_{\mathrm{hid}};\ \ln4{=}2\ln2\\
\mathrm{BM.36}\ (6,16,17,25);\ \text{ratios}\ (\tfrac13,8,\tfrac1{25});\ \text{Gram spectrum};\ \text{Yukawa-like coupling open} & \mathrm{BM.37}\ \text{routes}{=}\text{measure functionals};\ \text{alphabets disjoint};\ \text{trap closed} & \mathrm{BM.38}\ \text{the tier adjudication};\ \neg\text{Cal-BM parameter};\ \neg\text{a measured value}\\
\mathrm{Def\text{-}BM.11}\ c(i){=}\mathrm{Ad}(U^i)c(0);\ m_i^2{=}M^2(A_0{-}B_0\cos(\delta{-}\tfrac{2\pi i}3)) & \mathrm{BM.39}\ A_0{-}B_0{>}0\ (\text{SOS});\ \sum_im_i^2{=}3M^2A_0;\ \text{splitting}\ {\propto}\,gg' & \mathrm{BM.40}\ \sum_iW_i^2{=}2\Pi_{\mathcal R};\ \sigma\text{-vertex Casimir-neutral}\\
\mathrm{BM.41}\ \text{two-body impossible};\ \nu\text{-member forced};\ Q{=}{-}(m_\ell{+}m_\nu){<}0 & \mathrm{BM.42}\ \partial_\mu M^{\mu\nu\lambda}{=}0;\ p^2{=}p_r^2{+}L^2/r^2;\ 2^{\otimes4}{=}2{\cdot}0{\oplus}3{\cdot}1{\oplus}2 & \mathrm{BM.43}\ \nabla^2\Phi{=}4\pi G_N\rho;\ r_s{=}\tfrac{\mathfrak h E}2\\
\mathrm{BM.44}\ \Gamma_{\mathrm{tr}}\in[0,\tfrac13];\ \tau_{\mathrm{life}}\ge3\ \text{ticks};\ \neg\text{seconds} & \mathrm{BM.45}\ n\to p{+}e^-+\bar\nu_e;\ \text{grammar complete};\ \text{kinematics closed} & \neg\text{Cal-BM parameter};\ \neg\text{a measured value}\\
\mathrm{Def\text{-}BM.12}\ \partial_\mu F^{\mu\nu}{=}eJ^\nu_Q;\ E^{\mathrm{self}}{=}\tfrac{ce^2q^2}{4\pi\Delta x_{\mathrm B}};\ c\ \text{priced}\ (\tfrac12,\tfrac35) & \mathrm{BM.46}\ \Delta m_{\mathrm{EM}}{=}\tfrac{3ce^2\ln2}{2\pi\mathfrak h}{>}0;\ \text{the proton heavier} & \mathrm{BM.47}\ \text{UV: the envelope};\ \text{IR:}\ \mathrm{EM.11};\ f(k){=}1{-}\tfrac1k\\
\mathrm{BM.48}\ V{=}{-}\varepsilon_{\mathrm{EM}}/d;\ E_n{=}{-}\tfrac{\varepsilon_{\mathrm{EM}}^2m_r}{2n^2};\ \text{the bound channel}\ \neg\text{constructed} & \mathrm{BM.49}\ Q_{\mathrm{corr}}{<}0\ \text{deeper};\ E_{\max}{:}\langle E\rangle{=}2{:}1;\ \tfrac43\ \text{refuted} & \mathrm{BM.50}\ \text{the tier adjudication};\ \neg\text{a numerical}\ e\ [\mathrm{P5}]\\
\mathrm{BM.51}\ \bigl(\tfrac9{16}\bigr)^{6};\ V{=}8\Delta S_{\mathrm{pair}};\ \Delta m_{\mathrm{Lyapunov}}{=}0 & \mathrm{BM.52}\ \text{six sectors censused};\ \Delta m_{\mathrm{force}}{=}-\Delta m_{\mathrm{EM}} & \mathrm{BM.53}\ \Delta m_{\mathrm{metric}}{=}0;\ \Delta t\ge\pi\mathfrak h/(3ce^2\ln2)\\
\mathrm{BM.54}\ E^{\mathrm{info}}_{\mathrm{bind}}{=}0;\ \mathcal B^{(n)}{=}\mathcal B^{(p)};\ \text{the atom}\ \neg\text{constructed} & \mathrm{BM.55}\ \Delta m{=}-\tfrac{3ce^2\ln2}{2\pi\mathfrak h}{<}0;\ \tfrac{m_n}{m_p}{=}\tfrac{4\pi}{4\pi+ce^2} & \mathrm{BM.56}\ \text{the full-scope adjudication};\ \text{the order residue named},\ \neg\text{filled}\\
\mathrm{Def\text{-}BM.13}\ m_q{=}(2,1,1);\ W{=}(1,6,15,20,15,6,1);\ \Sigma_{\mathrm B}{=}\ln\tfrac43 & \mathrm{Def\text{-}BM.14}\ 3\ \text{floors}\ 2E{\in}\{12,24,60\};\ |\mathcal S|{=}3{\cdot}2{=}6;\ \text{the}\ D_S\ \text{row} & \mathrm{BM.57}\ |A_3|{=}3\ \text{vs}\ 1;\ \text{asymmetry}\ 2;\ \Delta m_{\mathrm{vertex}}{=}0\\
\mathrm{BM.58}\ \Delta m_{\mathrm{metric}}{=}0;\ \text{the seal asymmetry survives as}\ \Delta m_{\mathrm{EM}}{>}0 & \mathrm{BM.59}\ \Delta m_{\mathrm{transport}}{=}\ln(4/3)/\mathfrak h{>}0;\ \text{the residue filled on all four grounds} & \mathrm{BM.60}\ 1{:}2{:}5;\ D_S(\mathcal F_g,+){>}D_S(\mathcal F_g,-);\ \text{the inversion}\\
\mathrm{BM.61}\ \Sigma_{\mathrm B}\ \text{coupling-blind};\ \alpha_B\ \neg\text{derivable},\ (\mathrm O1)\text{--}(\mathrm O4) & \mathrm{BM.62}\ \Delta m{=}[\ln\tfrac43{-}6c\varepsilon_{\mathrm{EM}}\ln2]/\mathfrak h;\ \tfrac19{<}\varepsilon^{\mathrm{crit}}(\tfrac35){<}\tfrac18{<}\varepsilon^{\mathrm{crit}}(\tfrac12){<}\tfrac17 & \mathrm{BM.63}\ Q_\beta{<}-\ln3/\mathfrak h{<}0;\ \text{closed unconditionally}\\
\mathrm{BM.64}\ \text{the tier adjudication};\ m_n{>}m_p\ \text{sub-critically} & \neg\text{Cal-BM parameter};\ \neg\text{a new declaration beyond one name attachment} & \neg\text{a measured value}\ [\mathrm{P5}]\\
\mathrm{BM.65}\ \text{the baseline re-certified};\ \tfrac19{<}\varepsilon^{\mathrm{crit}}(\tfrac35){<}\tfrac18{<}\varepsilon^{\mathrm{crit}}(\tfrac12){<}\tfrac17 & \mathrm{BM.66}\ \operatorname{Tr}\widehat F{=}2c_0;\ I_F{=}4\widehat Q^2;\ \text{the operator rest input} & \mathrm{BM.67}\ U(\gamma){=}(-1)^{N_\pi};\ \operatorname{im}(\text{lift})\subseteq\{0,\pi\};\ \text{the obstruction record}\\
\mathrm{BM.68}\ \Omega{=}\Omega_{\mathrm{count}}\Omega_{\mathrm{charge}}\Omega_{\mathrm{geom}};\ \Omega_{\mathrm{geom}}^{\mathrm{stat}}{=}1;\ S_{\mathrm{joint}}{=}6\ln2{+}\ln W(Q) & \mathrm{BM.69}\ \tfrac{m^2_{\mathrm{pole}}}{m_{\mathrm{hid}}^2}{=}\rho_{\mathrm B}\tfrac{D_S}{4E^2};\ \text{Gap 7 minimal}:\ (\mathrm R1)\text{--}(\mathrm R3) & \mathrm{BM.70}\ F_{\ell\nu}{\ge}2\ln2;\ \text{formed}\implies\text{closed};\ \text{margins}\ \ln3,\ln12,\ln24\\
\mathrm{BM.71}\ \text{the five-route census};\ \Delta m_{\mathrm{bridge}}{=}0\ \text{in currency};\ \text{the decay closed} & \mathrm{BM.72}\ \text{the tier adjudication};\ \text{the partial bridge}\wedge\text{the no-go} & \neg\text{Cal-BM parameter};\ \neg\text{a new declaration};\ \neg\text{a measured value}\ [\mathrm{P5}]\\
\mathrm{BM.73}\ \text{A-BM.7 issued};\ (1,0)\ \text{occupancy};\ \hat\omega\hat\chi\hat\omega^{-1}{=}{-}\hat\chi & \mathrm{BM.74}\ \text{the ray-resolved spin census};\ j{=}|k|/2;\ \text{opposite rays at the vertex} & \mathrm{BM.75}\ \text{weight}\leftrightarrow\text{branch}\leftrightarrow\text{core unique};\ \Delta_g{=}(2,\sqrt2{-}1,2\cos\tfrac\pi{10}{-}\sqrt3)\hat\mu'^2\\
\mathrm{BM.76}\ \widehat V{=}E_+\otimes P_L;\ \widehat V^2{=}0;\ \operatorname{rk}\widehat V{=}1;\ [W_3\otimes I,\widehat V]{=}\widehat V & \mathrm{BM.77}\ \text{A-BM.8 issued};\ \Delta T{=}5\,m^2_B;\ \mathcal O_S\ne\varnothing;\ 828100{>}5819 & \mathrm{BM.78}\ \text{the tier adjudication};\ 1'\ \text{closed};\ \text{the decay allowed on}\ \mathcal O_S\\
\text{the decay closed on every pre-specification read} & \neg\text{Cal-BM parameter};\ \text{two addenda logged before use} & \neg\text{a measured value}\ [\mathrm{P5}]\\
\mathrm{BM.79}\ D(m){=}4\sin^2(m\delta_v/4){=}2{-}2T_{|m|}(\cos\tfrac{\delta_v}2);\ T_5(\mathsf c){=}0;\ r_{\mathrm{intra}}{=}\sin(k_f\delta_v/4)/\sin(\delta_v/4) & \mathrm{BM.80}\ 4z^3{-}24z^2{+}(21{-}12\gamma)z{-}(\gamma{-}1)^2;\ r^{(\ell)}_0{=}1{+}\gamma_q;\ q{<}\nu{<}\ell & \mathrm{BM.81}\ \Delta_1{>}\Delta_2{>}\Delta_3;\ \text{the doublet ratios};\ 3\ \neg\text{proportional ladders}\\
\mathrm{BM.82}\ \mathcal R_s{=}D_S/4E^2;\ \mathcal R_{\max}/\mathcal R_{\min}{=}50/(1{-}\mathsf c);\ 25{=}(30/6)^2 & \mathrm{BM.83}\ \rho_{\mathrm B},\sigma_\ell,\sigma_\nu;\ \ln1280/\ln960\notin\mathbb Q;\ t|_{m_n/m_p=4/3}{=}{-}\ln405/\ln16{<}0 & \mathrm{BM.84}\ m_n/m_p{<}21/20;\ \varepsilon^{\mathrm{crit}}\notin\mathbb Q;\ \text{the trichotomy}\\
\mathrm{BM.85}\ Q_\beta/m_n\le{-}\ln3/\ln1280;\ \text{two formed reads, disjoint alphabets} & \mathrm{BM.86}\ \text{the census adjudication};\ \neg\text{a declaration};\ \text{the ten open rows} & \neg\text{Cal-BM parameter};\ \neg\text{a measured value}\ [\mathrm{P5}]\\
\text{four currencies, zero conversions} & \text{item 7 stands with its residue named} & \text{the residue:}\ (\mathrm{R1})\text{--}(\mathrm{R3}),\ \sigma_\ell,\ \sigma_\nu,\ e\\
\mathrm{BM.87}\ c_{\mathrm{formed}}{=}\tfrac12\ \text{derived twice};\ \iota_{\mathrm{ball}}{=}\tfrac1{10}\ \text{void};\ \text{the selector, sign and monomial stand} & \mathrm{BM.88}\ e{=}e_0\sqrt{N\,f\,\mathsf g};\ f(k){=}1{-}\tfrac1k;\ G{=}2\pi E_{\mathrm{int}}^2\ln2;\ f(1){=}0 & \mathrm{BM.89}\ E^{\mathrm{self}}{=}\tfrac12Nw\mathsf g\,\varepsilon_{\mathrm{EM}}Q^2m_{\mathrm{hid}};\ w(0){=}1;\ w{=}1{-}4\varepsilon_q\ (\text{marked})\\
\mathrm{BM.90}\ w{=}w(\varepsilon_q)\ \text{branch-blind};\ \text{hypothesis (W) refuted on four grounds} & \mathrm{BM.91}\ \Delta m\,\mathfrak h{=}\ln\tfrac43{-}3Nw\mathsf g\,\varepsilon_{\mathrm{EM}}\ln2;\ Nw\mathsf g\,\varepsilon_{\mathrm{EM}}{\le}\tfrac18{\implies}m_n{>}m_p;\ m_n/m_p{<}\tfrac{10}9 & \mathrm{BM.92}\ Q_\beta\mathfrak h\le{-}\ln3\ \text{uniformly};\ \text{the specified read untouched}\\
\mathrm{BM.93}\ \text{the correction's adjudication};\ \neg\text{a declaration};\ (\mathrm{R4}),(\mathrm{R5})\ \text{named} & \neg\text{Cal-BM parameter};\ \neg\text{a measured coupling}\ [\mathrm{P5}] & \text{the residue:}\ (\mathrm{R1})\text{--}(\mathrm{R5}),\ e_0,\ \mathsf g,\ N\\
\mathrm{BM.94}\ w_{\mathrm{mat}}{=}1/(1{+}u{+}u^2)\in(\tfrac13,1];\ \langle E\rangle_{\mathrm{mat}}\in(2,6];\ \sigma/\langle E\rangle{<}\tfrac12 & \mathrm{BM.95}\ P(\text{fibre}\,|\,Q){=}\prod_im_{q_i}/W(Q);\ \ln\tfrac43{=}[\tfrac95\ln2{-}\ln3]{+}[\tfrac15\ln2] & \mathrm{BM.96}\ \delta_\nu{=}\tfrac18;\ G_2{>}G_1{>}A_0{>}G_0;\ m_\nu^{\mathrm{rest}}{>}0;\ (1,0);\ \neg\text{a seesaw-type mechanism}\\
\mathrm{BM.97}\ m^{\mathrm{rest}}{=}m_{\mathrm{pole}};\ \Delta_g^{\mathrm{crit}}{=}s(2M_0{+}s);\ \text{the specified read flicker-blind} & \mathrm{BM.98}\ \mathcal R{=}\tfrac{\pi\ln2}2\tfrac{\rho_{\mathrm B}D_S}{G}{=}\tfrac{\hat\mu'^2D_Sr^2}4(1{-}f)^2;\ \mathcal R_{\mathrm{mat}}/\mathcal R{=}D^2{<}9 & \mathrm{BM.99}\ N\mathsf g\,\varepsilon_{\mathrm{EM}}{\le}\tfrac18{\implies}\Delta m{>}0,\ {\ge}\tfrac5{12}{\implies}\Delta m{<}0;\ m_n/m_p{<}\ln80/\ln60{<}\tfrac{15}{14}\\
\mathrm{BM.100}\ \text{the tier's adjudication};\ \neg\text{a declaration};\ (\mathrm{R6})\ \text{named} & \neg\text{Cal-BM parameter};\ \neg\text{a measured value}\ [\mathrm{P5}] & \text{the residue:}\ (\mathrm{R1})\text{--}(\mathrm{R6}),\ e_0,\ \mathsf g,\ N,\ \rho_\ell,\ \rho_\nu\\
\mathrm{BM.101}\ Q^{\mathrm{inv}}{=}{-}Q^{\mathrm{fwd}};\ \varepsilon_\ast{=}(2{-}\log_23)/(3w)\in(\tfrac18,\tfrac5{12});\ P_{\max}{=}1{-}\tfrac3{2x_\ast} & \mathrm{BM.102}\ [\mathfrak h]{=}\mathfrak q^{-1/2},\ [\rho_{\mathrm B}]{=}\mathfrak q^0;\ \text{A-BM.9 declared, value}\ \neg\text{issued};\ \text{item 7 closed as an equation} & \mathrm{BM.103}\ W(Q){=}\binom6{Q+3};\ 4z^3{-}24z^2{+}(21{-}12\gamma)z{-}(\gamma{-}1)^2;\ \det(\lambda I{-}M)\ \neg\text{formed}\\
\mathrm{BM.104}\ G_{\mathrm{corr}}{=}G(1{+}a\mathsf g);\ \mathsf g_{\mathrm{fp}}{=}(1{+}G)/(1{-}aG);\ aG{=}\pi Nw\,\varepsilon_{\mathrm{EM}}Q^2E^2\ln2;\ \mathsf g(G)\ \neg\text{derived} & \mathrm{BM.105}\ w_{\mathrm{mat}}{=}\tfrac13\sum_ia_i\in[\tfrac13,1]\ \text{measure-free};\ S/(3(1{-}R));\ \text{hypothesis (W6) refuted} & \mathrm{BM.106}\ j{=}|\kappa|/2\ \text{with}\ j{=}\tfrac12,1\ \text{on one kind};\ \neg\text{an exchange operator};\ \text{deficit 11 stands};\ (\mathrm{R7})\ \text{named}\\
\mathrm{BM.107}\ \text{the six adjudications};\ \text{one declaration};\ (\mathrm{R4}'),(\mathrm{R7})\ \text{named} & \neg\text{a new dimensionful scale}\ [\mathrm{P1}];\ \neg\text{a measured value}\ [\mathrm{P5}] & \text{the residue:}\ (\mathrm{R1})\text{--}(\mathrm{R3}),\ (\mathrm{R4}'),\ (\mathrm{R5})\text{--}(\mathrm{R7}),\ \mathsf g,\ e_0,\ N\\
\mathrm{Def\text{-}BM.15}\ \mathrm{CP}:=\mathrm C\circ\mathrm P;\ c^{\mathrm{CP}};\ \text{the CP-pair};\ \text{the formal-mirror scope} & \mathrm{BM.108}\ \mu_{\mathrm{CP}}=e^{i\pi\Theta}=1\ \text{on}\ \Theta\in2\mathbb Z;\ \text{the forced conjugate};\ \neg\exists\tau^{-1} & \mathrm{BM.109}\ \mathcal W_{\text{pair}}=4\varepsilon_q(1{-}4\varepsilon_q)\le\tfrac14;\ Q_{\text{pair}}=0;\ \mathcal B_{\mathrm{info}}=0\\
\mathrm{Def\text{-}BM.16}\ \mathcal H_2=\mathcal H_\omega\otimes\Psi\ \text{(two chains)};\ \dim=4(|k_1|{+}1)(|k_2|{+}1);\ \hat P_{12}=\hat\mu_{\mathrm{CP}}\hat S_{12} & \mathrm{BM.110}\ g(2)=1:\ \text{the}\ \delta\text{-seal is unique};\ W_2^{\mathrm{CP}}(0)=W_2(0)=6;\ (\mathrm{R7})\ \text{discharged at}\ n=2 & \mathrm{BM.111}\ m_{\mathrm{hid}}^{\text{meson}}=2m_{\mathrm{hid}};\ \mathcal B_{\mathrm{info}}=\mathcal B_{\mathrm{pole}}=0;\ (\mathrm{R8})\ \text{named}\\
\mathrm{BM.112}\ \hat P_{12}^2=\mathrm{id};\ \operatorname{spec}\subseteq\{\pm1\};\ \operatorname{Tr}=d\,\mu_{\mathrm{CP}};\ \text{the orientation block is statistics-blind} & \mathrm{BM.113}\ \varepsilon=(-1)^{|k|}=(-1)^{2j};\ \textbf{deficit 11 closed};\ \mathrm{BM.6}(2)\ \text{discharged} & \neg\text{a spin-parity theorem}\ [\mathrm{BM.106}(1)\ \text{stands}]\\
\mathrm{Def\text{-}BM.17}\ \mathcal R=(n,\sigma,Q,\mathbf s,k);\ \mathcal T=\mathcal T_0/N;\ \Gamma=\Gamma_{\mathrm{tr}}f\,\tfrac{\varepsilon_{\mathrm{eff}}}{\varepsilon_{\mathrm{EM}}}\mathcal N\chi;\ \neg\text{a cosmological coordinate} & \mathrm{BM.114}\ (2{+}z{+}z^{-1})^n=(z{+}1)^{2n}/z^n;\ W_n(Q)=\binom{2n}{Q{+}n};\ \mathbb M,\ [\mathbb M:\mathbb Q]=16;\ 55\ \text{eigenvalues},\ 1\ \text{collision} & \mathrm{BM.115}\ \mathcal T_0=E_{\mathrm{int}}\mathfrak h/\ln2;\ \text{ladders}\ 2n..60n,\ \text{span}\ 30;\ \mathcal T=2r_s/N;\ \neg\text{a forbidden state}\\
\mathrm{BM.116}\ \Gamma_{\mathrm{tr}}\in[0,\tfrac13];\ (\mathrm Z1)\ \varepsilon_q{=}0,\ (\mathrm Z2)\ f{=}0,\ (\mathrm Z3)\ \mathcal N{=}0;\ \tau_{\mathrm{life}}\ge3\ \text{ticks} & \mathrm{BM.117}\ 20=\binom63\ \text{an orbit count},\ \neg\text{a decuplet};\ \lambda^{(n)}_{\max}=4n;\ \text{the octet}\ \neg\text{formed} & \mathrm{BM.118}\ \text{the tier adjudication};\ \neg\text{a declaration};\ (\mathrm{R8}),(\mathrm{R9})\ \text{named}\\
\text{four objectives, four derived verdicts} & \neg\text{a new dimensionful scale}\ [\mathrm{P1}];\ \neg\text{a measured value}\ [\mathrm{P5}] & \text{the residue:}\ (\mathrm{R1})\text{--}(\mathrm{R6}),\ (\mathrm{R7})\ \text{at}\ n\ge3,\ (\mathrm{R8}),(\mathrm{R9}),\ \mathsf g,\ e_0,\ N\\
\hline
\mathrm{Def\text{-}BM.18}\ \mathcal J=\operatorname{span}\{|\mathcal V\rangle,|\mathcal R\rangle\};\ \hat J_\pm;\ Q=x_+(\hat J_+{-}P_{\mathcal V})+x_-(\hat J_-{-}P_{\mathcal R});\ \mathcal T_{\mathcal J}:\ \dim=4(|k|{+}1) & \mathrm{BM.119}\ c{=}1\ \text{a convention};\ 1{-}v=\tfrac{m^2}{E(E+p)}{>}0;\ \text{the boundary is}\ \mathcal R/\mathcal V & \mathrm{BM.120}\ \hat J_\pm^2{=}0,\ \{\hat J_+,\hat J_-\}{=}\mathrm{id};\ \operatorname{spec}Q{=}\{0,-(x_+{+}x_-)\};\ \text{detailed balance derived}\\
\mathrm{BM.121}\ (\hat U^{2m})_{\varnothing\varnothing}{=}(x_+x_-)^m,\ (\hat U^{2m+1})_{\varnothing\varnothing}{=}0;\ \tfrac1{1-w^2x_+x_-};\ \text{pole}\ |w|\ge2 & \mathrm{BM.122}\ Q{=}0;\ \mathcal T_0{=}4\mathfrak h/\ln2\ \text{at the floor};\ \mathcal W{=}x_+^2\chi;\ \widehat V^2{=}0;\ \text{the two-photon bar} & \mathrm{BM.123}\ \mathcal W_{\mathrm{ann}}{=}x_-^2\chi;\ \tfrac{\mathcal W_{\mathrm{ann}}}{\mathcal W_{\mathrm{ch}}}{=}1\iff\varepsilon_q{=}\tfrac18;\ (\mathrm{R9})\ \text{closed at the channel level}\\
\mathrm{Def\text{-}BM.19}\ \Theta\ \text{vs}\ \Theta_R;\ \Delta(m){=}4\sin^2\tfrac{m\Theta_R}4{=}2{-}2T_{|m|}(\cos\tfrac{\Theta_R}2);\ \Delta(0){=}0\ \text{for even}\ k & \mathrm{BM.124}\ \Theta{=}2k{=}4j;\ \#\{m\}{=}2j{+}1;\ \Theta_R{=}2\pi k\ \text{impossible on dimension} & \mathrm{BM.125}\ 0\le m^2_{\mathrm{pole}}(m)\le4\hat\mu'^2;\ \hat\mu'^2\tfrac{\Theta_R^2}4{=}\hat\mu^2\ell_c^2G(c);\ \text{the product form}\ \neg\text{formed}\\
\mathrm{BM.126}\ \text{(W7) refuted on}\ (\mathrm O1)\text{--}(\mathrm O4);\ \text{the crossing is spin-blind};\ \text{the domain row is bound/transit} & \mathrm{BM.127}\ \varepsilon{=}(-1)^{|k|}\ \text{a ledger read};\ \text{the kind read separated} & \mathrm{BM.128}\ \dim\mathcal H_3{=}8\prod(|k_a|{+}1);\ \chi(\sigma){=}d^{c(\sigma)};\ \dim\wedge^3\mathbb C^2{=}0\ \text{vs}\ \dim\wedge^3\mathbb C^3{=}1;\ (\mathrm{R7})\ \text{closed at}\ n{=}3\\
\mathrm{BM.129}\ \varepsilon_M{=}{+}1\ \text{unconditionally};\ \varepsilon_B{=}(-1)^{k_1+k_2+k_3};\ \varepsilon(\Gamma){=}\prod_a\varepsilon(c_a) & \mathrm{BM.130}\ \text{the tier adjudication};\ \neg\text{a declaration};\ (\mathrm{R7}'),(\mathrm{R9}'),(\mathrm{R9}''),(\mathrm{R10})\ \text{named} & \text{the residue:}\ (\mathrm{R1})\text{--}(\mathrm{R6}),\ (\mathrm{R7}'),(\mathrm{R8}),(\mathrm{R9}'),(\mathrm{R9}''),(\mathrm{R10}),\ \mathsf g,\ e_0,\ N\\
\hline
\mathrm{BM.131}\ \mathrm{tw}(\rho){=}\tfrac{\mathfrak T(\Pi_{\mathcal R}\rho)}3\Pi_{\mathcal R}{+}\mathfrak T(\Pi_{\mathcal V}\rho)\Pi_{\mathcal V};\ \delta K{=}\gamma\rho{+}\log Z_\gamma I;\ \text{the common fixed point}\ \{\tfrac14I\};\ V(s){=}(1{-}\tfrac\gamma4)^{2s/N}V_0 & \mathrm{BM.132}\ N(k,E){=}\sqrt{1{-}\tfrac{G}{4\pi k}};\ \text{formed}\iff k\ge\tfrac{G}{4\pi}{=}\tfrac{E_{\mathrm{int}}^2\ln2}2;\ \text{the envelope}\ \neg\text{formed};\ N{\to}1 & \mathrm{BM.133}\ \mathsf g(G){=}\tfrac1{1+G}\ \text{selected};\ e{=}e_0[Nf/(1{+}G)]^{1/2};\ \text{the loop stays excluded}\\
\mathrm{Def\text{-}BM.20}\ R(z;\mathbf p);\ \text{the embedding and Fourier rows declared (A-BM.12)};\ \mathcal V{=}1\ \text{minimal} & \mathrm{BM.134}\ \text{the pole locus is the dispersion};\ \text{form-unique};\ d\mu_{\mathrm{shell}}{=}\tfrac{d^3p}{2E};\ \Phi_2{=}\tfrac{\sqrt{z^2-4N^2m^2}}{8\pi N^3z};\ \text{item 17 closed} & \mathrm{BM.135}\ |Q|{=}2:\ 12{+}2\ \text{states};\ |Q|{=}3:\ 2;\ \partial_\mu F^{\mu\nu}{=}eJ_Q^\nu;\ \text{Q10 closed};\ \text{deficit 10 closed}\\
\mathrm{BM.136}\ \varepsilon_{\sigma(i)\sigma(j)\sigma(k)}{=}\mathrm{sgn}(\sigma)\varepsilon_{ijk};\ \mathcal H_n{=}2^n\prod(|k_a|{+}1),\ \chi(\sigma){=}d^{c(\sigma)}\ \text{at every}\ n;\ \dim\operatorname{Inv}((\mathbb C^3)^{\otimes4}){=}0;\ (1,1,1){:}\ \varepsilon_B{=}{-}1,\ j_B{=}\tfrac32;\ 112/104\ \text{in B} & \mathrm{BM.137}\ w^{(n)}_{\mathrm{mat}}{=}\tfrac1{D_n};\ \mathcal B^{(2)}_{\mathrm{ens}}{=}{-}\tfrac{4u\ln2}{(1{+}u)\mathfrak h}{<}0;\ m_2^2{\equiv}0\ \text{by the pincer};\ 21\ \lambda^{(2)},\ \neg\text{a collision};\ (\mathrm{R8})\ \text{closed} & \mathrm{BM.138}\ \Psi{=}\Pi_k(\mathcal H_{\mathsf v}^{\otimes k});\ \operatorname{rk}\Pi_k{=}k{+}1;\ \text{the connector tensors alike};\ \text{item 19 closed}\\
\mathrm{BM.139}\ \text{A-BM.10}\ (\aleph\text{-domain, pointwise: ensembles literal, junctions E2});\ \text{A-BM.11}\ (\text{variant B});\ \text{A-BM.13}\ (\vee^2\ \text{locus});\ \text{item 16 closed} & \mathrm{BM.140}\ \delta_v^{\mathrm{dodeca}}{=}\tfrac\pi5\ \text{minimal, unique};\ x^4{-}48x^3{+}304x^2{-}512x{+}256;\ \text{five currencies censused};\ \text{deficit 12 adjudicated} & \mathrm{BM.141}\ J_3{=}\tfrac{p(1-p)(9-10p+2p^2)}{2D^2};\ J_2{\le}\tfrac18\ \text{at}\ p{=}\tfrac23;\ \sup J_3\in(\tfrac{93}{500},\tfrac3{16});\ (\mathrm{R6})\ \text{closed}\\
\mathrm{BM.142}\ (\hat J_+\otimes\hat J_+)^2{=}0;\ \mathcal M{=}e^2q^2x_+^2\chi;\ \sigma_{\mathrm{pair}},\sigma_{\mathrm{ann}}\ \text{exact symbolic},\ [\sigma]{=}\mathfrak q^{-1};\ \tfrac{\sigma_{\mathrm{ann}}}{\sigma_{\mathrm{pair}}}{=}(\tfrac{x_-}{x_+})^4\tfrac{z^2}{z^2-4N^2m^2};\ (\mathrm{R9}'),(\mathrm{R9}'')\ \text{closed} & \mathrm{BM.143}\ (\mathrm{R5})\ \text{void};\ (\mathrm{R2}),(\mathrm{R10})\ \text{registration data};\ (\mathrm{R4}')\ \text{void for BM};\ \text{C-BM.1--3 consolidated};\ (\mathrm{P5.1})\text{--}(\mathrm{P5.10})\ \text{named} & \mathrm{BM.144}\ \text{four declarations};\ \text{zero P4 additions};\ \neg\text{a retraction in}\ \mathrm{BM.1}\text{--}\mathrm{BM.130};\ \textbf{the era closed}\\
\text{the residue: every row dispositioned}\ [\mathrm{BM.143}(6)] & \text{for Part 5: the EM binding channel first}\ [\mathrm{BM.143}(7)] & \textbf{the Baryomesonic era is closed; the books are shelved}\\
\hline
\mathrm{BM.145}\ (3{+}3{+}5{+}3{+}1{+}1){=}16{=}C_{\mathrm{cap}};\ \dim\mathcal K_n{=}(6,16,17,25);\ \hat{\mathcal F}:\ (\Gamma_{\mathcal R},\text{registrations})\mapsto(M,g_{\mu\nu},T_{\mu\nu},\operatorname{spec});\ O(n){=}\max(0,2^{n+1}{-}16) & \mathrm{BM.146}\ \mathcal C(E,r){=}(g,N,\mathcal T,w_{\mathrm{mat}},\Gamma)\ \text{closed forms};\ \text{formed}\iff E\ge2\wedge\varkappa\|q\|{<}1\wedge E\le2r/\mathfrak h\wedge n\le3;\ \#\text{edges/cell}\ge2sE\ln2/\mathfrak h\to\infty & \mathrm{BM.147}\ g_{\mu\nu}{=}\eta_{\mu\nu}{+}\tfrac{\mathfrak hE}{2r}\mathbb 1_4{+}O(\cdot)\ \text{pure trace};\ \Omega^2{=}1{-}\tfrac{\mathfrak hE}{2r};\ m{=}\tfrac1{\kappa^2}\int R\,dV\\
\mathrm{BM.148}\ E_{\mathrm{crit}}(r){=}2r/\mathfrak h;\ S_{\mathrm{crit}}{=}2r\ln2/\mathfrak h;\ k_{\mathrm{crit}}{=}E^2\ln2/2;\ E_{\mathrm{crit}}(\mathfrak h){=}2;\ \text{no cell self-formed at its envelope} & \mathrm{BM.149}\ \langle\rho_{D^-}\rangle(1){=}(448,896,2976,2240,24384);\ \mathcal P_M^{(n)}F_{D^-}\mathcal P_M^{(n)}{=}0\ (n\le3);\ \text{what the metric cannot read, the ledger retains} & \mathrm{BM.150}\ \text{one operator, one arena, one action};\ \text{three fingerprints}\ (T\text{-channel},\ \text{mask slot},\ \text{coupling profile});\ \text{the pp-wave witness is source-free}\\
\mathrm{BM.151}\ |r_a|\le\sqrt3\ \text{exact};\ |\Delta|\le2\sqrt3;\ \mathrm{TV}\le2\sqrt3 m;\ \mathrm{CT2.6}(\mathrm{ii}),(\mathrm{iii})\ \textbf{reclassified as derived};\ \textbf{the extra-axiomatic list is empty} & \mathrm{BM.152}\ \text{zero declarations, one retired};\ \neg\text{a retraction in}\ \mathrm{BM.1}\text{--}\mathrm{BM.144};\ \textbf{Era 1 closed} & \textbf{the foundational era is closed; the corpus stands on its axioms alone}\\
\end{array}
\]

$\blacksquare$

